% Auto-generated from data/segments.json + layout.json (+ transforms) — do not edit by hand.
% volume: 14Sam03
% mode: printing
% segments: 2829
% transform_rules: 8
% toc_marks: 636

% Volume layout defaults (layout.json ``layout``).
\csromanlayoutapply{1}{1.5}{21.6pt}{8pt}{8pt}{65pt}{2.5em}%

\setcsromanpitaka{สุตฺตนฺตปิฏก}
\setcsromannikaya{สํยุตฺตนิกาย}
\csromanheader{2}{\textbf{สํยุตฺตนิกาย}}
\csromanmatikaanchor{mtk.0}
\csromantocmarknopage{chapter}{มหาวคฺคสํยุตฺตปาฬิ}{mtk.0}
\bookmark[dest={mtk.0},level=0]{มหาวคฺคสํยุตฺตปาฬิ}
\gambhira{มหาวคฺค\-สํยุตฺต\-ปาฬิ}
\namakkaram{นโม ตสฺส ภควโต อรหโต สมฺมา\-สมฺพุทฺธสฺส.}
\csromanheader{2}{1. \textbf{มคฺคสํยุตฺต}}
\chapterhead{1. \textbf{อวิชฺชาวคฺค}}
\csromanheader{2}{1. \textbf{อวิชฺชาสุตฺต}}
\proseitem{1}{\csromanfolio{1}เอวํ เม สุตํ\csromandash{}เอกํ สมยํ ภควา สาวตฺถิยํ วิหรติ เชตวเน อนาถ\-ปิณฺฑิกสฺส อาราเม. ตตฺร โข ภควา ภิกฺขู อามนฺเตสิ “ภิกฺขโว”ติ. “ภทนฺเต”ติ เต ภิกฺขู ภควโต ปจฺจสฺโสสุํ. ภควา เอตทโวจ\csromandash{}}

\prose{อวิชฺชา ภิกฺขเว ปุพฺพงฺคมา อกุสลานํ ธมฺมานํ สมาปตฺติยา, อนฺวเทว\csromansharedfootnote{f0e6a939c1847ebd}{อนุเทว (สี, อิ, ก)} อหิริกํ อโนตฺตปฺปํ. อวิชฺชาคตสฺส ภิกฺขเว อวิทฺทสุโน มิจฺฉาทิฏฺฐิ ปโหติ, มิจฺฉา\-ทิฏฺฐิสฺส มิจฺฉาสงฺกปฺโป ปโหติ, มิจฺฉา\-สงฺกปฺปสฺส มิจฺฉาวาจา ปโหติ, มิจฺฉาวาจสฺส มิจฺฉากมฺมนฺโต ปโหติ, มิจฺฉา\-กมฺมนฺตสฺส มิจฺฉาอาชีโว ปโหติ, มิจฺฉา อาชีวสฺส มิจฺฉาวายาโม ปโหติ, มิจฺฉา\-วายามสฺส มิจฺฉาสติ ปโหติ, มิจฺฉาสติสฺส มิจฺฉาสมาธิ ปโหติ.}

\prose{\csromanfolio{2}วิชฺชา จ โข ภิกฺขเว ปุพฺพงฺคมา กุสลานํ ธมฺมานํ สมาปตฺติยา, อนฺวเทว หิโรตฺตปฺปํ. วิชฺชาคตสฺส ภิกฺขเว วิทฺทสุโน สมฺมาทิฏฺฐิ ปโหติ, สมฺมาทิฏฺฐิสฺส สมฺมาสงฺกปฺโป ปโหติ, สมฺมา\-สงฺกปฺปสฺส สมฺมาวาจา ปโหติ, สมฺมาวาจสฺส สมฺมากมฺมนฺโต ปโหติ, สมฺมา\-กมฺมนฺตสฺส สมฺมาอาชีโว ปโหติ, สมฺมาอาชีวสฺส สมฺมาวายาโม ปโหติ, สมฺมาวายามสฺส สมฺมาสติ ปโหติ, สมฺมาสติสฺส สมฺมาสมาธิ ปโหตีติ.  ปฐมํ.}

\csromanheader{2}{2. \textbf{อุปฑฺฒสุตฺต}}
\proseitem{2}{เอวํ เม สุตํ\csromandash{}เอกํ สมยํ ภควา สเกฺยสุ วิหรติ นครกํ นาม\csromansharedfootnote{a679bcadf6666454}{นาครกํ นาม (สี), สกฺกรํ นาม (สฺยา, ก)} สกฺยานํ นิคโม. อถ โข อายสฺมา อานนฺโท เยน ภควา เตนุปสงฺกมิ, อุปสงฺกมิตฺวา ภควนฺตํ อภิวาเทตฺวา เอกมนฺตํ นิสีทิ, เอกมนฺตํ นิสินฺโน โข อายสฺมา อานนฺโท ภควนฺตํ เอตทโวจ “อุปฑฺฒมิทํ ภนฺเต พฺรหฺมจริยํ ยทิทํ กลฺยาณมิตฺตตา กลฺยาณ\-สหายตา กลฺยาณ\-สมฺปวงฺกตา”ติ.}

\prose{มา เหวํ อานนฺท, มา เหวํ อานนฺท, สกลเมวิทํ อานนฺท พฺรหฺมจริยํ ยทิทํ กลฺยาณมิตฺตตา กลฺยาณ\-สหายตา กลฺยาณ\-สมฺปวงฺกตา, กลฺยาณมิตฺตสฺเสตํ อานนฺท ภิกฺขุโน ปาฏิกงฺขํ กลฺยาณ\-สหายสฺส กลฺยาณ\-สมฺปวงฺกสฺส, อริยํ อฏฺฐงฺคิกํ มคฺคํ ภาเวสฺสติ, อริยํ อฏฺฐงฺคิกํ มคฺคํ พหุลีกริสฺสติ.}

\prose{กถญฺจานนฺท ภิกฺขุ กลฺยาณมิตฺโต กลฺยาณสหาโย กลฺยาณ\-สมฺปวงฺโก อริยํ อฏฺฐงฺคิกํ มคฺคํ ภาเวติ, อริยํ อฏฺฐงฺคิกํ มคฺคํ พหุลีกโรติ. อิธานนฺท ภิกฺขุ สมฺมาทิฏฺฐิํ ภาเวติ วิเวกนิสฺสิตํ วิราคนิสฺสิตํ นิโรธ\-นิสฺสิตํ โวสฺสคฺค\-ปริณามิํ, สมฺมาสงฺกปฺปํ ภาเวติ วิเวกนิสฺสิตํ ฯเปฯ สมฺมาวาจํ ภาเวติ ฯเปฯ สมฺมากมฺมนฺตํ ภาเวติ ฯเปฯ สมฺมาอาชีวํ ภาเวติ ฯเปฯ สมฺมาวายามํ ภาเวติ ฯเปฯ สมฺมาสติํ ภาเวติ ฯเปฯ สมฺมาสมาธิํ ภาเวติ วิเวกนิสฺสิตํ วิราคนิสฺสิตํ นิโรธ\-นิสฺสิตํ โวสฺสคฺค\-ปริณามิํ. เอวํ โข อานนฺท ภิกฺขุ กลฺยาณมิตฺโต กลฺยาณสหาโย กลฺยาณ\-สมฺปวงฺโก อริยํ อฏฺฐงฺคิกํ มคฺคํ ภาเวติ, อริยํ อฏฺฐงฺคิกํ มคฺคํ พหุลีกโรติ.}

\prose{\csromanfolio{3}ตทมินาเปตํ อานนฺท ปริยาเยน เวทิตพฺพํ, ยถา สกลเมวิทํ พฺรหฺมจริยํ ยทิทํ กลฺยาณมิตฺตตา กลฺยาณ\-สหายตา กลฺยาณสมฺป- วงฺกตา. มมํ หิ อานนฺท กลฺยาณมิตฺตํ อาคมฺม ชาติธมฺมา สตฺตา ชาติยา ปริมุจฺจนฺติ, ชราธมฺมา สตฺตา ชราย ปริมุจฺจนฺติ, มรณธมฺมา สตฺตา มรเณน ปริมุจฺจนฺติ, โสก\-ปริเทว\-ทุกฺข\-โทมนสฺสุ\-ปายาส\-ธมฺมา สตฺตา โสก\-ปริเทว\-ทุกฺข\-โทมนสฺสุปายาเสหิ ปริมุจฺจนฺติ. อิมินา โข เอตํ อานนฺท ปริยาเยน เวทิตพฺพํ, ยถา สกลเมวิทํ พฺรหฺมจริยํ ยทิทํ กลฺยาณมิตฺตตา กลฺยาณ\-สหายตา กลฺยาณ\-สมฺปวงฺกตาติ.  ทุติยํ.}

\csromanheader{2}{3. \textbf{สาริปุตฺตสุตฺต}}
\proseitem{3}{สาวตฺถิ\-นิทานํ. อถ โข อายสฺมา สาริปุตฺโต เยน ภควา เตนุปสงฺกมิ, อุปสงฺกมิตฺวา ภควนฺตํ อภิวาเทตฺวา เอกมนฺตํ นิสีทิ, เอกมนฺตํ นิสินฺโน โข อายสฺมา สาริปุตฺโต ภควนฺตํ เอตทโวจ “สกลมิทํ ภนฺเต พฺรหฺมจริยํ ยทิทํ กลฺยาณมิตฺตตา กลฺยาณ\-สหายตา กลฺยาณ\-สมฺปวงฺกตา”ติ.}

\prose{สาธุ สาธุ สาริปุตฺต, สกลมิทํ สาริปุตฺต พฺรหฺมจริยํ ยทิทํ กลฺยาณมิตฺตตา กลฺยาณ\-สหายตา กลฺยาณ\-สมฺปวงฺกตา. กลฺยาณมิตฺตสฺเสตํ สาริปุตฺต ภิกฺขุโน ปาฏิกงฺขํ กลฺยาณ\-สหายสฺส กลฺยาณ\-สมฺปวงฺกสฺส, อริยํ อฏฺฐงฺคิกํ มคฺคํ ภาเวสฺสติ, อริยํ อฏฺฐงฺคิกํ มคฺคํ พหุลีกริสฺสติ. กถญฺจ สาริปุตฺต ภิกฺขุ กลฺยาณมิตฺโต กลฺยาณสหาโย กลฺยาณ\-สมฺปวงฺโก อริยํ อฏฺฐงฺคิกํ มคฺคํ ภาเวติ, อริยํ อฏฺฐงฺคิกํ มคฺคํ พหุลีกโรติ.}

\prose{อิธ สาริปุตฺต ภิกฺขุ สมฺมาทิฏฺฐิํ ภาเวติ วิเวกนิสฺสิตํ วิราคนิสฺสิตํ นิโรธ\-นิสฺสิตํ โวสฺสคฺค\-ปริณามิํ ฯเปฯ สมฺมาสมาธิํ ภาเวติ วิเวกนิสฺสิตํ วิราคนิสฺสิตํ นิโรธ\-นิสฺสิตํ โวสฺสคฺค\-ปริณามิํ. เอวํ โข สาริปุตฺต ภิกฺขุ กลฺยาณมิตฺโต กลฺยาณสหาโย กลฺยาณ\-สมฺปวงฺโก อริยํ อฏฺฐงฺคิกํ มคฺคํ ภาเวติ, อริยํ อฏฺฐงฺคิกํ มคฺคํ พหุลีกโรติ.}

\prose{ตทมินาเปตํ สาริปุตฺต ปริยาเยน เวทิตพฺพํ, ยถา สกลมิทํ พฺรหฺมจริยํ ยทิทํ กลฺยาณมิตฺตตา กลฺยาณ\-สหายตา กลฺยาณ\-สมฺปวงฺกตา. มมํ หิ สาริปุตฺต กลฺยาณมิตฺตํ อาคมฺม ชาติธมฺมา สตฺตา \csromanfolio{4}ชาติยา ปริมุจฺจนฺติ, ชราธมฺมา สตฺตา ชราย ปริมุจฺจนฺติ, มรณธมฺมา สตฺตา มรเณน ปริมุจฺจนฺติ, โสก\-ปริเทว\-ทุกฺข\-โทมนสฺสุ\-ปายาส\-ธมฺมา สตฺตา โสก\-ปริเทว\-ทุกฺข\-โทมนสฺสุปายาเสหิ ปริมุจฺจนฺติ. อิมินา โข เอตํ สาริปุตฺต ปริยาเยน เวทิตพฺพํ, ยถา สกลมิทํ พฺรหฺมจริยํ ยทิทํ กลฺยาณมิตฺตตา กลฺยาณ\-สหายตา กลฺยาณ\-สมฺปวงฺกตาติ.  ตติยํ.}

\csromanheader{2}{4. ชาณุสฺโสณิ\-พฺราหฺมณ\-สุตฺต}
\proseitem{3}{สาวตฺถิ\-นิทานํ. อถ โข อายสฺมา อานนฺโท ปุพฺพณฺหสมยํ นิวาเสตฺวา ปตฺต\-จีวรมา\-ทาย สาวตฺถิํ ปิณฺฑาย ปาวิสิ. อทฺทสา โข อายสฺมา อานนฺโท ชาณุสฺโสณิํ พฺราหฺมณํ สพฺพเสเตน วฬวาภิ\-รเถน\csromansharedfootnote{7197fa05b4674517}{วฬภีรเถน (สี)} สาวตฺถิยา นิยฺยายนฺตํ, เสตา สุทํ อสฺสา ยุตฺตา โหนฺติ เสตาลงฺการา เสโต รโถ เสตปริวาโร เสตา รสฺมิโย เสตา ปโตทลฏฺฐิ เสตํ ฉตฺตํ เสตํ อุณฺหีสํ เสตานิ วตฺถานิ เสตา อุปาหนา เสตาย สุทํ วาลพีชนิยา พีชียติ. ตเมนํ ชโน ทิสฺวา เอวมาห “พฺรหฺมํ วต โภ ยานํ, พฺรหฺม\-ยานรูปํ วต โภ”ติ.}

\prose{อถ โข อายสฺมา อานนฺโท สาวตฺถิยํ ปิณฺฑาย จริตฺวา ปจฺฉาภตฺตํ ปิณฺฑปาต\-ปฏิกฺกนฺโต เยน ภควา เตนุปสงฺกมิ, อุปสงฺกมิตฺวา ภควนฺตํ อภิวาเทตฺวา เอกมนฺตํ นิสีทิ, เอกมนฺตํ นิสินฺโน โข อายสฺมา อานนฺโท ภควนฺตํ เอตทโวจ\csromandash{}}

\prose{อิธาหํ ภนฺเต ปุพฺพณฺหสมยํ นิวาเสตฺวา ปตฺต\-จีวรมา\-ทาย สาวตฺถิํ ปิณฺฑาย ปาวิสิํ, อทฺทสํ ขฺวาหํ ภนฺเต ชาณุสฺโสณิํ พฺราหฺมณํ สพฺพเสเตน วฬวาภิ\-รเถน สาวตฺถิยา นิยฺยายนฺตํ, เสตา สุทํ อสฺสา ยุตฺตา โหนฺติ เสตาลงฺการา เสโต รโถ เสตปริวาโร เสตา รสฺมิโย เสตา ปโตทลฏฺฐิ เสตํ ฉตฺตํ เสตํ อุณฺหีสํ เสตานิ วตฺถานิ เสตา อุปาหนา เสตาย สุทํ วาลพีชนิยา พีชียติ. ตเมนํ ชโน ทิสฺวา เอวมาห “พฺรหฺมํ วต โภ ยานํ, พฺรหฺม\-ยานรูปํ วต โภ”ติ. สกฺกา นุ โข ภนฺเต อิมสฺมิํ ธมฺมวินเย พฺรหฺมยานํ ปญฺญาเปตุนฺติ.}

\prose{\csromanfolio{5}“สกฺกา อานนฺทา”ติ ภควา อโวจ. อิมสฺเสว โข เอตํ อานนฺท อริยสฺส อฏฺฐงฺคิกสฺส มคฺคสฺส อธิวจนํ พฺรหฺมยานํ อิติปิ ธมฺมยานํ อิติปิ อนุตฺตโร สงฺคามวิชโย อิติปีติ.}

\prose{สมฺมาทิฏฺฐิ อานนฺท ภาวิตา พหุลีกตา ราค\-วินย\-ปริโยสานา โหติ, โทส\-วินย\-ปริโยสานา โหติ, โมห\-วินย\-ปริโยสานา โหติ. สมฺมาสงฺกปฺโป อานนฺท ภาวิโต พหุลีกโต ราค\-วินย\-ปริโยสาโน โหติ, โทสวินย\-ปริโยสาโน โหติ, โมหวินย\-ปริโยสาโน โหติ. สมฺมาวาจา อานนฺท ภาวิตา พหุลีกตา ราค\-วินย\-ปริโยสานา โหติ, โทส ฯเปฯ โมห\-วินย\-ปริโยสานา โหติ. สมฺมากมฺมนฺโต อานนฺท ภาวิโต พหุลีกโต ราค\-วินย\-ปริโยสาโน โหติ, โทส ฯเปฯ โมหวินย\-ปริโยสาโน โหติ. สมฺมาอาชีโว อานนฺท ภาวิโต พหุลีกโต ราค\-วินย\-ปริโยสาโน โหติ, โทส ฯเปฯ โมหวินย\-ปริโยสาโน โหติ. สมฺมาวายาโม อานนฺท ภาวิโต พหุลีกโต ราค\-วินย\-ปริโยสาโน โหติ, โทส ฯเปฯ โมหวินย\-ปริโยสาโน โหติ. สมฺมาสติ อานนฺท ภาวิตา พหุลีกตา ราค\-วินย\-ปริโยสานา โหติ, โทส ฯเปฯ โมห\-วินย\-ปริโยสานา โหติ. สมฺมาสมาธิ อานนฺท ภาวิโต พหุลีกโต ราค\-วินย\-ปริโยสาโน โหติ, โทส ฯเปฯ โมหวินย\-ปริโยสาโน โหติ.}

\prose{อิมินา โข เอตํ อานนฺท ปริยาเยน เวทิตพฺพํ, ยถา อิมสฺเสเวตํ อริยสฺส อฏฺฐงฺคิกสฺส มคฺคสฺส อธิวจนํ พฺรหฺมยานํ อิติปิ ธมฺมยานํ อิติปิ อนุตฺตโร สงฺคามวิชโย อิติปีติ. อิทมโวจ ภควา, อิทํ วตฺวาน สุคโต อถาปรํ เอตทโวจ สตฺถา\csromandash{}}

\setlength{\csromangathapagewidth}{0pt}
\setlength{\csromangathawakwidth}{0pt}
\csromangathameasuregroup{“ยสฺส สทฺธา จ ปญฺญา จ, \\ หิรี อีสา มโนโยตฺตํ, \\ รโถ สีลปริกฺขาโร, \\ อุเปกฺขา ธุรสมาธิ, \\ อพฺยาปาโท อวิหิํสา, \\ ติติกฺขาจมฺมสนฺนาโห\textsuperscript{1}, \\ เอตทตฺตนิ สมฺภูตํ,}{\csromangathabat{“ยสฺส สทฺธา จ ปญฺญา จ,}{ธมฺมา ยุตฺตา สทา ธุรํ.} \\ \csromangathabat{หิรี อีสา มโนโยตฺตํ,}{สติ อารกฺขสารถิ.} \\ \csromangathabat{รโถ สีลปริกฺขาโร,}{ฌานกฺโข จกฺกวีริโย.} \\ \csromangathabat{อุเปกฺขา ธุรสมาธิ,}{อนิจฺฉา ปริวารณํ.} \\ \csromangathabat{อพฺยาปาโท อวิหิํสา,}{วิเวโก ยสฺส อาวุธํ.} \\ \csromangathabat{ติติกฺขาจมฺมสนฺนาโห\textsuperscript{1},}{โยคกฺเขมาย วตฺตติ.} \\ \csromangathabat{เอตทตฺตนิ สมฺภูตํ,}{พฺรหฺมยานํ อนุตฺตรํ.}}
\csromangathasetleft{“ยสฺส สทฺธา จ ปญฺญา จ, \\ หิรี อีสา มโนโยตฺตํ, \\ รโถ สีลปริกฺขาโร, \\ อุเปกฺขา ธุรสมาธิ, \\ อพฺยาปาโท อวิหิํสา, \\ ติติกฺขาจมฺมสนฺนาโห\textsuperscript{1}, \\ เอตทตฺตนิ สมฺภูตํ,}
\csromangathagroup{\csromangathastanza{\csromangathabat{“ยสฺส สทฺธา จ ปญฺญา จ,}{ธมฺมา ยุตฺตา สทา ธุรํ.} \\ \csromangathabat{หิรี อีสา มโนโยตฺตํ,}{สติ อารกฺขสารถิ.}}\csromangathastanza{\csromangathabat{รโถ สีลปริกฺขาโร,}{ฌานกฺโข จกฺกวีริโย.} \\ \csromangathabat{อุเปกฺขา ธุรสมาธิ,}{อนิจฺฉา ปริวารณํ.}}\csromangathastanza{\csromangathabat{อพฺยาปาโท อวิหิํสา,}{วิเวโก ยสฺส อาวุธํ.} \\ \csromangathabat{ติติกฺขาจมฺมสนฺนาโห\csromansharedfootnote{240127230455fc85}{วมฺมสนฺนาโห (สี)},}{โยคกฺเขมาย วตฺตติ.} \\ \csromangathabat{เอตทตฺตนิ สมฺภูตํ,}{พฺรหฺมยานํ อนุตฺตรํ.}}}

\vspace{\tipitakaparskip}
\csromancenter{นิยฺยนฺติ ธีรา โลกมฺหา, อญฺญทตฺถุ ชยํ ชยนฺ”ติ. จตุตฺถํ.}
\csromanheader{2}{5. กิมตฺถิย\-สุตฺต}
\proseitem{5}{\csromanfolio{6}สาวตฺถิ\-นิทานํ. อถ โข สมฺพหุลา ภิกฺขู เยน ภควา เตนุ\-ปสงฺกมิํสุ ฯเปฯ เอกมนฺตํ นิสีทิํสุ, เอกมนฺตํ นิสินฺนา โข เต ภิกฺขู ภควนฺตํ เอตทโวจุํ\csromandash{}}

\prose{อิธ โน ภนฺเต อญฺญติตฺถิยา ปริพฺพาชกา อเมฺห เอวํ ปุจฺฉนฺติ “กิมตฺถิยํ อาวุโส สมเณ โคตเม พฺรหฺมจริยํ วุสฺสตี”ติ. เอวํ ปุฏฺฐา มยํ ภนฺเต เตสํ อญฺญ\-ติตฺถิยานํ ปริพฺพาชกานํ เอวํ พฺยากโรม “ทุกฺขสฺส โข อาวุโส ปริญฺญตฺถํ ภควติ พฺรหฺมจริยํ วุสฺสตี”ติ. กจฺจิ มยํ ภนฺเต เอวํ ปุฏฺฐา เอวํ พฺยากรมานา วุตฺตวาทิโน เจว ภควโต โหม, น จ ภควนฺตํ อภูเตน อพฺภาจิกฺขาม, ธมฺมสฺส จานุธมฺมํ พฺยากโรม, น จ โกจิ สหธมฺมิโก วาทานุวาโท คารยฺหํ ฐานํ อาคจฺฉตีติ.}

\prose{ตคฺฆ ตุเมฺห ภิกฺขเว เอวํ ปุฏฺฐา เอวํ พฺยากรมานา วุตฺตวาทิโน เจว เม โหถ, น จ มํ อภูเตน อพฺภาจิกฺขถ, ธมฺมสฺส จานุธมฺมํ พฺยากโรถ, น จ โกจิ สหธมฺมิโก วาทานุวาโท คารยฺหํ ฐานํ อาคจฺฉติ. ทุกฺขสฺส หิ ปริญฺญตฺถํ มยิ พฺรหฺมจริยํ วุสฺสติ. สเจ โว ภิกฺขเว อญฺญติตฺถิยา ปริพฺพาชกา เอวํ ปุจฺเฉยฺยุํ “อตฺถิ ปนาวุโส มคฺโค, อตฺถิ ปฏิปทา เอตสฺส ทุกฺขสฺส ปริญฺญายา”ติ. เอวํ ปุฏฺฐา ตุเมฺห ภิกฺขเว เตสํ อญฺญ\-ติตฺถิยานํ ปริพฺพาชกานํ เอวํ พฺยากเรยฺยาถ “อตฺถิ โข อาวุโส มคฺโค, อตฺถิ ปฏิปทา เอตสฺส ทุกฺขสฺส ปริญฺญายา”ติ.}

\prose{กตโม จ ภิกฺขเว มคฺโค, กตมา ปฏิปทา เอตสฺส ทุกฺขสฺส ปริญฺญายาติ. อยเมว อริโย อฏฺฐงฺคิโก มคฺโค. เสยฺยถิทํ, สมฺมาทิฏฺฐิ ฯเปฯ สมฺมาสมาธิ. อยํ ภิกฺขเว มคฺโค, อยํ ปฏิปทา เอตสฺส ทุกฺขสฺส ปริญฺญายาติ. เอวํ ปุฏฺฐา ตุเมฺห ภิกฺขเว เตสํ อญฺญ\-ติตฺถิยานํ ปริพฺพาชกานํ เอวํ พฺยากเรยฺยาถาติ.  ปญฺจมํ.}

\csromanmatikaanchor{mtk.1}
\csromanmatikaanchor{mtk.308}
\csromantocmarknopage{chapter}{4. นีวรณวคฺค}{mtk.1}
\bookmark[dest={mtk.1},level=0]{4. นีวรณวคฺค}
\csromanheader{6}{6. ปฐม\-อญฺญตรภิกฺขุ\-สุตฺต}
\proseitem{6}{\csromanfolio{7}สาวตฺถิ\-นิทานํ. อถ โข อญฺญตโร ภิกฺขุ เยน ภควา เตนุปสงฺกมิ ฯเปฯ เอกมนฺตํ นิสินฺโน โข โส ภิกฺขุ ภควนฺตํ เอตทโวจ\csromandash{}“พฺรหฺมจริยํ พฺรหฺม\-จริยนฺติ ภนฺเต วุจฺจติ, กตมํ นุ โข ภนฺเต พฺรหฺมจริยํ, กตมํ พฺรหฺมจริยปริโยสานนฺ”ติ.}

\prose{อยเมว โข ภิกฺขุ อริโย อฏฺฐงฺคิโก มคฺโค พฺรหฺมจริยํ. เสยฺยถิทํ, สมฺมาทิฏฺฐิ ฯเปฯ สมฺมาสมาธิ. โย โข ภิกฺขุ ราคกฺขโย โทสกฺขโย โมหกฺขโย, อิทํ พฺรหฺมจริย\-ปริโยสานนฺติ.  ฉฏฺฐํ.}

\csromanmatikaanchor{mtk.2}
\csromanmatikaanchor{mtk.310}
\csromantocmarkref{subsection}{1-2. กุสลสุตฺตานิ}{mtk.2}
\bookmark[dest={mtk.2},level=2]{1-2. กุสลสุตฺตานิ}
\csromanheader{6}{7. ทุติย\-อญฺญตรภิกฺขุ\-สุตฺต}
\proseitem{7}{สาวตฺถิ\-นิทานํ. อถ โข อญฺญตโร ภิกฺขุ เยน ภควา เตนุปสงฺกมิ ฯเปฯ เอกมนฺตํ นิสินฺโน โข โส ภิกฺขุ ภควนฺตํ เอตทโวจ\csromandash{}}

\prose{“ราควินโย โทสวินโย โมหวินโย”ติ ภนฺเต วุจฺจติ, กิสฺส นุ โข เอตํ ภนฺเต อธิวจนํ “ราควินโย โทสวินโย โมหวินโย”ติ. นิพฺพานธาตุยา โข เอตํ ภิกฺขุ อธิวจนํ ราควินโย โทสวินโย โมหวินโยติ, อาสวานํ ขโย เตน วุจฺจตีติ.}

\prose{เอวํ วุตฺเต โส ภิกฺขุ ภควนฺตํ เอตทโวจ “อมตํ อมตนฺ”ติ ภนฺเต วุจฺจติ, กตมํ นุ โข ภนฺเต อมตํ, กตโม อมตคามิ\-มคฺโคติ. โย โข ภิกฺขุ ราคกฺขโย โทสกฺขโย โมหกฺขโย, อิทํ วุจฺจติ อมตํ. อยเมว อริโย อฏฺฐงฺคิโก มคฺโค อมตคามิมคฺโค. เสยฺยถิทํ, สมฺมาทิฏฺฐิ ฯเปฯ สมฺมาสมาธีติ.  สตฺตมํ.}

\csromanheader{2}{8. \textbf{วิภงฺคสุตฺต}}
\proseitem{8}{สาวตฺถิ\-นิทานํ. อริยํ โว ภิกฺขเว อฏฺฐงฺคิกํ มคฺคํ เทเสสฺสามิ วิภชิสฺสามิ, ตํ สุณาถ สาธุกํ มนสิ กโรถ, ภาสิสฺสามีติ. “เอวํ ภนฺเต”ติ โข เต ภิกฺขู ภควโต ปจฺจสฺโสสุํ. ภควา เอตทโวจ\csromandash{}}

\prose{กตโม จ ภิกฺขเว อริโย อฏฺฐงฺคิโก มคฺโค, เสยฺยถิทํ, สมฺมาทิฏฺฐิ ฯเปฯ สมฺมาสมาธิ. กตมา จ ภิกฺขเว สมฺมาทิฏฺฐิ, ยํ โข ภิกฺขเว \csromanfolio{8}ทุกฺเข ญาณํ ทุกฺขสมุทเย ญาณํ ทุกฺขนิโรเธ ญาณํ ทุกฺขนิโรธ\-คามินิยา ปฏิปทาย ญาณํ. อยํ วุจฺจติ ภิกฺขเว สมฺมาทิฏฺฐิ.}

\prose{กตโม จ ภิกฺขเว สมฺมาสงฺกปฺโป, โย โข ภิกฺขเว เนกฺขมฺม\-สงฺกปฺโป อพฺยาปาทสงฺกปฺโป อวิหิํสาสงฺกปฺโป. อยํ วุจฺจติ ภิกฺขเว สมฺมาสงฺกปฺโป.}

\prose{กตมา จ ภิกฺขเว สมฺมาวาจา, ยา โข ภิกฺขเว มุสาวาทา เวรมณี ปิสุณาย วาจาย เวรมณี ผรุสาย วาจาย เวรมณี สมฺผปฺปลาปา เวรมณี. อยํ วุจฺจติ ภิกฺขเว สมฺมาวาจา.}

\prose{กตโม จ ภิกฺขเว สมฺมากมฺมนฺโต, ยา โข ภิกฺขเว ปาณาติปาตา เวรมณี อทินฺนาทานา เวรมณี อพฺรหฺมจริยา เวรมณี. อยํ วุจฺจติ ภิกฺขเว สมฺมากมฺมนฺโต.}

\prose{กตโม จ ภิกฺขเว สมฺมาอาชีโว, อิธ ภิกฺขเว อริยสาวโก มิจฺฉาอาชีวํ ปหาย สมฺมาอาชีเวน ชีวิตํ กปฺเปติ. อยํ วุจฺจติ ภิกฺขเว สมฺมาอาชีโว.}

\prose{กตโม จ ภิกฺขเว สมฺมาวายาโม, อิธ ภิกฺขเว ภิกฺขุ อนุปฺปนฺนานํ ปาปกานํ อกุสลานํ ธมฺมานํ อนุปฺปาทาย ฉนฺทํ ชเนติ วายมติ วีริยํ อารภติ, จิตฺตํ ปคฺคณฺหาติ ปทหติ, อุปฺปนฺนานํ ปาปกานํ อกุสลานํ ธมฺมานํ ปหานาย ฉนฺทํ ชเนติ ฯเปฯ อนุปฺปนฺนานํ กุสลานํ ธมฺมานํ อุปฺปาทาย ฉนฺทํ ชเนติ ฯเปฯ อุปฺปนฺนานํ กุสลานํ ธมฺมานํ ฐิติยา อสมฺโมสาย ภิยฺโยภาวาย เวปุลฺลาย ภาวนาย ปาริปูริยา ฉนฺทํ ชเนติ วายมติ วีริยํ อารภติ, จิตฺตํ ปคฺคณฺหาติ ปทหติ. อยํ วุจฺจติ ภิกฺขเว สมฺมาวายาโม.}

\prose{กตมา จ ภิกฺขเว สมฺมาสติ, อิธ ภิกฺขเว ภิกฺขุ กาเย กายานุปสฺสี วิหรติ อาตาปี สมฺปชาโน สติมา วิเนยฺย โลเก อภิชฺฌา\-โทมนสฺสํ, เวทนาสุ เวทนานุปสฺสี วิหรติ อาตาปี สมฺปชาโน สติมา วิเนยฺย โลเก อภิชฺฌา\-โทมนสฺสํ, จิตฺเต จิตฺตานุปสฺสี วิหรติ อาตาปี สมฺปชาโน สติมา วิเนยฺย โลเก อภิชฺฌา\-โทมนสฺสํ, ธมฺเมสุ ธมฺมานุปสฺสี วิหรติ อาตาปี สมฺปชาโน สติมา วิเนยฺย โลเก อภิชฺฌา\-โทมนสฺสํ. อยํ วุจฺจติ ภิกฺขเว สมฺมาสติ.}

\prose{\csromanfolio{9}กตโม จ ภิกฺขเว สมฺมาสมาธิ, อิธ ภิกฺขเว ภิกฺขุ วิวิจฺเจว กาเมหิ วิวิจฺจ อกุสเลหิ ธมฺเมหิ สวิตกฺกํ สวิจารํ วิเวกชํ ปีติสุขํ ปฐมํ ฌานํ อุปสมฺปชฺช วิหรติ, วิตกฺก\-วิจารานํ วูปสมา อชฺฌตฺตํ สมฺปสาทนํ เจตโส เอโกทิภาวํ อวิตกฺกํ อวิจารํ สมาธิชํ ปีติสุขํ ทุติยํ ฌานํ อุปสมฺปชฺช วิหรติ, ปีติยา จ วิราคา อุเปกฺขโก จ วิหรติ, สโต จ สมฺปชาโน สุขญฺจ กาเยน ปฏิสํเวเทติ. ยํ ตํ อริยา อาจิกฺขนฺติ “อุเปกฺขโก สติมา สุขวิหารี”ติ ตติยํ ฌานํ อุปสมฺปชฺช วิหรติ. สุขสฺส จ ปหานา ทุกฺขสฺส จ ปหานา ปุพฺเพว โสมนสฺส\-โทมนสฺสานํ อตฺถงฺคมา อทุกฺขมสุขํ อุเปกฺขา\-สติ\-ปาริสุทฺธิํ จตุตฺถํ ฌานํ อุปสมฺปชฺช วิหรติ. อยํ วุจฺจติ ภิกฺขเว สมฺมาสมาธีติ.  อฏฺฐมํ.}

\csromanheader{2}{9. \textbf{สูกสุตฺต}}
\proseitem{9}{สาวตฺถิ\-นิทานํ. เสยฺยถาปิ ภิกฺขเว สาลิสูกํ วา ยวสูกํ วา มิจฺฉา\-ปณิหิตํ หตฺเถน วา ปาเทน วา อกฺกนฺตํ หตฺถํ วา ปาทํ วา ภินฺทิสฺสติ\csromansharedfootnote{c56aa245e417025c}{เภจฺฉติ (ก)}, โลหิตํ วา อุปฺปาเทสฺสตีติ เนตํ ฐานํ วิชฺชติ. ตํ กิสฺส เหตุ, มิจฺฉา\-ปณิหิตตฺตา ภิกฺขเว สูกสฺส. เอวเมว โข ภิกฺขเว โส วต ภิกฺขุ มิจฺฉา\-ปณิหิตาย ทิฏฺฐิยา มิจฺฉา\-ปณิหิตาย มคฺคภาวนาย อวิชฺชํ ภินฺทิสฺสติ, วิชฺชํ อุปฺปาเทสฺสติ, นิพฺพานํ สจฺฉิกริสฺสตีติ เนตํ ฐานํ วิชฺชติ. ตํ กิสฺส เหตุ, มิจฺฉา\-ปณิหิตตฺตา ภิกฺขเว ทิฏฺฐิยา.}

\prose{เสยฺยถาปิ ภิกฺขเว สาลิสูกํ วา ยวสูกํ วา สมฺมาปณิหิตํ หตฺเถน วา ปาเทน วา อกฺกนฺตํ หตฺถํ วา ปาทํ วา ภินฺทิสฺสติ, โลหิตํ วา อุปฺปาเทสฺสตีติ ฐานเมตํ วิชฺชติ. ตํ กิสฺส เหตุ, สมฺมา\-ปณิหิตตฺตา ภิกฺขเว สูกสฺส. เอวเมว โข ภิกฺขเว โส วต ภิกฺขุ สมฺมา\-ปณิหิตาย ทิฏฺฐิยา สมฺมา\-ปณิหิตาย มคฺคภาวนาย อวิชฺชํ ภินฺทิสฺสติ, วิชฺชํ อุปฺปาเทสฺสติ, นิพฺพานํ สจฺฉิกริสฺสตีติ ฐานเมตํ วิชฺชติ. ตํ กิสฺส เหตุ, สมฺมา\-ปณิหิตตฺตา ภิกฺขเว ทิฏฺฐิยา.}

\prose{กถญฺจ ภิกฺขเว ภิกฺขุ สมฺมา\-ปณิหิตาย ทิฏฺฐิยา สมฺมา\-ปณิหิตาย มคฺคภาวนาย อวิชฺชํ ภินฺทติ, วิชฺชํ อุปฺปาเทติ, นิพฺพานํ สจฺฉิกโรตีติ. อิธ ภิกฺขเว ภิกฺขุ สมฺมาทิฏฺฐิํ ภาเวติ วิเวกนิสฺสิตํ วิราคนิสฺสิตํ \csromanfolio{10}นิโรธ\-นิสฺสิตํ โวสฺสคฺค\-ปริณามิํ ฯเปฯ สมฺมาสมาธิํ ภาเวติ วิเวกนิสฺสิตํ วิราคนิสฺสิตํ นิโรธ\-นิสฺสิตํ โวสฺสคฺค\-ปริณามิํ. เอวํ โข ภิกฺขเว ภิกฺขุ สมฺมา\-ปณิหิตาย ทิฏฺฐิยา สมฺมา\-ปณิหิตาย มคฺคภาวนาย อวิชฺชํ ภินฺทติ, วิชฺชํ อุปฺปาเทติ, นิพฺพานํ สจฺฉิกโรตีติ.  นวมํ.}

\csromanheader{2}{10. \textbf{นนฺทิยสุตฺต}}
\proseitem{9}{สาวตฺถิ\-นิทานํ. อถ โข นนฺทิโย ปริพฺพาชโก เยน ภควา เตนุปสงฺกมิ, อุปสงฺกมิตฺวา ภควตา สทฺธิํ สมฺโมทิ, สมฺโมทนียํ กถํ สารณียํ วีติสาเรตฺวา เอกมนฺตํ นิสีทิ, เอกมนฺตํ นิสินฺโน โข นนฺทิโย ปริพฺพาชโก ภควนฺตํ เอตทโวจ “กติ นุ โข โภ โคตม ธมฺมา ภาวิตา พหุลีกตา นิพฺพานงฺคมา โหนฺติ นิพฺพาน\-ปรายนา นิพฺพาน\-ปริโยสานา”ติ.}

\prosewithclosermidruled{อฏฺฐิเม โข นนฺทิย ธมฺมา ภาวิตา พหุลีกตา นิพฺพานงฺคมา โหนฺติ นิพฺพาน\-ปรายนา นิพฺพาน\-ปริโยสานา. กตเม อฏฺฐ, เสยฺยถิทํ, สมฺมาทิฏฺฐิ ฯเปฯ สมฺมาสมาธิ. อิเม โข นนฺทิย อฏฺฐ ธมฺมา ภาวิตา พหุลีกตา นิพฺพานงฺคมา โหนฺติ นิพฺพาน\-ปรายนา นิพฺพาน\-ปริโยสานาติ. เอวํ วุตฺเต นนฺทิโย ปริพฺพาชโก ภควนฺตํ เอตทโวจ “อภิกฺกนฺตํ โภ โคตม, อภิกฺกนฺตํ โภ โคตม ฯเปฯ อุปาสกํ มํ ภวํ โคตโม ธาเรตุ อชฺชตคฺเค ปาณุเปตํ สรณํ คตนฺ”ติ.  ทสมํ.}{อวิชฺชาวคฺโค ปฐโม.}
\csromancenter{\textbf{ตสฺสุทฺทานํ}}
\setlength{\csromangathapagewidth}{0pt}
\setlength{\csromangathawakwidth}{0pt}
\csromangathameasuregroup{อวิชฺชญฺจ อุปฑฺฒญฺจ, \\ กิมตฺถิโย จ เทฺว ภิกฺขู,}{\csromangathabat{อวิชฺชญฺจ อุปฑฺฒญฺจ,}{สาริปุตฺโต จ พฺราหฺมโณ.} \\ \csromangathabat{กิมตฺถิโย จ เทฺว ภิกฺขู,}{วิภงฺโค สูกนนฺทิยาติ.}}
\csromangathasetleft{อวิชฺชญฺจ อุปฑฺฒญฺจ, \\ กิมตฺถิโย จ เทฺว ภิกฺขู,}
\csromangathagroup{\csromangathastanza{\csromangathabat{อวิชฺชญฺจ อุปฑฺฒญฺจ,}{สาริปุตฺโต จ พฺราหฺมโณ.} \\ \csromangathabat{กิมตฺถิโย จ เทฺว ภิกฺขู,}{วิภงฺโค สูกนนฺทิยาติ.}}}

\chapterhead{2. \textbf{วิหารวคฺค}}
\csromanheader{2}{1. ปฐม\-วิหาร\-สุตฺต}
\proseitem{9}{สาวตฺถิ\-นิทานํ. อิจฺฉามหํ ภิกฺขเว อฑฺฒมาสํ ปฏิสลฺลิยิตุํ, นมฺหิ เกนจิ อุปสงฺกมิตพฺโพ อญฺญตฺร เอเกน ปิณฺฑปาตนีหารเกนาติ. “เอวํ}

\prose{\csromanfolio{11}ภนฺเต”ติ โข เต นฺหิกฺขู ภควโต ปฏิสฺสุตฺวา นาสฺสุธ โกจิ ภควนฺตํ อุปสงฺกมติ อญฺญตฺร เอเกน ปิณฺฑปาต\-นีหารเกน.}

\prose{อถ โข ภควา ตสฺส อฑฺฒมาสสฺส อจฺจเยน ปฏิสลฺลานา วุฏฺฐิโต ภิกฺขู อามนฺเตสิ “เยน สฺวาหํ ภิกฺขเว วิหาเรน ปฐมา\-ภิสมฺพุทฺโธ วิหรามิ, ตสฺส ปเทเสน วิหาสิํ, โส เอวํ ปชานามิ “มิจฺฉาทิฏฺฐิ\-ปจฺจยาปิ เวทยิตํ, สมฺมาทิฏฺฐิ\-ปจฺจยาปิ เวทยิตํ ฯเปฯ มิจฺฉาสมาธิ\-ปจฺจยาปิ เวทยิตํ, สมฺมาสมาธิ\-ปจฺจยาปิ เวทยิตํ, ฉนฺทปจฺจยาปิ เวทยิตํ, วิตกฺกปจฺจยาปิ เวทยิตํ, สญฺญาปจฺจยาปิ เวทยิตํ. ฉนฺโท จ อวูปสนฺโต โหติ, วิตกฺโก จ อวูปสนฺโต โหติ, สญฺญา จ อวูปสนฺตา โหติ, ตปฺปจฺจยาปิ เวทยิตํ. ฉนฺโท จ วูปสนฺโต โหติ, วิตกฺโก จ วูปสนฺโต โหติ, สญฺญา จ วูปสนฺตา โหติ, ตปฺปจฺจายาปิ เวทยิตํ. อปฺปตฺตสฺส ปตฺติยา อตฺถิ อายามํ\csromansharedfootnote{62d44272e3004926}{วายามํ (สี, สฺยา)}, ตสฺมิมฺปิ ฐาเน อนุปฺปตฺเต ตปฺปจฺจยาปิ เวทยิตนฺ’ติ”.  ปฐมํ.}

\csromanheader{2}{2. ทุติย\-วิหาร\-สุตฺต}
\proseitem{12}{สาวตฺถิ\-นิทานํ. อิจฺฉามหํ ภิกฺขเว เตมาสํ ปฏิสลฺลิยิตุํ, นมฺหิ เกนจิ อุปสงฺกมิตพฺโพ อญฺญตฺร เอเกน ปิณฺฑปาต\-นีหารเกนาติ. “เอวํ ภนฺเต”ติ โข เต ภิกฺขู ภควโต ปฏิสฺสุตฺวา นาสฺสุธ โกจิ ภควนฺตํ อุปสงฺกมติ อญฺญตฺร เอเกน ปิณฺฑปาต\-นีหารเกน.}

\prose{อถ โข ภควา ตสฺส เตมาสสฺส อจฺจเยน ปฏิสลฺลานา วุฏฺฐิโต ภิกฺขู อามนฺเตสิ “เยน สฺวาหํ ภิกฺขเว วิหาเรน ปฐมา\-ภิสมฺพุทฺโธ วิหรามิ, ตสฺส ปเทเสน วิหาสิํ. โส เอวํ ปชานามิ ‘มิจฺฉาทิฏฺฐิ\-ปจฺจยาปิ เวทยิตํ, มิจฺฉาทิฏฺฐิวูปสมปจฺจยาปิ เวทยิตํ, สมฺมาทิฏฺฐิ\-ปจฺจยาปิ เวทยิตํ, สมฺมาทิฏฺฐิ\-วูปสม\-ปจฺจยาปิ เวทยิตํ ฯเปฯ มิจฺฉาสมาธิ\-ปจฺจยาปิ เวทยิตํ, มิจฺฉาสมาธิ\-วูปสม\-ปจฺจยาปิ เวทยิตํ, สมฺมาสมาธิ\-ปจฺจยาปิ เวทยิตํ, สมฺมาสมาธิ\-วูปสม\-ปจฺจยาปิ เวทยิตํ, ฉนฺทปจฺจยาปิ เวทยิตํ, ฉนฺท\-วูปสม\-ปจฺจยาปิ เวทยิตํ, วิตกฺกปจฺจยาปิ เวทยิตํ, วิตกฺก\-วูปสม\-ปจฺจยาปิ เวทยิตํ, สญฺญาปจฺจยาปิ เวทยิตํ, สญฺญา\-วูปสม\-ปจฺจยาปิ เวทยิตํ. ฉนฺโท จ อวูปสนฺโต โหติ, วิตกฺโก จ อวูปสนฺโต โหติ, สญฺญา จ อวูปสนฺตา โหติ, ตปฺปจฺจยาปิ เวทยิตํ. ฉนฺโท จ วูปสนฺโต \csromanfolio{12}โหติ, วิตกฺโก จ วูปสนฺโต โหติ, สญฺญา จ วูปสนฺตา โหติ, ตปฺปจฺจยาปิ เวทยิตํ. อปฺปตฺตสฺส ปตฺติยา อตฺถิ อายามํ\csromansharedfootnote{62d44272e3004926}{วายามํ (สี, สฺยา)}, ตสฺมิมฺปิ ฐาเน อนุปฺปตฺเต ตปฺปจฺจยาปิ เวทยิตนฺ’ติ”.  ทุติยํ.}

\csromanheader{2}{3. \textbf{เสกฺขสุตฺต}}
\proseitem{13}{สาวตฺถิ\-นิทานํ. อถ โข อญฺญตโร ภิกฺขุ เยน ภควา เตนุปสงฺกมิ ฯเปฯ เอกมนฺตํ นิสินฺโน โข โส ภิกฺขุ ภควนฺตํ เอตทโวจ\csromandash{}“เสกฺโข เสกฺโข”ติ ภนฺเต วุจฺจติ. กิตฺตาวตา นุ โข ภนฺเต เสกฺโข โหตีติ. อิธ ภิกฺขุ เสกฺขาย สมฺมาทิฏฺฐิยา สมนฺนาคโต โหติ ฯเปฯ เสกฺเขน สมฺมาสมาธินา สมนฺนาคโต โหติ. เอตฺตาวตา โข ภิกฺขุ เสกฺโข โหตีติ.  ตติยํ.}

\prose{4. ปฐมอุปฺปาทสุตฺต}

\proseitem{14}{สาวตฺถิ\-นิทานํ. อฏฺฐิเม ภิกฺขเว ธมฺมา ภาวิตา พหุลีกตา อนุปฺปนฺนา อุปฺปชฺชนฺติ นาญฺญตฺร ตถาคตสฺส ปาตุภาวา อรหโต สมฺมา\-สมฺพุทฺธสฺส. กตเม อฏฺฐ, เสยฺยถิทํ, สมฺมาทิฏฺฐิ ฯเปฯ สมฺมาสมาธิ. อิเม โข ภิกฺขเว อฏฺฐ ธมฺมา ภาวิตา พหุลีกตา อนุปฺปนฺนา อุปฺปชฺชนฺติ นาญฺญตฺร ตถาคตสฺส ปาตุภาวา อรหโต สมฺมา\-สมฺพุทฺธสฺสาติ.  จตุตฺถํ.}

\prose{5. ทุติยอุปฺปาทสุตฺต}

\proseitem{15}{สาวตฺถิ\-นิทานํ. อฏฺฐิเม ภิกฺขเว ธมฺมา ภาวิตา พหุลีกตา อนุปฺปนฺนา อุปฺปชฺชนฺติ นาญฺญตฺร สุคตวินยา. กตเม อฏฺฐ, เสยฺยถิทํ, สมฺมาทิฏฺฐิ ฯเปฯ สมฺมาสมาธิ. อิเม โข ภิกฺขเว อฏฺฐ ธมฺมา ภาวิตา พหุลีกตา อนุปฺปนฺนา อุปฺปชฺชนฺติ นาญฺญตฺร สุคตวินยาติ.  ปญฺจมํ.}

\csromanheader{2}{6. ปฐม\-ปริสุทฺธ\-สุตฺต}
\proseitem{16}{สาวตฺถิ\-นิทานํ. อฏฺฐิเม ภิกฺขเว ธมฺมา ปริสุทฺธา ปริโยทาตา อนงฺคณา วิคตู\-ปกฺกิเลสา อนุปฺปนฺนา อุปฺปชฺชนฺติ นาญฺญตฺร ตถาคตสฺส ปาตุภาวา อรหโต สมฺมา\-สมฺพุทฺธสฺส. กตเม อฏฺฐ, เสยฺยถิทํ, สมฺมาทิฏฺฐิ ฯเปฯ สมฺมาสมาธิ. อิเม โข ภิกฺขเว อฏฺฐ ธมฺมา ปริสุทฺธา}

\prose{\csromanfolio{13}ปริโยทาตา อนงฺคณา วิคตู\-ปกฺกิเลสา อนุปฺปนฺนา อุปฺปชฺชนฺติ นาญฺญตฺร ตถาคตสฺส ปาตุภาวา อรหโต สมฺมาสมฺพุทฺธสฺสาติ.  ฉฏฺฐํ.}

\csromanheader{2}{7. ทุติย\-ปริสุทฺธ\-สุตฺต}
\proseitem{17}{สาวตฺถิ\-นิทานํ. อฏฺฐิเม ภิกฺขเว ธมฺมา ปริสุทฺธา ปริโยทาตา อนงฺคณา วิคตู\-ปกฺกิเลสา อนุปฺปนฺนา อุปฺปชฺชนฺติ นาญฺญตฺร สุคตวินยา. กตเม อฏฺฐ, เสยฺยถิทํ, สมฺมาทิฏฺฐิ ฯเปฯ สมฺมาสมาธิ. อิเม โข ภิกฺขเว อฏฺฐ ธมฺมา ปริสุทฺธา ปริโยทาตา อนงฺคณา วิคตู\-ปกฺกิเลสา อนุปฺปนฺนา อุปฺปชฺชนฺติ นาญฺญตฺร สุคตวินยาติ.  สตฺตมํ.}

\csromanheader{2}{8. ปฐม\-กุกฺกุฏาราม\-สุตฺต}
\proseitem{18}{เอวํ เม สุตํ\csromandash{}เอกํ สมยํ อายสฺมา จ อานนฺโท อายสฺมา จ ภทฺโท ปาฏลิปุตฺเต วิหรนฺติ กุกฺกุฏาราเม. อถ โข อายสฺมา ภทฺโท สายนฺหสมยํ ปฏิสลฺลานา วุฏฺฐิโต เยนายสฺมา อานนฺโท เตนุปสงฺกมิ, อุปสงฺกมิตฺวา อายสฺมตา อานนฺเทน สทฺธิํ สมฺโมทิ, สมฺโมทนียํ กถํ สารณียํ วีติสาเรตฺวา เอกมนฺตํ นิสีทิ, เอกมนฺตํ นิสินฺโน โข อายสฺมา ภทฺโท อายสฺมนฺตํ อานนฺทํ เอตทโวจ\csromandash{}}

\prose{“อพฺรหฺมจริยํ อพฺรหฺมจริยนฺ”ติ อาวุโส อานนฺท วุจฺจติ, กตมํ นุ โข อาวุโส อพฺรหฺมจริยนฺติ. สาธุ สาธุ อาวุโส ภทฺท, ภทฺทโก โข เต อาวุโส ภทฺท อุมฺมงฺโค, ภทฺทกํ ปฏิภานํ, กลฺยาณี ปริปุจฺฉา, เอวํ หิ ตฺวํ อาวุโส ภทฺท ปุจฺฉสิ “อพฺรหฺมจริยํ อพฺรหฺมจริยนฺติ อาวุโส อานนฺท วุจฺจติ, กตมํ นุ โข อาวุโส อพฺรหฺมจริยนฺ”ติ. เอวมาวุโสติ. อยเมว โข อาวุโส อฏฺฐงฺคิโก มิจฺฉามคฺโค อพฺรหฺมจริยํ. เสยฺยถิทํ, มิจฺฉาทิฏฺฐิ ฯเปฯ มิจฺฉาสมาธีติ.  อฏฺฐมํ.}

\csromanheader{2}{9. ทุติย\-กุกฺกุฏาราม\-สุตฺต}
\proseitem{19}{ปาฏลิปุตฺต\-นิทานํ. “พฺรหฺมจริยํ พฺรหฺมจริยนฺ”ติ อาวุโส อานนฺท วุจฺจติ, กตมํ นุ โข อาวุโส พฺรหฺมจริยํ, กตมํ พฺรหฺมจริย\-ปริโยสานนฺติ. สาธุ สาธุ อาวุโส ภทฺท, ภทฺทโก โข เต อาวุโส ภทฺท อุมฺมงฺโค, ภทฺทกํ ปฏิภานํ, กลฺยาณี ปริปุจฺฉา, เอวํ หิ ตฺวํ อาวุโส ภทฺท ปุจฺฉสิ “พฺรหฺมจริยํ พฺรหฺม\-จริยนฺติ อาวุโส อานนฺท วุจฺจติ, กตมํ นุ โข \csromanfolio{14}อาวุโส พฺรหฺมจริยํ, กตมํ พฺรหฺมจริยปริโยสานนฺ”ติ. เอวมาวุโสติ. อยเมว โข อาวุโส อริโย อฏฺฐงฺคิโก มคฺโค พฺรหฺมจริยํ. เสยฺยถิทํ, สมฺมาทิฏฺฐิ ฯเปฯ สมฺมาสมาธิ. โย โข อาวุโส ราคกฺขโย โทสกฺขโย โมหกฺขโย, อิทํ พฺรหฺมจริย\-ปริโยสานนฺติ.  นวมํ.}

\csromanmatikaanchor{mtk.3}
\csromanmatikaanchor{mtk.312}
\csromantocmarkref{subsection}{3-7. อุปกฺกิเลสสุตฺตาทีนิ}{mtk.3}
\bookmark[dest={mtk.3},level=2]{3-7. อุปกฺกิเลสสุตฺตาทีนิ}
\csromanheader{6}{10. ตติย\-กุกฺกุฏาราม\-สุตฺต}
\proseitemwithclosermidruled{20}{ปาฏลิปุตฺต\-นิทานํ. “พฺรหฺมจริยํ พฺรหฺมจริยนฺ”ติ อาวุโส อานนฺท วุจฺจติ, กตมํ นุ โข อาวุโส พฺรหฺมจริยํ, กตโม พฺรหฺมจารี, กตมํ พฺรหฺมจริย\-ปริโยสานนฺติ. สาธุ สาธุ อาวุโส ภทฺท, ภทฺทโก โข เต อาวุโส ภทฺท อุมฺมงฺโค, ภทฺทกํ ปฏิภานํ, กลฺยาณี ปริปุจฺฉา, เอวํ หิ ตฺวํ อาวุโส ภทฺท ปุจฺฉสิ “พฺรหฺมจริยํ พฺรหฺม\-จริยนฺติ อาวุโส อานนฺท วุจฺจติ, กตมํ นุ โข อาวุโส พฺรหฺมจริยํ, กตโม พฺรหฺมจารี, กตมํ พฺรหฺมจริยปริโยสานนฺ”ติ. เอวมาวุโสติ. อยเมว โข อาวุโส อริโย อฏฺฐงฺคิโก มคฺโค พฺรหฺมจริยํ. เสยฺยถิทํ, สมฺมาทิฏฺฐิ ฯเปฯ สมฺมาสมาธิ. โย โข อาวุโส อิมินา อริเยน อฏฺฐงฺคิเกน มคฺเคน สมนฺนาคโต, อยํ วุจฺจติ พฺรหฺมจารี. โย โข อาวุโส ราคกฺขโย โทสกฺขโย โมหกฺขโย, อิทํ พฺรหฺมจริย\-ปริโยสานนฺติ.  ทสมํ. ตีณิ สุตฺตนฺตานิ เอกนิทานานิ.}{วิหารวคฺโค ทุติโย.}
\csromancenter{\textbf{ตสฺสุทฺทานํ}}
\setlength{\csromangathapagewidth}{0pt}
\setlength{\csromangathawakwidth}{0pt}
\csromangathameasuregroup{เทฺว วิหารา จ เสกฺโข จ, \\ ปริสุทฺเธน เทฺว วุตฺตา,}{\csromangathabat{เทฺว วิหารา จ เสกฺโข จ,}{อุปฺปาทา อปเร ทุเว.} \\ \csromangathabat{ปริสุทฺเธน เทฺว วุตฺตา,}{กุกฺกุฏาราเมน ตโยติ.}}
\csromangathasetleft{เทฺว วิหารา จ เสกฺโข จ, \\ ปริสุทฺเธน เทฺว วุตฺตา,}
\csromangathagroup{\csromangathastanza{\csromangathabat{เทฺว วิหารา จ เสกฺโข จ,}{อุปฺปาทา อปเร ทุเว.} \\ \csromangathabat{ปริสุทฺเธน เทฺว วุตฺตา,}{กุกฺกุฏาราเมน ตโยติ.}}}

\chapterhead{3. \textbf{มิจฺฉตฺตวคฺค}}
\csromanheader{2}{1. \textbf{มิจฺฉตฺตสุตฺต}}
\proseitem{21}{สาวตฺถิ\-นิทานํ. มิจฺฉตฺตญฺจ โว ภิกฺขเว เทเสสฺสามิ สมฺมตฺตญฺจ, ตํ สุณาถ. กตมญฺจ ภิกฺขเว มิจฺฉตฺตํ, เสยฺยถิทํ, มิจฺฉทิฏฺฐิ ฯเปฯ}

\prose{\csromanfolio{15}มิจฺฉาสมาธิ. อิทํ วุจฺจติ ภิกฺขเว มิจฺฉตฺตํ. กตมญฺจ ภิกฺขเว สมฺมตฺตํ, เสยฺยถิทํ, สมฺมาทิฏฺฐิ (ป) สมฺมาสมาธิ. อิทํ วุจฺจติ ภิกฺขเว สมฺมตฺตนฺติ.  ปฐมํ.}

\csromanheader{2}{2. อกุสลธมฺม\-สุตฺต}
\proseitem{22}{สาวตฺถิ\-นิทานํ. อกุสเล จ โข ภิกฺขเว ธมฺเม เทเสสฺสามิ กุสเล จ ธมฺเม, ตํ สุณาถ. กตเม จ ภิกฺขเว อกุสลา ธมฺมา, เสยฺยถิทํ, มิจฺฉาทิฏฺฐิ ฯเปฯ มิจฺฉาสมาธิ. อิเม วุจฺจนฺติ ภิกฺขเว อกุสลา ธมฺมา. กตเม จ ภิกฺขเว กุสลา ธมฺมา, เสยฺยถิทํ, สมฺมาทิฏฺฐิ ฯเปฯ สมฺมาสมาธิ. อิเม วุจฺจนฺติ ภิกฺขเว กุสลา ธมฺมาติ.  ทุติยํ.}

\csromanheader{2}{3. ปฐม\-ปฏิปทา\-สุตฺต}
\proseitem{23}{สาวตฺถิ\-นิทานํ. มิจฺฉาปฏิปทญฺจ โว ภิกฺขเว เทเสสฺสามิ สมฺมาปฏิปทญฺจ, ตํ สุณาถ. กตมา จ ภิกฺขเว มิจฺฉาปฏิปทา, เสยฺยถิทํ, มิจฺฉาทิฏฺฐิ ฯเปฯ มิจฺฉาสมาธิ. อยํ วุจฺจติ ภิกฺขเว มิจฺฉาปฏิปทา. กตมา จ ภิกฺขเว สมฺมาปฏิปทา, เสยฺยถิทํ, สมฺมาทิฏฺฐิ ฯเปฯ สมฺมาสมาธิ. อยํ วุจฺจติ ภิกฺขเว สมฺมาปฏิปทาติ.  ตติยํ.}

\csromanheader{2}{4. ทุติย\-ปฏิปทา\-สุตฺต}
\proseitem{24}{สาวตฺถิ\-นิทานํ. คิหิโน วาหํ ภิกฺขเว ปพฺพชิตสฺส วา มิจฺฉา\-ปฏิปทํ น วณฺเณมิ, คิหิ วา ภิกฺขเว ปพฺพชิโต วา มิจฺฉา\-ปฏิปนฺโน มิจฺฉาปฏิปตฺตา\-ธิกรณ\-เหตุ นาราธโก โหติ ญายํ ธมฺมํ กุสลํ. กตมา จ ภิกฺขเว มิจฺฉาปฏิปทา, เสยฺยถิทํ, มิจฺฉาทิฏฺฐิ (ป) มิจฺฉาสมาธิ. อยํ วุจฺจติ ภิกฺขเว มิจฺฉาปฏิปทา. คิหิโน วาหํ ภิกฺขเว ปพฺพชิตสฺส วา มิจฺฉา\-ปฏิปทํ น วณฺเณมิ, คิหิ วา ภิกฺขเว ปพฺพชิโต วา มิจฺฉา\-ปฏิปนฺโน มิจฺฉาปฏิปตฺตา\-ธิกรณ\-เหตุ นาราธโก โหติ ญายํ ธมฺมํ กุสลํ.}

\prose{คิหิโน วาหํ ภิกฺขเว ปพฺพชิตสฺส วา สมฺมาปฏิปทํ วณฺเณมิ, คิหิ วา ภิกฺขเว ปพฺพชิโต วา สมฺมาปฏิปนฺโน สมฺมาปฏิปตฺตา\-ธิกรณ\-เหตุ อาราธโก โหติ ญายํ ธมฺมํ กุสลํ. กตมา จ ภิกฺขเว สมฺมาปฏิปทา, เสยฺยถิทํ, สมฺมาทิฏฺฐิ (ป) สมฺมาสมาธิ. อยํ วุจฺจติ ภิกฺขเว สมฺมาปฏิปทา. คิหิโน วาหํ ภิกฺขเว ปพฺพชิตสฺส วา สมฺมาปฏิปทํ วณฺเณมิ, คิหิ วา \csromanfolio{16}ภิกฺขเว ปพฺพชิโต วา สมฺมาปฏิปนฺโน สมฺมาปฏิปตฺตา\-ธิกรณ\-เหตุ อาราธโก โหติ ญายํ ธมฺมํ กุสลนฺติ.  จตุตฺถํ.}

\prose{5. ปฐม\-อสปฺปุริส\-สุตฺต}

\proseitem{25}{สาวตฺถิ\-นิทานํ. อสปฺปุริสญฺจ โว ภิกฺขเว เทเสสฺสามิ สปฺปุริสญฺจ, ตํ สุณาถ. กตโม จ ภิกฺขเว อสปฺปุริโส, อิธ ภิกฺขเว เอกจฺโจ มิจฺฉาทิฏฺฐิโก โหติ มิจฺฉาสงฺกปฺโป มิจฺฉาวาโจ มิจฺฉากมฺมนฺโต มิจฺฉาอาชีโว มิจฺฉาวายาโม มิจฺฉาสติ มิจฺฉาสมาธิ. อยํ วุจฺจติ ภิกฺขเว อสปฺปุริโส.}

\prose{กตโม จ ภิกฺขเว สปฺปุริโส, อิธ ภิกฺขเว เอกจฺโจ สมฺมาทิฏฺฐิโก โหติ สมฺมาสงฺกปฺโป สมฺมาวาโจ สมฺมากมฺมนฺโต สมฺมาอาชีโว สมฺมาวายาโม สมฺมาสติ สมฺมาสมาธิ. อยํ วุจฺจติ ภิกฺขเว สปฺปุริโสติ.  ปญฺจมํ.}

\prose{6. ทุติย\-อสปฺปุริส\-สุตฺต}

\proseitem{26}{สาวตฺถิ\-นิทานํ. อสปฺปุริสญฺจ โว ภิกฺขเว เทเสสฺสามิ อสปฺปุริเสน อสปฺปุริสตรญฺจ, สปฺปุริสญฺจ โว ภิกฺขเว เทเสสฺสามิ สปฺปุริเสน สปฺปุริสตรญฺจ, ตํ สุณาถ. กตโม จ ภิกฺขเว อสปฺปุริโส, อิธ ภิกฺขเว เอกจฺโจ มิจฺฉาทิฏฺฐิโก โหติ ฯเปฯ มิจฺฉาสมาธิ. อยํ วุจฺจติ ภิกฺขเว อสปฺปุริโส.}

\prose{กตโม จ ภิกฺขเว อสปฺปุริเสน อสปฺปุริสตโร, อิธ ภิกฺขเว เอกจฺโจ มิจฺฉาทิฏฺฐิโก โหติ ฯเปฯ มิจฺฉาสมาธิ มิจฺฉาญาณี มิจฺฉาวิมุตฺติ. อยํ วุจฺจติ ภิกฺขเว อสปฺปุริเสน อสปฺปุริสตโร.}

\prose{กตโม จ ภิกฺขเว สปฺปุริโส, อิธ ภิกฺขเว เอกจฺโจ สมฺมาทิฏฺฐิโก โหติ ฯเปฯ สมฺมาสมาธิ. อยํ วุจฺจติ ภิกฺขเว สปฺปุริโส.}

\prose{กตโม จ ภิกฺขเว สปฺปุริเสน สปฺปุริสตโร, อิธ ภิกฺขเว เอกจฺโจ สมฺมาทิฏฺฐิโก โหติ ฯเปฯ สมฺมาสมาธิ สมฺมาญาณี สมฺมาวิมุตฺติ. อยํ วุจฺจติ ภิกฺขเว สปฺปุริเสน สปฺปุริสตโรติ.  ฉฏฺฐํ.}

\csromanheader{2}{7. \textbf{กุมฺภสุตฺต}}
\proseitem{27}{\csromanfolio{17}สาวตฺถิ\-นิทานํ. เสยฺยถาปิ ภิกฺขเว กุมฺโภ อนาธาโร สุปฺปวตฺติโย โหติ, สาธาโร ทุปฺปวตฺติโย โหติ. เอวเมว โข ภิกฺขเว จิตฺตํ อนาธารํ สุปฺปวตฺติยํ โหติ, สาธารํ ทุปฺปวตฺติยํ โหติ. โก จ ภิกฺขเว จิตฺตสฺส อาธาโร, อยเมว อริโย อฏฺฐงฺคิโก มคฺโค. เสยฺยถิทํ, สมฺมาทิฏฺฐิ ฯเปฯ สมฺมาสมาธิ. อยํ จิตฺตสฺส อาธาโร. เสยฺยถาปิ ภิกฺขเว กุมฺโภ อนาธาโร สุปฺปวตฺติโย โหติ, สาธาโร ทุปฺปวตฺติโย โหติ. เอวเมว โข ภิกฺขเว จิตฺตํ อนาธารํ สุปฺปวตฺติยํ โหติ, สาธารํ ทุปฺปวตฺติยํ โหตีติ.  สตฺตมํ.}

\csromanheader{2}{8. \textbf{สมาธิสุตฺต}}
\proseitem{28}{สาวตฺถิ\-นิทานํ. อริยํ โว ภิกฺขเว สมฺมาสมาธิํ เทเสสฺสามิ สอุปนิสํ สปริกฺขารํ, ตํ สุณาถ. กตโม จ ภิกฺขเว อริโย สมฺมาสมาธิ สอุปนิโส สปริกฺขาโร. เสยฺยถิทํ, สมฺมาทิฏฺฐิ ฯเปฯ สมฺมาสติ\csromansharedfootnote{987f3bced3e82d13}{สมฺมาสมาธิ (สี, สฺยา, กํ, ก)}. ยา โข ภิกฺขเว อิเมหิ สตฺตหงฺเคหิ จิตฺตสฺส เอกคฺคตา สปริกฺขารตา\csromansharedfootnote{26afc8e6d7cfb107}{สปริกฺขตา (สี, อิ)}. อยํ วุจฺจติ ภิกฺขเว อริโย สมฺมาสมาธิ “สอุปนิโส” อิติปิ “สปริกฺขาโร” อิติปีติ.  อฏฺฐมํ.}

\csromanheader{2}{9. \textbf{เวทนาสุตฺต}}
\proseitem{29}{สาวตฺถิ\-นิทานํ. ติสฺโส อิมา ภิกฺขเว เวทนา. กตมา ติสฺโส, สุขา เวทนา ทุกฺขา เวทนา อทุกฺขมสุขา เวทนา. อิมา โข ภิกฺขเว ติสฺโส เวทนา. อิมาสํ โข ภิกฺขเว ติสฺสนฺนํ เวทนานํ ปริญฺญาย อริโย อฏฺฐงฺคิโก มคฺโค ภาเวตพฺโพ. กตโม อริโย อฏฺฐงฺคิโก มคฺโค, เสยฺยถิทํ, สมฺมาทิฏฺฐิ ฯเปฯ สมฺมาสมาธิ. อิมาสํ โข ภิกฺขเว ติสฺสนฺนํ เวทนานํ ปริญฺญาย อริโย อฏฺฐงฺคิโก มคฺโค ภาเวตพฺโพติ.  นวมํ.}

\csromanheader{2}{10. \textbf{อุตฺติยสุตฺต}}
\proseitemwithclosermidruled{30}{สาวตฺถิ\-นิทานํ. อถ โข อายสฺมา อุตฺติโย เยน ภควา เตนุปสงฺกมิ ฯเปฯ เอกมนฺตํ นิสินฺโน โข อายสฺมา อุตฺติโย ภควนฺตํ เอตทโวจ “อิธ มยฺหํ ภนฺเต รโหคตสฺส ปฏิสลฺลีนสฺส เอวํ เจตโส \csromanfolio{18}ปริวิตกฺโก อุทปาทิ ‘ปญฺจ กามคุณา วุตฺตา ภควตา, กตเม นุ โข ปญฺจ กามคุณา วุตฺตา ภควตา’ติ”. สาธุ สาธุ อุตฺติย, ปญฺจิเม โข อุตฺติย กามคุณา วุตฺตา มยา. กตเม ปญฺจ, จกฺขุวิญฺเญยฺยา รูปา อิฏฺฐา กนฺตา มนาปา ปิยรูปา กามูปสํหิตา รชนียา. โสตวิญฺเญยฺยา สทฺทา ฯเปฯ. ฆานวิญฺเญยฺยา คนฺธา ฯเปฯ. ชิวฺหาวิญฺเญยฺยา รสา ฯเปฯ. กายวิญฺเญยฺยา โผฏฺฐพฺพา อิฏฺฐา กนฺตา มนาปา ปิยรูปา กามูปสํหิตา รชนียา. อิเม โข อุตฺติย ปญฺจ กามคุณา วุตฺตา มยา. อิเมสํ โข อุตฺติย ปญฺจนฺนํ กามคุณานํ ปหานาย อริโย อฏฺฐงฺคิโก มคฺโค ภาเวตพฺโพ. กตโม อริโย อฏฺฐงฺคิโก มคฺโค, เสยฺยถิทํ, สมฺมาทิฏฺฐิ ฯเปฯ สมฺมาสมาธิ. อิเมสํ โข อุตฺติย ปญฺจนฺนํ กามคุณานํ ปหานาย อยํ อริโย อฏฺฐงฺคิโก มคฺโค ภาเวตพฺโพติ.  ทสมํ.}{มิจฺฉตฺตวคฺโค ตติโย.}
\csromancenter{\textbf{ตสฺสุทฺทานํ}}
\setlength{\csromangathapagewidth}{0pt}
\setlength{\csromangathawakwidth}{0pt}
\csromangathameasuregroup{มิจฺฉตฺตํ อกุสลํ ธมฺมํ, \\ อสปฺปุริเสน เทฺว กุมฺโภ,}{\csromangathabat{มิจฺฉตฺตํ อกุสลํ ธมฺมํ,}{ทุเว ปฏิปทาปิ จ.} \\ \csromangathabat{อสปฺปุริเสน เทฺว กุมฺโภ,}{สมาธิ เวทนุตฺติเยนาติ.}}
\csromangathasetleft{มิจฺฉตฺตํ อกุสลํ ธมฺมํ, \\ อสปฺปุริเสน เทฺว กุมฺโภ,}
\csromangathagroup{\csromangathastanza{\csromangathabat{มิจฺฉตฺตํ อกุสลํ ธมฺมํ,}{ทุเว ปฏิปทาปิ จ.} \\ \csromangathabat{อสปฺปุริเสน เทฺว กุมฺโภ,}{สมาธิ เวทนุตฺติเยนาติ.}}}

\chapterhead{4. \textbf{ปฏิปตฺติวคฺค}}
\csromanheader{2}{1. ปฐม\-ปฏิปตฺติ\-สุตฺต}
\proseitem{31}{สาวตฺถิ\-นิทานํ. มิจฺฉา\-ปฏิปตฺติญฺจ โว ภิกฺขเว เทเสสฺสามิ สมฺมาปฏิปตฺติญฺจ, ตํ สุณาถ. กตมา จ ภิกฺขเว มิจฺฉา\-ปฏิปตฺติ, เสยฺยถิทํ, มิจฺฉาทิฏฺฐิ ฯเปฯ มิจฺฉาสมาธิ. อยํ วุจฺจติ ภิกฺขเว มิจฺฉา\-ปฏิปตฺติ. กตมา จ ภิกฺขเว สมฺมาปฏิปตฺติ, เสยฺยถิทํ, สมฺมาทิฏฺฐิ ฯเปฯ สมฺมาสมาธิ. อยํ วุจฺจติ ภิกฺขเว สมฺมาปฏิปตฺตีติ.  ปฐมํ.}

\csromanheader{2}{2. ทุติย\-ปฏิปตฺติ\-สุตฺต}
\proseitem{32}{สาวตฺถิ\-นิทานํ. มิจฺฉาปฏิปนฺนญฺจ โว ภิกฺขเว เทเสสฺสามิ สมฺมาปฏิปนฺนญฺจ, ตํ สุณาถ. กตโม จ ภิกฺขเว มิจฺฉา\-ปฏิปนฺโน, อิธ ภิกฺขเว เอกจฺโจ มิจฺฉาทิฏฺฐิโก โหติ ฯเปฯ มิจฺฉาสมาธิ. อยํ วุจฺจติ}

\prose{\csromanfolio{19}ภิกฺขเว มิจฺฉา\-ปฏิปนฺโน. กตโม จ ภิกฺขเว สมฺมาปฏิปนฺโน, อิธ ภิกฺขเว เอกจฺโจ สมฺมาทิฏฺฐิโก โหติ ฯเปฯ สมฺมาสมาธิ. อยํ วุจฺจติ ภิกฺขเว สมฺมา\-ปฏิปนฺโนติ.  ทุติยํ.}

\csromanheader{2}{3. \textbf{วิรทฺธสุตฺต}}
\proseitem{33}{สาวตฺถิ\-นิทานํ. เยสํ เกสญฺจิ ภิกฺขเว อริโย อฏฺฐงฺคิโก มคฺโค วิรทฺโธ, วิรทฺโธ เตสํ อริโย อฏฺฐงฺคิโก มคฺโค สมฺมา ทุกฺข\-กฺขย\-คามี. เยสํ เกสญฺจิ ภิกฺขเว อริโย อฏฺฐงฺคิโก มคฺโค อารทฺโธ, อารทฺโธ เตสํ อริโย อฏฺฐงฺคิโก มคฺโค สมฺมา ทุกฺข\-กฺขย\-คามี. กตโม จ ภิกฺขเว อริโย อฏฺฐงฺคิโก มคฺโค, เสยฺยถิทํ, สมฺมาทิฏฺฐิ ฯเปฯ สมฺมาสมาธิ. เยสํ เกสญฺจิ ภิกฺขเว อยํ อริโย อฏฺฐงฺคิโก มคฺโค วิรทฺโธ, วิรทฺโธ เตสํ อริโย อฏฺฐงฺคิโก มคฺโค สมฺมา ทุกฺข\-กฺขย\-คามี. เยสํ เกสญฺจิ ภิกฺขเว อยํ อริโย อฏฺฐงฺคิโก มคฺโค อารทฺโธ, อารทฺโธ เตสํ อริโย อฏฺฐงฺคิโก มคฺโค สมฺมา ทุกฺข\-กฺขย\-คามีติ.  ตติยํ.}

\csromanheader{2}{4. \textbf{ปารํคมสุตฺต}}
\proseitem{34}{สาวตฺถิ\-นิทานํ. อฏฺฐิเม ภิกฺขเว ธมฺมา ภาวิตา พหุลีกตา อปารา ปารํ คมนาย สํวตฺตนฺติ. กตเม อฏฺฐ? เสยฺยถิทํ, สมฺมาทิฏฺฐิ ฯเปฯ สมฺมาสมาธิ. อิเม โข ภิกฺขเว อฏฺฐ ธมฺมา ภาวิตา พหุลีกตา อปารา ปารํ คมนาย สํวตฺตนฺตีติ.}

\prose{อิทมโวจ ภควา, อิทํ วตฺวาน สุคโต อถาปรํ เอตทโวจ สตฺถา\csromandash{}}

\setlength{\csromangathapagewidth}{0pt}
\setlength{\csromangathawakwidth}{0pt}
\csromangathameasuregroup{“อปฺปกา เต มนุสฺเสสุ, \\ อถายํ อิตรา ปชา, \\ เย จ โข สมฺมทกฺขาเต, \\ เต ชนา ปารเมสฺสนฺติ, \\ กณฺหํ ธมฺมํ วิปฺปหาย, \\ โอกา อโนกมาคมฺม, \\ ตตฺราภิรติมิจฺเฉยฺย, \\ ปริโยทเปยฺย อตฺตานํ, \\ เยสํ สมฺโพธิยงฺเคสุ, \\ อาทานปฏินิสฺสคฺเค,}{\csromangathabat{“อปฺปกา เต มนุสฺเสสุ,}{เย ชนา ปารคามิโน.} \\ \csromangathabat{อถายํ อิตรา ปชา,}{ตีรเมวานุธาวติ.} \\ \csromangathabat{เย จ โข สมฺมทกฺขาเต,}{ธมฺเม ธมฺมานุวตฺติโน.} \\ \csromangathabat{เต ชนา ปารเมสฺสนฺติ,}{มจฺจุเธยฺยํ สุทุตฺตรํ.} \\ \csromangathabat{กณฺหํ ธมฺมํ วิปฺปหาย,}{สุกฺกํ ภาเวถ ปณฺฑิโต.} \\ \csromangathabat{โอกา อโนกมาคมฺม,}{วิเวเก ยตฺถ ทูรมํ.} \\ \csromangathabat{ตตฺราภิรติมิจฺเฉยฺย,}{หิตฺวา กาเม อกิญฺจโน.} \\ \csromangathabat{ปริโยทเปยฺย อตฺตานํ,}{จิตฺตเกฺลเสหิ ปณฺฑิโต.} \\ \csromangathabat{เยสํ สมฺโพธิยงฺเคสุ,}{สมฺมา จิตฺตํ สุภาวิตํ.} \\ \csromangathabat{อาทานปฏินิสฺสคฺเค,}{อนุปาทาย เย รตา.}}
\csromangathasetleft{“อปฺปกา เต มนุสฺเสสุ, \\ อถายํ อิตรา ปชา, \\ เย จ โข สมฺมทกฺขาเต, \\ เต ชนา ปารเมสฺสนฺติ, \\ กณฺหํ ธมฺมํ วิปฺปหาย, \\ โอกา อโนกมาคมฺม, \\ ตตฺราภิรติมิจฺเฉยฺย, \\ ปริโยทเปยฺย อตฺตานํ, \\ เยสํ สมฺโพธิยงฺเคสุ, \\ อาทานปฏินิสฺสคฺเค,}
\csromangathagroup{\csromangathastanza{\csromangathabat{“อปฺปกา เต มนุสฺเสสุ,}{เย ชนา ปารคามิโน.} \\ \csromangathabat{อถายํ อิตรา ปชา,}{ตีรเมวานุธาวติ.}}\csromangathastanza{\csromangathabat{เย จ โข สมฺมทกฺขาเต,}{ธมฺเม ธมฺมานุวตฺติโน.} \\ \csromangathabat{เต ชนา ปารเมสฺสนฺติ,}{มจฺจุเธยฺยํ สุทุตฺตรํ.}}\csromangathastanza{\csromangathabat{กณฺหํ ธมฺมํ วิปฺปหาย,}{สุกฺกํ ภาเวถ ปณฺฑิโต.} \\ \csromangathabat{โอกา อโนกมาคมฺม,}{วิเวเก ยตฺถ ทูรมํ.}}\csromangathastanza{\csromangathabat{ตตฺราภิรติมิจฺเฉยฺย,}{หิตฺวา กาเม อกิญฺจโน.} \\ \csromangathabat{ปริโยทเปยฺย อตฺตานํ,}{จิตฺตเกฺลเสหิ ปณฺฑิโต.}}\csromangathastanza{\csromanfolio{20}\csromangathabat{เยสํ สมฺโพธิยงฺเคสุ,}{สมฺมา จิตฺตํ สุภาวิตํ.} \\ \csromangathabat{อาทานปฏินิสฺสคฺเค,}{อนุปาทาย เย รตา.}}}

\prose{ขีณาสวา ชุติมนฺโต, เต โลเก ปรินิพฺพุตา”ติ. จตุตฺถํ.}

\csromanheader{2}{5. ปฐม\-สามญฺญ\-สุตฺต}
\proseitem{35}{สาวตฺถิ\-นิทานํ. สามญฺญญฺจ โว ภิกฺขเว เทเสสฺสามิ สามญฺญผลานิ จ, ตํ สุณาถ. กตมญฺจ ภิกฺขเว สามญฺญํ, อยเมว อริโย อฏฺฐงฺคิโก มคฺโค, เสยฺยถิทํ, สมฺมาทิฏฺฐิ ฯเปฯ สมฺมาสมาธิ. อิทํ วุจฺจติ ภิกฺขเว สามญฺญํ. กตมานิ จ ภิกฺขเว สามญฺญผลานิ, โสตาปตฺติผลํ สกทาคามิผลํ อนาคามิผลํ อรหตฺตผลํ. อิมานิ วุจฺจนฺติ ภิกฺขเว สามญฺญผลานีติ.  ปญฺจมํ.}

\csromanheader{2}{6. ทุติย\-สามญฺญ\-สุตฺต}
\proseitem{36}{สาวตฺถิ\-นิทานํ. สามญฺญญฺจ โว ภิกฺขเว เทเสสฺสามิ สามญฺญตฺถญฺจ, ตํ สุณาถ. กตมญฺจ โข ภิกฺขเว สามญฺญํ, อยเมว อริโย อฏฺฐงฺคิโก มคฺโค, เสยฺยถิทํ, สมฺมาทิฏฺฐิ ฯเปฯ สมฺมาสมาธิ. อิทํ วุจฺจติ ภิกฺขเว สามญฺญํ. กตโม จ ภิกฺขเว สามญฺญตฺโถ, โย โข ภิกฺขเว ราคกฺขโย โทสกฺขโย โมหกฺขโย. อยํ วุจฺจติ ภิกฺขเว สามญฺญตฺโถติ.  ฉฏฺฐํ.}

\csromanheader{2}{7. ปฐม\-พฺรหฺมญฺญ\-สุตฺต}
\proseitem{37}{สาวตฺถิ\-นิทานํ. พฺรหฺมญฺญญฺจ โว ภิกฺขเว เทเสสฺสามิ พฺรหฺมญฺญ\-ผลานิ จ, ตํ สุณาถ. กตมญฺจ โข ภิกฺขเว พฺรหฺมญฺญํ, อยเมว อริโย อฏฺฐงฺคิโก มคฺโค, เสยฺยถิทํ, สมฺมาทิฏฺฐิ ฯเปฯ สมฺมาสมาธิ. อิทํ วุจฺจติ ภิกฺขเว พฺรหฺมญฺญํ. กตมานิ จ ภิกฺขเว พฺรหฺมญฺญ\-ผลานิ, โสตาปตฺติผลํ สกทาคามิผลํ อนาคามิผลํ อรหตฺตผลํ. อิมานิ วุจฺจนฺติ ภิกฺขเว พฺรหฺมญฺญผลานีติ.  สตฺตมํ.}

\csromanheader{2}{8. ทุติย\-พฺรหฺมญฺญ\-สุตฺต}
\proseitem{38}{สาวตฺถิ\-นิทานํ. พฺรหฺมญฺญญฺจ โว ภิกฺขเว เทเสสฺสามิ พฺรหฺมญฺญตฺถญฺจ, ตํ สุณาถ. กตมญฺจ ภิกฺขเว พฺรหฺมญฺญํ, อยเมว อริโย อฏฺฐงฺคิโก มคฺโค, เสยฺยถิทํ, สมฺมาทิฏฺฐิ ฯเปฯ สมฺมาสมาธิ. อิทํ วุจฺจติ ภิกฺขเว พฺรหฺมญฺญํ.}

\prose{\csromanfolio{21}กตโม จ ภิกฺขเว พฺรหฺมญฺญตฺโถ, โย โข ภิกฺขเว ราคกฺขโย โทสกฺขโย โมหกฺขโย. อยํ วุจฺจติ ภิกฺขเว พฺรหฺมญฺญ\-ตฺโถติ.  อฏฺฐมํ.}

\csromanheader{2}{9. ปฐม\-พฺรหฺมจริย\-สุตฺต}
\proseitem{39}{สาวตฺถิ\-นิทานํ. พฺรหฺมจริยญฺจ โว ภิกฺขเว เทเสสฺสามิ พฺรหฺมจริย\-ผลานิ จ, ตํ สุณาถ. กตมญฺจ ภิกฺขเว พฺรหฺมจริยํ. อยเมว อริโย อฏฺฐงฺคิโก มคฺโค, เสยฺยถิทํ, สมฺมาทิฏฺฐิ ฯเปฯ สมฺมาสมาธิ. อิทํ วุจฺจติ ภิกฺขเว พฺรหฺมจริยํ. กตมานิ จ ภิกฺขเว พฺรหฺมจริย\-ผลานิ, โสตาปตฺติผลํ สกทาคามิผลํ อนาคามิผลํ อรหตฺตผลํ. อิมานิ วุจฺจนฺติ ภิกฺขเว พฺรหฺมจริยผลานีติ.  นวมํ.}

\csromanheader{2}{10. ทุติย\-พฺรหฺมจริย\-สุตฺต}
\proseitemwithclosermidruled{40}{สาวตฺถิ\-นิทานํ. พฺรหฺมจริยญฺจ โว ภิกฺขเว เทเสสฺสามิ พฺรหฺมจริยตฺถญฺจ, ตํ สุณาถ. กตมญฺจ ภิกฺขเว พฺรหฺมจริยํ, อยเมว อริโย อฏฺฐงฺคิโก มคฺโค, เสยฺยถิทํ, สมฺมาทิฏฺฐิ ฯเปฯ สมฺมาสมาธิ. อิทํ วุจฺจติ ภิกฺขเว พฺรหฺมจริยํ. กตโม จ ภิกฺขเว พฺรหฺมจริย\-ตฺโถ, โย โข ภิกฺขเว ราคกฺขโย โทสกฺขโย โมหกฺขโย. อยํ วุจฺจติ ภิกฺขเว พฺรหฺมจริย\-ตฺโถติ.  ทสมํ.}{ปฏิปตฺติวคฺโค จตุตฺโถ.}
\csromancenter{\textbf{ตสฺสุทฺทานํ}}
\setlength{\csromangathapagewidth}{0pt}
\setlength{\csromangathawakwidth}{0pt}
\csromangathameasuregroup{ปฏิปตฺติ ปฏิปนฺโน จ, \\ สามญฺเญน จ เทฺว วุตฺตา,}{\csromangathabat{ปฏิปตฺติ ปฏิปนฺโน จ,}{วิรทฺธญฺจ ปารํคมา.} \\ \csromangathabat{สามญฺเญน จ เทฺว วุตฺตา,}{พฺรหฺมญฺญา อปเร ทุเว.} \\ \textbf{พฺรหฺมจริเยน เทฺว วุตฺตา, วคฺโค เตน ปวุจฺจตีติ.}}
\csromangathameasurewak{\textbf{พฺรหฺมจริเยน เทฺว วุตฺตา, วคฺโค เตน ปวุจฺจตีติ.}}
\csromangathasetleft{ปฏิปตฺติ ปฏิปนฺโน จ, \\ สามญฺเญน จ เทฺว วุตฺตา,}
\csromangathagroup{\csromangathastanza{\csromangathabat{ปฏิปตฺติ ปฏิปนฺโน จ,}{วิรทฺธญฺจ ปารํคมา.} \\ \csromangathabat{สามญฺเญน จ เทฺว วุตฺตา,}{พฺรหฺมญฺญา อปเร ทุเว.}}\csromangathastanza{\csromangathawak{\textbf{พฺรหฺมจริเยน เทฺว วุตฺตา, วคฺโค เตน ปวุจฺจตีติ.}}}}

\chapterhead{5. อญฺญติตฺถิย\-เปยฺยาล\-วคฺค}
\csromanheader{2}{1. ราควิราค\-สุตฺต}
\proseitem{41}{สาวตฺถิ\-นิทานํ. สเจ โว ภิกฺขเว อญฺญติตฺถิยา ปริพฺพาชกา เอวํ ปุจฺเฉยฺยุํ “กิมตฺถิยํ อาวุโส สมเณ โคตเม พฺรหฺมจริยํ วุสฺสตี”ติ. \csromanfolio{22}เอวํ ปุฏฺฐา ตุเมฺห ภิกฺขเว เตสํ อญฺญ\-ติตฺถิยานํ ปริพฺพาชกานํ เอวํ พฺยากเรยฺยาถ “ราควิราคตฺถํ โข อาวุโส ภควติ พฺรหฺมจริยํ วุสฺสตี”ติ. สเจ ปน โว ภิกฺขเว อญฺญติตฺถิยา ปริพฺพาชกา เอวํ ปุจฺเฉยฺยุํ “อตฺถิ ปนาวุโส มคฺโค อตฺถิ ปฏิปทา ราควิราคายา”ติ. เอวํ ปุฏฺฐา ตุเมฺห ภิกฺขเว เตสํ อญฺญ\-ติตฺถิยานํ ปริพฺพาชกานํ เอวํ พฺยากเรยฺยาถ “อตฺถิ โข อาวุโส มคฺโค อตฺถิ ปฏิปทา ราควิราคายา”ติ. กตโม จ ภิกฺขเว มคฺโค กตมา จ ปฏิปทา ราควิราคาย. อยเมว อริโย อฏฺฐงฺคิโก มคฺโค, เสยฺยถิทํ, สมฺมาทิฏฺฐิ ฯเปฯ สมฺมาสมาธิ. อยํ ภิกฺขเว มคฺโค อยํ ปฏิปทา ราควิราคายาติ. เอวํ ปุฏฺฐา ตุเมฺห ภิกฺขเว เตสํ อญฺญ\-ติตฺถิยานํ ปริพฺพาชกานํ เอวํ พฺยากเรยฺยาถาติ.  ปฐมํ.}

\csromanheader{2}{2-7. สํโยชนปฺปหานา\-ทิ\-สุตฺต\-ฉกฺก}
\proseitem{42-47}{สเจ โว ภิกฺขเว อญฺญติตฺถิยา ปริพฺพาชกา เอวํ ปุจฺเฉยฺยุํ “กิมตฺถิยํ อาวุโส สมเณ โคตเม พฺรหฺมจริยํ วุสฺสตี”ติ. เอวํ ปุฏฺฐา ตุเมฺห ภิกฺขเว เตสํ อญฺญ\-ติตฺถิยานํ ปริพฺพาชกานํ เอวํ พฺยากเรยฺยาถ “สํโยชนปฺปหานตฺถํ โข อาวุโส ภควติ พฺรหฺมจริยํ วุสฺสตี”ติ ฯเปฯ “อนุสยสมุคฺฆาตนตฺถํ โข อาวุโส ภควติ พฺรหฺมจริยํ วุสฺสตี”ติ ฯเปฯ “อทฺธานปริญฺญตฺถํ โข อาวุโส ภควติ พฺรหฺมจริยํ วุสฺสตี”ติ ฯเปฯ “อาสวานํ ขยตฺถํ โข อาวุโส ภควติ พฺรหฺมจริยํ วุสฺสตี”ติ ฯเปฯ “วิชฺชาวิมุตฺติผล- สจฺฉิกิริยตฺถํ โข อาวุโส ภควติ พฺรหฺมจริยํ วุสฺสตี”ติ ฯเปฯ “ญาณทสฺสนตฺถํ โข อาวุโส ภควติ พฺรหฺมจริยํ วุสฺสตี”ติ ฯเปฯ.  สตฺตมํ.}

\csromanheader{2}{8. อนุปาทาปรินิพฺพาน\-สุตฺต}
\proseitem{48}{สาวตฺถิ\-นิทานํ. สเจ โว ภิกฺขเว อญฺญติตฺถิยา ปริพฺพาชกา เอวํ ปุจฺเฉยฺยุํ “กิมตฺถิยํ อาวุโส สมเณ โคตเม พฺรหฺมจริยํ วุสฺสตี”ติ. เอวํ ปุฏฺฐา ตุเมฺห ภิกฺขเว เตสํ อญฺญ\-ติตฺถิยานํ ปริพฺพาชกานํ เอวํ พฺยากเรยฺยาถ “อนุปาทา\-ปรินิพฺพานตฺถํ โข อาวุโส ภควติ พฺรหฺมจริยํ วุสฺสตี”ติ. สเจ ปน โว ภิกฺขเว อญฺญติตฺถิยา ปริพฺพาชกา เอวํ ปุจฺเฉยฺยุํ “อตฺถิ ปนาวุโส มคฺโค อตฺถิ ปฏิปทา อนุปาทาปรินิพฺพานายา”ติ. เอวํ ปุฏฺฐา ตุเมฺห ภิกฺขเว เตสํ อญฺญ\-ติตฺถิยานํ ปริพฺพาชกานํ เอวํ}

\prose{\csromanfolio{23}พฺยากเรยฺยาถ “อตฺถิ โข อาวุโส มคฺโค อตฺถิ ปฏิปทา อนุปาทา\-ปรินิพฺพานายา”ติ. กตโม จ ภิกฺขเว มคฺโค กตมา จ ปฏิปทา อนุปาทา\-ปรินิพฺพานาย, อยเมว อริโย อฏฺฐงฺคิโก มคฺโค, เสยฺยถิทํ, สมฺมาทิฏฺฐิ ฯเปฯ สมฺมาสมาธิ. อยํ ภิกฺขเว มคฺโค อยํ ปฏิปทา อนุปาทา\-ปรินิพฺพานายาติ. เอวํ ปุฏฺฐา ตุเมฺห ภิกฺขเว เตสํ อญฺญ\-ติตฺถิยานํ ปริพฺพาชกานํ เอวํ พฺยากเรยฺยาถาติ.  อฏฺฐมํ. อญฺญติตฺถิย\-เปยฺยาลํ.}

\vspace{\tipitakaparskip}
\csromancenter{\textbf{ตสฺสุทฺทานํ}}
\setlength{\csromangathapagewidth}{0pt}
\setlength{\csromangathawakwidth}{0pt}
\csromangathameasuregroup{วิราคสํโยชนํ อนุสยํ, \\ วิชฺชา วิมุตฺติ ญาณญฺจ,}{\csromangathabat{วิราคสํโยชนํ อนุสยํ,}{อทฺธานํ อาสวา ขยา.} \\ \csromangathabat{วิชฺชา วิมุตฺติ ญาณญฺจ,}{อนุปาทาย อฏฺฐมี.}}
\csromangathasetleft{วิราคสํโยชนํ อนุสยํ, \\ วิชฺชา วิมุตฺติ ญาณญฺจ,}
\csromangathagroup{\csromangathastanza{\csromangathabat{วิราคสํโยชนํ อนุสยํ,}{อทฺธานํ อาสวา ขยา.} \\ \csromangathabat{วิชฺชา วิมุตฺติ ญาณญฺจ,}{อนุปาทาย อฏฺฐมี.}}}

\chapterhead{6. สูริย\-เปยฺยาล\-วคฺค}
\csromanheader{2}{1. กลฺยาณมิตฺต\-สุตฺต}
\proseitem{49}{สาวตฺถิ\-นิทานํ. สูริยสฺส ภิกฺขเว อุทยโต เอตํ ปุพฺพงฺคมํ เอตํ ปุพฺพนิมิตฺตํ, ยทิทํ อรุณุคฺคํ. เอวเมว โข ภิกฺขเว ภิกฺขุโน อริยสฺส อฏฺฐงฺคิกสฺส มคฺคสฺส อุปฺปาทาย เอตํ ปุพฺพงฺคมํ เอตํ ปุพฺพนิมิตฺตํ, ยทิทํ กลฺยาณมิตฺตตา. กลฺยาณมิตฺตสฺเสตํ ภิกฺขเว ภิกฺขุโน ปาฏิกงฺขํ อริยํ อฏฺฐงฺคิกํ มคฺคํ ภาเวสฺสติ, อริยํ อฏฺฐงฺคิกํ มคฺคํ พหุลีกริสฺสติ. กถญฺจ ภิกฺขเว ภิกฺขุ กลฺยาณมิตฺโต อริยํ อฏฺฐงฺคิกํ มคฺคํ ภาเวติ, อริยํ อฏฺฐงฺคิกํ มคฺคํ พหุลีกโรติ, อิธ ภิกฺขเว ภิกฺขุ สมฺมาทิฏฺฐิํ ภาเวติ วิเวกนิสฺสิตํ วิราคนิสฺสิตํ นิโรธ\-นิสฺสิตํ โวสฺสคฺค\-ปริณามิํ ฯเปฯ สมฺมาสมาธิํ ภาเวติ วิเวกนิสฺสิตํ วิราคนิสฺสิตํ นิโรธ\-นิสฺสิตํ โวสฺสคฺค\-ปริณามิํ. เอวํ โข ภิกฺขเว ภิกฺขุ กลฺยาณมิตฺโต อริยํ อฏฺฐงฺคิกํ มคฺคํ ภาเวติ, อริยํ อฏฺฐงฺคิกํ มคฺคํ พหุลีกโรตีติ.  ปฐมํ.}

\csromanheader{2}{2-6. สีลสมฺปทา\-ทิ\-สุตฺต\-ปญฺจก}
\proseitem{50-54}{สูริยสฺส ภิกฺขเว อุทยโต เอตํ ปุพฺพงฺคมํ เอตํ ปุพฺพนิมิตฺตํ, ยทิทํ อรุณุคฺคํ. เอวเมว โข ภิกฺขเว ภิกฺขุโน อริยสฺส อฏฺฐงฺคิกสฺส มคฺคสฺส \csromanfolio{24}อุปฺปาทาย เอตํ ปุพฺพงฺคมํ เอตํ ปุพฺพนิมิตฺตํ, ยทิทํ สีลสมฺปทา. สีลสมฺปนฺนสฺเสตํ ภิกฺขเว ภิกฺขุโน ปาฏิกงฺขํ ฯเปฯ ยทิทํ ฉนฺทสมฺปทา ฯเปฯ ยทิทํ อตฺตสมฺปทา ฯเปฯ ยทิทํ ทิฏฺฐิสมฺปทา ฯเปฯ ยทิทํ อปฺปมาท\-สมฺปทา ฯเปฯ.  ฉฏฺฐํ.}

\csromanheader{2}{7. โยนิโสมนสิการ\-สมฺปทา\-สุตฺต}
\proseitem{55}{สูริยสฺส ภิกฺขเว อุทยโต เอตํ ปุพฺพงฺคมํ เอตํ ปุพฺพนิมิตฺตํ, ยทิทํ อรุณุคฺคํ. เอวเมว โข ภิกฺขเว ภิกฺขุโน อริยสฺส อฏฺฐงฺคิกสฺส มคฺคสฺส อุปฺปาทาย เอตํ ปุพฺพงฺคมํ เอตํ ปุพฺพนิมิตฺตํ, ยทิทํ โยนิโสมนสิการ\-สมฺปทา. โยนิโสมนสิการสมฺปนฺนสฺเสตํ ภิกฺขเว ภิกฺขุโน ปาฏิกงฺขํ อริยํ อฏฺฐงฺคิกํ มคฺคํ ภาเวสฺสติ, อริยํ อฏฺฐงฺคิกํ มคฺคํ พหุลีกริสฺสติ. กถญฺจ ภิกฺขเว ภิกฺขุ โยนิโสมนสิการ\-สมฺปนฺโน อริยํ อฏฺฐงฺคิกํ มคฺคํ ภาเวติ, อริยํ อฏฺฐงฺคิกํ มคฺคํ พหุลีกโรติ, อิธ ภิกฺขเว ภิกฺขุ สมฺมาทิฏฺฐิํ ภาเวติ วิเวกนิสฺสิตํ วิราคนิสฺสิตํ นิโรธ\-นิสฺสิตํ โวสฺสคฺค\-ปริณามิํ ฯเปฯ สมฺมาสมาธิํ ภาเวติ วิเวกนิสฺสิตํ วิราคนิสฺสิตํ นิโรธ\-นิสฺสิตํ โวสฺสคฺค\-ปริณามิํ. เอวํ โข ภิกฺขเว ภิกฺขุ โยนิโสมนสิการ\-สมฺปนฺโน อริยํ อฏฺฐงฺคิกํ มคฺคํ ภาเวติ, อริยํ อฏฺฐงฺคิกํ มคฺคํ พหุลีกโรตีติ.  สตฺตมํ.}

\csromanheader{2}{1. กลฺยาณมิตฺต\-สุตฺต}
\proseitem{56}{สูริยสฺส ภิกฺขเว อุทยโต เอตํ ปุพฺพงฺคมํ เอตํ ปุพฺพนิมิตฺตํ, ยทิทํ อรุณุคฺคํ. เอวเมว โข ภิกฺขเว ภิกฺขุโน อริยสฺส อฏฺฐงฺคิกสฺส มคฺคสฺส อุปฺปาทาย เอตํ ปุพฺพงฺคมํ เอตํ ปุพฺพนิมิตฺตํ, ยทิทํ กลฺยาณมิตฺตตา. กลฺยาณมิตฺตสฺเสตํ ภิกฺขเว ภิกฺขุโน ปาฏิกงฺขํ อริยํ อฏฺฐงฺคิกํ มคฺคํ ภาเวสฺสติ, อริยํ อฏฺฐงฺคิกํ มคฺคํ พหุลีกริสฺสติ. กถญฺจ ภิกฺขเว ภิกฺขุ กลฺยาณมิตฺโต อริยํ อฏฺฐงฺคิกํ มคฺคํ ภาเวติ, อริยํ อฏฺฐงฺคิกํ มคฺคํ พหุลีกโรติ, อิธ หฺหิกฺขเว ภิกฺขุ สมฺมาทิฏฺฐิํ ภาเวติ ราค\-วินย\-ปริโยสานํ โทส\-วินย\-ปริโยสานํ โมห\-วินย\-ปริโยสานํ ฯเปฯ สมฺมาสมาธิํ ภาเวติ ราค\-วินย\-ปริโยสานํ โทส\-วินย\-ปริโยสานํ โมห\-วินย\-ปริโยสานํ. เอวํ โข ภิกฺขเว หฺหิกฺขุ กลฺยาณมิตฺโต อริยํ อฏฺฐงฺคิกํ มคฺคํ ภาเวติ, อริยํ อฏฺฐงฺคิกํ มคฺคํ พหุลีกโรตีติ.  ปฐมํ.}

\csromanheader{2}{2-6. สีลสมฺปทา\-ทิ\-สุตฺต\-ปญฺจก}
\proseitem{57-61}{\csromanfolio{25}สูริยสฺส ภิกฺขเว อุทยโต เอตํ ปุพฺพงฺคมํ เอตํ ปุพฺพนิมิตฺตํ, ยทิทํ อรุณุคฺคํ. เอวเมว โข ภิกฺขเว ภิกฺขุโน อริยสฺส อฏฺฐงฺคิกสฺส มคฺคสฺส อุปฺปาทาย เอตํ ปุพฺพงฺคมํ เอตํ ปุพฺพนิมิตฺตํ, ยทิทํ สีลสมฺปทา ฯเปฯ ยทิทํ ฉนฺทสมฺปทา ฯเปฯ ยทิทํ อตฺตสมฺปทา ฯเปฯ ยทิทํ ทิฏฺฐิสมฺปทา ฯเปฯ ยทิทํ อปฺปมาท\-สมฺปทา ฯเปฯ.  ฉฏฺฐํ.}

\csromanheader{2}{7. โยนิโสมนสิการ\-สมฺปทา\-สุตฺต}
\proseitem{62}{ยทิทํ โยนิโสมนสิการ\-สมฺปทา. โยนิโสมนสิการสมฺปนฺนสฺเสตํ ภิกฺขเว ภิกฺขุโน ปาฏิกงฺขํ อริยํ อฏฺฐงฺคิกํ มคฺคํ ภาเวสฺสติ, อริยํ อฏฺฐงฺคิกํ มคฺคํ พหุลีกริสฺสติ. กถญฺจ ภิกฺขเว ภิกฺขุ โยนิโสมนสิการ\-สมฺปนฺโน อริยํ อฏฺฐงฺคิกํ มคฺคํ ภาเวติ, อริยํ อฏฺฐงฺคิกํ มคฺคํ พหุลีกโรติ, อิธ ภิกฺขเว ภิกฺขุ สมฺมาทิฏฺฐิํ ภาเวติ ราค\-วินย\-ปริโยสานํ โทส\-วินย\-ปริโยสานํ โมห\-วินย\-ปริโยสานํ ฯเปฯ สมฺมาสมาธิํ ภาเวติ ราค\-วินย\-ปริโยสานํ โทส\-วินย\-ปริโยสานํ โมห\-วินย\-ปริโยสานํ. เอวํ โข ภิกฺขเว ภิกฺขุ โยนิโสมนสิการ\-สมฺปนฺโน อริยํ อฏฺฐงฺคิกํ มคฺคํ ภาเวติ, อริยํ อฏฺฐงฺคิกํ มคฺคํ พหุลีกโรตีติ.  สตฺตมํ. สูริยเปยฺยาลํ.}

\vspace{\tipitakaparskip}
\csromancenter{\textbf{ตสฺสุทฺทานํ}}
\setlength{\csromangathapagewidth}{0pt}
\setlength{\csromangathawakwidth}{0pt}
\csromangathameasuregroup{กลฺยาณมิตฺตํ สีลญฺจ, \\ ทิฏฺฐิ จ อปฺปมาโท จ,}{\csromangathabat{กลฺยาณมิตฺตํ สีลญฺจ,}{ฉนฺโท จ อตฺตสมฺปทา.} \\ \csromangathabat{ทิฏฺฐิ จ อปฺปมาโท จ,}{โยนิโส ภวติ สตฺตมํ.}}
\csromangathasetleft{กลฺยาณมิตฺตํ สีลญฺจ, \\ ทิฏฺฐิ จ อปฺปมาโท จ,}
\csromangathagroup{\csromangathastanza{\csromangathabat{กลฺยาณมิตฺตํ สีลญฺจ,}{ฉนฺโท จ อตฺตสมฺปทา.} \\ \csromangathabat{ทิฏฺฐิ จ อปฺปมาโท จ,}{โยนิโส ภวติ สตฺตมํ.}}}

\chapterhead{7. เอกธมฺม\-เปยฺยาล\-วคฺค}
\csromanheader{2}{1. กลฺยาณมิตฺต\-สุตฺต}
\proseitem{63}{สาวตฺถิ\-นิทานํ. เอกธมฺโม ภิกฺขเว พหูปกาโร อริยสฺส อฏฺฐงฺคิกสฺส มคฺคสฺส อุปฺปาทาย. กตโม เอกธมฺโม, ยทิทํ กลฺยาณมิตฺตตา. \csromanfolio{26}กลฺยาณมิตฺตสฺเสตํ ภิกฺขเว ภิกฺขุโน ปาฏิกงฺขํ อริยํ อฏฺฐงฺคิกํ มคฺคํ ภาเวสฺสติ, อริยํ อฏฺฐงฺคิกํ มคฺคํ พหุลีกริสฺสติ. กถญฺจ ภิกฺขเว ภิกฺขุ กลฺยาณมิตฺโต อริยํ อฏฺฐงฺคิกํ มคฺคํ ภาเวติ, อริยํ อฏฺฐงฺคิกํ มคฺคํ พหุลีกโรติ, อิธ ภิกฺขเว ภิกฺขุ สมฺมาทิฏฺฐิํ ภาเวติ วิเวกนิสฺสิตํ วิราคนิสฺสิตํ นิโรธ\-นิสฺสิตํ โวสฺสคฺค\-ปริณามิํ ฯเปฯ สมฺมาสมาธิํ ภาเวติ วิเวกนิสฺสิตํ วิราคนิสฺสิตํ นิโรธ\-นิสฺสิตํ โวสฺสคฺค\-ปริณามิํ. เอวํ โข ภิกฺขเว ภิกฺขุ กลฺยาณมิตฺโต อริยํ อฏฺฐงฺคิกํ มคฺคํ ภาเวติ, อริยํ อฏฺฐงฺคิกํ มคฺคํ พหุลีกโรตีติ.  ปฐมํ.}

\csromanheader{2}{2-6. สีลสมฺปทา\-ทิ\-สุตฺต\-ปญฺจก}
\proseitem{64-68}{เอกธมฺโม ภิกฺขเว พหูปกาโร อริยสฺส อฏฺฐงฺคิกสฺส มคฺคสฺส อุปฺปาทาย. กตโม เอกธมฺโม, ยทิทํ สีลสมฺปทา ฯเปฯ ยทิทํ ฉนฺทสมฺปทา ฯเปฯ ยทิทํ อตฺตสมฺปทา ฯเปฯ ยทิทํ ทิฏฺฐิสมฺปทา ฯเปฯ ยทิทํ อปฺปมาท\-สมฺปทา ฯเปฯ.  ฉฏฺฐํ.}

\csromanheader{2}{7. โยนิโสมนสิการ\-สมฺปทา\-สุตฺต}
\proseitem{69}{ยทิทํ โยนิโสมนสิการ\-สมฺปทา. โยนิโสมนสิการสมฺปนฺนสฺเสตํ ภิกฺขเว ภิกฺขุโน ปาฏิกงฺขํ อริยํ อฏฺฐงฺคิกํ มคฺคํ ภาเวสฺสติ, อริยํ อฏฺฐงฺคิกํ มคฺคํ พหุลีกริสฺสติ. กถญฺจ ภิกฺขเว ภิกฺขุ โยนิโสมนสิการ\-สมฺปนฺโน อริยํ อฏฺฐงฺคิกํ มคฺคํ ภาเวติ, อริยํ อฏฺฐงฺคิกํ มคฺคํ พหุลีกโรติ, อิธ ภิกฺขเว ภิกฺขุ สมฺมาทิฏฺฐิํ ภาเวติ วิเวกนิสฺสิตํ วิราคนิสฺสิตํ นิโรธ\-นิสฺสิตํ โวสฺสคฺค\-ปริณามิํ ฯเปฯ สมฺมาสมาธิํ ภาเวติ วิเวกนิสฺสิตํ วิราคนิสฺสิตํ นิโรธ\-นิสฺสิตํ โวสฺสคฺค\-ปริณามิํ. เอวํ โข ภิกฺขเว ภิกฺขุ โยนิโสมนสิการ\-สมฺปนฺโน อริยํ อฏฺฐงฺคิกํ มคฺคํ ภาเวติ, อริยํ อฏฺฐงฺคิกํ มคฺคํ พหุลีกโรตีติ.  สตฺตมํ.}

\csromanheader{2}{1. กลฺยาณมิตฺต\-สุตฺต}
\proseitem{70}{สาวตฺถิ\-นิทานํ. เอกธมฺโม ภิกฺขเว พหูปกาโร อริยสฺส อฏฺฐงฺคิกสฺส มคฺคสฺส อุปฺปาทาย. กตโม เอกธมฺโม, ยทิทํ กลฺยาณมิตฺตตา.}

\prose{\csromanfolio{27}กลฺยาณมิตฺตสฺเสตํ ภิกฺขเว ภิกฺขุโน ปาฏิกงฺขํ อริยํ อฏฺฐงฺคิกํ มคฺคํ ภาเวสฺสติ, อริยํ อฏฺฐงฺคิกํ มคฺคํ พหุลีกริสฺสติ. กถญฺจ ภิกฺขเว ภิกฺขุ กลฺยาณมิตฺโต อริยํ อฏฺฐงฺคิกํ มคฺคํ ภาเวติ, อริยํ อฏฺฐงฺคิกํ มคฺคํ พหุลีกโรติ, อิธ ภิกฺขเว ภิกฺขุ สมฺมาทิฏฺฐิํ ภาเวติ ราค\-วินย\-ปริโยสานํ โทส\-วินย\-ปริโยสานํ โมห\-วินย\-ปริโยสานํ ฯเปฯ สมฺมาสมาธิํ ภาเวติ ราค\-วินย\-ปริโยสานํ โทส\-วินย\-ปริโยสานํ โมห\-วินย\-ปริโยสานํ. เอวํ โข ภิกฺขเว ภิกฺขุ กลฺยาณมิตฺโต อริยํ อฏฺฐงฺคิกํ มคฺคํ ภาเวติ, อริยํ อฏฺฐงฺคิกํ มคฺคํ พหุลีกโรตีติ.  ปฐมํ.}

\csromanheader{2}{2-6. สีลสมฺปทา\-ทิ\-สุตฺต\-ปญฺจก}
\proseitem{71-75}{สาวตฺถิ\-นิทานํ. เอกธมฺโม ภิกฺขเว พหูปกาโร อริยสฺส อฏฺฐงฺคิกสฺส มคฺคสฺส อุปฺปาทาย. กตโม เอกธมฺโม, ยทิทํ สีลสมฺปทา ฯเปฯ ยทิทํ ฉนฺทสมฺปทา ฯเปฯ ยทิทํ อตฺตสมฺปทา ฯเปฯ ยทิทํ ทิฏฺฐิสมฺปทา ฯเปฯ ยทิทํ อปฺปมาท\-สมฺปทา ฯเปฯ.  ฉฏฺฐํ.}

\csromanheader{2}{7. โยนิโสมนสิการ\-สมฺปทา\-สุตฺต}
\proseitem{76}{ยทิทํ โยนิโสมนสิการ\-สมฺปทา. โยนิโสมนสิการสมฺปนฺนสฺเสตํ ภิกฺขเว ภิกฺขุโน ปาฏิกงฺขํ อริยํ อฏฺฐงฺคิกํ มคฺคํ ภาเวสฺสติ, อริยํ อฏฺฐงฺคิกํ มคฺคํ พหุลีกริสฺสติ. กถญฺจ ภิกฺขเว ภิกฺขุ โยนิโสมนสิการ\-สมฺปนฺโน อริยํ อฏฺฐงฺคิกํ มคฺคํ ภาเวติ, อริยํ อฏฺฐงฺคิกํ มคฺคํ พหุลีกโรติ, อิธ ภิกฺขเว ภิกฺขุ สมฺมาทิฏฺฐิํ ภาเวติ ฯเปฯ สมฺมาสมาธิํ ภาเวติ ราค\-วินย\-ปริโยสานํ โทส\-วินย\-ปริโยสานํ โมห\-วินย\-ปริโยสานํ. เอวํ โข ภิกฺขเว ภิกฺขุ โยนิโสมนสิการ\-สมฺปนฺโน อริยํ อฏฺฐงฺคิกํ มคฺคํ ภาเวติ, อริยํ อฏฺฐงฺคิกํ มคฺคํ พหุลีกโรตีติ.  สตฺตมํ. เอกธมฺม\-เปยฺยาลํ.}

\vspace{\tipitakaparskip}
\csromancenter{\textbf{ตสฺสุทฺทานํ}}
\setlength{\csromangathapagewidth}{0pt}
\setlength{\csromangathawakwidth}{0pt}
\csromangathameasuregroup{กลฺยาณมิตฺตํ สีลญฺจ, \\ ทิฏฺฐิ จ อปฺปมาโท จ,}{\csromangathabat{กลฺยาณมิตฺตํ สีลญฺจ,}{ฉนฺโท จ อตฺตสมฺปทา.} \\ \csromangathabat{ทิฏฺฐิ จ อปฺปมาโท จ,}{โยนิโส ภวติ สตฺตมํ.}}
\csromangathasetleft{กลฺยาณมิตฺตํ สีลญฺจ, \\ ทิฏฺฐิ จ อปฺปมาโท จ,}
\csromangathagroup{\csromangathastanza{\csromangathabat{กลฺยาณมิตฺตํ สีลญฺจ,}{ฉนฺโท จ อตฺตสมฺปทา.} \\ \csromangathabat{ทิฏฺฐิ จ อปฺปมาโท จ,}{โยนิโส ภวติ สตฺตมํ.}}}

\chapterheadpageread{8. ทุติยเอกธมฺมเปยฺยาลวคฺค}
\csromanheader{2}{1. กลฺยาณมิตฺต\-สุตฺต}
\proseitem{77}{\csromanfolio{28}สาวตฺถิ\-นิทานํ. นาหํ ภิกฺขเว อญฺญํ เอกธมฺมมฺปิ สมนุปสฺสามิ, เยน อนุปฺปนฺโน วา อริโย อฏฺฐงฺคิโก มคฺโค อุปฺปชฺชติ, อุปฺปนฺโน วา อริโย อฏฺฐงฺคิโก มคฺโค ภาวนา\-ปาริปูริํ คจฺฉติ, ยถยิทํ ภิกฺขเว กลฺยาณมิตฺตตา. กลฺยาณมิตฺตสฺเสตํ ภิกฺขเว ภิกฺขุโน ปาฏิกงฺขํ อริยํ อฏฺฐงฺคิกํ มคฺคํ ภาเวสฺสติ, อริยํ อฏฺฐงฺคิกํ มคฺคํ พหุลีกริสฺสติ. กถญฺจ ภิกฺขเว ภิกฺขุ กลฺยาณมิตฺโต อริยํ อฏฺฐงฺคิกํ มคฺคํ ภาเวติ, อริยํ อฏฺฐงฺคิกํ มคฺคํ พหุลีกโรติ. อิธ ภิกฺขเว ภิกฺขุ สมฺมาทิฏฺฐิํ ภาเวติ วิเวกนิสฺสิตํ ฯเปฯ สมฺมาสมาธิํ ภาเวติ วิเวกนิสฺสิตํ วิราคนิสฺสิตํ นิโรธ\-นิสฺสิตํ โวสฺสคฺค\-ปริณามิํ. เอวํ โข ภิกฺขเว ภิกฺขุ กลฺยาณมิตฺโต อริยํ อฏฺฐงฺคิกํ มคฺคํ ภาเวติ, อริยํ อฏฺฐงฺคิกํ มคฺคํ พหุลีกโรตีติ.  ปฐมํ.}

\csromanheader{2}{2-6. สีลสมฺปทา\-ทิ\-สุตฺต\-ปญฺจก}
\proseitem{78-82}{นาหํ ภิกฺขเว อญฺญํ เอกธมฺมมฺปิ สมนุปสฺสามิ, เยน อนุปฺปนฺโน วา อริโย อฏฺฐงฺคิโก มคฺโค อุปฺปชฺชติ, อุปฺปนฺโน วา อริโย อฏฺฐงฺคิโก มคฺโค ภาวนา\-ปาริปูริํ คจฺฉติ, ยถยิทํ ภิกฺขเว สีลสมฺปทา ฯเปฯ ยถยิทํ ภิกฺขเว ฉนฺทสมฺปทา ฯเปฯ ยถยิทํ ภิกฺขเว อตฺตสมฺปทา ฯเปฯ ยถยิทํ ภิกฺขเว ทิฏฺฐิสมฺปทา ฯเปฯ ยถยิทํ ภิกฺขเว อปฺปมาท\-สมฺปทา ฯเปฯ.  ฉฏฺฐํ.}

\csromanheader{2}{7. โยนิโสมนสิการ\-สมฺปทา\-สุตฺต}
\proseitem{83}{ยถยิทํ ภิกฺขเว โยนิโสมนสิการ\-สมฺปทา. โยนิโสมนสิการสมฺปนฺนสฺเสตํ ภิกฺขเว ภิกฺขุโน ปาฏิกงฺขํ อริยํ อฏฺฐงฺคิกํ มคฺคํ ภาเวสฺสติ, อริยํ อฏฺฐงฺคิกํ มคฺคํ พหุลีกริสฺสติ. กถญฺจ ภิกฺขเว ภิกฺขุ โยนิโสมนสิการ\-สมฺปนฺโน อริยํ อฏฺฐงฺคิกํ มคฺคํ ภาเวติ, อริยํ อฏฺฐงฺคิกํ มคฺคํ พหุลีกโรติ, อิธ ภิกฺขเว ภิกฺขุ สมฺมาทิฏฺฐิํ ภาเวติ วิเวกนิสฺสิตํ ฯเปฯ สมฺมาสมาธิํ ภาเวติ วิเวกนิสฺสิตํ วิราคนิสฺสิตํ นิโรธ\-นิสฺสิตํ โวสฺสคฺค\-ปริณามิํ. เอวํ โข ภิกฺขเว}

\prose{\csromanfolio{29}ภิกฺขุ โยนิโสมนสิการ\-สมฺปนฺโน อริยํ อฏฺฐงฺคิกํ มคฺคํ ภาเวติ, อริยํ อฏฺฐงฺคิกํ มคฺคํ พหุลีกโรตีติ.  สตฺตมํ.}

\csromanheader{2}{1. กลฺยาณมิตฺต\-สุตฺต}
\proseitem{84}{นาหํ ภิกฺขเว อญฺญํ เอกธมฺมมฺปิ สมนุปสฺสามิ, เยน อนุปฺปนฺโน วา อริโย อฏฺฐงฺคิโก มคฺโค อุปฺปชฺชติ, อุปฺปนฺโน วา อริโย อฏฺฐงฺคิโก มคฺโค ภาวนา\-ปาริปูริํ คจฺฉติ, ยถยิทํ ภิกฺขเว กลฺยาณมิตฺตตา. กลฺยาณมิตฺตสฺเสตํ ภิกฺขเว ภิกฺขุโน ปาฏิกงฺขํ อริยํ อฏฺฐงฺคิกํ มคฺคํ ภาเวสฺสติ, อริยํ อฏฺฐงฺคิกํ มคฺคํ พหุลีกริสฺสติ. กถญฺจ ภิกฺขเว ภิกฺขุ กลฺยาณมิตฺโต อริยํ อฏฺฐงฺคิกํ มคฺคํ ภาเวติ, อริยํ อฏฺฐงฺคิกํ มคฺคํ พหุลีกโรติ, อิธ ภิกฺขเว ภิกฺขุ สมฺมาทิฏฺฐิํ ภาเวติ ราค\-วินย\-ปริโยสานํ โทส\-วินย\-ปริโยสานํ โมห\-วินย\-ปริโยสานํ ฯเปฯ สมฺมาสมาธิํ ภาเวติ ราค\-วินย\-ปริโยสานํ โทส\-วินย\-ปริโยสานํ โมห\-วินย\-ปริโยสานํ. เอวํ โข ภิกฺขเว ภิกฺขุ กลฺยาณมิตฺโต อริยํ อฏฺฐงฺคิกํ มคฺคํ ภาเวติ, อริยํ อฏฺฐงฺคิกํ มคฺคํ พหุลีกโรตีติ.  ปฐมํ.}

\csromanheader{2}{2-6. สีลสมฺปทา\-ทิ\-สุตฺต\-ปญฺจก}
\proseitem{85-89}{นาหํ ภิกฺขเว อญฺญํ เอกธมฺมมฺปิ สมนุปสฺสามิ, เยน อนุปฺปนฺโน วา อริโย อฏฺฐงฺคิโก มคฺโค อุปฺปชฺชติ, อุปฺปนฺโน วา อริโย อฏฺฐงฺคิโก มคฺโค ภาวนา\-ปาริปูริํ คจฺฉติ, ยถยิทํ ภิกฺขเว สีลสมฺปทา ฯเปฯ ยถยิทํ ภิกฺขเว ฉนฺทสมฺปทา ฯเปฯ ยถยิทํ ภิกฺขเว อตฺตสมฺปทา ฯเปฯ ยถยิทํ ภิกฺขเว ทิฏฺฐิสมฺปทา ฯเปฯ ยถยิทํ ภิกฺขเว อปฺปมาท\-สมฺปทา ฯเปฯ.  ฉฏฺฐํ.}

\csromanheader{2}{7. โยนิโสมนสิการ\-สมฺปทา\-สุตฺต}
\proseitem{90}{ยถยิทํ ภิกฺขเว โยนิโสมนสิการ\-สมฺปทา. โยนิโสมนสิการสมฺปนฺนสฺเสตํ ภิกฺขเว ภิกฺขุโน ปาฏิกงฺขํ อริยํ อฏฺฐงฺคิกํ มคฺคํ ภาเวสฺสติ, อริยํ อฏฺฐงฺคิกํ มคฺคํ พหุลีกริสฺสติ. กถญฺจ ภิกฺขเว ภิกฺขุ โยนิโสมนสิการ\-สมฺปนฺโน อริยํ อฏฺฐงฺคิกํ มคฺคํ ภาเวติ, อริยํ อฏฺฐงฺคิกํ มคฺคํ พหุลีกโรติ, อิธ ภิกฺขเว ภิกฺขุ สมฺมาทิฏฺฐิํ ภาเวติ ราค\-วินย\-ปริโยสานํ โทส\-วินย\-ปริโยสานํ โมห\-วินย\-ปริโยสานํ ฯเปฯ สมฺมาสมาธิํ ภาเวติ ราค\-วินย\-ปริโยสานํ โทส\-วินย\-ปริโยสานํ โมห\-วินย\-ปริโยสานํ. เอวํ \csromanfolio{30}โข ภิกฺขเว ภิกฺขุ โยนิโสมนสิการ\-สมฺปนฺโน อริยํ อฏฺฐงฺคิกํ มคฺคํ ภาเวติ, อริยํ อฏฺฐงฺคิกํ มคฺคํ พหุลีกโรตีติ.  สตฺตมํ. เอกธมฺม\-เปยฺยาลํ.}

\vspace{\tipitakaparskip}
\csromancenter{\textbf{ตสฺสุทฺทานํ}}
\setlength{\csromangathapagewidth}{0pt}
\setlength{\csromangathawakwidth}{0pt}
\csromangathameasuregroup{กลฺยาณมิตฺตํ สีลญฺจ, \\ ทิฏฺฐิ จ อปฺปมาโท จ,}{\csromangathabat{กลฺยาณมิตฺตํ สีลญฺจ,}{ฉนฺโท จ อตฺตสมฺปทา.} \\ \csromangathabat{ทิฏฺฐิ จ อปฺปมาโท จ,}{โยนิโส ภวติ สตฺตมํ.}}
\csromangathasetleft{กลฺยาณมิตฺตํ สีลญฺจ, \\ ทิฏฺฐิ จ อปฺปมาโท จ,}
\csromangathagroup{\csromangathastanza{\csromangathabat{กลฺยาณมิตฺตํ สีลญฺจ,}{ฉนฺโท จ อตฺตสมฺปทา.} \\ \csromangathabat{ทิฏฺฐิ จ อปฺปมาโท จ,}{โยนิโส ภวติ สตฺตมํ.}}}

\csromanmatikaanchor{mtk.4}
\csromanmatikaanchor{mtk.23}
\csromantocmarkref{section}{8. อาวรณนีวรณสุตฺต}{mtk.4}
\bookmark[dest={mtk.4},level=1]{8. อาวรณนีวรณสุตฺต}
\csromanheader{4}{1. คงฺคา\-เปยฺยาล\-วคฺค}
\csromanheader{2}{1. ปฐม\-ปาจีนนินฺน\-สุตฺต}
\proseitem{91}{สาวตฺถิ\-นิทานํ. เสยฺยถาปิ ภิกฺขเว คงฺคา นที ปาจีนนินฺนา ปาจีนโปณา ปาจีนปพฺภารา. เอวเมว โข ภิกฺขเว ภิกฺขุ อริยํ อฏฺฐงฺคิกํ มคฺคํ ภาเวนฺโต อริยํ อฏฺฐงฺคิกํ มคฺคํ พหุลีกโรนฺโต นิพฺพานนินฺโน โหติ นิพฺพานโปโณ นิพฺพาน\-ปพฺภาโร. กถญฺจ ภิกฺขเว ภิกฺขุ อริยํ อฏฺฐงฺคิกํ มคฺคํ ภาเวนฺโต อริยํ อฏฺฐงฺคิกํ มคฺคํ พหุลีกโรนฺโต นิพฺพานนินฺโน โหติ นิพฺพานโปโณ นิพฺพาน\-ปพฺภาโร, อิธ ภิกฺขเว ภิกฺขุ สมฺมาทิฏฺฐิํ ภาเวติ วิเวกนิสฺสิตํ วิราคนิสฺสิตํ นิโรธ\-นิสฺสิตํ โวสฺสคฺค\-ปริณามิํ ฯเปฯ สมฺมาสมาธิํ ภาเวติ วิเวกนิสฺสิตํ วิราคนิสฺสิตํ นิโรธ\-นิสฺสิตํ โวสฺสคฺค\-ปริณามิํ. เอวํ โข ภิกฺขเว ภิกฺขุ อริยํ อฏฺฐงฺคิกํ มคฺคํ ภาเวนฺโต อริยํ อฏฺฐงฺคิกํ มคฺคํ พหุลีกโรนฺโต นิพฺพานนินฺโน โหติ นิพฺพานโปโณ นิพฺพาน\-ปพฺภาโรติ.  อฐมํ.}

\csromanheader{2}{2-5. ทุติยา\-ทิ\-ปาจีนนินฺน\-สุตฺต\-จตุกฺก}
\proseitem{92-95}{เสยฺยถาปิ ภิกฺขเว ยมุนา นที ปาจีนนินฺนา ปาจีนโปณา ปาจีนปพฺภารา. เอวเมว โข ภิกฺขเว ฯเปฯ. เสยฺยถาปิ ภิกฺขเว อจิรวตี นที ปาจีนนินฺนา ปาจีนโปณา ปาจีนปพฺภารา. เอวเมว โข ภิกฺขเว ฯเปฯ. เสยฺยถาปิ ภิกฺขเว สรภู นที ปาจีนนินฺนา ปาจีนโปณา ปาจีนปพฺภารา. เอวเมว โข ภิกฺขเว ฯเปฯ. เสยฺยถาปิ ภิกฺขเว มหี นที ปาจีนนินฺนา ปาจีนโปณา ปาจีนปพฺภารา. เอวเมว โข ภิกฺขเว ฯเปฯ.  ปญฺจมํ.}

\csromanheader{2}{6. ฉฏฺฐ\-ปาจีนนินฺน\-สุตฺต}
\proseitem{96}{\csromanfolio{31}เสยฺยถาปิ ภิกฺขเว ยา กาจิมา มหานทิโย, เสยฺยถิทํ. คงฺคา ยมุนา อจิรวตี สรภู มหี, สพฺพา ตา ปาจีนนินฺนา ปาจีนโปณา ปาจีนปพฺภารา. เอวเมว โข ภิกฺขเว ภิกฺขุ อริยํ อฏฺฐงฺคิกํ มคฺคํ ภาเวนฺโต อริยํ อฏฺฐงฺคิกํ มคฺคํ พหุลีกโรนฺโต นิพฺพานนินฺโน โหติ นิพฺพานโปโณ นิพฺพาน\-ปพฺภาโร. กถญฺจ ภิกฺขเว ภิกฺขุ อริยํ อฏฺฐงฺคิกํ มคฺคํ ภาเวนฺโต อริยํ อฏฺฐงฺคิกํ มคฺคํ พหุลีกโรนฺโต นิพฺพานนินฺโน โหติ นิพฺพานโปโณ นิพฺพาน\-ปพฺภาโร, อิธ ภิกฺขเว ภิกฺขุ สมฺมาทิฏฺฐิํ ภาเวติ วิเวกนิสฺสิตํ ฯเปฯ สมฺมาสมาธิํ ภาเวติ วิเวกนิสฺสิตํ วิราคนิสฺสิตํ นิโรธ\-นิสฺสิตํ โวสฺสคฺค\-ปริณามิํ. เอวํ โข ภิกฺขเว ภิกฺขุ อริยํ อฏฺฐงฺคิกํ มคฺคํ ภาเวนฺโต อริยํ อฏฺฐงฺคิกํ มคฺคํ พหุลีกโรนฺโต นิพฺพานนินฺโน โหติ นิพฺพานโปโณ นิพฺพาน\-ปพฺภาโรติ.  ฉฏฺฐํ.}

\csromanheader{2}{1. ปฐม\-สมุทฺทนินฺน\-สุตฺต}
\proseitem{97}{เสยฺยถาปิ ภิกฺขเว คงฺคา นที สมุทฺทนินฺนา สมุทฺทโปณา สมุทฺท\-ปพฺภารา. เอวเมว โข ภิกฺขเว ภิกฺขุ อริยํ อฏฺฐงฺคิกํ มคฺคํ ภาเวนฺโต อริยํ อฏฺฐงฺคิกํ มคฺคํ พหุลีกโรนฺโต นิพฺพานนินฺโน โหติ นิพฺพานโปโณ นิพฺพาน\-ปพฺภาโร. กถญฺจ ภิกฺขเว ภิกฺขุ อริยํ อฏฺฐงฺคิกํ มคฺคํ ภาเวนฺโต อริยํ อฏฺฐงฺคิกํ มคฺคํ พหุลีกโรนฺโต นิพฺพานนินฺโน โหติ นิพฺพานโปโณ นิพฺพาน\-ปพฺภาโร, อิธ ภิกฺขเว ภิกฺขุ สมฺมาทิฏฺฐิํ ภาเวติ วิเวกนิสฺสิตํ ฯเปฯ สมฺมาสมาธิํ ภาเวติ วิเวกนิสฺสิตํ วิราคนิสฺสิตํ นิโรธ\-นิสฺสิตํ โวสฺสคฺค\-ปริณามิํ. เอวํ โข ภิกฺขเว ภิกฺขุ อริยํ อฏฺฐงฺคิกํ มคฺคํ ภาเวนฺโต อริยํ อฏฺฐงฺคิกํ มคฺคํ พหุลีกโรนฺโต นิพฺพานนินฺโน โหติ นิพฺพานโปโณ นิพฺพาน\-ปพฺภาโรติ.  ปฐมํ.}

\csromanheader{2}{2-6. ทุติยา\-ทิ\-สมุทฺทนินฺน\-สุตฺต\-ปญฺจก}
\proseitem{98-102}{เสยฺยถาปิ ภิกฺขเว ยมุนา นที สมุทฺทนินฺนา สมุทฺทโปณา สมุทฺท\-ปพฺภารา. เอวเมว โข ภิกฺขเว ภิกฺขุ ฯเปฯ. เสยฺยถาปิ ภิกฺขเว อจิรวตี นที สมุทฺทนินฺนา สมุทฺทโปณา สมุทฺท\-ปพฺภารา. เอวเมว โข ภิกฺขเว ภิกฺขุ ฯเปฯ. เสยฺยถาปิ ภิกฺขเว สรภู นที สมุทฺทนินฺนา สมุทฺทโปณา สมุทฺทปพฺภรา. เอวเมว โข ภิกฺขเว ภิกฺขุ ฯเปฯ. เสยฺยถาปิ ภิกฺขเว \csromanfolio{32}มหี นที สมุทฺทนินฺนา สมุทฺทโปณา สมุทฺท\-ปพฺภารา. เอวเมว โข ภิกฺขเว ภิกฺขุ ฯเปฯ. เสยฺยถาปิ ภิกฺขเว ยา กาจิมา มหานทิโย, เสยฺยถิทํ, คงฺคา ยมุนา อจิรวตี สรภู มหี, สพฺพา ตา สมุทฺทนินฺนา สมุทฺทโปณา สมุทฺท\-ปพฺภารา. เอวเมว โข ภิกฺขเว ภิกฺขุ อริยํ อฏฺฐงฺคิกํ มคฺคํ ภาเวนฺโต อริยํ อฏฺฐงฺคิกํ มคฺคํ พหุลีกโรนฺโต นิพฺพานนินฺโน โหติ นิพฺพานโปโณ นิพฺพาน\-ปพฺภาโร. กถญฺจ ภิกฺขเว ภิกฺขุ อริยํ อฏฺฐงฺคิกํ มคฺคํ ภาเวนฺโต อริยํ อฏฺฐงฺคิกํ มคฺคํ พหุลีกโรนฺโต นิพฺพานนินฺโน โหติ นิพฺพานโปโณ นิพฺพาน\-ปพฺภาโร, อิธ ภิกฺขเว ภิกฺขุ สมฺมาทิฏฺฐิํ ภาเวติ วิเวกนิสฺสิตํ วิราคนิสฺสิตํ นิโรธ\-นิสฺสิตํ โวสฺสคฺค\-ปริณามิํ ฯเปฯ สมฺมาสมาธิํ ภาเวติ วิเวกนิสฺสตํ วิราคนิสฺสิตํ นิโรธ\-นิสฺสิตํ โวสฺสคฺค\-ปริณามิํ. เอวํ โข ภิกฺขเว ภิกฺขุ อริยํ อฏฺฐงฺคิกํ มคฺคํ ภาเวนฺโต อริยํ อฏฺฐงฺคิกํ มคฺคํ พหุลีกโรนฺโต นิพฺพานนินฺโน โหติ นิพฺพานโปโณ นิพฺพาน\-ปพฺภาโรติ.  ฉฏฺฐํ. คงฺคาเปยฺยาลํ.}

\vspace{\tipitakaparskip}
\csromancenter{\textbf{ตสฺสุทฺทานํ}}
\setlength{\csromangathapagewidth}{0pt}
\setlength{\csromangathawakwidth}{0pt}
\csromangathameasuregroup{ฉ ปาจีนโต นินฺนา, \\ เอเต เทฺว ฉ ทฺวาทส โหนฺติ,}{\csromangathabat{ฉ ปาจีนโต นินฺนา,}{ฉ นินฺนา จ สมุทฺทโต.} \\ \csromangathabat{เอเต เทฺว ฉ ทฺวาทส โหนฺติ,}{วคฺโค เตน ปวุจฺจตีติ.} \\ คงฺคาเปยฺยาลี ปาจินนินฺนวาจนมคฺคี,}
\csromangathameasurewak{คงฺคาเปยฺยาลี ปาจินนินฺนวาจนมคฺคี,}
\csromangathasetleft{ฉ ปาจีนโต นินฺนา, \\ เอเต เทฺว ฉ ทฺวาทส โหนฺติ,}
\csromangathagroup{\csromangathastanza{\csromangathabat{ฉ ปาจีนโต นินฺนา,}{ฉ นินฺนา จ สมุทฺทโต.} \\ \csromangathabat{เอเต เทฺว ฉ ทฺวาทส โหนฺติ,}{วคฺโค เตน ปวุจฺจตีติ.}}\csromangathastanza{\csromangathawak{คงฺคาเปยฺยาลี ปาจินนินฺนวาจนมคฺคี,}}}

\prose{วิเวกนิสฺสิตํ ทฺวาทสกี ปฐมกี.}
\csromansectionrule

\chapterhead{2. ทุติย\-คงฺคาเปยฺยาล\-วคฺค}
\csromanheader{2}{1. ปฐม\-ปาจีนนินฺน\-สุตฺต}
\proseitem{103}{เสยฺยถาปิ ภิกฺขเว คงฺคา นที ปาจีนนินฺนา ปาจีนโปณา ปาจีนปพฺภารา. เอวเมว โข ภิกฺขเว ภิกฺขุ อริยํ อฏฺฐงฺคิกํ มคฺคํ ภาเวนฺโต อริยํ อฏฺฐงฺคิกํ มคฺคํ พหุลีกโรนฺโต นิพฺพานนินฺโน โหติ นิพฺพานโปโณ นิพฺพาน\-ปพฺภาโร. กถญฺจ ภิกฺขเว ภิกฺขุ อริยํ อฏฺฐงฺคิกํ มคฺคํ ภาเวนฺโต อริยํ อฏฺฐงฺคิกํ มคฺคํ พหุลีกโรนฺโต นิพฺพานนินฺโน โหติ นิพฺพานโปโณ นิพฺพาน\-ปพฺภาโร, อิธ ภิกฺขเว ภิกฺขุ สมฺมาทิฏฺฐิํ ภาเวติ}

\prose{\csromanfolio{33}ราค\-วินย\-ปริโยสานํ โทส\-วินย\-ปริโยสานํ โมห\-วินย\-ปริโยสานํ ฯเปฯ สมฺมาสมาธิํ ภาเวติ ราค\-วินย\-ปริโยสานํ โทส\-วินย\-ปริโยสานํ โมห\-วินย\-ปริโยสานํ. เอวํ โข ภิกฺขเว ภิกฺขุ อริยํ อฏฺฐงฺคิกํ มคฺคํ ภาเวนฺโต อริยํ อฏฺฐงฺคิกํ มคฺคํ พหุลีกโรนฺโต นิพฺพานนินฺโน โหติ นิพฺพานโปโณ นิพฺพาน\-ปพฺภาโรติ.  ปฐมํ.}

\csromanheader{2}{2-6. ทุติยา\-ทิ\-ปาจีนนินฺน\-สุตฺต\-ปญฺจก}
\proseitem{104-108}{เสยฺยถาปิ ภิกฺขเว ยมุนา นที ปาจีนนินฺนา ปาจีนโปณา ปาจีนปพฺภารา. เอวเมว โข ภิกฺขเว ภิกฺขุ ฯเปฯ. เสยฺยถาปิ ภิกฺขเว อจิรวตี นที ปาจีนนินฺนา ปาจีนโปณา ปาจีนปพฺภารา. เอวเมว โข ภิกฺขเว ภิกฺขุ ฯเปฯ. เสยฺยถาปิ ภิกฺขเว สรภู นที ปาจีนนินฺนา ปาจีนโปณา ปาจีนปพฺภารา. เอวเมว โข ภิกฺขเว ภิกฺขุ ฯเปฯ. เสยฺยถาปิ ภิกฺขเว มหี นที ปาจีนนินฺนา ปาจีนโปณา ปาจีนปพฺภารา. เอวเมว โข ภิกฺขเว ภิกฺขุ ฯเปฯ. เสยฺยถาปิ ภิกฺขเว ยา กาจิมา มหานทิโย. เสยฺยถิทํ, คงฺคา ยมุนา อจิรวตี สรภู มหี, สพฺพา ตา ปาจีนนินฺนา ปาจีนโปณา ปาจีนปพฺภารา. เอวเมว โข ภิกฺขเว ภิกฺขุ ฯเปฯ. ฉฏฺฐํ.}

\csromanheader{2}{7. ปฐม\-สมุทฺทนินฺน\-สุตฺต}
\proseitem{109}{เสยฺยถาปิ ภิกฺขเว คงฺคา นที สมุทฺทนินฺนา สมุทฺทโปณา สมุทฺท\-ปพฺภารา. เอวเมว โข ภิกฺขเว ภิกฺขุ อริยํ อฏฺฐงฺคิกํ มคฺคํ ภาเวนฺโต อริยํ อฏฺฐงฺคิกํ มคฺคํ พหุลีกโรนฺโต นิพฺพานนินฺโน โหติ นิพฺพานโปโณ นิพฺพาน\-ปพฺภาโร. กถญฺจ ภิกฺขเว ภิกฺขุ อริยํ อฏฺฐงฺคิกํ มคฺคํ ภาเวนฺโต อริยํ อฏฺฐงฺคิกํ มคฺคํ พหุลีกโรนฺโต นิพฺพานนินฺโน โหติ นิพฺพานโปโณ นิพฺพาน\-ปพฺภาโร, อิธ ภิกฺขเว ภิกฺขุ สมฺมาทิฏฺฐิํ ภาเวติ ราค\-วินย\-ปริโยสานํ โทส\-วินย\-ปริโยสานํ โมห\-วินย\-ปริโยสานํ ฯเปฯ สมฺมาสมาธิํ ภาเวติ ราค\-วินย\-ปริโยสานํ โทส\-วินย\-ปริโยสานํ โมห\-วินย\-ปริโยสานํ. เอวํ โข ภิกฺขเว ภิกฺขุ อริยํ อฏฺฐงฺคิกํ มคฺคํ ภาเวนฺโต อริยํ อฏฺฐงฺคิกํ มคฺคํ พหุลีกโรนฺโต นิพฺพานนินฺโน โหติ นิพฺพานโปโณ นิพฺพาน\-ปพฺภาโรติ.  ปฐมํ.}

\csromanheader{2}{8-12. ทุติยา\-ทิ\-สมุทฺทนินฺน\-สุตฺต}
\proseitem{110-114}{เสยฺยถาปิ ภิกฺขเว ยมุนา นที สมุทฺทนินฺนา สมุทฺทโปณา สมุทฺทปพฺภรา. เอวเมว โข ภิกฺขเว ภิกฺขุ ฯเปฯ. เสยฺยถาปิ ภิกฺขเว \csromanfolio{34}อจิรวตี นที สมุทฺทนินฺนา สมุทฺทโปณา สมุทฺท\-ปพฺภารา. เอวเมว โข ภิกฺขเว ภิกฺขุ ฯเปฯ. เสยฺยถาปิ ภิกฺขเว สรภู นที สมุทฺทนินฺนา สมุทฺทโปณา สมุทฺท\-ปพฺภารา. เอวเมว โข ภิกฺขเว ภิกฺขุ ฯเปฯ. เสยฺยถาปิ ภิกฺขเว มหี นที สมุทฺทนินฺนา สมุทฺทโปณา สมุทฺท\-ปพฺภารา. เอวเมว โข ภิกฺขเว ภิกฺขุ ฯเปฯ. เสยฺยถาปิ ภิกฺขเว ยา กาจิมา มหานทิโย. เสยฺยถิทํ, คงฺคา ยมุนา อจิรวตี สรภู มหี, สพฺพา ตา สมุทฺทนินฺนา สมุทฺทโปณา สมุทฺท\-ปพฺภารา. เอวเมว โข ภิกฺขเว ภิกฺขุ อริยํ อฏฺฐงฺคิกํ มคฺคํ ภาเวนฺโต อริยํ อฏฺฐงฺคิกํ มคฺคํ พหุลีกโรนฺโต นิพฺพานนินฺโน โหติ นิพฺพานโปโณ นิพฺพาน\-ปพฺภาโร. กถญฺจ ภิกฺขเว ภิกฺขุ อริยํ อฏฺฐงฺคิกํ มคฺคํ ภาเวนฺโต อริยํ อฏฺฐงฺคิกํ มคฺคํ พหุลีกโรนฺโต นิพฺพานนินฺโน โหติ นิพฺพานโปโณ นิพฺพาน\-ปพฺภาโร. อิธ ภิกฺขเว ภิกฺขุ สมฺมาทิฏฺฐิํ ภาเวติ ราค\-วินย\-ปริโยสานํ โทส\-วินย\-ปริโยสานํ โมห\-วินย\-ปริโยสานํ ฯเปฯ สมฺมาสมาธิํ ภาเวติ ราค\-วินย\-ปริโยสานํ โทส\-วินย\-ปริโยสานํ โมห\-วินย\-ปริโยสานํ. เอวํ โข ภิกฺขเว ภิกฺขุ อริยํ อฏฺฐงฺคิกํ มคฺคํ ภาเวนฺโต อริยํ อฏฺฐงฺคิกํ มคฺคํ พหุลีกโรนฺโต นิพฺพานนินฺโน โหติ นิพฺพานโปโณ นิพฺพาน\-ปพฺภาโรติ.  ฉฏฺฐํ.}

\vspace{\tipitakaparskip}
\csromancenter{(ราควินย\-ทฺวาทสกี ทุติยกี สมุทฺท\-นินฺนนฺติ.)}
\csromanheader{2}{13. ปฐม\-ปาจีนนินฺน\-สุตฺต}
\proseitem{115}{เสยฺยถาปิ ภิกฺขเว คงฺคา นที ปาจีนนินฺนา ปาจีนโปณา ปาจีนปพฺภารา. เอวเมว โข ภิกฺขเว ภิกฺขุ อริยํ อฏฺฐงฺคิกํ มคฺคํ ภาเวนฺโต อริยํ อฏฺฐงฺคิกํ มคฺคํ พหุลีกโรนฺโต นิพฺพานนินฺโน โหติ นิพฺพานโปโณ นิพฺพาน\-ปพฺภาโร. กถญฺจ ภิกฺขเว ภิกฺขุ อริยํ อฏฺฐงฺคิกํ มคฺคํ ภาเวนฺโต อริยํ อฏฺฐงฺคิกํ มคฺคํ พหุลีกโรนฺโต นิพฺพานนินฺโน โหติ นิพฺพานโปโณ นิพฺพาน\-ปพฺภาโร, อิธ ภิกฺขเว ภิกฺขุ สมฺมาทิฏฺฐิํ ภาเวติ อมโตคธํ อมตปรายนํ อมต\-ปริโยสานํ ฯเปฯ สมฺมาสมาธิํ ภาเวติ อมโตคธํ อมตปรายนํ อมต\-ปริโยสานํ. เอวํ โข ภิกฺขเว ภิกฺขุ อริยํ อฏฺฐงฺคิกํ มคฺคํ ภาเวนฺโต อริยํ อฏฺฐงฺคิกํ มคฺคํ พหุลีกโรนฺโต นิพฺพานนินฺโน โหติ นิพฺพานโปโณ นิพฺพาน\-ปพฺภาโรติ.  ปฐมํ.}

\csromanheader{2}{14-18. ทุติยา\-ทิ\-ปาจีนนินฺน\-สุตฺต}
\proseitem{116-120}{\csromanfolio{35}เสยฺยถาปิ ภิกฺขเว ยมุนา นที ปาจีนนินฺนา ปาจีนโปณา ปาจีนปพฺภารา. เอวเมว โข ภิกฺขเว ภิกฺขุ ฯเปฯ. เสยฺยถาปิ ภิกฺขเว อจิรวตี นที ปาจีนนินฺนา ปาจีนโปณา ปาจีนปพฺภารา. เอวเมว โข ภิกฺขเว ภิกฺขุ ฯเปฯ. เสยฺยถาปิ ภิกฺขเว สรภู นที ปาจีนนินฺนา ปาจีนโปณา ปาจีนปพฺภารา. เอวเมว โข ภิกฺขเว ภิกฺขุ ฯเปฯ. เสยฺยถาปิ ภิกฺขเว มหี นที ปาจีนนินฺนา ปาจีนโปณา ปาจีนปพฺภารา. เอวเมว โข ภิกฺขเว ภิกฺขุ ฯเปฯ. เสยฺยถาปิ ภิกฺขเว ยา กาจิมา มหานทิโย. เสยฺยถิทํ, คงฺคา ยมุนา อจิรวตี สรภู มหี, สพฺพา ตา ปาจีนนินฺนา ปาจีนโปณา ปาจีนปพฺภารา. เอวเมว โข ภิกฺขเว ภิกฺขุ ฯเปฯ.  ฉฏฺฐํ.}

\csromanheader{2}{19. ปฐม\-สมุทฺทนินฺน\-สุตฺต}
\proseitem{121}{เสยฺยถาปิ ภิกฺขเว คงฺคา นที สมุทฺทนินฺนา สมุทฺทโปณา สมุทฺท\-ปพฺภารา. เอวเมว โข ภิกฺขเว ภิกฺขุ อริยํ อฏฺฐงฺคิกํ มคฺคํ ภาเวนฺโต อริยํ อฏฺฐงฺคิกํ มคฺคํ พหุลีกโรนฺโต นิพฺพานนินฺโน โหติ นิพฺพานโปโณ นิพฺพาน\-ปพฺภาโร. กถญฺจ ภิกฺขเว ภิกฺขุ อริยํ อฏฺฐงฺคิกํ มคฺคํ ภาเวนฺโต อริยํ อฏฺฐงฺคิกํ มคฺคํ พหุลีกโรนฺโต นิพฺพานนินฺโน โหติ นิพฺพานโปโณ นิพฺพาน\-ปพฺภาโร, อิธ ภิกฺขเว ภิกฺขุ สมฺมาทิฏฺฐิํ ภาเวติ อมโตคธํ อมตปรายนํ อมต\-ปริโยสานํ ฯเปฯ สมฺมาสมาธิํ ภาเวติ อมโตคธํ อมตปรายนํ อมต\-ปริโยสานํ. เอวํ โข ภิกฺขเว ภิกฺขุ อริยํ อฏฺฐงฺคิกํ มคฺคํ ภาเวนฺโต อริยํ อฏฺฐงฺคิกํ มคฺคํ พหุลีกโรนฺโต นิพฺพานนินฺโน โหติ นิพฺพานโปโณ นิพฺพาน\-ปพฺภาโรติ.  ปฐมํ.}

\csromanheader{2}{20-24. ทุติยา\-ทิ\-สมุทฺทนินฺน\-สุตฺต}
\proseitem{122-126}{เสยฺยถาปิ ภิกฺขเว ยมุนา นที สมุทฺทนินฺนา สมุทฺทโปณา สมุทฺทาปพฺภารา. เอวเมว โข ภิกฺขเว ภิกฺขุ ฯเปฯ. เสยฺยถาปิ ภิกฺขเว อจิรวตี นที สมุทฺทนินฺนา สมุทฺทโปณา สมุทฺท\-ปพฺภารา. เอวเมว โข ภิกฺขเว ภิกฺขุ ฯเปฯ. เสยฺยถาปิ ภิกฺขเว สรภู นที สมุทฺทนินฺนา สมุทฺทโปณา สมุทฺท\-ปพฺภารา. เอวเมว โข ภิกฺขเว ภิกฺขุ ฯเปฯ. เสยฺยถาปิ ภิกฺขเว มหี นที สมุทฺทนินฺนา สมุทฺทโปณา สมุทฺท\-ปพฺภารา. เอวเมว โข ภิกฺขเว ภิกฺขุ ฯเปฯ. เสยฺยถาปิ ภิกฺขเว ยา กาจิมา มหานทิโย. เสยฺยถิทํ, \csromanfolio{36}คงฺคา ยมุนา อจิรวตี สรภู มหี, สพฺพา ตา สมุทฺทนินฺนา สมุทฺทโปณา สมุทฺท\-ปพฺภารา. เอวเมว โข ภิกฺขเว ภิกฺขุ อริยํ อฏฺฐงฺคิกํ มคฺคํ ภาเวนฺโต อริยํ อฏฺฐงฺคิกํ มคฺคํ พหุลีกโรนฺโต นิพฺพานนินฺโน โหติ นิพฺพานโปโณ นิพฺพาน\-ปพฺภาโร. กถญฺจ ภิกฺขเว ภิกฺขุ อริยํ อฏฺฐงฺคิกํ มคฺคํ ภาเวนฺโต อริยํ อฏฺฐงฺคิกํ มคฺคํ พหุลีกโรนฺโต นิพฺพานนินฺโน โหติ นิพฺพานโปโณ นิพฺพาน\-ปพฺภาโร, อิธ ภิกฺขเว ภิกฺขุ สมฺมาทิฏฺฐิํ ภาเวติ อมโตคธํ อมตปรายนํ อมต\-ปริโยสานํ ฯเปฯ สมฺมาสมาธิํ ภาเวติ อมโตคธํ อมตปรายนํ อมต\-ปริโยสานํ. เอวํ โข ภิกฺขเว ภิกฺขุ อริยํ อฏฺฐงฺคิกํ มคฺคํ ภาเวนฺโต อริยํ อฏฺฐงฺคิกํ มคฺคํ พหุลีกโรนฺโต นิพฺพานนินฺโน โหติ นิพฺพานโปโณ นิพฺพาน\-ปพฺภาโรติ.  ฉฏฺฐํ. (อมโตคธ\-ทฺวาทสกี ตติยกี.)}

\csromanheader{2}{25. ปฐม\-ปาจีนนินฺน\-สุตฺต}
\proseitem{127}{เสยฺยถาปิ ภิกฺขเว คงฺคา นที ปาจีนนินฺนา ปาจีนโปณา ปาจีนปพฺภารา. เอวเมว โข ภิกฺขเว ภิกฺขุ อริยํ อฏฺฐงฺคิกํ มคฺคํ ภาเวนฺโต อริยํ อฏฺฐงฺคิกํ มคฺคํ พหุลีกโรนฺโต นิพฺพานนินฺโน โหติ นิพฺพานโปโณ นิพฺพาน\-ปพฺภาโร. กถญฺจ ภิกฺขเว ภิกฺขุ อริยํ อฏฺฐงฺคิกํ มคฺคํ ภาเวนฺโต อริยํ อฏฺฐงฺคิกํ มคฺคํ พหุลีกโรนฺโต นิพฺพานนินฺโน โหติ นิพฺพานโปโณ นิพฺพาน\-ปพฺภาโร, อิธ ภิกฺขเว ภิกฺขุ สมฺมาทิฏฺฐิํ ภาเวติ นิพฺพานนินฺนํ นิพฺพานโปณํ นิพฺพาน\-ปพฺภารํ ฯเปฯ สมฺมาสมาธิํ ภาเวติ นิพฺพานนินฺนํ นิพฺพานโปณํ นิพฺพาน\-ปพฺภารํ. เอวํ โข ภิกฺขเว ภิกฺขุ อริยํ อฏฺฐงฺคิกํ มคฺคํ ภาเวนฺโต อริยํ อฏฺฐงฺคิกํ มคฺคํ พหุลีกโรนฺโต นิพฺพานนินฺโน โหติ นิพฺพานโปโณ นิพฺพาน\-ปพฺภาโรติ.  ปฐมํ.}

\csromanheader{2}{26-30. ทุติยา\-ทิ\-ปาจีนนินฺน\-สุตฺต}
\proseitem{128-132}{เสยฺยถาปิ ภิกฺขเว ยมุนา นที ปาจีนนินฺนา ปาจีนโปณา ปาจีนปพฺภารา. เอวเมว โข ภิกฺขเว ภิกฺขุ ฯเปฯ. เสยฺยถาปิ ภิกฺขเว อจิรวตี นที ปาจีนนินฺนา ปาจีนโปณา ปาจีนปพฺภารา. เอวเมว โข ภิกฺขเว ภิกฺขุ ฯเปฯ. เสยฺยถาปิ ภิกฺขเว สรภู นที ปาจีนนินฺนา ปาจีนโปณา ปาจีนปพฺภารา. เอวเมว โข ภิกฺขเว ภิกฺขุ ฯเปฯ. เสยฺยถาปิ ภิกฺขเว มหี นที ปาจีนนินฺนา}

\prose{\csromanfolio{37}ปาจีนโปณา ปาจีนปพฺพารา. เอวเมว โข ภิกฺขเว ภิกฺขุ ฯเปฯ. เสยฺยถาปิ ภิกฺขเว ยา กาจิมา มหานทิโย. เสยฺยถิทํ, คงฺคา ยมุนา อจิรวตี สรภู มหี, สพฺพา ตา ปาจีนนินฺนา ปาจีนโปณา ปาจีนปพฺภารา. เอวเมว โข ภิกฺขเว ภิกฺขุ อริยํ อฏฺฐงฺคิกํ มคฺคํ ภาเวนฺโต อริยํ อฏฺฐงฺคิกํ มคฺคํ พหุลีกโรนฺโต นิพฺพานนินฺโน โหติ นิพฺพานโปโณ นิพฺพาน\-ปพฺภาโร. กถญฺจ ภิกฺขเว ภิกฺขุ อริยํ อฏฺฐงฺคิกํ มคฺคํ ภาเวนฺโต อริยํ อฏฺฐงฺคิกํ มคฺคํ พหุลีกโรนฺโต นิพฺพานนินฺโน โหติ นิพฺพานโปโณ นิพฺพาน\-ปพฺภาโร, อิธ ภิกฺขเว ภิกฺขุ สมฺมาทิฏฺฐิํ ภาเวติ นิพฺพานนินฺนํ นิพฺพานโปณํ นิพฺพาน\-ปพฺภารํ ฯเปฯ สมฺมาสมาธิํ ภาเวติ นิพฺพานนินฺนํ นิพฺพานโปณํ นิพฺพาน\-ปพฺภารํ. เอวํ โข ภิกฺขเว ภิกฺขุ อริยํ อฏฺฐงฺคิกํ มคฺคํ ภาเวนฺโต อริยํ อฏฺฐงฺคิกํ มคฺคํ พหุลีกโรนฺโต นิพฺพานนินฺโน โหติ นิพฺพานโปโณ นิพฺพาน\-ปพฺภาโรติ.  ฉฏฺฐํ.}

\csromanheader{2}{31. ปฐม\-สมุทฺทนินฺน\-สุตฺต}
\proseitem{133}{เสยฺยถาปิ ภิกฺขเว คงฺคา นที สมุทฺทนินฺนา สมุทฺทโปณา สมุทฺท\-ปพฺภารา. เอวเมว โข ภิกฺขเว ภิกฺขุ อริยํ อฏฺฐงฺคิกํ มคฺคํ ภาเวนฺโต อริยํ อฏฺฐงฺคิกํ มคฺคํ พหุลีกโรนฺโต นิพฺพานนินฺโน โหติ นิพฺพานโปโณ นิพฺพาน\-ปพฺภาโร. กถญฺจ ภิกฺขเว ภิกฺขุ อริยํ อฏฺฐงฺคิกํ มคฺคํ ภาเวนฺโต อริยํ อฏฺฐงฺคิกํ มคฺคํ พหุลีกโรนฺโต นิพฺพานนินฺโน โหติ นิพฺพานโปโณ นิพฺพาน\-ปพฺภาโร, อิธ ภิกฺขเว ภิกฺขุ สมฺมาทิฏฺฐิํ ภาเวติ นิพฺพานนินฺนํ นิพฺพานโปณํ นิพฺพาน\-ปพฺภารํ ฯเปฯ สมฺมาสมาธิํ ภาเวติ นิพฺพานนินฺนํ นิพฺพานโปณํ นิพฺพาน\-ปพฺภารํ. เอวํ โข ภิกฺขเว ภิกฺขุ อริยํ อฏฺฐงฺคิกํ มคฺคํ ภาเวนฺโต อริยํ อฏฺฐงฺคิกํ มคฺคํ พหุลีกโรนฺโต นิพฺพานนินฺโน โหติ นิพฺพานโปโณ นิพฺพาน\-ปพฺภาโรติ.  ปฐมํ.}

\csromanheader{2}{32-36. ทุติยา\-ทิ\-สมุทฺทนินฺน\-สุตฺต}
\proseitem{134-138}{เสยฺยถาปิ ภิกฺขเว ยมุนา นที สมุทฺทนินฺนา สมุทฺทโปณา สมุทฺท\-ปพฺภารา. เอวเมว โข ภิกฺขเว ภิกฺขุ ฯเปฯ. เสยฺยถาปิ ภิกฺขเว อจิรวตี นที สมุทฺทนินฺนา สมุทฺทโปณา สมุทฺท\-ปพฺภารา. เอวเมว โข ภิกฺขเว ภิกฺขุ ฯเปฯ. เสยฺยถาปิ ภิกฺขเว สรภู นที สมุทฺทนินฺนา สมุทฺทโปณา สมุทฺท\-ปพฺภารา. เอวเมว โข ภิกฺขเว ภิกฺขุ ฯเปฯ. เสยฺยถาปิ ภิกฺขเว มหี นที \csromanfolio{38}สมุทฺทนินฺนา สมุทฺทโปณา สมุทฺท\-ปพฺภารา. เอวเมว โข ภิกฺขเว ภิกฺขุ ฯเปฯ. เสยฺยถาปิ ภิกฺขเว ยา กาจิมา มหานทิโย. เสยฺยถิทํ, คงฺคา ยมุนา อจิรวตี สรภู มหี, สพฺพา ตา สมุทฺทนินฺนา สมุทฺทโปณา สมุทฺท\-ปพฺภารา. เอวเมว โข ภิกฺขเว ภิกฺขุ อริยํ อฏฺฐงฺคิกํ มคฺคํ ภาเวนฺโต อริยํ อฏฺฐงฺคิกํ มคฺคํ พหุลีกโรนฺโต นิพฺพานนินฺโน โหติ นิพฺพานโปโณ นิพฺพาน\-ปพฺภาโร. กถญฺจ ภิกฺขเว ภิกฺขุ อริยํ อฏฺฐงฺคิกํ มคฺคํ ภาเวนฺโต อริยํ อฏฺฐงฺคิกํ มคฺคํ พหุลีกโรนฺโต นิพฺพานนินฺโน โหติ นิพฺพานโปโณ นิพฺพาน\-ปพฺภาโร, อิธ ภิกฺขเว ภิกฺขุ สมฺมาทิฏฺฐิํ ภาเวติ นิพฺพานนินฺนํ นิพฺพานโปณํ นิพฺพาน\-ปพฺภารํ ฯเปฯ สมฺมาสมาธิํ ภาเวติ นิพฺพานนินฺนํ นิพฺพานโปณํ นิพฺพาน\-ปพฺภารํ. เอวํ โข ภิกฺขเว ภิกฺขุ อริยํ อฏฺฐงฺคิกํ มคฺคํ ภาเวนฺโต อริยํ อฏฺฐงฺคิกํ มคฺคํ พหุลีกโรนฺโต นิพฺพานนินฺโน โหติ นิพฺพานโปโณ นิพฺพาน\-ปพฺภาโรติ.  ฉฏฺฐํ. (คงฺคาเปยฺยาลี.)}

\vspace{\tipitakaparskip}
\csromancenter{\textbf{ตสฺสุทฺทานํ}}
\setlength{\csromangathapagewidth}{0pt}
\setlength{\csromangathawakwidth}{0pt}
\csromangathameasuregroup{ฉ ปาจีนโต นินฺนา, \\ เอเต เทฺว ฉ ทฺวาทส โหนฺติ,}{\csromangathabat{ฉ ปาจีนโต นินฺนา,}{ฉ นินฺนา จ สมุทฺทโต.} \\ \csromangathabat{เอเต เทฺว ฉ ทฺวาทส โหนฺติ,}{วคฺโค เตน ปวุจฺจตีติ.} \\ \textbf{นิพฺพานนินฺโน ทฺวาทสกี, จตุตฺถกี ฉฏฺฐา นวกี.}}
\csromangathameasurewak{\textbf{นิพฺพานนินฺโน ทฺวาทสกี, จตุตฺถกี ฉฏฺฐา นวกี.}}
\csromangathasetleft{ฉ ปาจีนโต นินฺนา, \\ เอเต เทฺว ฉ ทฺวาทส โหนฺติ,}
\csromangathagroup{\csromangathastanza{\csromangathabat{ฉ ปาจีนโต นินฺนา,}{ฉ นินฺนา จ สมุทฺทโต.} \\ \csromangathabat{เอเต เทฺว ฉ ทฺวาทส โหนฺติ,}{วคฺโค เตน ปวุจฺจตีติ.}}\csromangathastanza{\csromangathawak{\textbf{นิพฺพานนินฺโน ทฺวาทสกี, จตุตฺถกี ฉฏฺฐา นวกี.}}}}

\chapterhead{5. อปฺปมาท\-เปยฺยาล\-วคฺค}
\csromanheader{2}{1. \textbf{ตถาคตสุตฺต}}
\proseitem{139}{สาวตฺถิ\-นิทานํ. ยาวตา ภิกฺขเว สตฺตา อปทา วา ทฺวิปทา\csromansharedfootnote{0efc5171f6e27c02}{ทิปทา (สี)} วา จตุปฺปทา วา พหุปฺปทา\csromansharedfootnote{48c460fd12314009}{พหุปทา (?)} วา รูปิโน วา อรูปิโน วา สญฺญิโน วา อสญฺญิโน วา เนว\-สญฺญี\-นา\-สญฺญิโน วา, ตถาคโต เตสํ อคฺคมกฺขายติ อรหํ สมฺมาสมฺพุทฺโธ. เอวเมว โข ภิกฺขเว เย เกจิ กุสลา ธมฺมา สพฺเพ เต อปฺปมาทมูลกา อปฺปมาท\-สโมสรณา, อปฺปมาโท เตสํ ธมฺมานํ อคฺคมกฺขายติ. อปฺปมตฺตสฺเสตํ ภิกฺขเว ภิกฺขุโน ปาฏิกงฺขํ อริยํ อฏฺฐงฺคิกํ มคฺคํ ภาเวสฺสติ อริยํ อฏฺฐงฺคิกํ}

\prose{\csromanfolio{39}มคฺคํ พหุลีกริสฺสติ. กถญฺจ ภิกฺขเว ภิกฺขุ อปฺปมตฺโต อริยํ อฏฺฐงฺคิกํ มคฺคํ ภาเวติ อริยํ อฏฺฐงฺคิกํ มคฺคํ พหุลีกโรติ, อิธ ภิกฺขเว ภิกฺขุ สมฺมาทิฏฺฐิํ ภาเวติ วิเวกนิสฺสิตํ วิราคนิสฺสิตํ นิโรธ\-นิสฺสิตํ โวสฺสคฺค\-ปริณามิํ ฯเปฯ สมฺมาสมาธิํ ภาเวติ วิเวกนิสฺสิตํ วิราคนิสฺสิตํ นิโรธ\-นิสฺสิตํ โวสฺสคฺค\-ปริณามิํ. เอวํ โข ภิกฺขเว ภิกฺขุ อปฺปมตฺโต อริยํ อฏฺฐงฺคิกํ มคฺคํ ภาเวติ อริยํ อฏฺฐงฺคิกํ มคฺคํ พหุลีกโรตีติ.}

\prose{ยาวตา ภิกฺขเว สตฺตา อปทา วา ทฺวิปทา วา จตุปฺปทา วา พหุปฺปทา วา รูปิโน วา อรูปิโน วา สญฺญิโน วา อสญฺญิโน วา เนว\-สญฺญี\-นา\-สญฺญิโน วา, ตถาคโต เตสํ อคฺคมกฺขายติ อรหํ สมฺมาสมฺพุทฺโธ. เอวเมว โข ภิกฺขเว เย เกจิ กุสลา ธมฺมา สพฺเพ เต อปฺปมาทมูลกา อปฺปมาท\-สโมสรณา, อปฺปมาโท เตสํ ธมฺมานํ อคฺคมกฺขายติ. อปฺปมตฺตสฺเสตํ ภิกฺขเว ภิกฺขุโน ปาฏิกงฺขํ อริยํ อฏฺฐงฺคิกํ มคฺคํ ภาเวสฺสติ อริยํ อฏฺฐงฺคิกํ มคฺคํ พหุลีกริสฺสติ. กถญฺจ ภิกฺขเว ภิกฺขุ อปฺปมตฺโต อริยํ อฏฺฐงฺคิกํ มคฺคํ ภาเวติ, อริยํ อฏฺฐงฺคิกํ มคฺคํ พหุลีกโรติ, อิธ ภิกฺขเว ภิกฺขุ สมฺมาทิฏฺฐิํ ภาเวติ ราค\-วินย\-ปริโยสานํ โทส\-วินย\-ปริโยสานํ โมห\-วินย\-ปริโยสานํ ฯเปฯ สมฺมาสมาธิํ ภาเวติ ราค\-วินย\-ปริโยสานํ โทส\-วินย\-ปริโยสานํ โมห\-วินย\-ปริโยสานํ. เอวํ โข ภิกฺขเว ภิกฺขุ อปฺปมตฺโต อริยํ อฏฺฐงฺคิกํ มคฺคํ ภาเวติ อริยํ อฏฺฐงฺคิกํ มคฺคํ พหุลีกโรตีติ.}

\prose{ยาวตา ภิกฺขเว สตฺตา อปทา วา ทฺวิปทา วา จตุปฺปทา วา พหุปฺปทา วา รูปิโน วา อรูปิโน วา สญฺญิโน วา อสญฺญิโน วา เนว\-สญฺญี\-นา\-สญฺญิโน วา, ตถาคโต เตสํ อคฺคมกฺขายติ อรหํ สมฺมาสมฺพุทฺโธ. เอวเมว โข ภิกฺขเว เย เกจิ กุสลา ธมฺมา สพฺเพ เต อปฺปมาทมูลกา อปฺปมาท\-สโมสรณา, อปฺปมาโท เตสํ ธมฺมานํ อคฺคมกฺขายติ. อปฺปมตฺตสฺเสตํ ภิกฺขเว ภิกฺขุโน ปาฏิกงฺขํ อริยํ อฏฺฐงฺคิกํ มคฺคํ ภาเวสฺสติ อริยํ อฏฺฐงฺคิกํ มคฺคํ พหุลีกริสฺสติ. กถญฺจ ภิกฺขเว ภิกฺขุ อปฺปมตฺโต อริยํ อฏฺฐงฺคิกํ มคฺคํ ภาเวติ อริยํ อฏฺฐงฺคิกํ มคฺคํ พหุลีกโรติ, อิธ ภิกฺขเว ภิกฺขุ สมฺมาทิฏฺฐิํ ภาเวติ อมโตคธํ อมตปรายนํ อมต\-ปริโยสานํ ฯเปฯ สมฺมาสมาธิํ ภาเวติ อมโตคธํ อมตปรายนํ อมต\-ปริโยสานํ. เอวํ โข ภิกฺขเว}

\prose{\csromanfolio{40}ภิกฺขุ อปฺปมตฺโต อริยํ อฏฺฐงฺคิกํ มคฺคํ ภาเวติ อริยํ อฏฺฐงฺคิกํ มคฺคํ พหุลีกโรตีติ.}

\prose{ยาวตา ภิกฺขเว สตฺตา อปทา วา ทฺวิปทา วา จตุปฺปทา วา พหุปฺปทา วา รูปิโน วา อรูปิโน วา สญฺญิโน วา อสญฺญิโน วา เนว\-สญฺญี\-นา\-สญฺญิโน วา, ตถาคโต เตสํ อคฺคมกฺขายติ อรหํ สมฺมาสมฺพุทฺโธ. เอวเมว โข ภิกฺขเว เย เกจิ กุสลา ธมฺมา สพฺเพ เต อปฺปมาทมูลกา อปฺปมาท\-สโมสรณา, อปฺปมาโท เตสํ ธมฺมานํ อคฺคมกฺขายติ. อปฺปมตฺตสฺเสตํ ภิกฺขเว ภิกฺขุโน ปาฏิกงฺขํ อริยํ อฏฺฐงฺคิกํ มคฺคํ ภาเวสฺสติ อริยํ อฏฺฐงฺคิกํ มคฺคํ พหุลีกริสฺสติ. กถญฺจ ภิกฺขเว ภิกฺขุ อปฺปมตฺโต อริยํ อฏฺฐงฺคิกํ มคฺคํ ภาเวติ อริยํ อฏฺฐงฺคิกํ มคฺคํ พหุลีกโรติ, อิธ ภิกฺขเว ภิกฺขุ สมฺมาทิฏฺฐิํ ภาเวติ นิพฺพานนินฺนํ นิพฺพานโปณํ นิพฺพาน\-ปพฺภารํ ฯเปฯ สมฺมาสมาธิํ ภาเวติ นิพฺพานนินฺนํ นิพฺพานโปณํ นิพฺพาน\-ปพฺภารํ. เอวํ โข ภิกฺขเว ภิกฺขุ อปฺปมตฺโต อริยํ อฏฺฐงฺคิกํ มคฺคํ ภาเวติ อริยํ อฏฺฐงฺคิกํ มคฺคํ พหุลีกโรตีติ.  ปฐมํ.}

\csromanheader{2}{2. \textbf{ปทสุตฺต}}
\proseitem{140}{เสยฺยถาปิ ภิกฺขเว ยานิ กานิจิ ชงฺคลานํ ปาณานํ ปทชาตานิ สพฺพานิ ตานิ หตฺถิปเท สโมธานํ คจฺฉนฺติ, หตฺถิปทํ เตสํ อคฺคมกฺขายติ ยทิทํ มหนฺตตฺเตน. เอวเมว โข ภิกฺขเว เย เกจิ กุสลา ธมฺมา สพฺเพ เต อปฺปมาทมูลกา อปฺปมาท\-สโมสรณา, อปฺปมาโท เตสํ ธมฺมานํ อคฺคมกฺขายติ. อปฺปมตฺตสฺเสตํ ภิกฺขเว ภิกฺขุโน ปาฏิกงฺขํ อริยํ อฏฺฐงฺคิกํ มคฺคํ ภาเวสฺสติ อริยํ อฏฺฐงฺคิกํ มคฺคํ พหุลีกริสฺสติ. กถญฺจ ภิกฺขเว ภิกฺขุ อปฺปมตฺโต อริยํ อฏฺฐงฺคิกํ มคฺคํ ภาเวติ อริยํ อฏฺฐงฺคิกํ มคฺคํ พหุลีกโรติ, อิธ ภิกฺขเว ภิกฺขุ สมฺมาทิฏฺฐิํ ภาเวติ วิเวกนิสฺสิตํ วิราคนิสฺสิตํ นิโรธ\-นิสฺสิตํ โวสฺสคฺค\-ปริณามิํ ฯเปฯ สมฺมาสมาธิํ ภาเวติ วิเวกนิสฺสิตํ วิราคนิสฺสิตํ นิโรธ\-นิสฺสิตํ โวสฺสคฺค\-ปริณามิํ. เอวํ โข ภิกฺขเว ภิกฺขุ อปฺปมตฺโต อริยํ อฏฺฐงฺคิกํ มคฺคํ ภาเวติ อริยํ อฏฺฐงฺคิกํ มคฺคํ พหุลีกโรตีติ.  ทุติยํ.}

\csromanheader{2}{3-7. กูฏาทิสุตฺต}
\proseitem{141-145}{เสยฺยถาปิ ภิกฺขเว กูฏาคารสฺส ยา กาจิ โคปานสิโย สพฺพา ตา กูฏงฺคมา กูฏนินฺนา กูฏสโมสรณา, กูฏํ ตาสํ}

\prose{\csromanfolio{41}อคฺคมกฺขายติ. เอวเมว โข ภิกฺขเว. (ยถา เหฏฺฐิม\-สุตฺตนฺตํ, เอวํ วิตฺถาเรตพฺพํ.) เสยฺยถาปิ ภิกฺขเว เย เกจิ มูลคนฺธา, กาฬานุสาริยํ เตสํ อคฺคมกฺขายติ. เอวเมว โข ภิกฺขเว ฯเปฯ. เสยฺยถาปิ ภิกฺขเว เย เกจิ สารคนฺธา, โลหิตจนฺทนํ เตสํ อคฺคมกฺขายติ. เอวเมว โข ภิกฺขเว ฯเปฯ. เสยฺยถาปิ ภิกฺขเว เย เกจิ ปุปฺผคนฺธา, วสฺสิกํ เตสํ อคฺคมกฺขายติ. เอวเมว โข ภิกฺขเว ฯเปฯ. เสยฺยถาปิ ภิกฺขเว เย เกจิ กุฏฺฏราชาโน สพฺเพ เต รญฺโญ จกฺกวตฺติสฺส อนุยนฺตา ภวนฺติ, ราชา เตสํ จกฺกวตฺติ อคฺคมกฺขายติ. เอวเมว โข ภิกฺขเว ฯเปฯ.  สตฺตมํ.}

\csromanheader{2}{8-10. จนฺทิมาทิสุตฺต}
\proseitemwithclosermidruled{146-148}{เสยฺยถาปิ ภิกฺขเว ยา กาจิ ตารกรูปานํ ปภา, สพฺพา ตา จนฺทิม\-ปฺปภาย\csromansharedfootnote{2234edc2fa11ebba}{จนฺทิมาปภาย (สฺยา, ก)} กลํ นาคฺฆนฺติ โสฬสิํ, จนฺทปฺปภา ตาสํ อคฺคมกฺขายติ. เอวเมว โข ภิกฺขเว ฯเปฯ. เสยฺยถาปิ ภิกฺขเว สรทสมเย วิทฺเธ วิคตวลาหเก เทเว อาทิจฺโจ นภํ อพฺภุสฺสกฺกมาโน สพฺพํ อากาสคตํ ตมคตํ อภิวิหจฺจ ภาสเต จ ตปเต จ วิโรจติ จ. เอวเมว โข ภิกฺขเว ฯเปฯ. เสยฺยถาปิ ภิกฺขเว ยานิ กานิจิ ตนฺตาวุตานํ วตฺถานํ, กาสิกวตฺถํ เตสํ อคฺคมกฺขายติ. เอวเมว โข ภิกฺขเว เย เกจิ กุสลา ธมฺมา สพฺเพ เต อปฺปมาทมูลกา อปฺปมาท\-สโมสรณา, อปฺปมาโท เตสํ ธมฺมานํ อคฺคมกฺขายติ. อปฺปมตฺตสฺเสตํ ภิกฺขเว ภิกฺขุโน ปาฏิกงฺขํ อริยํ อฏฺฐงฺคิกํ มคฺคํ ภาเวสฺสติ อริยํ อฏฺฐงฺคิกํ มคฺคํ พหุลีกริสฺสติ. กถญฺจ ภิกฺขเว ภิกฺขุ อปฺปมตฺโต อริยํ อฏฺฐงฺคิกํ มคฺคํ ภาเวติ อริยํ อฏฺฐงฺคิกํ มคฺคํ พหุลีกโรติ, อิธ ภิกฺขเว ภิกฺขุ สมฺมาทิฏฺฐิํ ภาเวติ วิเวกนิสฺสิตํ วิราคนิสฺสิตํ นิโรธ\-นิสฺสิตํ โวสฺสคฺค\-ปริณามิํ ฯเปฯ สมฺมาสมาธิํ ภาเวติ วิเวกนิสฺสิตํ วิราคนิสฺสิตํ นิโรธ\-นิสฺสิตํ โวสฺสคฺค\-ปริณามิํ. เอวํ โข ภิกฺขเว ภิกฺขุ อปฺปมตฺโต อริยํ อฏฺฐงฺคิกํ มคฺคํ ภาเวติ อริยํ อฏฺฐงฺคิกํ มคฺคํ พหุลีกโรตีติ.  ทสมํ.}{(ยทปิ ตถาคตํ, ตทปิ วิตฺถาเรตพฺพํ.) อปฺปมาทวคฺโค ปญฺจโม.}
\csromancenter{\textbf{ตสฺสุทฺทานํ}}
\setlength{\csromangathapagewidth}{0pt}
\setlength{\csromangathawakwidth}{0pt}
\csromangathameasuregroup{ตถาคตํ ปทํ กูฏํ, \\ ราชา จนฺทิมสูริยา จ,}{\csromangathabat{ตถาคตํ ปทํ กูฏํ,}{มูลํ สาโร จ วสฺสิกํ.} \\ \csromangathabat{ราชา จนฺทิมสูริยา จ,}{วตฺเถน ทสมํ ปทํ.}}
\csromangathasetleft{ตถาคตํ ปทํ กูฏํ, \\ ราชา จนฺทิมสูริยา จ,}
\csromangathagroup{\csromangathastanza{\csromanfolio{42}\csromangathabat{ตถาคตํ ปทํ กูฏํ,}{มูลํ สาโร จ วสฺสิกํ.} \\ \csromangathabat{ราชา จนฺทิมสูริยา จ,}{วตฺเถน ทสมํ ปทํ.}}}

\csromanmatikaanchor{mtk.5}
\csromanmatikaanchor{mtk.27}
\csromantocmarkref{section}{9. รุกฺขสุตฺต}{mtk.5}
\bookmark[dest={mtk.5},level=1]{9. รุกฺขสุตฺต}
\csromanheader{4}{6. พลกรณีย\-วคฺค}
\csromanheader{2}{1. \textbf{พลสุตฺต}}
\proseitem{149}{สาวตฺถิ\-นิทานํ. เสยฺยถาปิ ภิกฺขเว เย เกจิ พลกรณียา กมฺมนฺตา กรียนฺติ, สพฺเพ เต ปถวิํ นิสฺสาย ปถวิยํ ปติฏฺฐาย เอวเมเต พลกรณียา กมฺมนฺตา กรียนฺติ. เอวเมว โข ภิกฺขเว ภิกฺขุ สีลํ นิสฺสาย สีเล ปติฏฺฐาย อริยํ อฏฺฐงฺคิกํ มคฺคํ ภาเวติ อริยํ อฏฺฐงฺคิกํ มคฺคํ พหุลีกโรติ. กถญฺจ ภิกฺขเว ภิกฺขุ สีลํ นิสฺสาย สีเล ปติฏฺฐาย อริยํ อฏฺฐงฺคิกํ มคฺคํ ภาเวติ อริยํ อฏฺฐงฺคิกํ มคฺคํ พหุลีกโรติ, อิธ ภิกฺขเว ภิกฺขุ สมฺมาทิฏฺฐิํ ภาเวติ วิเวกนิสฺสิตํ วิราคนิสฺสิตํ นิโรธ\-นิสฺสิตํ โวสฺสคฺค\-ปริณามิํ. เอวํ โข ภิกฺขเว ภิกฺขุ สีลํ นิสฺสาย สีเล ปติฏฺฐาย อริยํ อฏฺฐงฺคิกํ มคฺคํ ภาเวติ อริยํ อฏฺฐงฺคิกํ มคฺคํ พหุลีกโรตีติ.}

\vspace{\tipitakaparskip}
\csromancenter{(ปรคงฺคา\-เปยฺยาลี\-วณฺณิยโต ปริปุณฺณสุตฺตนฺติ วิตฺถารมคฺคี.)}
\prose{เสยฺยถาปิ ภิกฺขเว เย เกจิ พลกรณียา กมฺมนฺตา กรียนฺติ, สพฺเพ เต ปถวิํ นิสฺสาย ปถวิยํ ปติฏฺฐาย เอวเมเต พลกรณียา กมฺมนฺตา กรียนฺติ. เอวเมว โข ภิกฺขเว ภิกฺขุ สีลํ นิสฺสาย สีเล ปติฏฺฐาย อริยํ อฏฺฐงฺคิกํ มคฺคํ ภาเวติ อริยํ อฏฺฐงฺคิกํ มคฺคํ พหุลีกโรติ. กถญฺจ ภิกฺขเว ภิกฺขุ สีลํ นิสฺสาย สีเล ปติฏฺฐาย อริยํ อฏฺฐงฺคิกํ มคฺคํ ภาเวติ อริยํ อฏฺฐงฺคิกํ มคฺคํ พหุลีกโรติ, อิธ ภิกฺขเว ภิกฺขุ สมฺมาทิฏฺฐิํ ภาเวติ ราค\-วินย\-ปริโยสานํ โทส\-วินย\-ปริโยสานํ โมห\-วินย\-ปริโยสานํ ฯเปฯ สมฺมาสมาธิํ ภาเวติ ราค\-วินย\-ปริโยสานํ โทส\-วินย\-ปริโยสานํ โมห\-วินย\-ปริโยสานํ. เอวํ โข ภิกฺขเว ภิกฺขุ สีลํ นิสฺสาย สีเล ปติฏฺฐาย อริยํ อฏฺฐงฺคิกํ มคฺคํ ภาเวติ อริยํ อฏฺฐงฺคิกํ มคฺคํ พหุลีกโรตีติ.}

\prose{\csromanfolio{43}เสยฺยถาปิ ภิกฺขเว เย เกจิ พลกรณียา กมฺมนฺตา กรียนฺติ, สพฺเพ เต ปถวิํ นิสฺสาย ปถวิยํ ปติฏฺฐาย เอวเมเต พลกรณียา กมฺมนฺตา กรียนฺติ. เอวเมว โข ภิกฺขเว ภิกฺขุ สีลํ นิสฺสาย สีเล ปติฏฺฐาย อริยํ อฏฺฐงฺคิกํ มคฺคํ ภาเวติ อริยํ อฏฺฐงฺคิกํ มคฺคํ พหุลีกโรติ. กถญฺจ ภิกฺขเว ภิกฺขุ สีลํ นิสฺสาย สีเล ปติฏฺฐาย อริยํ อฏฺฐงฺคิกํ มคฺคํ ภาเวติ อริยํ อฏฺฐงฺคิกํ มคฺคํ พหุลีกโรติ. อิธ ภิกฺขเว ภิกฺขุ สมฺมาทิฏฺฐิํ ภาเวติ อมโตคธํ อมตปรายนํ อมต\-ปริโยสานํ ฯเปฯ สมฺมาสมาธิํ ภาเวติ อมโตคธํ อมตปรายนํ อมต\-ปริโยสานํ. เอวํ โข ภิกฺขเว ภิกฺขุ สีลํ นิสฺสาย สีเล ปติฏฺฐาย อริยํ อฏฺฐงฺคิกํ มคฺคํ ภาเวติ อริยํ อฏฺฐงฺคิกํ มคฺคํ พหุลีกโรตีติ.}

\prose{เสยฺยถาปิ ภิกฺขเว เย เกจิ พลกรณียา กมฺมนฺตา กรียนฺติ, สพฺเพ เต ปถวิํ นิสฺสาย ปถวิยํ ปติฏฺฐาย เอวเมเต พลกรณียา กมฺมนฺตา กรียนฺติ. เอวเมว โข ภิกฺขเว ภิกฺขุ สีลํ นิสฺสาย สีเล ปติฏฺฐาย อริยํ อฏฺฐงฺคิกํ มคฺคํ ภาเวติ อริยํ อฏฺฐงฺคิกํ มคฺคํ พหุลีกโรติ. กถญฺจ ภิกฺขเว ภิกฺขุ สีลํ นิสฺสาย สีเล ปติฏฺฐาย อริยํ อฏฺฐงฺคิกํ มคฺคํ ภาเวติ อริยํ อฏฺฐงฺคิกํ มคฺคํ พหุลีกโรติ, อิธ ภิกฺขเว ภิกฺขุ สมฺมาทิฏฺฐิํ ภาเวติ นิพฺพานนินฺนํ นิพฺพานโปณํ นิพฺพาน\-ปพฺภารํ ฯเปฯ สมฺมาสมาธิํ ภาเวติ นิพฺพานนินฺนํ นิพฺพานโปณํ นิพฺพาน\-ปพฺภารํ. เอวํ โข ภิกฺขเว ภิกฺขุ สีลํ นิสฺสาย สีเล ปติฏฺฐาย อริยํ อฏฺฐงฺคิกํ มคฺคํ ภาเวติ อริยํ อฏฺฐงฺคิกํ มคฺคํ พหุลีกโรตีติ. ปฐมํ.}

\csromanheader{2}{2. \textbf{พีชสุตฺต}}
\proseitem{150}{เสยฺยถาปิ ภิกฺขเว เย เกจิเม พีชคาม\-ภูตคามา วุฑฺฒิํ วิรูฬฺหิํ เวปุลฺลํ อาปชฺชนฺติ, สพฺเพ เต ปถวิํ นิสฺสาย ปถวิยํ ปติฏฺฐาย เอวเมเต พีชคาม\-ภูตคามา วุฑฺฒิํ วิรูฬฺหิํ เวปุลฺลํ อาปชฺชนฺติ. เอวเมว โข ภิกฺขเว ภิกฺขุ สีลํ นิสฺสาย สีเล ปติฏฺฐาย อริยํ อฏฺฐงฺคิกํ มคฺคํ ภาเวนฺโต อริยํ อฏฺฐงฺคิกํ มคฺคํ พหุลีกโรนฺโต วุฑฺฒิํ วิรูฬฺหิํ เวปุลฺลํ ปาปุณาติ ธมฺเมสุ. กถญฺจ ภิกฺขเว ภิกฺขุ สีลํ นิสฺสาย สีเล ปติฏฺฐาย อริยํ อฏฺฐงฺคิกํ มคฺคํ ภาเวนฺโต อริยํ อฏฺฐงฺคิกํ มคฺคํ พหุลีกโรนฺโต \csromanfolio{44}วุฑฺฒิํ วิรูฬฺหิํ เวปุลฺลํ ปาปุณาติ ธมฺเมสุ, อิธ ภิกฺขเว ภิกฺขุ สมฺมาทิฏฺฐิํ ภาเวติ วิเวกนิสฺสิตํ ฯเปฯ สมฺมาสมาธิํ ภาเวติ วิเวกนิสฺสิตํ วิราคนิสฺสิตํ นิโรธ\-นิสฺสิตํ โวสฺสคฺค\-ปริณามิํ. เอวํ โข ภิกฺขเว ภิกฺขุ สีลํ นิสฺสาย สีเล ปติฏฺฐาย อริยํ อฏฺฐงฺคิกํ มคฺคํ ภาเวนฺโต อริยํ อฏฺฐงฺคิกํ มคฺคํ พหุลีกโรนฺโต วุฑฺฒิํ วิรูฬฺหิํ เวปุลฺลํ ปาปุณาติ ธมฺเมสูติ.  ทุติยํ.}

\csromanheader{2}{3. \textbf{นาคสุตฺต}}
\proseitem{151}{เสยฺยถาปิ ภิกฺขเว หิมวนฺตํ ปพฺพตราชํ นิสฺสาย นาคา กายํ วฑฺเฒนฺติ, พลํ คาเหนฺติ. เต ตตฺถ กายํ วฑฺเฒตฺวา พลํ คาเหตฺวา กุโสพฺเภ โอตรนฺติ, กุโสพฺเภ\csromansharedfootnote{a495e922962eba3b}{กุสฺสุพฺเภ (สี, สฺยา), กุสุพฺเภ (อิ, ก)} โอตริตฺวา มหาโสพฺเภ โอตรนฺติ, มหาโสพฺเภ โอตริตฺวา กุนฺนทิโย โอตรนฺติ, กุนฺนทิโย โอตริตฺวา มหานทิโย โอตรนฺติ, มหานทิโย โอตริตฺวา มหาสมุทฺทํ\csromansharedfootnote{dfca408eecbd6db5}{มหาสมุทฺทสาครํ (สพฺพตฺถ) สํ 1. 269 ปิฏฺเฐปิ.} โอตรนฺติ. เต ตตฺถ มหนฺตตฺตํ เวปุลฺลตฺตํ อาปชฺชนฺติ กาเยน. เอวเมว โข ภิกฺขเว ภิกฺขุ สีลํ นิสฺสาย สีเล ปติฏฺฐาย อริยํ อฏฺฐงฺคิกํ มคฺคํ ภาเวนฺโต อริยํ อฏฺฐงฺคิกํ มคฺคํ พหุลีกโรนฺโต มหนฺตตฺตํ เวปุลฺลตฺตํ ปาปุณาติ ธมฺเมสุ. กถญฺจ ภิกฺขเว ภิกฺขุ สีลํ นิสฺสาย สีเล ปติฏฺฐาย อริยํ อฏฺฐงฺคิกํ มคฺคํ ภาเวนฺโต อริยํ อฏฺฐงฺคิกํ มคฺคํ พหุลีกโรนฺโต มหนฺตตฺตํ เวปุลฺลตฺตํ ปาปุณาติ ธมฺเมสุ, อิธ ภิกฺขเว ภิกฺขุ สมฺมาทิฏฺฐิํ ภาเวติ วิเวกนิสฺสิตํ วิราคนิสฺสิตํ นิโรธ\-นิสฺสิตํ โวสฺสคฺค\-ปริณามิํ ฯเปฯ สมฺมาสมาธิํ ภาเวติ วิเวกนิสฺสิตํ วิราคนิสฺสิตํ นิโรธ\-นิสฺสิตํ โวสฺสคฺค\-ปริณามิํ. เอวํ โข ภิกฺขเว ภิกฺขุ สีลํ นิสฺสาย สีเล ปติฏฺฐาย อริยํ อฏฺฐงฺคิกํ มคฺคํ ภาเวนฺโต อริยํ อฏฺฐงฺคิกํ มคฺคํ พหุลีกโรนฺโต มหนฺตตฺตํ เวปุลฺลตฺตํ ปาปุณาติ ธมฺเมสูติ.  ตติยํ.}

\csromanheader{2}{4. \textbf{รุกฺขสุตฺต}}
\proseitem{152}{เสยฺยถาปิ ภิกฺขเว รุกฺโข ปาจีนนินฺโน ปาจีนโปโณ ปาจีนปพฺภาโร, โส มูลจฺฉินฺโน\csromansharedfootnote{bd99f6a617705d96}{มูลจฺฉินฺเท กเต (สฺยา)} กตเมน ปปเตยฺยาติ. เยน ภนฺเต นินฺโน เยน โปโณ เยน ปพฺภาโรติ. เอวเมว โข ภิกฺขเว ภิกฺขุ อริยํ อฏฺฐงฺคิกํ มคฺคํ ภาเวนฺโต อริยํ อฏฺฐงฺคิกํ มคฺคํ พหุลีกโรนฺโต}

\prose{\csromanfolio{45}นิพฺพานนินฺโน โหติ นิพฺพานโปโณ นิพฺพาน\-ปพฺภาโร. กถญฺจ ภิกฺขเว ภิกฺขุ อริยํ อฏฺฐงฺคิกํ มคฺคํ ภาเวนฺโต อริยํ อฏฺฐงฺคิกํ มคฺคํ พหุลีกโรนฺโต นิพฺพานนินฺโน โหติ นิพฺพานโปโณ นิพฺพาน\-ปพฺภาโร, อิธ ภิกฺขเว ภิกฺขุ สมฺมาทิฏฺฐิํ ภาเวติ วิเวกนิสฺสิตํ วิราคนิสฺสิตํ นิโรธ\-นิสฺสิตํ โวสฺสคฺค\-ปริณามิํ ฯเปฯ สมฺมาสมาธิํ ภาเวติ วิเวกนิสฺสิตํ วิราคนิสฺสิตํ นิโรธ\-นิสฺสิตํ โวสฺสคฺค\-ปริณามิํ. เอวํ โข ภิกฺขเว ภิกฺขุ อริยํ อฏฺฐงฺคิกํ มคฺคํ ภาเวนฺโต อริยํ อฏฺฐงฺคิกํ มคฺคํ พหุลีกโรนฺโต นิพฺพานนินฺโน โหติ นิพฺพานโปโณ นิพฺพาน\-ปพฺภาโรติ.  จตุตฺถํ.}

\csromanheader{2}{5. \textbf{กุมฺภสุตฺต}}
\proseitem{153}{เสยฺยถาปิ ภิกฺขเว กุมฺโภ นิกฺกุชฺโช วมเตว อุทกํ โน ปจฺจาวมติ. เอวเมว โข ภิกฺขเว ภิกฺขุ อริยํ อฏฺฐงฺคิกํ มคฺคํ ภาเวนฺโต อริยํ อฏฺฐงฺคิกํ มคฺคํ พหุลีกโรนฺโต วมเตว ปาปเก อกุสเล ธมฺเม โน ปจฺจาวมติ. กถญฺจ ภิกฺขเว ภิกฺขุ อริยํ อฏฺฐงฺคิกํ มคฺคํ ภาเวนฺโต อริยํ อฏฺฐงฺคิกํ มคฺคํ พหุลีกโรนฺโต วมเตว ปาปเก อกุสเล ธมฺเม โน ปจฺจาวมติ, อิธ ภิกฺขเว ภิกฺขุ สมฺมาทิฏฺฐิํ ภาเวติ วิเวกนิสฺสิตํ วิราคนิสฺสิตํ นิโรธ\-นิสฺสิตํ โวสฺสคฺค\-ปริณามิํ ฯเปฯ สมฺมาสมาธิํ ภาเวติ วิเวกนิสฺสิตํ วิราคนิสฺสิตํ นิโรธ\-นิสฺสิตํ โวสฺสคฺค\-ปริณามิํ. เอวํ โข ภิกฺขเว ภิกฺขุ อริยํ อฏฺฐงฺคิกํ มคฺคํ ภาเวนฺโต อริยํ อฏฺฐงฺคิกํ มคฺคํ พหุลีกโรนฺโต วมเตว ปาปเก อกุสเล ธมฺเม โน ปจฺจาวมตีติ.  ปญฺจมํ.}

\csromanheader{2}{6. \textbf{สูกสุตฺต}}
\proseitem{154}{เสยฺยถาปิ ภิกฺขเว สาลิสูกํ วา ยวสูกํ วา สมฺมาปณิหิตํ หตฺเถน วา ปาเทน วา อกฺกนฺตํ หตฺถํ วา ปาทํ วา ภินฺทิสฺสติ โลหิตํ วา อุปฺปาเทสฺสตีติ ฐานเมตํ วิชฺชติ. ตํ กิสฺส เหตุ, สมฺมา\-ปณิหิตตฺตา ภิกฺขเว สูกสฺส. เอวเมว โข ภิกฺขเว ภิกฺขุ สมฺมา\-ปณิหิตาย ทิฏฺฐิยา สมฺมา\-ปณิหิตาย มคฺคภาวนาย อวิชฺชํ ภินฺทิสฺสติ วิชฺชํ อุปฺปาเทสฺสติ นิพฺพานํ สจฺฉิกริสฺสตีติ ฐานเมตํ วิชฺชติ. ตํ กิสฺส เหตุ, สมฺมา\-ปณิหิตตฺตา ภิกฺขเว ทิฏฺฐิยา. กถญฺจ ภิกฺขเว ภิกฺขุ สมฺมา\-ปณิหิตาย ทิฏฺฐิยา สมฺมา\-ปณิหิตาย มคฺคภาวนาย อวิชฺชํ ภินฺทติ วิชฺชํ อุปฺปาเทติ นิพฺพานํ สจฺฉิกโรติ, อิธ ภิกฺขเว ภิกฺขุ สมฺมาทิฏฺฐิํ ภาเวติ วิเวกนิสฺสิตํ ฯเปฯ \csromanfolio{46}สมฺมาสมาธิํ ภาเวติ วิเวกนิสฺสิตํ วิราคนิสฺสิตํ นิโรธ\-นิสฺสิตํ โวสฺสคฺค\-ปริณามิํ. เอวํ โข ภิกฺขเว ภิกฺขุ สมฺมา\-ปณิหิตาย ทิฏฺฐิยา สมฺมา\-ปณิหิตาย มคฺคภาวนาย อวิชฺชํ ภินฺทติ วิชฺชํ อุปฺปาเทติ นิพฺพานํ สจฺฉิกโรตีติ.  ฉฏฺฐํ.}

\csromanheader{2}{7. \textbf{อากาสสุตฺต}}
\proseitem{155}{เสยฺยถาปิ ภิกฺขเว อากาเส วิวิธา วาตา วายนฺติ, ปุรตฺถิมาปิ วาตา วายนฺติ, ปจฺฉิมาปิ วาตา วายนฺติ, อุตฺตราปิ วาตา วายนฺติ, ทกฺขิณาปิ วาตา วายนฺติ, สรชาปิ วาตา วายนฺติ, อรชาปิ วาตา วายนฺติ, สีตาปิ วาตา วายนฺติ, อุณฺหาปิ วาตา วายนฺติ, ปริตฺตาปิ วาตา วายนฺติ, อธิมตฺตาปิ วาตา วายนฺติ. เอวเมว โข ภิกฺขเว ภิกฺขุโน อริยํ อฏฺฐงฺคิกํ มคฺคํ ภาวยโต อริยํ อฏฺฐงฺคิกํ มคฺคํ พหุลีกโรโต จตฺตาโรปิ สติปฏฺฐานา ภาวนา\-ปาริปูริํ คจฺฉนฺติ, จตฺตาโรปิ สมฺมปฺปธานา ภาวนา\-ปาริปูริํ คจฺฉนฺติ, จตฺตาโรปิ อิทฺธิปาทา ภาวนา\-ปาริปูริํ คจฺฉนฺติ, ปญฺจปิ อินฺทฺริยานิ ภาวนา\-ปาริปูริํ คจฺฉนฺติ, ปญฺจปิ พลานิ ภาวนา\-ปาริปูริํ คจฺฉนฺติ, สตฺตปิ โพชฺฌงฺคา ภาวนา\-ปาริปูริํ คจฺฉนฺติ. กถญฺจ ภิกฺขเว ภิกฺขุโน อริยํ อฏฺฐงฺคิกํ มคฺคํ ภาวยโต อริยํ อฏฺฐงฺคิกํ มคฺคํ พหุลีกโรโต จตฺตาโรปิ สติปฏฺฐานา ภาวนา\-ปาริปูริํ คจฺฉนฺติ, จตฺตาโรปิ สมฺมปฺปธานา ภาวนา\-ปาริปูริํ คจฺฉนฺติ, จตฺตาโรปิ อิทฺธิปาทา ภาวนา\-ปาริปูริํ คจฺฉนฺติ, ปญฺจปิ อินฺทฺริยานิ ภาวนา\-ปาริปูริํ คจฺฉนฺติ, ปญฺจปิ พลานิ ภาวนา\-ปาริปูริํ คจฺฉนฺติ, สตฺตปิ โพชฺฌงฺคา ภาวนา\-ปาริปูริํ คจฺฉนฺติ, อิธ ภิกฺขเว ภิกฺขุ สมฺมาทิฏฺฐิํ ภาเวติ ฯเปฯ สมฺมาสมาธิํ ภาเวติ วิเวกนิสฺสิตํ วิราคนิสฺสิตํ นิโรธ\-นิสฺสิตํ โวสฺสคฺค\-ปริณามิํ. เอวํ โข ภิกฺขเว ภิกฺขุโน อริยํ อฏฺฐงฺคิกํ มคฺคํ ภาวยโต อริยํ อฏฺฐงฺคิกํ มคฺคํ พหุลีกโรโต จตฺตาโรปิ สติปฏฺฐานา ภาวนา\-ปาริปูริํ คจฺฉนฺติ, จตฺตาโรปิ สมฺมปฺปธานา ภาวนา\-ปาริปูริํ คจฺฉนฺติ, จตฺตาโรปิ อิทฺธิปาทา ภาวนา\-ปาริปูริํ คจฺฉนฺติ, ปญฺจปิ อินฺทฺริยานิ ภาวนา\-ปาริปูริํ คจฺฉนฺติ, ปญฺจปิ พลานิ ภาวนา\-ปาริปูริํ คจฺฉนฺติ, สตฺตปิ โพชฺฌงฺคา ภาวนา\-ปาริปูริํ คจฺฉนฺตีติ.  สตฺตมํ.}

\csromanheader{2}{8. ปฐม\-เมฆ\-สุตฺต}
\proseitem{156}{เสยฺยถาปิ ภิกฺขเว คิมฺหานํ ปจฺฉิเม มาเส อูหตํ รโชชลฺลํ, ตเมนํ มหาอกาลเมโฆ ฐานโส อนฺตรธาเปติ วูปสเมติ.}

\prose{\csromanfolio{47}เอวเมว โข ภิกฺขเว ภิกฺขุ อริยํ อฏฺฐงฺคิกํ มคฺคํ ภาเวนฺโต อริยํ อฏฺฐงฺคิกํ มคฺคํ พหุลีกโรนฺโต อุปฺปนฺนุปฺปนฺเน ปาปเก อกุสเล ธมฺเม ฐานโส อนฺตรธาเปติ วูปสเมติ. กถญฺจ ภิกฺขเว ภิกฺขุ อริยํ อฏฺฐงฺคิกํ มคฺคํ ภาเวนฺโต อริยํ อฏฺฐงฺคิกํ มคฺคํ พหุลีกโรนฺโต อุปฺปนฺนุปฺปนฺเน ปาปเก อกุสเล ธมฺเม ฐานโส อนฺตรธาเปติ วูปสเมติ, อิธ ภิกฺขเว ภิกฺขุ สมฺมาทิฏฺฐิํ ภาเวติ ฯเปฯ สมฺมาสมาธิํ ภาเวติ วิเวกนิสฺสิตํ วิราคนิสฺสิตํ นิโรธ\-นิสฺสิตํ โวสฺสคฺค\-ปริณามิํ. เอวํ โข ภิกฺขเว ภิกฺขุ อริยํ อฏฺฐงฺคิกํ มคฺคํ ภาเวนฺโต อริยํ อฏฺฐงฺคิกํ มคฺคํ พหุลีกโรนฺโต อุปฺปนฺนุปฺปนฺเน ปาปเก อกุสเล ธมฺเม ฐานโส อนฺตรธาเปติ วูปสเมตีติ.  อฏฺฐมํ.}

\csromanheader{2}{9. ทุติย\-เมฆ\-สุตฺต}
\proseitem{157}{เสยฺยถาปิ ภิกฺขเว อุปฺปนฺนํ มหาเมฆํ, ตเมนํ มหาวาโต อนฺตราเยว อนฺตรธาเปติ วูปสเมติ. เอวเมว โข ภิกฺขเว ภิกฺขุ อริยํ อฏฺฐงฺคิกํ มคฺคํ ภาเวนฺโต อริยํ อฏฺฐงฺคิกํ มคฺคํ พหุลีกโรนฺโต อุปฺปนฺนุปฺปนฺเน ปาปเก อกุสเล ธมฺเม อนฺตราเยว อนฺตรธาเปติ วูปสเมติ. กถญฺจ ภิกฺขเว ภิกฺขุ อริยํ อฏฺฐงฺคิกํ มคฺคํ ภาเวนฺโต อริยํ อฏฺฐงฺคิกํ มคฺคํ พหุลีกโรนฺโต อุปฺปนฺนุปฺปนฺเน ปาปเก อกุสเล ธมฺเม อนฺตราเยว อนฺตรธาเปติ วูปสเมติ, อิธ ภิกฺขเว ภิกฺขุ สมฺมาทิฏฺฐิํ ภาเวติ ฯเปฯ สมฺมาสมาธิํ ภาเวติ วิเวกนิสฺสิตํ วิราคนิสฺสิตํ นิโรธ\-นิสฺสิตํ โวสฺสคฺค\-ปริณามิํ. เอวํ โข ภิกฺขเว ภิกฺขุ อริยํ อฏฺฐงฺคิกํ มคฺคํ ภาเวนฺโต อริยํ อฏฺฐงฺคิกํ มคฺคํ พหุลีกโรนฺโต อุปฺปนฺนุปฺปนฺเน ปาปเก อกุสเล ธมฺเม อนฺตราเยว อนฺตรธาเปติ วูปสเมตีติ.  นวมํ.}

\csromanheader{2}{10. \textbf{นาวาสุตฺต}}
\proseitem{158}{เสยฺยถาปิ ภิกฺขเว สามุทฺทิกาย นาวาย เวตฺตพนฺธนพนฺธาย ฉ มาสานิ อุทเก ปริยาทาย\csromansharedfootnote{10fbf0efc4b5037c}{ปริยาตาย (ก), ปริยาหตาย (?)} เหมนฺติเกน ถลํ อุกฺขิตฺตาย วาตาตป\-ปเรตานิ พนฺธนานิ ตานิ ปาวุสฺสเกน เมเฆน อภิปฺปวุฏฺฐานิ อปฺปกสิเรเนว ปฏิปฺปสฺสมฺภนฺติ ปูติกานิ ภวนฺติ. เอวเมว โข ภิกฺขเว ภิกฺขุโน อริยํ อฏฺฐงฺคิกํ มคฺคํ ภาวยโต อริยํ อฏฺฐงฺคิกํ มคฺคํ พหุลีกโรโต อปฺปกสิเรเนว สํโยชนานิ ปฏิปฺปสฺสมฺภนฺติ ปูติกานิ \csromanfolio{48}ภวนฺติ. กถญฺจ ภิกฺขเว ภิกฺขุโน อริยํ อฏฺฐงฺคิกํ มคฺคํ ภาวยโต อริยํ อฏฺฐงฺคิกํ มคฺคํ พหุลีกโรโต อปฺปกสิเรเนว สํโยชนานิ ปฏิปฺปสฺสมฺภนฺติ ปูติกานิ ภวนฺติ, อิธ ภิกฺขเว ภิกฺขุ สมฺมาทิฏฺฐิํ ภาเวติ ฯเปฯ สมฺมาสมาธิํ ภาเวติ วิเวกนิสฺสิตํ วิราคนิสฺสิตํ นิโรธ\-นิสฺสิตํ โวสฺสคฺค\-ปริณามิํ. เอวํ โข ภิกฺขเว ภิกฺขุโน อริยํ อฏฺฐงฺคิกํ มคฺคํ ภาวยโต อริยํ อฏฺฐงฺคิกํ มคฺคํ พหุลีกโรโต อปฺปกสิเรเนว สํโยชนานิ ปฏิปฺปสฺสมฺภนฺติ ปูติกานิ ภวนฺตีติ.  ทสมํ.}

\csromanheader{2}{11. \textbf{อาคนฺตุกสุตฺต}}
\proseitem{159}{เสยฺยถาปิ ภิกฺขเว อาคนฺตุกาคารํ, ตตฺถ ปุรตฺถิมายปิ ทิสาย อาคนฺตฺวา วาสํ กปฺเปนฺติ, ปจฺฉิมายปิ ทิสาย อาคนฺตฺวา วาสํ กปฺเปนฺติ, อุตฺตรายปิ ทิสาย อาคนฺตฺวา วาสํ กปฺเปนฺติ, ทกฺขิณายปิ ทิสาย อาคนฺตฺวา วาสํ กปฺเปนฺติ, ขตฺติยาปิ อาคนฺตฺวา วาสํ กปฺเปนฺติ, พฺราหฺมณาปิ อาคนฺตฺวา วาสํ กปฺเปนฺติ, เวสฺสาปิ อาคนฺตฺวา วาสํ กปฺเปนฺติ, สุทฺทาปิ อาคนฺตฺวา วาสํ กปฺเปนฺติ. เอวเมว โข ภิกฺขเว ภิกฺขุ อริยํ อฏฺฐงฺคิกํ มคฺคํ ภาเวนฺโต อริยํ อฏฺฐงฺคิกํ มคฺคํ พหุลีกโรนฺโต เย ธมฺมา อภิญฺญา ปริญฺเญยฺยา, เต ธมฺเม อภิญฺญา ปริชานาติ ฯเปฯ เย ธมฺมา อภิญฺญา ปหาตพฺพา, เต ธมฺเม อภิญฺญา ปชหติ. เย ธมฺมา อภิญฺญา สจฺฉิกาตพฺพา, เต ธมฺเม อภิญฺญา สจฺฉิกโรติ. เย ธมฺมา อภิญฺญา ภาเวตพฺพา, เต ธมฺเม อภิญฺญา ภาเวติ.}

\prose{กตเม จ ภิกฺขเว ธมฺมา อภิญฺญา ปริญฺเญยฺยา, ปญฺจุ\-ปาทาน\-กฺขนฺ\-ธาติสฺส วจนียํ. กตเม ปญฺจ, เสยฺยถิทํ, รูปุปาทาน\-กฺขนฺโธ ฯเปฯ วิญฺญาณุ\-ปาทาน\-กฺขนฺโธ, อิเม ภิกฺขเว ธมฺมา อภิญฺญา ปริญฺเญยฺยา. กตเม จ ภิกฺขเว ธมฺมา อภิญฺญา ปหาตพฺพา, อวิชฺชา จ ภวตณฺหา จ, อิเม ภิกฺขเว ธมฺมา อภิญฺญา ปหาตพฺพา. กตเม จ ภิกฺขเว ธมฺมา อภิญฺญา สจฺฉิกาตพฺพา, วิชฺชา จ วิมุตฺติ จ, อิเม ภิกฺขเว ธมฺมา อภิญฺญา สจฺฉิกาตพฺพา. กตเม จ ภิกฺขเว ธมฺมา อภิญฺญา ภาเวตพฺพา, สมโถ จ วิปสฺสนา จ, อิเม ภิกฺขเว ธมฺมา อภิญฺญา ภาเวตพฺพา. กถญฺจ ภิกฺขเว ภิกฺขุ อริยํ อฏฺฐงฺคิกํ มคฺคํ ภาเวนฺโต อริยํ อฏฺฐงฺคิกํ มคฺคํ พหุลีกโรนฺโต เย ธมฺมา อภิญฺญา ปริญฺเญยฺยา, เต ธมฺเม อภิญฺญา ปริชานาติ.}

\prose{\csromanfolio{49}เย ธมฺมา อภิญฺญา ภาเวตพฺพา, เต ธมฺเม อภิญฺญา ภาเวติ, อิธ ภิกฺขเว ภิกฺขุ สมฺมาทิฏฺฐิํ ภาเวติ ฯเปฯ สมฺมาสมาธิํ ภาเวติ วิเวกนิสฺสิตํ วิราคนิสฺสิตํ นิโรธ\-นิสฺสิตํ โวสฺสคฺค\-ปริณามิํ. เอวํ โข ภิกฺขเว ภิกฺขุ อริยํ อฏฺฐงฺคิกํ มคฺคํ ภาเวนฺโต อริยํ อฏฺฐงฺคิกํ มคฺคํ พหุลีกโรนฺโต เย ธมฺมา อภิญฺญา ปริญฺเญยฺยา, เต ธมฺเม อภิญฺญา ปริชานาติ. เย ธมฺมา อภิญฺญา ปหาตพฺพา, เต ธมฺเม อภิญฺญา ปชหติ. เย ธมฺมา อภิญฺญา สจฺฉิกาตพฺพา, เต ธมฺเม อภิญฺญา สจฺฉิกโรติ. เย ธมฺมา อภิญฺญา ภาเวตพฺพา, เต ธมฺเม อภิญฺญา ภาเวตีติ.  เอกาทสมํ.}

\csromanheader{2}{12. \textbf{นทีสุตฺต}}
\proseitem{160}{เสยฺยถาปิ ภิกฺขเว คงฺคา นที ปาจีนนินฺนา ปาจีนโปณา ปาจีนปพฺภารา, อถ มหาชนกาโย อาคจฺเฉยฺย กุทาล ปิฏกํ อาทาย มยํ อิมํ คงฺคํ นทิํ ปจฺฉานินฺนํ กริสฺสาม ปจฺฉาโปณํ ปจฺฉา\-ปพฺภารนฺติ. ตํ กิํ มญฺญถ ภิกฺขเว, อปิ นุ โส มหาชนกาโย คงฺคํ นทิํ ปจฺฉานินฺนํ กเรยฺย ปจฺฉาโปณํ ปจฺฉา\-ปพฺภารนฺติ. โน เหตํ ภนฺเต. ตํ กิสฺส เหตุ, คงฺคา ภนฺเต นที ปาจีนนินฺนา ปาจีนโปณา ปาจีนปพฺภารา. สา น สุกรา ปจฺฉานินฺนํ กาตุํ ปจฺฉาโปณํ ปจฺฉา\-ปพฺภารํ, ยาวเทว ปน โส มหาชนกาโย กิลมถสฺส วิฆาตสฺส ภาคี อสฺสาติ. เอวเมว โข ภิกฺขเว ภิกฺขุํ อริยํ อฏฺฐงฺคิกํ มคฺคํ ภาเวนฺตํ อริยํ อฏฺฐงฺคิกํ มคฺคํ พหุลีกโรนฺตํ ราชาโน วา ราชมหามตฺตา วา มิตฺตา วา อมจฺจา วา ญาตี วา ญาติสาโลหิตา วา โภเคหิ อภิหฏฺฐุํ ปวาเรยฺยุํ “เอหมฺโภ ปุริส กิํ เต อิเม กาสาวา อนุทหนฺติ, กิํ มุณฺโฑ กปาลมนุสํจรสิ, เอหิ หีนายาวตฺติตฺวา โภเค จ ภุญฺชสฺสุ, ปุญฺญานิ จ กโรหี”ติ. โส วต ภิกฺขเว ภิกฺขุ อริยํ อฏฺฐงฺคิกํ มคฺคํ ภาเวนฺโต อริยํ อฏฺฐงฺคิกํ มคฺคํ พหุลีกโรนฺโต สิกฺขํ ปจฺจกฺขาย หีนายาวตฺติสฺสตีติ เนตํ ฐานํ วิชฺชติ. ตํ กิสฺส เหตุ, ยํ หิ ตํ ภิกฺขเว จิตฺตํ ทีฆรตฺตํ วิเวกนินฺนํ วิเวกโปณํ วิเวก\-ปพฺภารํ, ตํ วต หีนายาวตฺติสฺสตีติ เนตํ ฐานํ วิชฺชติ. กถญฺจ ภิกฺขเว ภิกฺขุ อริยํ อฏฺฐงฺคิกํ มคฺคํ ภาเวติ อริยํ อฏฺฐงฺคิกํ มคฺคํ พหุลีกโรติ, อิธ ภิกฺขเว ภิกฺขุ สมฺมาทิฏฺฐิํ ภาเวติ วิเวกนิสฺสิตํ ฯเปฯ สมฺมาสมาธิํ ภาเวติ วิเวกนิสฺสิตํ วิราคนิสฺสิตํ นิโรธ\-นิสฺสิตํ โวสฺสคฺค\-ปริณามิํ. เอวํ โข ภิกฺขเว \csromanfolio{50}ภิกฺขุ อริยํ อฏฺฐงฺคิกํ มคฺคํ ภาเวติ อริยํ อฏฺฐงฺคิกํ มคฺคํ พหุลีกโรตีติ. (ยทปิ พลกรณียํ, ตทปิ วิตฺถาเรตพฺพํ.). ทฺวาทสมํ. พลกรณีย\-วคฺโค ฉฏฺโฐ.}

\vspace{\tipitakaparskip}
\csromancenter{\textbf{ตสฺสุทฺทานํ}}
\setlength{\csromangathapagewidth}{0pt}
\setlength{\csromangathawakwidth}{0pt}
\csromangathameasuregroup{พลํ พีชญฺจ นาโค จ, \\ อากาเสน จ เทฺว เมฆา,}{\csromangathabat{พลํ พีชญฺจ นาโค จ,}{รุกฺโข กุมฺเภน สูกิยา.} \\ \csromangathabat{อากาเสน จ เทฺว เมฆา,}{นาวา อาคนฺตุกา นทีติ.}}
\csromangathasetleft{พลํ พีชญฺจ นาโค จ, \\ อากาเสน จ เทฺว เมฆา,}
\csromangathagroup{\csromangathastanza{\csromangathabat{พลํ พีชญฺจ นาโค จ,}{รุกฺโข กุมฺเภน สูกิยา.} \\ \csromangathabat{อากาเสน จ เทฺว เมฆา,}{นาวา อาคนฺตุกา นทีติ.}}}

\csromanmatikaanchor{mtk.6}
\csromanmatikaanchor{mtk.29}
\csromantocmarkref{section}{10. นีวรณสุตฺต}{mtk.6}
\bookmark[dest={mtk.6},level=1]{10. นีวรณสุตฺต}
\csromanheader{4}{7. \textbf{เอสนาวคฺค}}
\csromanheader{2}{1. \textbf{เอสนาสุตฺต}}
\proseitem{161}{สาวตฺถิ\-นิทานํ. ติสฺโส อิมา ภิกฺขเว เอสนา. กตมา ติสฺโส, กาเมสนา ภเวสนา พฺรหฺมจริเย\-สนา, อิมา โข ภิกฺขเว ติสฺโส เอสนา. อิมาสํ โข ภิกฺขเว ติสฺสนฺนํ เอสนานํ อภิญฺญาย อริโย อฏฺฐงฺคิโก มคฺโค ภาเวตพฺโพ. กตโม อริโย อฏฺฐงฺคิโก มคฺโค, อิธ ภิกฺขเว ภิกฺขุ สมฺมาทิฏฺฐิํ ภาเวติ วิเวกนิสฺสิตํ ฯเปฯ สมฺมาสมาธิํ ภาเวติ วิเวกนิสฺสิตํ วิราคนิสฺสิตํ นิโรธ\-นิสฺสิตํ โวสฺสคฺค\-ปริณามิํ. อิมาสํ โข ภิกฺขเว ติสฺสนฺนํ เอสนานํ อภิญฺญาย อยํ อริโย อฏฺฐงฺคิโก มคฺโค ภาเวตพฺโพติ.}

\prose{ติสฺโส อิมา โข ภิกฺขเว เอสนา. กตมา ติสฺโส, กาเมสนา ภเวสนา พฺรหฺมจริเย\-สนา, อิมา โข ภิกฺขเว ติสฺโส เอสนา. อิมาสํ โข ภิกฺขเว ติสฺสนฺนํ เอสนานํ อภิญฺญาย อริโย อฏฺฐงฺคิโก มคฺโค ภาเวตพฺโพ. กตโม อริโย อฏฺฐงฺคิโก มคฺโค, อิธ ภิกฺขเว ภิกฺขุ สมฺมาทิฏฺฐิํ ภาเวติ ฯเปฯ สมฺมาสมาธิํ ภาเวติ ราค\-วินย\-ปริโยสานํ โทส\-วินย\-ปริโยสานํ โมห\-วินย\-ปริโยสานํ. อิมาสํ โข ภิกฺขเว ติสฺสนฺนํ เอสนานํ อภิญฺญาย อยํ อริโย อฏฺฐงฺคิโก มคฺโค ภาเวตพฺโพติ.}

\prose{ติสฺโส อิมา โข ภิกฺขเว เอสนา. กตมา ติสฺโส, กาเมสนา ภเวสนา พฺรหฺมจริเย\-สนา, อิมา โข ภิกฺขเว ติสฺโส เอสนา. อิมาสํ}

\prose{\csromanfolio{51}โข ภิกฺขเว ติสฺสนฺนํ เอสนานํ อภิญฺญาย อริโย อฏฺฐงฺคิโก มคฺโค ภาเวตพฺโพ. กตโม อริโย อฏฺฐงฺคิโก มคฺโค, อิธ ภิกฺขเว ภิกฺขุ สมฺมาทิฏฺฐิํ ภาเวติ ฯเปฯ สมฺมาสมาธิํ ภาเวติ อมโตคธํ อมตปรายนํ อมต\-ปริโยสานํ. อิมาสํ โข ภิกฺขเว ติสฺสนฺนํ เอสนานํ อภิญฺญาย อยํ อริโย อฏฺฐงฺคิโก มคฺโค ภาเวตพฺโพติ.}

\prose{ติสฺโส อิมา โข ภิกฺขเว เอสนา. กตมา ติสฺโส, กาเมสนา ภเวสนา พฺรหฺมจริเย\-สนา, อิมา โข ภิกฺขเว ติสฺโส เอสนา. อิมาสํ โข ภิกฺขเว ติสฺสนฺนํ เอสนานํ อภิญฺญาย อริโย อฏฺฐงฺคิโก มคฺโค ภาเวตพฺโพ. กตโม อริโย อฏฺฐงฺคิโก มคฺโค, อิธ ภิกฺขเว ภิกฺขุ สมฺมาทิฏฺฐิํ ภาเวติ ฯเปฯ สมฺมาสมาธิํ ภาเวติ นิพฺพานนินฺนํ นิพฺพานโปณํ นิพฺพาน\-ปพฺภารํ. อิมาสํ โข ภิกฺขเว ติสฺสนฺนํ เอสนานํ อภิญฺญาย อยํ อริโย อฏฺฐงฺคิโก มคฺโค ภาเวตพฺโพติ.}

\prose{ติสฺโส อิมา ภิกฺขเว เอสนา. กตมา ติสฺโส, กาเมสนา ภเวสนา พฺรหฺมจริเย\-สนา, อิมา โข ภิกฺขเว ติสฺโส เอสนา. อิมาสํ โข ภิกฺขเว ติสฺสนฺนํ เอสนานํ ปริญฺญาย ฯเปฯ อยํ อริโย อฏฺฐงฺคิโก มคฺโค ภาเวตพฺโพติ.}

\vspace{\tipitakaparskip}
\csromancenter{(ยทปิ อภิญฺญา, ตทปิ ปริญฺญาย วิตฺถาเรตพฺพํ.)}
\prose{ติสฺโส อิมา ภิกฺขเว เอสนา. กตมา ติสฺโส, กาเมสนา ภเวสนา พฺรหฺมจริเย\-สนา, อิมา โข ภิกฺขเว ติสฺโส เอสนา. อิมาสํ โข ภิกฺขเว ติสฺสนฺนํ เอสนานํ ปริกฺขยาย ฯเปฯ อยํ อริโย อฏฺฐงฺคิโก มคฺโค ภาเวตพฺโพติ.}

\vspace{\tipitakaparskip}
\csromancenter{(ยทปิ อภิญฺญา, ตทปิ ปริกฺขยาย วิตฺถาเรตพฺพํ.)}
\prose{ติสฺโส อิมา ภิกฺขเว เอสนา. กตมา ติสฺโส, กาเมสนา ภเวสนา พฺรหฺมจริเย\-สนา, อิมา โข ภิกฺขเว ติสฺโส เอสนา. อิมาสํ โข ภิกฺขเว ติสฺสนฺนํ เอสนานํ ปหานาย อริโย อฏฺฐงฺคิโก มคฺโค ภาเวตพฺโพ. กตโม อริโย อฏฺฐงฺคิโก มคฺโค, อิธ ภิกฺขเว ภิกฺขุ สมฺมาทิฏฺฐิํ ภาเวติ ฯเปฯ สมฺมาสมาธิํ ภาเวติ วิเวกนิสฺสิตํ วิราคนิสฺสิตํ นิโรธ\-นิสฺสิตํ โวสฺสคฺค\-ปริณามิํ. อิมาสํ โข ภิกฺขเว ติสฺสนฺนํ เอสนานํ ปหานาย อยํ อริโย อฏฺฐงฺคิโก มคฺโค ภาเวตพฺโพติ. (ยทปิ อภิญฺญา, ตทปิ ปหานาย วิตฺถาเรตพฺพํ.). ปฐมํ.}

\csromanheader{2}{2. \textbf{วิธาสุตฺต}}
\proseitem{162}{\csromanfolio{52}ติสฺโส อิมา ภิกฺขเว วิธา. กตมา ติสฺโส, เสยฺโยหมสฺมีติวิธา สทิโสหมสฺมีติวิธา หีโนหมสฺมีติวิธา, อิมา โข ภิกฺขเว ติสฺโส วิธา. อิมาสํ โข ภิกฺขเว ติสฺสนฺนํ วิธานํ อภิญฺญาย ปริญฺญาย ปริกฺขยาย ปหานาย อริโย อฏฺฐงฺคิโก มคฺโค ภาเวตพฺโพ. กตโม อริโย อฏฺฐงฺคิโก มคฺโค, อิธ ภิกฺขเว ภิกฺขุ สมฺมาทิฏฺฐิํ ภาเวติ ฯเปฯ สมฺมาสมาธิํ ภาเวติ วิเวกนิสฺสิตํ วิราคนิสฺสิตํ นิโรธ\-นิสฺสิตํ โวสฺสคฺค\-ปริณามิํ. อิมาสํ โข ภิกฺขเว ติสฺสนฺนํ วิธานํ อภิญฺญาย ปริญฺญาย ปริกฺขยาย ปหานาย อยํ อริโย อฏฺฐงฺคิโก มคฺโค ภาเวตพฺโพติ. (ยถา เอสนา, เอวํ วิตฺถาเรตพฺพํ.). ทุติยํ.}

\csromanheader{2}{3. \textbf{อาสวสุตฺต}}
\proseitem{163}{ตโยเม ภิกฺขเว อาสวา. กตเม ตโย, กามาสโว ภวาสโว อวิชฺชาสโว, อิเม โข ภิกฺขเว ตโย อาสวา. อิเมสํ โข ภิกฺขเว ติณฺณนฺนํ อาสวานํ อภิญฺญาย ปริญฺญาย ปริกฺขยาย ปหานาย ฯเปฯ อยํ อริโย อฏฺฐงฺคิโก มคฺโค ภาเวตพฺโพติ.  ตติยํ.}

\csromanheader{2}{4. \textbf{ภวสุตฺต}}
\proseitem{164}{ตโยเม ภิกฺขเว ภวา. กตเม ตโย, กามภโว รูปภโว อรูปภโว, อิเม โข ภิกฺขเว ตโย ภวา. อิเมสํ โข ภิกฺขเว ติณฺณนฺนํ ภวานํ อภิญฺญาย ปริญฺญาย ปริกฺขยาย ปหานาย ฯเปฯ อยํ อริโย อฏฺฐงฺคิโก มคฺโค ภาเวตพฺโพติ.  จตุตฺถํ.}

\csromanheader{2}{5. \textbf{ทุกฺขตาสุตฺต}}
\proseitem{165}{ติสฺโส อิมา ภิกฺขเว ทุกฺขตา. กตมา ติสฺโส, ทุกฺขทุกฺขตา สงฺขาร\-ทุกฺขตา วิปริณาม\-ทุกฺขตา, อิมา โข ภิกฺขเว ติสฺโส ทุกฺขตา. อิมาสํ โข ภิกฺขเว ติสฺสนฺนํ ทุกฺขตานํ อภิญฺญาย ปริญฺญาย ปริกฺขยาย ปหานาย ฯเปฯ อยํ อริโย อฏฺฐงฺคิโก มคฺโค ภาเวตพฺโพติ.  ปญฺจมํ.}

\csromanheader{2}{6. \textbf{ขิลสุตฺต}}
\proseitem{166}{ตโยเม ภิกฺขเว ขิลา. กตเม ตโย, ราโค ขิโล โทโส ขิโล โมโห ขิโล, อิเม โข ภิกฺขเว ตโย ขิลา.}

\prose{\csromanfolio{53}อิเมสํ โข ภิกฺขเว ติณฺณนฺนํ ขิลานํ อภิญฺญาย ปริญฺญาย ปริกฺขยาย ปหานาย ฯเปฯ อยํ อริโย อฏฺฐงฺคิโก มคฺโค ภาเวตพฺโพติ.  ฉฏฺฐํ.}

\csromanheader{2}{7. \textbf{มลสุตฺต}}
\proseitem{167}{ตีณิมานิ ภิกฺขเว มลานิ. กตมานิ ตีณิ, ราโค มลํ โทโส มลํ โมโห มลํ, อิมานิ โข ภิกฺขเว ตีณิ มลานิ. อิเมสํ โข ภิกฺขเว ติณฺณนฺนํ มลานํ อภิญฺญาย ปริญฺญาย ปริกฺขยาย ปหานาย ฯเปฯ อยํ อริโย อฏฺฐงฺคิโก มคฺโค ภาเวตพฺโพติ.  สตฺตมํ.}

\csromanheader{2}{8. \textbf{นีฆสุตฺต}}
\proseitem{168}{ตโยเม ภิกฺขเว นีฆา. กตเม ตโย, ราโค นีโฆ โทโส นีโฆ โมโห นีโฆ, อิเม โข ภิกฺขเว ตโย นีฆา. อิเมสํ โข ภิกฺขเว ติณฺณนฺนํ นีฆานํ อภิญฺญาย ปริญฺญาย ปริกฺขยาย ปหานาย ฯเปฯ อยํ อริโย อฏฺฐงฺคิโก มคฺโค ภาเวตพฺโพติ.  อฏฺฐมํ.}

\csromanheader{2}{9. \textbf{เวทนาสุตฺต}}
\proseitem{169}{ติสฺโส อิมา ภิกฺขเว เวทนา. กตมา ติสฺโส, สุขา เวทนา ทุกฺขา เวทนา อทุกฺขมสุขา เวทนา, อิมา โข ภิกฺขเว ติสฺโส เวทนา. อิมาสํ โข ภิกฺขเว ติสฺสนฺนํ เวทนานํ อภิญฺญาย ปริญฺญาย ปริกฺขยาย ปหานาย ฯเปฯ อยํ อริโย อฏฺฐงฺคิโก มคฺโค ภาเวตพฺโพติ.  นวมํ.}

\csromanheader{2}{10. \textbf{ตณฺหาสุตฺต}}
\proseitem{170}{ติสฺโส อิมา ภิกฺขเว ตณฺหา. กตมา ติสฺโส, กามตณฺหา ภวตณฺหา วิภวตณฺหา, อิมา โข ภิกฺขเว ติสฺโส ตณฺหา. อิมาสํ โข ภิกฺขเว ติสฺสนฺนํ ตณฺหานํ อภิญฺญาย ปริญฺญาย ปริกฺขยาย ปหานาย ฯเปฯ อยํ อริโย อฏฺฐงฺคิโก มคฺโค ภาเวตพฺโพ. กตโม อริโย อฏฺฐงฺคิโก มคฺโค, อิธ ภิกฺขเว ภิกฺขุ สมฺมาทิฏฺฐิํ ภาเวติ วิเวกนิสฺสิตํ วิราคนิสฺสิตํ นิโรธ\-นิสฺสิตํ โวสฺสคฺค\-ปริณามิํ ฯเปฯ สมฺมาสมาธิํ ภาเวติ วิเวกนิสฺสิตํ วิราคนิสฺสิตํ นิโรธ\-นิสฺสิตํ โวสฺสคฺค\-ปริณามิํ. อิมาสํ โข ภิกฺขเว ติสฺสนฺนํ ตณฺหานํ อภิญฺญาย ปริญฺญาย ปริกฺขยาย ปหานาย ฯเปฯ อยํ อริโย อฏฺฐงฺคิโก มคฺโค ภาเวตพฺโพติ.  ทสมํ.}

\csromanheader{2}{11. \textbf{ตสินาสุตฺต}}
\proseitemwithclosermidruled{171}{\csromanfolio{54}ติสฺโส อิมา ภิกฺขเว ตสินา. กตมา ติสฺโส, กามตสินา ภวตสินา วิภวตสินา. อิมาสํ โข ภิกฺขเว ติสฺสนฺนํ ตสินานํ อภิญฺญาย ปริญฺญาย ปริกฺขยาย ปหานาย ฯเปฯ ราค\-วินย\-ปริโยสานํ โทส\-วินย\-ปริโยสานํ โมห\-วินย\-ปริโยสานํ ฯเปฯ อมโตคธํ อมตปรายนํ อมต\-ปริโยสานํ ฯเปฯ นิพฺพานนินฺนํ นิพฺพานโปณํ นิพฺพาน\-ปพฺภารํ. อิมาสํ โข ภิกฺขเว ติสฺสนฺนํ ตสินานํ อภิญฺญาย ปริญฺญาย ปริกฺขายาย ปหานาย ฯเปฯ อยํ อริโย อฏฺฐงฺคิโก มคฺโค ภาเวตพฺโพติ.  เอกาทสมํ.}{เอสนาวคฺโค สตฺตโม.}
\csromancenter{\textbf{ตสฺสุทฺทานํ}}
\setlength{\csromangathapagewidth}{0pt}
\setlength{\csromangathawakwidth}{0pt}
\csromangathameasuregroup{เอสนา วิธา อาสโว, \\ มลํ นีโฆ จ เวทนา,}{\csromangathabat{เอสนา วิธา อาสโว,}{ภโว จ ทุกฺขตา ขิลา.} \\ \csromangathabat{มลํ นีโฆ จ เวทนา,}{เทฺว ตณฺหา ตสินาย จาติ.}}
\csromangathasetleft{เอสนา วิธา อาสโว, \\ มลํ นีโฆ จ เวทนา,}
\csromangathagroup{\csromangathastanza{\csromangathabat{เอสนา วิธา อาสโว,}{ภโว จ ทุกฺขตา ขิลา.} \\ \csromangathabat{มลํ นีโฆ จ เวทนา,}{เทฺว ตณฺหา ตสินาย จาติ.}}}

\csromanmatikaanchor{mtk.7}
\csromanmatikaanchor{mtk.31}
\csromantocmarknopage{chapter}{5. จกฺกวตฺติวคฺค}{mtk.7}
\bookmark[dest={mtk.7},level=0]{5. จกฺกวตฺติวคฺค}
\csromanheader{4}{8. \textbf{โอฆวคฺค}}
\csromanheader{2}{1. \textbf{โอฆสุตฺต}}
\proseitem{172}{สาวตฺถิ\-นิทานํ. จตฺตาโรเม ภิกฺขเว โอฆา. กตเม จตฺตาโร, กาโมโฆ ภโวโฆ ทิฏฺโฐโฆ อวิชฺโชโฆ, อิเม โข ภิกฺขเว จตฺตาโร โอฆา. อิเมสํ โข ภิกฺขเว จตุนฺนํ โอฆานํ อภิญฺญาย ปริญฺญาย ปริกฺขยาย ปหานาย ฯเปฯ อยํ อริโย อฏฺฐงฺคิโก มคฺโค ภาเวตพฺโพติ. (ยถา เอสนา, เอวํ วิตฺถาเรตพฺพํ.). ปฐมํ.}

\csromanheader{2}{2. \textbf{โยคสุตฺต}}
\proseitem{173}{จตฺตาโรเม ภิกฺขเว โยคา. กตเม จตฺตาโร, กามโยโค ภวโยโค ทิฏฺฐิโยโค อวิชฺชาโยโค, อิเม โข ภิกฺขเว จตฺตาโร โยคา. อิเมสํ โข ภิกฺขเว จตุนฺนํ โยคานํ อภิญฺญาย ปริญฺญาย ปริกฺขยาย ปหานาย ฯเปฯ อยํ อริโย อฏฺฐงฺคิโก มคฺโค ภาเวตพฺโพติ.  ทุติยํ.}

\csromanheader{2}{3. \textbf{อุปาทานสุตฺต}}
\proseitem{174}{\csromanfolio{55}จตฺตาริมานิ ภิกฺขเว อุปาทานานิ. กตมานิ จตฺตาริ, กามุปาทานํ ทิฏฺฐุปาทานํ สีลพฺพตุ\-ปาทานํ อตฺตวาทุ\-ปาทานํ, อิมานิ โข ภิกฺขเว จตฺตาริ อุปาทานานิ. อิเมสํ โข ภิกฺขเว จตุนฺนํ อุปาทานานํ อภิญฺญาย ปริญฺญาย ปริกฺขยาย ปหานาย ฯเปฯ อยํ อริโย อฏฺฐงฺคิโก มคฺโค ภาเวตพฺโพติ.  ตติยํ.}

\csromanheader{2}{4. \textbf{คนฺถสุตฺต}}
\proseitem{175}{จตฺตาโรเม ภิกฺขเว คนฺถา. กตเม จตฺตาโร, อภิชฺฌา กายคนฺโถ พฺยาปาโท กายคนฺโถ สีลพฺพต\-ปรามาโส กายคนฺโถ อิทํสจฺจา\-ภินิเวโส กายคนฺโถ, อิเม โข ภิกฺขเว จตฺตาโร คนฺถา. อิเมสํ โข ภิกฺขเว จตุนฺนํ คนฺถานํ อภิญฺญาย ปริญฺญาย ปริกฺขยาย ปหานาย ฯเปฯ อยํ อริโย อฏฺฐงฺคิโก มคฺโค ภาเวตพฺโพติ.  จตุตฺถํ.}

\csromanheader{2}{5. \textbf{อนุสยสุตฺต}}
\proseitem{176}{สตฺติเม ภิกฺขเว อนุสยา. กตเม สตฺต, กามราคานุสโย ปฏิฆานุสโย ทิฏฺฐานุสโย วิจิกิจฺฉา\-นุสโย มานานุสโย ภวราคา\-นุสโย อวิชฺชานุสโย, อิเม โข ภิกฺขเว สตฺตานุสยา. อิเมสํ โข ภิกฺขเว สตฺตนฺนํ อนุสยานํ อภิญฺญาย ปริญฺญาย ปริกฺขยาย ปหานาย ฯเปฯ อยํ อริโย อฏฺฐงฺคิโก มคฺโค ภาเวตพฺโพติ.  ปญฺจมํ.}

\csromanheader{2}{6. \textbf{กามคุณสุตฺต}}
\proseitem{177}{ปญฺจิเม ภิกฺขเว กามคุณา. กตเม ปญฺจ, จกฺขุวิญฺเญยฺยา รูปา อิฏฺฐา กนฺตา มนาปา ปิยรูปา กามูปสํหิตา รชนียา, โสตวิญฺเญยฺยา สทฺทา ฯเปฯ ฆานวิญฺเญยฺยา คนฺธา ฯเปฯ ชิวฺหาวิญฺเญยฺยา รสา ฯเปฯ กายวิญฺเญยฺยา โผฏฺฐพฺพา อิฏฺฐา กนฺตา มนาปา ปิยรูปา กามูปสํหิตา รชนียา, อิเม โข ภิกฺขเว ปญฺจ กามคุณา. อิเมสํ โข ภิกฺขเว ปญฺจนฺนํ กามคุณานํ อภิญฺญาย ปริญฺญาย ปริกฺขยาย ปหานาย ฯเปฯ อยํ อริโย อฏฺฐงฺคิโก มคฺโค ภาเวตพฺโพติ.  ฉฏฺฐํ.}

\csromanheader{2}{7. \textbf{นีวรณสุตฺต}}
\proseitem{178}{\csromanfolio{56}ปญฺจิมานิ ภิกฺขเว นีวรณานิ. กตมานิ ปญฺจ, กามจฺฉนฺท\-นีวรณํ พฺยาปาท\-นีวรณํ ถินมิทฺธ\-นีวรณํ อุทฺธจฺจ\-กุกฺกุจฺจ\-นีวรณํ วิจิกิจฺฉา\-นีวรณํ, อิมานิ โข ภิกฺขเว ปญฺจ นีวรณานิ. อิเมสํ โข ภิกฺขเว ปญฺจนฺนํ นีวรณานํ อภิญฺญาย ปริญฺญาย ปริกฺขยาย ปหานาย ฯเปฯ อยํ อริโย อฏฺฐงฺคิโก มคฺโค ภาเวตพฺโพติ.  สตฺตมํ.}

\csromanheader{2}{8. อุปาทานกฺขนฺธ\-สุตฺต}
\proseitem{179}{ปญฺจิเม ภิกฺขเว อุปาทาน\-กฺขนฺธา. กตเม ปญฺจ, เสยฺยถิทํ, รูปุปาทาน\-กฺขนฺโธ เวทนุ\-ปาทาน\-กฺขนฺโธ สญฺญุ\-ปาทาน\-กฺขนฺโธ สงฺขารุ\-ปาทาน\-กฺขนฺโธ วิญฺญาณุ\-ปาทาน\-กฺขนฺโธ, อิเม โข ภิกฺขเว ปญฺจุ\-ปาทาน\-กฺขนฺธา. อิเมสํ โข ภิกฺขเว ปญฺจนฺนํ อุปาทาน\-กฺขนฺธานํ อภิญฺญาย ปริญฺญาย ปริกฺขยาย ปหานาย ฯเปฯ อยํ อริโย อฏฺฐงฺคิโก มคฺโค ภาเวตพฺโพติ.  อฏฺฐมํ.}

\csromanheader{2}{9. โอรมฺภาคิย\-สุตฺต}
\proseitem{180}{ปญฺจิมานิ ภิกฺขเว โอรมฺภาคิยานิ สํโยชนานิ. กตมานิ ปญฺจ, สกฺกายทิฏฺฐิ วิจิกิจฺฉา สีลพฺพต\-ปรามาโส กามจฺฉนฺโท พฺยาปาโท, อิมานิ โข ภิกฺขเว ปญฺโจ\-รมฺภาคิยานิ สํโยชนานิ. อิมาสํ โข ภิกฺขเว ปญฺจนฺนํ โอรมฺภาคิยานํ สํโยชนานํ อภิญฺญาย ปริญฺญาย ปริกฺขยาย ปหานาย ฯเปฯ อยํ อริโย อฏฺฐงฺคิโก มคฺโค ภาเวตพฺโพติ.  นวมํ.}

\csromanheader{2}{10. อุทฺธมฺภาคิย\-สุตฺต}
\proseitem{181}{ปญฺจิมานิ ภิกฺขเว อุทฺธมฺ\-ภาคิยานิ สํโยชนานิ. กตมานิ ปญฺจ, รูปราโค อรูปราโค มาโน อุทฺธจฺจํ อวิชฺชา, อิมานิ โข ภิกฺขเว ปญฺจุ\-ทฺธมฺภาคิยานิ สํโยชนานิ. อิเมสํ โข ภิกฺขเว ปญฺจนฺนํ อุทฺธมฺ\-ภาคิยานํ สํโยชนานํ อภิญฺญาย ปริญฺญาย ปริกฺขยาย ปหานาย อริโย อฏฺฐงฺคิโก มคฺโค ภาเวตพฺโพ. กตโม อริโย อฏฺฐงฺคิโก มคฺโค, อิธ ภิกฺขเว ภิกฺขุ สมฺมาทิฏฺฐิํ ภาเวติ วิเวกนิสฺสิตํ ฯเปฯ สมฺมาสมาธิํ ภาเวติ วิเวกนิสฺสิตํ วิราคนิสฺสิตํ นิโรธ\-นิสฺสิตํ โวสฺสคฺค\-ปริณามิํ. อิเมสํ โข ภิกฺขเว ปญฺจนฺนํ อุทฺธมฺ\-ภาคิยานํ สํโยชนานํ อภิญฺญาย}

\prose{\csromanfolio{57}ปริญฺญาย ปริกฺขยาย ปหานาย อยํ อริโย อฏฺฐงฺคิโก มคฺโค ภาเวตพฺโพติ.}

\prose{ปญฺจิมานิ ภิกฺขเว อุทฺธมฺ\-ภาคิยานิ สํโยชนานิ. กตมานิ ปญฺจ, รูปราโค อรูปราโค มาโน อุทฺธจฺจํ อวิชฺชา, อิมานิ โข ภิกฺขเว ปญฺจุ\-ทฺธมฺภาคิยานิ สํโยชนานิ. อิเมสํ โข ภิกฺขเว ปญฺจนฺนํ อุทฺธมฺ\-ภาคิยานํ สํโยชนานํ อภิญฺญาย ปริญฺญาย ปริกฺขยาย ปหานาย อริโย อฏฺฐงฺคิโก มคฺโค ภาเวตพฺโพ. กตโม อริโย อฏฺฐงฺคิโก มคฺโค, อิธ ภิกฺขเว ภิกฺขุ สมฺมาทิฏฺฐิํ ภาเวติ ฯเปฯ สมฺมาสมาธิํ ภาเวติ ราค\-วินย\-ปริโยสานํ โทส\-วินย\-ปริโยสานํ โมห\-วินย\-ปริโยสานํ. อมโตคธํ อมตปรายนํ อมต\-ปริโยสานํ. นิพฺพานนินฺนํ นิพฺพานโปณํ นิพฺพาน\-ปพฺภารํ. อิเมสํ โข ภิกฺขเว ปญฺจนฺนํ อุทฺธมฺ\-ภาคิยานํ สํโยชนานํ อภิญฺญาย ปริญฺญาย ปริกฺขยาย ปหานาย อยํ อริโย อฏฺฐงฺคิโก มคฺโค ภาเวตพฺโพติ.  ทสมํ. โอฆวคฺโค อฏฺฐโม.}

\vspace{\tipitakaparskip}
\csromancenter{\textbf{ตสฺสุทฺทานํ}}
\setlength{\csromangathapagewidth}{0pt}
\setlength{\csromangathawakwidth}{0pt}
\csromangathameasuregroup{โอโฆ โยโค อุปาทานํ, \\ กามคุณา นีวรณํ,}{\csromangathabat{โอโฆ โยโค อุปาทานํ,}{คนฺถํ อนุสเยน จ.} \\ \csromangathabat{กามคุณา นีวรณํ,}{ขนฺธา โอรุทฺธมฺภาคิยาติ.}}
\csromangathasetleft{โอโฆ โยโค อุปาทานํ, \\ กามคุณา นีวรณํ,}
\csromangathagroup{\csromangathastanza{\csromangathabat{โอโฆ โยโค อุปาทานํ,}{คนฺถํ อนุสเยน จ.} \\ \csromangathabat{กามคุณา นีวรณํ,}{ขนฺธา โอรุทฺธมฺภาคิยาติ.}}}

\csromanheader{2}{\textbf{วคฺคุทฺทานํ}}
\setlength{\csromangathapagewidth}{0pt}
\setlength{\csromangathawakwidth}{0pt}
\csromangathameasuregroup{อวิชฺชาวคฺโค ปฐโม, \\ มิจฺฉตฺตํ ตติโย วคฺโค, \\ ติตฺถิยํ ปญฺจโม วคฺโค, \\ พหุกเต สตฺตโม วคฺโค, \\ ทิวสวคฺโค นวโม, \\ เอกาทสพลวคฺโค,}{\csromangathabat{อวิชฺชาวคฺโค ปฐโม,}{ทุติยํ วิหารํ วุจฺจติ.} \\ \csromangathabat{มิจฺฉตฺตํ ตติโย วคฺโค,}{จตุตฺถํ ปฏิปนฺเนเนว.} \\ \csromangathabat{ติตฺถิยํ ปญฺจโม วคฺโค,}{จตสฺโส สูริเยน จ.} \\ \csromangathabat{พหุกเต สตฺตโม วคฺโค,}{อุปฺปาโท อฏฺฐเมน จ.} \\ \csromangathabat{ทิวสวคฺโค นวโม,}{ทสโม อปฺปมาเทน จ.} \\ \csromangathabat{เอกาทสพลวคฺโค,}{ทฺวาทส เอสนา ปาฬิยํ.} \\ โอฆวคฺโค ภวติ เตรสาติ. \\ มคฺคสํยุตฺตํ ปฐมํ.}
\csromangathameasurewak{โอฆวคฺโค ภวติ เตรสาติ. \\ มคฺคสํยุตฺตํ ปฐมํ.}
\csromangathasetleft{อวิชฺชาวคฺโค ปฐโม, \\ มิจฺฉตฺตํ ตติโย วคฺโค, \\ ติตฺถิยํ ปญฺจโม วคฺโค, \\ พหุกเต สตฺตโม วคฺโค, \\ ทิวสวคฺโค นวโม, \\ เอกาทสพลวคฺโค,}
\csromangathagroup{\csromangathastanza{\csromangathabat{อวิชฺชาวคฺโค ปฐโม,}{ทุติยํ วิหารํ วุจฺจติ.} \\ \csromangathabat{มิจฺฉตฺตํ ตติโย วคฺโค,}{จตุตฺถํ ปฏิปนฺเนเนว.}}\csromangathastanza{\csromangathabat{ติตฺถิยํ ปญฺจโม วคฺโค,}{จตสฺโส สูริเยน จ.} \\ \csromangathabat{พหุกเต สตฺตโม วคฺโค,}{อุปฺปาโท อฏฺฐเมน จ.}}\csromangathastanza{\csromangathabat{ทิวสวคฺโค นวโม,}{ทสโม อปฺปมาเทน จ.} \\ \csromangathabat{เอกาทสพลวคฺโค,}{ทฺวาทส เอสนา ปาฬิยํ.}}\csromangathastanza{\csromangathawak{โอฆวคฺโค ภวติ เตรสาติ. \\ มคฺคสํยุตฺตํ ปฐมํ.}}}

\csromanheader{2}{2. โพชฺฌงฺค\-สํยุตฺต}
\chapterhead{1. \textbf{ปพฺพตวคฺค}}
\csromanheader{2}{1. \textbf{หิมวนฺตสุตฺต}}
\proseitem{182}{\csromanfolio{58}สาวตฺถิ\-นิทานํ. เสยฺยถาปิ ภิกฺขเว หิมวนฺตํ ปพฺพตราชานํ นิสฺสาย นาคา กายํ วฑฺเฒนฺติ พลํ คาเหนฺติ, เต ตตฺถ กายํ วฑฺเฒตฺวา พลํ คาเหตฺวา กุโสพฺเภ โอตรนฺติ, กุโสพฺเภ โอตริตฺวา มหาโสพฺเภ โอตรนฺติ, มหาโสพฺเภ โอตริตฺวา กุนฺนทิโย โอตรนฺติ, กุนฺนทิโย โอตริตฺวา มหานทิโย โอตรนฺติ, มหานทิโย โอตริตฺวา มหา\-สมุทฺทสาครํ โอตรนฺติ, เต ตตฺถ มหนฺตตฺตํ เวปุลฺลตฺตํ อาปชฺชนฺติ กาเยน. เอวเมว โข ภิกฺขเว ภิกฺขุ สีลํ นิสฺสาย สีเล ปติฏฺฐาย สตฺต โพชฺฌงฺเค ภาเวนฺโต สตฺต โพชฺฌงฺเค พหุลีกโรนฺโต มหนฺตตฺตํ เวปุลฺลตฺตํ ปาปุณาติ ธมฺเมสุ. กถญฺจ ภิกฺขเว ภิกฺขุ สีลํ นิสฺสาย สีเล ปติฏฺฐาย สตฺต โพชฺฌงฺเค ภาเวนฺโต สตฺต โพชฺฌงฺเค พหุลีกโรนฺโต มหนฺตตฺตํ เวปุลฺลตฺตํ ปาปุณาติ ธมฺเมสูติ. อิธ ภิกฺขเว ภิกฺขุ สติสมฺโพชฺฌงฺคํ ภาเวติ วิเวกนิสฺสิตํ วิราคนิสฺสิตํ นิโรธ\-นิสฺสิตํ โวสฺสคฺค\-ปริณามิํ. ธมฺมวิจย\-สมฺโพชฺฌงฺคํ ภาเวติ ฯเปฯ. วีริย\-สมฺโพชฺฌงฺคํ ภาเวติ ฯเปฯ. ปีติสมฺโพชฺฌงฺคํ ภาเวติ ฯเปฯ. ปสฺสทฺธิ\-สมฺโพชฺฌงฺคํ ภาเวติ ฯเปฯ. สมาธิ\-สมฺโพชฺฌงฺคํ ภาเวติ ฯเปฯ. อุเปกฺขา\-สมฺโพชฺฌงฺคํ ภาเวติ วิเวกนิสฺสิตํ วิราคนิสฺสิตํ นิโรธ\-นิสฺสิตํ โวสฺสคฺค\-ปริณามิํ. เอวํ โข ภิกฺขเว ภิกฺขุ สีลํ นิสฺสาย สีเล ปติฏฺฐาย สตฺต โพชฺฌงฺเค ภาเวนฺโต สตฺต โพชฺฌางฺเค พหุลีกโรนฺโต มหนฺตตฺตํ เวปุลฺลตฺตํ ปาปุณาติ ธมฺเมสูติ.  ปฐมํ.}

\csromanheader{2}{2. \textbf{กายสุตฺต}}
\proseitem{183}{สาวตฺถิ\-นิทานํ. เสยฺยถาปิ ภิกฺขเว อยํ กาโย อาหารฏฺฐิติโก อาหารํ ปฏิจฺจ ติฏฺฐติ, อนาหาโร โน ติฏฺฐติ. เอวเมว โข ภิกฺขเว ปญฺจ นีวรณา อาหารฏฺฐิติกา อาหารํ ปฏิจฺจ ติฏฺฐนฺติ, อนาหารา โน ติฏฺฐนฺติ.}

\prose{\csromanfolio{59}โก จ ภิกฺขเว อาหาโร อนุปฺปนฺนสฺส วา กามจฺฉนฺทสฺส อุปฺปาทาย อุปฺปนฺนสฺส วา กามจฺฉนฺทสฺส ภิยฺโยภาวาย เวปุลฺลาย, อตฺถิ ภิกฺขเว สุภนิมิตฺตํ, ตตฺถ อโยนิโสมนสิการ\-พหุลีกาโร, อยมาหาโร อนุปฺปนฺนสฺส วา กามจฺฉนฺทสฺส อุปฺปาทาย อุปฺปนฺนสฺส วา กามจฺฉนฺทสฺส ภิยฺโยภาวาย เวปุลฺลาย.}

\prose{โก จ ภิกฺขเว อาหาโร อนุปฺปนฺนสฺส วา พฺยาปาทสฺส อุปฺปาทาย อุปฺปนฺนสฺส วา พฺยาปาทสฺส ภิยฺโยภาวาย เวปุลฺลาย, อตฺถิ ภิกฺขเว ปฏิฆ\-นิมิตฺตํ, ตตฺถ อโยนิโสมนสิการ\-พหุลีกาโร, อยมาหาโร อนุปฺปนฺนสฺส วา พฺยาปาทสฺส อุปฺปาทาย อุปฺปนฺนสฺส วา พฺยาปาทสฺส ภิยฺโยภาวาย เวปุลฺลาย.}

\prose{โก จ ภิกฺขเว อาหาโร อนุปฺปนฺนสฺส วา ถินมิทฺธสฺส อุปฺปาทาย อุปฺปนฺนสฺส วา ถินมิทฺธสฺส ภิยฺโยภาวาย เวปุลฺลาย, อตฺถิ ภิกฺขเว อรติ ตนฺทิ วิชมฺภิตา ภตฺตสมฺมโท เจตโส จ ลีนตฺตํ, ตตฺถ อโยนิโสมนสิการ\-พหุลีกาโร, อยมาหาโร อนุปฺปนฺนสฺส วา ถินมิทฺธสฺส อุปฺปาทาย อุปฺปนฺนสฺส วา ถินมิทฺธสฺส ภิยฺโยภาวาย เวปุลฺลาย.}

\prose{โก จ ภิกฺขเว อาหาโร อนุปฺปนฺนสฺส วา อุทฺธจฺจ\-กุกฺกุจฺจสฺส อุปฺปาทาย อุปฺปนฺนสฺส วา อุทฺธจฺจ\-กุกฺกุจฺจสฺส ภิยฺโยภาวาย เวปุลฺลาย, อตฺถิ ภิกฺขเว เจตโส อวูปสโม, ตตฺถ อโยนิโสมนสิการ\-พหุลีกาโร, อยมาหาโร อนุปฺปนฺนสฺส วา อุทฺธจฺจ\-กุกฺกุจฺจสฺส อุปฺปาทาย อุปฺปนฺนสฺส วา อุทฺธจฺจ\-กุกฺกุจฺจสฺส ภิยฺโยภาวาย เวปุลฺลาย.}

\prose{โก จ ภิกฺขเว อาหาโร อนุปฺปนฺนาย วา วิจิกิจฺฉาย อุปฺปาทาย อุปฺปนฺนาย วา วิจิกิจฺฉาย ภิยฺโยภาวาย เวปุลฺลาย, อตฺถิ ภิกฺขเว วิจิกิจฺฉา\-ฏฺฐานียา ธมฺมา, ตตฺถ อโยนิโสมนสิการ\-พหุลีกาโร, อยมาหาโร อนุปฺปนฺนาย วา วิจิกิจฺฉาย อุปฺปาทาย อุปฺปนฺนาย วา วิจิกิจฺฉาย ภิยฺโยภาวาย เวปุลฺลาย.}

\prose{เสยฺยถาปิ ภิกฺขเว อยํ กาโย อาหารฏฺฐิติโก อาหารํ ปฏิจฺจ ติฏฺฐติ, อนาหาโร โน ติฏฺฐติ. เอวเมว โข ภิกฺขเว อิเม ปญฺจ นีวรณา อาหารฏฺฐิติกา อาหารํ ปฏิจฺจ ติฏฺฐนฺติ, อนาหารา โน ติฏฺฐนฺติ.}

\prose{\csromanfolio{60}เสยฺยถาปิ ภิกฺขเว อยํ กาโย อาหารฏฺฐิติโก อาหารํ ปฏิจฺจ ติฏฺฐติ, อนาหาโร โน ติฏฺฐติ. เอวเมว โข ภิกฺขเว สตฺต โพชฺฌงฺคา อาหารฏฺฐิติกา อาหารํ ปฏิจฺจ ติฏฺฐนฺติ, อนาหารา โน ติฏฺฐนฺติ.}

\prose{โก จ ภิกฺขเว อาหาโร อนุปฺปนฺนสฺส วา สติสมฺโพชฺฌงฺคสฺส อุปฺปาทาย อุปฺปนฺนสฺส วา สติสมฺโพชฺฌงฺคสฺส ภาวนาย ปาริปูริยา, อตฺถิ ภิกฺขเว สติสมฺโพชฺฌงฺค\-ฏฺฐานียา ธมฺมา, ตตฺถ โยนิโสมนสิการ\-พหุลีกาโร, อยมาหาโร อนุปฺปนฺนสฺส วา สติสมฺโพชฺฌงฺคสฺส อุปฺปาทาย อุปฺปนฺนสฺส วา สติสมฺโพชฺฌงฺคสฺส ภาวนาย ปาริปูริยา.}

\prose{โก จ ภิกฺขเว อาหาโร อนุปฺปนฺนสฺส วา ธมฺมวิจย\-สมฺโพชฺฌงฺคสฺส อุปฺปาทาย อุปฺปนฺนสฺส วา ธมฺมวิจย\-สมฺโพชฺฌงฺคสฺส ภาวนาย ปาริปูริยา, อตฺถิ ภิกฺขเว กุสลากุสลา ธมฺมา สาวชฺชานวชฺชา ธมฺมา หีนปณีตา ธมฺมา กณฺห\-สุกฺก\-สปฺปฏิภาคา ธมฺมา, ตตฺถ โยนิโสมนสิการ\-พหุลีกาโร, อยมาหาโร อนุปฺปนฺนสฺส วา ธมฺมวิจย\-สมฺโพชฺฌงฺคสฺส อุปฺปาทาย อุปฺปนฺนสฺส วา ธมฺมวิจย\-สมฺโพชฺฌงฺคสฺส ภาวนาย ปาริปูริยา.}

\prose{โก จ ภิกฺขเว อาหาโร อนุปฺปนฺนสฺส วา วีริย\-สมฺโพชฺฌงฺคสฺส อุปฺปาทาย อุปฺปนฺนสฺส วา วีริย\-สมฺโพชฺฌงฺคสฺส ภาวนาย ปาริปูริยา, อตฺถิ ภิกฺขเว อารมฺภธาตุ\csromansharedfootnote{b0c6d6a96c7c3ec2}{อารพฺภธาตุ (สฺยา, ก)} นิกฺกมธาตุ ปรกฺกมธาตุ, ตตฺถ โยนิโสมนสิการ\-พหุลีกาโร, อยมาหาโร อนุปฺปนฺนสฺส วา วีริย\-สมฺโพชฺฌงฺคสฺส อุปฺปาทาย อุปฺปนฺนสฺส วา วีริย\-สมฺโพชฺฌงฺคสฺส ภาวนาย ปาริปูริยา.}

\prose{โก จ ภิกฺขเว อาหาโร อนุปฺปนฺนสฺส วา ปีติสมฺโพชฺฌงฺคสฺส อุปฺปาทาย อุปฺปนฺนสฺส วา ปีติสมฺโพชฺฌงฺคสฺส ภาวนาย ปาริปูริยา, อตฺถิ ภิกฺขเว ปีติสมฺโพชฺฌงฺค\-ฏฺฐานียา ธมฺมา, ตตฺถ โยนิโสมนสิการ\-พหุลีกาโร, อยมาหาโร อนุปฺปนฺนสฺส วา ปีติสมฺโพชฺฌงฺคสฺส อุปฺปาทาย อุปฺปนฺนสฺส วา ปีติสมฺโพชฺฌงฺคาสฺส ภาวนาย ปาริปูริยา.}

\prose{โก จ ภิกฺขเว อาหาโร อนุปฺปนฺนสฺส วา ปสฺสทฺธิ\-สมฺโพชฺฌงฺคสฺส อุปฺปาทาย อุปฺปนฺนสฺส วา ปสฺสทฺธิ\-สมฺโพชฺฌงฺคสฺส ภาวนาย ปาริปูริยา, อตฺถิ}

\prose{\csromanfolio{61}ภิกฺขเว กายปสฺสทฺธิ จิตฺตปสฺสทฺธิ, ตตฺถ โยนิโสมนสิการ\-พหุลีกาโร, อยมาหาโร อนุปฺปนฺนสฺส วา ปสฺสทฺธิ\-สมฺโพชฺฌงฺคสฺส อุปฺปาทาย อุปฺปนฺนสฺส วา ปสฺสทฺธิ\-สมฺโพชฺฌงฺคสฺส ภาวนาย ปาริปูริยา.}

\prose{โก จ ภิกฺขเว อาหาโร อนุปฺปนฺนสฺส วา สมาธิ\-สมฺโพชฺฌงฺคสฺส อุปฺปาทาย อุปฺปนฺนสฺส วา สมาธิ\-สมฺโพชฺฌงฺคสฺส ภาวนาย ปาริปูริยา, อตฺถิ ภิกฺขเว สมถ\-นิมิตฺตํ\csromansharedfootnote{c98c6d9bea9d0f49}{สมาธินิมิตฺตํ (สฺยา)} อพฺยคฺค\-นิมิตฺตํ, ตตฺถ โยนิโสมนสิการ\-พหุลีกาโร, อยมาหาโร อนุปฺปนฺนสฺส วา สมาธิ\-สมฺโพชฺฌงฺคสฺส อุปฺปาทาย อุปฺปนฺนสฺส วา สมาธิ\-สมฺโพชฺฌงฺคสฺส ภาวนาย ปาริปูริยา.}

\prose{โก จ ภิกฺขเว อาหาโร อนุปฺปนฺนสฺส วา อุเปกฺขา\-สมฺโพชฺฌงฺคสฺส อุปฺปาทาย อุปฺปนฺนสฺส วา อุเปกฺขา\-สมฺโพชฺฌงฺคสฺส ภาวนาย ปาริปูริยา, อตฺถิ ภิกฺขเว อุเปกฺขา\-สมฺโพชฺฌงฺค\-ฏฺฐานียา ธมฺมา, ตตฺถ โยนิโสมนสิการ\-พหุลีกาโร, อยมาหาโร อนุปฺปนฺนสฺส วา อุเปกฺขาสมฺโพชฺชงฺคสฺส อุปฺปาทาย อุปฺปนฺนสฺส วา อุเปกฺขา\-สมฺโพชฺฌงฺคสฺส ภาวนาย ปาริปูริยา.}

\prose{เสยฺยถาปิ ภิกฺขเว อยํ กาโย อาหารฏฺฐิติโก อาหารํ ปฏิจฺจ ติฏฺฐติ, อนาหาโร โน ติฏฺฐติ. เอวเมว โข ภิกฺขเว อิเม สตฺต โพชฺฌงฺคา อาหารฏฺฐิติกา อาหารํ ปฏิจฺจ ติฏฺฐนฺติ, อนาหารา โน ติฏฺฐนฺตีติ.  ทุติยํ.}

\csromanheader{2}{3. \textbf{สีลสุตฺต}}
\proseitem{184}{เย เต ภิกฺขเว ภิกฺขู สีลสมฺปนฺนา สมาธิ\-สมฺปนฺนา ญาณสมฺปนฺนา วิมุตฺติ\-สมฺปนฺนา วิมุตฺติ\-ญาณ\-ทสฺสน\-สมฺปนฺนา, ทสฺสนมฺปาหํ ภิกฺขเว เตสํ ภิกฺขูนํ พหุการํ\csromansharedfootnote{5512350a5ae45dae}{พหูปการํ (สฺยา)} วทามิ, สวนมฺปาหํ ภิกฺขเว เตสํ ภิกฺขูนํ พหุการํ วทามิ, อุปสงฺกมนมฺปาหํ ภิกฺขเว เตสํ ภิกฺขูนํ พหุการํ วทามิ, ปยิรุปาสนมฺปาหํ ภิกฺขเว เตสํ ภิกฺขูนํ พหุการํ วทามิ, อนุสฺสติมฺปาหํ ภิกฺขเว เตสํ ภิกฺขูนํ พหุการํ วทามิ, อนุปพฺพชฺชมฺปาหํ ภิกฺขเว เตสํ ภิกฺขูนํ พหุการํ วทามิ. ตํ กิสฺส เหตุ, ตถารูปานํ ภิกฺขเว ภิกฺขูนํ ธมฺมํ สุตฺวา ทฺวเยน วูปกาเสน วูปกฏฺโฐ\csromansharedfootnote{1a0a171eba2fb7ed}{ทฺวเยน วูปกฏฺโฐ (สี, สฺยา)} วิหรติ กายวูปกาเสน จ จิตฺต\-วูปกาเสน จ, โส ตถา วูปกฏฺโฐ วิหรนฺโต ตํ ธมฺมํ อนุสฺสรติ อนุวิตกฺเกติ.}

\prose{\csromanfolio{62}ยสฺมิํ สมเย ภิกฺขเว ภิกฺขุ ตถา วูปกฏฺโฐ วิหรนฺโต ตํ ธมฺมํ อนุสฺสรติ อนุวิตกฺเกติ, สติสมฺโพชฺฌงฺโค ตสฺมิํ สมเย ภิกฺขุโน อารทฺโธ โหติ, สติสมฺโพชฺฌงฺคํ ตสฺมิํ สมเย ภิกฺขุ ภาเวติ, สติสมฺโพชฺฌงฺโค ตสฺมิํ สมเย ภิกฺขุโน ภาวนา\-ปาริปูริํ คจฺฉติ, โส ตถา สโต วิหรนฺโต ตํ ธมฺมํ ปญฺญาย ปวิจินติ ปวิจรติ ปริวีมํสมา\-ปชฺชติ.}

\prose{ยสฺมิํ สมเย ภิกฺขเว ภิกฺขุ ตถา สโต วิหรนฺโต ตํ ธมฺมํ ปญฺญาย ปวิจินติ ปวิจรติ ปริวีมํสมา\-ปชฺชติ, ธมฺมวิจย\-สมฺโพชฺฌงฺโค ตสฺมิํ สมเย ภิกฺขุโน อารทฺโธ โหติ, ธมฺมวิจย\-สมฺโพชฺฌงฺคํ ตสฺมิํ สมเย ภิกฺขุ ภาเวติ, ธมฺมวิจย\-สมฺโพชฺฌงฺโค ตสฺมิํ สมเย ภิกฺขุโน ภาวนา\-ปาริปูริํ คจฺฉติ. ตสฺส ตํ ธมฺมํ ปญฺญาย ปวิจินโต ปวิจรโต ปริวีมํสมาปชฺชโต อารทฺธํ โหติ วีริยํ อสลฺลีนํ.}

\prose{ยสฺมิํ สมเย ภิกฺขเว ภิกฺขุโน ตํ ธมฺมํ ปญฺญาย ปวิจินโต ปวิจรโต ปริวีมํสมาปชฺชโต อารทฺธํ โหติ วีริยํ อสลฺลีนํ, วีริย\-สมฺโพชฺฌงฺโค ตสฺมิํ สมเย ภิกฺขุโน อารทฺโธ โหติ. วีริย\-สมฺโพชฺฌงฺคํ ตสฺมิํ สมเย ภิกฺขุ ภาเวติ, วีริย\-สมฺโพชฺฌงฺโค ตสฺมิํ สมเย ภิกฺขุโน ภาวนา\-ปาริปูริํ คจฺฉติ, อารทฺธ\-วีริยสฺส อุปฺปชฺชติ ปีติ นิรามิสา.}

\prose{ยสฺมิํ สมเย ภิกฺขเว ภิกฺขุโน อารทฺธ\-วีริยสฺส อุปฺปชฺชติ ปีติ นิรามิสา, ปีติสมฺโพชฺฌงฺโค ตสฺมิํ สมเย ภิกฺขุโน อารทฺโธ โหติ, ปีติสมฺโพชฺฌงฺคํ ตสฺมิํ สมเย ภิกฺขุ ภาเวติ, ปีติสมฺโพชฺฌงฺโค ตสฺมิํ สมเย ภิกฺขุโน ภาวนา\-ปาริปูริํ คจฺฉติ, ปีติมนสฺส กาโยปิ ปสฺสมฺภติ จิตฺตมฺปิ ปสฺสมฺภติ, ยสฺมิํ สมเย ภิกฺขเว ภิกฺขุโน ปีติมนสฺส กาโยปิ ปสฺสมฺภติ จิตฺตมฺปิ ปสฺสมฺภติ. ปสฺสทฺธิ\-สมฺโพชฺฌงฺโค ตสฺมิํ สมเย ภิกฺขุโน อารทฺโธ โหติ, ปสฺสทฺธิ\-สมฺโพชฺฌงฺคํ ตสฺมิํ สมเย ภิกฺขุ ภาเวติ, ปสฺสทฺธิ\-สมฺโพชฺฌงฺโค ตสฺมิํ สมเย ภิกฺขุโน ภาวนา\-ปาริปูริํ คจฺฉติ, ปสฺสทฺธ\-กายสฺส สุขิโน จิตฺตํ สมาธิยติ.}

\prose{ยสฺมิํ สมเย ภิกฺขเว ภิกฺขุโน ปสฺสทฺธ\-กายสฺส สุขิโน จิตฺตํ สมาธิยติ, สมาธิ\-สมฺโพชฺฌงฺโค ตสฺมิํ สมเย ภิกฺขุโน อารทฺโธ}

\prose{\csromanfolio{63}โหติ, สมาธิ\-สมฺโพชฺฌงฺคํ ตสฺมิํ สมเย ภิกฺขุ ภาเวติ, สมาธิ\-สมฺโพชฺฌงฺโค ตสฺมิํ สมเย ภิกฺขุโน ภาวนา\-ปาริปูริํ คจฺฉติ, โส ตถาสมาหิตํ จิตฺตํ สาธุกํ อชฺฌุเปกฺขิตา โหติ.}

\prose{ยสฺมิํ สมเย ภิกฺขเว ภิกฺขุ ตถาสมาหิตํ จิตฺตํ สาธุกํ อชฺฌุเปกฺขิตา โหติ, อุเปกฺขา\-สมฺโพชฺฌงฺโค ตสฺมิํ สมเย ภิกฺขุโน อารทฺโธ โหติ, อุเปกฺขา\-สมฺโพชฺฌงฺคํ ตสฺมิํ สมเย ภิกฺขุ ภาเวติ, อุเปกฺขา\-สมฺโพชฺฌงฺโค ตสฺมิํ สมเย ภิกฺขุโน ภาวนา\-ปาริปูริํ คจฺฉติ.}

\prose{เอวํ ภาวิเตสุ โข ภิกฺขเว สตฺตสุ สมฺโพชฺฌงฺเคสุ เอวํ พหุลีกเตสุ สตฺต ผลา สตฺตานิสํสา ปาฏิกงฺขา. กตเม สตฺต ผลา สตฺตานิสํสา, ทิฏฺเฐว ธมฺเม ปฏิกจฺจ อญฺญํ อาราเธติ. โน เจ ทิฏฺเฐว ธมฺเม ปฏิกจฺจ อญฺญํ อาราเธติ, อถ มรณกาเล อญฺญํ อาราเธติ. โน เจ ทิฏฺเฐว ธมฺเม ปฏิกจฺจ อญฺญํ อาราเธติ, โน เจ มรณกาเล อญฺญํ อาราเธติ, อถ ปญฺจนฺนํ โอรมฺภาคิยานํ สํโยชนานํ ปริกฺขยา อนฺตรา\-ปรินิพฺพายี โหติ. โน เจ ทิฏฺเฐว ธมฺเม ปฏิกจฺจ อญฺญํ อาราเธติ, โน เจ มรณกาเล อญฺญํ อาราเธติ, โน เจ ปญฺจนฺนํ โอรมฺภาคิยานํ สํโยชนานํ ปริกฺขยา อนฺตรา\-ปรินิพฺพายี โหติ, อถ ปญฺจนฺนํ โอรมฺภาคิยานํ สํโยชนานํ ปริกฺขยา อุปหจฺจ\-ปรินิพฺพายี โหติ. โน เจ ทิฏฺเฐว ธมฺเม ปฏิกจฺจ อญฺญํ อาราเธติ, โน เจ มรณกาเล อญฺญํ อาราเธติ, โน เจ ปญฺจนฺนํ โอรมฺภาคิยานํ สํโยชนานํ ปริกฺขยา อนฺตรา\-ปรินิพฺพายี โหติ, โน เจ ปญฺจนฺนํ โอรมฺภาคิยานํ สํโยชนานํ ปริกฺขยา อุปหจฺจ\-ปรินิพฺพายี โหติ, อถ ปญฺจนฺนํ โอรมฺภาคิยานํ สํโยชนานํ ปริกฺขยา อสงฺขารปรินิพฺพายี โหติ. โน เจ ทิฏฺเฐว ธมฺเม ปฏิกจฺจ อญฺญํ อารเธติ, โน เจ มรณกาเล อญฺญํ อาราเธติ, โน เจ ปญฺจนฺนํ โอรมฺภาคิยานํ สํโยชนานํ ปริกฺขยา อนฺตรา\-ปรินิพฺพายี โหติ, โน เจ ปญฺจนฺนํ โอรมฺภาคิยานํ สํโยชนานํ ปริกฺขยา อุปหจฺจ\-ปรินิพฺพายี โหติ, โน เจ ปญฺจนฺนํ โอรมฺภาคิยานํ สํโยชนานํ ปริกฺขยา อสงฺขารปรินิพฺพายี โหติ, อถ ปญฺจนฺนํ โอรมฺภาคิยานํ สํโยชนานํ ปริกฺขยา สสงฺขาร\-ปรินิพฺพายี โหติ. โน เจ ทิฏฺเฐว ธมฺเม ปฏิกจฺจ อญฺญํ อาราเธติ, โน เจ มรณกาเล อญฺญํ อาราเธติ, โน เจ ปญฺจนฺนํ โอรมฺภาคิยานํ สํโยชนานํ ปริกฺขยา อนฺตรา\-ปรินิพฺพายี}

\prose{\csromanfolio{64}โหติ, โน เจ ปญฺจนฺนํ โอรมฺภาคิยานํ สํโยชนานํ ปริกฺขยา อุปหจฺจ\-ปรินิพฺพายี โหติ, โน เจ ปญฺจนฺนํ โอรมฺภาคิยานํ สํโยชนานํ ปริกฺขยา อสงฺขารปรินิพฺพายี โหติ, โน เจ ปญฺจนฺนํ โอรมฺภาคิยานํ สํโยชนานํ ปริกฺขยา สสงฺขาร\-ปรินิพฺพายี โหติ, อถ ปญฺจนฺนํ โอรมฺภาคิยานํ สํโยชนานํ ปริกฺขยา อุทฺธํโสโต โหติ อกนิฏฺฐคามี, เอวํ ภาวิเตสุ โข ภิกฺขเว สตฺตสุ โพชฺฌงฺเคสุ เอวํ พหุลีกเตสุ อิเม สตฺต ผลา สตฺตานิสํสา ปาฏิกงฺขาติ.  ตติยํ.}

\csromanheader{2}{4. \textbf{วตฺถสุตฺต}}
\proseitem{185}{เอกํ สมยํ อายสฺมา สาริปุตฺโต สาวตฺถิยํ วิหรติ เชตวเน อนาถ\-ปิณฺฑิกสฺส อาราเม. ตตฺร โข อายสฺมา สาริปุตฺโต ภิกฺขู อามนฺเตสิ “อาวุโส ภิกฺขโว”ติ. “อาวุโส”ติ โข เต ภิกฺขู อายสฺมโต สาริปุตฺตสฺส ปจฺจสฺโสสุํ. อายสฺมา สาริปุตฺโต เอตทโวจ\csromandash{}}

\prose{สตฺติเม อาวุโส โพชฺฌงฺคา. กตเม สตฺต, สติสมฺโพชฺฌงฺโค ธมฺมวิจย\-สมฺโพชฺฌงฺโค วีริย\-สมฺโพชฺฌงฺโค ปีติสมฺโพชฺฌงฺโค ปสฺสทฺธิ\-สมฺโพชฺฌงฺโค สมาธิ\-สมฺโพชฺฌงฺโค อุเปกฺขา\-สมฺโพชฺฌงฺโค, อิเม โข อาวุโส สตฺต โพชฺฌงฺคา. อิเมสํ ขฺวาหํ อาวุโส สตฺตนฺนํ โพชฺฌงฺคานํ เยน เยน โพชฺฌงฺเคน อากงฺขามิ ปุพฺพณฺหสมยํ วิหริตุํ, เตน เตน โพชฺฌงฺเคน ปุพฺพณฺหสมยํ วิหรามิ. เยน เยน โพชฺฌงฺเคน อากงฺขามิ มชฺฌนฺหิกํ สมยํ วิหริตุํ, เตน เตน โพชฺฌงฺเคน มชฺฌนฺหิกํ สมยํ วิหรามิ. เยน เยน โพชฺฌงฺเคน อากงฺขามิ สายนฺหสมยํ วิหริตุํ, เตน เตน โพชฺฌงฺเคน สายนฺหสมยํ วิหรามิ. สติสมฺโพชฺฌงฺโค อิติ เจ เม อาวุโส โหติ, อปฺปมาโณติ เม โหติ, สุสมารทฺโธติ เม โหติ, ติฏฺฐนฺตํ จ นํ ติฏฺฐตีติ ปชานามิ, สเจปิ เม จวติ, อิทปฺปจฺจยา เม จวตีติ ปชานามิ ฯเปฯ. อุเปกฺขา\-สมฺโพชฺฌงฺโค อิติ เจ เม อาวุโส โหติ, อปฺปมาโณติ เม โหติ, สุสมารทฺโธติ เม โหติ, ติฏฺฐนฺตํ จ นํ ติฏฺฐตีติ ปชานามิ, สเจปิ เม จวติ, อิทปฺปจฺจยา เม จวตีติ ปชานามิ.}

\prose{เสยฺยถาปิ อาวุโส รญฺโญ วา ราช\-มหามตฺตสฺส วา นานารตฺตานํ ทุสฺสานํ ทุสฺสกรณฺฑโก ปูโร อสฺส, โส ยญฺญเทว ทุสฺสยุคํ}

\prose{\csromanfolio{65}อากงฺเขยฺย ปุพฺพณฺหสมยํ ปารุปิตุํ, ตํ ตเทว ทุสฺสยุคํ ปุพฺพณฺหสมยํ ปารุเปยฺย. ยญฺญเทว ทุสฺสยุคํ อากงฺเขยฺย มชฺฌนฺหิกํ สมยํ ปารุปิตุํ, ตํ ตเทว ทุสฺสยุคํ มชฺฌนฺหิกํ สมยํ ปารุเปยฺย. ยญฺญเทว ทุสฺสยุคํ อากงฺเขยฺย สายนฺหสมยํ ปารุปิตุํ, ตํ ตเทว ทุสฺสยุคํ สายนฺหสมยํ ปารุเปยฺย. เอวเมว ขฺวาหํ อาวุโส อิเมสํ สตฺตนฺนํ โพชฺฌงฺคานํ เยน เยน โพชฺฌงฺเคน อากงฺขามิ ปุพฺพณฺหสมยํ วิหริตุํ, เตน เตน โพชฺฌงฺเคน ปุพฺพณฺหสมยํ วิหรามิ. เยน เยน โพชฺฌงฺเคน อากงฺขามิ มชฺฌนฺหิกํ สมยํ วิหริตุํ, เตน เตน โพชฺฌงฺเคน มชฺฌนฺหิกํ สมยํ วิหรามิ. เยน เยน โพชฺฌงฺเคน อากงฺขามิ สายนฺหสมยํ วิหริตุํ, เตน เตน โพชฺฌงฺเคน สายนฺหสมยํ วิหรามิ. สติสมฺโพชฺฌงฺโค อิติ เจ เม อาวุโส โหติ, อปฺปมาโณติ เม โหติ, สุสมารทฺโธติ เม โหติ. ติฏฺฐนฺตํ จ นํ ติฏฺฐตีติ ปชานามิ, สเจปิ เม จวติ, อิทปฺปจฺจยา เม จวตีติ ปชานามิ ฯเปฯ. อุเปกฺขา\-สมฺโพชฺฌงฺโค อิติ เจ เม อาวุโส โหติ, อปฺปมาโณติ เม โหติ, สุสมารทฺโธติ เม โหติ, ติฏฺฐนฺตํ จ นํ ติฏฺฐตีติ ปชานามิ, สเจปิ เม จวติ, อิทปฺปจฺจยา เม จวตีติ ปชานามี”ติ.  จตุตฺถํ.}

\csromanheader{2}{5. \textbf{ภิกฺขุสุตฺต}}
\proseitem{186}{สาวตฺถิ\-นิทานํ. อถ โข อญฺญตโร ภิกฺขุ เยน ภควา เตนุปสงฺกมิ ฯเปฯ เอกมนฺตํ นิสินฺโน โข โส ภิกฺขุ ภควนฺตํ เอตทโวจ “โพชฺฌงฺคา โพชฺฌงฺคาติ ภนฺเต วุจฺจนฺติ, กิตฺตาวตา นุ โข ภนฺเต โพชฺฌงฺคาติ วุจฺจนฺตี”ติ. โพธาย สํวตฺตนฺตีติ โข ภิกฺขุ ตสฺมา โพชฺฌงฺคาติ วุจฺจนฺติ. อิธ ภิกฺขุ สติสมฺโพชฺฌงฺคํ ภาเวติ วิเวกนิสฺสิตํ วิราคนิสฺสิตํ นิโรธ\-นิสฺสิตํ โวสฺสคฺค\-ปริณามิํ ฯเปฯ. อุเปกฺขา\-สมฺโพชฺฌงฺคํ ภาเวติ วิเวกนิสฺสิตํ วิราคนิสฺสิตํ นิโรธ\-นิสฺสิตํ โวสฺสคฺค\-ปริณามิํ. ตสฺสิเม สตฺต โพชฺฌงฺเค ภาวยโต กามาสวาปิ จิตฺตํ วิมุจฺจติ, ภวาสวาปิ จิตฺตํ วิมุจฺจติ, อวิชฺชาสวาปิ จิตฺตํ วิมุจฺจติ, วิมุตฺตสฺมิํ “วิมุตฺตมฺ”อิติ ญาณํ โหติ, “ขีณา ชาติ, วุสิตํ พฺรหฺมจริยํ, กตํ กรณียํ, นาปรํ อิตฺถตฺตายา”ติ ปชานาติ. โพธาย สํวตฺตนฺตีติ ภิกฺขุ ตสฺมา โพชฺฌงฺคาติ วุจฺจนฺตีติ.  ปญฺจมํ.}

\csromanheader{2}{6. \textbf{กุณฺฑลิยสุตฺต}}
\proseitem{187}{\csromanfolio{66}เอกํ สมยํ ภควา สาเกเต วิหรติ อญฺจนวเน มิคทาเย. อถ โข กุณฺฑลิโย ปริพฺพาชโก เยน ภควา เตนุปสงฺกมิ, อุปสงฺกมิตฺวา ภควตา สทฺธิํ สมฺโมทิ, สมฺโมทนียํ กถํ สารณียํ วีติสาเรตฺวา เอกมนฺตํ นิสีทิ, เอกมนฺตํ นิสินฺโน โข กุณฺฑลิโย ปริพฺพาชโก ภควนฺตํ เอตทโวจ “อหมสฺมิ โภ โคตม อารามนิสฺสิยี\csromansharedfootnote{5b16dc44270c3b63}{อารามนิสาที (สี), อารามนิยาที (สฺยา)} ปริสาวจโร, ตสฺส มยฺหํ โภ โคตม ปจฺฉาภตฺตํ ภุตฺต\-ปาตราสสฺส อยมาจาโร\csromansharedfootnote{26a84cd2e0bf017c}{อยมาหาโร (สฺยา, ก)} โหติ ‘อาราเมน อารามํ อุยฺยาเนน อุยฺยานํ อนุจงฺกมามิ อนุวิจรามิ, โส ตตฺถ ปสฺสามิ เอเก สมณพฺราหฺมเณ อิติวาทปฺปโมกฺขานิสํส\-ญฺเจว กถํ กเถนฺเต อุปารมฺภานิสํสญฺจ, ภวํ ปน โคตโม กิมานิสํโส วิหรตี”ติ. วิชฺชาวิมุตฺติ\-ผลา\-นิสํโส โข กุณฺฑลิย ตถาคโต วิหรตีติ.}

\prose{กตเม ปน โภ โคตม ธมฺมา ภาวิตา พหุลีกตา วิชฺชาวิมุตฺติํ ปริปูเรนฺตีติ. สตฺต โข กุณฺฑลิย โพชฺฌงฺคา ภาวิตา พหุลีกตา วิชฺชาวิมุตฺติํ ปริปูเรนฺตีติ. กตเม ปน โภ โคตม ธมฺมา ภาวิตา พหุลีกตา สตฺต โพชฺฌงฺเค ปริปูเรนฺตีติ, จตฺตาโร โข กุณฺฑลิย สติปฏฺฐานา ภาวิตา พหุลีกตา สตฺต โพชฺฌงฺเค ปริปูเรนฺตีติ. กตเม ปน โภ โคตม ธมฺมา ภาวิตา พหุลีกตา จตฺตาโร สติปฏฺฐาเน ปริปูเรนฺตีติ, ตีณิ โข กุณฺฑลิย สุจริตานิ ภาวิตานิ พหุลีกตานิ จตฺตาโร สติปฏฺฐาเน ปริปูเรนฺตีติ. กตเม ปน โภ โคตม ธมฺมา ภาวิตา พหุลีกตา ตีณิ สุจริตานิ ปริปูเรนฺตีติ, อินฺทฺริยสํวโร โข กุณฺฑลิย ภาวิโต พหุลีกโต ตีณิ สุจริตานิ ปริปูเรตีติ.}

\prose{กถํ ภาวิโต จ กุณฺฑลิย อินฺทฺริยสํวโร กถํ พหุลีกโต ตีณิ สุจริตานิ ปริปูเรนฺตีติ, อิธ กุณฺฑลิย ภิกฺขุ จกฺขุนา รูปํ ทิสฺวา มนาปํ นาภิชฺฌติ นาภิหํสติ น ราคํ ชเนติ, ตสฺส ฐิโต จ กาโย โหติ, ฐิตํ จิตฺตํ อชฺฌตฺตํ สุสณฺฐิตํ สุวิมุตฺตํ. จกฺขุนา โข ปเนว รูปํ ทิสฺวา อมนาปํ น มงฺกุ โหติ อปฺปติฏฺฐิต\-จิตฺโต อทีนมานโส อพฺยาปนฺน\-เจตโส. ตสฺส ฐิโต จ กาโย โหติ, ฐิตํ จิตฺตํ อชฺฌตฺตํ สุสณฺฐิตํ สุวิมุตฺตํ.}

\prose{\csromanfolio{67}ปุน จปรํ กุณฺฑลิย ภิกฺขุ โสเตน สทฺทํ สุตฺวา ฯเปฯ ฆาเนน คนฺธํ ฆายิตฺวา ฯเปฯ ชิวฺหาย รสํ สายิตฺวา ฯเปฯ กาเยน โผฏฺฐพฺพํ ผุสิตฺวา ฯเปฯ มนสา ธมฺมํ วิญฺญาย มนาปํ นาภิชฺฌติ นาภิหํสติ น ราคํ ชเนติ. ตสฺส ฐิโต จ กาโย โหติ ฐิตํ จิตฺตํ อชฺฌตฺตํ สุสณฺฐิตํ สุวิมุตฺตํ. มนสา โข ปเนว ธมฺมํ วิญฺญาย อมนาปํ น มงฺกุ โหติ อปฺปติฏฺฐิต\-จิตฺโต อทีนมานโส อพฺยาปนฺน\-เจตโส. ตสฺส ฐิโต จ กาโย โหติ, ฐิตํ จิตฺตํ อชฺฌตฺตํ สุสณฺฐิตํ สุวิมุตฺตํ.}

\prose{ยโต โข กุณฺฑลิย ภิกฺขุโน จกฺขุนา รูปํ ทิสฺวา มนาปามนาเปสุ รูเปสุ ฐิโต จ กาโย โหติ, ฐิตํ จิตฺตํ อชฺฌตฺตํ สุสณฺฐิตํ สุวิมุตฺตํ. โสเตน สทฺทํ สุตฺวา ฯเปฯ. ฆาเนน คนฺธํ ฆายิตฺวา ฯเปฯ. ชิวฺหาย รสํ สายิตฺวา ฯเปฯ. กาเยน โผฏฺฐพฺพํ ผุสิตฺวา ฯเปฯ. มนสา ธมฺมํ วิญฺญาย มนาปามนาเปสุ ธมฺเมสุ ฐิโต จ กาโย โหติ, ฐิตํ จิตฺตํ อชฺฌตฺตํ สุสณฺฐิตํ สุวิมุตฺตํ. เอวํ ภาวิโต โข กุณฺฑลิย อินฺทฺริยสํวโร เอวํ พหุลีกโต ตีณิ สุจริตานิ ปริปูเรติ.}

\prose{กถํ ภาวิตานิ จ กุณฺฑลิย ตีณิ สุจริตานิ กถํ พหุลีกตานิ จตฺตาโร สติปฏฺฐาเน ปริปูเรนฺติ, อิธ กุณฺฑลิย ภิกฺขุ กายทุจฺจริตํ ปหาย กายสุจริตํ ภาเวติ, วจีทุจฺจริตํ ปหาย วจีสุจริตํ ภาเวติ, มโนทุจฺจริตํ ปหาย มโนสุจริตํ ภาเวติ. เอวํ ภาวิตานิ โข กุณฺฑลิย ตีณิ สุจริตานิ เอวํ พหุลีกตานิ จตฺตาโร สติปฏฺฐาเน ปริปูเรนฺติ.}

\prose{กถํ ภาวิตา จ กุณฺฑลิย จตฺตาโร สติปฏฺฐานา กถํ พหุลีกตา สตฺต โพชฺฌงฺเค ปริปูเรนฺติ, อิธ กุณฺฑลิย ภิกฺขุ กาเย กายานุปสฺสี วิหรติ อาตาปี สมฺปชาโน สติมา วิเนยฺย โลเก อภิชฺฌา\-โทมนสฺสํ. เวทนาสุ ฯเปฯ. จิตฺเต ฯเปฯ. ธมฺเมสุ ธมฺมานุปสฺสี วิหรติ อาตาปี สมฺปชาโน สติมา วิเนยฺย โลเก อภิชฺฌา\-โทมนสฺสํ. เอวํ ภาวิตา โข กุณฺฑลิย จตฺตาโร สติปฏฺฐานา เอวํ พหุลีกตา สตฺต โพชฺฌงฺเค ปริปูเรนฺติ.}

\prose{กถํ ภาวิตา จ กุณฺฑลิย สตฺต โพชฺฌงฺคา กถํ พหุลีกตา วิชฺชาวิมุตฺติํ ปริปูเรนฺติ, อิธ กุณฺฑลิย ภิกฺขุ สติสมฺโพชฺฌงฺคํ ภาเวติ วิเวกนิสฺสิตํ วิราคนิสฺสิตํ นิโรธ\-นิสฺสิตํ โวสฺสคฺค\-ปริณามิํ ฯเปฯ.}

\prose{\csromanfolio{68}อุเปกฺขาสมฺโพชฺฌางฺคํ ภาเวติ วิเวกนิสฺสิตํ วิราคนิสฺสิตํ นิโรธ\-นิสฺสิตํ โวสฺสคฺค\-ปริณามิํ. เอวํ ภาวิตา โข กุณฺฑลิย สตฺต โพชฺฌงฺคา เอวํ พหุลีกตา วิชฺชาวิมุตฺติํ ปริปูเรนฺตีติ.}

\prose{เอวํ วุตฺเต กุณฺฑลิโย ปริพฺพาชโก ภควนฺตํ เอตทโวจ “อภิกฺกนฺตํ โภ โคตม, อภิกฺกนฺตํ โภ โคตม, เสยฺยถาปิ โภ โคตม นิกฺกุชฺชิตํ วา อุกฺกุชฺเชยฺย, ปฏิจฺฉนฺนํ วา วิจเรยฺย, มูฬฺหสฺส วา มคฺคํ อาจิกฺเขยฺย, อนฺธกาเร วา เตลปชฺโชตํ ธาเรยฺย ‘จกฺขุมนฺโต รูปานิ ทกฺขนฺตี’ติ, เอวเมว โภตา โคตเมน อเนก\-ปริยาเยน ธมฺโม ปกาสิโต, เอสาหํ ภวนฺตํ โคตมํ สรณํ คจฺฉามิ ธมฺมญฺจ ภิกฺขุสํฆญฺจ, อุปาสกํ มํ ภวํ โคตโม ธาเรตุ อชฺชตคฺเค ปาณุเปตํ สรณํ คตนฺ”ติ.  ฉฏฺฐํ.}

\csromanheader{2}{7. \textbf{กูฏาคารสุตฺต}}
\proseitem{188}{เสยฺยถาปิ ภิกฺขเว กูฏาคารสฺส ยา กาจิ โคปานสิโย, สพฺพา ตา กูฏนินฺนา กูฏโปณา กูฏปพฺภารา. เอวเมว โข ภิกฺขเว ภิกฺขุ สตฺต โพชฺฌงฺเค ภาเวนฺโต สตฺต โพชฺฌงฺเค พหุลีกโรนฺโต นิพฺพานนินฺโน โหติ นิพฺพานโปโณ นิพฺพาน\-ปพฺภาโร.}

\prose{กถญฺจ ภิกฺขเว ภิกฺขุ สตฺต โพชฺฌงฺเค ภาเวนฺโต สตฺต โพชฺฌงฺเค พหุลีกโรนฺโต นิพฺพานนินฺโน โหติ นิพฺพานโปโณ นิพฺพาน\-ปพฺภาโร, อิธ ภิกฺขเว ภิกฺขุ สติสมฺโพชฺฌงฺคํ ภาเวติ วิเวกนิสฺสิตํ วิราคนิสฺสิตํ นิโรธ\-นิสฺสิตํ โวสฺสคฺค\-ปริณามิํ ฯเปฯ. อุเปกฺขา\-สมฺโพชฺฌงฺคํ ภาเวติ วิเวกนิสฺสิตํ วิราคนิสฺสิตํ นิโรธ\-นิสฺสิตํ โวสฺสคฺค\-ปริณามิํ. เอวํ โข ภิกฺขเว ภิกฺขุ สตฺต โพชฺฌงฺเค ภาเวนฺโต สตฺต โพชฺฌงฺเค พหุลีกโรนฺโต นิพฺพานนินฺโน โหติ นิพฺพานโปโณ นิพฺพาน\-ปพฺภาโรติ.  สตฺตมํ.}

\csromanheader{2}{8. \textbf{อุปวานสุตฺต}}
\proseitem{189}{เอกํ สมยํ อายสฺมา จ อุปวาโน อายสฺมา จ สาริปุตฺโต โกสมฺพิยํ วิหรนฺติ โฆสิตาราเม. อถ โข อายสฺมา สาริปุตฺโต สายนฺหสมยํ ปฏิสลฺลานา วุฏฺฐิโต เยนายสฺมา อุปวาโน เตนุปสงฺกมิ, อุปสงฺกมิตฺวา อายสฺมตา อุปวาเนน สทฺธิํ สมฺโมทิ, สมฺโมทนียํ กถํ สารณียํ วีติสาเรตฺวา เอกมนฺตํ นิสีทิ, เอกมนฺตํ นิสินฺโน โข อายสฺมา สาริปุตฺโต อายสฺมนฺตํ อุปวานํ เอตทโวจ\csromandash{}}

\prose{\csromanfolio{69}“ชาเนยฺย นุ โข อาวุโส อุปวาน ภิกฺขุ ปจฺจตฺตํ โยนิโส\-มนสิการา เอวํ สุสมารทฺธา เม สตฺต โพชฺฌงฺคา ผาสุวิหาราย สํวตฺตนฺตี”ติ. ชาเนยฺย โข อาวุโส สาริปุตฺต ภิกฺขุ ปจฺจตฺตํ โยนิโส\-มนสิการา เอวํ สุสมารทฺธา เม สตฺต โพชฺฌงฺคา ผาสุวิหาราย สํวตฺตนฺตีติ.}

\prose{สติสมฺโพชฺฌงฺคํ โข อาวุโส ภิกฺขุ อารพฺภมาโน ปชานาติ “จิตฺตญฺจ เม สุวิมุตฺตํ, ถินมิทฺธญฺจ เม สุสมูหตํ, อุทฺธจฺจ\-กุกฺกุจฺจญฺจ เม สุปฺปฏิวินีตํ, อารทฺธญฺจ เม วีริยํ, อฏฺฐิํกตฺวา มนสิ กโรมิ โน จ ลีนนฺ”ติ ฯเปฯ. อุเปกฺขา\-สมฺโพชฺฌงฺคํ อาวุโส ภิกฺขุ อารพฺภมาโน ปชานาติ “จิตฺตญฺจ เม สุวิมุตฺตํ, ถินมิทฺธญฺจ เม สุสมูหตํ, อุทฺธจฺจ\-กุกฺกุจฺจญฺจ เม สุปฺปฏิวินีตํ, อารทฺธญฺจ เม วีริยํ, อฏฺฐิํกตฺวา มนสิ กโรมิ โน จ ลีนนฺ”ติ. เอวํ โข อาวุโส สาริปุตฺต ภิกฺขุ ชาเนยฺย ปจฺจตฺตํ โยนิโส\-มนสิการา เอวํ สุสมารทฺธา เม สตฺต โพชฺฌงฺคา ผาสุวิหาราย สํวตฺตนฺตีติ.  อฏฺฐมํ.}

\prose{9. ปฐมอุปฺปนฺนสุตฺต}

\proseitem{190}{สตฺติเม ภิกฺขเว โพชฺฌงฺคา ภาวิตา พหุลีกตา อนุปฺปนฺนา อุปฺปชฺชนฺติ นาญฺญตฺร ตถาคตสฺส ปาตุภาวา อรหโต สมฺมา\-สมฺพุทฺธสฺส. กตเม สตฺต, สติสมฺโพชฺฌงฺโค ฯเปฯ อุเปกฺขา\-สมฺโพชฺฌงฺโค. อิเม โข ภิกฺขเว สตฺต โพชฺฌงฺคา ภาวิตา พหุลีกตา อนุปฺปนฺนา อุปฺปชฺชนฺติ นาญฺญตฺร ตถาคตสฺส ปาตุภาวา อรหโต สมฺมา\-สมฺพุทฺธสฺสาติ.  นวมํ.}

\prose{10. ทุติยอุปฺปนฺนสุตฺต}

\proseitemwithclosermidruled{191}{สตฺติเม ภิกฺขเว โพชฺฌงฺคา ภาวิตา พหุลีกตา อนุปฺปนฺนา อุปฺปชฺชนฺติ นาญฺญตฺร สุคตวินยา. กตเม สตฺต, สติสมฺโพชฺฌงฺโค ฯเปฯ อุเปกฺขา\-สมฺโพชฺฌงฺโค. อิเม โข ภิกฺขเว สตฺต โพชฺฌงฺคา ภาวิตา พหุลีกตา อนุปฺปนฺนา อุปฺปชฺชนฺติ นาญฺญตฺร สุคตวินยาติ.  ทสมํ.}{ปพฺพตวคฺโค ปฐโม.}
\csromancenter{\textbf{ตสฺสุทฺทานํ}}
\setlength{\csromangathapagewidth}{0pt}
\setlength{\csromangathawakwidth}{0pt}
\csromangathameasuregroup{หิมวนฺตํ กายํ สีลํ, \\ กูฏญฺจ อุปวานญฺจ,}{\csromangathabat{หิมวนฺตํ กายํ สีลํ,}{วตฺถํ ภิกฺขุ จ กุณฺฑลิ.} \\ \csromangathabat{กูฏญฺจ อุปวานญฺจ,}{อุปฺปนฺนา อปเร ทุเวติ.}}
\csromangathasetleft{หิมวนฺตํ กายํ สีลํ, \\ กูฏญฺจ อุปวานญฺจ,}
\csromangathagroup{\csromangathastanza{\csromangathabat{หิมวนฺตํ กายํ สีลํ,}{วตฺถํ ภิกฺขุ จ กุณฺฑลิ.} \\ \csromangathabat{กูฏญฺจ อุปวานญฺจ,}{อุปฺปนฺนา อปเร ทุเวติ.}}}

\chapterheadpageread{2. \textbf{คิลานวคฺค}}
\csromanheader{2}{1. \textbf{ปาณสุตฺต}}
\proseitem{192}{\csromanfolio{70}เสยฺยถาปิ ภิกฺขเว เย เกจิ ปาณา จตฺตาโร อิริยาปเถ กปฺเปนฺติ กาเลน คมนํ กาเลน ฐานํ กาเลน นิสชฺชํ กาเลน เสยฺยํ, สพฺเพ เต ปถวิํ นิสฺสาย, ปถวิยํ ปติฏฺฐาย, เอวเมเต จตฺตาโร อิริยาปเถ กปฺเปนฺติ. เอวเมว โข ภิกฺขเว ภิกฺขุ สีลํ นิสฺสาย สีเล ปติฏฺฐาย สตฺต โพชฺฌงฺเค ภาเวติ สตฺต โพชฺฌงฺเค พหุลีกโรติ.}

\prose{กถญฺจ ภิกฺขเว ภิกฺขุ สีลํ นิสฺสาย สีเล ปติฏฺฐาย สตฺต โพชฺฌงฺเค ภาเวติ สตฺต โพชฺฌงฺเค พหุลีกโรติ อิธ ภิกฺขเว ภิกฺขุ สติสมฺโพชฺฌงฺคํ ภาเวติ วิเวกนิสฺสิตํ วิราคนิสฺสิตํ นิโรธ\-นิสฺสิตํ โวสฺสคฺค\-ปริณามิํ ฯเปฯ. อุเปกฺขา\-สมฺโพชฺฌงฺคํ ภาเวติ วิเวกนิสฺสิตํ วิราคนิสฺสิตํ นิโรธ\-นิสฺสิตํ โวสฺสคฺค\-ปริณามิํ. เอวํ โข ภิกฺขเว ภิกฺขุ สีลํ นิสฺสาย สีเล ปติฏฺฐาย สตฺต โพชฺฌงฺเค ภาเวติ สตฺต โพชฺฌงฺเค พหุลีกโรตีติ.  ปฐมํ.}

\csromanheader{2}{2. ปฐม\-สูริยูปม\-สุตฺต}
\proseitem{193}{สูริยสฺส ภิกฺขเว อุทยโต เอตํ ปุพฺพงฺคมํ เอตํ ปุพฺพนิมิตฺตํ, ยทิทํ อรุณุคฺคํ. เอวเมว โข ภิกฺขเว ภิกฺขุโน สตฺตนฺนํ โพชฺฌงฺคานํ อุปฺปาทาย เอตํ ปุพฺพงฺคมํ เอตํ นิมิตฺตํ, ยทิทํ กลฺยาณมิตฺตตา. กลฺยาณมิตฺตสฺเสตํ ภิกฺขเว ภิกฺขุโน ปาฏิกงฺขํ สตฺต โพชฺฌงฺเค ภาเวสฺสติ สตฺต โพชฺฌงฺเค พหุลีกริสฺสติ.}

\prose{กถญฺจ ภิกฺขเว ภิกฺขุ กลฺยาณมิตฺโต สตฺต โพชฺฌงฺเค ภาเวติ สตฺต โพชฺฌงฺเค พหุลีกโรติ, อิธ ภิกฺขเว ภิกฺขุ สติสมฺโพชฺฌงฺคํ ภาเวติ วิเวกนิสฺสิตํ ฯเปฯ. อุเปกฺขา\-สมฺโพชฺฌงฺคํ ภาเวติ วิเวกนิสฺสิตํ วิราคนิสฺสิตํ นิโรธ\-นิสฺสิตํ โวสฺสคฺค\-ปริณามิํ. เอวํ โข ภิกฺขเว ภิกฺขุ กลฺยาณมิตฺโต สตฺต โพชฺฌงฺเค ภาเวติ สตฺต โพชฺฌงฺเค พหุลีกโรตีติ.  ทุติยํ.}

\csromanheader{2}{3. ทุติย\-สูริยูปม\-สุตฺต}
\proseitem{194}{\csromanfolio{71}สูริยสฺส ภิกฺขเว อุทยโต เอตํ ปุพฺพงฺคมํ เอตํ ปุพฺพนิมิตฺตํ, ยทิทํ อรุณุคฺคํ. เอวเมว โข ภิกฺขเว ภิกฺขุโน สตฺตนฺนํ โพชฺฌงฺคานํ อุปฺปาทาย เอตํ ปุพฺพงฺคมํ เอตํ ปุพฺพนิมิตฺตํ, ยทิทํ โยนิโส\-มนสิกาโร. โยนิโสมนสิการสมฺปนฺนสฺเสตํ ภิกฺขเว ภิกฺขุโน ปาฏิกงฺขํ สตฺต โพชฺฌงฺเค ภาเวสฺสติ, สตฺต โพชฺฌงฺเค พหุลีกริสฺสติ.}

\prose{กถญฺจ ภิกฺขเว ภิกฺขุ โยนิโสมนสิการ\-สมฺปนฺโน สตฺต โพชฺฌงฺเค ภาเวติ สตฺต โพชฺฌงฺเค พหุลีกโรติ, อิธ ภิกฺขเว ภิกฺขุ สติสมฺโพชฺฌงฺคํ ภาเวติ วิเวกนิสฺสิตํ ฯเปฯ. อุเปกฺขาสมฺโพชฺฌางฺคํ ภาเวติ วิเวกนิสฺสิตํ วิราคนิสฺสิตํ นิโรธ\-นิสฺสิตํ โวสฺสคฺค\-ปริณามิํ. เอวํ โข ภิกฺขเว ภิกฺขุ โยนิโสมนสิการ\-สมฺปนฺโน สตฺต โพชฺฌงฺเค ภาเวติ สตฺต โพชฺฌงฺเค พหุลีกโรตีติ.  ตติยํ.}

\csromanheader{2}{4. ปฐม\-คิลาน\-สุตฺต}
\proseitem{195}{เอกํ สมยํ ภควา ราชคเห วิหรติ เวฬุวเน กลนฺทก\-นิวาเป. เตน โข ปน สมเยน อายสฺมา มหากสฺสโป ปิปฺปลิคุหายํ\csromansharedfootnote{d951e99acb15b902}{ปิปฺผลิคุหายํ (สี)} วิหรติ อาพาธิโก ทุกฺขิโต พาฬฺหคิลาโน. อถ โข ภควา สายนฺหสมยํ ปฏิสลฺลานา วุฏฺฐิโต เยนายสฺมา มหากสฺสโป เตนุปสงฺกมิ, อุปสงฺกมิตฺวา ปญฺญตฺเต อาสเน นิสีทิ, นิสชฺช โข ภควา อายสฺมนฺตํ มหากสฺสปํ เอตทโวจ\csromandash{}}

\prose{“กจฺจิ เต กสฺสป ขมนียํ, กจฺจิ ยาปนียํ, กจฺจิ ทุกฺขา เวทนา ปฏิกฺกมนฺติ โน อภิกฺกมนฺติ, ปฏิกฺกโมสานํ ปญฺญายติ โน อภิกฺกโม”ติ. น เม ภนฺเต ขมนียํ, น ยาปนียํ, พฬฺหา เม ทุกฺขา เวทนา อภิกฺกมนฺติ โน ปฏิกฺกมนฺติ, อภิกฺกโมสานํ ปญฺญายติ โน ปฏิกฺกโมติ.}

\prose{สตฺติเม กสฺสป โพชฺฌงฺคา มยา สมฺมทกฺขาตา ภาวิตา พหุลีกตา อภิญฺญาย สมฺโพธาย นิพฺพานาย สํวตฺตนฺติ. กตเม สตฺต, สติสมฺโพชฺฌงฺโค โข กสฺสป มยา สมฺมทกฺขาโต ภาวิโต พหุลีกโต อภิญฺญาย สมฺโพธาย นิพฺพานาย สํวตฺตติ ฯเปฯ. อุเปกฺขา\-สมฺโพชฺฌงฺโค โข กสฺสป มยา สมฺมทกฺขาโต ภาวิโต พหุลีกโต อภิญฺญาย \csromanfolio{72}สมฺโพธาย นิพฺพานาย สํวตฺตติ. อิเม โข กสฺสป สตฺต โพชฺฌงฺคา มยา สมฺมทกฺขาตา ภาวิตา พหุลีกตา อภิญฺญาย สมฺโพธาย นิพฺพานาย สํวตฺตนฺตีติ. ตคฺฆ ภควา โพชฺฌงฺคา, ตคฺฆ สุคต โพชฺฌงฺคาติ.}

\prose{อิทมโวจ ภควา. อตฺตมโน อายสฺมา มหากสฺสโป ภควโต ภาสิตํ อภินนฺทิ. วุฏฺฐหิ จายสฺมา มหากสฺสโป ตมฺหา อาพาธา, ตถาปหีโน จายสฺมโต มหากสฺสปสฺส โส อาพาโธ อโหสีติ.  จตุตฺถํ.}

\csromanheader{2}{5. ทุติย\-คิลาน\-สุตฺต}
\proseitem{196}{เอกํ สมยํ ภควา ราชคเห วิหรติ เวฬุวเน กลนฺทก\-นิวาเป. เตน โข ปน สมเยน อายสฺมา มหาโมคฺคลฺลาโน คิชฺฌกูเฏ ปพฺพเต วิหรติ อาพาธิโก ทุกฺขิโต พาฬฺหคิลาโน. อถ โข ภควา สายนฺหสมยํ ปฏิสลฺลานา วุฏฺฐิโต เยนายสฺมา มหาโมคฺคลฺลาโน เตนุปสงฺกมิ, อุปสงฺกมิตฺวา ปญฺญตฺเต อาสเน นิสีทิ, นิสชฺช โข ภควา อายสฺมนฺตํ มหา\-โมคฺคลฺลานํ เอตทโวจ\csromandash{}}

\prose{“กจฺจิ เต โมคฺคลฺลาน ขมนียํ, กจฺจิ ยาปนียํ, กจฺจิ ทุกฺขา เวทนา ปฏิกฺกมนฺติ โน อภิกฺกมนฺติ, ปฏิกฺกโมสานํ ปญฺญายติ โน อภิกฺกโม”ติ. น เม ภนฺเต ขมนียํ, น ยาปนียํ, พาฬฺหา เม ทุกฺขา เวทนา อภิกฺกมนฺติ โน ปฏิกฺกมนฺติ, อภิกฺกโมสานํ ปญฺญายติ โน ปฏิกฺกโมติ.}

\prose{สตฺติเม โมคฺคลฺลาน โพชฺฌงฺคา มยา สมฺมทกฺขาตา ภาวิตา พหุลีกตา อภิญฺญาย สมฺโพธาย นิพฺพานาย สํวตฺตนฺติ. กตเม สตฺต, สติสมฺโพชฺฌงฺโค โข โมคฺคลฺลาน มยา สมฺมทกฺขาโต ภาวิโต พหุลีกโต อภิญฺญาย สมฺโพธาย นิพฺพานาย สํวตฺตติ ฯเปฯ. อุเปกฺขา\-สมฺโพชฺฌงฺโค โข โมคฺคลฺลาน มยา สมฺมทกฺขาโต ภาวิโต พหุลีกโต อภิญฺญาย สมฺโพธาย นิพฺพานาย สํวตฺตติ. อิเม โข โมคฺคลฺลาน สตฺต โพชฺฌงฺคา มยา สมฺมทกฺขาตา ภาวิตา พหุลีกตา อภิญฺญาย สมฺโพธาย นิพฺพานาย สํวตฺตนฺตีติ. ตคฺฆ ภควา โพชฺฌงฺคา, ตคฺฆ สุคต โพชฺฌงฺคาติ.}

\prose{\csromanfolio{73}อิทมโวจ ภควา. อตฺตมโน อายสฺมา มหาโมคฺคลฺลาโน ภควโต ภาสิตํ อภินนฺทิ. วุฏฺฐหิ จายสฺมา มหาโมคฺคลฺลาโน ตมฺหา อาพาธา, ตถาปหีโน จายสฺมโต มหา\-โมคฺคลฺลานสฺส โส อาพาโธ อโหสีติ.  ปญฺจมํ.}

\csromanheader{2}{6. ตติย\-คิลาน\-สุตฺต}
\proseitem{197}{เอกํ สมยํ ภควา ราชคเห วิหรติ เวฬุวเน กลนฺทก\-นิวาเป. เตน โข ปน สมเยน ภควา อาพาธิโก โหติ ทุกฺขิโต พาฬฺหคิลาโน. อถ โข อายสฺมา มหาจุนฺโท เยน ภควา เตนุปสงฺกมิ, อุปสงฺกมิตฺวา ภควนฺตํ อภิวาเทตฺวา เอกมนฺตํ นิสีทิ, เอกมนฺตํ นิสินฺนํ โข อายสฺมนฺตํ มหาจุนฺทํ ภควา เอตทโวจ “ปฏิภนฺตุ ตํ จุนฺท โพชฺฌงฺคา”ติ.}

\prose{สตฺติเม ภนฺเต โพชฺฌงฺคา ภควตา สมฺมทกฺขาตา ภาวิตา พหุลีกตา อภิญฺญาย สมฺโพธาย นิพฺพานาย สํวตฺตนฺติ. กตเม สตฺต, สติสมฺโพชฺฌงฺโค โข ภนฺเต ภควตา สมฺมทกฺขาโต ภาวิโต พหุลีกโต อภิญฺญาย สมฺโพธาย นิพฺพานาย สํวตฺตติ ฯเปฯ. อุเปกฺขา\-สมฺโพชฺฌงฺโค โข ภนฺเต ภควตา สมฺมทกฺขาโต ภาวิโต พหุลีกโต อภิญฺญาย สมฺโพธาย นิพฺพานาย สํวตฺตติ. อิเม โข ภนฺเต สตฺต โพชฺฌงฺคา ภควตา สมฺมทกฺขาตา ภาวิตา พหุลีกตา อภิญฺญาย สมฺโพธาย นิพฺพานาย สํวตฺตนฺตีติ. ตคฺฆ จุนฺท โพชฺฌงฺคา, ตคฺฆ จุนฺท โพชฺฌงฺคาติ.}

\prose{อิทมโวจายสฺมา จุนฺโท. สมนุญฺโญ สตฺถา อโหสิ. วุฏฺฐหิ จ ภควา ตมฺหา อาพาธา, ตถาปหีโน จ ภควโต โส อาพาโธ อโหสีติ.  ฉฏฺฐํ.}

\csromanheader{2}{7. \textbf{ปารงฺคมสุตฺต}}
\proseitem{198}{สตฺติเม ภิกฺขเว โพชฺฌงฺคา ภาวิตา พหุลีกตา อปารา ปารํ คมนาย สํวตฺตนฺติ. กตเม สตฺต, สติสมฺโพชฺฌงฺโค ฯเปฯ อุเปกฺขา\-สมฺโพชฺฌงฺโค. อิเม โข ภิกฺขเว สตฺต โพชฺฌงฺคา ภาวิตา พหุลีกตา อปารา ปารํ คมนาย สํวตฺตนฺตีติ.}

\setlength{\csromangathapagewidth}{0pt}
\setlength{\csromangathawakwidth}{0pt}
\csromangathameasuregroup{อปฺปกา เต มนุสฺเสสุ, \\ อถายํ อิตรา ปชา, \\ เย จ โข สมฺมทกฺขาเต, \\ เต ชนา ปารเมสฺสนฺติ, \\ กณฺหํ ธมฺมํ วิปฺปหาย, \\ โอกา อโนกมาคมฺม, \\ ตตฺราภิรติมิจฺเฉยฺย, \\ ปริโยทเปยฺย อตฺตานํ, \\ เยสํ สมฺโพธิยงฺเคสุ, \\ อาทานปฺปฏินิสฺสคฺเค, \\ ขีณาสวา ชุติมนฺโต,}{\csromangathabat{อปฺปกา เต มนุสฺเสสุ,}{เย ชนา ปารคามิโน.} \\ \csromangathabat{อถายํ อิตรา ปชา,}{ตีรเมวานุธาวติ.} \\ \csromangathabat{เย จ โข สมฺมทกฺขาเต,}{ธมฺเม ธมฺมานุวตฺติโน.} \\ \csromangathabat{เต ชนา ปารเมสฺสนฺติ,}{มจฺจุเธยฺยํ สุทุตฺตรํ.} \\ \csromangathabat{กณฺหํ ธมฺมํ วิปฺปหาย,}{สุกฺกํ ภาเวถ ปณฺฑิโต.} \\ \csromangathabat{โอกา อโนกมาคมฺม,}{วิเวเก ยตฺถ ทูรมํ.} \\ \csromangathabat{ตตฺราภิรติมิจฺเฉยฺย,}{หิตฺวา กาเม อกิญฺจโน.} \\ \csromangathabat{ปริโยทเปยฺย อตฺตานํ,}{จิตฺตเกฺลเสหิ ปณฺฑิโต.} \\ \csromangathabat{เยสํ สมฺโพธิยงฺเคสุ,}{สมฺมา จิตฺตํ สุภาวิตํ.} \\ \csromangathabat{อาทานปฺปฏินิสฺสคฺเค,}{อนุปาทาย เย รตา.} \\ \csromangathabat{ขีณาสวา ชุติมนฺโต,}{เต โลเก ปรินิพฺพุตาติ.} \\ สตฺตมํ.}
\csromangathasetleft{อปฺปกา เต มนุสฺเสสุ, \\ อถายํ อิตรา ปชา, \\ เย จ โข สมฺมทกฺขาเต, \\ เต ชนา ปารเมสฺสนฺติ, \\ กณฺหํ ธมฺมํ วิปฺปหาย, \\ โอกา อโนกมาคมฺม, \\ ตตฺราภิรติมิจฺเฉยฺย, \\ ปริโยทเปยฺย อตฺตานํ, \\ เยสํ สมฺโพธิยงฺเคสุ, \\ อาทานปฺปฏินิสฺสคฺเค, \\ ขีณาสวา ชุติมนฺโต,}
\csromangathagroup{\csromangathastanza{\csromanfolio{74}\csromangathabat{อปฺปกา เต มนุสฺเสสุ,}{เย ชนา ปารคามิโน.} \\ \csromangathabat{อถายํ อิตรา ปชา,}{ตีรเมวานุธาวติ.}}\csromangathastanza{\csromangathabat{เย จ โข สมฺมทกฺขาเต,}{ธมฺเม ธมฺมานุวตฺติโน.} \\ \csromangathabat{เต ชนา ปารเมสฺสนฺติ,}{มจฺจุเธยฺยํ สุทุตฺตรํ.}}\csromangathastanza{\csromangathabat{กณฺหํ ธมฺมํ วิปฺปหาย,}{สุกฺกํ ภาเวถ ปณฺฑิโต.} \\ \csromangathabat{โอกา อโนกมาคมฺม,}{วิเวเก ยตฺถ ทูรมํ.}}\csromangathastanza{\csromangathabat{ตตฺราภิรติมิจฺเฉยฺย,}{หิตฺวา กาเม อกิญฺจโน.} \\ \csromangathabat{ปริโยทเปยฺย อตฺตานํ,}{จิตฺตเกฺลเสหิ ปณฺฑิโต.}}\csromangathastanza{\csromangathabat{เยสํ สมฺโพธิยงฺเคสุ,}{สมฺมา จิตฺตํ สุภาวิตํ.} \\ \csromangathabat{อาทานปฺปฏินิสฺสคฺเค,}{อนุปาทาย เย รตา.}}\csromangathastanza{\csromangathabat{ขีณาสวา ชุติมนฺโต,}{เต โลเก ปรินิพฺพุตาติ.} \\ สตฺตมํ.}}

\csromanheader{2}{8. \textbf{วิรทฺธสุตฺต}}
\proseitem{199}{เยสํ เกสญฺจิ ภิกฺขเว สตฺต โพชฺฌงฺคา วิรทฺธา, วิรทฺโธ เตสํ อริโย มคฺโค สมฺมา ทุกฺข\-กฺขย\-คามี. เยสํ เกสญฺจิ ภิกฺขเว สตฺต โพชฺฌงฺคา อารทฺธา, อารทฺโธ เตสํ อริโย มคฺโค สมฺมา ทุกฺข\-กฺขย\-คามี. กตเม สตฺต, สติสมฺโพชฺฌงฺโค ฯเปฯ อุเปกฺขา\-สมฺโพชฺฌงฺโค. เยสํ เกสญฺจิ ภิกฺขเว อิเม สตฺต โพชฺฌงฺคา วิรทฺธา, วิรทฺโธ เตสํ อริโย มคฺโค สมฺมา ทุกฺข\-กฺขย\-คามี. เยสํ เกสญฺจิ ภิกฺขเว อิเม สตฺต โพชฺฌงฺคา อารทฺธา, อารทฺโธ เตสํ อริโย มคฺโค สมฺมา ทุกฺข\-กฺขย\-คามีติ.  อฏฺฐมํ.}

\csromanheader{2}{9. \textbf{อริยสุตฺต}}
\proseitem{200}{สตฺติเม ภิกฺขเว โพชฺฌงฺคา ภาวิตา พหุลีกตา อริยา นิยฺยานิกา นียนฺติ ตกฺกรสฺส สมฺมา ทุกฺขกฺขยาย. กตเม สตฺต, สติสมฺโพชฺฌงฺโค ฯเปฯ อุเปกฺขา\-สมฺโพชฺฌงฺโค. อิเม โข ภิกฺขเว สตฺต โพชฺฌงฺคา ภาวิตา พหุลีกตา อริยา นิยฺยานิกา นียนฺติ ตกฺกรสฺส สมฺมา ทุกฺขกฺขยายาติ.  นวมํ.}

\csromanheader{2}{10. \textbf{นิพฺพิทาสุตฺต}}
\proseitemwithclosermidruled{201}{\csromanfolio{75}สตฺติเม ภิกฺขเว โพชฺฌงฺคา ภาวิตา พหุลีกตา เอกนฺต\-นิพฺพิทาย วิราคาย นิโรธาย อุปสมาย อภิญฺญาย สมฺโพธาย นิพฺพานาย สํวตฺตนฺติ. กตเม สตฺต, สติสมฺโพชฺฌงฺโค ฯเปฯ อุเปกฺขา\-สมฺโพชฺฌงฺโค. อิเม โข ภิกฺขเว สตฺต โพชฺฌงฺคา ภาวิตา พหุลีกตา เอกนฺต\-นิพฺพิทาย วิราคาย นิโรธาย อุปสมาย อภิญฺญาย สมฺโพธาย นิพฺพานาย สํวตฺตนฺตีติ.  ทสมํ.}{คิลานวคฺโค ทุติโย.}
\csromancenter{\textbf{ตสฺสุทฺทานํ}}
\setlength{\csromangathapagewidth}{0pt}
\setlength{\csromangathawakwidth}{0pt}
\csromangathameasuregroup{ปาณา สูริยูปมา เทฺว, \\ ปารํคามี วิรทฺโธ จ,}{\csromangathabat{ปาณา สูริยูปมา เทฺว,}{คิลานา อปเร ตโย.} \\ \csromangathabat{ปารํคามี วิรทฺโธ จ,}{อริโย นิพฺพิทาย จาติ.}}
\csromangathasetleft{ปาณา สูริยูปมา เทฺว, \\ ปารํคามี วิรทฺโธ จ,}
\csromangathagroup{\csromangathastanza{\csromangathabat{ปาณา สูริยูปมา เทฺว,}{คิลานา อปเร ตโย.} \\ \csromangathabat{ปารํคามี วิรทฺโธ จ,}{อริโย นิพฺพิทาย จาติ.}}}

\chapterhead{3. \textbf{อุทายิวคฺค}}
\csromanheader{2}{1. \textbf{โพธายสุตฺต}}
\proseitem{202}{อถ โข อญฺญตโร ภิกฺขุ เยน ภควา เตนุปสงฺกมิ ฯเปฯ เอกมนฺตํ นิสินฺโน โข โส ภิกฺขุ ภควนฺตํ เอตทโวจ\csromandash{}}

\prose{“โพชฺฌงฺคา โพชฺฌงฺคา”ติ ภนฺเต วุจฺจนฺติ, กิตฺตาวตา นุ โข ภนฺเต “โพชฺฌงฺคา”ติ วุจฺจนฺตีติ. โพธาย สํวตฺตนฺตีติ โข ภิกฺขุ, ตสฺมา “โพชฺฌงฺคา”ติ วุจฺจนฺติ, อิธ ภิกฺขุ สติสมฺโพชฺฌงฺคํ ภาเวติ วิเวกนิสฺสิตํ ฯเปฯ. อุเปกฺขา\-สมฺโพชฺฌงฺคํ ภาเวติ วิเวกนิสฺสิตํ วิราคนิสฺสิตํ นิโรธ\-นิสฺสิตํ โวสฺสคฺค\-ปริณามิํ. โพธาย สํวตฺตนฺตีติ โข ภิกฺขุ, ตสฺมา “โพชฺฌงฺคา”ติ วุจฺจนฺตีติ.  ปฐมํ.}

\csromanheader{2}{2. โพชฺฌงฺค\-เทสนา\-สุตฺต}
\proseitem{203}{สตฺต โว ภิกฺขเว โพชฺฌงฺเค เทเสสฺสามิ ตํ สุณาถ. กตเม จ ภิกฺขเว สตฺต โพชฺฌงฺคา, สติสมฺโพชฺฌงฺโค ฯเปฯ อุเปกฺขา\-สมฺโพชฺฌงฺโค. อิเม โข ภิกฺขเว สตฺต โพชฺฌงฺคาติ.  ทุติยํ.}

\csromanheader{2}{3. \textbf{ฐานิยสุตฺต}}
\proseitem{204}{\csromanfolio{76}กามราค\-ฏฺฐานิยานํ\csromansharedfootnote{ae0c90eac8770e7f}{กามราคฏฺฐานียานํ (สี)} ภิกฺขเว ธมฺมานํ มนสิการ\-พหุลีการา อนุปฺปนฺโน เจว กามจฺฉนฺโท อุปฺปชฺชติ, อุปฺปนฺโน จ กามจฺฉนฺโท ภิยฺโยภาวาย เวปุลฺลาย สํวตฺตติ. พฺยาปาท\-ฏฺฐานิยานํ ภิกฺขเว ธมฺมานํ มนสิการ\-พหุลีการา อนุปฺปนฺโน เจว พฺยาปาโท อุปฺปชฺชติ, อุปฺปนฺโน จ พฺยาปาโท ภิยฺโยภาวาย เวปุลฺลาย สํวตฺตติ. ถินมิทฺธ\-ฏฺฐานิยานํ ภิกฺขเว ธมฺมานํ มนสิการ\-พหุลีการา อนุปฺปนฺนญฺเจว ถินมิทฺธํ อุปฺปชฺชติ, อุปฺปนฺนญฺจ ถินมิทฺธํ ภิยฺโยภาวาย เวปุลฺลาย สํวตฺตติ. อุทฺธจฺจกุกฺกุจฺจ\-ฏฺฐานิยานํ ภิกฺขเว ธมฺมานํ มนสิการ\-พหุลีการา อนุปฺปนฺนญฺเจว อุทฺธจฺจ\-กุกฺกุจฺจํ อุปฺปชฺชติ, อุปฺปนฺนญฺจ อุทฺธจฺจ\-กุกฺกุจฺจํ ภิยฺโยภาวาย เวปุลฺลาย สํวตฺตติ. วิจิกิจฺฉา\-ฏฺฐานิยานํ ภิกฺขเว ธมฺมานํ มนสิการ\-พหุลีการา อนุปฺปนฺนา เจว วิจิกิจฺฉา อุปฺปชฺชติ, อุปฺปนฺนา จ วิจิกิจฺฉา ภิยฺโยภาวาย เวปุลฺลาย สํวตฺตติ.}

\prose{สติสมฺโพชฺฌงฺค\-ฏฺฐานิยานํ ภิกฺขเว ธมฺมานํ มนสิการ\-พหุลีการา อนุปฺปนฺโน เจว สติสมฺโพชฺฌงฺโค อุปฺปชฺชติ, อุปฺปนฺโน จ สติสมฺโพชฺฌงฺโค ภาวนา\-ปาริปูริํ คจฺฉติ ฯเปฯ. อุเปกฺขาสมฺโพชฺฌงฺค\-ฏฺฐานิยานํ ภิกฺขเว ธมฺมานํ มนสิการ\-พหุลีการา อนุปฺปนฺโน เจว อุเปกฺขา\-สมฺโพชฺฌงฺโค อุปฺปชฺชติ, อุปฺปนฺโน จ อุเปกฺขา\-สมฺโพชฺฌงฺโค ภาวนา\-ปาริปูริํ คจฺฉตีติ.  ตติยํ.}

\csromanheader{2}{4. อโยนิโสมนสิการ\-สุตฺต}
\proseitem{205}{อโยนิโส ภิกฺขเว มนสิกโรโต อนุปฺปนฺโน เจว กามจฺฉนฺโท อุปฺปชฺชติ, อุปฺปนฺโน จ กามจฺฉนฺโท ภิยฺโยภาวาย เวปุลฺลาย สํวตฺตติ. อนุปฺปนฺโน เจว พฺยาปาโท อุปฺปชฺชติ, อุปฺปนฺโน จ พฺยาปาโท ภิยฺโยภาวาย เวปุลฺลาย สํวตฺตติ. อนุปฺปนฺนญฺเจว ถินมิทฺธํ อุปฺปชฺชติ, อุปฺปนฺนญฺจ ถินมิทฺธํ ภิยฺโยภาวาย เวปุลฺลาย สํวตฺตติ. อนุปฺปนฺนญฺเจว อุทฺธจฺจ\-กุกฺกุจฺจํ อุปฺปชฺชติ, อุปฺปนฺนญฺจ อุทฺธจฺจ\-กุกฺกุจฺจํ ภิยฺโยภาวาย เวปุลฺลาย สํวตฺตติ. อนุปฺปนฺนา เจว วิจิกิจฺฉา อุปฺปชฺชติ, อุปฺปนฺนา จ วิจิกิจฺฉา ภิยฺโยภาวาย เวปุลฺลาย สํวตฺตติ. อนุปฺปนฺโน เจว สติสมฺโพชฺฌงฺโค นุปฺปชฺชติ, อุปฺปนฺโน จ สติสมฺโพชฺฌงฺโค นิรุชฺฌติ ฯเปฯ. อนุปฺปนฺโน เจว อุเปกฺขา\-สมฺโพชฺฌงฺโค นุปฺปชฺชติ, อุปฺปนฺโน จ อุเปกฺขา\-สมฺโพชฺฌงฺโค นิรุชฺฌติ.}

\prose{\csromanfolio{77}โยนิโส จ โข ภิกฺขเว มนสิกโรโต อนุปฺปนฺโน เจว กามจฺฉนฺโท นุปฺปชฺชติ, อุปฺปนฺโน จ กามจฺฉนฺโท ปหียติ. อนุปฺปนฺโน เจว พฺยาปาโท นุปฺปชฺชติ, อุปฺปนฺโน จ พฺยาปาโท ปหียติ. อนุปฺปนฺนญฺเจว ถินมิทฺธํ นุปฺปชฺชติ, อุปฺปนฺนญฺจ ถินมิทฺธํ ปหียติ. อนุปฺปนฺนญฺเจว อุทฺธจฺจ\-กุกฺกุจฺจํ นุปฺปชฺชติ, อุปฺปนฺนญฺจ อุทฺธจฺจ\-กุกฺกุจฺจํ ปหียติ. อนุปฺปนฺนา เจว วิจิกิจฺฉา นุปฺปชฺชติ, อุปฺปนฺนา จ วิจิกิจฺฉา ปหียติ.}

\prose{อนุปฺปนฺโน เจว สติสมฺโพชฺฌงฺโค อุปฺปชฺชติ, อุปฺปนฺโน จ สติสมฺโพชฺฌงฺโค ภาวนา\-ปาริปูริํ คจฺฉติ ฯเปฯ. อนุปฺปนฺโน เจว อุเปกฺขา\-สมฺโพชฺฌงฺโค อุปฺปชฺชติ, อุปฺปนฺโน จ อุเปกฺขา\-สมฺโพชฺฌงฺโค ภาวนา\-ปาริปูริํ คจฺฉตีติ.  จตุตฺถํ.}

\csromanheader{2}{5. อปริหานิย\-สุตฺต}
\proseitem{206}{สตฺต โว ภิกฺขเว อปริหานิเย ธมฺเม เทเสสฺสามิ ตํ สุณาถ. กตเม จ ภิกฺขเว สตฺต อปริหานิยา ธมฺมา, ยทิทํ สตฺต โพชฺฌงฺคา. กตเม สตฺต, สติสมฺโพชฺฌงฺโค ฯเปฯ อุเปกฺขา\-สมฺโพชฺฌงฺโค. อิเม โข ภิกฺขเว สตฺต อปริหานิยา ธมฺมาติ.  ปญฺจมํ.}

\csromanheader{2}{6. ตณฺหกฺขย\-สุตฺต}
\proseitem{207}{โย ภิกฺขเว มคฺโค. ยา ปฏิปทา ตณฺหกฺขยาย สํวตฺตติ, ตํ มคฺคํ ตํ ปฏิปทํ ภาเวถ. กตโม จ ภิกฺขเว มคฺโค กตมา จ ปฏิปทา ตณฺหกฺขยาย สํวตฺตติ, ยทิทํ สตฺต โพชฺฌงฺคา. กตเม สตฺต, สติสมฺโพชฺฌงฺโค ฯเปฯ อุเปกฺขา\-สมฺโพชฺฌงฺโคติ. เอวํ วุตฺเต อายสฺมา อุทายี ภควนฺตํ เอตทโวจ “กถํ ภาวิตา นุ โข ภนฺเต สตฺต โพชฺฌงฺคา กถํ พหุลีกตา ตณฺหกฺขยาย สํวตฺตนฺตี”ติ.}

\prose{อิธ อุทายิ ภิกฺขุ สติสมฺโพชฺฌงฺคํ ภาเวติ วิเวกนิสฺสิตํ วิราคนิสฺสิตํ นิโรธ\-นิสฺสิตํ โวสฺสคฺค\-ปริณามิํ วิปุลํ มหคฺคตํ อปฺปมาณํ อพฺยาปชฺชํ\csromansharedfootnote{8057e252681ed807}{อพฺยาปชฺฌํ (สี, สฺยา, อิ)}. ตสฺส สติสมฺโพชฺฌงฺคํ ภาวยโต วิเวกนิสฺสิตํ วิราคนิสฺสิตํ นิโรธ\-นิสฺสิตํ โวสฺสคฺค\-ปริณามิํ วิปุลํ มหคฺคตํ อปฺปมาณํ อพฺยาปชฺชํ ตณฺหา ปหียติ, ตณฺหาย ปหานา กมฺมํ ปหียติ, กมฺมสฺส ปหานา ทุกฺขํ ปหียติ ฯเปฯ. อุเปกฺขา\-สมฺโพชฺฌงฺคํ ภาเวติ วิเวกนิสฺสิตํ \csromanfolio{78}วิราคนิสฺสิตํ นิโรธ\-นิสฺสิตํ โวสฺสคฺค\-ปริณามิํ วิปุลํ มหคฺคตํ อปฺปมาณํ อพฺยาปชฺชํ. ตสฺส อุเปกฺขา\-สมฺโพชฺฌงฺคํ ภาวยโต วิเวกนิสฺสิตํ วิราคนิสฺสิตํ นิโรธ\-นิสฺสิตํ โวสฺสคฺค\-ปริณามิํ วิปุลํ มหคฺคตํ อปฺปมาณํ อพฺยาปชฺชํ ตณฺหา ปหียติ, ตณฺหาย ปหานา กมฺมํ ปหียติ, กมฺมสฺส ปหานา ทุกฺขํ ปหียติ. อิติ โข อุทายิ ตณฺหกฺขยา กมฺมกฺขโย, กมฺมกฺขยา ทุกฺขกฺขโยติ.  ฉฏฺฐํ.}

\csromanheader{2}{7. ตณฺหานิโรธ\-สุตฺต}
\proseitem{208}{โย ภิกฺขเว มคฺโค ยา ปฏิปทา ตณฺหานิโรธาย สํวตฺตติ, ตํ มคฺคํ ตํ ปฏิปทํ ภาเวถ. กตโม จ ภิกฺขเว มคฺโค กตมา จ ปฏิปทา ตณฺหานิโรธาย สํวตฺตติ, ยทิทํ สตฺต โพชฺฌงฺคา. กตเม สตฺต, สติสมฺโพชฺฌงฺโค ฯเปฯ อุเปกฺขา\-สมฺโพชฺฌงฺโค. กถํ ภาวิตา จ ภิกฺขเว สตฺต โพชฺฌงฺคา กถํ พหุลีกตา ตณฺหานิโรธาย สํวตฺตนฺติ.}

\prose{อิธ ภิกฺขเว ภิกฺขุ สติสมฺโพชฺฌงฺคํ ภาเวติ วิเวกนิสฺสิตํ วิราคนิสฺสิตํ นิโรธ\-นิสฺสิตํ โวสฺสคฺค\-ปริณามิํ ฯเปฯ. อุเปกฺขา\-สมฺโพชฺฌงฺคํ ภาเวติ วิเวกนิสฺสิตํ วิราคนิสฺสิตํ นิโรธ\-นิสฺสิตํ โวสฺสคฺค\-ปริณามิํ. เอวํ ภาวิตา โข ภิกฺขเว สตฺต โพชฺฌงฺคา เอวํ พหุลีกตา ตณฺหานิโรธาย สํวตฺตนฺตีติ.  สตฺตมํ.}

\csromanheader{2}{8. นิพฺเพธภาคิย\-สุตฺต}
\proseitem{209}{นิพฺเพธ\-ภาคิยํ โว ภิกฺขเว มคฺคํ เทเสสฺสามิ ตํ สุณาถ. กตโม จ ภิกฺขเว นิพฺเพธ\-ภาคิโย มคฺโค, ยทิทํ สตฺต โพชฺฌงฺคา. กตเม สตฺต, สติสมฺโพชฺฌงฺโค ฯเปฯ อุเปกฺขา\-สมฺโพชฺฌงฺโคติ. เอวํ วุตฺเต อายสฺมา อุทายี ภควนฺตํ เอตทโวจ “กถํ ภาวิตา นุ โข ภนฺเต สตฺต โพชฺฌงฺคา กถํ พหุลีกตา นิพฺเพธาย สํวตฺตนฺตี”ติ.}

\prose{อิธ อุทายิ ภิกฺขุ สติสมฺโพชฺฌงฺคํ ภาเวติ วิเวกนิสฺสิตํ วิราคนิสฺสิตํ นิโรธ\-นิสฺสิตํ โวสฺสคฺค\-ปริณามิํ วิปุลํ มหคฺคตํ อปฺปมาณํ อพฺยาปชฺชํ. โส สติสมฺโพชฺฌงฺคํ ภาวิเตน จิตฺเตน อนิพฺพิทฺธปุพฺพํ อปฺปทาลิตปุพฺพํ โลภกฺขนฺธํ นิพฺพิชฺฌติ ปทาเลติ, อนิพฺพิทฺธปุพฺพํ อปฺปทาลิตปุพฺพํ โทสกฺขนฺธํ นิพฺพิชฺฌติ ปทาเลติ, อนิพฺพิทฺธปุพฺพํ อปฺปทาลิตปุพฺพํ โมหกฺขนฺธํ นิพฺพิชฺฌติ ปทาเลติ ฯเปฯ. อุเปกฺขา\-สมฺโพชฺฌงฺคํ ภาเวติ วิเวกนิสฺสิตํ วิราคนิสฺสิตํ}

\prose{\csromanfolio{79}นิโรธ\-นิสฺสิตํ โวสฺสคฺค\-ปริณามิํ วิปุลํ มหคฺคตํ อปฺปมาณํ อพฺยาปชฺชํ. โส อุเปกฺขา\-สมฺโพชฺฌงฺคํ ภาวิเตน จิตฺเตน อนิพฺพิทฺธปุพฺพํ อปฺปทาลิตปุพฺพํ โลภกฺขนฺธํ นิพฺพิชฺฌติ ปทาเลติ, อนิพฺพิทฺธปุพฺพํ อปฺปทาลิตปุพฺพํ โทสกฺขนฺธํ นิพฺพิชฺฌติ ปทาเลติ, อนิพฺพิทฺธปุพฺพํ อปฺปทาลิตปุพฺพํ โมหกฺขนฺธํ นิพฺพิชฺฌติ ปทาเลติ. เอวํ ภาวิตา โข อุทายิ สตฺต โพชฺฌงฺคา เอวํ พหุลีกตา นิพฺเพธาย สํวตฺตนฺตีติ.  อฏฺฐมํ.}

\csromanheader{2}{9. \textbf{เอกธมฺมสุตฺต}}
\proseitem{210}{นาหํ ภิกฺขเว อญฺญํ เอกธมฺมมฺปิ สมนุปสฺสามิ. โย เอวํ ภาวิโต พหุลีกโต สํโยชนียานํ ธมฺมานํ ปหานาย สํวตฺตติ, ยถยิทํ ภิกฺขเว สตฺต โพชฺฌงฺคา. กตเม สตฺต, สติสมฺโพชฺฌงฺโค ฯเปฯ อุเปกฺขา\-สมฺโพชฺฌงฺโค. กถํ ภาวิตา จ ภิกฺขเว สตฺต โพชฺฌงฺคา กถํ พหุลีกตา สํโยชนียานํ ธมฺมานํ ปหานาย สํวตฺตนฺติ.}

\prose{อิธ ภิกฺขเว ภิกฺขุ สติสมฺโพชฺฌงฺคํ ภาเวติ วิเวกนิสฺสิตํ ฯเปฯ. อุเปกฺขา\-สมฺโพชฺฌงฺคํ ภาเวติ วิเวกนิสฺสิตํ วิราคนิสฺสิตํ นิโรธ\-นิสฺสิตํ โวสฺสคฺค\-ปริณามิํ. เอวํ ภาวิตา โข ภิกฺขเว สตฺต โพชฺฌงฺคา เอวํ พหุลีกตา สํโยชนียานํ ธมฺมานํ ปหานาย สํวตฺตนฺติ.}

\prose{กตเม จ ภิกฺขเว สํโยชนียา ธมฺมา, จกฺขุ ภิกฺขเว สํโยชนีโย ธมฺโม, เอตฺเถเต อุปฺปชฺชนฺติ สํโยชน\-วินิพนฺธา อชฺโฌสานา ฯเปฯ ชิวฺหา สํโยชนียา ธมฺมา, เอตฺเถเต อุปฺปชฺชนฺติ สํโยชน\-วินิพนฺธา อชฺโฌสานา ฯเปฯ มโน สํโยชนีโย ธมฺโม, เอตฺเถเต อุปฺปชฺชนฺติ สํโยชน\-วินิพนฺธา อชฺโฌสานา. อิเม วุจฺจนฺติ ภิกฺขเว สํโยชนียา ธมฺมาติ.  นวมํ.}

\csromanheader{2}{10. \textbf{อุทายิสุตฺต}}
\proseitem{211}{เอกํ สมยํ ภควา สุมฺเภสุ วิหรติ เสตกํ นาม สุมฺภานํ นิคโม. อถ โข อายสฺมา อุทายี เยน ภควา เตนุปสงฺกมิ ฯเปฯ เอกมนฺตํ นิสินฺโน โข อายสฺมา อุทายี ภควนฺตํ เอตทโวจ\csromandash{}}

\prose{อจฺฉริยํ ภนฺเต, อพฺภุตํ ภนฺเต, ยาว พหุกตญฺจ เม ภนฺเต ภควติ เปมญฺจ คารโว จ หิรี จ โอตฺตปฺปญฺจ. อหํ หิ ภนฺเต ปุพฺเพ อคาริกภูโต \csromanfolio{80}สมาโน อพหุกโต อโหสิํ ธมฺเมน\csromansharedfootnote{89e1e123c2d51c29}{ธมฺเม (?)}, อพหุกโต สํเฆน. โส ขฺวาหํ ภควติ เปมญฺจ คารวญฺจ หิริํ จ โอตฺตปฺปญฺจ สมฺปสฺสมาโน อคารสฺมา อนคาริยํ ปพฺพชิโต, ตสฺส เม ภควา ธมฺมํ เทเสสิ “อิติ รูปํ อิติ รูปสฺส สมุทโย อิติ รูปสฺส อตฺถงฺคโม. อิติ เวทนา ฯเปฯ. อิติ สญฺญา. อิติ สงฺขารา. อิติ วิญฺญาณํ. อิติ วิญฺญาณสฺส สมุทโย อิติ วิญฺญาณสฺส อตฺถงฺคโม”ติ.}

\prose{โส ขฺวาหํ ภนฺเต สุญฺญาคารคโต อิเมสํ ปญฺจุ\-ปาทาน\-กฺขนฺธานํ อุกฺกุชฺชา\-วกุชฺชํ สมฺปริ\-วตฺเต\-นฺโต “อิทํ ทุกฺขนฺ”ติ ยถาภูตํ อพฺภญฺญาสิํ, “อยํ ทุกฺขสมุทโย”ติ ยถาภูตํ อพฺภญฺญาสิํ, “อยํ ทุกฺขนิโรโธ”ติ ยถาภูตํ อพฺภญฺญาสิํ, “อยํ ทุกฺขนิโรธ\-คามินี ปฏิปทา”ติ ยถาภูตํ อพฺภญฺญาสิํ. ธมฺโม จ เม ภนฺเต อภิสมิโต, มคฺโค จ เม ปฏิลทฺโธ. โย เม ภาวิโต พหุลีกโต, ตถา ตถา วิหรนฺตํ ตถตฺตาย อุปเนสฺสติ, ยถาหํ “ขีณา ชาติ, วุสิตํ พฺรหฺมจริยํ, กตํ กรณียํ, นาปรํ อิตฺถตฺตายา”ติ ปชานิสฺสามิ.}

\prose{สติสมฺโพชฺฌงฺโค เม ภนฺเต ปฏิลทฺโธ, โย เม ภาวิโต พหุลีกโต, ตถา ตถา วิหรนฺตํ ตถตฺตาย อุปเนสฺสติ, ยถาหํ “ขีณา ชาติ, วุสิตํ พฺรหฺมจริยํ, กตํ กรณียํ, นาปรํ อิตฺถตฺตายา”ติ ปชานิสฺสามิ ฯเปฯ. อุเปกฺขา\-สมฺโพชฺฌงฺโค เม ภนฺเต ปฏิลทฺโธ. โย เม ภาวิโต พหุลีกโต, ตถา ตถา วิหรนฺตํ ตถตฺตาย อุปเนสฺสติ, ยถาหํ “ขีณา ชาติ, วุสิตํ พฺรหฺมจริยํ, กตํ กรณียํ, นาปรํ อิตฺถตฺตายา”ติ ปชานิสฺสามิ. อยํ โข เม ภนฺเต มคฺโค ปฏิลทฺโธ. โย เม ภาวิโต พหุลีกโต, ตถา ตถา วิหรนฺตํ ตถตฺตาย อุปเนสฺสติ, ยถาหํ “ขีณา ชาติ, วุสิตํ พฺรหฺมจริยํ, กตํ กรณียํ, นาปรํ อิตฺถตฺตายา”ติ ปชานิสฺสามีติ.}

\prosewithclosermidruled{สาธุ สาธุ อุทายิ, เอโส หิ เต อุทายิ มคฺโค ปฏิลทฺโธ. โย เต ภาวิโต พหุลีกโต, ตถา ตถา วิหรนฺตํ ตถตฺตาย อุปเนสฺสติ, ยถา ตฺวํ “ขีณา ชาติ, วุสิตํ พฺรหฺมจริยํ, กตํ กรณียํ, นาปรํ อิตฺถตฺตายา”ติ ปชานิสฺสสีติ.  ทสมํ.}{อุทายิวคฺโค ตติโย.}
\csromancenter{\textbf{ตสฺสุทฺทานํ}}
\setlength{\csromangathapagewidth}{0pt}
\setlength{\csromangathawakwidth}{0pt}
\csromangathameasuregroup{โพธาย เทสนา ฐานา, \\ ขโย นิโรโธ นิพฺเพโธ,}{\csromangathabat{โพธาย เทสนา ฐานา,}{อโยนิโส จาปริหานี.} \\ \csromangathabat{ขโย นิโรโธ นิพฺเพโธ,}{เอกธมฺโม อุทายินาติ.}}
\csromangathasetleft{โพธาย เทสนา ฐานา, \\ ขโย นิโรโธ นิพฺเพโธ,}
\csromangathagroup{\csromangathastanza{\csromanfolio{81}\csromangathabat{โพธาย เทสนา ฐานา,}{อโยนิโส จาปริหานี.} \\ \csromangathabat{ขโย นิโรโธ นิพฺเพโธ,}{เอกธมฺโม อุทายินาติ.}}}

\chapterhead{4. \textbf{นีวรณวคฺค}}
\csromanheader{2}{1. ปฐม\-กุสล\-สุตฺต}
\proseitem{212}{เย เกจิ ภิกฺขเว ธมฺมา กุสลา กุสลภาคิยา กุสลปกฺขิกา, สพฺเพ เต อปฺปมาทมูลกา อปฺปมาท\-สโมสรณา. อปฺปมาโท เตสํ ธมฺมานํ อคฺคมกฺขายติ. อปฺปมตฺตสฺเสตํ ภิกฺขเว ภิกฺขุโน ปาฏิกงฺขํ สตฺต โพชฺฌงฺเค ภาเวสฺสติ สตฺต โพชฺฌงฺเค พหุลีกริสฺสติ.}

\prose{กถญฺจ ภิกฺขเว ภิกฺขุ อปฺปมตฺโต สตฺต โพชฺฌงฺเค ภาเวติ, สตฺตโพชฺฌงฺเค พหุลีกโรติ, อิธ ภิกฺขเว ภิกฺขุ สติสมฺโพชฺฌงฺคํ ภาเวติ วิเวกนิสฺสิตํ ฯเปฯ. อุเปกฺขา\-สมฺโพชฺฌงฺคํ ภาเวติ วิเวกนิสฺสิตํ วิราคนิสฺสิตํ นิโรธ\-นิสฺสิตํ โวสฺสคฺค\-ปริณามิํ. เอวํ โข ภิกฺขเว ภิกฺขุ อปฺปมตฺโต สตฺต โพชฺฌงฺเค ภาเวติ สตฺต โพชฺฌงฺเค พหุลีกโรตีติ.  ปฐมํ.}

\csromanheader{2}{2. ทุติย\-กุสล\-สุตฺต}
\proseitem{213}{เย เกจิ ภิกฺขเว ธมฺมา กุสลา กุสลภาคิยา กุสลปกฺขิกา, สพฺเพ เต โยนิโสมนสิการ\-มูลกา โยนิโสมนสิการ\-สโมสรณา. โยนิโส\-มนสิกาโร เตสํ ธมฺมานํ อคฺคมกฺขายติ. โยนิโสมนสิการสมฺปนฺนสฺเสตํ ภิกฺขเว ภิกฺขุโน ปาฏิกงฺขํ สตฺต โพชฺฌงฺเค ภาเวสฺสติ สตฺต โพชฺฌงฺเค พหุลีกริสฺสติ.}

\prose{กถญฺจ ภิกฺขเว ภิกฺขุ โยนิโสมนสิการ\-สมฺปนฺโน สตฺต โพชฺฌงฺเค ภาเวติ สตฺต โพชฺฌงฺเค พหุลีกโรติ, อิธ ภิกฺขเว ภิกฺขุ สติสมฺโพชฺฌงฺคํ ภาเวติ วิเวกนิสฺสิตํ ฯเปฯ. อุเปกฺขา\-สมฺโพชฺฌงฺคํ ภาเวติ วิเวกนิสฺสิตํ วิราคนิสฺสิตํ นิโรธ\-นิสฺสิตํ โวสฺสคฺค\-ปริณามิํ. เอวํ โข \csromanfolio{82}ภิกฺขเว ภิกฺขุ โยนิโสมนสิการ\-สมฺปนฺโน สตฺต โพชฺฌงฺเค ภาเวติ สตฺต โพชฺฌงฺเค พหุลีกโรตีติ.  ทุติยํ.}

\csromanheader{2}{3. อุปกฺกิเลส\-สุตฺต}
\proseitem{214}{ปญฺจิเม ภิกฺขเว ชาตรูปสฺส อุปกฺกิเลสา, เยหิ อุปกฺกิเลเสหิ อุปกฺกิลิฏฺฐํ ชาตรูปํ น เจว มุทุ โหติ น จ กมฺมนิยํ น จ ปภสฺสรํ ปภงฺคุ จ, น จ สมฺมา อุเปติ กมฺมาย. กตเม ปญฺจ, อโย ภิกฺขเว ชาตรูปสฺส อุปกฺกิเลโส, เยน อุปกฺกิเลเสน อุปกฺกิลิฏฺฐํ ชาตรูปํ น เจว มุทุ โหติ น จ กมฺมนิยํ น จ ปภสฺสรํ ปภงฺคุ จ, น จ สมฺมา อุเปติ กมฺมาย. โลหํ ภิกฺขเว ชาตรูปสฺส อุปกฺกิเลโส, เยน อุปกฺกิเลเสน อุปกฺกิลิฏฺฐํ ชาตรูปํ ฯเปฯ. ติปุ ภิกฺขเว ชาตรูปสฺส อุปกฺกิเลโส ฯเปฯ. สีสํ ภิกฺขเว ชาตรูปสฺส อุปกฺกิเลโส ฯเปฯ. สชฺฌุ ภิกฺขเว ชาตรูปสฺส อุปกฺกิเลโส, เยน อุปกฺกิเลเสน อุปกฺกิลิฏฺฐํ ชาตรูปํ น เจว มุทุ โหติ น จ กมฺมนิยํ น จ ปภสฺสรํ ปภงฺคุ จ, น จ สมฺมา อุเปติ กมฺมาย. อิเม โข ภิกฺขเว ปญฺจ ชาตรูปสฺส อุปกฺกิเลสา, เยหิ อุปกฺกิเลเสหิ อุปกฺกิลิฏฺฐํ ชาตรูปํ น เจว มุทุ โหติ น จ กมฺมนิยํ น จ ปภสฺสรํ ปภงฺคุ จ, น จ สมฺมา อุเปติ กมฺมาย.}

\prose{เอวเมว โข ภิกฺขเว ปญฺจิเม จิตฺตสฺส อุปกฺกิเลสา, เยหิ อุปกฺกิเลเสหิ อุปกฺกิลิฏฺฐํ จิตฺตํ น เจว มุทุ โหติ น จ กมฺมนิยํ น จ ปภสฺสรํ ปภงฺคุ จ, น จ สมฺมา สมาธิยติ อาสวานํ ขยาย. กตเม ปญฺจ, กามจฺฉนฺโท ภิกฺขเว จิตฺตสฺส อุปกฺกิเลโส, เยน อุปกฺกิเลเสน อุปกฺกิลิฏฺฐํ จิตฺตํ น เจว มุทุ โหติ น จ กมฺมนิยํ น จ ปภสฺสรํ ปภงฺคุ จ, น จ สมฺมา สมาธิยติ อาสวานํ ขยาย ฯเปฯ. อิเม โข ภิกฺขเว ปญฺจ จิตฺตสฺส อุปกฺกิเลสา, เยหิ อุปกฺกิเลเสหิ อุปกฺกิลิฏฺฐํ จิตฺตํ น เจว มุทุ โหติ น จ กมฺมนิยํ น จ ปภสฺสรํ ปภงฺคุ จ, น จ สมฺมา สมาธิยติ อาสวานํ ขยายาติ.  ตติยํ.}

\csromanheader{2}{4. อนุปกฺกิเลส\-สุตฺต}
\proseitem{215}{สตฺติเม ภิกฺขเว โพชฺฌงฺคา อนาวรณา อนีวรณา เจตโส อนุปกฺกิเลสา ภาวิตา พหุลีกตา วิชฺชา\-วิมุตฺติ\-ผล\-สจฺฉิกิริยาย สํวตฺตนฺติ. กตเม สตฺต, สติสมฺโพชฺฌงฺโค ภิกฺขเว อนาวรโณ}

\prose{\csromanfolio{83}อนีวรโณ เจตโส อนุปกฺกิเลโส ภาวิโต พหุลีกโต วิชฺชา\-วิมุตฺติ\-ผล\-สจฺฉิกิริยาย สํวตฺตติ ฯเปฯ. อุเปกฺขา\-สมฺโพชฺฌงฺโค ภิกฺขเว อนาวรโณ อนีวรโณ เจตโส อนุปกฺกิเลโส ภาวิโต พหุลีกโต วิชฺชา\-วิมุตฺติ\-ผล\-สจฺฉิกิริยาย สํวตฺตติ. อิเม โข ภิกฺขเว สตฺต โพชฺฌงฺคา อนาวรณา อนีวรณา เจตโส อนุปกฺกิเลสา ภาวิตา พหุลีกตา วิชฺชา\-วิมุตฺติ\-ผล\-สจฺฉิกิริยาย สํวตฺตนฺตีติ.  จตุตฺถํ.}

\csromanheader{2}{5. อโยนิโสมนสิการ\-สุตฺต}
\proseitem{216}{อโยนิโส ภิกฺขเว มนสิกโรโต อนุปฺปนฺโน เจว กามจฺฉนฺโท อุปฺปชฺชติ, อุปฺปนฺโน จ กามจฺฉนฺโท ภิยฺโยภาวาย เวปุลฺลาย สํวตฺตติ. อนุปฺปนฺโน เจว พฺยาปาโท อุปฺปชฺชติ, อุปฺปนฺโน จ พฺยาปาโท ภิยฺโยภาวาย เวปุลฺลาย สํวตฺตติ. อนุปฺปนฺนญฺเจว ถินมิทฺธํ อุปฺปชฺชติ, อุปฺปนฺนญฺจ ถินมิทฺธํ ภิยฺโยภาวาย เวปุลฺลาย สํวตฺตติ. อนุปฺปนฺนญฺเจว อุทฺธจฺจ\-กุกฺกุจฺจํ อุปฺปชฺชติ, อุปฺปนฺนญฺจ อุทฺธจฺจ\-กุกฺกุจฺจํ ภิยฺโยภาวาย เวปุลฺลาย สํวตฺตติ. อนุปฺปนฺนา เจว วิจิกิจฺฉา อุปฺปชฺชติ, อุปฺปนฺนา จ วิจิกิจฺฉา ภิยฺโยภาวาย เวปุลฺลาย สํวตฺตตีติ.  ปญฺจมํ.}

\csromanheader{2}{6. โยนิโสมนสิการ\-สุตฺต}
\proseitem{217}{โยนิโส จ โข ภิกฺขเว มนสิกโรโต อนุปฺปนฺโน เจว สติสมฺโพชฺฌงฺโค อุปฺปชฺชติ, อุปฺปนฺโน จ สติสมฺโพชฺฌงฺโค ภาวนา\-ปาริปูริํ คจฺฉติ ฯเปฯ อนุปฺปนฺโน เจว อุเปกฺขา\-สมฺโพชฺฌงฺโค อุปฺปชฺชติ, อุปฺปนฺโน จ อุเปกฺขา\-สมฺโพชฺฌงฺโค ภาวนา\-ปาริปูริํ คจฺฉตีติ.  ฉฏฺฐํ.}

\csromanheader{2}{7. \textbf{พุทฺธิสุตฺต}}
\proseitem{218}{สตฺติเม ภิกฺขเว โพชฺฌงฺคา ภาวิตา พหุลีกตา พุทฺธิยา อปริหานาย สํวตฺตนฺติ. กตเม สตฺต, สติสมฺโพชฺฌงฺโค ฯเปฯ อุเปกฺขา\-สมฺโพชฺฌงฺโค. อิเม โข ภิกฺขเว สตฺต โพชฺฌงฺคา ภาวิตา พหุลีกตา พุทฺธิยา อปริหานาย สํวตฺตนฺตีติ.  สตฺตมํ.}

\csromanheader{1}{8. อาวรณ\-นีวรณ\-สุตฺต}
\proseitem{219}{ปญฺจิเม ภิกฺขเว อาวรณา นีวรณา เจตโส อุปกฺกิเลสา ปญฺญาย ทุพฺพลีกรณา. กตเม ปญฺจ, กามจฺฉนฺโท ภิกฺขเว อาวรโณ \csromanfolio{84}นีวรโณ เจตโส อุปกฺกิเลโส ปญฺญาย ทุพฺพลีกรโณ, พฺยาปาโท ภิกฺขเว อาวรโณ นีวรโณ เจตโส อุปกฺกิเลโส ปญฺญาย ทุพฺพลีกรโณ, ถินมิทฺธํ ภิกฺขเว อาวรณํ นีวรณํ เจตโส อุปกฺกิเลสํ ปญฺญาย ทุพฺพลีกรณํ, อุทฺธจฺจ\-กุกฺกุจฺจํ ภิกฺขเว อาวรณํ นีวรณํ เจตโส อุปกฺกิเลสํ ปญฺญาย ทุพฺพลีกรณํ, วิจิกิจฺฉา ภิกฺขเว อาวรณา นีวรณา เจตโส อุปกฺกิเลสา ปญฺญาย ทุพฺพลีกรณา. อิเม โข ภิกฺขเว ปญฺจ อาวรณา นีวรณา เจตโส อุปกฺกิเลสา ปญฺญาย ทุพฺพลีกรณา.}

\prose{สตฺติเม ภิกฺขเว โพชฺฌงฺคา อนาวรณา อนีวรณา เจตโส อนุปกฺกิเลสา ภาวิตา พหุลีกตา วิชฺชา\-วิมุตฺติ\-ผล\-สจฺฉิกิริยาย สํวตฺตนฺติ. กตเม สตฺต, สติสมฺโพชฺฌงฺโค ภิกฺขเว อนาวรโณ อนีวรโณ เจตโส อนุปกฺกิเลโส ภาวิโต พหุลีกโต วิชฺชา\-วิมุตฺติ\-ผล\-สจฺฉิกิริยาย สํวตฺตติ ฯเปฯ. อุเปกฺขา\-สมฺโพชฺฌงฺโค ภิกฺขเว อนาวรโณ อนีวรโณ เจตโส อนุปกฺกิเลโส ภาวิโต พหุลีกโต วิชฺชาวิมุตฺติผลสจฺฉิ กิริยาย สํวตฺตติ. อิเม โข ภิกฺขเว สตฺต โพชฺฌงฺคา อนาวรณา อนีวรณา เจตโส อนุปกฺกิเลสา ภาวิตา พหุลีกตา วิชฺชา\-วิมุตฺติ\-ผล\-สจฺฉิกิริยาย สํวตฺตนฺตีติ.}

\prose{ยสฺมิํ ภิกฺขเว สมเย อริยสาวโก อฏฺฐิํ กตฺวา มนสิ กตฺวา สพฺพํ เจตโส สมนฺนาหริตฺวา โอหิตโสโต ธมฺมํ สุณาติ, อิมสฺส ปญฺจ นีวรณา ตสฺมิํ สมเย น โหนฺติ, สตฺต โพชฺฌงฺคา ตสฺมิํ สมเย ภาวนา\-ปาริปูริํ คจฺฉนฺติ.}

\prose{กตเม ปญฺจ นีวรณา ตสฺมิํ สมเย น โหนฺติ, กามจฺฉนฺท\-นีวรณํ ตสฺมิํ สมเย น โหติ, พฺยาปาท\-นีวรณํ ตสฺมิํ สมเย น โหติ, ถินมิทฺธ\-นีวรณํ ตสฺมิํ สมเย น โหติ, อุทฺธจฺจ\-กุกฺกุจฺจ\-นีวรณํ ตสฺมิํ สมเย น โหติ, วิจิกิจฺฉา\-นีวรณํ ตสฺมิํ สมเย น โหติ. อิมสฺส ปญฺจ นีวรณา ตสฺมิํ สมเย น โหนฺติ.}

\prose{กตเม สตฺต โพชฺฌงฺคา ตสฺมิํ สมเย ภาวนา\-ปาริปูริํ คจฺฉนฺติ, สติสมฺโพชฺฌงฺโค ตสฺมิํ สมเย ภาวนา\-ปาริปูริํ คจฺฉติ ฯเปฯ อุเปกฺขา\-สมฺโพชฺฌงฺโค ตสฺมิํ สมเย ภาวนา\-ปาริปูริํ คจฺฉติ. อิเม สตฺต โพชฺฌงฺคา ตสฺมิํ สมเย ภาวนา\-ปาริปูริํ คจฺฉนฺติ. ยสฺมิํ ภิกฺขเว สมเย อริยสาวโก อฏฺฐิํ กตฺวา มนสิ กตฺวา สพฺพํ เจตโส สมนฺนาหริตฺวา}

\prose{\csromanfolio{85}โอหิตโสโต ธมฺมํ สุณาติ, อิมสฺส ปญฺจ นีวรณา ตสฺมิํ สมเย น โหนฺติ. อิเม สตฺต โพชฺฌงฺคา ตสฺมิํ สมเย ภาวนา\-ปาริปูริํ คจฺฉนฺตีติ.  อฏฺฐมํ.}

\csromanheader{1}{9. \textbf{รุกฺขสุตฺต}}
\proseitem{220}{สนฺติ ภิกฺขเว มหารุกฺขา อณุพีชา มหากายา รุกฺขานํ อชฺฌารุหา, เยหิ รุกฺขา อชฺฌารูฬฺหา โอภคฺค\-วิภคฺคา วิปติตา เสนฺติ. กตเม จ เต ภิกฺขเว มหารุกฺขา อณุพีชา มหากายา รุกฺขานํ อชฺฌารุหา, เยหิ รุกฺขา อชฺฌารูฬฺหา โอภคฺค\-วิภคฺคา วิปติตา เสนฺติ\csromansharedfootnote{bca9087ba498fbef}{เสนฺติ. เสยฺยถิทํ (กตฺถจิ)}, อสฺสตฺโถ นิโคฺรโธ ปิลกฺโข อุทุมฺพโร กจฺฉโก กปิตฺถโน. อิเม โข เต ภิกฺขเว มหารุกฺขา อณุพีชา มหากายา รุกฺขานํ อชฺฌารุหา, เยหิ รุกฺขา อชฺฌารูฬฺหา โอภคฺค\-วิภคฺคา วิปติตา เสนฺติ. เอวเมว โข ภิกฺขเว อิเธกจฺโจ กุลปุตฺโต ยาทิสเก กาเม โอหาย อคารสฺมา อนคาริยํ ปพฺพชิโต โหติ, โส ตาทิสเกหิ กาเมหิ ตโต วา ปาปิฏฺฐตเรหิ โอภคฺค\-วิภคฺโค วิปติโต เสติ.}

\prose{ปญฺจิเม ภิกฺขเว อาวรณา นีวรณา เจตโส อชฺฌารุหา ปญฺญาย ทุพฺพลีกรณา. กตเม ปญฺจ, กามจฺฉนฺโท ภิกฺขเว อาวรโณ นีวรโณ เจตโส อชฺฌารุโห ปญฺญาย ทุพฺพลีกรโณ, พฺยาปาโท ภิกฺขเว อาวรโณ นีวรโณ เจตโส อชฺฌารุโห ปญฺญาย ทุพฺพลีกรโณ, ถินมิทฺธํ ภิกฺขเว อาวรณํ นีวรณํ เจตโส อชฺฌารุหํ ปญฺญาย ทุพฺพลีกรณํ, อุทฺธจฺจ\-กุกฺกุจฺจํ ภิกฺขเว อาวรณํ นีวรณํ เจตโส อชฺฌารุหํ ปญฺญาย ทุพฺพลีกรณํ, วิจิกิจฺฉา ภิกฺขเว อาวรณา นีวรณา เจตโส อชฺฌารุหา ปญฺญาย ทุพฺพลีกรณา. อิเม โข ภิกฺขเว ปญฺจ อาวรณา นีวรณา เจตโส อชฺฌารุหา ปญฺญาย ทุพฺพลีกรณา.}

\prose{สตฺติเม ภิกฺขเว โพชฺฌงฺคา อนาวรณา อนีวรณา เจตโส อนชฺฌารุหา ภาวิตา พหุลีกตา วิชฺชา\-วิมุตฺติ\-ผล\-สจฺฉิกิริยาย สํวตฺตนฺติ. กตเม สตฺต, สติสมฺโพชฺฌงฺโค ภิกฺขเว อนาวรโณ อนีวรโณ เจตโส อนชฺฌารุโห ภาวิโต พหุลีกโต วิชฺชา\-วิมุตฺติ\-ผล\-สจฺฉิกิริยาย สํวตฺตติ ฯเปฯ. อุเปกฺขา\-สมฺโพชฺฌงฺโค ภิกฺขเว \csromanfolio{86}อนาวรโณ อนีวรโณ เจตโส อนชฺฌารุโห ภาวิโต พหุลีกโต วิชฺชาวิมุตฺติผลสจฺฉิ กิริยาย สํวตฺตติ. อิเม โข ภิกฺขเว สตฺต โพชฺฌงฺคา อนาวรณา อนีวรณา เจตโส อนชฺฌารุหา ภาวิตา พหุลีกตา วิชฺชา\-วิมุตฺติ\-ผล\-สจฺฉิกิริยาย สํวตฺตนฺตีติ.  นวมํ.}

\csromanheader{1}{10. \textbf{นีวรณสุตฺต}}
\proseitem{221}{ปญฺจิเม ภิกฺขเว นีวรณา อนฺธกรณา อจกฺขุกรณา อญฺญาณกรณา ปญฺญานิโรธิกา วิฆาต\-ปกฺขิยา อนิพฺพานสํวตฺตนิกา. กตเม ปญฺจ, กามจฺฉนฺท\-นีวรณํ ภิกฺขเว อนฺธกรณํ อจกฺขุกรณํ อญฺญาณกรณํ ปญฺญา\-นิโรธิกํ วิฆาต\-ปกฺขิยํ อนิพฺพานสํวตฺตนิกํ, พฺยาปาท\-นีวรณํ ภิกฺขเว ฯเปฯ ถินมิทฺธ\-นีวรณํ ภิกฺขเว ฯเปฯ อุทฺธจฺจ\-กุกฺกุจฺจ\-นีวรณํ ภิกฺขเว ฯเปฯ วิจิกิจฺฉา\-นีวรณํ ภิกฺขเว อนฺธกรณํ อจกฺขุกรณํ อญฺญาณกรณํ ปญฺญา\-นิโรธิกํ วิฆาต\-ปกฺขิยํ อนิพฺพานสํวตฺตนิกํ. อิเม โข ภิกฺขเว ปญฺจ นีวรณา อนฺธกรณา อจกฺขุกรณา อญฺญาณกรณา ปญฺญานิโรธิกา วิฆาต\-ปกฺขิยา อนิพฺพานสํวตฺตนิกา.}

\prosewithclosermidruled{สตฺติเม ภิกฺขเว โพชฺฌงฺคา จกฺขุกรณา ญาณกรณา ปญฺญาพุทฺธิยา อวิฆาต\-ปกฺขิยา นิพฺพาน\-สํวตฺตนิกา. กตเม สตฺต, สติสมฺโพชฺฌงฺโค ภิกฺขเว จกฺขุกรโณ ญาณกรโณ ปญฺญาพุทฺธิโย อวิฆาต\-ปกฺขิโย นิพฺพาน\-สํวตฺตนิโก ฯเปฯ. อุเปกฺขา\-สมฺโพชฺฌงฺโค ภิกฺขเว จกฺขุกรโณ ญาณกรโณ ปญฺญาพุทฺธิโย อวิฆาต\-ปกฺขิโย นิพฺพาน\-สํวตฺตนิโก. อิเม โข ภิกฺขเว สตฺต โพชฺฌงฺคา จกฺขุกรณา ญาณกรณา ปญฺญาพุทฺธิยา อวิฆาต\-ปกฺขิยา นิพฺพาน\-สํวตฺตนิกาติ.  ทสมํ.}{นีวรณวคฺโค จตุตฺโถ.}
\csromancenter{\textbf{ตสฺสุทฺทานํ}}
\setlength{\csromangathapagewidth}{0pt}
\setlength{\csromangathawakwidth}{0pt}
\csromangathameasuregroup{เทฺว กุสลา กิเลสา จ, \\ อาวรณา นีวรณา รุกฺขํ,}{\csromangathabat{เทฺว กุสลา กิเลสา จ,}{เทฺว โยนิโส จ พุทฺธิ จ.} \\ \csromangathabat{อาวรณา นีวรณา รุกฺขํ,}{นีวรณญฺจ เต ทสาติ.}}
\csromangathasetleft{เทฺว กุสลา กิเลสา จ, \\ อาวรณา นีวรณา รุกฺขํ,}
\csromangathagroup{\csromangathastanza{\csromangathabat{เทฺว กุสลา กิเลสา จ,}{เทฺว โยนิโส จ พุทฺธิ จ.} \\ \csromangathabat{อาวรณา นีวรณา รุกฺขํ,}{นีวรณญฺจ เต ทสาติ.}}}

\chapterheadpageread{5. จกฺกวตฺติ\-วคฺค}
\csromanmatikaanchor{mtk.8}
\csromantocmarkref{section}{1. วิธาสุตฺต}{mtk.8}
\bookmark[dest={mtk.8},level=1]{1. วิธาสุตฺต}
\csromanheader{1}{1. \textbf{วิธาสุตฺต}}
\csromanmatikaanchor{mtk.9}
\csromantocmarkref{subsection}{2-8. จกฺกวตฺติสุตฺตาทีนิ}{mtk.9}
\bookmark[dest={mtk.9},level=2]{2-8. จกฺกวตฺติสุตฺตาทีนิ}
\proseitem{222}{\csromanfolio{87}สาวตฺถิ\-นิทานํ. เย หิ เกจิ ภิกฺขเว อตีตมทฺธานํ สมณา วา พฺราหฺมณา วา ติสฺโส วิธา ปชหิํสุ, สพฺเพ เต สตฺตนฺนํ โพชฺฌงฺคานํ ภาวิตตฺตา พหุลีกตตฺตา. เย หิ เกจิ ภิกฺขเว อนาคตมทฺธานํ สมณา วา พฺราหฺมณา วา ติสฺโส วิธา ปชหิสฺสนฺติ, สพฺเพ เต สตฺตนฺนํ โพชฺฌงฺคานํ ภาวิตตฺตา พหุลีกตตฺตา. เย หิ เกจิ ภิกฺขเว เอตรหิ สมณา วา พฺราหฺมณา วา ติสฺโส วิธา ปชหนฺติ, สพฺเพ เต สตฺตนฺนํ โพชฺฌงฺคานํ ภาวิตตฺตา พหุลีกตตฺตา. กตเมสํ สตฺตนฺนํ โพชฺฌงฺคานํ, สติสมฺโพชฺฌงฺคสฺส ฯเปฯ อุเปกฺขา\-สมฺโพชฺฌงฺคสฺส. เย หิ เกจิ ภิกฺขเว อตีตมทฺธานํ สมณา วา พฺราหฺมณา วา ติสฺโส วิธา ปชหิํสุ ฯเปฯ ปชหิสฺสนฺติ ฯเปฯ ปชหนฺติ, สพฺเพ เต อิเมสํเยว สตฺตนฺนํ โพชฺฌงฺคานํ ภาวิตตฺตา พหุลีกตตฺตาติ.  ปฐมํ.}

\csromanheader{2}{2. จกฺกวตฺติ\-สุตฺต}
\proseitem{223}{รญฺโญ ภิกฺขเว จกฺกวตฺติสฺส ปาตุภาวา สตฺตนฺนํ รตนานํ ปาตุภาโว โหติ. กตเมสํ สตฺตนฺนํ, จกฺกรตนสฺส ปาตุภาโว โหติ, หตฺถิ\-รตนสฺส ปาตุภาโว โหติ, อสฺสรตนสฺส ปาตุภาโว โหติ, มณิรตนสฺส ปาตุภาโว โหติ, อิตฺถิรตนสฺส ปาตุภาโว โหติ, คหปติ\-รตนสฺส ปาตุภาโว โหติ, ปริณายก\-รตนสฺส ปาตุภาโว โหติ. รญฺโญ ภิกฺขเว จกฺกวตฺติสฺส ปาตุภาวา อิเมสํ สตฺตนฺนํ รตนานํ ปาตุภาโว โหติ.}

\prose{ตถาคตสฺส ภิกฺขเว ปาตุภาวา อรหโต สมฺมา\-สมฺพุทฺธสฺส สตฺตนฺนํ โพชฺฌงฺค\-รตนานํ ปาตุภาโว โหติ. กตเมสํ สตฺตนฺนํ, สติสมฺโพชฺฌงฺคสฺส รตนสฺส ปาตุภาโว โหติ ฯเปฯ อุเปกฺขา\-สมฺโพชฺฌงฺคสฺส รตนสฺส ปาตุภาโว โหติ. ตถาคตสฺส ภิกฺขเว ปาตุภาวา อรหโต สมฺมา\-สมฺพุทฺธสฺส อิเมสํ สตฺตนฺนํ โพชฺฌงฺค\-รตนานํ ปาตุภาโว โหตีติ.  ทุติยํ.}

\csromanheader{2}{3. \textbf{มารสุตฺต}}
\proseitem{224}{\csromanfolio{88}มารเสน\-ปฺปมทฺทนํ โว ภิกฺขเว มคฺคํ เทเสสฺสามิ, ตํ สุณาถ. กตโม จ ภิกฺขเว มารเสน\-ปฺปมทฺทโน มคฺโค, ยทิทํ สตฺต โพชฺฌงฺคา. กตเม สตฺต, สติสมฺโพชฺฌงฺโค ฯเปฯ อุเปกฺขา\-สมฺโพชฺฌงฺโค. อยํ โข ภิกฺขเว มารเสน\-ปฺปมทฺทโน มคฺโคติ.  ตติยํ.}

\csromanheader{2}{4. \textbf{ทุปฺปญฺญสุตฺต}}
\proseitem{225}{อถ โข อญฺญตโร ภิกฺขุ เยน ภควา เตนุปสงฺกมิ ฯเปฯ เอกมนฺตํ นิสินฺโน โข โส ภิกฺขุ ภควนฺตํ เอตทโวจ “ทุปฺปญฺโญ เอฬมูโค ทุปฺปญฺโญ เอฬมูโค”ติ ภนฺเต วุจฺจติ, กิตฺตาวตา นุ โข ภนฺเต “ทุปฺปญฺโญ เอฬมูโค”ติ วุจฺจตีติ. สตฺตนฺนํ โข ภิกฺขุ โพชฺฌงฺคานํ อภาวิตตฺตา อพหุลีกตตฺตา “ทุปฺปญฺโญ เอฬมูโค”ติ วุจฺจติ. กตเมสํ สตฺตนฺนํ, สติสมฺโพชฺฌงฺคสฺส ฯเปฯ อุเปกฺขา\-สมฺโพชฺฌงฺคสฺส. อิเมสํ โข ภิกฺขุ สตฺตนฺนํ โพชฺฌงฺคานํ อภาวิตตฺตา อพหุลีกตตฺตา “ทุปฺปญฺโญ เอฬมูโค”ติ วุจฺจตีติ.  จตุตฺถํ.}

\csromanheader{2}{5. ปญฺญวนฺต\-สุตฺต}
\proseitem{226}{“ปญฺญวา อเนฬมูโค ปญฺญวา อเนฬมูโค”ติ ภนฺเต วุจฺจติ, กิตฺตาวตา นุ โข ภนฺเต “ปญฺญวา อเนฬมูโค”ติ วุจฺจตีติ. สตฺตนฺนํ โข ภิกฺขุ โพชฺฌงฺคานํ ภาวิตตฺตา พหุลีกตตฺตา “ปญฺญวา อเนฬมูโค”ติ วุจฺจติ. กตเมสํ สตฺตนฺนํ, สติสมฺโพชฺฌงฺคสฺส ฯเปฯ อุเปกฺขา\-สมฺโพชฺฌงฺคสฺส. อิเมสํ โข ภิกฺขุ สตฺตนฺนํ โพชฺฌงฺคานํ ภาวิตตฺตา พหุลีกตตฺตา “ปญฺญวา อเนฬมูโค”ติ วุจฺจตีติ.  ปญฺจมํ.}

\csromanheader{2}{6. \textbf{ทลิทฺทสุตฺต}}
\proseitem{227}{“ทลิทฺโท ทลิทฺโท”ติ ภนฺเต วุจฺจติ, กิตฺตาวตา นุ โข ภนฺเต “ทลิทฺโท”ติ วุจฺจตีติ. สตฺตนฺนํ โข ภิกฺขุ โพชฺฌงฺคานํ อภาวิตตฺตา อพหุลีกตตฺตา “ทลิทฺโท”ติ วุจฺจติ. กตเมสํ สตฺตนฺนํ, สติสมฺโพชฺฌงฺคสฺส ฯเปฯ อุเปกฺขา\-สมฺโพชฺฌงฺคสฺส. อิเมสํ โข ภิกฺขุ สตฺตนฺนํ โพชฺฌงฺคานํ อภาวิตตฺตา อพหุลีกตตฺตา “ทลิทฺโท”ติ วุจฺจตีติ.  ฉฏฺฐํ.}

\csromanheader{2}{7. \textbf{อทลิทฺทสุตฺต}}
\proseitem{228}{\csromanfolio{89}“อทลิทฺโท อทลิทฺโท”ติ ภนฺเต วุจฺจติ, กิตฺตาวตา นุ โข ภนฺเต “อทลิทฺโท”ติ วุจฺจตีติ. สตฺตนฺนํ โข ภิกฺขุ โพชฺฌงฺคานํ ภาวิตตฺตา พหุลีกตตฺตา “อทลิทฺโท”ติ วุจฺจติ. กตเมสํ สตฺตนฺนํ, สติสมฺโพชฺฌงฺคสฺส ฯเปฯ อุเปกฺขา\-สมฺโพชฺฌงฺคสฺส. อิเมสํ โข ภิกฺขุ สตฺตนฺนํ โพชฺฌงฺคานํ ภาวิตตฺตา พหุลีกตตฺตา “อทลิทฺโท”ติ วุจฺจตีติ.  สตฺตมํ.}

\csromanheader{2}{8. \textbf{อาทิจฺจสุตฺต}}
\proseitem{229}{อาทิจฺจสฺส ภิกฺขเว อุทยโต เอตํ ปุพฺพงฺคมํ เอตํ ปุพฺพนิมิตฺตํ, ยทิทํ อรุณุคฺคํ. เอวเมว โข ภิกฺขเว ภิกฺขุโน สตฺตนฺนํ โพชฺฌงฺคานํ อุปฺปาทาย เอตํ ปุพฺพงฺคมํ เอตํ ปุพฺพนิมิตฺตํ, ยทิทํ กลฺยาณมิตฺตตา. กลฺยาณมิตฺตสฺเสตํ ภิกฺขเว ภิกฺขุโน ปาฏิกงฺขํ สตฺต โพชฺฌงฺเค ภาเวสฺสติ สตฺต โพชฺฌงฺเค พหุลีกริสฺสติ. กถญฺจ ภิกฺขเว ภิกฺขุ กลฺยาณมิตฺโต สตฺต โพชฺฌงฺเค ภาเวติ สตฺต โพชฺฌงฺเค พหุลีกโรติ, อิธ ภิกฺขเว ภิกฺขุ สติสมฺโพชฺฌงฺคํ ภาเวติ วิเวกนิสฺสิตํ ฯเปฯ อุเปกฺขา\-สมฺโพชฺฌงฺคํ ภาเวติ วิเวกนิสฺสิตํ วิราคนิสฺสิตํ นิโรธ\-นิสฺสิตํ โวสฺสคฺค\-ปริณามิํ. เอวํ โข ภิกฺขเว ภิกฺขุ กลฺยาณมิตฺโต สตฺต โพชฺฌงฺเค ภาเวติ สตฺต โพชฺฌงฺเค พหุลีกโรตีติ.  อฏฺฐมํ.}

\csromanmatikaanchor{mtk.10}
\csromantocmarkref{subsection}{9. อชฺฌตฺติกงฺคสุตฺต}{mtk.10}
\bookmark[dest={mtk.10},level=2]{9. อชฺฌตฺติกงฺคสุตฺต}
\csromanheader{2}{9. อชฺฌตฺติก\-งฺค\-สุตฺต}
\proseitem{230}{“อชฺฌตฺติกํ ภิกฺขเว องฺคนฺ”ติ กริตฺวา นาญฺญํ เอกงฺคมฺปิ สมนุปสฺสามิ สตฺตนฺนํ โพชฺฌงฺคานํ อุปฺปาทาย, ยถยิทํ ภิกฺขเว โยนิโส\-มนสิกาโร. โยนิโสมนสิการสมฺปนฺนสฺเสตํ ภิกฺขเว ภิกฺขุโน ปาฏิกงฺขํ สตฺต โพชฺฌงฺเค ภาเวสฺสติ สตฺต โพชฺฌงฺเค พหุลีกริสฺสติ. กถญฺจ ภิกฺขเว ภิกฺขุ โยนิโสมนสิการ\-สมฺปนฺโน สตฺต โพชฺฌงฺเค ภาเวติ สตฺต โพชฺฌงฺเค พหุลีกโรติ, อิธ ภิกฺขเว ภิกฺขุ สติสมฺโพชฺฌงฺคํ ภาเวติ วิเวกนิสฺสิตํ ฯเปฯ อุเปกฺขา\-สมฺโพชฺฌงฺคํ ภาเวติ วิเวกนิสฺสิตํ วิราคนิสฺสิตํ นิโรธ\-นิสฺสิตํ โวสฺสคฺค\-ปริณามิํ. เอวํ โข ภิกฺขเว ภิกฺขุ โยนิโสมนสิการ\-สมฺปนฺโน สตฺต โพชฺฌงฺเค ภาเวติ สตฺต โพชฺฌงฺเค พหุลีกโรตีติ.  นวมํ.}

\csromanmatikaanchor{mtk.11}
\csromantocmarkref{subsection}{10. พาหิรงฺคสุตฺต}{mtk.11}
\bookmark[dest={mtk.11},level=2]{10. พาหิรงฺคสุตฺต}
\csromanheader{2}{10. \textbf{พาหิรงฺคสุตฺต}}
\proseitemwithclosermidruled{231}{\csromanfolio{90}“พาหิรํ ภิกฺขเว องฺคนฺ”ติ กริตฺวา นาญฺญํ เอกงฺคมฺปิ สมนุปสฺสามิ สตฺตนฺนํ โพชฺฌงฺคานํ อุปฺปาทาย, ยถยิทํ ภิกฺขเว กลฺยาณมิตฺตตา. กลฺยาณมิตฺตสฺเสตํ ภิกฺขเว ภิกฺขุโน ปาฏิกงฺขํ สตฺต โพชฺฌงฺเค ภาเวสฺสติ สตฺต โพชฺฌงฺเค พหุลีกริสฺสติ. กถญฺจ ภิกฺขเว ภิกฺขุ กลฺยาณมิตฺโต สตฺต โพชฺฌงฺเค ภาเวติ สตฺต โพชฺฌงฺเค พหุลีกโรติ, อิธ ภิกฺขเว ภิกฺขุ สติสมฺโพชฺฌงฺคํ ภาเวติ วิเวกนิสฺสิตํ ฯเปฯ อุเปกฺขา\-สมฺโพชฺฌงฺคํ ภาเวติ วิเวกนิสฺสิตํ วิราคนิสฺสิตํ นิโรธ\-นิสฺสิตํ โวสฺสคฺค\-ปริณามิํ. เอวํ โข ภิกฺขเว ภิกฺขุ กลฺยาณมิตฺโต สตฺต โพชฺฌงฺเค ภาเวติ สตฺต โพชฺฌงฺเค พหุลีกโรตีติ.}{จกฺกวตฺติ\-วคฺโค ปญฺจโม.}
\csromancenter{\textbf{ตสฺสุทฺทานํ}}
\setlength{\csromangathapagewidth}{0pt}
\setlength{\csromangathawakwidth}{0pt}
\csromangathameasuregroup{วิธา จกฺกวตฺติ มาโร, \\ ทลิทฺโท อทลิทฺโท จ,}{\csromangathabat{วิธา จกฺกวตฺติ มาโร,}{ทุปฺปญฺโญ ปญฺญเวน จ.} \\ \csromangathabat{ทลิทฺโท อทลิทฺโท จ,}{อาทิจฺจงฺเคน เต ทสาติ.}}
\csromangathasetleft{วิธา จกฺกวตฺติ มาโร, \\ ทลิทฺโท อทลิทฺโท จ,}
\csromangathagroup{\csromangathastanza{\csromangathabat{วิธา จกฺกวตฺติ มาโร,}{ทุปฺปญฺโญ ปญฺญเวน จ.} \\ \csromangathabat{ทลิทฺโท อทลิทฺโท จ,}{อาทิจฺจงฺเคน เต ทสาติ.}}}

\csromanmatikaanchor{mtk.12}
\csromantocmarknopage{subsection}{6. สากจฺฉวคฺค}{mtk.12}
\bookmark[dest={mtk.12},level=2]{6. สากจฺฉวคฺค}
\csromanheader{2}{6. \textbf{สากจฺฉวคฺค}}
\csromanmatikaanchor{mtk.13}
\csromantocmarkref{paragraph}{1. อาหารสุตฺต}{mtk.13}
\bookmark[dest={mtk.13},level=4]{1. อาหารสุตฺต}
\csromanheader{4}{1. \textbf{อาหารสุตฺต}}
\proseitem{232}{สาวตฺถิ\-นิทานํ. ปญฺจนฺนญฺจ ภิกฺขเว นีวรณานํ สตฺตนฺนญฺจ โพชฺฌงฺคานํ อาหารญฺจ อนาหารญฺจ เทเสสฺสามิ, ตํ สุณาถ. โก จ ภิกฺขเว อาหาโร อนุปฺปนฺนสฺส วา กามจฺฉนฺทสฺส อุปฺปาทาย อุปฺปนฺนสฺส วา กามจฺฉนฺทสฺส ภิยฺโยภาวาย เวปุลฺลาย, อตฺถิ ภิกฺขเว สุภนิมิตฺตํ, ตตฺถ อโยนิโสมนสิการ\-พหุลีกาโร, อยมาหาโร อนุปฺปนฺนสฺส วา กามจฺฉนฺทสฺส อุปฺปาทาย อุปฺปนฺนสฺส วา กามจฺฉนฺทสฺส ภิยฺโยภาวาย เวปุลฺลาย.}

\prose{โก จ ภิกฺขเว อาหาโร อนุปฺปนฺนสฺส วา พฺยาปาทสฺส อุปฺปาทาย อุปฺปนฺนสฺส วา พฺยาปาทสฺส ภิยฺโยภาวาย เวปุลฺลาย, อตฺถิ ภิกฺขเว}

\prose{\csromanfolio{91}ปฏิฆ\-นิมิตฺตํ, ตตฺถ อโยนิโสมนสิการ\-พหุลีกาโร, อยมาหาโร อนุปฺปนฺนสฺส วา พฺยาปาทสฺส อุปฺปาทาย อุปฺปนฺนสฺส วา พฺยาปาทสฺส ภิยฺโยภาวาย เวปุลฺลาย.}

\prose{โก จ ภิกฺขเว อาหาโร อนุปฺปนฺนสฺส วา ถินมิทฺธสฺส อุปฺปาทาย อุปฺปนฺนสฺส วา ถินมิทฺธสฺส ภิยฺโยภาวาย เวปุลฺลาย, อตฺถิ ภิกฺขเว อรติ ตนฺทิ วิชมฺภิตา ภตฺตสมฺมโท เจตโส จ ลีนตฺตํ, ตตฺถ อโยนิโสมนสิการ\-พหุลีกาโร, อยมาหาโร อนุปฺปนฺนสฺส วา ถินมิทฺธสฺส อุปฺปาทาย อุปฺปนฺนสฺส วา ถินมิทฺธสฺส ภิยฺโยภาวาย เวปุลฺลาย.}

\prose{โก จ ภิกฺขเว อาหาโร อนุปฺปนฺนสฺส วา อุทฺธจฺจ\-กุกฺกุจฺจสฺส อุปฺปาทาย อุปฺปนฺนสฺส วา อุทฺธจฺจ\-กุกฺกุจฺจสฺส ภิยฺโยภาวาย เวปุลฺลาย, อตฺถิ ภิกฺขเว เจตโส อวูปสโม, ตตฺถ อโยนิโสมนสิการ\-พหุลีกาโร, อยมาหาโร อนุปฺปนฺนสฺส วา อุทฺธจฺจ\-กุกฺกุจฺจสฺส อุปฺปาทาย อุปฺปนฺนสฺส วา อุทฺธจฺจ\-กุกฺกุจฺจสฺส ภิยฺโยภาวาย เวปุลฺลาย.}

\prose{โก จ ภิกฺขเว อาหาโร อนุปฺปนฺนาย วา วิจิกิจฺฉาย อุปฺปาทาย อุปฺปนฺนาย วา วิจิกิจฺฉาย ภิยฺโยภาวาย เวปุลฺลาย, อตฺถิ ภิกฺขเว วิจิกิจฺฉา\-ฏฺฐานียา ธมฺมา, ตตฺถ อโยนิโสมนสิการ\-พหุลีกาโร, อยมาหาโร อนุปฺปนฺนาย วา วิจิกิจฺฉาย อุปฺปาทาย อุปฺปนฺนาย วา วิจิกิจฺฉาย ภิยฺโยภาวาย เวปุลฺลาย.}

\prose{โก จ ภิกฺขเว อาหาโร อนุปฺปนฺนสฺส วา สติสมฺโพชฺฌงฺคสฺส อุปฺปาทาย อุปฺปนฺนสฺส วา สติสมฺโพชฺฌงฺคสฺส ภาวนาย ปาริปูริยา, อตฺถิ ภิกฺขเว สติสมฺโพชฺฌงฺค\-ฏฺฐานียา ธมฺมา, ตตฺถ โยนิโสมนสิการ\-พหุลีกาโร, อยมาหาโร อนุปฺปนฺนสฺส วา สติสมฺโพชฺฌงฺคสฺส อุปฺปาทาย อุปฺปนฺนสฺส วา สติสมฺโพชฺฌงฺคสฺส ภาวนาย ปาริปูริยา.}

\prose{โก จ ภิกฺขเว อาหาโร อนุปฺปนฺนสฺส วา ธมฺมวิจย\-สมฺโพชฺฌงฺคสฺส อุปฺปาทาย อุปฺปนฺนสฺส วา ธมฺมวิจย\-สมฺโพชฺฌงฺคสฺส ภาวนาย ปาริปูริยา, อตฺถิ ภิกฺขเว กุสลากุสลา ธมฺมา สาวชฺชานวชฺชา ธมฺมา หีนปณีตา ธมฺมา กณฺห\-สุกฺก\-สปฺปฏิภาคา ธมฺมา, ตตฺถ โยนิโสมนสิการ\-พหุลีกาโร, อยมาหาโร อนุปฺปนฺนสฺส วา ธมฺมวิจย\-สมฺโพชฺฌงฺคสฺส อุปฺปาทาย อุปฺปนฺนสฺส วา ธมฺมวิจย\-สมฺโพชฺฌงฺคสฺส ภาวนาย ปาริปูริยา.}

\prose{\csromanfolio{92}โก จ ภิกฺขเว อาหาโร อนุปฺปนฺนสฺส วา วีริย\-สมฺโพชฺฌงฺคสฺส อุปฺปาทาย อุปฺปนฺนสฺส วา วีริย\-สมฺโพชฺฌงฺคสฺส ภาวนาย ปาริปูริยา, อตฺถิ ภิกฺขเว อารมฺภธาตุ นิกฺกมธาตุ ปรกฺกมธาตุ, ตตฺถ โยนิโสมนสิการ\-พหุลีกาโร, อยมาหาโร อนุปฺปนฺนสฺส วา วีริย\-สมฺโพชฺฌงฺคสฺส อุปฺปาทาย อุปฺปนฺนสฺส วา วีริย\-สมฺโพชฺฌงฺคสฺส ภาวนาย ปาริปูริยา.}

\prose{โก จ ภิกฺขเว อาหาโร อนุปฺปนฺนสฺส วา ปีติสมฺโพชฺฌงฺคสฺส อุปฺปาทาย อุปฺปนฺนสฺส วา ปีติสมฺโพชฺฌงฺคสฺส ภาวนาย ปาริปูริยา, อตฺถิ ภิกฺขเว ปีติสมฺโพชฺฌงฺค\-ฏฺฐานียา ธมฺมา, ตตฺถ โยนิโสมนสิการ\-พหุลีกาโร, อยมาหาโร อนุปฺปนฺนสฺส วา ปีติสมฺโพชฺฌงฺคสฺส อุปฺปาทาย อุปฺปนฺนสฺส วา ปีติสมฺโพชฺฌงฺคสฺส ภาวนาย ปาริปูริยา.}

\prose{โก จ ภิกฺขเว อาหาโร อนุปฺปนฺนสฺส วา ปสฺสทฺธิ\-สมฺโพชฺฌงฺคสฺส อุปฺปาทาย อุปฺปนฺนสฺส วา ปสฺสทฺธิ\-สมฺโพชฺฌงฺคสฺส ภาวนาย ปาริปูริยา, อตฺถิ ภิกฺขเว กายปฺปสฺสทฺธิ จิตฺต\-ปฺปสฺสทฺธิ, ตตฺถ โยนิโสมนสิการ\-พหุลีกาโร, อยมาหาโร อนุปฺปนฺนสฺส วา ปสฺสทฺธิ\-สมฺโพชฺฌงฺคสฺส อุปฺปาทาย อุปฺปนฺนสฺส วา ปสฺสทฺธิ\-สมฺโพชฺฌงฺคสฺส ภาวนาย ปาริปูริยา.}

\prose{โก จ ภิกฺขเว อาหาโร อนุปฺปนฺนสฺส วา สมาธิ\-สมฺโพชฺฌงฺคสฺส อุปฺปาทาย อุปฺปนฺนสฺส วา สมาธิ\-สมฺโพชฺฌงฺคสฺส ภาวนาย ปาริปูริยา, อตฺถิ ภิกฺขเว สมถ\-นิมิตฺตํ อพฺยคฺค\-นิมิตฺตํ, ตตฺถ โยนิโสมนสิการ\-พหุลีกาโร, อยมาหาโร อนุปฺปนฺนสฺส วา สมาธิ\-สมฺโพชฺฌงฺคสฺส อุปฺปาทาย อุปฺปนฺนสฺส วา สมาธิ\-สมฺโพชฺฌงฺคสฺส ภาวนาย ปาริปูริยา.}

\prose{โก จ ภิกฺขเว อาหาโร อนุปฺปนฺนสฺส วา อุเปกฺขา\-สมฺโพชฺฌงฺคสฺส อุปฺปาทาย อุปฺปนฺนสฺส วา อุเปกฺขา\-สมฺโพชฺฌงฺคสฺส ภาวนาย ปาริปูริยา, อตฺถิ ภิกฺขเว อุเปกฺขา\-สมฺโพชฺฌงฺค\-ฏฺฐานียา ธมฺมา, ตตฺถ โยนิโสมนสิการ\-พหุลีกาโร, อยมาหาโร อนุปฺปนฺนสฺส วา อุเปกฺขา\-สมฺโพชฺฌงฺคสฺส อุปฺปาทาย อุปฺปนฺนสฺส วา อุเปกฺขา\-สมฺโพชฺฌงฺคสฺส ภาวนาย ปาริปูริยา.}

\prose{โก จ ภิกฺขเว อนาหาโร อนุปฺปนฺนสฺส วา กามจฺฉนฺทสฺส อุปฺปาทาย อุปฺปนฺนสฺส วา กามจฺฉนฺทสฺส ภิยฺโยภาวาย เวปุลฺลาย, อตฺถิ ภิกฺขเว}

\prose{\csromanfolio{93}อสุภนิมิตฺตํ, ตตฺถ โยนิโสมนสิการ\-พหุลีกาโร, อยมนาหาโร อนุปฺปนฺนสฺส วา กามจฺฉนฺทสฺส อุปฺปาทาย อุปฺปนฺนสฺส วา กามจฺฉนฺทสฺส ภิยฺโยภาวาย เวปุลฺลาย.}

\prose{โก จ ภิกฺขเว อนาหาโร อนุปฺปนฺนสฺส วา พฺยาปาทสฺส อุปฺปาทาย อุปฺปนฺนสฺส วา พฺยาปาทสฺส ภิยฺโยภาวาย เวปุลฺลาย, อตฺถิ ภิกฺขเว เมตฺตา\-เจโตวิมุตฺติ, ตตฺถ โยนิโสมนสิการ\-พหุลีกาโร, อยมนาหาโร อนุปฺปนฺนสฺส วา พฺยาปาทสฺส อุปฺปาทาย อุปฺปนฺนสฺส วา พฺยาปาทสฺส ภิยฺโยภาวาย เวปุลฺลาย.}

\prose{โก จ ภิกฺขเว อนาหาโร อนุปฺปนฺนสฺส วา ถินมิทฺธสฺส อุปฺปาทาย อุปฺปนฺนสฺส วา ถินมิทฺธสฺส ภิยฺโยภาวาย เวปุลฺลาย, อตฺถิ ภิกฺขเว อารมฺภธาตุ นิกฺกมธาตุ ปรกฺกมธาตุ, ตตฺถ โยนิโสมนสิการ\-พหุลีกาโร, อยมนาหาโร อนุปฺปนฺนสฺส วา ถินมิทฺธสฺส อุปฺปาทาย อุปฺปนฺนสฺส วา ถินมิทฺธสฺส ภิยฺโยภาวาย เวปุลฺลาย.}

\prose{โก จ ภิกฺขเว อนาหาโร อนุปฺปนฺนสฺส วา อุทฺธจฺจ\-กุกฺกุจฺจสฺส อุปฺปาทาย อุปฺปนฺนสฺส วา อุทฺธจฺจ\-กุกฺกุจฺจสฺส ภิยฺโยภาวาย เวปุลฺลาย, อตฺถิ ภิกฺขเว เจตโส วูปสโม, ตตฺถ โยนิโสมนสิการ\-พหุลีกาโร, อยมนาหาโร อนุปฺปนฺนสฺส วา อุทฺธจฺจ\-กุกฺกุจฺจสฺส อุปฺปาทาย อุปฺปนฺนสฺส วา อุทฺธจฺจ\-กุกฺกุจฺจสฺส ภิยฺโยภาวาย เวปุลฺลาย.}

\prose{โก จ ภิกฺขเว อนาหาโร อนุปฺปนฺนาย วา วิจิกิจฺฉาย อุปฺปาทาย อุปฺปนฺนาย วา วิจิกิจฺฉาย ภิยฺโยภาวาย เวปุลฺลาย, อตฺถิ ภิกฺขเว กุสลากุสลา ธมฺมา สาวชฺชานวชฺชา ธมฺมา หีนปณีตา ธมฺมา กณฺห\-สุกฺก\-สปฺปฏิภาคา ธมฺมา, ตตฺถ โยนิโสมนสิการ\-พหุลีกาโร, อยมนาหาโร อนุปฺปนฺนาย วา วิจิกิจฺฉาย อุปฺปาทาย อุปฺปนฺนาย วา วิจิกิจฺฉาย ภิยฺโยภาวาย เวปุลฺลาย.}

\prose{โก จ ภิกฺขเว อนาหาโร อนุปฺปนฺนสฺส วา สติสมฺโพชฺฌงฺคสฺส อุปฺปาทาย อุปฺปนฺนสฺส วา สติสมฺโพชฺฌงฺคสฺส ภาวนาย ปาริปูริยา, อตฺถิ ภิกฺขเว สติสมฺโพชฺฌงฺค\-ฏฺฐานียา ธมฺมา, ตตฺถ อมนสิการพหุลีกาโร, อยมนาหาโร อนุปฺปนฺนสฺส วา สติสมฺโพชฺฌงฺคสฺส อุปฺปาทาย อุปฺปนฺนสฺส วา สติสมฺโพชฺฌงฺคสฺส ภาวนาย ปาริปูริยา.}

\prose{\csromanfolio{94}โก จ ภิกฺขเว อนาหาโร อนุปฺปนฺนสฺส วา ธมฺมวิจย\-สมฺโพชฺฌงฺคสฺส อุปฺปาทาย อุปฺปนฺนสฺส วา ธมฺมวิจย\-สมฺโพชฺฌงฺคสฺส ภาวนาย ปาริปูริยา, อตฺถิ ภิกฺขเว กุสลากุสลา ธมฺมา สาวชฺชานวชฺชา ธมฺมา หีนปณีตา ธมฺมา กณฺห\-สุกฺก\-สปฺปฏิภาคา ธมฺมา, ตตฺถ อมนสิการพหุลีกาโร, อยมนาหาโร อนุปฺปนฺนสฺส วา ธมฺมวิจย\-สมฺโพชฺฌงฺคสฺส อุปฺปาทาย อุปฺปนฺนสฺส วา ธมฺมวิจย\-สมฺโพชฺฌงฺคสฺส ภาวนาย ปาริปูริยา.}

\prose{โก จ ภิกฺขเว อนาหาโร อนุปฺปนฺนสฺส วา วีริย\-สมฺโพชฺฌงฺคสฺส อุปฺปาทาย อุปฺปนฺนสฺส วา วีริย\-สมฺโพชฺฌงฺคสฺส ภาวนาย ปาริปูริยา, อตฺถิ ภิกฺขเว อารมฺภธาตุ นิกฺกมธาตุ ปรกฺกมธาตุ, ตตฺถ อมนสิการพหุลีกาโร, อยมนาหาโร อนุปฺปนฺนสฺส วา วีริย\-สมฺโพชฺฌงฺคสฺส อุปฺปาทาย อุปฺปนฺนสฺส วา วีริย\-สมฺโพชฺฌงฺคสฺส ภาวนาย ปาริปูริยา.}

\prose{โก จ ภิกฺขเว อนาหาโร อนุปฺปนฺนสฺส วา ปีติสมฺโพชฺฌงฺคสฺส อุปฺปาทาย อุปฺปนฺนสฺส วา ปีติสมฺโพชฺฌงฺคสฺส ภาวนาย ปาริปูริยา, อตฺถิ ภิกฺขเว ปีติสมฺโพชฺฌงฺค\-ฏฺฐานียา ธมฺมา, ตตฺถ อมนสิการพหุลีกาโร, อยมนาหาโร อนุปฺปนฺนสฺส วา ปีติสมฺโพชฺฌงฺคสฺส อุปฺปาทาย อุปฺปนฺนสฺส วา ปีติสมฺโพชฺฌงฺคสฺส ภาวนาย ปาริปูริยา.}

\prose{โก จ ภิกฺขเว อนาหาโร อนุปฺปนฺนสฺส วา ปสฺสทฺธิ\-สมฺโพชฺฌงฺคสฺส อุปฺปาทาย อุปฺปนฺนสฺส วา ปสฺสทฺธิ\-สมฺโพชฺฌงฺคสฺส ภาวนาย ปาริปูริยา, อตฺถิ ภิกฺขเว กายปฺปสฺสทฺธิ จิตฺต\-ปฺปสฺสทฺธิ, ตตฺถ อมนสิการพหุลีกาโร อยมนาหาโร อนุปฺปนฺนสฺส วา ปสฺสทฺธิ\-สมฺโพชฺฌงฺคสฺส อุปฺปาทาย อุปฺปนฺนสฺส วา ปสฺสทฺธิ\-สมฺโพชฺฌงฺคสฺส ภาวนาย ปาริปูริยา.}

\prose{โก จ ภิกฺขเว อนาหาโร อนุปฺปนฺนสฺส วา สมาธิ\-สมฺโพชฺฌงฺคสฺส อุปฺปาทาย อุปฺปนฺนสฺส วา สมาธิ\-สมฺโพชฺฌงฺคสฺส ภาวนาย ปาริปูริยา, อตฺถิ ภิกฺขเว สมถ\-นิมิตฺตํ อพฺยคฺค\-นิมิตฺตํ, ตตฺถ อมนสิการพหุลีกาโร อยมนาหาโร อนุปฺปนฺนสฺส วา สมาธิ\-สมฺโพชฺฌงฺคสฺส อุปฺปาทาย อุปฺปนฺนสฺส วา สมาธิสมฺโพงฺคสฺส ภาวนาย ปาริปูริยา.}

\prose{โก จ ภิกฺขเว อนาหาโร อนุปฺปนฺนสฺส วา อุเปกฺขา\-สมฺโพชฺฌงฺคสฺส อุปฺปาทาย อุปฺปนฺนสฺส วา อุเปกฺขา\-สมฺโพชฺฌงฺคสฺส ภาวนาย ปาริปูริยา, อตฺถิ ภิกฺขเว อุเปกฺขา\-สมฺโพชฺฌงฺค\-ฏฺฐานียา ธมฺมา, ตตฺถ อมนสิการพหุลีกาโร, อยมนาหาโร อนุปฺปนฺนสฺส วา อุเปกฺขา\-สมฺโพชฺฌงฺคสฺส}

\vspace{\tipitakaparskip}
\csromancenter{อุปฺปาทาย อุปฺปนฺนสฺส วา อุเปกฺขา\-สมฺโพชฺฌงฺคสฺส ภาวนาย ปาริปูริยาติ.  ปฐมํ.}
\csromanmatikaanchor{mtk.14}
\csromantocmarkref{paragraph}{2. ปริยายสุตฺต}{mtk.14}
\bookmark[dest={mtk.14},level=4]{2. ปริยายสุตฺต}
\csromanheader{4}{2. \textbf{ปริยายสุตฺต}}
\proseitem{233}{\csromanfolio{95}อถ โข สมฺพหุลา ภิกฺขู ปุพฺพณฺหสมยํ นิวาเสตฺวา ปตฺต\-จีวรมา\-ทาย สวตฺถิํ ปิณฺฑาย ปวิสิํสุ. อถ โข เตสํ ภิกฺขูนํ เอตทโหสิ “อติปฺปโค โข ตาว สาวตฺถิยํ ปิณฺฑาย จริตุํ, ยํนูน มยํ เยน อญฺญ\-ติตฺถิยานํ ปริพฺพาชกานํ อาราโม เตนุปสงฺกเมยฺยามา”ติ.}

\prose{อถ โข เต ภิกฺขู เยน อญฺญ\-ติตฺถิยานํ ปริพฺพาชกานํ อาราโม เตนุ\-ปสงฺกมิํสุ, อุปสงฺกมิตฺวา เตหิ อญฺญติตฺถิเยหิ ปริพฺพาชเกหิ สทฺธิํ สมฺโมทิํสุ, สมฺโมทนียํ กถํ สารณียํ วีติสาเรตฺวา เอกมนฺตํ นิสีทิํสุ, เอกมนฺตํ นิสินฺเน โข เต ภิกฺขู อญฺญติตฺถิยา ปริพฺพาชกา เอตทโวจุํ\csromandash{}}

\prose{สมโณ อาวุโส โคตโม สาวกานํ เอวํ ธมฺมํ เทเสติ “เอถ ตุเมฺห ภิกฺขเว ปญฺจ นีวรเณ ปหาย เจตโส อุปกฺกิเลเส ปญฺญาย ทุพฺพลีกรเณ สตฺต โพชฺฌงฺเค ยถาภูตํ ภาเวถา”ติ. มยมฺปิ โข อาวุโส สาวกานํ เอวํ ธมฺมํ เทเสม “เอถ ตุเมฺห อาวุโส ปญฺจ นีวรเณ ปหาย เจตโส อุปกฺกิเลเส ปญฺญาย ทุพฺพลีกรเณ สตฺต โพชฺฌงฺเค ยถาภูตํ ภาเวถา”ติ. อิธ โน อาวุโส โก วิเสโส, โก อธิปฺปยาโส, กิํ นานากรณํ สมณสฺส วา โคตมสฺส อมฺหากํ วา, ยทิทํ ธมฺมเทสนาย วา ธมฺมเทสนํ อนุสาสนิยา วา อนุสาสนินฺติ.}

\prose{อถ โข เต ภิกฺขู เตสํ อญฺญ\-ติตฺถิยานํ ปริพฺพาชกานํ ภาสิตํ เนว อภินนฺทิํสุ นปฺปฏิกฺโกสิํสุ, อนภินนฺทิตฺวา อปฺปฏิกฺโกสิตฺวา อุฏฺฐายาสนา ปกฺกมิํสุ “ภควโต สนฺติเก เอตสฺส ภาสิตสฺส อตฺถํ อาชานิสฺสามา”ติ. อถ โข เต ภิกฺขู สาวตฺถิํ ปิณฺฑาย จริตฺวา ปจฺฉาภตฺตํ ปิณฺฑปาต\-ปฏิกฺกนฺตา เยน ภควา เตนุ\-ปสงฺกมิํสุ, อุปสงฺกมิตฺวา ภควนฺตํ อภิวาเทตฺวา เอกมนฺตํ นิสีทิํสุ, เอกมนฺตํ นิสินฺนา โข เต ภิกฺขู ภควนฺตํ เอตทโวจุํ\csromandash{}}

\prose{\csromanfolio{96}อิธ มยํ ภนฺเต ปุพฺพณฺหสมยํ นิวาเสตฺวา ปตฺต\-จีวรมา\-ทาย สาวตฺถิํ ปิณฺฑาย ปวิสิมฺห, เตสํ โน ภนฺเต อมฺหากํ เอตทโหสิ “อติปฺปโค โข ตาว สาวตฺถิยํ ปิณฺฑาย จริตุํ, ยํนูน มยํ เยน อญฺญ\-ติตฺถิยานํ ปริพฺพาชกานํ อาราโม เตนุปสงฺกเมยฺยามา”ติ. อถ โข มยํ ภนฺเต เยน อญฺญ\-ติตฺถิยานํ ปริพฺพาชกานํ อาราโม เตนุ\-ปสงฺกมิมฺห, อุปสงฺกมิตฺวา เตหิ อญฺญติตฺถิเยหิ ปริพฺพาชเกหิ สทฺธิํ สมฺโมทิมฺห, สมฺโมทนียํ กถํ สารณียํ วีติสาเรตฺวา เอกมนฺตํ นิสีทิมฺห, เอกมนฺตํ นิสินฺเน โข อเมฺห ภนฺเต อญฺญติตฺถิยา ปริพฺพาชกา เอตทโวจุํ\csromandash{}}

\prose{สมโณ อาวุโส โคตโม สาวกานํ เอวํ ธมฺมํ เทเสติ “เอถ ตุเมฺห ภิกฺขเว ปญฺจ นีวรเณ ปหาย เจตโส อุปกฺกิเลเส ปญฺญาย ทุพฺพลีกรเณ สตฺต โพชฺฌงฺเค ยถาภูตํ ภาเวถา”ติ. มยมฺปิ โข อาวุโส สาวกานํ เอวํ ธมฺมํ เทเสม “เอถ ตุเมฺห อาวุโส ปญฺจ นีวรเณ ปหาย เจตโส อุปกฺกิเลเส ปญฺญาย ทุพฺพลีกรเณ สตฺต โพชฺฌงฺเค ยถาภูตํ ภาเวถา”ติ, อิธ โน อาวุโส โก วิเสโส, โก อธิปฺปยาโส, กิํ นานากรณํ สมณสฺส วา โคตมสฺส อมฺหากํ วา, ยทิทํ ธมฺมเทสนาย วา ธมฺมเทสนํ อนุสาสนิยา วา อนุสาสนินฺ”ติ.}

\prose{อถ โข มยํ ภนฺเต เตสํ อญฺญ\-ติตฺถิยานํ ปริพฺพาชกานํ ภาสิตํ เนว อภินนฺทิมฺห นปฺปฏิกฺโกสิมฺห, อนภินนฺทิตฺวา อปฺปฏิกฺโกสิตฺวา อุฏฺฐายาสนา ปกฺกมิมฺห “ภควโต สนฺติเก เอตสฺส ภาสิตสฺส อตฺถํ อาชานิสฺสามา”ติ.}

\prose{เอวํวาทิโน ภิกฺขเว อญฺญติตฺถิยา ปริพฺพาชกา เอวมสฺสุวจนียา “อตฺถิ ปนาวุโส ปริยาโย, ยํ ปริยายํ อาคมฺม ปญฺจ นีวรณา ทส โหนฺติ, สตฺต โพชฺฌงฺคา จตุทฺทสา”ติ. เอวํ ปุฏฺฐา ภิกฺขเว อญฺญติตฺถิยา ปริพฺพาชกา น เจว สมฺปายิสฺสนฺติ, อุตฺตริญฺจ วิฆาตํ อาปชฺชิสฺสนฺติ. ตํ กิสฺส เหตุ, ยถา ตํ ภิกฺขเว อวิสยสฺมิํ. นาหํ ตํ ภิกฺขเว ปสฺสามิ สเทวเก โลเก สมารเก สพฺรหฺมเก สสฺสมณ\-พฺราหฺมณิยา ปชาย สเทว\-มนุสฺสาย, โย อิเมสํ ปญฺหานํ เวยฺยากรเณน จิตฺตํ อาราเธยฺย อญฺญตฺร ตถาคเตน วา ตถาคต\-สาวเกน วา อิโต วา ปน สุตฺวา.}

\prose{\csromanfolio{97}กตโม จ ภิกฺขเว ปริยาโย, ยํ ปริยายํ อาคมฺม ปญฺจ นีวรณา ทส โหนฺติ. ยทปิ ภิกฺขเว อชฺฌตฺตํ กามจฺฉนฺโท, ตทปิ นีวรณํ, ยทปิ พหิทฺธา กามจฺฉนฺโท, ตทปิ นีวรณํ กามจฺฉนฺท\-นีวรณนฺติ อิติ หิทํ อุทฺเทสํ คจฺฉติ, ตทมินาเปตํ ปริยาเยน ทฺวยํ โหติ. ยทปิ ภิกฺขเว อชฺฌตฺตํ พฺยาปาโท, ตทปิ นีวรณํ, ยทปิ พหิทฺธา พฺยาปาโท, ตทปิ นีวรณํ, พฺยาปาท\-นีวรณนฺติ อิติ หิทํ อุทฺเทสํ คจฺฉติ, ตทมินาเปตํ ปริยาเยน ทฺวยํ โหติ. ยทปิ ภิกฺขเว ถินํ, ตทปิ นีวรณํ, ยทปิ มิทฺธํ, ตทปิ นีวรณํ, ถินมิทฺธ\-นีวรณนฺติ อิติ หิทํ อุทฺเทสํ คจฺฉติ, ตทมินาเปตํ ปริยาเยน ทฺวยํ โหติ. ยทปิ ภิกฺขเว อุทฺธจฺจํ, ตทปิ นีวรณํ, ยทปิ กุกฺกุจฺจํ, ตทปิ นีวรณํ, อุทฺธจฺจกุกฺกุจฺจ\-นีวรณนฺติ อิติ หิทํ อุทฺเทสํ คจฺฉติ, ตทมินาเปตํ ปริยาเยน ทฺวยํ โหติ. ยทปิ ภิกฺขเว อชฺฌตฺตํ ธมฺเมสุ วิจิกิจฺฉา, ตทปิ นีวรณํ, ยทปิ พหิทฺธา ธมฺเมสุ วิจิกิจฺฉา, ตทปิ นีวรณํ, วิจิกิจฺฉา\-นีวรณนฺติ อิติ หิทํ อุทฺเทสํ คจฺฉติ, ตทมินาเปตํ ปริยาเยน ทฺวยํ โหติ. อยํ โข ภิกฺขเว ปริยาโย, ยํ ปริยายํ อาคมฺม ปญฺจ นีวรณา ทส โหนฺติ.}

\prose{กตโม จ ภิกฺขเว ปริยาโย, ยํ ปริยายํ อาคมฺม สตฺต โพชฺฌงฺคา จตุทฺทส โหนฺติ. ยทปิ ภิกฺขเว อชฺฌตฺตํ ธมฺเมสุ สติ, ตทปิ สติสมฺโพชฺฌงฺโค, ยทปิ พหิทฺธา ธมฺเมสุ สติ, ตทปิ สติสมฺโพชฺฌงฺโค, สติสมฺโพชฺฌงฺโคติ อิติ หิทํ อุทฺเทสํ คจฺฉติ, ตทมินาเปตํ ปริยาเยน ทฺวยํ โหติ.}

\prose{ยทปิ ภิกฺขเว อชฺฌตฺตํ ธมฺเมสุ ปญฺญาย ปวิจินติ\csromansharedfootnote{5ef2b40ec808c89f}{ปวิจินาติ (ก)} ปวิจรติ ปริวีมํสมา\-ปชฺชติ, ตทปิ ธมฺมวิจย\-สมฺโพชฺฌงฺโค. ยทปิ พหิทฺธา ธมฺเมสุ ปญฺญาย ปวิจินติ ปวิจรติ ปริวีมํสมา\-ปชฺชติ, ตทปิ ธมฺมวิจย\-สมฺโพชฺฌงฺโค, ธมฺมวิจย\-สมฺโพชฺฌงฺโคติ อิติ หิทํ อุทฺเทสํ คจฺฉติ, ตทมินาเปตํ ปริยาเยน ทฺวยํ โหติ.}

\prose{ยทปิ ภิกฺขเว กายิกํ วีริยํ, ตทปิ วีริย\-สมฺโพชฺฌงฺโค, ยทปิ เจตสิกํ วีริยํ, ตทปิ วีริย\-สมฺโพชฺฌงฺโค, วีริย\-สมฺโพชฺฌงฺโคติ อิติ หิทํ อุทฺเทสํ คจฺฉติ, ตทมินาเปตํ ปริยาเยน ทฺวยํ โหติ.}

\prose{\csromanfolio{98}ยทปิ ภิกฺขเว สวิตกฺก\-สวิจารา ปีติ, ตทปิ ปีติสมฺโพชฺฌงฺโค, ยทปิ อวิตกฺก\-อวิจารา ปีติ, ตทปิ ปีติสมฺโพชฺฌงฺโค, ปีติสมฺโพชฺฌงฺโคติ อิติ หิทํ อุทฺเทสํ คจฺฉติ, ตทมินาเปตํ ปริยาเยน ทฺวยํ โหติ.}

\prose{ยทปิ ภิกฺขเว กายปฺปสฺสทฺธิ, ตทปิ ปสฺสทฺธิ\-สมฺโพชฺฌงฺโค, ยทปิ จิตฺต\-ปฺปสฺสทฺธิ, ตทปิ ปสฺสทฺธิ\-สมฺโพชฺฌงฺโค, ปสฺสทฺธิ\-สมฺโพชฺฌงฺโคติ อิติ หิทํ อุทฺเทสํ คจฺฉติ, ตทมินาเปตํ ปริยาเยน ทฺวยํ โหติ.}

\prose{ยทปิ ภิกฺขเว สวิตกฺโก สวิจาโร สมาธิ, ตทปิ สมาธิ\-สมฺโพชฺฌงฺโค, ยทปิ อวิตกฺก\-อวิจาโร สมาธิ, ตทปิ สมาธิ\-สมฺโพชฺฌงฺโค, สมาธิ\-สมฺโพชฺฌงฺโคติ อิติ หิทํ อุทฺเทสํ คจฺฉติ, ตทมินาเปตํ ปริยาเยน ทฺวยํ โหติ.}

\prose{ยทปิ ภิกฺขเว อชฺฌตฺตํ ธมฺเมสุ อุเปกฺขา, ตทปิ อุเปกฺขา\-สมฺโพชฺฌงฺโค, ยทปิ พหิทฺธา ธมฺเมสุ อุเปกฺขา, ตทปิ อุเปกฺขา\-สมฺโพชฺฌงฺโค, อุเปกฺขา\-สมฺโพชฺฌงฺโคติ อิติ หิทํ อุทฺเทสํ คจฺฉติ, ตทมินาเปตํ ปริยาเยน ทฺวยํ โหติ. อยํ โข ภิกฺขเว ปริยาโย, ยํ ปริยายํ อาคมฺม สตฺต โพชฺฌงฺคา จตุทฺทสาติ.  ทุติยํ.}

\csromanmatikaanchor{mtk.15}
\csromantocmarkref{paragraph}{3. อคฺคิสุตฺต}{mtk.15}
\bookmark[dest={mtk.15},level=4]{3. อคฺคิสุตฺต}
\csromanheader{4}{3. \textbf{อคฺคิสุตฺต}}
\proseitem{234}{อถ โข สมฺพหุลา ภิกฺขู ปุพฺพณฺหสมยํ นิวาเสตฺวา ปตฺต\-จีวรมา\-ทาย สาวตฺถิํ ปิณฺฑาย ปวิสิํสุ. (ปริยายสุตฺต\-สทิสํ.)}

\prose{เอวํวาทิโน ภิกฺขเว อญฺญติตฺถิยา ปริพฺพาชกา เอวมสฺสุ วจนียา “ยสฺมิํ อาวุโส สมเย ลีนํ จิตฺตํ โหติ, กตเมสํ ตสฺมิํ สมเย โพชฺฌงฺคานํ อกาโล ภาวนาย, กตเมสํ ตสฺมิํ สมเย โพชฺฌงฺคานํ กาโล ภาวนาย. ยสฺมิํ ปนาวุโส สมเย อุทฺธตํ จิตฺตํ โหติ, กตเมสํ ตสฺมิํ สมเย โพชฺฌงฺคานํ อกาโล ภาวนาย, กตเมสํ ตสฺมิํ สมเย โพชฺฌงฺคานํ กาโล ภาวนายา”ติ. เอวํ ปุฏฺฐา ภิกฺขเว อญฺญติตฺถิยา ปริพฺพาชกา น เจว สมฺปายิสฺสนฺติ, อุตฺตริญฺจ วิฆาตํ อาปชฺชิสฺสนฺติ. ตํ กิสฺส เหตุ, ยถา ตํ ภิกฺขเว อวิสยสฺมิํ.}

\prose{นาหํ ตํ ภิกฺขเว ปสฺสามิ สเทวเก โลเก สมารเก สพฺรหฺมเก สสฺสมณ\-พฺราหฺมณิยา ปชาย สเทว\-มนุสฺสาย, โย อิเมสํ ปญฺหานํ}

\prose{\csromanfolio{99}เวยฺยากรเณน จิตฺตํ อาราเธยฺย อญฺญตฺร ตถาคเตน วา ตถาคต\-สาวเกน วา อิโต วา ปน สุตฺวา.}

\prose{ยสฺมิํ ภิกฺขเว สมเย ลีนํ จิตฺตํ โหติ, อกาโล ตสฺมิํ สมเย ปสฺสทฺธิ\-สมฺโพชฺฌงฺคสฺส ภาวนาย, อกาโล สมาธิ\-สมฺโพชฺฌงฺคสฺส ภาวนาย, อกาโล อุเปกฺขา\-สมฺโพชฺฌงฺคสฺส ภาวนาย. ตํ กิสฺส เหตุ, ลีนํ ภิกฺขเว จิตฺตํ, ตํ เอเตหิ ธมฺเมหิ ทุสฺสมุฏฺฐาปยํ โหติ.}

\prose{เสยฺยถาปิ ภิกฺขเว ปุริโส ปริตฺตํ อคฺคิํ อุชฺชาเลตุกาโม อสฺส, โส ตตฺถ อลฺลานิ เจว ติณานิ ปกฺขิเปยฺย, อลฺลานิ จ โคมยานิ ปกฺขิเปยฺย, อลฺลานิ จ กฏฺฐานิ ปกฺขิเปยฺย, อุทกวาตญฺจ ทเทยฺย, ปํสุเกน จ โอกิเรยฺย, ภพฺโพ นุ โข โส ปุริโส ปริตฺตํ อคฺคิํ อุชฺชาเลตุนฺติ. โน เหตํ ภนฺเต.}

\prose{เอวํเมว โข ภิกฺขเว ยสฺมิํ สมเย ลีนํ จิตฺตํ โหติ, อกาโล ตสฺมิํ สมเย ปสฺสทฺธิ\-สมฺโพชฺฌงฺคสฺส ภาวนาย, อกาโล สมาธิ\-สมฺโพชฺฌงฺคสฺส ภาวนาย, อกาโล อุเปกฺขา\-สมฺโพชฺฌงฺคสฺส ภาวนาย. ตํ กิสฺส เหตุ, ลีนํ ภิกฺขเว จิตฺตํ, ตํ เอเตหิ ธมฺเมหิ ทุสฺสมุฏฺฐาปยํ โหติ.}

\prose{ยสฺมิํ จ โข ภิกฺขเว สมเย ลีนํ จิตฺตํ โหติ, กาโล ตสฺมิํ สมเย ธมฺมวิจย\-สมฺโพชฺฌงฺคสฺส ภาวนาย, กาโล วีริย\-สมฺโพชฺฌงฺคสฺส ภาวนาย, กาโล ปีติสมฺโพชฺฌงฺคสฺส ภาวนาย. ตํ กิสฺส เหตุ, ลีนํ ภิกฺขเว จิตฺตํ, ตํ เอเตหิ ธมฺเมหิ สุสมุฏฺฐาปยํ โหติ.}

\prose{เสยฺยถาปิ ภิกฺขเว ปุริโส ปริตฺตํ อคฺคิํ อุชฺชาเลตุกาโม อสฺส, โส ตตฺถ สุกฺขานิ เจว ติณานิ ปกฺขิเปยฺย, สุกฺขานิ โคมยานิ ปกฺขิเปยฺย, สุกฺขานิ กฏฺฐานิ ปกฺขิเปยฺย, มุขวาตญฺจ ทเทยฺย, น จ ปํสุเกน โอกิเรยฺย, ภพฺโพ นุ โข โส ปุริโส ปริตฺตํ อคฺคิํ อุชฺชาเลตุนฺติ. เอวํ ภนฺเต.}

\prose{เอวเมว โข ภิกฺขเว ยสฺมิํ สมเย ลีนํ จิตฺตํ โหติ, กาโล ตสฺมิํ สมเย ธมฺมวิจย\-สมฺโพชฺฌงฺคสฺส ภาวนาย, กาโล วิริย\-สมฺโพชฺฌงฺคสฺส ภาวนาย, กาโล ปีติสมฺโพชฺฌงฺคสฺส ภาวนาย. ตํ กิสฺส เหตุ, ลีนํ ภิกฺขเว จิตฺตํ, ตํ เอเตหิ ธมฺเมหิ สุสมุฏฺฐาปยํ โหติ.}

\prose{\csromanfolio{100}ยสฺมิํ ภิกฺขเว สมเย อุทฺธตํ จิตฺตํ โหติ, อกาโล ตสฺมิํ สมเย ธมฺมวิจย\-สมฺโพชฺฌงฺคสฺส ภาวนาย, อกาโล วีริย\-สมฺโพชฺฌงฺคสฺส ภาวนาย, อกาโล ปีติสมฺโพชฺฌงฺคสฺส ภาวนาย. ตํ กิสฺส เหตุ, อุทฺธตํ ภิกฺขเว จิตฺตํ, ตํ เอเตหิ ธมฺเมหิ ทุวูปสมยํ โหติ.}

\prose{เสยฺยถาปิ ภิกฺขเว ปุริโส มหนฺตํ อคฺคิกฺขนฺธํ นิพฺพาเปตุกาโม อสฺส, โส ตตฺถ สุกฺขานิ เจว ติณานิ ปกฺขิเปยฺย, สุกฺขานิ จ โคมยานิ ปกฺขิเปยฺย, สุกฺขานิ จ กฏฺฐานิ ปกฺขิเปยฺย, มุขวาตญฺจ ทเทยฺย, น จ ปํสุเกน โอกิเรยฺย, ภพฺโพ นุ โข โส ปุริโส มหนฺตํ อคฺคิกฺขนฺธํ นิพฺพาเปตุนฺติ. โน เหตํ ภนฺเต.}

\prose{เอวเมว โข ภิกฺขเว ยสฺมิํ สมเย อุทฺธตํ จิตฺตํ โหติ, อกาโล ตสฺมิํ สมเย ธมฺมวิจย\-สมฺโพชฺฌงฺคสฺส ภาวนาย, อกาโล วีริย\-สมฺโพชฺฌงฺคสฺส ภาวนาย, อกาโล ปีติสมฺโพชฺฌงฺคสฺส ภาวนาย. ตํ กิสฺส เหตุ, อุทฺธตํ ภิกฺขเว จิตฺตํ, ตํ เอเตหิ ธมฺเมหิ ทุวูปสมยํ โหติ.}

\prose{ยสฺมิํ จ โข ภิกฺขเว สมเย อุทฺธตํ จิตฺตํ โหติ, กาโล ตสฺมิํ สมเย ปสฺสทฺธิ\-สมฺโพชฺฌงฺคสฺส ภาวนาย, กาโล สมาธิ\-สมฺโพชฺฌงฺคสฺส ภาวนาย, กาโล อุเปกฺขา\-สมฺโพชฺฌงฺคสฺส ภาวนาย. ตํ กิสฺส เหตุ, อุทฺธตํ ภิกฺขเว จิตฺตํ, ตํ เอเตหิ ธมฺเมหิ สุวูปสมยํ โหติ.}

\prose{เสยฺยถาปิ ภิกฺขเว ปุริโส มหนฺตํ อคฺคิกฺขนฺธํ นิพฺพาเปตุกาโม อสฺส, โส ตตฺถ อลฺลานิ เจว ติณานิ ปกฺขิเปยฺย, อลฺลานิ จ โคมยานิ ปกฺขิเปยฺย, อลฺลานิ จ กฏฺฐานิ ปกฺขิเปยฺย, อุทกวาตญฺจ ทเทยฺย, ปํสุเกน จ โอกิเรยฺย, ภพฺโพ นุ โข ปุริโส มหนฺตํ อคฺคิกฺขนฺธํ นิพฺพาเปตุนฺติ. เอวํ ภนฺเต.}

\prose{เอวเมว โข ภิกฺขเว ยสฺมิํ สมเย อุทฺธตํ จิตฺตํ โหติ, กาโล ตสฺมิํ สมเย ปสฺสทฺธิ\-สมฺโพชฺฌงฺคสฺส ภาวนาย, กาโล สมาธิ\-สมฺโพชฺฌงฺคสฺส ภาวนาย, กาโล อุเปกฺขา\-สมฺโพชฺฌงฺคสฺส ภาวนาย. ตํ กิสฺส เหตุ, อุทฺธตํ ภิกฺขเว จิตฺตํ, ตํ เอเตหิ ธมฺเมหิ สุวูปสมยํ โหติ. สติํ จ ขฺวาหํ ภิกฺขเว สพฺพตฺถิกํ วทามีติ.  ตติยํ.}

\csromanmatikaanchor{mtk.16}
\csromantocmarkref{paragraph}{4. เมตฺตาสหคตสุตฺต (หลิทฺทวสนสุตฺต)}{mtk.16}
\bookmark[dest={mtk.16},level=4]{4. เมตฺตาสหคตสุตฺต (หลิทฺทวสนสุตฺต)}
\csromanheader{4}{4. เมตฺตาสหคต\-สุตฺต}
\proseitem{235}{\csromanfolio{101}เอกํ สมยํ ภควา โกลิเยสุ วิหรติ หลิทฺทวสนํ นาม โกลิยานํ นิคโม. อถ โข สมฺพหุลา ภิกฺขู ปุพฺพณฺหสมยํ นิวาเสตฺวา ปตฺต\-จีวรมา\-ทาย หลิทฺทวสนํ ปิณฺฑาย ปวิสิํสุ. อถ โข เตสํ ภิกฺขูนํ เอตทโหสิ “อติปฺปโค โข ตาว หลิทฺทวสเน ปิณฺฑาย จริตุํ, ยํนูน มยํ เยน อญฺญ\-ติตฺถิยานํ ปริพฺพาชกานํ อาราโม เตนุปสงฺกเมยฺยามา”ติ.}

\prose{อถ โข เต ภิกฺขู เยน อญฺญ\-ติตฺถิยานํ ปริพฺพาชกานํ อาราโม เตนุ\-ปสงฺกมิํสุ, อุปสงฺกมิตฺวา เตหิ อญฺญติตฺถิเยหิ ปริพฺพาชเกหิ สทฺธิํ สมฺโมทิํสุ, สมฺโมทนียํ กถํ สารณียํ วีติสาเรตฺวา เอกมนฺตํ นิสีทิํสุ, เอกมนฺตํ นิสินฺเน โข เต ภิกฺขู อญฺญติตฺถิยา ปริพฺพาชกา เอตทโวจุํ\csromandash{}}

\prose{“สมโณ อาวุโส โคตโม สาวกานํ เอวํ ธมฺมํ เทเสติ ‘เอถ ตุเมฺห ภิกฺขเว ปญฺจ นีวรเณ ปหาย เจตโส อุปกฺกิเลเส ปญฺญาย ทุพฺพลีกรเณ เมตฺตา\-สหคเตน เจตสา เอกํ ทิสํ ผริตฺวา วิหรถ. ตถา ทุติยํ. ตถา ตติยํ. ตถา จตุตฺถํ. อิติ อุทฺธมโธ ติริยํ สพฺพธิ สพฺพตฺตตาย สพฺพาวนฺตํ โลกํ เมตฺตา\-สหคเตน เจตสา วิปุเลน มหคฺคเตน อปฺปมาเณน อเวเรน อพฺยาปชฺเชน ผริตฺวา วิหรถ. กรุณา\-สหคเตน เจตสา เอกํ ทิสํ ผริตฺวา วิหรถ. ตถา ทุติยํ. ตถา ตติยํ. ตถา จตุตฺถํ. อิติ อุทฺธมโธ ติริยํ สพฺพธิ สพฺพตฺตตาย สพฺพาวนฺตํ โลกํ กรุณา\-สหคเตน เจตสา วิปุเลน มหคฺคเตน อปฺปมาเณน อเวเรน อพฺยาปชฺเชน ผริตฺวา วิหรถ. มุทิตา\-สหคเตน เจตสา เอกํ ทิสํ ผริตฺวา วิหรถ. ตถา ทุติยํ. ตถา ตติยํ. ตถา จตุตฺถํ. อิติ อุทฺธมโธ ติริยํ สพฺพธิ สพฺพตฺตตาย สพฺพาวนฺตํ โลกํ มุทิตา\-สหคเตน เจตสา วิปุเลน มหคฺคเตน อปฺปมาเณน อเวเรน อพฺยาปชฺเชน ผริตฺวา วิหรถ. อุเปกฺขา\-สหคเตน เจตสา เอกํ ทิสํ ผริตฺวา วิหรถ. ตถา ทุติยํ. ตถา ตติยํ. ตถา จตุตฺถํ. อิติ อุทฺธมโธ ติริยํ สพฺพธิ สพฺพตฺตตาย สพฺพาวนฺตํ โลกํ อุเปกฺขา\-สหคเตน เจตสา วิปุเลน มหคฺคเตน อปฺปมาเณน อเวเรน อพฺยาปชฺเชน ผริตฺวา วิหรถา’ติ.}

\prose{\csromanfolio{102}มยมฺปิ โข อาวุโส สาวกานํ เอวํ ธมฺมํ เทเสม ‘เอถ ตุเมฺห อาวุโส ปญฺจ นีวรเณ ปหาย เจตโส อุปกฺกิเลเส ปญฺญาย ทุพฺพลีกรเณ เมตฺตา\-สหคเตน เจตสา เอกํ ทิสํ ผริตฺวา วิหรถ ฯเปฯ กรุณา\-สหคเตน เจตสา. มุทิตา\-สหคเตน เจตสา. อุเปกฺขา\-สหคเตน เจตสา เอกํ ทิสํ ผริตฺวา วิหรถ. ตถา ทุติยํ. ตถา ตติยํ. ตถา จตุตฺถํ. อิติ อุทฺธมโธ ติริยํ สพฺพธิ สพฺพตฺตตาย สพฺพาวนฺตํ โลกํ อุเปกฺขา\-สหคเตน เจตสา วิปุเลน มหคฺคเตน อปฺปมาเณน อเวเรน อพฺยาปชฺเชน ผริตฺวา วิหรถา’ติ. อิธ โน อาวุโส โก วิเสโส, โก อธิปฺปยาโส, กิํ นานากรณํ สมณสฺส วา โคตมสฺส อมฺหากํ วา, ยทิทํ ธมฺมเทสนาย วา ธมฺมเทสนํ, อนุสาสนิยา วา อนุสาสนินฺ”ติ.}

\prose{อถ โข เต ภิกฺขู เตสํ อญฺญ\-ติตฺถิยานํ ปริพฺพาชกานํ ภาสิตํ เนว อภินนฺทิํสุ นปฺปฏิกฺโกสิํสุ, อนภินนฺทิตฺวา อปฺปฏิกฺโกสิตฺวา อุฏฺฐายาสนา ปกฺกมิํสุ “ภควโต สนฺติเก เอตสฺส ภาสิตสฺส อตฺถํ อาชานิสฺสามา”ติ. อถ โข เต ภิกฺขู หลิทฺทวสเน ปิณฺฑาย จริตฺวา ปจฺฉาภตฺตํ ปิณฺฑปาต\-ปฏิกฺกนฺตา เยน ภควา เตนุ\-ปสงฺกมิํสุ, อุปสงฺกมิตฺวา ภควนฺตํ อภิวาเทตฺวา เอกมนฺตํ นิสินฺนา โข เต ภิกฺขู ภควนฺตํ เอตทโวจุํ\csromandash{}}

\prose{“อิธ มยํ ภนฺเต ปุพฺพณฺหสมยํ นิวาเสตฺวา ปตฺต\-จีวรมา\-ทาย หลิทฺทวสเน ปิณฺฑาย ปวิสิมฺห, เตสํ โน ภนฺเต อมฺหากํ เอตทโหสิ ‘อติปฺปโค โข ตาว หลิทฺทวสเน ปิณฺฑาย จริตุํ, ยํนูน มยํ เยน อญฺญ\-ติตฺถิยานํ ปริพฺพาชกานํ อาราโม เตนุปสงฺกเมยฺยามา’ติ.}

\prose{อถ โข มยํ ภนฺเต เยน อญฺญ\-ติตฺถิยานํ ปริพฺพาชกานํ อาราโม เตนุ\-ปสงฺกมิมฺห, อุปสงฺกมิตฺวา เตหิ อญฺญติตฺถิเยหิ ปริพฺพาชเกหิ สทฺธิํ สมฺโมทิมฺห, สมฺโมทนียํ กถํ สารณียํ วีติสาเรตฺวา เอกมนฺตํ นิสีทิมฺห, เอกมนฺตํ นิสินฺเน โข อเมฺห ภนฺเต เต อญฺญติตฺถิยา ปริพฺพาชกา เอตทโวจุํ\csromandash{}}

\prose{สมโณ อาวุโส โคตโม สาวกานํ เอวํ ธมฺมํ เทเสติ ‘เอถ ตุเมฺห ภิกฺขเว ปญฺจ นีวรเณ ปหาย เจตโส อุปกฺกิเลเส ปญฺญาย ทุพฺพลีกรเณ เมตฺตา\-สหคเตน เจตสา เอกํ ทิสํ}

\prose{\csromanfolio{103}ผริตฺวา วิหรถ ฯเปฯ กรุณา\-สหคเตน เจตสา ฯเปฯ มุทิตา\-สหคเตน เจตสา ฯเปฯ อุเปกฺขา\-สหคเตน เจตสา เอกํ ทิสํ ผริตฺวา วิหรถ. ตถา ทุติยํ. ตถา ตติยํ. ตถา จตุตฺถํ. อิติ อุทฺธมโธ ติริยํ สพฺพธิ สพฺพตฺตตาย สพฺพาวนฺตํ โลกํ อุเปกฺขา\-สหคเตน เจตสา วิปุเลน มหคฺคเตน อปฺปมาเณน อเวเรน อพฺยาปชฺเชน ผริตฺวา วิหรถา’ติ.}

\prose{มยมฺปิ โข อาวุโส สาวกานํ เอวํ ธมฺมํ เทเสม ‘เอถ ตุเมฺห อาวุโส ปญฺจ นีวรเณ ปหาย เจตโส อุปกฺกิเลเส ปญฺญาย ทุพฺพลีกรเณ เมตฺตา\-สหคเตน เจตสา เอกํ ทิสํ ผริตฺวา วิหรถ ฯเปฯ กรุณา\-สหคเตน เจตสา ฯเปฯ มุทิตา\-สหคเตน เจตสา ฯเปฯ อุเปกฺขา\-สหคเตน เจตสา เอกํ ทิสํ ผริตฺวา วิหรถ. ตถา ทุติยํ. ตถ ตติยํ. ตถา จตุตฺถํ. อิติ อุทฺธมโธ ติริยํ สพฺพธิ สพฺพตฺตตาย สพฺพาวนฺตํ โลกํ อุเปกฺขา\-สหคเตน เจตสา วิปุเลน มหคฺคเตน อปฺปมาเณน อเวเรน อพฺยาปชฺเชน ผริตฺวา วิหรถา’ติ. อิธ โน อาวุโส โก วิเสโส, โก อธิปฺปยาโส, กิํ นานากรณํ สมณสฺส วา โคตมสฺส อมฺหากํ วา, ยทิทํ ธมฺมเทสนาย วา ธมฺมเทสนํ อนุสาสนิยา วา อนุสาสนินฺ’ติ.}

\prose{อถ โข มยํ ภนฺเต เตสํ อญฺญ\-ติตฺถิยานํ ปริพฺพาชกานํ ภาสิตํ เนว อภินนฺทิมฺห นปฺปฏิกฺโกสิมฺห, อนภินนฺทิตฺวา อปฺปฏิกฺโกสิตฺวา อุฏฺฐายาสนา ปกฺกมิมฺห ‘ภควโต สนฺติเก เอตสฺส ภาสิตสฺส อตฺถํ อาชานิสฺสามา’ติ”.}

\prose{เอวํวาทิโน ภิกฺขเว อญฺญติตฺถิยา ปริพฺพาชกา เอวมสฺสุ วจนียา “กถํ ภาวิตา ปนาวุโส เมตฺตา\-เจโตวิมุตฺติ กิํคติกา โหติ กิํปรมา กิํผลา กิํปริโยสานา, กถํ ภาวิตา ปนาวุโส กรุณา\-เจโตวิมุตฺติ กิํคติกา โหติ กิํปรมา กิํผลา กิํปริโยสานา, กถํ ภาวิตา ปนาวุโส มุทิตา\-เจโตวิมุตฺติ กิํคติกา โหติ กิํปรมา กิํผลา กิํปริโยสานา, กถํ ภาวิตา ปนาวุโส อุเปกฺขา\-เจโตวิมุตฺติ กิํคติกา โหติ กิํปรมา กิํผลา กิํปริโยสานา”ติ. เอวํ ปุฏฺฐา ภิกฺขเว อญฺญติตฺถิยา ปริพฺพาชกา น เจว สมฺปายิสฺสนฺติ, อุตฺตริญฺจ วิฆาตํ อาปชฺชิสฺสนฺติ. ตํ กิสฺส เหตุ, ยถา ตํ ภิกฺขเว อวิสยสฺมิํ. นาหํ ตํ ภิกฺขเว ปสฺสามิ}

\prose{\csromanfolio{104}สเทวเก โลเก สมารเก สพฺรหฺมเก สสฺสมณ\-พฺราหฺมณิยา ปชาย สเทว\-มนุสฺสาย, โย อิเมสํ ปญฺหานํ เวยฺยากรเณน จิตฺตํ อาราเธยฺย อญฺญตฺร ตถาคเตน วา ตถาคต\-สาวเกน วา อิโต วา ปน สุตฺวา.}

\prose{กถํ ภาวิตา จ ภิกฺขเว เมตฺตา\-เจโตวิมุตฺติ กิํคติกา โหติ กิํปรมา กิํผลา กิํปริโยสานา, อิธ ภิกฺขเว ภิกฺขุ เมตฺตาสหคตํ สติสมฺโพชฺฌงฺคํ ภาเวติ ฯเปฯ เมตฺตาสหคตํ อุเปกฺขา\-สมฺโพชฺฌงฺคํ ภาเวติ วิเวกนิสฺสิตํ วิราคนิสฺสิตํ นิโรธ\-นิสฺสิตํ โวสฺสคฺค\-ปริณามิํ. โส สเจ อากงฺขติ “อปฺปฏิกูเล ปฏิกูลสญฺญี วิหเรยฺยนฺ”ติ, ปฏิกูลสญฺญี ตตฺถ วิหรติ. สเจ อากงฺขติ “ปฏิกูเล อปฺปฏิกูลสญฺญี วิหเรยฺยนฺ”ติ, อปฺปฏิกูลสญฺญี ตตฺถ วิหรติ. สเจ อากงฺขติ “อปฺปฏิกูเล จ ปฏิกูเล จ ปฏิกูลสญฺญี วิหเรยฺยนฺ”ติ, ปฏิกูลสญฺญี ตตฺถ วิหรติ. สเจ อากงฺขติ “ปฏิกูเล จ อปฺปฏิกูเล จ อปฺปฏิกูลสญฺญี วิหเรยฺยนฺ”ติ, อปฺปฏิกูลสญฺญี ตตฺถ วิหรติ. สเจ อากงฺขติ “อปฺปฏิกูลญฺจ ปฏิกูลญฺจ ตทุภยํ อภินิวชฺเชตฺวา อุเปกฺขโก วิหเรยฺยํ สโต สมฺปชาโน”ติ, อุเปกฺขโก จ ตตฺถ วิหรติ สโต สมฺปชาโน. สุภํ วา โข ปน วิโมกฺขํ อุปสมฺปชฺช วิหรติ, สุภปรมาหํ ภิกฺขเว เมตฺตา\-เจโตวิมุตฺติํ วทามิ อิธปญฺญสฺส ภิกฺขุโน อุตฺตริวิมุตฺติํ อปฺปฏิวิชฺฌโต.}

\prose{กถํ ภาวิตา จ ภิกฺขเว กรุณา\-เจโตวิมุตฺติ กิํคติกา โหติ กิํปรมา กิํผลา กิํปริโยสานา, อิธ ภิกฺขเว ภิกฺขุ กรุณา\-สหคตํ สติสมฺโพชฺฌงฺคํ ภาเวติ ฯเปฯ กรุณา\-สหคตํ อุเปกฺขา\-สมฺโพชฺฌงฺคํ ภาเวติ วิเวกนิสฺสิตํ วิราคนิสฺสิตํ นิโรธ\-นิสฺสิตํ โวสฺสคฺค\-ปริณามิํ. โส สเจ อากงฺขติ “อปฺปฏิกูเล ปฏิกูลสญฺญี วิหเรยฺยนฺ”ติ, ปฏิกูลสญฺญี ตตฺถ วิหรติ ฯเปฯ. สเจ อากงฺขติ “อปฺปฏิกูลญฺจ ปฏิกูลญฺจ ตทุภยํ อภินิวชฺเชตฺวา อุเปกฺขโก วิหเรยฺยํ สโต สมฺปชาโน”ติ, อุเปกฺขโก ตตฺถ วิหรติ สโต สมฺปชาโน. สพฺพโส วา ปน รูปสญฺญานํ สมติกฺกมา ปฏิฆ\-สญฺญานํ อตฺถงฺคมา นานตฺต\-สญฺญานํ อมนสิการา “อนนฺโต อากาโส”ติ อากาสา\-นญฺจา\-ยตนํ อุปสมฺปชฺช วิหรติ, อากาสานญฺจายตนปรมาหํ ภิกฺขเว}

\prose{\csromanfolio{105}กรุณา\-เจโตวิมุตฺติํ วทามิ อิธปญฺญสฺส ภิกฺขุโน อุตฺตริวิมุตฺติํ อปฺปฏิวิชฺฌโต.}

\prose{กถํ ภาวิตา จ ภิกฺขเว มุทิตา\-เจโตวิมุตฺติ กิํคติกา โหติ กิํปรมา กิํผลา กิํปริโยสานา, อิธ ภิกฺขเว ภิกฺขุ มุทิตา\-สหคตํ สติสมฺโพชฺฌงฺคํ ภาเวติ ฯเปฯ มุทิตา\-สหคตํ อุเปกฺขา\-สมฺโพชฺฌงฺคํ ภาเวติ วิเวกนิสฺสิตํ วิราคนิสฺสิตํ นิโรธ\-นิสฺสิตํ โวสฺสคฺค\-ปริณามิํ. โส สเจ อากงฺขติ “อปฺปฏิกูเล ปฏิกูลสญฺญี วิหเรยฺยนฺ”ติ, ปฏิกูลสญฺญี ตตฺถ วิหรติ ฯเปฯ. สเจ อากงฺขติ “อปฺปฏิกูลญฺจ ปฏิกูลญฺจ ตทุภยํ อภินิวชฺเชตฺวา อุเปกฺขโก วิหเรยฺยํ สโต สมฺปชาโน”ติ, อุเปกฺขโก ตตฺถ วิหรติ สโต สมฺปชาโน. สพฺพโส วา ปน อากาสา\-นญฺจา\-ยตนํ สมติกฺกมฺม “อนนฺตํ วิญฺญาณนฺ”ติ วิญฺญาณญฺจา\-ยตนํ อุปสมฺปชฺช วิหรติ, วิญฺญาณญฺจายตนปรมาหํ ภิกฺขเว มุทิตา\-เจโตวิมุตฺติํ วทามิ อิธปญฺญสฺส ภิกฺขุโน อุตฺตริวิมุตฺติํ อปฺปฏิวิชฺฌโต.}

\prose{กถํ ภาวิตา จ ภิกฺขเว อุเปกฺขา\-เจโตวิมุตฺติ กิํคติกา โหติ กิํปรมา กิํผลา กิํปริโยสานา, อิธ ภิกฺขเว ภิกฺขุ อุเปกฺขา\-สหคตํ สติสมฺโพชฺฌงฺคํ ภาเวติ วิเวกนิสฺสิตํ วิราคนิสฺสิตํ นิโรธ\-นิสฺสิตํ โวสฺสคฺค\-ปริณามิํ ฯเปฯ อุเปกฺขา\-สหคตํ อุเปกฺขา\-สมฺโพชฺฌงฺคํ ภาเวติ วิเวกนิสฺสิตํ วิราคนิสฺสิตํ นิโรธ\-นิสฺสิตํ โวสฺสคฺค\-ปริณามิํ. โส สเจ อากงฺขติ “อปฺปฏิกูเล ปฏิกูลสญฺญี วิหเรยฺยนฺ”ติ, ปฏิกูลสญฺญี ตตฺถ วิหรติ. สเจ อากงฺขติ “ปฏิกูเล อปฺปฏิกูลสญฺญี วิหเรยฺยนฺ”ติ, อปฺปฏิกูลสญฺญี ตตฺถ วิหรติ. สเจ อากงฺขติ “อปฺปฏิกูเล จ ปฏิกูเล จ ปฏิกูลสญฺญี วิหเรยฺยนฺ”ติ, ปฏิกูลสญฺญี ตตฺถ วิหรติ. สเจ อากงฺขติ “ปฏิกูเล จ อปฺปฏิกูเล จ อปฺปฏิกูลสญฺญี วิหเรยฺยนฺ”ติ, อปฺปฏิกูลสญฺญี ตตฺถ วิหรติ. สเจ อากงฺขติ “อปฺปฏิกูลญฺจ ปฏิกูลญฺจ ตทุภยํ อภินิวชฺเชตฺวา อุเปกฺขโก วิหเรยฺยํ สโต สมฺปชาโน”ติ, อุเปกฺขโก ตตฺถ วิหรติ สโต สมฺผชาโน. สพฺพโส วา ปน วิญฺญาณญฺจา\-ยตนํ สมติกฺกมฺม “นตฺถิ กิญฺจี”ติ อากิญฺจญฺญา\-ยตนํ อุปสมฺปชฺช วิหรติ, อากิญฺจญฺญายตนปรมาหํ ภิกฺขเว อุเปกฺขา\-เจโตวิมุตฺติํ วทามิ อิธปญฺญสฺส ภิกฺขุโน อุตฺตริวิมุตฺตํ อปฺปฏิวิชฺฌโตติ.  จตุตฺถํ.}

\csromanmatikaanchor{mtk.17}
\csromantocmarkref{paragraph}{5. สงฺคารวสุตฺต}{mtk.17}
\bookmark[dest={mtk.17},level=4]{5. สงฺคารวสุตฺต}
\csromanheader{4}{5. \textbf{สงฺคารวสุตฺต}}
\proseitem{236}{\csromanfolio{106}สาวตฺถิ\-นิทานํ. อถ โข สงฺคารโว พฺราหฺมโณ เยน ภควา เตนุปสงฺกมิ, อุปสงฺกมิตฺวา ภควตา สทฺธิํ สมฺโมทิ, สมฺโมทนียํ กถํ สารณียํ วีติสาเรตฺวา เอกมนฺตํ นิสีทิ, เอกมนฺตํ นิสินฺโน โข สงฺคารโว พฺราหฺมโณ ภควนฺตํ เอตทโวจ\csromandash{}}

\prose{“โก นุ โข โภ โคตม เหตุ โก ปจฺจโย, เยเนกทา ทีฆรตฺตํ สชฺฌายกตาปิ มนฺตา นปฺปฏิภนฺติ, ปเคว อสชฺฌายกตา. โก ปน โภ โคตม เหตุ โก ปจฺจโย, เยเนกทา ทีฆรตฺตํ อสชฺฌายกตาปิ มนฺตา ปฏิภนฺติ, ปเคว สชฺฌายกตา”ติ.}

\prose{ยสฺมิํ โข พฺราหฺมณ สมเย กาม\-ราค\-ปริยุฏฺฐิเตน เจตสา วิหรติ กามราค\-ปเรเตน, อุปฺปนฺนสฺส จ กามราคสฺส นิสฺสรณํ ยถาภูตํ นปฺปชานาติ. อตฺตตฺถมฺปิ ตสฺมิํ สมเย ยถาภูตํ น ชานาติ น ปสฺสติ, ปรตฺถมฺปิ ตสฺมิํ สมเย ยถาภูตํ น ชานาติ น ปสฺสติ, อุภยตฺถมฺปิ ตสฺมิํ สมเย ยถาภูตํ น ชานาติ น ปสฺสติ, ทีฆรตฺตํ สชฺฌายกตาปิ มนฺตา นปฺปฏิภนฺติ, ปเคว อสชฺฌายกตา. เสยฺยถาปิ พฺราหฺมณ อุทปตฺโต สํสฏฺโฐ ลาขาย วา หลิทฺทิยา วา นีลิยา วา มญฺชิฏฺฐาย วา ตตฺถ จกฺขุมา ปุริโส สกํ มุขนิมิตฺตํ ปจฺจเวกฺขมาโน ยถาภูตํ น ชาเนยฺย น ปสฺเสยฺย. เอวเมว โข พฺราหฺมณ ยสฺมิํ สมเย กาม\-ราค\-ปริยุฏฺฐิเตน เจตสา วิหรติ กามราค\-ปเรเตน, อุปฺปนฺนสฺส จ กามราคสฺส นิสฺสรณํ ยถาภูตํ นปฺปชานาติ. อตฺตตฺถมฺปิ ตสฺมิํ สมเย ยถาภูตํ น ชานาติ น ปสฺสติ, ปรตฺถมฺปิ ฯเปฯ อุภยตฺถมฺปิ ตสฺมิํ สมเย ยถาภูตํ น ชานาติ น ปสฺสติ, ทีฆรตฺตํ สชฺฌายกตาปิ มนฺตา นปฺปฏิภนฺติ, ปเคว อสชฺฌายกตา.}

\prose{ปุน จปรํ พฺราหฺมณ ยสฺมิํ สมเย พฺยาปาท\-ปริยุฏฺฐิเตน เจตสา วิหรติ พฺยาปาท\-ปเรเตน, อุปฺปนฺนสฺส จ พฺยาปาทสฺส นิสฺสรณํ ยถาภูตํ นปฺปชานาติ. อตฺตตฺถมฺปิ ตสฺมิํ สมเย ยถาภูตํ น ชานาติ น ปสฺสติ, ปรตฺถมฺปิ ฯเปฯ อุภยตฺถมฺปิ ตสฺมิํ สมเย ยถาภูตํ น ชานาติ น ปสฺสติ, ทีฆรตฺตํ สชฺฌายกตาปิ มนฺตา นปฺปฏิภนฺติ, ปเคว อสชฺฌายกตา. เสยฺยถาปิ พฺรหฺมณ อุทปตฺโต อคฺคินา สนฺตตฺโต ปกฺกุถิโต\csromansharedfootnote{3dcb801d925d3a5a}{ปกฺกุธิโต (ก), อุกฺกฏฺฐิโต (สี), อุกฺกุฏฺฐิโต (สฺยา)}}

\prose{\csromanfolio{107}อุสฺมุทกชาโต\csromansharedfootnote{23ce4fc0ce6697c3}{อุสฺสทกชาโต (สี), อุสฺมาทกชาโต (สฺยา)}, ตตฺถ จกฺขุมา ปุริโส สกํ มุขนิมิตฺตํ ปจฺจเวกฺขมาโน ยถาภูตํ น ชาเนยฺย น ปสฺเสยฺย. เอวเมว โข พฺราหฺมณ ยสฺมิํ สมเย พฺยาปาท\-ปริยุฏฺฐิเตน เจตสา วิหรติ พฺยาปาท\-ปเรเตน. อุปฺปนฺนสฺส จ พฺยาปาทสฺส นิสฺสรณํ ยถาภูตํ นปฺปชานาติ. อตฺตตฺถมฺปิ ตสฺมิํ สมเย ยถาภูตํ น ชานาติ น ปสฺสติ, ปรตฺถมฺปิ ตสฺมิํ สมเย ฯเปฯ อุภยตฺถมฺปิ ตสฺมิํ สมเย ยถาภูตํ น ชานาติ น ปสฺสติ, ทีฆรตฺตํ สชฺฌายกตาปิ มนฺตา นปฺปฏิภนฺติ, ปเคว อสชฺฌายกตา.}

\prose{ปุน จปรํ พฺราหฺมณ ยสฺมิํ สมเย ถินมิทฺธ\-ปริยุฏฺฐิเตน เจตสา วิหรติ ถินมิทฺธ\-ปเรเตน, อุปฺปนฺนสฺส จ ถินมิทฺธสฺส นิสฺสรณํ ยถาภูตํ นปฺปชานาติ. อตฺตตฺถมฺปิ ตสฺมิํ สมเย ยถาภูตํ น ชานาติ น ปสฺสติ, ปรตฺถมฺปิ ฯเปฯ อุภยตฺถมฺปิ ตสฺมิํ สมเย ยถาภูตํ น ชานาติ น ปสฺสติ, ทีฆรตฺตํ สชฺฌายกตาปิ มนฺตา นปฺปฏิภนฺติ, ปเคว อสชฺฌายกตา.}

\prose{เสยฺยถาปิ พฺราหฺมณ อุทปตฺโต เสวาล\-ปณก\-ปริโยนทฺโธ, ตตฺถ จกฺขุมา ปุริโส สกํ มุขนิมิตฺตํ ปจฺจเวกฺขมาโน ยถาภูตํ น ชาเนยฺย น ปสฺเสยฺย. เอวเมว โข พฺราหฺมณ ยสฺมิํ สมเย ถินมิทฺธ\-ปริยุฏฺฐิเตน เจตสา วิหรติ ถินมิทฺธ\-ปเรเตน, อุปฺปนฺนสฺส จ ถินมิทฺธสฺส นิสฺสรณํ ยถาภูตํ นปฺปชานาติ. อตฺตตฺถมฺปิ ตสฺมิํ สมเย ยถาภูตํ น ชานาติ น ปสฺสติ, ปรตฺถมฺปิ ฯเปฯ อุภยตฺถมฺปิ ตสฺมิํ สมเย ยถาภูตํ น ชานาติ น ปสฺสติ, ทีฆรตฺตํ สชฺฌายกตาปิ มนฺตา นปฺปฏิภนฺติ, ปเคว อสชฺฌายกตา.}

\prose{ปุน จปรํ พฺราหฺมณ ยสฺมิํ สมเย อุทฺธจฺจ\-กุกฺกุจฺจ\-ปริยุฏฺฐิเตน เจตสา วิหรติ อุทฺธจฺจ\-กุกฺกุจฺจ\-ปเรเตน, อุปฺปนฺนสฺส จ อุทฺธจฺจ\-กุกฺกุจฺจสฺส นิสฺสรณํ ยถาภูตํ นปฺปชานาติ. อตฺตตฺถมฺปิ ตสฺมิํ สมเย ยถาภูตํ น ชานาติ น ปสฺสติ, ปรตฺถมฺปิ ฯเปฯ อุภยตฺถมฺปิ ตสฺมิํ สเมเย ยถาภูตํ น ชานาติ น ปสฺสติ, ทีฆรตฺตํ สชฺฌายกตาปิ มนฺตา นปฺปฏิภนฺติ, ปเคว อสชฺฌายกตา.}

\prose{เสยฺยถาปิ พฺราหฺมณ อุทปตฺโต วาเตริโต จลิโต ภนฺโต อูมิชาโต, ตตฺถ จกฺขุมา ปุริโส สกํ มุขนิมิตฺตํ ปจฺจเวกฺขมาโน}

\prose{\csromanfolio{108}ยถาภูตํ น ชาเนยฺย น ปสฺเสยฺย. เอวเมว โข พฺราหฺมณ ยสฺมิํ สมเย อุทฺธจฺจ\-กุกฺกุจฺจ\-ปริยุฏฺฐิเตน เจตสา วิหรติ อุทฺธจฺจ\-กุกฺกุจฺจ\-ปเรเตน, อุปฺปนฺนสฺส จ อุทฺธจฺจ\-กุกฺกุจฺจสฺส นิสฺสรณํ ยถาภูตํ นปฺปชานาติ. อตฺตตฺถมฺปิ ตสฺมิํ สมเย ยถาภูตํ น ชานาติ น ปสฺสติ, ปรตฺถมฺปิ ฯเปฯ อุภยตฺถมฺปิ ตสฺมิํ สมเย ยถาภูตํ น ชานาติ น ปสฺสติ, ทีฆรตฺตํ สชฺฌายกตาปิ มนฺตา นปฺปฏิภนฺติ, ปเคว อสชฺฌายกตา.}

\prose{ปุน จปรํ พฺราหฺมณ ยสฺมิํ สมเย วิจิกิจฺฉา\-ปริยุฏฺฐิเตน เจตสา วิหรติ วิจิกิจฺฉา\-ปเรเตน, อุปฺปนฺนาย จ วิจิกิจฺฉาย นิสฺสรณํ ยถาภูตํ นปฺปชานาติ. อตฺตตฺถมฺปิ ตสฺมิํ สมเย ยถาภูตํ น ชานาติ น ปสฺสติ, ปรตฺถมฺปิ ฯเปฯ อุภยตฺถมฺปิ ฯเปฯ ทีฆรตฺตํ สชฺฌายกตาปิ มนฺตา นปฺปฏิภนฺติ, ปเคว อสชฺฌยกตา.}

\prose{เสยฺยถาปิ พฺราหฺมณ อุทปตฺโต อาวิโล ลุฬิโต กลลีภูโต อนฺธกาเร นิกฺขิตฺโต, ตตฺถ จกฺขุมา ปุริโส สกํ มุขนิมิตฺตํ ปจฺจเวกฺขมาโน ยถาภูตํ น ชาเนยฺย น ปสฺเสยฺย. เอวเมว โข พฺราหฺมณ ยสฺมิํ สมเย วิจิกิจฺฉา\-ปริยุฏฺฐิเตน เจตสา วิหรติ วิจิกิจฺฉา\-ปเรเตน, อุปฺปนฺนาย จ วิจิกิจฺฉาย นิสฺสรณํ ยถาภูตํ นปฺปชานาติ. อตฺตตฺถมฺปิ ตสฺมิํ สมเย ยถาภูตํ น ชานาติ น ปสฺสติ, ปรตฺถมฺปิ ตสฺมิํ สมเย ยถาภูตํ น ชานาติ น ปสฺสติ, อุภยตฺถมฺปิ ตสฺมิํ สมเย ยถาภูตํ น ชานาติ น ปสฺสติ, ทีฆรตฺตํ สชฺฌายกตาปิ มนฺตา นปฺปฏิภนฺติ, ปเคน อสชฺฌายกตา. อยํ โข พฺราหฺมณ เหตุ อยํ ปจฺจโย, เยเนกทา ทีฆรตฺตํ สชฺฌายกตาปิ มนฺตา นปฺปฏิภนฺติ, ปเคว อสชฺฌายกตา.}

\prose{ยสฺมิํ จ โข พฺราหฺมณ สมเย น กาม\-ราค\-ปริยุฏฺฐิเตน เจตสา วิหรติ น กามราค\-ปเรเตน, อุปฺปนฺนสฺส จ กามราคสฺส นิสฺสรณํ ยถาภูตํ ปชานาติ. อตฺตตฺถมฺปิ ตสฺมิํ สมเย ยถาภูตํ ชานาติ ปสฺสติ, ปรตฺถมฺปิ ตสฺมิํ สมเย ยถาภูตํ ชานาติ ปสฺสติ, อุภยตฺถมฺปิ ตสฺมิํ สมเย ยถาภูตํ ชานาติ ปสฺสติ, ทีฆรตฺตํ อสชฺฌายกตาปิ มนฺตา ปฏิภนฺติ, ปเคว สชฺฌายกตา.}

\prose{\csromanfolio{109}เสยฺยถาปิ พฺราหฺมณ อุทปตฺโต อสํสฏฺโฐ ลาขาย วา หลิทฺทิยา วา นีลิยา วา มญฺชิฏฺฐาย วา, ตตฺถ จกฺขุมา ปุริโส สกํ มุขนิมิตฺตํ ปจฺจเวกฺขมาโน ยถาภูตํ ชาเนยฺย ปสฺเสยฺย. เอวเมว โข พฺราหฺมณ ยสฺมิํ สมเย น กาม\-ราค\-ปริยุฏฺฐิเตน เจตสา วิหรติ น กามราค\-ปเรเตน, อุปฺปนฺนสฺส จ กามราคสฺส นิสฺสรณํ ยถาภูตํ ปชานาติ ฯเปฯ.}

\prose{ปุน จปรํ พฺราหฺมณ ยสฺมิํ สมเย น พฺยาปาท\-ปริยุฏฺฐิเตน เจตสา วิหรติ น พฺยาปาท\-ปเรเตน, อุปฺปนฺนสฺส จ พฺยาปาทสฺส นิสฺสรณํ ยถาภูตํ ปชานาติ. อตฺตตฺถมฺปิ ตสฺมิํ สมเย ยถาภูตํ ชานาติ ปสฺสติ, ปรตฺถมฺปิ ฯเปฯ อุภยตฺถมฺปิ ฯเปฯ ทีฆรตฺตํ อสชฺฌายกตาปิ มนฺตา ปฏิภนฺติ, ปเคว สชฺฌายกตา.}

\prose{เสยฺยถาปิ พฺราหฺมณ อุทปตฺโต น อคฺคินา สนฺตตฺโต น ปกฺกุถิโต น อุสฺมุทกชาโต, ตตฺถ จกฺขุมา ปุริโส สกํ มุขนิมิตฺตํ ปจฺจเวกฺขมาโน ยถาภูตํ ชาเนยฺย ปสฺเสยฺย. เอวเมว โข พฺราหฺมณ ยสฺมิํ สมเย น พฺยาปาท\-ปริยุฏฺฐิเตน เจตสา วิหรติ น พฺยาปาท\-ปเรเตน, อุปฺปนฺนสฺส จ พฺยาปาทสฺส นิสฺสรณํ ยถาภูตํ ปชานาติ. อตฺตตฺถมฺปิ ตสฺมิํ สมเย ยถาภูตํ ชานาติ ปสฺสติ, ปรตฺถมฺปิ ฯเปฯ อุภยตฺถมฺปิ ฯเปฯ ทีฆรตฺตํ อสชฺฌายกตาปิ มนฺตา ปฏิภนฺติ, ปเคว สชฺฌายกตา.}

\prose{ปุน จปรํ พฺราหฺมณ ยสฺมิํ สมเย น ถินมิทฺธ\-ปริยุฏฺฐิเตน เจตสา วิหรติ น ถินมิทฺธ\-ปเรเตน, อุปฺปนฺนสฺส จ ถินมิทฺธสฺส นิสฺสรณํ ยถาภูตํ ปชานาติ. อตฺตตฺถมฺปิ ตสฺมิํ สมเย ยถาภูตํ ชานาติ ปสฺสติ, ปรตฺถมฺปิ ฯเปฯ อุภยตฺถมฺปิ ฯเปฯ ทีฆรตฺตํ อสชฺฌายกตาปิ มนฺตา ปฏิภนฺติ, ปเคว สชฺฌายกตา.}

\prose{เสยฺยถาปิ พฺราหฺมณ อุทปตฺโต น เสวาล\-ปณก\-ปริโยนทฺโธ, ตตฺถ จกฺขุมา ปุริโส สกํ มุขนิมิตฺตํ ปจฺจเวกฺขมาโน ยถาภูตํ ชาเนยฺย ปสฺเสยฺย. เอวเมว โข พฺราหฺมณ ยสฺมิํ สมเย น ถินมิทฺธ\-ปริยุฏฺฐิเตน เจตสา วิหรติ น ถินมิทฺธ\-ปเรเตน, อุปฺปนฺนสฺส จ ถินมิทฺธสฺส นิสฺสรณํ ยถาภูตํ ปชานาติ. อตฺตตฺถมฺปิ ตสฺมิํ สมเย ยถาภูตํ ชานาติ ปสฺสติ, ปรตฺถมฺปิ ฯเปฯ อุภยตฺถมฺปิ ฯเปฯ ทีฆรตฺตํ อสชฺฌายกตาปิ มนฺตา ปฏิภนฺติ, ปเคว สชฺฌายกตา.}

\prose{\csromanfolio{110}ปุน จปรํ พฺราหฺมณ ยสฺมิํ สมเย น อุทฺธจฺจ\-กุกฺกุจฺจ\-ปริยุฏฺฐิเตน เจตสา วิหรติ น อุทฺธจฺจ\-กุกฺกุจฺจ\-ปเรเตน, อุปฺปนฺนสฺส จ อุทฺธจฺจ\-กุกฺกุจฺจสฺส นิสฺสรณํ ยถาภูตํ ปชานาติ. อตฺตตฺถมฺปิ ตสฺมิํ สมเย ยถาภูตํ ชานาติ ปสฺสติ, ปรตฺถมฺปิ ฯเปฯ อุภยตฺถมฺปิ ฯเปฯ ทีฆรตฺตํ อสชฺฌายกตาปิ มนฺตา ปฏิภนฺติ, ปเคว สชฺฌายกตา.}

\prose{เสยฺยถาปิ พฺราหฺมณ อุทปตฺโต น วาเตริโต น จลิโต น ภนฺโต น อูมิชาโต, ตตฺถ จกฺขุมา ปุริโส สกํ มุขนิมิตฺตํ ปจฺจเวกฺขมาโน ยถาภูตํ ชาเนยฺย ปสฺเสยฺย. เอวเมว โข พฺราหฺมณ ยสฺมิํ สมเย น อุทฺธจฺจ\-กุกฺกุจฺจ\-ปริยุฏฺฐิเตน เจตสา วิหรติ น อุทฺธจฺจ\-กุกฺกุจฺจ\-ปเรเตน, อุปฺปนฺนสฺส จ อุทฺธจฺจ\-กุกฺกุจฺจสฺส นิสฺสรณํ ยถาภูตํ ปชานาติ. อตฺตตฺถมฺปิ ตสฺมิํ สมเย ยถาภูตํ ชานาติ ปสฺสติ, ปรตฺถมฺปิ ฯเปฯ อุภยตฺถมฺปิ ฯเปฯ ทีฆรตฺตํ อสชฺฌายกตาปิ มนฺตา ปฏิภนฺติ, ปเคว สชฺฌายกตา.}

\prose{ปุน จปรํ พฺราหฺมณ ยสฺมิํ สมเย น วิจิกิจฺฉา\-ปริยุฏฺฐิเตน เจตสา วิหรติ น วิจิกิจฺฉา\-ปเรเตน, อุปฺปนฺนาย จ วิจิกิจฺฉาย นิสฺสรณํ ยถาภูตํ ปชานาติ\csromansharedfootnote{d98d6df0feaf4024}{ปชานาติ ปสฺสติ (สฺยา)}. อตฺตตฺถมฺปิ ตสฺมิํ สมเย ยถาภูตํ ชานาติ ปสฺสติ, ปรตฺถมฺปิ ตสฺมิํ สมเย ยถาภูตํ ชานาติ ปสฺสติ, อุภยตฺถมฺปิ ตสฺมิํ สมเย ยถาภูตํ ชานาติ ปสฺสติ, ทีฆรตฺตํ อสชฺฌายกตาปิ มนฺตา ปฏิภนฺติ, ปเคว สชฺฌายกตา.}

\prose{เสยฺยถาปิ พฺราหฺมณ อุทปตฺโต อจฺโฉ วิปฺปสนฺโน อนาวิโล อาโลเก นิกฺขิตฺโต, ตตฺถ จกฺขุมา ปุริโส สกํ มุขนิมิตฺตํ ปจฺจเวกฺขมาโน ยถาภูตํ ชาเนยฺย ปสฺเสยฺย. เอวเมว โข พฺราหฺมณ ยสฺมิํ สมเย น วิจิกิจฺฉา\-ปริยุฏฺฐิเตน เจตสา วิหรติ น วิจิกิจฺฉา\-ปเรเตน, อุปฺปนฺนาย จ วิจิกิจฺฉาย นิสฺสรณํ ยถาภูตํ ปชานาติ, อตฺตตฺถมฺปิ ตสฺมิํ สมเย ยถาภูตํ ชานาติ ปสฺสติ, ปรตฺถมฺปิ ตสฺมิํ สมเย ยถาภูตํ ชานาติ ปสฺสติ, อุภยตฺถมฺปิ ตสฺมิํ สมเย ยถาภูตํ ชานาติ ปสฺสติ, ทีฆรตฺตํ อสชฺฌายกตาปิ มนฺตา ปฏิภนฺติ, ปเคว สชฺฌายกตา. อยํ โข พฺราหฺมณ เหตุ อยํ ปจฺจโย, เยเนกทา ทีฆรตฺตํ อสชฺฌายกตาปิ มนฺตา ปฏิภนฺติ, ปเคว สชฺฌายกตา.}

\prose{\csromanfolio{111}สตฺติเม พฺราหฺมณ โพชฺฌงฺคา อนาวรณา อนีวรณา เจตโส อนุปกฺกิเลสา ภาวิตา พหุลีกตา วิชฺชา\-วิมุตฺติ\-ผล\-สจฺฉิกิริยาย สํวตฺตนฺติ. กตเม สตฺต, สติสมฺโพชฺฌงฺโค โข พฺราหฺมณ อนาวรโณ อนีวรโณ เจตโส อนุปกฺกิเลโส ภาวิโต พหุลีกโต วิชฺชา\-วิมุตฺติ\-ผล\-สจฺฉิกิริยาย สํวตฺตติ ฯเปฯ อุเปกฺขา\-สมฺโพชฺฌงฺโค โข พฺราหฺมณ อนาวรโณ อนีวรโณ เจตโส อนุปกฺกิเลโส ภาวิโต พหุลีกโต วิชฺชาวิมุตฺติผลสจฺฉิ กิริยาย สํวตฺตติ. อิเม โข พฺราหฺมณ สตฺต โพชฺฌงฺคา อนาวรณา อนีวรณา เจตโส อนุปกฺกิเลสา ภาวิตา พหุลีกตา วิชฺชา\-วิมุตฺติ\-ผล\-สจฺฉิกิริยาย สํวตฺตนฺตีติ. เอวํ วุตฺเต สงฺคารโว พฺราหฺมโณ ภควนฺตํ เอตทโวจ “อภิกฺกนฺตํ โภ โคตม ฯเปฯ อุปาสกํ มํ ภวํ โคตโม ธาเรตุ อชฺชตคฺเค ปาณุเปตํ สรณํ คตนฺ”ติ. ปญฺจมํ.}

\csromanmatikaanchor{mtk.18}
\csromantocmarkref{paragraph}{6. อภยสุตฺต}{mtk.18}
\bookmark[dest={mtk.18},level=4]{6. อภยสุตฺต}
\csromanheader{4}{6. \textbf{อภยสุตฺต}}
\proseitem{237}{เอวํ เม สุตํ\csromandash{}เอกํ สมยํ ภควา ราชคเห วิหรติ คิชฺฌกูเฏ ปพฺพเต. อถ โข อภโย ราชกุมาโร เยน ภควา เตนุปสงฺกมิ, อุปสงฺกมิตฺวา ภควนฺตํ อภิวาเทตฺวา เอกมนฺตํ นิสีทิ, เอกมนฺตํ นิสินฺโน โข อภโย ราชกุมาโร ภควนฺตํ เอตทโวจ “ปูรโณ ภนฺเต กสฺสโป เอวมาห นตฺถิ เหตุ นตฺถิ ปจฺจโย อญฺญาณาย อทสฺสนาย, อเหตุ อปฺปจฺจโย\csromansharedfootnote{f277cb41dc503f29}{อปฺปจฺจยา (สี), อปฺปจฺจยํ (?)} อญฺญาณํ อทสฺสนํ โหติ. นตฺถิ เหตุ นตฺถิ ปจฺจโย ญาณาย ทสฺสนาย, อเหตุ อปฺปจฺจโย ญาณํ ทสฺสนํ โหตี”ติ. อิธ ภควา กิมาหาติ. อตฺถิ ราชกุมาร เหตุ อตฺถิ ปจฺจโย อญฺญาณาย อทสฺสนาย, สเหตุ สปฺปจฺจโย\csromansharedfootnote{c07a2928f8386d1d}{สปฺปจฺจยา (สี), สปฺปจฺจยํ (?)} อญฺญาณํ อทสฺสนํ โหติ. อตฺถิ ราชกุมาร เหตุ อตฺถิ ปจฺจโย ญาณาย ทสฺสนาย, สเหตุ สปฺปจฺจโย ญาณํ ทสฺสนํ โหตีติ.}

\prose{กตโม ปน ภนฺเต เหตุ กตโม ปจฺจโย อญฺญาณาย อทสฺสนาย, กถํ สเหตุ สปฺปจฺจโย อญฺญาณํ อทสฺสนํ โหตีติ. ยสฺมิํ โข ราชกุมาร สมเย กาม\-ราค\-ปริยุฏฺฐิเตน เจตสา วิหรติ กามราค\-ปเรเตน, อุปฺปนฺนสฺส จ กามราคสฺส นิสฺสรณํ ยถาภูตํ น ชานาติ น ปสฺสติ. อยมฺปิ โข ราชกุมาร เหตุ \csromanfolio{112}อยํ ปจฺจโย อญฺญาณาย อทสฺสนาย, เอวมฺปิ สเหตุ สปฺปจฺจโย อญฺญาณํ อทสฺสนํ โหติ.}

\prose{ปุน จปรํ ราชกุมาร ยสฺมิํ สมเย พฺยาปาท\-ปริยุฏฺฐิเตน เจตสา วิหรติ พฺยาปาท\-ปเรเตน ฯเปฯ ถินมิทฺธ\-ปริยุฏฺฐิเตน ฯเปฯ อุทฺธจฺจ\-กุกฺกุจฺจ\-ปริยุฏฺฐิเตน ฯเปฯ วิจิกิจฺฉา\-ปริยุฏฺฐิเตน เจตสา วิหรติ วิจิกิจฺฉา\-ปเรเตน, อุปฺปนฺนาย จ วิจิกิจฺฉาย นิสฺสรณํ ยถาภูตํ น ชานาติ น ปสฺสติ. อยมฺปิ โข ราชกุมาร เหตุ อยํ ปจฺจโย อญฺญาณาย อทสฺสนาย, เอวมฺปิ สเหตุ สปฺปจฺจโย อญฺญาณํ อทสฺสนํ โหตีติ.}

\prose{โก นามายํ ภนฺเต ธมฺม\-ปริยาโยติ. นีวรณา นาเมเต ราชกุมาราติ. ตคฺฆ ภควา นีวรณา, ตคฺฆ สุคต นีวรณา, เอกเมเกนปิ โข ภนฺเต นีวรเณน อภิภูโต ยถาภูตํ น ชาเนยฺย น ปสฺเสยฺย, โก ปน วาโท ปญฺจหิ นีวรเณหิ.}

\prose{กตโม ปน ภนฺเต เหตุ กตโม ปจฺจโย ญาณาย ทสฺสนาย, กถํ สเหตุ สปฺปจฺจโย ญาณํ ทสฺสนํ โหตีติ. อิธ ราชกุมาร ภิกฺขุ สติสมฺโพชฺฌงฺคํ ภาเวติ วิเวกนิสฺสิตํ วิราคนิสฺสิตํ นิโรธ\-นิสฺสิตํ โวสฺสคฺค\-ปริณามิํ, โส สติสมฺโพชฺฌงฺคํ ภาวิเตน จิตฺเตน ยถาภูตํ ชานาติ ปสฺสติ. อยมฺปิ โข ราชกุมาร เหตุ อยํ ปจฺจโย ญาณาย ทสฺสนาย, เอวมฺปิ สเหตุ สปฺปจฺจโย ญาณํ ทสฺสนํ โหติ.}

\prose{ปุน จปรํ ราชกุมาร ภิกฺขุ ฯเปฯ อุเปกฺขา\-สมฺโพชฺฌงฺคํ ภาเวติ วิเวกนิสฺสิตํ วิราคนิสฺสิตํ นิโรธ\-นิสฺสิตํ โวสฺสคฺค\-ปริณามิํ, โส อุเปกฺขา\-สมฺโพชฺฌงฺคํ ภาวิเตน จิตฺเตน ยถาภูตํ ชานาติ ปสฺสติ. อยมฺปิ โข ราชกุมาร เหตุ อยํ ปจฺจโย ญาณาย ทสฺสนาย, เอวํ สเหตุ สปฺปจฺจโย ญาณํ ทสฺสนํ โหตีติ.}

\prose{โก นามายํ ภนฺเต ธมฺม\-ปริยาโยติ. โพชฺฌงฺคา นาเมเต ราชกุมาราติ. ตคฺฆ ภควา โพชฺฌงฺคา, ตคฺฆ สุคต โพชฺฌงฺคา, เอกเมเกนปิ โข ภนฺเต โพชฺฌงฺเคน สมนฺนาคโต ยถาภูตํ ชาเนยฺย ปสฺเสยฺย, โก ปน วาโท สตฺตหิ โพชฺฌงฺเคหิ. โยปิ เม ภนฺเต คิชฺฌกูฏํ ปพฺพตํ}

\csromanmatikaanchor{mtk.19}
\csromanmatikaanchor{mtk.20}
\csromantocmarknopage{paragraph}{7. อานาปานวคฺค}{mtk.19}
\bookmark[dest={mtk.19},level=4]{7. อานาปานวคฺค}
\csromantocmarkref{subparagraph}{1-10. อฏฺฐิกมหปฺผลสุตฺตาทีนิ}{mtk.20}
\bookmark[dest={mtk.20},level=5]{1-10. อฏฺฐิกมหปฺผลสุตฺตาทีนิ}
\prose{\csromanfolio{113}อาโรหนฺตสฺส กายกิลมโถ จิตฺตกิลมโถ, โสปิ เม ปฏิปฺปสฺสทฺโธ, ธมฺโม จ เม อภิสมิโตติ.  ฉฏฺฐํ. สากจฺฉวคฺโค ฉฏฺโฐ.}

\vspace{\tipitakaparskip}
\csromancenter{\textbf{ตสฺสุทฺทานํ}}
\setlength{\csromangathapagewidth}{0pt}
\setlength{\csromangathawakwidth}{0pt}
\csromangathameasuregroup{อาหารา ปริยายมคฺคิ, \\ อภโย ปุจฺฉิโต ปญฺหํ,}{\csromangathabat{อาหารา ปริยายมคฺคิ,}{เมตฺตํ สงฺคารเวน จ.} \\ \csromangathabat{อภโย ปุจฺฉิโต ปญฺหํ,}{คิชฺฌกูฏมฺหิ ปพฺพเตติ.}}
\csromangathasetleft{อาหารา ปริยายมคฺคิ, \\ อภโย ปุจฺฉิโต ปญฺหํ,}
\csromangathagroup{\csromangathastanza{\csromangathabat{อาหารา ปริยายมคฺคิ,}{เมตฺตํ สงฺคารเวน จ.} \\ \csromangathabat{อภโย ปุจฺฉิโต ปญฺหํ,}{คิชฺฌกูฏมฺหิ ปพฺพเตติ.}}}

\csromanheader{4}{7. \textbf{อานาปานวคฺค}}
\csromanheader{2}{1. อฏฺฐิก\-มหปฺผล\-สุตฺต}
\proseitem{238}{สาวตฺถิ\-นิทานํ. อฏฺฐิกสญฺญา ภิกฺขเว ภาวิตา พหุลีกตา มหปฺผลา โหติ มหานิสํสา. กถํ ภาวิตา จ ภิกฺขเว อฏฺฐิกสญฺญา กถํ พหุลีกตา มหปฺผลา โหติ มหานิสํสา, อิธ ภิกฺขเว ภิกฺขุ อฏฺฐิกสญฺญา\-สหคตํ สติสมฺโพชฺฌงฺคํ ภาเวติ วิเวกนิสฺสิตํ วิราคนิสฺสิตํ นิโรธ\-นิสฺสิตํ โวสฺสคฺค\-ปริณามิํ ฯเปฯ อฏฺฐิกสญฺญา\-สหคตํ อุเปกฺขา\-สมฺโพชฺฌงฺคํ ภาเวติ วิเวกนิสฺสิตํ วิราคนิสฺสิตํ นิโรธ\-นิสฺสิตํ โวสฺสคฺค\-ปริณามิํ. เอวํ ภาวิตา โข ภิกฺขเว อฏฺฐิกสญฺญา เอวํ พหุลีกตา มหปฺผลา โหติ มหานิสํสาติ.}

\csromanheader{2}{อญฺญตร\-ผล\-สุตฺต}
\prose{อฏฺฐิกสญฺญาย ภิกฺขเว ภาวิตาย พหุลีกตาย ทฺวินฺนํ ผลานํ อญฺญตรํ ผลํ ปาฏิกงฺขํ, ทิฏฺเฐว ธมฺเม อญฺญา, สติ วา อุปาทิเสเส อนาคามิตา. กถํ ภาวิตาย จ โข ภิกฺขเว อฏฺฐิกสญฺญาย กถํ พหุลีกตาย ทฺวินฺนํ ผลานํ อญฺญตรํ ผลํ ปาฏิกงฺขํ ทิฏฺเฐว ธมฺเม อญฺญา, สติ วา อุปาทิเสเส อนาคามิตา, อิธ ภิกฺขเว ภิกฺขุ อฏฺฐิกสญฺญา\-สหคตํ สติสมฺโพชฺฌงฺคํ ภาเวติ ฯเปฯ อฏฺฐิกสญฺญา\-สหคตํ อุเปกฺขา\-สมฺโพชฺฌงฺคํ ภาเวติ วิเวกนิสฺสิตํ วิราคนิสฺสิตํ นิโรธ\-นิสฺสิตํ โวสฺสคฺค\-ปริณามิํ. เอวํ ภาวิตาย โข ภิกฺขเว อฏฺฐิกสญฺญาย เอวํ}

\prose{\csromanfolio{114}พหุลีกตาย ทฺวินฺนํ ผลานํ อญฺญตรํ ผลํ ปาฏิกงฺขํ ทิฏฺเฐว ธมฺเม อญฺญา, สติ วา อุปาทิเสเส อนาคามิตาติ.}

\csromanheader{2}{\textbf{มหตฺถสุตฺต}}
\prose{อฏฺฐิกสญฺญา ภิกฺขเว ภาวิตา พหุลีกตา มหโต อตฺถาย สํวตฺตติ. กถํ ภาวิตา จ ภิกฺขเว อฏฺฐิกสญฺญา กถํ พหุลีกตา มหโต อตฺถาย สํวตฺตติ, อิธ ภิกฺขเว ภิกฺขุ อฏฺฐิกสญฺญา\-สหคตํ สติสมฺโพชฺฌงฺคํ ภาเวติ ฯเปฯ อฏฺฐิกสญฺญา\-สหคตํ อุเปกฺขา\-สมฺโพชฺฌงฺคํ ภาเวติ วิเวกนิสฺสิตํ วิราคนิสฺสิตํ นิโรธ\-นิสฺสิตํ โวสฺสคฺค\-ปริณามิํ. เอวํ ภาวิตา โข ภิกฺขเว อฏฺฐิกสญฺญา เอวํ พหุลีกตา มหโต อตฺถาย สํวตฺตตีติ.}

\csromanheader{2}{โยคกฺเขม\-สุตฺต}
\prose{อฏฺฐิกสญฺญา ภิกฺขเว ภาวิตา พหุลีกตา มหโต โยคกฺเขมาย สํวตฺตติ. กถํ ภาวิตา จ ภิกฺขเว อฏฺฐิกสญฺญา กถํ พหุลีกตา มหโต โยคกฺเขมาย สํวตฺตติ, อิธ ภิกฺขเว ภิกฺขุ อฏฺฐิกสญฺญา\-สหคตํ สติสมฺโพชฺฌงฺคํ ภาเวติ ฯเปฯ อฏฺฐิกสญฺญา\-สหคตํ อุเปกฺขา\-สมฺโพชฺฌงฺคํ ภาเวติ วิเวกนิสฺสิตํ วิราคนิสฺสิตํ นิโรธ\-นิสฺสิตํ โวสฺสคฺค\-ปริณามิํ. เอวํ ภาวิตา โข ภิกฺขเว อฏฺฐิกสญฺญา เอวํ พหุลีกตา มหโต โยคกฺเขมาย สํวตฺตตีติ.}

\csromanheader{2}{\textbf{สํเวคสุตฺต}}
\prose{อฏฺฐิกสญฺญา ภิกฺขเว ภาวิตา พหุลีกตา มหโต สํเวคาย สํวตฺตติ. กถํ ภาวิตา จ ภิกฺขเว อฏฺฐิกสญฺญา กถํ พหุลีกตา มหโต สํเวคาย สํวตฺตติ, อิธ ภิกฺขเว ภิกฺขุ อฏฺฐิกสญฺญา\-สหคตํ สติสมฺโพชฺฌงฺคํ ภาเวติ ฯเปฯ อฏฺฐิกสญฺญา\-สหคตํ อุเปกฺขา\-สมฺโพชฺฌงฺคํ ภาเวติ วิเวกนิสฺสิตํ วิราคนิสฺสิตํ นิโรธ\-นิสฺสิตํ โวสฺสคฺค\-ปริณามิํ. เอวํ ภาวิตา โข ภิกฺขเว อฏฺฐิกสญฺญา เอวํ พหุลีกตา มหโต สํเวคาย สํวตฺตตีติ.}

\csromanheader{2}{ผาสุวิหาร\-สุตฺต}
\prose{\csromanfolio{115}อฏฺฐิกสญฺญา ภิกฺขเว ภาวิตา พหุลีกตา มหโต ผาสุวิหาราย สํวตฺตติ. กถํ ภาวิตา จ ภิกฺขเว อฏฺฐิกสญฺญา กถํ พหุลีกตา มหโต ผาสุวิหาราย สํวตฺตติ, อิธ ภิกฺขเว ภิกฺขุ อฏฺฐิกสญฺญา\-สหคตํ สติสมฺโพชฺฌงฺคํ ภาเวติ ฯเปฯ อฏฺฐิกสญฺญา\-สหคตํ อุเปกฺขา\-สมฺโพชฺฌงฺคํ ภาเวติ วิเวกนิสฺสิตํ วิราคนิสฺสิตํ นิโรธ\-นิสฺสิตํ โวสฺสคฺค\-ปริณามิํ. เอวํ ภาวิตา โข ภิกฺขเว อฏฺฐิกสญฺญา เอวํ พหุลีกตา มหโต ผาสุวิหาราย สํวตฺตตีติ.  ปฐมํ.}

\csromanheader{2}{2. \textbf{ปุฬวกสุตฺต}}
\vspace{\tipitakaparskip}
\csromancenter{ปุฬวกสญฺญา\csromansharedfootnote{41bd4becab01424a}{ปุฬุวกสญฺญา (ก)} ภิกฺขเว ภาวิตา ฯเปฯ.  ทุติยํ.}
\csromanheader{2}{3. \textbf{วินีลกสุตฺต}}
\vspace{\tipitakaparskip}
\csromancenter{วินีลกสญฺญา ภิกฺขเว ฯเปฯ. ตติยํ.}
\csromanheader{2}{4. วิจฺฉิทฺทก\-สุตฺต}
\vspace{\tipitakaparskip}
\csromancenter{วิจฺฉิทฺทก\-สญฺญา ภิกฺขเว ฯเปฯ.  จตุตฺถํ.}
\csromanheader{2}{5. อุทฺธุมาตก\-สุตฺต}
\vspace{\tipitakaparskip}
\csromancenter{อุทฺธุมาตก\-สญฺญา ภิกฺขเว ฯเปฯ.  ปญฺจมํ.}
\csromanheader{2}{6. \textbf{เมตฺตาสุตฺต}}
\vspace{\tipitakaparskip}
\csromancenter{เมตฺตา ภิกฺขเว ภาวิตา ฯเปฯ.  ฉฏฺฐํ.}
\csromanheader{2}{7. \textbf{กรุณาสุตฺต}}
\vspace{\tipitakaparskip}
\csromancenter{กรุณา ภิกฺขเว ภาวิตา ฯเปฯ.  สตฺตมํ.}
\csromanheader{2}{8. \textbf{มุทิตาสุตฺต}}
\vspace{\tipitakaparskip}
\csromancenter{มุทิตา ภิกฺขเว ภาวิตา ฯเปฯ.  อฏฺฐมํ.}
\csromanheader{2}{9. \textbf{อุเปกฺขาสุตฺต}}
\vspace{\tipitakaparskip}
\csromancenter{อุเปกฺขา ภิกฺขเว ภาวิตา ฯเปฯ.  นวมํ.}
\csromanheader{2}{10. \textbf{อานาปานสุตฺต}}
\nitthitammidruled{247. อานาปานสฺสติ ภิกฺขเว ภาวิตา ฯเปฯ. . ทสมํ. อานาปานวคฺโค สตฺตโม.}
\csromancenter{\textbf{ตสฺสุทฺทานํ}}
\csromanmatikaanchor{mtk.21}
\csromanmatikaanchor{mtk.22}
\csromantocmarknopage{paragraph}{8. นิโรธวคฺค}{mtk.21}
\bookmark[dest={mtk.21},level=4]{8. นิโรธวคฺค}
\csromantocmarkref{subparagraph}{1-10. อสุภสุตฺตาทีนิ}{mtk.22}
\bookmark[dest={mtk.22},level=5]{1-10. อสุภสุตฺตาทีนิ}
\csromantocmarknopage{paragraph}{9. คงฺคาเปยฺยาลวคฺค}{mtk.23}
\bookmark[dest={mtk.23},level=4]{9. คงฺคาเปยฺยาลวคฺค}
\setlength{\csromangathapagewidth}{0pt}
\setlength{\csromangathawakwidth}{0pt}
\csromangathameasuregroup{อฏฺฐิก ปุฬวกํ วินีลกํ, \\ เมตฺตา กรุณา มุทิตา,}{\csromangathabat{อฏฺฐิก ปุฬวกํ วินีลกํ,}{วิจฺฉิทฺทกํ อุทฺธุมาเตน ปญฺจมํ.} \\ \csromangathabat{เมตฺตา กรุณา มุทิตา,}{อุเปกฺขา อานาปาเนน เต ทสาติ.}}
\csromangathasetleft{อฏฺฐิก ปุฬวกํ วินีลกํ, \\ เมตฺตา กรุณา มุทิตา,}
\csromangathagroup{\csromangathastanza{\csromanfolio{116}\csromangathabat{อฏฺฐิก ปุฬวกํ วินีลกํ,}{วิจฺฉิทฺทกํ อุทฺธุมาเตน ปญฺจมํ.} \\ \csromangathabat{เมตฺตา กรุณา มุทิตา,}{อุเปกฺขา อานาปาเนน เต ทสาติ.}}}

\csromanheader{4}{8. \textbf{นิโรธวคฺค}}
\csromanheader{2}{1. \textbf{อสุภสุตฺต}}
\vspace{\tipitakaparskip}
\csromancenter{อสุภสญฺญา ภิกฺขเว ฯเปฯ.  ปฐมํ.}
\csromanheader{2}{2. \textbf{มรณสุตฺต}}
\vspace{\tipitakaparskip}
\csromancenter{มรณสญฺญา ภิกฺขเว ฯเปฯ.  ทุติยํ.}
\csromanheader{2}{3. อาหาเรปฏิกูล\-สุตฺต}
\vspace{\tipitakaparskip}
\csromancenter{อาหาเร ปฏิกูลสญฺญา ภิกฺขเว ฯเปฯ. ตติยํ.}
\csromanheader{2}{4. อนภิรติ\-สุตฺต}
\vspace{\tipitakaparskip}
\csromancenter{สพฺพโลเก อนภิรติ\-สญฺญา ภิกฺขเว ฯเปฯ.  จตุตฺถํ.}
\csromanheader{2}{5. \textbf{อนิจฺจสุตฺต}}
\vspace{\tipitakaparskip}
\csromancenter{อนิจฺจสญฺญา ภิกฺขเว ฯเปฯ.  ปญฺจมํ.}
\csromanheader{2}{6. \textbf{ทุกฺขสุตฺต}}
\vspace{\tipitakaparskip}
\csromancenter{อนิจฺเจ ทุกฺขสญฺญา ภิกฺขเว ฯเปฯ.  ฉฏฺฐํ.}
\csromanheader{2}{7. \textbf{อนตฺตสุตฺต}}
\vspace{\tipitakaparskip}
\csromancenter{ทุกฺเข อนตฺตสญฺญา ภิกฺขเว ฯเปฯ.  สตฺตมํ.}
\csromanheader{2}{8. \textbf{ปหานสุตฺต}}
\vspace{\tipitakaparskip}
\csromancenter{ปหานสญฺญา ภิกฺขเว ฯเปฯ.  อฏฺฐมํ.}
\csromanheader{2}{9. \textbf{วิราคสุตฺต}}
\vspace{\tipitakaparskip}
\csromancenter{วิราคสญฺญา ภิกฺขเว ฯเปฯ.  นวมํ.}
\csromanheader{2}{10. \textbf{นิโรธสุตฺต}}
\proseitem{257}{\csromanfolio{117}นิโรธสญฺญา ภิกฺขเว ภาวิตา พหุลีกตา มหปฺผลา โหติ มหานิสํสา. กถํ ภาวิตา จ ภิกฺขเว นิโรธสญฺญา กถํ พหุลีกตา มหปฺผลา โหติ มหานิสํสา, อิธ ภิกฺขเว ภิกฺขุ นิโรธสญฺญา\-สหคตํ สติสมฺโพชฺฌงฺคํ ภาเวติ ฯเปฯ นิโรธสญฺญา\-สหคตํ อุเปกฺขา\-สมฺโพชฺฌงฺคํ ภาเวติ วิเวกนิสฺสิตํ วิราคนิสฺสิตํ นิโรธ\-นิสฺสิตํ โวสฺสคฺค\-ปริณามิํ. เอวํ ภาวิตา โข ภิกฺขเว นิโรธสญฺญา เอวํ พหุลีกตา มหปฺผลา โหติ มหานิสํสาติ.}

\prose{นิโรธสญฺญาย ภิกฺขเว ภาวิตาย พหุลีกตาย ทฺวินฺนํ ผลานํ อญฺญตรํ ผลํ ปาฏิกงฺขํ ทิฏฺเฐว ธมฺเม อญฺญา, สติ วา อุปาทิเสเส อนาคามิตา. กถํ ภาวิตาย ภิกฺขเว นิโรธสญฺญาย กถํ พหุลีกตาย ทฺวินฺนํ ผลานํ อญฺญตรํ ผลํ ปาฏิกงฺขํ ทิฏฺเฐว ธมฺเม อญฺญา, สติ วา อุปาทิเสเส อนาคามิตา, อิธ ภิกฺขเว ภิกฺขุ นิโรธสญฺญา\-สหคตํ สติสมฺโพชฺฌงฺคํ ภาเวติ ฯเปฯ นิโรธสญฺญา\-สหคตํ อุเปกฺขา\-สมฺโพชฺฌงฺคํ ภาเวติ วิเวกนิสฺสิตํ วิราคนิสฺสิตํ นิโรธ\-นิสฺสิตํ โวสฺสคฺค\-ปริณามิํ. เอวํ ภาวิตาย โข ภิกฺขเว นิโรธสญฺญาย เอวํ พหุลีกตาย ทฺวินฺนํ ผลานํ อญฺญตรํ ผลํ ปาฏิกงฺขํ ทิฏฺเฐว ธมฺเม อญฺญา, สติ วา อุปาทิเสเส อนาคามิตาติ.}

\csromanmatikaanchor{mtk.24}
\csromantocmarkref{subparagraph}{1-12. คงฺคานทีอาทิสุตฺต}{mtk.24}
\bookmark[dest={mtk.24},level=5]{1-12. คงฺคานทีอาทิสุตฺต}
\prose{นิโรธสญฺญา ภิกฺขเว ภาวิตา พหุลีกตา มหโต อตฺถาย สํวตฺตติ, มหโต โยคกฺเขมาย สํวตฺตติ, มหโต สํเวคาย สํวตฺตติ, มหโต ผาสุวิหาราย สํวตฺตติ. กถํ ภาวิตา จ ภิกฺขเว นิโรธสญฺญา กถํ พหุลีกตา มหโต อตฺถาย สํวตฺตติ, มหโต \csromanfolio{118}โยคกฺเขมาย สํวตฺตติ, มหโต สํเวคาย สํวตฺตติ, มหโต ผาสุวิหาราย สํวตฺตติ. อิธ ภิกฺขเว ภิกฺขุ นิโรธสญฺญา\-สหคตํ สติสมฺโพชฺฌงฺคํ ภาเวติ ฯเปฯ นิโรธสญฺญา\-สหคตํ อุเปกฺขา\-สมฺโพชฺฌงฺคํ ภาเวติ วิเวกนิสฺสิตํ วิราคนิสฺสิตํ นิโรธ\-นิสฺสิตํ โวสฺสคฺค\-ปริณามิํ. เอวํ ภาวิตา โข ภิกฺขเว นิโรธสญฺญา เอวํ พหุลีกตา มหโต อตฺถาย สํวตฺตติ, มหโต โยคกฺเขมาย สํวตฺตติ, มหโต สํเวคาย สํวตฺตติ, มหโต ผาสุวิหาราย สํวตฺตตีติ.  ทสมํ. นิโรธวคฺโค อฏฺฐโม.}

\vspace{\tipitakaparskip}
\csromancenter{\textbf{ตสฺสุทฺทานํ}}
\setlength{\csromangathapagewidth}{0pt}
\setlength{\csromangathawakwidth}{0pt}
\csromangathameasuregroup{อสุภมรณอาหาเร, \\ อนิจฺจทุกฺขอนตฺตปหานํ,}{\csromangathabat{อสุภมรณอาหาเร,}{ปฏิกูลอนภิรเตน\textsuperscript{1}.} \\ \csromangathabat{อนิจฺจทุกฺขอนตฺตปหานํ,}{วิราคนิโรเธน เต ทสาติ.}}
\csromangathasetleft{อสุภมรณอาหาเร, \\ อนิจฺจทุกฺขอนตฺตปหานํ,}
\csromangathagroup{\csromangathastanza{\csromangathabat{อสุภมรณอาหาเร,}{ปฏิกูลอนภิรเตน\csromansharedfootnote{338dd9d64fd7ee1e}{[มิสฺสินฺคฺ โนเต 0]}.} \\ \csromangathabat{อนิจฺจทุกฺขอนตฺตปหานํ,}{วิราคนิโรเธน เต ทสาติ.}}}

\chapterhead{9. คงฺคา\-เปยฺยาล\-วคฺค}
\prose{1-12. คงฺคานที\-อาทิ\-สุตฺต}

\proseitemwithclosermidruled{258-269}{เสยฺยถาปิ ภิกฺขเว คงฺคา นที ปาจีนนินฺนา ปาจีนโปณา ปาจีนปพฺภารา. เอวเมว โข ภิกฺขเว ภิกฺขุ สตฺต โพชฺฌงฺเค ภาเวนฺโต สตฺต โพชฺฌงฺเค พหุลีกโรนฺโต นิพฺพานนินฺโน โหติ นิพฺพานโปโณ นิพฺพาน\-ปพฺภาโร. กถญฺจ ภิกฺขเว ภิกฺขุ สตฺต โพชฺฌงฺเค ภาเวนฺโต สตฺต โพชฺฌงฺเค พหุลีกโรนฺโต นิพฺพานนินฺโน โหติ นิพฺพานโปโณ นิพฺพาน\-ปพฺภาโร. อิธ ภิกฺขเว ภิกฺขุ สติสมฺโพชฺฌงฺคํ ภาเวติ ฯเปฯ อุเปกฺขา\-สมฺโพชฺฌงฺคํ ภาเวติ วิเวกนิสฺสิตํ วิราคนิสฺสิตํ นิโรธ\-นิสฺสิตํ โวสฺสคฺค\-ปริณามิํ. เอวํ โข ภิกฺขเว ภิกฺขุ สตฺต โพชฺฌงฺเค ภาเวนฺโต สตฺต โพชฺฌงฺเค พหุลีกโรนฺโต นิพฺพานนินฺโน โหติ นิพฺพานโปโณ นิพฺพาน\-ปพฺภาโรติ.}{(ยาว เอสนา ปาฬิ วิตฺถาเรตพฺพา.) คงฺคา\-เปยฺยาล\-วคฺโค นวโม.}
\csromancenter{\textbf{ตสฺสุทฺทานํ}}
\setlength{\csromangathapagewidth}{0pt}
\setlength{\csromangathawakwidth}{0pt}
\csromangathameasuregroup{ฉ ปาจีนโต นินฺนา, \\ เทฺว เต ฉ ทฺวาทส โหนฺติ,}{\csromangathabat{ฉ ปาจีนโต นินฺนา,}{ฉ นินฺนา จ สมุทฺทโต.} \\ \csromangathabat{เทฺว เต ฉ ทฺวาทส โหนฺติ,}{วคฺโค เตน ปวุจฺจตีติ.}}
\csromangathasetleft{ฉ ปาจีนโต นินฺนา, \\ เทฺว เต ฉ ทฺวาทส โหนฺติ,}
\csromangathagroup{\csromangathastanza{\csromanfolio{119}\csromangathabat{ฉ ปาจีนโต นินฺนา,}{ฉ นินฺนา จ สมุทฺทโต.} \\ \csromangathabat{เทฺว เต ฉ ทฺวาทส โหนฺติ,}{วคฺโค เตน ปวุจฺจตีติ.}}}

\csromanmatikaanchor{mtk.25}
\csromantocmarknopage{paragraph}{10. อปฺปมาทวคฺค}{mtk.25}
\bookmark[dest={mtk.25},level=4]{10. อปฺปมาทวคฺค}
\csromanheader{4}{10. \textbf{อปฺปมาทวคฺค}}
\csromanmatikaanchor{mtk.26}
\csromantocmarkref{subparagraph}{1-10. ตถาคตาทิสุตฺต}{mtk.26}
\bookmark[dest={mtk.26},level=5]{1-10. ตถาคตาทิสุตฺต}
\csromantocmarknopage{paragraph}{11. พลกรณียวคฺค}{mtk.27}
\bookmark[dest={mtk.27},level=4]{11. พลกรณียวคฺค}
\csromanheader{5}{1-10. ตถาคตา\-ทิ\-สุตฺต}
\proseitemwithclosermidruled{270-279}{ยาวตา ภิกฺขเว สตฺตา อปทา วา ทฺวิปทา วา จตุปฺปทา วา พหุปฺปทา วาติ วิตฺถาเรตพฺพํ.}{อปฺปมาทวคฺโค ทสโม.}
\csromancenter{\textbf{ตสฺสุทฺทานํ}}
\setlength{\csromangathapagewidth}{0pt}
\setlength{\csromangathawakwidth}{0pt}
\csromangathameasuregroup{ตถาคตํ ปทํ กูฏํ, \\ ราชา จนฺทิมสูริยา จ,}{\csromangathabat{ตถาคตํ ปทํ กูฏํ,}{มูลํ สาเรน วสฺสิกํ.} \\ \csromangathabat{ราชา จนฺทิมสูริยา จ,}{วตฺเถน ทสมํ ปทนฺติ.}}
\csromangathasetleft{ตถาคตํ ปทํ กูฏํ, \\ ราชา จนฺทิมสูริยา จ,}
\csromangathagroup{\csromangathastanza{\csromangathabat{ตถาคตํ ปทํ กูฏํ,}{มูลํ สาเรน วสฺสิกํ.} \\ \csromangathabat{ราชา จนฺทิมสูริยา จ,}{วตฺเถน ทสมํ ปทนฺติ.}}}

\vspace{\tipitakaparskip}
\csromancenter{(อปฺปมาทวคฺโค โพชฺฌงฺค\-สํยุตฺตสฺส โพชฺฌงฺควเสน วิตฺถาเรตพฺโพ.)}
\csromansectionrule
\chapterhead{11. พลกรณีย\-วคฺค}
\csromanmatikaanchor{mtk.28}
\csromantocmarkref{subparagraph}{1-12. พลาทิสุตฺต}{mtk.28}
\bookmark[dest={mtk.28},level=5]{1-12. พลาทิสุตฺต}
\csromantocmarknopage{paragraph}{12. เอสนาวคฺค}{mtk.29}
\bookmark[dest={mtk.29},level=4]{12. เอสนาวคฺค}
\csromanheader{5}{1-12. พลาทิสุตฺต}
\vspace{\tipitakaparskip}
\csromancenter{เสยฺยถาปิ ภิกฺขเว เย เกจิ พลกรณียา กมฺมนฺตา กรียนฺตีติ วิตฺถาเรตพฺพํ.}
\nitthitammidruled{พลกรณีย\-วคฺโค เอกาทสโม.}
\csromancenter{\textbf{ตสฺสุทฺทานํ}}
\setlength{\csromangathapagewidth}{0pt}
\setlength{\csromangathawakwidth}{0pt}
\csromangathameasuregroup{พลํ พีชญฺจ นาโค จ, \\ อากาเสน จ เทฺว เมฆา,}{\csromangathabat{พลํ พีชญฺจ นาโค จ,}{รุกฺโข กุมฺเภน สูกิยา.} \\ \csromangathabat{อากาเสน จ เทฺว เมฆา,}{นาวา อาคนฺตุกา นทีติ.}}
\csromangathasetleft{พลํ พีชญฺจ นาโค จ, \\ อากาเสน จ เทฺว เมฆา,}
\csromangathagroup{\csromangathastanza{\csromangathabat{พลํ พีชญฺจ นาโค จ,}{รุกฺโข กุมฺเภน สูกิยา.} \\ \csromangathabat{อากาเสน จ เทฺว เมฆา,}{นาวา อาคนฺตุกา นทีติ.}}}

\vspace{\tipitakaparskip}
\csromancenter{(พลกรณีย\-วคฺโค โพชฺฌงฺค\-สํยุตฺตสฺส โพชฺฌงฺควเสน วิตฺถาเรตพฺโพ.)}
\csromanmatikaanchor{mtk.30}
\csromanmatikaanchor{mtk.40}
\csromantocmarkref{subparagraph}{1-10. เอสนาทิสุตฺต}{mtk.30}
\bookmark[dest={mtk.30},level=5]{1-10. เอสนาทิสุตฺต}
\csromantocmarknopage{paragraph}{13. โอฆวคฺค}{mtk.31}
\bookmark[dest={mtk.31},level=4]{13. โอฆวคฺค}
\csromanheader{5}{12. \textbf{เอสนาวคฺค}}
\csromanheader{5}{1-10. เอสนาทิสุตฺต}
\proseitemwithclosermidruled{292-301}{\csromanfolio{120}ติสฺโส อิมา ภิกฺขเว เอสนา. กตมา ติสฺโส, กาเมสนา ภเวสนา พฺรหฺมจริเย\-สนาติ วิตฺถาเรตพฺพํ.}{เอสนาวคฺโค ทฺวาทสโม.}
\csromancenter{\textbf{ตสฺสุทฺทานํ}}
\setlength{\csromangathapagewidth}{0pt}
\setlength{\csromangathawakwidth}{0pt}
\csromangathameasuregroup{เอสนา วิธา อาสโว, \\ ขิลํ มลญฺจ นีโฆ จ,}{\csromangathabat{เอสนา วิธา อาสโว,}{ภโว จ ทุกฺขตา ติสฺโส.} \\ \csromangathabat{ขิลํ มลญฺจ นีโฆ จ,}{เวทนา ตณฺหา ตสินาย จาติ.}}
\csromangathasetleft{เอสนา วิธา อาสโว, \\ ขิลํ มลญฺจ นีโฆ จ,}
\csromangathagroup{\csromangathastanza{\csromangathabat{เอสนา วิธา อาสโว,}{ภโว จ ทุกฺขตา ติสฺโส.} \\ \csromangathabat{ขิลํ มลญฺจ นีโฆ จ,}{เวทนา ตณฺหา ตสินาย จาติ.}}}

\vspace{\tipitakaparskip}
\csromancenter{(โพชฺฌงฺค\-สํยุตฺตสฺส เอสนาเปยฺยาลํ วิเวกนิสฺสิตโต วิตฺถาเรตพฺพํ.)}
\csromansectionrule
\csromanmatikaanchor{mtk.32}
\csromanmatikaanchor{mtk.42}
\csromantocmarkref{subparagraph}{1-9. โอฆาทิสุตฺต}{mtk.32}
\bookmark[dest={mtk.32},level=5]{1-9. โอฆาทิสุตฺต}
\csromanheader{5}{13. \textbf{โอฆวคฺค}}
\csromanheader{5}{1-9. โอฆาทิสุตฺต}
\proseitem{302-310}{จตฺตาโรเม ภิกฺขเว โอฆา. กตเม จตฺตาโร, กาโมโฆ ภโวโฆ ทิฏฺโฐโฆ อวิชฺโชโฆติ วิตฺถาเรตพฺพํ.}

\csromanmatikaanchor{mtk.33}
\csromantocmarkref{paragraph}{10. อุทฺธมฺภาคิยสุตฺต}{mtk.33}
\bookmark[dest={mtk.33},level=4]{10. อุทฺธมฺภาคิยสุตฺต}
\csromanheader{4}{10. อุทฺธมฺภาคิย\-สุตฺต}
\proseitem{311}{สาวตฺถิ\-นิทานํ. ปญฺจิมานิ ภิกฺขเว อุทฺธมฺ\-ภาคิยานิ สํโยชนานิ. กตมานิ ปญฺจ, รูปราโค อรูปราโค มาโน อุทฺธจฺจํ อวิชฺชา. อิมานิ โข ภิกฺขเว ปญฺจุ\-ทฺธมฺภาคิยานิ สํโยชนานิ. อิเมสํ โข ภิกฺขเว ปญฺจนฺนํ อุทฺธมฺ\-ภาคิยานํ สํโยชนานํ อภิญฺญาย ปริญฺญาย ปริกฺขยาย ปหานาย สตฺต โพชฺฌงฺคา ภาเวตพฺพา. กตเม สตฺต, อิธ ภิกฺขเว ภิกฺขุ สติสมฺโพชฺฌงฺคํ ภาเวติ วิเวกนิสฺสิตํ วิราคนิสฺสิตํ นิโรธ\-นิสฺสิตํ โวสฺสคฺค\-ปริณามิํ ฯเปฯ อุเปกฺขา\-สมฺโพชฺฌงฺคํ ภาเวติ ราค\-วินย\-ปริโยสานํ โทส\-วินย\-ปริโยสานํ โมห\-วินย\-ปริโยสานํ. อมโตคธํ อมตปรายนํ อมต\-ปริโยสานํ. นิพฺพานนินฺนํ นิพฺพานโปณํ นิพฺพาน\-ปพฺภารํ. อิเมสํ โข ภิกฺขเว ภิกฺขุ ปญฺจนฺนํ อุทฺธมฺ\-ภาคิยานํ สํโยชนานํ อภิญฺญาย ปริญฺญาย ปริกฺขยาย ปหานาย อิเม สตฺต โพชฺฌงฺคา ภาเวตพฺพาติ.  ทสมํ. โอฆวคฺโค เตรสโม.}

\vspace{\tipitakaparskip}
\csromancenter{\textbf{ตสฺสุทฺทานํ}}
\csromanmatikaanchor{mtk.34}
\csromanmatikaanchor{mtk.35}
\csromanmatikaanchor{mtk.37}
\csromanmatikaanchor{mtk.39}
\csromantocmarknopage{subparagraph}{14. ปุนคงฺคาเปยฺยาลวคฺค}{mtk.34}
\bookmark[dest={mtk.34},level=5]{14. ปุนคงฺคาเปยฺยาลวคฺค}
\csromantocmarkref{subparagraph}{312-323. ปุนคงฺคานทีอาทิสุตฺต}{mtk.35}
\bookmark[dest={mtk.35},level=5]{312-323. ปุนคงฺคานทีอาทิสุตฺต}
\setlength{\csromangathapagewidth}{0pt}
\setlength{\csromangathawakwidth}{0pt}
\csromangathameasuregroup{โอโฆ โยโค อุปาทานํ, \\ กามคุณา นีวรณา,}{\csromangathabat{โอโฆ โยโค อุปาทานํ,}{คนฺถา อนุสเยน จ.} \\ \csromangathabat{กามคุณา นีวรณา,}{ขนฺธา โอรุทฺธมฺภาคิยานีติ.}}
\csromangathasetleft{โอโฆ โยโค อุปาทานํ, \\ กามคุณา นีวรณา,}
\csromangathagroup{\csromangathastanza{\csromanfolio{121}\csromangathabat{โอโฆ โยโค อุปาทานํ,}{คนฺถา อนุสเยน จ.} \\ \csromangathabat{กามคุณา นีวรณา,}{ขนฺธา โอรุทฺธมฺภาคิยานีติ.}}}

\csromanheader{5}{14. ปุนคงฺคาเปยฺยาล\-วคฺค}
\nitthitammidruled{ปุนคงฺคานที\-อาทิ\-สุตฺต วคฺโค จุทฺทสโม.}
\csromanheader{2}{\textbf{อุทฺทานํ}}
\setlength{\csromangathapagewidth}{0pt}
\setlength{\csromangathawakwidth}{0pt}
\csromangathameasuregroup{ฉ ปาจีนโต นินฺนา, \\ เทฺว เต ฉ ทฺวาทส โหนฺติ,}{\csromangathabat{ฉ ปาจีนโต นินฺนา,}{ฉ นินฺนา จ สมุทฺทโต.} \\ \csromangathabat{เทฺว เต ฉ ทฺวาทส โหนฺติ,}{วคฺโค เตน ปวุจฺจตีติ.}}
\csromangathasetleft{ฉ ปาจีนโต นินฺนา, \\ เทฺว เต ฉ ทฺวาทส โหนฺติ,}
\csromangathagroup{\csromangathastanza{\csromangathabat{ฉ ปาจีนโต นินฺนา,}{ฉ นินฺนา จ สมุทฺทโต.} \\ \csromangathabat{เทฺว เต ฉ ทฺวาทส โหนฺติ,}{วคฺโค เตน ปวุจฺจตีติ.}}}

\vspace{\tipitakaparskip}
\csromancenter{(โพชฺฌงฺค\-สํยุตฺตสฺส คงฺคาเปยฺยาลํ ราควเสน วิตฺถาเรตพฺพํ.)}
\csromansectionrule
\chapterhead{15. ปุนอปฺปมาท\-วคฺค}
\vspace{\tipitakaparskip}
\csromancenter{ตถาคตา\-ทิ\-สุตฺต ปนฺนรสโม.}
\csromanheader{2}{\textbf{อุทฺทานํ}}
\setlength{\csromangathapagewidth}{0pt}
\setlength{\csromangathawakwidth}{0pt}
\csromangathameasuregroup{ตถาคตํ ปทํ กูฏํ,}{\csromangathabat{ตถาคตํ ปทํ กูฏํ,}{มูลํ สาเรน วสฺสิกํ.}}
\csromangathasetleft{ตถาคตํ ปทํ กูฏํ,}
\csromangathagroup{\csromangathastanza{\csromangathabat{ตถาคตํ ปทํ กูฏํ,}{มูลํ สาเรน วสฺสิกํ.}}}

\vspace{\tipitakaparskip}
\csromancenter{ราชา จนฺทิมสูริยา จ, วตฺเถน ทสมํ ปทนฺติ. \textbf{(อปฺปมาทวคฺโค ราควเสน วิตฺถาเรตพฺโพ.)}}
\csromansectionrule
\csromanmatikaanchor{mtk.36}
\csromanmatikaanchor{mtk.38}
\csromantocmarknopage{subparagraph}{15. ปุนอปฺปมาทวคฺค}{mtk.36}
\bookmark[dest={mtk.36},level=5]{15. ปุนอปฺปมาทวคฺค}
\csromantocmarkref{subparagraph}{324-333. ตถาคตาทิสุตฺต}{mtk.37}
\bookmark[dest={mtk.37},level=5]{324-333. ตถาคตาทิสุตฺต}
\csromantocmarknopage{subparagraph}{16. ปุนพลกรณียวคฺค}{mtk.38}
\bookmark[dest={mtk.38},level=5]{16. ปุนพลกรณียวคฺค}
\csromantocmarkref{subparagraph}{334-356. ปุนพลาทิสุตฺต}{mtk.39}
\bookmark[dest={mtk.39},level=5]{334-356. ปุนพลาทิสุตฺต}
\csromantocmarknopage{subparagraph}{17. ปุนเอสนาวคฺค}{mtk.40}
\bookmark[dest={mtk.40},level=5]{17. ปุนเอสนาวคฺค}
\csromanheader{5}{16. ปุนพลกรณีย\-วคฺค}
\vspace{\tipitakaparskip}
\csromancenter{ปุนพลา\-ทิ\-สุตฺต โสฬสโม.}
\csromanheader{2}{\textbf{อุทฺทานํ}}
\csromanmatikaanchor{mtk.41}
\csromanmatikaanchor{mtk.43}
\csromantocmarkref{subparagraph}{346-356. ปุนเอสนาทิสุตฺต}{mtk.41}
\bookmark[dest={mtk.41},level=5]{346-356. ปุนเอสนาทิสุตฺต}
\csromantocmarknopage{subparagraph}{18. ปุนโอฆวคฺค}{mtk.42}
\bookmark[dest={mtk.42},level=5]{18. ปุนโอฆวคฺค}
\csromantocmarkref{subparagraph}{357-366. ปุนโอฆาทิสุตฺต}{mtk.43}
\bookmark[dest={mtk.43},level=5]{357-366. ปุนโอฆาทิสุตฺต}
\setlength{\csromangathapagewidth}{0pt}
\setlength{\csromangathawakwidth}{0pt}
\csromangathameasuregroup{พลํ พีชญฺจ นาโค จ, \\ อากาเสน จ เทฺว เมฆา,}{\csromangathabat{พลํ พีชญฺจ นาโค จ,}{รุกฺโข กุมฺเภน สูกิยา.} \\ \csromangathabat{อากาเสน จ เทฺว เมฆา,}{นาวา อาคนฺตุกา นทีติ.}}
\csromangathasetleft{พลํ พีชญฺจ นาโค จ, \\ อากาเสน จ เทฺว เมฆา,}
\csromangathagroup{\csromangathastanza{\csromanfolio{122}\csromangathabat{พลํ พีชญฺจ นาโค จ,}{รุกฺโข กุมฺเภน สูกิยา.} \\ \csromangathabat{อากาเสน จ เทฺว เมฆา,}{นาวา อาคนฺตุกา นทีติ.}}}

\vspace{\tipitakaparskip}
\csromancenter{(โพชฺฌงฺค\-สํยุตฺตสฺส พลกรณีย\-วคฺโค ราควเสน วิตฺถาเรตพฺโพ.)}
\csromansectionrule
\chapterhead{17. ปุนเอสนาวคฺค}
\proseitem{346-356}{ปุนเอสนาทิสุตฺต ปุนเอสนาวคฺโค สตฺตรสโม.}

\csromanheader{2}{\textbf{อุทฺทานํ}}
\setlength{\csromangathapagewidth}{0pt}
\setlength{\csromangathawakwidth}{0pt}
\csromangathameasuregroup{เอสนา วิธา อาสโว, \\ ขิลํ มลญฺจ นีโฆ จ,}{\csromangathabat{เอสนา วิธา อาสโว,}{ภโว จ ทุกฺขตา ติสฺโส.} \\ \csromangathabat{ขิลํ มลญฺจ นีโฆ จ,}{เวทนา ตณฺหา ตสินาย จาติ.}}
\csromangathasetleft{เอสนา วิธา อาสโว, \\ ขิลํ มลญฺจ นีโฆ จ,}
\csromangathagroup{\csromangathastanza{\csromangathabat{เอสนา วิธา อาสโว,}{ภโว จ ทุกฺขตา ติสฺโส.} \\ \csromangathabat{ขิลํ มลญฺจ นีโฆ จ,}{เวทนา ตณฺหา ตสินาย จาติ.}}}

\chapterhead{18. ปุนโอฆวคฺค}
\prose{357-366. ปุนโอฆาทิสุตฺต}

\prose{โพชฺฌงฺค\-สํยุตฺตสฺส ปุนโอฆวคฺโค อฏฺฐารสโม.}

\csromanheader{2}{\textbf{อุทฺทานํ}}
\setlength{\csromangathapagewidth}{0pt}
\setlength{\csromangathawakwidth}{0pt}
\csromangathameasuregroup{โอโฆ โยโค อุปาทานํ, \\ กามคุณา นีวรณา,}{\csromangathabat{โอโฆ โยโค อุปาทานํ,}{คนฺถา อนุสเยน จ.} \\ \csromangathabat{กามคุณา นีวรณา,}{ขนฺธา โอรุทฺธมฺภาคิยานีติ.}}
\csromangathasetleft{โอโฆ โยโค อุปาทานํ, \\ กามคุณา นีวรณา,}
\csromangathagroup{\csromangathastanza{\csromangathabat{โอโฆ โยโค อุปาทานํ,}{คนฺถา อนุสเยน จ.} \\ \csromangathabat{กามคุณา นีวรณา,}{ขนฺธา โอรุทฺธมฺภาคิยานีติ.}}}

\vspace{\tipitakaparskip}
\csromancenter{(ราควินยปริโยสานโทสวินยปริโยสานโมหวินยปริโยสาน- วคฺโค วิตฺถาเรตพฺโพ.) (ยทปิ มคฺคสํยุตฺตํ วิตฺถาเรตพฺพํ, ตทปิ โพชฺฌงฺค\-สํยุตฺตํ วิตฺถาเรตพฺพํ.)}
\csromancenter{โพชฺฌงฺค\-สํยุตฺตํ ทุติยํ.}
\csromanmatikaanchor{mtk.44}
\csromanmatikaanchor{mtk.45}
\csromantocmarknopage{subsection}{3. สติปฏฺฐานสํยุตฺต}{mtk.44}
\bookmark[dest={mtk.44},level=2]{3. สติปฏฺฐานสํยุตฺต}
\csromantocmarknopage{paragraph}{1. อมฺพปาลิวคฺค}{mtk.45}
\bookmark[dest={mtk.45},level=4]{1. อมฺพปาลิวคฺค}
\csromanheader{4}{1. \textbf{อมฺพปาลิวคฺค}}
\csromanmatikaanchor{mtk.46}
\csromantocmarkref{subparagraph}{1. อมฺพปาลิสุตฺต}{mtk.46}
\bookmark[dest={mtk.46},level=5]{1. อมฺพปาลิสุตฺต}
\csromanheader{5}{1. \textbf{อมฺพปาลิสุตฺต}}
\proseitem{367}{\csromanfolio{123}เอวํ เม สุตํ\csromandash{}เอกํ สมยํ ภควา เวสาลิยํ วิหรติ อมฺพปาลิวเน. ตตฺร โข ภควา ภิกฺขู อามนฺเตสิ “ภิกฺขโว”ติ. “ภทนฺเต”ติ เต ภิกฺขู ภควโต ปจฺจสฺโสสุํ. ภควา เอตทโวจ\csromandash{}}

\prose{เอกายนฺวายํ ภิกฺขเว มคฺโค สตฺตานํ วิสุทฺธิยา โสก\-ปริเทวานํ สมติกฺกมาย ทุกฺข\-โทมนสฺสานํ อตฺถงฺคมาย ญายสฺส อธิคมาย นิพฺพานสฺส สจฺฉิกิริยาย, ยทิทํ จตฺตาโร สติปฏฺฐานา. กตเม จตฺตาโร, อิธ ภิกฺขเว ภิกฺขุ กาเย กายานุปสฺสี วิหรติ อาตาปี สมฺปชาโน สติมา วิเนยฺย โลเก อภิชฺฌา\-โทมนสฺสํ, เวทนาสุ เวทนานุปสฺสี วิหรติ อาตาปี สมฺปชาโน สติมา วิเนยฺย โลเก อภิชฺฌา\-โทมนสฺสํ, จิตฺเต จิตฺตานุปสฺสี วิหรติ อาตาปี สมฺปชาโน สติมา วิเนยฺย โลเก อภิชฺฌา\-โทมนสฺสํ, ธมฺเมสุ ธมฺมานุปสฺสี วิหรติ อาตาปี สมฺปชาโน สติมา วิเนยฺย โลเก อภิชฺฌา\-โทมนสฺสํ. เอกายนฺวายํ ภิกฺขเว มคฺโค สตฺตานํ วิสุทฺธิยา โสก\-ปริเทวานํ สมติกฺกมาย ทุกฺข\-โทมนสฺสานํ อตฺถงฺคมาย ญายสฺส อธิคมาย นิพฺพานสฺส สจฺฉิกิริยาย, ยทิทํ จตฺตาโร สติปฏฺฐานาติ.}

\prose{อิทมโวจ ภควา. อตฺตมนา เต ภิกฺขู ภควโต ภาสิตํ อภินนฺทุนฺติ.  ปฐมํ.}

\csromanmatikaanchor{mtk.47}
\csromantocmarkref{subparagraph}{2. สติสุตฺต}{mtk.47}
\bookmark[dest={mtk.47},level=5]{2. สติสุตฺต}
\csromanheader{5}{2. \textbf{สติสุตฺต}}
\proseitem{368}{เอกํ สมยํ ภควา เวสาลิยํ วิหรติ อมฺพปาลิวเน. ตตฺร โข ภควา ภิกฺขู อามนฺเตสิ “ภิกฺขโว”ติ. “ภทนฺเต”ติ เต ภิกฺขู ภควโต ปจฺจสฺโสสุํ. ภควา เอตทโวจ\csromandash{}}

\prose{สโต ภิกฺขเว ภิกฺขุ วิหเรยฺย สมฺปชาโน อยํ โว อมฺหากํ อนุสาสนี. กถญฺจ ภิกฺขเว ภิกฺขุ สโต โหติ, อิธ ภิกฺขเว ภิกฺขุ กาเย กายานุปสฺสี วิหรติ อาตาปี สมฺปชาโน สติมา \csromanfolio{124}วิเนยฺย โลเก อภิชฺฌา\-โทมนสฺสํ, เวทนาสุ ฯเปฯ จิตฺเต ฯเปฯ ธมฺเมสุ ธมฺมานุปสฺสี วิหรติ อาตาปี สมฺปชาโน สติมา วิเนยฺย โลเก อภิชฺฌา\-โทมนสฺสํ. เอวํ โข ภิกฺขเว ภิกฺขุ สโต โหติ.}

\prose{กถญฺจ ภิกฺขเว ภิกฺขุ สมฺปชาโน โหติ, อิธ ภิกฺขเว ภิกฺขุ อภิกฺกนฺเต ปฏิกฺกนฺเต สมฺปชานการี โหติ, อาโลกิเต วิโลกิเต สมฺปชานการี โหติ, สมิญฺชิเต ปสาริเต สมฺปชานการี โหติ, สงฺฆาฏิ\-ปตฺต\-จีวร\-ธารเณ สมฺปชานการี โหติ, อสิเต ปีเต ขายิเต สายิเต สมฺปชานการี โหติ, อุจฺจาร\-ปสฺสาว\-กมฺเม สมฺปชานการี โหติ, คเต ฐิเต นิสินฺเน สุตฺเต ชาคริเต ภาสิเต ตุณฺหีภาเว สมฺปชานการี โหติ. เอวํ โข ภิกฺขเว ภิกฺขุ สมฺปชานการี โหติ, สโต ภิกฺขเว ภิกฺขุ วิหเรยฺย สมฺปชาโน. อยํ โว อมฺหากํ อนุสาสนีติ.  ทุติยํ.}

\csromanmatikaanchor{mtk.48}
\csromantocmarkref{subparagraph}{3. ภิกฺขุสุตฺต}{mtk.48}
\bookmark[dest={mtk.48},level=5]{3. ภิกฺขุสุตฺต}
\csromanheader{5}{3. \textbf{ภิกฺขุสุตฺต}}
\proseitem{369}{เอกํ สมยํ ภควา สาวตฺถิยํ วิหรติ เชตวเน อนาถ\-ปิณฺฑิกสฺส อาราเม. อถ โข อญฺญตโร ภิกฺขุ เยน ภควา เตนุปสงฺกมิ, อุปสงฺกมิตฺวา ภควนฺตํ อภิวาเทตฺวา เอกมนฺตํ นิสีทิ, เอกมนฺตํ นิสินฺโน โข โส ภิกฺขุ ภควนฺตํ เอตทโวจ “สาธุ เม ภนฺเต ภควา สํขิตฺเตน ธมฺมํ เทเสตุ, ยมหํ ภควโต ธมฺมํ สุตฺวา เอโก วูปกฏฺโฐ อปฺปมตฺโต อาตาปี ปหิตตฺโต วิหเรยฺยนฺ”ติ, เอวเมว ปนิเธกจฺเจ โมฆปุริสา มญฺเจว\csromansharedfootnote{088b1786965fcd82}{มเมว (สี)} อชฺเฌสนฺติ, ธมฺเม จ ภาสิเต มเมว อนุพนฺธิตพฺพํ มญฺญนฺตีติ. เทเสตุ เม ภนฺเต ภควา สํขิตฺเตน ธมฺมํ, เทเสตุ สุคโต สํขิตฺเตน ธมฺมํ, อปฺเปว นามาหํ ภควโต ภาสิตสฺส อตฺถํ ชาเนยฺยํ. อปฺเปว นามาหํ ภควโต ภาสิตสฺส ทายาโท อสฺสนฺติ. ตสฺมาติห ตฺวํ ภิกฺขุ อาทิเมว วิโสเธหิ กุสเลสุ ธมฺเมสุ. โก จาทิ กุสลานํ ธมฺมานํ, สีลญฺจ สุวิสุทฺธํ ทิฏฺฐิ จ อุชุกา. ยโต โข เต ภิกฺขุ สีลญฺจ สุวิสุทฺธํ ภวิสฺสติ, ทิฏฺฐิ จ อุชุกา. ตโต ตฺวํ ภิกฺขุ สีลํ นิสฺสาย สีเล ปติฏฺฐาย จตฺตาโร สติปฏฺฐาเน ติวิเธน ภาเวยฺยาสิ.}

\prose{\csromanfolio{125}กตเม จตฺตาโร, อิธ ตฺวํ ภิกฺขุ อชฺฌตฺตํ วา กาเย กายานุปสฺสี วิหราหิ อาตาปี สมฺปชาโน สติมา วิเนยฺย โลเก อภิชฺฌา\-โทมนสฺสํ. พหิทฺธา วา กาเย กายานุปสฺสี วิหราหิ อาตาปี สมฺปชาโน สติมา วิเนยฺย โลเก อภิชฺฌา\-โทมนสฺสํ. อชฺฌตฺต\-พหิทฺธา วา กาเย กายานุปสฺสี วิหราหิ อาตาปี สมฺปชาโน สติมา วิเนยฺย โลเก อภิชฺฌา\-โทมนสฺสํ. อชฺฌตฺตํ วา เวทนาสุ. พหิทฺธา วา เวทนาสุ. อชฺฌตฺต\-พหิทฺธา วา เวทนาสุ เวทนานุปสฺสี วิหราหิ อาตาปี สมฺปชาโน สติมา วิเนยฺย โลเก อภิชฺฌา\-โทมนสฺสํ. อชฺฌตฺตํ วา จิตฺเต. พหิทฺธา วา จิตฺเต. อชฺฌตฺต\-พหิทฺธา วา จิตฺเต จิตฺตานุปสฺสี วิหราหิ อาตาปี สมฺปชาโน สติมา วิเนยฺย โลเก อภิชฺฌา\-โทมนสฺสํ. อชฺฌตฺตํ วา ธมฺเมสุ. พหิทฺธา วา ธมฺเมสุ. อชฺฌตฺต\-พหิทฺธา วา ธมฺเมสุ ธมฺมานุปสฺสี วิหราหิ อาตาปี สมฺปชาโน สติมา วิเนยฺย โลเก อภิชฺฌา\-โทมนสฺสํ. ยโต โข ตฺวํ ภิกฺขุ สีลํ นิสฺสาย สีเล ปติฏฺฐาย อิเม จตฺตาโร สติปฏฺฐาเน เอวํ ติวิเธน ภาเวสฺสสิ, ตโต ตุยฺหํ ภิกฺขุ ยา รตฺติ วา ทิวโส วา อาคมิสฺสติ, วุทฺธิเยว ปาฏิกงฺขา กุสเลสุ ธมฺเมสุ โน ปริหานีติ.}

\prose{อถ โข โส ภิกฺขเว ภควโต ภาสิตํ อภินนฺทิตฺวา อนุโมทิตฺวา อุฏฺฐายาสนา ภควนฺตํ อภิวาเทตฺวา ปทกฺขิณํ กตฺวา ปกฺกามิ. อถ โข โส ภิกฺขุ เอโก วูปกฏฺโฐ อปฺปมตฺโต อาตาปี ปหิตตฺโต วิหรนฺโต นจิรสฺเสว ยสฺสตฺถาย กุลปุตฺตา สมฺมเทว อคารสฺมา อนคาริยํ ปพฺพชนฺติ, ตทนุตฺตรํ พฺรหฺมจริย\-ปริโยสานํ ทิฏฺเฐว ธมฺเม สยํ อภิญฺญา สจฺฉิกตฺวา อุปสมฺปชฺช วิหรติ, “ขีณา ชาติ, วุสิตํ พฺรหฺมจริยํ, กตํ กรณียํ, นาปรํ อิตฺถตฺตายา”ติ อพฺภญฺญาสิ. อญฺญตโร จ ปน โส ภิกฺขุ อรหตํ อโหสีติ.  ตติยํ.}

\csromanmatikaanchor{mtk.49}
\csromantocmarkref{subparagraph}{4. สาลสุตฺต}{mtk.49}
\bookmark[dest={mtk.49},level=5]{4. สาลสุตฺต}
\csromanheader{5}{4. \textbf{สาลสุตฺต}}
\proseitem{370}{เอกํ สมยํ ภควา โกสเลสุ วิหรติ สาลาย พฺราหฺมณคาเม. ตตฺร โข ภควา ภิกฺขู อามนฺเตสิ ฯเปฯ เอตทโวจ\csromandash{}}

\prose{เย เต ภิกฺขเว ภิกฺขู นวา อจิรปพฺพชิตา อธุนาคตา อิมํ ธมฺมวินยํ, เต โว ภิกฺขเว ภิกฺขู จตุนฺนํ สติปฏฺฐานานํ ภาวนาย สมาทเปตพฺพา นิเวเสตพฺพา ปติฏฺฐาเป\-ตพฺพา. กตเมสํ จตุนฺนํ, เอถ ตุเมฺห อาวุโส \csromanfolio{126}กาเย กายานุปสฺสิโน วิหรถ อาตาปิโน สมฺปชานา เอโกทิภูตา วิปฺปสนฺน\-จิตฺตา สมาหิตา เอกคฺคจิตฺตา กายสฺส ยถาภูตํ ญาณาย. เวทนาสุ เวทนา\-นุปสฺสิโน วิหรถ อาตาปิโน สมฺปชานา เอโกทิภูตา วิปฺปสนฺน\-จิตฺตา สมาหิตา เอกคฺคจิตฺตา เวทนานํ ยถาภูตํ ญาณาย. จิตฺเต จิตฺตานุปสฺสิโน วิหรถ อาตาปิโน สมฺปชานา เอโกทิภูตา วิปฺปสนฺน\-จิตฺตา สมาหิตา เอกคฺคจิตฺตา จิตฺตสฺส ยถาภูตํ ญาณาย. ธมฺเมสุ ธมฺมา\-นุปสฺสิโน วิหรถ อาตาปิโน สมฺปชานา เอโกทิภูตา วิปฺปสนฺน\-จิตฺตา สมาหิตา เอกคฺคจิตฺตา ธมฺมานํ ยถาภูตํ ญาณาย. เยปิ เต ภิกฺขเว ภิกฺขู เสขา อปฺปตฺตมานสา อนุตฺตรํ โยคกฺเขมํ ปตฺถยมานา วิหรนฺติ, เตปิ กาเย กายานุปสฺสิโน วิหรนฺติ อาตาปิโน สมฺปชานา เอโกทิภูตา วิปฺปสนฺน\-จิตฺตา สมาหิตา เอกคฺคจิตฺตา กายสฺส ปริญฺญาย. เวทนาสุ เวทนา\-นุปสฺสิโน วิหรนฺติ อาตาปิโน สมฺปชานา เอโกทิภูตา วิปฺปสนฺน\-จิตฺตา สมาหิตา เอกคฺคจิตฺตา เวทนานํ ปริญฺญาย. จิตฺเต จิตฺตานุปสฺสิโน วิหรนฺติ อาตาปิโน สมฺปชานา เอโกทิภูตา วิปฺปสนฺน\-จิตฺตา สมาหิตา เอกคฺคจิตฺตา จิตฺตสฺส ปริญฺญาย. ธมฺเมสุ ธมฺมา\-นุปสฺสิโน วิหรนฺติ อาตาปิโน สมฺปชานา เอโกทิภูตา วิปฺปสนฺน\-จิตฺตา สมาหิตา เอกคฺคจิตฺตา ธมฺมานํ ปริญฺญาย.}

\prose{เยปิ เต ภิกฺขเว ภิกฺขู อรหนฺโต ขีณาสวา วุสิตวนฺโต กตกรณียา โอหิตภารา อนุปฺปตฺต\-สทตฺถา ปริกฺขีณ\-ภวสํโยชนา สมฺมทญฺญา วิมุตฺตา, เตปิ กาเย กายานุปสฺสิโน วิหรนฺติ อาตาวิโน สมฺปชานา เอโกทิภูตา วิปฺปสนฺน\-จิตฺตา สมาหิตา เอกคฺคจิตฺตา กาเยน วิสํยุตฺตา. เวทนาสุ เวทนา\-นุปสฺสิโน วิหรนฺติ อาตาปิโน สมฺปชานา เอโกทิภูตา วิปฺปสนฺน\-จิตฺตา สมาหิตา เอกคฺคจิตฺตา เวทนาหิ วิสํยุตฺตา. จิตฺเต จิตฺตานุปสฺสิโน วิหรนฺติ อาตาปิโน สมฺปชานา เอโกทิภูตา วิปฺปสนฺน\-จิตฺตา สมาหิตา เอกคฺคจิตฺตา จิตฺเตน วิสํยุตฺตา. ธมฺเมสุ ธมฺมา\-นุปสฺสิโน วิหรนฺติ อาตาปิโน สมฺปชานา เอโกทิภูตา วิปฺปสนฺน\-จิตฺตา สมาหิตา เอกคฺคจิตฺตา ธมฺเมหิ วิสํยุตฺตา.}

\prose{เยปิ เต ภิกฺขเว ภิกฺขู นวา อจิรปพฺพชิตา อธุนาคตา อิมํ ธมฺมวินยํ, เต โว ภิกฺขเว ภิกฺขู อิเมสํ จตุนฺนํ สติปฏฺฐานานํ ภาวนาย สมาทเปตพฺพา นิเวเสตพฺพา ปติฏฺฐาเป\-ตพฺพาติ.  จตุตฺถํ.}

\csromanheader{2}{3. สติปฏฺฐาน\-สํยุตฺต}
\csromanmatikaanchor{mtk.50}
\csromantocmarkref{subparagraph}{5. อกุสลราสิสุตฺต}{mtk.50}
\bookmark[dest={mtk.50},level=5]{5. อกุสลราสิสุตฺต}
\csromanheader{5}{5. อกุสลราสิ\-สุตฺต}
\proseitem{371}{\csromanfolio{127}สาวตฺถิ\-นิทานํ. ตตฺร โข ภควา เอตทโวจ\csromandash{}“อกุสลราสี”ติ ภิกฺขเว วทมาโน ปญฺจ นีวรเณ สมฺมา วทมาโน วเทยฺย. เกวโลหายํ ภิกฺขเว อกุสลราสิ ยทิทํ ปญฺจ นีวรณา. กตเม ปญฺจ, กามจฺฉนฺท\-นีวรณํ พฺยาปาท\-นีวรณํ ถินมิทฺธ\-นีวรณํ อุทฺธจฺจ\-กุกฺกุจฺจ\-นีวรณํ วิจิกิจฺฉา\-นีวรณํ. “อกุสลราสี”ติ ภิกฺขเว วทมาโน อิเม ปญฺจ นีวรเณ สมฺมา วทมาโน วเทยฺย. เกวโล หายํ ภิกฺขเว อกุสลราสิ ยทิทํ ปญฺจ นีวรณา.}

\prose{“กุสลราสี”ติ ภิกฺขเว วทมาโน จตฺตาโร สติปฏฺฐาเน สมฺมา วทมาโน วเทยฺย. เกวโล หายํ ภิกฺขเว กุสลราสิ ยทิทํ จตฺตาโร สติปฏฺฐานา. กตเม จตฺตาโร, อิธ ภิกฺขเว ภิกฺขุ กาเย กายานุปสฺสี วิหรติ อาตาปี สมฺปชาโน สติมา วิเนยฺย โลเก อภิชฺฌา\-โทมนสฺสํ. เวทนาสุ ฯเปฯ. จิตฺเต ฯเปฯ. ธมฺเมสุ ธมฺมานุปสฺสี วิหรติ อาตาปี สมฺปชาโน สติมา วิเนยฺย โลเก อภิชฺฌา\-โทมนสฺสํ. “กุสลราสี”ติ ภิกฺขเว วทมาโน อิเม จตฺตาโร สติปฏฺฐาเน สมฺมา วทมาโน วเทยฺย. เกวโล หายํ ภิกฺขเว กุสลราสิ ยทิทํ จตฺตาโร สติปฏฺฐานาติ.  ปญฺจมํ.}

\csromanmatikaanchor{mtk.51}
\csromantocmarkref{subparagraph}{6. สกุณคฺฆิสุตฺต}{mtk.51}
\bookmark[dest={mtk.51},level=5]{6. สกุณคฺฆิสุตฺต}
\csromanheader{5}{6. สกุณคฺฆิ\-สุตฺต}
\proseitem{372}{ภูตปุพฺพํ ภิกฺขเว สกุณคฺฆิ ลาปํ สกุณํ สหสา อชฺฌปฺปตฺตา อคฺคเหสิ. อถ โข ภิกฺขเว ลาโป สกุโณ สกุณคฺฆิยา หริยมาโน เอวํ ปริเทวสิ “มยเมวมฺห\csromansharedfootnote{fa6f255890ce7493}{มยเมวามฺห (ก)} อลกฺขิกา, มยํ อปฺปปุญฺญา, เย มยํ อโคจเร จริมฺห ปรวิสเย, สเจชฺช มยํ โคจเร จเรยฺยาม สเก เปตฺติเก วิสเย, น มฺยายํ\csromansharedfootnote{c6327765a3eab82b}{น จายํ (สี)} สกุณคฺฆิ อลํ อภวิสฺส ยทิทํ ยุทฺธายา”ติ. โก ปน เต ลาป โคจโร สโก เปตฺติโก วิสโยติ. ยทิทํ นงฺคล\-กฏฺฐ\-กรณํ เลฑฺฑุฏฺฐานนฺติ. อถ โข ภิกฺขเว สกุณคฺฆิ สเก พเล อปตฺถทฺธา สเก พเล อสํวทมานา\csromansharedfootnote{9485592d8a5c2e12}{อวจมานา (สี)} ลาปํ สกุณํ ปมุญฺจิ “คจฺฉ โข ตฺวํ ลาป ตตฺรปิ เม คนฺตฺวา น โมกฺขสี”ติ.}

\prose{\csromanfolio{128}อถ โข ภิกฺขเว ลาโป สกุโณ นงฺคล\-กฏฺฐ\-กรณํ โลฑฺฑุฏฺฐานํ คนฺตฺวา มหนฺตํ เลฑฺฑุํ อภิรุหิตฺวา สกุณคฺฆิํ วทมาโน อฏฺฐาสิ “เอหิ โข ทานิ เม สกุณคฺฆิ, เอหิ โข ทานิ เม สกุณคฺฆี”ติ. อถ โข สา ภิกฺขเว สกุณคฺฆิ สเก พเล อปตฺถทฺธา สเก พเล อสํวทมานา อุโภ ปกฺเข สนฺนยฺห\csromansharedfootnote{2884ee88bda843a6}{สนฺธาย (สี, สฺยา)} ลาปํ สกุณํ สหสา อชฺฌปฺปตฺตา. ยทา โข ภิกฺขเว อญฺญาสิ ลาโป สกุโณ พหุอาคโต โข มฺยายํ สกุณคฺฆีติ, อถ ตสฺเสว เลฑฺฑุสฺส อนฺตรํ ปจฺจุปาทิ. อถ โข ภิกฺขเว สกุณคฺฆิ ตตฺเถว อุรํ ปจฺจตาเฬสิ. เอวํ หิ ตํ\csromansharedfootnote{84f22dc3b8c31328}{เอวํ เหตํ (สี)} ภิกฺขเว โหติ โย อโคจเร จรติ ปรวิสเย.}

\prose{ตสฺมาติห ภิกฺขเว มา อโคจเร จริตฺถ ปรวิสเย, อโคจเร ภิกฺขเว จรตํ ปรวิสเย ลจฺฉติ มาโร โอตารํ, ลจฺฉติ มาโร อารมฺมณํ. โก จ ภิกฺขเว ภิกฺขุโน อโคจโร ปรวิสโย, ยทิทํ ปญฺจ กามคุณา. กตเม ปญฺจ, จกฺขุวิญฺเญยฺยา รูปา อิฏฺฐา กนฺตา มนาปา ปิยรูปา กามูปสํหิตา รชนียา. โสตวิญฺเญยฺยา สทฺทา ฯเปฯ. ฆานวิญฺเญยฺยา คนฺธา ฯเปฯ. ชิวฺหาวิญฺเญยฺยา รสา ฯเปฯ. กายวิญฺเญยฺยา โผฏฺฐพฺพา อิฏฺฐา กนฺตา มนาปา ปิยรูปา กามูปสํหิตา รชนียา. อยํ ภิกฺขเว ภิกฺขุโน อโคจโร ปรวิสโย.}

\prose{โคจเร ภิกฺขเว จรถ สเก เปตฺติเก วิสเย, โคจเร ภิกฺขเว จรตํ สเก เปตฺติเก วิสเย น ลจฺฉติ มาโร โอตารํ, น ลจฺฉติ มาโร อารมฺมณํ. โก จ ภิกฺขเว ภิกฺขุโน โคจโร สโก เปตฺติโก วิสโย, ยทิทํ จตฺตาโร สติปฏฺฐานา. กตเม จตฺตาโร, อิธ ภิกฺขเว ภิกฺขุ กาเย กายานุปสฺสี วิหรติ อาตาปี สมฺปชาโน สติมา วิเนยฺย โลเก อภิชฺฌา\-โทมนสฺสํ. เวทนาสุ ฯเปฯ. จิตฺเต ฯเปฯ. ธมฺเมสุ ธมฺมานุปสฺสี วิหรติ อาตาปี สมฺปชาโน สติมา วิเนยฺย โลเก อภิชฺฌา\-โทมนสฺสํ. อยํ ภิกฺขเว ภิกฺขุโน โคจโร สโก เปตฺติโก วิสโยติ.  ฉฏฺฐํ.}

\csromanmatikaanchor{mtk.52}
\csromantocmarkref{subparagraph}{7. มกฺกฏสุตฺต}{mtk.52}
\bookmark[dest={mtk.52},level=5]{7. มกฺกฏสุตฺต}
\csromanheader{5}{7. \textbf{มกฺกฏสุตฺต}}
\proseitem{373}{อตฺถิ ภิกฺขเว หิมวโต ปพฺพตราชสฺส ทุคฺคา วิสมา เทสา, ยตฺถ เนว มกฺกฏานํ จารี น มนุสฺสานํ. อตฺถิ ภิกฺขเว หิมวโต}

\prose{\csromanfolio{129}ปพฺพตราชสฺส ทุคฺคา วิสมา เทสา, ยตฺถ มกฺกฏานํ หิ โข จารี น มนุสฺสานํ. อตฺถิ ภิกฺขเว หิมวโต ปพฺพตราชสฺส สมา ภูมิภาคา รมณียา, ยตฺถ มกฺกฏานญฺเจว จารี มนุสฺสานญฺจ. ตตฺร ภิกฺขเว ลุทฺทา มกฺกฏวีถีสุ เลปํ โอฑฺเฑนฺติ มกฺกฏานํ พาธนาย.}

\prose{ตตฺร ภิกฺขเว เย เต มกฺกฏา อพาลชาติกา อโลลชาติกา, เต ตํ เลปํ ทิสฺวา อารกา ปริวชฺชนฺติ. โย ปน โส โหติ มกฺกโฏ พาลชาติโก โลลชาติโก, โส ตํ เลปํ อุปสงฺกมิตฺวา หตฺเถน คณฺหาติ, โส ตตฺถ พชฺฌติ. “หตฺถํ โมเจสฺสามี”ติ ทุติเยน หตฺเถน คณฺหาติ, โส ตตฺถ พชฺฌติ. “อุโภ หตฺเถ โมเจสฺสามี”ติ ปาเทน คณฺหาติ, โส ตตฺถ พชฺฌติ. “อุโภ หตฺเถ โมเจสฺสามิ ปาทญฺจา”ติ ทุติเยน ปาเทน คณฺหาติ, โส ตตฺถ พชฺฌติ. “อุโภ หตฺเถ โมเจสฺสามิ ปาเท จา”ติ ตุณฺเฑน คณฺหาติ, โส ตตฺถ พชฺฌติ. เอวํ หิ โส ภิกฺขเว มกฺกโฏ ปญฺโจฑฺฑิโต ถุนํ เสติ อนยํ อาปนฺโน, พฺยสนํ อาปนฺโน, ยถา\-กาม\-กรณีโย ลุทฺทสฺส. ตเมนํ ภิกฺขเว ลุทฺโท วิชฺฌิตฺวา ตสฺมิํ เยว กฏฺฐ\-กต\-งฺคาเร\csromansharedfootnote{fbc7f7d7b1c4e837}{ตสฺมิํ เยว มกฺกฏํ อุทฺธริตฺวา (สี, สฺยา)} อวสฺสชฺเชตฺวา เยน กามํ ปกฺกมติ. เอวํ โส ตํ ภิกฺขเว โหติ, โย อโคจเร จรติ ปรวิสเย.}

\prose{ตสฺมาติห ภิกฺขเว มา อโคจเร จริตฺถ ปรวิสเย, อโคจเร ภิกฺขเว จรตํ ปรวิสเย ลจฺฉติ มาโร โอตารํ, ลจฺฉติ มาโร อารมฺมณํ. โก จ ภิกฺขเว ภิกฺขุโน อโคจโร ปรวิสโย, ยทิทํ ปญฺจ กามคุณา. กตเม ปญฺจ, จกฺขุวิญฺเญยฺยา รูปา อิฏฺฐา กนฺตา มนาปา ปิยรูปา กามูปสํหิตา รชนียา. โสตวิญฺเญยฺยา สทฺทา ฯเปฯ. ฆานวิญฺเญยฺยา คนฺธา ฯเปฯ. ชิวฺหาวิญฺเญยฺยา รสา ฯเปฯ. กายวิญฺเญยฺยา โผฏฺฐพฺพา อิฏฺฐา กนฺตา มนาปา ปิยรูปา กามูปสํหิตา รชนียา. อยํ ภิกฺขเว ภิกฺขุโน อโคจโร ปรวิสโย.}

\prose{โคจเร ภิกฺขเว จรถ สเก เปตฺติเก วิสเย, โคจเร ภิกฺขเว จรตํ สเก เปตฺติเก วิสเย น ลจฺฉติ มาโร โอตารํ, น ลจฺฉติ มาโร อารมฺมณํ. โก จ ภิกฺขเว ภิกฺขุโน โคจโร สโก เปตฺติโก วิสโย, ยทิทํ จตฺตาโร สติปฏฺฐานา. กตเม}

\prose{\csromanfolio{130}จตฺตาโร, อิธ ภิกฺขเว ภิกฺขุ กาเย กายานุปสฺสี วิหรติ อาตาปี สมฺปชาโน สติมา วิเนยฺย โลเก อภิชฺฌา\-โทมนสฺสํ. เวทนาสุ ฯเปฯ. จิตฺเต ฯเปฯ. ธมฺเมสุ ธมฺมานุปสฺสี วิหรติ อาตาปี สมฺปชาโน สติมา วิเนยฺย โลเก อภิชฺฌา\-โทมนสฺสํ. อยํ ภิกฺขเว ภิกฺขุโน โคจโร สโก เปตฺติโก วิสโยติ.  สตฺตมํ.}

\csromanmatikaanchor{mtk.53}
\csromantocmarkref{subparagraph}{8. สูทสุตฺต}{mtk.53}
\bookmark[dest={mtk.53},level=5]{8. สูทสุตฺต}
\csromanheader{5}{8. \textbf{สูทสุตฺต}}
\proseitem{374}{เสยฺยถาปิ ภิกฺขเว พาโล อพฺยตฺโต อกุสโล สูโท ราชานํ วา ราชมหามตฺตํ วา\csromansharedfootnote{6b6d082f8335933a}{ราชมหามตฺตานํ วา (สี)} นานจฺจเยหิ สูเปหิ ปจฺจุปฏฺฐิโต อสฺส อมฺพิลคฺเคหิปิ ติตฺตกคฺเคหิปิ กฏุกคฺเคหิปิ มธุรคฺเคหิปิ ขาริเกหิปิ อขาริเกหิปิ โลณิเกหิปิ อโลณิเกหิปิ.}

\prose{ส โข โส ภิกฺขเว พาโล อพฺยตฺโต อกุสโล สูโท สกสฺส ภตฺตุ นิมิตฺตํ น อุคฺคณฺหาติ “อิทํ วา เม อชฺช ภตฺตุ สูเปยฺยํ รุจฺจติ, อิมสฺส วา อภิหรติ. อิมสฺส วา พหุํ คณฺหาติ, อิมสฺส วา วณฺณํ ภาสติ. อมฺพิลคฺคํ วา เม อชฺช ภตฺตุ สูเปยฺยํ รุจฺจติ, อมฺพิลคฺคสฺส วา อภิหรติ, อมฺพิลคฺคสฺส วา พหุํ คณฺหาติ, อมฺภิลคฺคสฺส วา วณฺณํ ภาสติ. ติตฺตกคฺคํ วา เม อชฺช. กฏุกคฺคํ วา เม อชฺช. มธุรคฺคํ วา เม อชฺช. ขาริกํ วา เม อชฺช. อขาริกํ วา เม อชฺช. โลณิกํ วา เม อชฺช. อโลณิกํ วา เม อชฺช ภตฺตุ สูเปยฺยํ รุจฺจติ, อโลณิกสฺส วา อภิหรติ, อโลณิกสฺส วา พหุํ คณฺหาติ, อโลณิกสฺส วา วณฺณํ ภาสตี”ติ.}

\prose{ส โข โส ภิกฺขเว พาโล อพฺยตฺโต อกุสโล สูโท น เจว ลาภี โหติ อจฺฉาทนสฺส, น ลาภี เวตนสฺส, น ลาภี อภิหารานํ. ตํ กิสฺส เหตุ, ตถา หิ โส ภิกฺขเว พาโล อพฺยตฺโต อกุสโล สูโท สกสฺส ภตฺตุ นิมิตฺตํ น อุคฺคณฺหาติ. เอวเมว โข ภิกฺขเว อิเธกจฺโจ พาโล อพฺยตฺโต อกุสโล ภิกฺขุ กาเย กายานุปสฺสี วิหรติ อาตาปี สมฺปชาโน สติมา วิเนยฺย โลเก อภิชฺฌา\-โทมนสฺสํ, ตสฺส กาเย กายานุปสฺสิโน วิหรโต จิตฺตํ น สมาธิยติ, อุปกฺกิเลสา น ปหียนฺติ, โส ตํ นิมิตฺตํ น อุคฺคณฺหาติ. เวทนาสุ เวทนานุปสฺสี วิหรติ ฯเปฯ. จิตฺเต จิตฺตานุปสฺสี}

\prose{\csromanfolio{131}วิหรติ ฯเปฯ. ธมฺเมสุ ธมฺมานุปสฺสี วิหรติ อาตาปี สมฺปชาโน สติมา วิเนยฺย โลเก อภิชฺฌา\-โทมนสฺสํ, ตสฺส ธมฺเมสุ ธมฺมา\-นุปสฺสิโน วิหรโต จิตฺตํ น สมาธิยติ, อุปกฺกิเลสา น ปหียนฺติ, โส ตํ นิมิตฺตํ น อุคฺคณฺหาติ.}

\prose{ส โข โส ภิกฺขเว พาโล อพฺยตฺโต อกุสโล ภิกฺขุ น เจว ลาภี โหติ ทิฏฺเฐว ธมฺเม สุขวิหารานํ, น ลาภี สติสมฺปชญฺญสฺส. ตํ กิสฺส เหตุ, ตถา หิ โส ภิกฺขเว พาโล อพฺยตฺโต อกุสโล ภิกฺขุ สกสฺส จิตฺตสฺส นิมิตฺตํ น อุคฺคณฺหาติ.}

\prose{เสยฺยถาปิ ภิกฺขเว ปณฺฑิโต พฺยตฺโต กุสโล สูโท ราชานํ วา ราชมหามตฺตํ วา นานจฺจเยหิ สูเปหิ ปจฺจุปฏฺฐิโต อสฺส อมฺพิลคฺเคหิปิ ติตฺตกคฺเคหิปิ กฏุกคฺเคหิปิ มธุรคฺเคหิปิ ขาริเกหิปิ อขาริเกหิปิ โลณิเกหิปิ อโลณิเกหิปิ.}

\prose{ส โข โส ภิกฺขเว ปณฺฑิโต พฺยตฺโต กุสโล สูโท สกสฺส ภตฺตุ นิมิตฺตํ อุคฺคณฺหาติ “อิทํ วา เม อชฺช ภตฺตุ สูเปยฺยํ รุจฺจติ, อิมสฺส วา อภิหรติ, อิมสฺส วา พหุํ คณฺหาติ, อิมสฺส วา วณฺณํ ภาสติ. อมฺพิลคฺคํ วา เม อชฺช ภตฺตุ สูเปยฺยํ รุจฺจติ, อมฺพิลคฺคสฺส วา อภิหรติ, อมฺพิลคฺคสฺส วา พหุํ คณฺหาติ, อมฺพิลคฺคสฺส วา วณฺณํ ภาสติ. ติตฺตกคฺคํ วา เม อชฺช. กฏุกคฺคํ วา เม อชฺช. มธุรคฺคํ วา เม อชฺช. ขาริกํ วา เม อชฺช. อขาริกํ วา เม อชฺช. โลณิกํ วา เม อชฺช. อโลณิกํ วา เม อชฺช ภตฺตุ สูเปยฺยํ รุจฺจติ. อโลณิกสฺส วา อภิหรติ, อโลณิกสฺส วา พหุํ คณฺหาติ. อโลณิกสฺส วา วณฺณํ ภาสตี”ติ.}

\prose{ส โข โส ภิกฺขเว ปณฺฑิโต พฺยตฺโต กุสโล สูโท ลาภี เจว โหติ อจฺฉาทนสฺส, ลาภี เวตนสฺส, ลาภี อภิหารานํ. ตํ กิสฺส เหตุ, ตถา หิ โส ภิกฺขเว ปณฺฑิโต พฺยตฺโต กุสโล สูโท สกสฺส ภตฺตุ นิมิตฺตํ อุคฺคณฺหาติ. เอวเมว โข ภิกฺขเว อิเธกจฺโจ ปณฺฑิโต พฺยตฺโต กุสโล ภิกฺขุ กาเย กายานุปสฺสี วิหรติ อาตาปี สมฺปชาโน สติมา วิเนยฺย โลเก อภิชฺฌา\-โทมนสฺสํ, ตสฺส กาเย กายานุปสฺสิโน วิหรโต จิตฺตํ สมาธิยติ, อุปกฺกิเลสา ปหียนฺติ, โส ตํ นิมิตฺตํ อุคฺคณฺหาติ. เวทนาสุ เวทนานุปสฺสี วิหรติ ฯเปฯ. จิตฺเต จิตฺตานุปสฺสี วิหรติ ฯเปฯ. ธมฺเมสุ ธมฺมานุปสฺสี}

\prose{\csromanfolio{132}วิหรติ อาตาปี สมฺปชาโน สติมา วิเนยฺย โลเก อภิชฺฌา\-โทมนสฺสํ, ตสฺส ธมฺเมสุ ธมฺมา\-นุปสฺสิโน วิหรโต จิตฺตํ สมาธิยติ, อุปกฺกิเลสา ปหียนฺติ, โส ตํ นิมิตฺตํ อุคฺคณฺหาติ.}

\prose{ส โข โส ภิกฺขเว ปณฺฑิโต พฺยตฺโต กุสโล ภิกฺขุ ลาภี เจว โหติ ทิฏฺเฐว ธมฺเม สุขวิหารานํ, ลาภี โหติ สติสมฺปชญฺญสฺส. ตํ กิสฺส เหตุ, ตถา หิ โส ภิกฺขเว ปณฺฑิโต พฺยตฺโต กุสโล ภิกฺขุ สกสฺส จิตฺตสฺส นิมิตฺตํ อุคฺคณฺหาตีติ.  อฏฺฐมํ.}

\csromanmatikaanchor{mtk.54}
\csromantocmarkref{subparagraph}{9. คิลานสุตฺต}{mtk.54}
\bookmark[dest={mtk.54},level=5]{9. คิลานสุตฺต}
\csromanheader{5}{9. \textbf{คิลานสุตฺต}}
\proseitem{375}{เอวํ เม สุตํ\csromandash{}เอกํ สมยํ ภควา เวสาลิยํ วิหรติ เวฬุวคามเก\csromansharedfootnote{dfe80a66b4b645fe}{เพลุว คามเก (สี, สฺยา, กํ, อิ)}. ตตฺร โข ภควา ภิกฺขู อมนฺเตสิ “เอถ ตุเมฺห ภิกฺขเว สมนฺตา เวสาลิยา ยถามิตฺตํ ยถาสนฺทิฏฺฐํ ยถาสมฺภตฺตํ วสฺสํ อุเปถ, อิเธวาหํ เวฬุวคามเก วสฺสํ อุปคจฺฉามี”ติ. “เอวํ ภนฺเต”ติ โข เต ภิกฺขู ภควโต ปฏิสฺสุตฺวา สมนฺตา เวสาลิยา ยถามิตฺตํ ยถาสนฺทิฏฺฐํ ยถาสมฺภตฺตํ วสฺสํ อุปคจฺฉุํ, ภควา ปน ตตฺเถว เวฬุวคามเก วสฺสํ อุปคจฺฉิ\csromansharedfootnote{721f9c106aa11a61}{อุปคญฺฉิ (สี, อิ)}.}

\prose{อถ โข ภควโต วสฺสูปคตสฺส ขโร อาพาโธ อุปฺปชฺชิ, พาฬฺหา เวทนา วตฺตนฺติ มารณนฺติกา, ตตฺร สุทํ ภควา สโต สมฺปชาโน อธิวาเสสิ อวิหญฺญมาโน. อถ โข ภควโต เอตทโหสิ “น โข เม ตํ ปติรูปํ, โยหํ อนามนฺเตตฺวา อุปฏฺฐาเก อนปโลเกตฺวา ภิกฺขุสํฆํ ปรินิพฺพาเยยฺยํ, ยํนูนาหํ อิมํ อาพาธํ วีริเยน ปฏิปณาเมตฺวา ชีวิต\-สงฺขารํ อธิฏฺฐาย วิหเรยฺยนฺ”ติ. อถ โข ภควา ตํ อาพาธํ วีริเยน ปฏิปณาเมตฺวา ชีวิต\-สงฺขารํ อธิฏฺฐาย วิหาสิ. (อถ โข ภควโต โส อาพาโธ ปฏิปฺปสฺสมฺภิ.)\csromansharedfootnote{8f90839b00cc02f1}{( ) ที 2. 84 ปิฏฺเฐ ทิสฺสติ.}}

\prose{อถ โข ภควา คิลานา วุฏฺฐิโต\csromansharedfootnote{7ad0dbd929e80867}{คิลานวุฏฺฐิโต (สทฺทนีติ)} อจิรวุฏฺฐิโต เคลญฺญา วิหารา นิกฺขมิตฺวา วิหาร\-ปจฺฉายายํ\csromansharedfootnote{8eed76fa099b969a}{วิหารปจฺฉาฉายายํ (พหูสุ)} ปญฺญตฺเต อาสเน นิสีทิ. อถ โข อายสฺมา อานนฺโท เยน ภควา เตนุปสงฺกมิ, อุปสงฺกมิตฺวา ภควนฺตํ อภิวาเทตฺวา เอกมนฺตํ นิสีทิ, เอกมนฺตํ นิสินฺโน โข อายสฺมา อานนฺโท}

\prose{\csromanfolio{133}ภควนฺตํ เอตทโวจ “ทิฏฺโฐ เม ภนฺเต ภควโต ผาสุ, ทิฏฺฐํ ภนฺเต ภควโต ขมนียํ, ทิฏฺฐํ ภนฺเต ภควโต ยาปนียํ, อปิ จ เม ภนฺเต มธุรกชาโต วิย กาโย, ทิสาปิ เม น ปกฺขายนฺติ, ธมฺมาปิ มํ นปฺปฏิภนฺติ ภควโต เคลญฺเญน, อปิ จ เม ภนฺเต อโหสิ กาจิเทว อสฺสาสมตฺตา, น ตาว ภควา ปรินิพฺพายิสฺสติ, น ยาว ภควา ภิกฺขุสํฆํ อารพฺภ กิญฺจิเทว อุทาหรตี”ติ.}

\prose{กิํ ปน ทานิ อานนฺท ภิกฺขุสํโฆ มยิ ปจฺจาสีสติ\csromansharedfootnote{2dfb335c1531753a}{ปจฺจาสิํ สติ (สี, สฺยา, กํ, อิ)}, เทสิโต อานนฺท มยา ธมฺโม อนนฺตรํ อพาหิรํ กริตฺวา, นตฺถานนฺท ตถาคตสฺส ธมฺเมสุ อาจริยมุฏฺฐิ, ยสฺส นูน อานนฺท เอวมสฺส “อหํ ภิกฺขุสํฆํ ปริหริสฺสามี”ติ วา “มมุทฺเทสิโก ภิกฺขุสํโฆ”ติ วา, โส นูน อานนฺท ภิกฺขุสํฆํ อารพฺภ กิญฺจิเทว อุทาหเรยฺย. ตถาคตสฺส โข อานนฺท น เอวํ โหติ “อหํ ภิกฺขุสํฆํ ปริหริสฺสามี”ติ วา “มมุทฺเทสิโก ภิกฺขุสํโฆ”ติ วา, ส กิํ\csromansharedfootnote{b3af893a070cb4ed}{โส นูน (สี, อิ)} อานนฺท ตถาคโต ภิกฺขุสํฆํ อารพฺภ กิญฺจิเทว อุทาหริสฺสติ, เอตรหิ โข ปนาหํ อานนฺท ชิณฺโณ วุทฺโธ มหลฺลโก อทฺธคโต วโยอนุปฺปตฺโต, อาสีติโก เม วโย วตฺตติ. เสยฺยถาปิ อานนฺท ชชฺชรสกฏํ\csromansharedfootnote{1a5a6a8a804113a8}{ชรสกฏํ (สพฺพตฺถ) 4. เวคมิสฺสเกน (สี), เวฬุมิสฺสเกน (สฺยา, กํ), เวธมิสฺสเกน (อิ, ก), เวขมิสฺสเกน (ก) 5. ผาสุตรํ (สพฺพตฺถ)} เวฐมิสฺสเกน\csromansharedfootnote{39ff17ed4cb8af0f}{เวคมิสฺสเกน (สี), เวฬุมิสฺสเกน (สฺยา, กํ), เวธมิสฺสเกน (อิ, ก), เวขมิสฺสเกน (ก) 5. ผาสุตรํ (สพฺพตฺถ)} ยาเปติ. เอวเมว โข อานนฺท เวขมิสฺสเกน มญฺเญ ตถาคตสฺส กาโย ยาเปติ.}

\prose{ยสฺมิํ อานนฺท สมเย ตถาคโต สพฺพ\-นิมิตฺตานํ อมนสิการา เอกจฺจานํ เวทนานํ นิโรธา อนิมิตฺตํ เจโตสมาธิํ อุปสมฺปชฺช วิหรติ. ผาสุตโร5 อานนฺท ตสฺมิํ สมเย ตถาคตสฺส กาโย โหติ\csromansharedfootnote{01401eb4e439502e}{ตถาคตสฺส โหติ (พหูสุ)}. ตสฺมาติหานนฺท อตฺตทีปา วิหรถ อตฺตสรณา อนญฺญสรณา ธมฺมทีปา ธมฺมสรณา อนญฺญสรณา.}

\prose{กถญฺจานนฺท ภิกฺขุ อตฺตทีโป วิหรติ อตฺตสรโณ อนญฺญสรโณ ธมฺมทีโป ธมฺมสรโณ อนญฺญสรโณ, อิธานนฺท ภิกฺขุ กาเย กายานุปสฺสี วิหรติ อาตาปี สมฺปชาโน สติมา วิเนยฺย โลเก อภิชฺฌา\-โทมนสฺสํ. เวทนาสุ ฯเปฯ. จิตฺเต ฯเปฯ. ธมฺเมสุ ธมฺมานุปสฺสี วิหรติ อาตาปี}

\prose{\csromanfolio{134}สมฺปชาโน สติมา วิเนยฺย โลเก อภิชฺฌา\-โทมนสฺสํ. เอวํ โข อานนฺท ภิกฺขุ อตฺตทีโป วิหรติ อตฺตสรโณ อนญฺญสรโณ ธมฺมทีโป ธมฺมสรโณ อนญฺญสรโณ. เย หิ เกจิ อานนฺท เอตรหิ วา มมจฺจเย วา อตฺตทีปา วิหริสฺสนฺติ อตฺตสรณา อนญฺญสรณา ธมฺมทีปา ธมฺมสรณา อนญฺญสรณา. ตมตคฺเค เมเต อานนฺท ภิกฺขู ภวิสฺสนฺติ เย เกจิ สิกฺขากามาติ.  นวมํ.}

\csromanmatikaanchor{mtk.55}
\csromantocmarkref{subparagraph}{10. ภิกฺขุนุปสฺสยสุตฺต}{mtk.55}
\bookmark[dest={mtk.55},level=5]{10. ภิกฺขุนุปสฺสยสุตฺต}
\csromanheader{5}{10. ภิกฺขุนุปสฺสย\-สุตฺต}
\proseitem{376}{อถ โข อายสฺมา อานนฺโท ปุพฺพณฺหสมยํ นิวาเสตฺวา ปตฺต\-จีวรมา\-ทาย เยน อญฺญตโร ภิกฺขุนุ\-ปสฺสโย เตนุปสงฺกมิ, อุปสงฺกมิตฺวา ปญฺญตฺเต อาสเน นิสีทิ. อถ โข สมฺพหุลา ภิกฺขุนิโย เยนายสฺมา อานนฺโท เตนุ\-ปสงฺกมิํสุ, อุปสงฺกมิตฺวา อายสฺมนฺตํ อานนฺทํ อภิวาเทตฺวา เอกมนฺตํ นิสีทิํสุ, เอกมนฺตํ นิสินฺโน โข ตา ภิกฺขุนิโย อายสฺมนฺตํ อานนฺทํ เอตทโวจุํ\csromandash{}}

\prose{อิธ ภนฺเต อานนฺท สมฺพหุลา ภิกฺขุนิโย จตูสุ สติปฏฺฐาเนสุ สุปฺปติฏฺฐิต\-จิตฺตา\csromansharedfootnote{223667655bec1a82}{สุปฏฺฐิตจิตฺตา (สี, อิ, ก)} วิหรนฺติโย อุฬารํ ปุพฺเพนาปรํ วิเสสํ สญฺชานนฺตีติ\csromansharedfootnote{4a7b3820996021f6}{สมฺปชานนฺตีติ (ก)}. เอวเมตํ ภคินิโย, เอวเมตํ ภคินิโย, โย หิ โกจิ ภคินิโย ภิกฺขุ วา ภิกฺขุนี วา จตูสุ สติปฏฺฐาเนสุ สุปฺปติฏฺฐิต\-จิตฺโต วิหรติ, ตสฺเสตํ ปาฏิกงฺขํ “อุฬารํ ปุพฺเพนาปรํ วิเสสํ สญฺชานิสฺสตี”ติ.}

\prose{อถ โข อายสฺมา อานนฺโท ตา ภิกฺขุนิโย ธมฺมิยา กถาย สนฺทสฺเสตฺวา สมาทเปตฺวา สมุตฺเตเชตฺวา สมฺปหํเสตฺวา อุฏฺฐายาสนา ปกฺกามิ. อถ โข อายสฺมา อานนฺโท สาวตฺถิยํ ปิณฺฑาย จริตฺวา ปจฺฉาภตฺตํ ปิณฺฑปาต\-ปฏิกฺกนฺโต เยน ภควา เตนุปสงฺกมิ, อุปสงฺกมิตฺวา ภควนฺตํ อภิวาเทตฺวา เอกมนฺตํ นิสีทิ, เอกมนฺตํ นิสินฺโน โข อายสฺมา อานนฺโท ภควนฺตํ เอตทโวจ\csromandash{}}

\prose{อิธาหํ ภนฺเต ปุพฺพณฺหสมยํ นิวาเสตฺวา ปตฺต\-จีวรมา\-ทาย เยน อญฺญตโร ภิกฺขุนุ\-ปสฺสโย เตนุปสงฺกมิํ, อุปสงฺกมิตฺวา ปญฺญตฺเต อาสเน นิสีทิ. อถ โข ภนฺเต สมฺพหุลา ภิกฺขุนิโย เยนาหํ เตนุ\-ปสงฺกมิํสุ, อุปสงฺกมิตฺวา มํ อภิวาเทตฺวา เอกมนฺตํ นิสีทิํสุ,}

\prose{\csromanfolio{135}เอกมนฺตํ นิสินฺนา โข ภนฺเต ตา ภิกฺขุนิโย มํ เอตทโวจุํ “อิธ ภนฺเต อานนฺท สมฺพหุลา ภิกฺขุนิโย จตูสุ สติปฏฺฐาเนสุ สุปฺปติฏฺฐิต\-จิตฺตา วิหรนฺติโย อุฬารํ ปุพฺเพนาปรํ วิเสสํ สญฺชานนฺตี”ติ. เอวํ วุตฺตาหํ ภนฺเต ตา ภิกฺขุนิโย เอตทโวจํ “เอวเมตํ ภคินิโย, เอวเมตํ ภคินิโย, โย หิ โกจิ ภคินิโย ภิกฺขุ วา ภิกฺขุนี วา จตูสุ สติปฏฺฐาเนสุ สุปฺปติฏฺฐิต\-จิตฺโต วิหรติ, ตสฺเสตํ ปาฏิกงฺขํ อุฬารํ ปุพฺเพนาปรํ วิเสสํ สญฺชานิสฺสตี”ติ.}

\prose{เอวเมตํ อานนฺท, เอวเมตํ อานนฺท, โย หิ โกจิ อานนฺท ภิกฺขุ วา ภิกฺขุนี วา จตูสุ สติปฏฺฐาเนสุ สุปฺปติฏฺฐิต\-จิตฺโต วิหรติ, ตสฺเสตํ ปาฏิกงฺขํ อุฬารํ ปุพฺเพนาปรํ วิเสสํ สญฺชานิสฺสติ\csromansharedfootnote{6ee5b42501ee250c}{สญฺชานิสฺสตีติ (พหูสุ)}.}

\prose{กตเมสุ จตูสุ, อิธานนฺท ภิกฺขุ กาเย กายานุปสฺสี วิหรติ อาตาปี สมฺปชาโน สติมา วิเนยฺย โลเก อภิชฺฌา\-โทมนสฺสํ, ตสฺส กาเย กายานุปสฺสิโน วิหรโต กายารมฺมโณ วา อุปฺปชฺชติ กายสฺมิํ ปริฬาโห เจตโส วา ลีนตฺตํ, พหิทฺธา วา จิตฺตํ วิกฺขิปติ. เตนานนฺท\csromansharedfootnote{96602003cf0f048f}{เตนหานนฺท (สี)} ภิกฺขุนา กิสฺมิญฺจิเทว ปสาทนีเย นิมิตฺเต จิตฺตํ ปณิทหิตพฺพํ, ตสฺส กิสฺมิญฺจิเทว ปสาทนีเย นิมิตฺเต จิตฺตํ ปณิทหโต ปาโมชฺชํ ชายติ, ปมุทิตสฺส ปีติ ชายติ, ปีติมนสฺส กาโย ปสฺสมฺภติ, ปสฺสทฺธกาโย สุขํ เวทยติ\csromansharedfootnote{8f9d689ff21b5bf2}{เวทิยติ (สี)}, สุขิโน จิตฺตํ สมาธิยติ. โส อิติ ปฏิสญฺจิกฺขติ “ยสฺส ขฺวาหํ อตฺถาย จิตฺตํ ปณิทหิํ, โส เม อตฺโถ อภินิปฺผนฺโน, หนฺท ทานิ ปฏิสํหรามี”ติ. โส ปฏิสํหรติ เจว น จ วิตกฺเกติ น จ วิจาเรติ, “อวิตกฺโกมฺหิ อวิจาโร อชฺฌตฺตํ สติมา สุขมสฺมี”ติ ปชานาติ.}

\prose{ปุน จปรํ อานนฺท ภิกฺขุ เวทนาสุ ฯเปฯ. จิตฺเต ฯเปฯ. ธมฺเมสุ ธมฺมานุปสฺสี วิหรติ อาตาปี สมฺปชาโน สติมา วิเนยฺย โลเก อภิชฺฌา\-โทมนสฺสํ, ตสฺส ธมฺเมสุ ธมฺมา\-นุปสฺสิโน วิหรโต ธมฺมารมฺมโณ วา อุปฺปชฺชติ กายสฺมิํ ปริฬาโห เจตโส วา ลีนตฺตํ, พหิทฺธา วา จิตฺตํ วิกฺขิปติ. เตนานนฺท ภิกฺขุนา กิสฺมิญฺจิเทว ปสาทนีเย นิมิตฺเต จิตฺตํ ปณิทหิตพฺพํ, ตสฺส กิสฺมิญฺจิ เทว ปสาทนีเย นิมิตฺเต จิตฺตํ ปณิทหโต ปาโมชฺชํ}

\prose{\csromanfolio{136}ชายติ, ปมุทิตสฺส ปีติ ชายติ, ปีติมนสฺส กาโย ปสฺสมฺภติ, ปสฺสทฺธกาโย สุขํ เวทยติ, สุขิโน จิตฺตํ สมาธิยติ. โส อิติ ปฏิสญฺจิกฺขติ “ยสฺส ขฺวาหํ อตฺถาย จิตฺตํ ปณิทหิํ, โส เม อตฺโถ อภินิปฺผนฺโน, หนฺท ทานิ ปฏิสํหรามี”ติ. โส ปฏิสํหรติ เจว น จ วิตกฺเกติ น จ วิจาเรติ, “อวิตกฺโกมฺหิ อวิจาโร อชฺฌตฺตํ สติมา สุขมสฺมี”ติ ปชานาติ. เอวํ โข อานนฺท ปณิธาย ภาวนา โหติ.}

\prose{กถญฺจานนฺท อปฺปณิธาย ภาวนา โหติ, พหิทฺธา อานนฺท ภิกฺขุ จิตฺตํ อปฺปณิธาย “อปฺปณิหิตํ เม พหิทฺธา จิตฺตนฺ”ติ ปชานาติ. อถ “ปจฺฉาปุเร อสํขิตฺตํ วิมุตฺตํ อปฺปณิหิตนฺ”ติ ปชานาติ. อถ จ ปน “กาเย กายานุปสฺสี วิหรามิ, อาตาปี สมฺปชาโน สติมา สุขมสฺมี”ติ ปชานาติ. พหิทฺธา อานนฺท ภิกฺขุ จิตฺตํ อปฺปณิธาย “อปฺปณิหิตํ เม พหิทฺธา จิตฺตนฺ”ติ ปชานาติ. อถ “ปจฺฉาปุเร อสํขิตฺตํ วิมุตฺตํ อปฺปณิหิตนฺ”ติ ปชานาติ. อถ จ ปน “เวทนาสุ เวทนานุปสฺสี วิหรามิ, อาตาปี สมฺปชาโน สติมา สุขมสฺมี”ติ ปชานาติ. พหิทฺธา อานนฺท ภิกฺขุ จิตฺตํ อปฺปณิธาย “อปฺปณิหิตํ เม พหิทฺธา จิตฺตนฺ”ติ ปชานาติ. อถ “ปจฺฉา ปุเร อสํขิตฺตํ วิมุตฺตํ อปฺปณิหิตนฺ”ติ ปชานาติ. อถ จ ปน “จิตฺเต จิตฺตานุปสฺสี วิหรามิ, อาตาปี สมฺปชาโน สติมา สุขมสฺมี”ติ ปชานาติ. พหิทฺธา อานนฺท ภิกฺขุ จิตฺตํ อปฺปณิธาย “อปฺปณิหิตํ เม พหิทฺธา จิตฺตนฺ”ติ ปชานาติ. อถ “ปจฺฉา ปุเร อสํขิตฺตํ วิมุตฺตํ อปฺปณิหิตนฺ”ติ ปชานาติ. อถ จ ปน “ธมฺเมสุ ธมฺมานุปสฺสี วิหรามิ, อาตาปี สมฺปชาโน สติมา สุขมสฺมี”ติ ปชานาติ. เอวํ โข อานนฺท อปฺปณิธาย ภาวนา โหติ.}

\prose{อิติ โข อานนฺท เทสิตา มยา ปณิธาย ภาวนา, เทสิตา อปฺปณิธาย ภาวนา. ยํ อานนฺท สตฺถารา กรณียํ สาวกานํ หิเตสินา อนุกมฺปเกน อนุกมฺปํ อุปาทาย, กตํ โว ตํ มยา. เอตานิ อานนฺท รุกฺขมูลานิ, เอตานิ สุญฺญาคารานิ. ฌายถานนฺท มา ปมาทตฺถ, มา ปจฺฉา วิปฺปฏิสาริโน อหุวตฺถ, อยํ โว อมฺหากํ อนุสาสนีติ.}

\prosewithclosermidruled{อิทมโวจ ภควา. อตฺตมโน อายสฺมา อานนฺโท ภควโต ภาสิตํ อภินนฺทีติ.  ทสมํ.}{อมฺพปาลิวคฺโค ปฐโม.}
\csromancenter{\textbf{ตสฺสุทฺทานํ}}
\setlength{\csromangathapagewidth}{0pt}
\setlength{\csromangathawakwidth}{0pt}
\csromangathameasuregroup{อมฺพปาลิ สโต ภิกฺขุ, \\ สกุณคฺฆิ มกฺกโฏ สูโท,}{\csromangathabat{อมฺพปาลิ สโต ภิกฺขุ,}{สาลา’กุสลราสิ จ.} \\ \csromangathabat{สกุณคฺฆิ มกฺกโฏ สูโท,}{คิลาโน ภิกฺขุนุปสฺสโยติ.}}
\csromangathasetleft{อมฺพปาลิ สโต ภิกฺขุ, \\ สกุณคฺฆิ มกฺกโฏ สูโท,}
\csromangathagroup{\csromangathastanza{\csromanfolio{137}\csromangathabat{อมฺพปาลิ สโต ภิกฺขุ,}{สาลา’กุสลราสิ จ.} \\ \csromangathabat{สกุณคฺฆิ มกฺกโฏ สูโท,}{คิลาโน ภิกฺขุนุปสฺสโยติ.}}}

\csromanmatikaanchor{mtk.56}
\csromantocmarknopage{paragraph}{2. นาลนฺทวคฺค}{mtk.56}
\bookmark[dest={mtk.56},level=4]{2. นาลนฺทวคฺค}
\csromanheader{4}{2. \textbf{นาลนฺทวคฺค}}
\csromanmatikaanchor{mtk.57}
\csromantocmarkref{subparagraph}{1. มหาปุริสสุตฺต}{mtk.57}
\bookmark[dest={mtk.57},level=5]{1. มหาปุริสสุตฺต}
\csromanheader{5}{1. มหาปุริส\-สุตฺต}
\proseitem{377}{สาวตฺถิ\-นิทานํ. อถ โข อายสฺมา สาริปุตฺโต เยน ภควา เตนุปสงฺกมิ, อุปสงฺกมิตฺวา ภควนฺตํ อภิวาเทตฺวา เอกมนฺตํ นิสีทิ, เอกมนฺตํ นิสินฺโน โข อายสฺมา สาริปุตฺโต ภควนฺตํ เอตทโวจ “มหาปุริโส มหาปุริโสติ ภนฺเต วุจฺจติ, กิตฺตาวตา นุ โข ภนฺเต มหาปุริโส โหตี”ติ. วิมุตฺต\-จิตฺตตฺตา ขฺวาหํ สาริปุตฺต “มหาปุริโส”ติ วทามิ, อวิมุตฺต\-จิตฺตตฺตา “โน มหาปุริโส”ติ วทามิ.}

\prose{กถญฺจ สาริปุตฺต วิมุตฺตจิตฺโต โหติ, อิธ สาริปุตฺต ภิกฺขุ กาเย กายานุปสฺสี วิหรติ อาตาปี สมฺปชาโน สติมา วิเนยฺย โลเก อภิชฺฌา\-โทมนสฺสํ, ตสฺส กาเย กายานุปสฺสิโน วิหรโต จิตฺตํ วิรชฺชติ วิมุจฺจติ อนุปาทาย อาสเวหิ. เวทนาสุ ฯเปฯ. จิตฺเต ฯเปฯ. ธมฺเมสุ ธมฺมานุปสฺสี วิหรติ อาตาปี สมฺปชาโน สติมา วิเนยฺย โลเก อภิชฺฌา\-โทมนสฺสํ, ตสฺส ธมฺเมสุ ธมฺมา\-นุปสฺสิโน วิหรโต จิตฺตํ วิรชฺชติ วิมุจฺจติ อนุปาทาย อาสเวหิ. เอวํ โข สาริปุตฺต วิมุตฺตจิตฺโต โหติ, วิมุตฺต\-จิตฺตตฺตา ขฺวาหํ สาริปุตฺต “มหาปุริโส”ติ วทามิ, อวิมุตฺต\-จิตฺตตฺตา “โน มหาปุริโส”ติ วทามีติ.  ปฐมํ.}

\csromanmatikaanchor{mtk.58}
\csromantocmarkref{subparagraph}{2. นาลนฺทสุตฺต}{mtk.58}
\bookmark[dest={mtk.58},level=5]{2. นาลนฺทสุตฺต}
\csromanheader{5}{2. \textbf{นาลนฺทสุตฺต}}
\proseitem{378}{เอกํ สมยํ ภควา นาลนฺทายํ วิหรติ ปาวาริก\-มฺพวเน. อถ โข อายสฺมา สาริปุตฺโต เยน ภควา เตนุปสงฺกมิ, อุปสงฺกมิตฺวา ภควนฺตํ อภิวาเทตฺวา เอกมนฺตํ นิสีทิ, เอกมนฺตํ นิสินฺโน โข อายสฺมา สาริปุตฺโต ภควนฺตํ เอตทโวจ “เอวํปสนฺโน อหํ ภนฺเต ภควติ, น จาหุ \csromanfolio{138}น จ ภวิสฺสติ น เจตรหิ วิชฺชติ อญฺโญ สมโณ วา พฺราหฺมโณ วา ภควตา ภิยฺโยภิญฺญตโร ยทิทํ สมฺโพธิยนฺ”ติ. อุฬารา โข ตฺยายํ สาริปุตฺต อาสภี วาจา ภาสิตา, เอกํโส คหิโต, สีหนาโท นทิโต “เอวํปสนฺโน อหํ ภนฺเต ภควติ, น จาหุ น จ ภวิสฺสติ น เจตรหิ วิชฺชติ อญฺโญ สมโณ วา พฺราหฺมโณ วา ภควตา ภิยฺโยภิญฺญตโร ยทิทํ สมฺโพธิยนฺ”ติ.}

\prose{กิํ นุ เต สาริปุตฺต เยเต อเหสุํ อตีตมทฺธานํ อรหนฺโต สมฺมาสมฺพุทฺธา, สพฺเพเต ภควนฺโต เจตสา เจโต ปริจฺจ วิทิตา “เอวํสีลา เต ภควนฺโต อเหสุํ อิติ วา, เอวํธมฺมา เต ภควนฺโต อเหสุํ อิติ วา, เอวํปญฺญา เต ภควนฺโต อเหสุํ อิติ วา, เอวํวิหาริโน เต ภควนฺโต อเหสุํ อิติ วา, เอวํวิมุตฺตา เต ภควนฺโต อเหสุํ อิติ วา”ติ. โน เหตํ ภนฺเต.}

\prose{กิํ ปน เต สาริปุตฺต เย เต ภวิสฺสนฺติ อนาคตมทฺธานํ อรหนฺโต สมฺมาสมฺพุทฺธา, สพฺเพเต ภควนฺโต เจตสา เจโต ปริจฺจ วิทิตา “เอวํสีลา เต ภควนฺโต ภวิสฺสนฺติ อิติ วา, เอวํธมฺมา เต ภควนฺโต ภวิสฺสนฺติ อิติ วา, เอวํปญฺญา เต ภควนฺโต ภวิสฺสนฺติ อิติ วา, เอวํวิหาริโน เต ภควนฺโต ภวิสฺสนฺติ อิติ วา, เอวํวิมุตฺตา เต ภควนฺโต ภวิสฺสนฺติ อิติ วา”ติ. โน เหตํ ภนฺเต.}

\prose{กิํ ปน ตฺยาหํ\csromansharedfootnote{73ccba92ff890a6f}{กิํ ปน เต (สี)} สาริปุตฺต เอตรหิ อรหํ สมฺมาสมฺพุทฺโธ เจตสา เจโต ปริจฺจ วิทิโต “เอวํสีโล ภควา อิติ วา, เอวํธมฺโม ภควา อิติ วา, เอวํปญฺโญ ภควา อิติ วา, เอวํวิหารี ภควา อิติ วา, เอวํวิมุตฺโต ภควา อิติ วา”ติ. โน เหตํ ภนฺเต.}

\prose{เอตฺถ จ เต สาริปุตฺต อตีตา\-นาคต\-ปจฺจุปฺปนฺเนสุ อรหนฺเตสุ สมฺมา\-สมฺพุทฺเธสุ เจโตปริย\-ญาณํ\csromansharedfootnote{9a07c8b765f68538}{เจโตปริยายญาณํ (พหูสุ)} นตฺถิ, อถ กิญฺจรหิ ตฺยายํ สาริปุตฺต อุฬารา อาสภี วาจา ภาสิตา, เอกํโส คหิโต, สีหนาโท นทิโต “เอวํปสนฺโน อหํ ภนฺเต ภควติ, น จาหุ น จ ภวิสฺสติ น เจตรหิ วิชฺชติ อญฺโญ สมโณ วา พฺราหฺมโณ วา ภควตา ภิยฺโยภิญฺญตโร ยทิทํ สมฺโพธิยนฺ”ติ.}

\prose{\csromanfolio{139}น โข เม\csromansharedfootnote{3a0c709fa3543b4c}{น โข เม ตํ (สฺยา, กํ, ก)} ภนฺเต อตีตา\-นาคต\-ปจฺจุปฺปนฺเนสุ อรหนฺเตสุ สมฺมา\-สมฺพุทฺเธสุ เจโตปริย\-ญาณํ อตฺถิ, อปิ จ เม ธมฺมนฺวโย วิทิโต. เสยฺยถาปิ ภนฺเต รญฺโญ ปจฺจนฺติมํ นครํ ทฬฺหุทฺธาปํ\csromansharedfootnote{bf1bda518944d9ed}{ทฬฺหุทฺทาปํ (สี, อิ, ก), ทฬฺหทฺธาปํ (สฺยา, กํ)} ทฬฺห\-ปาการ\-โตรณํ เอกทฺวารํ, ตตฺรสฺส โทวาริโก ปณฺฑิโต พฺยตฺโต เมธาวี อญฺญาตานํ นิวาเรตา ญาตานํ ปเวเสตา. โส ตสฺส นครสฺส สมนฺตา อนุปริยายปถํ อนุกฺกมมาโน น ปสฺเสยฺย ปาการสนฺธิํ วา ปาการวิวรํ วา อนฺตมโส พิฬาร\-นิกฺขมน\-มตฺตมฺปิ. ตสฺส เอวมสฺส “เย โข เกจิ โอฬาริกา ปาณา อิมํ นครํ ปวิสนฺติ วา นิกฺขมนฺติ วา, สพฺเพเต อิมินาว ทฺวาเรน ปวิสนฺติ วา นิกฺขมนฺติ วา’ติ. เอวเมว โข เม ภนฺเต ธมฺมนฺวโย วิทิโต, เยปิ เต ภนฺเต อเหสุํ อตีตมทฺธานํ อรหนฺโต สมฺมาสมฺพุทฺธา, สพฺเพเต ภควนฺโต ปญฺจ นีวรเณ ปหาย เจตโส อุปกฺกิเลเส ปญฺญาย ทุพฺพลีกรเณ จตูสุ สติปฏฺฐาเนสุ สุปฺปติฏฺฐิต\-จิตฺตา สตฺต โพชฺฌงฺเค ยถาภูตํ ภาเวตฺวา อนุตฺตรํ สมฺมาสมฺโพธิํ อภิสมฺ\-พุชฺฌิํสุ. เยปิ เต ภนฺเต ภวิสฺสนฺติ อนาคตมทฺธานํ อรหนฺโต สมฺมาสมฺพุทฺธา, สพฺเพเต ภควนฺโต ปญฺจ นีวรเณ ปหาย เจตโส อุปกฺกิเลเส ปญฺญาย ทุพฺพลีกรเณ จตูสุ สติปฏฺฐาเนสุ สุปฺปติฏฺฐิต\-จิตฺตา สตฺต โพชฺฌงฺเค ยถาภูตํ ภาเวตฺวา อนุตฺตรํ สมฺมาสมฺโพธิํ อภิสมฺพุชฺฌิสฺสนฺติ. ภควาปิ ภนฺเต เอตรหิ อรหํ สมฺมาสมฺพุทฺโธ ปญฺจ นีวรเณ ปหาย เจตโส อุปกฺกิเลเส ปญฺญาย ทุพฺพลีกรเณ จตูสุ สติปฏฺฐาเนสุ สุปฺปติฏฺฐิต\-จิตฺโต สตฺต โพชฺฌงฺเค ยถาภูตํ ภาเวตฺวา อนุตฺตรํ สมฺมาสมฺโพธิํ อภิสมฺพุทฺโธติ.}

\prose{สาธุ สาธุ สาริปุตฺต, ตสฺมาติห ตฺวํ สาริปุตฺต อิมํ ธมฺม\-ปริยายํ อภิกฺขณํ ภาเสยฺยาสิ ภิกฺขูนํ ภิกฺขุนีนํ อุปาสกานํ อุปาสิกานํ. เยสมฺปิ หิ สาริปุตฺต โมฆปุริสานํ ภวิสฺสติ ตถาคเต กงฺขา วา วิมติ วา, เตสมฺปิมํ ธมฺม\-ปริยายํ สุตฺวา ยา ตถาคเต กงฺขา วา วิมติ วา, สา ปหียิสฺสตีติ.  ทุติยํ.}

\csromanmatikaanchor{mtk.59}
\csromantocmarkref{subparagraph}{3. จุนฺทสุตฺต}{mtk.59}
\bookmark[dest={mtk.59},level=5]{3. จุนฺทสุตฺต}
\csromanheader{5}{3. \textbf{จุนฺทสุตฺต}}
\proseitem{379}{เอกํ สมยํ ภควา สาวตฺถิยํ วิหรติ เชตวเน อนาถ\-ปิณฺฑิกสฺส อาราเม. เตน โข ปน สมเยน อายสฺมา สาริปุตฺโต \csromanfolio{140}มคเธสุ วิหรติ นาลกคามเก อาพาธิโก ทุกฺขิโต พาฬฺหคิลาโน. จุนฺโท จ สมณุทฺเทโส อายสฺมโต สาริปุตฺตสฺส อุปฏฺฐาโก โหติ.}

\prose{อถ โข อายสฺมา สาริปุตฺโต เตเนว อาพาเธน ปรินิพฺพายิ. อถ โข จุนฺโท สมณุทฺเทโส อายสฺมโต สาริปุตฺตสฺส ปตฺต\-จีวรมา\-ทาย เยน สาวตฺถิ เชตวนํ อนาถ\-ปิณฺฑิกสฺส อาราโม, เยนายสฺมา อานนฺโท เตนุปสงฺกมิ, อุปสงฺกมิตฺวา อายสฺมนฺตํ อานนฺทํ อภิวาเทตฺวา เอกมนฺตํ นิสีทิ, เอกมนฺตํ นิสินฺโน โข จุนฺโท สมณุทฺเทโส อายสฺมนฺตํ อานนฺทํ เอตทโวจ “อายสฺมา ภนฺเต สาริปุตฺโต ปรินิพฺพุโต, อิทมสฺส ปตฺตจีวรนฺ”ติ. อตฺถิ โข อิทํ อาวุโส จุนฺท กถาปาภตํ ภควนฺตํ ทสฺสนาย, อายามาวุโส จุนฺท, เยน ภควา เตนุปสงฺกมิสฺสาม, อุปสงฺกมิตฺวา ภควโต เอตมตฺถํ อาโรเจสฺสามาติ. “เอวํ ภนฺเต”ติ โข จุนฺโท สมณุทฺเทโส อายสฺมโต อานนฺทสฺส ปจฺจสฺโสสิ.}

\prose{อถ โข อายสฺมา จ อานนฺโท จุนฺโท จ สมณุทฺเทโส เยน ภควา เตนุ\-ปสงฺกมิํสุ, อุปสงฺกมิตฺวา ภควนฺตํ อภิวาเทตฺวา เอกมนฺตํ นิสีทิํสุ, เอกมนฺตํ นิสินฺโน โข อายสฺมา อานนฺโท ภควนฺตํ เอตทโวจ “อยํ ภนฺเต จุนฺโท สมณุทฺเทโส เอวมาห ‘อยสฺมา ภนฺเต สาริปุตฺโต ปรินิพฺพุโต, อิทมสฺส ปตฺตจีวรนฺ’ติ. อปิ จ เม ภนฺเต มธุรกชาโต วิย กาโย, ทิสาปิ เม น ปกฺขายนฺติ, ธมฺมาปิ มํ นปฺปฏิภนฺติ ‘อายสฺมา สาริปุตฺโต ปรินิพฺพุโต’ติ สุตฺวา”.}

\prose{กิํ นุ โข เต อานนฺท สาริปุตฺโต สีลกฺขนฺธํ วา อาทาย ปรินิพฺพุโต, สมาธิ\-กฺขนฺธํ วา อาทาย ปรินิพฺพุโต, ปญฺญากฺขนฺธํ วา อาทาย ปรินิพฺพุโต, วิมุตฺติ\-กฺขนฺธํ วา อาทาย ปรินิพฺพุโต, วิมุตฺติ\-ญาณ\-ทสฺสนกฺขนฺธํ วา อาทาย ปรินิพฺพุโตติ. น จ โข เม ภนฺเต อายสฺมา สาริปุตฺโต สีลกฺขนฺธํ วา อาทาย ปรินิพฺพุโต, สมาธิ\-กฺขนฺธํ วา ฯเปฯ ปญฺญากฺขนฺธํ วา ฯเปฯ วิมุตฺติ\-กฺขนฺธํ วา ฯเปฯ วิมุตฺติ\-ญาณ\-ทสฺสนกฺขนฺธํ วา อาทาย ปรินิพฺพุโต, อปิ จ เม ภนฺเต อายสฺมา สาริปุตฺโต โอวาทโก อโหสิ โอติณฺโณ วิญฺญาปโก สนฺทสฺสโก สมาทปโก สมุตฺเตชโก สมฺปหํสโก, อกิลาสุ ธมฺมเทสนาย, อนุคฺคาหโก สพฺรหฺมจารีนํ, ตํ มยํ อายสฺมโต สาริปุตฺตสฺส ธมฺโมชํ ธมฺมโภคํ ธมฺมานุคฺคหํ อนุสฺสรามาติ.}

\prose{\csromanfolio{141}นนุ ตํ อานนฺท มยา ปฏิกจฺเจว\csromansharedfootnote{aa443bc1a45a8a11}{ปฏิคจฺเจว (สี, อิ)} อกฺขาตํ, สพฺเพหิ ปิเยหิ มนาเปหิ นานาภาโว วินาภาโว อญฺญถาภาโว, ตํ กุเตตฺถ อานนฺท ลพฺภา “ยํ ตํ ชาตํ ภูตํ สงฺขตํ ปโลกธมฺมํ, ตํ วต มา ปลุชฺชี”ติ เนตํ ฐานํ วิชฺชติ. เสยฺยถาปิ อานนฺท มหโต รุกฺขสฺส ติฏฺฐโต สารวโต โย มหนฺตตโร ขนฺโธ โส ปลุชฺเชยฺย. เอวเมว โข อานนฺท มหโต ภิกฺขุ\-สํฆสฺส ติฏฺฐโต สารวโต สาริปุตฺโต ปรินิพฺพุโต, ตํ กุเตตฺถ อานนฺท ลพฺภา “ยํ ตํ ชาตํ ภูตํ สงฺขตํ ปโลกธมฺมํ, ตํ วต มา ปลุชฺชี”ติ เนตํ ฐานํ วิชฺชติ. ตสฺมาติหานนฺท อตฺตทีปา วิหรถ อตฺตสรณา อนญฺญสรณา ธมฺมทีปา ธมฺมสรณา อนญฺญสรณา.}

\prose{กถญฺจานนฺท ภิกฺขุ อตฺตทีโป วิหรติ อตฺตสรโณ อนญฺญสรโณ ธมฺมทีโป ธมฺมสรโณ อนญฺญสรโณ, อิธานนฺท ภิกฺขุ กาเย กายานุปสฺสี วิหรติ อาตาปี สมฺปชาโน สติมา วิเนยฺย โลเก อภิชฺฌา\-โทมนสฺสํ. เวทนาสุ ฯเปฯ. จิตฺเต ฯเปฯ. ธมฺเมสุ ธมฺมานุปสฺสี วิหรติ อาตาปี สมฺปชาโน สติมา วิเนยฺย โลเก อภิชฺฌา\-โทมนสฺสํ. เอวํ โข อานนฺท ภิกฺขุ อตฺตทีโป วิหรติ อตฺตสรโณ อนญฺญสรโณ ธมฺมทีโป ธมฺมสรโณ อนญฺญสรโณ. เย หิ เกจิ อานนฺท เอตรหิ วา มมจฺจเย วา อตฺตทีปา วิหริสฺสนฺติ อตฺตสรณา อนญฺญสรณา ธมฺมทีปา ธมฺมสรณา อนญฺญสรณา, ตมตคฺเค เมเต อานนฺท ภิกฺขู ภวิสฺสนฺติ, เย เกจิ สิกฺขากามาติ.  ตติยํ.}

\csromanmatikaanchor{mtk.60}
\csromantocmarkref{subparagraph}{4. อุกฺกเจลสุตฺต}{mtk.60}
\bookmark[dest={mtk.60},level=5]{4. อุกฺกเจลสุตฺต}
\csromanheader{5}{4. \textbf{อุกฺกเจลสุตฺต}}
\proseitem{380}{เอกํ สมยํ ภควา วชฺชีสุ วิหรติ อุกฺกเจลายํ คงฺคาย นทิยา ตีเร มหตา ภิกฺขุ\-สํเฆน สทฺธิํ อจิร\-ปรินิพฺพุเตสุ สาริปุตฺต\-โมคฺคลฺลาเนสุ. เตน โข ปน สมเยน ภควา ภิกฺขุ\-สํฆ\-ปริวุโต อชฺโฌกาเส นิสินฺโน โหติ.}

\prose{อถ โข ภควา ตุณฺหีภูตํ ภิกฺขุสํฆํ อนุวิโลเกตฺวา ภิกฺขู อามนฺเตสิ “อปิ มฺยายํ ภิกฺขเว ปริสา สุญฺญา วิย ขายติ ปรินิพฺพุเตสุ สาริปุตฺต\-โมคฺคลฺลาเนสุ. อสุญฺญา เม ภิกฺขเว ปริสา \csromanfolio{142}โหติ, อนเปกฺขา ตสฺสํ ทิสายํ โหติ, ยสฺสํ ทิสายํ สาริปุตฺต\-โมคฺคลฺลานา วิหรนฺติ. เย หิ เต ภิกฺขเว อเหสุํ อตีตมทฺธานํ อรหนฺโต สมฺมาสมฺพุทฺธา, เตสมฺปิ ภควนฺตานํ เอตปฺปรมํเยว\csromansharedfootnote{2b3c45c63d829f26}{เอตปรมํเยว (สี, สฺยา, กํ, อิ)} สาวกฺยุคํ อโหสิ, เสยฺยถาปิ มยฺหํ สาริปุตฺต\-โมคฺคลฺลานา. เยปิ เต ภิกฺขเว ภวิสฺสนฺติ อนาคตมทฺธานํ อรหนฺโต สมฺมาสมฺพุทฺธา, เตสมฺปิ ภควนฺตานํ เอตปฺปรมํเยว สาวกฺยุคํ ภวิสฺสติ, เสยฺยถาปิ มยฺหํ สาริปุตฺต\-โมคฺคลฺลานา. อจฺฉริยํ ภิกฺขเว สาวกานํ, อพฺภุตํ ภิกฺขเว สาวกานํ, สตฺถุ จ นาม สาสนกรา ภวิสฺสนฺติ โอวาทปฺปฏิกรา, จตุนฺนญฺจ ปริสานํ ปิยา ภวิสฺสนฺติ มนาปา ครุภาวนียา จ. อจฺฉริยํ ภิกฺขเว ตถาคตสฺส, อพฺภุตํ ภิกฺขเว ตถคตสฺส, เอวรูเปปิ นาม สาวกฺยุเค ปรินิพฺพุเต นตฺถิ ตถาคตสฺส โสโก วา ปริเทโว วา, ตํ กุเตตฺถ ภิกฺขเว ลพฺภา “ยํ ตํ ชาตํ ภูตํ สงฺขตํ ปโลกธมฺมํ, ตํ วต มา ปลุชฺชี”ติ เนตํ ฐานํ วิชฺชติ. เสยฺยถาปิ ภิกฺขเว มหโต รุกฺขสฺส ติฏฺฐโต สารวโต เย มหนฺตตรา ขนฺธา, เต ปลุชฺเชยฺยุํ. เอวเมว โข ภิกฺขเว มหโต ภิกฺขุ\-สํฆสฺส ติฏฺฐโต สารวโต สาริปุตฺต\-โมคฺคลฺลานา ปรินิพฺพุตา, ตํ กุเตตฺถ ภิกฺขเว ลพฺภา “ยํ ตํ ชาตํ ภูตํ สงฺขตํ ปโลกธมฺมํ, ตํ วต มา ปลุชฺชี”ติ เนตํ ฐานํ วิชฺชติ. ตสฺมาติห ภิกฺขเว อตฺตทีปา วิหรถ อตฺตสรณา อนญฺญสรณา ธมฺมทีปา ธมฺมสรณา อนญฺญสรณา.}

\prose{กถญฺจ ภิกฺขเว ภิกฺขุ อตฺตทีโป วิหรติ อตฺตสรโณ อนญฺญสรโณ ธมฺมทีโป ธมฺมสรโณ อนญฺญสรโณ, อิธ ภิกฺขเว ภิกฺขุ กาเย กายานุปสฺสี วิหรติ อาตาปี สมฺปชาโน สติมา วิเนยฺย โลเก อภิชฺฌา\-โทมนสฺสํ. เวทนาสุ ฯเปฯ. จิตฺเต ฯเปฯ. ธมฺเมสุ ธมฺมานุปสฺสี วิหรติ อาตาปี สมฺปชาโน สติมา วิเนยฺย โลเก อภิชฺฌา\-โทมนสฺสํ. เอวํ โข ภิกฺขเว ภิกฺขุ อตฺตทีโป วิหรติ อตฺตสรโณ อนญฺญสรโณ ธมฺมทีโป ธมฺมสรโณ อนญฺญสรโณ. เย หิ เกจิ ภิกฺขเว เอตรหิ วา มมจฺจเย วา อตฺตทีปา วิหริสฺสนฺติ อตฺตสรณา อนญฺญสรณา ธมฺมทีปา ธมฺมสรณา อนญฺญสรณา. ตมตคฺเค เมเต ภิกฺขเว ภิกฺขู ภวิสฺสนฺติ, เย เกจิ สิกฺขากามาติ.  จตุตฺถํ.}

\csromanheader{2}{5. \textbf{พาหิยสุตฺต}}
\csromanmatikaanchor{mtk.61}
\csromantocmarkref{subparagraph}{5. ภาหิยสุตฺต}{mtk.61}
\bookmark[dest={mtk.61},level=5]{5. ภาหิยสุตฺต}
\proseitem{381}{\csromanfolio{143}สาวตฺถิ\-นิทานํ. อถ โข อายสฺมา พาหิโย เยน ภควา เตนุปสงฺกมิ, อุปสงฺกมิตฺวา ภควนฺตํ อภิวาเทตฺวา เอกมนฺตํ นิสีทิ, เอกมนฺตํ นิสินฺโน โข อายสฺมา พาหิโย ภควนฺตํ เอตทโวจ “สาธุ เม ภนฺเต ภควา สํขิตฺเตน ธมฺมํ เทเสตุ, ยมหํ ภควโต ธมฺมํ สุตฺวา เอโก วูปกฏฺโฐ อปฺปมตฺโต อาตาปี ปหิตตฺโต วิหเรยฺยนฺ”ติ. ตสฺมาติห ตฺวํ พาหิย อาทิเมว วิโสเธหิ กุสเลสุ ธมฺเมสุ. โก จาทิ กุสลานํ ธมฺมานํ, สีลญฺจ สุวิสุทฺธํ ทิฏฺฐิ จ อุชุกา. ยโต จ โข เต พาหิย สีลญฺจ สุวิสุทฺธํ ภวิสฺสติ ทิฏฺฐิ จ อุชุกา, ตโต ตฺวํ พาหิย สีลํ นิสฺสาย สีเล ปติฏฺฐาย จตฺตาโร สติปฏฺฐาเน ภาเวยฺยาสิ.}

\prose{กตเม จตฺตาโร, อิธ ตฺวํ พาหิย กาเย กายานุปสฺสี วิหราหิ อาตาปี สมฺปชาโน สติมา วิเนยฺย โลเก อภิชฺฌา\-โทมนสฺสํ. เวทนาสุ ฯเปฯ. จิตฺเต ฯเปฯ. ธมฺเมสุ ธมฺมานุปสฺสี วิหราหิ อาตาปี สมฺปชาโน สติมา วิเนยฺย โลเก อภิชฺฌา\-โทมนสฺสํ. ยโต โข ตฺวํ พาหิย สีลํ นิสฺสาย สีเล ปติฏฺฐาย อิเม จตฺตาโร สติปฏฺฐาเน เอวํ ภาเวสฺสสิ, ตโต ตุยฺหํ พาหิย ยา รตฺติ วา ทิวโส วา อาคมิสฺสติ, วุทฺธิเยว ปาฏิกงฺขา กุสเลสฺสุ ทมฺเมสุ, โน ปริหานีติ.}

\prose{อถ โข อายสฺมา พาหิโย ภควโต ภาสิตํ อภินนฺทิตฺวา อนุโมทิตฺวา อุฏฺฐายาสนา ภควนฺตํ อภิวาเทตฺวา ปทกฺขิณํ กตฺวา ปกฺกามิ. อถ โข อายสฺมา พาหิโย เอโก วูปกฏฺโฐ อปฺปมตฺโต อาตาปี ปหิตตฺโต วิหรนฺโต นจิรสฺเสว ยสฺสตฺถาย กุลปุตฺตา สมฺมเทว อคารสฺมา อนคาริยํ ปพฺพชนฺติ, ตทนุตฺตรํ พฺรหฺมจริย\-ปริโยสานํ ทิฏฺเฐว ธมฺเม สยํ อภิญฺญา สจฺฉิกตฺวา อุปสมฺปชฺช วิหาสิ, “ขีณา ชาติ, วุสิตํ พฺรหฺมจริยํ, กตํ กรณียํ, นาปรํ อิตฺถตฺตายา”ติ อพฺภญฺญาสิ. อญฺญตโร จ ปนายสฺมา พาหิโย อรหตํ อโหสีติ.  ปญฺจมํ.}

\csromanmatikaanchor{mtk.62}
\csromantocmarkref{subparagraph}{6. อุตฺติยสุตฺต}{mtk.62}
\bookmark[dest={mtk.62},level=5]{6. อุตฺติยสุตฺต}
\csromanheader{5}{6. \textbf{อุตฺติยสุตฺต}}
\proseitem{382}{สาวตฺถิ\-นิทานํ. อถ โข อายสฺมา อุตฺติโย เยน ภควา เตนุปสงฺกมิ ฯเปฯ เอกมนฺตํ นิสินฺโน โข อายสฺมา อุตฺติโย ภควนฺตํ เอตทโวจ “สาธุ เม ภนฺเต ภควา สํขิตฺเตน ธมฺมํ เทเสตุ, ยมหํ \csromanfolio{144}ภควโต ธมฺมํ สุตฺวา เอโก วูปกฏฺโฐ อปฺปมตฺโต อาตาปี ปหิตตฺโต วิหเรยฺยนฺ”ติ. ตสฺมาติห ตฺวํ อุตฺติย อาทิเมว วิโสเธหิ กุสเลสุ ธมฺเมสุ. โก จาทิ กุสลานํ ธมฺมานํ, สีลญฺจ สุวิสุทฺธํ ทิฏฺฐิ จ อุชุกา. ยโต จ โข เต อุตฺติย สีลญฺจ สุวิสุทฺธํ ภวิสฺสติ ทิฏฺฐิ จ อุชุกา, ตโต ตฺวํ อุตฺติย สีลํ นิสฺสาย สีเล ปติฏฺฐาย จตฺตาโร สติปฏฺฐาเน ภาเวยฺยาสิ.}

\prose{กตเม จตฺตาโร, อิธ ตฺวํ อุตฺติย กาเย กายานุปสฺสี วิหราหิ อาตาปี สมฺปชาโน สติมา วิเนยฺย โลเก อภิชฺฌา\-โทมนสฺสํ. เวทนาสุ ฯเปฯ. จิตฺเต ฯเปฯ. ธมฺเมสุ ธมฺมานุปสฺสี วิหราหิ อาตาปี สมฺปชาโน สติมา วิเนยฺย โลเก อภิชฺฌา\-โทมนสฺสํ. ยโต โข ตฺวํ อุตฺติย สีลํ นิสฺสาย สีเล ปติฏฺฐาย อิเม จตฺตาโร สติปฏฺฐาเน เอวํ ภาเวสฺสสิ, ตโต ตฺวํ อุตฺติย คมิสฺสสิ มจฺจุเธยฺยสฺส ปารนฺติ.}

\prose{อถ โข อายสฺมา อุตฺติโย ภควโต ภาสิตํ อภินนฺทิตฺวา อนุโมทิตฺวา อุฏฺฐายาสนา ภควนฺตํ อภิวาเทตฺวา ปทกฺขิณํ กตฺวา ปกฺกามิ. อถ โข อายสฺมา อุตฺติโย เอโก วูปกฏฺโฐ อปฺปมตฺโต อาตาปี ปหิตตฺโต วิหรนฺโต นจิรสฺเสว ยสฺสตฺถาย กุลปุตฺตา สมฺมเทว อคารสฺมา อนคาริยํ ปพฺพชนฺติ, ตทนุตฺตรํ พฺรหฺมจริย\-ปริโยสานํ ทิฏฺเฐว ธมฺเม สยํ อภิญฺญา สจฺฉิกตฺวา อุปสมฺปชฺช วิหาสิ, “ขีณา ชาติ, วุสิตํ พฺรหฺมจริยํ, กตํ กรณียํ, นาปรํ อิตฺถตฺตายา”ติ อพฺภญฺญาสิ. อญฺญตโร จ ปนายสฺมา อุตฺติโย อรหตํ อโหสีติ.  ฉฏฺฐํ.}

\csromanmatikaanchor{mtk.63}
\csromantocmarkref{subparagraph}{7. อริยสุตฺต}{mtk.63}
\bookmark[dest={mtk.63},level=5]{7. อริยสุตฺต}
\csromanheader{5}{7. \textbf{อริยสุตฺต}}
\proseitem{383}{จตฺตาโรเม ภิกฺขเว สติปฏฺฐานา ภาวิตา พหุลีกตา อริยา นิยฺยานิกา นิยฺยนฺติ ตกฺกรสฺส สมฺมา ทุกฺขกฺขยาย. กตเม จตฺตาโร, อิธ ภิกฺขเว ภิกฺขุ กาเย กายานุปสฺสี วิหรติ อาตาปี สมฺปชาโน สติมา วิเนยฺย โลเก อภิชฺฌา\-โทมนสฺสํ. เวทนาสุ ฯเปฯ. จิตฺเต ฯเปฯ. ธมฺเมสุ ธมฺมานุปสฺสี วิหรติ อาตาปี สมฺปชาโน สติมา วิเนยฺย โลเก อภิชฺฌา\-โทมนสฺสํ. อิเม โข ภิกฺขเว จตฺตาโร สติปฏฺฐานา ภาวิตา พหุลีกตา อริยา นิยฺยานิกา นิยฺยนฺติ ตกฺกรสฺส สมฺมา ทุกฺขกฺขยายาติ.  สตฺตมํ.}

\csromanmatikaanchor{mtk.64}
\csromantocmarkref{subparagraph}{8. พฺรหฺมสุตฺต}{mtk.64}
\bookmark[dest={mtk.64},level=5]{8. พฺรหฺมสุตฺต}
\csromanheader{5}{8. \textbf{พฺรหฺมสุตฺต}}
\proseitem{384}{\csromanfolio{145}เอกํ สมยํ ภควา อุรุเวลายํ วิหรติ นชฺชา เนรญฺชราย ตีเร อชปาล\-นิโคฺรเธ ปฐมา\-ภิสมฺพุทฺโธ. อถ โข ภควโต รโหคตสฺส ปฏิสลฺลีนสฺส เอวํ เจตโส ปริวิตกฺโก อุทปาทิ\csromandash{}เอกายนฺวายํ มคฺโค สตฺตานํ วิสุทฺธิยา โสก\-ปริเทวานํ สมติกฺกมาย ทุกฺข\-โทมนสฺสานํ อตฺถงฺคมาย ญายสฺส อธิคมาย นิพฺพานสฺส สจฺฉิกิริยาย, ยทิทํ จตฺตาโร สติปฏฺฐานา.}

\prose{กตเม จตฺตาโร, กาเย วา ภิกฺขุ กายานุปสฺสี วิหเรยฺย อาตาปี สมฺปชาโน สติมา วิเนยฺย โลเก อภิชฺฌา\-โทมนสฺสํ. เวทนาสุ วา ภิกฺขุ ฯเปฯ. จิตฺเต วา ภิกฺขุ ฯเปฯ. ธมฺเมสุ วา ภิกฺขุ ธมฺมานุปสฺสี วิหเรยฺย อาตาปี สมฺปชาโน สติมา วิเนยฺย โลเก อภิชฺฌา\-โทมนสฺสํ. เอกายนฺวายํ มคฺโค สตฺตานํ วิสุทฺธิยา โสก\-ปริเทวานํ สมติกฺกมาย ทุกฺข\-โทมนสฺสานํ อตฺถงฺคมาย ญายสฺส อธิคมาย นิพฺพานสฺส สจฺฉิกิริยาย, ยทิทํ จตฺตาโร สติปฏฺฐานาติ.}

\prose{อถ โข พฺรหฺมา สหมฺปติ ภควโต เจตสา เจโต\-ปริวิตกฺกม\-ญฺญาย เสยฺยถาปิ นาม พลวา ปุริโส สมิญฺชิตํ วา พาหํ ปสาเรยฺย, ปสาริตํ วา พาหํ สมิญฺเชยฺย. เอวเมว โข พฺรหฺมโลเก อนฺตรหิโต ภควโต ปุรโต ปาตุรโหสิ. อถ โข พฺรหฺมา สหมฺปติ เอกํสํ อุตฺตราสงฺคํ กริตฺวา เยน ภควา เตนญฺชลิํ ปณาเมตฺวา ภควนฺตํ เอตทโวจ\csromandash{} เอวเมตํ ภควา, เอวเมตํ สุคต, เอกายนฺวายํ ภนฺเต มคฺโค สตฺตานํ วิสุทฺธิยา โสก\-ปริเทวานํ สมติกฺกมาย ทุกฺข\-โทมนสฺสานํ อตฺถงฺคมาย ญายสฺส อธิคมาย นิพฺพานสฺส สจฺฉิกิริยาย, ยทิทํ จตฺตาโร สติปฏฺฐานา.}

\prose{กตเม จตฺตาโร, กาเย วา ภนฺเต ภิกฺขุ กายานุปสฺสี วิหเรยฺย อาตาปี สมฺปชาโน สติมา วิเนยฺย โลเก อภิชฺฌา\-โทมนสฺสํ. เวทนาสุ วา ภนฺเต ภิกฺขุ ฯเปฯ. จิตฺเต วา ภนฺเต ภิกฺขุ ฯเปฯ. ธมฺเมสุ วา ภนฺเต ภิกฺขุ ธมฺมานุปสฺสี วิหเรยฺย อาตาปี สมฺปชาโน สติมา วิเนยฺย โลเก อภิชฺฌา\-โทมนสฺสํ. เอกายนฺวายํ ภนฺเต มคฺโค สตฺตานํ วิสุทฺธิยา โสก\-ปริเทวานํ สมติกฺกมาย ทุกฺข\-โทมนสฺสานํ \csromanfolio{146}อตฺถงฺคมาย ญายสฺส อธิคมาย นิพฺพานสฺส สจฺฉิกิริยาย, ยทิทํ จตฺตาโร สติปฏฺฐานาติ.}

\prose{อิทมโวจ พฺรหฺมา สหมฺปติ, อิทํ วตฺวา อถาปรํ เอตทโวจ\csromandash{}}

\setlength{\csromangathapagewidth}{0pt}
\setlength{\csromangathawakwidth}{0pt}
\csromangathameasuregroup{“เอกายนํ ชาติขยนฺตทสฺสี,}{\csromangathabat{“เอกายนํ ชาติขยนฺตทสฺสี,}{มคฺคํ ปชานาติ หิตานุกมฺปี.}}
\csromangathasetleft{“เอกายนํ ชาติขยนฺตทสฺสี,}
\csromangathagroup{\csromangathastanza{\csromangathabat{“เอกายนํ ชาติ\-ขยนฺต\-ทสฺสี,}{มคฺคํ ปชานาติ หิตานุกมฺปี.}}}

\prose{เอเตน มคฺเคน ตริํสุ ปุพฺเพ, ตริสฺสนฺติ เย จ ตรนฺติ โอฆนฺ”ติ. อฏฺฐมํ.}

\csromanmatikaanchor{mtk.65}
\csromantocmarkref{subparagraph}{9. เสทกสุตฺต}{mtk.65}
\bookmark[dest={mtk.65},level=5]{9. เสทกสุตฺต}
\csromanheader{5}{9. \textbf{เสทกสุตฺต}}
\proseitem{385}{เอกํ สมยํ ภควา สุมฺเภสุ วิหรติ เสทกํ นาม สุมฺภานํ นิคโม. ตตฺร โข ภควา ภิกฺขู อามนฺเตสิ “ภูตปุพฺพํ ภิกฺขเว จณฺฑาลวํสิโก จณฺฑาลวํสํ อุสฺสาเปตฺวา เมทกถาลิกํ อนฺเตวาสิํ อามนฺเตสิ ‘เอหิ ตฺวํ สมฺม เมทกถาลิเก จณฺฑาลวํสํ อภิรุหิตฺวา มม อุปริขนฺเธ ติฏฺฐาหี’ติ”. “เอวํ อาจริยา”ติ โข ภิกฺขเว เมทกถาลิกา อนฺเตวาสี จณฺฑาล\-วํสิ\-กสฺส ปฏิสฺสุตฺวา จณฺฑาลวํสํ อภิรุหิตฺวา อาจริยสฺส อุปริขนฺเธ อฏฺฐาสิ. อถ โข ภิกฺขเว จณฺฑาลวํสิโก เมทกถาลิกํ อนฺเตวาสิํ เอตทโวจ “ตฺวํ สมฺม เมทกถาลิเก มมํ รกฺข, อหํ ตํ รกฺขิสฺสามิ, เอวํ มยํ อญฺญมญฺญํ คุตฺตา อญฺญมญฺญํ รกฺขิตา สิปฺปานิ เจว ทสฺเสสฺสาม, ลาภญฺจ\csromansharedfootnote{e362e75fe64cef4a}{ลาเภ จ (สี)} ลจฺฉาม, โสตฺถินา จ จณฺฑาลวํสา โอโรหิสฺสามา”ติ. เอวํ วุตฺเต ภิกฺขเว เมทกถาลิกา อนฺเตวาสี จณฺฑาล\-วํสิกํ เอตทโวจ “น โข ปเนตํ อาจริย เอวํ ภวิสฺสติ, ตฺวํ อาจริย อตฺตานํ รกฺข, อหํ อตฺตานํ รกฺขิสฺสามิ, เอวํ มยํ อตฺตคุตฺตา อตฺตรกฺขิตา สิปฺปานิ เจว ทสฺเสสฺสาม, ลาภญฺจ ลจฺฉาม, โสตฺถินา จ จณฺฑาลวํสา โอโรหิสฺสามา”ติ.}

\prose{“โส ตตฺถ ญาโย”ติ ภควา เอตทโวจ, ยถา เมทกถาลิกา อนฺเตวาสี อาจริยํ อโวจ “อตฺตานํ ภิกฺขเว รกฺขิสฺสามี”ติ สติปฏฺฐานํ เสวิตพฺพํ, “ปรํ รกฺขิสฺสามี”ติ สติปฏฺฐานํ เสวิตพฺพํ. อตฺตานํ ภิกฺขเว รกฺขนฺโต ปรํ รกฺขติ, ปรํ รกฺขนฺโต อตฺตานํ รกฺขติ.}

\prose{\csromanfolio{147}กถญฺจ ภิกฺขเว อตฺตานํ รกฺขนฺโต ปรํ รกฺขติ, อาเสวนาย ภาวนาย พหุลีกมฺเมน. เอวํ โข ภิกฺขเว อตฺตานํ รกฺขนฺโต ปรํ รกฺขติ. กถญฺจ ภิกฺขเว ปรํ รกฺขนฺโต อตฺตานํ รกฺขติ, ขนฺติยา อวิหิํสาย เมตฺตจิตฺตตาย อนุทยตาย. เอวํ โข ภิกฺขเว ปรํ รกฺขนฺโต อตฺตานํ รกฺขติ. “อตฺตานํ ภิกฺขเว รกฺขิสฺสามี”ติ สติปฏฺฐานํ เสวิตพฺพํ, “ปรํ รกฺขิสฺสามี”ติ สติปฏฺฐานํ เสวิตพฺพํ, อตฺตานํ ภิกฺขเว รกฺขนฺโต ปรํ รกฺขติ, ปรํ รกฺขนฺโต อตฺตานํ รกฺขตีติ.  นวมํ.}

\csromanmatikaanchor{mtk.66}
\csromantocmarkref{subparagraph}{10. ชนปทกลฺยาณีสุตฺต}{mtk.66}
\bookmark[dest={mtk.66},level=5]{10. ชนปทกลฺยาณีสุตฺต}
\csromanheader{5}{10. ชนปท\-กลฺยาณีสุตฺต}
\proseitem{386}{เอวํ เม สุตํ\csromandash{}เอกํ สมยํ ภควา สุมฺเภสุ วิหรติ เสทกํ นาม สุมฺภานํ นิคโม. ตตฺร โข ภควา ภิกฺขู อามนฺเตสิ “ภิกฺขโว”ติ. “ภทนฺเต”ติ เต ภิกฺขู ภควโต ปจฺจสฺโสสุํ. ภควา เอตทโวจ\csromandash{}}

\prose{เสยฺยถาปิ ภิกฺขเว “ชนปท\-กลฺยาณี ชนปท\-กลฺยาณี”ติ โข ภิกฺขเว มหาชนกาโย สนฺนิปเตยฺย, สา โข ปนสฺส ชนปท\-กลฺยาณี ปรมปาสาวินี นจฺเจ, ปรมปาสาวินี คีเต. “ชนปท\-กลฺยาณี นจฺจติ คายตี”ติ โข ภิกฺขเว ภิยฺโยโส มตฺตาย มหาชนกาโย สนฺนิปเตยฺย, อถ ปุริโส อาคจฺเฉยฺย ชีวิตุกาโม อมริตุกาโม สุขกาโม ทุกฺข\-ปฺปฏิกูโล. ตเมนํ เอวํ วเทยฺย “อยํ เต อมฺโภ ปุริส สมติตฺติโก เตลปตฺโต อนฺตเรน จ มหาสมชฺชํ อนฺตเรน จ ชนปท\-กลฺยาณิยา ปริหริตพฺโพ, ปุริโส จ เต อุกฺขิตฺตาสิโก ปิฏฺฐิโต ปิฏฺฐิโต อนุพนฺธิสฺสติ “ยตฺเถว นํ โถกมฺปิ ฉฑฺเฑสฺสติ, ตตฺเถว เต สิโร ปาเตสฺสตี”ติ. ตํ กิํ มญฺญถ ภิกฺขเว, อปิ นุ โส ปุริโส อมุํ เตลปตฺตํ อมนสิกริตฺวา พหิทฺธา ปมาทํ อาหเรยฺยาติ. โน เหตํ ภนฺเต.}

\prosewithclosermidruled{อุปมา โข มฺยายํ ภิกฺขเว กตา อตฺถสฺส วิญฺญาปนาย, อยํ เจเวตฺต อตฺโถ, “สมติตฺติโก เตลปตฺโต”ติ โข ภิกฺขเว กายคตาย เอตํ สติยา อธิวจนํ. ตสฺมาติห ภิกฺขเว เอวํ สิกฺขิตพฺพํ “กายคตา สติ โน ภาวิตา ภวิสฺสติ พหุลีกตา ยานีกตา \csromanfolio{148}วตฺถุกตา อนุฏฺฐิตา ปริจิตา สุสมารทฺธา”ติ. เอวํ หิ โข ภิกฺขเว สิกฺขิตพฺพนฺติ.  ทสมํ.}{นาลนฺทวคฺโค ทุติโย.}
\csromancenter{\textbf{ตสฺสุทฺทานํ}}
\setlength{\csromangathapagewidth}{0pt}
\setlength{\csromangathawakwidth}{0pt}
\csromangathameasuregroup{มหาปุริโส นาลนฺทํ, \\ อุตฺติโย อริโย พฺรหฺมา,}{\csromangathabat{มหาปุริโส นาลนฺทํ,}{จุนฺโท เจลญฺจ พาหิโย.} \\ \csromangathabat{อุตฺติโย อริโย พฺรหฺมา,}{เสทกํ ชนปเทน จาติ.}}
\csromangathasetleft{มหาปุริโส นาลนฺทํ, \\ อุตฺติโย อริโย พฺรหฺมา,}
\csromangathagroup{\csromangathastanza{\csromangathabat{มหาปุริโส นาลนฺทํ,}{จุนฺโท เจลญฺจ พาหิโย.} \\ \csromangathabat{อุตฺติโย อริโย พฺรหฺมา,}{เสทกํ ชนปเทน จาติ.}}}

\csromanmatikaanchor{mtk.67}
\csromantocmarknopage{paragraph}{3. สีลฏฺฐิติวคฺค}{mtk.67}
\bookmark[dest={mtk.67},level=4]{3. สีลฏฺฐิติวคฺค}
\csromanheader{4}{3. สีลฏฺฐิติ\-วคฺค}
\csromanmatikaanchor{mtk.68}
\csromantocmarkref{subparagraph}{1. สีลสุตฺต}{mtk.68}
\bookmark[dest={mtk.68},level=5]{1. สีลสุตฺต}
\csromanheader{5}{1. \textbf{สีลสุตฺต}}
\proseitem{387}{เอวํ เม สุตํ\csromandash{}เอกํ สมยํ อายสฺมา จ อานนฺโท อายสฺมา จ ภทฺโท ปาฏลิปุตฺเต วิหรนฺติ กุกฺกุฏาราเม. อถ โข อายสฺมา ภทฺโท สายนฺหสมยํ ปฏิสลฺลานา วุฏฺฐิโต เยนายสฺมา อานนฺโท เตนุปสงฺกมิ, อุปสงฺกมิตฺวา อายสฺมตา อานนฺเทน สทฺธิํ สมฺโมทิ, สมฺโมทนียํ กถํ สารณียํ วีติสาเรตฺวา เอกมนฺตํ นิสีทิ, เอกมนฺตํ นิสินฺโน โข อายสฺมา ภทฺโท อายสฺมนฺตํ อานนฺทํ เอตทโวจ “ยานิมานิ อาวุโส อานนฺท กุสลานิ สีลานิ วุตฺตานิ ภควตา, อิมานิ กุสลานิ สีลานิ กิมตฺถิยานิ วุตฺตานิ ภควตา”ติ.}

\prose{สาธุ สาธุ อาวุโส ภทฺท, ภทฺทโก โข เต อาวุโส ภทฺท อุมฺมงฺโค\csromansharedfootnote{c305010ce3dddcba}{อุมฺมคฺโค (สี, สฺยา, กํ)}, ภทฺทกํ ปฏิภานํ, กลฺยาณี ปริปุจฺฉา, เอวํ หิ ตฺวํ อาวุโส ภทฺท ปุจฺฉสิ “ยานิมานิ อาวุโส อานนฺท กุสลานิ สีลานิ วุตฺตานิ ภควตา, อิมานิ กุสลานิ สีลานิ กิมตฺถิยานิ วุตฺตานิ ภควตา”ติ. เอวมาวุโสติ. ยานิมานิ อาวุโส ภทฺท กุสลานิ สีลานิ วุตฺตานิ ภควตา, อิมานิ กุสลานิ สีลานิ ยาวเทว จตุนฺนํ สติปฏฺฐานานํ ภาวนาย วุตฺตานิ ภควตา.}

\prose{กตเมสํ จตุนฺนํ, อิธาวุโส ภิกฺขุ กาเย กายานุปสฺสี วิหรติ อาตาปี สมฺปชาโน สติมา วิเนยฺย โลเก อภิชฺฌา\-โทมนสฺสํ.}

\prose{\csromanfolio{149}เวทนาสุ ฯเปฯ. จิตฺเต ฯเปฯ. ธมฺเมสุ ธมฺมานุปสฺสี วิหรติ อาตาปี สมฺปชาโน สติมา วิเนยฺย โลเก อภิชฺฌา\-โทมนสฺสํ. ยานิมานิ อาวุโส ภทฺท กุสลานิ สีลานิ วุตฺตานิ ภควตา, อิมานิ กุสลานิ สีลานิ ยาวเทว อิเมสํ จตุนฺนํ สติปฏฺฐานานํ ภาวนาย วุตฺตานิ ภควตาติ.  ปฐมํ.}

\csromanmatikaanchor{mtk.69}
\csromantocmarkref{subparagraph}{2. จิรฏฺฐิติสุตฺต}{mtk.69}
\bookmark[dest={mtk.69},level=5]{2. จิรฏฺฐิติสุตฺต}
\csromanheader{5}{2. จิรฏฺฐิติ\-สุตฺต}
\proseitem{388}{ตํเยว นิทานํ. เอกมนฺตํ นิสินฺโน โข อายสฺมา ภทฺโท อายสฺมนฺตํ อานนฺทํ เอตทโวจ “โก นุ โข อาวุโส อานนฺท เหตุ โก ปจฺจโย, เยน ตถาคเต ปรินิพฺพุเต สทฺธมฺโม น จิรฏฺฐิติโก โหติ, โก ปนาวุโส อานนฺท เหตุ โก ปจฺจโย, เยน ตถาคเต ปรินิพฺพุเต สทฺธมฺโม จิรฏฺฐิติโก โหตี”ติ.}

\prose{สาธุ สาธุ อาวุโส ภทฺท, ภทฺทโก โข เต อาวุโส ภทฺท อุมฺมงฺโค, ภทฺทกํ ปฏิภานํ, กลฺยาณี ปริปุจฺฉา, เอวํ หิ ตฺวํ อาวุโส ภทฺท ปุจฺฉสิ “โก นุ โข อาวุโส อานนฺท เหตุ โก ปจฺจโย, เยน ตถาคเต ปรินิพฺพุเต สทฺธมฺโม น จิรฏฺฐิติโก โหติ, โก ปนาวุโส อานนฺท เหตุ โก ปจฺจโย, เยน ตถาคเต ปรินิพฺพุเต สทฺธมฺโม จิรฏฺฐิติโก โหตี”ติ. เอวมาวุโสติ. จตุนฺนํ โข อาวุโส สติปฏฺฐานานํ อภาวิตตฺตา อพหุลีกตตฺตา ตถาคเต ปรินิพฺพุเต สทฺธมฺโม น จิรฏฺฐิติโก โหติ. จตุนฺนญฺจ โข อาวุโส สติปฏฺฐานานํ ภาวิตตฺตา พหุลีกตตฺตา ตถาคเต ปรินิพฺพุเต สทฺธมฺโม จิรฏฺฐิติโก โหติ.}

\prose{กตเมสํ จตุนฺนํ, อิธาวุโส ภิกฺขุ กาเย กายานุปสฺสี วิหรติ อาตาปี สมฺปชาโน สติมา วิเนยฺย โลเก อภิชฺฌา\-โทมนสฺสํ. เวทนาสุ ฯเปฯ. จิตฺเต ฯเปฯ. ธมฺเมสุ ธมฺมานุปสฺสี วิหรติ อาตาปี สมฺปชาโน สติมา วิเนยฺย โลเก อภิชฺฌา\-โทมนสฺสํ. อิเมสํ โข อาวุโส จตุนฺนํ สติปฏฺฐานานํ อภาวิตตฺตา อพหุลีกตตฺตา ตถาคเต ปรินิพฺพุเต สทฺธมฺโม น จิรฏฺฐิติโก โหติ. อิเมสญฺจ โข อาวุโส จตุนฺนํ สติปฏฺฐานานํ ภาวิตตฺตา พหุลีกตตฺตา ตถาคเต ปรินิพฺพุเต สทฺธมฺโม จิรฏฺฐิติโก โหตีติ.  ทุติยํ.}

\csromanheader{2}{3. \textbf{ปริหานสุตฺต}}
\csromanmatikaanchor{mtk.70}
\csromantocmarkref{subparagraph}{3-4. ปริหานสุตฺตาทีนิ}{mtk.70}
\bookmark[dest={mtk.70},level=5]{3-4. ปริหานสุตฺตาทีนิ}
\proseitem{389}{\csromanfolio{150}เอกํ สมยํ อายสฺมา จ อานนฺโท อายสฺมา จ ภทฺโท ปาฏลิปุตฺเต วิหรนฺติ กุกฺกุฏาราเม. อถ โข อายสฺมา ภทฺโท สายนฺหสมยํ ปฏิสลฺลานา วุฏฺฐิโต เยนายสฺมา อานนฺโท เตนุปสงฺกมิ, อุปสงฺกมิตฺวา อายสฺมตา อานนฺเทน สทฺธิํ สมฺโมทิ, สมฺโมทนียํ กถํ สารณียํ วีติสาเรตฺวา เอกมนฺตํ นิสีทิ, เอกมนฺตํ นิสินฺโน โข อายสฺมา ภทฺโท อายสฺมนฺตํ อานนฺทํ เอตทโวจ “โก นุ โข อาวุโส อานนฺท เหตุ โก ปจฺจโย, เยน สทฺธมฺม\-ปริหานํ โหติ, โก นุ โข อาวุโส อานนฺท เหตุ โก ปจฺจโย, เยน สทฺธมฺม\-อปริหานํ โหตี”ติ.}

\prose{สาธุ สาธุ อาวุโส ภทฺท, ภทฺทโก โข เต อาวุโส ภทฺท อุมฺมงฺโค, ภทฺทกํ ปฏิภานํ, กลฺยาณี ปริปุจฺฉา, เอวํ หิ ตฺวํ อาวุโส ภทฺท ปุจฺฉสิ “โก นุ โข อาวุโส อานนฺท เหตุ โก ปจฺจโย, เยน สทฺธมฺม\-ปริหานํ โหติ, โก ปนาวุโส อานนฺท เหตุ โก ปจฺจโย, เยน สทฺธมฺม\-อปริหานํ โหตี”ติ. เอวมาวุโสติ. จตุนฺนํ โข อาวุโส สติปฏฺฐานานํ อภาวิตตฺตา อพหุลีกตตฺตา สทฺธมฺม\-ปริหานํ โหติ. จตุนฺนญฺจ โข อาวุโส สติปฏฺฐานานํ ภาวิตตฺตา พหุลีกตตฺตา สทฺธมฺม- อปริหานํ โหติ.}

\prose{กตเมสํ จตุนฺนํ, อิธาวุโส ภิกฺขุ กาเย กายานุปสฺสี วิหรติ อาตาปี สมฺปชาโน สติมา วิเนยฺย โลเก อภิชฺฌา\-โทมนสฺสํ. เวทนาสุ ฯเปฯ. จิตฺเต ฯเปฯ. ธมฺเมสุ ธมฺมานุปสฺสี วิหรติ อาตาปี สมฺปชาโน สติมา วิเนยฺย โลเก อภิชฺฌา\-โทมนสฺสํ. อิเมสํ โข อาวุโส จตุนฺนํ สติปฏฺฐานานํ อภาวิตตฺตา อพหุลีกตตฺตา สทฺธมฺม\-ปริหานํ โหติ. อิเมสญฺจ โข อาวุโส จตุนฺนํ สติปฏฺฐานานํ ภาวิตตฺตา พหุลีกตตฺตา สทฺธมฺม\-อปริหานํ โหตีติ.  ตติยํ.}

\csromanheader{2}{4. \textbf{สุทฺธสุตฺต}}
\proseitem{390}{สาวตฺถิ\-นิทานํ. จตฺตาโรเม ภิกฺขเว สติปฏฺฐานา. กตเม จตฺตาโร, อิธ ภิกฺขเว ภิกฺขุ กาเย กายานุปสฺสี วิหรติ อาตาปี สมฺปชาโน สติมา วิเนยฺย โลเก อภิชฺฌา\-โทมนสฺสํ. เวทนาสุ ฯเปฯ.}

\prose{\csromanfolio{151}จิตฺเต ฯเปฯ. ธมฺเมสุ ธมฺมานุปสฺสี วิหรติ อาตาปี สมฺปชาโน สติมา วิเนยฺย โลเก อภิชฺฌา\-โทมนสฺสํ. อิเม โข ภิกฺขเว จตฺตาโร สติปฏฺฐานาติ.  จตุตฺถํ.}

\csromanmatikaanchor{mtk.71}
\csromantocmarkref{subparagraph}{5. อญฺญตรพฺราหฺมณสุตฺต}{mtk.71}
\bookmark[dest={mtk.71},level=5]{5. อญฺญตรพฺราหฺมณสุตฺต}
\csromanheader{6}{5. อญฺญตร\-พฺราหฺมณ\-สุตฺต}
\proseitem{391}{เอวํ เม สุตํ\csromandash{}เอกํ สมยํ ภควา สาวตฺถิยํ วิหรติ เชตวเน อนาถ\-ปิณฺฑิกสฺส อาราเม. อถ โข อญฺญตโร พฺราหฺมโณ เยน ภควา เตนุปสงฺกมิ, อุปสงฺกมิตฺวา ภควตา สทฺธิํ สมฺโมทิ, สมฺโมทนียํ กถํ สารณียํ วีติสาเรตฺวา เอกมนฺตํ นิสีทิ, เอกมนฺตํ นิสินฺโน โข โส พฺราหฺมโณ ภควนฺตํ เอตทโวจ “โก นุ โข โภ โคตม เหตุ โก ปจฺจโย, เยน ตถาคเต ปรินิพฺพุเต สทฺธมฺโม น จิรฏฺฐิติโก โหติ, โก ปน โภ โคตม เหตุ โก ปจฺจโย, เยน ตถาคเต ปรินิพฺพุเต สทฺธมฺโม จิรฏฺฐิติโก โหตี”ติ.}

\prose{จตุนฺนํ โข พฺราหฺมณ สติปฏฺฐานานํ อภาวิตตฺตา อพหุลีกตตฺตา ตถาคเต ปรินิพฺพุเต สทฺธมฺโม น จิรฏฺฐิติโก โหติ, จตุนฺนญฺจ โข พฺราหฺมณ สติปฏฺฐานานํ ภาวิตตฺตา พหุลีกตตฺตา ตถาคเต ปรินิพฺพุเต สทฺธมฺโม จิรฏฺฐิติโก โหติ.}

\prose{กตเมสํ จตุนฺนํ, อิธ พฺราหฺมณ ภิกฺขุ กาเย กายานุปสฺสี วิหรติ อาตาปี สมฺปชาโน สติมา วิเนยฺย โลเก อภิชฺฌา\-โทมนสฺสํ. เวทนาสุ ฯเปฯ. จิตฺเต ฯเปฯ. ธมฺเมสุ ธมฺมานุปสฺสี วิหรติ อาตาปี สมฺปชาโน สติมา วิเนยฺย โลเก อภิชฺฌา\-โทมนสฺสํ. อิเมสํ โข พฺราหฺมณ จตุนฺนํ สติปฏฺฐานานํ อภาวิตตฺตา อพหุลีกตตฺตา ตถาคเต ปรินิพฺพุเต สทฺธมฺโม น จิรฏฺฐิติโก โหติ. อิเมสญฺจ โข พฺราหฺมณ จตุนฺนํ สติปฏฺฐานานํ ภาวิตตฺตา พหุลีกตตฺตา ตถาคเต ปรินิพฺพุเต สทฺธมฺโม จิรฏฺฐิติโก โหตีติ.}

\prose{เอวํ วุตฺเต โส พฺราหฺมโณ ภควนฺตํ เอตทโวจ “อภิกฺกนฺตํ โภ โคตม ฯเปฯ อุปาสกํ มํ ภวํ โคตโม ธาเรตุ อชฺชตคฺเค ปาณุเปตํ สรณํ คตนฺ”ติ.  ปญฺจมํ.}

\csromanmatikaanchor{mtk.72}
\csromantocmarkref{subparagraph}{6. ปเทสสุตฺต}{mtk.72}
\bookmark[dest={mtk.72},level=5]{6. ปเทสสุตฺต}
\csromanheader{6}{6. \textbf{ปเทสสุตฺต}}
\proseitem{392}{\csromanfolio{152}เอกํ สมยํ อายสฺมา จ สาริปุตฺโต อายสฺมา จ มหาโมคฺคลฺลาโน อายสฺมา จ อนุรุทฺโธ สาเกเต วิหรนฺติ กณฺฑกีวเน\csromansharedfootnote{6eaf2ad64b51c798}{กณฺฑกีวเน (สี, สฺยา, กํ, อิ)}. อถ โข อายสฺมา จ สาริปุตฺโต อายสฺมา จ มหาโมคฺคลฺลาโน สายนฺหสมยํ ปฏิสลฺลานา วุฏฺฐิตา เยนายสฺมา อนุรุทฺโธ เตนุ\-ปสงฺกมิํสุ, อุปสงฺกมิตฺวา อายสฺมตา อนุรุทฺเธน สทฺธิํ สมฺโมทิํสุ, สมฺโมทนียํ กถํ สารณียํ วีติสาเรตฺวา เอกมนฺตํ นิสีทิํสุ, เอกมนฺตํ นิสินฺโน โข อายสฺมา สาริปุตฺโต อายสฺมนฺตํ อนุรุทฺธํ เอตทโวจ “เสโข เสโขติ\csromansharedfootnote{a5b316b840cf8108}{เสกฺโข เสกฺโขติ (สฺยา, กํ)} อาวุโส อนุรุทฺธ วุจฺจติ, กิตฺตาวตา นุ โข อาวุโส เสโข โหตี”ติ, จตุนฺนํ โข อาวุโส สติปฏฺฐานานํ ปเทสํ ภาวิตตฺตา เสโข โหติ.}

\prose{กตเมสํ จตุนฺนํ, อิธาวุโส ภิกฺขุ กาเย กายานุปสฺสี วิหรติ อาตาปี สมฺปชาโน สติมา วิเนยฺย โลเก อภิชฺฌา\-โทมนสฺสํ. เวทนาสุ ฯเปฯ. จิตฺเต ฯเปฯ. ธมฺเมสุ ธมฺมานุปสฺสี วิหรติ อาตาปี สมฺปชาโน สติมา วิเนยฺย โลเก อภิชฺฌา\-โทมนสฺสํ. อิเมสํ โข อาวุโส จตุนฺนํ สติปฏฺฐานานํ ปเทสํ ภาวิตตฺตา เสโข โหตีติ.  ฉฏฺฐํ.}

\csromanmatikaanchor{mtk.73}
\csromantocmarkref{subparagraph}{7. สมตฺตสุตฺต}{mtk.73}
\bookmark[dest={mtk.73},level=5]{7. สมตฺตสุตฺต}
\csromanheader{6}{7. \textbf{สมตฺตสุตฺต}}
\proseitem{393}{ตํเยว นิทานํ. เอกมนฺตํ นิสินฺโน โข อายสฺมา สาริปุตฺโต อายสฺมนฺตํ อนุรุทฺธํ เอตทโวจ “อเสโข อเสโขติ อาวุโส อนุรุทฺธ วุจฺจติ, กิตฺตาวตา นุ โข อาวุโส อเสโข โหตี”ติ. จตุนฺนํ โข อาวุโส สติปฏฺฐานานํ สมตฺตํ ภาวิตตฺตา อเสโข โหติ.}

\prose{กตเมสํ จตุนฺนํ, อิธาวุโส ภิกฺขุ กาเย กายานุปสฺสี วิหรติ อาตาปี สมฺปชาโน สติมา วิเนยฺย โลเก อภิชฺฌา\-โทมนสฺสํ. เวทนาสุ ฯเปฯ. จิตฺเต ฯเปฯ. ธมฺเมสุ ธมฺมานุปสฺสี วิหรติ อาตาปี สมฺปชาโน สติมา วิเนยฺย โลเก อภิชฺฌา\-โทมนสฺสํ. อิเมสํ}

\vspace{\tipitakaparskip}
\csromancenter{โข อาวุโส จตุนฺนํ สติปฏฺฐานานํ สมตฺตํ ภาวิตตฺตา อเสโข โหตีติ.  สตฺตมํ.}
\csromanmatikaanchor{mtk.74}
\csromantocmarkref{subparagraph}{8. โลกสุตฺต}{mtk.74}
\bookmark[dest={mtk.74},level=5]{8. โลกสุตฺต}
\csromanheader{6}{8. \textbf{โลกสุตฺต}}
\proseitem{394}{\csromanfolio{153}ตํเยว นิทานํ. เอกมนฺตํ นิสินฺโน โข อายสฺมา สาริปุตฺโต อายสฺมนฺตํ อนุรุทฺธํ เอตทโวจ “กตเมสํ อาวุโส อนุรุทฺธ ธมฺมานํ ภาวิตตฺตา พหุลีกตตฺตา มหาภิญฺญตํ\csromansharedfootnote{65d22a72d0bc27ae}{มหาภิญฺญาตํ (อิ)} ปตฺโต”ติ. จตุนฺนํ อาวุโส สติปฏฺฐานานํ ภาวิตตฺตา พหุลีกตตฺตา มหาภิญฺญตํ ปตฺโต.}

\prose{กตเมสํ จตุนฺนํ, อิธาหํ อาวุโส กาเย กายานุปสฺสี วิหรามิ อาตาปี สมฺปชาโน สติมา วิเนยฺย โลเก อภิชฺฌา\-โทมนสฺสํ. เวทนาสุ ฯเปฯ. จิตฺเต ฯเปฯ. ธมฺเมสุ ธมฺมานุปสฺสี วิหรามิ อาตาปี สมฺปชาโน สติมา วิเนยฺย โลเก อภิชฺฌา\-โทมนสฺสํ. อิเมสํ ขฺวาหํ อาวุโส จตุนฺนํ สติปฏฺฐานานํ ภาวิตตฺตา พหุลีกตตฺตา มหาภิญฺญตํ ปตฺโต. อิเมสญฺจ ปนาหํ อาวุโส จตุนฺนํ สติปฏฺฐานานํ ภาวิตตฺตา พหุลีกตตฺตา สหสฺสํ โลกํ อภิชานามีติ.  อฏฺฐมํ.}

\csromanmatikaanchor{mtk.75}
\csromantocmarkref{subparagraph}{9. สิริวฑฺฒสุตฺต}{mtk.75}
\bookmark[dest={mtk.75},level=5]{9. สิริวฑฺฒสุตฺต}
\csromanheader{6}{9. สิริวฑฺฒ\-สุตฺต}
\proseitem{395}{เอกํ สมยํ อายสฺมา อานนฺโท ราชคเห วิหรติ เวฬุวเน กลนฺทก\-นิวาเป. เตน โข ปน สมเยน สิริวฑฺโฒ\csromansharedfootnote{fe9b13f168320571}{สิริวฑฺโฒ (ก)} คหปติ อาพาธิโก โหติ ทุกฺขิโต พาฬฺหคิลาโน. อถ โข สิริวฑฺโฒ คหปติ อญฺญตรํ ปุริสํ อามนฺเตสิ “เอหิ ตฺวํ อมฺโภ ปุริส เยนายสฺมา อานนฺโท เตนุปสงฺกมิ, อุปสงฺกมิตฺวา มม วจเนน อายสฺมโต อานนฺทสฺส ปาเท สิรสา วนฺท ‘สิริวฑฺโฒ ภนฺเต คหปติ อาพาธิโก ทุกฺขิโต พาฬฺหคิลาโน, โส อายสฺมโต อานนฺทสฺส ปาเท สิรสา วนฺทตี’ติ, เอวญฺจ วเทหิ ‘สาธุ กิร ภนฺเต อายสฺมา อานนฺโท เยน สิริวฑฺฒสฺส คหปติสฺส นิเวสนํ เตนุ\-ปสงฺกมตุ อนุกมฺปํ อุปาทายา’ติ”. “เอวํ ภนฺเต”ติ โข โส ปุริโส สิริวฑฺฒสฺส คหปติสฺส ปฏิสฺสุตฺวา เยนายสฺมา อานนฺโท เตนุปสงฺกมิ, อุปสงฺกมิตฺวา อายสฺมนฺตํ อานนฺทํ อภิวาเทตฺวา เอกมนฺตํ นิสีทิ, เอกมนฺตํ นิสินฺโน \csromanfolio{154}โข โส ปุริโส อายสฺมนฺตํ อานนฺทํ เอตทโวจ “สิริวฑฺโฒ ภนฺเต คหปติ อาพาธิโก ทุกฺขิโต พาฬฺหคิลาโน, โส อายสฺมโต อานนฺทสฺส ปาเท สิรสา วนฺทติ, เอวญฺจ วเทติ ‘สาธุ กิร ภนฺเต อายสฺมา อานนฺโท เยน สิริวฑฺฒสฺส คหปติสฺส นิเวสนํ เตนุ\-ปสงฺกมตุ อนุกมฺปํ อุปาทายา’ติ”. อธิวาเสสิ โข อายสฺมา อานนฺโท ตุณฺหีภาเวน.}

\prose{อถ โข อายสฺมา อานนฺโท ปุพฺพณฺหสมยํ นิวาเสตฺวา ปตฺต\-จีวรมา\-ทาย เยน สิริวฑฺฒสฺส คหปติสฺส นิเวสนํ เตนุปสงฺกมิ, อุปสงฺกมิตฺวา ปญฺญตฺเต อาสเน นิสีทิ, นิสชฺช โข อายสฺมา อานนฺโท สิริวฑฺฒํ คหปติํ เอตทโวจ “กจฺจิ เต คหปติ ขมนียํ, กจฺจิ ยาปนียํ, กจฺจิ ทุกฺขา เวทนา ปฏิกฺกมนฺติ โน อภิกฺกมนฺติ, ปฏิกฺกโมสานํ ปญฺญายติ โน อภิกฺกโม”ติ. น เม ภนฺเต ขมนียํ, น ยาปนียํ, พาฬฺหา เม ทุกฺขา เวทนา อภิกฺกมนฺติ โน ปฏิกฺกมนฺติ, อภิกฺกโมสานํ ปญฺญายติ โน ปฏิกฺกโมติ.}

\prose{ตสฺมาติห เต คหปติ เอวํ สิกฺขิตพฺพํ “กาเย กายานุปสฺสี วิหริสฺสามิ อาตาปี สมฺปชาโน สติมา วิเนยฺย โลเก อภิชฺฌา\-โทมนสฺสํ. เวทนาสุ ฯเปฯ. จิตฺเต ฯเปฯ. ธมฺเมสุ ธมฺมานุปสฺสี วิหริสฺสามิ อาตาปี สมฺปชาโน สติมา วิเนยฺย โลเก อภิชฺฌาโทมนสฺสนฺ”ติ. เอวํ หิ เต คหปติ สิกฺขิตพฺพนฺติ.}

\prose{เยเม ภนฺเต ภควตา จตฺตาโร สติปฏฺฐานา เทสิตา, สํวิชฺชนฺติ เต ธมฺมา\csromansharedfootnote{9f8dfb548282e0e0}{สํวิชฺชนฺเต รตนธมฺมา (สี)} มยิ, อหญฺจ เตสุ ธมฺเมสุ สนฺทิสฺสามิ. อหํ หิ ภนฺเต กาเย กายานุปสฺสี วิหรามิ อาตาปี สมฺปชาโน สติมา วิเนยฺย โลเก อภิชฺฌา\-โทมนสฺสํ. เวทนาสุ ฯเปฯ. จิตฺเต ฯเปฯ. ธมฺเมสุ ธมฺมานุปสฺสี วิหรามิ อาตาปี สมฺปชาโน สติมา วิเนยฺย โลเก อภิชฺฌา\-โทมนสฺสํ. ยานิ จิมานิ ภนฺเต ภควตา ปญฺโจ\-รมฺภาคิยานิ สํโยชนานิ เทสิตานิ, นาหํ ภนฺเต เตสํ กิญฺจิ อตฺตนิ อปฺปหีนํ สมนุปสฺสามีติ. ลาภา เต คหปติ, สุลทฺธํ เต คหปติ, อนาคามิผลํ ตยา คหปติ พฺยากตนฺติ.  นวมํ.}

\csromanmatikaanchor{mtk.76}
\csromantocmarkref{subparagraph}{10. มานทินฺนสุตฺต}{mtk.76}
\bookmark[dest={mtk.76},level=5]{10. มานทินฺนสุตฺต}
\csromanheader{6}{10. \textbf{มานทินฺนสุตฺต}}
\proseitemwithclosermidruled{396}{\csromanfolio{155}ตํเยว นิทานํ. เตน โข ปน สมเยน มานทินฺโน คหปติ อาพาธิโก โหติ ทุกฺขิโต พาฬฺหคิลาโน. อถ โข มานทินฺโน คหปติ อญฺญตรํ ปุริสํ อามนฺเตสิ “เอหิ ตฺวํ อมฺโภ ปุริส ฯเปฯ” น เม ภนฺเต ขมนียํ, น ยาปนียํ, พาฬฺหา เม ทุกฺขา เวทนา อภิกฺกมนฺติ โน ปฏิกฺกมนฺติ, อภิกฺกโมสานํ ปญฺญายติ โน ปฏิกฺกโมติ. เอวรูปาย จาหํ ภนฺเต ทุกฺขาย เวทนาย ผุฏฺโฐ สมาโน กาเย กายานุปสฺสี วิหรามิ อาตาปี สมฺปชาโน สติมา วิเนยฺย โลเก อภิชฺฌา\-โทมนสฺสํ. เวทนาสุ ฯเปฯ. จิตฺเต ฯเปฯ. ธมฺเมสุ ธมฺมานุปสฺสี วิหรามิ อาตาปี สมฺปชาโน สติมา วิเนยฺย โลเก อภิชฺฌา\-โทมนสฺสํ. ยานิ จิมานิ ภนฺเต ภควตา ปญฺโจ\-รมฺภาคิยานิ สํโยชนานิ เทสิตานิ, นาหํ ภนฺเต เตสํ กิญฺจิ อตฺตนิ อปฺปหีนํ สมนุปสฺสามีติ. ลาภา เต คหปติ, สุลทฺธํ เต คหปติ, อนาคามิผลํ ตยา คหปติ พฺยากตนฺติ.  ทสมํ.}{สีลฏฺฐิติ\-วคฺโค ตติโย.}
\csromancenter{\textbf{ตสฺสุทฺทานํ}}
\setlength{\csromangathapagewidth}{0pt}
\setlength{\csromangathawakwidth}{0pt}
\csromangathameasuregroup{สีลํ ฐิติ ปริหานํ, \\ สมตฺตํ โลโก สิริวฑฺโฒ,}{\csromangathabat{สีลํ ฐิติ ปริหานํ,}{สุทฺธํ พฺราหฺมณปเทสํ.} \\ \csromangathabat{สมตฺตํ โลโก สิริวฑฺโฒ,}{มานทินฺเนน เต ทสาติ.}}
\csromangathasetleft{สีลํ ฐิติ ปริหานํ, \\ สมตฺตํ โลโก สิริวฑฺโฒ,}
\csromangathagroup{\csromangathastanza{\csromangathabat{สีลํ ฐิติ ปริหานํ,}{สุทฺธํ พฺราหฺมณปเทสํ.} \\ \csromangathabat{สมตฺตํ โลโก สิริวฑฺโฒ,}{มานทินฺเนน เต ทสาติ.}}}

\csromanmatikaanchor{mtk.77}
\csromantocmarknopage{subparagraph}{4. อนนุสฺสุตวคฺค}{mtk.77}
\bookmark[dest={mtk.77},level=5]{4. อนนุสฺสุตวคฺค}
\csromanheader{6}{4. อนนุสฺสุต\-วคฺค}
\csromanmatikaanchor{mtk.78}
\csromantocmarkref{subparagraph}{1. อนนุสฺสุตสุตฺต}{mtk.78}
\bookmark[dest={mtk.78},level=5]{1. อนนุสฺสุตสุตฺต}
\csromanheader{6}{1. อนนุสฺสุต\-สุตฺต}
\proseitem{397}{สาวตฺถิ\-นิทานํ. อยํ กาเย กายานุปสฺสนาติ เม ภิกฺขเว ปุพฺเพ อนนุสฺสุเตสุ ธมฺเมสุ จกฺขุํ อุทปาทิ, ญาณํ อุทปาทิ, ปญฺญา อุทปาทิ, วิชฺชา อุทปาทิ, อาโลโก อุทปาทิ. สา โข ปนายํ กาเย กายานุปสฺสนา ภาเวตพฺพาติ เม ภิกฺขเว ฯเปฯ ภาวิตาติ เม ภิกฺขเว ปุพฺเพ อนนุสฺสุเตสุ ธมฺเมสุ จกฺขุํ อุทปาทิ, ญาณํ อุทปาทิ, ปญฺญา อุทปาทิ, วิชฺชา อุทปาทิ, อาโลโก อุทปาทิ.}

\prose{\csromanfolio{156}อยํ เวทนาสุ เวทนา\-นุปสฺสนาติ เม ภิกฺขเว ปุพฺเพ อนนุสฺสุเตสุ ธมฺเมสุ จกฺขุํ อุทปาทิ, ญาณํ อุทปาทิ, ปญฺญา อุทปาทิ, วิชฺชา อุทปาทิ, อาโลโก อุทปาทิ. สา โข ปนายํ เวทนาสุ เวทนา\-นุปสฺสนา ภาเวตพฺพาติ เม ภิกฺขเว ฯเปฯ ภาวิตาติ เม ภิกฺขเว ปุพฺเพ อนนุสฺสุเตสุ ธมฺเมสุ จกฺขุํ อุทปาทิ, ญาณํ อุทปาทิ, ปญฺญา อุทปาทิ, วิชฺชา อุทปาทิ, อาโลโก อุทปาทิ.}

\prose{อยํ จิตฺเต จิตฺตานุปสฺสนาติ เม ภิกฺขเว ปุพฺเพ อนนุสฺสุเตสุ ธมฺเมสุ จกฺขุํ อุทปาทิ, ญาณํ อุทปาทิ, ปญฺญา อุทปาทิ, วิชฺชา อุทปาทิ, อาโลโก อุทปาทิ. สา โข ปนายํ จิตฺเต จิตฺตานุปสฺสนา ภาเวตพฺพาติ เม ภิกฺขเว ฯเปฯ ภาวิตาติ เม ภิกฺขเว ปุพฺเพ อนนุสฺสุเตสุ ธมฺเมสุ จกฺขุํ อุทปาทิ, ญาณํ อุทปาทิ, ปญฺญา อุทปาทิ, วิชฺชา อุทปาทิ, อาโลโก อุทปาทิ.}

\prose{อยํ ธมฺเมสุ ธมฺมา\-นุปสฺสนาติ เม ภิกฺขเว ปุพฺเพ อนนุสฺสุเตสุ ธมฺเมสุ จกฺขุํ อุทปาทิ, ญาณํ อุทปาทิ, ปญฺญา อุทปาทิ, วิชฺชา อุทปาทิ, อาโลโก อุทปาทิ. สา โข ปนายํ ธมฺเมสุ ธมฺมา\-นุปสฺสนา ภาเวตพฺพาติ เม ภิกฺขเว ฯเปฯ ภาวิตาติ เม ภิกฺขเว ปุพฺเพ อนนุสฺสุเตสุ ธมฺเมสุ จกฺขุํ อุทปาทิ, ญาณํ อุทปาทิ, ปญฺญา อุทปาทิ, วิชฺชา อุทปาทิ, อาโลโก อุทปาทีติ.  ปฐมํ.}

\csromanmatikaanchor{mtk.79}
\csromantocmarkref{subparagraph}{2. วิราคสุตฺต}{mtk.79}
\bookmark[dest={mtk.79},level=5]{2. วิราคสุตฺต}
\csromanheader{6}{2. \textbf{วิราคสุตฺต}}
\proseitem{398}{จตฺตาโรเม ภิกฺขเว สติปฏฺฐานา ภาวิตา พหุลีกตา เอกนฺต\-นิพฺพิทาย วิราคาย นิโรธาย อุปสมาย อภิญฺญาย สมฺโพธาย นิพฺพานาย สํวตฺตนฺติ.}

\prose{กตเม จตฺตาโร, อิธ ภิกฺขเว ภิกฺขุ กาเย กายานุปสฺสี วิหรติ อาตาปี สมฺปชาโน สติมา วิเนยฺย โลเก อภิชฺฌา\-โทมนสฺสํ. เวทนาสุ ฯเปฯ. จิตฺเต ฯเปฯ. ธมฺเมสุ ธมฺมานุปสฺสี วิหรติ อาตาปี สมฺปชาโน สติมา วิเนยฺย โลเก อภิชฺฌา\-โทมนสฺสํ. อิเม โข ภิกฺขเว จตฺตาโร สติปฏฺฐานา ภาวิตา พหุลีกตา เอกนฺต\-นิพฺพิทาย วิราคาย นิโรธาย อุปสมาย อภิญฺญาย สมฺโพธาย นิพฺพานาย สํวตฺตนฺตีติ.  ทุติยํ.}

\csromanheader{2}{3. \textbf{วิรทฺธสุตฺต}}
\csromanmatikaanchor{mtk.80}
\csromanmatikaanchor{mtk.81}
\csromantocmarkref{subparagraph}{3-4. วิรทฺธสุตฺตาทีนิ}{mtk.80}
\bookmark[dest={mtk.80},level=5]{3-4. วิรทฺธสุตฺตาทีนิ}
\csromantocmarkref{subparagraph}{5-9. สติสุตฺตาทีนิ}{mtk.81}
\bookmark[dest={mtk.81},level=5]{5-9. สติสุตฺตาทีนิ}
\proseitem{399}{\csromanfolio{157}เยสํ เกสญฺจิ ภิกฺขเว จตฺตาโร สติปฏฺฐานา วิรทฺธา, วิรทฺโธ เตสํ อริโย มคฺโค\csromansharedfootnote{5cdab4cbe2162f9d}{อริโย อฏฺฐงฺคิโก มคฺโค (ก) อิมสฺมิํ เยว สุตฺเต ทิสฺสติ อฏฺฐงฺคิโกติปทํ, น ปนาญฺญตฺถ อิทฺธิปาท อนุรุทฺธาทีสุ.} สมฺมา ทุกฺข\-กฺขย\-คามี. เยสํ เกสญฺจิ ภิกฺขเว จตฺตาโร สติปฏฺฐานา อารทฺธา, อารทฺโธ เตสํ อริโย มคฺโค สมฺมา ทุกฺข\-กฺขย\-คามี.}

\prose{กตเม จตฺตาโร, อิธ ภิกฺขเว ภิกฺขุ กาเย กายานุปสฺสี วิหรติ อาตาปี สมฺปชาโน สติมา วิเนยฺย โลเก อภิชฺฌา\-โทมนสฺสํ. เวทนาสุ ฯเปฯ. จิตฺเต ฯเปฯ. ธมฺเมสุ ธมฺมานุปสฺสี วิหรติ อาตาปี สมฺปชาโน สติมา วิเนยฺย โลเก อภิชฺฌา\-โทมนสฺสํ. เยสํ เกสญฺจิ ภิกฺขเว อิเม จตฺตาโร สติปฏฺฐานา วิรทฺธา, วิรทฺโธ เตสํ อริโย มคฺโค สมฺมา ทุกฺข\-กฺขย\-คามี. เยสํ เกสญฺจิ ภิกฺขเว อิเม จตฺตาโร สติปฏฺฐานา อารทฺธา, อารทฺโธ เตสํ อริโย มคฺโค สมฺมา ทุกฺข\-กฺขย\-คามีติ.  ตติยํ.}

\csromanheader{2}{4. \textbf{ภาวิตสุตฺต}}
\proseitem{400}{จตฺตาโรเม ภิกฺขเว สติปฏฺฐานา ภาวิตา พหุลีกตา อปารา ปารํ คมนาย สํวตฺตนฺติ.}

\prose{กตเม จตฺตาโร, อิธ ภิกฺขเว ภิกฺขุ กาเย กายานุปสฺสี วิหรติ อาตาปี สมฺปชาโน สติมา วิเนยฺย โลเก อภิชฺฌา\-โทมนสฺสํ. เวทนาสุ ฯเปฯ. จิตฺเต ฯเปฯ. ธมฺเมสุ ธมฺมานุปสฺสี วิหรติ อาตาปี สมฺปชาโน สติมา วิเนยฺย โลเก อภิชฺฌา\-โทมนสฺสํ. อิเม โข ภิกฺขเว จตฺตาโร สติปฏฺฐานา ภาวิตา พหุลีกตา อปารา ปารํ คมนาย สํวตฺตนฺตีติ.  จตุตฺถํ.}

\csromanheader{2}{5. \textbf{สติสุตฺต}}
\proseitem{401}{สาวตฺถิ\-นิทานํ. สโต ภิกฺขเว ภิกฺขุ วิหเรยฺย สมฺปชาโน, อยํ โว อมฺหากํ อนุสาสนี.}

\prose{กถญฺจ ภิกฺขเว ภิกฺขุ สโต โหติ, อิธ ภิกฺขเว ภิกฺขุ กาเย กายานุปสฺสี วิหรติ อาตาปี สมฺปชาโน สติมา วิเนยฺย โลเก \csromanfolio{158}อภิชฺฌา\-โทมนสฺสํ. เวทนาสุ ฯเปฯ. จิตฺเต ฯเปฯ. ธมฺเมสุ ธมฺมานุปสฺสี วิหรติ อาตาปี สมฺปชาโน สติมา วิเนยฺย โลเก อภิชฺฌา\-โทมนสฺสํ. เอวํ โข ภิกฺขเว ภิกฺขุ สโต โหติ.}

\prose{กถญฺจ ภิกฺขเว ภิกฺขุ สมฺปชาโน โหติ, อิธ ภิกฺขเว ภิกฺขุโน วิทิตา เวทนา อุปฺปชฺชนฺติ, วิทิตา อุปฏฺฐหนฺติ, วิทิตา อพฺภตฺถํ คจฺฉนฺติ. วิทิตา วิตกฺกา อุปฺปชฺชนฺติ, วิทิตา อุปฏฺฐหนฺติ, วิทิตา อพฺภตฺถํ คจฺฉนฺติ. วิทิตา สญฺญา อุปฺปชฺชนฺติ, วิทิตา อุปฏฺฐหนฺติ, วิทิตา อพฺภตฺถํ คจฺฉนฺติ. เอวํ โข ภิกฺขเว ภิกฺขุ สมฺปชาโน โหติ, สโต ภิกฺขเว ภิกฺขุ วิหเรยฺย สมฺปชาโน, อยํ โว อมฺหากํ อนุสาสนีติ.  ปญฺจมํ.}

\csromanheader{2}{6. \textbf{อญฺญาสุตฺต}}
\proseitem{402}{จตฺตาโรเม ภิกฺขเว สติปฏฺฐานา. กตเม จตฺตาโร, อิธ ภิกฺขเว ภิกฺขุ กาเย กายานุปสฺสี วิหรติ อาตาปี สมฺปชาโน สติมา วิเนยฺย โลเก อภิชฺฌา\-โทมนสฺสํ. เวทนาสุ ฯเปฯ. จิตฺเต ฯเปฯ. ธมฺเมสุ ธมฺมานุปสฺสี วิหรติ อาตาปี สมฺปชาโน สติมา วิเนยฺย โลเก อภิชฺฌา\-โทมนสฺสํ. อิเม โข ภิกฺขเว จตฺตาโร สติปฏฺฐานา. อิเมสํ โข ภิกฺขเว จตุนฺนํ สติปฏฺฐานานํ ภาวิตตฺตา พหุลีกตตฺตา ทฺวินฺนํ ผลานํ อญฺญตรํ ผลํ ปาฏิกงฺขํ, ทิฏฺเฐว ธมฺเม อญฺญา, สติ วา อุปาทิเสเส อนาคามิตาติ.  ฉฏฺฐํ.}

\csromanheader{2}{7. \textbf{ฉนฺทสุตฺต}}
\proseitem{403}{จตฺตาโรเม ภิกฺขเว สติปฏฺฐานา. กตเม จตฺตาโร, อิธ ภิกฺขเว ภิกฺขุ กาเย กายานุปสฺสี วิหรติ อาตาปี สมฺปชาโน สติมา วิเนยฺย โลเก อภิชฺฌา\-โทมนสฺสํ, ตสฺส กาเย กายานุปสฺสิโน วิหรโต โย กายสฺมิํ ฉนฺโท, โส ปหียติ, ฉนฺทสฺส ปหานา อมตํ สจฺฉิกตํ โหติ.}

\prose{เวทนาสุ เวทนานุปสฺสี วิหรติ อาตาปี สมฺปชาโน สติมา วิเนยฺย โลเก อภิชฺฌา\-โทมนสฺสํ, ตสฺส เวทนาสุ เวทนา\-นุปสฺสิโน วิหรโต โย เวทนาสุ ฉนฺโท, โส ปหียติ, ฉนฺทสฺส ปหานา อมตํ สจฺฉิกตํ โหติ.}

\prose{\csromanfolio{159}จิตฺเต จิตฺตานุปสฺสี วิหรติ อาตาปี สมฺปชาโน สติมา วิเนยฺย โลเก อภิชฺฌา\-โทมนสฺสํ, ตสฺส จิตฺเต จิตฺตานุปสฺสิโน วิหรโต โย จิตฺตมฺหิ ฉนฺโท, โส ปหียติ, ฉนฺทสฺส ปหานา อมตํ สจฺฉิกตํ โหติ.}

\prose{ธมฺเมสุ ธมฺมานุปสฺสี วิหรติ อาตาปี สมฺปชาโน สติมา วิเนยฺย โลเก อภิชฺฌา\-โทมนสฺสํ, ตสฺส ธมฺเมสุ ธมฺมา\-นุปสฺสิโน วิหรโต โย ธมฺเมสุ ฉนฺโท, โส ปหียติ, ฉนฺทสฺส ปหานา อมตํ สจฺฉิกตํ โหตีติ.  สตฺตมํ.}

\csromanheader{2}{8. \textbf{ปริญฺญาตสุตฺต}}
\proseitem{404}{จตฺตาโรเม ภิกฺขเว สติปฏฺฐานา. กตเม จตฺตาโร, อิธ ภิกฺขเว ภิกฺขุ กาเย กายานุปสฺสี วิหรติ อาตาปี สมฺปชาโน สติมา วิเนยฺย โลเก อภิชฺฌา\-โทมนสฺสํ, ตสฺส กาเย กายานุปสฺสิโน วิหรโต กาโย ปริญฺญาโต โหติ, กายสฺส ปริญฺญาตตฺตา อมตํ สจฺฉิกตํ โหติ.}

\prose{เวทนาสุ เวทนานุปสฺสี วิหรติ อาตาปี สมฺปชาโน สติมา วิเนยฺย โลเก อภิชฺฌา\-โทมนสฺสํ, ตสฺส เวทนาสุ เวทนา\-นุปสฺสิโน วิหรโต เวทนา ปริญฺญาตา โหนฺติ, เวทนานํ ปริญฺญาตตฺตา อมตํ สจฺฉิกตํ โหติ.}

\prose{จิตฺเต จิตฺตานุปสฺสี วิหรติ อาตาปี สมฺปชาโน สติมา วิเนยฺย โลเก อภิชฺฌา\-โทมนสฺสํ, ตสฺส จิตฺเต จิตฺตานุปสฺสิโน วิหรโต จิตฺตํ ปริญฺญาตํ โหติ, จิตฺตสฺส ปริญฺญาตตฺตา อมตํ สจฺฉิกตํ โหติ.}

\prose{ธมฺเมสุ ธมฺมานุปสฺสี วิหรติ อาตาปี สมฺปชาโน สติมา วิเนยฺย โลเก อภิชฺฌา\-โทมนสฺสํ, ตสฺส ธมฺเมสุ ธมฺมานุปสฺสีโน วิหรโต ธมฺมา ปริญฺญาตา โหนฺติ, ธมฺมานํ ปริญฺญาตตฺตา อมตํ สจฺฉิกตํ โหตีติ.  อฏฺฐมํ.}

\csromanheader{2}{9. \textbf{ภาวนาสุตฺต}}
\proseitem{405}{\csromanfolio{160}จตุนฺนํ ภิกฺขเว สติปฏฺฐานานํ ภาวนํ เทเสสฺสามิ, ตํ สุณาถ. กตมา ภิกฺขเว จตุนฺนํ สติปฏฺฐานานํ ภาวนา, อิธ ภิกฺขเว ภิกฺขุ กาเย กายานุปสฺสี วิหรติ อาตาปี สมฺปชาโน สติมา วิเนยฺย โลเก อภิชฺฌา\-โทมนสฺสํ. เวทนาสุ ฯเปฯ. จิตฺเต ฯเปฯ. ธมฺเมสุ ธมฺมานุปสฺสี วิหรติ อาตาปี สมฺปชาโน สติมา วิเนยฺย โลเก อภิชฺฌา\-โทมนสฺสํ. อยํ โข ภิกฺขเว จตุนฺนํ สติปฏฺฐานานํ ภาวนาติ.  นวมํ.}

\csromanmatikaanchor{mtk.82}
\csromantocmarkref{subparagraph}{10. วิภงฺคสุตฺต}{mtk.82}
\bookmark[dest={mtk.82},level=5]{10. วิภงฺคสุตฺต}
\csromanheader{6}{10. \textbf{วิภงฺคสุตฺต}}
\proseitem{406}{สติปฏฺฐานญฺจ โว ภิกฺขเว เทเสสฺสามิ สติปฏฺฐานภาวนญฺจ สติปฏฺฐานภาวนาคามินิญฺจ ปฏิปทํ, ตํ สุณาถ. กตมญฺจ ภิกฺขเว สติปฏฺฐานํ, อิธ ภิกฺขเว ภิกฺขุ กาเย กายานุปสฺสี วิหรติ อาตาปี สมฺปชาโน สติมา วิเนยฺย โลเก อภิชฺฌา\-โทมนสฺสํ. เวทนาสุ เวทนานุปสฺสี วิหรติ ฯเปฯ. จิตฺเต จิตฺตานุปสฺสี วิหรติ ฯเปฯ. ธมฺเมสุ ธมฺมานุปสฺสี วิหรติ อาตาปี สมฺปชาโน สติมา วิเนยฺย โลเก อภิชฺฌา\-โทมนสฺสํ. อิทํ วุจฺจติ ภิกฺขเว สติปฏฺฐานํ.}

\prose{กตมา จ ภิกฺขเว สติปฏฺฐาน\-ภาวนา, อิธ ภิกฺขเว ภิกฺขุ สมุทยธมฺมา\-นุปสฺสี กายสฺมิํ วิหรติ, วยธมฺมา\-นุปสฺสี กายสฺมิํ วิหรติ, สมุทย\-วย\-ธมฺมานุปสฺสี กายสฺมิํ วิหรติ อาตาปี สมฺปชาโน สติมา วิเนยฺย โลเก อภิชฺฌา\-โทมนสฺสํ. สมุทยธมฺมา\-นุปสฺสี เวทนาสุ วิหรติ ฯเปฯ. สมุทยธมฺมา\-นุปสฺสี จิตฺเต วิหรติ ฯเปฯ. สมุทยธมฺมา\-นุปสฺสี ธมฺเมสุ วิหรติ, วยธมฺมา\-นุปสฺสี ธมฺเมสุ วิหรติ, สมุทย\-วย\-ธมฺมานุปสฺสี ธมฺเมสุ วิหรติ อาตาปี สมฺปชาโน สติมา วิเนยฺย โลเก อภิชฺฌา\-โทมนสฺสํ. อยํ วุจฺจติ ภิกฺขเว สติปฏฺฐาน\-ภาวนา.}

\prosewithclosermidruled{กตมา จ ภิกฺขเว สติปฏฺฐาน\-ภาวนาคามินี ปฏิปทา, อยเมว อริโย อฏฺฐงฺคิโก มคฺโค. เสยฺยถิทํ, สมฺมาทิฏฺฐิ สมฺมาสงฺกปฺโป สมฺมาวาจา สมฺมากมฺมนฺโต สมฺมาอาชีโว สมฺมาวายาโม สมฺมาสติ สมฺมาสมาธิ. อยํ วุจฺจติ ภิกฺขเว สติปฏฺฐาน\-ภาวนาคามินี ปฏิปทาติ.  ทสมํ.}{อนนุสฺสุต\-วคฺโค จตุตฺโถ.}
\csromancenter{\textbf{ตสฺสุทฺทานํ}}
\csromanmatikaanchor{mtk.83}
\csromanmatikaanchor{mtk.84}
\csromanmatikaanchor{mtk.85}
\csromantocmarknopage{subparagraph}{5. อมตวคฺค}{mtk.83}
\bookmark[dest={mtk.83},level=5]{5. อมตวคฺค}
\csromantocmarkref{subparagraph}{1-2. อมตสุตฺตาทีนิ}{mtk.84}
\bookmark[dest={mtk.84},level=5]{1-2. อมตสุตฺตาทีนิ}
\csromantocmarkref{subparagraph}{3-4. มคฺคสุตฺตาทีนิ}{mtk.85}
\bookmark[dest={mtk.85},level=5]{3-4. มคฺคสุตฺตาทีนิ}
\setlength{\csromangathapagewidth}{0pt}
\setlength{\csromangathawakwidth}{0pt}
\csromangathameasuregroup{อนนุสฺสุตํ วิราโค, \\ อญฺญา ฉนฺทํ ปริญฺญาย,}{\csromangathabat{อนนุสฺสุตํ วิราโค,}{วิรทฺโธ ภาวนา สติ.} \\ \csromangathabat{อญฺญา ฉนฺทํ ปริญฺญาย,}{ภาวนา วิภงฺเคน จาติ.}}
\csromangathasetleft{อนนุสฺสุตํ วิราโค, \\ อญฺญา ฉนฺทํ ปริญฺญาย,}
\csromangathagroup{\csromangathastanza{\csromanfolio{161}\csromangathabat{อนนุสฺสุตํ วิราโค,}{วิรทฺโธ ภาวนา สติ.} \\ \csromangathabat{อญฺญา ฉนฺทํ ปริญฺญาย,}{ภาวนา วิภงฺเคน จาติ.}}}

\csromanheader{6}{5. \textbf{อมตวคฺค}}
\csromanheader{2}{1. \textbf{อมตสุตฺต}}
\proseitem{407}{สาวตฺถิ\-นิทานํ. จตูสุ ภิกฺขเว สติปฏฺฐาเนสุ สุปฺปติฏฺฐิต\-จิตฺตา วิหรถ, มา โว อมตํ ปนสฺส. กตเมสุ จตูสุ, อิธ ภิกฺขเว ภิกฺขุ กาเย กายานุปสฺสี วิหรติ อาตาปี สมฺปชาโน สติมา วิเนยฺย โลเก อภิชฺฌา\-โทมนสฺสํ. เวทนาสุ ฯเปฯ. จิตฺเต ฯเปฯ. ธมฺเมสุ ธมฺมานุปสฺสี วิหรติ อาตาปี สมฺปชาโน สติมา วิเนยฺย โลเก อภิชฺฌา\-โทมนสฺสํ. อิเมสุ ภิกฺขเว จตูสุ สติปฏฺฐาเนสุ สุปฺปติฏฺฐิต\-จิตฺตา วิหรถ, มา โว อมตํ ปนสฺสาติ.  ปฐมํ.}

\csromanheader{2}{2. \textbf{สมุทยสุตฺต}}
\proseitem{408}{จตุนฺนํ ภิกฺขเว สติปฏฺฐานานํ สมุทยญฺจ อตฺถงฺคมญฺจ เทเสสฺสามิ, ตํ สุณาถ. โก จ ภิกฺขเว กายสฺส สมุทโย, อาหารสมุทยา กายสฺส สมุทโย, อาหารนิโรธา กายสฺส อตฺถงฺคโม. ผสฺสสมุทยา เวทนานํ สมุทโย, ผสฺสนิโรธา เวทนานํ อตฺถงฺคโม. นามรูป\-สมุทยา จิตฺตสฺส สมุทโย, นามรูป\-นิโรธา จิตฺตสฺส อตฺถงฺคโม. มนสิการ\-สมุทยา ธมฺมานํ สมุทโย, มนสิการ\-นิโรธา ธมฺมานํ อตฺถงฺคโมติ.  ทุติยํ.}

\csromanheader{2}{3. \textbf{มคฺคสุตฺต}}
\proseitem{409}{สาวตฺถิ\-นิทานํ. ตตฺร โข ภควา ภิกฺขู อามนฺเตสิ\csromandash{} เอกมิทาหํ ภิกฺขเว สมยํ อุรุเวลายํ วิหรามิ นชฺชา เนรญฺชราย ตีเร อชปาล\-นิโคฺรเธ ปฐมา\-ภิสมฺพุทฺโธ. ตสฺส มยฺหํ ภิกฺขเว รโหคตสฺส ปฏิสลฺลีนสฺส เอวํ เจตโส ปริวิตกฺโก อุทปาทิ เอกายนฺวายํ มคฺโค \csromanfolio{162}สตฺตานํ วิสุทฺธิยา โสก\-ปริเทวานํ สมติกฺกมาย ทุกฺข\-โทมนสฺสานํ อตฺถงฺคมาย ญายสฺส อธิคมาย นิพฺพานสฺส สจฺฉิกิริยาย ยทิทํ จตฺตาโร สติปฏฺฐานา.}

\prose{กตเม จตฺตาโร, กาเย วา ภิกฺขุ กายานุปสฺสี วิหเรยฺย อาตาปี สมฺปชาโน สติมา วิเนยฺย โลเก อภิชฺฌา\-โทมนสฺสํ. เวทนาสุ วา ภิกฺขุ เวทนานุปสฺสี วิหเรยฺย ฯเปฯ. จิตฺเต วา ภิกฺขุ จิตฺตานุปสฺสี วิหเรยฺย ฯเปฯ. ธมฺเมสุ วา ภิกฺขุ ธมฺมานุปสฺสี วิหเรยฺย อาตาปี สมฺปชาโน สติมา วิเนยฺย โลเก อภิชฺฌา\-โทมนสฺสํ. เอกายนฺวายํ มคฺโค สตฺตานํ วิสุทฺธิยา โสก\-ปริเทวานํ สมติกฺกมาย ทุกฺข\-โทมนสฺสานํ อตฺถงฺคมาย ญายสฺส อธิคมาย นิพฺพานสฺส สจฺฉิกิริยาย ยทิทํ จตฺตาโร สติปฏฺฐานาติ.}

\prose{อถ โข ภิกฺขเว พฺรหฺมา สหมฺปติ มม เจตสา เจโต\-ปริวิตกฺกม\-ญฺญาย, เสยฺยถาปิ นาม พลวา ปุริโส สมิญฺชิตํ วา พาหํ ปสาเรยฺย, ปสาริตํ วา พาหํ สมิญฺเชยฺย, เอวเมว พฺรหฺมโลเก อนฺตรหิโต มม ปุรโต ปาตุรโหสิ. อถ โข ภิกฺขเว พฺรหฺมา สหมฺปติ เอกํสํ อุตฺตราสงฺคํ กริตฺวา เยนาหํ เตนญฺชลิํ ปณาเมตฺวา มํ เอตทโวจ\csromandash{}เอวเมตํ ภควา, เอวเมตํ สุคต, เอกายนฺวายํ ภนฺเต มคฺโค สตฺตานํ วิสุทฺธิยา โสก\-ปริเทวานํ สมติกฺกมาย ทุกฺข\-โทมนสฺสานํ อตฺถงฺคมาย ญายสฺส อธิคมาย นิพฺพานสฺส สจฺฉิกิริยาย ยทิทํ จตฺตาโร สติปฏฺฐานา.}

\prose{กตเม จตฺตาโร, กาเย วา ภนฺเต ภิกฺขุ กายานุปสฺสี วิหเรยฺย อาตาปี สมฺปชาโน สติมา วิเนยฺย โลเก อภิชฺฌา\-โทมนสฺสํ. เวทนาสุ วา ฯเปฯ. จิตฺเต วา ฯเปฯ. ธมฺเมสุ วา ภนฺเต ภิกฺขุ ธมฺมานุปสฺสี วิหเรยฺย อาตาปี สมฺปชาโน สติมา วิเนยฺย โลเก อภิชฺฌา\-โทมนสฺสํ. เอกายนฺวายํ ภนฺเต มคฺโค สตฺตานํ วิสุทฺธิยา โสก\-ปริเทวานํ สมติกฺกมาย ทุกฺข\-โทมนสฺสานํ อตฺถงฺคมาย ญายสฺส อธิคมาย นิพฺพานสฺส สจฺฉิกิริยาย ยทิทํ จตฺตาโร สติปฏฺฐานาติ.}

\prose{อิทมโวจ ภิกฺขเว พฺรหฺมา สหมฺปติ, อิทํ วตฺวา อถาปรํ เอตทโวจ\csromandash{}}

\setlength{\csromangathapagewidth}{0pt}
\setlength{\csromangathawakwidth}{0pt}
\csromangathameasuregroup{“เอกายนํ ชาติขยนฺตทสฺสี,}{\csromangathabat{“เอกายนํ ชาติขยนฺตทสฺสี,}{มคฺคํ ปชานาติ หิตานุกมฺปี.}}
\csromangathasetleft{“เอกายนํ ชาติขยนฺตทสฺสี,}
\csromangathagroup{\csromangathastanza{\csromangathabat{“เอกายนํ ชาติ\-ขยนฺต\-ทสฺสี,}{มคฺคํ ปชานาติ หิตานุกมฺปี.}}}

\prose{เอเตน มคฺเคน ตริํสุ ปุพฺเพ, ตริสฺสนฺติ เย จ ตรนฺติ โอฆนฺ”ติ. ตติยํ.}

\csromanheader{2}{4. \textbf{สติสุตฺต}}
\proseitem{410}{\csromanfolio{163}สโต ภิกฺขเว ภิกฺขุ วิหเรยฺย, อยํ โว อมฺหากํ อนุสาสนี. กถญฺจ ภิกฺขเว ภิกฺขุ สโต โหติ, อิธ ภิกฺขเว ภิกฺขุ กาเย กายานุปสฺสี วิหรติ อาตาปี สมฺปชาโน สติมา วิเนยฺย โลเก อภิชฺฌา\-โทมนสฺสํ. เวทนาสุ ฯเปฯ. จิตฺเต ฯเปฯ. ธมฺเมสุ ธมฺมานุปสฺสี วิหรติ อาตาปี สมฺปชาโน สติมา วิเนยฺย โลเก อภิชฺฌา\-โทมนสฺสํ. เอวํ โข ภิกฺขเว ภิกฺขุ สโต โหติ, สโต ภิกฺขเว ภิกฺขุ วิหเรยฺย, อยํ โว อมฺหากํ อนุสาสนีติ.  จตุตฺถํ.}

\csromanmatikaanchor{mtk.86}
\csromantocmarkref{subparagraph}{5. กุสลราสิสุตฺต}{mtk.86}
\bookmark[dest={mtk.86},level=5]{5. กุสลราสิสุตฺต}
\csromanheader{6}{5. กุสลราสิ\-สุตฺต}
\proseitem{411}{“กุสลราสี”ติ ภิกฺขเว วทมาโน จตฺตาโร สติปฏฺฐาเน สมฺมา วทมาโน วเทยฺย. เกวโล หายํ ภิกฺขเว กุสลราสิ ยทิทํ จตฺตาโร สติปฏฺฐานา.}

\prose{กตเม จตฺตาโร, อิธ ภิกฺขเว ภิกฺขุ กาเย กายานุปสฺสี วิหรติ อาตาปี สมฺปชาโน สติมา วิเนยฺย โลเก อภิชฺฌา\-โทมนสฺสํ. เวทนาสุ ฯเปฯ. จิตฺเต ฯเปฯ. ธมฺเมสุ ธมฺมานุปสฺสี วิหรติ อาตาปี สมฺปชาโน สติมา วิเนยฺย โลเก อภิชฺฌา\-โทมนสฺสํ. เวทนาสุ ฯเปฯ. จิตฺเต ฯเปฯ. ธมฺเมสุ ธมฺมานุปสฺสี วิหรติ อาตาปิ สมฺปชาโน สติมา วิเนยฺย โลเก อภิชฺฌา\-โทมนสฺสํ. “กุสลราสี”ติ ภิกฺขเว วทมาโน อิเม จตฺตาโร สติปฏฺฐาเน สมฺมา วทมาโน วเทยฺย. เกวโล หายํ ภิกฺขเว กุสลราสิ ยทิทํ จตฺตาโร สติปฏฺฐานาติ.  ปญฺจมํ.}

\csromanmatikaanchor{mtk.87}
\csromantocmarkref{subparagraph}{6. ปาติโมกฺขสํวรสุตฺต}{mtk.87}
\bookmark[dest={mtk.87},level=5]{6. ปาติโมกฺขสํวรสุตฺต}
\csromanheader{6}{6. ปาติโมกฺข\-สํวรสุตฺต}
\proseitem{412}{อถ โข อญฺญตโร ภิกฺขุ เยน ภควา เตนุปสงฺกมิ ฯเปฯ เอกมนฺตํ นิสินฺโน โข โส ภิกฺขุ ภควนฺตํ เอตทโวจ\csromandash{}}

\prose{สาธุ เม ภนฺเต ภควา สํขิตฺเตน ธมฺมํ เทเสตุ, ยมหํ ภควโต ธมฺมํ สุตฺวา เอโก วูปกฏฺโฐ อปฺปมตฺโต อาตาปี ปหิตตฺโต วิหเรยฺยนฺติ. ตสฺมาติห ตฺวํ ภิกฺขุ อาทิเมว วิโสเธหิ กุสเลสุ ธมฺเมสุ. โก จาทิ กุสลานํ ธมฺมานํ, อิธ ตฺวํ ภิกฺขุ ปาติโมกฺข\-สํวร\-สํวุโต วิหราหิ อาจารโคจร\-สมฺปนฺโน อณุมตฺเตสุ วชฺเชสุ ภยทสฺสาวี, สมาทาย สิกฺขสฺสุ สิกฺขาปเทสุ, ยโต โข ตฺวํ ภิกฺขุ ปาติโมกฺข\-สํวร\-สํวุโต วิหริสฺสสิ อาจารโคจร\-สมฺปนฺโน อณุมตฺเตสุ \csromanfolio{164}วชฺเชสุ ภยทสฺสาวี, สมาทาย สิกฺขสฺสุ สิกฺขาปเทสุ, ตโต ตฺวํ ภิกฺขุ สีลํ นิสฺสาย สีเล ปติฏฺฐาย จตฺตาโร สติปฏฺฐาเน ภาเวยฺยาสิ.}

\prose{กตเม จตฺตาโร, อิธ ตฺวํ ภิกฺขุ กาเย กายานุปสฺสี วิหราหิ อาตาปี สมฺปชาโน สติมา วิเนยฺย โลเก อภิชฺฌา\-โทมนสฺสํ. เวทนาสุ ฯเปฯ. จิตฺเต ฯเปฯ. ธมฺเมสุ ธมฺมานุปสฺสี วิหราหิ อาตาปี สมฺปชาโน สติมา วิเนยฺย โลเก อภิชฺฌา\-โทมนสฺสํ. ยโต โข ตฺวํ ภิกฺขุ สีลํ นิสฺสาย สีเล ปติฏฺฐาย อิเม จตฺตาโร สติปฏฺฐาเน เอวํ ภาเวสฺสสิ, ตโต ตุยฺหํ ภิกฺขุ ยา รตฺติ วา ทิวโส วา อาคมิสฺสติ, วุทฺธิเยว ปาฏิกงฺขา กุสเลสุ ธมฺเมสุ, โน ปริหานีติ. อถ โข โส ภิกฺขุ ภควโต ภาสิตํ อภินนฺทิตฺวา อนุโมทิตฺวา อุฏฺฐายาสนา ภควนฺตํ อภิวาเทตฺวา ปทกฺขิณํ กตฺวา ปกฺกามิ.}

\prose{อถ โข โส ภิกฺขุ เอโก วูปกฏฺโฐ อปฺปมตฺโต อาตาปี ปหิตตฺโต วิหรนฺโต นจิรสฺเสว ยสฺสตฺถาย กุลปุตฺตา สมฺมเทว อคารสฺมา อนคาริยํ ปพฺพชนฺติ, ตทนุตฺตรํ พฺรหฺมจริย\-ปริโยสานํ ทิฏฺเฐว ธมฺเม สยํ อภิญฺญา สจฺฉิกตฺวา อุปสมฺปชฺช วิหรติ, “ขีณา ชาติ, วุสิตํ พฺรหฺมจริยํ, กตํ กรณียํ, นาปรํ อิตฺถตฺตายา”ติ อพฺภญฺญาสิ, อญฺญตโร จ ปน โส ภิกฺขุ อรหตํ อโหสีติ.  ฉฏฺฐํ.}

\csromanmatikaanchor{mtk.88}
\csromantocmarkref{subparagraph}{7. ทุจฺจริตสุตฺต}{mtk.88}
\bookmark[dest={mtk.88},level=5]{7. ทุจฺจริตสุตฺต}
\csromanheader{6}{7. \textbf{ทุจฺจริตสุตฺต}}
\proseitem{413}{อถ โข อญฺญตโร ภิกฺขุ เยน ภควา เตนุปสงฺกมิ ฯเปฯ สาธุ เม ภนฺเต ภควา สํขิตฺเตน ธมฺมํ เทเสตุ, ยมหํ ภควโต ธมฺมํ สุตฺวา เอโก วูปกฏฺโฐ อปฺปมตฺโต อาตาปี ปหิตตฺโต วิหเรยฺยนฺติ. ตสฺมาติห ตฺวํ ภิกฺขุ อาทิเมว วิโสเธหิ กุสเลสุ ธมฺเมสุ. โก จาทิ กุสลานํ ธมฺมานํ, อิธ ตฺวํ ภิกฺขุ กายทุจฺจริตํ ปหาย กายสุจริตํ ภาเวสฺสสิ, วจีทุจฺจริตํ ปหาย วจีสุจริตํ ภาเวสฺสสิ, มโนทุจฺจริตํ ปหาย มโนสุจริตํ ภาเวสฺสสิ, ยโต โข ตฺวํ ภิกฺขุ กายทุจฺจริตํ ปหาย กายสุจริตํ ภาเวสฺสสิ, วจีทุจฺจริตํ ปหาย วจีสุจริตํ ภาเวสฺสสิ, มโนทุจฺจริตํ ปหาย มโนสุจริตํ ภาเวสฺสสิ, ตโต ตฺวํ ภิกฺขุ สีลํ นิสฺสาย สีเล ปติฏฺฐาย จตฺตาโร สติปฏฺฐาเน ภาเวยฺยาสิ.}

\csromanmatikaanchor{mtk.89}
\csromantocmarkref{subparagraph}{8-10. มิตฺตสุตฺตาทีนิ}{mtk.89}
\bookmark[dest={mtk.89},level=5]{8-10. มิตฺตสุตฺตาทีนิ}
\prose{\csromanfolio{165}กตเม จตฺตาโร, อิธ ตฺวํ ภิกฺขุ กาเย กายานุปสฺสี วิหราหิ อาตาปี สมฺปชาโน สติมา วิเนยฺย โลเก อภิชฺฌา\-โทมนสฺสํ. เวทนาสุ ฯเปฯ. จิตฺเต ฯเปฯ. ธมฺเมสุ ธมฺมานุปสฺสี วิหราหิ อาตาปี สมฺปชาโน สติมา วิเนยฺย โลเก อภิชฺฌา\-โทมนสฺสํ. ยโต โข ตฺวํ ภิกฺขุ สีลํ นิสฺสาย สีเล ปติฏฺฐาย อิเม จตฺตาโร สติปฏฺฐาเน เอวํ ภาเวสฺสสิ, ตโต ตุยฺหํ ภิกฺขุ ยา รตฺติ วา ทิวโส วา อาคมิสฺสติ, วุทฺธิเยว ปาฏิกงฺขา กุสเลสุ ธมฺเมสุ, โน ปริหานีติ ฯเปฯ อญฺญตโร จ ปน โส ภิกฺขุ อรหตํ อโหสีติ.  สตฺตมํ.}

\csromanheader{2}{8. \textbf{มิตฺตสุตฺต}}
\proseitem{414}{เย ภิกฺขเว อนุกมฺเปยฺยาถ, เย จ โข โสตพฺพํ มญฺเญยฺยุํ มิตฺตา วา อมจฺจา วา ญาตี วา สาโลหิตา วา, เต โว ภิกฺขเว จตุนฺนํ สติปฏฺฐานานํ ภาวนาย สมาทเปตพฺพา นิเวเสตพฺพา ปติฏฺฐาเป\-ตพฺพา.}

\prose{กตเมสํ จตุนฺนํ, อิธ ภิกฺขเว ภิกฺขุ กาเย กายานุปสฺสี วิหรติ อาตาปี สมฺปชาโน สติมา วิเนยฺย โลเก อภิชฺฌา\-โทมนสฺสํ. เวทนาสุ ฯเปฯ. จิตฺเต ฯเปฯ. ธมฺเมสุ ธมฺมานุปสฺสี วิหรติ อาตาปี สมฺปชาโน สติมา วิเนยฺย โลเก อภิชฺฌา\-โทมนสฺสํ. เย ภิกฺขเว อนุกมฺเปยฺยาถ, เย จ โสตพฺพํ มญฺเญยฺยุํ มิตฺตา วา อมจฺจา วา ญาตี วา สาโลหิตา วา, เต โว ภิกฺขเว อิเมสํ จตุนฺนํ สติปฏฺฐานานํ ภาวนาย สมาทเปตพฺพา นิเวเสตพฺพา ปติฏฺฐาเป\-ตพฺพาติ.  อฏฺฐมํ.}

\csromanheader{2}{9. \textbf{เวทนาสุตฺต}}
\proseitem{415}{ติสฺโส อิมา ภิกฺขเว เวทนา. กตมา ติสฺโส, สุขา เวทนา ทุกฺขา เวทนา อทุกฺขมสุขา เวทนา, อิมา โข ภิกฺขเว ติสฺโส เวทนา. อิมาสํ โข ภิกฺขเว ติสฺสนฺนํ เวทนานํ ปริญฺญาย จตฺตาโร สติปฏฺฐานา ภาเวตพฺพา.}

\prose{กตเม จตฺตาโร, อิธ ภิกฺขเว ภิกฺขุ กาเย กายานุปสฺสี วิหรติ อาตาปี สมฺปชาโน สติมา วิเนยฺย โลเก อภิชฺฌา\-โทมนสฺสํ. เวทนาสุ ฯเปฯ. จิตฺเต ฯเปฯ. ธมฺเมสุ ธมฺมานุปสฺสี วิหรติ อาตาปี สมฺปชาโน สติมา วิเนยฺย โลเก อภิชฺฌา\-โทมนสฺสํ. อิมาสํ โข ภิกฺขเว ติสฺสนฺนํ เวทนานํ ปริญฺญาย อิเม จตฺตาโร สติปฏฺฐานา ภาเวตพฺพาติ.  นวมํ.}

\csromanheader{2}{10. \textbf{อาสวสุตฺต}}
\csromanmatikaanchor{mtk.90}
\csromanmatikaanchor{mtk.91}
\csromantocmarknopage{subparagraph}{6. คงฺคาเปยฺยาลวคฺค}{mtk.90}
\bookmark[dest={mtk.90},level=5]{6. คงฺคาเปยฺยาลวคฺค}
\csromantocmarkref{subparagraph}{1-12. คงฺคานทีอาทิสุตฺตทฺวาทสก}{mtk.91}
\bookmark[dest={mtk.91},level=5]{1-12. คงฺคานทีอาทิสุตฺตทฺวาทสก}
\proseitemwithclosermidruled{416}{\csromanfolio{166}ตโยเม ภิกฺขเว อาสวา. กตเม ตโย, กามาสโว ภวาสโว อวิชฺชาสโว, อิเม โข ภิกฺขเว ตโย อาสวา. อิเมสํ โข ภิกฺขเว ติณฺณนฺนํ อาสวานํ ปหานาย จตฺตาโร สติปฏฺฐานา ภาเวตพฺพา. กตเม จตฺตาโร, อิธ ภิกฺขเว ภิกฺขุ กาเย กายานุปสฺสี วิหรติ อาตาปี สมฺปชาโน สติมา วิเนยฺย โลเก อภิชฺฌา\-โทมนสฺสํ. เวทนาสุ ฯเปฯ. จิตฺเต ฯเปฯ. ธมฺเมสุ ธมฺมานุปสฺสี วิหรติ อาตาปี สมฺปชาโน สติมา วิเนยฺย โลเก อภิชฺฌา\-โทมนสฺสํ. อิเมสํ โข ภิกฺขเว ติณฺณนฺนํ อาสวานํ ปหานาย อิเม จตฺตาโร สติปฏฺฐานา ภาเวตพฺพาติ.  ทสมํ.}{อมตวคฺโค ปญฺจโม.}
\csromancenter{\textbf{ตสฺสุทฺทานํ}}
\setlength{\csromangathapagewidth}{0pt}
\setlength{\csromangathawakwidth}{0pt}
\csromangathameasuregroup{อมตํ สมุทโย มคฺโค, \\ ปาติโมกฺขํ ทุจฺจริตํ,}{\csromangathabat{อมตํ สมุทโย มคฺโค,}{สติ กุสลราสิ จ.} \\ \csromangathabat{ปาติโมกฺขํ ทุจฺจริตํ,}{มิตฺตเวทนา อาสเวน จาติ.}}
\csromangathasetleft{อมตํ สมุทโย มคฺโค, \\ ปาติโมกฺขํ ทุจฺจริตํ,}
\csromangathagroup{\csromangathastanza{\csromangathabat{อมตํ สมุทโย มคฺโค,}{สติ กุสลราสิ จ.} \\ \csromangathabat{ปาติโมกฺขํ ทุจฺจริตํ,}{มิตฺตเวทนา อาสเวน จาติ.}}}

\chapterhead{6. คงฺคา\-เปยฺยาล\-วคฺค}
\prose{1-12. คงฺคานที\-อาทิ\-สุตฺต\-ทฺวาทสก}

\proseitem{417-428}{เสยฺยถาปิ ภิกฺขเว คงฺคา นที ปาจีนนินฺนา ปาจีนโปณา ปาจีนปพฺภารา, เอวเมว โข ภิกฺขเว ภิกฺขุ จตฺตาโร สติปฏฺฐาเน ภาเวนฺโต จตฺตาโร สติปฏฺฐาเน พหุลีกโรนฺโต นิพฺพานนินฺโน โหติ นิพฺพานโปโณ นิพฺพาน\-ปพฺภาโร.}

\prose{กถญฺจ ภิกฺขเว ภิกฺขุ จตฺตาโร สติปฏฺฐาเน ภาเวนฺโต จตฺตาโร สติปฏฺฐาเน พหุลีกโรนฺโต นิพฺพานนินฺโน โหติ นิพฺพานโปโณ นิพฺพาน\-ปพฺภาโร, อิธ ภิกฺขเว ภิกฺขุ กาเย กายานุปสฺสี วิหรติ อาตาปี สมฺปชาโน สติมา วิเนยฺย โลเก อภิชฺฌา\-โทมนสฺสํ.}

\prose{\csromanfolio{167}เวทนาสุ ฯเปฯ. จิตฺเต ฯเปฯ. ธมฺเมสุ ธมฺมานุปสฺสี วิหรติ อาตาปี สมฺปชาโน สติมา วิเนยฺย โลเก อภิชฺฌา\-โทมนสฺสํ. เอวํ โข ภิกฺขเว ภิกฺขุ จตฺตาโร สติปฏฺฐาเน ภาเวนฺโต จตฺตาโร สติปฏฺฐาเน พหุลีกโรนฺโต นิพฺพานนินฺโน โหติ นิพฺพานโปโณ นิพฺพาน\-ปพฺภาโรติ วิตฺถาเรตพฺพํ. คงฺคา\-เปยฺยาล\-วคฺโค ฉฏฺโฐ.}

\vspace{\tipitakaparskip}
\csromancenter{\textbf{ตสฺสุทฺทานํ}}
\setlength{\csromangathapagewidth}{0pt}
\setlength{\csromangathawakwidth}{0pt}
\csromangathameasuregroup{ฉ ปาจีนโต นินฺนา, \\ เอเต เทฺว ฉ ทฺวาทส โหนฺติ,}{\csromangathabat{ฉ ปาจีนโต นินฺนา,}{ฉ นินฺนา จ สมุทฺทโต.} \\ \csromangathabat{เอเต เทฺว ฉ ทฺวาทส โหนฺติ,}{วคฺโค เตน ปวุจฺจตีติ.}}
\csromangathasetleft{ฉ ปาจีนโต นินฺนา, \\ เอเต เทฺว ฉ ทฺวาทส โหนฺติ,}
\csromangathagroup{\csromangathastanza{\csromangathabat{ฉ ปาจีนโต นินฺนา,}{ฉ นินฺนา จ สมุทฺทโต.} \\ \csromangathabat{เอเต เทฺว ฉ ทฺวาทส โหนฺติ,}{วคฺโค เตน ปวุจฺจตีติ.}}}

\csromanheader{5}{7. \textbf{อปฺปมาทวคฺค}}
\csromanmatikaanchor{mtk.92}
\csromanmatikaanchor{mtk.93}
\csromantocmarknopage{subparagraph}{7. อปฺปมาทวคฺค}{mtk.92}
\bookmark[dest={mtk.92},level=5]{7. อปฺปมาทวคฺค}
\csromantocmarkref{subparagraph}{1-10. ตถาคตาทิสุตฺต}{mtk.93}
\bookmark[dest={mtk.93},level=5]{1-10. ตถาคตาทิสุตฺต}
\csromanheader{6}{1-10. ตถาคตา\-ทิ\-สุตฺต}
\proseitemwithclosermidruled{429-438}{ยาวตา ภิกฺขเว สตฺตา อปทา วา ทฺวิปทา วา จตุปฺปทา วา พหุปฺปทา วาติ วิตฺถาเรตพฺพํ.}{อปฺปมาทวคฺโค สตฺตโม.}
\csromancenter{\textbf{ตสฺสุทฺทานํ}}
\setlength{\csromangathapagewidth}{0pt}
\setlength{\csromangathawakwidth}{0pt}
\csromangathameasuregroup{ตถาคต ปทํ กูฏํ, \\ ราชา จนฺทิมสูริยา,}{\csromangathabat{ตถาคต ปทํ กูฏํ,}{มูลํ สาโร จ วสฺสิกํ.} \\ \csromangathabat{ราชา จนฺทิมสูริยา,}{วตฺเถน ทสมํ ปทนฺติ.}}
\csromangathasetleft{ตถาคต ปทํ กูฏํ, \\ ราชา จนฺทิมสูริยา,}
\csromangathagroup{\csromangathastanza{\csromangathabat{ตถาคต ปทํ กูฏํ,}{มูลํ สาโร จ วสฺสิกํ.} \\ \csromangathabat{ราชา จนฺทิมสูริยา,}{วตฺเถน ทสมํ ปทนฺติ.}}}

\chapterhead{8. พลกรณีย\-วคฺค}
\csromanmatikaanchor{mtk.94}
\csromanmatikaanchor{mtk.95}
\csromantocmarknopage{subparagraph}{8. พลกรณียวคฺค}{mtk.94}
\bookmark[dest={mtk.94},level=5]{8. พลกรณียวคฺค}
\csromantocmarkref{subparagraph}{1-12. พลาทิสุตฺต}{mtk.95}
\bookmark[dest={mtk.95},level=5]{1-12. พลาทิสุตฺต}
\csromanheader{6}{1-12. พลาทิสุตฺต}
\proseitem{439-450}{เสยฺยถาปิ ภิกฺขเว เย เกจิ พลกรณียา กมฺมนฺตา กรียนฺตีติ วิตฺถาเรตพฺพํ. พลกรณีย\-วคฺโค อฏฺฐโม.}
\csromansectionrule

\vspace{\tipitakaparskip}
\csromancenter{\textbf{ตสฺสุทฺทานํ}}
\setlength{\csromangathapagewidth}{0pt}
\setlength{\csromangathawakwidth}{0pt}
\csromangathameasuregroup{พลํ พีชญฺจ นาโค จ, \\ อากาเสน จ เทฺว เมฆา,}{\csromangathabat{พลํ พีชญฺจ นาโค จ,}{รุกฺโข กุมฺเภน สูกิยา.} \\ \csromangathabat{อากาเสน จ เทฺว เมฆา,}{นาวา อาคนฺตุกา นทีติ.}}
\csromangathasetleft{พลํ พีชญฺจ นาโค จ, \\ อากาเสน จ เทฺว เมฆา,}
\csromangathagroup{\csromangathastanza{\csromanfolio{168}\csromangathabat{พลํ พีชญฺจ นาโค จ,}{รุกฺโข กุมฺเภน สูกิยา.} \\ \csromangathabat{อากาเสน จ เทฺว เมฆา,}{นาวา อาคนฺตุกา นทีติ.}}}

\chapterhead{9. \textbf{เอสนาวคฺค}}
\csromanmatikaanchor{mtk.96}
\csromanmatikaanchor{mtk.97}
\csromantocmarknopage{subparagraph}{9. เอสนาวคฺค}{mtk.96}
\bookmark[dest={mtk.96},level=5]{9. เอสนาวคฺค}
\csromantocmarkref{subparagraph}{1-10. เอสนาทิสุตฺต}{mtk.97}
\bookmark[dest={mtk.97},level=5]{1-10. เอสนาทิสุตฺต}
\csromanheader{6}{1-10. เอสนาทิสุตฺต}
\proseitemwithclosermidruled{451-460}{ติสฺโส อิมา ภิกฺขเว เอสนา. กตมา ติสฺโส, กาเมสนา ภเวสนา พฺรหฺมจริเย\-สนาติ วิตฺถาเรตพฺพํ.}{เอสนาวคฺโค นวโม.}
\csromancenter{\textbf{ตสฺสุทฺทานํ}}
\setlength{\csromangathapagewidth}{0pt}
\setlength{\csromangathawakwidth}{0pt}
\csromangathameasuregroup{เอสนา วิธา อาสโว, \\ ขิลํ มลญฺจ นีโฆ จ,}{\csromangathabat{เอสนา วิธา อาสโว,}{ภโว จ ทุกฺขตา ติสฺโส.} \\ \csromangathabat{ขิลํ มลญฺจ นีโฆ จ,}{เวทนา ตณฺหา ตสินาย จาติ.}}
\csromangathasetleft{เอสนา วิธา อาสโว, \\ ขิลํ มลญฺจ นีโฆ จ,}
\csromangathagroup{\csromangathastanza{\csromangathabat{เอสนา วิธา อาสโว,}{ภโว จ ทุกฺขตา ติสฺโส.} \\ \csromangathabat{ขิลํ มลญฺจ นีโฆ จ,}{เวทนา ตณฺหา ตสินาย จาติ.}}}

\chapterhead{10. \textbf{โอฆวคฺค}}
\csromanmatikaanchor{mtk.98}
\csromanmatikaanchor{mtk.99}
\csromantocmarknopage{subparagraph}{10. โอฆวคฺค}{mtk.98}
\bookmark[dest={mtk.98},level=5]{10. โอฆวคฺค}
\csromantocmarkref{subparagraph}{1-10. อุทฺธมฺภาคิยาทิสุตฺต}{mtk.99}
\bookmark[dest={mtk.99},level=5]{1-10. อุทฺธมฺภาคิยาทิสุตฺต}
\csromanheader{6}{1-10. อุทฺธมฺภาคิยา\-ทิ\-สุตฺต}
\proseitem{461-470}{ปญฺจิมานิ ภิกฺขเว อุทฺธมฺ\-ภาคิยานิ สํโยชนานิ. กตมานิ ปญฺจ, รูปราโค อรูปราโค มาโน อุทฺธจฺจํ อวิชฺชา, อิมานิ โข ภิกฺขเว ปญฺจุ\-ทฺธมฺภาคิยานิ สํโยชนานิ. อิเมสํ โข ภิกฺขเว ปญฺจนฺนํ อุทฺธมฺ\-ภาคิยานํ สํโยชนานํ อภิญฺญาย ปริญฺญาย ปริกฺขยาย ปหานาย จตฺตาโร สติปฏฺฐานา ภาเวตพฺพา.}

\prosewithclosermidruled{\csromanfolio{169}กตเม จตฺตาโร, อิธ ภิกฺขเว ภิกฺขุ กาเย กายานุปสฺสี วิหรติ อาตาปี สมฺปชาโน สติมา วิเนยฺย โลเก อภิชฺฌา\-โทมนสฺสํ. เวทนาสุ ฯเปฯ. จิตฺเต ฯเปฯ. ธมฺเมสุ ธมฺมานุปสฺสี วิหรติ อาตาปี สมฺปชาโน สติมา วิเนยฺย โลเก อภิชฺฌา\-โทมนสฺสํ. อิเมสํ โข ภิกฺขเว ปญฺจนฺนํ อุทฺธมฺ\-ภาคิยานํ สํโยชนานํ อภิญฺญาย ปริญฺญาย ปริกฺขยาย ปหานาย อิเม จตฺตาโร สติปฏฺฐานา ภาเวตพฺพาติ.  ทสมํ.}{(ยถา มคฺคสํยุตฺตํ, ตถา สติปฏฺฐาน\-สํยุตฺตํ วิตฺถาเรตพฺพํ.) โอฆวคฺโค ทสโม.}
\csromancenter{\textbf{ตสฺสุทฺทานํ}}
\setlength{\csromangathapagewidth}{0pt}
\setlength{\csromangathawakwidth}{0pt}
\csromangathameasuregroup{โอโฆ โยโค อุปาทานํ, \\ กามคุณา นีวรณา,}{\csromangathabat{โอโฆ โยโค อุปาทานํ,}{คนฺถา อนุสเยน จ.} \\ \csromangathabat{กามคุณา นีวรณา,}{ขนฺธา โอรุทฺธมฺภาคิยาติ.} \\ สติปฏฺฐานสํยุตฺตํ ตติยํ.}
\csromangathameasurewak{สติปฏฺฐานสํยุตฺตํ ตติยํ.}
\csromangathasetleft{โอโฆ โยโค อุปาทานํ, \\ กามคุณา นีวรณา,}
\csromangathagroup{\csromangathastanza{\csromangathabat{โอโฆ โยโค อุปาทานํ,}{คนฺถา อนุสเยน จ.} \\ \csromangathabat{กามคุณา นีวรณา,}{ขนฺธา โอรุทฺธมฺภาคิยาติ.}}\csromangathastanza{\csromangathawak{สติปฏฺฐาน\-สํยุตฺตํ ตติยํ.}}}

\csromanmatikaanchor{mtk.100}
\csromanmatikaanchor{mtk.101}
\csromantocmarknopage{subparagraph}{4. อินฺทฺริยสุํยุตฺต}{mtk.100}
\bookmark[dest={mtk.100},level=5]{4. อินฺทฺริยสุํยุตฺต}
\csromantocmarknopage{subparagraph}{1. สุทฺธิกวคฺค}{mtk.101}
\bookmark[dest={mtk.101},level=5]{1. สุทฺธิกวคฺค}
\csromanheader{6}{1. \textbf{สุทฺธิกวคฺค}}
\csromanheader{2}{1. \textbf{สุทฺธิกสุตฺต}}
\csromanmatikaanchor{mtk.102}
\csromantocmarkref{subparagraph}{1-7. สุทฺธิกสุตฺตาทีนิ}{mtk.102}
\bookmark[dest={mtk.102},level=5]{1-7. สุทฺธิกสุตฺตาทีนิ}
\proseitem{471}{\csromanfolio{170}สาวตฺถิ\-นิทานํ. ตตฺร โข ภควา เอตทโวจ\csromandash{}ปญฺจิมานิ ภิกฺขเว อินฺทฺริยานิ. กตมานิ ปญฺจ, สทฺธินฺทฺริยํ วีริยินฺทฺริยํ สตินฺทฺริยํ สมาธินฺทฺริยํ ปญฺญินฺทฺริยํ. อิมานิ โข ภิกฺขเว ปญฺจิ\-นฺทฺริยานีติ.  ปฐมํ.}

\csromanheader{2}{2. ปฐม\-โสตาปนฺน\-สุตฺต}
\proseitem{472}{ปญฺจิมานิ ภิกฺขเว อินฺทฺริยานิ. กตมานิ ปญฺจ, สทฺธินฺทฺริยํ วีริยินฺทฺริยํ สตินฺทฺริยํ สมาธินฺทฺริยํ ปญฺญินฺทฺริยํ. ยโต โข ภิกฺขเว อริยสาวโก อิเมสํ ปญฺจนฺนํ อินฺทฺริยานํ อสฺสาทญฺจ\csromansharedfootnote{b921b536220ec6e9}{สมุทยญฺจ อตฺถงฺคมญฺจ อสฺสาทญฺจ (สฺยา, กํ, อิ, ก) สํ 1. 433 ปิฏฺเฐปิ.} อาทีนวญฺจ นิสฺสรณญฺจ ยถาภูตํ ปชานาติ. อยํ วุจฺจติ ภิกฺขเว อริยสาวโก โสตาปนฺโน อวินิปาต\-ธมฺโม นิยโต สมฺโพธิ\-ปรายโณติ.  ทุติยํ.}

\csromanheader{2}{3. ทุติย\-โสตาปนฺน\-สุตฺต}
\proseitem{473}{ปญฺจิมานิ ภิกฺขเว อินฺทฺริยานิ. กตมานิ ปญฺจ, สทฺธินฺทฺริยํ วีริยินฺทฺริยํ สตินฺทฺริยํ สมาธินฺทฺริยํ ปญฺญินฺทฺริยํ. ยโต โข ภิกฺขเว อริยสาวโก อิเมสํ ปญฺจนฺนํ อินฺทฺริยานํ สมุทยญฺจ อตฺถงฺคมญฺจ อสฺสาทญฺจ อาทีนวญฺจ นิสฺสรณญฺจ ยถาภูตํ ปชานาติ. อยํ วุจฺจติ ภิกฺขเว อริยสาวโก โสตาปนฺโน อวินิปาต\-ธมฺโม นิยโต สมฺโพธิ\-ปรายโณติ.  ตติยํ.}

\prose{4. ปฐม\-อรหนฺต\-สุตฺต}

\proseitem{474}{ปญฺจิมานิ ภิกฺขเว อินฺทฺริยานิ. กตมานิ ปญฺจ, สทฺธินฺทฺริยํ วีริยินฺทฺริยํ สตินฺทฺริยํ สมาธินฺทฺริยํ ปญฺญินฺทฺริยํ. ยโต โข ภิกฺขเว อริยสาวโก อิเมสํ ปญฺจนฺนํ อินฺทฺริยานํ อสฺสาทญฺจ\csromansharedfootnote{b921b536220ec6e9}{สมุทยญฺจ อตฺถงฺคมญฺจ อสฺสาทญฺจ (สฺยา, กํ, อิ, ก) สํ 1. 433 ปิฏฺเฐปิ.} อาทีนวญฺจ นิสฺสรณญฺจ ยถาภูตํ}

\prose{\csromanfolio{171}วิทิตฺวา อนุปาทาวิมุตฺโต โหติ. อยํ วุจฺจติ ภิกฺขเว ภิกฺขุ อรหํ ขีณาสโว วุสิตวา กตกรณีโย โอหิตภาโร อนุปฺปตฺต\-สทตฺโถ ปริกฺขีณ\-ภวสํโยชโน สมฺมทญฺญา\-วิมุตฺโตติ.  จตุตฺถํ.}

\prose{5. ทุติย\-อรหนฺต\-สุตฺต}

\proseitem{475}{ปญฺจิมานิ ภิกฺขเว อินฺทฺริยานิ. กตมานิ ปญฺจ, สทฺธินฺทฺริยํ วีริยินฺทฺริยํ สตินฺทฺริยํ สมาธินฺทฺริยํ ปญฺญินฺทฺริยํ. ยโต โข ภิกฺขเว ภิกฺขุ อิเมสํ ปญฺจนฺนํ อินฺทฺริยานํ สมุทยญฺจ อตฺถงฺคมญฺจ อสฺสาทญฺจ อาทีนวญฺจ นิสฺสรณญฺจ ยถาภูตํ วิทิตฺวา อนุปาทาวิมุตฺโต โหติ, อยํ วุจฺจติ ภิกฺขเว ภิกฺขุ อรหํ ขีณาสโว วุสิตวา กตกรณีโย โอหิตภาโร อนุปฺปตฺต\-สทตฺโถ ปริกฺขีณ\-ภวสํโยชโน สมฺมทญฺญา\-วิมุตฺโตติ.  ปญฺจมํ.}

\csromanheader{2}{6. ปฐม\-สมณพฺราหฺมณ\-สุตฺต}
\proseitem{476}{ปญฺจิมานิ ภิกฺขเว อินฺทฺริยานิ. กตมานิ ปญฺจ สทฺธินฺทฺริยํ ฯเปฯ ปญฺญินฺทฺริยํ. เย หิ เกจิ ภิกฺขเว สมณา วา พฺราหฺมณา วา อิเมสํ ปญฺจนฺนํ อินฺทฺริยานํ อสฺสาทญฺจ อาทีนวญฺจ นิสฺสรณญฺจ ยถาภูตํ นปฺปชานนฺติ, น เม เต ภิกฺขเว สมณา วา พฺราหฺมณา วา สมเณสุ วา สมณสมฺมตา, พฺราหฺมเณสุ วา พฺราหฺมณ\-สมฺมตา, น จ ปเนเต อายสฺมนฺโต สามญฺญตฺถํ วา พฺรหฺมญฺญตฺถํ วา ทิฏฺเฐว ธมฺเม สยํ อภิญฺญา สจฺฉิกตฺวา อุปสมฺปชฺช วิหรนฺติ.}

\prose{เย จ โข เกจิ\csromansharedfootnote{a47e3426384ac83e}{เย จ โข เต (สฺยา, กํ, ก) สํ 1. 432 ปิฏฺเฐปิ.} ภิกฺขเว สมณา วา พฺราหฺมณา วา อิเมสํ ปญฺจนฺนํ อินฺทฺริยานํ สมุทยญฺจ อตฺถงฺคมญฺจ อสฺสาทญฺจ อาทีนวญฺจ นิสฺสรณญฺจ ยถาภูตํ ปชานนฺติ, เต โข เม ภิกฺขเว สมณา วา พฺราหฺมณา วา สมเณสุ เจว สมณสมฺมตา, พฺราหฺมเณสุ จ พฺราหฺมณ\-สมฺมตา, เต จ ปนายสฺมนฺโต สามญฺญตฺถญฺจ พฺรหฺมญฺญตฺถญฺจ ทิฏฺเฐว ธมฺเม สยํ อภิญฺญา สจฺฉิกตฺวา อุปสมฺปชฺช วิหรนฺตีติ.  ฉฏฺฐํ.}

\csromanheader{2}{7. ทุติย\-สมณพฺราหฺมณ\-สุตฺต}
\csromanmatikaanchor{mtk.103}
\csromanmatikaanchor{mtk.104}
\csromantocmarkref{subparagraph}{8. ทฏฺฐพฺพสุตฺต}{mtk.103}
\bookmark[dest={mtk.103},level=5]{8. ทฏฺฐพฺพสุตฺต}
\csromantocmarkref{subparagraph}{9-10. วิภงฺคสุตฺตานิ}{mtk.104}
\bookmark[dest={mtk.104},level=5]{9-10. วิภงฺคสุตฺตานิ}
\proseitem{477}{เย หิ เกจิ ภิกฺขเว สมณา วา พฺราหฺมณา วา สทฺธินฺทฺริยํ นปฺปชานนฺติ, สทฺธินฺทฺริย\-สมุทยํ นปฺปชานนฺติ, สทฺธินฺทฺริย\-นิโรธํ นปฺปชานนฺติ, สทฺธินฺทฺริยนิโรธ\-คามินิํ ปฏิปทํ นปฺปชานนฺติ. วีริยินฺทฺริยํ นปฺปชานนฺติ ฯเปฯ. สตินฺทฺริยํ \csromanfolio{172}นปฺปชานนฺติ ฯเปฯ. สมาธินฺทฺริยํ นปฺปชานนฺติ ฯเปฯ. ปญฺญินฺทฺริยํ นปฺปชานนฺติ, ปญฺญินฺทฺริย\-สมุทยํ นปฺปชานนฺติ, ปญฺญินฺทฺริย\-นิโรธํ นปฺปชานนฺติ, ปญฺญินฺทฺริยนิโรธ\-คามินิํ ปฏิปทํ นปฺปชานนฺติ. น เม เต ภิกฺขเว สมณา วา พฺราหฺมณา วา สมเณสุ วา สมณสมฺมตา, พฺราหฺมเณสุ วา พฺราหฺมณ\-สมฺมตา, น จ ปเนเต อายสฺมนฺโต สามญฺญตฺถํ วา พฺราหฺมญฺญตฺถํ วา ทิฏฺเฐว ธมฺเม สยํ อภิญฺญา สจฺฉิกตฺวา อุปสมฺปชฺช วิหรนฺติ.}

\prose{เย จ โข เกจิ ภิกฺขเว สมณา วา พฺราหฺมณา วา สทฺธินฺทฺริยํ ปชานนฺติ, สทฺธินฺทฺริย\-สมุทยํ ปชานนฺติ, สทฺธินฺทฺริย\-นิโรธํ ปชานนฺติ, สทฺธินฺทฺริยนิโรธ\-คามินิํ ปฏิปทํ ปชานนฺติ. วีริยินฺทฺริยํ ปชานนฺติ, วีริยินฺทฺริย\-สมุทยํ ปชานนฺติ, วีริยินฺทฺริย\-นิโรธํ ปชานนฺติ, วีริยินฺทฺริยนิโรธ\-คามินิํ ปฏิปทํ ปชานนฺติ. สตินฺทฺริยํ ปชานนฺติ ฯเปฯ. สมาธินฺทฺริยํ ปชานนฺติ ฯเปฯ. ปญฺญินฺทฺริยํ ปชานนฺติ, ปญฺญินฺทฺริย\-สมุทยํ ปชานนฺติ, ปญฺญินฺทฺริย\-นิโรธํ ปชานนฺติ, ปญฺญินฺทฺริยนิโรธ\-คามินิํ ปฏิปทํ ปชานนฺติ. เต โข เม ภิกฺขเว สมณา วา พฺราหฺมณา วา สมเณสุ เจว สมณสมฺมตา, พฺราหฺมเณสุ จ พฺราหฺมณ\-สมฺมตา, เต จ ปนายสฺมนฺโต สามญฺญตฺถญฺจ พฺราหฺมญฺญตฺถญฺจ ทิฏฺเฐว ธมฺเม สยํ อภิญฺญา สจฺฉิกตฺวา อุปสมฺปชฺช วิหรนฺตีติ.  สตฺตมํ.}

\csromanheader{6}{8. \textbf{ทฏฺฐพฺพสุตฺต}}
\proseitem{478}{ปญฺจิมานิ ภิกฺขเว อินฺทฺริยานิ. กตมานิ ปญฺจ, สทฺธินฺทฺริยํ ฯเปฯ ปญฺญินฺทฺริยํ. กตฺถ จ ภิกฺขเว สทฺธินฺทฺริยํ ทฏฺฐพฺพํ, จตูสุ โสตาปตฺติยงฺเคสุ เอตฺถ สทฺธินฺทฺริยํ ทฏฺฐพฺพํ. กตฺถ จ ภิกฺขเว วีริยินฺทฺริยํ ทฏฺฐพฺพํ, จตูสุ สมฺม\-ปฺปธาเนสุ เอตฺถ วีริยินฺทฺริยํ ทฏฺฐพฺพํ. กตฺถ จ ภิกฺขเว สตินฺทฺริยํ ทฏฺฐพฺพํ, จตูสุ สติปฏฺฐาเนสุ เอตฺถ สตินฺทฺริยํ ทฏฺฐพฺพํ. กตฺถ จ ภิกฺขเว สมาธินฺทฺริยํ ทฏฺฐพฺพํ, จตูสุ ฌาเนสุ เอตฺถ สมาธินฺทฺริยํ ทฏฺฐพฺพํ. กตฺถ จ ภิกฺขเว ปญฺญินฺทฺริยํ ทฏฺฐพฺพํ, จตูสุ อริยสจฺเจสุ เอตฺถ ปญฺญินฺทฺริยํ ทฏฺฐพฺพํ. อิมานิ โข ภิกฺขเว ปญฺจิ\-นฺทฺริยานีติ.  อฏฺฐมํ.}

\csromanheader{2}{9. ปฐม\-วิภงฺค\-สุตฺต}
\proseitem{479}{ปญฺจิมานิ ภิกฺขเว อินฺทฺริยานิ. กตมานิ ปญฺจ, สทฺธินฺทฺริยํ ฯเปฯ ปญฺญินฺทฺริยํ. กตมญฺจ ภิกฺขเว สทฺธินฺทฺริยํ, อิธ ภิกฺขเว อริยสาวโก สทฺโธ โหติ, สทฺทหติ ตถาคตสฺส โพธิํ “อิติปิ โส ภควา อรหํ สมฺมาสมฺพุทฺโธ}

\prose{\csromanfolio{173}วิชฺชา\-จรณ\-สมฺปนฺโน สุคโต โลกวิทู อนุตฺตโร ปุริส\-ทมฺม\-สารถิ สตฺถา เทวมนุสฺสานํ พุทฺโธ ภควา”ติ. อิทํ วุจฺจติ ภิกฺขเว สทฺธินฺทฺริยํ.}

\prose{กตมญฺจ ภิกฺขเว วีริยินฺทฺริยํ, อิธ ภิกฺขเว อริยสาวโก อารทฺธวีริโย วิหรติ อกุสลานํ ธมฺมานํ ปหานาย กุสลานํ ธมฺมานํ อุปสมฺปทาย ถามวา ทฬฺหปรกฺกโม อนิกฺขิตฺตธุโร กุสเลสุ ธมฺเมสุ. อิทํ วุจฺจติ ภิกฺขเว วีริยินฺทฺริยํ.}

\prose{กตมญฺจ ภิกฺขเว สตินฺทฺริยํ, อิธ ภิกฺขเว อริยสาวโก สติมา โหติ ปรเมน สติเนปกฺเกน สมนฺนาคโต, จิรกตมฺปิ จิรภาสิตมฺปิ สริตา อนุสฺสริตา. อิทํ วุจฺจติ ภิกฺขเว สตินฺทฺริยํ.}

\prose{กตมญฺจ ภิกฺขเว สมาธินฺทฺริยํ, อิธ ภิกฺขเว อริยสาวโก โวสฺสคฺคา\-รมฺมณํ กริตฺวา ลภติ สมาธิํ, ลภติ จิตฺตสฺส เอกคฺคตํ. อิทํ วุจฺจติ ภิกฺขเว สมาธินฺทฺริยํ.}

\prose{กตมญฺจ ภิกฺขเว ปญฺญินฺทฺริยํ, อิธ ภิกฺขเว อริยสาวโก ปญฺญวา โหติ อุทย\-ตฺถคามินิยา ปญฺญาย สมนฺนาคโต อริยาย นิพฺเพธิกาย สมฺมา ทุกฺข\-กฺขย\-คามินิยา. อิทํ วุจฺจติ ภิกฺขเว ปญฺญินฺทฺริยํ. อิมานิ โข ภิกฺขเว ปญฺจิ\-นฺทฺริยานีติ.  นวมํ.}

\csromanheader{2}{10. ทุติย\-วิภงฺค\-สุตฺต}
\proseitem{480}{ปญฺจิมานิ ภิกฺขเว อินฺทฺริยานิ. กตมานิ ปญฺจ, สทฺธินฺทฺริยํ ฯเปฯ ปญฺญินฺทฺริยํ. กตมญฺจ ภิกฺขเว สทฺธินฺทฺริยํ, อิธ ภิกฺขเว อริยสาวโก สทฺโธ โหติ, สทฺทหติ ตถาคตสฺส โพธิํ “อิติปิ โส ภควา อรหํ สมฺมาสมฺพุทฺโธ วิชฺชา\-จรณ\-สมฺปนฺโน สุคโต โลกวิทู อนุตฺตโร ปุริส\-ทมฺม\-สารถิ สตฺถา เทวมนุสฺสานํ พุทฺโธ ภควา”ติ. อิทํ วุจฺจติ ภิกฺขเว สทฺธินฺทฺริยํ.}

\prose{กตมญฺจ ภิกฺขเว วีริยินฺทฺริยํ, อิธ ภิกฺขเว อริยสาวโก อารทฺธวีริโย วิหรติ อกุสลานํ ธมฺมานํ ปหานาย กุสลานํ ธมฺมานํ อุปสมฺปทาย ถามวา ทฬฺหปรกฺกโม อนิกฺขิตฺตธุโร กุสเลสุ ธมฺเมสุ, โส อนุปฺปนฺนานํ ปาปกานํ อกุสลานํ ธมฺมานํ อนุปฺปาทาย ฉนฺทํ ชเนติ วายมติ วีริยํ อารภติ จิตฺตํ ปคฺคณฺหาติ ปทหติ, อุปฺปนฺนานํ ปาปกานํ อกุสลานํ ธมฺมานํ ปหานาย ฉนฺทํ ชเนติ วายมติ \csromanfolio{174}วีริยํ อารภติ จิตฺตํ ปคฺคณฺหาติ ปทหติ, อนุปฺปนฺนานํ กุสลานํ ธมฺมานํ อุปฺปาทาย ฉนฺทํ ชเนติ วายมติ วีริยํ อารภติ จิตฺตํ ปคฺคณฺหาติ ปทหติ, อุปฺปนฺนานํ กุสลานํ ธมฺมานํ ฐิติยา\csromansharedfootnote{6e8c04efcbf4f73c}{สมาปตฺติยา (สฺยา, กํ, ก)} อสมฺโมสาย ภิยฺโยภาวาย เวปุลฺลาย ภาวนาย ปาริปูริยา ฉนฺทํ ชเนติ วายมติ วีริยํ อารภติ จิตฺตํ ปคฺคณฺหาติ ปทหติ. อิทํ วุจฺจติ ภิกฺขเว วีริยินฺทฺริยํ.}

\prose{กตมญฺจ ภิกฺขเว สตินฺทฺริยํ, อิธ ภิกฺขเว อริยสาวโก สติมา โหติ ปรเมน สติเนปกฺเกน สมนฺนาคโต, จิรกตมฺปิ จิรภาสิตมฺปิ สริตา อนุสฺสริตา, โส กาเย กายานุปสฺสี วิหรติ อาตาปี สมฺปชาโน สติมา วิเนยฺย โลเก อภิชฺฌา\-โทมนสฺสํ. เวทนาสุ ฯเปฯ. จิตฺเต ฯเปฯ. ธมฺเมสุ ธมฺมานุปสฺสี วิหรติ อาตาปี สมฺปชาโน สติมา วิเนยฺย โลเก อภิชฺฌา\-โทมนสฺสํ. อิทํ วุจฺจติ ภิกฺขเว สตินฺทฺริยํ.}

\prose{กตมญฺจ ภิกฺขเว สมาธินฺทฺริยํ, อิธ ภิกฺขเว อริยสาวโก โวสฺสคฺคา\-รมฺมณํ กริตฺวา ลภติ สมาธิํ, ลภติ จิตฺตสฺส เอกคฺคตํ, โส วิวิจฺเจว กาเมหิ วิวิจฺจ อกุสเลหิ ธมฺเมหิ สวิตกฺกํ สวิจารํ วิเวกชํ ปีติสุขํ ปฐมํ ฌานํ อุปสมฺปชฺช วิหรติ, วิตกฺก\-วิจารานํ วูปสมา อชฺฌตฺตํ สมฺปสาทนํ เจตโส เอโกทิภาวํ อวิตกฺกํ อวิจารํ สมาธิชํ ปีติสุขํ ทุติยํ ฌานํ อุปสมฺปชฺช วิหรติ, ปีติยา จ วิราคา อุเปกฺขโก จ วิหรติ สโต จ สมฺปชาโน สุขญฺจ กาเยน ปฏิสํเวเทติ, ยํ ตํ อริยา อาจิกฺขนฺติ “อุเปกฺขโก สติมา สุขวิหารี”ติ, ตติยํ ฌานํ อุปสมฺปชฺช วิหรติ. สุขสฺส จ ปหานา ทุกฺขสฺส จ ปหานา ปุพฺเพว โสมนสฺส\-โทมนสฺสานํ อตฺถงฺคมา อทุกฺขมสุขํ อุเปกฺขา\-สติ\-ปาริสุทฺธิํ จตุตฺถํ ฌานํ อุปสมฺปชฺช วิหรติ. อิทํ วุจฺจติ ภิกฺขเว สมาธินฺทฺริยํ.}

\prose{กตมญฺจ ภิกฺขเว ปญฺญินฺทฺริยํ, อิธ ภิกฺขเว อริยสาวโก ปญฺญวา โหติ อุทย\-ตฺถคามินิยา ปญฺญาย สมนฺนาคโต อริยาย นิพฺเพธิกาย สมฺมา ทุกฺข\-กฺขย\-คามินิยา, โส “อิทํ ทุกฺขนฺ”ติ ยถาภูตํ ปชานาติ, “อยํ ทุกฺขสมุทโย”ติ ยถาภูตํ ปชานาติ, “อยํ ทุกฺขนิโรโธ”ติ ยถาภูตํ ปชานาติ. “อยํ ทุกฺขนิโรธ\-คามินี ปฏิปทา”ติ}

\prosewithclosermidruled{\csromanfolio{175}ยถาภูตํ ปชานาติ. อิทํ วุจฺจติ ภิกฺขเว ปญฺญินฺทฺริยํ. อิมานิ โข ภิกฺขเว ปญฺจิ\-นฺทฺริยานีติ.  ทสมํ.}{สุทฺธิกวคฺโค ปฐโม.}
\csromancenter{\textbf{ตสฺสุทฺทานํ}}
\setlength{\csromangathapagewidth}{0pt}
\setlength{\csromangathawakwidth}{0pt}
\csromangathameasuregroup{สุทฺธิกญฺเจว เทฺว โสตา, \\ สมณพฺราหฺมณา ทฏฺฐพฺพํ,}{\csromangathabat{สุทฺธิกญฺเจว เทฺว โสตา,}{อรหนฺตา อปเร ทุเว.} \\ \csromangathabat{สมณพฺราหฺมณา ทฏฺฐพฺพํ,}{วิภงฺคา อปเร ทุเวติ.}}
\csromangathasetleft{สุทฺธิกญฺเจว เทฺว โสตา, \\ สมณพฺราหฺมณา ทฏฺฐพฺพํ,}
\csromangathagroup{\csromangathastanza{\csromangathabat{สุทฺธิกญฺเจว เทฺว โสตา,}{อรหนฺตา อปเร ทุเว.} \\ \csromangathabat{สมณพฺราหฺมณา ทฏฺฐพฺพํ,}{วิภงฺคา อปเร ทุเวติ.}}}

\csromanmatikaanchor{mtk.105}
\csromantocmarknopage{subparagraph}{2. มุทุตรวคฺค}{mtk.105}
\bookmark[dest={mtk.105},level=5]{2. มุทุตรวคฺค}
\csromanheader{6}{2. \textbf{มุทุตรวคฺค}}
\csromanmatikaanchor{mtk.106}
\csromantocmarkref{subparagraph}{1. ปฏิลาภสุตฺต}{mtk.106}
\bookmark[dest={mtk.106},level=5]{1. ปฏิลาภสุตฺต}
\csromanheader{6}{1. \textbf{ปฏิลาภสุตฺต}}
\proseitem{481}{ปญฺจิมานิ ภิกฺขเว อินฺทฺริยานิ. กตมานิ ปญฺจ, สทฺธินฺทฺริยํ ฯเปฯ ปญฺญินฺทฺริยํ. กตมญฺจ ภิกฺขเว สทฺธินฺทฺริยํ, อิธ ภิกฺขเว อริยสาวโก สทฺโธ โหติ, สทฺทหติ ตถาคตสฺส โพธิํ “อิติปิ โส ภควา อรหํ สมฺมาสมฺพุทฺโธ วิชฺชา\-จรณ\-สมฺปนฺโน สุคโต โลกวิทู อนุตฺตโร ปุริส\-ทมฺม\-สารถิ สตฺถา เทวมนุสฺสานํ พุทฺโธ ภควา”ติ. อิทํ วุจฺจติ ภิกฺขเว สทฺธินฺทฺริยํ.}

\prose{กตมญฺจ ภิกฺขเว วีริยินฺทฺริยํ, ยํ โข ภิกฺขเว จตฺตาโร สมฺมปฺปธาเน อารพฺภ วีริยํ ปฏิลภติ. อิทํ วุจฺจติ ภิกฺขเว วีริยินฺทฺริยํ.}

\prose{กตมญฺจ ภิกฺขเว สตินฺทฺริยํ, ยํ โข ภิกฺขเว จตฺตาโร สติปฏฺฐาเน อารพฺภ สติํ ปฏิลภติ. อิทํ วุจฺจติ ภิกฺขเว สตินฺทฺริยํ.}

\prose{กตมญฺจ ภิกฺขเว สมาธินฺทฺริยํ, อิธ ภิกฺขเว อริยสาวโก โวสฺสคฺคา\-รมฺมณํ กริตฺวา ลภติ สมาธิํ, ลภติ จิตฺตสฺส เอกคฺคตํ. อิทํ วุจฺจติ ภิกฺขเว สมาธินฺทฺริยํ.}

\prose{กตมญฺจ ภิกฺขเว ปญฺญินฺทฺริยํ, อิธ ภิกฺขเว อริยสาวโก ปญฺญวา โหติ อุทย\-ตฺถคามินิยา ปญฺญาย สมนฺนาคโต อริยาย นิพฺเพธิกาย สมฺมา ทุกฺข\-กฺขย\-คามินิยา. อิทํ วุจฺจติ ภิกฺขเว ปญฺญินฺทฺริยํ. อิมานิ โข ภิกฺขเว ปญฺจิ\-นฺทฺริยานีติ.  ปฐมํ.}

\csromanheader{2}{2. ปฐม\-สํขิตฺต\-สุตฺต}
\csromanmatikaanchor{mtk.107}
\csromanmatikaanchor{mtk.108}
\csromantocmarkref{subparagraph}{2-4. สํขิตฺตสุตฺตานิ}{mtk.107}
\bookmark[dest={mtk.107},level=5]{2-4. สํขิตฺตสุตฺตานิ}
\csromantocmarkref{subparagraph}{5-7. วิตฺถารสุตฺตานิ}{mtk.108}
\bookmark[dest={mtk.108},level=5]{5-7. วิตฺถารสุตฺตานิ}
\proseitem{482}{\csromanfolio{176}ปญฺจิมานิ ภิกฺขเว อินฺทฺริยานิ. กตมานิ ปญฺจ, สทฺธินฺทฺริยํ ฯเปฯ ปญฺญินฺทฺริยํ. อิมานิ โข ภิกฺขเว ปญฺจินฺทฺริยานิ. อิเมสํ โข ภิกฺขเว ปญฺจนฺนํ อินฺทฺริยานํ สมตฺตา ปริปูรตฺตา อรหํ โหติ, ตโต มุทุตเรหิ อนาคามี โหติ, ตโต มุทุตเรหิ สกทาคามี โหติ, ตโต มุทุตเรหิ โสตาปนฺโน โหติ, ตโต มุทุตเรหิ ธมฺมานุสารี โหติ, ตโต มุทุตเรหิ สทฺธานุสารี โหตีติ.  ทุติยํ.}

\csromanheader{2}{3. ทุติย\-สํขิตฺต\-สุตฺต}
\proseitem{483}{ปญฺจิมานิ ภิกฺขเว อินฺทฺริยานิ. กตมานิ ปญฺจ, สทฺธินฺทฺริยํ ฯเปฯ ปญฺญินฺทฺริยํ. อิมานิ โข ภิกฺขเว ปญฺจินฺทฺริยานิ. อิเมสํ โข ภิกฺขเว ปญฺจนฺนํ อินฺทฺริยานํ สมตฺตา ปริปูรตฺตา อรหํ โหติ, ตโต มุทุตเรหิ อนาคามี โหติ, ตโต มุทุตเรหิ สกทาคามี โหติ, ตโต มุทุตเรหิ โสตาปนฺโน โหติ, ตโต มุทุตเรหิ ธมฺมานุสารี โหติ. ตโต มุทุตเรหิ สทฺธานุสารี โหติ. อิติ โข ภิกฺขเว อินฺทฺริย\-เวมตฺตตา ผลเวมตฺตตา โหติ, ผลเวมตฺตตา ปุคฺคล\-เวมตฺตตาติ.  ตติยํ.}

\csromanheader{2}{4. ตติย\-สํขิตฺต\-สุตฺต}
\proseitem{484}{ปญฺจิมานิ ภิกฺขเว อินฺทฺริยานิ. กตมานิ ปญฺจ, สทฺธินฺทฺริยํ ฯเปฯ ปญฺญินฺทฺริยํ. อิมานิ โข ภิกฺขเว ปญฺจินฺทฺริยานิ. อิเมสํ โข ภิกฺขเว ปญฺจนฺนํ อินฺทฺริยานํ สมตฺตา ปริปูรตฺตา อรหํ โหติ, ตโต มุทุตเรหิ อนาคามี โหติ, ตโต มุทุตเรหิ สกทาคามี โหติ, ตโต มุทุตเรหิ โสตาปนฺโน โหติ, ตโต มุทุตเรหิ ธมฺมานุสารี โหติ, ตโต มุทุตเรหิ สทฺธานุสารี โหติ. อิติ โข ภิกฺขเว ปริปูรํ ปริปูรการี อาราเธติ, ปเทสํ ปเทสการี อาราเธติ, อวญฺฌานิเตฺววาหํ ภิกฺขเว ปญฺจนฺทฺริยานีติ วทามีติ.  จตุตฺถํ.}

\csromanheader{2}{5. ปฐม\-วิตฺถาร\-สุตฺต}
\proseitem{485}{ปญฺจิมานิ ภิกฺขเว อินฺทฺริยานิ. กตมานิ ปญฺจ, สทฺธินฺทฺริยํ ฯเปฯ ปญฺญินฺทฺริยํ. อิมานิ โข ภิกฺขเว ปญฺจินฺทฺริยานิ. อิเมสํ โข ภิกฺขเว ปญฺจนฺนํ อินฺทฺริยานํ}

\prose{\csromanfolio{177}สมตฺตา ปริปูรตฺตา อรหํ โหติ, ตโต มุทุตเรหิ อนฺตรา\-ปรินิพฺพายี โหติ, ตโต มุทุตเรหิ อุปหจฺจ\-ปรินิพฺพายี โหติ, ตโต มุทุตเรหิ อสงฺขารปรินิพฺพายี โหติ, ตโต มุทุตเรหิ สสงฺขาร\-ปรินิพฺพายี โหติ, ตโต มุทุตเรหิ อุทฺธํโสโต โหติ อกนิฏฺฐคามี, ตโต มุทุตเรหิ สกทาคามี โหติ, ตโต มุทุตเรหิ โสตาปนฺโน โหติ, ตโต มุทุตเรหิ ธมฺมานุสารี โหติ, ตโต มุทุตเรหิ สทฺธานุสารี โหตีติ.  ปญฺจมํ.}

\csromanheader{2}{6. ทุติย\-วิตฺถาร\-สุตฺต}
\proseitem{486}{ปญฺจิมานิ ภิกฺขเว อินฺทฺริยานิ. กตมานิ ปญฺจ, สทฺธินฺทฺริยํ ฯเปฯ ปญฺญินฺทฺริยํ. อิมานิ โข ภิกฺขเว ปญฺจินฺทฺริยานิ. อิเมสํ โข ภิกฺขเว ปญฺจนฺนํ อินฺทฺริยานํ สมตฺตา ปริปูรตฺตา อรหํ โหติ, ตโต มุทุตเรหิ อนฺตรา\-ปรินิพฺพายี โหติ, ตโต มุทุตเรหิ อุปหจฺจ\-ปรินิพฺพายี โหติ, ตโต มุทุตเรหิ อสงฺขารปรินิพฺพายี โหติ, ตโต มุทุตเรหิ สสงฺขาร\-ปรินิพฺพายี โหติ, ตโต มุทุตเรหิ อุทฺธํโสโต โหติ อกนิฏฺฐคามี, ตโต มุทุตเรหิ สกทาคามี โหติ, ตโต มุทุตเรหิ โสตาปนฺโน โหติ, ตโต มุทุตเรหิ ธมฺมานุสารี โหติ, ตโต มุทุตเรหิ สทฺธานุสารี โหติ. อิติ โข ภิกฺขเว อินฺทฺริย\-เวมตฺตตา ผลเวมตฺตตา โหติ, ผลเวมตฺตตา ปุคฺคล\-เวมตฺตตา โหตีติ.  ฉฏฺฐํ.}

\csromanheader{2}{7. ตติย\-วิตฺถาร\-สุตฺต}
\csromanmatikaanchor{mtk.109}
\csromantocmarkref{subparagraph}{8-10. ปฏิปนฺนสุตฺตาทีนิ}{mtk.109}
\bookmark[dest={mtk.109},level=5]{8-10. ปฏิปนฺนสุตฺตาทีนิ}
\proseitem{487}{ปญฺจิมานิ ภิกฺขเว อินฺทฺริยานิ. กตมานิ ปญฺจ, สทฺธินฺทฺริยํ ฯเปฯ ปญฺญินฺทฺริยํ. อิมานิ โข ภิกฺขเว ปญฺจินฺทฺริยานิ. อิเมสํ โข ภิกฺขเว ปญฺจนฺนํ อินฺทฺริยานํ สมตฺตา ปริปูรตฺตา อรหํ โหติ, ตโต มุทุตเรหิ อนฺตรา\-ปรินิพฺพายี โหติ, ตโต มุทุตเรหิ อุปหจฺจ\-ปรินิพฺพายี โหติ, ตโต มุทุตเรหิ อสงฺขารปรินิพฺพายี โหติ, ตโต มุทุตเรหิ สสงฺขาร\-ปรินิพฺพายี โหติ, ตโต มุทุตเรหิ อุทฺธํโสโต โหติ อกนิฏฺฐคามี, ตโต มุทุตเรหิ สกทาคามี โหติ, ตโต มุทุตเรหิ โสตาปนฺโน โหติ, ตโต มุทุตเรหิ ธมฺมานุสารี โหติ, ตโต มุทุตเรหิ สทฺธานุสารี โหติ. อิติ โข ภิกฺขเว \csromanfolio{178}ปริปูรํ ปริปูรการี อาราเธติ, ปเทสํ ปเทสการี อาราเธติ, อวญฺฌานิ เตฺววาหํ ภิกฺขเว ปญฺจิ\-นฺทฺริยานีติ วทามีติ.  สตฺตมํ.}

\csromanheader{2}{8. \textbf{ปฏิปนฺนสุตฺต}}
\proseitem{488}{ปญฺจิมานิ ภิกฺขเว อินฺทฺริยานิ. กตมานิ ปญฺจ, สทฺธินฺทฺริยํ ฯเปฯ ปญฺญินฺทฺริยํ. อิมานิ โข ภิกฺขเว ปญฺจินฺทฺริยานิ. อิเมสํ โข ภิกฺขเว ปญฺจนฺนํ อินฺทฺริยานํ สมตฺตา ปริปูรตฺตา อรหํ โหติ, ตโต มุทุตเรหิ อรหตฺตผลสจฺฉิ กิริยาย ปฏิปนฺโน โหติ, ตโต มุทุตเรหิ อนาคามี โหติ, ตโต มุทุตเรหิ อนาคามิ\-ผล\-สจฺฉิกิริยาย ปฏิปนฺโน โหติ, ตโต มุทุตเรหิ สกทาคามี โหติ, ตโต มุทุตเรหิ สกทาคามิ\-ผล\-สจฺฉิกิริยาย ปฏิปนฺโน โหติ, ตโต มุทุตเรหิ โสตาปนฺโน โหติ, ตโต มุทุตเรหิ โสตาปตฺติผล\-สจฺฉิกิริยาย ปฏิปนฺโน โหติ. ยสฺส โข ภิกฺขเว อิมานิ ปญฺจินฺทฺริยานิ สพฺเพน สพฺพํ สพฺพถา สพฺพํ นตฺถิ, ตมหํ “พาหิโร ปุถุชฺชน\-ปกฺเข ฐิโต”ติ วทามีติ.  อฏฺฐมํ.}

\csromanheader{2}{9. \textbf{สมฺปนฺนสุตฺต}}
\proseitem{489}{อถ โข อญฺญตโร ภิกฺขุ เยน ภควา เตนุปสงฺกมิ, อุปสงฺกมิตฺวา ภควนฺตํ อภิวาเทตฺวา เอกมนฺตํ นิสีทิ, เอกมนฺตํ นิสินฺโน โข โส ภิกฺขุ ภควนฺตํ เอตทโวจ\csromandash{}}

\prose{“อินฺทฺริย\-สมฺปนฺโน อินฺทฺริย\-สมฺปนฺโนติ ภนฺเต วุจฺจติ, กิตฺตาวตา นุ โข ภนฺเต อินฺทฺริย\-สมฺปนฺโน โหตี”ติ. อิธ ภิกฺขุ ภิกฺขุ สทฺธินฺทฺริยํ ภาเวติ อุปสมคามิํ สมฺโพธคามิํ, วีริยินฺทฺริยํ ภาเวติ อุปสมคามิํ สมฺโพธคามิํ, สตินฺทฺริยํ ภาเวติ อุปสมคามิํ สมฺโพธคามิํ, สมาธินฺทฺริยํ ภาเวติ อุปสมคามิํ สมฺโพธคามิํ, ปญฺญินฺทฺริยํ ภาเวติ อุปสมคามิํ สมฺโพธคามิํ. เอตฺตาวตา โข ภิกฺขุ ภิกฺขุ อินฺทฺริย\-สมฺปนฺโน โหตีติ.  นวมํ.}

\csromanheader{2}{10. อาสวกฺขย\-สุตฺต}
\proseitem{490}{ปญฺจิมานิ ภิกฺขเว อินฺทฺริยานิ. กตมานิ ปญฺจ, สทฺธินฺทฺริยํ ฯเปฯ ปญฺญินฺทฺริยํ. อิมานิ โข ภิกฺขเว ปญฺจินฺทฺริยานิ. อิเมสํ โข ภิกฺขเว ปญฺจนฺนํ อินฺทฺริยานํ}

\csromanmatikaanchor{mtk.110}
\csromanmatikaanchor{mtk.111}
\csromantocmarknopage{subparagraph}{3. ฉฬินฺทฺริยวคฺค}{mtk.110}
\bookmark[dest={mtk.110},level=5]{3. ฉฬินฺทฺริยวคฺค}
\csromantocmarkref{subparagraph}{1-10. ปุนพฺภวสุตฺตาทีนิ}{mtk.111}
\bookmark[dest={mtk.111},level=5]{1-10. ปุนพฺภวสุตฺตาทีนิ}
\prosewithclosermidruled{\csromanfolio{179}ภาวิตตฺตา พหุลีกตตฺตา ภิกฺขุ อาสวานํ ขยา อนาสวํ เจโตวิมุตฺติํ ปญฺญาวิมุตฺติํ ทิฏฺเฐว ธมฺเม สยํ อภิญฺญา สจฺฉิกตฺวา อุปสมฺปชฺช วิหรตีติ.  ทสมํ.}{มุทุตรวคฺโค ทุติโย.}
\csromancenter{\textbf{ตสฺสุทฺทานํ}}
\setlength{\csromangathapagewidth}{0pt}
\setlength{\csromangathawakwidth}{0pt}
\csromangathameasuregroup{ปฏิลาโภ ตโย สํขิตฺตา, \\ ปฏิปนฺโน จ สมฺปนฺโน\textsuperscript{1},}{\csromangathabat{ปฏิลาโภ ตโย สํขิตฺตา,}{วิตฺถารา อปเร ตโย.} \\ \csromangathabat{ปฏิปนฺโน จ สมฺปนฺโน\textsuperscript{1},}{ทสมํ อาสวกฺขยนฺติ.}}
\csromangathasetleft{ปฏิลาโภ ตโย สํขิตฺตา, \\ ปฏิปนฺโน จ สมฺปนฺโน\textsuperscript{1},}
\csromangathagroup{\csromangathastanza{\csromangathabat{ปฏิลาโภ ตโย สํขิตฺตา,}{วิตฺถารา อปเร ตโย.} \\ \csromangathabat{ปฏิปนฺโน จ สมฺปนฺโน\csromansharedfootnote{338dd9d64fd7ee1e}{[มิสฺสินฺคฺ โนเต 0]},}{ทสมํ อาสวกฺขยนฺติ.}}}

\csromanheader{6}{3. ฉฬินฺทฺริย\-วคฺค}
\csromanheader{2}{1. ปุนพฺภว\-สุตฺต}
\proseitem{491}{ปญฺจิมานิ ภิกฺขเว อินฺทฺริยานิ. กตมานิ ปญฺจ, สทฺธินฺทฺริยํ ฯเปฯ ปญฺญินฺทฺริยํ. ยาวกีวญฺจาหํ ภิกฺขเว อิเมสํ ปญฺจนฺนํ อินฺทฺริยานํ สมุทยญฺจ อตฺถงฺคมญฺจ อสฺสาทญฺจ อาทีนวญฺจ นิสฺสรณญฺจ ยถาภูตํ นาพฺภญฺญาสิํ, เนว ตาวาหํ ภิกฺขเว สเทวเก โลเก สมารเก สพฺรหฺมเก สสฺสมณ\-พฺราหฺมณิยา ปชาย สเทว\-มนุสฺสาย อนุตฺตรํ สมฺมาสมฺโพธิํ อภิสมฺพุทฺโธติ ปจฺจญฺญาสิํ\csromansharedfootnote{27ebb5cb172b84b2}{อภิสมฺพุทฺโธ ปจฺจญฺญาสิํ (สี, สฺยา, กํ)}. ยโต จ ขฺวาหํ ภิกฺขเว อิเมสํ ปญฺจนฺนํ อินฺทฺริยานํ สมุทยญฺจ อตฺถงฺคมญฺจ อสฺสาทญฺจ อาทีนวญฺจ นิสฺสรณญฺจ ยถาภูตํ อพฺภญฺญาสิํ, อถาหํ ภิกฺขเว สเทวเก โลเก สมารเก สพฺรหฺมเก สสฺสมณ\-พฺราหฺมณิยา ปชาย สเทว\-มนุสฺสาย อนุตฺตรํ สมฺมาสมฺโพธิํ อภิสมฺพุทฺโธติ ปจฺจญฺญาสิํ, ญาณญฺจ ปน เม ทสฺสนํ อุทปาทิ, อกุปฺปา เม วิมุตฺติ\csromansharedfootnote{4a253074c6d5b4b9}{เจโตวิมุตฺติ (สี, อิ, ก)}, อยมนฺติมา ชาติ, นตฺถิทานิ ปุนพฺภโวติ.  ปฐมํ.}

\csromanheader{2}{2. ชีวิตินฺทฺริย\-สุตฺต}
\proseitem{492}{ตีณิมานิ ภิกฺขเว อินฺทฺริยานิ. กตมานิ ตีณิ, อิตฺถินฺทฺริยํ ปุริสินฺทฺริยํ ชีวิตินฺทฺริยํ. อิมานิ โข ภิกฺขเว ตีณิ อินฺทฺริยานีติ.  ทุติยํ.}

\csromanheader{2}{3. อญฺญินฺทฺริย\-สุตฺต}
\proseitem{493}{\csromanfolio{180}ตีณิมานิ ภิกฺขเว อินฺทฺริยานิ. กตมานิ ตีณิ, อนญฺญาต\-ญฺญสฺสามี\-ตินฺทฺริยํ อญฺญินฺทฺริยํ อญฺญาตาวิ\-นฺทฺริยํ. อิมานิ โข ภิกฺขเว ตีณิ อินฺทฺริยานีติ.  ตติยํ.}

\csromanheader{2}{4. \textbf{เอกพีชีสุตฺต}}
\proseitem{494}{ปญฺจิมานิ ภิกฺขเว อินฺทฺริยานิ. กตมานิ ปญฺจ, สทฺธินฺทฺริยํ ฯเปฯ ปญฺญินฺทฺริยํ. อิมานิ โข ภิกฺขเว ปญฺจินฺทฺริยานิ. อิเมสํ โข ภิกฺขเว ปญฺจนฺนํ อินฺทฺริยานํ สมตฺตา ปริปูรตฺตา อรหํ โหติ, ตโต มุทุตเรหิ อนฺตรา\-ปรินิพฺพายี โหติ, ตโต มุทุตเรหิ อุปหจฺจ\-ปรินิพฺพายี โหติ, ตโต มุทุตเรหิ อสงฺขารปรินิพฺพายี โหติ, ตโต มุทุตเรหิ สสงฺขาร\-ปรินิพฺพายี โหติ, ตโต มุทุตเรหิ อุทฺธํโสโต โหติ อกนิฏฺฐคามี, ตโต มุทุตเรหิ สกทาคามี โหติ, ตโต มุทุตเรหิ เอกพีชี\csromansharedfootnote{65bf893b428c1c56}{เอกพีชิ (ก)} โหติ, ตโต มุทุตเรหิ โกลํโกโล โหติ, ตโต มุทุตเรหิ สตฺตกฺขตฺตุ\-ปรโม โหติ, ตโต มุทุตเรหิ ธมฺมานุสารี โหติ, ตโต มุทุตเรหิ สทฺธานุสารี โหตีติ.  จตุตฺถํ.}

\csromanheader{2}{5. \textbf{สุทฺธกสุตฺต}}
\proseitem{495}{ฉยิมานิ ภิกฺขเว อินฺทฺริยานิ. กตมานิ ฉ, จกฺขุนฺทฺริยํ โสตินฺทฺริยํ ฆานินฺทฺริยํ ชิวฺหินฺทฺริยํ กายินฺทฺริยํ มนินฺทฺริยํ. อิมานิ โข ภิกฺขเว ฉ อินฺทฺริยานีติ.  ปญฺจมํ.}

\csromanheader{2}{6. \textbf{โสตาปนฺนสุตฺต}}
\proseitem{496}{ฉยิมานิ ภิกฺขเว อินฺทฺริยานิ. กตมานิ ฉ, จกฺขุนฺทฺริยํ ฯเปฯ มนินฺทฺริยํ. ยโต โข ภิกฺขเว อริยสาวโก อิเมสํ ฉนฺนํ อินฺทฺริยานํ สมุทยญฺจ อตฺถงฺคมญฺจ อสฺสาทญฺจ อาทีนวญฺจ นิสฺสรณญฺจ ยถาภูตํ ปชานาติ. อยํ วุจฺจติ ภิกฺขเว อริยสาวโก โสตาปนฺโน อวินิปาต\-ธมฺโม นิยโต สมฺโพธิ\-ปรายโณติ.  ฉฏฺฐํ.}

\csromanheader{2}{4. อินฺทฺริย\-สํยุตฺต}
\csromanheader{2}{7. \textbf{อรหนฺตสุตฺต}}
\proseitem{497}{\csromanfolio{181}ฉยิมานิ ภิกฺขเว อินฺทฺริยานิ. กตมานิ ฉ, จกฺขุนฺทฺริยํ โสตินฺทฺริยํ ฆานินฺทฺริยํ ชิวฺหินฺทฺริยํ กายินฺทฺริยํ มนินฺทฺริยํ. ยโต โข ภิกฺขเว ภิกฺขุ อิเมสํ ฉนฺนํ อินฺทฺริยานํ สมุทยญฺจ อตฺถงฺคมญฺจ อสฺสาทญฺจ อาทีนวญฺจ นิสฺสรณญฺจ ยถาภูตํ วิทิตฺวา อนุปาทาวิมุตฺโต โหติ. อยํ วุจฺจติ ภิกฺขเว ภิกฺขุ อรหํ ขีณาสโว วุสิตวา กตกรณีโย โอหิตภาโร อนุปฺปตฺต\-สทตฺโถ ปริกฺขีณ\-ภวสํโยชโน สมฺมทญฺญา\-วิมุตฺโตติ.  สตฺตมํ.}

\csromanheader{2}{8. \textbf{สมฺพุทฺธสุตฺต}}
\proseitem{498}{ฉยิมานิ ภิกฺขเว อินฺทฺริยานิ. กตมานิ ฉ, จกฺขุนฺทฺริยํ โสตินฺทฺริยํ ฆานินฺทฺริยํ ชิวฺหินฺทฺริยํ กายินฺทฺริยํ มนินฺทฺริยํ. ยาวกีวญฺจาหํ ภิกฺขเว อิเมสํ ฉนฺนํ อินฺทฺริยานํ สมุทยญฺจ อตฺถงฺคมญฺจ อสฺสาทญฺจ อาทีนวญฺจ นิสฺสรณญฺจ ยถาภูตํ นาพฺภญฺญาสิํ, เนว ตาวาหํ ภิกฺขเว สเทวเก โลเก สมารเก สพฺรหฺมเก สสฺสมณ\-พฺราหฺมณิยา ปชาย สเทว\-มนุสฺสาย อนุตฺตรํ สมฺมาสมฺโพธิํ อภิสมฺพุทฺโธติ ปจฺจญฺญาสิํ. ยโต จ ขฺวาหํ ภิกฺขเว อิเมสํ ฉนฺนํ อินฺทฺริยานํ สมุทยญฺจ อตฺถงฺคมญฺจ อสฺสาทญฺจ อาทีนวญฺจ นิสฺสรณญฺจ ยถาภูตํ อพฺภญฺญาสิํ, อถาหํ ภิกฺขเว สเทวเก โลเก สมารเก สพฺรหฺมเก สสฺสมณ\-พฺราหฺมณิยา ปชาย สเทว\-มนุสฺสาย อนุตฺตรํ สมฺมาสมฺโพธิํ อภิสมฺพุทฺโธติ ปจฺจญฺญาสิํ, ญาณญฺจ ปน เม ทสฺสนํ อุทปาทิ, อกุปฺปา เม วิมุตฺติ, อยมนฺติมา ชาติ, นตฺถิทานิ ปุนพฺภโวติ.  อฏฺฐมํ.}

\csromanheader{2}{9. ปฐม\-สมณพฺราหฺมณ\-สุตฺต}
\proseitem{499}{ฉยิมานิ ภิกฺขเว อินฺทฺริยานิ. กตมานิ ฉ, จกฺขุนฺทฺริยํ โสตินฺทฺริยํ ฆานินฺทฺริยํ ชิวฺหินฺทฺริยํ กายินฺทฺริยํ มนินฺทฺริยํ. เย หิ เกจิ ภิกฺขเว สมณา วา พฺราหฺมณา วา อิเมสํ ฉนฺนํ อินฺทฺริยานํ สมุทยญฺจ อตฺถงฺคมญฺจ อสฺสาทญฺจ อาทีนวญฺจ นิสฺสรณญฺจ ยถาภูตํ นปฺปชานนฺติ, น เม เต ภิกฺขเว สมณา วา พฺราหฺมณา วา สมเณสุ วา สมณสมฺมตา, พฺราหฺมเณสุ วา พฺราหฺมณ\-สมฺมตา, น จ ปเนเต อายสฺมนฺโต สามญฺญตฺถํ วา พฺรหฺมญฺญตฺถํ \csromanfolio{182}วา ทิฏฺเฐว ธมฺเม สยํ อภิญฺญา สจฺฉิกตฺวา อุปสมฺปชฺช วิหรนฺติ. เย จ โข เกจิ ภิกฺขเว สมณา วา พฺราหฺมณา วา อิเมสํ ฉนฺนํ อินฺทฺริยานํ สมุทยญฺจ อตฺถงฺคมญฺจ อสฺสาทญฺจ อาทีนวญฺจ นิสฺสรณญฺจ ยถาภูตํ ปชานนฺติ, เต โข เม ภิกฺขเว สมณา วา พฺราหฺมณา วา สมเณสุ เจว สมณสมฺมตา, พฺราหฺมเณสุ จ พฺราหฺมณ\-สมฺมตา, เต จ ปนายสฺมนฺโต สามญฺญตฺถญฺจ พฺรหฺมญฺญตฺถญฺจ ทิฏฺเฐว ธมฺเม สยํ อภิญฺญา สจฺฉิกตฺวา อุปสมฺปชฺช วิหรนฺตีติ.  นวมํ.}

\csromanheader{2}{10. ทุติย\-สมณพฺราหฺมณ\-สุตฺต}
\proseitem{500}{เย หิ เกจิ ภิกฺขเว สมณา วา พฺราหฺมณา วา จกฺขุนฺทฺริยํ นปฺปชานนฺติ, จกฺขุนฺทฺริย\-สมุทยํ นปฺปชานนฺติ, จกฺขุนฺทฺริย\-นิโรธํ นปฺปชานนฺติ, จกฺขุนฺทฺริยนิโรธ\-คามินิํ ปฏิปทํ นปฺปชานนฺติ. โสตินฺทฺริยํ ฯเปฯ. ฆานินฺทฺริยํ ฯเปฯ. ชิวฺหานฺทฺริยํ ฯเปฯ. กายินฺทฺริยํ ฯเปฯ. มนินฺทฺริยํ นปฺปชานนฺติ, มนินฺทฺริย\-สมุทยํ นปฺปชานนฺติ, มนินฺทฺริย\-นิโรธํ นปฺปชานนฺติ. มนินฺทฺริยนิโรธ\-คามินิํ ปฏิปทํ นปฺปชานนฺติ. น เม เต ภิกฺขเว ฯเปฯ สยํ อภิญฺญา สจฺฉิกตฺวา อุปสมฺปชฺช วิหรนฺติ.}

\prosewithclosermidruled{เย จ โข เกจิ ภิกฺขเว สมณา วา พฺราหฺมณา วา จกฺขุนฺทฺริยํ ปชานนฺติ, จกฺขุนฺทฺริย\-สมุทยํ ปชานนฺติ, จกฺขุนฺทฺริย\-นิโรธํ ปชานนฺติ, จกฺขุนฺทฺริยนิโรธ\-คามินิํ ปฏิปทํ ปชานนฺติ. โสตินฺทฺริยํ ฯเปฯ. ฆานินฺทฺริยํ ฯเปฯ. ชิวฺหินฺทฺริยํ ฯเปฯ. กายินฺทฺริยํ ฯเปฯ. มนินฺทฺริยํ ปชานนฺติ, มนินฺทฺริย\-สมุทยํ ปชานนฺติ, มนินฺทฺริย\-นิโรธํ ปชานนฺติ, มนินฺทฺริยนิโรธ\-คามินิํ ปฏิปทํ ปชานนฺติ, เต โข เม ภิกฺขเว สมณา วา พฺราหฺมณา วา สมเณสุ เจว สมณสมฺมตา, พฺราหฺมเณสุ จ พฺราหฺมณ\-สมฺมตา, เต จ ปนายสฺมนฺโต สามญฺญตฺถญฺจ พฺรหฺมญฺญตฺถญฺจ ทิฏฺเฐว ธมฺเม สยํ อภิญฺญา สจฺฉิกตฺวา อุปสมฺปชฺช วิหรนฺตีติ.  ทสมํ.}{ฉฬินฺทิยวคฺโค ตติโย.}
\csromancenter{\textbf{ตสฺสุทฺทานํ}}
\setlength{\csromangathapagewidth}{0pt}
\setlength{\csromangathawakwidth}{0pt}
\csromangathameasuregroup{ปุนพฺภโว ชีวิตญฺญาย, \\ โสโต อรหสมฺพุทฺโธ,}{\csromangathabat{ปุนพฺภโว ชีวิตญฺญาย,}{เอกพีชี จ สุทฺธกํ.} \\ \csromangathabat{โสโต อรหสมฺพุทฺโธ,}{เทฺว จ สมณพฺราหฺมณาติ.}}
\csromangathasetleft{ปุนพฺภโว ชีวิตญฺญาย, \\ โสโต อรหสมฺพุทฺโธ,}
\csromangathagroup{\csromangathastanza{\csromangathabat{ปุนพฺภโว ชีวิตญฺญาย,}{เอกพีชี จ สุทฺธกํ.} \\ \csromangathabat{โสโต อรหสมฺพุทฺโธ,}{เทฺว จ สมณ\-พฺราหฺมณาติ.}}}

\csromanmatikaanchor{mtk.112}
\csromantocmarknopage{subparagraph}{4. สุขินฺทฺริยวคฺค}{mtk.112}
\bookmark[dest={mtk.112},level=5]{4. สุขินฺทฺริยวคฺค}
\csromanheader{6}{4. สุขินฺทฺริย\-วคฺค}
\csromanheader{2}{1. \textbf{สุทฺธิกสุตฺต}}
\csromanmatikaanchor{mtk.113}
\csromanmatikaanchor{mtk.114}
\csromantocmarkref{subparagraph}{1-3. สุทฺธิกสุตฺตาทีนิ}{mtk.113}
\bookmark[dest={mtk.113},level=5]{1-3. สุทฺธิกสุตฺตาทีนิ}
\csromantocmarkref{subparagraph}{4-8. ปฐมสมณพฺราหฺมณสุตฺตาทีนิ}{mtk.114}
\bookmark[dest={mtk.114},level=5]{4-8. ปฐมสมณพฺราหฺมณสุตฺตาทีนิ}
\proseitem{501}{\csromanfolio{183}ปญฺจิมานิ ภิกฺขเว อินฺทฺริยานิ. กตมานิ ปญฺจ, สุขินฺทฺริยํ ทุกฺขินฺทฺริยํ โสมนสฺสิ\-นฺทฺริยํ โทมนสฺสิ\-นฺทฺริยํ อุเปกฺขินฺทฺริยํ. อิมานิ โข ภิกฺขเว ปญฺจิ\-นฺทฺริยานีติ.  ปฐมํ.}

\csromanheader{2}{2. \textbf{โสตาปนฺนสุตฺต}}
\proseitem{502}{ปญฺจิมานิ ภิกฺขเว อินฺทฺริยานิ. กตมานิ ปญฺจ, สุขินฺทฺริยํ ทุกฺขินฺทฺริยํ โสมนสฺสิ\-นฺทฺริยํ โทมนสฺสิ\-นฺทฺริยํ อุเปกฺขินฺทฺริยํ. ยโต โข ภิกฺขเว อริยสาวโก อิเมสํ ปญฺจนฺนํ อินฺทฺริยานํ สมุทยญฺจ อตฺถงฺคมญฺจ อสฺสาทญฺจ อาทีนวญฺจ นิสฺสรณญฺจ ยถาภูตํ ปชานาติ. อยํ วุจฺจติ ภิกฺขเว อริยสาวโก โสตาปนฺโน อวินิปาต\-ธมฺโม นิยโต สมฺโพธิ\-ปรายโณติ.  ทุติยํ.}

\csromanheader{2}{3. \textbf{อรหนฺตสุตฺต}}
\proseitem{503}{ปญฺจิมานิ ภิกฺขเว อินฺทฺริยานิ. กตมานิ ปญฺจ, สุขินฺทฺริยํ ทุกฺขินฺทฺริยํ โสมนสฺสิ\-นฺทฺริยํ โทมนสฺสิ\-นฺทฺริยํ อุเปกฺขินฺทฺริยํ. ยโต โข ภิกฺขเว ภิกฺขุ อิเมสํ ปญฺจนฺนํ อินฺทฺริยานํ สมุทยญฺจ อตฺถงฺคมญฺจ อสฺสาทญฺจ อาทีนวญฺจ นิสฺสรณญฺจ ยถาภูตํ วิทิตฺวา อนุปาทาวิมุตฺโต โหติ. อยํ วุจฺจติ ภิกฺขเว ภิกฺขุ อรหํ ขีณาสโว วุสิตวา กตกรณีโย โอหิตภาโร อนุปฺปตฺต\-สทตฺโถ ปริกฺขีณ\-ภวสํโยชโน สมฺมทญฺญา\-วิมุตฺโตติ.  ตติยํ.}

\csromanheader{2}{4. ปฐม\-สมณพฺราหฺมณ\-สุตฺต}
\proseitem{504}{ปญฺจิมานิ ภิกฺขเว อินฺทฺริยานิ. กตมานิ ปญฺจ, สุขินฺทฺริยํ ทุกฺขินฺทฺริยํ โสมนสฺสิ\-นฺทฺริยํ โทมนสฺสิ\-นฺทฺริยํ อุเปกฺขินฺทฺริยํ. เย หิ เกจิ ภิกฺขเว สมณา วา พฺราหฺมณา วา อิเมสํ ปญฺจนฺนํ อินฺทฺริยานํ สมุทยญฺจ อตฺถงฺคมญฺจ อสฺสาทญฺจ อาทีนวญฺจ นิสฺสรณญฺจ ยถาภูตํ นปฺปชานนฺติ, น เม เต ภิกฺขเว สมณา วา พฺราหฺมณา วา สมเณสุ วา สมณสมฺมตา, พฺราหฺมเณสุ วา พฺราหฺมณ\-สมฺมตา, น จ ปเนเต อายสฺมนฺโต สามญฺญตฺถํ วา พฺรหฺมญฺญตฺถํ วา ทิฏฺเฐว ธมฺเม สยํ อภิญฺญา สจฺฉิกตฺวา อุปสมฺปชฺช วิหรนฺติ.}

\prose{\csromanfolio{184}เย จ โข เกจิ ภิกฺขเว สมณา วา พฺราหฺมณา วา อิเมสํ ปญฺจนฺนํ อินฺทฺริยานํ สมุทยญฺจ อตฺถงฺคมญฺจ อสฺสาทญฺจ อาทีนวญฺจ นิสฺสรณญฺจ ยถาภูตํ ปชานนฺติ, เต โข เม ภิกฺขเว สมณา วา พฺราหฺมณา วา สมเณสุ เจว สมณสมฺมตา, พฺราหฺมเณสุ จ พฺราหฺมณ\-สมฺมตา, เต จ ปนายสฺมนฺโต สามญฺญตฺถญฺจ พฺรหฺมญฺญตฺถญฺจ ทิฏฺเฐว ธมฺเม สยํ อภิญฺญา สจฺฉิกตฺวา อุปสมฺปชฺช วิหรนฺตีติ.  จตุตฺถํ.}

\csromanheader{2}{5. ทุติย\-สมณพฺราหฺมณ\-สุตฺต}
\proseitem{505}{ปญฺจิมานิ ภิกฺขเว อินฺทฺริยานิ. กตมานิ ปญฺจ, สุขินฺทฺริยํ ทุกฺขินฺทฺริยํ โสมนสฺสิ\-นฺทฺริยํ โทมนสฺสิ\-นฺทฺริยํ อุเปกฺขินฺทฺริยํ. เย หิ เกจิ ภิกฺขเว สมณา วา พฺราหฺมณา วา สุขินฺทฺริยํ นปฺปชานนฺติ, สุขินฺทฺริย\-สมุทยํ นปฺปชานนฺติ, สุขินฺทฺริย\-นิโรธํ นปฺปชานนฺติ, สุขินฺทฺริยนิโรธ\-คามินิํ ปฏิปทํ นปฺปชานนฺติ. ทุกฺขินฺทฺริยํ นปฺปชานนฺติ ฯเปฯ. โสมนสฺสิ\-นฺทฺริยํ นปฺปชานนฺติ ฯเปฯ. โทมนสฺสิ\-นฺทฺริยํ นปฺปชานนฺติ ฯเปฯ. อุเปกฺขินฺทฺริยํ นปฺปชานนฺติ, อุเปกฺขินฺทฺริย\-สมุทยํ นปฺปชานนฺติ, อุเปกฺขินฺทฺริย\-นิโรธํ นปฺปชานนฺติ, อุเปกฺขินฺทฺริยนิโรธ\-คามินิํ ปฏิปทํ นปฺปชานนฺติ. น เม เต ภิกฺขเว สมณา วา พฺราหฺมณา วา สมเณสุ วา สมณสมฺมตา, พฺราหฺมเณสุ วา พฺราหฺมณ\-สมฺมตา, น จ ปเนเต อายสฺมนฺโต สามญฺญตฺถํ วา พฺราหฺมญฺญตฺถํ วา ทิฏฺเฐว ธมฺเม สยํ อภิญฺญา สจฺฉิกตฺวา อุปสมฺปชฺช วิหรนฺติ.}

\prose{เย จ โข เกจิ ภิกฺขเว สมณา วา พฺราหฺมณา วา สุขินฺทฺริยํ ปชานนฺติ, สุขินฺทฺริย\-สมุทยํ ปชานนฺติ, สุขินฺทฺริย\-นิโรธํ ปชานนฺติ, สุขินฺทฺริยนิโรธ\-คามินิํ ปฏิปทํ ปชานนฺติ. ทุกฺขินฺทฺริยํ ปชานนฺติ ฯเปฯ. โสมนสฺสิ\-นฺทฺริยํ ปชานนฺติ ฯเปฯ. โทมนสฺสิ\-นฺทฺริยํ ปชานนฺติ ฯเปฯ. อุเปกฺขินฺทฺริยํ ปชานนฺติ, อุเปกฺขินฺทฺริย\-สมุทยํ ปชานนฺติ, อุเปกฺขินฺทฺริย\-นิโรธํ ปชานนฺติ, อุเปกฺขินฺทฺริยนิโรธ\-คามินิํ ปฏิปทํ ปชานนฺติ. เต จ โข เม ภิกฺขเว สมณา วา พฺราหฺมณา วา สมเณสุ เจว สมณสมฺมตา, พฺราหฺมเณสุ จ พฺราหฺมณ\-สมฺมตา, เต จ ปนายสฺมนฺโต สามญฺญตฺถญฺจ พฺรหฺมญฺญตฺถญฺจ ทิฏฺเฐว ธมฺเม สยํ อภิญฺญา สจฺฉิกตฺวา อุปสมฺปชฺช วิหรนฺตีติ.  ปญฺจมํ.}

\csromanheader{2}{6. ปฐม\-วิภงฺค\-สุตฺต}
\proseitem{506}{ปญฺจิมานิ ภิกฺขเว อินฺทฺริยานิ. กตมานิ ปญฺจ, สุขินฺทฺริยํ ทุกฺขินฺทฺริยํ โสมนสฺสิ\-นฺทฺริยํ โทมนสฺสิ\-นฺทฺริยํ อุเปกฺขินฺทฺริยํ.}

\prose{\csromanfolio{185}กตมญฺจ ภิกฺขเว สุขินฺทฺริยํ, ยํ โข ภิกฺขเว กายิกํ สุขํ กายิกํ สาตํ กาย\-สมฺผสฺสชํ สุขํ สาตํ เวทยิตํ. อิทํ วุจฺจติ ภิกฺขเว สุขินฺทฺริยํ.}

\prose{กตมญฺจ ภิกฺขเว ทุกฺขินฺทฺริยํ, ยํ โข ภิกฺขเว กายิกํ ทุกฺขํ กายิกํ อสาตํ กาย\-สมฺผสฺสชํ ทุกฺขํ อสาตํ เวทยิตํ. อิทํ วุจฺจติ ภิกฺขเว ทุกฺขินฺทฺริยํ.}

\prose{กตมญฺจ ภิกฺขเว โสมนสฺสิ\-นฺทฺริยํ, ยํ โข ภิกฺขเว เจตสิกํ สุขํ เจตสิกํ สาตํ มโน\-สมฺผสฺสชํ สุขํ สาตํ เวทยิตํ. อิทํ วุจฺจติ ภิกฺขเว โสมนสฺสิ\-นฺทฺริยํ.}

\prose{กตมญฺจ ภิกฺขเว โทมนสฺสิ\-นฺทฺริยํ, ยํ โข ภิกฺขเว เจตสิกํ ทุกฺขํ เจตสิกํ อสาตํ มโน\-สมฺผสฺสชํ ทุกฺขํ อสาตํ เวทยิตํ. อิทํ วุจฺจติ ภิกฺขเว โทมนสฺสิ\-นฺทฺริยํ.}

\prose{กตมญฺจ ภิกฺขเว อุเปกฺขินฺทฺริยํ, ยํ โข ภิกฺขเว กายิกํ วา เจตสิกํ วา เนวสาตํ นาสาตํ เวทยิตํ. อิทํ วุจฺจติ ภิกฺขเว อุเปกฺขินฺทฺริยํ. อิมานิ โข ภิกฺขเว ปญฺจิ\-นฺทฺริยานีติ.  ฉฏฺฐํ.}

\csromanheader{2}{7. ทุติย\-วิภงฺค\-สุตฺต}
\proseitem{507}{ปญฺจิมานิ ภิกฺขเว อินฺทฺริยานิ. กตมานิ ปญฺจ, สุขินฺทฺริยํ ทุกฺขินฺทฺริยํ โสมนสฺสิ\-นฺทฺริยํ โทมนสฺสิ\-นฺทฺริยํ อุเปกฺขินฺทฺริยํ.}

\prose{กตมญฺจ ภิกฺขเว สุขินฺทฺริยํ, ยํ โข ภิกฺขเว กายิกํ สุขํ กายิกํ สาตํ กาย\-สมฺผสฺสชํ สุขํ สาตํ เวทยิตํ. อิทํ วุจฺจติ ภิกฺขเว สุขินฺทฺริยํ.}

\prose{กตมญฺจ ภิกฺขเว ทุกฺขินฺทฺริยํ, ยํ โข ภิกฺขเว กายิกํ ทุกฺขํ กายิกํ อสาตํ กาย\-สมฺผสฺสชํ ทุกฺขํ อสาตํ เวทยิตํ. อิทํ วุจฺจติ ภิกฺขเว ทุกฺขินฺทฺริยํ.}

\prose{กตมญฺจ ภิกฺขเว โสมนสฺสิ\-นฺทฺริยํ, ยํ โข ภิกฺขเว เจตสิกํ สุขํ เจตสิกํ สาตํ มโน\-สมฺผสฺสชํ สุขํ สาตํ เวทยิตํ. อิทํ วุจฺจติ ภิกฺขเว โสมนสฺสิ\-นฺทฺริยํ.}

\prose{\csromanfolio{186}กตมญฺจ ภิกฺขเว โทมนสฺสิ\-นฺทฺริยํ, ยํ โข ภิกฺขเว เจตสิกํ ทุกฺขํ เจตสิกํ อสาตํ มโน\-สมฺผสฺสชํ ทุกฺขํ อสาตํ เวทยิตํ. อิทํ วุจฺจติ ภิกฺขเว โทมนสฺสิ\-นฺทฺริยํ.}

\prose{กตมญฺจ ภิกฺขเว อุเปกฺขินฺทฺริยํ, ยํ โข ภิกฺขเว กายิกํ วา เจตสิกํ วา เนวสาตํ นาสาตํ เวทยิตํ. อิทํ วุจฺจติ ภิกฺขเว อุเปกฺขินฺทฺริยํ.}

\prose{ตตฺร ภิกฺขเว ยญฺจ สุขินฺทฺริยํ ยญฺจ โสมนสฺสิ\-นฺทฺริยํ, สุขา สา เวทนา ทฏฺฐพฺพา. ตตฺร ภิกฺขเว ยญฺจ ทุกฺขินฺทฺริยํ ยญฺจ โทมนสฺสิ\-นฺทฺริยํ, ทุกฺขา สา เวทนา ทฏฺฐพฺพา. ตตฺร ภิกฺขเว ยทิทํ อุเปกฺขินฺทฺริยํ, อทุกฺขมสุขา สา เวทนา ทฏฺฐพฺพา. อิมานิ โข ภิกฺขเว ปญฺจิ\-นฺทฺริยานีติ.  สตฺตมํ.}

\csromanheader{2}{8. ตติย\-วิภงฺค\-สุตฺต}
\proseitem{508}{ปญฺจิมานิ ภิกฺขเว อินฺทฺริยานิ. กตมานิ ปญฺจ, สุขินฺทฺริยํ ทุกฺขินฺทฺริยํ โสมนสฺสิ\-นฺทฺริยํ โทมนสฺสิ\-นฺทฺริยํ อุเปกฺขินฺทฺริยํ.}

\prose{กตมญฺจ ภิกฺขเว สุขินฺทฺริยํ, ยํ โข ภิกฺขเว กายิกํ สุขํ กายิกํ สาตํ กาย\-สมฺผสฺสชํ สุขํ สาตํ เวทยิตํ. อิทํ วุจฺจติ ภิกฺขเว สุขินฺทฺริยํ.}

\prose{กตมญฺจ ภิกฺขเว ทุกฺขินฺทฺริยํ, ยํ โข ภิกฺขเว กายิกํ ทุกฺขํ กายิกํ อสาตํ กาย\-สมฺผสฺสชํ ทุกฺขํ อสาตํ เวทยิตํ. อิทํ วุจฺจติ ภิกฺขเว ทุกฺขินฺทฺริยํ.}

\prose{กตมญฺจ ภิกฺขเว โสมนสฺสิ\-นฺทฺริยํ, ยํ โข ภิกฺขเว เจตสิกํ สุขํ เจตสิกํ สาตํ มโน\-สมฺผสฺสชํ สุขํ สาตํ เวทยิตํ. อิทํ วุจฺจติ ภิกฺขเว โสมนสฺสิ\-นฺทฺริยํ.}

\prose{กตมญฺจ ภิกฺขเว โทมนสฺสิ\-นฺทฺริยํ, ยํ โข ภิกฺขเว เจตสิกํ ทุกฺขํ เจตสิกํ อสาตํ มโน\-สมฺผสฺสชํ ทุกฺขํ อสาตํ เวทยิตํ. อิทํ วุจฺจติ ภิกฺขเว โทมนสฺสิ\-นฺทฺริยํ.}

\prose{กตมญฺจ ภิกฺขเว อุเปกฺขินฺทฺริยํ, ยํ โข ภิกฺขเว กายิกํ วา เจตสิกํ วา เนวสาตํ นาสาตํ เวทยิตํ. อิทํ วุจฺจติ ภิกฺขเว อุเปกฺขินฺทฺริยํ.}

\prose{ตตฺร ภิกฺขเว ยญฺจ สุขินฺทฺริยํ ยญฺจ โสมนสฺสิ\-นฺทฺริยํ, สุขา สา เวทนา ทฏฺฐพฺพา. ตตฺร ภิกฺขเว ยญฺจ ทุกฺขินฺทฺริยํ ยญฺจ โทมนสฺสิ\-นฺทฺริยํ, ทุกฺขา สา เวทนา ทฏฺฐพฺพา. ตตฺร ภิกฺขเว ยทิทํ อุเปกฺขินฺทฺริยํ. อทุกฺขมสุขา สา เวทนา}

\prose{\csromanfolio{187}ทฏฺฐพฺพา. อิติ โข ภิกฺขเว อิมานิ ปญฺจินฺทฺริยานิ ปญฺจ หุตฺวา ตีณิ โหนฺติ, ตีณิ หุตฺวา ปญฺจ โหนฺติ ปริยาเยนาติ.  อฏฺฐมํ.}

\csromanmatikaanchor{mtk.115}
\csromantocmarkref{subparagraph}{9. กฏฺโฐปมสุตฺต}{mtk.115}
\bookmark[dest={mtk.115},level=5]{9. กฏฺโฐปมสุตฺต}
\csromanheader{6}{9. กฏฺโฐ\-ปม\-สุตฺต}
\proseitem{509}{ปญฺจิมานิ ภิกฺขเว อินฺทฺริยานิ. กตมานิ ปญฺจ, สุขินฺทฺริยํ ทุกฺขินฺทฺริยํ โสมนสฺสิ\-นฺทฺริยํ โทมนสฺสิ\-นฺทฺริยํ อุเปกฺขินฺทฺริยํ. สุขเวทนิยํ ภิกฺขเว ผสฺสํ ปฏิจฺจ อุปฺปชฺชติ สุขินฺทฺริยํ, โส สุขิโตว สมาโน “สุขิโตสฺมี”ติ ปชานาติ, ตสฺเสว สุขเวทนิยสฺส ผสฺสสฺส นิโรธา ยํ ตชฺชํ เวทยิตํ สุขเวทนิยํ ผสฺสํ ปฏิจฺจ อุปฺปนฺนํ สุขินฺทฺริยํ, ตํ นิรุชฺฌติ, ตํ วูปสมฺมตีติ ปชานาติ.}

\prose{ทุกฺข\-เวทนิยํ ภิกฺขเว ผสฺสํ ปฏิจฺจ อุปฺปชฺชติ ทุกฺขินฺทฺริยํ, โส ทุกฺขิโตว สมาโน “ทุกฺขิโตสฺมี”ติ ปชานาติ, ตสฺเสว ทุกฺข\-เวทนิยสฺส ผสฺสสฺส นิโรธา ยํ ตชฺชํ เวทยิตํ ทุกฺข\-เวทนิยํ ผสฺสํ ปฏิจฺจ อุปฺปนฺนํ ทุกฺขินฺทฺริยํ, ตํ นิรุชฺฌติ, ตํ วูปสมฺมตีติ ปชานาติ.}

\prose{โสมนสฺส\-เวทนิยํ ภิกฺขเว ผสฺสํ ปฏิจฺจ อุปฺปชฺชติ โสมนสฺสิ\-นฺทฺริยํ, โส สุมโนว สมาโน “สุมโนสฺมี”ติ ปชานาติ, ตสฺเสว โสมนสฺส\-เวทนิยสฺส ผสฺสสฺส นิโรธา ยํ ตชฺชํ เวทยิตํ โสมนสฺส\-เวทนิยํ ผสฺสํ ปฏิจฺจ อุปฺปนฺนํ โสมนสฺสิ\-นฺทฺริยํ, ตํ นิรุชฺฌติ, ตํ วูปสมฺมตีติ ปชานาติ.}

\prose{โทมนสฺส\-เวทนิยํ ภิกฺขเว ผสฺสํ ปฏิจฺจ อุปฺปชฺชติ โทมนสฺสิ\-นฺทฺริยํ, โส ทุมฺมโนว สมาโน “ทุมฺมโนสฺมี”ติ ปชานาติ, ตสฺเสว โทมนสฺส\-เวทนิยสฺส ผสฺสสฺสฺส นิโรธา ยํ ตชฺชํ เวทยิตํ โทมนสฺส\-เวทนิยํ ผสฺสํ ปฏิจฺจ อุปฺปนฺนํ โทมนสฺสิ\-นฺทฺริยํ, ตํ นิรุชฺฌติ, ตํ วูปสมฺมตีติ ปชานาติ.}

\prose{อุเปกฺขา\-เวทนิยํ ภิกฺขเว ผสฺสํ ปฏิจฺจ อุปฺปชฺชติ อุเปกฺขินฺทฺริยํ, โส อุเปกฺขโกว สมาโน “อุเปกฺขโกสฺมี”ติ ปชานาติ, ตสฺเสว อุเปกฺขา\-เวทนิยสฺส ผสฺสสฺส นิโรธา ยํ ตชฺชํ เวทยิตํ อุเปกฺขา\-เวทนิยํ ผสฺสํ ปฏิจฺจ อุปฺปนฺนํ อุเปกฺขินฺทฺริยํ, ตํ นิรุชฺฌติ, ตํ วูปสมฺมตีติ ปชานาติ.}

\prose{เสยฺยถาปิ ภิกฺขเว ทฺวินฺนํ กฏฺฐานํ สํฆฏฺฏนสโมธานา\csromansharedfootnote{8c75b2f320fa67fa}{สํฆฏฺฏนา สโมธานา (อิ, ก), สํฆฏนสโมธานา (สฺยา, กํ)} อุสฺมา ชายติ, เตโช อภินิพฺพตฺตติ, เตสํเยว กฏฺฐานํ นานาภาวา\-วินิกฺเขปา\csromansharedfootnote{81a8dcc7ff810e71}{นานาภาวนิกฺเขปา (สฺยา, กํ, อิ, ก)} \csromanfolio{188}ยา ตชฺชา อุสฺมา, สา นิรุชฺฌติ, สา วูปสมฺมติ. เอวเมว โข ภิกฺขเว สุขเวทนิยํ ผสฺสํ ปฏิจฺจ อุปฺปชฺชติ สุขินฺทฺริยํ, โส สุขิโตว สมาโน “สุขิโตสฺมี”ติ ปชานาติ, ตสฺเสว สุขเวทนิยสฺส ผสฺสสฺส นิโรธา ยํ ตชฺชํ เวทยิตํ สุขเวทนิยํ ผสฺสํ ปฏิจฺจ อุปฺปชฺชติ สุขินฺทฺริยํ, ตํ นิรุชฺฌติ, ตํ วูปสมฺมตีติ ปชานาติ. ทุกฺข\-เวทนิยํ ภิกฺขเว ผสฺสํ ปฏิจฺจ ฯเปฯ. โสมนสฺส\-เวทนิยํ ภิกฺขเว ผสฺสํ ปฏิจฺจ ฯเปฯ. โทมนสฺส\-เวทนิยํ ภิกฺขเว ผสฺสํ ปฏิจฺจ ฯเปฯ. อุเปกฺขา\-เวทนิยํ ภิกฺขเว ผสฺสํ ปฏิจฺจ อุปฺปชฺชติ อุเปกฺขินฺทฺริยํ, โส อุเปกฺขโกว สมาโน “อุเปกฺขโกสฺมี”ติ ปชานาติ, ตสฺเสว อุเปกฺขา\-เวทนิยสฺส ผสฺสสฺส นิโรธา ยํ ตชฺชํ เวทยิตํ อุเปกฺขา\-เวทนิยํ ผสฺสํ ปฏิจฺจ อุปฺปชฺชติ อุเปกฺขินฺทฺริยํ. ตํ นิรุชฺฌติ, ตํ วูปสมฺมตีติ ปชานาติ.  นวมํ.}

\csromanmatikaanchor{mtk.116}
\csromantocmarkref{subparagraph}{10. อุปฺปฏิปาฏิกสุตฺต}{mtk.116}
\bookmark[dest={mtk.116},level=5]{10. อุปฺปฏิปาฏิกสุตฺต}
\csromanheader{6}{10. อุปฺปฏิปาฏิก\-สุตฺต}
\proseitem{510}{ปญฺจิมานิ ภิกฺขเว อินฺทฺริยานิ. กตมานิ ปญฺจ, ทุกฺขินฺทฺริยํ โทมนสฺสิ\-นฺทฺริยํ สุขินฺทฺริยํ โสมนสฺสิ\-นฺทฺริยํ อุเปกฺขินฺทฺริยํ. อิธ ภิกฺขเว ภิกฺขุโน อปฺปมตฺตสฺส อาตาปิโน ปหิตตฺตสฺส วิหรโต อุปฺปชฺชติ ทุกฺขินฺทฺริยํ, โส เอวํ ปชานาติ “อุปฺปนฺนํ โข เม อิทํ ทุกฺขินฺทฺริยํ, ตญฺจ โข สนิมิตฺตํ สนิทานํ สสงฺขารํ สปฺปจฺจยํ, ตญฺจ อนิมิตฺตํ อนิทานํ อสงฺขารํ อปฺปจฺจยํ ทุกฺขินฺทฺริยํ อุปฺปชฺชิสฺสตี”ติ เนตํ ฐานํ วิชฺชติ. โส ทุกฺขินฺทฺริยญฺจ ปชานาติ, ทุกฺขินฺทฺริยสมุทยญฺจ ปชานาติ, ทุกฺขินฺทฺริยนิโรธญฺจ ปชานาติ. ยตฺถ จุปฺปนฺนํ ทุกฺขินฺทฺริยํ อปริเสสํ นิรุชฺฌติ, ตญฺจ ปชานาติ. กตฺถ จุปฺปนฺนํ ทุกฺขินฺทฺริยํ อปริเสสํ นิรุชฺฌติ, อิธ ภิกฺขเว ภิกฺขุ วิวิจฺเจว กาเมหิ วิวิจฺจ อกุสเลหิ ธมฺเมหิ สวิตกฺกํ สวิจารํ วิเวกชํ ปีติสุขํ ปฐมํ ฌานํ อุปสมฺปชฺช วิหรติ, เอตฺถ จุปฺปนฺนํ ทุกฺขินฺทฺริยํ อปริเสสํ นิรุชฺฌติ. อยํ วุจฺจติ ภิกฺขเว ภิกฺขุ อญฺญาสิ ทุกฺขิ\-นฺทฺริยสฺส นิโรธํ, ตทตฺถาย จิตฺตํ อุปสํหรติ.}

\prose{อิธ ปน ภิกฺขเว ภิกฺขุโน อปฺปมตฺตสฺส อาตาปิโน ปหิตตฺตสฺส วิหรโต อุปฺปชฺชติ โทมนสฺสิ\-นฺทฺริยํ, โส เอวํ ปชานาติ “อุปฺปนฺนํ โข เม อิทํ โทมนสฺสิ\-นฺทฺริยํ, ตญฺจ โข สนิมิตฺตํ สนิทานํ สสงฺขารํ สปฺปจฺจยํ, ตญฺจ อนิมิตฺตํ อนิทานํ อสงฺขารํ อปฺปจฺจยํ โทมนสฺสิ\-นฺทฺริยํ อุปฺปชฺชิสฺสตี”ติ เนตํ ฐานํ วิชฺชติ. โส โทมนสฺสินฺทฺริยญฺจ ปชานาติ, โทมนสฺสินฺทฺริยสมุทยญฺจ ปชานาติ, โทมนสฺสินฺทฺริยนิโรธญฺจ ปชานาติ. ยตฺถ จุปฺปนฺนํ โทมนสฺสิ\-นฺทฺริยํ}

\prose{\csromanfolio{189}อปริเสสํ นิรุชฺฌติ, ตญฺจ ปชานาติ. กตฺถ จุปฺปนฺนํ โทมนสฺสิ\-นฺทฺริยํ อปริเสสํ นิรุชฺฌติ, อิธ ภิกฺขเว ภิกฺขุ วิตกฺก\-วิจารานํ วูปสมา อชฺฌตฺตํ สมฺปสาทนํ เจตโส เอโกทิภาวํ อวิตกฺกํ อวิจารํ สมาธิชํ ปีติสุขํ ทุติยํ ฌานํ อุปสมฺปชฺช วิหรติ, เอตฺถ จุปฺปนฺนํ โทมนสฺสิ\-นฺทฺริยํ อปริเสสํ นิรุชฺฌติ. อยํ วุจฺจติ ภิกฺขเว ภิกฺขุ อญฺญาสิ โทมนสฺสิ\-นฺทฺริยสฺส นิโรธํ, ตทตฺถาย จิตฺตํ อุปสํหรติ.}

\prose{อิธ ปน ภิกฺขเว ภิกฺขุโน อปฺปมตฺตสฺส อาตาปิโน ปหิตตฺตสฺส วิหรโต อุปฺปชฺชติ สุขินฺทฺริยํ, โส เอวํ ปชานาติ “อุปฺปนฺนํ โข เม อิทํ สุขินฺทฺริยํ, ตญฺจ โข สนิมิตฺตํ สนิทานํ สสงฺขารํ สปฺปจฺจยํ, ตญฺจ อนิมิตฺตํ อนิทานํ อสงฺขารํ อปฺปจฺจยํ สุขินฺทฺริยํ อุปฺปชฺชิสฺสตี”ติ เนตํ ฐานํ วิชฺชติ. โส สุขินฺทฺริยญฺจ ปชานาติ สุขินฺทฺริยสมุทยญฺจ ปชานาติ, สุขินฺทฺริยนิโรธญฺจ ปชานาติ. ยตฺถ จุปฺปนฺนํ สุขินฺทฺริยํ อปริเสสํ นิรุชฺฌติ, ตญฺจ ปชานาติ. กตฺถ จุปฺปนฺนํ สุขินฺทฺริยํ อปริเสสํ นิรุชฺฌติ, อิธ ภิกฺขเว ภิกฺขุ ปีติยา จ วิราคา อุเปกฺขโก จ วิหรติ สโต จ สมฺปชาโน สุขญฺจ กาเยน ปฏิสํเวเทติ. ยํ ตํ อริยา อาจิกฺขนฺติ “อุเปกฺขโก สติมา สุขวิหารี”ติ, ตติยํ ฌานํ อุปสมฺปชฺช วิหรติ, เอตฺถ จุปฺปนฺนํ สุขินฺทฺริยํ อปริเสสํ นิรุชฺฌติ. อยํ วุจฺจติ ภิกฺขเว ภิกฺขุ อญฺญาสิ สุขินฺทฺริยสฺส นิโรธํ, ตทตฺถาย จิตฺตํ อุปสํหรติ.}

\prose{อิธ ปน ภิกฺขเว ภิกฺขุโน อปฺปมตฺตสฺส อาตาปิโน ปหิตตฺตสฺส วิหรโต อุปฺปชฺชติ โสมนสฺสิ\-นฺทฺริยํ, โส เอวํ ปชานาติ “อุปฺปนฺนํ โข เม อิทํ โสมนสฺสิ\-นฺทฺริยํ, ตญฺจ โข สนิมิตฺตํ สนิทานํ สสงฺขารํ สปฺปจฺจยํ, ตญฺจ อนิมิตฺตํ อนิทานํ อสงฺขารํ อปฺปจฺจยํ โสมนสฺสิ\-นฺทฺริยํ อุปฺปชฺชิสฺสตี”ติ เนตํ ฐานํ วิชฺชติ. โส โสมนสฺสินฺทฺริยญฺจ ปชานาติ, โสมนสฺสินฺทฺริยสมุทยญฺจ ปชานาติ, โสมนสฺสินฺทิยนิโรธญฺจ ปชานาติ. ยตฺถ จุปฺปนฺนํ โสมนสฺสิ\-นฺทฺริยํ อปริเสสํ นิรุชฺฌติ, ตญฺจ ปชานาติ. กตฺถ จุปฺปนฺนํ โสมนสฺสิ\-นฺทฺริยํ อปริเสสํ นิรุชฺฌติ, อิธ ภิกฺขเว ภิกฺขุ สุขสฺส จ ปหานา ทุกฺขสฺส จ ปหานา ปุพฺเพว โสมนสฺส\-โทมนสฺสานํ อตฺถงฺคมา อทุกฺขมสุขํ อุเปกฺขา\-สติ\-ปาริสุทฺธิํ จตุตฺถํ ฌานํ อุปสมฺปชฺช วิหรติ, เอตฺถ จุปฺปนฺนํ โสมนสฺสิ\-นฺทฺริยํ อปริเสสํ นิรุชฺฌติ. อยํ วุจฺจติ ภิกฺขเว ภิกฺขุ อญฺญาสิ โสมนสฺสิ\-นฺทฺริยสฺส นิโรธํ, ตทตฺถาย จิตฺตํ อุปสํหรติ.}

\prosewithclosermidruled{\csromanfolio{190}อิธ ปน ภิกฺขเว ภิกฺขุโน อปฺปมตฺตสฺส อาตาปิโน ปหิตตฺตสฺส วิหรโต อุปฺปชฺชติ อุเปกฺขินฺทฺริยํ, โส เอวํ ปชานาติ “อุปฺปนฺนํ โข เม อิทํ อุเปกฺขินฺทฺริยํ, ตญฺจ โข สนิมิตฺตํ สนิทานํ สสงฺขารํ สปฺปจฺจยํ, ตญฺจ อนิมิตฺตํ อนิทานํ อสงฺขารํ อปฺปจฺจยํ อุเปกฺขินฺทฺริยํ อุปฺปชฺชิสฺสตี”ติ เนตํ ฐานํ วิชฺชติ. โส อุเปกฺขินฺทฺริยญฺจ ปชานาติ, อุเปกฺขินฺทฺริยสมุทยญฺจ ปชานาติ, อุเปกฺขินฺทฺริยนิโรธญฺจ ปชานาติ. ยตฺถ จุปฺปนฺนํ อุเปกฺขินฺทฺริยํ อปริเสสํ นิรุชฺฌติ, ตญฺจ ปชานาติ. กตฺถ จุปฺปนฺนํ อุเปกฺขินฺทฺริยํ อปริเสสํ นิรุชฺฌติ, อิธ ภิกฺขเว ภิกฺขุ สพฺพโส เนว\-สญฺญา\-นา\-สญฺญา\-ยตนํ สมติกฺกมฺม สญฺญาเวทยิต\-นิโรธํ อุปสมฺปชฺช วิหรติ, เอตฺถ จุปฺปนฺนํ อุเปกฺขินฺทฺริยํ อปริเสสํ นิรุชฺฌติ. อยํ วุจฺจติ ภิกฺขเว ภิกฺขุ อญฺญาสิ อุเปกฺขิ\-นฺทฺริยสฺส นิโรธํ, ตทตฺถาย จิตฺตํ อุปสํหรตีติ.  ทสมํ.}{สุขินฺทฺริย\-วคฺโค จตุตฺโถ.}
\csromancenter{\textbf{ตสฺสุทฺทานํ}}
\setlength{\csromangathapagewidth}{0pt}
\setlength{\csromangathawakwidth}{0pt}
\csromangathameasuregroup{สุทฺธิกญฺจ โสโต อรหา, \\ วิภงฺเคน ตโย วุตฺตา,}{\csromangathabat{สุทฺธิกญฺจ โสโต อรหา,}{ทุเว สมณพฺราหฺมณา.} \\ \csromangathabat{วิภงฺเคน ตโย วุตฺตา,}{กฏฺโฐ อุปฺปฏิปาฏิกนฺติ.}}
\csromangathasetleft{สุทฺธิกญฺจ โสโต อรหา, \\ วิภงฺเคน ตโย วุตฺตา,}
\csromangathagroup{\csromangathastanza{\csromangathabat{สุทฺธิกญฺจ โสโต อรหา,}{ทุเว สมณพฺราหฺมณา.} \\ \csromangathabat{วิภงฺเคน ตโย วุตฺตา,}{กฏฺโฐ อุปฺปฏิปาฏิกนฺติ.}}}

\csromanmatikaanchor{mtk.117}
\csromantocmarknopage{subparagraph}{5. ชราวคฺค}{mtk.117}
\bookmark[dest={mtk.117},level=5]{5. ชราวคฺค}
\csromanheader{6}{5. \textbf{ชราวคฺค}}
\csromanmatikaanchor{mtk.118}
\csromantocmarkref{subparagraph}{1. ชราธมฺมสุตฺต}{mtk.118}
\bookmark[dest={mtk.118},level=5]{1. ชราธมฺมสุตฺต}
\csromanheader{6}{1. ชราธมฺม\-สุตฺต}
\proseitem{511}{เอวํ เม สุตํ\csromandash{}เอกํ สมยํ ภควา สาวตฺถิยํ วิหรติ ปุพฺพาราเม มิคาร\-มาตุ\-ปาสาเท. เตน โข ปน สมเยน ภควา สายนฺหสมยํ ปฏิสลฺลานา วุฏฺฐิโต ปจฺฉาตเป นิสินฺโน โหติ ปิฏฺฐิํ โอตาปยมาโน.}

\prose{อถ โข อายสฺมา อานนฺโท เยน ภควา เตนุปสงฺกมิ, อุปสงฺกมิตฺวา ภควนฺตํ อภิวาเทตฺวา ภควโต คตฺตานิ ปาณินา อโนมชฺชนฺโต ภควนฺตํ เอตทโวจ “อจฺฉริยํ ภนฺเต, อพฺภุตํ ภนฺเต, น เจวํ ทานิ ภนฺเต ภควโต ตาว ปริสุทฺโธ ฉวิวณฺโณ ปริโยทาโต, สิถิลานิ จ คตฺตานิ สพฺพานิ วลิยชาตานิ, ปุรโต}

\prose{\csromanfolio{191}ปพฺภาโร จ กาโย, ทิสฺสติ จ อินฺทฺริยานํ อญฺญถตฺตํ จกฺขุนฺทฺริยสฺส โสตินฺทฺริยสฺส ฆานินฺทฺริยสฺส ชิวฺหินฺทฺริยสฺส กายินฺทฺริยสฺสา”ติ.}

\prose{เอวํ เหตํ อานนฺท โหติ ชราธมฺโม โยพฺพญฺเญ, พฺยาธิธมฺโม อาโรเคฺย, มรณธมฺโม ชีวิเต, น เจว ตาว ปริสุทฺโธ โหติ ฉวิวณฺโณ ปริโยทาโต, สิถิลานิ จ โหนฺติ คตฺตานิ สพฺพานิ วลิยชาตานิ, ปุรโต ปพฺภาโร จ กาโย, ทิสฺสติ จ อินฺทฺริยานํ อญฺญถตฺตํ จกฺขุนฺทฺริยสฺส โสตินฺทฺริยสฺส ฆานินฺทฺริยสฺส ชิวฺหินฺทฺริยสฺส กายินฺทฺริยสฺสาติ.}

\prose{อิทมโวจ ภควา, อิทํ วตฺวา จ พุคโต อถาปรํ เอตทโวจ สตฺถา\csromandash{}}

\setlength{\csromangathapagewidth}{0pt}
\setlength{\csromangathawakwidth}{0pt}
\csromangathameasuregroup{“ธี ตํ ชมฺมิ ชเร อตฺถุ, \\ ตาว มโนรมํ พิมฺพํ, \\ โยปิ วสฺสสตํ ชีเว,}{\csromangathabat{“ธี ตํ ชมฺมิ ชเร อตฺถุ,}{ทุพฺพณฺณกรณี ชเร.} \\ \csromangathabat{ตาว มโนรมํ พิมฺพํ,}{ชราย อภิมทฺทิตํ.} \\ \csromangathabat{โยปิ วสฺสสตํ ชีเว,}{โสปิ มจฺจุปรายโณ\textsuperscript{1}.}}
\csromangathasetleft{“ธี ตํ ชมฺมิ ชเร อตฺถุ, \\ ตาว มโนรมํ พิมฺพํ, \\ โยปิ วสฺสสตํ ชีเว,}
\csromangathagroup{\csromangathastanza{\csromangathabat{“ธี ตํ ชมฺมิ ชเร อตฺถุ,}{ทุพฺพณฺณกรณี ชเร.} \\ \csromangathabat{ตาว มโนรมํ พิมฺพํ,}{ชราย อภิมทฺทิตํ.} \\ \csromangathabat{โยปิ วสฺสสตํ ชีเว,}{โสปิ มจฺจุปรายโณ\csromansharedfootnote{4cb0daef375b1edc}{สพฺเพ มจฺจุปรายนา (สฺยา, กํ, ก)}.}}}

\vspace{\tipitakaparskip}
\csromancenter{น กิญฺจิ ปริวชฺเชติ, สพฺพเมวาภิมทฺทตี”ติ. ปฐมํ.}
\csromanmatikaanchor{mtk.119}
\csromantocmarkref{subparagraph}{2. อุณฺณาภพฺราหฺมณสุตฺต}{mtk.119}
\bookmark[dest={mtk.119},level=5]{2. อุณฺณาภพฺราหฺมณสุตฺต}
\csromanheader{6}{2. อุณฺณาภ\-พฺราหฺมณ\-สุตฺต}
\proseitem{512}{สาวตฺถิ\-นิทานํ. อถ โข อุณฺณาโภ พฺราหฺมโณ เยน ภควา เตนุปสงฺกมิ, อุปสงฺกมิตฺวา ภควตา สทฺธิํ สมฺโมทิ, สมฺโมทนียํ กถํ สารณียํ วีติสาเรตฺวา เอกมนฺตํ นิสีทิ, เอกมนฺตํ นิสินฺโน โข อุณฺณาโภ พฺราหฺมโณ ภควนฺตํ เอตทโวจ\csromandash{}}

\prose{ปญฺจิมานิ โภ โคตม อินฺทฺริยานิ นานาวิสยานิ นานาโคจรานิ น อญฺญมญฺญสฺส โคจรวิสยํ ปจฺจนุโภนฺติ. กตมานิ ปญฺจ, จกฺขุนฺทฺริยํ โสตินฺทฺริยํ ฆานินฺทฺริยํ ชิวฺหินฺทฺริยํ กายินฺทฺริยํ. อิเมสํ นุ โข โภ โคตม ปญฺจนฺนํ อินฺทฺริยานํ นานาวิสยานํ นานาโคจรานํ น อญฺญมญฺญสฺส โคจรวิสยํ ปจฺจนุโภนฺตานํ กิํ ปฏิสรณํ, โก จ เนสํ โคจรวิสยํ ปจฺจนุโภตีติ.}

\prose{ปญฺจิมานิ พฺราหฺมณ อินฺทฺริยานิ นานาวิสยานิ นานาโคจรานิ น อญฺญมญฺญสฺส โคจรวิสยํ ปจฺจนุโภนฺติ. กตมานิ ปญฺจ, จกฺขุนฺทฺริยํ \csromanfolio{192}โสตินฺทฺริยํ ฆานินฺทฺริยํ ชิวฺหินฺทฺริยํ กายินฺทฺริยํ. อิเมสํ โข พฺราหฺมณ ปญฺจนฺนํ อินฺทฺริยานํ นานาวิสยานํ นานาโคจรานํ น อญฺญมญฺญสฺส โคจรวิสยํ ปจฺจนุโภนฺตานํ มโน ปฏิสรณํ, มโนว เนสํ โคจรวิสยํ ปจฺจนุโภตีติ.}

\prose{มนสฺส ปน โภ โคตม กิํ ปฏิสรณนฺติ, มนสฺส โข พฺราหฺมณ สติ ปฏิสรณนฺติ. สติยา ปน โภ โคตม กิํ ปฏิสรณนฺติ, สติยา โข พฺราหฺมณ วิมุตฺติ ปฏิสรณนฺติ. วิมุตฺติยา ปน โภ โคตม กิํ ปฏิสรณนฺติ, วิมุตฺติยา โข พฺราหฺมณ นิพฺพานํ ปฏิสรณนฺติ. นิพฺพานสฺส ปน โภ โคตม กิํ ปฏิสรณนฺติ, อจฺจยาสิ\csromansharedfootnote{9f0e58dfbdb576ee}{อจฺจสรา (สี, สฺยา, กํ), อชฺฌปรํ (อิ, ก)} พฺราหฺมณ ปญฺหํ นาสกฺขิ ปญฺหสฺส ปริยนฺตํ คเหตุํ, นิพฺพาโนคธํ หิ พฺราหฺมณ พฺรหฺมจริยํ วุสฺสติ นิพฺพาน\-ปรายณํ นิพฺพาน\-ปริโยสานนฺติ.}

\prose{อถ โข อุณฺณาโภ พฺราหฺมโณ ภควโต ภาสิตํ อภินนฺทิตฺวา อนุโมทิตฺวา อุฏฺฐายาสนา ภควนฺตํ อภิวาเทตฺวา ปทกฺขิณํ กตฺวา ปกฺกามิ.}

\prose{อถ โข ภควา อจิรปกฺกนฺเต อุณฺณาเภ พฺราหฺมเณ ภิกฺขู อามนฺเตสิ “เสยฺยถาปิ ภิกฺขเว กูฏาคาเร วา กูฏาคาร\-สาลายํ วา\csromansharedfootnote{2719e88633900ef9}{กูฏาคารํ วา กูฏาคารสาลํ วา อุตฺตราย (กสี)} ปาจีนวาตปานา สูริเย อุคฺคจฺฉนฺเต วาตปาเนน รสฺมิ\csromansharedfootnote{2d3cd8de69fa061a}{รสฺมิโย (สฺยา, ก)} ปวิสิตฺวา กฺวาสฺส\csromansharedfootnote{0d14269988f31815}{กาย (สฺยา, ก)} ปติฏฺฐิตา”ติ. ปจฺฉิมายํ ภนฺเต ภิตฺติยนฺติ. เอวเมว โข ภิกฺขเว อุณฺณาภสฺส พฺราหฺมณสฺส ตถาคเต สทฺธา นิวิฏฺฐา, มูลชาตา ปติฏฺฐิตา, ทฬฺหา อสํหาริยา สมเณน วา พฺราหฺมเณน วา เทเวน วา มาเรน วา พฺรหฺมุนา วา เกนจิ วา โลกสฺมิํ. อิมมฺหิ เจ ภิกฺขเว สมเย อุณฺณาโภ พฺราหฺมโณ กาลํ กเรยฺย, นตฺถิ ตํ สํโยชนํ, เยน สํโยชเนน สํยุตฺโต อุณฺณาโภ พฺราหฺมโณ ปุน อิมํ โลกํ อาคจฺเฉยฺยาติ.  ทุติยํ.}

\csromanmatikaanchor{mtk.120}
\csromantocmarkref{subparagraph}{3. สาเกตสุตฺต}{mtk.120}
\bookmark[dest={mtk.120},level=5]{3. สาเกตสุตฺต}
\csromanheader{6}{3. \textbf{สาเกตสุตฺต}}
\proseitem{513}{เอวํ เม สุตํ\csromandash{}เอกํ สมยํ ภควา สาเกเต วิหรติ อญฺชนวเน มิคทาเย. ตตฺร โข ภควา ภิกฺขู อามนฺเตสิ “อตฺถิ นุ โข}

\prose{\csromanfolio{193}ภิกฺขเว ปริยาโย, ยํ ปริยายํ อาคมฺม ยานิ ปญฺจินฺทฺริยานิ, ตานิ ปญฺจ พลานิ โหนฺติ. ยานิ ปญฺจ พลานิ, ตานิ ปญฺจินฺทฺริยานิ โหนฺตี”ติ.}

\prose{ภควํมูลกา โน ภนฺเต ธมฺมา ภควํ\-เนตฺติกา ภควํ\-ปฏิสรณา, สาธุ วต ภนฺเต, ภควนฺตํเยว ปฏิภาตุ เอตสฺส ภาสิตสฺส อตฺโถ, ภควโต สุตฺวา ภิกฺขู ธาเรสฺสนฺตีติ. อตฺถิ ภิกฺขเว ปริยาโย, ยํ ปริยายํ อาคมฺม ยานิ ปญฺจินฺทฺริยานิ, ตานิ ปญฺจ พลานิ, โหนฺติ. ยานิ ปญฺจ พลานิ, ตานิ ปญฺจินฺทฺริยานิ โหนฺติ.}

\prose{กตโม จ ภิกฺขเว ปริยาโย, ยํ ปริยายํ อาคมฺม ยานิ ปญฺจินฺทฺริยานิ, ตานิ ปญฺจ พลานิ โหนฺติ. ยานิ ปญฺจ พลานิ, ตานิ ปญฺจินฺทฺริยานิ โหนฺติ. ยํ ภิกฺขเว สทฺธินฺทฺริยํ, ตํ สทฺธาพลํ, ยํ สทฺธาพลํ, ตํ สทฺธินฺทฺริยํ, ยํ วีริยินฺทฺริยํ, ตํ วีริยพลํ, ยํ วีริยพลํ, ตํ วีริยินฺทฺริยํ, ยํ สตินฺทฺริยํ, ตํ สติพลํ, ยํ สติพลํ, ตํ สตินฺทฺริยํ, ยํ สมาธินฺทฺริยํ, ตํ สมาธิพลํ, ยํ สมาธิพลํ, ตํ สมาธินฺทฺริยํ, ยํ ปญฺญินฺทฺริยํ, ตํ ปญฺญาพลํ, ยํ ปญฺญาพลํ, ตํ ปญฺญินฺทฺริยํ. เสยฺยถาปิ ภิกฺขเว นที ปาจีนนินฺนา ปาจีนโปณา ปาจีนปพฺภารา, ตสฺส มชฺเฌ ทีโป, อตฺถิ ภิกฺขเว ปริยาโย, ยํ ปริยายํ อาคมฺม ตสฺสา นทิยา เอโก โสโตเตฺวว สงฺขฺยํ\csromansharedfootnote{95a4bc23b28b0157}{สงฺขํ (สี, สฺยา, กํ)} คจฺฉติ, อตฺถิ ปน ภิกฺขเว ปริยาโย, ยํ ปริยายํ อาคมฺม ตสฺสา นทิยา เทฺว โสตานิเตฺวว สงฺขฺยํ คจฺฉติ.}

\prose{กตโม จ ภิกฺขเว ปริยาโย, ยํ ปริยายํ อาคมฺม ตสฺสา นทิยา เอโก โสโตเตฺวว สงฺขฺยํ คจฺฉติ. ยญฺจ ภิกฺขเว ตสฺส ทีปสฺส ปุริมนฺเต\csromansharedfootnote{960415d4271bbaee}{ปุรตฺถิมนฺเต (สี, สฺยา, กํ, อิ)} อุทกํ, ยญฺจ ปจฺฉิมนฺเต อุทกํ, อยํ โข ภิกฺขเว ปริยาโย, ยํ ปริยายํ อาคมฺม ตสฺสา นทิยา เอโก โสโตเตฺวว สงฺขฺยํ คจฺฉติ.}

\prose{กตโม จ ภิกฺขเว ปริยาโย, ยํ ปริยายํ อาคมฺม ตสฺสา นทิยา เทฺว โสตานิเตฺวว สงฺขฺยํ คจฺฉนฺติ. ยญฺจ ภิกฺขเว ตสฺส ทีปสฺส อุตฺตรนฺเต อุทกํ, ยญฺจ ทกฺขิณนฺเต อุทกํ, อยํ โข ภิกฺขเว ปริยาโย, ยํ ปริยายํ อาคมฺม ตสฺสา นทิยา เทฺว โสตานิเตฺวว สงฺขฺยํ คจฺฉนฺติ. เอวเมว โข ภิกฺขเว ยํ สทฺธินฺทฺริยํ, ตํ สทฺธาพลํ, ยํ สทฺธาพลํ, ตํ สทฺธินฺทฺริยํ, ยํ วีริยินฺทฺริยํ, ตํ วีริยพลํ, ยํ วีริยพลํ, ตํ วีริยินฺทฺริยํ, ยํ สตินฺทฺริยํ, ตํ}

\prose{\csromanfolio{194}สติพลํ, ยํ สติพลํ, ตํ สตินฺทฺริยํ, ยํ สมาธินฺทฺริยํ, ตํ สมาธิพลํ, ยํ สมาธิพลํ, ตํ สมาธินฺทฺริยํ, ยํ ปญฺญินฺทฺริยํ, ตํ ปญฺญาพลํ, ยํ ปญฺญาพลํ, ตํ ปญฺญินฺทฺริยํ. ปญฺจนฺนํ ภิกฺขเว อินฺทฺริยานํ ภาวิตตฺตา พหุลีกตตฺตา ภิกฺขุ อาสวานํ ขยา อนาสวํ เจโตวิมุตฺติํ ปญฺญาวิมุตฺติํ ทิฏฺเฐว ธมฺเม สยํ อภิญฺญา สจฺฉิกตฺวา อุปสมฺปชฺช วิหรตีติ.  ตติยํ.}

\csromanmatikaanchor{mtk.121}
\csromantocmarkref{subparagraph}{4. ปุพฺพโกฏฺฐกสุตฺต}{mtk.121}
\bookmark[dest={mtk.121},level=5]{4. ปุพฺพโกฏฺฐกสุตฺต}
\csromanheader{6}{4. ปุพฺพโกฏฺฐก\-สุตฺต}
\proseitem{514}{เอวํ เม สุตํ\csromandash{}เอกํ สมยํ ภควา สาวตฺถิยํ วิหรติ ปุพฺพโกฏฺฐเก. ตตฺร โข ภควา อายสฺมนฺตํ สาริปุตฺตํ อามนฺเตสิ “สทฺทหสิ\csromansharedfootnote{a1bb3c55216ade0e}{สทฺทหาสิ (สี, อิ)} ตฺวํ สาริปุตฺต สทฺธินฺทฺริยํ ภาวิตํ พหุลีกตํ อมโตคธํ โหติ อมตปรายณํ อมต\-ปริโยสานํ ฯเปฯ ปญฺญินฺทฺริยํ ภาวิตํ พหุลีกตํ อมโตคธํ โหติ อมตปรายณํ อมตปริโยสานนฺ”ติ.}

\prose{น ขฺวาหํ เอตฺถ ภนฺเต ภควโต สทฺธาย คจฺฉามิ สทฺธินฺทฺริยํ ฯเปฯ ปญฺญินฺทฺริยํ ภาวิตํ พหุลีกตํ อมโตคธํ โหติ อมตปรายณํ อมต\-ปริโยสานํ. เยสํ เหตํ ภนฺเต อญฺญาตํ อสฺส อทิฏฺฐํ อวิทิตํ อสจฺฉิกตํ อผสฺสิตํ\csromansharedfootnote{86f3ea9aa01a30a6}{อปสฺสิตํ (สี, สฺยา, กํ, ก), อผุสิตํ (อิ)} ปญฺญาย, เต ตตฺถ ปเรสํ สทฺธาย คจฺเฉยฺยุํ. สทฺธินฺทฺริยํ ฯเปฯ ปญฺญินฺทฺริยํ ภาวิตํ พหุลีกตํ อมโตคธํ โหติ อมตปรายณํ อมต\-ปริโยสานํ. เยสญฺจ โข เอตํ ภนฺเต ญาตํ ทิฏฺฐํ วิทิตํ สจฺฉิกตํ ผสฺสิตํ ปญฺญาย, นิกฺกงฺขา เต ตตฺถ นิพฺพิจิกิจฺฉา สทฺธินฺทฺริยํ ฯเปฯ ปญฺญินฺทฺริยํ ภาวิตํ พหุลีกตํ อมโตคธํ โหติ อมตปรายณํ อมต\-ปริโยสานํ. มยฺหญฺจ โข เอตํ ภนฺเต ญาตํ ทิฏฺฐํ วิทิตํ สจฺฉิกตํ ผสฺสิตํ ปญฺญาย, นิกฺกงฺขฺวาหํ ตตฺถ นิพฺพิจิกิจฺโฉ สทฺธินฺทฺริยํ ฯเปฯ ปญฺญินฺทฺริยํ ภาวิตํ พหุลีกตํ อมโตคธํ โหติ อมตปรายณํ อมต\-ปริโยสานนฺติ.}

\prose{สาธุ สาธุ สาริปุตฺต, เยสํ เหตํ สาริปุตฺต อญฺญาตํ อสฺส อทิฏฺฐํ อวิทิตํ อสจฺฉิกตํ อผสฺสิตํ ปญฺญาย, เต ตตฺถ ปเรสํ สทฺธาย คจฺเฉยฺยุํ. สทฺธินฺทฺริยํ ภาวิตํ พหุลีกตํ อมโตคธํ โหติ อมตปรายณํ อมต\-ปริโยสานํ ฯเปฯ ปญฺญินฺทฺริยํ ภาวิตํ พหุลีกตํ อมโตคธํ โหติ อมตปรายณํ อมต\-ปริโยสานํ. เยสญฺจ โข เอตํ สาริปุตฺต ญาตํ ทิฏฺฐํ วิทิตํ สจฺฉิกตํ ผสฺสิตํ ปญฺญาย, นิกฺกงฺขา เต ตตฺถ}

\csromanmatikaanchor{mtk.122}
\csromantocmarkref{subparagraph}{5-8. ปุพฺพารามสุตฺตานิ}{mtk.122}
\bookmark[dest={mtk.122},level=5]{5-8. ปุพฺพารามสุตฺตานิ}
\prose{\csromanfolio{195}นิพฺพิจิกิจฺฉา สทฺธินฺทฺริยํ ภาวิตํ พหุลีกตํ อมโตคธํ โหติ อมตปรายณํ อมต\-ปริโยสานํ ฯเปฯ ปญฺญินฺทฺริยํ ภาวิตํ พหุลีกตํ อมโตคธํ โหติ อมตปรายณํ อมต\-ปริโยสานนฺติ.  จตุตฺถํ.}

\csromanheader{2}{5. ปฐม\-ปุพฺพาราม\-สุตฺต}
\proseitem{515}{เอวํ เม สุตํ\csromandash{}เอกํ สมยํ ภควา สาวตฺถิยํ วิหรติ ปุพฺพาราเม มิคาร\-มาตุ\-ปาสาเท. ตตฺร โข ภวคา ภิกฺขู อมนฺเตสิ “กตินํ นุ โข ภิกฺขเว อินฺทฺริยานํ ภาวิตตฺตา พหุลีกตตฺตา ขีณาสโว ภิกฺขุ อญฺญํ พฺยากโรติ ‘ขีณา ชาติ, วุสิตํ พฺรหฺมจริยํ, กตํ กรณียํ, นาปรํ อิตฺถตฺตายาติ ปชานามี’ติ”.}

\prose{ภควํมูลกา โน ภนฺเต ธมฺมา ฯเปฯ. เอกสฺส โข ภิกฺขเว อินฺทฺริยสฺส ภาวิตตฺตา พหุลีกตตฺตา ขีณาสโว ภิกฺขุ อญฺญํ พฺยากโรติ “ขีณา ชาติ, วุสิตํ พฺรหฺมจริยํ, กตํ กรณียํ, นาปรํ อิตฺถตฺตายาติ ปชานามี”ติ. กตมสฺส เอกสฺส, ปญฺญินฺทฺริยสฺส. ปญฺญวโต ภิกฺขเว อริยสาวกสฺส ตทนฺวยา สทฺธา สณฺฐาติ, ตทนฺวยํ วีริยํ สณฺฐาติ, ตทนฺวยา สติ สณฺฐาติ, ตทนฺวโย สมาธิ สณฺฐาติ. อิมสฺส โข ภิกฺขเว เอกสฺส อินฺทฺริยสฺส ภาวิตตฺตา พหุลีกตตฺตา ขีณาสโว ภิกฺขุ อญฺญํ พฺยากโรติ “ขีณา ชาติ, วุสิตํ พฺรหฺมจริยํ, กตํ กรณียํ, นาปรํ อิตฺถตฺตายาติ ปชานามี”ติ.  ปญฺจมํ.}

\csromanheader{2}{6. ทุติย\-ปุพฺพาราม\-สุตฺต}
\proseitem{516}{ตํเยว นิทานํ. กตินํ นุ โข ภิกฺขเว อินฺทฺริยานํ ภาวิตตฺตา พหุลีกตตฺตา ขีณาสโว ภิกฺขุ อญฺญํ พฺยากโรติ “ขีณา ชาติ, วุสิตํ พฺรหฺมจริยํ, กตํ กรณียํ, นาปรํ อิตฺถตฺตายาติ ปชานามี”ติ. ภควํมูลกา โน ภนฺเต ธมฺมา ฯเปฯ. ทฺวินฺนํ โข ภิกฺขเว อินฺทฺริยานํ ภาวิตตฺตา พหุลีกตตฺตา ขีณาสโว ภิกฺขุ อญฺญํ พฺยากโรติ “ขีณา ชาติ, วุสิตํ พฺรหฺมจริยํ, กตํ กรณียํ, นาปรํ อิตฺถตฺตายาติ ปชานามี”ติ. กตเมสํ ทฺวินฺนํ, อริยาย จ ปญฺญาย อริยาย จ วิมุตฺติยา. ยา หิสฺส ภิกฺขเว อริยา ปญฺญา, ตทสฺส ปญฺญินฺทฺริยํ. ยา หิสฺส ภิกฺขเว อริยา วิมุตฺติ, ตทสฺส สมาธินฺทฺริยํ. อิเมสํ โข ภิกฺขเว ทฺวินฺนํ อินฺทฺริยานํ ภาวิตตฺตา พหุลีกตตฺตา ขีณาสโว ภิกฺขุ อญฺญํ พฺยากโรติ “ขีณา ชาติ, \csromanfolio{196}วุสิตํ พฺรหฺมจริยํ, กตํ กรณียํ, นาปรํ อิตฺถตฺตายาติ ปชานามี”ติ.  ฉฏฺฐํ.}

\csromanheader{2}{7. ตติย\-ปุพฺพาราม\-สุตฺต}
\proseitem{517}{ตํเยว นิทานํ. กตินํ นุ โข ภิกฺขเว อินฺทฺริยานํ ภาวิตตฺตา พหุลีกตตฺตา ขีณาสโว ภิกฺขุ อญฺญํ พฺยากโรติ “ขีณา ชาติ, วุสิตํ พฺรหฺมจริยํ, กตํ กรณียํ, นาปรํ อิตฺถตฺตายาติ ปชานามี”ติ. ภควํมูลกา โน ภนฺเต ธมฺมา ฯเปฯ. จตุนฺนํ โข ภิกฺขเว อินฺทฺริยานํ ภาวิตตฺตา พหุลีกตตฺตา ขีณาสโว ภิกฺขุ อญฺญํ พฺยากโรติ “ขีณา ชาติ, วุสิตํ พฺรหฺมจริยํ, กตํ กรณียํ, นาปรํ อิตฺถตฺตายาติ ปชานามี”ติ. กตเมสํ จตุนฺนํ, วีริยิ\-นฺทฺริยสฺส สตินฺทฺริยสฺส สมาธิ\-นฺทฺริยสฺส ปญฺญินฺทฺริยสฺส. อิเมสํ โข ภิกฺขเว จตุนฺนํ อินฺทฺริยานํ ภาวิตตฺตา พหุลีกตตฺตา ขีณาสโว ภิกฺขุ อญฺญํ พฺยากโรติ “ขีณา ชาติ, วุสิตํ พฺรหฺมจริยํ, กตํ กรณียํ, นาปรํ อิตฺถตฺตายาติ ปชานามี”ติ.  สตฺตมํ.}

\csromanheader{2}{8. จตุตฺถ\-ปุพฺพาราม\-สุตฺต}
\proseitem{518}{ตํเยว นิทานํ. กตินํ นุ โข ภิกฺขเว อินฺทฺริยานํ ภาวิตตฺตา พหุลีกตตฺตา ขีณาสโว ภิกฺขุ อญฺญํ พฺยากโรติ “ขีณา ชาติ, วุสิตํ พฺรหฺมจริยํ, กตํ กรณียํ, นาปรํ อิตฺถตฺตายาติ ปชานามี”ติ. ภควํมูลกา โน ภนฺเต ธมฺมา ฯเปฯ. ปญฺจนฺนํ โข ภิกฺขเว อินฺทฺริยานํ ภาวิตตฺตา พหุลีกตตฺตา ขีณาสโว ภิกฺขุ อญฺญํ พฺยากโรติ “ขีณา ชาติ, วุสิตํ พฺรหฺมจริยํ, กตํ กรณียํ, นาปรํ อิตฺถตฺตายาติ ปชานามี”ติ. กตเมสํ ปญฺจนฺนํ, สทฺธิ\-นฺทฺริยสฺส วีริยิ\-นฺทฺริยสฺส สตินฺทฺริยสฺส สมาธิ\-นฺทฺริยสฺส ปญฺญินฺทฺริยสฺส. อิเมสํ โข ภิกฺขเว ปญฺจนฺนํ อินฺทฺริยานํ ภาวิตตฺตา พหุลีกตตฺตา ขีณาสโว ภิกฺขุ อญฺญํ พฺยากโรติ “ขีณา ชาติ, วุสิตํ พฺรหฺมจริยํ, กตํ กรณียํ, นาปรํ อิตฺถตฺตายาติ ปชานามี”ติ.  อฏฺฐมํ.}

\csromanmatikaanchor{mtk.123}
\csromantocmarkref{subparagraph}{9. ปิณฺโฑลภารทฺวาชสุตฺต}{mtk.123}
\bookmark[dest={mtk.123},level=5]{9. ปิณฺโฑลภารทฺวาชสุตฺต}
\csromanheader{6}{9. ปิณฺโฑลภารทฺวาช\-สุตฺต}
\proseitem{519}{เอวํ เม สุตํ\csromandash{}เอกํ สมยํ ภควา โกสมฺพิยํ วิหรติ โฆสิตาราเม. เตน โข ปน สมเยน อายสฺมตา ปิณฺโฑล\-ภารทฺวาเชน อญฺญา พฺยากตา โหติ “ขีณา ชาติ, วุสิตํ พฺรหฺมจริยํ, กตํ}

\prose{\csromanfolio{197}กรณียํ, นาปรํ อิตฺถตฺตายาติ ปชานามี”ติ. อถ โข สมฺพหุลา ภิกฺขู เยน ภควา เตนุ\-ปสงฺกมิํสุ, อุปสงฺกมิตฺวา ภควนฺตํ อภิวาเทตฺวา เอกมนฺตํ นิสีทิํสุ, เอกมนฺตํ นิสินฺนา โข เต ภิกฺขู ภควนฺตํ เอตทโวจุํ\csromandash{}}

\prose{อายสฺมตา ภนฺเต ปิณฺโฑล\-ภารทฺวาเชน อญฺญา พฺยากตา “ขีณา ชาติ, วุสิตํ พฺรหฺมจริยํ, กตํ กรณียํ, นาปรํ อิตฺถตฺตายาติ ปชานามี”ติ. กิํ นุ โข ภนฺเต อตฺถวสํ สมฺปสฺสมาเนน อายสฺมตา ปิณฺโฑล\-ภารทฺวาเชน อญฺญา พฺยากตา “ขีณา ชาติ, วุสิตํ พฺรหฺมจริยํ, กตํ กรณียํ, นาปรํ อิตฺถตฺตายาติ ปชานามี”ติ.}

\prose{ติณฺณนฺนํ โข ภิกฺขเว อินฺทฺริยานํ ภาวิตตฺตา พหุลีกตตฺตา ปิณฺโฑล\-ภารทฺวาเชน ภิกฺขุนา อญฺญา พฺยากตา “ขีณา ชาติ, วุสิตํ พฺรหฺมจริยํ, กตํ กรณียํ, นาปรํ อิตฺถตฺตายาติ ปชานามี”ติ. กตเมสํ ติณฺณนฺนํ, สตินฺทฺริยสฺส สมาธิ\-นฺทฺริยสฺส ปญฺญินฺทฺริยสฺส. อิเมสํ โข ภิกฺขเว ติณฺณนฺนํ อินฺทฺริยานํ ภาวิตตฺตา พหุลีกตตฺตา ปิณฺโฑล\-ภารทฺวาเชน ภิกฺขุนา อญฺญา พฺยากตา “ขีณา ชาติ, วุสิตํ พฺรหฺมจริยํ, กตํ กรณียํ, นาปรํ อิตฺถตฺตายาติ ปชานามี”ติ. อิมานิ จ ภิกฺขเว ตีณินฺทฺริยานิ กิมนฺตานิ, ขยนฺตานิ. กิสฺส ขยนฺตานิ, ชาติ\-ชรา\-มรณสฺส, “ชาติ\-ชรา\-มรณํ ขยนฺ”ติ โข ภิกฺขเว สมฺปสฺสมาเนน ปิณฺโฑล\-ภารทฺวาเชน ภิกฺขุนา อญฺญา พฺยากตา “ขีณา ชาติ, วุสิตํ พฺรหฺมจริยํ, กตํ กรณียํ, นาปรํ อิตฺถตฺตายาติ ปชานามี”ติ.  นวมํ.}

\csromanmatikaanchor{mtk.124}
\csromantocmarkref{subparagraph}{10. อาปณสุตฺต}{mtk.124}
\bookmark[dest={mtk.124},level=5]{10. อาปณสุตฺต}
\csromanheader{6}{10. \textbf{อาปณสุตฺต}}
\proseitem{520}{เอวํ เม สุตํ\csromandash{}เอกํ สมยํ ภควา องฺเคสุ วิหรติ อาปณํ นาม องฺคานํ นิคโม. ตตฺร โข ภควา อายสฺมนฺตํ สาริปุตฺตํ อามนฺเตสิ “โย โส สาริปุตฺต อริยสาวโก ตถาคเต เอกนฺตคโต\csromansharedfootnote{60393ca57bec9e9f}{เอกนฺติคโต (สี)} อภิปฺปสนฺโน, น โส ตถาคเต วา ตถาคต\-สาสเน วา กงฺเขยฺย วา วิจิกิจฺเฉยฺย วา”ติ.}

\prose{โย โส ภนฺเต อริยสาวโก ตถาคเต เอกนฺตคโต อภิปฺปสนฺโน, น โส ตถาคเต วา ตถาคต\-สาสเน วา กงฺเขยฺย วา วิจิกิจฺเฉยฺย วา. สทฺธสฺส หิ ภนฺเต อริยสาวกสฺส เอตํ ปาฏิกงฺขํ \csromanfolio{198}ยํ อารทฺธวีริโย วิหริสฺสติ อกุสลานํ ธมฺมานํ ปหานาย กุสลานํ ธมฺมานํ อุปสมฺปทาย ถามวา ทฬฺหปรกฺกโม อนิกฺขิตฺตธุโร กุสเลสุ ธมฺเมสุ. ยํ หิสฺส ภนฺเต วีริยํ, ตทสฺส วีริยินฺทฺริยํ.}

\prose{สทฺธสฺส หิ ภนฺเต อริยสาวกสฺส อารทฺธ\-วีริยสฺส เอตํ ปาฏิกงฺขํ ยํ สติมา ภวิสฺสติ ปรเมน สติเนปกฺเกน สมนฺนาคโต จิรกตมฺปิ จิรภาสิตมฺปิ สริตา อนุสฺสริตา. ยา หิสฺส ภนฺเต สติ, ตทสฺส สตินฺทฺริยํ.}

\prose{สทฺธสฺส หิ ภนฺเต อริยสาวกสฺส อารทฺธ\-วีริยสฺส อุปฏฺฐิต\-สฺสติโน เอตํ ปาฏิกงฺขํ ยํ โวสฺสคฺคา\-รมฺมณํ กริตฺวา ลภิสฺสติ สมาธิํ, ลภิสฺสติ จิตฺตสฺส เอกคฺคตํ. โย หิสฺส ภนฺเต สมาธิ, ตทสฺส สมาธินฺทฺริยํ.}

\prose{สทฺธสฺส หิ ภนฺเต อริยสาวกสฺส อารทฺธ\-วีริยสฺส อุปฏฺฐิต\-สฺสติโน สมาหิต\-จิตฺตสฺส เอตํ ปาฏิกงฺขํ ยํ เอวํ ปชานิสฺสติ “อนมตคฺโค โข สํสาโร, ปุพฺพโกฏิ น ปญฺญายติ อวิชฺชา\-นีวรณานํ สตฺตานํ ตณฺหา\-สํโยชนานํ สนฺธาวตํ สํสรตํ, อวิชฺชาย เตฺวว ตโมกายสฺส อเสส\-วิราค\-นิโรโธ สนฺตเมตํ ปทํ, ปณีตเมตํ ปทํ, ยทิทํ สพฺพ\-สงฺขาร\-สมโถ สพฺพู\-ปธิ\-ปฏินิสฺสคฺโค ตณฺหากฺขโย วิราโค นิโรโธ นิพฺพานํ”\csromansharedfootnote{a23eab1915399b26}{นิพฺพานนฺติ (?)}. ยา หิสฺส ภนฺเต ปญฺญา, ตทสฺส ปญฺญินฺทฺริยํ.}

\prose{สทฺโธ โส\csromansharedfootnote{22dc53cf91daf7e7}{ส โข โส (สี, สฺยา, กํ)} ภนฺเต อริยสาวโก เอวํ ปทหิตฺวา ปทหิตฺวา เอวํ สริตฺวา สริตฺวา เอวํ สมาทหิตฺวา สมาทหิตฺวา เอวํ ปชานิตฺวา ปชานิตฺวา เอวํ อภิสทฺทหติ “อิเม โข เต ธมฺมา, เย เม ปุพฺเพ สุตวา อเหสุํ, เตนาหํ เอตรหิ กาเยน จ ผุสิตฺวา วิหรามิ, ปญฺญาย จ อติวิชฺฌ\csromansharedfootnote{83c48de37c50b8ec}{ปฏิวิชฺฌ (สี, ก) ตทฏฺฐกถาสุ ปน อติวิชฺฌิตฺวาติ วณฺณิตํ.} ปสฺสามี”ติ. ยา หิสฺส ภนฺเต สทฺธา, ตทสฺส สทฺธิ\-นฺทฺริยนฺติ.}

\prose{สาธุ สาธุ สาริปุตฺต, โย โส สาริปุตฺต อริยสาวโก ตถาคเต เอกนฺตคโต อภิปฺปสนฺโน, น โส ตถาคเต วา ตถาคต\-สาสเน วา กงฺเขยฺย วา วิจิกิจฺเฉยฺย วา. สทฺธสฺส หิ สาริปุตฺต อริยสาวกสฺส เอตํ ปาฏิกงฺขํ ยํ อารทฺธวีริโย วิหริสฺสติ}

\prose{\csromanfolio{199}อกุสลานํ ธมฺมานํ ปหานาย กุสลานํ ธมฺมานํ อุปสมฺปทาย ถามวา ทฬฺหปรกฺกโม อนิกฺขิตฺตธุโร กุสเลสุ ธมฺเมสุ. ยํ หิสฺส สาริปุตฺต วีริยํ, ตทสฺส วีริยินฺทฺริยํ.}

\prose{สทฺธสฺส หิ สาริปุตฺต อริยสาวกสฺส อารทฺธ\-วีริยสฺส เอตํ ปาฏิกงฺขํ ยํ สติมา ภวิสฺสติ ปรเมน สติเนปกฺเกน สมนฺนาคโต จิรกตมฺปิ จิรภาสิตมฺปิ สริตา อนุสฺสริตา. ยา หิสฺส สาริปุตฺต สติ, ตทสฺส สตินฺทฺริยํ.}

\prose{สทฺธสฺส หิ สาริปุตฺต อริยสาวกสฺส อารทฺธ\-วีริยสฺส อุปฏฺฐิต\-สฺสติโน เอตํ ปาฏิกงฺขํ ยํ โวสฺสคฺคา\-รมฺมณํ กริตฺวา ลภิสฺสติ สมาธิํ, ลภิสฺสติ จิตฺตสฺส เอกคฺคตํ. โย หิสฺส สาริปุตฺต สมาธิ, ตทสฺส สมาธินฺทฺริยํ.}

\prose{สทฺธสฺส หิ สาริปุตฺต อริยสาวกสฺส อารทฺธ\-วีริยสฺส อุปฏฺฐิต\-สฺสติโน สมาหิต\-จิตฺตสฺส เอตํ ปาฏิกงฺขํ ยํ เอวํ ปชานิสฺสติ “อนมตคฺโค โข สํสาโร, ปุพฺพโกฏิ น ปญฺญายติ อวิชฺชา\-นีวรณานํ สตฺตานํ ตณฺหา\-สํโยชนานํ สนฺธาวตํ สํสรตํ, อวิชฺชาย เตฺวว ตโมกายสฺส อเสส\-วิราค\-นิโรโธ สนฺตเมตํ ปทํ, ปณีตเมตํ ปทํ, ยทิทํ สพฺพ\-สงฺขาร\-สมโถ สพฺพู\-ปธิ\-ปฏินิสฺสคฺโค ตณฺหากฺขโย วิราโค นิโรโธ นิพฺพานํ”. ยา หิสฺส สาริปุตฺต ปญฺญา, ตทสฺส ปญฺญินฺทฺริยํ.}

\prosewithclosermidruled{สทฺโธ โส\csromansharedfootnote{88f0f3006a76bd88}{ส โข โส (สี, สฺยา, กํ, อิ)} สาริปุตฺต อริยสาวโก เอวํ ปทหิตฺวา ปทหิตฺวา เอวํ สริตฺวา สริตฺวา เอวํ สมาทหิตฺวา สมาทหิตฺวา เอวํ ปชานิตฺวา ปชานิตฺวา เอวํ อภิสทฺทหติ “อิเม โข เต ธมฺมา, เย เม ปุพฺเพ สุตวา อเหสุํ, เตนาหํ เอตรหิ กาเยน จ ผุสิตฺวา วิหรามิ, ปญฺญาย จ อติวิชฺฌ\csromansharedfootnote{da23fc3175dc4f14}{ปฏิวิชฺฌ (กสี, ก)} ปสฺสามี”ติ. ยา หิสฺส สาริปุตฺต สทฺธา, ตทสฺส สทฺธิ\-นฺทฺริยนฺติ.  ทสมํ.}{ชราวคฺโค ปญฺจโม.}
\csromancenter{\textbf{ตสฺสุทฺทานํ}}
\setlength{\csromangathapagewidth}{0pt}
\setlength{\csromangathawakwidth}{0pt}
\csromangathameasuregroup{ชรา อุณฺณาโภ พฺราหฺมโณ, \\ ปุพฺพาราเม จ จตฺตาริ,}{\csromangathabat{ชรา อุณฺณาโภ พฺราหฺมโณ,}{สาเกโต ปุพฺพโกฏฺฐโก.} \\ \csromangathabat{ปุพฺพาราเม จ จตฺตาริ,}{ปิณฺโฑโล อาปเณน จาติ\textsuperscript{1}.}}
\csromangathasetleft{ชรา อุณฺณาโภ พฺราหฺมโณ, \\ ปุพฺพาราเม จ จตฺตาริ,}
\csromangathagroup{\csromangathastanza{\csromangathabat{ชรา อุณฺณาโภ พฺราหฺมโณ,}{สาเกโต ปุพฺพโกฏฺฐโก.} \\ \csromangathabat{ปุพฺพาราเม จ จตฺตาริ,}{ปิณฺโฑโล อาปเณน จาติ\csromansharedfootnote{cf5a8fa8849f5b37}{สทฺเธน เต ทสาติ (สฺยา, กํ, ก)}.}}}

\csromanmatikaanchor{mtk.125}
\csromantocmarknopage{subparagraph}{6. สูกรขตวคฺค}{mtk.125}
\bookmark[dest={mtk.125},level=5]{6. สูกรขตวคฺค}
\csromanheader{6}{6. สูกรขต\-วคฺค}
\csromanmatikaanchor{mtk.126}
\csromantocmarkref{subparagraph}{1. สาลสุตฺต}{mtk.126}
\bookmark[dest={mtk.126},level=5]{1. สาลสุตฺต}
\csromanheader{6}{1. \textbf{สาลสุตฺต}}
\proseitem{521}{\csromanfolio{200}เอวํ เม สุตํ\csromandash{}เอกํ สมยํ ภควา โกสเลสุ วิหรติ สาลาย พฺราหฺมณคาเม. ตตฺร โข ภควา ภิกฺขู อามนฺเตสิ\csromandash{}เสยฺยถาปิ ภิกฺขเว เย เกจิ ติรจฺฉานคตา ปาณา, สีโห มิคราชา เตสํ อคฺคมกฺขายติ ยทิทํ ถาเมน ชเวน สูเรน\csromansharedfootnote{2b269e66527cb533}{สูริเยน (สี, สฺยา, กํ)}. เอวเมว โข ภิกฺขเว เย เกจิ โพธิปกฺขิยา ธมฺมา, ปญฺญินฺทฺริยํ เตสํ อคฺคมกฺขายติ ยทิทํ โพธาย.}

\prose{กตเม จ ภิกฺขเว โพธิปกฺขิยา ธมฺมา, สทฺธินฺทฺริยํ ภิกฺขเว โพธิปกฺขิโย ธมฺโม, ตํ โพธาย สํวตฺตติ. วีริยินฺทฺริยํ โพธิปกฺขิโย ธมฺโม, ตํ โพธาย สํวตฺตติ. สตินฺทฺริยํ โพธิปกฺขิโย ธมฺโม, ตํ โพธาย สํวตฺตติ. สมาธินฺทฺริยํ โพธิปกฺขิโย ธมฺโม, ตํ โพธาย สํวตฺตติ, ปญฺญินฺทฺริยํ โพธิปกฺขิโย ธมฺโม, ตํ โพธาย สํวตฺตติ. เสยฺยถาปิ ภิกฺขเว เย เกจิ ติรจฺฉานคตา ปาณา, สีโห มิคราชา เตสํ อคฺคมกฺขายติ ยทิทํ ถาเมน ชเวน สูเรน. เอวเมว โข ภิกฺขเว เย เกจิ โพธิปกฺขิยา ธมฺมา, ปญฺญินฺทฺริยํ เตสํ อคฺคมกฺขายติ ยทิทํ โพธายาติ.  ปฐมํ.}

\csromanmatikaanchor{mtk.127}
\csromantocmarkref{subparagraph}{2. มลฺลิกสุตฺต}{mtk.127}
\bookmark[dest={mtk.127},level=5]{2. มลฺลิกสุตฺต}
\csromanheader{6}{2. \textbf{มลฺลิกสุตฺต}}
\proseitem{522}{เอวํ เม สุตํ\csromandash{}เอกํ สมยํ ภควา มลฺเลสุ\csromansharedfootnote{ffb43debbc49766a}{มลฺลเกสุ (สี, สฺยา, กํ), มลฺลิเกสุ (ก)} วิหรติ อุรุเวลกปฺปํ นาม มลฺลานํ นิคโม. ตตฺร โข ภควา ภิกฺขู อามนฺเตสิ\csromandash{}ยาวกีวญฺจ ภิกฺขเว อริยสาวกสฺส อริยญาณํ น อุปฺปนฺนํ โหติ, เนว ตาว จตุนฺนํ อินฺทฺริยานํ สณฺฐิติ โหติ, เนว ตาว จตุนฺนํ อินฺทฺริยานํ อวฏฺฐิติ โหติ. ยโต จ โข ภิกฺขเว อริยสาวกสฺส อริยญาณํ อุปฺปนฺนํ โหติ, อถ จตุนฺนํ อินฺทฺริยานํ สณฺฐิติ โหติ, อถ จตุนฺนํ อินฺทฺริยานํ อวฏฺฐิติ โหติ.}

\prose{เสยฺยถาปิ ภิกฺขเว ยาวกีวญฺจ กูฏาคารสฺส กูฏํ น อุสฺสิตํ โหติ, เนว ตาว โคปานสีนํ สณฺฐิติ โหติ, เนว ตาว โคปานสีนํ อวฏฺฐิติ โหติ. ยโต จ โข ภิกฺขเว กูฏาคารสฺส กูฏํ}

\prose{\csromanfolio{201}อุสฺสิตํ โหติ, อถ โคปานสีนํ สณฺฐิติ โหติ, อถ โคปานสีนํ อวฏฺฐิติ โหติ. เอวเมว โข ภิกฺขเว ยาวกีวญฺจ อริยสาวกสฺส อริยญาณํ น อุปฺปนฺนํ โหติ, เนว ตาว จตุนฺนํ อินฺทฺริยานํ สณฺฐิติ โหติ, เนว ตาว จตุนฺนํ อินฺทฺริยานํ อวฏฺฐิติ โหติ. ยโต จ โข ภิกฺขเว อริยสาวกสฺส อริยญาณํ อุปฺปนฺนํ โหติ, อถ จตุนฺนํ อินฺทฺริยานํ ฯเปฯ อวฏฺฐิติ โหติ.}

\prose{กตเมสํ จตุนฺนํ, สทฺธิ\-นฺทฺริยสฺส วีริยิ\-นฺทฺริยสฺส สตินฺทฺริยสฺส สมาธิ\-นฺทฺริยสฺส. ปญฺญวโต ภิกฺขเว อริยสาวกสฺส ตทนฺวยา สทฺธา สณฺฐาติ, ตทนฺวยํ วีริยํ สณฺฐาติ, ตทนฺวยา สติ สณฺฐาติ, ตทนฺวโย สมาธิ สณฺฐาตีติ.  ทุติยํ.}

\csromanmatikaanchor{mtk.128}
\csromantocmarkref{subparagraph}{3. เสขสุตฺต}{mtk.128}
\bookmark[dest={mtk.128},level=5]{3. เสขสุตฺต}
\csromanheader{6}{3. \textbf{เสขสุตฺต}}
\proseitem{523}{เอวํ เม สุตํ\csromandash{}เอกํ สมยํ ภควา โกสมฺพิยํ วิหรติ โฆสิตาราเม. ตตฺร โข ภควา ภิกฺขู อามนฺเตสิ\csromandash{}อตฺถิ นุ โข ภิกฺขเว ปริยาโย, ยํ ปริยายํ อาคมฺม เสโข ภิกฺขุ เสขภูมิยํ ฐิโต “เสโขสฺมี”ติ ปชาเนยฺย, อเสโข ภิกฺขุ อเสขภูมิยํ ฐิโต “อเสโขสฺมี”ติ ปชาเนยฺยาติ.}

\prose{ภควํมูลกา โน ภนฺเต ธมฺมา ฯเปฯ. อตฺถิ ภิกฺขเว ปริยาโย, ยํ ปริยายํ อาคมฺม เสโข ภิกฺขุ เสขภูมิยํ ฐิโต “เสโขสฺมี”ติ ปชาเนยฺย, อเสโข ภิกฺขุ อเสขภูมิยํ ฐิโต “อเสโขสฺมี”ติ ปชาเนยฺย.}

\prose{กตโม จ ภิกฺขเว ปริยาโย, ยํ ปริยายํ อาคมฺม เสโข ภิกฺขุ เสขภูมิยํ ฐิโต “เสโขสฺมี”ติ ปชานาติ, อิธ ภิกฺขเว เสโข ภิกฺขุ “อิทํ ทุกฺขนฺ”ติ ยถาภูตํ ปชานาติ, “อยํ ทุกฺขสมุทโย”ติ ยถาภูตํ ปชานาติ, “อยํ ทุกฺขนิโรโธ”ติ ยถาภูตํ ปชานาติ, “อยํ ทุกฺขนิโรธ\-คามินี ปฏิปทา”ติ ยถาภูตํ ปชานาติ. อยมฺปิ โข ภิกฺขเว ปริยาโย, ยํ ปริยายํ อาคมฺม เสโข ภิกฺขุ เสขภูมิยํ ฐิโต “เสโขสฺมี”ติ ปชานาติ.}

\csromanmatikaanchor{mtk.129}
\csromantocmarkref{subparagraph}{4-6. ปทสุตฺตาทีนิ}{mtk.129}
\bookmark[dest={mtk.129},level=5]{4-6. ปทสุตฺตาทีนิ}
\prose{ปุน จปรํ ภิกฺขเว เสโข ภิกฺขุ อิติ ปฏิสญฺจิกฺขติ “อตฺถิ นุ โข อิโต พหิทฺธา อญฺโญ สมโณ วา พฺราหฺมโณ วา, โย เอวํ ภูตํ ตจฺฉํ ตถํ \csromanfolio{202}ธมฺมํ เทเสติ ยถา ภควา”ติ. โส เอวํ ปชานาติ “นตฺถิ โข อิโต พหิทฺธา อญฺโญ สมโณ วา พฺราหฺมโณ วา, โย เอวํ ภูตํ ตจฺฉํ ตถํ ธมฺมํ เทเสติ ยถา ภควา”ติ. อยมฺปิ โข ภิกฺขเว ปริยาโย, ยํ ปริยายํ อาคมฺม เสโข ภิกฺขุ เสขภูมิยํ ฐิโต “เสโขสฺมี”ติ ปชานาติ.}

\prose{ปุน จปรํ ภิกฺขเว เสโข ภิกฺขุ ปญฺจินฺทฺริยานิ ปชานาติ, สทฺธินฺทฺริยํ วีริยินฺทฺริยํ สตินฺทฺริยํ สมาธินฺทฺริยํ ปญฺญินฺทฺริยํ. ยํคติกานิ ยํปรมานิ ยํผลานิ ยํปริโยสานานิ น เหว โข กาเยน ผุสิตฺวา วิหรติ, ปญฺญาย จ อติวิชฺฌ ปสฺสติ. อยมฺปิ โข ภิกฺขเว ปริยาโย, ยํ ปริยายํ อาคมฺม เสโข ภิกฺขุ เสขภูมิยํ ฐิโต “เสโขสฺมี”ติ ปชานาติ.}

\prose{กตโม จ ภิกฺขเว ปริยาโย, ยํ ปริยายํ อาคมฺม อเสโข ภิกฺขุ อเสขภูมิยํ “ฐิโต อเสโขสฺมี”ติ ปชานาติ, อิธ ภิกฺขเว อเสโข ภิกฺขุ ปญฺจินฺทฺริยานิ ปชานาติ สทฺธินฺทฺริยํ วีริยินฺทฺริยํ สตินฺทฺริยํ สมาธินฺทฺริยํ ปญฺญินฺทฺริยํ. ยํคติกานิ ยํปรมานิ ยํผลานิ ยํปริโยสานานิ กาเยน จ ผุสิตฺวา วิหรติ, ปญฺญาย จ อติวิชฺฌ ปสฺสติ. อยมฺปิ โข ภิกฺขเว ปริยาโย, ยํ ปริยายํ อาคมฺม อเสโข ภิกฺขุ อเสขภูมิยํ ฐิโต “อเสโขสฺมี”ติ ปชานาติ.}

\prose{ปุน จปรํ ภิกฺขเว อเสโข ภิกฺขุ ฉ อินฺทฺริยานิ ปชานาติ, จกฺขุนฺทฺริยํ โสตินฺทฺริยํ ฆานินฺทฺริยํ ชิวฺหินฺทฺริยํ กายินฺทฺริยํ มนินฺทฺริยํ. อิมานิ โข ฉ อินฺทฺริยานิ สพฺเพน สพฺพํ สพฺพถา สพฺพํ อปริเสสํ นิรุชฺฌิสฺสนฺติ, อญฺญานิ จ ฉ อินฺทฺริยานิ “น กุหิญฺจิ กิสฺมิญฺจิ อุปฺปชฺชิสฺสนฺตี”ติ ปชานาติ. อยมฺปิ โข ภิกฺขเว ปริยาโย, ยํ ปริยายํ อาคมฺม อเสโข ภิกฺขุ อเสขภูมิยํ ฐิโต “อเสโขสฺมี”ติ ปชานาตีติ.  ตติยํ.}

\csromanheader{2}{4. \textbf{ปทสุตฺต}}
\proseitem{524}{เสยฺยถาปิ ภิกฺขเว ยานิ กานิจิ ชงฺคลานํ\csromansharedfootnote{5f11cdf58c1c21ba}{ชงฺคมานํ (สี, อิ)} ปาณานํ ปทชาตานิ, สพฺพานิ ตานิ หตฺถิปเท สโมธานํ คจฺฉนฺติ, หตฺถิปทํ เตสํ อคฺคมกฺขายติ ยทิทํ มหนฺตตฺเตน. เอวเมว โข ภิกฺขเว ยานิ กานิจิ ปทานิ โพธาย}

\prose{\csromanfolio{203}สํวตฺตนฺติ, ปญฺญินฺทฺริยํ ปทํ เตสํ อคฺคมกฺขายติ ยทิทํ โพธาย. กตมานิ จ ภิกฺขเว ปทานิ โพธาย สํวตฺตนฺติ, สทฺธินฺทฺริยํ ภิกฺขเว ปทํ, ตํ โพธาย สํวตฺตติ. วีริยินฺทฺริยํ ปทํ, ตํ โพธาย สํวตฺตติ. สตินฺทฺริยํ ปทํ, ตํ โพธาย สํวตฺตติ. สมาธินฺทฺริยํ ปทํ, ตํ โพธาย สํวตฺตติ. ปญฺญินฺทฺริยํ ปทํ, ตํ โพธาย สํวตฺตติ. เสยฺยถาปิ ภิกฺขเว ยานิ กานิจิ ชงฺคลานํ ปาณานํ ปทชาตานิ, สพฺพานิ ตานิ หตฺถิปเท สโมธานํ คจฺฉนฺติ, หตฺถิปทํ เตสํ อคฺคมกฺขายติ ยทิทํ มหนฺตตฺเตน. เอวเมว โข ภิกฺขเว ยานิ กานิจิ ปทานิ โพธาย สํวตฺตนฺติ, ปญฺญินฺทฺริยํ ปทํ เตสํ อคฺคมกฺขายติ ยทิทํ โพธายาติ.  จตุตฺถํ.}

\csromanheader{2}{5. \textbf{สารสุตฺต}}
\proseitem{525}{เสยฺยถาปิ ภิกฺขเว เย เกจิ สารคนฺธา, โลหิตจนฺทนํ เตสํ อคฺคมกฺขายติ. เอวเมว โข ภิกฺขเว เย เกจิ โพธิปกฺขิยา ธมฺมา, ปญฺญินฺทฺริยํ เตสํ อคฺคมกฺขายติ ยทิทํ โพธาย. กตเม จ ภิกฺขเว โพธิปกฺขิยา ธมฺมา, สทฺธินฺทฺริยํ ภิกฺขเว โพธิปกฺขิโย ธมฺโม, ตํ โพธาย สํวตฺตติ. วีริยินฺทฺริยํ ฯเปฯ. สตินฺทฺริยํ ฯเปฯ. สมาธินฺทฺริยํ ฯเปฯ. ปญฺญินฺทฺริยํ โพธิปกฺขิโย ธมฺโม, ตํ โพธาย สํวตฺตติ. เสยฺยถาปิ ภิกฺขเว เย เกจิ สารคนฺธา, โลหิตจนฺทนํ เตสํ อคฺคมกฺขายติ. เอวเมว โข ภิกฺขเว เย เกจิ โพธิปกฺขิยา ธมฺมา, ปญฺญินฺทฺริยํ เตสํ อคฺคมกฺขายติ ยทิทํ โพธายาติ.  ปญฺจมํ.}

\csromanheader{2}{6. ปติฏฺฐิต\-สุตฺต}
\proseitem{526}{เอกธมฺเม ปติฏฺฐิตสฺส ภิกฺขเว ภิกฺขุโน ปญฺจินฺทฺริยานิ ภาวิตานิ โหนฺติ สุภาวิตานิ. กตมสฺมิํ เอกธมฺเม อปฺปมาเท. กตโม จ ภิกฺขเว อปฺปมาโท, อิธ ภิกฺขเว ภิกฺขุ จิตฺตํ รกฺขติ อาสเวสุ จ สาสเวสุ จ ธมฺเมสุ, ตสฺส จิตฺตํ รกฺขโต อาสเวสุ จ สาสเวสุ จ ธมฺเมสุ สทฺธิ\-นฺทฺริยมฺปิ ภาวนา\-ปาริปูริํ คจฺฉติ, วีริยิ\-นฺทฺริยมฺปิ ภาวนา\-ปาริปูริํ คจฺฉติ, สตินฺทฺริยมฺปิ ภาวนา\-ปาริปูริํ คจฺฉติ, สมาธิ\-นฺทฺริยมฺปิ ภาวนา\-ปาริปูริํ คจฺฉติ, ปญฺญินฺทฺริยมฺปิ ภาวนา\-ปาริปูริํ คจฺฉติ. เอวมฺปิ โข ภิกฺขเว เอกธมฺเม ปติฏฺฐิตสฺส ภิกฺขุโน ปญฺจินฺทฺริยานิ ภาวิตานิ โหนฺติ สุภาวิตานีติ.  ฉฏฺฐํ.}

\csromanmatikaanchor{mtk.130}
\csromantocmarkref{subparagraph}{7. สหมฺปติพฺรหฺมสุตฺต}{mtk.130}
\bookmark[dest={mtk.130},level=5]{7. สหมฺปติพฺรหฺมสุตฺต}
\csromanheader{6}{7. สหมฺปติพฺรหฺม\-สุตฺต}
\proseitem{527}{\csromanfolio{204}เอกํ สมยํ ภควา อุรุเวลายํ วิหรติ นชฺชา เนรญฺชราย ตีเร อชปาล\-นิโคฺรเธ ปฐมา\-ภิสมฺพุทฺโธ. อถ โข ภควโต รโหคตสฺส ปฏิสลฺลีนสฺส เอวํ เจตโส ปริวิตกฺโก อุทปาทิ “ปญฺจินฺทฺริยานิ ภาวิตานิ พหุลีกตานิ อมโตคธานิ โหนฺติ อมต\-ปรายณานิ อมต\-ปริโยสานานิ. กตมานิ ปญฺจ, สทฺธินฺทฺริยํ ภาวิตํ พหุลีกตํ อมโตคธํ โหติ อมตปรายณํ อมต\-ปริโยสานํ. วีริยินฺทฺริยํ ฯเปฯ. สตินฺทฺริยํ ฯเปฯ. สมาธินฺทฺริยํ ฯเปฯ. ปญฺญินฺทฺริยํ ภาวิตํ พหุลีกตํ อมโตคธํ โหติ อมตปรายณํ อมต\-ปริโยสานํ. อิมานิ ปญฺจินฺทฺริยานิ ภาวิตานิ พหุลีกตานิ อมโตคธานิ โหนฺติ อมต\-ปรายณานิ อมตปริโยสานานี”ติ.}

\prose{อถ โข พฺรหฺมา สหมฺปติ ภควโต เจตสา เจโต\-ปริวิตกฺกม\-ญฺญาย เสยฺยถาปิ นาม พลวา ปุริโส สมิญฺชิตํ วา พาหํ ปสาเรยฺย, ปสาริตํ วา พาหํ สมิญฺเชยฺย, เอวเมว พฺรหฺมโลเก อนฺตรหิโต ภควโต ปุรโต ปาตุรโหสิ. อถ โข พฺรหฺมา สหมฺปติ เอกํสํ อุตฺตราสงฺคํ กริตฺวา เยน ภควา เตนญฺชลิํ ปณาเมตฺวา ภควนฺตํ เอตทโวจ\csromandash{}เอวเมตํ ภควา, เอวเมตํ สุคต, ปญฺจินฺทฺริยานิ ภาวิตานิ พหุลีกตานิ อมโตคธานิ โหนฺติ อมต\-ปรายณานิ อมต\-ปริโยสานานิ. กตมานิ ปญฺจ, สทฺธินฺทฺริยํ ภาวิตํ พหุลีกตํ อมโตคธํ โหติ อมตปรายณํ อมต\-ปริโยสานํ ฯเปฯ ปญฺญินฺทฺริยํ ภาวิตํ พหุลีกตํ อมโตคธํ โหติ อมตปรายณํ อมต\-ปริโยสานํ. อิมานิ ปญฺจินฺทฺริยานิ ภาวิตานิ พหุลีกตานิ อมโตคธานิ โหนฺติ อมต\-ปรายณานิ อมต\-ปริโยสานานิ.}

\prose{ภูตปุพฺพาหํ ภนฺเต กสฺสเป สมฺมาสมฺพุทฺเธ พฺรหฺมจริยํ อจริํ, ตตฺรปิ มํ เอวํ ชานนฺติ “สหโก ภิกฺขุ สหโก ภิกฺขู”ติ. โส ขฺวาหํ ภนฺเต อิเมสํเยว ปญฺจนฺนํ อินฺทฺริยานํ ภาวิตตฺตา พหุลีกตตฺตา กาเมสุ กามจฺฉนฺทํ วิราเชตฺวา กายสฺส เภทา ปรํ มรณา สุคติํ พฺรหฺมโลกํ อุปปนฺโน, ตตฺรปิ มํ เอวํ ชานนฺติ “พฺรหฺมา สหมฺปติ พฺรหฺมา สหมฺปตี”ติ. เอวเมตํ ภควา, เอวเมตํ สุคต, อหเมตํ ชานามิ, อหเมตํ ปสฺสามิ, ยถา อิมานิ ปญฺจินฺทฺริยานิ ภาวิตานิ พหุลีกตานิ อมโตคธานิ โหนฺติ อมต\-ปรายณานิ อมตปริโยสานานีติ.  สตฺตมํ.}

\csromanheader{2}{8. สูกรขต\-สุตฺต}
\csromanmatikaanchor{mtk.131}
\csromantocmarkref{subparagraph}{8-10. สูกรขตสุตฺตาทีนิ}{mtk.131}
\bookmark[dest={mtk.131},level=5]{8-10. สูกรขตสุตฺตาทีนิ}
\proseitem{528}{\csromanfolio{205}เอกํ สมยํ ภควา ราชคเห วิหรติ คิชฺฌกูเฏ ปพฺพเต สูกรขตายํ. ตตฺร โข ภควา อายสฺมนฺตํ สาริปุตฺตํ อามนฺเตสิ “กิํ นุ โข สาริปุตฺต อตฺถวสํ สมฺปสฺสมาโน ขีณาสโว ภิกฺขุ ตถาคเต วา ตถาคต\-สาสเน วา ปรม\-นิปจฺจ\-การํ ปวตฺตมาโน ปวตฺตตี”ติ\csromansharedfootnote{668d05165946a34a}{ปวตฺเตตีติ (สี)}. อนุตฺตรํ หิ ภนฺเต โยคกฺเขมํ สมฺปสฺสมาโน ขีณาสโว ภิกฺขุ ตถาคเต วา ตถาคต\-สาสเน วา ปรม\-นิปจฺจ\-การํ ปวตฺตมาโน ปวตฺตตีติ. สาธุ สาธุ สาริปุตฺต, อนุตฺตรํ หิ สาริปุตฺต โยคกฺเขมํ สมฺปสฺสมาโน ขีณาสโว ภิกฺขุ ตถาคเต วา ตถาคต\-สาสเน วา ปรม\-นิปจฺจ\-การํ ปวตฺตมาโน ปวตฺตติ.}

\prose{กตโม จ สาริปุตฺต อนุตฺตโร โยคกฺเขโม, ยํ สมฺปสฺสมาโน ขีณาสโว ภิกฺขุ ตถาคเต วา ตถาคต\-สาสเน วา ปรม\-นิปจฺจ\-การํ ปวตฺตมาโน ปวตฺตตีติ. อิธ ภนฺเต ขีณาสโว ภิกฺขุ สทฺธินฺทฺริยํ ภาเวติ อุปสมคามิํ สมฺโพธคามิํ. วีริยินฺทฺริยํ ภาเวติ ฯเปฯ. สตินฺทฺริยํ ภาเวติ. สมาธินฺทฺริยํ ภาเวติ. ปญฺญินฺทฺริยํ ภาเวติ อุปสมคามิํ สมฺโพธคามิํ. อยํ โข ภนฺเต อนุตฺตโร โยคกฺเขโม, ยํ สมฺปสฺสมาโน ขีณาสโว ภิกฺขุ ตถาคเต วา ตถาคต\-สาสเน วา ปรม\-นิปจฺจ\-การํ ปวตฺตมาโน ปวตฺตตีติ. สาธุ สาธุ สาริปุตฺต, เอโส หิ สาริปุตฺต อนุตฺตโร โยคกฺเขโม, ยํ สมฺปสฺสมาโน ขีณาสโว ภิกฺขุ ตถาคเต วา ตถาคต\-สาสเน วา ปรม\-นิปจฺจ\-การํ ปวตฺตมาโน ปวตฺตตีติ.}

\csromanmatikaanchor{mtk.132}
\csromanmatikaanchor{mtk.133}
\csromantocmarknopage{subparagraph}{7. โพธิปกฺขิยวคฺค}{mtk.132}
\bookmark[dest={mtk.132},level=5]{7. โพธิปกฺขิยวคฺค}
\csromantocmarkref{subparagraph}{1-6. สํโยชนสุตฺตาทีนิ}{mtk.133}
\bookmark[dest={mtk.133},level=5]{1-6. สํโยชนสุตฺตาทีนิ}
\prose{กตโม จ สาริปุตฺต ปรม\-นิปจฺจ\-กาโร, ยํ ขีณาสโว ภิกฺขุ ตถาคเต วา ตถาคต\-สาสเน วา ปรม\-นิปจฺจ\-การํ ปวตฺตมาโน ปวตฺตตีติ. อิธ ภนฺเต ขีณาสโว ภิกฺขุ สตฺถริ สคารโว วิหรติ สปฺปติสฺโส\csromansharedfootnote{66705d72f7774f41}{สปฺปฏิสฺโส (สฺยา, กํ, ก)}, ธมฺเม สคารโว วิหรติ สปฺปติสฺโส, สํเฆ สคารโว วิหรติ สปฺปติสฺโส, สิกฺขาย สคารโว วิหรติ สปฺปติสฺโส, สมาธิสฺมิํ สคารโว วิหรติ สปฺปติสฺโส. อยํ โข ภนฺเต ปรม\-นิปจฺจ\-กาโร, ยํ ขีณาสโว ภิกฺขุ ตถาคเต วา ตถาคต\-สาสเน วา ปรม\-นิปจฺจ\-การํ ปวตฺตมาโน ปวตฺตตีติ. สาธุ สาธุ สาริปุตฺต, เอโส \csromanfolio{206}หิ สาริปุตฺต ปรม\-นิปจฺจ\-กาโร, ยํ ขีณาสโว ภิกฺขุ ตถาคเต วา ตถาคต\-สาสเน วา ปรม\-นิปจฺจ\-การํ ปวตฺตมาโน ปวตฺตตีติ.  อฏฺฐมํ.}

\prose{9. ปฐมอุปฺปาทสุตฺต}

\proseitem{529}{สาวตฺถิ\-นิทานํ. ปญฺจิมานิ ภิกฺขเว อินฺทฺริยานิ ภาวิตานิ พหุลีกตานิ อนุปฺปนฺนานิ อุปฺปชฺชนฺติ นาญฺญตฺร ตถาคตสฺส ปาตุภาวา อรหโต สมฺมา\-สมฺพุทฺธสฺส. กตมานิ ปญฺจ, สทฺธินฺทฺริยํ วีริยินฺทฺริยํ สตินฺทฺริยํ สมาธินฺทฺริยํ ปญฺญินฺทฺริยํ. อิมานิ โข ภิกฺขเว ปญฺจินฺทฺริยานิ ภาวิตานิ พหุลีกตานิ อนุปฺปนฺนานิ อุปฺปชฺชนฺติ นาญฺญตฺร ตถาคตสฺส ปาตุภาวา อรหโต สมฺมา\-สมฺพุทฺธสฺสาติ.  นวมํ.}

\prose{10. ทุติยอุปฺปาทสุตฺต}

\proseitem{530}{ปญฺจิมานิ ภิกฺขเว อินฺทฺริยานิ ภาวิตานิ พหุลีกตานิ อนุปฺปนฺนานิ อุปฺปชฺชนฺติ นาญฺญตฺร สุคตวินยา. กตมานิ ปญฺจ, สทฺธินฺทฺริยํ วีริยินฺทฺริยํ สตินฺทฺริยํ สมาธินฺทฺริยํ ปญฺญินฺทฺริยํ. อิมานิ โข ภิกฺขเว ปญฺจินฺทฺริยานิ ภาวิตานิ พหุลีกตานิ อนุปฺปนฺนานิ อุปฺปชฺชนฺติ นาญฺญตฺร สุคตวินยาติ.  ทสมํ. สูกรขต\-วคฺโค ฉฏฺโฐ.}

\vspace{\tipitakaparskip}
\csromancenter{\textbf{ตสฺสุทฺทานํ}}
\setlength{\csromangathapagewidth}{0pt}
\setlength{\csromangathawakwidth}{0pt}
\csromangathameasuregroup{สาลํ มลฺลิกํ เสโข จ, \\ พฺรหฺมสูกรขตาโย,}{\csromangathabat{สาลํ มลฺลิกํ เสโข จ,}{ปทํ สารํ ปติฏฺฐิตํ.} \\ \csromangathabat{พฺรหฺมสูกรขตาโย,}{อุปฺปาทา อปเร ทุเวติ.}}
\csromangathasetleft{สาลํ มลฺลิกํ เสโข จ, \\ พฺรหฺมสูกรขตาโย,}
\csromangathagroup{\csromangathastanza{\csromangathabat{สาลํ มลฺลิกํ เสโข จ,}{ปทํ สารํ ปติฏฺฐิตํ.} \\ \csromangathabat{พฺรหฺมสูกรขตาโย,}{อุปฺปาทา อปเร ทุเวติ.}}}

\csromanheader{6}{7. โพธิปกฺขิย\-วคฺค}
\csromanheader{2}{1. \textbf{สํโยชนสุตฺต}}
\proseitem{531}{สาวตฺถิ\-นิทานํ. ปญฺจิมานิ ภิกฺขเว อินฺทฺริยานิ ภาวิตานิ พหุลีกตานิ สํโยชน\-ปฺปหานาย\csromansharedfootnote{bfec27c49162467d}{สํโยชนานํ ปหานาย (สฺยา, ก)} สํวตฺตนฺติ. กตมานิ ปญฺจ, สทฺธินฺทฺริยํ ฯเปฯ ปญฺญินฺทฺริยํ. อิมานิ โข ภิกฺขเว ปญฺจินฺทฺริยานิ ภาวิตานิ พหุลีกตานิ สํโยชน\-ปฺปหานาย สํวตฺตนฺตีติ.  ปฐมํ.}

\csromanheader{2}{2. \textbf{อนุสยสุตฺต}}
\proseitem{532}{\csromanfolio{207}ปญฺจิมานิ ภิกฺขเว อินฺทฺริยานิ ภาวิตานิ พหุลีกตานิ อนุสย\-สมุคฺฆาตาย สํวตฺตนฺติ. กตมานิ ปญฺจ, สทฺธินฺทฺริยํ ฯเปฯ ปญฺญินฺทฺริยํ. อิมานิ โข ภิกฺขเว ปญฺจินฺทฺริยานิ ภาวิตานิ พหุลีกตานิ อนุสย\-สมุคฺฆาตาย สํวตฺตนฺตีติ.  ทุติยํ.}

\csromanheader{2}{3. \textbf{ปริญฺญาสุตฺต}}
\proseitem{533}{ปญฺจิมานิ ภิกฺขเว อินฺทฺริยานิ ภาวิตานิ พหุลีกตานิ อทฺธาน\-ปริญฺญาย สํวตฺตนฺติ. กตมานิ ปญฺจ, สทฺธินฺทฺริยํ ฯเปฯ ปญฺญินฺทฺริยํ. อิมานิ โข ภิกฺขเว ปญฺจินฺทฺริยานิ ภาวิตานิ พหุลีกตานิ อทฺธาน\-ปริญฺญาย สํวตฺตนฺตีติ.  ตติยํ.}

\csromanheader{2}{4. อาสวกฺขย\-สุตฺต}
\proseitem{534}{ปญฺจิมานิ ภิกฺขเว อินฺทฺริยานิ ภาวิตานิ พหุลีกตานิ อาสวานํ ขยาย สํวตฺตนฺติ. กตมานิ ปญฺจ, สทฺธินฺทฺริยํ ฯเปฯ ปญฺญินฺทฺริยํ. อิมานิ โข ภิกฺขเว ปญฺจินฺทฺริยานิ ภาวิตานิ พหุลีกตานิ อาสวานํ ขยาย สํวตฺตนฺตีติ.}

\prose{ปญฺจิมานิ ภิกฺขเว อินฺทฺริยานิ ภาวิตานิ พหุลีกตานิ สํโยชน\-ปฺปหานาย สํวตฺตนฺติ, อนุสย\-สมุคฺฆาตาย สํวตฺตนฺติ, อทฺธาน\-ปริญฺญาย สํวตฺตนฺติ, อาสวานํ ขยาย สํวตฺตนฺติ. กตมานิ ปญฺจ, สทฺธินฺทฺริยํ ฯเปฯ ปญฺญินฺทฺริยํ. อิมานิ โข ภิกฺขเว ปญฺจินฺทฺริยานิ ภาวิตานิ พหุลีกตานิ สํโยชน\-ปฺปหานาย สํวตฺตนฺติ, อนุสย\-สมุคฺฆาตาย สํวตฺตนฺติ, อทฺธาน\-ปริญฺญาย สํวตฺตนฺติ, อาสวานํ ขยาย สํวตฺตนฺตีติ.  จตุตฺถํ.}

\csromanheader{2}{5. ปฐม\-ผล\-สุตฺต}
\proseitem{535}{ปญฺจิมานิ ภิกฺขเว อินฺทฺริยานิ. กตมานิ ปญฺจ, สทฺธินฺทฺริยํ ฯเปฯ ปญฺญินฺทฺริยํ. อิมานิ โข ภิกฺขเว ปญฺจินฺทฺริยานิ. อิเมสํ โข ภิกฺขเว ปญฺจนฺนํ อินฺทฺริยานํ ภาวิตตฺตา พหุลีกตตฺตา ทฺวินฺนํ ผลานํ อญฺญตรํ ผลํ ปาฏิกงฺขํ, ทิฏฺเฐว ธมฺเม อญฺญา, สติ วา อุปาทิเสเส อนาคามิตาติ.  ปญฺจมํ.}

\csromanheader{2}{6. ทุติย\-ผล\-สุตฺต}
\csromanmatikaanchor{mtk.134}
\csromantocmarkref{subparagraph}{7-10. รุกฺขสุตฺต}{mtk.134}
\bookmark[dest={mtk.134},level=5]{7-10. รุกฺขสุตฺต}
\proseitem{536}{\csromanfolio{208}ปญฺจิมานิ ภิกฺขเว อินฺทฺริยานิ. กตมานิ ปญฺจ, สทฺธินฺทฺริยํ ฯเปฯ ปญฺญินฺทฺริยํ. อิมานิ โข ภิกฺขเว ปญฺจินฺทฺริยานิ. อิเมสํ โข ภิกฺขเว ปญฺจนฺนํ อินฺทฺริยานํ ภาวิตตฺตา พหุลีกตตฺตา สตฺต ผลา สตฺตานิสํสา ปาฏิกงฺขา. กตเม สตฺต ผลา สตฺตานิสํสา, ทิฏฺเฐว ธมฺเม ปฏิกจฺจ อญฺญํ อาราเธติ. โน เจ ทิฏฺเฐว ธมฺเม ปฏิกจฺจ อญฺญํ อาราเธติ, อถ มรณกาเล อญฺญํ อาราเธติ. โน เจ ทิฏฺเฐว ธมฺเม อญฺญํ อาราเธติ, โน เจ มรณกาเล อญฺญํ อาราเธติ, อถ ปญฺจนฺนํ โอรมฺภาคิยานํ สํโยชนานํ ปริกฺขยา อนฺตรา\-ปรินิพฺพายี โหติ. อุปหจฺจ\-ปรินิพฺพายี โหติ. อสงฺขารปรินิพฺพายี โหติ. สสงฺขาร\-ปรินิพฺพายี โหติ. อุทฺธํโสโต โหติ อกนิฏฺฐคามี. อิเมสํ โข ภิกฺขเว ปญฺจนฺนํ อินฺทฺริยานํ ภาวิตตฺตา พหุลีกตตฺตา อิเม สตฺต ผลา สตฺตานิสํสา ปาฏิกงฺขาติ.  ฉฏฺฐํ.}

\csromanheader{2}{7. ปฐม\-รุกฺข\-สุตฺต}
\proseitem{537}{เสยฺยถาปิ ภิกฺขเว เย เกจิ ชมฺพุทีปกา รุกฺขา, ชมฺพู เตสํ อคฺคมกฺขายติ. เอวเมว โข ภิกฺขเว เย เกจิ โพธิปกฺขิยา ธมฺมา, ปญฺญินฺทฺริยํ เตสํ อคฺคมกฺขายติ ยทิทํ โพธาย. กตเม จ ภิกฺขเว โพธิปกฺขิยา ธมฺมา, สทฺธินฺทฺริยํ ภิกฺขเว โพธิปกฺขิโย ธมฺโม, ตํ โพธาย สํวตฺตติ. วีริยินฺทฺริยํ ฯเปฯ. สตินฺทฺริยํ. สมาธินฺทฺริยํ. ปญฺญินฺทฺริยํ โพธิปกฺขิโย ธมฺโม, ตํ โพธาย สํวตฺตติ. เสยฺยถาปิ ภิกฺขเว เย เกจิ ชมฺพุทีปกา รุกฺขา, ชมฺพู เตสํ อคฺคมกฺขายติ. เอวเมว โข ภิกฺขเว เย เกจิ โพธิปกฺขิยา ธมฺมา, ปญฺญินฺทฺริยํ เตสํ อคฺคมกฺขายติ ยทิทํ โพธายาติ.  สตฺตมํ.}

\csromanheader{2}{8. ทุติย\-รุกฺข\-สุตฺต}
\proseitem{538}{เสยฺยถาปิ ภิกฺขเว เย เกจิ เทวานํ ตาวติํสานํ รุกฺขา, ปาริฉตฺตโก เตสํ อคฺคมกฺขายติ. เอวเมว โข ภิกฺขเว เย เกจิ โพธิปกฺขิยา ธมฺมา, ปญฺญินฺทฺริยํ เตสํ อคฺคมกฺขายติ ยทิทํ โพธาย. กตเม จ ภิกฺขเว โพธิปกฺขิยา ธมฺมา, สทฺธินฺทฺริยํ ภิกฺขเว โพธิปกฺขิโย ธมฺโม, ตํ โพธาย สํวตฺตติ. วีริยินฺทฺริยํ ฯเปฯ. สตินฺทฺริยํ. สมาธินฺทฺริยํ. ปญฺญินฺทฺริยํ โพธิปกฺขิโย ธมฺโม, ตํ โพธาย สํวตฺตติ. เสยฺยถาปิ}

\prose{\csromanfolio{209}ภิกฺขเว เย เกจิ เทวานํ ตาวติํสานํ รุกฺขา, ปาริฉตฺตโก เตสํ อคฺคมกฺขายติ. เอวเมว โข ภิกฺขเว เย เกจิ โพธิปกฺขิยา ธมฺมา, ปญฺญินฺทฺริยํ เตสํ อคฺคมกฺขายติ ยทิทํ โพธายาติ.  อฏฺฐมํ.}

\csromanheader{2}{9. ตติย\-รุกฺข\-สุตฺต}
\proseitem{539}{เสยฺยถาปิ ภิกฺขเว เย เกจิ อสุรานํ รุกฺขา, จิตฺตปาฏลิ เตสํ อคฺคมกฺขายติ. เอวเมว โข ภิกฺขเว เย เกจิ โพธิปกฺขิยา ธมฺมา, ปญฺญินฺทฺริยํ เตสํ อคฺคมกฺขายติ ยทิทํ โพธาย. กตเม จ ภิกฺขเว โพธิปกฺขิยา ธมฺมา, สทฺธินฺทฺริยํ ภิกฺขเว โพธิปกฺขิโย ธมฺโม, ตํ โพธาย สํวตฺตติ ฯเปฯ. ปญฺญินฺทฺริยํ โพธิปกฺขิโย ธมฺโม, ตํ โพธาย สํวตฺตติ. เสยฺยถาปิ ภิกฺขเว เย เกจิ อสุรานํ รุกฺขา, จิตฺตปาฏลิ เตสํ อคฺคมกฺขายติ. เอวเมว โข ภิกฺขเว เย เกจิ โพธิปกฺขิยา ธมฺมา, ปญฺญินฺทฺริยํ เตสํ อคฺคมกฺขายติ ยทิทํ โพธายาติ.  นวมํ.}

\csromanheader{2}{10. จตุตฺถ\-รุกฺข\-สุตฺต}
\proseitemwithclosermidruled{540}{เสยฺยถาปิ ภิกฺขเว เย เกจิ สุปณฺณานํ รุกฺขา, กูฏสิมฺพลี\csromansharedfootnote{5a9cf4b28b820fe4}{โกฏสิมฺพลิ (สฺยา, กํ)} เตสํ อคฺคมกฺขายติ. เอวเมว โข ภิกฺขเว เย เกจิ โพธิปกฺขิยา ธมฺมา, ปญฺญินฺทฺริยํ เตสํ อคฺคมกฺขายติ ยทิทํ โพธาย. กตเม จ ภิกฺขเว โพธิปกฺขิยา ธมฺมา, สทฺธินฺทฺริยํ ภิกฺขเว โพธิปกฺขิโย ธมฺโม, ตํ โพธาย สํวตฺตติ ฯเปฯ. ปญฺญินฺทฺริยํ โพธิปกฺขิโย ธมฺโม, ตํ โพธาย สํวตฺตติ. เสยฺยถาปิ ภิกฺขเว เย เกจิ สุปณฺณานํ รุกฺขา, กูฏสิมฺพลี เตสํ อคฺคมกฺขายติ. เอวเมว โข ภิกฺขเว เย เกจิ โพธิปกฺขิยา ธมฺมา, ปญฺญินฺทฺริยํ เตสํ อคฺคมกฺขายติ ยทิทํ โพธายาติ.  ทสมํ.}{โพธิปกฺขิย\-วคฺโค สตฺตโม.}
\csromancenter{\textbf{ตสฺสุทฺทานํ}}
\setlength{\csromangathapagewidth}{0pt}
\setlength{\csromangathawakwidth}{0pt}
\csromangathameasuregroup{สํโยชนา อนุสยา, \\ เทฺว ผลา จตุโร รุกฺขา,}{\csromangathabat{สํโยชนา อนุสยา,}{ปริญฺญา อาสวกฺขยา.} \\ \csromangathabat{เทฺว ผลา จตุโร รุกฺขา,}{วคฺโค เตน ปวุจฺจตีติ.}}
\csromangathasetleft{สํโยชนา อนุสยา, \\ เทฺว ผลา จตุโร รุกฺขา,}
\csromangathagroup{\csromangathastanza{\csromangathabat{สํโยชนา อนุสยา,}{ปริญฺญา อาสวกฺขยา.} \\ \csromangathabat{เทฺว ผลา จตุโร รุกฺขา,}{วคฺโค เตน ปวุจฺจตีติ.}}}

\chapterheadpageread{8. คงฺคา\-เปยฺยาล\-วคฺค}
\csromanmatikaanchor{mtk.135}
\csromanmatikaanchor{mtk.136}
\csromantocmarknopage{subparagraph}{8. คงฺคาเปยฺยาลวคฺค}{mtk.135}
\bookmark[dest={mtk.135},level=5]{8. คงฺคาเปยฺยาลวคฺค}
\csromantocmarkref{subparagraph}{1-12. ปาจีนาทิสุตฺต}{mtk.136}
\bookmark[dest={mtk.136},level=5]{1-12. ปาจีนาทิสุตฺต}
\csromanheader{6}{1-12. ปาจีนาทิสุตฺต}
\proseitem{541-552}{\csromanfolio{210}เสยฺยถาปิ ภิกฺขเว คงฺคา นที ปาจีนนินฺนา ปาจีนโปณา ปาจีนปพฺภารา. เอวเมว โข ภิกฺขเว ภิกฺขุ ปญฺจินฺทฺริยานิ ภาเวนฺโต ปญฺจินฺทฺริยานิ พหุลีกโรนฺโต นิพฺพานนินฺโน โหติ นิพฺพานโปโณ นิพฺพาน\-ปพฺภาโร. กถญฺจ ภิกฺขเว ภิกฺขุ ปญฺจินฺทฺริยานิ ภาเวนฺโต ปญฺจินฺทฺริยานิ พหุลีกโรนฺโต นิพฺพานนินฺโน โหติ นิพฺพานโปโณ นิพฺพาน\-ปพฺภาโร, อิธ ภิกฺขเว ภิกฺขุ สทฺธินฺทฺริยํ ภาเวติ วิเวกนิสฺสิตํ วิราคนิสฺสิตํ นิโรธ\-นิสฺสิตํ โวสฺสคฺค\-ปริณามิํ. วีริยินฺทฺริยํ ฯเปฯ. สตินฺทฺริยํ. สมาธินฺทฺริยํ. ปญฺญินฺทฺริยํ ภาเวติ วิเวกนิสฺสิตํ วิราคนิสฺสิตํ นิโรธ\-นิสฺสิตํ โวสฺสคฺค\-ปริณามิํ. เอวํ โข ภิกฺขเว ภิกฺขุ ปญฺจินฺทฺริยานิ ภาเวนฺโต ปญฺจินฺทฺริยานิ พหุลีกโรนฺโต นิพฺพานนินฺโน โหติ นิพฺพานโปโณ นิพฺพาน\-ปพฺภาโรติ.  ทฺวาทสมํ. คงฺคา\-เปยฺยาล\-วคฺโค อฏฺฐโม.}

\vspace{\tipitakaparskip}
\csromancenter{\textbf{ตสฺสุทฺทานํ}}
\setlength{\csromangathapagewidth}{0pt}
\setlength{\csromangathawakwidth}{0pt}
\csromangathameasuregroup{ฉ ปาจีนโต นินฺนา,}{\csromangathabat{ฉ ปาจีนโต นินฺนา,}{ฉ นินฺนา จ สมุทฺทโต.} \\ เทฺวเต ฉ ทฺวาทส โหนฺติ, วคฺโค เตน ปวุจฺจตีติ. \\ อปฺปมาทวคฺโค วิตฺถาเรตพฺโพ.}
\csromangathasetleft{ฉ ปาจีนโต นินฺนา,}
\csromangathagroup{\csromangathastanza{\csromangathabat{ฉ ปาจีนโต นินฺนา,}{ฉ นินฺนา จ สมุทฺทโต.} \\ เทฺวเต ฉ ทฺวาทส โหนฺติ, วคฺโค เตน ปวุจฺจตีติ. \\ อปฺปมาทวคฺโค วิตฺถาเรตพฺโพ.}}

\vspace{\tipitakaparskip}
\csromancenter{\textbf{ตสฺสุทฺทานํ}}
\setlength{\csromangathapagewidth}{0pt}
\setlength{\csromangathawakwidth}{0pt}
\csromangathameasuregroup{ตถาคตํ ปทํ กูฏํ,}{\csromangathabat{ตถาคตํ ปทํ กูฏํ,}{มูลํ สาเรน วสฺสิกํ.} \\ ราชา จนฺทิมสูริยา, วตฺเถน ทสมํ ปทนฺติ. \\ พลกรณียวคฺโค วิตฺถาเรตพฺโพ.}
\csromangathasetleft{ตถาคตํ ปทํ กูฏํ,}
\csromangathagroup{\csromangathastanza{\csromangathabat{ตถาคตํ ปทํ กูฏํ,}{มูลํ สาเรน วสฺสิกํ.} \\ ราชา จนฺทิมสูริยา, วตฺเถน ทสมํ ปทนฺติ. \\ พลกรณีย\-วคฺโค วิตฺถาเรตพฺโพ.}}

\vspace{\tipitakaparskip}
\csromancenter{\textbf{ตสฺสุทฺทานํ}}
\setlength{\csromangathapagewidth}{0pt}
\setlength{\csromangathawakwidth}{0pt}
\csromangathameasuregroup{พลํ พีชญฺจ นาโค จ, \\ อากาเสน จ เทฺว เมฆา,}{\csromangathabat{พลํ พีชญฺจ นาโค จ,}{รุกฺโข กุมฺเภน สูกิยา.} \\ \csromangathabat{อากาเสน จ เทฺว เมฆา,}{นาวา อาคนฺตุกา นทีติ.}}
\csromangathasetleft{พลํ พีชญฺจ นาโค จ, \\ อากาเสน จ เทฺว เมฆา,}
\csromangathagroup{\csromangathastanza{\csromangathabat{พลํ พีชญฺจ นาโค จ,}{รุกฺโข กุมฺเภน สูกิยา.} \\ \csromangathabat{อากาเสน จ เทฺว เมฆา,}{นาวา อาคนฺตุกา นทีติ.}}}

\vspace{\tipitakaparskip}
\csromancenter{เอสนาวคฺโค วิตฺถาเรตพฺโพ.}
\csromancenter{\textbf{ตสฺสุทฺทานํ}}
\setlength{\csromangathapagewidth}{0pt}
\setlength{\csromangathawakwidth}{0pt}
\csromangathameasuregroup{เอสนา วิธา อาสโว, \\ ขิลํ มลญฺจ นีโฆ จ,}{\csromangathabat{เอสนา วิธา อาสโว,}{ภโว จ ทุกฺขตา ติสฺโส.} \\ \csromangathabat{ขิลํ มลญฺจ นีโฆ จ,}{เวทนา ตณฺหา ตสินา จาติ.}}
\csromangathasetleft{เอสนา วิธา อาสโว, \\ ขิลํ มลญฺจ นีโฆ จ,}
\csromangathagroup{\csromangathastanza{\csromanfolio{211}\csromangathabat{เอสนา วิธา อาสโว,}{ภโว จ ทุกฺขตา ติสฺโส.} \\ \csromangathabat{ขิลํ มลญฺจ นีโฆ จ,}{เวทนา ตณฺหา ตสินา จาติ.}}}

\chapterhead{12. \textbf{โอฆวคฺค}}
\csromanmatikaanchor{mtk.137}
\csromanmatikaanchor{mtk.138}
\csromantocmarknopage{subparagraph}{12. โอฆวคฺค}{mtk.137}
\bookmark[dest={mtk.137},level=5]{12. โอฆวคฺค}
\csromantocmarkref{subparagraph}{1-10. โอฆาทิสุตฺต}{mtk.138}
\bookmark[dest={mtk.138},level=5]{1-10. โอฆาทิสุตฺต}
\csromanheader{6}{1-10. โอฆาทิสุตฺต}
\proseitem{587-596}{ปญฺจิมานิ ภิกฺขเว อุทฺธมฺ\-ภาคิยานิ สํโยชนานิ. กตมานิ ปญฺจ, รูปราโค อรูปราโค มาโน อุทฺธจฺจํ อวิชฺชา, อิมานิ โข ภิกฺขเว ปญฺจุ\-ทฺธมฺภาคิยานิ สํโยชนานิ. อิเมสํ โข ภิกฺขเว ปญฺจนฺนํ อุทฺธมฺ\-ภาคิยานํ สํโยชนานํ อภิญฺญาย ปริญฺญาย ปริกฺขยาย ปหานาย ปญฺจินฺทฺริยานิ ภาเวตพฺพานิ. กตมานิ ปญฺจ, อิธ ภิกฺขเว ภิกฺขุ สทฺธินฺทฺริยํ ภาเวติ วิเวกนิสฺสิตํ วิราคนิสฺสิตํ นิโรธ\-นิสฺสิตํ โวสฺสคฺค\-ปริณามิํ ฯเปฯ. ปญฺญินฺทฺริยํ ภาเวติ วิเวกนิสฺสิตํ วิราคนิสฺสิตํ นิโรธ\-นิสฺสิตํ โวสฺสคฺค\-ปริณามิํ. อิเมสํ โข ภิกฺขเว ปญฺจนฺนํ อุทฺธมฺ\-ภาคิยานํ สํโยชนานํ อภิญฺญาย ปริญฺญาย ปริกฺขยาย ปหานาย อิมานิ ปญฺจินฺทฺริยานิ ภาเวตพฺพานีติ.  ทสมํ.  (ยถา มคฺคสํยุตฺตํ, ตถา วิตฺถาเรตพฺพํ.) โอฆวคฺโค ทฺวาทสโม.}

\vspace{\tipitakaparskip}
\csromancenter{\textbf{ตสฺสุทฺทานํ}}
\setlength{\csromangathapagewidth}{0pt}
\setlength{\csromangathawakwidth}{0pt}
\csromangathameasuregroup{โอโฆ โยโค อุปาทานํ, \\ กามคุณา นีวรณา,}{\csromangathabat{โอโฆ โยโค อุปาทานํ,}{คนฺถา อนุสเยน จ.} \\ \csromangathabat{กามคุณา นีวรณา,}{ขนฺธา โอรุทฺธมฺภาคิยาติ.}}
\csromangathasetleft{โอโฆ โยโค อุปาทานํ, \\ กามคุณา นีวรณา,}
\csromangathagroup{\csromangathastanza{\csromangathabat{โอโฆ โยโค อุปาทานํ,}{คนฺถา อนุสเยน จ.} \\ \csromangathabat{กามคุณา นีวรณา,}{ขนฺธา โอรุทฺธมฺภาคิยาติ.}}}

\chapterheadpageread{13. คงฺคา\-เปยฺยาล\-วคฺค}
\csromanmatikaanchor{mtk.139}
\csromanmatikaanchor{mtk.140}
\csromantocmarknopage{subparagraph}{13. คงฺคาเปยฺยาลวคฺค}{mtk.139}
\bookmark[dest={mtk.139},level=5]{13. คงฺคาเปยฺยาลวคฺค}
\csromantocmarkref{subparagraph}{1-12. ปาจีนาทิสุตฺต}{mtk.140}
\bookmark[dest={mtk.140},level=5]{1-12. ปาจีนาทิสุตฺต}
\csromanheader{6}{1-12. ปาจีนาทิสุตฺต}
\proseitem{597-608}{\csromanfolio{212}เสยฺยถาปิ ภิกฺขเว คงฺคา นที ปาจีนนินฺนา ปาจีนโปณา ปาจีนปพฺภารา. เอวเมว โข ภิกฺขเว ภิกฺขุ ปญฺจินฺทฺริยานิ ภาเวนฺโต ปญฺจินฺทฺริยานิ พหุลีกโรนฺโต นิพฺพานนินฺโน โหติ นิพฺพานโปโณ นิพฺพาน\-ปพฺภาโร. กถญฺจ ภิกฺขเว ภิกฺขุ ปญฺจินฺทฺริยานิ ภาเวนฺโต ปญฺจินฺทฺริยานิ พหุลีกโรนฺโต นิพฺพานนินฺโน โหติ นิพฺพานโปโณ นิพฺพาน\-ปพฺภาโร, อิธ ภิกฺขเว ภิกฺขุ สทฺธินฺทฺริยํ ภาเวติ ราค\-วินย\-ปริโยสานํ โทส\-วินย\-ปริโยสานํ โมห\-วินย\-ปริโยสานํ ฯเปฯ. ปญฺญินฺทฺริยํ ภาเวติ ราค\-วินย\-ปริโยสานํ โทส\-วินย\-ปริโยสานํ โมห\-วินย\-ปริโยสานํ. เอวํ โข ภิกฺขเว ภิกฺขุ ปญฺจินฺทฺริยานิ ภาเวนฺโต ปญฺจินฺทฺริยานิ พหุลีกโรนฺโต นิพฺพานนินฺโน โหติ นิพฺพานโปโณ นิพฺพาน\-ปพฺภาโรติ.  ทฺวาทสมํ. คงฺคา\-เปยฺยาล\-วคฺโค เตรสโม.}

\vspace{\tipitakaparskip}
\csromancenter{\textbf{ตสฺสุทฺทานํ}}
\setlength{\csromangathapagewidth}{0pt}
\setlength{\csromangathawakwidth}{0pt}
\csromangathameasuregroup{ฉ ปาจีนโต นินฺนา, \\ เทฺวเต ฉ ทฺวาทส โหนฺติ,}{\csromangathabat{ฉ ปาจีนโต นินฺนา,}{ฉ นินฺนา จ สมุทฺทโต.} \\ \csromangathabat{เทฺวเต ฉ ทฺวาทส โหนฺติ,}{วคฺโค เตน ปวุจฺจตีติ.}}
\csromangathasetleft{ฉ ปาจีนโต นินฺนา, \\ เทฺวเต ฉ ทฺวาทส โหนฺติ,}
\csromangathagroup{\csromangathastanza{\csromangathabat{ฉ ปาจีนโต นินฺนา,}{ฉ นินฺนา จ สมุทฺทโต.} \\ \csromangathabat{เทฺวเต ฉ ทฺวาทส โหนฺติ,}{วคฺโค เตน ปวุจฺจตีติ.}}}

\vspace{\tipitakaparskip}
\csromancenter{(อปฺปมาทวคฺค พลกรณีย\-วคฺค เอสนาวคฺคา วิตฺถาเรตพฺพา.)}
\csromansectionrule
\chapterhead{17. \textbf{โอฆวคฺค}}
\csromanmatikaanchor{mtk.141}
\csromanmatikaanchor{mtk.142}
\csromantocmarknopage{subparagraph}{17. โอฆวคฺค}{mtk.141}
\bookmark[dest={mtk.141},level=5]{17. โอฆวคฺค}
\csromantocmarkref{subparagraph}{1-10. โอฆาทิสุตฺต}{mtk.142}
\bookmark[dest={mtk.142},level=5]{1-10. โอฆาทิสุตฺต}
\csromanheader{6}{1-10. โอฆาทิสุตฺต}
\proseitem{641-650}{ปญฺจิมานิ ภิกฺขเว อุทฺธมฺ\-ภาคิยานิ สํโยชนานิ. กตมานิ ปญฺจ, รูปราโค อรูปราโค มาโน อุทฺธจฺจํ อวิชฺชา, อิมานิ โข ภิกฺขเว ปญฺจุ\-ทฺธมฺภาคิยานิ สํโยชนานิ, อิเมสํ โข ภิกฺขเว ปญฺจนฺนํ อุทฺธมฺ\-ภาคิยานํ สํโยชนานํ อภิญฺญาย ปริญฺญาย ปริกฺขยาย ปหานาย ปญฺจินฺทฺริยานิ ภาเวตพฺพานิ. กตมานิ ปญฺจ, อิธ ภิกฺขเว ภิกฺขุ สทฺธินฺทฺริยํ ภาเวติ ราค\-วินย\-ปริโยสานํ โทส\-วินย\-ปริโยสานํ โมห\-วินย\-ปริโยสานํ.}

\prose{\csromanfolio{213}วีริยินฺทฺริยํ ฯเปฯ. สตินฺทฺริยํ. สมาธินฺทฺริยํ. ปญฺญินฺทฺริยํ ภาเวติ ราค\-วินย\-ปริโยสานํ โทส\-วินย\-ปริโยสานํ โมห\-วินย\-ปริโยสานํ. อิเมสํ โข ภิกฺขเว ปญฺจนฺนํ อุทฺธมฺ\-ภาคิยานํ สํโยชนานํ อภิญฺญาย ปริญฺญาย ปริกฺขยาย ปหานาย อิมานิ ปญฺจินฺทฺริยานิ ภาเวตพฺพานีติ. โอฆวคฺโค สตฺตรสโม.}

\vspace{\tipitakaparskip}
\csromancenter{\textbf{ตสฺสุทฺทานํ}}
\setlength{\csromangathapagewidth}{0pt}
\setlength{\csromangathawakwidth}{0pt}
\csromangathameasuregroup{โอโฆ โยโค อุปาทานํ, \\ กามคุณา นีวรณา,}{\csromangathabat{โอโฆ โยโค อุปาทานํ,}{คนฺถา อนุสเยน จ.} \\ \csromangathabat{กามคุณา นีวรณา,}{ขนฺธา โอรุทฺธมฺภาคิยาติ.} \\ อินฺทฺริยสํยุตฺตํ จตุตฺถํ.}
\csromangathameasurewak{อินฺทฺริยสํยุตฺตํ จตุตฺถํ.}
\csromangathasetleft{โอโฆ โยโค อุปาทานํ, \\ กามคุณา นีวรณา,}
\csromangathagroup{\csromangathastanza{\csromangathabat{โอโฆ โยโค อุปาทานํ,}{คนฺถา อนุสเยน จ.} \\ \csromangathabat{กามคุณา นีวรณา,}{ขนฺธา โอรุทฺธมฺภาคิยาติ.}}\csromangathastanza{\csromangathawak{อินฺทฺริยสํยุตฺตํ จตุตฺถํ.}}}

\csromanmatikaanchor{mtk.143}
\csromantocmarknopage{subparagraph}{5. สมฺมปฺปธานสํยุตฺต}{mtk.143}
\bookmark[dest={mtk.143},level=5]{5. สมฺมปฺปธานสํยุตฺต}
\csromanheader{6}{5. สมฺมปฺปธาน\-สํยุตฺต}
\chapterhead{1. คงฺคา\-เปยฺยาล\-วคฺค}
\csromanmatikaanchor{mtk.144}
\csromanmatikaanchor{mtk.145}
\csromantocmarknopage{subparagraph}{1. คงฺคาเปยฺยาลวคฺค}{mtk.144}
\bookmark[dest={mtk.144},level=5]{1. คงฺคาเปยฺยาลวคฺค}
\csromantocmarkref{subparagraph}{1-12. ปาจีนาทิสุตฺต}{mtk.145}
\bookmark[dest={mtk.145},level=5]{1-12. ปาจีนาทิสุตฺต}
\csromanheader{6}{1-12. ปาจีนาทิสุตฺต}
\proseitem{651-662}{\csromanfolio{214}สาวตฺถิ\-นิทานํ. ตตฺร โข ภควา เอตทโวจ\csromandash{}จตฺตาโรเม ภิกฺขเว สมฺมปฺปธานา. กตเม จตฺตาโร, อิธ ภิกฺขเว ภิกฺขุ อนุปฺปนฺนานํ ปาปกานํ อกุสลานํ ธมฺมานํ อนุปฺปาทาย ฉนฺทํ ชเนติ วายมติ, วีริยํ อารภติ, จิตฺตํ ปคฺคณฺหาติ ปทหติ, อุปฺปนฺนานํ ปาปกานํ อกุสลานํ ธมฺมานํ ปหานาย ฉนฺทํ ชเนติ วายมติ, วีริยํ อารภติ, จิตฺตํ ปคฺคณฺหาติ ปทหติ, อนุปฺปนฺนานํ กุสลานํ ธมฺมานํ อุปฺปาทาย ฉนฺทํ ชเนติ วายมติ, วีริยํ อารภติ, จิตฺตํ ปคฺคณฺหาติ ปทหติ, อุปฺปนฺนานํ กุสลานํ ธมฺมานํ ฐิติยา อสมฺโมสาย ภิยฺโยภาวาย เวปุลฺลาย ภาวนาย ปาริปูริยา ฉนฺทํ ชเนติ วายมติ, วีริยํ อารภติ, จิตฺตํ ปคฺคณฺหาติ ปทหติ. อิเม โข ภิกฺขเว จตฺตาโร สมฺมปฺปธานาติ.}

\prosewithclosermidruled{เสยฺยถาปิ ภิกฺขเว คงฺคา นที ปาจีนนินฺนา ปาจีนโปณา ปาจีนปพฺภารา. เอวเมว โข ภิกฺขเว ภิกฺขุ จตฺตาโร สมฺมปฺปธาเน ภาเวนฺโต จตฺตาโร สมฺมปฺปธาเน พหุลีกโรนฺโต นิพฺพานนินฺโน โหติ นิพฺพานโปโณ นิพฺพาน\-ปพฺภาโร. กถญฺจ ภิกฺขเว ภิกฺขุ จตฺตาโร สมฺมปฺปธาเน ภาเวนฺโต จตฺตาโร สมฺมปฺปธาเน พหุลีกโรนฺโต นิพฺพานนินฺโน โหติ นิพฺพานโปโณ นิพฺพาน\-ปพฺภาโร, อิธ ภิกฺขเว ภิกฺขุ อนุปฺปนฺนานํ ปาปกานํ อกุสลานํ ธมฺมานํ อนุปฺปาทาย ฉนฺทํ ชเนติ, วายมติ วีริยํ อารภติ, จิตฺตํ ปคฺคณฺหาติ ปทหติ, อุปฺปนฺนานํ ปาปกานํ อกุสลานํ ธมฺมานํ ปหานาย ฉนฺทํ ชเนติ วายมติ, วีริยํ อารภติ, จิตฺตํ ปคฺคณฺหาติ, ปทหติ, อนุปฺปนฺนานํ กุสลานํ ธมฺมานํ อุปฺปาทาย ฉนฺทํ ชเนติ วายมติ, วีริยํ อารภติ, จิตฺตํ ปคฺคณฺหาติ ปทหติ, อุปฺปนฺนานํ กุสลานํ ธมฺมานํ ฐิติยา อสมฺโมสาย ภิยฺโยภาวาย เวปุลฺลาย ภาวนาย ปาริปูริยา ฉนฺทํ ชเนติ วายมติ, วีริยํ อารภติ, จิตฺตํ ปคฺคณฺหาติ ปทหติ. เอวํ โข ภิกฺขเว ภิกฺขุ จตฺตาโร สมฺมปฺปธาเน ภาเวนฺโต จตฺตาโร สมฺมปฺปธาเน พหุลีกโรนฺโต นิพฺพานนินฺโน โหติ นิพฺพานโปโณ นิพฺพาน\-ปพฺภาโรติ.  ทฺวาทสมํ.}{คงฺคา\-เปยฺยาล\-วคฺโค ปฐโม.}
\csromancenter{\textbf{ตสฺสุทฺทานํ}}
\csromanmatikaanchor{mtk.146}
\csromanmatikaanchor{mtk.147}
\csromantocmarknopage{subparagraph}{2. อปฺปมาทวคฺค}{mtk.146}
\bookmark[dest={mtk.146},level=5]{2. อปฺปมาทวคฺค}
\csromantocmarkref{section}{อุทฺทานคาถา}{mtk.147}
\bookmark[dest={mtk.147},level=1]{อุทฺทานคาถา}
\setlength{\csromangathapagewidth}{0pt}
\setlength{\csromangathawakwidth}{0pt}
\csromangathameasuregroup{ฉ ปาจีนโต นินฺนา, \\ เทฺวเต ฉ ทฺวาทส โหนฺติ,}{\csromangathabat{ฉ ปาจีนโต นินฺนา,}{ฉ นินฺนา จ สมุทฺทโต.} \\ \csromangathabat{เทฺวเต ฉ ทฺวาทส โหนฺติ,}{วคฺโค เตน ปวุจฺจตีติ.}}
\csromangathasetleft{ฉ ปาจีนโต นินฺนา, \\ เทฺวเต ฉ ทฺวาทส โหนฺติ,}
\csromangathagroup{\csromangathastanza{\csromanfolio{215}\csromangathabat{ฉ ปาจีนโต นินฺนา,}{ฉ นินฺนา จ สมุทฺทโต.} \\ \csromangathabat{เทฺวเต ฉ ทฺวาทส โหนฺติ,}{วคฺโค เตน ปวุจฺจตีติ.}}}

\vspace{\tipitakaparskip}
\csromancenter{(สมฺมปฺปธานสํยุตฺตสฺส คงฺคาเปยฺยาลี สมฺมปฺปธาน\-วเสน วิตฺถาเรตพฺพา.)}
\csromansectionrule
\chapterhead{2. \textbf{อปฺปมาทวคฺค}}
\vspace{\tipitakaparskip}
\csromancenter{\textbf{ตสฺสุทฺทานํ}}
\setlength{\csromangathapagewidth}{0pt}
\setlength{\csromangathawakwidth}{0pt}
\csromangathameasuregroup{ตถาคตํ ปทํ กูฏํ, \\ ราชา จนฺทิมสูริยา,}{\csromangathabat{ตถาคตํ ปทํ กูฏํ,}{มูลํ สาเรน วสฺสิกํ.} \\ \csromangathabat{ราชา จนฺทิมสูริยา,}{วตฺเถน ทสมํ ปทนฺติ.}}
\csromangathasetleft{ตถาคตํ ปทํ กูฏํ, \\ ราชา จนฺทิมสูริยา,}
\csromangathagroup{\csromangathastanza{\csromangathabat{ตถาคตํ ปทํ กูฏํ,}{มูลํ สาเรน วสฺสิกํ.} \\ \csromangathabat{ราชา จนฺทิมสูริยา,}{วตฺเถน ทสมํ ปทนฺติ.}}}

\vspace{\tipitakaparskip}
\csromancenter{(อปฺปมาทวคฺโค สมฺมปฺปธาน\-วเสน วิตฺถาเรตพฺโพ.)}
\csromansectionrule
\chapterhead{3. พลกรณีย\-วคฺค}
\csromanmatikaanchor{mtk.148}
\csromanmatikaanchor{mtk.149}
\csromantocmarknopage{subparagraph}{3. พลกรณียวคฺค}{mtk.148}
\bookmark[dest={mtk.148},level=5]{3. พลกรณียวคฺค}
\csromantocmarkref{subparagraph}{1-12. พลกรณียาทิสุตฺต}{mtk.149}
\bookmark[dest={mtk.149},level=5]{1-12. พลกรณียาทิสุตฺต}
\csromanheader{6}{1-12. พลกรณียา\-ทิ\-สุตฺต}
\proseitemwithclosermidruled{673-684}{เสยฺยถาปิ ภิกฺขเว เย เกจิ พลกรณียา กมฺมนฺตา กยิรนฺติ, สพฺเพเต ปถวิํ นิสฺสาย ปถวิยํ ปติฏฺฐาย, เอวเมเต พลกรณียา กมฺมนฺตา กยิรนฺติ. เอวเมว โข ภิกฺขเว ภิกฺขุ สีลํ นิสฺสาย สีเล ปติฏฺฐาย จตฺตาโร สมฺมปฺปธาเน ภาเวติ จตฺตาโร สมฺมปฺปธาเน พหุลีกโรติ. กถญฺจ ภิกฺขเว ภิกฺขุ สีลํ นิสฺสาย สีเล ปติฏฺฐาย จตฺตาโร สมฺมปฺปธาเน ภาเวติ จตฺตาโร สมฺมปฺปธาเน พหุลีกโรติ, อิธ ภิกฺขเว ภิกฺขุ อนุปฺปนฺนานํ ปาปกานํ อกุสลานํ ธมฺมานํ อนุปฺปาทาย ฉนฺทํ ชเนติ วายมติ, วีริยํ อารภติ, จิตฺตํ ปคฺคณฺหาติ ปทหติ ฯเปฯ อุปฺปนฺนานํ กุสลานํ ธมฺมานํ ฐิติยา อสมฺโมสาย ภิยฺโยภาวาย เวปุลฺลาย ภาวนาย ปาริปูริยา ฉนฺทํ ชเนติ วายมติ, วีริยํ อารภติ, จิตฺตํ ปคฺคณฺหาติ ปทหติ. เอวํ โข ภิกฺขเว ภิกฺขุ สีลํ \csromanfolio{216}นิสฺสาย สีเล ปติฏฺฐาย จตฺตาโร สมฺมปฺปธาเน ภาเวติ จตฺตาโร สมฺมปฺปธาเน พหุลีกโรตีติ. (เอวํ พลกรณีย\-วคฺโค สมฺมปฺปธาน\-วเสน วิตฺถาเรตพฺโพ.). ทฺวาทสมํ.}{พลกรณีย\-วคฺโค ตติโย.}
\csromanheader{2}{\textbf{ตสฺสุทานํ}}
\setlength{\csromangathapagewidth}{0pt}
\setlength{\csromangathawakwidth}{0pt}
\csromangathameasuregroup{พลํ พีชญฺจ นาโค จ,}{\csromangathabat{พลํ พีชญฺจ นาโค จ,}{รุกฺโข กุมฺเภน สูกิยา.}}
\csromangathasetleft{พลํ พีชญฺจ นาโค จ,}
\csromangathagroup{\csromangathastanza{\csromangathabat{พลํ พีชญฺจ นาโค จ,}{รุกฺโข กุมฺเภน สูกิยา.}}}

\vspace{\tipitakaparskip}
\csromancenter{อากาเสน จ เทฺว เมฆา, นาวา อาคนฺตุกา นทีติ.}
\csromansectionrule
\chapterhead{4. \textbf{เอสนาวคฺค}}
\csromanmatikaanchor{mtk.150}
\csromanmatikaanchor{mtk.151}
\csromantocmarknopage{subparagraph}{4. เอสนาวคฺค}{mtk.150}
\bookmark[dest={mtk.150},level=5]{4. เอสนาวคฺค}
\csromantocmarkref{subparagraph}{1-10. เอสนาทิสุตฺตทสก}{mtk.151}
\bookmark[dest={mtk.151},level=5]{1-10. เอสนาทิสุตฺตทสก}
\csromanheader{6}{1-10. เอสนา\-ทิ\-สุตฺต\-ทสก}
\proseitemwithclosermidruled{685-694}{ติสฺโส อิมา ภิกฺขเว เอสนา. กตมา ติสฺโส, กาเมสนา ภเวสนา พฺรหฺมจริเย\-สนา. อิมา โข ภิกฺขเว ติสฺโส เอสนา. อิมาสํ โข ภิกฺขเว ติสฺสนฺนํ เอสนานํ อภิญฺญาย ปริญฺญาย ปริกฺขยาย ปหานาย จตฺตาโร สมฺมปฺปธานา ภาเวตพฺพา. กตเม จตฺตาโร, อิธ ภิกฺขเว ภิกฺขุ อนุปฺปนฺนานํ ฯเปฯ อุปฺปนฺนานํ กุสลานํ ธมฺมานํ ฐิติยา อสมฺโมสาย ภิยฺโยภาวาย เวปุลฺลาย ภาวนาย ปาริปูริยา ฉนฺทํ ชเนติ วายมติ, วีริยํ อารภติ, จิตฺตํ ปคฺคณฺหาติ ปทหติ. อิมาสํ โข ภิกฺขเว ติสฺสนฺนํ เอสนานํ อภิญฺญาย ปริญฺญาย ปริกฺขยาย ปหานาย อิเม จตฺตาโร สมฺมปฺปธานา ภาเวตพฺพาติ. วิตฺถาเรตพฺพํ.  ทสมํ.}{เอสนาวคฺโค จตุตฺโถ.}
\csromancenter{\textbf{ตสฺสุทฺทานํ}}
\setlength{\csromangathapagewidth}{0pt}
\setlength{\csromangathawakwidth}{0pt}
\csromangathameasuregroup{เอสนา วิธา อาสโว, \\ ขิลํ มลญฺจ นีโฆ จ,}{\csromangathabat{เอสนา วิธา อาสโว,}{ภโว จ ทุกฺขตา ติสฺโส.} \\ \csromangathabat{ขิลํ มลญฺจ นีโฆ จ,}{เวทนา ตณฺหา ตสินา จาติ.}}
\csromangathasetleft{เอสนา วิธา อาสโว, \\ ขิลํ มลญฺจ นีโฆ จ,}
\csromangathagroup{\csromangathastanza{\csromangathabat{เอสนา วิธา อาสโว,}{ภโว จ ทุกฺขตา ติสฺโส.} \\ \csromangathabat{ขิลํ มลญฺจ นีโฆ จ,}{เวทนา ตณฺหา ตสินา จาติ.}}}

\chapterheadpageread{5. \textbf{โอฆวคฺค}}
\csromanmatikaanchor{mtk.152}
\csromanmatikaanchor{mtk.153}
\csromantocmarknopage{subparagraph}{5. โอฆวคฺค}{mtk.152}
\bookmark[dest={mtk.152},level=5]{5. โอฆวคฺค}
\csromantocmarkref{subparagraph}{1-10. โอฆาทิสุตฺต}{mtk.153}
\bookmark[dest={mtk.153},level=5]{1-10. โอฆาทิสุตฺต}
\csromanheader{6}{1-10. โอฆาทิสุตฺต}
\proseitemwithclosermidruled{695-704}{\csromanfolio{217}ปญฺจิมานิ ภิกฺขเว อุทฺธมฺ\-ภาคิยานิ สํโยชนานิ. กตมานิ ปญฺจ, รูปราโค อรูปราโค มาโน อุทฺธจฺจํ อวิชฺชา, อิมานิ โข ภิกฺขเว ปญฺจุ\-ทฺธมฺภาคิยานิ สํโยชนานิ. อิเมสํ โข ภิกฺขเว ปญฺจนฺนํ อุทฺธมฺ\-ภาคิยานํ สํโยชนานํ อภิญฺญาย ปริญฺญาย ปริกฺขยาย ปหานาย จตฺตาโร สมฺมปฺปธานา ภาเวตพฺพา. กตเม จตฺตาโร, อิธ ภิกฺขเว ภิกฺขุ อนุปฺปนฺนานํ ฯเปฯ อุปฺปนฺนานํ กุสลานํ ธมฺมานํ ฐิติยา อสมฺโมสาย ภิยฺโยภาวาย เวปุลฺลาย ภาวนาย ปาริปูริยา ฉนฺทํ ชเนติ วายมติ, วีริยํ อารภติ, จิตฺตํ ปคฺคณฺหาติ ปทหติ. อิเมสํ โข ภิกฺขเว ปญฺจนฺนํ อุทฺธมฺ\-ภาคิยานํ สํโยชนานํ อภิญฺญาย ปริญฺญาย ปริกฺขยาย ปหานาย อิเม จตฺตาโร สมฺมปฺปธานา ภาเวตพฺพาติ. วิตฺถาเรตพฺพา. ทสมํ.}{โอฆวคฺโค ปญฺจโม.}
\csromancenter{\textbf{ตสฺสุทฺทานํ}}
\setlength{\csromangathapagewidth}{0pt}
\setlength{\csromangathawakwidth}{0pt}
\csromangathameasuregroup{โอโฆ โยโค อุปาทานํ, \\ กามคุณา นีวรณา,}{\csromangathabat{โอโฆ โยโค อุปาทานํ,}{คนฺถา อนุสเยน จ.} \\ \csromangathabat{กามคุณา นีวรณา,}{ขนฺธา โอรุทฺธมฺภาคิยาติ.} \\ สมฺมปฺปธานสํยุตฺตํ ปญฺจมํ.}
\csromangathameasurewak{สมฺมปฺปธานสํยุตฺตํ ปญฺจมํ.}
\csromangathasetleft{โอโฆ โยโค อุปาทานํ, \\ กามคุณา นีวรณา,}
\csromangathagroup{\csromangathastanza{\csromangathabat{โอโฆ โยโค อุปาทานํ,}{คนฺถา อนุสเยน จ.} \\ \csromangathabat{กามคุณา นีวรณา,}{ขนฺธา โอรุทฺธมฺภาคิยาติ.}}\csromangathastanza{\csromangathawak{สมฺมปฺปธานสํยุตฺตํ ปญฺจมํ.}}}

\csromanmatikaanchor{mtk.154}
\csromantocmarknopage{subparagraph}{6. พลสํยุตฺต}{mtk.154}
\bookmark[dest={mtk.154},level=5]{6. พลสํยุตฺต}
\csromanheader{6}{6. \textbf{พลสํยุตฺต}}
\chapterhead{1. คงฺคา\-เปยฺยาล\-วคฺค}
\csromanmatikaanchor{mtk.155}
\csromanmatikaanchor{mtk.156}
\csromantocmarknopage{subparagraph}{1. คงฺคาเปยฺยาลวคฺค}{mtk.155}
\bookmark[dest={mtk.155},level=5]{1. คงฺคาเปยฺยาลวคฺค}
\csromantocmarkref{subparagraph}{1-12. พลาทิสุตฺต}{mtk.156}
\bookmark[dest={mtk.156},level=5]{1-12. พลาทิสุตฺต}
\csromanheader{6}{1-12. พลาทิสุตฺต}
\proseitemwithclosermidruled{705-716}{\csromanfolio{218}ปญฺจิมานิ ภิกฺขเว พลานิ. กตมานิ ปญฺจ, สทฺธาพลํ วีริยพลํ สติพลํ สมาธิพลํ ปญฺญาพลํ. อิมานิ โข ภิกฺขเว ปญฺจ พลานีติ. เสยฺยถาปิ ภิกฺขเว คงฺคา นที ปาจีนนินฺนา ปาจีนโปณา ปาจีนปพฺภารา. เอวเมว โข ภิกฺขเว ภิกฺขุ ปญฺจ พลานิ ภาเวนฺโต ปญฺจพลานิ พหุลีกโรนฺโต นิพฺพานนินฺโน โหติ นิพฺพานโปโณ นิพฺพาน\-ปพฺภาโร. กถญฺจ ภิกฺขเว ภิกฺขุ ปญฺจ พลานิ ภาเวนฺโต ปญฺจ พลานิ พหุลีกโรนฺโต นิพฺพานนินฺโน โหติ นิพฺพานโปโณ นิพฺพาน\-ปพฺภาโร, อิธ ภิกฺขเว ภิกฺขุ สทฺธาพลํ ภาเวติ วิเวกนิสฺสิตํ วิราคนิสฺสิตํ นิโรธ\-นิสฺสิตํ โวสฺสคฺค\-ปริณามิํ. วีริยพลํ ฯเปฯ. สติพลํ. สมาธิพลํ. ปญฺญาพลํ ภาเวติ วิเวกนิสฺสิตํ วิราคนิสฺสิตํ นิโรธ\-นิสฺสิตํ โวสฺสคฺค\-ปริณามิํ. เอวํ โข ภิกฺขเว ภิกฺขุ ปญฺจ พลานิ ภาเวนฺโต ปญฺจ พลานิ พหุลีกโรนฺโต นิพฺพานนินฺโน โหติ นิพฺพานโปโณ นิพฺพาน\-ปพฺภาโรติ.  ทฺวาทสมํ.}{คงฺคา\-เปยฺยาล\-วคฺโค ปฐโม.}
\csromancenter{\textbf{ตสฺสุทฺทานํ}}
\setlength{\csromangathapagewidth}{0pt}
\setlength{\csromangathawakwidth}{0pt}
\csromangathameasuregroup{ฉ ปาจีนโต นินฺนา,}{\csromangathabat{ฉ ปาจีนโต นินฺนา,}{ฉ นินฺนา จ สมุทฺทโต.} \\ เทฺวเต ฉ ทฺวาทส โหนฺติ, วคฺโค เตน ปวุจฺจตีติ. \\ อปฺปมาทวคฺโค วิตฺถาเรตพฺโพ.}
\csromangathasetleft{ฉ ปาจีนโต นินฺนา,}
\csromangathagroup{\csromangathastanza{\csromangathabat{ฉ ปาจีนโต นินฺนา,}{ฉ นินฺนา จ สมุทฺทโต.} \\ เทฺวเต ฉ ทฺวาทส โหนฺติ, วคฺโค เตน ปวุจฺจตีติ. \\ อปฺปมาทวคฺโค วิตฺถาเรตพฺโพ.}}

\vspace{\tipitakaparskip}
\csromancenter{\textbf{ตสฺสุทฺทานํ}}
\setlength{\csromangathapagewidth}{0pt}
\setlength{\csromangathawakwidth}{0pt}
\csromangathameasuregroup{ตถาคตํ ปทํ กูฏํ, \\ ราชา จนฺทิมสูริยา,}{\csromangathabat{ตถาคตํ ปทํ กูฏํ,}{มูลํ สาเรน วสฺสิกํ.} \\ \csromangathabat{ราชา จนฺทิมสูริยา,}{วตฺเถน ทสมํ ปทนฺติ.}}
\csromangathasetleft{ตถาคตํ ปทํ กูฏํ, \\ ราชา จนฺทิมสูริยา,}
\csromangathagroup{\csromangathastanza{\csromangathabat{ตถาคตํ ปทํ กูฏํ,}{มูลํ สาเรน วสฺสิกํ.} \\ \csromangathabat{ราชา จนฺทิมสูริยา,}{วตฺเถน ทสมํ ปทนฺติ.}}}

\vspace{\tipitakaparskip}
\csromancenter{พลกรณีย\-วคฺโค วิตฺถาเรตพฺโพ.}
\csromancenter{\textbf{ตสฺสุทฺทานํ}}
\setlength{\csromangathapagewidth}{0pt}
\setlength{\csromangathawakwidth}{0pt}
\csromangathameasuregroup{พลํ พีชญฺจ นาโค จ,}{\csromangathabat{พลํ พีชญฺจ นาโค จ,}{รุกฺโข กุมฺเภน สูกิยา.} \\ อากาเสน จ เทฺว เมฆา, นาวา อาคนฺตุกา นทีติ. \\ เอสนาวคฺโค วิตฺถาเรตพฺโพ.}
\csromangathasetleft{พลํ พีชญฺจ นาโค จ,}
\csromangathagroup{\csromangathastanza{\csromanfolio{219}\csromangathabat{พลํ พีชญฺจ นาโค จ,}{รุกฺโข กุมฺเภน สูกิยา.} \\ อากาเสน จ เทฺว เมฆา, นาวา อาคนฺตุกา นทีติ. \\ เอสนาวคฺโค วิตฺถาเรตพฺโพ.}}

\vspace{\tipitakaparskip}
\csromancenter{\textbf{ตสฺสุทฺทานํ}}
\setlength{\csromangathapagewidth}{0pt}
\setlength{\csromangathawakwidth}{0pt}
\csromangathameasuregroup{เอสนา วิธา อาสโว, \\ ขิลํ มลญฺจ นีโฆ จ,}{\csromangathabat{เอสนา วิธา อาสโว,}{ภโว จ ทุกฺขตา ติสฺโส.} \\ \csromangathabat{ขิลํ มลญฺจ นีโฆ จ,}{เวทนา ตณฺหา ตสินา จาติ.}}
\csromangathasetleft{เอสนา วิธา อาสโว, \\ ขิลํ มลญฺจ นีโฆ จ,}
\csromangathagroup{\csromangathastanza{\csromangathabat{เอสนา วิธา อาสโว,}{ภโว จ ทุกฺขตา ติสฺโส.} \\ \csromangathabat{ขิลํ มลญฺจ นีโฆ จ,}{เวทนา ตณฺหา ตสินา จาติ.}}}

\chapterhead{5. \textbf{โอฆวคฺค}}
\csromanmatikaanchor{mtk.157}
\csromanmatikaanchor{mtk.158}
\csromantocmarknopage{subparagraph}{5. โอฆวคฺค}{mtk.157}
\bookmark[dest={mtk.157},level=5]{5. โอฆวคฺค}
\csromantocmarkref{subparagraph}{1-10. โอฆาทิสุตฺต}{mtk.158}
\bookmark[dest={mtk.158},level=5]{1-10. โอฆาทิสุตฺต}
\csromanheader{6}{1-10. โอฆาทิสุตฺต}
\proseitem{749-758}{ปญฺจิมานิ ภิกฺขเว อุทฺธมฺ\-ภาคิยานิ สํโยชนานิ. กตมานิ ปญฺจ, รูปราโค อรูปราโค มาโน อุทฺธจฺจํ อวิชฺชา, อิมานิ โข ภิกฺขเว ปญฺจุ\-ทฺธมฺภาคิยานิ สํโยชนานิ. อิเมสํ โข ภิกฺขเว ปญฺจนฺนํ อุทฺธมฺ\-ภาคิยานํ สํโยชนานํ อภิญฺญาย ปริญฺญาย ปริกฺขยาย ปหานาย ปญฺจ พลานิ ภาเวตพฺพานิ. กตมานิ ปญฺจ, อิธ ภิกฺขเว ภิกฺขุ สทฺธาพลํ ภาเวติ วิเวกนิสฺสิตํ วิราคนิสฺสิตํ นิโรธ\-นิสฺสิตํ โวสฺสคฺค\-ปริณามิํ. วีริยพลํ ฯเปฯ. สติพลํ สมาธิพลํ ปญฺญาพลํ ภาเวติ วิเวกนิสฺสิตํ วิราคนิสฺสิตํ นิโรธ\-นิสฺสิตํ โวสฺสคฺค\-ปริณามิํ. อิเมสํ โข ภิกฺขเว ปญฺจนฺนํ อุทฺธมฺ\-ภาคิยานํ สํโยชนานํ อภิญฺญาย ปริญฺญาย ปริกฺขยาย ปหานาย อิมานิ ปญฺจ พลานิ ภาเวตพฺพานีติ. (เอวํ วิตฺถาเรตพฺพา.)}

\chapterheadpageread{6. คงฺคา\-เปยฺยาล\-วคฺค}
\csromanmatikaanchor{mtk.159}
\csromanmatikaanchor{mtk.160}
\csromantocmarknopage{subparagraph}{6. คงฺคาเปยฺยาลวคฺค}{mtk.159}
\bookmark[dest={mtk.159},level=5]{6. คงฺคาเปยฺยาลวคฺค}
\csromantocmarkref{subparagraph}{1-12. ปาจีนาทิสุตฺต}{mtk.160}
\bookmark[dest={mtk.160},level=5]{1-12. ปาจีนาทิสุตฺต}
\csromanheader{6}{1-12. ปาจีนาทิสุตฺต}
\proseitem{759-770}{\csromanfolio{220}เสยฺยถาปิ ภิกฺขเว คงฺคา นที ปาจีนนินฺนา ปาจีนโปณา ปาจีนปพฺภารา. เอวเมว โข ภิกฺขเว ภิกฺขุ ปญฺจ พลานิ ภาเวนฺโต ปญฺจ พลานิ พหุลีกโรนฺโต นิพฺพานนินฺโน โหติ นิพฺพานโปโณ นิพฺพาน\-ปพฺภาโร. กถญฺจ ภิกฺขเว ภิกฺขุ ปญฺจ พลานิ ภาเวนฺโต ปญฺจ พลานิ พหุลีกโรนฺโต นิพฺพานนินฺโน โหติ นิพฺพานโปโณ นิพฺพาน\-ปพฺภาโร, อิธ ภิกฺขเว ภิกฺขุ สทฺธาพลํ ภาเวติ ราค\-วินย\-ปริโยสานํ โทส\-วินย\-ปริโยสานํ โมห\-วินย\-ปริโยสานํ. เอวํ โข ภิกฺขเว ภิกฺขุ ปญฺจ พลานิ ภาเวนฺโต ปญฺจ พลานิ พหุลีกโรนฺโต นิพฺพานนินฺโน โหติ นิพฺพานโปโณ นิพฺพาน\-ปพฺภาโรติ. วิตฺถาเรตพฺพา.  ทฺวาทสมํ. คงฺคา\-เปยฺยาล\-วคฺโค ฉฏฺโฐ.}

\vspace{\tipitakaparskip}
\csromancenter{\textbf{ตสฺสุทฺทานํ}}
\setlength{\csromangathapagewidth}{0pt}
\setlength{\csromangathawakwidth}{0pt}
\csromangathameasuregroup{ฉ ปาจีนโต นินฺนา, \\ เทฺวเต ฉ ทฺวาทส โหนฺติ,}{\csromangathabat{ฉ ปาจีนโต นินฺนา,}{ฉ นินฺนา จ สมุทฺทโต.} \\ \csromangathabat{เทฺวเต ฉ ทฺวาทส โหนฺติ,}{วคฺโค เตน ปวุจฺจตีติ.}}
\csromangathasetleft{ฉ ปาจีนโต นินฺนา, \\ เทฺวเต ฉ ทฺวาทส โหนฺติ,}
\csromangathagroup{\csromangathastanza{\csromangathabat{ฉ ปาจีนโต นินฺนา,}{ฉ นินฺนา จ สมุทฺทโต.} \\ \csromangathabat{เทฺวเต ฉ ทฺวาทส โหนฺติ,}{วคฺโค เตน ปวุจฺจตีติ.}}}

\vspace{\tipitakaparskip}
\csromancenter{(อปฺปมาท พลกรณียวคฺคา วิตฺถาเรตพฺพา.)}
\csromansectionrule
\chapterhead{9. \textbf{เอสนาวคฺค}}
\csromanmatikaanchor{mtk.161}
\csromanmatikaanchor{mtk.162}
\csromantocmarknopage{subparagraph}{9. เอสนาวคฺค}{mtk.161}
\bookmark[dest={mtk.161},level=5]{9. เอสนาวคฺค}
\csromantocmarkref{subparagraph}{1-12. เอสนาทิสุตฺต}{mtk.162}
\bookmark[dest={mtk.162},level=5]{1-12. เอสนาทิสุตฺต}
\csromanheader{6}{1-12. เอสนาทิสุตฺต}
\proseitemwithclosermidruled{792-802}{เอวํ เอสนา ปาฬิ วิตฺถาเรตพฺพา ราค\-วินย\-ปริโยสานํ โทส\-วินย\-ปริโยสานํ โมห\-วินย\-ปริโยสานํ.}{เอสนาวคฺโค นวโม.}
\csromancenter{\textbf{ตสฺสุทฺทานํ}}
\setlength{\csromangathapagewidth}{0pt}
\setlength{\csromangathawakwidth}{0pt}
\csromangathameasuregroup{เอสนา วิธา อาสโว, \\ ขิลํ มลญฺจ นีโฆ จ,}{\csromangathabat{เอสนา วิธา อาสโว,}{ภโว จ ทุกฺขตา ติสฺโส.} \\ \csromangathabat{ขิลํ มลญฺจ นีโฆ จ,}{เวทนา ตณฺหา ตสินา จาติ.}}
\csromangathasetleft{เอสนา วิธา อาสโว, \\ ขิลํ มลญฺจ นีโฆ จ,}
\csromangathagroup{\csromangathastanza{\csromangathabat{เอสนา วิธา อาสโว,}{ภโว จ ทุกฺขตา ติสฺโส.} \\ \csromangathabat{ขิลํ มลญฺจ นีโฆ จ,}{เวทนา ตณฺหา ตสินา จาติ.}}}

\chapterheadpageread{10. \textbf{โอฆวคฺค}}
\csromanmatikaanchor{mtk.163}
\csromanmatikaanchor{mtk.164}
\csromantocmarknopage{subparagraph}{10. โอฆวคฺค}{mtk.163}
\bookmark[dest={mtk.163},level=5]{10. โอฆวคฺค}
\csromantocmarkref{subparagraph}{1-10. โอฆาทิสุตฺต}{mtk.164}
\bookmark[dest={mtk.164},level=5]{1-10. โอฆาทิสุตฺต}
\csromanheader{6}{1-10. โอฆาทิสุตฺต}
\proseitemwithclosermidruled{803-821}{\csromanfolio{221}ปญฺจิมานิ ภิกฺขเว อุทฺธมฺ\-ภาคิยานิ สํโยชนานิ. กตมานิ ปญฺจ, รูปราโค อรูปราโค มาโน อุทฺธจฺจํ อวิชฺชา, อิมานิ โข ภิกฺขเว ปญฺจุ\-ทฺธมฺภาคิยานิ สํโยชนานิ. อิเมสํ โข ภิกฺขเว ปญฺจนฺนํ อุทฺธมฺ\-ภาคิยานํ สํโยชนานํ อภิญฺญาย ปริญฺญาย ปริกฺขยาย ปหานาย ปญฺจ พลานิ ภาเวตพฺพานิ. กตมานิ ปญฺจ, อิธ ภิกฺขเว ภิกฺขุ สทฺธาพลํ ภาเวติ ฯเปฯ ปญฺญาพลํ ภาเวติ ราค\-วินย\-ปริโยสานํ โทส\-วินย\-ปริโยสานํ โมห\-วินย\-ปริโยสานํ. อิเมสํ โข ภิกฺขเว ปญฺจนฺนํ อุทฺธมฺ\-ภาคิยานํ สํโยชนานํ อภิญฺญาย ปริญฺญาย ปริกฺขยาย ปหานาย อิมานิ ปญฺจ พลานิ ภาเวตพฺพานีติ.  ทสมํ.}{โอฆวคฺโค ทสโม.}
\csromancenter{\textbf{ตสฺสุทฺทานํ}}
\setlength{\csromangathapagewidth}{0pt}
\setlength{\csromangathawakwidth}{0pt}
\csromangathameasuregroup{โอโฆ โยโค อุปาทานํ, \\ กามคุณา นีวรณา,}{\csromangathabat{โอโฆ โยโค อุปาทานํ,}{คนฺถา อนุสเยน จ.} \\ \csromangathabat{กามคุณา นีวรณา,}{ขนฺธา โอรุทฺธมฺภาคิยาติ.} \\ พลสํยุตฺตํ ฉฏฺฐํ.}
\csromangathameasurewak{พลสํยุตฺตํ ฉฏฺฐํ.}
\csromangathasetleft{โอโฆ โยโค อุปาทานํ, \\ กามคุณา นีวรณา,}
\csromangathagroup{\csromangathastanza{\csromangathabat{โอโฆ โยโค อุปาทานํ,}{คนฺถา อนุสเยน จ.} \\ \csromangathabat{กามคุณา นีวรณา,}{ขนฺธา โอรุทฺธมฺภาคิยาติ.}}\csromangathastanza{\csromangathawak{พลสํยุตฺตํ ฉฏฺฐํ.}}}

\csromanmatikaanchor{mtk.165}
\csromanmatikaanchor{mtk.166}
\csromantocmarknopage{subparagraph}{7. อิทฺธิปาทสํยุตฺต}{mtk.165}
\bookmark[dest={mtk.165},level=5]{7. อิทฺธิปาทสํยุตฺต}
\csromantocmarknopage{subparagraph}{1. จาปาลวคฺค}{mtk.166}
\bookmark[dest={mtk.166},level=5]{1. จาปาลวคฺค}
\csromanheader{6}{1. \textbf{จาปาลวคฺค}}
\csromanmatikaanchor{mtk.167}
\csromantocmarkref{subparagraph}{1. อปารสุตฺต}{mtk.167}
\bookmark[dest={mtk.167},level=5]{1. อปารสุตฺต}
\csromanheader{6}{1. \textbf{อปารสุตฺต}}
\csromanmatikaanchor{mtk.168}
\csromantocmarkref{subparagraph}{2-4. วิรทฺธสุตฺตาทีนิ}{mtk.168}
\bookmark[dest={mtk.168},level=5]{2-4. วิรทฺธสุตฺตาทีนิ}
\proseitem{813}{\csromanfolio{222}จตฺตาโรเม ภิกฺขเว อิทฺธิปาทา ภาวิตา พหุลีกตา อปารา ปารํ คมนาย สํวตฺตนฺติ. กตเม จตฺตาโร, อิธ ภิกฺขเว ภิกฺขุ ฉนฺท\-สมาธิ\-ปฺปธาน\-สงฺขาร\-สมนฺนาคตํ อิทฺธิปาทํ ภาเวติ, วีริย\-สมาธิ\-ปฺปธาน\-สงฺขาร\-สมนฺนาคตํ อิทฺธิปาทํ ภาเวติ, จิตฺต\-สมาธิ\-ปฺปธาน\-สงฺขาร\-สมนฺนาคตํ อิทฺธิปาทํ ภาเวติ, วีมํสา\-สมาธิ\-ปฺปธาน\-สงฺขาร\-สมนฺนาคตํ อิทฺธิปาทํ ภาเวติ. อิเม โข ภิกฺขเว จตฺตาโร อิทฺธิปาทา ภาวิตา พหุลีกตา อปารา ปารํ คมนาย สํวตฺตนฺตีติ.  ปฐมํ.}

\csromanheader{2}{2. \textbf{วิรทฺธสุตฺต}}
\proseitem{814}{เยสํ เกสญฺจิ ภิกฺขเว จตฺตาโร อิทฺธิปาทา วิรทฺธา, วิรทฺโธ เตสํ อริโย มคฺโค สมฺมา ทุกฺข\-กฺขย\-คามี. เยสํ เกสญฺจิ ภิกฺขเว จตฺตาโร อิทฺธิปาทา อารทฺธา, อารทฺโธ เตสํ อริโย มคฺโค สมฺมา ทุกฺข\-กฺขย\-คามี. กตเม จตฺตาโร, อิธ ภิกฺขเว ภิกฺขุ ฉนฺท\-สมาธิ\-ปฺปธาน\-สงฺขาร\-สมนฺนาคตํ อิทฺธิปาทํ ภาเวติ, วีริยสมาธิ ฯเปฯ จิตฺตสมาธิ ฯเปฯ วีมํสา\-สมาธิ\-ปฺปธาน\-สงฺขาร\-สมนฺนาคตํ อิทฺธิปาทํ ภาเวติ. เยสํ เกสญฺจิ ภิกฺขเว อิเม จตฺตาโร อิทฺธิปาทา วิรทฺธา, วิรทฺโธ เตสํ อริโย มคฺโค สมฺมา ทุกฺข\-กฺขย\-คามี. เยสํ เกสญฺจิ ภิกฺขเว อิเม จตฺตาโร อิทฺธิปาทา อารทฺธา, อารทฺโธ เตสํ อริโย มคฺโค สมฺมา ทุกฺข\-กฺขย\-คามีติ.  ทุติยํ.}

\csromanheader{2}{3. \textbf{อริยสุตฺต}}
\proseitem{815}{จตฺตาโรเม ภิกฺขเว อิทฺธิปาทา ภาวิตา พหุลีกตา อริยา นิยฺยานิกา นิยฺยนฺติ ตกฺกรสฺส สมฺมา ทุกฺขกฺขยาย. กตเม จตฺตาโร, อิธ ภิกฺขเว ภิกฺขุ ฉนฺท\-สมาธิ\-ปฺปธาน\-สงฺขาร\-สมนฺนาคตํ อิทฺธิปาทํ ภาเวติ, วีริยสมาธิ ฯเปฯ จิตฺตสมาธิ ฯเปฯ วีมํสา\-สมาธิ\-ปฺปธาน\-สงฺขาร\-สมนฺนาคตํ อิทฺธิปาทํ ภาเวติ. อิเม โข ภิกฺขเว จตฺตาโร อิทฺธิปาทา}

\prose{\csromanfolio{223}ภาวิตา พหุลีกตา อริยา นิยฺยานิกา นิยฺยนฺติ ตกฺกรสฺส สมฺมา ทุกฺขกฺขยายาติ.  ตติยํ.}

\csromanheader{2}{4. \textbf{นิพฺพิทาสุตฺต}}
\proseitem{816}{จตฺตาโรเม ภิกฺขเว อิทฺธิปาทา ภาวิตา พหุลีกตา เอกนฺต\-นิพฺพิทาย วิราคาย นิโรธาย อุปสมาย อภิญฺญาย สมฺโพธาย นิพฺพานาย สํวตฺตนฺติ. กตเม จตฺตาโร, อิธ ภิกฺขเว ภิกฺขุ ฉนฺท\-สมาธิ\-ปฺปธาน\-สงฺขาร\-สมนฺนาคตํ อิทฺธิปาทํ ภาเวติ, วีริยสมาธิ ฯเปฯ จิตฺตสมาธิ ฯเปฯ วีมํสา\-สมาธิ\-ปฺปธาน\-สงฺขาร\-สมนฺนาคตํ อิทฺธิปาทํ ภาเวติ. อิเม โข ภิกฺขเว จตฺตาโร อิทฺธิปาทา ภาวิตา พหุลีกตา เอกนฺต\-นิพฺพิทาย วิราคาย นิโรธาย อุปสมาย อภิญฺญาย สมฺโพธาย นิพฺพานาย สํวตฺตนฺตีติ.  จตุตฺถํ.}

\csromanmatikaanchor{mtk.169}
\csromantocmarkref{subparagraph}{5. อิทฺธิปเทสสุตฺต}{mtk.169}
\bookmark[dest={mtk.169},level=5]{5. อิทฺธิปเทสสุตฺต}
\csromanheader{6}{5. อิทฺธิปเทส\-สุตฺต}
\proseitem{817}{เย หิ เกจิ ภิกฺขเว อตีตมทฺธานํ สมณา วา พฺราหฺมณา วา อิทฺธิปเทสํ อภินิ\-ปฺผาเทสุํ, สพฺเพ เต จตุนฺนํ อิทฺธิปาทานํ ภาวิตตฺตา พหุลีกตตฺตา. เย หิ เกจิ ภิกฺขเว อนาคตมทฺธานํ สมณา วา พฺราหฺมณา วา อิทฺธิปเทสํ อภินิปฺผาเทสฺสนฺติ, สพฺเพ เต จตุนฺนํ อิทฺธิปาทานํ ภาวิตตฺตา พหุลีกตตฺตา. เย หิ เกจิ ภิกฺขเว เอตรหิ สมณา วา พฺราหฺมณา วา อิทฺธิปเทสํ อภินิ\-ปฺผาเทนฺติ, สพฺเพ เต จตุนฺนํ อิทฺธิปาทานํ ภาวิตตฺตา พหุลีกตตฺตา.}

\prose{กตเมสํ จตุนฺนํ, อิธ ภิกฺขเว ภิกฺขุ ฉนฺท\-สมาธิ\-ปฺปธาน\-สงฺขาร\-สมนฺนาคตํ อิทฺธิปาทํ ภาเวติ, วีริยสมาธิ ฯเปฯ จิตฺตสมาธิ ฯเปฯ วีมํสา\-สมาธิ\-ปฺปธาน\-สงฺขาร\-สมนฺนาคตํ อิทฺธิปาทํ ภาเวติ. เย หิ เกจิ ภิกฺขเว อตีตมทฺธานํ สมณา วา พฺราหฺมณา วา อิทฺธิปเทสํ อภินิ\-ปฺผาเทสุํ, สพฺเพ เต อิเมสํเยว จตุนฺนํ อิทฺธิปาทานํ ภาวิตตฺตา พหุลีกตตฺตา. เย หิ เกจิ ภิกฺขเว อนาคตมทฺธานํ สมณา วา พฺราหฺมณา วา อิทฺธิปเทสํ อภินิปฺผาเทสฺสนฺติ, สพฺเพ เต อิเมสํเยว จตุนฺนํ อิทฺธิปาทานํ ภาวิตตฺตา พหุลีกตตฺตา. เย หิ เกจิ ภิกฺขเว เอตรหิ สมณา วา พฺราหฺมณา วา อิทฺธิปเทสํ อภินิ\-ปฺผาเทนฺติ, สพฺเพ เต อิเมสํเยว จตุนฺนํ อิทฺธิปาทานํ ภาวิตตฺตา พหุลีกตตฺตาติ.  ปญฺจมํ.}

\csromanmatikaanchor{mtk.170}
\csromantocmarkref{subparagraph}{6. สมตฺตสุตฺต}{mtk.170}
\bookmark[dest={mtk.170},level=5]{6. สมตฺตสุตฺต}
\csromanheader{6}{6. \textbf{สมตฺตสุตฺต}}
\csromanmatikaanchor{mtk.171}
\csromantocmarkref{subparagraph}{7-9. ภิกฺขุสุตฺตาทีนิ}{mtk.171}
\bookmark[dest={mtk.171},level=5]{7-9. ภิกฺขุสุตฺตาทีนิ}
\proseitem{818}{\csromanfolio{224}เย หิ เกจิ ภิกฺขเว อตีตมทฺธานํ สมณา วา พฺราหฺมณา วา สมตฺตํ อิทฺธิํ อภินิ\-ปฺผาเทสุํ, สพฺเพ เต จตุนฺนํ อิทฺธิปาทานํ ภาวิตตฺตา พหุลีกตตฺตา. เย หิ เกจิ ภิกฺขเว อนาคตมทฺธานํ สมณา วา พฺราหฺมณา วา สมตฺตํ อิทฺธิํ อภินิปฺผาเทสฺสนฺติ, สพฺเพ เต จตุนฺนํ อิทฺธิปาทานํ ภาวิตตฺตา พหุลีกตตฺตา. เย หิ เกจิ ภิกฺขเว เอตรหิ สมณา วา พฺราหฺมณา วา สมตฺตํ อิทฺธิํ อภินิ\-ปฺผาเทนฺติ, สพฺเพ เต จตุนฺนํ อิทฺธิปาทานํ ภาวิตตฺตา พหุลีกตตฺตา.}

\prose{กตเมสํ จตุนฺนํ, อิธ ภิกฺขเว ภิกฺขุ ฉนฺท\-สมาธิ\-ปฺปธาน\-สงฺขาร\-สมนฺนาคตํ อิทฺธิปาทํ ภาเวติ, วีริยสมาธิ ฯเปฯ จิตฺตสมาธิ ฯเปฯ วีมํสา\-สมาธิ\-ปฺปธาน\-สงฺขาร\-สมนฺนาคตํ อิทฺธิปาทํ ภาเวติ. เย หิ เกจิ ภิกฺขเว อตีตมทฺธานํ สมณา วา พฺราหฺมณา วา สมตฺตํ อิทฺธิํ อภินิ\-ปฺผาเทสุํ, สพฺเพ เต อิเมสํเยว จตุนฺนํ อิทฺธิปาทานํ ภาวิตตฺตา พหุลีกตตฺตา. เย หิ เกจิ ภิกฺขเว อนาคตมทฺธานํ สมณา วา พฺราหฺมณา วา สมตฺตํ อิทฺธิํ อภินิปฺผาเทสฺสนฺติ, สพฺเพ เต อิเมสํเยว จตุนฺนํ อิทฺธิปาทานํ ภาวิตตฺตา พหุลีกตตฺตา. เย หิ เกจิ ภิกฺขเว เอตรหิ สมณา วา พฺราหฺมณา วา สมตฺตํ อิทฺธิํ อภินิ\-ปฺผาเทนฺติ, สพฺเพ เต อิเมสํเยว จตุนฺนํ อิทฺธิปาทานํ ภาวิตตฺตา พหุลีกตตฺตาติ.  ฉฏฺฐํ.}

\csromanheader{2}{7. \textbf{ภิกฺขุสุตฺต}}
\proseitem{819}{เย หิ เกจิ ภิกฺขเว อตีตมทฺธานํ ภิกฺขู อาสวานํ ขยา อนาสวํ เจโตวิมุตฺติํ ปญฺญาวิมุตฺติํ ทิฏฺเฐว ธมฺเม สยํ อภิญฺญา สจฺฉิกตฺวา อุปสมฺปชฺช วิหริํสุ, สพฺเพ เต จตุนฺนํ อิทฺธิปาทานํ ภาวิตตฺตา พหุลีกตตฺตา. เย หิ เกจิ ภิกฺขเว อนาคตมทฺธานํ ภิกฺขู อาสวานํ ขยา อนาสวํ เจโตวิมุตฺติํ ปญฺญาวิมุตฺติํ ทิฏฺเฐว ธมฺเม สยํ อภิญฺญา สจฺฉิกตฺวา อุปสมฺปชฺช วิหริสฺสนฺติ, สพฺเพ เต จตุนฺนํ อิทฺธิปาทานํ ภาวิตตฺตา พหุลีกตตฺตา. เย หิ เกจิ ภิกฺขเว เอตรหิ ภิกฺขู อาสวานํ ขยา อนาสวํ เจโตวิมุตฺติํ ปญฺญาวิมุตฺติํ ทิฏฺเฐว ธมฺเม สยํ อภิญฺญา สจฺฉิกตฺวา อุปสมฺปชฺช วิหรนฺติ, สพฺเพ เต จตุนฺนํ อิทฺธิปาทานํ ภาวิตตฺตา พหุลีกตตฺตา.}

\prose{\csromanfolio{225}กตเมสํ จตุนฺนํ, อิธ ภิกฺขเว ภิกฺขุ ฉนฺท\-สมาธิ\-ปฺปธาน\-สงฺขาร\-สมนฺนาคตํ อิทฺธิปาทํ ภาเวติ, วีริยสมาธิ ฯเปฯ จิตฺตสมาธิ. วีมํสา\-สมาธิ\-ปฺปธาน\-สงฺขาร\-สมนฺนาคตํ อิทฺธิปาทํ ภาเวติ. เย หิ เกจิ ภิกฺขเว อตีตมทฺธานํ ภิกฺขู อาสวานํ ขยา อนาสวํ เจโตวิมุตฺติํ ปญฺญาวิมุตฺติํ ทิฏฺเฐว ธมฺเม สยํ อภิญฺญา สจฺฉิกตฺวา อุปสมฺปชฺช วิหริํสุ, สพฺเพ เต อิเมสํเยว จตุนฺนํ อิทฺธิปาทานํ ภาวิตตฺตา พหุลีกตตฺตา. เย หิ เกจิ ภิกฺขเว อนาคตมทฺธานํ ภิกฺขู อาสวานํ ขยา อนาสวํ เจโตวิมุตฺติํ ปญฺญาวิมุตฺติํ ทิฏฺเฐว ธมฺเม สยํ อภิญฺญา สจฺฉิกตฺวา อุปสมฺปชฺช วิหริสฺสนฺติ, สพฺเพ เต อิเมสํเยว จตุนฺนํ อิทฺธิปาทานํ ภาวิตตฺตา พหุลีกตตฺตา. เย หิ เกจิ ภิกฺขเว เอตรหิ ภิกฺขู อาสวานํ ขยา อนาสวํ เจโตวิมุตฺติํ ปญฺญาวิมุตฺติํ ทิฏฺเฐว ธมฺเม สยํ อภิญฺญา สจฺฉิกตฺวา อุปสมฺปชฺช วิหรนฺติ, สพฺเพ เต อิเมสํเยว จตุนฺนํ อิทฺธิปาทานํ ภาวิตตฺตา พหุลีกตตฺตาติ.  สตฺตมํ.}

\csromanheader{2}{8. \textbf{พุทฺธสุตฺต}}
\proseitem{820}{จตฺตาโรเม ภิกฺขเว อิทฺธิปาทา. กตเม จตฺตาโร, อิธ ภิกฺขเว ภิกฺขุ ฉนฺท\-สมาธิ\-ปฺปธาน\-สงฺขาร\-สมนฺนาคตํ อิทฺธิปาทํ ภาเวติ, วีริยสมาธิ ฯเปฯ จิตฺตสมาธิ. วีมํสา\-สมาธิ\-ปฺปธาน\-สงฺขาร\-สมนฺนาคตํ อิทฺธิปาทํ ภาเวติ. อิเม โข ภิกฺขเว จตฺตาโร อิทฺธิปาทา, อิเมสํ โข ภิกฺขเว จตุนฺนํ อิทฺธิปาทานํ ภาวิตตฺตา พหุลีกตตฺตา ตถาคโต อรหํ สมฺมา\-สมฺพุทฺโธติ วุจฺจตีติ.  อฏฺฐมํ.}

\csromanheader{2}{9. \textbf{ญาณสุตฺต}}
\proseitem{821}{“อยํ ฉนฺท\-สมาธิ\-ปฺปธาน\-สงฺขาร\-สมนฺนาคโต อิทฺธิปาโท”ติ เม ภิกฺขเว ปุพฺเพ อนนุสฺสุเตสุ ธมฺเมสุ จกฺขุํ อุทปาทิ, ญาณํ อุทปาทิ, ปญฺญา อุทปาทิ, วิชฺชา อุทปาทิ, อาโลโก อุทปาทิ. “โส โข ปนายํ ฉนฺท\-สมาธิ\-ปฺปธาน\-สงฺขาร\-สมนฺนาคโต อิทฺธิปาโท ภาเวตพฺโพ”ติ เม ภิกฺขเว. “ภาวิโต”ติ เม ภิกฺขเว ปุพฺเพ อนนุสฺสุเตสุ ธมฺเมสุ จกฺขุํ อุทปาทิ, ญาณํ อุทปาทิ, ปญฺญา อุทปาทิ, วิชฺชา อุทปาทิ, อาโลโก อุทปาทิ.}

\prose{\csromanfolio{226}“อยํ วีริย\-สมาธิ\-ปฺปธาน\-สงฺขาร\-สมนฺนาคโต อิทฺธิปาโท”ติ เม ภิกฺขเว ปุพฺเพ อนนุสฺสุเตสุ ธมฺเมสุ จกฺขุํ อุทปาทิ, ญาณํ อุทปาทิ, ปญฺญา อุทปาทิ, วิชฺชา อุทปาทิ, อาโลโก อุทปาทิ. “โส โข ปนายํ วีริย\-สมาธิ\-ปฺปธาน\-สงฺขาร\-สมนฺนาคโต อิทฺธิปาโท ภาเวตพฺโพ”ติ เม ภิกฺขเว. “ภาวิโต”ติ เม ภิกฺขเว ปุพฺเพ อนนุสฺสุเตสุ ธมฺเมสุ จกฺขุํ อุทปาทิ, ญาณํ อุทปาทิ, ปญฺญา อุทปาทิ, วิชฺชา อุทปาทิ, อาโลโก อุทปาทิ.}

\prose{“อยํ จิตฺต\-สมาธิ\-ปฺปธาน\-สงฺขาร\-สมนฺนาคโต อิทฺธิปาโท”ติ เม ภิกฺขเว ปุพฺเพ อนนุสฺสุเตสุ ธมฺเมสุ จกฺขุํ อุทปาทิ, ญาณํ อุทปาทิ, ปญฺญา อุทปาทิ, วิชฺชา อุทปาทิ, อาโลโก อุทปาทิ. “โส โข ปนายํ จิตฺต\-สมาธิ\-ปฺปธาน\-สงฺขาร\-สมนฺนาคโต อิทฺธิปาโท ภาเวตพฺโพ”ติ เม ภิกฺขเว. “ภาวิโต”ติ เม ภิกฺขเว ปุพฺเพ อนนุสฺสุเตสุ ธมฺเมสุ จกฺขุํ อุทปาทิ, ญาณํ อุทปาทิ, ปญฺญา อุทปาทิ, วิชฺชา อุทปาทิ, อาโลโก อุทปาทิ.}

\prose{“อยํ วีมํสา\-สมาธิ\-ปฺปธาน\-สงฺขาร\-สมนฺนาคโต อิทฺธิปาโท”ติ เม ภิกฺขเว ปุพฺเพ อนนุสฺสุเตสุ ธมฺเมสุ จกฺขุํ อุทปาทิ, ญาณํ อุทปาทิ, ปญฺญา อุทปาทิ, วิชฺชา อุทปาทิ, อาโลโก อุทปาทิ. “โส โข ปนายํ วีมํสา\-สมาธิ\-ปฺปธาน\-สงฺขาร\-สมนฺนาคโต อิทฺธิปาโท ภาเวตพฺโพ”ติ เม ภิกฺขเว. “ภาวิโต”ติ เม ภิกฺขเว ปุพฺเพ อนนุสฺสุเตสุ ธมฺเมสุ จกฺขุํ อุทปาทิ, ญาณํ อุทปาทิ, ปญฺญา อุทปาทิ, วิชฺชา อุทปาทิ, อาโลโก อุทปาทีติ.  นวมํ.}

\csromanmatikaanchor{mtk.172}
\csromantocmarkref{subparagraph}{10. เจติยสุตฺต}{mtk.172}
\bookmark[dest={mtk.172},level=5]{10. เจติยสุตฺต}
\csromanheader{6}{10. \textbf{เจติยสุตฺต}}
\proseitem{822}{เอวํ เม สุตํ\csromandash{}เอกํ สมยํ ภควา เวสาลิยํ วิหรติ มหาวเน กูฏาคาร\-สาลายํ. อถ โข ภควา ปุพฺพณฺหสมยํ นิวาเสตฺวา ปตฺต\-จีวรมา\-ทาย เวสาลิํ ปิณฺฑาย ปาวิสิ, เวสาลิยํ ปิณฺฑาย จริตฺวา ปจฺฉาภตฺตํ ปิณฺฑปาต\-ปฏิกฺกนฺโต อายสฺมนฺตํ อานนฺทํ อามนฺเตสิ “คณฺหาหิ อานนฺท นิสีทนํ, เยน จาปาลํ เจติยํ\csromansharedfootnote{82a7a0f63ea9c67a}{ปาวาลเจติยํ (สฺยา, กํ)} เตนุปสงฺกมิสฺสาม ทิวาวิหารายา”ติ. “เอวํ ภนฺเต”ติ โข อายสฺมา อานนฺโท ภควโต ปฏิสฺสุตฺวา นิสีทนํ อาทาย ภควนฺตํ ปิฏฺฐิโต ปิฏฺฐิโต อนุพนฺธิ. อถ โข ภควา เยน จาปาลํ เจติยํ เตนุปสงฺกมิ,}

\prose{\csromanfolio{227}อุปสงฺกมิตฺวา ปญฺญตฺเต อาสเน นิสีทิ. อายสฺมาปิ โข อานนฺโท ภควนฺตํ อภิวาเทตฺวา เอกมนฺตํ นิสีทิ, เอกมนฺตํ นิสินฺนํ โข อายสฺมนฺตํ อานนฺทํ ภควา เอตทโวจ\csromandash{}}

\prose{รมณียา อานนฺท เวสาลี, รมณียํ อุเทนํ เจติยํ, รมณียํ โคตมกํ เจติยํ, รมณียํ สตฺตมฺพํ เจติยํ, รมณียํ พหุปุตฺตํ เจติยํ\csromansharedfootnote{ad8f96adf804bf5c}{พหุปุตฺตกเจติยํ (สฺยา, กํ, อิ, ก)}, รมณียํ สารนฺททํ เจติยํ\csromansharedfootnote{c04ce288446d6038}{อานนฺทเจติยํ (ก), สานนฺทรํ (ก)}, รมณียํ จาปาลํ เจติยํ. ยสฺส กสฺสจิ อานนฺท จตฺตาโร อิทฺธิปาทา ภาวิตา พหุลีกตา ยานีกตา วตฺถุกตา อนุฏฺฐิตา ปริจิตา สุสมารทฺธา, โส อากงฺขมาโน กปฺปํ วา ติฏฺเฐยฺย กปฺปาวเสสํ วา, ตถาคตสฺส โข อานนฺท จตฺตาโร อิทฺธิปาทา ภาวิตา พหุลีกตา ยานีกตา วตฺถุกตา อนุฏฺฐิตา ปริจิตา สุสมารทฺธา, อากงฺขมาโน อานนฺท ตถาคโต กปฺปํ วา ติฏฺเฐยฺย กปฺปาวเสสํ วาติ.}

\prose{เอวมฺปิ โข อายสฺมา อานนฺโท ภควตา โอฬาริเก นิมิตฺเต กยิรมาเน โอฬาริเก โอภาเส กยิรมาเน นาสกฺขิ ปฏิวิชฺฌิตุํ, น ภควนฺตํ ยาจิ “ติฏฺฐตุ ภนฺเต ภควา กปฺปํ, ติฏฺฐตุ สุคโต กปฺปํ\csromansharedfootnote{d514043a049828af}{สุคโต กปฺปาวเสสํ (อิ, ก) ที 2. 86 ปิฏฺเฐปิ.} พหุชนหิตาย พหุชน\-สุขาย โลกานุกมฺปาย อตฺถาย หิตาย สุขาย เทวมนุสฺสานนฺ”ติ, ยถา ตํ มาเรน ปริยุฏฺฐิต\-จิตฺโต.}

\prose{ทุติยมฺปิ โข ภควา ฯเปฯ. ตติยมฺปิ โข ภควา อายสฺมนฺตํ อานนฺทํ อามนฺเตสิ “รมณียา อานนฺท เวสาลี, รมณียํ อุเทนํ เจติยํ, รมณียํ โคตมกํ เจติยํ, รมณียํ สตฺตพฺพํ เจติยํ, รมณียํ พหุปุตฺตํ เจติยํ, รมณียํ สารนฺททํ เจติยํ, รมณียํ จาปาลํ เจติยํ. ยสฺส กสฺสจิ อานนฺท จตฺตาโร อิทฺธิปาทา ภาวิตา พหุลีกตา ยานีกตา วตฺถุกตา อนุฏฺฐิตา ปริจิตา สุสมารทฺธา, โส อากงฺขมาโน กปฺปํ วา ติฏฺเฐยฺย กปฺปาวเสสํ วา, ตถาคตสฺส โข อานนฺท จตฺตาโร อิทฺธิปาทา ภาวิตา พหุลีกตา ยานีกตา วตฺถุกตา อนุฏฺฐิตา ปริจิตา สุสมารทฺธา, อากงฺขมาโน อานนฺท ตถาคโต กปฺปํ วา ติฏฺเฐยฺย กปฺปาวเสสํ วา”ติ.}

\prose{\csromanfolio{228}เอวมฺปิ โข อายสฺมา อานนฺโท ภควตา โอฬาริเก นิมิตฺเต กยิรมาเน โอฬาริเก โอภาเส กยิรมาเน นาสกฺขิ ปฏิวิชฺฌิตุํ, น ภควนฺตํ ยาจิ “ติฏฺฐตุ ภนฺเต ภควา กปฺปํ, ติฏฺฐตุ สุคโต กปฺปํ พหุชนหิตาย พหุชน\-สุขาย โลกานุกมฺปาย อตฺถาย หิตาย สุขาย เทวมนุสฺสานนฺ”ติ, ยถา ตํ มาเรน ปริยุฏฺฐิต\-จิตฺโต.}

\prose{อถ โข ภควา อายสฺมนฺตํ อานนฺทํ อามนฺเตสิ “คจฺฉ โข ตฺวํ อานนฺท, ยสฺสทานิ กาลํ มญฺญสี”ติ. “เอวํ ภนฺเต”ติ โข อายสฺมา อานนฺโท ภควโต ปฏิสฺสุตฺวา อุฏฺฐายาสนา ภควนฺตํ อภิวาเทตฺวา ปทกฺขิณํ กตฺวา อวิทูเร อญฺญตรสฺมิํ รุกฺขมูเล นิสีทิ. อถ โข มาโร ปาปิมา อจิรปกฺกนฺเต อายสฺมนฺเต อานนฺเท เยน ภควา เตนุปสงฺกมิ, อุปสงฺกมิตฺวา ภควนฺตํ เอตทโวจ “ปรินิพฺพาตุ\-ทานิ ภนฺเต ภควา, ปรินิพฺพาตุ\-ทานิ สุคโต\csromansharedfootnote{cfabaac9fcb09f92}{ปรินิพฺพาตุ สุคโต (สี, สฺยา, กํ) เอวมุปริปิ.}, ปรินิพฺพาน\-กาโลทานิ ภนฺเต ภควโต. ภาสิตา โข ปเนสา ภนฺเต ภควตา วาจา “น ตาวาหํ ปาปิม ปรินิพฺพายิสฺสามิ, ยาว เม ภิกฺขู น สาวกา ภวิสฺสนฺติ วิยตฺตา วินีตา วิสารทา พหุสฺสุตา ธมฺมธรา ธมฺมา\-นุธมฺม\-ปฺปฏิปนฺนา สามีจิ\-ปฺปฏิปนฺนา อนุธมฺม\-จาริโน สกํ อาจริยกํ อุคฺคเหตฺวา อาจิกฺขิสฺสนฺติ เทเสสฺสนฺติ ปญฺญเปสฺสนฺติ ปฏฺฐเปสฺสนฺติ วิวริสฺสนฺติ วิภชิสฺสนฺติ อุตฺตานีกริสฺสนฺติ\csromansharedfootnote{908a4c4396e358da}{อุตฺตานิํกริสฺสนฺติ (ก), อุตฺตานิกริสฺสนฺติ (อิ)}, อุปฺปนฺนํ ปรปฺปวาทํ สหธมฺเมน สุนิคฺคหิตํ นิคฺคเหตฺวา สปฺปาฏิหาริยํ ธมฺมํ เทเสสฺสนฺตี”ติ.}

\prose{สนฺติ โข ปน ภนฺเต เอตรหิ ภิกฺขู ภควโต สาวกา วิยตฺตา วินีตา วิสารทา พหุสฺสุตา ธมฺมธรา ธมฺมา\-นุธมฺม\-ปฺปฏิปนฺนา สามีจิ\-ปฺปฏิปนฺนา อนุธมฺม\-จาริโน สกํ อาจริยกํ อุคฺคเหตฺวา อาจิกฺขนฺติ เทเสนฺติ ปญฺญเปนฺติ ปฏฺฐเปนฺติ วิวรนฺติ วิภชนฺติ อุตฺตานีกโรนฺติ, อุปฺปนฺนํ ปรปฺปวาทํ สหธมฺเมน สุนิคฺคหิตํ นิคฺคเหตฺวา สปฺปาฏิหาริยํ ธมฺมํ เทเสนฺติ. ปรินิพฺพาตุ\-ทานิ ภนฺเต ภควา, ปรินิพฺพาตุ\-ทานิ สุคโต, ปรินิพฺพาน\-กาโลทานิ ภนฺเต ภควโต.}

\prose{ภาสิตา โข ปเนสา ภนฺเต ภควตา วาจา “น ตาวาหํ ปาปิม ปรินิพฺพายิสฺสามิ, ยาว เม ภิกฺขุนิโย น สาวิกา ภวิสฺสนฺติ วิยตฺตา}

\prose{\csromanfolio{229}วินีตา วิสารทา พหุสฺสุตา ธมฺมธรา ธมฺมา\-นุธมฺม\-ปฺปฏิปนฺนา สามีจิ\-ปฺปฏิปนฺนา อนุธมฺม\-จารินิโย สกํ อาจริยกํ อุคฺคเหตฺวา อาจิกฺขิสฺสนฺติ เทเสสฺสนฺติ ปญฺญเปสฺสนฺติ ปฏฺฐเปสฺสนฺติ วิวริสฺสนฺติ วิภชิสฺสนฺติ อุตฺตานีกริสฺสนฺติ, อุปฺปนฺนํ ปรปฺปวาทํ สหธมฺเมน สุนิคฺคหิตํ นิคฺคเหตฺวา สปฺปาฏิหาริยํ ธมฺมํ เทเสสฺสนฺตี”ติ.}

\prose{สนฺติ โข ปน ภนฺเต เอตรหิ ภิกฺขุนิโย ภควโต สาวิกา วิยตฺตา วินีตา วิสารทา พหุสฺสุตา ธมฺมธรา ธมฺมา\-นุธมฺม\-ปฺปฏิปนฺนา สามีจิ\-ปฺปฏิปนฺนา อนุธมฺม\-จารินิโย สกํ อาจริยกํ อุคฺคเหตฺวา อาจิกฺขนฺติ เทเสนฺติ ปญฺญเปนฺติ ปฏฺฐเปนฺติ วิวรนฺติ วิภชนฺติ อุตฺตานีกโรนฺติ, อุปฺปนฺนํ ปรปฺปวาทํ สหธมฺเมน สุนิคฺคหิตํ นิคฺคเหตฺวา สปฺปาฏิหาริยํ ธมฺมํ เทเสนฺติ, ปรินิพฺพาตุ\-ทานิ ภนฺเต ภควา, ปรินิพฺพาตุ\-ทานิ สุคโต, ปรินิพฺพาน\-กาโลทานิ ภนฺเต ภควโต.}

\prose{ภาสิตา โข ปเนสา ภนฺเต ภควตา วาจา “น ตาวาหํ ปาปิม ปรินิพฺพายิสฺสามิ, ยาว เม อุปาสกา ฯเปฯ ยาว เม อุปาสิกา น สาวิกา ภวิสฺสนฺติ วิยตฺตา วินีตา วิสารทา พหุสฺสุตา ธมฺมธรา ธมฺมา\-นุธมฺม\-ปฺปฏิปนฺนา สามีจิ\-ปฺปฏิปนฺนา อนุธมฺมจาริโย สกํ อาจริยกํ อุคฺคเหตฺวา อาจิกฺขิสฺสนฺติ เทเสสฺสนฺติ ปญฺญเปสฺสนฺติ ปฏฺฐเปสฺสนฺติ วิวริสฺสนฺติ วิภชิสฺสนฺติ อุตฺตานีกริสฺสนฺติ, อุปฺปนฺนํ ปรปฺปวาทํ สหธมฺเมน สุนิคฺคหิตํ นิคฺคเหตฺวา สปฺปาฏิหาริยํ ธมฺมํ เทเสสฺสนฺตี”ติ.}

\prose{สนฺติ โข ปน ภนฺเต เอตรหิ อุปาสกา อุปาสิกา ภควโต สาวิกา วิยตฺตา วินีตา วิสารทา พหุสฺสุตา ธมฺมธรา ธมฺมา\-นุธมฺม\-ปฺปฏิปนฺนา สามีจิ\-ปฺปฏิปนฺนา อนุธมฺม\-จารินิโย สกํ อาจริยกํ อุคฺคเหตฺวา อาจิกฺขนฺติ เทเสนฺติ ปญฺญเปนฺติ ปฏฺฐเปนฺติ วิวรนฺติ วิภชนฺติ อุตฺตานีกโรนฺติ, อุปฺปนฺนํ ปรปฺปวาทํ สหธมฺเมน สุนิคฺคหิตํ นิคฺคเหตฺวา สปฺปาฏิหาริยํ ธมฺมํ เทเสนฺติ. ปรินิพฺพาตุ\-ทานิ ภนฺเต ภควา, ปรินิพฺพาตุ\-ทานิ สุคโต, ปรินิพฺพาน\-กาโลทานิ ภนฺเต ภควโต.}

\prose{ภาสิตา โข ปเนสา ภนฺเต ภควตา วาจา “น ตาวาหํ ปาปิม ปรินิพฺพายิสฺสามิ, ยาว เม อิทํ พฺรหฺมจริยํ น อิทฺธญฺเจว ภวิสฺสติ ผีตญฺจ, วิตฺถาริตํ พาหุชญฺญํ ปุถุภูตํ ยาว เทวมนุสฺเสหิ สุปฺปกาสิตนฺ”ติ. ตยิทํ ภนฺเต ภควโต พฺรหฺมจริยํ อิทฺธญฺเจว ผีตญฺจ, วิตฺถาริตํ พาหุชญฺญํ}

\prose{\csromanfolio{230}ปุถุภูตํ ยาว เทวมนุสฺเสหิ สุปฺปกาสิตํ. ปรินิพฺพาตุ\-ทานิ ภนฺเต ภควา, ปรินิพฺพาตุ\-ทานิ สุคโต, ปรินิพฺพาน\-กาโลทานิ ภนฺเต ภควโตติ.}

\prose{เอวํ วุตฺเต ภควา มารํ ปาปิมนฺตํ เอตทโวจ “อปฺโปสฺสุกฺโก ตฺวํ ปาปิม โหหิ, น จิรํ\csromansharedfootnote{896f75dac2b6d212}{นจิรสฺเสว (ก)} ตถาคตสฺส ปรินิพฺพานํ ภวิสฺสติ, อิโต ติณฺณํ มาสานํ อจฺจเยน ตถาคโต ปรินิพฺพายิสฺสตี”ติ. อถ โข ภควา จาปาเล เจติเย สโต สมฺปชาโน อายุสงฺขารํ โอสฺสชิ, โอสฺสฏฺเฐ จ\csromansharedfootnote{93e006239c58a05a}{โอสชฺเช ปน (ก)} ภควตา อายุสงฺขาเร มหาภูมิจาโล อโหสิ ภิํสนโก โลมหํโส, เทวทุนฺทุภิโย\csromansharedfootnote{2e21bc2f8b7cd744}{เทวทุทฺรภิโย (ก)} จ ผลิํสุ. อถ โข ภควา เอตมตฺถํ วิทิตฺวา ตายํ เวลายํ อิมํ อุทานํ อุทาเนสิ\csromandash{}}

\setlength{\csromangathapagewidth}{0pt}
\setlength{\csromangathawakwidth}{0pt}
\csromangathameasuregroup{“ตุลมตุลญฺจ สมฺภวํ,}{\csromangathabat{“ตุลมตุลญฺจ สมฺภวํ,}{ภวสงฺขารมวสฺสชิ มุนิ.}}
\csromangathasetleft{“ตุลมตุลญฺจ สมฺภวํ,}
\csromangathagroupwithclosermidruled{\csromangathastanza{\csromangathabat{“ตุลมตุลญฺจ สมฺภวํ,}{ภวสงฺขารม\-วสฺสชิ มุนิ.}}}{อชฺฌตฺตรโต สมาหิโต, อภินฺทิ กวจมิวตฺตสมฺภวนฺ”ติ. ทสมํ. จาปาลวคฺโค ปฐโม.}
\csromancenter{\textbf{ตสฺสุทฺทานํ}}
\setlength{\csromangathapagewidth}{0pt}
\setlength{\csromangathawakwidth}{0pt}
\csromangathameasuregroup{อปาราปิ วิรทฺโธ จ, \\ ปเทสํ สมตฺตํ ภิกฺขุ,}{\csromangathabat{อปาราปิ วิรทฺโธ จ,}{อริยา นิพฺพิทาปิ จ.} \\ \csromangathabat{ปเทสํ สมตฺตํ ภิกฺขุ,}{พุทฺธํ ญาณญฺจ เจติยนฺติ.}}
\csromangathasetleft{อปาราปิ วิรทฺโธ จ, \\ ปเทสํ สมตฺตํ ภิกฺขุ,}
\csromangathagroup{\csromangathastanza{\csromangathabat{อปาราปิ วิรทฺโธ จ,}{อริยา นิพฺพิทาปิ จ.} \\ \csromangathabat{ปเทสํ สมตฺตํ ภิกฺขุ,}{พุทฺธํ ญาณญฺจ เจติยนฺติ.}}}

\csromanmatikaanchor{mtk.173}
\csromantocmarknopage{subparagraph}{2. ปาสาทกมฺปนวคฺค}{mtk.173}
\bookmark[dest={mtk.173},level=5]{2. ปาสาทกมฺปนวคฺค}
\csromanheader{6}{2. ปาสาทกมฺปน\-วคฺค}
\csromanmatikaanchor{mtk.174}
\csromantocmarkref{subparagraph}{1. ปุพฺพสุตฺต}{mtk.174}
\bookmark[dest={mtk.174},level=5]{1. ปุพฺพสุตฺต}
\csromanheader{6}{1. \textbf{ปุพฺพสุตฺต}}
\proseitem{823}{สาวตฺถิ\-นิทานํ. ปุพฺเพว เม ภิกฺขเว สมฺโพธา อนภิสมฺพุทฺธสฺส โพธิสตฺตสฺเสว สโต เอตทโหสิ “โก นุ โข เหตุ โก ปจฺจโย อิทฺธิปาทภาวนายา”ติ. ตสฺส มยฺหํ ภิกฺขเว เอตทโหสิ “อิธ ภิกฺขุ ฉนฺท\-สมาธิ\-ปฺปธาน\-สงฺขาร\-สมนฺนาคตํ อิทฺธิปาทํ ภาเวติ, อิติ เม ฉนฺโท น จ อติลีโน ภวิสฺสติ, น จ อติปฺปคฺคหิโต ภวิสฺสติ, น จ อชฺฌตฺตํ สํขิตฺโต ภวิสฺสติ. น จ พหิทฺธา วิกฺขิตฺโต ภวิสฺสติ, ปจฺฉา\-ปุเร\-สญฺญี จ วิหรติ, ยถา ปุเร ตถา ปจฺฉา, ยถา ปจฺฉา ตถา ปุเร,}

\prose{\csromanfolio{231}ยถา อโธ ตถา อุทฺธํ, ยถา อุทฺธํ ตถา อโธ, ยถา ทิวา ตถา รตฺติํ, ยถา รตฺติํ ตถา ทิวา, อิติ วิวเฏน เจตสา อปริโยนทฺเธน สปฺปภาสํ จิตฺตํ ภาเวติ.}

\prose{วีริย\-สมาธิ\-ปฺปธาน\-สงฺขาร\-สมนฺนาคตํ อิทฺธิปาทํ ภาเวติ, อิติ เม วีริยํ น จ อติลีนํ ภวิสฺสติ, น จ อติปฺปคฺคหิตํ ภวิสฺสติ, น จ อชฺฌตฺตํ สํขิตฺตํ ภวิสฺสติ, น จ พหิทฺธา วิกฺขิตฺตํ ภวิสฺสติ, ปจฺฉา\-ปุเร\-สญฺญี จ วิหรติ, ยถา ปุเร ตถา ปจฺฉา, ยถา ปจฺฉา ตถา ปุเร, ยถา อโธ ตถา อุทฺธํ, ยถา อุทฺธํ ตถา อโธ, ยถา ทิวา ตถา รตฺติํ, ยถา รตฺติํ ตถา ทิวา, อิติ วิวเฏน เจตสา อปริโยนทฺเธน สปฺปภาสํ จิตฺตํ ภาเวติ.}

\prose{จิตฺต\-สมาธิ\-ปฺปธาน\-สงฺขาร\-สมนฺนาคตํ อิทฺธิปาทํ ภาเวติ, อิติ เม จิตฺตํ น จ อติลีนํ ภวิสฺสติ, น จ อติปฺปคฺคหิตํ ภวิสฺสติ, น จ อชฺฌตฺตํ สํขิตฺตํ ภวิสฺสติ, น จ พหิทฺธา วิกฺขิตฺตํ ภวิสฺสติ, ปจฺฉา\-ปุเร\-สญฺญี จ วิหรติ, ยถา ปุเร ตถา ปจฺฉา, ยถา ปจฺฉา ตถา ปุเร, ยถา อโธ ตถา อุทฺธํ, ยถา อุทฺธํ ตถา อโธ, ยถา ทิวา ตถา รตฺติํ, ยถา รตฺติํ ตถา ทิวา, อิติ วิวเฏน เจตสา อปริโยนทฺเธน สปฺปภาสํ จิตฺตํ ภาเวติ.}

\prose{วีมํสา\-สมาธิ\-ปฺปธาน\-สงฺขาร\-สมนฺนาคตํ อิทฺธิปาทํ ภาเวติ, อิติ เม วีมํสา น จ อติลีนา ภวิสฺสติ, น จ อติปฺปคฺคหิตา ภวิสฺสติ, น จ อชฺฌตฺตํ สํขิตฺตา ภวิสฺสติ, น จ พหิทฺธา วิกฺขิตฺตา ภวิสฺสติ, ปจฺฉา\-ปุเร\-สญฺญี จ วิหรติ, ยถา ปุเร ตถา ปจฺฉา, ยถา ปจฺฉา ตถา ปุเร, ยถา อโธ ตถา อุทฺธํ, ยถา อุทฺธํ ตถา อโธ, ยถา ทิวา ตถา รตฺติํ, ยถา รตฺติํ ตถา ทิวา, อิติ วิวเฏน เจตสา อปริโยนทฺเธน สปฺปภาสํ จิตฺตํ ภาเวติ.}

\prose{เอวํ ภาวิเตสุ โข ภิกฺขุ จตูสุ อิทฺธิปาเทสุ เอวํ พหุลีกเตสุ อเนกวิหิตํ อิทฺธิวิธํ ปจฺจนุโภติ, เอโกปิ หุตฺวา พหุธา โหติ, พหุธาปิ หุตฺวา เอโก โหติ, อาวิภาวํ ติโรภาวํ ติโรกุฏฺฏํ ติโรปาการํ ติโรปพฺพตํ อสชฺชมาโน คจฺฉติ เสยฺยถาปิ อากาเส. ปถวิยาปิ อุมฺมุชฺช\-นิมุชฺชํ กโรติ เสยฺยถาปิ อุทเก. อุทเกปิ}

\prose{\csromanfolio{232}อภิชฺชมาเน\csromansharedfootnote{351f9c41fcc0a63b}{อภิชฺชมาโน (สี, อิ, ก)} คจฺฉติ เสยฺยถาปิ ปถวิยํ. อากาเสปิ ปลฺลงฺเกน กมติ เสยฺยถาปิ ปกฺขี สกุโณ. อิเมปิ จนฺทิมสูริเย เอวํมหิทฺธิเก เอวํ\-มหา\-นุภาเว ปาณินา ปริมสติ\csromansharedfootnote{2d0e1cfd11f58810}{ปรามสติ (สี, สฺยา, กํ)} ปริมชฺชติ, ยาว พฺรหฺมโลกาปิ กาเยน วสํ วตฺเตติ.}

\prose{เอวํ ภาวิเตสุ โข ภิกฺขุ จตูสุ อิทฺธิปาเทสุ เอวํ พหุลีกเตสุ ทิพฺพาย โสตธาตุยา วิสุทฺธาย อติกฺกนฺตมานุสิกาย อุโภ สทฺเท สุณาติ ทิพฺเพ จ มานุเส จ ทูเร สนฺติเก จาติ.}

\prose{เอวํ ภาวิเตสุ โข ภิกฺขุ จตูสุ อิทฺธิปาเทสุ เอวํ พหุลีกเตสุ ปรสตฺตานํ ปรปุคฺคลานํ เจตสา เจโต ปริจฺจ ปชานาติ, สราคํ วา จิตฺตํ “สราคํ จิตฺตนฺ”ติ ปชานาติ, วีตราคํ วา จิตฺตํ “วีตราคํ จิตฺตนฺ”ติ ปชานาติ, สโทสํ วา จิตฺตํ “สโทสํ จิตฺตนฺ”ติ ปชานาติ, วีตโทสํ วา จิตฺตํ “วีตโทสํ จิตฺตนฺ”ติ ปชานาติ, สโมหํ วา จิตฺตํ “สโมหํ จิตฺตนฺ”ติ ปชานาติ, วีตโมหํ วา จิตฺตํ “วีตโมหํ จิตฺตนฺ”ติ ปชานาติ, สํขิตฺตํ วา จิตฺตํ “สํขิตฺตํ จิตฺตนฺ”ติ ปชานาติ, วิกฺขิตฺตํ วา จิตฺตํ “วิกฺขิตฺตํ จิตฺตนฺ”ติ ปชานาติ, มหคฺคตํ วา จิตฺตํ “มหคฺคตํ จิตฺตํนฺ”ติ ปชานาติ, อมหคฺคตํ วา จิตฺตํ “อมหคฺคตํ จิตฺตนฺ”ติ ปชานาติ, สอุตฺตรํ วา จิตฺตํ “สอุตฺตรํ จิตฺตนฺ”ติ ปชานาติ, อนุตฺตรํ วา จิตฺตํ “อนุตฺตรํ จิตฺตนฺ”ติ ปชานาติ, สมาหิตํ วา จิตฺตํ “สมาหิตํ จิตฺตนฺ”ติ ปชานาติ, อสมาหิตํ วา จิตฺตํ “อสมาหิตํ จิตฺตนฺ”ติ ปชานาติ, วิมุตฺตํ วา จิตฺตํ “วิมุตฺตํ จิตฺตนฺ”ติ ปชานาติ, อวิมุตฺตํ วา จิตฺตํ “อวิมุตฺตํ จิตฺตนฺ”ติ ปชานาติ.}

\prose{เอวํ ภาวิเตสุ โข ภิกฺขุ จตูสุ อิทฺธิปาเทสุ เอวํ พหุลีกเตสุ อเนกวิหิตํ ปุพฺเพนิวาสํ อนุสฺสรติ. เสยฺยถิทํ, เอกมฺปิ ชาติํ เทฺวปิ ชาติโย ติสฺโสปิ ชาติโย จตสฺโสปิ ชาติโย ปญฺจปิ ชาติโย ทสปิ ชาติโย วีสมฺปิ ชาติโย ติํสมฺปิ ชาติโย จตฺตาลีสมฺปิ ชาติโย ปญฺญาสมฺปิ ชาติโย ชาติสตมฺปิ ชาติสหสฺสมฺปิ ชาติ\-สต\-สหสฺสมฺปิ อเนเกปิ สํวฏฺฏกปฺเป อเนเกปิ วิวฏฺฏกปฺเป อเนเกปิ สํวฏฺฏ\-วิวฏฺฏ\-กปฺเป, “อมุตฺราสิํ เอวํนาโม เอวํโคตฺโต เอวํวณฺโณ เอวมาหาโร เอวํ\-สุขทุกฺข\-ปฺปฏิสํเวที เอวมา\-ยุปริยนฺโต, โส ตโต จุโต อมุตฺร อุทปาทิํ\csromansharedfootnote{2c7ed2bc63e27f12}{อุปฺปาทิํ (สี)}, ตตฺราปาสิํ}

\prose{\csromanfolio{233}“เอวํนาโม เอวํโคตฺโต เอวํวณฺโณ เอวมาหาโร เอวํ\-สุขทุกฺข\-ปฺปฏิสํเวที เอวมา\-ยุปริยนฺโต, โส ตโต จุโต อิธูปปนฺโน”ติ. อิติ สาการํ สอุทฺเทสํ อเนกวิหิตํ ปุพฺเพนิวาสํ อนุสฺสรติ.}

\prose{เอวํ ภาวิเตสุ โข ภิกฺขุ จตูสุ อิทฺธิปาเทสุ เอวํ พหุลีกเตสุ ทิพฺเพน จกฺขุนา วิสุทฺเธน อติกฺกนฺต\-มานุสเกน สตฺเต ปสฺสติ จวมาเน อุปปชฺชมาเน หีเน ปณีเต สุวณฺเณ ทุพฺพณฺเณ สุคเต ทุคฺคเต, ยถากมฺมูปเค สตฺเต ปชานาติ “อิเม วต โภนฺโต สตฺตา กาย\-ทุจฺจริเตน สมนฺนาคตา วจีทุจฺจริเตน สมนฺนาคตา มโน\-ทุจฺจริเตน สมนฺนาคตา, อริยานํ อุปวาทกา มิจฺฉาทิฏฺฐิกา มิจฺฉาทิฏฺฐิ\-กมฺม\-สมาทานา, เต กายสฺส เภทา ปรํ มรณา อปายํ ทุคฺคติํ วินิปาตํ นิรยํ อุปปนฺนา. อิเม วา ปน โภนฺโต สตฺตา กายสุจริเตน สมนฺนาคตา วจีสุจริเตน สมนฺนาคตา, มโนสุจริเตน สมนฺนาคตา อริยานํ อนุปวาทกา สมฺมาทิฏฺฐิกา สมฺมา\-ทิฏฺฐิ\-กมฺม\-สมาทานา, เต กายสฺส เภทา ปรํ มรณา สุคติํ สคฺคํ โลกํ อุปปนฺนา”ติ. อิติ ทิพฺเพน จกฺขุนา วิสุทฺเธน อติกฺกนฺต\-มานุสเกน สตฺเต ปสฺสติ จวมาเน อุปปชฺชมาเน หีเน ปณีเต สุวณฺเณ ทุพฺพณฺเณ สุคเต ทุคฺคเต, ยถากมฺมูปเค สตฺเต ปชานาติ.}

\prose{เอวํ ภาวิเตสุ โข ภิกฺขุ จตูสุ อิทฺธิปาเทสุ เอวํ พหุลีกเตสุ อาสวานํ ขยา อนาสวํ เจโตวิมุตฺติํ ปญฺญาวิมุตฺติํ ทิฏฺเฐว ธมฺเม สยํ อภิญฺญา สจฺฉิกตฺวา อุปสมฺปชฺช วิหรตีติ.  ปฐมํ.}

\csromanmatikaanchor{mtk.175}
\csromantocmarkref{subparagraph}{2. มหปฺผลสุตฺต}{mtk.175}
\bookmark[dest={mtk.175},level=5]{2. มหปฺผลสุตฺต}
\csromanheader{6}{2. มหปฺผล\-สุตฺต}
\csromanmatikaanchor{mtk.176}
\csromantocmarkref{subparagraph}{3. ฉนฺทสมาทิสุตฺต}{mtk.176}
\bookmark[dest={mtk.176},level=5]{3. ฉนฺทสมาทิสุตฺต}
\proseitem{824}{จตฺตาโรเม ภิกฺขเว อิทฺธิปาทา ภาวิตา พหุลีกตา มหปฺผลา โหนฺติ มหานิสํสา. กถํ ภาวิตา จ ภิกฺขเว จตฺตาโร อิทฺธิปาทา กถํ พหุลีกตา มหปฺผลา โหนฺติ มหานิสํสา. อิธ ภิกฺขเว ภิกฺขุ ฉนฺท\-สมาธิ\-ปฺปธาน\-สงฺขาร\-สมนฺนาคตํ อิทฺธิปาทํ ภาเวติ, อิติ เม ฉนฺโท น จ อติลีโน ภวิสฺสติ, น จ อติปฺปคฺคหิโต ภวิสฺสติ, น จ อชฺฌตฺตํ สํขิตฺโต ภวิสฺสติ, น จ พหิทฺธา วิกฺขิตฺโต ภวิสฺสติ, ปจฺฉา\-ปุเร\-สญฺญี จ วิหรติ, ยถา ปุเร ตถา ปจฺฉา, ยถา ปจฺฉา ตถา ปุเร, ยถา อโธ ตถา อุทฺธํ, ยถา อุทฺธํ ตถา อโธ, ยถา ทิวา ตถา รตฺติํ, ยถา รตฺติํ \csromanfolio{234}ตถา ทิวา, อิติ วิวเฏน เจตสา อปริโยนทฺเธน สปฺปภาสํ จิตฺตํ ภาเวติ.}

\prose{วีริยสมาธิ ฯเปฯ. จิตฺตสมาธิ. วีมํสา\-สมาธิ\-ปฺปธาน\-สงฺขาร\-สมนฺนาคตํ อิทฺธิปาทํ ภาเวติ, อิติ เม วีมํสา น จ อติลีนา ภวิสฺสติ, น จ อติปฺปคฺคหิตา ภวิสฺสติ, น จ อชฺฌตฺตํ สํขิตฺตา ภวิสฺสติ, น จ พหิทฺธา วิกฺขิตฺตา ภวิสฺสติ, ปจฺฉา\-ปุเร\-สญฺญี จ วิหรติ, ยถา ปุเร ตถา ปจฺฉา, ยถา ปจฺฉา ตถา ปุเร, ยถา อโธ ตถา อุทฺธํ, ยถา อุทฺธํ ตถา อโธ, ยถา ทิวา ตถา รตฺติํ, ยถา รตฺติํ ตถา ทิวา, อิติ วิวเฏน เจตสา อปริโยนทฺเธน สปฺปภาสํ จิตฺตํ ภาเวติ. เอวํ ภาวิตา โข ภิกฺขเว จตฺตาโร อิทฺธิปาทา เอวํ พหุลีกตา มหปฺผลา โหนฺติ มหานิสํสา.}

\prose{เอวํ ภาวิเตสุ โข ภิกฺขเว ภิกฺขุ จตูสุ อิทฺธิปาเทสุ เอวํ พหุลีกเตสุ อเนกวิหิตํ อิทฺธิวิธํ ปจฺจนุโภติ, เอโกปิ หุตฺวา พหุธา โหติ ฯเปฯ ยาว พฺรหฺมโลกาปิ กาเยน วสํ วตฺเตติ.}

\prose{เอวํ ภาวิเตสุ โข ภิกฺขเว ภิกฺขุ จตูสุ อิทฺธิปาเทสุ เอวํ พหุลีกเตสุ อาสวานํ ขยา อนาสวํ เจโตวิมุตฺติํ ปญฺญาวิมุตฺติํ ทิฏฺเฐว ธมฺเม สยํ อภิญฺญา สจฺฉิกตฺวา อุปสมฺปชฺช วิหรตีติ.  ทุติยํ.}

\csromanheader{2}{3. ฉนฺทสมาธิ\-สุตฺต}
\proseitem{825}{ฉนฺทํ เจ ภิกฺขเว ภิกฺขุ นิสฺสาย ลภติ สมาธิํ, ลภติ จิตฺตสฺส เอกคฺคตํ, อยํ วุจฺจติ ฉนฺทสมาธิ. โส อนุปฺปนฺนานํ ปาปกานํ อกุสลานํ ธมฺมานํ อนุปฺปาทาย ฉนฺทํ ชเนติ วายมติ, วีริยํ อารภติ, จิตฺตํ ปคฺคณฺหาติ ปทหติ. อุปฺปนฺนานํ ปาปกานํ อกุสลานํ ธมฺมานํ ปหานาย ฉนฺทํ ชเนติ วายมติ, วีริยํ อารภติ, จิตฺตํ ปคฺคณฺหาติ ปทหติ. อนุปฺปนฺนานํ กุสลานํ ธมฺมานํ อุปฺปาทาย ฉนฺทํ ชเนติ วายมติ, วีริยํ อารภติ, จิตฺตํ ปคฺคณฺหาติ ปทหติ. อุปฺปนฺนานํ กุสลานํ ธมฺมานํ ฐิติยา อสมฺโมสาย ภิยฺโยภาวาย เวปุลฺลาย ภาวนาย ปาริปูริยา ฉนฺทํ ชเนติ วายมติ, วีริยํ อารภติ, จิตฺตํ ปคฺคณฺหาติ ปทหติ. อิเม วุจฺจนฺติ ปธาน\-สงฺขาราติ. อิติ อยญฺจ ฉนฺโท อยญฺจ ฉนฺทสมาธิ อิเม จ ปธาน\-สงฺขารา, อยํ วุจฺจติ ภิกฺขเว ฉนฺท\-สมาธิ\-ปฺปธาน\-สงฺขาร\-สมนฺนาคโต อิทฺธิปาโท.}

\prose{\csromanfolio{235}วีริยํ เจ ภิกฺขเว ภิกฺขุ นิสฺสาย ลภติ สมาธิํ, ลภติ จิตฺตสฺส เอกคฺคตํ, อยํ วุจฺจติ วีริยสมาธิ. โส อนุปฺปนฺนานํ ฯเปฯ. อุปฺปนฺนานํ กุสลานํ ธมฺมานํ ฐิติยา อสมฺโมสาย ภิยฺโยภาวาย เวปุลฺลาย ภาวนาย ปาริปูริยา ฉนฺทํ ชเนติ วายมติ, วีริยํ อารภติ, จิตฺตํ ปคฺคณฺหาติ ปทหติ. อิเม วุจฺจนฺติ ปธาน\-สงฺขาราติ. อิติ อิทญฺจ วีริยํ อยญฺจ วีริยสมาธิ อิเม จ ปธาน\-สงฺขารา, อยํ วุจฺจติ ภิกฺขเว วีริย\-สมาธิ\-ปฺปธาน\-สงฺขาร\-สมนฺนาคโต อิทฺธิปาโท.}

\prose{จิตฺตํ เจ ภิกฺขเว ภิกฺขุ นิสฺสาย ลภติ สมาธิํ, ลภติ จิตฺตสฺส เอกคฺคตํ, อยํ วุจฺจติ จิตฺตสมาธิ. โส อนุปฺปนฺนานํ ปาปกานํ ฯเปฯ. อุปฺปนฺนานํ กุสลานํ ธมฺมานํ ฐิติยา อสมฺโมสาย ภิยฺโยภาวาย เวปุลฺลาย ภาวนาย ปาริปูริยา ฉนฺทํ ชเนติ วายมติ, วีริยํ อารภติ, จิตฺตํ ปคฺคณฺหาติ ปทหติ. อิเม วุจฺจนฺติ ปธาน\-สงฺขาราติ. อิติ อิทญฺจ จิตฺตํ อยญฺจ จิตฺตสมาธิ อิเม จ ปธาน\-สงฺขารา, อยํ วุจฺจติ ภิกฺขเว จิตฺต\-สมาธิ\-ปฺปธาน\-สงฺขาร\-สมนฺนาคโต อิทฺธิปาโท.}

\prose{วีมํสํ เจ ภิกฺขเว ภิกฺขุ นิสฺสาย ลภติ สมาธิํ, ลภติ จิตฺตสฺส เอกคฺคตํ, อยํ วุจฺจติ วีมํสาสมาธิ. โส อนุปฺปนฺนานํ ปาปกานํ อกุสลานํ ธมฺมานํ อนุปฺปาทาย ฉนฺทํ ชเนติ วายมติ, วีริยํ อารภติ, จิตฺตํ ปคฺคณฺหาติ ปทหติ ฯเปฯ. อุปฺปนฺนานํ กุสลานํ ธมฺมานํ ฐิติยา อสมฺโมสาย ภิยฺโยภาวาย เวปุลฺลาย ภาวนาย ปาริปูริยา ฉนฺทํ ชเนติ วายมติ, วีริยํ อารภติ, จิตฺตํ ปคฺคณฺหาติ ปทหติ. อิเม วุจฺจนฺติ ปธาน\-สงฺขารา. อิติ อยญฺจ วีมํสา อยญฺจ วีมํสาสมาธิ อิเม จ ปธาน\-สงฺขารา, อยํ วุจฺจติ ภิกฺขเว วีมํสา\-สมาธิ\-ปฺปธาน\-สงฺขาร\-สมนฺนาคโต อิทฺธิปาโทติ.  ตติยํ.}

\csromanmatikaanchor{mtk.177}
\csromantocmarkref{subparagraph}{4. โมคฺคลฺลานสุตฺต}{mtk.177}
\bookmark[dest={mtk.177},level=5]{4. โมคฺคลฺลานสุตฺต}
\csromanheader{6}{4. โมคฺคลฺลาน\-สุตฺต}
\proseitem{826}{เอวํ เม สุตํ\csromandash{}เอกํ สมยํ ภควา สาวตฺถิยํ วิหรติ ปุพฺพาราเม มิคาร\-มาตุ\-ปาสาเท. เตน โข ปน สมเยน สมฺพหุลา ภิกฺขู เหฏฺฐา\-มิคาร\-มาตุ\-ปาสาเท วิหรนฺติ อุทฺธตา อุนฺนฬา จปลา มุขรา วิกิณฺณวาจา มุฏฺฐสฺสติโน อสมฺปชานา อสมาหิตา ภนฺตจิตฺตา\csromansharedfootnote{c24b4112d59ffd6b}{วิพฺภนฺตจิตฺตา (สี, สฺยา, กํ)} ปากตินฺทฺริยา.}

\prose{\csromanfolio{236}อถ โข ภควา อายสฺมนฺตํ มหา\-โมคฺคลฺลานํ อามนฺเตสิ “เอเต โข โมคฺคลฺลาน สพฺรหฺมจาริโน เหฏฺฐา\-มิคาร\-มาตุ\-ปาสาเท วิหรนฺติ อุทฺธตา อุนฺนฬา จปลา มุขรา วิกิณฺณวาจา มุฏฺฐสฺสติโน อสมฺปชานา อสมาหิตา ภนฺตจิตฺตา ปากตินฺทฺริยา, คจฺฉ โมคฺคลฺลาน, เต ภิกฺขู สํเวเชหี”ติ.}

\prose{“เอวํ ภนฺเต”ติ โข อายสฺมา มหาโมคฺคลฺลาโน ภควโต ปฏิสฺสุตฺวา ตถารูปํ อิทฺธา\-ภิสงฺขารํ อภิสงฺขาสิ\csromansharedfootnote{b914b96a2662d788}{อภิสงฺขเรสิ (สฺยา, อิ, ก)}. ยถา ปาท\-งฺคุฏฺฐเกน มิคาร\-มาตุ\-ปาสาทํ สงฺกมฺเปสิ สมฺปกมฺเปสิ สมฺปจาเลสิ. อถ โข เต ภิกฺขู สํวิคฺคา โลมหฏฺฐชาตา เอกมนฺตํ อฏฺฐํสุ, อจฺฉริยํ วต โภ, อพฺภุตํ วต โภ, นิวาตญฺจ วต, อยญฺจ มิคาร\-มาตุ\-ปาสาโท คมฺภีรเนโม สุนิขาโต อจโล อสมฺปกมฺปี, อถ จ ปน สงฺกมฺปิโต สมฺปกมฺปิโต สมฺปจาลิโตติ.}

\prose{อถ โข ภควา เยน เต ภิกฺขู เตนุปสงฺกมิ, อุปสงฺกมิตฺวา เต ภิกฺขู ภควา เอตทโวจ “กิํ นุ ตุเมฺห ภิกฺขเว สํวคฺคา โลมหฏฺฐชาตา เอกมนฺตํ ฐิตา”ติ. อจฺฉริยํ ภนฺเต, อพฺภุตํ ภนฺเต, นิวาตญฺจ วต, อยญฺจ มิคาร\-มาตุ\-ปาสาโท คมฺภีรเนโม สุนิขาโต อจโล อสมฺปกมฺปี, อถ จ ปน สงฺกมฺปิโต สมฺปกมฺปิโต สมฺปจาลิโตติ. ตุเมฺหว โข ภิกฺขเว สํเวเชตุ\-กาเมน โมคฺคลฺลาเนน ภิกฺขุนา ปาท\-งฺคุฏฺฐเกน มิคาร\-มาตุ\-ปาสาโท สงฺกมฺปิโต สมฺปกมฺปิโต สมฺปจาลิโต. ตํ กิํ มญฺญถ ภิกฺขเว, กตเมสํ ธมฺมานํ ภาวิตตฺตา พหุลีกตตฺตา โมคฺคลฺลาโน ภิกฺขุ เอวํมหิทฺธิโก เอวํ\-มหา\-นุภาโวติ. ภควํมูลกา โน ภนฺเต ธมฺมา ภควํ\-เนตฺติกา ภควํ\-ปฏิสรณา, สาธุ วต ภนฺเต ภควนฺตํเยว ปฏิภาตุ เอตสฺส ภาสิตสฺส อตฺโถ, ภควโต สุตฺวา ภิกฺขู ธาเรสฺสนฺตีติ.}

\prose{เตน หิ ภิกฺขเว สุณาถ, จตุนฺนํ โข ภิกฺขเว อิทฺธิปาทานํ ภาวิตตฺตา พหุลีกตตฺตา โมคฺคลฺลาโน ภิกฺขุ เอวํมหิทฺธิโก เอวํ\-มหา\-นุภาโว. กตเมสํ จตุนฺนํ, อิธ ภิกฺขเว โมคฺคลฺลาโน ภิกฺขุ ฉนฺท\-สมาธิ\-ปฺปธาน\-สงฺขาร\-สมนฺนาคตํ อิทฺธิปาทํ ภาเวติ, วีริยสมาธิ ฯเปฯ จิตฺตสมาธิ, วีมํสา\-สมาธิ\-ปฺปธาน\-สงฺขาร\-สมนฺนาคตํ อิทฺธิปาทํ ภาเวติ, อิติ เม วีมํสา น}

\prose{\csromanfolio{237}จ อติลีนา ภวิสฺสติ, น จ อติปฺปคฺคหิตา ภวิสฺสติ น จ อชฺฌตฺตํ สํขิตฺตา ภวิสฺสติ, น จ พหิทฺธา วิกฺขิตฺตา ภวิสฺสติ. ปจฺฉา\-ปุเร\-สญฺญี จ วิหรติ, ยถา ปุเร ตถา ปจฺฉา, ยถา ปจฺฉา ตถา ปุเร, ยถา อโธ ตถา อุทฺธํ, ยถา อุทฺธํ ตถา อโธ, ยถา ทิวา ตถา รตฺติํ, ยถา รตฺติํ ตถา ทิวา, อิติ วิวเฏน เจตสา อปริโยนทฺเธน สปฺปภาสํ จิตฺตํ ภาเวติ. อิเมสํ โข ภิกฺขเว จตุนฺนํ อิทฺธิปาทานํ ภาวิตตฺตา พหุลีกตตฺตา โมคฺคลฺลาโน ภิกฺขุ เอวํมหิทฺธิโก เอวํ\-มหา\-นุภาโว, อิเมสญฺจ ปน ภิกฺขเว จตุนฺนํ อิทฺธิปาทานํ ภาวิตตฺตา พหุลีกตตฺตา โมคฺคลฺลาโน ภิกฺขุ อเนกวิหิตํ อิทฺธิวิธํ ปจฺจนุโภติ ฯเปฯ ยาว พฺรหฺมโลกาปิ กาเยน วสํ วตฺเตติ. อิเมสญฺจ ปน ภิกฺขเว จตุนฺนํ อิทฺธิปาทานํ ภาวิตตฺตา พหุลีกตตฺตา โมคฺคลฺลาโน ภิกฺขุ อาสวานํ ขยา อนาสวํ เจโตวิมุตฺติํ ปญฺญาวิมุตฺติํ ทิฏฺเฐว ธมฺเม สยํ อภิญฺญา สจฺฉิกตฺวา อุปสมฺปชฺช วิหรตีติ.  จตุตฺถํ.}

\csromanmatikaanchor{mtk.178}
\csromantocmarkref{subparagraph}{5. อุณฺณาภพฺราหฺมณสุตฺต}{mtk.178}
\bookmark[dest={mtk.178},level=5]{5. อุณฺณาภพฺราหฺมณสุตฺต}
\csromanheader{6}{5. อุณฺณาภ\-พฺราหฺมณ\-สุตฺต}
\proseitem{827}{เอวํ เม สุตํ\csromandash{}เอกํ สมยํ อายสฺมา อานนฺโท โกสมฺพิยํ วิหรติ โฆสิตาราเม. อถ โข อุณฺณาโภ พฺราหฺมโณ เยนายสฺมา อานนฺโท เตนุปสงฺกมิ, อุปสงฺกมิตฺวา อายสฺมตา อานนฺเทน สทฺธิํ สมฺโมทิ, สมฺโมทนียํ กถํ สารณียํ วีติสาเรตฺวา เอกมนฺตํ นิสีทิ, เอกมนฺตํ นิสินฺโน โข อุณฺณาโภ พฺราหฺมโณ อายสฺมนฺตํ อานนฺทํ เอตทโวจ “กิมตฺถิยํ นุ โข โภ อานนฺท สมเณ โคตเม พฺรหฺมจริยํ วุสฺสตี”ติ. ฉนฺท\-ปฺปหานตฺถํ โข พฺราหฺมณ ภควติ พฺรหฺมจริยํ วุสฺสตีติ.}

\prose{อตฺถิ ปน โภ อานนฺท มคฺโค, อตฺถิ ปฏิปทา เอตสฺส ฉนฺทสฺส ปหานายาติ. อตฺถิ โข พฺราหฺมณ มคฺโค, อตฺถิ ปฏิปทา เอตสฺส ฉนฺทสฺส ปหานายาติ.}

\prose{กตโม ปน โภ อานนฺท มคฺโค กตมา ปฏิปทา เอตสฺส ฉนฺทสฺส ปหานายาติ, อิธ พฺราหฺมณ ภิกฺขุ, ฉนฺท\-สมาธิ\-ปฺปธาน\-สงฺขาร\-สมนฺนาคตํ อิทฺธิปาทํ ภาเวติ, วีริยสมาธิ ฯเปฯ. จิตฺตสมาธิ. วีมํสา\-สมาธิ\-ปฺปธาน\-สงฺขาร\-สมนฺนาคตํ อิทฺธิปาทํ ภาเวติ. อยํ โข พฺราหฺมณ มคฺโค อยํ ปฏิปทา เอตสฺส ฉนฺทสฺส ปหานายาติ.}

\prose{\csromanfolio{238}เอวํ สนฺเต โภ อานนฺท สนฺตกํ โหติ โน อสนฺตกํ. ฉนฺเทเนว ฉนฺทํ ปชหิสฺสตีติ เนตํ ฐานํ วิชฺชติ. เตน หิ พฺราหฺมณ ตญฺเญเวตฺถ ปฏิปุจฺฉิสฺสามิ. ยถา เต ขเมยฺย, ตถา ตํ พฺยากเรยฺยาสิ. ตํ กิํ มญฺญสิ พฺราหฺมณ, อโหสิ เต ปุพฺเพ ฉนฺโท “อารามํ คมิสฺสามี”ติ, ตสฺส เต อารามคตสฺส โย ตชฺโช ฉนฺโท, โส ปฏิปฺปสฺสทฺโธติ. เอวํ โภ. อโหสิ เต ปุพฺเพ วีริยํ “อารามํ คมิสฺสามี”ติ, ตสฺส เต อารามคตสฺส ยํ ตชฺชํ วีริยํ, ตํ ปฏิปฺปสฺสทฺธนฺติ. เอวํ โภ. อโหสิ เต ปุพฺเพ จิตฺตํ “อารามํ คมิสฺสามี”ติ, ตสฺส เต อารามคตสฺส ยํ ตชฺชํ จิตฺตํ, ตํ ปฏิปฺปสฺสทฺธนฺติ. เอวํ โภ. อโหสิ เต ปุพฺเพ วีมํสา “อารามํ คมิสฺสามี”ติ, ตสฺส เต อารามคตสฺส ยา ตชฺชา วีมํสา, สา ปฏิปฺปสฺสทฺธาติ. เอวํ โภ.}

\prose{เอวเมว โข พฺราหฺมณ โย โส ภิกฺขุ อรหํ ขีณาสโว วุสิตวา กตกรณีโย โอหิตภาโร อนุปฺปตฺต\-สทตฺโถ ปริกฺขีณ\-ภวสํโยชโน สมฺมทญฺญา\-วิมุตฺโต, ตสฺส โย ปุพฺเพ ฉนฺโท อโหสิ อรหตฺต\-ปฺปตฺติยา, อรหตฺตปฺปตฺเต\csromansharedfootnote{6340b50e637a58e5}{อรหตฺเต ปตฺเต (สี, สฺยา, กํ)} โย ตชฺโช ฉนฺโท, โส ปฏิปฺปสฺสทฺโธ. ยํ ปุพฺเพ วีริยํ อโหสิ อรหตฺต\-ปฺปตฺติยา, อรหตฺตปฺปตฺเต ยํ ตชฺชํ วีริยํ, ตํ ปฏิปฺปสฺสทฺธํ. ยํ ปุพฺเพ จิตฺตํ อโหสิ อรหตฺต\-ปฺปตฺติยา, อรหตฺตปฺปตฺเต ยํ ตชฺชํ จิตฺตํ, ตํ ปฏิปฺปสฺสทฺธํ. ยา ปุพฺเพ วีมํสา อโหสิ อรหตฺต\-ปฺปตฺติยา, อรหตฺตปฺปตฺเต ยา ตชฺชา วีมํสา, สา ปฏิปฺปสฺสทฺธา. ตํ กิํ มญฺญสิ พฺราหฺมณ, อิติ เอวํ สนฺเต สนฺตกํ วา โหติ โน อสนฺตกํ วาติ.}

\prose{อทฺธา โภ อานนฺท เอวํ สนฺเต สนฺตกํ โหติ โน อสนฺตกํ. อภิกฺกนฺตํ โภ อานนฺท, อภิกฺกนฺตํ โภ อานนฺท, เสยฺยถาปิ โภ อานนฺท นิกฺกุชฺชิตํ วา อุกฺกุชฺเชยฺย, ปฏิจฺฉนฺนํ วา วิวเรยฺย, มูฬฺหสฺส วา มคฺคํ อาจิกฺเขยฺย, อนฺธกาเร วา เตลปชฺโชตํ ธาเรยฺย “จกฺขุมนฺโต รูปานิ ทกฺขนฺตี”ติ, เอวเมวํ โภตา อานนฺเทน อเนก\-ปริยาเยน ธมฺโม ปกาสิโต, เอสาหํ โภ อานนฺท ตํ ภวนฺตํ โคตมํ สรณํ คจฺฉามิ ธมฺมญฺจ ภิกฺขุสํฆญฺจ, อุปาสกํ มํ ภวํ อานนฺโท ธาเรตุ อชฺชตคฺเค ปาณุเปตํ สรณํ คตนฺติ.  ปญฺจมํ.}

\csromanheader{6}{7. อิทฺธิปาท\-สํยุตฺต}
\csromanheader{2}{6. ปฐม\-สมณพฺราหฺมณ\-สุตฺต}
\csromanmatikaanchor{mtk.179}
\csromantocmarkref{subparagraph}{6-9. ปฐมสมณพฺราหฺมณสุตฺตาทีนิ}{mtk.179}
\bookmark[dest={mtk.179},level=5]{6-9. ปฐมสมณพฺราหฺมณสุตฺตาทีนิ}
\proseitem{828}{\csromanfolio{239}เย หิ เกจิ ภิกฺขเว อตีตมทฺธานํ สมณา วา พฺราหฺมณา วา มหิทฺธิกา อเหสุํ มหานุภาวา, สพฺเพ เต จตุนฺนํ อิทฺธิปาทานํ ภาวิตตฺตา พหุลีกตตฺตา. เย หิ เกจิ ภิกฺขเว อนาคตมทฺธานํ สมณา วา พฺราหฺมณา วา มหิทฺธิกา ภวิสฺสนฺติ มหานุภาวา, สพฺเพ เต จตุนฺนํ อิทฺธิปาทานํ ภาวิตตฺตา พหุลีกตตฺตา. เย หิ เกจิ ภิกฺขเว เอตรหิ สมณา วา พฺราหฺมณา วา มหิทฺธิกา มหานุภาวา, สพฺเพ เต จตุนฺนํ อิทฺธิปาทานํ ภาวิตตฺตา พหุลีกตตฺตา.}

\prose{กตเมสํ จตุนฺนํ, อิธ ภิกฺขเว ภิกฺขุ ฉนฺท\-สมาธิ\-ปฺปธาน\-สงฺขาร\-สมนฺนาคตํ อิทฺธิปาทํ ภาเวติ. วีริยสมาธิ ฯเปฯ. จิตฺตสมาธิ. วีมํสา\-สมาธิ\-ปฺปธาน\-สงฺขาร\-สมนฺนาคตํ อิทฺธิปาทํ ภาเวติ. เย หิ เกจิ ภิกฺขเว อตีตมทฺธานํ สมณา วา พฺราหฺมณา วา มหิทฺธิกา อเหสุํ มหานุภาวา, สพฺเพ เต อิเมสํเยว จตุนฺนํ อิทฺธิปาทานํ ภาวิตตฺตา พหุลีกตตฺตา. เย หิ เกจิ ภิกฺขเว อนาคตมทฺธานํ สมณา วา พฺราหฺมณา วา มหิทฺธิกา ภวิสฺสนฺติ มหานุภาวา, สพฺเพ เต อิเมสํเยว จตุนฺนํ อิทฺธิปาทานํ ภาวิตตฺตา พหุลีกตตฺตา. เย หิ เกจิ ภิกฺขเว เอตรหิ สมณา วา พฺราหฺมณา วา มหิทฺธิกา มหานุภาวา, สพฺเพ เต อิเมสํเยว จตุนฺนํ อิทฺธิปาทานํ ภาวิตตฺตา พหุลีกตตฺตาติ.  ฉฏฺฐํ.}

\csromanheader{2}{7. ทุติย\-สมณพฺราหฺมณ\-สุตฺต}
\proseitem{829}{เย หิ เกจิ ภิกฺขเว อตีตมทฺธานํ สมณา วา พฺราหฺมณา วา อเนกวิหิตํ อิทฺธิวิธํ ปจฺจนุโภสุํ, เอโกปิ หุตฺวา พหุธา อเหสุํ, พหุธาปิ หุตฺวา เอโก อเหสุํ, อาวิภาวํ ติโรภาวํ ติโรกุฏฺฏํ ติโรปาการํ ติโรปพฺพตํ อสชฺชมานา อคมํสุ เสยฺยถาปิ อากาเส. ปถวิยาปิ อุมฺมุชฺช\-นิมุชฺชํ อกํสุ เสยฺยถาปิ อุทเก. อุทเกปิ อภิชฺชมาเน\csromansharedfootnote{c510b00b9dfaed74}{อภิชฺชมานา (สี, อิ, ก)} อคมํสุ เสยฺยถาปิ ปถวิยํ. อากาเสปิ ปลฺลงฺเกน กมิํสุ เสยฺยถาปิ ปกฺขี สกุโณ. อิเมปิ จนฺทิมสูริเย เอวํมหิทฺธิเก เอวํ\-มหา\-นุภาเว ปาณินา ปริมสิํสุ\csromansharedfootnote{88ac0283ba022da6}{ปรามสิํสุ (สฺยา, กํ, ก)} ปริมชฺชิํสุ, ยาว พฺรหฺมโลกาปิ กาเยน วสํ วตฺเตสุํ, สพฺเพ เต จตุนฺนํ อิทฺธิปาทานํ ภาวิตตฺตา พหุลีกตตฺตา.}

\prose{\csromanfolio{240}เย หิ เกจิ ภิกฺขเว อนาคตมทฺธานํ สมณา วา พฺราหฺมณา วา อเนกวิหิตํ อิทฺธิวิธํ ปจฺจนุโภสฺสนฺติ, เอโกปิ หุตฺวา พหุธา ภวิสฺสนฺติ, พหุธาปิ หุตฺวา เอโก ภวิสฺสนฺติ, อาวิภาวํ ติโรภาวํ ติโรกุฏฺฏํ ติโรปาการํ ติโรปพฺพตํ อสชฺชมานา คมิสฺสนฺติ เสยฺยถาปิ อากาเส. ปถวิยาปิ อุมฺมุชฺช\-นิมุชฺชํ กริสฺสนฺติ เสยฺยถาปิ อุทเก. อุทเกปิ อภิชฺชมาเน คมิสฺสนฺติ เสยฺยถาปิ ปถวิยํ. อากาเสปิ ปลฺลงฺเกน กมิสฺสนฺติ เสยฺยถาปิ ปกฺขี สกุโณ. อิเมปิ จนฺทิมสูริเย เอวํมหิทฺธิเก เอวํ\-มหา\-นุภาเว ปาณินา ปริมสิสฺสนฺติ ปริมชฺชิสฺสนฺติ, ยาว พฺรหฺมโลกาปิ กาเยน วสํ วตฺติสฺสนฺติ, สพฺเพ เต จตุนฺนํ อิทฺธิปาทานํ ภาวิตตฺตา พหุลีกตตฺตา.}

\prose{เย หิ เกจิ ภิกฺขเว เอตรหิ สมณา วา พฺราหฺมณา วา อเนกวิหิตํ อิทฺธิวิธํ ปจฺจนุโภนฺติ, เอโกปิ หุตฺวา พหุธา โหนฺติ, พหุธาปิ หุตฺวา เอโก โหนฺติ, อาวิภาวํ ติโรภาวํ ติโรกุฏฺฏํ ติโรปาการํ ติโรปพฺพตํ อสชฺชมานา คจฺฉนฺติ เสยฺยถาปิ อากาเส. ปถวิยาปิ อุมฺมุชฺช\-นิมุชฺชํ กโรนฺติ เสยฺยถาปิ อุทเก. อุทเกปิ อภิชฺชมาเน คจฺฉนฺติ เสยฺยถาปิ ปถวิยํ. อากาเสปิ ปลฺลงฺเกน กมนฺติ เสยฺยถาปิ ปกฺขี สกุโณ. อิเมปิ จนฺทิมสูริเย เอวํมหิทฺธิเก เอวํ\-มหา\-นุภาเว ปาณินา ปริมสนฺติ\csromansharedfootnote{bab39c8b672347b9}{ปรามสนฺติ (สฺยา, กํ)} ปริมชฺชนฺติ, ยาว พฺรหฺมโลกาปิ กาเยน วสํ วตฺเตนฺติ, สพฺเพ เต จตุนฺนํ อิทฺธิปาทานํ ภาวิตตฺตา พหุลีกตตฺตาติ.}

\prose{กตเมสํ จตุนฺนํ, อิธ ภิกฺขเว ภิกฺขุ ฉนฺท\-สมาธิ\-ปฺปธาน\-สงฺขาร\-สมนฺนาคตํ อิทฺธิปาทํ ภาเวติ. วีริยสมาธิ ฯเปฯ. จิตฺตสมาธิ. วีมํสา\-สมาธิ\-ปฺปธาน\-สงฺขาร\-สมนฺนาคตํ อิทฺธิปาทํ ภาเวติ. เย หิ เกจิ ภิกฺขเว อตีตมทฺธานํ สมณา วา พฺราหฺมณา วา อเนกวิหิตํ อิทฺธิวิธํ ปจฺจนุโภสุํ, เอโกปิ หุตฺวา พหุธา อเหสุํ ฯเปฯ ยาว พฺรหฺมโลกาปิ กาเยน วสํ วตฺเตสุํ, สพฺเพ เต อิเมสํเยว จตุนฺนํ อิทฺธิปาทานํ ภาวิตตฺตา พหุลีกตตฺตา.}

\prose{เย หิ เกจิ ภิกฺขเว อนาคตมทฺธานํ สมณา วา พฺราหฺมณา วา อเนกวิหิตํ อิทฺธิวิธํ ปจฺจนุโภสฺสนฺติ, เอโกปิ หุตฺวา พหุธา ภวิสฺสนฺติ ฯเปฯ ยาว พฺรหฺมโลกาปิ กาเยน วสํ วตฺติสฺสนฺติ, สพฺเพ เต อิเมสํเยว จตุนฺนํ อิทฺธิปาทานํ ภาวิตตฺตา พหุลีกตตฺตา.}

\prose{\csromanfolio{241}เย หิ เกจิ ภิกฺขเว เอตรหิ สมณา วา พฺราหฺมณา วา อเนกวิหิตํ อิทฺธิวิธํ ปจฺจนุโภนฺติ, เอโกปิ หุตฺวา พหุธา โหนฺติ ฯเปฯ ยาว พฺรหฺมโลกาปิ กาเยน วสํ วตฺเตนฺติ, สพฺเพ เต อิเมสํเยว จตุนฺนํ อิทฺธิปาทานํ ภาวิตตฺตา พหุลีกตตฺตาติ.  สตฺตมํ.}

\csromanheader{2}{8. \textbf{ภิกฺขุสุตฺต}}
\proseitem{830}{จตุนฺนํ ภิกฺขเว อิทฺธิปาทานํ ภาวิตตฺตา พหุลีกตตฺตา ภิกฺขุ อาสวานํ ขยา อนาสวํ เจโตวิมุตฺติํ ปญฺญาวิมุตฺติํ ทิฏฺเฐว ธมฺเม สยํ อภิญฺญา สจฺฉิกตฺวา อุปสมฺปชฺช วิหรติ.}

\prose{กตเมสํ จตุนฺนํ, อิธ ภิกฺขเว ภิกฺขุ ฉนฺท\-สมาธิ\-ปฺปธาน\-สงฺขาร\-สมนฺนาคตํ อิทฺธิปาทํ ภาเวติ. วีริยสมาธิ ฯเปฯ. จิตฺตสมาธิ. วีมํสา\-สมาธิ\-ปฺปธาน\-สงฺขาร\-สมนฺนาคตํ อิทฺธิปาทํ ภาเวติ. อิเมสํ โข ภิกฺขเว จตุนฺนํ อิทฺธิปาทานํ ภาวิตตฺตา พหุลีกตตฺตา ภิกฺขุ อาสวานํ ขยา อนาสวํ เจโตวิมุตฺติํ ปญฺญาวิมุตฺติํ ทิฏฺเฐว ธมฺเม สยํ อภิญฺญา สจฺฉิกตฺวา อุปสมฺปชฺช วิหรตีติ.  อฏฺฐมํ.}

\csromanheader{2}{9. อิทฺธา\-ทิ\-เทสนา\-สุตฺต}
\proseitem{831}{อิทฺธิํ โว ภิกฺขเว เทเสสฺสามิ อิทฺธิปาทญฺจ อิทฺธิปาทภาวนญฺจ อิทฺธิปาทภาวนาคามินิญฺจ ปฏิปทํ, ตํ สุณาถ.}

\prose{กตมา จ ภิกฺขเว อิทฺธิ, อิธ ภิกฺขเว ภิกฺขุ อเนกวิหิตํ อิทฺธิวิธํ ปจฺจนุโภติ, เอโกปิ หุตฺวา พหุธา โหติ, พหุธาปิ หุตฺวา เอโก โหติ ฯเปฯ ยาว พฺรหฺมโลกาปิ กาเยน วสํ วตฺเตติ, อยํ วุจฺจติ ภิกฺขเว อิทฺธิ.}

\prose{กตโม จ ภิกฺขเว อิทฺธิปาโท, โย โส ภิกฺขเว มคฺโค ยา ปฏิปทา อิทฺธิลาภาย อิทฺธิ\-ปฏิลาภาย สํวตฺตติ, อยํ วุจฺจติ ภิกฺขเว อิทฺธิปาโท.}

\prose{กตมา จ ภิกฺขเว อิทฺธิปาท\-ภาวนา, อิธ ภิกฺขเว ภิกฺขุ ฉนฺท\-สมาธิ\-ปฺปธาน\-สงฺขาร\-สมนฺนาคตํ อิทฺธิปาทํ ภาเวติ. วีริยสมาธิ ฯเปฯ. จิตฺตสมาธิ. วีมํสา\-สมาธิ\-ปฺปธาน\-สงฺขาร\-สมนฺนาคตํ อิทฺธิปาทํ ภาเวติ, อยํ วุจฺจติ ภิกฺขเว อิทฺธิปาท\-ภาวนา.}

\prose{\csromanfolio{242}กตมา จ ภิกฺขเว อิทฺธิปาท\-ภาวนาคามินี ปฏิปทา, อยเมว อริโย อฏฺฐงฺคิโก มคฺโค. เสยฺยถิทํ, สมฺมาทิฏฺฐิ สมฺมาสงฺกปฺโป สมฺมาวาจา สมฺมากมฺมนฺโต สมฺมาอาชีโว สมฺมาวายาโม สมฺมาสติ สมฺมาสมาธิ. อยํ วุจฺจติ ภิกฺขเว อิทฺธิปาท\-ภาวนาคามินี ปฏิปทาติ.  นวมํ.}

\csromanmatikaanchor{mtk.180}
\csromantocmarkref{subparagraph}{10. วิภงฺคสุตฺต}{mtk.180}
\bookmark[dest={mtk.180},level=5]{10. วิภงฺคสุตฺต}
\csromanheader{6}{10. \textbf{วิภงฺคสุตฺต}}
\proseitem{832}{จตฺตาโรเม ภิกฺขเว อิทฺธิปาทา ภาวิตา พหุลีกตา มหปฺผลา โหนฺติ มหานิสํสา.}

\prose{กถํ ภาวิตา จ ภิกฺขเว จตฺตาโร อิทฺธิปาทา กถํ พหุลีกตา มหปฺผลา โหนฺติ มหานิสํสา, อิธ ภิกฺขเว ภิกฺขุ ฉนฺท\-สมาธิ\-ปฺปธาน\-สงฺขาร\-สมนฺนาคตํ อิทฺธิปาทํ ภาเวติ, อิติ เม ฉนฺโท น จ อติลีโน ภวิสฺสติ, น จ อติปฺปคฺคหิโต ภวิสฺสติ, น จ อชฺฌตฺตํ สํขิตฺโต ภวิสฺสติ, น จ พหิทฺธา วิกฺขิตฺโต ภวิสฺสติ. ปจฺฉา\-ปุเร\-สญฺญี จ วิหรติ, ยถา ปุเร ตถา ปจฺฉา, ยถา ปจฺฉา ตถา ปุเร. ยถา อโธ ตถา อุทฺธํ, ยถา อุทฺธํ ตถา อโธ. ยถา ทิวา ตถา รตฺติํ, ยถา รตฺติํ ตถา ทิวา. อิติ วิวเฏน เจตสา อปริโยนทฺเธน สปฺปภาสํ จิตฺตํ ภาเวติ. วีริยสมาธิ ฯเปฯ. จิตฺตสมาธิ. วีมํสา\-สมาธิ\-ปฺปธาน\-สงฺขาร\-สมนฺนาคตํ อิทฺธิปาทํ ภาเวติ, อิติ เม วีมํสา น จ อติลีนา ภวิสฺสติ, น จ อติปฺปคฺคหิตา ภวิสฺสติ, น จ อชฺฌตฺตํ สํขิตฺตา ภวิสฺสติ, น จ พหิทฺธา วิกฺขิตฺตา ภวิสฺสติ. ปจฺฉา\-ปุเร\-สญฺญี จ วิหรติ, ยถา ปุเร ตถา ปจฺฉา, ยถา ปจฺฉา ตถา ปุเร. ยถา อโธ ตถา อุทฺธํ, ยถา อุทฺธํ ตถา อโธ. ยถา ทิวา ตถา รตฺติํ, ยถา รตฺติํ ตถา ทิวา. อิติ วิวเฏน เจตสา อปริโยนทฺเธน สปฺปภาสํ จิตฺตํ ภาเวติ.}

\prose{กตโม จ ภิกฺขเว อติลีโน ฉนฺโท, โย ภิกฺขเว ฉนฺโท โกสชฺช\-สหคโต โกสชฺช\-สมฺปยุตฺโต, อยํ วุจฺจติ ภิกฺขเว อติลีโน ฉนฺโท.}

\prose{กตโม จ ภิกฺขเว อติปฺปคฺคหิโต ฉนฺโท, โย ภิกฺขเว ฉนฺโท อุทฺธจฺจ\-สหคโต อุทฺธจฺจ\-สมฺปยุตฺโต, อยํ วุจฺจติ ภิกฺขเว อติปฺปคฺคหิโต ฉนฺโท.}

\prose{\csromanfolio{243}กตโม จ ภิกฺขเว อชฺฌตฺตํ สํขิตฺโต ฉนฺโท, โย ภิกฺขเว ฉนฺโท ถินมิทฺธ\-สหคโต ถินมิทฺธ\-สมฺปยุตฺโต, อยํ วุจฺจติ ภิกฺขเว อชฺฌตฺตํ สํขิตฺโต ฉนฺโท.}

\prose{กตโม จ ภิกฺขเว พหิทฺธา วิกฺขิตฺโต ฉนฺโท, โย ภิกฺขเว ฉนฺโท พหิทฺธา ปญฺจ กามคุเณ อารพฺภ อนุวิกฺขิตฺโต อนุวิสโฏ, อยํ วุจฺจติ ภิกฺขเว พหิทฺธา วิกฺขิตฺโต ฉนฺโท.}

\prose{กถญฺจ ภิกฺขเว ภิกฺขุ ปจฺฉา\-ปุเร\-สญฺญี จ วิหรติ, ยถา ปุเร ตถา ปจฺฉา, ยถา ปจฺฉา ตถา ปุเร, อิธ ภิกฺขเว ภิกฺขุโน ปจฺฉาปุเร\-สญฺญา สุคฺคหิตา โหติ สุมนสิกตา สูปธาริตา สุปฺปฏิวิทฺธา ปญฺญาย. เอวํ โข ภิกฺขเว ภิกฺขุ ปจฺฉา\-ปุเร\-สญฺญี จ วิหรติ, ยถา ปุเร ตถา ปจฺฉา, ยถา ปจฺฉา ตถา ปุเร.}

\prose{กถญฺจ ภิกฺขเว ภิกฺขุ ยถา อโธ ตถา อุทฺธํ, ยถา อุทฺธํ ตถา อโธ วิหรติ, อิธ ภิกฺขเว ภิกฺขุ อิมเมว กายํ อุทฺธํ ปาทตลา อโธ เกสมตฺถกา ตจปริยนฺตํ ปูรํ นานปฺปการสฺส อสุจิโน ปจฺจเวกฺขติ “อตฺถิ อิมสฺมิํ กาเย เกสา โลมา นขา ทนฺตา ตโจ, มํสํ นฺหารุ อฏฺฐิ อฏฺฐิมิญฺชํ วกฺกํ, หทยํ ยกนํ กิโลมกํ ปิหกํ ปปฺผาสํ, อนฺตํ อนฺตคุณํ อุทริยํ กรีสํ, ปิตฺตํ เสมฺหํ ปุพฺโพ โลหิตํ เสโท เมโท, อสฺสุ วสา เขโฬ สิงฺฆาณิกา ลสิกา มุตฺตนฺ”ติ. เอวํ โข ภิกฺขเว ภิกฺขุ ยถา อโธ ตถา อุทฺธํ, ยถา อุทฺธํ ตถา อโธ วิหรติ.}

\prose{กถญฺจ ภิกฺขเว ภิกฺขุ ยถา ทิวา ตถา รตฺติํ, ยถา รตฺติํ ตถา ทิวา วิหรติ, อิธ ภิกฺขเว ภิกฺขุ เยหิ อากาเรหิ เยหิ ลิงฺเคหิ เยหิ นิมิตฺเตหิ ทิวา ฉนฺท\-สมาธิ\-ปฺปธาน\-สงฺขาร\-สมนฺนาคตํ อิทฺธิปาทํ ภาเวติ, โส เตหิ อากาเรหิ เตหิ ลิงฺเคหิ เตหิ นิมิตฺเตหิ รตฺติํ ฉนฺท\-สมาธิ\-ปฺปธาน\-สงฺขาร\-สมนฺนาคตํ อิทฺธิปาทํ ภาเวติ. เยหิ วา ปน อากาเรหิ เยหิ ลิงฺเคหิ เยหิ นิมิตฺเตหิ รตฺติํ ฉนฺท\-สมาธิ\-ปฺปธาน\-สงฺขาร\-สมนฺนาคตํ อิทฺธิปาทํ ภาเวติ, โส เตหิ อากาเรหิ เตหิ ลิงฺเคหิ เตหิ นิมิตฺเตหิ ทิวา ฉนฺท\-สมาธิ\-ปฺปธาน\-สงฺขาร\-สมนฺนาคตํ อิทฺธิปาทํ ภาเวติ. เอวํ โข ภิกฺขเว ภิกฺขุ ยถา ทิวา ตถา รตฺติํ, ยถา รตฺติํ ตถา ทิวา วิหรติ.}

\prose{\csromanfolio{244}กถญฺจ ภิกฺขเว ภิกฺขุ วิวเฏน เจตสา อปริโยนทฺเธน สปฺปภาสํ จิตฺตํ ภาเวติ, อิธ ภิกฺขเว ภิกฺขุโน อาโลกสญฺญา สุคฺคหิตา โหติ, ทิวาสญฺญา สฺวาธิฏฺฐิตา. เอวํ โข ภิกฺขเว ภิกฺขุ วิวเฏน เจตสา อปริโยนทฺเธน สปฺปภาสํ จิตฺตํ ภาเวติ.}

\prose{กถมญฺจ ภิกฺขเว อติลีนํ วีริยํ, ยํ ภิกฺขเว วีริยํ โกสชฺช\-สหคตํ โกสชฺช\-สมฺปยุตฺตํ, อิทํ วุจฺจติ ภิกฺขเว อติลีนํ วีริยํ.}

\prose{กตมญฺจ ภิกฺขเว อติปฺปคฺคหิตํ วีริยํ, ยํ ภิกฺขเว วีริยํ อุทฺธจฺจ\-สหคตํ อุทฺธจฺจ\-สมฺปยุตฺตํ, อิทํ วุจฺจติ ภิกฺขเว อติปฺปคฺคหิตํ วีริยํ.}

\prose{กตมญฺจ ภิกฺขเว อชฺฌตฺตํ สํขิตฺตํ วีริยํ, ยํ ภิกฺขเว วีริยํ ถินมิทฺธ\-สหคตํ ถินมิทฺธ\-สมฺปยุตฺตํ, อิทํ วุจฺจติ ภิกฺขเว อชฺฌตฺตํ สํขิตฺตํ วีริยํ.}

\prose{กตมญฺจ ภิกฺขเว พหิทฺธา วิกฺขิตฺตํ วีริยํ, ยํ ภิกฺขเว วีริยํ พหิทฺธา ปญฺจ กามคุเณ อารพฺภ อนุวิกฺขิตฺตํ อนุวิสฏํ, อิทํ วุจฺจติ ภิกฺขเว พหิทฺธา วิกฺขิตฺตํ วีริยํ ฯเปฯ.}

\prose{กถญฺจ ภิกฺขเว ภิกฺขุ วิวเฏน เจตสา อปริโยนทฺเธน สปฺปภาสํ จิตฺตํ ภาเวติ, อิธ ภิกฺขเว ภิกฺขุโน อาโลกสญฺญา สุคฺคหิตา โหติ, ทิวาสญฺญา สฺวาธิฏฺฐิตา. เอวํ โข ภิกฺขเว ภิกฺขุ วิวเฏน เจตสา อปริโยนทฺเธน สปฺปภาสํ จิตฺตํ ภาเวติ.}

\prose{กถมญฺจ ภิกฺขเว อติลีนํ จิตฺตํ, ยํ ภิกฺขเว จิตฺตํ โกสชฺช\-สหคตํ โกสชฺช\-สมฺปยุตฺตํ, อิทํ วุจฺจติ ภิกฺขเว อติลีนํ จิตฺตํ.}

\prose{กตมญฺจ ภิกฺขเว อติปฺปคฺคหิตํ จิตฺตํ, ยํ ภิกฺขเว จิตฺตํ อุทฺธจฺจ\-สหคตํ อุทฺธจฺจ\-สมฺปยุตฺตํ, อิทํ วุจฺจติ ภิกฺขเว อติปฺปคฺคหิตํ จิตฺตํ.}

\prose{กตมญฺจ ภิกฺขเว อชฺฌตฺตํ สํขิตฺตํ จิตฺตํ, ยํ ภิกฺขเว จิตฺตํ ถินมิทฺธ\-สหคตํ ถินมิทฺธ\-สมฺปยุตฺตํ, อิทํ วุจฺจติ ภิกฺขเว อชฺฌตฺตํ สํขิตฺตํ จิตฺตํ.}

\prose{กตมญฺจ ภิกฺขเว พหิทฺธา วิกฺขิตฺตํ จิตฺตํ, ยํ ภิกฺขเว จิตฺตํ พหิทฺธา ปญฺจ กามคุเณ อารพฺภ อนุวิกฺขิตฺตํ อนุวิสฏํ, อิทํ วุจฺจติ ภิกฺขเว พหิทฺธา วิกฺขิตฺตํ จิตฺตํ ฯเปฯ. เอวํ โข ภิกฺขเว ภิกฺขุ วิวเฏน เจตสา อปริโยนทฺเธน สปฺปภาสํ จิตฺตํ ภาเวติ.}

\prose{กตมา จ ภิกฺขเว อติลีนา วีมํสา, ยา ภิกฺขเว วีมํสา โกสชฺช\-สหคตา โกสชฺช\-สมฺปยุตฺตา, อยํ วุจฺจติ ภิกฺขเว อติลีนา วิมํสา.}

\prose{\csromanfolio{245}กตมา จ ภิกฺขเว อติปฺปคฺคหิตา วีมํสา, ยา ภิกฺขเว วีมํสา อุทฺธจฺจ\-สหคตา อุทฺธจฺจ\-สมฺปยุตฺตา, อยํ วุจฺจติ ภิกฺขเว อติปฺปคฺคหิตา วีมํสา.}

\prose{กตมา จ ภิกฺขเว อชฺฌตฺตํ สํขิตฺตา วีมํสา, ยา ภิกฺขเว วีมํสา ถินมิทฺธ\-สหคตา ถินมิทฺธ\-สมฺปยุตฺตา, อยํ วุจฺจติ ภิกฺขเว อชฺฌตฺตํ สํขิตฺตา วีมํสา.}

\prose{กตมา จ ภิกฺขเว พหิทฺธา วิกฺขิตฺตา วีมํสา, ยา ภิกฺขเว วีมํสา พหิทฺธา ปญฺจ กามคุเณ อารพฺภ อนุวิกฺขิตฺตา อนุวิสฏา, อยํ วุจฺจติ ภิกฺขเว พหิทฺธา วิกฺขิตฺตา วีมํสา ฯเปฯ. เอวํ โข ภิกฺขเว ภิกฺขุ วิวเฏน เจตสา อปริโยนทฺเธน สปฺปภาสํ จิตฺตํ ภาเวติ, เอวํ ภาวิตา โข ภิกฺขเว จตฺตาโร อิทฺธิปาทา เอวํ พหุลีกตา มหปฺผลา โหนฺติ มหานิสํสา.}

\prosewithclosermidruled{เอวํ ภาวิเตสุ โข ภิกฺขเว ภิกฺขุ จตูสุ อิทฺธิปาเทสุ เอวํ พหุลีกเตสุ อเนกวิหิตํ อิทฺธิวิธํ ปจฺจนุโภติ, เอโกปิ หุตฺวา พหุธา โหติ, พหุธาปิ หุตฺวา เอโก โหติ ฯเปฯ ยาว พฺรหฺมโลกาปิ กาเยน วสํ วตฺเตติ. เอวํ ภาวิเตสุ โข ภิกฺขเว ภิกฺขุ จตูสุ อิทฺธิปาเทสุ เอวํ พหุลีกเตสุ อาสวานํ ขยา อนาสวํ เจโตวิมุตฺติํ ปญฺญาวิมุตฺติํ ทิฏฺเฐว ธมฺเม สยํ อภิญฺญา สจฺฉิกตฺวา อุปสมฺปชฺช วิหรตีติ.  ทสมํ.}{ปาสาทกมฺปน\-วคฺโค ทุติโย.}
\csromancenter{\textbf{ตสฺสุทฺทานํ}}
\setlength{\csromangathapagewidth}{0pt}
\setlength{\csromangathawakwidth}{0pt}
\csromangathameasuregroup{ปุพฺพํ มหปฺผลํ ฉนฺทํ, \\ เทฺว สมณพฺราหฺมณา ภิกฺขุ,}{\csromangathabat{ปุพฺพํ มหปฺผลํ ฉนฺทํ,}{โมคฺคลฺลานญฺจ อุณฺณาภํ.} \\ \csromangathabat{เทฺว สมณพฺราหฺมณา ภิกฺขุ,}{เทสนา วิภงฺเคน จาติ.}}
\csromangathasetleft{ปุพฺพํ มหปฺผลํ ฉนฺทํ, \\ เทฺว สมณพฺราหฺมณา ภิกฺขุ,}
\csromangathagroup{\csromangathastanza{\csromangathabat{ปุพฺพํ มหปฺผลํ ฉนฺทํ,}{โมคฺคลฺลานญฺจ อุณฺณาภํ.} \\ \csromangathabat{เทฺว สมณพฺราหฺมณา ภิกฺขุ,}{เทสนา วิภงฺเคน จาติ.}}}

\csromanmatikaanchor{mtk.181}
\csromantocmarknopage{subparagraph}{3. อโยคุฬวคฺค}{mtk.181}
\bookmark[dest={mtk.181},level=5]{3. อโยคุฬวคฺค}
\csromanheader{6}{3. \textbf{อโยคุฬวคฺค}}
\csromanmatikaanchor{mtk.182}
\csromantocmarkref{subparagraph}{1. มคฺคสุตฺต}{mtk.182}
\bookmark[dest={mtk.182},level=5]{1. มคฺคสุตฺต}
\csromanheader{6}{1. \textbf{มคฺคสุตฺต}}
\csromanmatikaanchor{mtk.183}
\csromantocmarkref{subparagraph}{2-10. อโยคุฬสุตฺตาทีนิ}{mtk.183}
\bookmark[dest={mtk.183},level=5]{2-10. อโยคุฬสุตฺตาทีนิ}
\proseitem{833}{สาวตฺถิ\-นิทานํ. ปุพฺเพว เม ภิกฺขเว สมฺโพธา อนภิสมฺพุทฺธสฺส โพธิสตฺตสฺเสว สโต เอตทโหสิ “โก นุ โข มคฺโค กา ปฏิปทา \csromanfolio{246}อิทฺธิปาทภาวนายา”ติ. ตสฺส มยฺหํ ภิกฺขเว เอตทโหสิ\csromandash{}อิธ ภิกฺขุ ฉนฺท\-สมาธิ\-ปฺปธาน\-สงฺขาร\-สมนฺนาคตํ อิทฺธิปาทํ ภาเวติ, อิติ เม ฉนฺโท น จ อติลีโน ภวิสฺสติ, น จ อติปฺปคฺคหิโต ภวิสฺสติ, น จ อชฺฌตฺตํ สํขิตฺโต ภวิสฺสติ, น จ พหิทฺธา วิกฺขิตฺโต ภวิสฺสติ. ปจฺฉา\-ปุเร\-สญฺญี จ วิหรติ, ยถา ปุเร ตถา ปจฺฉา, ยถา ปจฺฉา ตถา ปุเร. ยถา อโธ ตถา อุทฺธํ, ยถา อุทฺธํ ตถา อโธ. ยถา ทิวา ตถา รตฺติํ ยถา รตฺติํ ตถา ทิวา. อิติ วิวเฏน เจตสา อปริโยนทฺเธน สปฺปภาสํ จิตฺตํ ภาเวติ. วีริยสมาธิ ฯเปฯ. จิตฺตสมาธิ. วีมํสา\-สมาธิ\-ปฺปธาน\-สงฺขาร\-สมนฺนาคตํ อิทฺธิปาทํ ภาเวติ, อิติ เม วีมํสา น จ อติลีนา ภวิสฺสติ, น จ อติปฺปคฺคหิตา ภวิสฺสติ, น จ อชฺฌตฺตํ สํขิตฺตา ภวิสฺสติ, น จ พหิทฺธา วิกฺขิตฺตา ภวิสฺสติ. ปจฺฉา\-ปุเร\-สญฺญี จ วิหรติ, ยถา ปุเร ตถา ปจฺฉา, ยถา ปจฺฉา ตถา ปุเร. ยถา อโธ ตถา อุทฺธํ, ยถา อุทฺธํ ตถา อโธ. ยถา ทิวา ตถา รตฺติํ, ยถา รตฺติํ ตถา ทิวา. อิติ วิวเฏน เจตสา อปริโยนทฺเธน สปฺปภาสํ จิตฺตํ ภาเวติ.}

\prose{เอวํ ภาวิเตสุ โข ภิกฺขเว ภิกฺขุ จตูสุ อิทฺธิปาเทสุ เอวํ พหุลีกเตสุ อเนกวิหิตํ อิทฺธิวิธํ ปจฺจนุโภติ, เอโกปิ หุตฺวา พหุธา โหติ, พหุธาปิ หุตฺวา เอโก โหติ ฯเปฯ ยาว พฺรหฺมโลกาปิ กาเยน วสํ วตฺเตติ. เอวํ ภาวิเตสุ โข ภิกฺขเว ภิกฺขุ จตูสุ อิทฺธิปาเทสุ เอวํ พหุลีกเตสุ อาสวานํ ขยา อนาสวํ เจโตวิมุตฺติํ ปญฺญาวิมุตฺติํ ทิฏฺเฐว ธมฺเม สยํ อภิญฺญา สจฺฉิกตฺวา อุปสมฺปชฺช วิหรตีติ.  ปฐมํ. (ฉปิ อภิญฺญาโย วิตฺถาเรตพฺพา.)}

\csromanheader{2}{2. \textbf{อโยคุฬสุตฺต}}
\proseitem{834}{สาวตฺถิ\-นิทานํ. อถ โข อายสฺมา อานนฺโท เยน ภควา เตนุปสงฺกมิ, อุปสงฺกมิตฺวา ภควนฺตํ อภิวาเทตฺวา เอกมนฺตํ นิสีทิ, เอกมนฺตํ นิสินฺโน โข อายสฺมา อานนฺโท ภควนฺตํ เอตทโวจ “อภิชานาติ นุ โข ภนฺเต ภควา อิทฺธิยา มโนมเยน กาเยน พฺรหฺมโลกํ อุปสงฺกมิตา”ติ. อภิชานามิ ขฺวาหํ อานนฺท อิทฺธิยา มโนมเยน กาเยน พฺรหฺมโลกํ อุปสงฺกมิตาติ. อภิชานาติ ปน ภนฺเต ภควา อิมินา จาตุมหาภูติเกน กาเยน อิทฺธิยา พฺรหฺมโลกํ}

\prose{\csromanfolio{247}อุปสงฺกมิตาติ. อภิชานามิ ขฺวาหํ อานนฺท อิมินา จาตุมหาภูติเกน\csromansharedfootnote{018a0f94a7a795aa}{จาตุมฺมหาภูติเกน (สี, สฺยา, กํ)} กาเยน อิทฺธิยา พฺรหฺมโลกํ อุปสงฺกมิตาติ.}

\prose{ยญฺจ โข โอมาติ ภนฺเต ภควา อิทฺธิยา มโนมเยน กาเยน พฺรหฺมโลกํ อุปสงฺกมิตุํ, ยญฺจ โข อภิชานาติ ภนฺเต ภควา อิมินา จาตุมหาภูติเกน กาเยน อิทฺธิยา พฺรหฺมโลกํ อุปสงฺกมิตา, ตยิทํ ภนฺเต ภควโต อจฺฉริยํ เจว อพฺภุตํ จาติ. อจฺฉริยา เจว อานนฺท ตถาคตา อจฺฉริยธมฺม\-สมนฺนาคตา จ, อพฺภุตา เจว อานนฺท ตถาคตา อพฺภุตธมฺม\-สมนฺนาคตา จ.}

\prose{ยสฺมิํ อานนฺท สมเย ตถาคโต กายมฺปิ จิตฺเต สโมทหติ\csromansharedfootnote{fdc1ae4ab9eefb9e}{สมาทหติ (สี, สฺยา, อิ)}, จิตฺตมฺปิ กาเย สโมทหติ, สุขสญฺญญฺจ ลหุสญฺญญฺจ กาเย โอกฺกมิตฺวา วิหรติ. ตสฺมิํ อานนฺท สมเย ตถาคตสฺส กาโย ลหุตโร เจว โหติ มุทุตโร จ กมฺมนิยตโร จ ปภสฺสรตโร จ.}

\prose{เสยฺยถาปิ อานนฺท อโยคุโฬ ทิวสํ สนฺตตฺโต ลหุตโร เจว โหติ มุทุตโร จ กมฺมนิยตโร จ ปภสฺสรตโร จ. เอวเมว โข อานนฺท ยสฺมิํ สมเย ตถาคโต กายมฺปิ จิตฺเต สโมทหติ, จิตฺตมฺปิ กาเย สโมทหติ, สุขสญฺญญฺจ ลหุสญฺญญฺจ กาเย โอกฺกมิตฺวา วิหรติ, ตสฺมิํ อานนฺท สมเย ตถาคตสฺส กาโย ลหุตโร เจว โหติ มุทุตโร จ กมฺมนิยตโร จ ปภสฺสรตโร จ.}

\prose{ยสฺมิํ อานนฺท สมเย ตถาคโต กายมฺปิ จิตฺเต สโมทหติ, จิตฺตมฺปิ กาเย สโมทหติ, สุขสญฺญญฺจ ลหุสญฺญญฺจ กาเย โอกฺกมิตฺวา วิหรติ, ตสฺมิํ อานนฺท สมเย ตถาคตสฺส กาโย อปฺปกสิเรเนว ปถวิยา เวหาสํ อพฺภุคฺคจฺฉติ. โส อเนกวิหิตํ อิทฺธิวิธํ ปจฺจนุโภติ, เอโกปิ หุตฺวา พหุธา โหติ, พหุธาปิ หุตฺวา เอโก โหติ ฯเปฯ ยาว พฺรหฺมโลกาปิ กาเยน วสํ วตฺเตติ.}

\prose{เสยฺยถาปิ อานนฺท ตูลปิจุ วา กปฺปาสปิจุ วา ลหุโก วาตูปาทาโน อปฺปกสิเรเนว ปถวิยา เวหาสํ อพฺภุคฺคจฺฉติ. เอวเมว โข อานนฺท ยสฺมิํ สมเย ตถาคโต กายมฺปิ จิตฺเต สโมทหติ,}

\prose{\csromanfolio{248}จิตฺตมฺปิ กาเย สโมทหติ, สุขสญฺญญฺจ ลหุสญฺญญฺจ กาเย โอกฺกมิตฺวา วิหรติ, ตสฺมิํ อานนฺท สมเย ตถาคตสฺส กาโย อปฺปกสิเรเนว ปถวิยา เวหาสํ อพฺภุคฺคจฺฉติ. โส อเนกวิหิตํ อิทฺธิวิธํ ปจฺจนุโภติ. เอโกปิ หุตฺวา พหุธา โหติ ฯเปฯ ยาว พฺรหฺมโลกาปิ กาเยน วสํ วตฺเตตีติ.  ทุติยํ.}

\csromanheader{2}{3. \textbf{ภิกฺขุสุตฺต}}
\proseitem{835}{จตฺตาโรเม ภิกฺขเว อิทฺธิปาทา. กตเม จตฺตาโร, อิธ ภิกฺขเว ภิกฺขุ ฉนฺท\-สมาธิ\-ปฺปธาน\-สงฺขาร\-สมนฺนาคตํ อิทฺธิปาทํ ภาเวติ. วีริยสมาธิ ฯเปฯ. จิตฺตสมาธิ. วีมํสา\-สมาธิ\-ปฺปธาน\-สงฺขาร\-สมนฺนาคตํ อิทฺธิปาทํ ภาเวติ. อิเม โข ภิกฺขเว จตฺตาโร อิทฺธิปาทา, อิเมสํ โข ภิกฺขเว จตุนฺนํ อิทฺธิปาทานํ ภาวิตตฺตา พหุลีกตตฺตา ภิกฺขุ อาสวานํ ขยา อนาสวํ เจโตวิมุตฺติํ ปญฺญาวิมุตฺติํ ทิฏฺเฐว ธมฺเม สยํ อภิญฺญา สจฺฉิกตฺวา อุปสมฺปชฺช วิหรตีติ.  ตติยํ.}

\csromanheader{2}{4. \textbf{สุทฺธิกสุตฺต}}
\proseitem{836}{จตฺตาโรเม ภิกฺขเว อิทฺธิปาทา. กตเม จตฺตาโร, อิธ ภิกฺขเว ภิกฺขุ ฉนฺท\-สมาธิ\-ปฺปธาน\-สงฺขาร\-สมนฺนาคตํ อิทฺธิปาทํ ภาเวติ. วีริยสมาธิ ฯเปฯ. จิตฺตสมาธิ. วีมํสา\-สมาธิ\-ปฺปธาน\-สงฺขาร\-สมนฺนาคตํ อิทฺธิปาทํ ภาเวติ. อิเม โข ภิกฺขเว จตฺตาโร อิทฺธิปาทาติ.  จตุตฺถํ.}

\csromanheader{2}{5. ปฐม\-ผล\-สุตฺต}
\proseitem{837}{จตฺตาโรเม ภิกฺขเว อิทฺธิปาทา. กตเม จตฺตาโร, อิธ ภิกฺขเว ภิกฺขุ ฉนฺท\-สมาธิ\-ปฺปธาน\-สงฺขาร\-สมนฺนาคตํ อิทฺธิปาทํ ภาเวติ. วีริยสมาธิ ฯเปฯ. จิตฺตสมาธิ. วีมํสา\-สมาธิ\-ปฺปธาน\-สงฺขาร\-สมนฺนาคตํ อิทฺธิปาทํ ภาเวติ. อิเม โข ภิกฺขเว จตฺตาโร อิทฺธิปาทา, อิเมสํ โข ภิกฺขเว จตุนฺนํ อิทฺธิปาทานํ ภาวิตตฺตา พหุลีกตตฺตา ภิกฺขุนา ทฺวินฺนํ ผลานํ อญฺญตรํ ผลํ ปาฏิกงฺขํ, ทิฏฺเฐว ธมฺเม อญฺญา, สติ วา อุปาทิเสเส อนาคามิตาติ.  ปญฺจมํ.}

\csromanheader{2}{6. ทุติย\-ผล\-สุตฺต}
\proseitem{838}{\csromanfolio{249}จตฺตาโรเม ภิกฺขเว อิทฺธิปาทา. กตเม จตฺตาโร, อิธ ภิกฺขเว ภิกฺขุ ฉนฺท\-สมาธิ\-ปฺปธาน\-สงฺขาร\-สมนฺนาคตํ อิทฺธิปาทํ ภาเวติ. วีริยสมาธิ ฯเปฯ. จิตฺตสมาธิ. วีมํสา\-สมาธิ\-ปฺปธาน\-สงฺขาร\-สมนฺนาคตํ อิทฺธิปาทํ ภาเวติ. อิเม โข ภิกฺขเว จตฺตาโร อิทฺธิปาทา, อิเมสํ โข ภิกฺขเว จตุนฺนํ อิทฺธิปาทานํ ภาวิตตฺตา พหุลีกตตฺตา สตฺต ผลา สตฺตานิสํสา ปาฏิกงฺขา.}

\prose{กตเม สตฺต ผลา สตฺตานิสํสา, ทิฏฺเฐว ธมฺเม ปฏิกจฺจ\csromansharedfootnote{98ea22c719a3a772}{ปฏิคจฺจ (สี), ปฏิหจฺจ (อิ)} อญฺญํ อาราเธติ. โน เจ ทิฏฺเฐว ธมฺเม ปฏิกจฺจ อญฺญํ อาราเธติ, อถ มรณกาเล อญฺญํ อาราเธติ. โน เจ ทิฏฺเฐว ธมฺเม ปฏิกจฺจ อญฺญํ อาราเธติ, โน เจ มรณกาเล อญฺญํ อาราเธติ, อถ ปญฺจนฺนํ โอรมฺภาคิยานํ สํโยชนานํ ปริกฺขยา อนฺตรา\-ปรินิพฺพายี โหติ, อุปหจฺจ\-ปรินิพฺพายี โหติ, อสงฺขารปรินิพฺพายี โหติ, สสงฺขาร\-ปรินิพฺพายี โหติ, อุทฺธํโสโต โหติ อกนิฏฺฐคามี. อิเมสํ โข ภิกฺขเว จตุนฺนํ อิทฺธิปาทานํ ภาวิตตฺตา พหุลีกตตฺตา อิเม สตฺต ผลา สตฺตานิสํสา ปาฏิกงฺขาติ.  ฉฏฺฐํ.}

\prose{7. ปฐม\-อานนฺท\-สุตฺต}

\proseitem{839}{สาวตฺถิ\-นิทานํ. อถ โข อายสฺมา อานนฺโท เยน ภควา เตนุปสงฺกมิ, อุปสงฺกมิตฺวา ภควนฺตํ อภิวาเทตฺวา เอกมนฺตํ นิสีทิ, เอกมนฺตํ นิสินฺโน โข อายสฺมา อานนฺโท ภควนฺตํ เอตทโวจ\csromandash{}}

\prose{กตมา นุ โข ภนฺเต อิทฺธิ, กตโม อิทฺธิปาโท, กตมา อิทฺธิปาท\-ภาวนา, กตมา อิทฺธิปาท\-ภาวนาคามินี ปฏิปทาติ. อิธานนฺท ภิกฺขุ อเนกวิหิตํ อิทฺธิวิธํ ปจฺจนุโภติ, เอโกปิ หุตฺวา พหุธา โหติ, พหุธาปิ หุตฺวา เอโก โหติ ฯเปฯ ยาว พฺรหฺมโลกาปิ กาเยน วสํ วตฺเตติ, อยํ วุจฺจตานนฺท อิทฺธิ.}

\prose{กตโม จานนฺท อิทฺธิปาโท, โย อานนฺท มคฺโค ยา ปฏิปทา อิทฺธิลาภาย อิทฺธิ\-ปฏิลาภาย สํวตฺตติ, อยํ วุจฺจตานนฺท อิทฺธิปาโท.}

\prose{\csromanfolio{250}กตมา จานนฺท อิทฺธิปาท\-ภาวนา, อิธานนฺท ภิกฺขุ ฉนฺท\-สมาธิ\-ปฺปธาน\-สงฺขาร\-สมนฺนาคตํ อิทฺธิปาทํ ภาเวติ. วีริยสมาธิ ฯเปฯ. จิตฺตสมาธิ. วีมํสา\-สมาธิ\-ปฺปธาน\-สงฺขาร\-สมนฺนาคตํ อิทฺธิปาทํ ภาเวติ, อยํ วุจฺจตานนฺท อิทฺธิปาท\-ภาวนา.}

\prose{กตมา จานนฺท อิทฺธิปาท\-ภาวนาคามินี ปฏิปทา, อยเมว อริโย อฏฺฐงฺคิโก มคฺโค, เสยฺยถิทํ, สมฺมาทิฏฺฐิ ฯเปฯ สมฺมาสมาธิ, อยํ วุจฺจตานนฺท อิทฺธิปาท\-ภาวนาคามินี ปฏิปทาติ.  สตฺตมํ.}

\prose{8. ทุติย\-อานนฺท\-สุตฺต}

\proseitem{840}{เอกมนฺตํ นิสินฺนํ โข อายสฺมนฺตํ อานนฺทํ ภควา เอตทโวจ “กตมา นุ โข อานนฺท อิทฺธิ, กตโม อิทฺธิปาโท, กตมา อิทฺธิปาท\-ภาวนา, กตมา อิทฺธิปาท\-ภาวนาคามินี ปฏิปทา”ติ. ภควํมูลกา โน ภนฺเต ธมฺมา, ภควํ\-เนตฺติกา ฯเปฯ.}

\prose{อิธานนฺท ภิกฺขุ อเนกวิหิตํ อิทฺธิวิธํ ปจฺจนุโภติ, เอโกปิ หุตฺวา พหุธา โหติ ฯเปฯ ยาว พฺรหฺมโลกาปิ กาเยน วสํ วตฺเตติ, อยํ วุจฺจตานนฺท อิทฺธิ.}

\prose{กตโม จานนฺท อิทฺธิปาโท, โย อานนฺท มคฺโค ยา ปฏิปทา อิทฺธิลาภาย อิทฺธิ\-ปฏิลาภาย สํวตฺตติ, อยํ วุจฺจตานนฺท อิทฺธิปาโท.}

\prose{กตมา จานนฺท อิทฺธิปาท\-ภาวนา, อิธานนฺท ภิกฺขุ ฉนฺท\-สมาธิ\-ปฺปธาน\-สงฺขาร\-สมนฺนาคตํ อิทฺธิปาทํ ภาเวติ. วีริยสมาธิ ฯเปฯ. จิตฺตสมาธิ. วีมํสา\-สมาธิ\-ปฺปธาน\-สงฺขาร\-สมนฺนาคตํ อิทฺธิปาทํ ภาเวติ, อยํ วุจฺจตานนฺท อิทฺธิปาท\-ภาวนา.}

\prose{กตมา จานนฺท อิทฺธิปาท\-ภาวนาคามินี ปฏิปทา, อยเมว อริโย อฏฺฐงฺคิโก มคฺโค. เสยฺยถิทํ, สมฺมาทิฏฺฐิ ฯเปฯ สมฺมาสมาธิ, อยํ วุจฺจตานนฺท อิทฺธิปาท\-ภาวนาคามินี ปฏิปทาติ.  อฏฺฐมํ.}

\csromanheader{2}{9. ปฐม\-ภิกฺขุ\-สุตฺต}
\proseitem{841}{อถ โข สมฺพหุลา ภิกฺขู เยน ภควา เตนุ\-ปสงฺกมิํสุ, อุปสงฺกมิตฺวา ภควนฺตํ อภิวาเทตฺวา เอกมนฺตํ นิสีทิํสุ, เอกมนฺตํ นิสินฺนา}

\prose{\csromanfolio{251}โข เต ภิกฺขู ภควนฺตํ เอตทโวจุํ “กตมา นุ โข ภนฺเต อิทฺธิ, กตโม อิทฺธิปาโท, กตมา อิทฺธิปาท\-ภาวนา, กตมา อิทฺธิปาท\-ภาวนาคามินี ปฏิปทา”ติ.}

\prose{อิธ ภิกฺขเว ภิกฺขุ อเนกวิหิตํ อิทฺธิวิธํ ปจฺจนุโภติ, เอโกปิ หุตฺวา พหุธา โหติ ฯเปฯ ยาว พฺรหฺมโลกาปิ กาเยน วสํ วตฺเตติ, อยํ วุจฺจติ ภิกฺขเว อิทฺธิ.}

\prose{กตโม จ ภิกฺขเว อิทฺธิปาโท, โย ภิกฺขเว มคฺโค ยา ปฏิปทา อิทฺธิลาภาย อิทฺธิ\-ปฏิลาภาย สํวตฺตติ, อยํ วุจฺจติ ภิกฺขเว อิทฺธิปาโท.}

\prose{กตมา จ ภิกฺขเว อิทฺธิปาท\-ภาวนา, อิธ ภิกฺขเว ภิกฺขุ ฉนฺท\-สมาธิ\-ปฺปธาน\-สงฺขาร\-สมนฺนาคตํ อิทฺธิปาทํ ภาเวติ. วีริยสมาธิ ฯเปฯ. จิตฺตสมาธิ. วีมํสา\-สมาธิ\-ปฺปธาน\-สงฺขาร\-สมนฺนาคตํ อิทฺธิปาทํ ภาเวติ, อยํ วุจฺจติ ภิกฺขเว อิทฺธิปาท\-ภาวนา.}

\prose{กตมา จ ภิกฺขเว อิทฺธิปาท\-ภาวนาคามินี ปฏิปทา, อยเมว อริโย อฏฺฐงฺคิโก มคฺโค. เสยฺยถิทํ, สมฺมาทิฏฺฐิ ฯเปฯ สมฺมาสมาธิ, อยํ วุจฺจติ ภิกฺขเว อิทฺธิปาท\-ภาวนาคามินี ปฏิปทาติ.  นวมํ.}

\csromanheader{2}{10. ทุติย\-ภิกฺขุ\-สุตฺต}
\proseitem{842}{อถ โข สมฺพหุลา ภิกฺขู เยน ภควา เตนุ\-ปสงฺกมิํสุ ฯเปฯ เอกมนฺตํ นิสินฺเน โข เต ภิกฺขู ภควา เอตทโวจ “กตมา นุ โข ภิกฺขเว อิทฺธิ, กตโม อิทฺธิปาโท, กตมา อิทฺธิปาท\-ภาวนา, กตมา อิทฺธิปาท\-ภาวนาคามินี ปฏิปทา”ติ. ภควํมูลกา โน ภนฺเต ธมฺมา, ภควํ\-เนตฺติกา ฯเปฯ.}

\prose{กตมา จ ภิกฺขเว อิทฺธิ, อิธ ภิกฺขเว ภิกฺขุ อเนกวิหิตํ อิทฺธิวิธํ ปจฺจนุโภติ, เอโกปิ หุตฺวา พหุธา โหติ ฯเปฯ ยาว พฺรหฺมโลกาปิ กาเยน วสํ วตฺเตติ, อยํ วุจฺจติ ภิกฺขเว อิทฺธิ.}

\prose{กตโม จ ภิกฺขเว อิทฺธิปาโท, โย ภิกฺขเว มคฺโค ยา ปฏิปทา อิทฺธิลาภาย อิทฺธิ\-ปฏิลาภาย สํวตฺตติ, อยํ วุจฺจติ ภิกฺขเว อิทฺธิปาโท.}

\prose{กตมา จ ภิกฺขเว อิทฺธิปาท\-ภาวนา, อิธ ภิกฺขเว ภิกฺขุ ฉนฺท\-สมาธิ\-ปฺปธาน\-สงฺขาร\-สมนฺนาคตํ อิทฺธิปาทํ ภาเวติ. วีริยสมาธิ ฯเปฯ. จิตฺตสมาธิ. \csromanfolio{252}วีมํสา\-สมาธิ\-ปฺปธาน\-สงฺขาร\-สมนฺนาคตํ อิทฺธิปาทํ ภาเวติ, อยํ วุจฺจติ ภิกฺขเว อิทฺธิปาท\-ภาวนา.}

\prose{กตมา จ ภิกฺขเว อิทฺธิปาท\-ภาวนาคามินี ปฏิปทา, อยเมว อริโย อฏฺฐงฺคิโก มคฺโค. เสยฺยถิทํ, สมฺมาทิฏฺฐิ ฯเปฯ สมฺมาสมาธิ, อยํ วุจฺจติ ภิกฺขเว อิทฺธิปาท\-ภาวนาคามินี ปฏิปทาติ.  ทสมํ.}

\csromanmatikaanchor{mtk.184}
\csromantocmarkref{subparagraph}{11. โมคฺคลฺลานสุตฺต}{mtk.184}
\bookmark[dest={mtk.184},level=5]{11. โมคฺคลฺลานสุตฺต}
\csromanheader{6}{11. โมคฺคลฺลาน\-สุตฺต}
\proseitem{843}{ตตฺร โข ภควา ภิกฺขู อามนฺเตสิ “ตํ กิํ มญฺญถ ภิกฺขเว, กตเมสํ ธมฺมานํ ภาวิตตฺตา พหุลีกตตฺตา โมคฺคลฺลาโน ภิกฺขุ เอวํมหิทฺธิโก เอวํ\-มหา\-นุภาโว”ติ. ภควํมูลกา โน ภนฺเต ธมฺมา, ภควํ\-เนตฺติกา ฯเปฯ. จตุนฺนํ โข ภิกฺขเว อิทฺธิปาทานํ ภาวิตตฺตา พหุลีกตตฺตา โมคฺคลฺลาโน ภิกฺขุ เอวํมหิทฺธิโก เอวํ\-มหา\-นุภาโว.}

\prose{กตเมสํ จตุนฺนํ, อิธ ภิกฺขเว โมคฺคลฺลาโน ภิกฺขุ ฉนฺท\-สมาธิ\-ปฺปธาน\-สงฺขาร\-สมนฺนาคตํ อิทฺธิปาทํ ภาเวติ, อิติ เม ฉนฺโท น จ อติลีโน ภวิสฺสติ, น จ อติปฺปคฺคหิโต ภวิสฺสติ, น จ อชฺฌตฺตํ สํขิตฺโต ภวิสฺสติ, น จ พหิทฺธา วิกฺขิตฺโต ภวิสฺสติ. ปจฺฉา\-ปุเร\-สญฺญี จ วิหรติ, ยถา ปุเร ตถา ปจฺฉา, ยถา ปจฺฉา ตถา ปุเร. ยถา อโธ ตถา อุทฺธํ, ยถา อุทฺธํ ตถา อโธ. ยถา ทิวา ตถา รตฺติํ, ยถา รตฺติํ ตถา ทิวา. อิติ วิวเฏน เจตสา อปริโยนทฺเธน สปฺปภาสํ จิตฺตํ ภาเวติ. วีริยสมาธิ ฯเปฯ. จิตฺตสมาธิ. วีมํสา\-สมาธิ\-ปฺปธาน\-สงฺขาร\-สมนฺนาคตํ อิทฺธิปาทํ ภาเวติ, อิติ เม วีมํสา น จ อติลีนา ภวิสฺสติ, น จ อติปฺปคฺคหิตา ภวิสฺสติ, น จ อชฺฌตฺตํ สํขิตฺตา ภวิสฺสติ, น จ พหิทฺธา วิกฺขิตฺตา ภวิสฺสติ ฯเปฯ. อิติ วิวเฏน เจตสา อปริโยนทฺเธน สปฺปภาสํ จิตฺตํ ภาเวติ. อิเมสํ โข ภิกฺขเว จตุนฺนํ อิทฺธิปาทานํ ภาวิตตฺตา พหุลีกตตฺตา โมคฺคลฺลาโน ภิกฺขุ เอวํมหิทฺธิโก เอวํ\-มหา\-นุภาโว.}

\prose{อิเมสญฺจ ปน ภิกฺขเว จตุนฺนํ อิทฺธิปาทานํ ภาวิตตฺตา พหุลีกตตฺตา โมคฺคลฺลาโน ภิกฺขุ เอวํ อเนกวิหิตํ อิทฺธิวิธํ ปจฺจนุโภติ, เอโกปิ หุตฺวา พหุธา โหติ, พหุธาปิ หุตฺวา เอโก โหติ ฯเปฯ.}

\prose{\csromanfolio{253}ยาว พฺรหฺมโลกาปิ กาเยน วสํ วตฺเตติ. อิเมสญฺจ ปน ภิกฺขเว จตุนฺนํ อิทฺธิปาทานํ ภาวิตตฺตา พหุลีกตตฺตา โมคฺคลฺลาโน ภิกฺขุ อาสวานํ ขยา อนาสวํ เจโตวิมุตฺติํ ปญฺญาวิมุตฺติํ ทิฏฺเฐว ธมฺเม สยํ อภิญฺญา สจฺฉิกตฺวา อุปสมฺปชฺช วิหรตีติ.  เอกาทสมํ.}

\csromanmatikaanchor{mtk.185}
\csromantocmarkref{subparagraph}{12. ตถาคตสุตฺต}{mtk.185}
\bookmark[dest={mtk.185},level=5]{12. ตถาคตสุตฺต}
\csromanheader{6}{12. \textbf{ตถาคตสุตฺต}}
\proseitem{844}{ตตฺร โข ภควา ภิกฺขู อามนฺเตสิ “ตํ กิํ มญฺญถ ภิกฺขเว, กตเมสํ ธมฺมานํ ภาวิตตฺตา พหุลีกตตฺตา ตถาคโต เอวํมหิทฺธิโก เอวํ\-มหา\-นุภาโว”ติ. ภควํมูลกา โน ภนฺเต ธมฺมา ฯเปฯ. จตุนฺนํ โข ภิกฺขเว อิทฺธิปาทานํ ภาวิตตฺตา พหุลีกตตฺตา ตถาคโต เอวํมหิทฺธิโก เอวํ\-มหา\-นุภาโว.}

\prose{กตเมสํ จตุนฺนํ, อิธ ภิกฺขเว ตถาคโต ฉนฺท\-สมาธิ\-ปฺปธาน\-สงฺขาร\-สมนฺนาคตํ อิทฺธิปาทํ ภาเวติ, อิติ เม ฉนฺโท น จ อติลีโน ภวิสฺสติ, น จ อติปฺปคฺคหิโต ภวิสฺสติ, น จ อชฺฌตฺตํ สํขิตฺโต ภวิสฺสติ, น จ พหิทฺธา วิกฺขิตฺโต ภวิสฺสติ. ปจฺฉา\-ปุเร\-สญฺญี จ วิหรติ, ยถา ปุเร ตถา ปจฺฉา, ยถา ปจฺฉา ตถา ปุเร. ยถา อโธ ตถา อุทฺธํ, ยถา อุทฺธํ ตถา อโธ. ยถา ทิวา ตถา รตฺติํ, ยถา รตฺติํ ตถา ทิวา. อิติ วิวเฏน เจตสา อปริโยนทฺเธน สปฺปภาสํ จิตฺตํ ภาเวติ. วีริยสมาธิ ฯเปฯ. จิตฺตสมาธิ. วีมํสา\-สมาธิ\-ปฺปธาน\-สงฺขาร\-สมนฺนาคตํ อิทฺธิปาทํ ภาเวติ, อิติ เม วีมํสา น จ อติลีนา ภวิสฺสติ, น จ อติปฺปคฺคหิตา ภวิสฺสติ, น จ อชฺฌตฺตํ สํขิตฺตา ภวิสฺสติ, น จ พหิทฺธา วิกฺขิตฺตา ภวิสฺสติ. ปจฺฉา\-ปุเร\-สญฺญี จ วิหรติ, ยถา ปุเร ตถา ปจฺฉา, ยถา ปจฺฉา ตถา ปุเร. ยถา อโธ ตถา อุทฺธํ, ยถา อุทฺธํ ตถา อโธ. ยถา ทิวา ตถา รตฺติํ, ยถา รตฺติํ ตถา ทิวา. อิติ วิวเฏน เจตสา อปริโยนทฺเธน สปฺปภาสํ จิตฺตํ ภาเวติ. อิเมสํ โข ภิกฺขเว จตุนฺนํ อิทฺธิปาทานํ ภาวิตตฺตา พหุลีกตตฺตา ตถาคโต เอวํมหิทฺธิโก เอวํ\-มหา\-นุภาโว.}

\csromanmatikaanchor{mtk.186}
\csromanmatikaanchor{mtk.187}
\csromantocmarknopage{subparagraph}{4. คงฺคาเปยฺยาลวคฺค}{mtk.186}
\bookmark[dest={mtk.186},level=5]{4. คงฺคาเปยฺยาลวคฺค}
\csromantocmarkref{subparagraph}{1-12. คงฺคานทีอาทิสุตฺต}{mtk.187}
\bookmark[dest={mtk.187},level=5]{1-12. คงฺคานทีอาทิสุตฺต}
\prosewithclosermidruled{อิเมสญฺจ ปน ภิกฺขเว จตุนฺนํ อิทฺธิปาทานํ ภาวิตตฺตา พหุลีกตตฺตา ตถาคโต อเนกวิหิตํ อิทฺธิวิธํ ปจฺจนุโภติ, เอโกปิ หุตฺวา พหุธา โหติ ฯเปฯ ยาว พฺรหฺมโลกาปิ กาเยน วสํ วตฺเตติ. อิเมสญฺจ ปน ภิกฺขเว จตุนฺนํ อิทฺธิปาทานํ ภาวิตตฺตา พหุลีกตตฺตา ตถาคโต \csromanfolio{254}อาสวานํ ขยา อนาสวํ เจโตวิมุตฺติํ ปญฺญาวิมุตฺติํ ทิฏฺเฐว ธมฺเม สยํ อภิญฺญา สจฺฉิกตฺวา อุปสมฺปชฺช วิหรตีติ.  ทฺวาทสมํ.}{(ฉปิ อภิญฺญาโย วิตฺถาเรตพฺพา.) อโยคุฬวคฺโค ตติโย.}
\csromancenter{\textbf{ตสฺสุทฺทานํ}}
\setlength{\csromangathapagewidth}{0pt}
\setlength{\csromangathawakwidth}{0pt}
\csromangathameasuregroup{มคฺโค อโยคุโฬ ภิกฺขุ, \\ เทฺว จานนฺทา ทุเว ภิกฺขู,}{\csromangathabat{มคฺโค อโยคุโฬ ภิกฺขุ,}{สุทฺธิกญฺจาปิ เทฺว ผลา.} \\ \csromangathabat{เทฺว จานนฺทา ทุเว ภิกฺขู,}{โมคฺคลฺลาโน ตถาคโตติ.}}
\csromangathasetleft{มคฺโค อโยคุโฬ ภิกฺขุ, \\ เทฺว จานนฺทา ทุเว ภิกฺขู,}
\csromangathagroup{\csromangathastanza{\csromangathabat{มคฺโค อโยคุโฬ ภิกฺขุ,}{สุทฺธิกญฺจาปิ เทฺว ผลา.} \\ \csromangathabat{เทฺว จานนฺทา ทุเว ภิกฺขู,}{โมคฺคลฺลาโน ตถาคโตติ.}}}

\chapterhead{4. คงฺคา\-เปยฺยาล\-วคฺค}
\prose{1-12. คงฺคานที\-อาทิ\-สุตฺต}

\proseitem{845-856}{เสยฺยถาปิ ภิกฺขเว คงฺคา นที ปาจีนนินฺนา ปาจีนโปณา ปาจีนปพฺภารา, เอวเมว โข ภิกฺขเว ภิกฺขุ จตฺตาโร อิทฺธิปาเท ภาเวนฺโต จตฺตาโร อิทฺธิปาเท พหุลีกโรนฺโต นิพฺพานนินฺโน โหติ นิพฺพานโปโณ นิพฺพาน\-ปพฺภาโร. กถญฺจ ภิกฺขเว ภิกฺขุ จตฺตาโร อิทฺธิปาเท ภาเวนฺโต จตฺตาโร อิทฺธิปาเท พหุลีกโรนฺโต นิพฺพานนินฺโน โหติ นิพฺพานโปโณ นิพฺพาน\-ปพฺภาโร, อิธ ภิกฺขเว ภิกฺขุ ฉนฺท\-สมาธิ\-ปฺปธาน\-สงฺขาร\-สมนฺนาคตํ อิทฺธิปาทํ ภาเวติ. วีริยสมาธิ ฯเปฯ. จิตฺตสมาธิ. วีมํสา\-สมาธิ\-ปฺปธาน\-สงฺขาร\-สมนฺนาคตํ อิทฺธิปาทํ ภาเวติ.}

\prosewithclosermidruled{เอวํ โข ภิกฺขเว ภิกฺขุ จตฺตาโร อิทฺธิปาเท ภาเวนฺโต จตฺตาโร อิทฺธิปาเท พหุลีกโรนฺโต นิพฺพานนินฺโน โหติ นิพฺพานโปโณ นิพฺพาน\-ปพฺภาโรติ.  ทฺวาทสมํ.}{คงฺคา\-เปยฺยาล\-วคฺโค จตุตฺโถ.}
\csromancenter{\textbf{ตสฺสุทฺทานํ}}
\setlength{\csromangathapagewidth}{0pt}
\setlength{\csromangathawakwidth}{0pt}
\csromangathameasuregroup{ฉ ปาจีนโต นินฺนา, \\ เทฺวเต ฉ ทฺวาทส โหนฺติ,}{\csromangathabat{ฉ ปาจีนโต นินฺนา,}{ฉ นินฺนา จ สมุทฺทโต.} \\ \csromangathabat{เทฺวเต ฉ ทฺวาทส โหนฺติ,}{วคฺโค เตน ปวุจฺจตีติ.}}
\csromangathasetleft{ฉ ปาจีนโต นินฺนา, \\ เทฺวเต ฉ ทฺวาทส โหนฺติ,}
\csromangathagroup{\csromangathastanza{\csromangathabat{ฉ ปาจีนโต นินฺนา,}{ฉ นินฺนา จ สมุทฺทโต.} \\ \csromangathabat{เทฺวเต ฉ ทฺวาทส โหนฺติ,}{วคฺโค เตน ปวุจฺจตีติ.}}}

\vspace{\tipitakaparskip}
\csromancenter{อปฺปมาทวคฺโค วิตฺถาเรตพฺโพ.}
\csromancenter{\textbf{ตสฺสุทฺทานํ}}
\setlength{\csromangathapagewidth}{0pt}
\setlength{\csromangathawakwidth}{0pt}
\csromangathameasuregroup{ตถาคตํ ปทํ กูฏํ,}{\csromangathabat{ตถาคตํ ปทํ กูฏํ,}{มูลํ สาโร จ วสฺสิกํ.} \\ ราชา จนฺทิมสูริยา, วตฺเถน ทสมํ ปทนฺติ. \\ พลกรณียวคฺโค วิตฺถาเรตพฺโพ.}
\csromangathasetleft{ตถาคตํ ปทํ กูฏํ,}
\csromangathagroup{\csromangathastanza{\csromanfolio{255}\csromangathabat{ตถาคตํ ปทํ กูฏํ,}{มูลํ สาโร จ วสฺสิกํ.} \\ ราชา จนฺทิมสูริยา, วตฺเถน ทสมํ ปทนฺติ. \\ พลกรณีย\-วคฺโค วิตฺถาเรตพฺโพ.}}

\vspace{\tipitakaparskip}
\csromancenter{\textbf{ตสฺสุทฺทานํ}}
\setlength{\csromangathapagewidth}{0pt}
\setlength{\csromangathawakwidth}{0pt}
\csromangathameasuregroup{พลํ พีชญฺจ นาโค จ,}{\csromangathabat{พลํ พีชญฺจ นาโค จ,}{รุกฺโข กุมฺเภน สูกิยา.} \\ อากาเสน จ เทฺว เมฆา, นาวา อาคนฺตุกา นทีติ. \\ เอสนาวคฺโค วิตฺถาเรตพฺโพ.}
\csromangathasetleft{พลํ พีชญฺจ นาโค จ,}
\csromangathagroup{\csromangathastanza{\csromangathabat{พลํ พีชญฺจ นาโค จ,}{รุกฺโข กุมฺเภน สูกิยา.} \\ อากาเสน จ เทฺว เมฆา, นาวา อาคนฺตุกา นทีติ. \\ เอสนาวคฺโค วิตฺถาเรตพฺโพ.}}

\vspace{\tipitakaparskip}
\csromancenter{\textbf{ตสฺสุทฺทานํ}}
\setlength{\csromangathapagewidth}{0pt}
\setlength{\csromangathawakwidth}{0pt}
\csromangathameasuregroup{เอสนา วิธา อาสโว, \\ ขิลํ มลญฺจ นีโฆ จ,}{\csromangathabat{เอสนา วิธา อาสโว,}{ภโว จ ทุกฺขตา ติสฺโส.} \\ \csromangathabat{ขิลํ มลญฺจ นีโฆ จ,}{เวทนา ตณฺหา ตสินา จาติ.}}
\csromangathasetleft{เอสนา วิธา อาสโว, \\ ขิลํ มลญฺจ นีโฆ จ,}
\csromangathagroup{\csromangathastanza{\csromangathabat{เอสนา วิธา อาสโว,}{ภโว จ ทุกฺขตา ติสฺโส.} \\ \csromangathabat{ขิลํ มลญฺจ นีโฆ จ,}{เวทนา ตณฺหา ตสินา จาติ.}}}

\chapterhead{8. \textbf{โอฆวคฺค}}
\csromanmatikaanchor{mtk.188}
\csromanmatikaanchor{mtk.189}
\csromantocmarknopage{subparagraph}{8. โอฆวคฺค}{mtk.188}
\bookmark[dest={mtk.188},level=5]{8. โอฆวคฺค}
\csromantocmarkref{subparagraph}{1-10. โอฆาทิสุตฺต}{mtk.189}
\bookmark[dest={mtk.189},level=5]{1-10. โอฆาทิสุตฺต}
\csromanheader{6}{1-10. โอฆาทิสุตฺต}
\proseitem{889-898}{ปญฺจิมานิ ภิกฺขเว อุทฺธมฺ\-ภาคิยานิ สํโยชนานิ. กตมานิ ปญฺจ, รูปราโค อรูปราโค มาโน อุทฺธจฺจํ อวิชฺชา, อิมานิ โข ภิกฺขเว ปญฺจุ\-ทฺธมฺภาคิยานิ สํโยชนานิ. อิเมสํ โข ภิกฺขเว ปญฺจนฺนํ อุทฺธมฺ\-ภาคิยานํ สํโยชนานํ อภิญฺญาย ปริญฺญาย ปริกฺขยาย ปหานาย จตฺตาโร อิทฺธิปาทา ภาเวตพฺพา. กตเม จตฺตาโร, อิธ ภิกฺขเว ภิกฺขุ ฉนฺท\-สมาธิ\-ปฺปธาน\-สงฺขาร\-สมนฺนาคตํ อิทฺธิปาทํ ภาเวติ. วีริยสมาธิ ฯเปฯ. \csromanfolio{256}จิตฺตสมาธิ. วีมํสา\-สมาธิ\-ปฺปธาน\-สงฺขาร\-สมนฺนาคตํ อิทฺธิปาทํ ภาเวติ. อิเมสํ โข ภิกฺขเว ปญฺจนฺนํ อุทฺธมฺ\-ภาคิยานํ สํโยชนานํ อภิญฺญาย ปริญฺญาย ปริกฺขยาย ปหานาย อิเม จตฺตาโร อิทฺธิปาทา ภาเวตพฺพาติ. (ยถา มคฺคสํยุตฺตํ, ตถา วิตฺถาเรตพฺพํ.) โอฆวคฺโค อฏฺฐโม.}

\vspace{\tipitakaparskip}
\csromancenter{\textbf{ตสฺสุทฺทานํ}}
\setlength{\csromangathapagewidth}{0pt}
\setlength{\csromangathawakwidth}{0pt}
\csromangathameasuregroup{โอโฆ โยโค อุปาทานํ, \\ กามคุณา นีวรณา,}{\csromangathabat{โอโฆ โยโค อุปาทานํ,}{คนฺถา อนุสเยน จ.} \\ \csromangathabat{กามคุณา นีวรณา,}{ขนฺธา โอรุทฺธมฺภาคิยาติ.} \\ \textbf{อิทฺธิปาทสํยุตฺตํ สตฺตมํ.}}
\csromangathameasurewak{\textbf{อิทฺธิปาทสํยุตฺตํ สตฺตมํ.}}
\csromangathasetleft{โอโฆ โยโค อุปาทานํ, \\ กามคุณา นีวรณา,}
\csromangathagroup{\csromangathastanza{\csromangathabat{โอโฆ โยโค อุปาทานํ,}{คนฺถา อนุสเยน จ.} \\ \csromangathabat{กามคุณา นีวรณา,}{ขนฺธา โอรุทฺธมฺภาคิยาติ.}}\csromangathastanza{\csromangathawak{\textbf{อิทฺธิปาทสํยุตฺตํ สตฺตมํ.}}}}

\csromanmatikaanchor{mtk.190}
\csromanmatikaanchor{mtk.191}
\csromantocmarknopage{subparagraph}{8. อนุรุทฺธสํยุตฺต}{mtk.190}
\bookmark[dest={mtk.190},level=5]{8. อนุรุทฺธสํยุตฺต}
\csromantocmarknopage{subparagraph}{1. รโหคตวคฺค}{mtk.191}
\bookmark[dest={mtk.191},level=5]{1. รโหคตวคฺค}
\csromanheader{6}{1. \textbf{รโหคตวคฺค}}
\csromanmatikaanchor{mtk.192}
\csromantocmarkref{subparagraph}{1. ปฐมรโหคตสุตฺต}{mtk.192}
\bookmark[dest={mtk.192},level=5]{1. ปฐมรโหคตสุตฺต}
\csromanheader{6}{1. ปฐม\-รโหคต\-สุตฺต}
\proseitem{899}{\csromanfolio{257}เอวํ เม สุตํ\csromandash{}เอกํ สมยํ อายสฺมา อนุรุทฺโธ สาวตฺถิยํ วิหรติ เชตวเน อนาถ\-ปิณฺฑิกสฺส อาราเม. อถ โข อายสฺมโต อนุรุทฺธสฺส รโหคตสฺส ปฏิสลฺลีนสฺส เอวํ เจตโส ปริวิตกฺโก อุทปาทิ “เยสํ เกสญฺจิ จตฺถาโร สติปฏฺฐานา วิรทฺธา, วิรทฺโธ เตสํ อริโย มคฺโค สมฺมา ทุกฺข\-กฺขย\-คามี. เยสํ เกสญฺจิ จตฺตาโร สติปฏฺฐานา อารทฺธา, อารทฺโธ เตสํ อริโย มคฺโค สมฺมา ทุกฺข\-กฺขย\-คามี”ติ.}

\prose{อถ โข อายสฺมา มหาโมคฺคลฺลาโน อายสฺมโต อนุรุทฺธสฺส เจตสา เจโต\-ปริวิตกฺกม\-ญฺญาย เสยฺยถาปิ นาม พลวา ปุริโส สมิญฺชิตํ วา พาหํ ปสาเรยฺย, ปสาริตํ วา พาหํ สมิญฺเชยฺย, เอวเมว อายสฺมโต อนุรุทฺธสฺส สมฺมุเข ปาตุรโหสิ. อถ โข อายสฺมา มหาโมคฺคลฺลาโน อายสฺมนฺตํ อนุรุทฺธํ เอตทโวจ “กิตฺตาวตา นุ โข อาวุโส อนุรุทฺธ ภิกฺขุโน จตฺตาโร สติปฏฺฐานา อารทฺธา โหนฺตี”ติ.}

\prose{อิธาวุโส ภิกฺขุ อชฺฌตฺตํ กาเย สมุทยธมฺมา\-นุปสฺสี วิหรติ. อชฺฌตฺตํ กาเย วยธมฺมา\-นุปสฺสี วิหรติ. อชฺฌตฺตํ กาเย สมุทย\-วย\-ธมฺมานุปสฺสี วิหรติ อาตาปี สมฺปชาโน สติมา วิเนยฺย โลเก อภิชฺฌา\-โทมนสฺสํ. พหิทฺธา กาเย สมุทยธมฺมา\-นุปสฺสี วิหรติ. พหิทฺธา กาเย วยธมฺมา\-นุปสฺสี วิหรติ. พหิทฺธา กาเย สมุทย\-วย\-ธมฺมานุปสฺสี วิหรติ อาตาปี สมฺปชาโน สติมา วิเนยฺย โลเก อภิชฺฌา\-โทมนสฺสํ. อชฺฌตฺต\-พหิทฺธา กาเย สมุทยธมฺมา\-นุปสฺสี วิหรติ. อชฺฌตฺต\-พหิทฺธา กาเย วยธมฺมา\-นุปสฺสี วิหรติ. อชฺฌตฺต\-พหิทฺธา กาเย สมุทย\-วย\-ธมฺมานุปสฺสี วิหรติ อาตาปี สมฺปชาโน สติมา วิเนยฺย โลเก อภิชฺฌา\-โทมนสฺสํ.}

\prose{\csromanfolio{258}โส สเจ อากงฺขติ “อปฺปฏิกูเล ปฏิกูลสญฺญี วิหเรยฺยนฺ”ติ, ปฏิกูลสญฺญี ตตฺถ วิหรติ. สเจ อากงฺขติ “ปฏิกูเล อปฺปฏิกูลสญฺญี วิหเรยฺยนฺ”ติ, อปฺปฏิกูลสญฺญี ตตฺถ วิหรติ. สเจ อากงฺขติ “อปฺปฏิกูเล จ ปฏิกูเล จ ปฏิกูลสญฺญี วิหเรยฺยนฺ”ติ, ปฏิกูลสญฺญี ตตฺถ วิหรติ. สเจ อากงฺขติ “ปฏิกูเล จ อปฺปฏิกูเล จ อปฺปฏิกูลสญฺญี วิหเรยฺยนฺ”ติ, อปฺปฏิกูลสญฺญี ตตฺถ วิหรติ. สเจ อากงฺขติ “อปฺปฏิกูลญฺจ ปฏิกูลญฺจ ตทุภยํ อภินิวชฺเชตฺวา อุเปกฺขโก วิหเรยฺยํ สโต สมฺปชาโน”ติ, อุเปกฺขโก ตตฺถ วิหรติ สโต สมฺปชาโน.}

\prose{อชฺฌตฺตํ เวทนาสุ สมุทยธมฺมา\-นุปสฺสี วิหรติ. อชฺฌตฺตํ เวทนาสุ วยธมฺมา\-นุปสฺสี วิหรติ. อชฺฌตฺตํ เวทนาสุ สมุทย\-วย\-ธมฺมานุปสฺสี วิหรติ อาตาปี สมฺปชาโน สติมา วิเนยฺย โลเก อภิชฺฌา\-โทมนสฺสํ. พหิทฺธา เวทนาสุ สมุทยธมฺมา\-นุปสฺสี วิหรติ. พหิทฺธา เวทนาสุ วยธมฺมา\-นุปสฺสี วิหรติ. พหิทฺธา เวทนาสุ สมุทย\-วย\-ธมฺมานุปสฺสี วิหรติ อาตาปี สมฺปชาโน สติมา วิเนยฺย โลเก อภิชฺฌา\-โทมนสฺสํ. อชฺฌตฺต\-พหิทฺธา เวทนาสุ สมุทยธมฺมา\-นุปสฺสี วิหรติ. อชฺฌตฺต\-พหิทฺธา เวทนาสุ วยธมฺมา\-นุปสฺสี วิหรติ. อชฺฌตฺต\-พหิทฺธา เวทนาสุ สมุทย\-วย\-ธมฺมานุปสฺสี วิหรติ อาตาปี สมฺปชาโน สติมา วิเนยฺย โลเก อภิชฺฌา\-โทมนสฺสํ.}

\prose{โส สเจ อากงฺขติ “อปฺปฏิกูเล ปฏิกูลสญฺญี วิหเรยฺยนฺ”ติ, ปฏิกูลสญฺญี ตตฺถ วิหรติ. สเจ อากงฺขติ “ปฏิกูเล อปฺปฏิกูลสญฺญี วิหเรยฺยนฺ”ติ, อปฺปฏิกูลสญฺญี ตตฺถ วิหรติ. สเจ อากงฺขติ “อปฺปฏิกูเล จ ปฏิกูเล จ ปฏิกูลสญฺญี วิหเรยฺยนฺ”ติ, ปฏิกูลสญฺญี ตตฺถ วิหรติ. สเจ อากงฺขติ “ปฏิกูเล จ อปฺปฏิกูเล จ อปฺปฏิกูลสญฺญี วิหเรยฺยนฺ”ติ, อปฺปฏิกูลสญฺญี ตตฺถ วิหรติ. สเจ อากงฺขติ “อปฺปฏิกูลญฺจ ปฏิกูลญฺจ ตทุภยํ อภินิวชฺเชตฺวา อุเปกฺขโก วิหเรยฺยํ สโต สมฺปชาโน”ติ, อุเปกฺขโก ตตฺถ วิหรติ สโต สมฺปชาโน.}

\prose{อชฺฌตฺตํ จิตฺเต ฯเปฯ. พหิทฺธา จิตฺเต ฯเปฯ. อชฺฌตฺต\-พหิทฺธา จิตฺเต สมุทยธมฺมา\-นุปสฺสี วิหรติ. อชฺฌตฺต\-พหิทฺธา จิตฺเต วยธมฺมา\-นุปสฺสี วิหรติ. อชฺฌตฺต\-พหิทฺธา จิตฺเต สมุทย\-วย\-ธมฺมานุปสฺสี วิหรติ อาตาปี ฯเปฯ อภิชฺฌา\-โทมนสฺสํ. โส สเจ อากงฺขติ “อปฺปฏิกูเล ปฏิกูลสญฺญี}

\prose{\csromanfolio{259}วิหเรยฺยนฺ”ติ, ปฏิกูลสญฺญี ตตฺถ วิหรติ ฯเปฯ อุเปกฺขโก ตตฺถ วิหรติ สโต สมฺปชาโน.}

\prose{อชฺฌตฺตํ ธมฺเมสุ ฯเปฯ. พหิทฺธา ธมฺเมสุ ฯเปฯ. อชฺฌตฺต\-พหิทฺธา ธมฺเมสุ สมุทยธมฺมา\-นุปสฺสี วิหรติ. อชฺฌตฺต\-พหิทฺธา ธมฺเมสุ วยธมฺมา\-นุปสฺสี วิหรติ. อชฺฌตฺต\-พหิทฺธา ธมฺเมสุ สมุทย\-วย\-ธมฺมานุปสฺสี วิหรติ อาตาปี สมฺปชาโน สติมา วิเนยฺย โลเก อภิชฺฌา\-โทมนสฺสํ. โส สเจ อากงฺขติ “อปฺปฏิกูเล ปฏิกูลสญฺญี วิหเรยฺยนฺ”ติ, ปฏิกูลสญฺญี ตตฺถ วิหรติ ฯเปฯ อุเปกฺขโก ตตฺถ วิหรติ สโต สมฺปชาโน. เอตฺตาวตา โข อาวุโส ภิกฺขุโน จตฺตาโร สติปฏฺฐานา อารทฺธา โหนฺตีติ.  ปฐมํ.}

\csromanmatikaanchor{mtk.193}
\csromantocmarkref{subparagraph}{2. ทุติยรโหคตสุตฺต}{mtk.193}
\bookmark[dest={mtk.193},level=5]{2. ทุติยรโหคตสุตฺต}
\csromanheader{6}{2. ทุติย\-รโหคต\-สุตฺต}
\proseitem{900}{สาวตฺถิ\-นิทานํ. อถ โข อายสฺมโต อนุรุทฺธสฺส รโหคตสฺส ปฏิสลฺลีนสฺส เอวํ เจตโส ปริวิตกฺโก อุทปาทิ “เยสํ เกสญฺจิ จตฺตาโร สติปฏฺฐานา วิรทฺธา, วิรทฺโธ เตสํ อริโย มคฺโค สมฺมา ทุกฺข\-กฺขย\-คามี. เยสํ เกสญฺจิ จตฺตาโร สติปฏฺฐานา อารทฺธา, อารทฺโธ เตสํ อริโย มคฺโค สมฺมา ทุกฺข\-กฺขย\-คามี”ติ.}

\prose{อถ โข อายสฺมา มหาโมคฺคลฺลาโน อายสฺมโต อนุรุทฺธสฺส เจตสา เจโต\-ปริวิตกฺกม\-ญฺญาย เสยฺยถาปิ นาม พลวา ปุริโส สมิญฺชิตํ วา พาหํ ปสาเรยฺย, ปสาริตํ วา พาหํ สมิญฺเชยฺย, เอวเมว อายสฺมโต อนุรุทฺธสฺส สมฺมุเข ปาตุรโหสิ.}

\prose{อถ โข อายสฺมา มหาโมคฺคลฺลาโน อายสฺมนฺตํ อนุรุทฺธํ เอตทโวจ “กิตฺตาวตา นุ โข อาวุโส อนุรุทฺธ ภิกฺขุโน จตฺตาโร สติปฏฺฐานา อารทฺธา โหนฺตี”ติ. อิธาวุโส ภิกฺขุ อชฺฌตฺตํ กาเย กายานุปสฺสี วิหรติ อาตาปี สมฺปชาโน สติมา วิเนยฺย โลเก อภิชฺฌา\-โทมนสฺสํ. พหิทฺธา กาเย กายานุปสฺสี วิหรติ ฯเปฯ. อชฺฌตฺต\-พหิทฺธา กาเย กายานุปสฺสี วิหรติ อาตาปี สมฺปชาโน สติมา วิเนยฺย โลเก อภิชฺฌา\-โทมนสฺสํ.}

\csromanmatikaanchor{mtk.194}
\csromanmatikaanchor{mtk.195}
\csromantocmarkref{subparagraph}{3. สุตนุสุตฺต}{mtk.194}
\bookmark[dest={mtk.194},level=5]{3. สุตนุสุตฺต}
\csromantocmarkref{subparagraph}{4-8. ปฐมกณฺฑกีสุตฺตาทีนิ}{mtk.195}
\bookmark[dest={mtk.195},level=5]{4-8. ปฐมกณฺฑกีสุตฺตาทีนิ}
\prose{อชฺฌตฺตํ เวทนาสุ เวทนานุปสฺสี วิหรติ อาตาปี สมฺปชาโน สติมา วิเนยฺย โลเก อภิชฺฌา\-โทมนสฺสํ. พหิทฺธา เวทนาสุ \csromanfolio{260}เวทนานุปสฺสี วิหรติ อาตาปี สมฺปชาโน สติมา วิเนยฺย โลเก อภิชฺฌา\-โทมนสฺสํ. อชฺฌตฺต\-พหิทฺธา เวทนาสุ เวทนานุปสฺสี วิหรติ อาตาปี สมฺปชาโน สติมา วิเนยฺย โลเก อภิชฺฌา\-โทมนสฺสํ.}

\prose{อชฺฌตฺตํ จิตฺเต ฯเปฯ. พหิทฺธา จิตฺเต ฯเปฯ. อชฺฌตฺต\-พหิทฺธา จิตฺเต จิตฺตานุปสฺสี วิหรติ อาตาปี สมฺปชาโน สติมา วิเนยฺย โลเก อภิชฺฌา\-โทมนสฺสํ.}

\prose{อชฺฌตฺตํ ธมฺเมสุ ฯเปฯ. พหิทฺธา ธมฺเมสุ ฯเปฯ. อชฺฌตฺต\-พหิทฺธา ธมฺเมสุ ธมฺมานุปสฺสี วิหรติ อาตาปี สมฺปชาโน สติมา วิเนยฺย โลเก อภิชฺฌา\-โทมนสฺสํ. เอตฺตาวตา โข อาวุโส ภิกฺขุโน จตฺตาโร สติปฏฺฐานา อารทฺธา โหนฺตีติ.  ทุติยํ.}

\csromanheader{6}{3. \textbf{สุตนุสุตฺต}}
\proseitem{901}{เอกํ สมยํ อายสฺมา อนุรุทฺโธ สาวตฺถิยํ วิหรติ สุตนุตีเร. อถ โข สมฺพหุลา ภิกฺขู เยนายสฺมา อนุรุทฺโธ เตนุ\-ปสงฺกมิํสุ, อุปสงฺกมิตฺวา อายสฺมตา อนุรุทฺเธน สทฺธิํ สมฺโมทิํสุ, สมฺโมทนียํ กถํ สารณียํ วีติสาเรตฺวา เอกมนฺตํ นิสีทิํสุ, เอกมนฺตํ นิสินฺนา โข เต ภิกฺขู อายสฺมนฺตํ อนุรุทฺธํ เอตทโวจุํ “กตเมสํ อายสฺมา อนุรุทฺโธ ธมฺมานํ ภาวิตตฺตา พหุลีกตตฺตา มหาภิญฺญตํ ปตฺโต”ติ.}

\prose{จตุนฺนํ ขฺวาหํ อาวุโส สติปฏฺฐานานํ ภาวิตตฺตา พหุลีกตตฺตา มหาภิญฺญตํ ปตฺโต. กตเมสํ จตุนฺนํ, อิธาหํ อาวุโส กาเย กายานุปสฺสี วิหรามิ อาตาปี สมฺปชาโน สติมา วิเนยฺย โลเก อภิชฺฌา\-โทมนสฺสํ. เวทนาสุ ฯเปฯ. จิตฺเต ฯเปฯ. ธมฺเมสุ ธมฺมานุปสฺสี วิหรามิ อาตาปี สมฺปชาโน สติมา วิเนยฺย โลเก อภิชฺฌา\-โทมนสฺสํ. อิเมสํ ขฺวาหํ อาวุโส จตุนฺนํ สติปฏฺฐานานํ ภาวิตตฺตา พหุลีกตตฺตา มหาภิญฺญตํ ปตฺโต. อิเมสญฺจ ปนาหํ อาวุโส จตุนฺนํ สติปฏฺฐานานํ ภาวิตตฺตา พหุลีกตตฺตา หีนํ ธมฺมํ หีนโต อพฺภญฺญาสิํ, มชฺฌิมํ ธมฺมํ มชฺฌิมโต อพฺภญฺญาสิํ, ปณีตํ ธมฺมํ ปณีตโต อพฺภญฺญาสินฺติ.  ตติยํ.}

\csromanheader{2}{4. ปฐม\-กณฺฑกี\-สุตฺต}
\proseitem{902}{เอกํ สมยํ อายสฺมา จ อนุรุทฺโธ อายสฺมา จ สาริปุตฺโต อายสฺมา จ มหาโมคฺคลฺลาโน สาเกเต วิหรนฺติ กณฺฑกีวเน.}

\prose{\csromanfolio{261}อถ โข อายสฺมา จ สาริปุตฺโต อายสฺมา จ มหาโมคฺคลฺลาโน สายนฺหสมยํ ปฏิสลฺลานา วุฏฺฐิตา เยนายสฺมา อนุรุทฺโธ เตนุ\-ปสงฺกมิํสุ, อุปสงฺกมิตฺวา อายสฺมตา อนุรุทฺเธน สทฺธิํ สมฺโมทิํสุ, สมฺโมทนียํ กถํ สารณียํ วีติสาเรตฺวา เอกมนฺตํ นิสีทิํสุ, เอกมนฺตํ นิสินฺโน โข อายสฺมา สาริปุตฺโต อายสฺมนฺตํ อนุรุทฺธํ เอตทโวจ “เสเขนาวุโส อนุรุทฺธ ภิกฺขุนา กตเม ธมฺมา อุปสมฺปชฺช วิหาตพฺพา”ติ.}

\prose{เสเขนาวุโส สาริปุตฺต ภิกฺขุนา จตฺตาโร สติปฏฺฐานา อุปสมฺปชฺช วิหาตพฺพา. กตเม จตฺตาโร, อิธาวุโส ภิกฺขุ กาเย กายานุปสฺสี วิหรติ อาตาปี สมฺปชาโน สติมา วิเนยฺย โลเก อภิชฺฌา\-โทมนสฺสํ. เวทนาสุ ฯเปฯ. จิตฺเต ฯเปฯ. ธมฺเมสุ ธมฺมานุปสฺสี วิหรติ อาตาปี สมฺปชาโน สติมา วิเนยฺย โลเก อภิชฺฌา\-โทมนสฺสํ. เสเขนาวุโส สาริปุตฺต ภิกฺขุนา อิเม จตฺตาโร สติปฏฺฐานา อุปสมฺปชฺช วิหาตพฺพาติ.  จตุตฺถํ.}

\csromanheader{2}{5. ทุติย\-กณฺฑกี\-สุตฺต}
\proseitem{903}{สาเกตนิทานํ. เอกมนฺตํ นิสินฺโน โข อายสฺมา สาริปุตฺโต อายสฺมนฺตํ อนุรุทฺธํ เอตทโวจ “อเสเขนาวุโส อนุรุทฺธ ภิกฺขุนา กตเม ธมฺมา อุปสมฺปชฺช วิหาตพฺพา”ติ. อเสเขนาวุโส สาริปุตฺต ภิกฺขุนา จตฺตาโร สติปฏฺฐานา อุปสมฺปชฺช วิหาตพฺพา. กตเม จตฺตาโร, อิธาวุโส ภิกฺขุ กาเย กายานุปสฺสี วิหรติ อาตาปี สมฺปชาโน สติมา วิเนยฺย โลเก อภิชฺฌา\-โทมนสฺสํ. เวทนาสุ ฯเปฯ. จิตฺเต ฯเปฯ. ธมฺเมสุ ธมฺมานุปสฺสี วิหรติ อาตาปี สมฺปชาโน สติมา วิเนยฺย โลเก อภิชฺฌา\-โทมนสฺสํ. อเสเขนาวุโส สาริปุตฺต ภิกฺขุนา อิเม จตฺตาโร สติปฏฺฐานา อุปสมฺปชฺช วิหาตพฺพาติ.  ปญฺจมํ.}

\csromanheader{2}{6. ตติย\-กณฺฑกี\-สุตฺต}
\proseitem{904}{สาเกตนิทานํ. เอกมนฺตํ นิสินฺโน โข อายสฺมา สาริปุตฺโต อายสฺมนฺตํ อนุรุทฺธํ เอตทโวจ “กตเมสํ อายสฺมา อนุรุทฺโธ ธมฺมานํ ภาวิตตฺตา พหุลีกตตฺตา มหาภิญฺญตํ ปตฺโต”ติ. จตุนฺนํ ขฺวาหํ อาวุโส สติปฏฺฐานานํ ภาวิตตฺตา พหุลีกตตฺตา มหาภิญฺญตํ \csromanfolio{262}ปตฺโต. กตเมสํ จตุนฺนํ, อิธาหํ อาวุโส กาเย กายานุปสฺสี วิหรามิ อาตาปี สมฺปชาโน สติมา วิเนยฺย โลเก อภิชฺฌา\-โทมนสฺสํ. เวทนาสุ ฯเปฯ. จิตฺเต ฯเปฯ. ธมฺเมสุ ธมฺมานุปสฺสี วิหรามิ อาตาปี สมฺปชาโน สติมา วิเนยฺย โลเก อภิชฺฌา\-โทมนสฺสํ. อิเมสํ ขฺวาหํ อาวุโส จตุนฺนํ สติปฏฺฐานานํ ภาวิตตฺตา พหุลีกตตฺตา มหาภิญฺญตํ ปตฺโต. อิเมสญฺจ ปนาหํ อาวุโส จตุนฺนํ สติปฏฺฐานานํ ภาวิตตฺตา พหุลีกตตฺตา สหสฺสํ โลกํ อภิชานามีติ.  ฉฏฺฐํ.}

\csromanheader{2}{7. ตณฺหากฺขย\-สุตฺต}
\proseitem{905}{สาวตฺถิ\-นิทานํ. ตตฺร โข อายสฺมา อนุรุทฺโธ ภิกฺขู อามนฺเตสิ “อาวุโส ภิกฺขโว”ติ. “อาวุโส”ติ โข เต ภิกฺขู อายสฺมโต อนุรุทฺธสฺส ปจฺจสฺโสสุํ. อายสฺมา อนุรุทฺโธ เอตทโวจ\csromandash{}}

\prose{จตฺตาโรเม อาวุโส สติปฏฺฐานา ภาวิตา พหุลีกตา ตณฺหากฺขยาย สํวตฺตนฺติ. กตเม จตฺตาโร, อิธาวุโส ภิกฺขุ กาเย กายานุปสฺสี วิหรติ ฯเปฯ. เวทนาสุ ฯเปฯ. จิตฺเต ฯเปฯ. ธมฺเมสุ ธมฺมานุปสฺสี วิหรติ อาตาปี สมฺปชาโน สติมา วิเนยฺย โลเก อภิชฺฌา\-โทมนสฺสํ. อิเม โข อาวุโส จตฺตาโร สติปฏฺฐานา ภาวิตา พหุลีกตา ตณฺหากฺขยาย สํวตฺตนฺตีติ.  สตฺตมํ.}

\csromanheader{2}{8. สลฬาคาร\-สุตฺต}
\proseitem{906}{เอกํ สมยํ อายสฺมา อนุรุทฺโธ สาวตฺถิยํ วิหรติ สลฬาคาเร. ตตฺร โข อายสฺมา อนุรุทฺโธ ภิกฺขู อามนฺเตสิ ฯเปฯ เอตทโวจ\csromandash{}เสยฺยถาปิ อาวุโส คงฺคา นที ปาจีนนินฺนา ปาจีนโปณา ปาจีนปพฺภารา, อถ มหาชนกาโย อาคจฺเฉยฺย กุทาลปิฏกํ\csromansharedfootnote{6ba53f5f1e69157e}{กุทฺทาลปิฏกํ (พหูสุ)} อาทาย “มยํ อิมํ คงฺคานทิํ ปจฺฉานินฺนํ กริสฺสาม ปจฺฉาโปณํ ปจฺฉาปพฺภารนฺ”ติ. ตํ กิํ มญฺญถาวุโส, อปิ นุ โส มหาชนกาโย คงฺคานทิํ ปจฺฉานินฺนํ กเรยฺย ปจฺฉาโปณํ ปจฺฉา\-ปพฺภารนฺติ. โน เหตํ อาวุโส. ตํ กิสฺส เหตุ, คงฺคา อาวุโส นที ปาจีนนินฺนา ปาจีนโปณา ปาจีนปพฺภารา. สา น สุกรา ปจฺฉานินฺนํ กาตุํ ปจฺฉาโปณํ ปจฺฉา\-ปพฺภารํ, ยาวเทว จ ปน โส มหาชนกาโย กิลมถสฺส วิฆาตสฺส ภาคี อสฺสาติ.}

\prose{\csromanfolio{263}เอวเมว โข อาวุโส ภิกฺขุํ จตฺตาโร สติปฏฺฐาเน ภาเวนฺตํ จตฺตาโร สติปฏฺฐาเน พหุลีกโรนฺตํ ราชาโน วา ราชมหามตฺตา วา มิตฺตา วา อมจฺจา วา ญาตี วา สาโลหิตา วา โภเคหิ อภิหฏฺฐุํ ปวาเรยฺยุํ “เอหมฺโภ ปุริส, กิํ เต อิเม กาสาวา อนุทหนฺติ, กิํ มุณฺโฑ กปาลมนุสญฺจรสิ, เอหิ หีนายาวตฺติตฺวา โภเค จ ภุญฺชสฺสุ, ปุญฺญานิ จ กโรหี”ติ.}

\prose{โส วต อาวุโส ภิกฺขุ จตฺตาโร สติปฏฺฐาเน ภาเวนฺโต จตฺตาโร สติปฏฺฐาเน พหุลีกโรนฺโต สิกฺขํ ปจฺจกฺขาย หีนายาวตฺติสฺสตีติ เนตํ ฐานํ วิชฺชติ. ตํ กิสฺส เหตุ, ยํ หิ ตํ อาวุโส จิตฺตํ ทีฆรตฺตํ วิเวกนินฺนํ วิเวกโปณํ วิเวก\-ปพฺภารํ, ตํ วต หีนายาวตฺติสฺสตีติ เนตํ ฐานํ วิชฺชติ. กถญฺจาวุโส ภิกฺขุ จตฺตาโร สติปฏฺฐาเน ภาเวติ, จตฺตาโร สติปฏฺฐาเน พหุลีกโรตีติ. อิธาวุโส ภิกฺขุ กาเย กายานุปสฺสี วิหรติ ฯเปฯ. เวทนาสุ ฯเปฯ. จิตฺเต ฯเปฯ. ธมฺเมสุ ธมฺมานุปสฺสี วิหรติ อาตาปี สมฺปชาโน สติมา วิเนยฺย โลเก อภิชฺฌา\-โทมนสฺสํ. เอวํ โข อาวุโส ภิกฺขุ จตฺตาโร สติปฏฺฐาเน ภาเวติ. จตฺตาโร สติปฏฺฐาเน พหุลีกโรตีติ.  อฏฺฐมํ.}

\csromanmatikaanchor{mtk.196}
\csromantocmarkref{subparagraph}{9. อมฺพปาลิวนสุตฺต}{mtk.196}
\bookmark[dest={mtk.196},level=5]{9. อมฺพปาลิวนสุตฺต}
\csromanheader{6}{9. อมฺพปาลิวน\-สุตฺต}
\proseitem{907}{เอกํ สมยํ อายสฺมา จ อนุรุทฺโธ อายสฺมา จ สาริปุตฺโต เวสาลิยํ วิหรนฺติ อมฺพปาลิวเน. อถ โข อายสฺมา สาริปุตฺโต สายนฺหสมยํ ปฏิสลฺลานา วุฏฺฐิโต ฯเปฯ เอกมนฺตํ นิสินฺโน โข อายสฺมา สาริปุตฺโต อายสฺมนฺตํ อนุรุทฺธํ เอตทโวจ\csromandash{}}

\prose{วิปฺปสนฺนานิ โข เต อาวุโส อนุรุทฺธ อินฺทฺริยานิ, ปริสุทฺโธ มุขวณฺโณ ปริโยทาโต, กตเมนายสฺมา อนุรุทฺโธ วิหาเรน เอตรหิ พหุลํ วิหรตีติ. จตูสุ ขฺวาหํ อาวุโส สติปฏฺฐาเนสุ สุปฺปติฏฺฐิต\-จิตฺโต เอตรหิ พหุลํ วิหรามิ. กตเมสุ จตูสุ, อิธาหํ อาวุโส กาเย กายานุปสฺสี วิหรามิ อาตาปี สมฺปชาโน สติมา วิเนยฺย โลเก อภิชฺฌา\-โทมนสฺสํ. เวทนาสุ ฯเปฯ. จิตฺเต ฯเปฯ. ธมฺเมสุ ธมฺมานุปสฺสี วิหรามิ \csromanfolio{264}อาตาปี สมฺปชาโน สติมา วิเนยฺย โลเก อภิชฺฌา\-โทมนสฺสํ. อิเมสุ ขฺวาหํ อาวุโส จตูสุ สติปฏฺฐาเนสุ สุปฺปติฏฺฐิต\-จิตฺโต เอตรหิ พหุลํ วิหรามิ. โย โส อาวุโส ภิกฺขุ อรหํ ขีณาสโว วุสิตวา กตกรณีโย โอหิตภาโร อนุปฺปตฺต\-สทตฺโถ ปริกฺขีณ\-ภวสํโยชโน สมฺมทญฺญา\-วิมุตฺโต, โส อิเมสุ จตูสุ สติปฏฺฐาเนสุ สุปฺปติฏฺฐิต\-จิตฺโต พหุลํ วิหรตีติ.}

\prose{ลาภา วต โน อาวุโส, สุลทฺธํ วต โน อาวุโส, เย มยํ อายสฺมโต อนุรุทฺธสฺส สมฺมุขาว อสฺสุมฺห อาสภิํ วาจํ ภาสมานสฺสาติ.  นวมํ.}

\csromanmatikaanchor{mtk.197}
\csromantocmarkref{subparagraph}{10. พาฬฺหคิลานสุตฺต}{mtk.197}
\bookmark[dest={mtk.197},level=5]{10. พาฬฺหคิลานสุตฺต}
\csromanheader{6}{10. พาฬฺหคิลาน\-สุตฺต}
\proseitem{908}{เอกํ สมยํ อายสฺมา อนุรุทฺโธ สาวตฺถิยํ วิหรติ อนฺธวนสฺมิํ อาพาธิโก ทุกฺขิโต พาฬฺหคิลาโน. อถ โข สมฺพหุลา ภิกฺขู เยนายสฺมา อนุรุทฺโธ เตนุ\-ปสงฺกมิํสุ, อุปสงฺกมิตฺวา อายสฺมนฺตํ อนุรุทฺธํ เอตทโวจุํ\csromandash{}}

\prosewithclosermidruled{กตเมนา\-ยสฺมโต อนุรุทฺธสฺส วิหาเรน วิหรโต อุปฺปนฺนา สารีริกา ทุกฺขา เวทนา จิตฺตํ น ปริยาทาย ติฏฺฐนฺตีติ. จตูสุ โข เม อาวุโส สติปฏฺฐาเนสุ สุปฺปติฏฺฐิต\-จิตฺตสฺส วิหรโต อุปฺปนฺนา สารีริกา ทุกฺขา เวทนา จิตฺตํ น ปริยาทาย ติฏฺฐนฺติ. กตเมสุ จตูสุ, อิธาหํ อาวุโส กาเย กายานุปสฺสี วิหรามิ ฯเปฯ. เวทนาสุ ฯเปฯ. จิตฺเต ฯเปฯ. ธมฺเมสุ ธมฺมานุปสฺสี วิหรามิ อาตาปี สมฺปชาโน สติมา วิเนยฺย โลเก อภิชฺฌา\-โทมนสฺสํ. อิเมสุ โข เม อาวุโส จตูสุ สติปฏฺฐาเนสุ สุปฺปติฏฺฐิต\-จิตฺตสฺส วิหรโต อุปฺปนฺนา สารีริกา ทุกฺขา เวทนา จิตฺตํ น ปริยาทาย ติฏฺฐนฺตีติ.  ทสมํ.}{รโหคตวคฺโค ปฐโม.}
\csromancenter{\textbf{ตสฺสุทฺทานํ}}
\setlength{\csromangathapagewidth}{0pt}
\setlength{\csromangathawakwidth}{0pt}
\csromangathameasuregroup{รโหคเตน เทฺว วุตฺตา, \\ ตณฺหากฺขยสลฬาคารํ,}{\csromangathabat{รโหคเตน เทฺว วุตฺตา,}{สุตนุ กณฺฑกี ตโย.} \\ \csromangathabat{ตณฺหากฺขยสลฬาคารํ,}{อมฺพปาลิ จ คิลานนฺติ.}}
\csromangathasetleft{รโหคเตน เทฺว วุตฺตา, \\ ตณฺหากฺขยสลฬาคารํ,}
\csromangathagroup{\csromangathastanza{\csromangathabat{รโหคเตน เทฺว วุตฺตา,}{สุตนุ กณฺฑกี ตโย.} \\ \csromangathabat{ตณฺหากฺขยสลฬาคารํ,}{อมฺพปาลิ จ คิลานนฺติ.}}}

\csromanmatikaanchor{mtk.198}
\csromantocmarknopage{subparagraph}{2. ทุติยวคฺค}{mtk.198}
\bookmark[dest={mtk.198},level=5]{2. ทุติยวคฺค}
\csromanheader{6}{2. \textbf{ทุติยวคฺค}}
\csromanmatikaanchor{mtk.199}
\csromantocmarkref{subparagraph}{1. กปฺปสหสฺสสุตฺต}{mtk.199}
\bookmark[dest={mtk.199},level=5]{1. กปฺปสหสฺสสุตฺต}
\csromanheader{6}{1. กปฺปสหสฺส\-สุตฺต}
\csromanmatikaanchor{mtk.200}
\csromantocmarkref{subparagraph}{2-14. อิทฺธิวิธสุตฺตาทีนิ}{mtk.200}
\bookmark[dest={mtk.200},level=5]{2-14. อิทฺธิวิธสุตฺตาทีนิ}
\proseitem{909}{\csromanfolio{265}เอกํ สมยํ อายสฺมา อนุรุทฺโธ สาวตฺถิยํ วิหรติ เชตวเน อนาถ\-ปิณฺฑิกสฺส อาราเม. อถ โข สมฺพหุลา ภิกฺขู เยนายสฺมา อนุรุทฺโธ เตนุ\-ปสงฺกมิํสุ, อุปสงฺกมิตฺวา อายสฺมตา อนุรุทฺเธน สทฺธิํ ฯเปฯ เอกมนฺตํ นิสินฺนา โข เต ภิกฺขู อายสฺมนฺตํ อนุรุทฺธํ เอตทโวจุํ\csromandash{}}

\prose{กตเมสํ อายสฺมา อนุรุทฺโธ ธมฺมานํ ภาวิตตฺตา พหุลีกตตฺตา มหาภิญฺญตํ ปตฺโตติ. จตุนฺนํ ขฺวาหํ อาวุโส สติปฏฺฐานานํ ภาวิตตฺตา พหุลีกตตฺตา มหาภิญฺญตํ ปตฺโต. กตเมสํ จตุนฺนํ, อิธาหํ อาวุโส กาเย กายานุปสฺสี วิหรามิ ฯเปฯ. เวทนาสุ ฯเปฯ. จิตฺเต ฯเปฯ. ธมฺเมสุ ธมฺมานุปสฺสี วิหรามิ อาตาปี สมฺปชาโน สติมา วิเนยฺย โลเก อภิชฺฌา\-โทมนสฺสํ. อิเมสํ ขฺวาหํ อาวุโส จตุนฺนํ สติปฏฺฐานานํ ภาวิตตฺตา พหุลีกตตฺตา มหาภิญฺญตํ ปตฺโต. อิเมสญฺจ ปนาหํ อาวุโส จตุนฺนํ สติปฏฺฐานานํ ภาวิตตฺตา พหุลีกตตฺตา กปฺปสหสฺสํ อนุสฺสรามีติ.  ปฐมํ.}

\csromanheader{2}{2. อิทฺธิวิธ\-สุตฺต}
\proseitem{910}{อิเมสญฺจ ปนาหํ อาวุโส จตุนฺนํ สติปฏฺฐานานํ ภาวิตตฺตา พหุลีกตตฺตา อเนกวิหิตํ อิทฺธิวิธํ ปจฺจนุโภมิ, เอโกปิ หุตฺวา พหุธา โหมิ ฯเปฯ ยาว พฺรหฺมโลกาปิ กาเยน วสํ วตฺเตมีติ.  ทุติยํ.}

\csromanheader{2}{3. \textbf{ทิพฺพโสตสุตฺต}}
\proseitem{911}{อิเมสญฺจ ปนาหํ อาวุโส จตุนฺนํ สติปฏฺฐานานํ ภาวิตตฺตา พหุลีกตตฺตา ทิพฺพาย โสตธาตุยา วิสุทฺธาย อติกฺกนฺตมานุสิกาย อุโภ สทฺเท สุณามิ ทิพฺเพ จ มานุเส จ เย ทูเร สนฺติเก จาติ.  ตติยํ.}

\csromanheader{2}{4. เจโตปริย\-สุตฺต}
\proseitem{912}{อิเมสญฺจ ปนาหํ อาวุโส จตุนฺนํ สติปฏฺฐานานํ ภาวิตตฺตา พหุลีกตตฺตา ปรสตฺตานํ ปรปุคฺคลานํ เจตสา เจโต ปริจฺจ \csromanfolio{266}ปชานามิ, สราคํ วา จิตฺตํ “สราคํ จิตฺตนฺ”ติ ปชานามิ ฯเปฯ อวิมุตฺตํ วา จิตฺตํ “อวิมุตฺตํ จิตฺตนฺ”ติ ปชานามีติ.  จตุตฺถํ.}

\csromanheader{2}{5. \textbf{ฐานสุตฺต}}
\proseitem{913}{อิเมสญฺจ ปนาหํ อาวุโส จตุนฺนํ สติปฏฺฐานานํ ภาวิตตฺตา พหุลีกตตฺตา ฐานญฺจ ฐานโต อฏฺฐานญฺจ อฏฺฐานโต ยถาภูตํ ปชานามีติ.  ปญฺจมํ.}

\csromanheader{2}{6. กมฺมสมาทาน\-สุตฺต}
\proseitem{914}{อิเมสญฺจ ปนาหํ อาวุโส จตุนฺนํ สติปฏฺฐานานํ ภาวิตตฺตา พหุลีกตตฺตา อตีตา\-นาคต\-ปจฺจุปฺปนฺนานํ กมฺม\-สมาทานานํ ฐานโส เหตุโส วิปากํ ยถาภูตํ ปชานามีติ.  ฉฏฺฐํ.}

\csromanheader{2}{7. สพฺพตฺถคามินิ\-สุตฺต}
\proseitem{915}{อิเมสญฺจ ปนาหํ อาวุโส จตุนฺนํ สติปฏฺฐานานํ ภาวิตตฺตา พหุลีกตตฺตา สพฺพตฺถคามินิ\-ปฺปฏิปทํ ยถาภูตํ ปชานามีติ.  สตฺตมํ.}

\csromanheader{2}{8. \textbf{นานาธาตุสุตฺต}}
\proseitem{916}{อิเมสญฺจ ปนาหํ อาวุโส จตุนฺนํ สติปฏฺฐานานํ ภาวิตตฺตา พหุลีกตตฺตา อเนกธาตุ\-นานาธาตุ\-โลกํ ยถาภูตํ ปชานามีติ.  อฏฺฐมํ.}

\csromanheader{2}{9. นานา\-ธิมุตฺติ\-สุตฺต}
\proseitem{917}{อิเมสญฺจ ปนาหํ อาวุโส จตุนฺนํ สติปฏฺฐานานํ ภาวิตตฺตา พหุลีกตตฺตา สตฺตานํ นานา\-ธิมุตฺติกตํ ยถาภูตํ ปชานามีติ.  นวมํ.}

\csromanheader{2}{10. อินฺทฺริยปโรปริยตฺต\-สุตฺต}
\proseitem{918}{อิเมสญฺจ ปนาหํ อาวุโส จตุนฺนํ สติปฏฺฐานานํ ภาวิตตฺตา พหุลีกตตฺตา ปรสตฺตานํ ปรปุคฺคลานํ อินฺทฺริย\-ปโรปริยตฺตํ ยถาภูตํ ปชานามีติ.  ทสมํ.}

\csromanheader{6}{8. อนุรุทฺธ\-สํยุตฺต}
\csromanheader{2}{11. \textbf{ฌานาทิสุตฺต}}
\proseitem{919}{\csromanfolio{267}อิเมสญฺจ ปนาหํ อาวุโส จตุนฺนํ สติปฏฺฐานานํ ภาวิตตฺตา พหุลีกตตฺตา ฌาน\-วิโมกฺข\-สมาธิ\-สมาปตฺตีนํ สํกิเลสํ โวทานํ วุฏฺฐานํ ยถาภูตํ ปชานามีติ.  เอกาทสมํ.}

\csromanheader{2}{12. ปุพฺเพนิวาส\-สุตฺต}
\proseitem{920}{อิเมสญฺจ ปนาหํ อาวุโส จตุนฺนํ สติปฏฺฐานานํ ภาวิตตฺตา พหุลีกตตฺตา อเนกวิหิตํ ปุพฺเพนิวาสํ อนุสฺสารามิ. เสยฺยถิทํ, เอกมฺปิ ชาติํ เทฺวปิ ชาติโย ฯเปฯ. อิติ สาการํ สอุทฺเทสํ อเนกวิหิตํ ปุพฺเพนิวาสํ อนุสฺสรามีติ.  ทฺวาทสมํ.}

\csromanheader{2}{13. ทิพฺพจกฺขุ\-สุตฺต}
\proseitem{921}{อิเมสญฺจ ปนาหํ อาวุโส จตุนฺนํ สติปฏฺฐานานํ ภาวิตตฺตา พหุลีกตตฺตา ทิพฺเพน จกฺขุนา วิสุทฺเธน อติกฺกนฺต\-มานุสเกน สตฺเต ปสฺสามิ จวมาเน อุปปชฺชมาเน ฯเปฯ. อิติ ทิพฺเพน จกฺขุนา วิสุทฺเธน อติกฺกนฺต\-มานุสเกน ยถากมฺมูปเค สตฺเต ปชานามีติ.  เตรสมํ.}

\csromanheader{2}{14. อาสวกฺขย\-สุตฺต}
\proseitem{922}{อิเมสญฺจ ปนาหํ อาวุโส จตุนฺนํ สติปฏฺฐานานํ ภาวิตตฺตา พหุลีกตฺตา อาสวานํ ขยา อนาสวํ เจโตวิมุตฺติํ ปญฺญาวิมุตฺติํ ทิฏฺเฐว ธมฺเม สยํ อภิญฺญา สจฺฉิกตฺวา อุปสมฺปชฺช วิหรามีติ.  จุทฺทสมํ. ทุติโย วคฺโค.}

\vspace{\tipitakaparskip}
\csromancenter{\textbf{ตสฺสุทฺทานํ}}
\setlength{\csromangathapagewidth}{0pt}
\setlength{\csromangathawakwidth}{0pt}
\csromangathameasuregroup{มหาภิญฺญํ อิทฺธิ ทิพฺพํ, \\ สพฺพตฺถธาตุธิมุตฺติ,}{\csromangathabat{มหาภิญฺญํ อิทฺธิ ทิพฺพํ,}{เจโตปริยํ ฐานํ กมฺมํ.} \\ \csromangathabat{สพฺพตฺถธาตุธิมุตฺติ,}{อินฺทฺริยํ ฌานํ ติสฺโส วิชฺชาติ.} \\ อนุรุทฺธสํยุตฺตํ อฏฺฐมํ.}
\csromangathameasurewak{อนุรุทฺธสํยุตฺตํ อฏฺฐมํ.}
\csromangathasetleft{มหาภิญฺญํ อิทฺธิ ทิพฺพํ, \\ สพฺพตฺถธาตุธิมุตฺติ,}
\csromangathagroup{\csromangathastanza{\csromangathabat{มหาภิญฺญํ อิทฺธิ ทิพฺพํ,}{เจโตปริยํ ฐานํ กมฺมํ.} \\ \csromangathabat{สพฺพตฺถธาตุธิมุตฺติ,}{อินฺทฺริยํ ฌานํ ติสฺโส วิชฺชาติ.}}\csromangathastanza{\csromangathawak{อนุรุทฺธสํยุตฺตํ อฏฺฐมํ.}}}

\csromanmatikaanchor{mtk.201}
\csromantocmarknopage{subparagraph}{9. ฌานสํยุตฺต}{mtk.201}
\bookmark[dest={mtk.201},level=5]{9. ฌานสํยุตฺต}
\csromanheader{6}{9. \textbf{ฌานสํยุตฺต}}
\chapterhead{1. คงฺคา\-เปยฺยาล\-วคฺค}
\csromanmatikaanchor{mtk.202}
\csromanmatikaanchor{mtk.203}
\csromantocmarknopage{subparagraph}{1. คงฺคาเปยฺยาลวคฺค}{mtk.202}
\bookmark[dest={mtk.202},level=5]{1. คงฺคาเปยฺยาลวคฺค}
\csromantocmarkref{subparagraph}{1-12. ฌานาทิสุตฺต}{mtk.203}
\bookmark[dest={mtk.203},level=5]{1-12. ฌานาทิสุตฺต}
\csromanheader{6}{1-12. ฌานาทิสุตฺต}
\proseitem{923-934}{\csromanfolio{268}สาวตฺถิ\-นิทานํ. ตตฺร โข ฯเปฯ จตฺตาโรเม ภิกฺขเว ฌานา. กตเมน จตฺตาโร, อิธ ภิกฺขเว ภิกฺขุ วิวิจฺเจว กาเมหิ วิวิจฺจ อกุสเลหิ ธมฺเมหิ สวิตกฺกํ สวิจารํ วิเวกชํ ปีติสุขํ ปฐมํ ฌานํ อุปสมฺปชฺช วิหรติ. วิตกฺก\-วิจารานํ วูปสมา อชฺฌตฺตํ สมฺปสาทนํ เจตโส เอโกทิภาวํ อวิตกฺกํ อวิจารํ สมาธิชํ ปีติสุขํ ทุติยํ ฌานํ อุปสมฺปชฺช วิหรติ, ปีติยา จ วิราคา อุเปกฺขโก จ วิหรติ สโต จ สมฺปชาโน สุขญฺจ กาเยน ปฏิสํเวเทติ, ยํ ตํ อริยา อาจิกฺขนฺติ “อุเปกฺขโก สติมา สุขวิหารี”ติ, ตติยํ ฌานํ อุปสมฺปชฺช วิหรติ, สุขสฺส จ ปหานา ทุกฺขสฺส จ ปหานา ปุพฺเพว โสมนสฺส\-โทมนสฺสานํ อตฺถงฺคมา อทุกฺขมสุขํ อุเปกฺขา\-สติ\-ปาริสุทฺธิํ จตุตฺถํ ฌานํ อุปสมฺปชฺช วิหรติ. อิเม โข ภิกฺขเว จตฺตาโร ฌานาติ.}

\prosewithclosermidruled{เสยฺยถาปิ ภิกฺขเว คงฺคา นที ปาจีนนินฺนา ปาจีนโปณา ปาจีนปพฺภารา, เอวเมว โข ภิกฺขเว ภิกฺขุ จตฺตาโร ฌาเน ภาเวนฺโต จตฺตาโร ฌาเน พหุลีกโรนฺโต นิพฺพานนินฺโน โหติ นิพฺพานโปโณ นิพฺพาน\-ปพฺภาโร. กถญฺจ ภิกฺขเว ภิกฺขุ จตฺตาโร ฌาเน ภาเวนฺโต จตฺตาโร ฌาเน พหุลีกโรนฺโต นิพฺพานนินฺโน โหติ นิพฺพานโปโณ นิพฺพาน\-ปพฺภาโร, อิธ ภิกฺขเว ภิกฺขุ วิวิจฺเจว กาเมหิ วิวิจฺจ อกุสเลหิ ธมฺเมหิ สวิตกฺกํ สวิจารํ วิเวกชํ ปีติสุขํ ปฐมํ ฌานํ อุปสมฺปชฺช วิหรติ, วิตกฺก\-วิจารานํ วูปสมา ฯเปฯ ทุติยํ ฌานํ. ตติยํ ฌานํ. จตุตฺถํ ฌานํ อุปสมฺปชฺช วิหรติ. เอวํ โข ภิกฺขเว ภิกฺขุ จตฺตาโร ฌาเน ภาเวนฺโต จตฺตาโร ฌาเน พหุลีกโรนฺโต นิพฺพานนินฺโน โหติ นิพฺพานโปโณ นิพฺพาน\-ปพฺภาโรติ.  ทฺวาทสมํ.}{คงฺคา\-เปยฺยาล\-วคฺโค ปฐโม.}
\csromanheader{2}{9. ฌานสํยุตฺต}
\vspace{\tipitakaparskip}
\csromancenter{\textbf{ตสฺสุทฺทานํ}}
\setlength{\csromangathapagewidth}{0pt}
\setlength{\csromangathawakwidth}{0pt}
\csromangathameasuregroup{ฉ ปาจีนโต นินฺนา,}{\csromangathabat{ฉ ปาจีนโต นินฺนา,}{ฉ นินฺนา จ สมุทฺทโต.} \\ เทฺวเต ฉ ทฺวาทส โหนฺติ, วคฺโค เตน ปวุจฺจตีติ. \\ อปฺปมาทวคฺโค วิตฺถาเรตพฺโพ.}
\csromangathasetleft{ฉ ปาจีนโต นินฺนา,}
\csromangathagroup{\csromangathastanza{\csromanfolio{269}\csromangathabat{ฉ ปาจีนโต นินฺนา,}{ฉ นินฺนา จ สมุทฺทโต.} \\ เทฺวเต ฉ ทฺวาทส โหนฺติ, วคฺโค เตน ปวุจฺจตีติ. \\ อปฺปมาทวคฺโค วิตฺถาเรตพฺโพ.}}

\vspace{\tipitakaparskip}
\csromancenter{\textbf{ตสฺสุทฺทานํ}}
\setlength{\csromangathapagewidth}{0pt}
\setlength{\csromangathawakwidth}{0pt}
\csromangathameasuregroup{ตถาคตํ ปทํ กูฏํ,}{\csromangathabat{ตถาคตํ ปทํ กูฏํ,}{มูลํ สาโร จ วสฺสิกํ.} \\ ราชา จนฺทิมสูริยา, วตฺเถน ทสมํ ปทนฺติ. \\ พลกรณียวคฺโค วิตฺถาเรตพฺโพ.}
\csromangathasetleft{ตถาคตํ ปทํ กูฏํ,}
\csromangathagroup{\csromangathastanza{\csromangathabat{ตถาคตํ ปทํ กูฏํ,}{มูลํ สาโร จ วสฺสิกํ.} \\ ราชา จนฺทิมสูริยา, วตฺเถน ทสมํ ปทนฺติ. \\ พลกรณีย\-วคฺโค วิตฺถาเรตพฺโพ.}}

\vspace{\tipitakaparskip}
\csromancenter{\textbf{ตสฺสุทฺทานํ}}
\setlength{\csromangathapagewidth}{0pt}
\setlength{\csromangathawakwidth}{0pt}
\csromangathameasuregroup{พลํ พีชญฺจ นาโค จ,}{\csromangathabat{พลํ พีชญฺจ นาโค จ,}{รุกฺโข กุมฺเภน สูกิยา.} \\ อากาเสน จ เทฺว เมฆา, นาวา อาคนฺตุกา นทีติ. \\ เอสนาวคฺโค วิตฺถาเรตพฺโพ.}
\csromangathasetleft{พลํ พีชญฺจ นาโค จ,}
\csromangathagroup{\csromangathastanza{\csromangathabat{พลํ พีชญฺจ นาโค จ,}{รุกฺโข กุมฺเภน สูกิยา.} \\ อากาเสน จ เทฺว เมฆา, นาวา อาคนฺตุกา นทีติ. \\ เอสนาวคฺโค วิตฺถาเรตพฺโพ.}}

\vspace{\tipitakaparskip}
\csromancenter{\textbf{ตสฺสุทฺทานํ}}
\setlength{\csromangathapagewidth}{0pt}
\setlength{\csromangathawakwidth}{0pt}
\csromangathameasuregroup{เอสนา วิธา อาสโว, \\ ขิลํ มลญฺจ นีโฆ จ,}{\csromangathabat{เอสนา วิธา อาสโว,}{ภโว จ ทุกฺขตา ติสฺโส.} \\ \csromangathabat{ขิลํ มลญฺจ นีโฆ จ,}{เวทนา ตณฺหา ตสินา จาติ.}}
\csromangathasetleft{เอสนา วิธา อาสโว, \\ ขิลํ มลญฺจ นีโฆ จ,}
\csromangathagroup{\csromangathastanza{\csromangathabat{เอสนา วิธา อาสโว,}{ภโว จ ทุกฺขตา ติสฺโส.} \\ \csromangathabat{ขิลํ มลญฺจ นีโฆ จ,}{เวทนา ตณฺหา ตสินา จาติ.}}}

\chapterhead{5. \textbf{โอฆวคฺค}}
\csromanmatikaanchor{mtk.204}
\csromanmatikaanchor{mtk.205}
\csromantocmarknopage{subparagraph}{5. โอฆวคฺค}{mtk.204}
\bookmark[dest={mtk.204},level=5]{5. โอฆวคฺค}
\csromantocmarkref{subparagraph}{1-10. โอฆาทิสุตฺต}{mtk.205}
\bookmark[dest={mtk.205},level=5]{1-10. โอฆาทิสุตฺต}
\csromanheader{6}{1-10. โอฆาทิสุตฺต}
\proseitemwithclosermidruled{967-976}{ปญฺจิมานิ ภิกฺขเว อุทฺธมฺ\-ภาคิยานิ สํโยชนานิ. กตมานิ ปญฺจ, รูปราโค อรูปราโค มาโน อุทฺธจฺจํ อวิชฺชา. อิมานิ โข ภิกฺขเว \csromanfolio{270}ปญฺจุ\-ทฺธมฺภาคิยานิ สํโยชนานิ. อิเมสํ โข ภิกฺขเว ปญฺจนฺนํ อุทฺธมฺ\-ภาคิยานํ สํโยชนานํ อภิญฺญาย ปริญฺญาย ปริกฺขยาย ปหานาย จตฺตาโร ฌาเน ภาเวตพฺพา. กตเม จตฺตาโร, อิธ ภิกฺขเว ภิกฺขุ วิวิจฺเจว กาเมหิ วิวิจฺจ อกุสเลหิ ธมฺเมหิ สวิตกฺกํ สวิจารํ วิเวกชํ ปีติสุขํ ปฐมํ ฌานํ อุปสมฺปชฺช วิหรติ, วิตกฺก\-วิจารานํ วูปสมา อชฺฌตฺตํ สมฺปสาทนํ เจตโส เอโกทิภาวํ อวิตกฺกํ อวิจารํ สมาธิชํ ปีติสุขํ ทุติยํ ฌานํ ฯเปฯ ตติยํ ฌานํ ฯเปฯ จตุตฺถํ ฌานํ อุปสมฺปชฺช วิหรติ, อิเมสํ โข ภิกฺขเว ปญฺจนฺนํ อุทฺธมฺ\-ภาคิยานํ สํโยชนานํ อภิญฺญาย ปริญฺญาย ปริกฺขยาย ปหานาย อิเม จตฺตาโร ฌานา ภาเวตพฺพาติ วิตฺถาเรตพฺพํ.}{(ยถา มคฺคสํยุตฺตํ, ตถา วิตฺถาเรตพฺพํ.) โอฆวคฺโค ปญฺจโม.}
\csromancenter{\textbf{ตสฺสุทฺทานํ}}
\setlength{\csromangathapagewidth}{0pt}
\setlength{\csromangathawakwidth}{0pt}
\csromangathameasuregroup{โอโฆ โยโค อุปาทานํ, \\ กามคุณา นีวรณา,}{\csromangathabat{โอโฆ โยโค อุปาทานํ,}{คนฺถา อนุสเยน จ.} \\ \csromangathabat{กามคุณา นีวรณา,}{ขนฺธา โอรุทฺธมฺภาคิยาติ.} \\ \textbf{ฌานสํยุตฺตํ นวมํ.}}
\csromangathameasurewak{\textbf{ฌานสํยุตฺตํ นวมํ.}}
\csromangathasetleft{โอโฆ โยโค อุปาทานํ, \\ กามคุณา นีวรณา,}
\csromangathagroup{\csromangathastanza{\csromangathabat{โอโฆ โยโค อุปาทานํ,}{คนฺถา อนุสเยน จ.} \\ \csromangathabat{กามคุณา นีวรณา,}{ขนฺธา โอรุทฺธมฺภาคิยาติ.}}\csromangathastanza{\csromangathawak{\textbf{ฌานสํยุตฺตํ นวมํ.}}}}

\csromanmatikaanchor{mtk.206}
\csromanmatikaanchor{mtk.207}
\csromantocmarknopage{subparagraph}{10. อานาปานสํยุตฺต}{mtk.206}
\bookmark[dest={mtk.206},level=5]{10. อานาปานสํยุตฺต}
\csromantocmarknopage{subparagraph}{1. เอกธมฺมวคฺค}{mtk.207}
\bookmark[dest={mtk.207},level=5]{1. เอกธมฺมวคฺค}
\csromanheader{6}{1. \textbf{เอกธมฺมวคฺค}}
\csromanmatikaanchor{mtk.208}
\csromantocmarkref{subparagraph}{1. เอกธมฺมสุตฺต}{mtk.208}
\bookmark[dest={mtk.208},level=5]{1. เอกธมฺมสุตฺต}
\csromanheader{6}{1. \textbf{เอกธมฺมสุตฺต}}
\csromanmatikaanchor{mtk.209}
\csromantocmarkref{subparagraph}{2-4. โพชฺฌงฺคสุตฺตาทีนิ}{mtk.209}
\bookmark[dest={mtk.209},level=5]{2-4. โพชฺฌงฺคสุตฺตาทีนิ}
\proseitem{977}{\csromanfolio{271}สาวตฺถิ\-นิทานํ. ตตฺร โข ฯเปฯ เอตทโวจ\csromandash{}เอกธมฺโม ภิกฺขเว ภาวิโต พหุลีกโต มหปฺผโล โหติ มหานิสํโส. กตโม เอกธมฺโม, อานาปานสฺสติ\csromansharedfootnote{87c20ac719670e67}{อานาปานสติ (สี, อิ)}. กถํ ภาวิตา จ ภิกฺขเว อานาปานสฺสติ กถํ พหุลีกตา มหปฺผลา โหติ มหานิสํสา. อิธ ภิกฺขเว ภิกฺขุ อรญฺญคโต วา รุกฺขมูลคโต วา สุญฺญาคารคโต วา นิสีทติ ปลฺลงฺกํ อาภุชิตฺวา อุชุํ กายํ ปณิธาย ปริมุขํ สติํ อุปฏฺฐเปตฺวา, โส สโตว อสฺสสติ, สโตว\csromansharedfootnote{736e57e54ee87350}{สโต (พหูสุ) ตติยปาราชิเกปิ.} ปสฺสสติ. ทีฆํ วา อสฺสสนฺโต “ทีฆํ อสฺสสามี”ติ ปชานาติ, ทีฆํ วา ปสฺสสนฺโต “ทีฆํ ปสฺสสามี”ติ ปชานาติ, รสฺสํ วา ปสฺสสนฺโต “รสฺสํ ปสฺสสามี”ติ ปชานาติ, รสฺสํ วา ปสฺสสนฺโต “รสฺสํ ปสฺสสามี”ติ ปชานาติ. สพฺพกาย\-ปฺปฏิสํเวที “อสฺสสิสฺสามี”ติ สิกฺขติ, สพฺพกาย\-ปฺปฏิสํเวที “ปสฺสสิสฺสามี”ติ สิกฺขติ, ปสฺสมฺภยํ กายสงฺขารํ อสฺสสิสฺสามี”ติ สิกฺขติ, ปสฺสมฺภยํ กายสงฺขารํ “ปสฺสสิสฺสามี”ติ สิกฺขติ, ปีติปฺปฏิสํเวที “อสฺสสิสฺสามี”ติ สิกฺขติ, ปีติปฺปฏิสํเวที “ปสฺสสิสฺสามี”ติ สิกฺขติ, สุขปฺปฏิสํเวที “อสฺสสิสฺสามี”ติ สิกฺขติ สุขปฺปฏิสํเวที “ปสฺสสิสฺสามี”ติ สิกฺขติ, จิตฺตสงฺขาร\-ปฺปฏิสํเวที “อสฺสสิสฺสามี”ติ สิกฺขติ, จิตฺตสงฺขาร\-ปฺปฏิสํเวที “ปสฺสสิสฺสามี”ติ สิกฺขติ, ปสฺสมฺภยํ จิตฺตสงฺขารํ “อสฺสสิสฺสามี”ติ สิกฺขติ, ปสฺสมฺภยํ จิตฺตสงฺขารํ “ปสสิสฺสามี”ติ สิกฺขติ, จิตฺต\-ปฺปฏิสํเวที “อสฺสสิสฺสามี”ติ สิกฺขติ, จิตฺต\-ปฺปฏิสํเวที “ปสฺสสิสฺสามี”ติ สิกฺขติ, อภิปฺปโมทยํ จิตฺตํ “อสฺสสิสฺสามี”ติ สิกฺขติ, อภิปฺปโมทยํ จิตฺตํ “ปสฺสสิสฺสามี”ติ สิกฺขติ, สมาทหํ จิตฺตํ “อสฺสสิสฺสามี”ติ สิกฺขติ, สมาทหํ จิตฺตํ “ปสฺสสิสฺสามี”ติ สิกฺขติ, วิโมจยํ จิตฺตํ “อสฺสสิสฺสามี”ติ สิกฺขติ, \csromanfolio{272}วิโมจยํ จิตฺตํ “ปสฺสสิสฺสามี”ติ สิกฺขติ, อนิจฺจานุปสฺสี “อสฺสสิสฺสามี”ติ สิกฺขติ, อนิจฺจานุปสฺสี “ปสฺสสิสฺสามี”ติ สิกฺขติ, วิราคานุปสฺสี “อสฺสสิสฺสามี”ติ สิกฺขติ, วิราคานุปสฺสี “ปสฺสสิสฺสามี”ติ สิกฺขติ, “นิโรธานุปสฺสี อสฺสสิสฺสามี”ติ สิกฺขติ, “นิโรธานุปสฺสี ปสฺสสิสฺสามี”ติ สิกฺขติ, “ปฏินิสฺสคฺคา\-นุปสฺสี อสฺสสิสฺสามี”ติ สิกฺขติ, “ปฏินิสฺสคฺคา\-นุปสฺสี ปสฺสสิสฺสามี”ติ สิกฺขติ. เอวํ ภาวิตา โข ภิกฺขเว อานาปานสฺสติ เอวํ พหุลีกตา มหปฺผลา โหติ มหานิสํสาติ.  ปฐมํ.}

\csromanheader{2}{2. \textbf{โพชฺฌงฺคสุตฺต}}
\proseitem{978}{อานาปานสฺสติ ภิกฺขเว ภาวิตา พหูลีกตา มหปฺผลา โหติ มหานิสํสา. กถํ ภาวิตา จ ภิกฺขเว อานาปานสฺสติ กถํ พหุลีกตา มหปฺผลา โหติ มหานิสํสา, อิธ ภิกฺขเว ภิกฺขุ อานาปานสฺสติ\-สหคตํ สติสมฺโพชฺฌงฺคํ ภาเวติ วิเวกนิสฺสิตํ วิราคนิสฺสิตํ นิโรธ\-นิสฺสิตํ โวสฺสคฺค\-ปริณามิํ. อานาปานสฺสติ\-สหคตํ ธมฺมวิจย\-สมฺโพชฺฌงฺคํ ภาเวติ ฯเปฯ. อานาปานสฺสติ\-สหคตํ อุเปกฺขา\-สมฺโพชฺฌงฺคํ ภาเวติ วิเวกนิสฺสิตํ วิราคนิสฺสิตํ นิโรธ\-นิสฺสิตํ โวสฺสคฺค\-ปริณามิํ. เอวํ ภาวิตา โข ภิกฺขเว อานาปานสฺสติ เอวํ พหุลีกตา มหปฺผลา โหติ มหานิสํสาติ.  ทุติยํ.}

\csromanheader{2}{3. \textbf{สุทฺธิกสุตฺต}}
\proseitem{979}{อานาปานสฺสติ ภิกฺขเว ภาวิตา พหุลีกตา มหปฺผลา โหติ มหานิสํสา. กถํ ภาวิตา จ ภิกฺขเว อานาปานสฺสติ กถํ พหุลีกตา มหปฺผลา โหติ มหานิสํสา, อิธ ภิกฺขเว ภิกฺขุ อรญฺญคโต วา รุกฺขมูลคโต วา สุญฺญาคารคโต วา นิสีทติ ปลฺลงฺกํ อาภุชิตฺวา อุชุํ กายํ ปณิธาย ปริมุขํ สติํ อุปฏฺฐเปตฺวา, โส สโตว อสฺสสติ, สโตว ปสฺสสติ ฯเปฯ “ปฏินิสฺสคฺคา\-นุปสฺสี อสฺสสิสฺสามี”ติ สิกฺขติ, “ปฏินิสฺสคฺคา\-นุปสฺสี ปสฺสสิสฺสามี”ติ สิกฺขติ. เอวํ ภาวิตา โข ภิกฺขเว อานาปานสฺสติ เอวํ พหุลีกตา มหปฺผลา โหติ มหานิสํสาติ.  ตติยํ.}

\csromanheader{6}{10. อานาปาน\-สํยุตฺต}
\csromanheader{2}{4. ปฐม\-ผล\-สุตฺต}
\csromanmatikaanchor{mtk.210}
\csromantocmarkref{subparagraph}{5-6. ทุติยผลสุตฺตาทีนิ}{mtk.210}
\bookmark[dest={mtk.210},level=5]{5-6. ทุติยผลสุตฺตาทีนิ}
\proseitem{980}{\csromanfolio{273}อานาปานสฺสติ ภิกฺขเว ภาวิตา พหุลีกตา มหปฺผลา โหติ มหานิสํสา. กถํ ภาวิตา จ ภิกฺขเว อานาปานสฺสติ กถํ พหุลีกตา มหปฺผลา โหติ มหานิสํสา, อิธ ภิกฺขเว ภิกฺขุ อรญฺญคโต วา รุกฺขมูลคโต วา สุญฺญาคารคโต วา นิสีทติ ปลฺลงฺกํ อาภุชิตฺวา อุชุํ กายํ ปณิธาย ปริมุขํ สติํ อุปฏฺฐเปตฺวา, โส สโตว อสฺสสติ, สโตว ปสฺสสติ ฯเปฯ “ปฏินิสฺสคฺคา\-นุปสฺสี อสฺสสิสฺสามี”ติ สิกฺขติ, “ปฏินิสฺสคฺคา\-นุปสฺสี ปสฺสสิสฺสามี”ติ สิกฺขติ. เอวํ ภาวิตา โข ภิกฺขเว อานาปานสฺสติ เอวํ พหุลีกตา มหปฺผลา โหติ มหานิสํสา. เอวํ ภาวิตาย โข ภิกฺขเว อานาปานสฺสติยา เอวํ พหุลีกตาย ทฺวินฺนํ ผลานํ อญฺญตรํ ผลํ ปาฏิกงฺขํ ทิฏฺเฐว ธมฺเม อญฺญา, สติ วา อุปาทิเสเส อนาคามิตาติ.  จตุตฺถํ.}

\csromanheader{2}{5. ทุติย\-ผล\-สุตฺต}
\proseitem{981}{อานาปานสฺสติ ภิกฺขเว ภาวิตา พหุลีกตา มหปฺผลา โหติ มหานิสํสา. กถํ ภาวิตา จ ภิกฺขเว อานาปานสฺสติ กถํ พหุลีกตา มหปฺผลา โหติ มหานิสํสา, อิธ ภิกฺขเว ภิกฺขุ อรญฺญคโต วา รุกฺขมูลคโต วา สุญฺญาคารคโต วา นิสีทติ ปลฺลงฺกํ อาภุชิตฺวา อุชุํ กายํ ปณิธาย ปริมุขํ สติํ อุปฏฺฐเปตฺวา, โส สโตว อสฺสสติ, สโตว ปสฺสสติ ฯเปฯ “ปฏินิสฺสคฺคา\-นุปสฺสี อสฺสสิสฺสามี”ติ สิกฺขติ, “ปฏินิสฺสคฺคา\-นุปสฺสี ปสฺสสิสฺสามี”ติ สิกฺขติ. เอวํ ภาวิตา โข ภิกฺขเว อานาปานสฺสติ เอวํ พหุลีกตา มหปฺผลา โหติ มหานิสํสา.}

\prose{เอวํ ภาวิตาย โข ภิกฺขเว อานาปานสฺสติยา เอวํ พหุลีกตาย สตฺต ผลา สตฺตานิสํสา ปาฏิกงฺขา. กตเม สตฺต ผลา สตฺตานิสํสา, ทิฏฺเฐว ธมฺเม ปฏิกจฺจ อญฺญํ อาราเธติ, โน เจ ทิฏฺเฐว ธมฺเม ปฏิกจฺจ อญฺญํ อาราเธติ, อถ มรณกาเล อญฺญํ อาราเธติ, โน เจ ทิฏฺเฐว ธมฺเม ปฏิกจฺจ อญฺญํ อาราเธติ, โน เจ มรณกาเล อญฺญํ อาราเธติ, อถ ปญฺจนฺนํ โอรมฺภาคิยานํ สํโยชนานํ ปริกฺขยา อนฺตรา\-ปรินิพฺพายี โหติ, อุปหจฺจ\-ปรินิพฺพายี โหติ, อสงฺขารปรินิพฺพายี โหติ, สสงฺขาร\-ปรินิพฺพายี โหติ, อุทฺธํโสโต โหติ อกนิฏฺฐคามี. เอวํ \csromanfolio{274}ภาวิตาย โข ภิกฺขเว อานาปานสฺสติยา เอวํ พหุลีกตาย อิเม สตฺต ผลา สตฺตานิสํสา ปาฏิกงฺขาติ.  ปญฺจมํ.}

\csromanheader{2}{6. \textbf{อริฏฺฐสุตฺต}}
\proseitem{982}{สาวตฺถิ\-นิทานํ. ตตฺร โข ภควา ฯเปฯ เอตทโวจ “ภาเวถ โน ตุเมฺห ภิกฺขเว อานาปานสฺสตินฺ”ติ. เอวํ วุตฺเต อายสฺมา อริฏฺโฐ ภควนฺตํ เอตทโวจ “อหํ โข ภนฺเต ภาเวมิ อานาปานสฺสตินฺ”ติ. ยถา กถํ ปน ตฺวํ อริฏฺฐ ภาเวสิ อานาปาน\-สฺสตินฺติ. อตีเตสุ เม ภนฺเต กาเมสุ กามจฺฉนฺโท ปหีโน, อนาคเตสุ เม กาเมสุ กามจฺฉนฺโท วิคโต, อชฺฌตฺต\-พหิทฺธา\csromansharedfootnote{3d908726a5b89517}{อชฺฌตฺตํ พหิทฺธา (สฺยา, กํ, อิ, ก)} จ เม ธมฺเมสุ ปฏิฆสญฺญา สุปฺปฏิวินีตา, โส\csromansharedfootnote{971bf57e2b920fab}{โสหํ (?)} สโตว อสฺสสิสฺสามิ, สโตว ปสฺสสิสฺสามิ. เอวํ ขฺวาหํ ภนฺเต ภาเวมิ อานาปาน\-สฺสตินฺติ.}

\prose{อตฺเถสา อริฏฺฐ อานาปานสฺสติ, เนสา นตฺถีติ วทามิ, อปิ จ อริฏฺฐ ยถา อานาปานสฺสติ วิตฺถาเรน ปริปุณฺณา โหติ, ตํ สุณาหิ สาธุกํ มนสิ กโรหิ, ภาสิสฺสามีติ. “เอวํ ภนฺเต”ติ โข อายสฺมา อริฏฺโฐ ภควโต ปจฺจสฺโสสิ. ภควา เอตทโวจ\csromandash{}}

\prose{กถญฺจ อริฏฺฐ อานาปานสฺสติ วิตฺถาเรน ปริปุณฺณา โหติ, อิธ อริฏฺฐ ภิกฺขุ อรญฺญคโต วา รุกฺขมูลคโต วา สุญฺญาคารคโต วา นิสีทติ ปลฺลงฺกํ อาภุชิตฺวา อุชุํ กายํ ปณิธาย ปริมุขํ สติํ อุปฏฺฐเปตฺวา, โส สโตว อสฺสสติ, สโตว ปสฺสสติ, ทีฆํ วา อสฺสสนฺโต “ทีฆํ อสฺสสามี”ติ ปชานาติ ฯเปฯ “ปฏินิสฺสคฺคา\-นุปสฺสี อสฺสสิสฺสามี”ติ สิกฺขติ, “ปฏินิสฺสคฺคา\-นุปสฺสี ปสฺสสิสฺสามี”ติ สิกฺขติ. เอวํ โข อริฏฺฐ อานาปานสฺสติ วิตฺถาเรน ปริปุณฺณา โหตีติ.  ฉฏฺฐํ.}

\csromanmatikaanchor{mtk.211}
\csromantocmarkref{subparagraph}{7. มหากปฺปินสุตฺต}{mtk.211}
\bookmark[dest={mtk.211},level=5]{7. มหากปฺปินสุตฺต}
\csromanheader{6}{7. มหากปฺปิน\-สุตฺต}
\proseitem{983}{สาวตฺถิ\-นิทานํ. เตน โข ปน สมเยน อายสฺมา มหากปฺปิโน ภควโต อวิทูเร นิสินฺโน โหติ ปลฺลงฺกํ อาภุชิตฺวา อุชุํ กายํ ปณิธาย ปริมุขํ สติํ อุปฏฺฐเปตฺวา. อทฺทสา โข ภควา อายสฺมนฺตํ มหากปฺปินํ อวิทูเร นิสินฺนํ ปลฺลงฺกํ อาภุชิตฺวา อุชุํ กายํ ปณิธาย ปริมุขํ สติํ อุปฏฺฐเปตฺวา, ทิสฺวาน ภิกฺขู อามนฺเตสิ\csromandash{}}

\prose{\csromanfolio{275}ปสฺสถ โน ตุเมฺห ภิกฺขเว เอตสฺส ภิกฺขุโน กายสฺส อิญฺชิตตฺตํ วา ผนฺทิตตฺตํ วาติ. ยทาปิ มยํ ภนฺเต ตํ อายสฺมนฺตํ ปสฺสาม สํฆมชฺเฌ วา นิสินฺนํ เอกํ วา รโห นิสินฺนํ, ตทาปิ มยํ ตสฺส อายสฺมโต น ปสฺสาม กายสฺส อิญฺชิตตฺตํ วา ผนฺทิตตฺตํ วาติ.}

\prose{ยสฺส ภิกฺขเว สมาธิสฺส ภาวิตตฺตา พหุลีกตตฺตา เนว กายสฺส อิญฺชิตตฺตํ วา โหติ ผนฺทิตตฺตํ วา, น จิตฺตสฺส อิญฺชิตตฺตํ วา โหติ ผนฺทิตตฺตํ วา, ตสฺส โส ภิกฺขเว ภิกฺขุ สมาธิสฺส นิกามลาภี อกิจฺฉลาภี อกสิรลาภี. กตมสฺส จ ภิกฺขเว สมาธิสฺส ภาวิตตฺตา พหุลีกตตฺตา เนว กายสฺส อิญฺชิตตฺตํ วา โหติ ผนฺทิตตฺตํ วา, น จิตฺตสฺส อิญฺชิตตฺตํ วา โหติ ผนฺทิตตฺตํ วา.}

\prose{อานาปาน\-สฺสติ\-สมาธิสฺส ภิกฺขเว ภาวิตตฺตา พหุลีกตตฺตา เนว กายสฺส อิญฺชิตตฺตํ วา โหติ ผนฺทิตตฺตํ วา, น จิตฺตสฺส อิญฺชิตตฺตํ วา โหติ ผนฺทิตตฺตํ วา. กถํ ภาวิเต จ ภิกฺขเว อานาปาน\-สฺสติ\-สมาธิมฺหิ กถํ พหุลีกเต เนว กายสฺส อิญฺชิตตฺตํ วา โหติ ผนฺทิตตฺตํ วา, น จิตฺตสฺส อิญฺชิตตฺตํ วา โหติ ผนฺทิตตฺตํ วา.}

\prose{อิธ ภิกฺขเว ภิกฺขุ อรญฺญคโต วา รุกฺขมูลคโต วา สุญฺญาคารคโต วา นิสีทติ ปลฺลงฺกํ อาภุชิตฺวา อุชุํ กายํ ปณิธาย ปริมุขํ สติํ อุปฏฺฐเปตฺวา, โส สโตว อสฺสสติ, สโตว ปสฺสสติ ฯเปฯ “ปฏินิสฺสคฺคา\-นุปสฺสี อสฺสสิสฺสามี”ติ สิกฺขติ, “ปฏินิสฺสคฺคา\-นุปสฺสี ปสฺสสิสฺสามี”ติ สิกฺขติ. เอวํ ภาวิเต จ โข ภิกฺขเว อานาปาน\-สฺสติ\-สมาธิมฺหิ เอวํ พหุลีกเต เนว กายสฺส อิญฺชิตตฺตํ วา โหติ ผนฺทิตตฺตํ วา, น จิตฺตสฺส อิญฺชิตตฺตํ วา โหติ ผนฺทิตตฺตํ วาติ.  สตฺตมํ.}

\csromanmatikaanchor{mtk.212}
\csromantocmarkref{subparagraph}{8. ปทีโปปมสุตฺต}{mtk.212}
\bookmark[dest={mtk.212},level=5]{8. ปทีโปปมสุตฺต}
\csromanheader{6}{8. ปทีโปปม\-สุตฺต}
\proseitem{984}{อานาปาน\-สฺสติ\-สมาธิ ภิกฺขเว ภาวิโต พหุลีกโต มหปฺผโล โหติ มหานิสํโส. กถํ ภาวิโต จ ภิกฺขเว อานาปาน\-สฺสติ\-สมาธิ กถํ พหุลีกโต มหปฺผโล โหติ มหานิสํโส.}

\prose{อิธ ภิกฺขเว ภิกฺขุ อรญฺญคโต วา รุกฺขมูลคโต วา สุญฺญาคารคโต วา นิสีทติ ปลฺลงฺกํ อาภุชิตฺวา อุชุํ กายํ ปณิธาย \csromanfolio{276}ปริมุขํ สติํ อุปฏฺฐเปตฺวา, โส สโตว อสฺสสติ, สโตว ปสฺสสติ, ทีฆํ วา อสฺสสนฺโต “ทีฆํ อสฺสสามี”ติ ปชานาติ ฯเปฯ “ปฏินิสฺสคฺคา\-นุปสฺสี อสฺสสิสฺสามี”ติ สิกฺขติ, “ปฏินิสฺสคฺคา\-นุปสฺสี ปสฺสสิสฺสามี”ติ สิกฺขติ. เอวํ ภาวิโต โข ภิกฺขเว อานาปาน\-สฺสติ\-สมาธิ เอวํ พหุลีกโต มหปฺผโล โหติ มหานิสํโส.}

\prose{อหมฺปิ สุทํ ภิกฺขเว ปุพฺเพว สมฺโพธา อนภิสมฺพุทฺโธ โพธิสตฺโตว สมาโน อิมินา วิหาเรน พหุลํ วิหรามิ, ตสฺส มยฺหํ ภิกฺขเว อิมินา วิหาเรน พหุลํ วิหรโต เนว กาโย กิลมติ น จกฺขูนิ, อนุปาทาย จ เม อาสเวหิ จิตฺตํ วิมุจฺจิ.}

\prose{ตสฺมาติห ภิกฺขเว ภิกฺขุ เจปิ อากงฺเขยฺย “เนว เม กาโย กิลเมยฺย น จกฺขูนิ, อนุปาทาย จ เม อาสเวหิ จิตฺตํ วิมุจฺเจยฺยา”ติ, อยเมว อานาปาน\-สฺสติ\-สมาธิ สาธุกํ มนสิ กาตพฺโพ. ตสฺมาติห ภิกฺขเว ภิกฺขุ เจปิ อากงฺเขยฺย “เย เม เคหสิตา สรสงฺกปฺปา, เต ปหีเยยฺยุนฺ”ติ, อยเมว อานาปาน\-สฺสติ\-สมาธิ สาธุกํ มนสิ กาตพฺโพ. ตสฺมาติห ภิกฺขเว ภิกฺขุ เจปิ อากงฺเขยฺย “อปฺปฏิกูเล ปฏิกูลสญฺญี วิหเรยฺยนฺ”ติ, อยเมว อานาปาน\-สฺสติ\-สมาธิ สาธุกํ มนสิ กาตพฺโพ. ตสฺมาติห ภิกฺขเว ภิกฺขุ เจปิ อากงฺเขยฺย “ปฏิกูเล อปฺปฏิกูลสญฺญี วิหเรยฺยนฺ”ติ, อยเมว อานาปาน\-สฺสติ\-สมาธิ สาธุกํ มนสิ กาตพฺโพ. ตสฺมาติห ภิกฺขเว ภิกฺขุ เจปิ อากงฺเขยฺย “ปฏิกูเล จ อปฺปฏิกูเล จ ปฏิกูลสญฺญี วิหเรยฺยนฺ”ติ, อยเมว อานาปาน\-สฺสติ\-สมาธิ สาธุกํ มนสิ กาตพฺโพ. ตสฺมาติห ภิกฺขเว ภิกฺขุ เจปิ อากงฺเขยฺย “ปฏิกูเล จ อปฺปฏิกูเล, จ อปฺปฏิกูลสญฺญี วิหเรยฺยนฺ”ติ, อยเมว อานาปาน\-สฺสติ\-สมาธิ สาธุกํ มนสิ กาตพฺโพ. ตสฺมาติห ภิกฺขเว ภิกฺขุ เจปิ อากงฺเขยฺย “อปฺปฏิกูลญฺจ ปฏิกูลญฺจ ตทุภยํ อภินิวชฺเชตฺวา อุเปกฺขโก วิหเรยฺยํ สโต สมฺปชาโน”ติ, อยเมว อานาปาน\-สฺสติ\-สมาธิ สาธุกํ มนสิ กาตพฺโพ.}

\prose{ตสฺมาติห ภิกฺขเว ภิกฺขุ เจปิ อากงฺเขยฺย “วิวิจฺเจว กาเมหิ วิวิจฺจ อกุสเลหิ ธมฺเมหิ สวิตกฺกํ สวิจารํ วิเวกชํ ปีติสุขํ ปฐมํ ฌานํ อุปสมฺปชฺช วิหเรยฺยนฺ”ติ, อยเมว อานาปาน\-สฺสติ\-สมาธิ สาธุกํ มนสิ กาตพฺโพ. ตสฺมาติห ภิกฺขเว ภิกฺขุ เจปิ อากงฺเขยฺย}

\prose{\csromanfolio{277}“วิตกฺก\-วิจารานํ วูปสมา อชฺฌตฺตํ สมฺปสาทนํ เจตโส เอโกทิภาวํ อวิตกฺกํ อวิจารํ สมาธิชํ ปีติสุขํ ทุติยํ ฌานํ อุปสมฺปชฺช วิหเรยฺยนฺ”ติ, อยเมว อานาปาน\-สฺสติ\-สมาธิ สาธุกํ มนสิ กาตพฺโพ. ตสฺมาติห ภิกฺขเว ภิกฺขุ เจปิ อากงฺเขยฺย “ปีติยา จ วิราคา อุเปกฺขโก จ วิหเรยฺยํ, สโต จ สมฺปชาโน สุขญฺจ กาเยน ปฏิสํเวเทยฺยํ, ยํ ตํ อริยา อาจิกฺกนฺติ ‘อุเปกฺขโก สติมา สุขวิหารี’ติ, ตติยํ ฌานํ อุปสมฺปชฺช วิหเรยฺยนฺ”ติ, อยเมว อานาปาน\-สฺสติ\-สมาธิ สาธุกํ มนสิ กาตพฺโพ. ตสฺมาติห ภิกฺขเว ภิกฺขุ เจปิ อากงฺเขยฺย “สุขสฺส จ ปหานา ทุกฺขสฺส จ ปหานา ปุพฺเพว โสมนสฺส\-โทมนสฺสานํ อตฺถงฺคมา อทุกฺขมสุขํ อุเปกฺขา\-สติ\-ปาริสุทฺธิํ จตุตฺถํ ฌานํ อุปสมฺปชฺช วิหเรยฺยนฺ”ติ, อยเมว อานาปาน\-สฺสติ\-สมาธิ สาธุกํ มนสิ กาตพฺโพ.}

\prose{ตสฺมาติห ภิกฺขเว ภิกฺขุ เจปิ อากงฺเขยฺย “สพฺพโส รูปสญฺญานํ สมติกฺกมา ปฏิฆ\-สญฺญานํ อตฺถงฺคมา นานตฺต\-สญฺญานํ อมนสิการา ‘อนนฺโต อากาโส’ติ อากาสา\-นญฺจา\-ยตนํ อุปสมฺปชฺช วิหเรยฺยนฺ”ติ, อยเมว อานาปาน\-สฺสติ\-สมาธิ สาธุกํ มนสิ กาตพฺโพ. ตสฺมาติห ภิกฺขเว ภิกฺขุ เจปิ อากงฺเขยฺย “สพฺพโส อากาสา\-นญฺจา\-ยตนํ สมติกฺกมฺม ‘อนนฺตํ วิญฺญาณนฺ’ติ วิญฺญาณญฺจา\-ยตนํ อุปสมฺปชฺช วิหเรยฺยนฺ”ติ, อยเมว อานาปาน\-สฺสติ\-สมาธิ สาธุกํ มนสิ กาตพฺโพ. ตสฺมาติห ภิกฺขเว ภิกฺขุ เจปิ อากงฺเขยฺย “สพฺพโส วิญฺญาณญฺจา\-ยตนํ สมติกฺกมฺม ‘นตฺถิ กิญฺจี’ติ อากิญฺจญฺญา\-ยตนํ อุปสมฺปชฺช วิหเรยฺยนฺ”ติ, อยเมว อานาปาน\-สฺสติ\-สมาธิ สาธุกํ มนสิ กาตพฺโพ. ตสฺมาติห ภิกฺขเว ภิกฺขุ เจปิ อากงฺเขยฺย “สพฺพโส อากิญฺจญฺญา\-ยตนํ สมติกฺกมฺม เนว\-สญฺญา\-นา\-สญฺญา\-ยตนํ อุปสมฺปชฺช วิหเรยฺยนฺ”ติ, อยเมว อานาปาน\-สฺสติ\-สมาธิ สาธุกํ มนสิ กาตพฺโพ. ตสฺมาติห ภิกฺขเว ภิกฺขุ เจปิ อากงฺเขยฺย “สพฺพโส เนว\-สญฺญา\-นา\-สญฺญา\-ยตนํ สมติกฺกมฺม สญฺญาเวทยิต\-นิโรธํ อุปสมฺปชฺช วิหเรยฺยนฺ”ติ, อยเมว อานาปาน\-สฺสติ\-สมาธิ สาธุกํ มนสิ กาตพฺโพ.}

\prose{เอวํ ภาวิเต โข ภิกฺขเว อานาปาน\-สฺสติ\-สมาธิมฺหิ เอวํ พหุลีกเต สุขํ เจ เวทนํ เวทยติ, สา “อนิจฺจา”ติ ปชานาติ, “อนชฺโฌสิตา”ติ ปชานาติ, “อนภินนฺทิตา”ติ ปชานาติ. ทุกฺขํ เจ เวทนํ เวทยติ, สา “อนิจฺจา”ติ ปชานาติ, “อนชฺโฌสิตา”ติ ปชานาติ,}

\prose{\csromanfolio{278}“อนภินนฺทิตา”ติ ปชานาติ. อทุกฺขมสุขํ เจ เวทนํ เวทยติ, สา “อนิจฺจา”ติ ปชานาติ, “อนชฺโฌสิตา”ติ ปชานาติ, “อนภินนฺทิตา”ติ ปชานาติ.}

\prose{สุขํ\csromansharedfootnote{8f65f4ea3fc7c163}{โส สุขํ (สี, สฺยา, กํ, อิ) ม 3. 288 ปิฏฺเฐปิ. อฏฺฐกถาฏีกา โอโลเกตพฺพา.} เจ เวทนํ เวทยติ, วิสํยุตฺโต นํ เวทยติ. ทุกฺขํ เจ เวทนํ เวทยติ, วิสํยุตฺโต นํ เวทยติ. อทุกฺขมสุขํ เจ เวทนํ เวทยติ, วิสํยุตฺโต นํ เวทยติ. โส กาย\-ปริยนฺติกํ เวทนํ เวทยมาโน “กาย\-ปริยนฺติกํ เวทนํ เวทยามี”ติ ปชานาติ, ชีวิต\-ปริยนฺติกํ เวทนํ เวทยมาโน “ชีวิต\-ปริยนฺติกํ เวทนํ เวทยามี”ติ ปชานาติ, กายสฺส เภทา อุทฺธํ ชีวิต\-ปริยาทานา “อิเธว สพฺพ\-เวทยิตานิ อนภินนฺทิตานิ สีตีภวิสฺสนฺตี”ติ ปชานาติ.}

\prose{เสยฺยถาปิ ภิกฺขเว เตลญฺจ ปฏิจฺจ วฏฺฏิญฺจ ปฏิจฺจ เตลปฺปทีโป ฌาเยยฺย, ตสฺเสว เตลสฺส จ วฏฺฏิยา จ ปริยาทานา อนาหาโร นิพฺพาเยยฺย. เอวเมว โข ภิกฺขเว ภิกฺขุ กาย\-ปริยนฺติกํ เวทนํ เวทยมาโน “กาย\-ปริยนฺติกํ เวทนํ เวทยามี”ติ ปชานาติ, ชีวิต\-ปริยนฺติกํ เวทนํ เวทยมาโน “ชีวิต\-ปริยนฺติกํ เวทนํ เวทยามี”ติ ปชานาติ, กายสฺส เภทา อุทฺธํ ชีวิต\-ปริยาทานา “อิเธว สพฺพ\-เวทยิตานิ อนภินนฺทิตานิ สีตีภวิสฺสนฺตี”ติ ปชานาตีติ.  อฏฺฐมํ.}

\csromanmatikaanchor{mtk.213}
\csromantocmarkref{subparagraph}{9. เวสาลีสุตฺต}{mtk.213}
\bookmark[dest={mtk.213},level=5]{9. เวสาลีสุตฺต}
\csromanheader{6}{9. \textbf{เวสาลีสุตฺต}}
\proseitem{985}{เอวํ เม สุตํ\csromandash{}เอกํ สมยํ ภควา เวสาลิยํ วิหรติ มหาวเน กูฏาคาร\-สาลายํ. เตน โข ปน สมเยน ภควา ภิกฺขูนํ อเนก\-ปริยาเยน อสุภกถํ กเถติ, อสุภาย วณฺณํ ภาสติ, อสุภ\-ภาวนาย วณฺณํ ภาสติ.}

\prose{อถ โข ภควา ภิกฺขู อามนฺเตสิ “อิจฺฉามหํ ภิกฺขเว อฑฺฒมาสํ ปฏิสลฺลียิตุํ, นามฺหิ เกนจิ อุปสงฺกมิตพฺโพ อญฺญตฺร เอเกน ปิณฺฑปาต\-นีหารเกนา”ติ. “เอวํ ภนฺเต”ติ โข เต ภิกฺขู ภควโต ปฏิสฺสุตฺวา นาสฺสุธ โกจิ ภควนฺตํ อุปสงฺกมติ อญฺญตฺร เอเกน ปิณฺฑปาต\-นีหารเกน.}

\prose{\csromanfolio{279}อถ โข เต ภิกฺขู “ภควา อเนก\-ปริยาเยน อสุภกถํ กเถติ, อสุภาย วณฺณํ ภาสติ, อสุภ\-ภาวนาย วณฺณํ ภาสตี”ติ อเนกา\-การ\-โวการํ อสุภ\-ภาวนา\-นุโยคม\-นุยุตฺตา วิหรนฺติ. เต อิมินา กาเยน อฏฺฏียมานา\csromansharedfootnote{6b4a27443befbaaa}{อฏฺฏิยมานา (สี, สฺยา, กํ, อิ, ก)} หรายมานา ชิคุจฺฉมานา สตฺถหารกํ ปริเยสนฺติ, ทสปิ ภิกฺขู เอกาเหน สตฺถํ อาหรนฺติ, วีสมฺปิ ฯเปฯ ติํสมฺปิ ภิกฺขู เอกาเหน สตฺถํ อาหรนฺติ.}

\prose{อถ โข ภควา ตสฺส อฑฺฒมาสสฺส อจฺจเยน ปฏิสลฺลานา วุฏฺฐิโต อายสฺมนฺตํ อานนฺทํ อามนฺเตสิ “กิํ นุ โข อานนฺท ตนุภูโต วิย ภิกฺขุสํโฆ”ติ. ตถา หิ ปน ภนฺเต “ภควา ภิกฺขูนํ อเนก\-ปริยาเยน อสุภกตํ กเถติ, อสุภาย วณฺณํ ภาสติ, อสุภ\-ภาวนาย วณฺณํ ภาสตี”ติ อเนกา\-การ\-โวการํ อสุภ\-ภาวนา\-นุโยคม\-นุยุตฺตา วิหรนฺติ. เต อิมินา กาเยน อฏฺฏียมานา หรายมานา ชิคุจฺฉมานา สตฺถหารกํ ปริเยสนฺติ, ทสปิ ภิกฺขู เอกาเหน สตฺถํ อาหรนฺติ, วีสมฺปิ ภิกฺขู. ติํสมฺปิ ภิกฺขู เอกาเหน สตฺถํ อาหรนฺติ, สาธุ ภนฺเต ภควา อญฺญํ ปริยายํ อาจิกฺขตุ, ยถายํ ภิกฺขุสํโฆ อญฺญาย สณฺฐเหยฺยาติ.}

\prose{เตนหานนฺท ยาวติกา ภิกฺขู เวสาลิํ อุปนิสฺสาย วิหรนฺติ, เต สพฺเพ อุปฏฺฐาน\-สาลายํ สนฺนิปาเตหีติ. “เอวํ ภนฺเต”ติ โข อายสฺมา อานนฺโท ภควโต ปฏิสฺสุตฺวา ยาวติกา ภิกฺขู เวสาลิํ อุปนิสฺสาย วิหรนฺติ, เต สพฺเพ อุปฏฺฐาน\-สาลายํ สนฺนิปาเตตฺวา เยน ภควา เตนุปสงฺกมิ, อุปสงฺกมิตฺวา ภควนฺตํ เอตทโวจ “สนฺนิปติโต\csromansharedfootnote{ba9bb5507bacb105}{สนฺนิปาติโต (สี)} ภนฺเต ภิกฺขุสํโฆ, ยสฺสทานิ ภนฺเต ภควา กาลํ มญฺญตี”ติ.}

\prose{อถ โข ภควา เยน อุปฏฺฐานสาลา เตนุปสงฺกมิ, อุปสงฺกมิตฺวา ปญฺญตฺเต อาสเน นิสีทิ, นิสฺสชฺช โข ภควา ภิกฺขู อามนฺเตสิ\csromandash{}อยมฺปิ โข ภิกฺขเว อานาปาน\-สฺสติ\-สมาธิ ภาวิโต พหุลีกโต สนฺโต เจว ปณีโต จ อเสจนโก จ สุโข จ วิหาโร, อุปฺปนฺนุปฺปนฺเน จ ปาปเก อกุสเล ธมฺเม ฐานโส อนฺตรธาเปติ วูปสเมติ.}

\prose{\csromanfolio{280}เสยฺยถาปิ ภิกฺขเว คิมฺหานํ ปจฺฉิเม มาเส อูหตํ รโชชลฺลํ, ตเมนํ มหาอกาลเมโฆ ฐานโส อนฺตรธาเปติ วูปสเมติ. เอวเมว โข ภิกฺขเว อานาปาน\-สฺสติ\-สมาธิ ภาวิโต พหุลีกโต สนฺโต เจว ปณีโต จ อเสจนโก จ สุโข จ วิหาโร, อุปฺปนฺนุปฺปนฺเน จ ปาปเก อกุสเล ธมฺเม ฐานโส อนฺตรธาเปติ วูปสเมติ. กถํ ภาวิโต จ ภิกฺขเว อานาปาน\-สฺสติ\-สมาธิ กถํ พหุลีกโต สนฺโต เจว ปณีโต จ อเสจนโก จ สุโข จ วิหาโร, อุปฺปนฺนุปฺปนฺเน จ ปาปเก อกุสเล ธมฺเม ฐานโส อนฺตรธาเปติ วูปสเมติ.}

\prose{อิธ ภิกฺขเว ภิกฺขุ อรญฺญคโต วา รุกฺขมูลคโต วา สุญฺญาคารคโต วา นิสีทติ ปลฺลงฺกํ อาภุชิตฺวา อุชุํ กายํ ปณิธาย ปริมุขํ สติํ อุปฏฺฐเปตฺวา, โส สโตว อสฺสสติ, สโตว ปสฺสสติ ฯเปฯ “ปฏินิสฺสคฺคา\-นุปสฺสี อสฺสสิสฺสามี”ติ สิกฺขติ, “ปฏินิสฺสคฺคา\-นุปสฺสี ปสฺสสิสฺสามี”ติ สิกฺขติ. เอวํ ภาวิโต โข ภิกฺขเว อานาปาน\-สฺสติ\-สมาธิ เอวํ พหุลีกโต สนฺโต เจว ปณีโต จ อเสจนโก จ สุโข จ วิหาโร, อุปฺปนฺนุปฺปนฺเน จ ปาปเก อกุสเล ธมฺเม ฐานโส อนฺตรธาเปติ วูปสเมตีติ.  นวมํ.}

\csromanmatikaanchor{mtk.214}
\csromantocmarkref{subparagraph}{10. กิมิลสุตฺต}{mtk.214}
\bookmark[dest={mtk.214},level=5]{10. กิมิลสุตฺต}
\csromanheader{6}{10. \textbf{กิมิลสุตฺต}}
\proseitem{986}{เอวํ เม สุตํ\csromandash{}เอกํ สมยํ ภควา กิมิลายํ\csromansharedfootnote{1bff7655608296ac}{กิมฺพิลายํ (สี, อิ)} วิหรติ เวฬุวเน. ตตฺร โข ภควา อายสฺมนฺตํ กิมิลํ อามนฺเตสิ “กถํ ภาวิโต นุ โข กิมิล อานาปาน\-สฺสติ\-สมาธิ กถํ พหุลีกโต มหปฺผโล โหติ มหานิสํโส”ติ.}

\prose{เอวํ วุตฺเต อายสฺมา กิมิโล ตุณฺหี อโหสิ. ทุติยมฺปิ โข ภควา ฯเปฯ. ตติยมฺปิ โข ภควา อายสฺมนฺตํ กิมิลํ อามนฺเตสิ “กถํ ภาวิโต นุ โข กิมิล อานาปาน\-สฺสติ\-สมาธิ กถํ พหุลีกโต มหปฺผโล โหติ มหานิสํโส”ติ. ตติยมฺปิ โข อายสฺมา กิมิโล ตุณฺหี อโหสิ.}

\prose{เอวํ วุตฺเต อายสฺมา อานนฺโท ภควนฺตํ เอตทโวจ “เอตสฺส ภควา กาโล, เอตสฺส สุคต กาโล, ยํ ภควา อานาปาน\-สฺสติ\-สมาธิํ}

\prose{\csromanfolio{281}ภาเสยฺย, ภควโต สุตฺวา ภิกฺขู ธาเรสฺสนฺตีติ. เตนหานนฺท สุณาหิ สาธุกํ มนสิ กโรหิ, ภาสิสฺสามีติ. “เอวํ ภนฺเต”ติ โข อายสฺมา อานนฺโท ภควโต ปจฺจสฺโสสิ. ภควา เอตทโวจ\csromandash{}}

\prose{กถํ ภาวิโต จ อานนฺท อานาปาน\-สฺสติ\-สมาธิ กถํ พหุลีกโต มหปฺผโล โหติ มหานิสํโส, อิธานนฺท ภิกฺขุ อรญฺญคโต วา รุกฺขมูลคโต วา สุญฺญาคารคโต ว นิสีทติ ปลฺลงฺกํ อาภุชิตฺวา อุชุํ กายํ ปณิธาย ปริมุขํ สติํ อุปฏฺฐเปตฺวา, โส สโตว อสฺสสติ, สโตว ปสฺสสติ ฯเปฯ “ปฏินิสฺสคฺคา\-นุปสฺสี อสฺสสิสฺสามี”ติ สิกฺขติ, “ปฏินิสฺสคฺคา\-นุปสฺสี ปสฺสสิสฺสามี”ติ สิกฺขติ. เอวํ ภาวิโต โข อานนฺท อานาปาน\-สฺสติ\-สมาธิ เอวํ พหุลีกโต มหปฺผโล โหติ มหานิสํโส.}

\prose{ยสฺมิํ สมเย อานนฺท ภิกฺขุ ทีฆํ วา อสฺสสนฺโต “ทีฆํ อสฺสสามี”ติ ปชานาติ, ทีฆํ วา ปสฺสสนฺโต “ทีฆํ ปสฺสสามี”ติ ปชานาติ, รสฺสํ วา อสฺสสนฺโต “รสฺสํ อสฺสสามี”ติ ปชานาติ, รสฺสํ วา ปสฺสสนฺโต “รสฺสํ ปสฺสสามี”ติ ปชานาติ, “สพฺพกาย\-ปฺปฏิสํเวที อสฺสสิสฺสามี”ติ สิกฺขติ, “สพฺพกาย\-ปฺปฏิสํเวที ปสฺสสิสฺสามี”ติ สิกฺขติ, “ปสฺสมฺภยํ กายสงฺขารํ อสฺสสิสฺสามี”ติ สิกฺขติ, “ปสฺสมฺภยํ กายสงฺขารํ ปสฺสสิสฺสามี”ติ สิกฺขติ, กาเย กายานุปสฺสี อานนฺท ภิกฺขุ ตสฺมิํ สมเย วิหรติ อาตาปี สมฺปชาโน สติมา วิเนยฺย โลเก อภิชฺฌา\-โทมนสฺสํ. ตํ กิสฺส เหตุ, กายญฺญตราหํ อานนฺท เอตํ วทามิ ยทิทํ อสฺสาสปสฺสาสํ, ตสฺมาติหานนฺท กาเย กายานุปสฺสี ภิกฺขุ ตสฺมิํ สมเย วิหรติ อาตาปี สมฺปชาโน สติมา วิเนยฺย โลเก อภิชฺฌา\-โทมนสฺสํ. (1)}

\prose{ยสฺมิํ สมเย อานนฺท ภิกฺขุ “ปีติปฺปฏิสํเวที อสฺสสิสฺสามี”ติ สิกฺขติ, “ปีติปฺปฏิสํเวที ปสฺสสิสฺสามี”ติ สิกฺขติ, “สุขปฺปฏิสํเวที อสฺสสิสฺสามี”ติ สิกฺขติ, “สุขปฺปฏิสํเวที ปสฺสสิสฺสามี”ติ สิกฺขติ, “จิตฺตสงฺขาร\-ปฺปฏิสํเวที อสฺสสิสฺสามี”ติ สิกฺขติ, “จิตฺตสงฺขาร\-ปฺปฏิสํเวที ปสฺสสิสฺสามี”ติ สิกฺขติ, “ปสฺสมฺภยํ จิตฺตสงฺขารํ ปสฺสสิสฺสามี”ติ สิกฺขติ, เวทนาสุ เวทนานุปสฺสี อานนฺท ภิกฺขุ ตสฺมิํ สมเย วิหรติ อาตาปี}

\prose{\csromanfolio{282}สมฺปชาโน สติมา วิเนยฺย โลเก อภิชฺฌา\-โทมนสฺสํ. ตํ กิสฺส เหตุ, เวทนาญฺญตราหํ อานนฺท เอตํ วทามิ ยทิทํ อสฺสาส\-ปสฺสาสานํ\csromansharedfootnote{0f2595c108b3fb3d}{อสฺสาสปสฺสาสํ (อิ, ก) ม 3. 126 ปิฏฺเฐปิ.} สาธุกํ มนสิการํ. ตสฺมาติหานนฺท เวทนาสุ เวทนานุปสฺสี ภิกฺขุ ตสฺมิํ สมเย วิหรติ อาตาปี สมฺปชาโน สติมา วิเนยฺย โลเก อภิชฺฌา\-โทมนสฺสํ. (2)}

\prose{ยสฺมิํ สมเย อานนฺท ภิกฺขุ “จิตฺต\-ปฺปฏิสํเวที อสฺสสิสฺสามี”ติ สิกฺขติ, “จิตฺต\-ปฺปฏิสํเวที ปสฺสสิสฺสามี”ติ สิกฺขติ, อภิปฺปโมทยํ จิตฺตํ ฯเปฯ สมาทหํ จิตฺตํ. “วิโมจยํ จิตฺตํ อสฺสสิสฺสามี”ติ สิกฺขติ, “วิโมจยํ จิตฺตํ ปสฺสสิสฺสามี”ติ สิกฺขติ. จิตฺเต จิตฺตานุปสฺสี อานนฺท ภิกฺขุ ตสฺมิํ สมเย วิหรติ อาตาปี สมฺปชาโน สติมา วิเนยฺย โลเก อภิชฺฌา\-โทมนสฺสํ. ตํ กิสฺส เหตุ, นาหํ อานนฺท มุฏฺฐสฺสติสฺส อสมฺปชานสฺส อานาปาน\-สฺสติ\-สมาธิภาวนํ วทามิ. ตสฺมาติหานนฺท จิตฺเต จิตฺตานุปสฺสี ภิกฺขุ ตสฺมิํ สมเย วิหรติ อาตาปี สมฺปชาโน สติมา วิเนยฺย โลเก อภิชฺฌา\-โทมนสฺสํ. (3)}

\prose{ยสฺมิํ สมเย อานนฺท ภิกฺขุ “อนิจฺจานุปสฺสี อสฺสสิสฺสามี”ติ สิกฺขติ ฯเปฯ วิราคานุปสฺสี. นิโรธานุปสฺสี. “ปฏินิสฺสคฺคา\-นุปสฺสี อสฺสสิสฺสามี”ติ สิกฺขติ. “ปฏินิสฺสคฺคา\-นุปสฺสี ปสฺสสิสฺสามี”ติ สิกฺขติ. ธมฺเมสุ ธมฺมานุปสฺสี อานนฺท ภิกฺขุ ตสฺมิํ สมเย วิหรติ อาตาปี สมฺปชาโน สติมา วิเนยฺย โลเก อภิชฺฌา\-โทมนสฺสํ. โส ยํ ตํ โหติ อภิชฺฌา\-โทมนสฺสานํ ปหานํ, ตํ ปญฺญาย ทิสฺวา สาธุกํ อชฺฌุเปกฺขิตา โหติ. ตสฺมาติหานนฺท ธมฺเมสุ ธมฺมานุปสฺสี ภิกฺขุ ตสฺมิํ สมเย วิหรติ อาตาปี สมฺปชาโน สติมา วิเนยฺย โลเก อภิชฺฌา\-โทมนสฺสํ. (4)}

\prose{เสยฺยถาปิ อานนฺท จตุมหาปเถ\csromansharedfootnote{49d7f3b18fabb048}{จาตุมฺมหาปเถ (สี, สฺยา, กํ)} มหาปํสุปุญฺโช ปุรตฺถิมาย เจปิ ทิสาย อาคจฺเฉยฺย สกฏํ วา รโถ วา, อุปหนเตว ตํ ปํสุปุญฺชํ. ปจฺฉิมาย เจปิ ทิสาย อาคจฺเฉยฺย. อุตฺตราย เจปิ ทิสาย. ทกฺขิณาย เจปิ ทิสาย อาคจฺเฉยฺย สกฏํ วา รโถ วา, อุปหนเตว ตํ ปํสุปุญฺชํ. เอวเมว โข อานนฺท ภิกฺขุ กาเย กายานุปสฺสี วิหรนฺโตปิ}

\prosewithclosermidruled{\csromanfolio{283}อุปหนเตว ปาปเก อกุสเล ธมฺเม. เวทนาสุ ฯเปฯ. จิตฺเต ฯเปฯ. ธมฺเมสุ ธมฺมานุปสฺสี วิหรนฺโตปิ อุปหนเตว ปาปเก อกุสเล ธมฺเมติ.  ทสมํ.}{เอกธมฺมวคฺโค ปฐโม.}
\csromancenter{\textbf{ตสฺสุทฺทานํ}}
\setlength{\csromangathapagewidth}{0pt}
\setlength{\csromangathawakwidth}{0pt}
\csromangathameasuregroup{เอกธมฺโม จ โพชฺฌงฺโค, \\ อริฏฺโฐ กปฺปิโน ทีโป,}{\csromangathabat{เอกธมฺโม จ โพชฺฌงฺโค,}{สุทฺธิกญฺจ ทุเว ผลา.} \\ \csromangathabat{อริฏฺโฐ กปฺปิโน ทีโป,}{เวสาลี กิมิเลน จาติ.}}
\csromangathasetleft{เอกธมฺโม จ โพชฺฌงฺโค, \\ อริฏฺโฐ กปฺปิโน ทีโป,}
\csromangathagroup{\csromangathastanza{\csromangathabat{เอกธมฺโม จ โพชฺฌงฺโค,}{สุทฺธิกญฺจ ทุเว ผลา.} \\ \csromangathabat{อริฏฺโฐ กปฺปิโน ทีโป,}{เวสาลี กิมิเลน จาติ.}}}

\chapterhead{2. \textbf{ทุติยวคฺค}}
\csromanmatikaanchor{mtk.215}
\csromanmatikaanchor{mtk.216}
\csromantocmarknopage{subparagraph}{2. ทุติยวคฺค}{mtk.215}
\bookmark[dest={mtk.215},level=5]{2. ทุติยวคฺค}
\csromantocmarkref{subparagraph}{1. อิจฺฉานงฺคลสุตฺต}{mtk.216}
\bookmark[dest={mtk.216},level=5]{1. อิจฺฉานงฺคลสุตฺต}
\csromanheader{6}{1. อิจฺฉานงฺคล\-สุตฺต}
\proseitem{987}{เอกํ สมยํ ภควา อิจฺฉานงฺคเล วิหรติ อิจฺฉานงฺคล\-วนสณฺเฑ. ตตฺร โข ภควา ภิกฺขู อามนฺเตสิ “อิจฺฉามหํ ภิกฺขเว เตมาสํ ปฏิสลฺลียิตุํ, นามฺหิ เกนจิ อุปสงฺกมิตพฺโพ อญฺญตฺร เอเกน ปิณฺฑปาต\-นีหารเกนา”ติ. “เอวํ ภนฺเต”ติ โข เต ภิกฺขู ภควโต ปฏิสฺสุตฺวา นาสฺสุธ โกจิ ภควนฺตํ อุปสงฺกมติ อญฺญตฺร เอเกน ปิณฺฑปาต\-นีหารเกน.}

\prose{อถ โข ภควา ตสฺส เตมาสสฺส อจฺจเยน ปฏิสลฺลานา วุฏฺฐิโต ภิกฺขู อามนฺเตสิ\csromandash{}สเจ โข ภิกฺขเว อญฺญติตฺถิยา ปริพฺพาชกา เอวํ ปุจฺเฉยฺยุํ “กตเมนาวุโส วิหาเรน สมโณ โคตโม วสฺสาวาสํ พหุลํ วิหาสี”ติ, เอวํ ปุฏฺฐา ตุเมฺห ภิกฺขเว เตสํ อญฺญ\-ติตฺถิยานํ ปริพฺพาชกานํ เอวํ พฺยากเรยฺยาถ “อานาปาน\-สฺสติ\-สมาธินา โข อาวุโส ภควา วสฺสาวาสํ พหุลํ วิหาสี”ติ. อิธาหํ ภิกฺขเว สโต อสฺสสามิ, สโต ปสฺสสามิ. ทีฆํ อสฺสสนฺโต “ทีฆํ อสฺสสามี”ติ ปชานามิ, ทีฆํ ปสฺสสนฺโต “ทีฆํ ปสฺสสามี”ติ ปชานามิ, รสฺสํ อสฺสสนฺโต “รสฺสํ อสฺสสามี”ติ ปชานามิ, รสฺสํ ปสฺสสนฺโต “รสฺสํ ปสฺสสามี”ติ ปชานามิ, “สพฺพกาย\-ปฺปฏิสํเวที อสฺสสิสฺสามี”ติ ปชานามิ ฯเปฯ “ปฏินิสฺสคฺคา\-นุปสฺสี อสฺสสิสฺสามี”ติ ปชานามิ, “ปฏินิสฺสคฺคา\-นุปสฺสี ปสฺสสิสฺสามี”ติ ปชานามิ.}

\prose{\csromanfolio{284}ยํ หิ ตํ ภิกฺขเว สมฺมา วทมาโน วเทยฺย อริยวิหาโร อิติปิ, พฺรหฺมวิหาโร อิติปิ, ตถาคต\-วิหาโร อิติปิ. อานาปาน\-สฺสติ\-สมาธิํ สมฺมา วทมาโน วเทยฺย อริยวิหาโร อิติปิ, พฺรหฺมวิหาโร อิติปิ, ตถาคต\-วิหาโร อิติปิ. เย เต ภิกฺขเว ภิกฺขู เสขา อปฺปตฺตมานสา อนุตฺตรํ โยคกฺเขมํ ปตฺถยมานา วิหรนฺติ, เตสํ อานาปาน\-สฺสติ\-สมาธิ ภาวิโต พหุลีกโต อาสวานํ ขยาย สํวตฺตติ, เย จ โข เต ภิกฺขเว ภิกฺขู อรหนฺโต ขีณาสวา วุสิตวนฺโต กตกรณียา โอหิตภารา อนุปฺปตฺต\-สทตฺถา ปริกฺขีณ\-ภวสํโยชนา สมฺมทญฺญา\-วิมุตฺตา, เตสํ อานาปาน\-สฺสติ\-สมาธิ ภาวิโต พหุลีกโต ทิฏฺฐ\-ธมฺม\-สุข\-วิหาราย เจว สํวตฺตติ สติสมฺปชญฺญาย จ.}

\prose{ยํ หิ ตํ ภิกฺขเว สมฺมา วทมาโน วเทยฺย อริยวิหาโร อิติปิ, พฺรหฺมวิหาโร อิติปิ, ตถาคต\-วิหาโร อิติปิ. อานาปาน\-สฺสติ\-สมาธิํ สมฺมา วทมาโน วเทยฺย อริยวิหาโร อิติปิ, พฺรหฺมวิหาโร อิติปิ, ตถาคต\-วิหาโร อิติปีติ.  ปฐมํ.}

\csromanmatikaanchor{mtk.217}
\csromantocmarkref{subparagraph}{2. กงฺเขยฺยสุตฺต}{mtk.217}
\bookmark[dest={mtk.217},level=5]{2. กงฺเขยฺยสุตฺต}
\csromanheader{6}{2. \textbf{กงฺเขยฺยสุตฺต}}
\proseitem{988}{เอกํ สมยํ อายสฺมา โลมสกํภิโย\csromansharedfootnote{d5c823dbdb812325}{โลมสวงฺคิโส (สี, อิ)} สกฺเกสุ วิหรติ กปิล\-วตฺถุสฺมิํ นิโคฺรธาราเม. อถ โข มาหานาโม สกฺโก เยนายสฺมา โลมสกํภิโย เตนุปสงฺกมิ, อุปสงฺกมิตฺวา อายสฺมนฺตํ โลมสกํภิยํ อภิวาเทตฺวา เอกมนฺตํ นิสีทิ, เอกมนฺตํ นิสินฺโน โข มหานาโม สกฺโก อายสฺมนฺตํ โลมสกํภิยํ เอตทโวจ “โส เอว นุ โข ภนฺเต เสโข วิหาโร โส ตถาคต\-วิหาโร, อุทาหุ อญฺโญว\csromansharedfootnote{a27447da1f7c5647}{อญฺโญ (สฺยา, กํ, อิ, ก)} เสโข วิหาโร อญฺโญ ตถาคต\-วิหาโร”ติ.}

\prose{น โข อาวุโส มหานาม เสฺวว เสโข วิหาโร, โส ตถาคต\-วิหาโร, อญฺโญ โข อาวุโส มหานาม เสโข วิหาโร, อญฺโญ ตถาคต\-วิหาโร. เย เต อาวุโส มหานาม ภิกฺขู เสขา อปฺปตฺตมานสา อนุตฺตรํ โยคกฺเขมํ ปตฺถยมานา วิหรนฺติ, เต ปญฺจ นีวรเณ ปหาย วิหรนฺติ. กตเม ปญฺจ, กามจฺฉนฺท\-นีวรณํ ปหาย วิหรนฺติ, พฺยาปาท\-นีวรณํ ฯเปฯ ถินมิทฺธ\-นีวรณํ. อุทฺธจฺจ\-กุกฺกุจฺจ\-นีวรณํ. วิจิกิจฺฉา\-นีวรณํ ปหาย วิหรนฺติ.}

\prose{\csromanfolio{285}เยปิ เต อาวุโส มหานาม ภิกฺขู เสขา อปฺปตฺตมานสา อนุตฺตรํ โยคกฺเขมํ ปตฺถยมานา วิหรนฺติ, เต อิเม ปญฺจ นีวรเณ ปหาย วิหรนฺติ.}

\prose{เย จ โข เต อาวุโส มหานาม ภิกฺขู อรหนฺโต ขีณาสวา วุสิตวนฺโต กตกรณียา โอหิตภารา อนุปฺปตฺต\-สทตฺถา ปริกฺขีณ\-ภวสํโยชนา สมฺมทญฺญา\-วิมุตฺตา, เตสํ ปญฺจ นีวรณา ปหีนา อุจฺฉินฺนมูลา ตาลาวตฺถุกตา อนภาวํกตา\csromansharedfootnote{fe0270e15c5a9573}{อนภาวกตา (สี, อิ)} อายติํ อนุปฺปาทธมฺมา. กตเม ปญฺจ, กามจฺฉนฺท\-นีวรณํ ปหีนํ อุจฺฉินฺนมูลํ ตาลา\-วตฺถุกตํ อนภาวํกตํ อายติํ อนุปฺปาท\-ธมฺมํ. พฺยาปาท\-นีวรณํ ปหีนํ ฯเปฯ ถินมิทฺธ\-นีวรณํ. อุทฺธจฺจ\-กุกฺกุจฺจ\-นีวรณํ. วิจิกิจฺฉา\-นีวรณํ ปหีนํ อุจฺฉินฺนมูลํ กาลาวตฺถุกตํ อนภาวํกตํ อายติํ อนุปฺปาท\-ธมฺมํ.}

\prose{เย เต อาวุโส มหานาม ภิกฺขู อรหนฺโต ขีณาสวา วุสิตวนฺโต กตกรณียา โอหิตภารา อนุปฺปตฺต\-สทตฺถา ปริกฺขีณ\-ภวสํโยชนา สมฺมทญฺญา\-วิมุตฺตา, เตสํ อิเม ปญฺจ นีวรณา ปหีนา อุจฺฉินฺนมูลา ตาลาวตฺถุกตา อนภาวํกตา อายติํ อนุปฺปาทธมฺมา. ตทมินาเปตํ อาวุโส มหานาม ปริยาเยน เวทิตพฺพํ, ยถา อญฺโญว เสโข วิหาโร, อญฺโญ ตถาคต\-วิหาโร.}

\prose{เอกมิทํ อาวุโส มหานาม สมยํ ภควา อิจฺฉานงฺคเล วิหรติ อิจฺฉานงฺคล\-วนสณฺเฑ. ตตฺร โข อาวุโส มหานาม ภควา ภิกฺขู อามนฺเตสิ “อิจฺฉามหํ ภิกฺขเว เตมาสํ ปฏิสลฺลียิตุํ, นามฺหิ เกนจิ อุปสงฺกมิตพฺโพ อญฺญตฺร เอเกน ปิณฺฑปาต\-นีหารเกนา”ติ. “เอวํ ภนฺเต”ติ โข อาวุโส มหานาม เต ภิกฺขู ภควโต ปฏิสฺสุตฺวา นาสฺสุธ โกจิ ภควนฺตํ อุปสงฺกมติ อญฺญตฺร เอเกน ปิณฺฑปาต\-นีหารเกน.}

\prose{อถ โข อาวุโส ภควา ตสฺส เตมาสสฺส อจฺจเยน ปฏิสลฺลานา วุฏฺฐิโต ภิกฺขู อามนฺเตสิ\csromandash{}สเจ โข ภิกฺขเว อญฺญติตฺถิยา ปริพฺพาชกา เอวํ ปุจฺเฉยฺยุํ “กตเมนาวุโส วิหาเรน สมโณ โคตโม วสฺสาวาสํ พหุลํ วิหาสี”ติ. เอวํ ปุฏฺฐา ตุเมฺห ภิกฺขเว}

\csromanmatikaanchor{mtk.218}
\csromantocmarkref{subparagraph}{3-4. อานนฺทสุตฺตานิ}{mtk.218}
\bookmark[dest={mtk.218},level=5]{3-4. อานนฺทสุตฺตานิ}
\prose{\csromanfolio{286}เตสํ อญฺญ\-ติตฺถิยานํ ปริพฺพาชกานํ เอวํ พฺยากเรยฺยาถ “อานาปาน\-สฺสติ\-สมาธินา โข อาวุโส ภควา วสฺสาวาสํ พหุลํ วิหาสี”ติ. อิธาหํ ภิกฺขเว สโต อสฺสสามิ, สโต ปสฺสสามิ, ทีฆํ อสฺสสนฺโต “ทีฆํ อสฺสสามี”ติ ปชานามิ, ทีฆํ ปสฺสสนฺโต “ทีฆํ ปสฺสสามี”ติ ปชานามิ ฯเปฯ “ปฏินิสฺสคฺคา\-นุปสฺสี อสฺสสิสฺสามี”ติ ปชานามิ, “ปฏินิสฺสคฺคา\-นุปสฺสี ปสฺสสิสฺสามี”ติ ปชานามิ.}

\prose{ยํ หิ ตํ ภิกฺขเว สมฺมา วทมาโน วเทยฺย อริยวิหาโร อิติปิ, พฺรหฺมวิหาโร อิติปิ, ตถาคต\-วิหาโร อิติปิ. อานาปาน\-สฺสติ\-สมาธิํ สมฺมา วทมาโน วเทยฺย อริยวิหาโร อิติปิ, พฺรหฺมวิหาโร อิติปิ, ตถาคต\-วิหาโร อิติปิ.}

\prose{เย เต ภิกฺขเว ภิกฺขู เสขา อปฺปตฺตมานสา อนุตฺตรํ โยคกฺเขมํ ปตฺถยมานา วิหรนฺติ, เตสํ อานาปาน\-สฺสติ\-สมาธิ ภาวิโต พหุลีกโต อาสวานํ ขยาย สํวตฺตติ.}

\prose{เย จ โข เต ภิกฺขเว ภิกฺขู อรหนฺโต ขีณาสวา วุสิตวนฺโต กตกรณียา โอหิตภารา อนุปฺปตฺต\-สทตฺถา ปริกฺขีณ\-ภวสํโยชนา สมฺมทญฺญา\-วิมุตฺตา, เตสํ อานาปาน\-สฺสติ\-สมาธิ ภาวิโต พหุลีกโต ทิฏฺเฐว ธมฺเม สุขวิหาราย เจว สํวตฺตติ สติสมฺปชญฺญาย จ.}

\prose{ยํ หิ ตํ ภิกฺขเว สมฺมา วทมาโน วเทยฺย อริยวิหาโร อิติปิ, พฺรหฺมวิหาโร อิติปิ, ตถาคต\-วิหาโร อิติปิ. อานาปาน\-สฺสติ\-สมาธิํ สมฺมา วทมาโน วเทยฺย อริยวิหาโร อิติปิ, พฺรหฺมวิหาโร อิติปิ, ตถาคต\-วิหาโร อิติปีติ. อิมินา โข เอตํ อาวุโส มหานาม ปริยาเยน เวทิตพฺพํ, ยถา อญฺโญว เสโข วิหาโร, อญฺโญ ตถาคต\-วิหาโรติ.}

\prose{3. ปฐม\-อานนฺท\-สุตฺต}

\proseitem{989}{สาวตฺถิ\-นิทานํ. อถ โข อายสฺมา อานนฺโท เยน ภควา เตนุปสงฺกมิ, อุปสงฺกมิตฺวา ภควนฺตํ อภิวาเทตฺวา เอกมนฺตํ นิสีทิ, เอกมนฺตํ นิสินฺโน โข อายสฺมา อานนฺโท ภควนฺตํ เอตทโวจ “อตฺถิ นุ โข}

\prose{\csromanfolio{287}ภนฺเต เอกธมฺโม\csromansharedfootnote{3af4af4157c14d09}{เอโก ธมฺโม (สี)} ภาวิโต พหุลีกโต จตฺตาโร ธมฺเม ปริปูเรติ, จตฺตาโร ธมฺมา ภาวิตา พหุลีกตา สตฺต ธมฺเม ปริปูเรนฺติ, สตฺต ธมฺมา ภาวิตา พหุลีกตา เทฺว ธมฺเม ปริปูเรนฺตี”ติ.}

\prose{อตฺถิ โข อานนฺท เอกธมฺโม ภาวิโต พหุลีกโต จตฺตาโร ธมฺเม ปริปูเรติ, จตฺตาโร ธมฺมา ภาวิตา พหุลีกตา สตฺต ธมฺเม ปริปูเรนฺติ. สตฺต ธมฺมา ภาวิตา พหุลีกตา เทฺว ธมฺเม ปริปูเรนฺตีติ.}

\prose{กตโม ปน ภนฺเต เอกธมฺโม\csromansharedfootnote{3af4af4157c14d09}{เอโก ธมฺโม (สี)} ภาวิโต พหุลีกโต จตฺตาโร ธมฺเม ปริปูเรติ, จตฺตาโร ธมฺมา ภาวิตา พหุลีกตา สตฺต ธมฺเม ปริปูเรนฺติ, สตฺต ธมฺมา ภาวิตา พหุลีกตา เทฺว ธมฺเม ปริปูเรนฺตีติ. อานาปาน\-สฺสติ\-สมาธิ โข อานนฺท เอกธมฺโม ภาวิโต พหุลีกโต จตฺตาโร สติปฏฺฐาเน ปริปูเรติ, จตฺตาโร สติปฏฺฐานา ภาวิตา พหุลีกตา สตฺต โพชฺฌงฺเค ปริปูเรนฺติ, สตฺต โพชฺฌงฺคา ภาวิตา พหุลีกตา วิชฺชาวิมุตฺติํ ปริปูเรนฺติ.}

\prose{กถํ ภาวิโต อานนฺท อานาปาน\-สฺสติ\-สมาธิ กถํ พหุลีกโต จตฺตาโร สติปฏฺฐาเน ปริปูเรติ. อิธานนฺท ภิกฺขุ อรญฺญคโต วา รุกฺขมูลคโต วา สุญฺญาคารคโต วา นิสีทติ ปลฺลงฺกํ อาภุชิตฺวา อุชุํ กายํ ปณิธาย ปริมุขํ สติํ อุปฏฺฐเปตฺวา, โส สโตว อสฺสสติ, สโตว ปสฺสสติ, ทีฆํ วา อสฺสสนฺโต “ทีฆํ อสฺสสามี”ติ ปชานาติ, ทีฆํ วา ปสฺสสนฺโต “ทีฆํ ปสฺสสามี”ติ ปชานาติ ฯเปฯ “ปฏินิสฺสคฺคา\-นุปสฺสี อสฺสสิสฺสามี”ติ สิกฺขติ, “ปฏินิสฺสคฺคา\-นุปสฺสี ปสฺสสิสฺสามี”ติ สิกฺขติ, ยสฺมิํ สมเย อานนฺท ภิกฺขุ ทีฆํ วา อสฺสสนฺโต “ทีฆํ อสฺสสามี”ติ ปชานาติ, ทีฆํ วา ปสฺสสนฺโต “ทีฆํ ปสฺสสามี”ติ ปชานาติ, รสฺสํ วา ฯเปฯ “ปสฺสมฺภยํ กายสงฺขารํ อสฺสสิสฺสามี”ติ สิกฺขติ, “ปสฺสมฺภยํ กายสงฺขารํ ปสฺสสิสฺสามี”ติ สิกฺขติ. กาเย กายานุปสฺสี อานนฺท ภิกฺขุ ตสฺมิํ สมเย วิหรติ อาตาปี สมฺปชาโน สติมา วิเนยฺย โลเก อภิชฺฌา\-โทมนสฺสํ. ตํ กิสฺส เหตุ, กายญฺญตราหํ อานนฺท เอตํ วทามิ ยทิทํ อสฺสาสปสฺสาสํ. ตสฺมาติหานนฺท กาเย กายานุปสฺสี ภิกฺขุ ตสฺมิํ สมเย วิหรติ อาตาปี สมฺปชาโน สติมา วิเนยฺย โลเก อภิชฺฌา\-โทมนสฺสํ. (1)}

\prose{\csromanfolio{288}ยสฺมิํ สมเย อานนฺท ภิกฺขุ “ปีติปฺปฏิสํเวที อสฺสสิสฺสามี”ติ สิกฺขติ, สุขปฺปฏิสํเวที ฯเปฯ จิตฺตสงฺขาร\-ปฺปฏิสํเวที. “ปสฺสมฺภยํ จิตฺตสงฺขารํ อสฺสสิสฺสามี”ติ สิกฺขติ, “ปสฺสมฺภยํ จิตฺตสงฺขารํ ปสฺสสิสฺสามี”ติ สิกฺขติ. เวทนาสุ เวทนานุปสฺสี อานนฺท ภิกฺขุ ตสฺมิํ สมเย วิหรติ อาตาปี สมฺปชาโน สติมา วิเนยฺย โลเก อภิชฺฌา\-โทมนสฺสํ. ตํ กิสฺส เหตุ, เวทนาญฺญตราหํ อานนฺท เอตํ วทามิ, ยทิทํ อสฺสาส\-ปสฺสาสานํ สาธุกํ มนสิการํ. ตสฺมาติหานนฺท เวทนาสุ เวทนานุปสฺสี ภิกฺขุ ตสฺมิํ สมเย วิหรติ อาตาปี สมฺปชาโน สติมา วิเนยฺย โลเก อภิชฺฌา\-โทมนสฺสํ. (2)}

\prose{ยสฺมิํ สมเย อานนฺท ภิกฺขุ “จิตฺต\-ปฺปฏิสํเวที อสฺสสิสฺสามี”ติ สิกฺขติ, “จิตฺต\-ปฺปฏิสํเวที ปสฺสสิสฺสามี”ติ สิกฺขติ, อภิปฺปโมทยํ จิตฺตํ ฯเปฯ สมาทหํ จิตฺตํ. “วิโมจยํ จิตฺตํ อสฺสสิสฺสามี”ติ สิกฺขติ, “วิโมจยํ จิตฺตํ ปสฺสสิสฺสามี”ติ สิกฺขติ. จิตฺเต จิตฺตานุปสฺสี อานนฺท ภิกฺขุ ตสฺมิํ สมเย วิหรติ อาตาปี สมฺปชาโน สติมา วิเนยฺย โลเก อภิชฺฌา\-โทมนสฺสํ. ตํ กิสฺส เหตุ, นาหํ อานนฺท มุฏฺฐสฺสติสฺส อสมฺปชานสฺส อานาปาน\-สฺสติ\-สมาธิภาวนํ วทามิ. ตสฺมา ติหานนฺท จิตฺเต จิตฺตานุปสฺสี ภิกฺขุ ตสฺมิํ สมเย วิหรติ อาตาปี สมฺปชาโน สติมา วิเนยฺย โลเก อภิชฺฌา\-โทมนสฺสํ. (3)}

\prose{ยสฺมิํ สมเย อานนฺท ภิกฺขุ อนิจฺจานุปสฺสี ฯเปฯ วิราคานุปสฺสี. นิโรธานุปสฺสี. “ปฏินิสฺสคฺคา\-นุปสฺสี อสฺสสิสฺสามี”ติ สิกฺขติ, “ปฏินิสฺสคฺคา\-นุปสฺสี ปสฺสสิสฺสามี”ติ สิกฺขติ. ธมฺเมสุ ธมฺมานุปสฺสี อานนฺท ภิกฺขุ ตสฺมิํ สมเย วิหรติ อาตาปี สมฺปชาโน สติมา วิเนยฺย โลเก อภิชฺฌา\-โทมนสฺสํ. โส ยํ ตํ โหติ อภิชฺฌา\-โทมนสฺสานํ ปหานํ, ตํ ปญฺญาย ทิสฺวา สาธุกํ อชฺฌุเปกฺขิตา โหติ. ตสฺมาติหานนฺท ธมฺเมสุ ธมฺมานุปสฺสี ภิกฺขุ ตสฺมิํ สมเย วิหรติ อาตาปี สมฺปชาโน สติมา วิเนยฺย โลเก อภิชฺฌา\-โทมนสฺสํ. (4)}

\prose{เอวํ ภาวิโต โข อานนฺท อานาปาน\-สฺสติ\-สมาธิ เอวํ พหุลีกโต จตฺตาโร สติปฏฺฐาเน ปริปูเรติ.}

\prose{กถํ ภาวิตา จานนฺท จตฺตาโร สติปฏฺฐานา กถํ พหุลีกตา สตฺต โพชฺฌงฺเค ปริปูเรนฺติ, ยสฺมิํ สมเย อานนฺท ภิกฺขุ กาเย}

\prose{\csromanfolio{289}กายานุปสฺสี วิหรติ, อุปฏฺฐิตาสฺส\csromansharedfootnote{59b917bf55a9ab57}{อุปฏฺฐิตสฺสติ (อิ, ก)} ตสฺมิํ สมเย ภิกฺขุโน\csromansharedfootnote{7a1e2b18e2defeea}{ตสฺมิํ สมเย อานนฺท ภิกฺขุโน (อิ, ก), ตสฺมิํ สมเย (ม 3. 127 ปิฏฺเฐ.)} สติ โหติ อสมฺมุฏฺฐา. ยสฺมิํ สมเย อานนฺท ภิกฺขุโน อุปฏฺฐิตา สติ โหติ อสมฺมุฏฺฐา, สติสมฺโพชฺฌงฺโค ตสฺมิํ สมเย ภิกฺขุโน อารทฺโธ โหติ, สติสมฺโพชฺฌงฺคํ ตสฺมิํ สมเย ภิกฺขุ ภาเวติ, สติสมฺโพชฺฌงฺโค ตสฺมิํ สมเย ภิกฺขุโน ภาวนา\-ปาริปูริํ คจฺฉติ. (1)}

\prose{โส ตถาสโต วิหรนฺโต ตํ ธมฺมํ ปญฺญาย ปวิจินติ ปวิจรติ ปริวีมํสมา\-ปชฺชติ. ยสฺมิํ สมเย อานนฺท ภิกฺขุ ตถาสโต วิหรนฺโต ตํ ธมฺมํ ปญฺญาย ปวิจินติ ปวิจรติ ปริวีมํสมา\-ปชฺชติ. ธมฺมวิจย\-สมฺโพชฺฌงฺโค ตสฺมิํ สมเย ภิกฺขุโน อารทฺโธ โหติ, ธมฺมวิจย\-สมฺโพชฺฌงฺคํ ตสฺมิํ สมเย ภิกฺขุ ภาเวติ, ธมฺมวิจย\-สมฺโพชฺฌงฺโค ตสฺมิํ สมเย ภิกฺขุโน ภาวนา\-ปาริปูริํ คจฺฉติ. (2)}

\prose{ตสฺส ตํ ธมฺมํ ปญฺญาย ปวิจินโต ปวิจรโต ปริวีมํสมาปชฺชโต อารทฺธํ โหติ วีริยํ อสลฺลีนํ. ยสฺมิํ สมเย อานนฺท ภิกฺขุโน ตํ ธมฺมํ ปญฺญาย ปวิจินโต ปวิจรโต ปริวีมํสมาปชฺชโต อารทฺธํ โหติ วีริยํ อสลฺลีนํ, วีริย\-สมฺโพชฺฌงฺโค ตสฺมิํ สมเย ภิกฺขุโน อารทฺโธ โหติ, วีริย\-สมฺโพชฺฌงฺคํ ตสฺมิํ สมเย ภิกฺขุ ภาเวติ, วีริย\-สมฺโพชฺฌงฺโค ตสฺมิํ สมเย ภิกฺขุโน ภาวนา\-ปาริปูริํ คจฺฉติ. (3)}

\prose{อารทฺธ\-วีริยสฺส อุปฺปชฺชติ ปีติ นิรามิสา. ยสฺมิํ สมเย อานนฺท ภิกฺขุโน อารทฺธ\-วีริยสฺส อุปฺปชฺชติ ปีติ นิรามิสา, ปีติสมฺโพชฺฌงฺโค ตสฺมิํ สมเย ภิกฺขุโน อารทฺโธ โหติ, ปีติสมฺโพชฺฌงฺคํ ตสฺมิํ สมเย ภิกฺขุ ภาเวติ, ปีติสมฺโพชฺฌงฺโค ตสฺมิํ สมเย ภิกฺขุโน ภาวนา\-ปาริปูริํ คจฺฉติ. (4)}

\prose{ปีติมนสฺส กาโยปิ ปสฺสมฺภติ, จิตฺตมฺปิ ปสฺสมฺภติ. ยสฺมิํ สมเย อานนฺท ภิกฺขุโน ปีติมนสฺส กาโยปิ ปสฺสมฺภติ, จิตฺตมฺปิ ปสฺสมฺภติ, ปสฺสทฺธิ\-สมฺโพชฺฌงฺโค ตสฺมิํ สมเย ภิกฺขุโน อารทฺโธ โหติ, ปสฺสทฺธิ\-สมฺโพชฺฌงฺคํ ตสฺมิํ สมเย ภิกฺขุ ภาเวติ, ปสฺสทฺธิ\-สมฺโพชฺฌงฺโค ตสฺมิํ สมเย ภิกฺขุโน ภาวนา\-ปาริปูริํ คจฺฉติ. (5)}

\prose{\csromanfolio{290}ปสฺสทฺธ\-กายสฺส สุขิโน จิตฺตํ สมาธิยติ. ยสฺมิํ สมเย อานนฺท ภิกฺขุโน ปสฺสทฺธ\-กายสฺส สุขิโน จิตฺตํ สมาธิยติ, สมาธิ\-สมฺโพชฺฌงฺโค ตสฺมิํ สมเย ภิกฺขุโน อารทฺโธ โหติ, สมาธิ\-สมฺโพชฺฌงฺคํ ตสฺมิํ สมเย ภิกฺขุ ภาเวติ, สมาธิ\-สมฺโพชฺฌงฺโค ตสฺมิํ สมเย ภิกฺขุโน ภาวนา\-ปาริปูริํ คจฺฉติ. (6)}

\prose{โส ตถาสมาหิตํ จิตฺตํ สาธุกํ อชฺฌุเปกฺขิตา โหติ. ยสฺมิํ สมเย อานนฺท ภิกฺขุ ตถาสมาหิตํ จิตฺตํ สาธุกํ อชฺฌุเปกฺขิตา โหติ, อุเปกฺขา\-สมฺโพชฺฌงฺโค ตสฺมิํ สมเย ภิกฺขุโน อารทฺโธ โหติ, อุเปกฺขา\-สมฺโพชฺฌงฺคํ ตสฺมิํ สมเย ภิกฺขุ ภาเวติ, อุเปกฺขา\-สมฺโพชฺฌงฺโค ตสฺมิํ สมเย ภิกฺขุโน ภาวนา\-ปาริปูริํ คจฺฉติ. (7)}

\prose{ยสฺมิํ สมเย อานนฺท ภิกฺขุ เวทนาสุ ฯเปฯ. จิตฺเต ฯเปฯ. ธมฺเมสุ ธมฺมานุปสฺสี วิหรติ, อุปฏฺฐิตาสฺส ตสฺมิํ สมเย ภิกฺขุโน สติ โหติ อสมฺมุฏฺฐา. ยสฺมิํ สมเย อานนฺท ภิกฺขุโน อุปฏฺฐิตา สติ โหติ อสมฺมุฏฺฐา, สติสมฺโพชฺฌงฺโค ตสฺมิํ สมเย ภิกฺขุโน อารทฺโธ โหติ, สติสมฺโพชฺฌงฺคํ ตสฺมิํ สมเย ภิกฺขุ ภาเวติ, สติสมฺโพชฺฌงฺโค ตสฺมิํ สมเย ภิกฺขุโน ภาวนา\-ปาริปูริํ คจฺฉติ.}

\vspace{\tipitakaparskip}
\csromancenter{(ยถา ปฐมํ สติปฏฺฐานํ, เอวํ วิตฺถาเรตพฺพํ.)}
\prose{โส ตถาสมาหิตํ จิตฺตํ สาธุกํ อชฺฌุเปกฺขิตา โหติ, ยสฺมิํ สมเย อานนฺท ภิกฺขุ ตถาสมาหิตํ จิตฺตํ สาธุกํ อชฺฌุเปกฺขิตา โหติ, อุเปกฺขา\-สมฺโพชฺฌงฺโค ตสฺมิํ สมเย ภิกฺขุโน อารทฺโธ โหติ, อุเปกฺขา\-สมฺโพชฺฌงฺคํ ตสฺมิํ สมเย ภิกฺขุ ภาเวติ, อุเปกฺขา\-สมฺโพชฺฌงฺโค ตสฺมิํ สมเย ภิกฺขุโน ภาวนา\-ปาริปูริํ คจฺฉติ. เอวํ ภาวิตา โข อานนฺท จตฺตาโร สติปฏฺฐานา เอวํ พหุลีกตา สตฺต โพชฺฌงฺเค ปริปูเรนฺติ.}

\prose{กถํ ภาวิตา อานนฺท สตฺต โพชฺฌงฺคา กถํ พหุลีกตา วิชฺชาวิมุตฺติํ ปริปูเรนฺติ. อิธานนฺท ภิกฺขุ สติสมฺโพชฺฌงฺคํ ภาเวติ วิเวกนิสฺสิตํ วิราคนิสฺสิตํ นิโรธ\-นิสฺสิตํ โวสฺสคฺค\-ปริณามิํ. ธมฺมวิจย\-สมฺโพชฺฌงฺคํ}

\csromanmatikaanchor{mtk.219}
\csromantocmarkref{subparagraph}{5-6. ภิกฺขุสุตฺตานิ}{mtk.219}
\bookmark[dest={mtk.219},level=5]{5-6. ภิกฺขุสุตฺตานิ}
\prose{\csromanfolio{291}ภาเวติ ฯเปฯ. อุเปกฺขา\-สมฺโพชฺฌงฺคํ ภาเวติ วิเวกนิสฺสิตํ วิราคนิสฺสิตํ นิโรธ\-นิสฺสิตํ โวสฺสคฺค\-ปริณามิํ. เอวํ ภาวิตา โข อานนฺท สตฺต โพชฺฌงฺคา เอวํ พหุลีกตา วิชฺชาวิมุตฺติํ ปริปูเรนฺตีติ.  ตติยํ.}

\prose{4. ทุติย\-อานนฺท\-สุตฺต}

\proseitem{990}{อถ โข อายสฺมา อานนฺโท เยน ภควา เตนุปสงฺกมิ, อุปสงฺกมิตฺวา ภควนฺตํ อภิวาเทตฺวา เอกมนฺตํ นิสีทิ, เอกมนฺตํ นิสินฺนํ โข อายสฺมนฺตํ อานนฺทํ ภควา เอตทโวจ “อตฺถิ นุ โข อานนฺท เอกธมฺโม ภาวิโต พหุลีกโต จตฺตาโร ธมฺเม ปริปูเรติ, จตฺตาโร ธมฺมา ภาวิตา พหุลีกตา สตฺต ธมฺเม ปริปูเรนฺติ, สตฺต ธมฺมา ภาวิตา พหุลีกตา เทฺว ธมฺเม ปริปูเรนฺตี”ติ. ภควํมูลกา โน ภนฺเต ธมฺมา ฯเปฯ. อตฺถานนฺท เอกธมฺโม ภาวิโต พหุลีกโต จตฺตาโร ธมฺเม ปริปูเรติ, จตฺตาโร ธมฺมา ภาวิตา พหุลีกตา สตฺต ธมฺเม ปริปูเรนฺติ, สตฺต ธมฺมา ภาวิตา พหุลีกตา เทฺว ธมฺเม ปริปูเรนฺติ.}

\prose{กตโม จานนฺท เอกธมฺโม ภาวิโต พหุลีกโต จตฺตาโร ธมฺเม ปริปูเรติ, จตฺตาโร ธมฺมา ภาวิตา พหุลีกตา สตฺต ธมฺเม ปริปูเรนฺติ, สตฺต ธมฺมา ภาวิตา พหุลีกตา เทฺว ธมฺเม ปริปูเรนฺติ. อานาปาน\-สฺสติ\-สมาธิ อานนฺท เอกธมฺโม ภาวิโต พหุลีกโต จตฺตาโร สติปฏฺฐาเน ปริปูเรติ, จตฺตาโร สติปฏฺฐานา ภาวิตา พหุลีกตา สตฺต โพชฺฌงฺเค ปริปูเรนฺติ, สตฺต โพชฺฌงฺคา ภาวิตา พหุลีกตา วิชฺชาวิมุตฺติํ ปริปูเรนฺตีติ. กถํ ภาวิโต จานนฺท อานาปาน\-สฺสติ\-สมาธิ กถํ พหุลีกโต จตฺตาโร สติปฏฺฐาเน ปริปูเรติ, อิธานนฺท ภิกฺขุ อรญฺญคโต วา ฯเปฯ. เอวํ ภาวิตา โข อานนฺท สตฺต โพชฺฌงฺคา เอวํ พหุลีกตา วิชฺชาวิมุตฺติํ ปริปูเรนฺตีติ.  จตุตฺถํ.}

\csromanheader{2}{5. ปฐม\-ภิกฺขุ\-สุตฺต}
\proseitem{991}{อถ โข สมฺพหุลา ภิกฺขู เยน ภควา เตนุ\-ปสงฺกมิํสุ, อุปสงฺกมิตฺวา ภควนฺตํ อภิวาเทตฺวา เอกมนฺตํ นิสีทิํสุ, เอกมนฺตํ นิสินฺนา โข เต ภิกฺขู ภควนฺตํ เอตทโวจุํ “อตฺถิ นุ โข ภนฺเต เอกธมฺโม ภาวิโต พหุลีกโต จตฺตาโร ธมฺเม ปริปูเรติ, จตฺตาโร ธมฺมา ภาวิตา พหุลีกตา สตฺต ธมฺเม ปริปูเรนฺติ, สตฺต ธมฺมา \csromanfolio{292}ภาวิตา พหุลีกตา เทฺว ธมฺเม ปริปูเรนฺตี”ติ. อตฺถิ โข ภิกฺขเว เอกธมฺโม ภาวิโต พหุลีกโต จตฺตาโร ธมฺเม ปริปูเรติ, จตฺตาโร ธมฺมา ภาวิตา พหุลีกตา สตฺต ธมฺเม ปริปูเรนฺติ, สตฺต ธมฺมา ภาวิตา พหุลีกตา เทฺว ธมฺเม ปริปูเรนฺตีติ.}

\prose{กตโม ปน ภนฺเต เอกธมฺโม ภาวิโต พหุลีกโต จตฺตาโร ธมฺเม ปริปูเรติ, จตฺตาโร ธมฺมา ภาวิตา พหุลีกตา สตฺต ธมฺเม ปริปูเรนฺติ, สตฺต ธมฺมา ภาวิตา พหุลีกตา เทฺว ธมฺเม ปริปูเรนฺตีติ. อานาปาน\-สฺสติ\-สมาธิ โข ภิกฺขเว เอกธมฺโม ภาวิโต พหุลีกโต จตฺตาโร สติปฏฺฐาเน ปริปูเรติ, จตฺตาโร สติปฏฺฐานา ภาวิตา พหุลีกตา สตฺต โพชฺฌงฺเค ปริปูเรนฺติ, สตฺต โพชฺฌงฺคา ภาวิตา พหุลีกตา วิชฺชาวิมุตฺติํ ปริปูเรนฺตีติ.}

\prose{กถํ ภาวิโต จ ภิกฺขเว อานาปาน\-สฺสติ\-สมาธิ กถํ พหุลีกโต จตฺตาโร สติปฏฺฐาเน ปริปูเรติ. อิธ ภิกฺขเว ภิกฺขุ อรญฺญคโต วา ฯเปฯ. เอวํ ภาวิตา โข ภิกฺขเว สตฺต โพชฺฌงฺคา เอวํ พหุลีกตา วิชฺชาวิมุตฺติํ ปริปูเรนฺตีติ.  ปญฺจมํ.}

\csromanheader{2}{6. ทุติย\-ภิกฺขุ\-สุตฺต}
\proseitem{992}{อถ โข สมฺพหุลา ภิกฺขู เยน ภควา เตนุ\-ปสงฺกมิํสุ, อุปสงฺกมิตฺวา ภควนฺตํ อภิวาเทตฺวา เอกมนฺตํ นิสีทิํสุ, เอกมนฺตํ นิสินฺเน โข เต ภิกฺขู ภควา เอตทโวจ “อตฺถิ นุ โข ภิกฺขเว เอกธมฺโม ภาวิโต พหุลีกโต จตฺตาโร ธมฺเม ปริปูเรติ, จตฺตาโร ธมฺมา ภาวิตา พหุลีกตา สตฺต ธมฺเม ปริปูเรนฺติ, สตฺต ธมฺมา ภาวิตา พหุลีกตา เทฺว ธมฺเม ปริปูเรนฺตี”ติ. ภควํมูลกา โน ภนฺเต ธมฺมา ฯเปฯ ภควโต สุตฺวา ภิกฺขู ธาเรสฺสนฺตีติ. อตฺถิ ภิกฺขเว เอกธมฺโม ภาวิโต พหุลีกโต จตฺตาโร ธมฺเม ปริปูเรติ, จตฺตาโร ธมฺมา ภาวิตา พหุลีกตา สตฺต ธมฺเม ปริปูเรนฺติ, สตฺต ธมฺมา ภาวิตา พหุลีกตา เทฺว ธมฺเม ปริปูเรนฺติ.}

\prose{กตโม จ ภิกฺขเว เอกธมฺโม ภาวิโต พหุลีกโต จตฺตาโร ธมฺเม ปริปูเรติ, จตฺตาโร ธมฺมา ภาวิตา พหุลีกตา สตฺต ธมฺเม ปริปูเรนฺติ, สตฺต ธมฺมา ภาวิตา พหุลีกตา เทฺว ธมฺเม ปริปูเรนฺติ. อานาปาน\-สฺสติ\-สมาธิ ภิกฺขเว เอกธมฺโม ภาวิโต พหุลีกโต}

\prose{\csromanfolio{293}จตฺตาโร สติปฏฺฐาเน ปริปูเรติ, จตฺตาโร สติปฏฺฐานา ภาวิตา พหุลีกตา สตฺต โพชฺฌงฺเค ปริปูเรนฺติ, สตฺต โพชฺฌงฺคา ภาวิตา พหุลีกตา วิชฺชาวิมุตฺติํ ปริปูเรนฺตีติ.}

\prose{กถํ ภาวิโต จ ภิกฺขเว อานาปาน\-สฺสติ\-สมาธิ กถํ พหุลีกโต จตฺตาโร สติปฏฺฐาเน ปริปูเรติ. อิธ ภิกฺขเว ภิกฺขุ อรญฺญคโต วา รุกฺขมูลคโต วา สุญฺญาคารคโต วา นิสีทติ, ปลฺลงฺกํ อาภุชิตฺวา อุชุํ กายํ ปณิธาย ปริมุขํ สติํ อุปฏฺฐเปตฺวา, โส สโตว อสฺสสติ, สโตว ปสฺสสติ ฯเปฯ “ปฏินิสฺสคฺคา\-นุปสฺสี อสฺสสิสฺสามี”ติ สิกฺขติ, “ปฏินิสฺสคฺคา\-นุปสฺสี ปสฺสสิสฺสามี”ติ สิกฺขติ.}

\prose{ยสฺมิํ สมเย ภิกฺขเว ภิกฺขุ ทีฆํ วา อสฺสสนฺโต “ทีฆํ อสฺสสามี”ติ ปชานาติ, ทีฆํ วา ปสฺสสนฺโต “ทีฆํ ปสฺสสามี”ติ ปชานาติ. รสฺสํ วา อสฺสสนฺโต “รสฺสํ อสฺสสามี”ติ ปชานาติ ฯเปฯ สพฺพกาย\-ปฺปฏิสํเวที ฯเปฯ “ปสฺสมฺภยํ กายสงฺขารํ อสฺสสิสฺสามี”ติ สิกฺขติ, “ปสฺสมฺภยํ กายสงฺขารํ ปสฺสสิสฺสามี”ติ สิกฺขติ. กาเย กายานุปสฺสี ภิกฺขเว ภิกฺขุ ตสฺมิํ สมเย วิหรติ อาตาปี สมฺปชาโน สติมา วิเนยฺย โลเก อภิชฺฌา\-โทมนสฺสํ. ตํ กิสฺส เหตุ, กายญฺญตราหํ ภิกฺขเว เอตํ วทามิ ยทิทํ อสฺสาสปสฺสาสํ. ตสฺมาติห ภิกฺขเว กาเย กายานุปสฺสี ภิกฺขุ ตสฺมิํ สมเย วิหรติ อาตาปี สมฺปชาโน สติมา วิเนยฺย โลเก อภิชฺฌา\-โทมนสฺสํ. (1)}

\prose{ยสฺมิํ สมเย ภิกฺขเว ภิกฺขุ ปีติปฺปฏิสํเวที ฯเปฯ สุขปฺปฏิสํเวที. จิตฺตสงฺขาร\-ปฺปฏิสํเวที. “ปสฺสมฺภยํ จิตฺตสงฺขารํ อสฺสสิสฺสามี”ติ สิกฺขติ, ปสฺสมฺภยํ จิตฺตสงฺขารํ “ปสฺสสิสฺสามี”ติ สิกฺขติ. เวทนาสุ เวทนานุปสฺสี ภิกฺขเว ภิกฺขุ ตสฺมิํ สมเย วิหรติ อาตาปี สมฺปชาโน สติมา วิเนยฺย โลเก อภิชฺฌา\-โทมนสฺสํ. ตํ กิสฺส เหตุ, เวทนาญฺญตราหํ ภิกฺขเว เอตํ วทามิ ยทิทํ อสฺสาส\-ปสฺสาสานํ สาธุกํ มนสิการํ. ตสฺมาติห ภิกฺขเว เวทนาสุ เวทนานุปสฺสี ภิกฺขุ ตสฺมิํ สมเย วิหรติ อาตาปี สมฺปชาโน สติมา วิเนยฺย โลเก อภิชฺฌา\-โทมนสฺสํ. (2)}

\prose{ยสฺมิํ สมเย ภิกฺขเว ภิกฺขุ จิตฺต\-ปฺปฏิสํเวที ฯเปฯ อภิปฺปโมทยํ จิตฺตํ ฯเปฯ “สมาทหํ จิตฺตํ อสฺสสิสฺสามี”ติ สิกฺขติ, “สมาทหํ จิตฺตํ}

\prose{\csromanfolio{294}ปสฺสสิสฺสามี”ติ สิกฺขติ, “วิโมจยํ จิตฺตํ อสฺสสิสฺสามี”ติ สิกฺขติ, “วิโมจยํ จิตฺตํ ปสฺสสิสฺสามี”ติ สิกฺขติ. จิตฺเต จิตฺตานุปสฺสี ภิกฺขเว ภิกฺขุ ตสฺมิํ สมเย วิหรติ อาตาปี สมฺปชาโน สติมา วิเนยฺย โลเก อภิชฺฌา\-โทมนสฺสํ. ตํ กิสฺส เหตุ, นาหํ ภิกฺขเว มุฏฺฐสฺสติสฺส อสมฺปชานสฺส อานาปาน\-สฺสติ\-สมาธิภาวนํ วทามิ. ตสฺมาติห ภิกฺขเว จิตฺเต จิตฺตานุปสฺสี ภิกฺขุ ตสฺมิํ สมเย วิหรติ อาตาปี สมฺปชาโน สติมา วิเนยฺย โลเก อภิชฺฌา\-โทมนสฺสํ. (3)}

\prose{ยสฺมิํ สมเย ภิกฺขเว ภิกฺขุ อนิจฺจานุปสฺสี ฯเปฯ วิราคานุปสฺสี ฯเปฯ นิโรธานุปสฺสี ฯเปฯ “ปฏินิสฺสคฺคา\-นุปสฺสี อสฺสสิสฺสามี”ติ สิกฺขติ, “ปฏินิสฺสคฺคา\-นุปสฺสี ปสฺสสิสฺสามี”ติ สิกฺขติ. ธมฺเมสุ ธมฺมานุปสฺสี ภิกฺขเว ภิกฺขุ ตสฺมิํ สมเย วิหรติ อาตาปี สมฺปชาโน สติมา วิเนยฺย โลเก อภิชฺฌา\-โทมนสฺสํ. โส ยํ ตํ โหติ อภิชฺฌา\-โทมนสฺสานํ ปหานํ, ตํ ปญฺญาย ทิสฺวา สาธุกํ อชฺฌุเปกฺขิตา โหติ. ตสฺมาติห ภิกฺขเว ธมฺเมสุ ธมฺมานุปสฺสี ภิกฺขุ ตสฺมิํ สมเย วิหรติ อาตาปี สมฺปชาโน สติมา วิเนยฺย โลเก อภิชฺฌา\-โทมนสฺสํ. (4)}

\prose{เอวํ ภาวิโต โข ภิกฺขเว อานาปาน\-สฺสติ\-สมาธิ เอวํ พหุลีกโต จตฺตาโร สติปฏฺฐาเน ปริปูเรติ.}

\prose{กถํ ภาวิตา จ ภิกฺขเว จตฺตาโร สติปฏฺฐานา กถํ พหุลีกตา สตฺต โพชฺฌงฺเค ปริปูเรนฺติ. ยสฺมิํ สมเย ภิกฺขเว ภิกฺขุ กาเย กายานุปสฺสี วิหรติ, อุปฏฺฐิตาสฺส ตสฺมิํ สมเย ภิกฺขุโน สติ โหติ อสมฺมุฏฺฐา. ยสฺมิํ สมเย ภิกฺขเว ภิกฺขุโน อุปฏฺฐิตา สติ อสมฺมุฏฺฐา, สติสมฺโพชฺฌงฺโค ตสฺมิํ สมเย ภิกฺขุโน อารทฺโธ โหติ, สติสมฺโพชฺฌงฺคํ ตสฺมิํ สมเย ภิกฺขุ ภาเวติ, สติสมฺโพชฺฌงฺโค ตสฺมิํ สมเย ภิกฺขุโน ภาวนา\-ปาริปูริํ คจฺฉติ. (1)}

\prose{โส ตถาสโต วิหรนฺโต ตํ ธมฺมํ ปญฺญาย ปวิจินติ ปวิจรติ ปริวีมํสมา\-ปชฺชติ. ยสฺมิํ สมเย ภิกฺขเว ภิกฺขุ ตถาสโต วิหรนฺโต ตํ ธมฺมํ ปญฺญาย ปวิจินติ ปวิจรติ ปริวีมํสมา\-ปชฺชติ, ธมฺมวิจย\-สมฺโพชฺฌงฺโค ตสฺมิํ สมเย ภิกฺขุโน อารทฺโธ โหติ, ธมฺมวิจย\-สมฺโพชฺฌงฺคํ ตสฺมิํ สมเย ภิกฺขุ ภาเวติ, ธมฺมวิจย\-สมฺโพชฺฌงฺโค ตสฺมิํ สมเย ภิกฺขุโน ภาวนา\-ปาริปูริํ คจฺฉติ. (2)}

\prose{\csromanfolio{295}ตสฺส ตํ ธมฺมํ ปญฺญาย ปวิจินโต ปวิจรโต ปริวีมํสมาปชฺชโต อารทฺธํ โหติ วีริยํ อสลฺลีนํ. ยสฺมิํ สมเย ภิกฺขเว ภิกฺขุโน ตํ ธมฺมํ ปญฺญาย ปวิจินโต ปวิจรโต ปริวีมํสมาปชฺชโต อารทฺธํ โหติ วีริยํ อสลฺลีนํ, วีริย\-สมฺโพชฺฌงฺโค ตสฺมิํ สมเย ภิกฺขุโน อารทฺโธ โหติ, วีริย\-สมฺโพชฺฌงฺคํ ตสฺมิํ สมเย ภิกฺขุ ภาเวติ, วีริย\-สมฺโพชฺฌงฺโค ตสฺมิํ สมเย ภิกฺขุโน ภาวนา\-ปาริปูริํ คจฺฉติ. (3)}

\prose{อารทฺธ\-วีริยสฺส อุปฺปชฺชติ ปีติ นิรามิสา. ยสฺมิํ สมเย ภิกฺขเว ภิกฺขุโน อารทฺธ\-วีริยสฺส อุปฺปชฺชติ ปีติ นิรามิสา, ปีติสมฺโพชฺฌงฺโค ตสฺมิํ สมเย ภิกฺขุโน อารทฺโธ โหติ, ปีติสมฺโพชฺฌงฺคํ ตสฺมิํ สมเย ภิกฺขุ ภาเวติ, ปีติสมฺโพชฺฌงฺโค ตสฺมิํ สมเย ภิกฺขุโน ภาวนา\-ปาริปูริํ คจฺฉติ. (4)}

\prose{ปีติมนสฺส กาโยปิ ปสฺสมฺภติ, จิตฺตมฺปิ ปสฺสมฺภติ. ยสฺมิํ สมเย ภิกฺขเว ภิกฺขุโน ปีติมนสฺส กาโยปิ ปสฺสมฺภติ, จิตฺตมฺปิ ปสฺสมฺภติ, ปสฺสทฺธิ\-สมฺโพชฺฌงฺโค ตสฺมิํ สมเย ภิกฺขุโน อารทฺโธ โหติ, ปสฺสทฺธิ\-สมฺโพชฺฌงฺคํ ตสฺมิํ สมเย ภิกฺขุ ภาเวติ, ปสฺสทฺธิ\-สมฺโพชฺฌงฺโค ตสฺมิํ สมเย ภิกฺขุโน ภาวนา\-ปาริปูริํ คจฺฉติ. (5)}

\prose{ปสฺสทฺธ\-กายสฺส สุขิโน จิตฺตํ สมาธิยติ. ยสฺมิํ สมเย ภิกฺขเว ภิกฺขุโน ปสฺสทฺธ\-กายสฺส สุขิโน จิตฺตํ สมาธิยติ, สมาธิ\-สมฺโพชฺฌงฺโค ตสฺมิํ สมเย ภิกฺขุโน อารทฺโธ โหติ, สมาธิ\-สมฺโพชฺฌงฺคํ ตสฺมิํ สมเย ภิกฺขุ ภาเวติ, สมาธิ\-สมฺโพชฺฌงฺโค ตสฺมิํ สมเย ภิกฺขุโน ภาวนา\-ปาริปูริํ คจฺฉติ. (6)}

\prose{โส ตถาสมาหิตํ จิตฺตํ สาธุกํ อชฺฌุเปกฺขิตา โหติ. ยสฺมิํ สมเย ภิกฺขเว ภิกฺขุ ตถาสมาหิตํ จิตฺตํ สาธุกํ อชฺฌุเปกฺขิตา โหติ, อุเปกฺขา\-สมฺโพชฺฌงฺโค ตสฺมิํ สมเย ภิกฺขุโน อารทฺโธ โหติ, อุเปกฺขา\-สมฺโพชฺฌงฺคํ ตสฺมิํ สมเย ภิกฺขุ ภาเวติ, อุเปกฺขา\-สมฺโพชฺฌงฺโค ตสฺมิํ สมเย ภิกฺขุโน ภาวนา\-ปาริปูริํ คจฺฉติ. (7)}

\prose{ยสฺมิํ สมเย ภิกฺขเว ภิกฺขุ เวทนาสุ ฯเปฯ. จิตฺเต ฯเปฯ. ธมฺเมสุ ธมฺมานุปสฺสี วิหรติ, อุปฏฺฐิตาสฺส ตสฺมิํ สมเย ภิกฺขุโน สติ}

\csromanmatikaanchor{mtk.220}
\csromantocmarkref{subparagraph}{7-10. สํโยชนปฺปหานสุตฺตาทีนิ}{mtk.220}
\bookmark[dest={mtk.220},level=5]{7-10. สํโยชนปฺปหานสุตฺตาทีนิ}
\prose{\csromanfolio{296}โหติ อสมฺมุฏฺฐา. ยสฺมิํ สมเย ภิกฺขเว ภิกฺขุโน อุปฏฺฐิตา สติ โหติ อสมฺมุฏฺฐา, สติสมฺโพชฺฌงฺโค ตสฺมิํ สมเย ภิกฺขุโน อารทฺโธ โหติ, สติสมฺโพชฺฌงฺคํ ตสฺมิํ สมเย ภิกฺขุ ภาเวติ, สติสมฺโพชฺฌงฺโค ตสฺมิํ สมเย ภิกฺขุโน ภาวนา\-ปาริปูริํ คจฺฉติ ฯเปฯ.}

\prose{โส ตถาสมาหิตํ จิตฺตํ สาธุกํ อชฺฌุเปกฺขิตา โหติ. ยสฺมิํ สมเย ภิกฺขเว ภิกฺขุ ตถาสมาหิตํ จิตฺตํ สาธุกํ อชฺฌุเปกฺขิตา โหติ, อุเปกฺขา\-สมฺโพชฺฌงฺโค ตสฺมิํ สมเย ภิกฺขุโน อารทฺโธ โหติ, อุเปกฺขา\-สมฺโพชฺฌงฺคํ ตสฺมิํ สมเย ภิกฺขุ ภาเวติ, อุเปกฺขา\-สมฺโพชฺฌงฺโค ตสฺมิํ สมเย ภิกฺขุโน ภาวนา\-ปาริปูริํ คจฺฉติ. เอวํ ภาวิตา โข ภิกฺขเว จตฺตาโร สติปฏฺฐานา เอวํ พหุลีกตา สตฺต โพชฺฌงฺเค ปริปูเรนฺติ.}

\prose{กถํ ภาวิตา จ ภิกฺขเว สตฺต โพชฺฌงฺคา กถํ พหุลีกตา วิชฺชาวิมุตฺติํ ปริปูเรนฺติ, อิธ ภิกฺขเว ภิกฺขุ สติสมฺโพชฺฌงฺคํ ภาเวติ วิเวกนิสฺสิตํ วิราคนิสฺสิตํ นิโรธ\-นิสฺสิตํ โวสฺสคฺค\-ปริณามิํ. ธมฺมวิจย\-สมฺโพชฺฌงฺคํ ภาเวติ วิเวกนิสฺสิตํ วิราคนิสฺสิตํ นิโรธ\-นิสฺสิตํ โวสฺสคฺค\-ปริณามิํ ฯเปฯ. อุเปกฺขา\-สมฺโพชฺฌงฺคํ ภาเวติ วิเวกนิสฺสิตํ วิราคนิสฺสิตํ นิโรธ\-นิสฺสิตํ โวสฺสคฺค\-ปริณามิํ. เอวํ ภาวิตา โข ภิกฺขเว สตฺต โพชฺฌงฺคา เอวํ พหุลีกตา วิชฺชาวิมุตฺติํ ปริปูเรนฺตีติ.  ฉฏฺฐํ.}

\csromanheader{2}{7. สํโยชน\-ปฺปหาน\-สุตฺต}
\proseitem{993}{อานาปาน\-สฺสติ\-สมาธิ ภิกฺขเว ภาวิโต พหุลีกโต สํโยชน\-ปฺปหานาย สํวตฺตติ ฯเปฯ.  สตฺตมํ.}

\csromanheader{2}{8. อนุสยสมุคฺฆาต\-สุตฺต}
\vspace{\tipitakaparskip}
\csromancenter{อนุสย\-สมุคฺฆาตาย สํวตฺตติ.  อฏฺฐมํ.}
\csromanheader{2}{9. อทฺธานปริญฺญา\-สุตฺต}
\vspace{\tipitakaparskip}
\csromancenter{อทฺธาน\-ปริญฺญาย สํวตฺตติ.  นวมํ.}
\csromanheader{2}{10. อาสวกฺขย\-สุตฺต}
\proseitem{996}{\csromanfolio{297}อาสวานํ ขยาย สํวตฺตติ. กถํ ภาวิโต จ ภิกฺขเว อานาปาน\-สฺสติ\-สมาธิ กถํ พหุลีกโต สํโยชน\-ปฺปหานาย สํวตฺตติ. อนุสย\-สมุคฺฆาตาย สํวตฺตติ. อทฺธาน\-ปริญฺญาย สํวตฺตติ. อาสวานํ ขยาย สํวตฺตติ, อิธ ภิกฺขเว ภิกฺขุ อรญฺญคโต วา รุกฺขมูลคโต วา ฯเปฯ “ปฏินิสฺสคฺคา\-นุปสฺสี อสฺสสิสฺสามี”ติ สิกฺขติ, “ปฏินิสฺสคฺคา\-นุปสฺสี ปสฺสสิสฺสามี”ติ สิกฺขติ. เอวํ ภาวิโต โข ภิกฺขเว อานาปาน\-สฺสติ\-สมาธิ เอวํ พหุลีกโต สํโยชน\-ปฺปหานาย สํวตฺตติ ฯเปฯ. อนุสย\-สมุคฺฆาตาย สํวตฺตติ ฯเปฯ. อทฺธาน\-ปริญฺญาย สํวตฺตติ ฯเปฯ. สาสวานํ ขยาย สํวตฺตตีติ.  ทสมํ. ทุติโย วคฺโค.}

\vspace{\tipitakaparskip}
\csromancenter{\textbf{ตสฺสุทฺทานํ}}
\setlength{\csromangathapagewidth}{0pt}
\setlength{\csromangathawakwidth}{0pt}
\csromangathameasuregroup{อิจฺฉานงฺคลํ กงฺเขยฺยํ, \\ ภิกฺขู สํโยชนานุสยา,}{\csromangathabat{อิจฺฉานงฺคลํ กงฺเขยฺยํ,}{อานนฺท อปเร ทุเว.} \\ \csromangathabat{ภิกฺขู สํโยชนานุสยา,}{อทฺธานํ อาสวกฺขยนฺติ.} \\ อานาปานสํยุตฺตํ ทสมํ.}
\csromangathameasurewak{อานาปานสํยุตฺตํ ทสมํ.}
\csromangathasetleft{อิจฺฉานงฺคลํ กงฺเขยฺยํ, \\ ภิกฺขู สํโยชนานุสยา,}
\csromangathagroup{\csromangathastanza{\csromangathabat{อิจฺฉานงฺคลํ กงฺเขยฺยํ,}{อานนฺท อปเร ทุเว.} \\ \csromangathabat{ภิกฺขู สํโยชนานุสยา,}{อทฺธานํ อาสวกฺขยนฺติ.}}\csromangathastanza{\csromangathawak{อานาปานสํยุตฺตํ ทสมํ.}}}

\csromanmatikaanchor{mtk.221}
\csromanmatikaanchor{mtk.222}
\csromantocmarknopage{subparagraph}{11. โสตาปตฺติสํยุตฺต}{mtk.221}
\bookmark[dest={mtk.221},level=5]{11. โสตาปตฺติสํยุตฺต}
\csromantocmarknopage{subparagraph}{1. เวฬุทฺวารวคฺค}{mtk.222}
\bookmark[dest={mtk.222},level=5]{1. เวฬุทฺวารวคฺค}
\csromanheader{6}{1. \textbf{เวฬุทฺวารวคฺค}}
\csromanmatikaanchor{mtk.223}
\csromantocmarkref{subparagraph}{1. จกฺกวตฺติราชสุตฺต}{mtk.223}
\bookmark[dest={mtk.223},level=5]{1. จกฺกวตฺติราชสุตฺต}
\csromanheader{6}{1. จกฺกวตฺติราช\-สุตฺต}
\proseitem{997}{\csromanfolio{298}สาวตฺถิ\-นิทานํ. ตตฺร โข ภควา ฯเปฯ เอตทโวจ “กิญฺจาปิ ภิกฺขเว ราชา จกฺกวตฺตี\csromansharedfootnote{2b364623350c6086}{จกฺกวตฺติ (สฺยา, กํ, อิ, ก)} จตุนฺนํ ทีปานํ อิสฺสริยาธิปจฺจํ รชฺชํ กาเรตฺวา กายสฺส เภทา ปรํ มรณา สุคติํ สคฺคํ โลกํ อุปปชฺชติ เทวานํ ตาวติํสานํ สหพฺยตํ, โส ตตฺถ นนฺทเน วเน อจฺฉรา\-สงฺฆ\-ปริวุโต ทิพฺเพหิ จ ปญฺจหิ กามคุเณหิ สมปฺปิโต สมงฺคีภูโต ปริจาเรติ โส จตูหิ ธมฺเมหิ อสมนฺนาคโต, อถ โข โส อปริมุตฺโตว\csromansharedfootnote{3a1dabebc645e0f9}{อปริมุตฺโต จ (สฺยา, กํ, ก)} นิรยา, อปริมุตฺโต ติรจฺฉาน\-โยนิยา, อปริมุตฺโต เปตฺติวิสยา, อปริมุตฺโต อปาย\-ทุคฺคติ\-วินิปาตา. กิญฺจาปิ ภิกฺขเว อริยสาวโก ปิณฺฑิยาโลเปน ยาเปติ, นนฺตกานิ จ ธาเรติ. โส จตูหิ ธมฺเมหิ สมนฺนาคโต, อถ โข โส ปริมุตฺโต\csromansharedfootnote{dd3dca71dbec9412}{ปริมุตฺโต จ (สฺยา, กํ, ก)} นิรยา, ปริมุตฺโต ติรจฺฉาน\-โยนิยา, ปริมุตฺโต เปตฺติวิสยา, ปริมุตฺโต อปาย\-ทุคฺคติ\-วินิปาตา.}

\prose{กตเมหิ จตูหิ, อิธ ภิกฺขเว อริยสาวโก พุทฺเธ อเวจฺจ\-ปฺปสาเทน สมนฺนาคโต โหติ “อิติปิ โส ภควา อรหํ สมฺมาสมฺพุทฺโธ วิชฺชา\-จรณ\-สมฺปนฺโน สุคโต โลกวิทู อนุตฺตโร ปุริส\-ทมฺม\-สารถิ สตฺถา เทวมนุสฺสานํ พุทฺโธ ภควา”ติ. ธมฺเม อเวจฺจ\-ปฺปสาเทน สมนฺนาคโต โหติ “สฺวากฺขาโต ภควตา ธมฺโม สนฺทิฏฺฐิโก อกาลิโก เอหิปสฺสิโก โอปเนยฺยิโก ปจฺจตฺตํ เวทิตพฺโพ วิญฺญูหี”ติ. สํเฆ อเวจฺจ\-ปฺปสาเทน สมนฺนาคโต โหติ “สุปฺปฏิปนฺโน ภควโต สาวกสํโฆ อุชุปฺปฏิปนฺโน ภควโต สาวกสํโฆ ญายปฺปฏิปนฺโน ภควโต สาวกสํโฆ สามีจิ\-ปฺปฏิปนฺโน ภควโต สาวกสํโฆ, ยทิทํ จตฺตาริ ปุริสยุคานิ อฏฺฐ ปุริสปุคฺคลา, เอส ภควโต สาวกสํโฆ อาหุเนยฺโย ปาหุเนยฺโย ทกฺขิเณยฺโย อญฺชลิกรณีโย อนุตฺตรํ ปุญฺญกฺเขตฺตํ โลกสฺสา”ติ. อริยกนฺเตหิ}

\prose{\csromanfolio{299}สีเลหิ สมนฺนาคโต โหติ อขณฺเฑหิ อจฺฉิทฺเทหิ อสพเลหิ อกมฺมาเสหิ ภุชิสฺเสหิ วิญฺญุ\-ปฺปสตฺเถหิ อปรามฏฺเฐหิ สมาธิ\-สํวตฺตนิเกหิ. อิเมหิ จตูหิ ธมฺเมหิ สมนฺนาคโต โหติ, โย จ ภิกฺขเว จตุนฺนํ ทีปานํ ปฏิลาโภ, โย จตุนฺนํ ธมฺมานํ ปฏิลาโภ, จตุนฺนํ ทีปานํ ปฏิลาโภ จตุนฺนํ ธมฺมานํ ปฏิลาภสฺส กลํ นาคฺฆติ โสฬสินฺติ.  ปฐมํ.}

\csromanmatikaanchor{mtk.224}
\csromantocmarkref{subparagraph}{2. พฺรหฺมจริโยคธสุตฺต}{mtk.224}
\bookmark[dest={mtk.224},level=5]{2. พฺรหฺมจริโยคธสุตฺต}
\csromanheader{6}{2. พฺรหฺมจริโยคธ\-สุตฺต}
\proseitem{998}{จตูหิ ภิกฺขเว ธมฺเมหิ สมนฺนาคโต อริยสาวโก โสตาปนฺโน โหติ อวินิปาต\-ธมฺโม นิยโต สมฺโพธิ\-ปรายโณ.}

\prose{กตเมหิ จตูหิ, อิธ ภิกฺขเว อริยสาวโก พุทฺเธ อเวจฺจ\-ปฺปสาเทน สมนฺนาคโต โหติ “อิติปิ โส ภควา อรหํ สมฺมาสมฺพุทฺโธ วิชฺชา\-จรณ\-สมฺปนฺโน สุคโต โลกวิทู อนุตฺตโร ปุริส\-ทมฺม\-สารถิ สตฺถา เทวมนุสฺสานํ พุทฺโธ ภควา”ติ. ธมฺเม ฯเปฯ. สํเฆ ฯเปฯ. อริยกนฺเตหิ สีเลหิ สมนฺนาคโต โหติ อขณฺเฑหิ ฯเปฯ สมาธิ\-สํวตฺตนิเกหิ. อิเมหิ โข ภิกฺขเว จตูหิ ธมฺเมหิ สมนฺนาคโต อริยสาวโก โสตาปนฺโน โหติ อวินิปาต\-ธมฺโม นิยโต สมฺโพธิ\-ปรายโณติ.}

\prose{อิทมโวจ ภควา, อิทํ วตฺวาน\csromansharedfootnote{f786312130b697a4}{วตฺวา (สี, อิ) เอวมีทิเสสุ ฐาเนสุ.} สุคโต, อถาปรํ เอตทโวจ สตฺถา\csromandash{}}

\setlength{\csromangathapagewidth}{0pt}
\setlength{\csromangathawakwidth}{0pt}
\csromangathameasuregroup{“เยสํ สทฺธา จ สีลญฺจ,}{\csromangathabat{“เยสํ สทฺธา จ สีลญฺจ,}{ปสาโท ธมฺมทสฺสนํ.}}
\csromangathasetleft{“เยสํ สทฺธา จ สีลญฺจ,}
\csromangathagroup{\csromangathastanza{\csromangathabat{“เยสํ สทฺธา จ สีลญฺจ,}{ปสาโท ธมฺมทสฺสนํ.}}}

\prose{เต เว กาเลน ปจฺเจนฺติ, พฺรหฺมจริโยคธํ สุขนฺ”ติ. ทุติยํ.}

\csromanmatikaanchor{mtk.225}
\csromantocmarkref{subparagraph}{3. ทีฆาวุอุปาสกสุตฺต}{mtk.225}
\bookmark[dest={mtk.225},level=5]{3. ทีฆาวุอุปาสกสุตฺต}
\csromanheader{6}{3. ทีฆาวุ\-อุปาสก\-สุตฺต}
\proseitem{999}{เอกํ สมยํ ภควา ราชคเห วิหรติ เวฬุวเน กลนฺทก\-นิวาเป. เตน โข ปน สมเยน ทีฆาวุ อุปาสโก อาพาธิโก โหติ ทุกฺขิโต พาฬฺหคิลาโน. อถ โข ทีฆาวุ อุปาสโก ปิตรํ โชติกํ คหปติํ อามนฺเตสิ “เอหิ ตฺวํ คหปติ เยน ภควา เตนุปสงฺกม, อุปสงฺกมิตฺวา มม วจเนน ภควโต ปาเท สิรสา วนฺท ‘ทีฆาวุ \csromanfolio{300}ภนฺเต อุปาสโก อาพาธิโก โหติ ทุกฺขิโต พาฬฺหคิลาโน, โส ภควโต ปาเท สิรสา วนฺทตี’ติ. เอวญฺจ วเทหิ ‘สาธุ กิร ภนฺเต ภควา เยน ทีฆาวุสฺส อุปาสกสฺส นิเวสนํ เตนุ\-ปสงฺกมตุ อนุกมฺปํ อุปาทายา’ติ”. “เอวํ ตาตา”ติ โข โชติโก คหปติ ทีฆาวุสฺส อุปาสกสฺส ปฏิสฺสุตฺวา เยน ภควา เตนุปสงฺกมิ, อุปสงฺกมิตฺวา ภควนฺตํ อภิวาเทตฺวา เอกมนฺตํ นิสีทิ, เอกมนฺตํ นิสินฺโน โข โชติโก คหปติ ภควนฺตํ เอตทโวจ “ทีฆาวุ ภนฺเต อุปาสโก อาพาธิโก โหติ ทุกฺขิโต พาฬฺหคิลาโน, โส ภควโต ปาเท สิรสา วนฺทติ, เอวญฺจ วเทติ ‘สาธุ กิร ภนฺเต ภควา เยน ทีฆาวุสฺส อุปาสกสฺส นิเวสนํ เตนุ\-ปสงฺกมตุ อนุกมฺปํ อุปาทายา’ติ”. อธิวาเสสิ ภควา ตุณฺหีภาเวน.}

\prose{อถ โข ภควา นิวาเสตฺวา ปตฺต\-จีวรมา\-ทาย เยน ทีฆาวุสฺส อุปาสกสฺส นิเวสนํ เตนุปสงฺกมิ, อุปสงฺกมิตฺวา ปญฺญตฺเต อาสเน นิสีทิ, นิสชฺช โข ภควา ทีฆาวุํ อุปาสกํ เอตทโวจ “กจฺจิ เต ทีฆาวุ ขมนียํ กจฺจิ ยาปนียํ, กจฺจิ ทุกฺขา เวทนา ปฏิกฺกมนฺติ โน อภิกฺกมนฺติ, ปฏิกฺกโมสานํ ปญฺญายติ โน อภิกฺกโม”ติ. น เม ภนฺเต ขมนียํ น ยาปนียํ, พาฬฺหา เม ทุกฺขา เวทนา อภิกฺกมนฺติ โน ปฏิกฺกมนฺติ, อภิกฺกโมสานํ ปญฺญายติ โน ปฏิกฺกโมติ. ตสฺมาติห เต ทีฆาวุ เอวํ สิกฺขิตพฺพํ “พุทฺเธ อเวจฺจ\-ปฺปสาเทน สมนฺนาคโต ภวิสฺสามิ, “อิติปิ โส ภควา อรหํ สมฺมาสมฺพุทฺโธ วิชฺชา\-จรณ\-สมฺปนฺโน สุคโต โลกวิทู อนุตฺตโร ปุริส\-ทมฺม\-สารถิ สตฺถา เทวมนุสฺสานํ พุทฺโธ ภควา”ติ. ธมฺเม ฯเปฯ. สํเฆ ฯเปฯ. อริยกนฺเตหิ สีเลหิ สมนฺนาคโต ภวิสฺสามิ อขณฺเฑหิ ฯเปฯ สมาธิ\-สํวตฺตนิเกหิ”. เอวํ หิ เต ทีฆาวุ สิกฺขิตพฺพนฺติ.}

\prose{ยานิมานิ ภนฺเต ภควตา จตฺตาริ โสตาปตฺติยงฺคานิ เทสิตานิ, สํวิชฺชนฺเตเต ธมฺมา มยิ, อหญฺจ เตสุ ธมฺเมสุ สนฺทิสฺสามิ. อหํ หิ ภนฺเต พุทฺเธ อเวจฺจ\-ปฺปสาเทน สมนฺนาคโต “อิติปิ โส ภควา ฯเปฯ สตฺถา เทวมนุสฺสานํ พุทฺโธ ภควา”ติ. ธมฺเม ฯเปฯ. สํเฆ ฯเปฯ. อริยกนฺเตหิ สีเลหิ สมนฺนาคโต อขณฺเฑหิ ฯเปฯ สมาธิสํวตฺตนิเกหีติ. ตสฺมาติห ตฺวํ ทีฆาวุ อิเมสุ จตูสุ โสตาปตฺติยงฺเคสุ ปติฏฺฐาย ฉ วิชฺชาภาคิเย ธมฺเม อุตฺตริ ภาเวยฺยาสิ “อิธ ตฺวํ ทีฆาวุ สพฺพ\-สงฺขาเรสุ}

\csromanmatikaanchor{mtk.226}
\csromantocmarkref{subparagraph}{4-5. สาริปุตฺตสุตฺตานิ}{mtk.226}
\bookmark[dest={mtk.226},level=5]{4-5. สาริปุตฺตสุตฺตานิ}
\prose{\csromanfolio{301}อนิจฺจานุปสฺสี วิหราหิ อนิจฺเจ ทุกฺขสญฺญี ทุกฺเข อนตฺตสญฺญี ปหานสญฺญี วิราคสญฺญี นิโรธสญฺญี”ติ. เอวํ หิ เต ทีฆาวุ สิกฺขิตพฺพนฺติ.}

\prose{เยเม ภนฺเต ภควตา ฉ วิชฺชาภาคิยา ธมฺมา เทสิตา, สํวิชฺชนฺเตเต ธมฺมา มยิ, อหญฺจ เตสุ ธมฺเมสุ สนฺทิสฺสามิ. อหํ หิ เภนฺเต สพฺพ\-สงฺขาเรสุ อนิจฺจานุปสฺสี วิหรามิ อนิจฺเจ ทุกฺขสญฺญิ ทุกฺเข อนตฺตสญฺญี ปหานสญฺญี วิราคสญฺญี นิโรธสญฺญี. อปิ จ เม ภนฺเต เอวํ โหติ “มา เหวายํ โชติโก คหปติ มมจฺจเยน วิฆาตํ อาปชฺชีติ\csromansharedfootnote{988f894a07091107}{อาปชฺชติ (ก)}, มา ตฺวํ ตาต ทีฆาวุ เอวํ มนสากาสิ, อิงฺฆ ตฺวํ ตาต ทีฆาวุ ยเทว เต ภควา อาห, ตเทว ตฺวํ สาธุกํ มนสิ กโรหี”ติ.}

\prose{อถ โข ภควา ทีฆาวุํ อุปาสกํ อิมินา โอวาเทน โอวทิตฺวา อุฏฺฐายาสนา ปกฺกามิ. อถ โข ทีฆาวุ อุปาสโก อจิร\-ปกฺกนฺตสฺส ภควโต กาลมกาสิ. อถ โข สมฺพหุลา ภิกฺขู เยน ภควา เตนุ\-ปสงฺกมิํสุ, อุปสงฺกมิตฺวา ภควนฺตํ อภิวาเทตฺวา เอกมนฺตํ นิสีทิํสุ, เอกมนฺตํ นิสินฺนา โข เต ภิกฺขู ภควนฺตํ เอตทโวจุํ “โย โส ภนฺเต ทีฆาวุ นาม อุปาสโก ภควตา สํขิตฺเตน โอวาเทน โอวทิโต, โส กาลงฺกโต, ตสฺส กา คติ, โก อภิสมฺปราโย”ติ. ปณฺฑิโต ภิกฺขเว ทีฆาวุ อุปาสโก ปจฺจปาทิ\csromansharedfootnote{cc7cbf9de28136f2}{อโหสิ สจฺจวาที (สฺยา, กํ, อิ, ก)} ธมฺมสฺสา\-นุธมฺมํ, น จ มํ ธมฺมา\-ธิกรณํ\csromansharedfootnote{d961578b7dcbc7af}{น จ ธมฺมาธิกรณํ (สฺยา, กํ, อิ, ก)} วิเหเสสิ\csromansharedfootnote{b7d575b3513fc175}{วิเหเฐสิ (อิติปิ อญฺญตฺถ)}, ทีฆาวุ ภิกฺขเว อุปาสโก ปญฺจนฺนํ โอรมฺภาคิยานํ สํโยชนานํ ปริกฺขยา โอปปาติโก ตตฺถ ปรินิพฺพายี อนาวตฺติธมฺโม ตสฺมา โลกาติ.  ตติยํ.}

\csromanheader{2}{4. ปฐม\-สาริปุตฺต\-สุตฺต}
\proseitem{1000}{เอกํ สมยํ อายสฺมา จ สาริปุตฺโต อายสฺมา จ อานนฺโท สาวตฺถิยํ วิหรนฺติ เชตวเน อนาถ\-ปิณฺฑิกสฺส อาราเม. อถ โข อายสฺมา อานนฺโท สายนฺหสมยํ ปฏิสลฺลานา วุฏฺฐิโต ฯเปฯ เอกมนฺตํ นิสินฺโน โข อายสฺมา อานนฺโท อายสฺมนฺตํ สาริปุตฺตํ เอตทโวจ “กตินํ นุ โข อาวุโส สาริปุตฺต ธมฺมานํ สมนฺนาคมน\-เหตุ เอวมยํ ปชา ภควตา พฺยากตา ‘โสตาปนฺนา อวินิปาต\-ธมฺมา นิยตา \csromanfolio{302}สมฺโพธิ\-ปรายณา’ติ”. จตุนฺนํ โข อาวุโส ธมฺมานํ สมนฺนาคมน\-เหตุ เอวมยํ ปชา ภควตา พฺยากตา “โสตาปนฺนา อวินิปาต\-ธมฺมา นิยตา สมฺโพธิ\-ปรายณา”.}

\prose{กตเมสํ จตุนฺนํ, อิธาวุโส อริยสาวโก พุทฺเธ อเวจฺจ\-ปฺปสาเทน สมนฺนาคโต โหติ “อิติปิ โส ภควา ฯเปฯ สตฺถา เทวมนุสฺสานํ พุทฺโธ ภควา”ติ. ธมฺเม ฯเปฯ. สํเฆ ฯเปฯ. อริยกนฺเตหิ สีเลหิ สมนฺนาคโต โหติ อขณฺเฑหิ ฯเปฯ สมาธิ\-สํวตฺตนิเกหิ. อิเมสํ โข อาวุโส จตุนฺนํ ธมฺมานํ สมนฺนาคมน\-เหตุ เอวมยํ ปชา ภควตา พฺยากตา “โสตาปนฺนา อวินิปาต\-ธมฺมา นิยตา สมฺโพธิ\-ปรายณา”ติ.  จตุตฺถํ.}

\csromanheader{2}{5. ทุติย\-สาริปุตฺต\-สุตฺต}
\proseitem{1001}{อถ โข อายสฺมา สาริปุตฺโต เยน ภควา เตนุปสงฺกมิ, อุปสงฺกมิตฺวา ภควนฺตํ อภิวาเทตฺวา เอกมนฺตํ นิสีทิ, เอกมนฺตํ นิสินฺนํ โข อายสฺมนฺตํ สาริปุตฺตํ ภควา เอตทโวจ “โสตาปตฺติยงฺคํ โสตาปตฺติยงฺคนฺติ หิทํ สาริปุตฺต วุจฺจติ, กตมํ นุ โข สาริปุตฺต โสตาปตฺติยงฺคนฺ”ติ. สปฺปุริส\-สํเสโว หิ ภนฺเต โสตาปตฺติยงฺคํ, สทฺธมฺม\-สฺสวนํ โสตาปตฺติยงฺคํ, โยนิโส\-มนสิกาโร โสตาปตฺติยงฺคํ, ธมฺมา\-นุธมฺม\-ปฺปฏิปตฺติ โสตาปตฺติยงฺคนฺติ. สาธุ สาธุ สาริปุตฺต, สปฺปุริส\-สํเสโว หิ สาริปุตฺต โสตาปตฺติยงฺคํ, สทฺธมฺม\-สฺสวนํ โสตาปตฺติยงฺคํ, โยนิโส\-มนสิกาโร โสตาปตฺติยงฺคํ, ธมฺมานธมฺมปฺปฏิปตฺติ โสตาปตฺติยงฺคํ.}

\prose{“โสโต โสโต”ติ หิทํ สาริปุตฺต วุจฺจติ, กตโม นุ โข สาริปุตฺต โสโตติ. อยเมว หิ ภนฺเต อริโย อฏฺฐงฺคิโก มคฺโค โสโต. เสยฺยถิทํ, สมฺมาทิฏฺฐิ สมฺมาสงฺกปฺโป สมฺมาวาจา สมฺมากมฺมนฺโต สมฺมาอาชีโว สมฺมาวายาโม สมฺมาสติ สมฺมาสมาธีติ. สาธุ สาธุ สาริปุตฺต, อยเมว หิ สาริปุตฺต อริโย อฏฺฐงฺคิโก มคฺโค โสโต. เสยฺยถิทํ, สมฺมาทิฏฺฐิ ฯเปฯ สมฺมาสมาธิ.}

\prose{“โสตาปนฺโน โสตาปนฺโน”ติ หิทํ สาริปุตฺต วุจฺจติ, กตโม นุ โข สาริปุตฺต โสตาปนฺโนติ. โย หิ ภนฺเต อิมินา อริเยน อฏฺฐงฺคิเกน มคฺคเน สมนฺนาคโต, อยํ วุจฺจติ โสตาปนฺโน, สฺวายํ}

\prose{\csromanfolio{303}อายสฺมา เอวํนาโม เอวํโคตฺโตติ. สาธุ สาธุ สาริปุตฺต, โย หิ สาริปุตฺต อิมินา อริเยน อฏฺฐงฺคิเกน มคฺเคน สมนฺนาคโต, อยํ วุจฺจติ โสตาปนฺโน, สฺวายํ อายสฺมา เอวํนาโม เอวํโคตฺโตติ.  ปญฺจมํ.}

\csromanmatikaanchor{mtk.227}
\csromantocmarkref{subparagraph}{6. ถปติสุตฺต}{mtk.227}
\bookmark[dest={mtk.227},level=5]{6. ถปติสุตฺต}
\csromanheader{6}{6. \textbf{ถปติสุตฺต}}
\proseitem{1002}{สาวตฺถิ\-นิทานํ. เตน โข ปน สมเยน สมฺพหุลา ภิกฺขู ภควโต จีวรกมฺมํ กโรนฺติ “นิฏฺฐิตจีวโร ภควา เตมาสจฺจเยน จาริกํ ปกฺกมิสฺสตี”ติ. เตน โข ปน สมเยน อิสิทตฺตปุราณา ถปตโย สาธุเก ปฏิวสนฺติ เกนจิเทว กรณีเยน. อสฺโสสุํ โข อิสิทตฺตปุราณา ถปตโย “สมฺพหุลา กิร ภิกฺขู ภควโต จีวรกมฺมํ กโรนฺติ “นิฏฺฐิตจีวโร ภควา เตมาสจฺจเยน จาริกํ ปกฺกมิสฺสตี”ติ.}

\prose{อถ โข อิสิทตฺตปุราณา ถปตโย มคฺเค ปุริสํ ฐเปสุํ “ยทา ตฺวํ อมฺโภ ปุริส ปสฺเสยฺยาสิ ภควนฺตํ อาคจฺฉนฺตํ อรหนฺตํ สมฺมา\-สมฺพุทฺธํ, อถ อมฺหากํ อาโรเจยฺยาสี”ติ. ทฺวีหตีหํ ฐิโต โข โส ปุริโส อทฺทส ภควนฺตํ ทูรโตว อาคจฺฉนฺตํ ทิสฺวาน เยน อิสิทตฺตปุราณา ถปตโย เตนุปสงฺกมิ, อุปสงฺกมิตฺวา อิสิทตฺตปุราเณ ถปตโย เอตทโวจ “อยํ โส ภนฺเต ภควา อาคจฺฉติ อรหํ สมฺมาสมฺพุทฺโธ, ยสฺสทานิ กาลํ มญฺญถา”ติ.}

\prose{อถ โข อิสิทตฺตปุราณา ถปตโย เยน ภควา เตนุ\-ปสงฺกมิํสุ, อุปสงฺกมิตฺวา ภควนฺตํ อภิวาเทตฺวา ภควนฺตํ ปิฏฺฐิโต ปิฏฺฐิโต อนุพนฺธิํสุ. อถ โข ภควา มคฺคา โอกฺกมฺม เยน อญฺญตรํ รุกฺขมูลํ เตนุปสงฺกมิ, อุปสงฺกมิตฺวา ปญฺญตฺเต อาสเน นิสีทิ, อิสิทตฺตปุราณา ถปตโย ภควนฺตํ อภิวาเทตฺวา เอกมนฺตํ นิสีทิํสุ, เอกมนฺตํ นิสินฺนา โข เต อิสิทตฺตปุราณา ถปตโย ภควนฺตํ เอตทโวจุํ\csromandash{}}

\prose{ยทา มยํ ภนฺเต ภควนฺตํ สุณาม “สาวตฺถิยา โกสเลสุ จาริกํ ปกฺกมิสฺสตี”ติ, โหติ โน ตสฺมิํ สมเย อนตฺตมนตา โหติ โทมนสฺสํ “ทูเร โน ภควา ภวิสฺสตี”ติ. ยทา ปน มยํ ภนฺเต ภควนฺตํ สุณาม “สาวตฺถิยา โกสเลสุ จาริกํ ปกฺกนฺโต”ติ, โหติ โน ตสฺมิํ สมเย อนตฺตมนตา โหติ โทมนสฺสํ, “ทูเร โน ภควา”ติ.}

\prose{\csromanfolio{304}ยทา ปน มยํ ภนฺเต ภควนฺตํ สุณาม “โกสเลหิ มลฺเลสุ จาริกํ ปกฺกมิสฺสตี”ติ, โหติ โน ตสฺมิํ สมเย อนตฺตมนตา โหติ โทมนสฺสํ “ทูเร โน ภควา ภวิสฺสตี”ติ. ยทา ปน มยํ ภนฺเต ภควนฺตํ สุณาม “โกสเลหิ มลฺเลสุ จาริกํ ปกฺกนฺโต”ติ, โหติ โน ตสฺมิํ สมเย อนตฺตมนตา โหติ โทมนสฺสํ “ทูเร โน ภควา”ติ.}

\prose{ยทา ปน มยํ ภนฺเต ภควนฺตํ สุณาม “มลฺเลหิ วชฺชีสุ จาริกํ ปกฺกมิสฺสตี”ติ, โหติ โน ตสฺมิํ สมเย อนตฺตมนตา โหติ โทมนสฺสํ “ทูเร โน ภควา ภวิสฺสตี”ติ. ยทา ปน มยํ ภนฺเต ภควนฺตํ สุณาม “มลฺเลหิ วชฺชีสุ จาริกํ ปกฺกนฺโต”ติ, โหติ โน ตสฺมิํ สมเย อนตฺตมนตา โหติ โทมนสฺสํ “ทูเร โน ภควา”ติ.}

\prose{ยทา ปน มยํ ภนฺเต ภควนฺตํ สุณาม “วชฺชีหิ กาสีสุ จาริกํ ปกฺกมิสฺสตี”ติ, โหติ โน ตสฺมิํ สมเย อนตฺตมนตา โหติ โทมนสฺสํ “ทูเร โน ภควา ภวิสฺสตี”ติ. ยทา ปน มยํ ภนฺเต ภควนฺตํ สุณาม “วชฺชีหิ กาสีสุ จาริกํ ปกฺกนฺโต”ติ, โหติ โน ตสฺมิํ สมเย อนตฺตมนตา โหติ โทมนสฺสํ “ทูเร โน ภควา”ติ.}

\prose{ยทา ปน มยํ ภนฺเต ภควนฺตํ สุณาม “กาสีหิ มาคเธ จาริกํ ปกฺกมิสฺสตี”ติ, โหติ โน ตสฺมิํ สมเย อนตฺตมนตา โหติ โทมนสฺสํ “ทูเร โน ภควา ภวิสฺสตี”ติ. ยทา ปน มยํ ภนฺเต ภควนฺตํ สุณาม “กาสีหิ มาคเธ จาริกํ ปกฺกนฺโต”ติ, โหติ อนปฺปกา โน ตสฺมิํ สมเย อนตฺตมนตา โหติ อนปฺปกํ โทมนสฺสํ “ทูเร โน ภควา”ติ.}

\prose{ยทา ปน มยํ ภนฺเต ภควนฺตํ สุณาม “มาคเธหิ กาสีสุ จาริกํ ปกฺกมิสฺสตี”ติ, โหติ โน ตสฺมิํ สมเย อตฺตมนตา โหติ โสมนสฺสํ “อาสนฺเน โน ภควา ภวิสฺสตี”ติ. ยทา ปน มยํ ภนฺเต ภควนฺตํ สุณาม “มาคเธหิ กาสีสุ จาริกํ ปกฺกนฺโต”ติ, โหติ โน ตสฺมิํ สมเย อตฺตมนตา โหติ โสมนสฺสํ “อาสนฺเน โน ภควา”ติ.}

\prose{ยทา ปน มยํ ภนฺเต ภควนฺตํ สุณาม “กาสีหิ วชฺชีสุ จาริกํ ปกฺกมิสฺสตี”ติ, โหติ โน ตสฺมิํ สมเย อตฺตมนตา, โหติ}

\prose{\csromanfolio{305}โสมนสฺสํ “อาสนฺเน โน ภควา ภวิสฺสตี”ติ. ยทา ปน มยํ ภนฺเต ภควนฺตํ สุณม “กาสีหิ วชฺชีสุ จาริกํ ปกฺกนฺโต”ติ, โหติ โน ตสฺมิํ สมเย อตฺตมนตา, โหติ โสมนสฺสํ “อาสนฺเน โน ภควา”ติ.}

\prose{ยทา ปน มยํ ภนฺเต ภควนฺตํ สุณาม “วชฺชีหิ มลฺเลสุ จาริกํ ปกฺกมิสฺสตี”ติ, โหติ โน ตสฺมิํ สมเย อตฺตมนตา, โหติ โสมนสฺสํ “อาสนฺเน โน ภควา ภวิสฺสตี”ติ. ยทา ปน มยํ ภนฺเต ภควนฺตํ สุณาม “วชฺชีหิ มลฺเลสุ จาริกํ ปกฺกนฺโต”ติ, โหติ โน ตสฺมิํ สมเย อตฺตมนตา, โหติ โสมนสฺสํ “อาสนฺเน โน ภควา”ติ.}

\prose{ยทา ปน มยํ ภนฺเต ภควนฺตํ สุณาม “มลฺเลหิ โกสเล จาริกํ ปกฺกมิสฺสตี”ติ, โหติ โน ตสฺมิํ สมเย อตฺตมนตา, โหติ โสมนสฺสํ “อาสนฺเน โน ภควา ภวิสฺสตี”ติ. ยทา ปน มยํ ภนฺเต ภควนฺตํ สุณาม “มลฺเลหิ โกสเล จาริกํ ปกฺกนฺโต”ติ, โหติ โน ตสฺมิํ สมเย อตฺตมนตา, โหติ โสมนสฺสํ “อาสนฺเน โน ภควา”ติ.}

\prose{ยทา ปน มยํ ภนฺเต ภควนฺตํ สุณาม “โกสเลหิ สาวตฺถิํ จาริกํ ปกฺกมิสฺสตี”ติ, โหติ โน ตสฺมิํ สมเย อตฺตมนตา, โหติ โสมนสฺสํ “อาสนฺเน โน ภควา ภวิสฺสตี”ติ. ยทา ปน มยํ ภนฺเต ภควนฺตํ สุณาม “สาวตฺถิยํ วิหรติ เชตวเน อนาถ\-ปิณฺฑิกสฺส อาราเม”ติ, โหติ อนปฺปกา โน ตสฺมิํ สมเย อตฺตมนตา, โหติ อนปฺปกํ โสมนสฺสํ “อาสนฺเน โน ภควา”ติ.}

\prose{ตสฺมาติห ถปตโย สมฺพาโธ ฆราวาโส รชาปโถ, อพฺโภกาโส ปพฺพชฺชา, อลญฺจ ปน โว ถปตโย อปฺปมาทายาติ. อตฺถิ โข โน ภนฺเต เอตมฺหา สมฺพาธา อญฺโญ สมฺพาโธ สมฺพาธตโร เจว สมฺพาธ\-สงฺขาต\-ตโร จาติ. กตโม ปน โว ถปตโย เอตมฺหา สมฺพาธา อญฺโญ สมฺพาโธ สมฺพาธตโร เจว สมฺพาธ\-สงฺขาต\-ตโร จาติ.}

\prose{อิธ มยํ ภนฺเต ยทา ราชา ปเสนทิ โกสโล อุยฺยานภูมิํ นิยฺยาตุกาโม โหติ, เย เต รญฺโญ ปเสนทิสฺส โกสลสฺส นาคา โอปวยฺหา, เต กปฺเปตฺวา, ยา ตา รญฺโญ ปเสนทิสฺส โกสลสฺส ปชาปติโย ปิยา มนาปา, ตา เอกํ ปุรโต}

\prose{\csromanfolio{306}เอกํ ปจฺฉโต นิสีทาเปม. ตาสํ โข ปน ภนฺเต ภคินีนํ เอวรูโป คนฺโธ โหติ, เสยฺยถาปิ นาม คนฺธ\-กรณฺฑกสฺส ตาวเทว วิวริย\-มานสฺส ยถา ตํ ราชกญฺญานํ คนฺเธน วิภูสิตานํ. ตาสํ โข ปน ภนฺเต ภคินีนํ เอวรูโป กายสมฺผสฺโส โหติ, เสยฺยถาปิ นาม ตูลปิจุโน วา กปฺปาสปิจุโน วา ยถา ตํ ราชกญฺญานํ สุเขธิตานํ. ตสฺมิํ โข ปน ภนฺเต สมเย นาโคปิ รกฺขิตพฺโพ โหติ, ตาปิ ภคินิโย รกฺขิตพฺพา โหนฺติ, อตฺตาปิ รกฺขิตพฺโพ โหติ, น โข ปน มยํ ภนฺเต อภิชานาม ตาสุ ภคินีสุ ปาปกํ จิตฺตํ อุปฺปาเทตา. อยํ โข โน ภนฺเต เอตมฺหา สมฺพาธา อญฺโญ สมฺพาโธ สมฺพาธตโร เจว สมฺพาธ\-สงฺขาต\-ตโร จาติ.}

\prose{ตสฺมาติห ถปตโย สมฺพาโธ ฆราวาโส รชาปโถ, อพฺโภกาโส ปพฺพชฺชา, อลญฺจ ปน โว ถปตโย อปฺปมาทาย. จตูหิ โข ถปตโย ธมฺเมหิ สมนฺนาคโต อริยสาวโก โสตาปนฺโน โหติ อวินิปาต\-ธมฺโม นิยโต สมฺโพธิ\-ปรายโณ.}

\prose{กตเมหิ จตูหิ, อิธ ถปตโย อริยสาวโก พุทฺเธ อเวจฺจ\-ปฺปสาเทน สมนฺนาคโต โหติ “อิติปิ โส ภควา ฯเปฯ สตฺถา เทวมนุสฺสานํ พุทฺโธ ภควา”ติ. ธมฺเม ฯเปฯ. สํเฆ ฯเปฯ. วิคต\-มล\-มจฺเฉเรน เจตสา อชฺฌาคารํ วสติ มุตฺตจาโค ปยตปาณิ โวสฺสคฺครโต ยาจโยโค ทาน\-สํวิภาค\-รโต, อิเมหิ โข ถปตโย จตูหิ ธมฺเมหิ สมนฺนาคโต อริยสาวโก โสตาปนฺโน โหติ อวินิปาต\-ธมฺโม นิยโต สมฺโพธิ\-ปรายโณ.}

\prose{ตุเมฺห โข ถปตโย พุทฺเธ อเวจฺจ\-ปฺปสาเทน สมนฺนาคตา “อิติปิ โส ภควา ฯเปฯ สตฺถา เทวมนุสฺสานํ พุทฺโธ ภควา”ติ. ธมฺเม ฯเปฯ. สํเฆ ฯเปฯ. ยํ โข ปน กิญฺจิ กุเล เทยฺยธมฺมํ, สพฺพํ ตํ อปฺปฏิวิภตฺตํ สีลวนฺเตหิ กลฺยาณ\-ธมฺเมหิ. ตํ กิํ มญฺญถ ถปตโย, กติวิธา เต โกสเลสุ มนุสฺสา เย ตุมฺหากํ สมสมา ยทิทํ ทาน\-สํวิภาเคติ. ลาภา โน ภนฺเต, สุลทฺธํ โน ภนฺเต, เยสํ โน ภควา เอวํ ปชานาตีติ.  ฉฏฺฐํ.}

\csromanmatikaanchor{mtk.228}
\csromantocmarkref{subparagraph}{7. เวฬุทฺวาเรยฺยสุตฺต}{mtk.228}
\bookmark[dest={mtk.228},level=5]{7. เวฬุทฺวาเรยฺยสุตฺต}
\csromanheader{6}{7. เวฬุทฺวาเรยฺย\-สุตฺต}
\proseitem{1003}{เอวํ เม สุตํ\csromandash{}เอกํ สมยํ ภควา โกสเลสุ จาริกํ จรมาโน มหตา ภิกฺขุ\-สํเฆน สทฺธิํ เยน เวฬุทฺวารํ นาม โกสลานํ}

\prose{\csromanfolio{307}พฺราหฺมณคาโม ตทวสริ. อสฺโสสุํ โข เต เวฬุทฺวาเรยฺยกา พฺราหฺมณ\-คหปติกา “สมโณ ขลุ โภ โคตโม สกฺยปุตฺโต สกฺยกุลา ปพฺพชิโต โกสเลสุ จาริกํ จรมาโน มหตา ภิกฺขุ\-สํเฆน สทฺธิํ เวฬุทฺวารํ อนุปฺปตฺโต, ตํ โข ปน ภวนฺตํ โคตมํ เอวํ กลฺยาโณ กิตฺติสทฺโท อพฺภุคฺคโต ‘อิติปิ โส ภควา อรหํ สมฺมาสมฺพุทฺโธ วิชฺชา\-จรณ\-สมฺปนฺโน สุคโต โลกวิทู อนุตฺตโร ปุริส\-ทมฺม\-สารถิ สตฺถา เทวมนุสฺสานํ พุทฺโธ ภควา’\csromansharedfootnote{524d89cf393f7c6e}{ภควาติ (สี, สฺยา, กํ, อิ)} โส อิมํ โลกํ สเทวกํ สมารกํ สพฺรหฺมกํ สสฺสมณ\-พฺราหฺมณิํ ปชํ สเทวมนุสฺสํ สยํ อภิญฺญา สจฺฉิกตฺวา ปเวเทติ, โส ธมฺมํ เทเสติ อาทิกลฺยาณํ มชฺเฌกลฺยาณํ ปริโยสาน\-กลฺยาณํ สาตฺถํ สพฺยญฺชนํ เกวล\-ปริปุณฺณํ ปริสุทฺธํ พฺรหฺมจริยํ ปกาเสติ, สาธุ โข ปน ตถารูปานํ อรหตํ ทสฺสนํ โหตี”ติ.}

\prose{อถ โข เต เวฬุทฺวาเรยฺยกา พฺราหฺมณ\-คหปติกา เยน ภควา เตนุ\-ปสงฺกมิํสุ, อุปสงฺกมิตฺวา อปฺเปกจฺเจ ภควนฺตํ อภิวาเทตฺวา เอกมนฺตํ นิสีทิํสุ, อปฺเปกจฺเจ ภควตา สทฺธิํ สมฺโมทิํสุ, สมฺโมทนียํ กถํ สารณียํ วีติสาเรตฺวา เอกมนฺตํ นิสีทิํสุ, อปฺเปกจฺเจ เยน ภควา เตนญฺชลิํ ปณาเมตฺวา เอกมนฺตํ นิสีทิํสุ, อปฺเปกจฺเจ ภควโต สนฺติเก นามโคตฺตํ สาเวตฺวา เอกมนฺตํ นิสีทิํสุ, อปฺเปกจฺเจ ตุณฺหีภูตา เอกมนฺตํ นิสีทิํสุ. เอกมนฺตํ นิสินฺนา โข เต เวฬุทฺวาเรยฺยกา พฺราหฺมณ\-คหปติกา ภควนฺตํ เอตทโวจุํ “มยํ โภ โคตม เอวํกามา เอวํฉนฺทา เอวํอธิปฺปายา ปุตฺต\-สมฺพาธ\-สยนํ อชฺฌาวเสยฺยาม, กาสิกจนฺทนํ ปจฺจนุภเวยฺยาม, มาลา\-คนฺธ\-วิเลปนํ ธาเรยฺยาม, ชาตรูป\-รชตํ สาทิเยยฺยาม, กายสฺส เภทา ปรํ มรณา สุคติํ สคฺคํ โลกํ อุปปชฺเชยฺยาม, เตสํ โน ภวํ โคตโม อมฺหากํ เอวํกามานํ เอวํฉนฺทานํ เอวํ\-อธิปฺปายานํ ตถา ธมฺมํ เทเสตุ, ยถา มยํ ปุตฺต\-สมฺพาธ\-สยนํ อชฺฌาวเสยฺยาม ฯเปฯ สุคติํ สคฺคํ โลกํ อุปปชฺเชยฺยามา”ติ.}

\prose{อตฺตูปนายิกํ โว คหปตโย ธมฺม\-ปริยายํ เทเสสฺสามิ, ตํ สุณาถ สาธุกํ มนสิ กโรถ, ภาสิสฺสามีติ. “เอวํ โภ”ติ โข เต เวฬุทฺวาเรยฺยกา พฺราหฺมณ\-คหปติกา ภควโต ปจฺจสฺโสสุํ. ภควา เอตทโวจ\csromandash{}}

\prose{\csromanfolio{308}กตโม จ คหปตโย อตฺตุปนายิโก ธมฺมปริยาโย, อิธ คหปตโย อริยสาวโก อิติ ปฏิสญฺจิกฺขติ “อหํ โขสฺมิ ชีวิตุกาโม อมริตุกาโม สุขกาโม ทุกฺข\-ปฺปฏิกูโล. โย โข มํ ชีวิตุกามํ อมริตุกามํ สุขกามํ ทุกฺข\-ปฺปฏิกูลํ ชีวิตา โวโรเปยฺย, น เมตํ อสฺส ปิยํ มนาปํ. อหญฺเจว โข ปน ปรํ ชีวิตุกามํ อมริตุกามํ สุขกามํ ทุกฺข\-ปฺปฏิกูลํ ชีวิตา โวโรเปยฺยํ, ปรสฺสปิ ตํ อสฺส อปฺปิยํ อมนาปํ. โย โข มฺยายํ ธมฺโม อปฺปิโย อมนาโป, ปรสฺสเปโส ธมฺโม อปฺปิโย อมนาโป. โย โข มฺยายํ ธมฺโม อปฺปิโย อมนาโป, กถาหํ ปรํ เตน สํโยเชยฺยนฺ”ติ. โส อิติ ปฏิสงฺขาย อตฺตนา จ ปาณาติปาตา ปฏิวิรโต โหติ, ปรญฺจ ปาณาติปาตา เวรมณิยา สมาทเปติ, ปาณาติปาตา เวรมณิยา จ วณฺณํ ภาสติ, เอวมสฺสายํ กายสมาจาโร ติโกฏิ\-ปริสุทฺโธ โหติ.}

\prose{ปุน จปรํ คหปตโย อริยสาวโก อิติ ปฏิสญฺจิกฺขติ “โย โข เม อทินฺนํ เถยฺยสงฺขาตํ อาทิเยยฺย, น เมตํ อสฺส ปิยํ มนาปํ. อหญฺเจว โข ปน ปรสฺส อทินฺนํ เถยฺยสงฺขาตํ อาทิยฺเยยฺยํ, ปรสฺสปิ ตํ อสฺส อปฺปิยํ อมนาปํ. โย โข มฺยายํ ธมฺโม อปฺปิโย อมนาโป, ปรสฺสเปโส ธมฺโม อปฺปิโย อมนาโป. โย โข มฺยายํ ธมฺโม อปฺปิโย อมนาโป, กถาหํ ปรํ เตน สํโยเชยฺยนฺ”ติ. โส อิติ ปฏิสงฺขาย อตฺตนา จ อทินฺนาทานํ ปฏิวิรโต โหติ, ปรญฺจ อทินฺนาทานา เวรมณิยา สมาทเปติ, อทินฺนาทานา เวรมณิยา จ วณฺณํ ภาสติ, เอวมสฺสายํ กายสมาจาโร ติโกฏิ\-ปริสุทฺโธ โหติ.}

\prose{ปุน จปรํ คหปตโย อริยสาวโก อิติ ปฏิสญฺจิกฺขติ “โย โข เม ทาเรสุ จาริตฺตํ อาปชฺเชยฺย, น เมตํ อสฺส ปิยํ มนาปํ. อยญฺเจว โข ปน ปรสฺส ทาเรสุ จาริตฺตํ อาปชฺเชยฺยํ, ปรสฺสปิ ตํ อสฺส อปฺปิยํ อมนาปํ. โย โข มฺยายํ ธมฺโม อปฺปิโย อมนาโป, ปรสฺสเปโส ธมฺโม อปฺปิโย อมนาโป. โย โข มฺยายํ ธมฺโม อปฺปิโย อมนาโป, กถาหํ ปรํ เตน สํโยเชยฺยนฺ”ติ. โส อิติ ปฏิสงฺขาย อตฺตนา จ กาเมสุ\-มิจฺฉาจารา ปฏิวิรโต โหติ, ปรญฺจ กาเมสุ\-มิจฺฉาจารา เวรมณิยา สมาทเปติ, กาเมสุ\-มิจฺฉาจารา เวรมณิยา จ}

\prose{\csromanfolio{309}วณฺณํ ภาสติ, เอวมสฺสายํ กายสมาจาโร ติโกฏิ\-ปริสุทฺโธ โหติ.}

\prose{ปุน จปรํ คหปตโย อริยสาวโก อิติ ปฏิสญฺจิกฺขติ “โย โข เม มุสาวาเทน อตฺถํ ภญฺเชยฺย, น เมตํ อสฺส ปิยํ มนาปํ. อหญฺเจว โข ปน ปรสฺส มุสาวาเทน อตฺถํ ภญฺเชยฺยํ, ปรสฺสปิ ตํ อสฺส อปฺปิยํ อมนาปํ. โย โข มฺยายํ ธมฺโม อปฺปิโย อมนาโป, ปรสฺสเปโส ธมฺโม อปฺปิโย อมนาโป. โย โข มฺยายํ ธมฺโม อปฺปิโย อมนาโป, กถาหํ ปรํ เตน สํโยเชยฺยนฺ”ติ. โส อิติ ปฏิสงฺขาย อตฺตนา จ มุสาวาทา ปฏิวิรโต โหติ, ปรญฺจ มุสาวาทา เวรมณิยา สมาทเปติ, มุสาวาทา เวรมณิยา จ วณฺณํ ภาสติ, เอวมสฺสายํ วจีสมาจาโร ติโกฏิ\-ปริสุทฺโธ โหติ.}

\prose{ปุน จปรํ คหปตโย อริยสาวโก อิติ ปฏิสญฺจิกฺขติ โย โข มํ ปิสุณาย วาจาย มิตฺเต ภินฺเทยฺย\csromansharedfootnote{a5417c846763610d}{มิตฺเตหิ เภเทยฺย (สฺยา, กํ, อิ, ก)}, น เมตํ อสฺส ปิยํ มนาปํ. อหญฺเจว โข ปน ปรํ ปิสุณาย วาจาย มิตฺเต ภินฺเทยฺยํ, ปรสฺสปิ ตํ อสฺส อปฺปิยํ อมนาปํ ฯเปฯ เอวมสฺสายํ วจีสมาจาโร ติโกฏิ\-ปริสุทฺโธ โหติ.}

\prose{ปุน จปรํ คหปตโย อริยสาวโก อิติ ปฏิสญฺจิกฺขติ โย โข มํ ผรุสาย วาจาย สมุทาจเรยฺย, น เมตํ อสฺส ปิยํ มนาปํ. อหญฺเจว โข ปน ปรํ ผรุสาย วาจาย สมุทาจเรยฺยํ, ปรสฺสปิ ตํ อสฺส อปฺปิยํ อมนาปํ. โย โข มฺยายํ ธมฺโม ฯเปฯ เอวมสฺสายํ วจีสมาจาโร ติโกฏิ\-ปริสุทฺโธ โหติ.}

\prose{ปุน จปรํ คหปตโย อริยสาวโก อิติ ปฏิสญฺจิกฺขติ “โย โข มํ สมฺผภาเสน สมฺผปฺปลาป\-ภาเสน สมุทาจเรยฺย, น เมตํ อสฺส ปิยํ มนาปํ. อหญฺเจว โข ปน ปรํ สมฺผภาเสน สมฺผปฺปลาป\-ภาเสน สมุทาจเรยฺยํ, ปรสฺสปิ ตํ อสฺส อปฺปิยํ อมนาปํ. โย โข มฺยายํ ธมฺโม อปฺปิโย อมนาโป, ปรสฺสเปโส ธมฺโม อปฺปิโย อมนาโป. โย โข มฺยายํ ธมฺโม อปฺปิโย อมนาโป, กถาหํ ปรํ เตน สโยเชยฺยนฺ”ติ. โส อิติ ปฏิสงฺขาย อตฺตนา จ สมฺผปฺปลาปา ปฏิวิรโต}

\csromanmatikaanchor{mtk.229}
\csromantocmarkref{subparagraph}{8-10. คิญฺชกาวสถสุตฺตานิ}{mtk.229}
\bookmark[dest={mtk.229},level=5]{8-10. คิญฺชกาวสถสุตฺตานิ}
\prose{\csromanfolio{310}โหติ, ปรญฺจ สมฺผปฺปลาปา เวรมณิยา สมาทเปติ, สมฺผปฺปลาปา เวรมณิยา จ วณฺณํ ภาสติ, เอวมสฺสายํ วจีสมาจาโร ติโกฏิ\-ปริสุทฺโธ โหติ.}

\prose{โส พุทฺเธ อเวจฺจ\-ปฺปสาเทน สมนฺนาคโต โหติ “อิติปิ โส ภควา ฯเปฯ สตฺถา เทวมนุสฺสานํ พุทฺโธ ภควา”ติ. ธมฺเม ฯเปฯ. สํเฆ อเวจฺจ\-ปฺปสาเทน สมนฺนาคโต โหติ “สุปฺปฏิปนฺโน ภควโต สาวกสํโฆ ฯเปฯ อนุตฺตรํ ปุญฺญกฺเขตฺตํ โลกสฺสา”ติ. อริยกนฺเตหิ สีเลหิ สมนฺนาคโต โหติ อขณฺเฑหิ ฯเปฯ สมาธิ\-สํวตฺตนิเกหิ. ยโต โข คหปตโย อริยสาวโก อิเมหิ สตฺตหิ สทฺธมฺเมหิ\csromansharedfootnote{a608102cd12125b3}{ธมฺเมหิ (สี)} สมนฺนาคโต โหติ อิเมหิ จตูหิ อากงฺขิเยหิ ฐาเนหิ, โส อากงฺขมาโน อตฺตนาว อตฺตานํ พฺยากเรยฺย “ขีณนิรโยมฺหิ ขีณติรจฺฉานโยนิ\csromansharedfootnote{3e29a9e1e77e05b8}{ขีณติรจฺฉานโยนิโย (สี, สฺยา, กํ, อิ), ขีณติรจฺฉานโยนิโก (ก)} ขีณเปตฺติวิสโย ขีณา\-ปายทุคฺคติวินิปาโต, โสตาปนฺโนหมสฺมิ อวินิปาต\-ธมฺโม นิยโต สมฺโพธิ\-ปรายโณ”ติ.}

\prose{เอวํ วุตฺเต เวฬุทฺวาเรยฺยกา พฺราหฺมณ\-คหปติกา ภควนฺตํ เอตทโวจุํ “อภิกฺกนฺตํ โภ โคตม ฯเปฯ เอเต มยํ ภวนฺตํ โคตมํ สรณํ คจฺฉาม ธมฺมญฺจ ภิกฺขุสํฆญฺจ, อุปาสเก โน ภวํ โคตโม ธาเรตุ อชฺชตคฺเค ปาณุเปเต\csromansharedfootnote{ebd455903f695188}{ปาณุเปตํ (ก)} สรณํ คเต”ติ. สตฺตมํ.}

\csromanheader{2}{8. ปฐม\-คิญฺชกาวสถ\-สุตฺต}
\proseitem{1004}{เอวํ เม สุตํ\csromandash{}เอกํ สมยํ ภควา ญาติเก วิหรติ คิญฺชกาวสเถ. อถ โข อายสฺมา อานนฺโท เยน ภควา เตนุปสงฺกมิ, อุปสงฺกมิตฺวา ภควนฺตํ อภิวาเทตฺวา เอกมนฺตํ นิสีทิ, เอกมนฺตํ นิสินฺโน โข อายสฺมา อานนฺโท ภควนฺตํ เอตทโวจ\csromandash{}}

\prose{สาโฬฺห นาม ภนฺเต ภิกฺขุ กาลงฺกโต, ตสฺส กา คติ, โก อภิสมฺปราโย. นนฺทา นาม ภนฺเต ภิกฺขุนี กาลงฺกตา, ตสฺสา กา คติ, โก อภิสมฺปราโย. สุทตฺโต นาม ภนฺเต อุปาสโก กาลงฺกโต, ตสฺส กา คติ, โก อภิสมฺปราโย. สุชาตา นาม ภนฺเต อุปาสิกา กาลงฺกตา, ตสฺสา กา คติ, โก อภิสมฺปราโยติ.}

\prose{\csromanfolio{311}สาโฬฺห อานนฺท ภิกฺขุ กาลงฺกโต อาสวานํ ขยา อนาสวํ เจโตวิมุตฺติํ ปญฺญาวิมุตฺติํ ทิฏฺเฐว ธมฺเม สยํ อภิญฺญา สจฺฉิกตฺวา อุปสมฺปชฺช วิหาสิ. นนฺทา อานนฺท ภิกฺขุนี กาลงฺกตา ปญฺจนฺนํ โอรมฺภาคิยานํ สํโยชนานํ ปริกฺขยา โอปปาติกา ตตฺถ ปรินิพฺพายินี อนาวตฺติธมฺมา ตสฺมา โลกา. สุทตฺโต อานนฺท อุปาสโก กาลงฺกโต ติณฺณํ สํโยชนานํ ปริกฺขยา ราค\-โทส\-โมหานํ ตนุตฺตา สกทาคามี สกิเทว อิมํ โลกํ อาคนฺตฺวา ทุกฺขสฺสนฺตํ กริสฺสติ. สุชาตา อานนฺท อุปาสิกา กาลงฺกตา ติณฺณํ สํโยชนานํ ปริกฺขยา โสตาปนฺนา อวินิปาต\-ธมฺมา นิยตา สมฺโพธิ\-ปรายณา.}

\prose{อนจฺฉริยํ โข ปเนตํ อานนฺท ยํ มนุสฺสภูโต กาลํ กเรยฺย, ตสฺมิํ ตสฺมิํ เจ มํ กาลงฺกเต อุปสงฺกมิตฺวา เอตมตฺถํ ปฏิปุจฺฉิสฺสถ, วิเหสาเปสา อานนฺท อสฺส ตถาคตสฺส, ตสฺมาติหานนฺท ธมฺมาทาสํ นาม ธมฺม\-ปริยายํ เทเสสฺสามิ, เยน สมนฺนาคโต อริยสาวโก อากงฺขมาโน อตฺตนาว อตฺตานํ พฺยากเรยฺย “ขีณนิรโยมฺหิ ขีณติรจฺฉานโยนิ ขีณเปตฺติวิสโย ขีณา\-ปายทุคฺคติวินิปาโต, โสตาปนฺโนหมสฺมิ อวินิปาต\-ธมฺโม นิยโต สมฺโพธิ\-ปรายโณ”.}

\prose{กตโม จ โส อานนฺท ธมฺมาทาโส ธมฺมปริยาโย, เยน สมนฺนาคโต อริยสาวโก อากงฺขมาโน อตฺตนาว อตฺตานํ พฺยากเรยฺย “ขีณนิรโยมฺหิ ขีณติรจฺฉานโยนิ ขีณเปตฺติวิสโย ขีณา\-ปายทุคฺคติวินิปาโต, โสตาปนฺโนหมสฺมิ อวินิปาต\-ธมฺโม นิยโต สมฺโพธิ\-ปรายโณ”.}

\prose{อิธ อานนฺท อริยสาวโก พุทฺเธ อเวจฺจ\-ปฺปสาเทน สมนฺนาคโต โหติ “อิติปิ โส ภควา ฯเปฯ สตฺถา เทวมนุสฺสานํ พุทฺโธ ภควา”ติ. ธมฺเม ฯเปฯ. สํเฆ ฯเปฯ. อริยกนฺเตหิ สีเลหิ สมนฺนาคโต โหติ อขณฺเฑหิ ฯเปฯ สมาธิ\-สํวตฺตนิเกหิ. อยํ โข โส อานนฺท ธมฺมฺทาโส ธมฺมปรุยาโย, เยน สมนฺนาคโต อริยสาวโก อากงฺขมาโน อตฺตนาว อตฺตานํ พฺยากเรยฺย “ขีณนิรโยมฺหิ ขีณติรจฺฉานโยนิ ขีณเปตฺติวิสโย ขีณา\-ปายทุคฺคติวินิปาโต, โสตาปนฺโนหมสฺมิ อวินิปาต\-ธมฺโม นิยโต สมฺโพธิ\-ปรายโณ”ติ. อฏฺฐมํ. (ตีณิปิ สุตฺตนฺตานิ เอกนิทานานิ.)}

\csromanheader{2}{9. ทุติย\-คิญฺชกาวสถ\-สุตฺต}
\proseitem{1005}{\csromanfolio{312}เอกมนฺตํ นิสินฺโน โข อายสฺมา อานนฺโท ภควนฺตํ เอตทโวจ “อโสโก นาม ภนฺเต ภิกฺขุ กาลงฺกโต, ตสฺส กา คติ, โก อภิสมฺปราโย. อโสกา นาม ภนฺเต ภิกฺขุนี กาลงฺกตา ฯเปฯ. อโสโก นาม ภนฺเต อุปาสโก กาลงฺกโต ฯเปฯ อโสกา นาม ภนฺเต อุปาสิกา กาลงฺกตา, ตสฺสา กา คติ, โก อภิสมฺปราโย”ติ.}

\prose{อโสโก อานนฺท ภิกฺขุ กาลงฺกโต อาสวานํ ขยา อนาสวํ เจโตวิมุตฺติํ ปญฺญาวิมุตฺติํ ทิฏฺเฐว ธมฺเม สยํ อภิญฺญา สจฺฉิกตฺวา อุปสมฺปชฺช วิหาสิ ฯเปฯ. (ปุริม\-เวยฺยากรเณน เอกนิทานํ.)}

\prose{อยํ โข โส อานนฺท ธมฺมาทาโส ธมฺมปริยาโย, เยน สมนฺนาคโต อริยสาวโก อากงฺขมาโน อตฺตนาว อตฺตานํ พฺยากเรยฺย “ขีณนิรโยมฺหิ ขีณติรจฺฉานโยนิ ขีณเปตฺติวิสโย ขิณาปายทุคฺคติวินิปาโต, โสตาปนฺโนหมสฺมิ อวินิปาต\-ธมฺโม นิยโต สมฺโพธิ\-ปรายโณ”ติ.  นวมํ.}

\csromanheader{2}{10. ตติย\-คิญฺชกาวสถ\-สุตฺต}
\proseitem{1006}{เอกมนฺตํ นิสินฺโน โข อายสฺมา อานนฺโท ภควนฺตํ เอตทโวจ “กกฺกโฏ นาม ภนฺเต ญาติเก อุปาสโก กาลงฺกโต, ตสฺส กา คติ, โก อภิสมฺปราโย. กฬิโภ นาม ภนฺเต ญาติเก อุปาสโก ฯเปฯ นิกโต นาม ภนฺเต ญาติเก อุปาสโก. กฏิสฺสโห นาม ภนฺเต ญาติเก อุปาสโก. ตุฏฺโฐ นาม ภนฺเต ญาติเก อุปาสโก. สนฺตุฏฺโฐ นาม ภนฺเต ญาติเก อุปาสโก. ภทฺโท นาม ภนฺเต ญาติเก อุปาสโก. สุภทฺโท นาม ภนฺเต ญาติเก อุปาสโก กาลงฺกโต, ตสฺส กา คติ, โก อภิสมฺปราโยติ.}

\prose{กกฺกโฏ อานนฺท อุปาสโก กาลงฺกโต ปญฺจนฺนํ โอรมฺภาคิยานํ สํโยชนานํ ปริกฺขยา โอปปาติโก ตตฺถ ปรินิพฺพายี อนาวตฺติธมฺโม ตสฺมา โลกา. กฬิโภ อานนฺท ฯเปฯ. นิกโต อานนฺท. กฏิสฺสโห}

\prose{\csromanfolio{313}อานนฺท. ตุฏฺโฐ อานนฺท. สนฺตุฏฺโฐ อานนฺท. ภทฺโท อานนฺท. สุภทฺโท อานนฺท อุปาสโก กาลงฺกโต ปญฺจนฺนํ โอรมฺภาคิยานํ สํโยชนานํ ปริกฺขยา โอปปาติโก ตตฺถ ปรินิพฺพายี อนาวตฺติธมฺโม ตสฺมา โลกา. (สพฺเพ เอกคติกา กาตพฺพา.)}

\prose{ปโรปญฺญาส อานนฺท ญาติเก อุปาสกา กาลงฺกตา ปญฺจนฺนํ โอรมฺภาคิยานํ สํโยชนานํ ปริกฺขยา โอปปาติกา ตตฺถ ปรินิพฺพายิโน อนาวตฺติธมฺมา ตสฺมา โลกา. สาธิกนวุติ อานนฺท ญาติเก อุปาสกา กาลงฺกตา ติณฺณํ สํโยชนานํ ปริกฺขยา ราค\-โทส\-โมหานํ ตนุตฺตา สกทาคามิโน สกิเทว อิมํ โลกํ อาคนฺตฺวา ทุกฺขสฺสนฺตํ กริสฺสนฺติ. ฉาติเรกานิ โข อานนฺท ปญฺจสตานิ ญาติเก อุปาสกา กาลงฺกตา ติณฺณํ สํโยชนานํ ปริกฺขยา โสตาปนฺนา อวินิปาต\-ธมฺมา นิยตา สมฺโพธิ\-ปรายณา.}

\prose{อนจฺฉริยํ โข ปเนตํ อานนฺท ยํ มนุสฺสภูโต กาลํ กเรยฺย, ตสฺมิํ ตสฺมิํ เจ มํ กาลงฺกเต อุปสงฺกมิตฺวา เอตมตฺถํ ปฏิปุจฺฉิสฺสถ, วิเหสาเปสา อาอานนฺท อสฺส ตถาคตสฺส. ตสฺมาติหานนฺท ธมฺมาทาสํ นาม ธมฺม\-ปริยายํ เทเสสฺสามิ, เยน สมนฺนาคโต อริยสาวโก อากงฺขมาโน อตฺตนาว อตฺตานํ พฺยากเรยฺย “ขีณนิรโยมฺหิ ขีณติรจฺฉานโยนิ ขีณเปตฺติวิสโย ขีณา\-ปายทุคฺคติวินิปาโต, โสตาปนฺโนหมสฺมิ อวินิปาต\-ธมฺโม นิยโต สมฺโพธิ\-ปรายโณ”.}

\prose{กตโม จ โส อานนฺท ธมฺมาทาโส ธมฺมปริยาโย เยน สมนฺนาคโต อริยสาวโก อากงฺขมาโน อตฺตนาว อตฺตานํ พฺยากเรยฺย “ขีณนิรโยมฺหิ ขีณติรจฺฉานโยนิ ขีณเปตฺติวิสโย ขีณา\-ปายทุคฺคติวินิปาโต, โสตาปนฺโนหมสฺมิ อวินิปาต\-ธมฺโม นิยโต สมฺโพธิ\-ปรายโณ”.}

\prose{อิธานนฺท อริยสาวโก พุทฺเธ อเวจฺจ\-ปฺปสาเทน สมนฺนาคโต โหติ “อิติปิ โส ภควา ฯเปฯ สตฺถา เทวมนุสฺสานํ พุทฺโธ ภควา”ติ. ธมฺเม ฯเปฯ. สํเฆ ฯเปฯ. อริยกนฺเตหิ สีเลหิ สมนฺนาคโต โหติ อขณฺเฑหิ ฯเปฯ สมาธิ\-สํวตฺตนิเกหิ. อยํ โข โส อานนฺท ธมฺมาทาโส ธมฺมปริยาโย เยน สมนฺนาคโต อริยสาวโก อากงฺขมาโน อตฺตนาว อตฺตานํ พฺยากเรยฺย “ขีณนิรโยมฺหิ ขีณติรจฺฉานโยนิ}

\prosewithclosermidruled{\csromanfolio{314}ขีณเปตฺติวิสโย ขีณา\-ปายทุคฺคติวินิปาโต, โสตาปนฺโนหมสฺมิ อวินิปาต\-ธมฺโม นิยโต สมฺโพธิ\-ปรายโณ”ติ.  ทสมํ.}{เวฬุทฺวารวคฺโค ปฐโม.}
\csromancenter{\textbf{ตสฺสุทฺทานํ}}
\setlength{\csromangathapagewidth}{0pt}
\setlength{\csromangathawakwidth}{0pt}
\csromangathameasuregroup{ราชา โอคธทีฆาวุ, \\ ถปตี เวฬุทฺวาเรยฺยา,}{\csromangathabat{ราชา โอคธทีฆาวุ,}{สาริปุตฺตาปเร ทุเว.} \\ \csromangathabat{ถปตี เวฬุทฺวาเรยฺยา,}{คิญฺชกาวสเถ ตโยติ.}}
\csromangathasetleft{ราชา โอคธทีฆาวุ, \\ ถปตี เวฬุทฺวาเรยฺยา,}
\csromangathagroup{\csromangathastanza{\csromangathabat{ราชา โอคธทีฆาวุ,}{สาริปุตฺตาปเร ทุเว.} \\ \csromangathabat{ถปตี เวฬุทฺวาเรยฺยา,}{คิญฺชกาวสเถ ตโยติ.}}}

\csromanmatikaanchor{mtk.230}
\csromantocmarknopage{subparagraph}{10. ราชการามวคฺค}{mtk.230}
\bookmark[dest={mtk.230},level=5]{10. ราชการามวคฺค}
\csromanheader{6}{2. ราชการาม\-วคฺค}
\csromanmatikaanchor{mtk.231}
\csromantocmarkref{subparagraph}{1. สหสฺสภิกฺขุนิสํฆสุตฺต}{mtk.231}
\bookmark[dest={mtk.231},level=5]{1. สหสฺสภิกฺขุนิสํฆสุตฺต}
\csromanheader{6}{1. \textbf{สหสฺสภิกฺขุนิสํฆสุตฺต}}
\proseitem{1007}{เอกํ สมยํ ภควา สาวตฺถิยํ วิหรติ ราชการาเม. อถ โข สหสฺส\-ภิกฺขุนิ\-สํโฆ\csromansharedfootnote{17b7f8c6f7edaa71}{สหสฺโส ภิกฺขุนิสํโฆ (สี)} เยน ภควา เตนุปสงฺกมิ, อุปสงฺกมิตฺวา ภควนฺตํ อภิวาเทตฺวา เอกมนฺตํ อฏฺฐาสิ, เอกมนฺตํ ฐิตา โข ตา ภิกฺขุนิโย ภควา เอตทโวจ\csromandash{}}

\prose{จตูหิ โข ภิกฺขุนิโย ธมฺเมหิ สมนฺนาคโต อริยสาวโก โสตาปนฺโน โหติ อวินิปาต\-ธมฺโม นิยโต สมฺโพธิ\-ปรายโณ. กตเมหิ จตูหิ, อิธ ภิกฺขุนิโย อริยสาวโก พุทฺเธ อเวจฺจ\-ปฺปสาเทน สมนฺนาคโต โหติ “อิติปิ โส ภควา ฯเปฯ สตฺถา เทวมนุสฺสานํ พุทฺโธ ภควา”ติ. ธมฺเม ฯเปฯ. สํเฆ ฯเปฯ. อริยกนฺเตหิ สีเลหิ สมนฺนาคโต โหติ อขณฺเฑหิ ฯเปฯ สมาธิ\-สํวตฺตนิเกหิ. อิเมหิ โข ภิกฺขุนิโย จตูหิ ธมฺเมหิ สมนฺนาคโต อริยสาวโก โสตาปนฺโน โหติ อวินิปาต\-ธมฺโม นิยโต สมฺโพธิ\-ปรายโณติ.  ปฐมํ.}

\csromanmatikaanchor{mtk.232}
\csromantocmarkref{subparagraph}{2. พฺราหฺมณสุตฺต}{mtk.232}
\bookmark[dest={mtk.232},level=5]{2. พฺราหฺมณสุตฺต}
\csromanheader{6}{2. \textbf{พฺราหฺมณสุตฺต}}
\proseitem{1008}{สาวตฺถิ\-นิทานํ. พฺราหฺมณา ภิกฺขเว อุทยคามินิํ นาม ปฏิปทํ ปญฺญเปนฺติ, เต สาวกํ เอวํ สมาทเปนฺติ “เอหิ ตฺวํ อมฺโภ ปุริส, กาลสฺเสว อุฏฺฐาย ปาจีนมุโข ยาหิ, โส ตฺวํ มา โสพฺภํ ปริวชฺเชหิ}

\csromanmatikaanchor{mtk.233}
\csromantocmarkref{subparagraph}{3-5. อานนฺทตฺเถรสุตฺตาทีนิ}{mtk.233}
\bookmark[dest={mtk.233},level=5]{3-5. อานนฺทตฺเถรสุตฺตาทีนิ}
\prose{\csromanfolio{315}มา ปปาตํ มา ขาณุํ มา กณฺฑกฐานํ\csromansharedfootnote{a0ca19e9507458c4}{กณฺฑกํ ฐานํ (อิ, ก)} มา จนฺทนิยํ มา โอฬิคลฺลํ, ยตฺถ\csromansharedfootnote{639130a2656df3b4}{ยตฺเถว (สฺยา, กํ), ยานิ วา (สี)} ปปเตยฺยาสิ, ตตฺเถว มรณํ อาคเมยฺยาสิ. เอวํ ตฺวํ อมฺโภ ปุริส กายสฺส เภทา ปรํ มรณา สุคติํ สคฺคํ โลกํ อุปปชฺชิสฺสสี”ติ.}

\prose{ตํ โข ปเนตํ ภิกฺขเว พฺราหฺมณานํ พาลคมนเมตํ\csromansharedfootnote{1aa56ea0e8aa264d}{พาลานํ คมนเมตํ (สี)} มูฬฺหคมนเมตํ, น นิพฺพิทาย น วิราคาย น นิโรธาย น อุปสมาย น อภิญฺญาย น สมฺโพธาย น นิพฺพานาย สํวตฺตติ. อหญฺจ โข ภิกฺขเว อริยสฺส วินเย อุทยคามินิํ ปฏิปทํ ปญฺญาเปมิ, ยา เอกนฺต\-นิพฺพิทาย วิราคาย นิโรธาย อุปสมาย อภิญฺญาย สมฺโพธาย นิพฺพานาย สํวตฺตติ.}

\prose{กตมา จ สา ภิกฺขเว อุทยคามินี ปฏิปทา, ยา เอกนฺต\-นิพฺพิทาย ฯเปฯ นิพฺพานาย สํวตฺตติ, อิธ ภิกฺขเว อริยสาวโก พุทฺเธ อเวจฺจ\-ปฺปสาเทน สมนฺนาคโต โหติ “อิติปิ โส ภควา อรหํ สมฺมาสมฺพุทฺโธ ฯเปฯ สตฺถา เทวมนุสฺสานํ พุทฺโธ ภควา”ติ. ธมฺเม ฯเปฯ. สํเฆ ฯเปฯ. อริยกนฺเตหิ สีเลหิ สมนฺนาคโต โหติ อขณฺเฑหิ ฯเปฯ สมาธิ\-สํวตฺตนิเกหิ. อยํ โข สา ภิกฺขเว อุทยคามินี ปฏิปทา เอกนฺต\-นิพฺพิทาย ฯเปฯ นิพฺพานาย สํวตฺตตีติ.  ทุติยํ.}

\csromanheader{2}{3. อานนฺทตฺเถร\-สุตฺต}
\proseitem{1009}{เอกํ สมยํ อายสฺมา จ อานนฺโท อายสฺมา จ สาริปุตฺโต สาวตฺถิยํ วิหรนฺติ เชตวเน อนาถ\-ปิณฺฑิกสฺส อาราเม. อถ โข อายสฺมา สาริปุตฺโต สายนฺหสมยํ ปฏิสลฺลานา วุฏฺฐิโต เยนายสฺมา อานนฺโท เตนุปสงฺกมิ, อุปสงฺกมิตฺวา อายสฺมตา อานนฺเทน สทฺธิํ สมฺโมทิ, สมฺโมทนียํ กถํ สารณียํ วีติสาเรตฺวา เอกมนฺตํ นิสีทิ, เอกมนฺตํ นิสินฺโน โข อายสฺมา สาริปุตฺโต อายสฺมนฺตํ อานนฺทํ เอตทโวจ “กตินํ โข อาวุโส อานนฺท ธมฺมานํ ปหานา, กตินํ ธมฺมานํ สมนฺนาคมน\-เหตุ เอวมยํ ปชา ภควตา พฺยากตา ‘โสตาปนฺนา อวินิปาต\-ธมฺมา นิยตา สมฺโพธิ\-ปรายณา’ติ”. จตุนฺนํ โข อาวุโส ธมฺมานํ ปหานา จตุนฺนํ ธมฺมานํ สมนฺนาคมน\-เหตุ เอวมยํ ปชา ภควตา พฺยากตา “โสตาปนฺนา อวินิปาต\-ธมฺมา นิยตา สมฺโพธิ\-ปรายณา”ติ.}

\prose{\csromanfolio{316}กตเมสํ จตุนฺนํ, ยถารูเปน โข อาวุโส พุทฺเธ อปฺปสาเทน สมนฺนาคโต อสฺสุตวา ปุถุชฺชโน กายสฺส เภทา ปรํ มรณา อปายํ ทุคฺคติํ วินิปาตํ นิรยํ อุปปชฺชติ, ตถารูปสฺส พุทฺเธ อปฺปสาโท น โหติ. ยถารูเปน จ โข อาวุโส พุทฺเธ อเวจฺจ\-ปฺปสาเทน สมนฺนาคโต สุตวา อริยสาวโก กายสฺส เภทา ปรํ มรณา สุคติํ สคฺคํ โลกํ อุปปชฺชติ, ตถารูปสฺส พุทฺเธ อเวจฺจปฺปสาโท โหติ “อิติปิ โส ภควา ฯเปฯ สตฺถา เทวมนุสฺสานํ พุทฺโธ ภควา”ติ.}

\prose{ยถารูเปน จ โข อาวุโส ธมฺเม อปฺปสาเทน สมนฺนาคโต อสฺสุตวา ปุถุชฺชโน กายสฺส เภทา ปรํ มรณา อปายํ ทุคฺคติํ วินิปาตํ นิรยํ อุปปชฺชติ, ตถารูปสฺส ธมฺเม อปฺปสาโท น โหติ. ยถารูเปน จ โข อาวุโส ธมฺเม อเวจฺจ\-ปฺปสาเทน สมนฺนาคโต สุตวา อริยสาวโก กายสฺส เภทา ปรํ มรณา สุคติํ สคฺคํ โลกํ อุปปชฺชติ, ตถารูปสฺส ธมฺเม อเวจฺจปฺปสาโท โหติ “สฺวากฺขาโต ภควตา ธมฺโม ฯเปฯ วิญฺญูหี”ติ.}

\prose{ยถารูเปน จ โข อาวุโส สํเฆ อปฺปสาเทน สมนฺนาคโต อสฺสุตวา ปุถุชฺชโน กายสฺส เภทา ปรํ มรณา อปายํ ทุคฺคติํ วินิปาตํ นิรยํ อุปปชฺชติ, ตถารูปสฺส สํเฆ อปฺปสาโท น โหติ. ยถารูเปน จ โข อาวุโส สํเฆ อเวจฺจ\-ปฺปสาเทน สมนฺนาคโต สุตวา อริยสาวโก กายสฺส เภทา ปรํ มรณา สุคติํ สคฺคํ โลกํ อุปปชฺชติ, ตถารูปสฺส สํเฆ อเวจฺจปฺปสาโท โหติ “สุปฺปฏิปนฺโน ภควโต สาวกสํโฆ ฯเปฯ อนุตฺตรํ ปุญฺญกฺเขตฺตํ โลกสฺสา”ติ.}

\prose{ยถารูเปน จ โข อาวุโส ทุสฺสีเลฺยน สมนฺนาคโต อสฺสุตวา ปุถุชฺชโน กายสฺส เภทา ปรํ มรณา อปายํ ทุคฺคติํ วินิปาตํ นิรยํ อุปปชฺชติ, ตถารูปสฺส ทุสฺสีลฺยํ น โหติ. ยถารูเปหิ จ โข อาวุโส อริยกนฺเตหิ สีเลหิ สมนฺนาคโต สุตวา อริยสาวโก กายสฺส เภทา ปรํ มรณา สุคติํ สคฺคํ โลกํ อุปปชฺชติ, ตถารูปานิ อริยกนฺตานิ สีลานิ โหนฺติ อขณฺฑานิ ฯเปฯ สมาธิ\-สํวตฺตนิกานิ. อิเมสํ โข อาวุโส จตุนฺนํ ธมฺมานํ ปหานา อิเมสํ}

\csromanmatikaanchor{mtk.234}
\csromantocmarkref{subparagraph}{6-7. มิตฺตามจฺจสุตฺตานิ}{mtk.234}
\bookmark[dest={mtk.234},level=5]{6-7. มิตฺตามจฺจสุตฺตานิ}
\prose{\csromanfolio{317}จตุนฺนํ ธมฺมานํ สมนฺนาคมน\-เหตุ เอวมยํ ปชา ภควตา พฺยากตา “โสตาปนฺนา อวินิปาต\-ธมฺมา นิยตา สมฺโพธิ\-ปรายณา”ติ.  ตติยํ.}

\csromanheader{2}{4. ทุคฺคติภย\-สุตฺต}
\proseitem{1010}{จตูหิ ภิกฺขเว ธมฺเมหิ สมนฺนาคโต อริยสาวโก สพฺพ\-ทุคฺคติภยํ สมติกฺกนฺโต โหติ. กตเมหิ จตูหิ, อิธ ภิกฺขเว อริยสาวโก พุทฺเธ อเวจฺจ\-ปฺปสาเทน สมนฺนาคโต โหติ “อิติปิ โส ภควา ฯเปฯ สตฺถา เทวมนุสฺสานํ พุทฺโธ ภควา”ติ. ธมฺเม ฯเปฯ. สํเฆ ฯเปฯ. อริยกนฺเตหิ สีเลหิ สมนฺนาคโต โหติ อขณฺเฑหิ ฯเปฯ สมาธิ\-สํวตฺตนิเกหิ. อิเมหิ โข ภิกฺขเว จตูหิ ธมฺเมหิ สมนฺนาคโต อริยสาวโก สพฺพ\-ทุคฺคติภยํ สมติกฺกนฺโต โหตีติ.  จตุตฺถํ.}

\csromanheader{2}{5. ทุคฺคติวินิปาตภย\-สุตฺต}
\proseitem{1011}{จตูหิ ภิกฺขเว ธมฺเมหิ สมนฺนาคโต อริยสาวโก สพฺพ\-ทุคฺคติ\-วินิปาตภยํ สมติกฺกนฺโต โหติ. กตเมหิ จตูหิ, อิธ ภิกฺขเว อริยสาวโก พุทฺเธ อเวจฺจ\-ปฺปสาเทน สมนฺนาคโต โหติ “อิติปิ โส ภควา ฯเปฯ สตฺถา เทวมนุสฺสานํ พุทฺโธ ภควา”ติ. ธมฺเม ฯเปฯ. สํเฆ ฯเปฯ. อริยกนฺเตหิ สีเลหิ สมนฺนาคโต โหติ อขณฺเฑหิ ฯเปฯ สมาธิ\-สํวตฺตนิเกหิ. อิเมหิ โข ภิกฺขเว จตูหิ ธมฺเมหิ สมนฺนาคโต อริยสาวโก สพฺพ\-ทุคฺคติ\-วินิปาตภยํ สมติกฺกนฺโต โหตีติ.  ปญฺจมํ.}

\csromanheader{2}{6. ปฐม\-มิตฺตามจฺจ\-สุตฺต}
\csromanmatikaanchor{mtk.235}
\csromantocmarkref{subparagraph}{8-10. เทวจริกสุตฺตานิ}{mtk.235}
\bookmark[dest={mtk.235},level=5]{8-10. เทวจริกสุตฺตานิ}
\proseitem{1012}{เย เต ภิกฺขเว อนุกมฺเปยฺยาถ, เย จ โสตพฺพํ มญฺเญยฺยุํ มิตฺตา วา อมจฺจา วา ญาตี วา สาโลหิตา วา, เต ภิกฺขเว จตูสุ โสตาปตฺติยงฺเคสุ สมาทเปตพฺพา นิเวเสตพฺพา ปติฏฺฐาเป\-ตพฺพา. กตเมสุ จตูสุ, พุทฺเธ อเวจฺจปฺปสาเท สมาทเปตพฺพา นิเวเสตพฺพา ปติฏฺฐาเป\-ตพฺพา “อิติปิ โส ภควา ฯเปฯ สตฺถา เทวมนุสฺสานํ พุทฺโธ ภควา”ติ. ธมฺเม ฯเปฯ. สํเฆ ฯเปฯ. อริยกนฺเตสุ สีเลสุ สมาทเปตพฺพา นิเวเสตพฺพา ปติฏฺฐาเป\-ตพฺพา อขณฺเฑสุ ฯเปฯ สมาธิ\-สํวตฺตนิเกสุ. เย เต ภิกฺขเว อนุกมฺเปยฺยาถ, เย จ โสตพฺพํ มญฺเญยฺยุํ มิตฺตา วา \csromanfolio{318}อมจฺจา วา ญาตี วา สาโลหิตา วา, เต ภิกฺขเว อิเมสุ จตูสุ โสตาปตฺติยงฺเคสุ สมาทเปตพฺพา นิเวเสตพฺพา ปติฏฺฐาเป\-ตพฺพาติ.  ฉฏฺฐํ.}

\csromanheader{2}{7. ทุติย\-มิตฺตามจฺจ\-สุตฺต}
\proseitem{1013}{เย เต ภิกฺขเว อนุกมฺเปยฺยาถ, เย จ โสตพฺพํ มญฺเญยฺยุํ มิตฺตา วา อมจฺจา วา ญาตี วา สาโลหิตา วา, เต ภิกฺขเว จตูสุ โสตาปตฺติยงฺเคสุ สมาทเปตพฺพา นิเวเสตพฺพา ปติฏฺฐาเป\-ตพฺพา. กตเมสุ จตูสุ, พุทฺเธ อเวจฺจปฺปสาเท สมาทเปตพฺพา นิเวเสตพฺพา ปติฏฺฐาเป\-ตพฺพา “อิติปิ โส ภควา ฯเปฯ สตฺถา เทวมนุสฺสานํ พุทฺโธ ภควา”ติ.}

\prose{สิยา ภิกฺขเว จตุนฺนํ มหาภูตานํ อญฺญถตฺตํ ปถวีธาตุยา อาโปธาตุยา เตโชธาตุยา วาโยธาตุยา, น เตฺวว พุทฺเธ อเวจฺจ\-ปฺปสาเทน สมนฺนาคตสฺส อริยสาวกสฺส สิยา อญฺญถตฺตํ. ตตฺริทํ อญฺญถตฺตํ, โส วต พุทฺเธ อเวจฺจ\-ปฺปสาเทน สมนฺนาคโต อริยสาวโก นิรยํ วา ติรจฺฉาน\-โยนิํ วา เปตฺติวิสยํ วา อุปปชฺชิสฺสตีติ เนตํ ฐานํ วิชฺชติ. ธมฺเม ฯเปฯ. สํเฆ ฯเปฯ. อริยกนฺเตสุ สีเลสุ สมาทเปตพฺพา นิเวเสตพฺพา ปติฏฺฐาเป\-ตพฺพา อขณฺเฑสุ ฯเปฯ สมาธิ\-สํวตฺตนิเกสุ. สิยา ภิกฺขเว จตุนฺนํ มหาภูตานํ อญฺญถตฺตํ ปถวีธาตุยา อาโปธาตุยา เตโชธาตุยา วาโยธาตุยา, น เตฺวว อริยกนฺเตหิ สีเลหิ สมนฺนาคตสฺส อริยสาวกสฺส สิยา อญฺญถตฺตํ. ตตฺริทํ อญฺญถตฺตํ, โส วต อริยกนฺเตหิ สีเลหิ สมนฺนาคโต อริยสาวโก นิรยํ วา ติรจฺฉาน\-โยนิํ วา เปตฺติวิสยํ วา อุปปชฺชิสฺสตีติ เนตํ ฐานํ วิชฺชติ. เย เต ภิกฺขเว อนุกมฺเปยฺยาถ, เย จ โสตพฺพํ มญฺเญยฺยุํ มิตฺตา วา อมจฺจา วา ญาตี วา สาโลหิตา วา, เต ภิกฺขเว อิเมสุ จตูสุ โสตาปตฺติยงฺเคสุ สมาทเปตพฺพา นิเวเสตพฺพา ปติฏฺฐาเป\-ตพฺพาติ.  สตฺตมํ.}

\csromanheader{2}{8. ปฐม\-เทวจาริก\-สุตฺต}
\proseitem{1014}{สาวตฺถิ\-นิทานํ. อถ โข อายสฺมา มหาโมคฺคลฺลาโน เสยฺยถาปิ นาม พลวา ปุริโส สมิญฺชิตํ วา พาหํ ปสาเรยฺย, ปสาริตํ}

\prose{\csromanfolio{319}วา พาหํ สมิญฺเชยฺย, เอวเมว เชตวเน อนฺตรหิโต เทเวสุ ตาวติํเสสุ ปาตุรโหสิ. อถ โข สมฺพหุลา ตาวติํส\-กายิกา เทวตาโย เยนายสฺมา มหาโมคฺคลฺลาโน เตนุ\-ปสงฺกมิํสุ, อุปสงฺกมิตฺวา อายสฺมนฺตํ มหา\-โมคฺคลฺลานํ อภิวาเทตฺวา เอกมนฺตํ อฏฺฐํสุ, เอกมนฺตํ ฐิตา โข ตา เทวตาโย อายสฺมา มหาโมคฺคลฺลาโน เอตทโวจ\csromandash{}}

\prose{สาธุ โข อาวุโส พุทฺเธ อเวจฺจ\-ปฺปสาเทน สมนฺนาคมนํ โหติ “อิติปิ โส ภควา ฯเปฯ สตฺถา เทวมนุสฺสานํ พุทฺโธ ภควา”ติ. พุทฺเธ อเวจฺจ\-ปฺปสาเทน สมนฺนาคมน\-เหตุ โข อาวุโส เอวมิเธกจฺเจ สตฺตา กายสฺส เภทา ปรํ มรณา สุคติํ สคฺคํ โลกํ อุปปชฺชนฺติ. สาธุ โข อาวุโส ธมฺเม ฯเปฯ. สํเฆ ฯเปฯ. สาธุ โข อาวุโส อริยกนฺเตหิ สีเลหิ สมนฺนาคมนํ โหติ อขณฺเฑหิ ฯเปฯ สมาธิ\-สํวตฺตนิเกหิ. อริยกนฺเตหิ สีเลหิ สมนฺนาคมน\-เหตุ โข อาวุโส เอวมิเธกจฺเจ สตฺตา กายสฺส เภทา ปรํ มรณา สุคติํ สคฺคํ โลกํ อุปปชฺชนฺตีติ.}

\prose{สาธุ โข มาริส โมคฺคลฺลาน, พุทฺเธ อเวจฺจ\-ปฺปสาเทน สมนฺนาคมนํ โหติ “อิติปิ โส ภควา ฯเปฯ สตฺถา เทวมนุสฺสานํ พุทฺโธ ภควา”ติ. พุทฺเธ อเวจฺจ\-ปฺปสาเทน สมนฺนาคมน\-เหตุ โข มาริส โมคฺคลฺลาน เอวมิเธกจฺเจ สตฺตา กายสฺส เภทา ปรํ มรณา สุคติํ สคฺคํ โลกํ อุปปชฺชนฺติ. สาธุ โข มาริส โมคฺคลฺลาน ธมฺเม ฯเปฯ. สํเฆ ฯเปฯ. อริยกนฺเตหิ สีเลหิ สมนฺนาคมนํ โหติ อขณฺเฑหิ ฯเปฯ สมาธิ\-สํวตฺตนิเกหิ. อริยกนฺเตหิ สีเลหิ สมนฺนาคมน\-เหตุ โข มาริส โมคฺคลฺลาน เอวมิเธกจฺเจ สตฺตา กายสฺส เภทา ปรํ มรณา สุคติํ สคฺคํ โลกํ อุปปชฺชนฺตีติ.  อฏฺฐมํ.}

\csromanheader{2}{9. ทุติย\-เทวจาริก\-สุตฺต}
\proseitem{1015}{สาวตฺถิ\-นิทานํ. อถ โข อายสฺมา มหาโมคฺคลฺลาโน เสยฺยถาปิ นาม พลวา ปุริโส สมิญฺชิตํ วา พาหํ ปสาเรยฺย, ปสาริตํ วา พาหํ สมิญฺเชยฺย, เอวเมว เชตวเน อนฺตรหิโต เทเวสุ ตาวติํเสสุ ปาตุรโหสิ. อถ โข สมฺพหุลา ตาวติํส\-กายิกา เทวตาโย เยนายสฺมา มหาโมคฺคลฺลาโน เตนุ\-ปสงฺกมิํสุ, อุปสงฺกมิตฺวา อายสฺมนฺตํ มหา\-โมคฺคลฺลานํ อภิวาเทตฺวา เอกมนฺตํ อฏฺฐํสุ, \csromanfolio{320}เอกมนฺตํ ฐิตา โข ตา เทวตาโย ยสฺมา มหาโมคฺคลฺลาโน เอตทโวจ\csromandash{}}

\prose{สาธุ โข อาวุโส พุทฺเธ อเวจฺจ\-ปฺปสาเทน สมนฺนาคมนํ โหติ “อิติปิ โส ภควา ฯเปฯ สตฺถา เทวมนุสฺสานํ พุทฺโธ ภควา”ติ. พุทฺเธ อเวจฺจ\-ปฺปสาเทน สมนฺนาคมน\-เหตุ โข อาวุโส เอวมิเธกจฺเจ สตฺตา กายสฺส เภทา ปรํ มรณา สุคติํ สคฺคํ โลกํ อุปปนฺนา. สาธุ โข อาวุโส ธมฺเม ฯเปฯ. สํเฆ ฯเปฯ. อริยกนฺเตหิ สีเลหิ สมนฺนาคมนํ โหติ อขณฺเฑหฺ ฯเปฯ สมาธิ\-สํวตฺตนิเกหิ. อริยกนฺเตหิ สีเลหิ สมนฺนาคมน\-เหตุ โข อาวุโส เอวมิเธกจฺเจ สตฺตา กายสฺส เภทา ปรํ มรณา สุคติํ สคฺคํ โลกํ อุปปนฺนาติ.}

\prose{สาธุ โข มาริส โมคฺคลฺลาน, พุทฺเธ อเวจฺจ\-ปฺปสาเทน สมนฺนาคมนํ โหติ “อิติปิ โส ภควา ฯเปฯ สตฺถา เทวมนุสฺสานํ พุทฺโธ ภควา”ติ. พุทฺเธ อเวจฺจ\-ปฺปสาเทน สมนฺนาคมน\-เหตุ โข มาริส โมคฺคลฺลาน เอวมิเธกจฺเจ สตฺตา กายสฺส เภทา ปรํ มรณา สุคติํ สคฺคํ โลกํ อุปปนฺนา. สาธุ โข มาริส โมคฺคลฺลาน ธมฺเม ฯเปฯ. สํเฆ ฯเปฯ. อริยกนฺเตหิ สีเลหิ สมนฺนาคมนํ โหติ อขณฺเฑหิ ฯเปฯ สมาธิ\-สํวตฺตนิเกหิ. อริยกนฺเตหิ สีเลหิ สมนฺนาคมน\-เหตุ โข มาริส โมคฺคลฺลาน เอวมิเธกจฺเจ สตฺตา กายสฺส เภทา ปรํ มรณา สุคติํ สคฺคํ โลกํ อุปปนฺนาติ.  นวมํ.}

\csromanheader{2}{10. ตติย\-เทวจาริก\-สุตฺต}
\proseitem{1016}{อถ โข ภควา เสยฺยถาปิ นาม พลวา ปุริโส สมิญฺชิตํ วา พาหํ ปสาเรยฺย, ปสาริตํ วา พาหํ สมิญฺเชยฺย, เอวเมว เชตวเน อนฺตรหิโต เทเวสุ ตาวติํเสสุ ปาตุรโหสิ. อถ โข สมฺพหุลา ตาวติํส\-กายิกา เทวตาโย เยน ภควา เตนุ\-ปสงฺกมิํสุ, อุปสงฺกมิตฺวา ภควนฺตํ อภิวาเทตฺวา เอกมนฺตํ อฏฺฐํสุ, เอกมนฺตํ ฐิตา โข ตา เทวตาโย ภควา เอตทโวจ\csromandash{}}

\prose{สาธุ โข อาวุโส พุทฺเธ อเวจฺจ\-ปฺปสาเทน สมนฺนาคมนํ โหติ “อิติปิ โส ภควา ฯเปฯ สตฺถา เทวมนุสฺสานํ พุทฺโธ ภควา”ติ, พุทฺเธ อเวจฺจ\-ปฺปสาเทน สมนฺนาคมน\-เหตุ โข อาวุโส เอวมิเธกจฺเจ สตฺตา โสตาปนฺนา อวินิปาต\-ธมฺมา นิยตา สมฺโพธิ\-ปรายณา. สาธุ โข}

\csromanmatikaanchor{mtk.236}
\csromanmatikaanchor{mtk.237}
\csromantocmarknopage{subparagraph}{3. สรณานิวคฺค}{mtk.236}
\bookmark[dest={mtk.236},level=5]{3. สรณานิวคฺค}
\csromantocmarkref{subparagraph}{1-2. มหานามสุตฺตานิ}{mtk.237}
\bookmark[dest={mtk.237},level=5]{1-2. มหานามสุตฺตานิ}
\prose{\csromanfolio{321}อาวุโส ธมฺเม ฯเปฯ. สํเฆ ฯเปฯ. อริยกนฺเตหิ สีเลหิ สมนฺนาคมนํ โหติ อขณฺเฑหิ ฯเปฯ สมาธิ\-สํวตฺตนิเกหิ. อริยกนฺเตหิ สีเลหิ สมนฺนาคมน\-เหตุ โข อาวุโส เอวมิเธกจฺเจ สตฺตา โสตาปนฺนา อวินิปาต\-ธมฺมา นิยตา สมฺโพธิ\-ปรายณาติ.}

\prosewithclosermidruled{สาธุ โข มาริส พุทฺเธ อเวจฺจ\-ปฺปสาเทน สมนฺนาคมนํ โหติ “อิติปิ โส ภควา ฯเปฯ สตฺถา เทวมนุสฺสานํ พุทฺโธ ภควา”ติ, พุทฺเธ อเวจฺจ\-ปฺปสาเทน สมนฺนาคมน\-เหตุ โข มาริส เอวมยํ ปชา โสตาปนฺนา อวินิปาต\-ธมฺมา นิยตา สมฺโพธิ\-ปรายณา. สาธุ โข มาริส ธมฺเม ฯเปฯ. สํเฆ ฯเปฯ. อริยกนฺเตหิ สีเลหิ สมนฺนาคมนํ โหติ อขณฺเฑหิ ฯเปฯ สมาธิ\-สํวตฺตนิเกหิ. อริยกนฺเตหิ สีเลหิ สมนฺนาคมน\-เหตุ โข มาริส เอวมยํ ปชา โสตาปนฺนา อวินิปาต\-ธมฺมา นิยตา สมฺโพธิ\-ปรายณาติ.  ทสมํ.}{ราชการาม\-วคฺโค ทุติโย.}
\csromancenter{\textbf{ตสฺสุทฺทานํ}}
\setlength{\csromangathapagewidth}{0pt}
\setlength{\csromangathawakwidth}{0pt}
\csromangathameasuregroup{สหสฺสพฺราหฺมณานนฺท, \\ มิตฺตามจฺจา ทุเว วุตฺตา,}{\csromangathabat{สหสฺสพฺราหฺมณานนฺท,}{ทุคฺคติ อปเร ทุเว.} \\ \csromangathabat{มิตฺตามจฺจา ทุเว วุตฺตา,}{ตโย จ เทวจาริกาติ.}}
\csromangathasetleft{สหสฺสพฺราหฺมณานนฺท, \\ มิตฺตามจฺจา ทุเว วุตฺตา,}
\csromangathagroup{\csromangathastanza{\csromangathabat{สหสฺสพฺราหฺมณานนฺท,}{ทุคฺคติ อปเร ทุเว.} \\ \csromangathabat{มิตฺตามจฺจา ทุเว วุตฺตา,}{ตโย จ เทวจาริกาติ.}}}

\csromanheader{6}{3. \textbf{สรณานิวคฺค}}
\csromanheader{2}{1. ปฐม\-มหานาม\-สุตฺต}
\proseitem{1017}{เอวํ เม สุตํ\csromandash{}เอกํ สมยํ ภควา สกฺเกสุ วิหรติ กปิล\-วตฺถุสฺมิํ นิโคฺรธาราเม. อถ โข มหานาโม สกฺโก เยน ภควา เตนุปสงฺกมิ, อุปสงฺกมิตฺวา ภควนฺตํ อภิวาเทตฺวา เอกมนฺตํ นิสีทิ, เอกมนฺตํ นิสินฺโน โข มหานาโม สกฺโก ภควนฺตํ เอตทโวจ “อิทํ ภนฺเต กปิลวตฺถุ อิทฺธญฺเจว ผีตญฺจ พาหุชญฺญํ อากิณฺณมนุสฺสํ สมฺพาธพฺยูหํ, โส ขฺวาหํ ภนฺเต ภควนฺตํ วา ปยิรุปาสิตฺวา มโนภาวนีเย วา ภิกฺขู สายนฺหสมยํ กปิลวตฺถุํ ปวิสนฺโต ภนฺเตนปิ\csromansharedfootnote{dcf9f688f0eff639}{วิพฺภนฺเตนปิ (สี), ภมนฺเตนปิ (ก)} หตฺถินา \csromanfolio{322}สมาคจฺฉามิ, ภนฺเตนปิ อสฺเสน สมาคจฺฉามิ, ภนฺเตนปิ รเถน สมาคจฺฉามิ, ภนฺเตนปิ สกเฏน สมาคจฺฉามิ, ภนฺเตนปิ ปุริเสน สมาคจฺฉามิ. ตสฺส มยฺหํ ภนฺเต ตสฺมิํ สมเย มุสฺสเตว\csromansharedfootnote{04031656badf30ec}{มุสเตว (?)} ภควนฺตํ อารพฺภ สติ, มุสฺสติ\csromansharedfootnote{591391576bd0cb06}{มุสติ (?)} ธมฺมํ อารพฺภ สติ, มุสฺสติ สํฆํ อารพฺภ สติ. ตสฺส มยฺหํ ภนฺเต เอวํ โหติ “อิมมฺหิ จาหํ สมเย กาลํ กเรยฺยํ, กา มยฺหํ คติ, โก อภิสมฺปราโย”ติ.}

\prose{มา ภายิ มหานาม, มา ภายิ มหานาม, อปาปกํ เต มรณํ ภวิสฺสติ อปาปิกา กาลํกิริยา\csromansharedfootnote{3238dad5e1bcc411}{กาลกิริยา (สี, สฺยา, กํ)}. ยสฺส กสฺสจิ มหานาม ทีฆรตฺตํ สทฺธา\-ปริภาวิตํ จิตฺตํ, สีลปริภาวิตํ จิตฺตํ, สุตปริภาวิตํ จิตฺตํ, จาค\-ปริภาวิตํ จิตฺตํ, ปญฺญา\-ปริภาวิตํ จิตฺตํ, ตสฺส โย หิ ขฺวายํ กาโย รูปี จาตุมหาภูติโก\csromansharedfootnote{e9d37cc6de879a5d}{จาตุมฺมหาภูติโก (สี, สฺยา, กํ) 5. สิคาลา (สี, สฺยา, กํ, อิ)} มาตา\-เปตฺติก\-สมฺภโว โอทน\-กุมฺมาสู\-ปจโย อนิจฺจุ\-จฺฉาทน\-ปริมทฺทน\-เภทน- วิทฺธํสน\-ธมฺโม. ตํ อิเธว กากา วา ขาทนฺติ, คิชฺฌา วา ขาทนฺติ, กุลลา วา ขาทนฺติ, สุนขา วา ขาทนฺติ, สิงฺคาลา\csromansharedfootnote{59d80dae1024887a}{สิคาลา (สี, สฺยา, กํ, อิ)} วา ขาทนฺติ, วิวิธา วา ปาณกชาตา ขาทนฺติ. ยญฺจ ขฺวสฺส จิตฺตํ ทีฆรตฺตํ สทฺธา\-ปริภาวิตํ ฯเปฯ ปญฺญา\-ปริภาวิตํ, ตํ อุทฺธคามิ โหติ วิเสสคามิ.}

\prose{เสยฺยถาปิ มหานาม ปุริโส สปฺปิกุมฺภํ วา เตลกุมฺภํ วา คมฺภีรํ อุทกรหทํ โอคาหิตฺวา ภินฺเทยฺย. ตตฺร ยา อสฺส สกฺขรา วา กฐลา\csromansharedfootnote{b8ae4d766087ea86}{กถลา (อิ, ก)} วา, สา อโธคามี อสฺส. ยญฺจ ขฺวสฺส ตตฺร สปฺปิ วา เตลํ วา, ตํ อุทฺธคามิ อสฺส วิเสสคามิ. เอวเมว โข มหานาม ยสฺส กสฺสจิ ทีฆรตฺตํ สทฺธา\-ปริภาวิตํ จิตฺตํ ฯเปฯ ปญฺญา\-ปริภาวิตํ จิตฺตํ. ตสฺส โย หิ ขฺวายํ กาโย รูปี จาตุมหาภูติโก มาตา\-เปตฺติก\-สมฺภโว โอทนกุมฺมาสูปวจโย อนิจฺจุ\-จฺฉาทน\-ปริมทฺทน\-เภทน- วิทฺธํสน\-ธมฺโม. ตํ อิเธว กากา วา ขาทนฺติ, คิชฺฌา วา ขาทนฺติ, กุลลา วา ขาทนฺติ, สุนขา วา ขาทนฺติ, สิงฺคาลา วา ขาทนฺติ, วิวิธา วา ปาณกชาตา ขาทนฺติ. ยญฺจ ขฺวสฺส จิตฺตํ ทีฆรตฺตํ สทฺธา\-ปริภาวิตํ ฯเปฯ ปญฺญา\-ปริภาวิตํ, ตํ อุทฺธคามิ โหติ วิเสสคามิ. ตุยฺหํ โข ปน มหานาม ทีฆรตฺตํ สทฺธา\-ปริภาวิตํ จิตฺตํ ฯเปฯ ปญฺญา\-ปริภาวิตํ จิตฺตํ, มา ภายิ}

\prose{\csromanfolio{323}มหานาม, มา ภายิ มหานาม, อปาปกํ เต มรณํ ภวิสฺสติ อปาปิกา กาลํกิริยาติ. ปฐมํ.}

\csromanheader{2}{2. ทุติย\-มหานาม\-สุตฺต}
\proseitem{1018}{เอวํ เม สุตํ\csromandash{}เอกํ สมยํ ภควา สกฺเกสุ วิหรติ กปิล\-วตฺถุสฺมิํ นิโคฺรธาราเม. อถ โข มหานาโม สกฺโก เยน ภควา เตนุปสงฺกมิ, อุปสงฺกมิตฺวา ภควนฺตํ อภิวาเทตฺวา เอกมนฺตํ นิสีทิ, เอกมนฺตํ นิสินฺโน โข มหานาโม สกฺโก ภควนฺตํ เอตทโวจ “อิทํ ภนฺเต กปิลวตฺถุ อิทฺธญฺเจว ผีตญฺจ พาหุชญฺญํ อากิณฺณมนุสฺสํ สมฺพาธพฺยูหํ, โส ขฺวาหํ ภนฺเต ภควนฺตํ วา ปยิรุปาสิตฺวา มโนภาวนีเย วา ภิกฺขู สายนฺหสมยํ กปิลวตฺถุํ ปวิสนฺโต ภนฺเตนปิ หตฺถินา สมาคจฺฉามิ, ภนฺเตนปิ อสฺเสน สมาคจฺฉามิ, ภนฺเตนปิ รเถน สมาคจฺฉามิ, ภนฺเตนปิ สกเฏน สมาคจฺฉามิ, ภนฺเตนปิ ปุริเสน สมาคจฺฉามิ, ตสฺส มยฺหํ ภนฺเต ตสฺมิํ สมเย มุสฺสเตว ภควนฺตํ อารพฺภ สติ, มุสฺสติ ธมฺมํ อารพฺภ สติ, มุสฺสติ สํฆํ อารพฺภ สติ. ตสฺส มยฺหํ ภนฺเต เอวํ โหติ “อิมมฺหิ จาหํ สมเย กาลํกเรยฺยํ, กา มยฺหํ คติ, โก อภิสมฺปราโย”ติ.}

\prose{มา ภายิ มหานาม, มา ภายิ มหานาม, อปาปกํ เต มรณํ ภวิสฺสติ อปาปิกา กาลํกิริยา. จตูหิ โข มหานาม ธมฺเมหิ สมนฺนาคโต อริยสาวโก นิพฺพานนินฺโน โหติ นิพฺพานโปโณ นิพฺพาน\-ปพฺภาโร. กตเมหิ จตูหิ, อิธ มหานาม อริยสาวโก พุทฺเธ อเวจฺจ\-ปฺปสาเทน สมนฺนาคโต โหติ “อิติปิ โส ภควา ฯเปฯ สตฺถา เทวมนุสฺสานํ พุทฺโธ ภควา”ติ. ธมฺเม ฯเปฯ. สํเฆ ฯเปฯ. อริยกนฺเตหิ สีเลหิ สมนฺนาคโต โหติ อขณฺเฑหิ ฯเปฯ สมาธิ\-สํวตฺตนิเกหิ.}

\prose{เสยฺยถาปิ มหานาม รุกฺโข ปาจีนนินฺโน ปาจีนโปโณ ปาจีนปพฺภาโร, โส มูลจฺฉินฺโน กตเมน ปปเตยฺยาติ. เยน ภนฺเต นินฺโน เยน โปโณ เยน ปพฺภาโรติ. เอวเมว โข มหานาม อิเมหิ จตูหิ ธมฺเมหิ สมนฺนาคโต อริยสาวโก นิพฺพานนินฺโน โหติ นิพฺพานโปโณ นิพฺพาน\-ปพฺภาโรติ. ทุติยํ.}

\csromanmatikaanchor{mtk.238}
\csromantocmarkref{subparagraph}{3. โคธสกฺกสุตฺต}{mtk.238}
\bookmark[dest={mtk.238},level=5]{3. โคธสกฺกสุตฺต}
\csromanheader{6}{3. โคธ\-สกฺก\-สุตฺต}
\proseitem{1019}{\csromanfolio{324}กปิลวตฺถุ\-นิทานํ. อถ โข มหานาโม สกฺโก เยน โคธา สกฺโก เตนุปสงฺกมิ, อุปสงฺกมิตฺวา โคธํ สกฺกํ เอตทโวจ “กติหิ\csromansharedfootnote{1b488688f32b7071}{กตีหิ (อิ, ก) รูปสิทฺธิ โอโลเกตพฺพา.} ตฺวํ โคเธ ธมฺเมหิ สมนฺนาคตํ โสตาปนฺน\-ปุคฺคลํ อาชานาสิ อวินิปาต\-ธมฺมํ นิยตํ สมฺโพธิ\-ปรายณนฺติ.}

\prose{ตีหิ ขฺวาหํ มหานาม ธมฺเมหิ สมนฺนาคตํ โสตาปนฺน\-ปุคฺคลํ อาชานามิ อวินิปาต\-ธมฺมํ นิยตํ สมฺโพธิ\-ปรายณํ. กตเมหิ ตีหิ, อิธ มหานาม อริยสาวโก พุทฺเธ อเวจฺจ\-ปฺปสาเทน สมนฺนาคโต โหติ “อิติปิ โส ภควา ฯเปฯ สตฺถา เทวมนุสฺสานํ พุทฺโธ ภควา”ติ. ธมฺเม ฯเปฯ. สํเฆ อเวจฺจ\-ปฺปสาเทน สมนฺนาคโต โหติ “สุปฺปฏิปนฺโน ภควโต สาวกสํโฆ ฯเปฯ อนุตฺตรํ ปุญฺญกฺเขตฺตํ โลกสฺสา”ติ. อิเมหิ ขฺวาหํ มหานาม ตีหิ ธมฺเมหิ สมนฺนาคตํ โสตาปนฺน\-ปุคฺคลํ อาชานามิ อวินิปาต\-ธมฺมํ นิยตํ สมฺโพธิ\-ปรายณํ.}

\prose{ตฺวํ ปน มหานาม กติหิ ธมฺเมหิ สมนฺนาคตํ โสตาปนฺน\-ปุคฺคลํ อาชานาสิ อวินิปาต\-ธมฺมํ นิยตํ สมฺโพธิ\-ปรายณนฺติ. จตูหิ ขฺวาหํ โคเธ ธมฺเมหิ สมนฺนาคตํ โสตาปนฺน\-ปุคฺคลํ อาชานามิ อวินิปาต\-ธมฺมํ นิยตํ สมฺโพธิ\-ปรายณํ. กตเมหิ จตูหิ, อิธ โคเธ อริยสาวโก พุทฺเธ อเวจฺจ\-ปฺปสาเทน สมนฺนาคโต โหติ “อิติปิ โส ภควา ฯเปฯ สตฺถา เทวมนุสฺสานํ พุทฺโธ ภควา”ติ. ธมฺเม ฯเปฯ. สํเฆ ฯเปฯ. อริยกนฺเตหิ สีเลหิ สมนฺนาคโต โหติ อขณฺเฑหิ ฯเปฯ สมาธิ\-สํวตฺตนิเกหิ. อิเมหิ ขฺวาหํ โคเธ จตูหิ ธมฺเมหิ สมนฺนาคตํ โสตาปนฺน\-ปุคฺคลํ อาชานามิ อวินิปาต\-ธมฺมํ นิยตํ สมฺโพธิ\-ปรายณนฺติ.}

\prose{อาคเมหิ ตฺวํ มหานาม, อาคเมหิ ตฺวํ มหานาม, ภควาว เอตํ ชาเนยฺย เอเตหิ ธมฺเมหิ สมนฺนาคตํ วา อสมนฺนาคตํ วาติ, อายาม โคเธ, เยน ภควา เตนุปสงฺกเมยฺยาม, อุปสงฺกมิตฺวา ภควโต เอตมตฺถํ อาโรเจสฺสามาติ. อถ โข มหานาโม สกฺโก โคธา จ สกฺโก เยน ภควา เตนุ\-ปสงฺกมิํสุ, อุปสงฺกมิตฺวา ภควนฺตํ อภิวาเทตฺวา เอกมนฺตํ นิสีทิํสุ, เอกมนฺตํ นิสินฺโน โข มหานาโม สกฺโก ภควนฺตํ เอตทโวจ\csromandash{}}

\prose{\csromanfolio{325}อิธาหํ ภนฺเต เยน โคธา สกฺโก เตนุปสงฺกมิํ, อุปสงฺกมิตฺวา โคธํ สกฺกํ เอตทโวจํ “กติหิ ตฺวํ โคเธ ธมฺเมหิ สมนฺนาคตํ โสตาปนฺน\-ปุคฺคลํ อาชานาสิ อวินิปาต\-ธมฺมํ นิยตํ สมฺโพธิ\-ปรายณํ”. เอวํ วุตฺเต พนฺเต โคธา สกฺโก มํ เอตทโวจ\csromandash{}}

\prose{ตีหิ ขฺวาหํ มหานาม ธมฺเมหิ สมนฺนาคตํ โสตาปนฺน\-ปุคฺคลํ อาชานามิ อวินิปาต\-ธมฺมํ นิยตํ สมฺโพธิ\-ปรายณํ. กตเมหิ ตีหิ, อิธ มหานาม อริยสาวโก พุทฺเธ อเวจฺจ\-ปฺปสาเทน สมนฺนาคโต โหติ “อิติปิ โส ภควา ฯเปฯ สตฺถา เทวมนุสฺสานํ พุโธ ภควา”ติ. ธมฺเม ฯเปฯ. สํเฆ อเวจฺจ\-ปฺปสาเทน สมนฺนาคโต โหติ “สุปฺปฏิปนฺโน ภควโต สาวกสํโฆ ฯเปฯ อนุตฺตรํ ปุญฺญกฺเขตฺตํ โลกสฺสา”ติ. อิเมหิ ขฺวาหํ มหานาม ตีหิ ธมฺเมหิ สมนฺนาคตํ โสตาปนฺน\-ปุคฺคลํ อาชานามิ อวินิปาต\-ธมฺมํ นิยตํ สมฺโพธิ\-ปรายณํ. ตฺวํ ปน มหานาม กตเมหิ ธมฺเมหิ สมนฺนาคตํ โสตาปนฺน\-ปุคฺคลํ อาชานาสิ อวินิปาต\-ธมฺมํ นิยตํ สมฺโพธิ\-ปรายณนฺติ.}

\prose{เอวํ วุตฺตาหํ ภนฺเต โคธํ สกฺกํ เอตทโวจํ “จตูหิ ขฺวาหํ โคเธ ธมฺเมหิ สมนฺนาคตํ โสตาปนฺน\-ปุคฺคลํ อาชานามิ อวินิปาต\-ธมฺมํ นิยตํ สมฺโพธิ\-ปรายณํ. กตเมหิ จตูหิ, อิธ โคเธ อริยสาวโก พุทฺเธ อเวจฺจ\-ปฺปสาเทน สมนฺนาคโต โหติ ‘อิติปิ โส ภควา ฯเปฯ สตฺถา เทวมนุสฺสานํ พุทฺโธ ภควา’ติ. ธมฺเม ฯเปฯ. สํเฆ ฯเปฯ. อริยกนฺเตหิ สีเลหิ สมนฺนาคโต โหติ อขณฺเฑหิ ฯเปฯ สมาธิ\-สํวตฺตนิเกหิ. อิเมหิ ขฺวาหํ โคเธ จตูหิ ธมฺเมหิ สมนฺนาคตํ โสตาปนฺน\-ปุคฺคลํ อาชานามิ อวินิปาต\-ธมฺมํ นิยตํ สมฺโพธิปรายณนฺ”ติ.}

\prose{เอวํ วุตฺเต ภนฺเต โคธา สกฺโก มํ เอตทโวจ “อาคเมหิ ตฺวํ มหานาม, อาคเมหิ ตฺวํ มหานาม, ภควาว เอตํ ชาเนยฺย เอเตหิ ธมฺเมหิ สมนฺนาคตํ วา อสมนฺนาคตํ วา”ติ, อิธ ภนฺเต โกจิเทว ธมฺโม สมุปฺปาโท อุปฺปชฺเชยฺย, เอกโต อสฺส ภควา เอกโต ภิกฺขุสํโฆ จ, เยเนว ภควา, เตเนวาหํ อสฺสํ, เอวํ ปสนฺนํ มํ ภนฺเต ภควา ธาเรตุ, อิธ ภนฺเต โกจิเทว ภมฺโม สมุปฺปาโท อุปฺปชฺเชยฺย, เอกโต อสฺส ภควา, เอกโต ภิกฺขุสํโฆ ภิกฺขุนิ\-สํโฆ จ, เยเนว ภควา, เตเนวาหํ อสฺสํ, เอวํ ปสนฺนํ มํ ภนฺเต ภควา ธาเรตุ. อิธ ภนฺเต โกจิเทว ธมฺโม สมุปฺปาโท อุปฺปชฺเชยฺย, เอกโต อสฺส ภควา, เอกโต}

\csromanmatikaanchor{mtk.239}
\csromantocmarkref{subparagraph}{4-5. สรณานิสกฺกสุตฺตานิ}{mtk.239}
\bookmark[dest={mtk.239},level=5]{4-5. สรณานิสกฺกสุตฺตานิ}
\prose{\csromanfolio{326}ภิกฺขุสํโฆ ภิกฺขุนิ\-สํโฆ จ อุปาสกา จ, เยเนว ภควา, เตเนวาหํ อสฺสํ, เอวํ ปสนฺนํ มํ ภนฺเต ภควา ธาเรตุ, อิธ ภนฺเต โกจิเทว ธมฺโม สมุปฺปาโท อุปฺปชฺเชยฺย, เอกโต อสฺส ภควา, เอกโต ภิกฺขุสํโฆ ภิกฺขุนิ\-สํโฆ อุปาสกา อุปาสิกาโย จ, เยเนว ภควา, เตเนวาหํ อสฺสํ, เอวํ ปสนฺนํ มํ ภนฺเต ภควา ธาเรตุ, อิธ ภนฺเต โกจิเทว ธมฺโม สมุปฺปาโท อุปฺปชฺเชยฺย, เอกโต อสฺส ภควา, เอกโต ภิกฺขุสํโฆ ภิกฺขุนิ\-สํโฆ อุปาสกา อุปาสิกาโย สเทวโก จ โลโก สมารโก สพฺรหฺมโก สสฺสมณ\-พฺราหฺมณี ปชา สเทวมนุสฺสา, เยเนว ภควา, เตเนวาหํ อสฺสํ, เอวํ ปสนฺนํ มํ ภนฺเต ภควา ธาเรตูติ. เอวํวาที ตฺวํ โคเธ มหานามํ สกฺกํ กิํ วเทสีติ. เอวํวาทาหํ ภนฺเต มหานามํ สกฺกํ น กิญฺจิ วทามิ อญฺญตฺร กลฺยาณา อญฺญตฺร กุสลาติ. ตติยํ.}

\csromanheader{2}{4. ปฐม\-สรณานิ\-สกฺก\-สุตฺต}
\proseitem{1020}{กปิลวตฺถุ\-นิทานํ. เตน โข ปน สมเยน สรณานิ\csromansharedfootnote{d08a4d9a35753d7a}{สรกานิ (สี, สฺยา, กํ, อิ)} สกฺโก กาลงฺกโต โหติ, โส ภควตา พฺยากโต “โสตาปนฺโน อวินิปาต\-ธมฺโม นิยโต สมฺโพธิ\-ปรายโณ”ติ. ตตฺร สุทํ สมฺพหุลา สกฺกา สงฺคมฺม สมาคมฺม อุชฺฌายนฺติ ขียนฺติ วิปาเจนฺติ “อจฺฉริยํ วต โภ, อพฺภุตํ วต โภ, เอตฺถ ทานิ โก น โสตาปนฺโน ภวิสฺสติ, ยตฺร หิ นาม สรณานิ สกฺโก กาลงฺกโต, โส ภควตา พฺยากโต ‘โสตาปนฺโน อวินิปาต\-ธมฺโม นิยโต สมฺโพธิ\-ปรายโณ’ติ. สรณานิ สกฺโก สิกฺขาทุพฺพลฺยมา\-ปาทิ, มชฺชปานํ อปายี”ติ.}

\prose{อถ โข มหานาโม สกฺโก เยน ภควา เตนุปสงฺกมิ, อุปสงฺกมิตฺวา ภควนฺตํ อภิวาเทตฺวา เอกมนฺตํ นิสีทิ, เอกมนฺตํ นิสินฺโน โข มหานาโม สกฺโก ภควนฺตํ เอตทโวจ\csromandash{}อิธ ภนฺเต สรณานิ สกฺโก กาลงฺกโต, โส ภควตา พฺยากโต “โสตาปนฺโน อวินิปาต\-ธมฺโม นิยโต สมฺโพธิ\-ปรายโณ”ติ. ตตฺร สุทํ ภนฺเต สมฺพหุลา สกฺกา สงฺคมฺม สมาคมฺม อุชฺฌายนฺติ ขียนฺติ วิปาเจนฺติ “อจฺฉริยํวต โภ, อพฺภุตํ วต โภ, เอตฺถ ทานิ โก น โสตาปนฺโน ภวิสฺสติ, ยตฺร หิ นาม สรณานิ สกฺโก กาลงฺกโต, โส ภควตา}

\prose{\csromanfolio{327}พฺยากโต ‘โสตาปนฺโน อวินิปาต\-ธมฺโม นิยโต สมฺโพธิ\-ปรายโณ’ติ, สรณานิ สกฺโก สิกฺขาทุพฺพลฺยมา\-ปาทิ, มชฺชปานํ อปายี”ติ.}

\prose{โย โส มหานาม ทีฆรตฺตํ อุปาสโก พุทฺธํ สรณํ คโต, ธมฺมํ สรณํ คโต, สํฆํ สรณํ คโต, โส กถํ วินิปาตํ คจฺเฉยฺย. ยํ หิ ตํ มหานาม สมฺมา วทมาโน วเทยฺย “ทีฆรตฺตํ อุปาสโก พุทฺธํ สรณํ คโต, ธมฺมํ สรณํ คโต, สํฆํ สรณํ คโต”ติ. สรณานิ สกฺกํ สมฺมา วทมาโน วเทยฺย, สรณานิ มหานาม สกฺโก ทีฆรตฺตํ อุปาสโก พุทฺธํ สรณํ คโต, ธมฺมํ สรณํ คโต, สํฆํ สรณํ คโต, โส กถํ วินิปาตํ คจฺเฉยฺย.}

\prose{อิธ มหานาม เอกจฺโจ ปุคฺคโล พุทฺเธ อเวจฺจ\-ปฺปสาเทน สมนฺนาคโต โหติ “อิติปิ โส ภควา ฯเปฯ สตฺถา เทวมนุสฺสานํ พุทฺโธ ภควา”ติ. ธมฺเม ฯเปฯ. สํเฆ ฯเปฯ หาสปญฺโญ, ชวนปญฺโญ, วิมุตฺติยา จ สมนฺนาคโต, โส อาสวานํ ขยา อนาสวํ เจโตวิมุตฺติํ ปญฺญาวิมุตฺติํ ทิฏฺเฐว ธมฺเม สยํ อภิญฺญา สจฺฉิกตฺวา อุปสมฺปชฺช วิหรติ. อยมฺปิ โข มหานาม ปุคฺคโล ปริมุตฺโต นิรยา, ปริมุตฺโต ติรจฺฉาน\-โยนิยา, ปริมุตฺโต เปตฺติวิสยา, ปริมุตฺโต อปาย\-ทุคฺคติ\-วินิปาตา.}

\prose{อิธ ปน มหานาม เอกจฺโจ ปุคฺคโล พุทฺเธ อเวจฺจ\-ปฺปสาเทน สมนฺนาคโต โหติ “อิติปิ โส ภควา ฯเปฯ สตฺถา เทวมนุสฺสานํ พุทฺโธ ภควา”ติ. ธมฺเม ฯเปฯ. สํเฆ ฯเปฯ หาสปญฺโญ, ชวนปญฺโญ, น จ วิมุตฺติยา สมนฺนาคโต, โส ปญฺจนฺนํ โอรมฺภาคิยานํ สํโยชนานํ ปริกฺขยา โอปปาติโก โหติ ตตฺถ ปรินิพฺพายี อนาวตฺติธมฺโม ตสฺมา\csromansharedfootnote{3dc8f1594152d076}{อสฺมา (สฺยา, กํ, อิ, ก)} โลกา, อยมฺปิ โข มหานาม ปุคฺคโล ปริมุตฺโต นิรยา, ปริมุตฺโต ติรจฺฉาน\-โยนิยา, ปริมุตฺโต เปตฺติวิสยา, ปริมุตฺโต อปาย\-ทุคฺคติ\-วินิปาตา.}

\prose{“อิธ ปน มหานาม เอกจฺโจ ปุคฺคโล พุทฺเธ อเวจฺจ\-ปฺปสาเทน สมนฺนาคโต โหติ อิติปิ โส ภควา ฯเปฯ สตฺถา เทวมนุสฺสานํ พุทฺโธ ภควา”ติ. ธมฺเม ฯเปฯ. สํเฆ ฯเปฯ น หาสปญฺโญ, น ชวนปญฺโญ, น จ วิมุตฺติยา สมนฺนาคโต, โส ติณฺณํ สํโยชนานํ ปริกฺขยา ราค\-โทส\-โมหานํ ตนุตฺตา สกทาคามี โหติ, สกิเทว อิมํ โลกํ อาคนฺตฺวา ทุกฺขสฺสนฺตํ}

\prose{\csromanfolio{328}กโรติ, อยมฺปิ โข มหานาม ปุคฺคโล ปริมุตฺโต นิรยา, ปริมุตฺโต ติรจฺฉาน\-โยนิยา, ปริมุตฺโต เปตฺติวิสยา, ปริมุตฺโต อปาย\-ทุคฺคติ\-วินิปาตา.}

\prose{อิธ ปน มหานาม เอกจฺโจ ปุคฺคโล พุทฺเธ อเวจฺจ\-ปฺปสาเทน สมนฺนาคโต โหติ “อิติปิ โส ภควา ฯเปฯ สตฺถา เทวมนุสฺสานํ พุทฺโธ ภควา”ติ. ธมฺเม ฯเปฯ. สํเฆ ฯเปฯ น หาสปญฺโญ, น ชวนปญฺโญ, น จ วิมุตฺติยา สมนฺนาคโต, โส ติณฺณํ สํโยชนานํ ปริกฺขยา โสตาปนฺโน โหติ อวินิปาต\-ธมฺโม นิยโต สมฺโพธิ\-ปรายโณติ. อยมฺปิ โข มหานาม ปุคฺคโล ปริมุตฺโต นิรยา, ปริมุตฺโต ติรจฺฉาน\-โยนิยา, ปริมุตฺโต เปตฺติวิสยา, ปริมุตฺโต อปาย\-ทุคฺคติ\-วินิปาตา.}

\prose{อิธ ปน มหานาม เอกจฺโจ ปุคฺคโล น เหว โข พุทฺเธ อเวจฺจ\-ปฺปสาเทน สมนฺนาคโต โหติ. น ธมฺเม ฯเปฯ. น สํเฆ ฯเปฯ น หาสปญฺโญ, น ชวนปญฺโญ, น จ วิมุตฺติยา สมนฺนาคโต. อปิ จสฺส อิเม ธมฺมา โหนฺติ สทฺธินฺทฺริยํ วีริยินฺทฺริยํ สตินฺทฺริยํ สมาธินฺทฺริยํ ปญฺญินฺทฺริยํ, ตถาคต\-ปฺปเวทิตา จสฺส ธมฺมา ปญฺญาย มตฺตโส นิชฺฌานํ ขมนฺติ, อยมฺปิ โข มหานาม ปุคฺคโล อคนฺตา นิรยํ อคนฺตา ติรจฺฉาน\-โยนิํ อคนฺตา เปตฺติวิสยํ อคนฺตา อปายํ ทุคฺคติํ วินิปาตํ.}

\prose{อิธ ปน มหานาม เอกจฺโจ ปุคฺคโล น เหว โข พุทฺเธ อเวจฺจ\-ปฺปสาเทน สมนฺนาคโต โหติ. น ธมฺเม ฯเปฯ. น สํเฆ ฯเปฯ น หาสปญฺโญ, น ชวนปญฺโญ, น จ วิมุตฺติยา สมนฺนาคโต. อปิ จสฺส อิเม ธมฺมา โหนฺติ สทฺธินฺทฺริยํ ฯเปฯ ปญฺญินฺทฺริยํ. ตถาคเต จสฺส สทฺธามตฺตํ โหติ เปมมตฺตํ. อยมฺปิ โข มหานาม ปุคฺคโล อคนฺตา นิรยํ อคนฺตา ติรจฺฉาน\-โยนิํ อคนฺตา เปตฺติวิสยํ อคนฺตา อปายํ ทุคฺคติํ วินิปาตํ. อิเม เจปิ มหานาม มหาสาลา สุภาสิตํ ทุพฺภาสิตํ อาชาเนยฺยุํ, อิเมจาหํ\csromansharedfootnote{a11db8a486754553}{อิเมวาหํ (สฺยา, กํ), อิเมสาหํ (ก)} มหาสาเล พฺยากเรยฺยํ “โสตาปนฺนา อวินิปาต\-ธมฺมา นิยตา สมฺโพธิ\-ปรายณา”ติ, กิมงฺคํ\csromansharedfootnote{cbf4b6dfaa995165}{กิมงฺค (สี, สฺยา, กํ, อิ)} ปน สรณานิํ สกฺกํ, สรณานิ มหานาม สกฺโก มรณกาเล สิกฺขํ สมาทิยีติ.  จตุตฺถํ.}

\csromanheader{6}{11. โสตาปตฺติ\-สํยุตฺต}
\csromanheader{2}{5. ทุติย\-สรณานิ\-สกฺก\-สุตฺต}
\proseitem{1021}{\csromanfolio{329}กปิลวตฺถุ\-นิทานํ. เตน โข ปน สมเยน สรณานิ สกฺโก กาลงฺกโต โหติ. โส ภควตา พฺยากโต “โสตาปนฺโน อวินิปาต\-ธมฺโม นิยโต สมฺโพธิ\-ปรายโณ”ติ. ตตฺร สุทํ สมฺพหุลา สกฺกา สงฺคมฺม สมาคมฺม อุชฺฌายนฺติ ขียนฺติ วิปาเจนฺติ “อจฺฉริยํ วต โภ, อพฺภุตํ วต โภ, เอตฺถ ทานิ โก น โสตาปนฺโน ภวิสฺสติ, ยตฺร หิ นาม สรณานิ สกฺโก กาลงฺกโต, โส ภควตา พฺยากโต ‘โสตาปนฺโน อวินิปาต\-ธมฺโม นิยโต สมฺโพธิ\-ปรายโณ’ติ. สรณานิ สกฺโก สิกฺขาย อปริปูรการี อโหสี”ติ. อถ โข มหานาโม สกฺโก เยน ภควา เตนุปสงฺกมิ, อุปสงฺกมิตฺวา ภควนฺตํ อภิวาเทตฺวา เอกมนฺตํ นิสีทิ, เอกมนฺตํ นิสินฺโน โข มหานาโม สกฺโก ภควนฺตํ เอตทโวจ\csromandash{}}

\prose{“อิธ ภนฺเต สรณานิ สกฺโก กาลงฺกโต, โส ภควตา พฺยากโต ‘โสตาปนฺโน อวินิปาต\-ธมฺโม นิยโต สมฺโพธิ\-ปรายโณ’ติ. ตตฺร สุทํ ภนฺเต สมฺพหุลา สกฺกา สงฺคมฺม สมาคมฺม อุชฺฌายนฺติ ขียนฺติ วิปาเจนฺติ ‘อจฺฉริยํ วต โภ, อพฺภุตํ วต โภ, เอตฺถ ทานิ โก น โสตาปนฺโน ภวิสฺสติ, ยตฺร หิ นาม สรณานิ สกฺโก กาลงฺกโต, โส ภควตา พฺยากโต ‘โสตาปนฺโน อวินิปาต\-ธมฺโม นิยโต สมฺโพธิ\-ปรายโณ’ติ, สรณานิ สกฺโก สิกฺขาย อปริปูรการี อโหสี”ติ.}

\prose{โย โส มหานาม ทีฆรตฺตํ อุปาสโก พุทฺธํ สรณํ คโต, ธมฺมํ สรณํ คโต, สํฆํ สรณํ คโต, โส กถํ วินิปาตํ คจฺเฉยฺย. ยํ หิ ตํ มหานาม สมฺมา วทมาโน วเทยฺย “ทีฆรตฺตํ อุปาสโก พุทฺธํ สรณํ คโต, ธมฺมํ สรณํ คโต, สํฆํ สรณํ คโต”. สรณานิํ สกฺกํ สมฺมา วทมาโน วเทยฺย “สรณานิ มหานาม สกฺโก ทีฆรตฺตํ อุปาสโก พุทฺธํ สรณํ คโต, ธมฺมํ สรณํ คโต, สํฆํ สรณํ คโต, โส กถํ วินิปาตํ คจฺเฉยฺย”.}

\prose{อิธ มหานาม เอกจฺโจ ปุคฺคโล พุทฺเธ เอกนฺตคโต โหติ อภิปฺปสนฺโน “อิติปิ โส ภควา ฯเปฯ สตฺถา เทวมนุสฺสานํ พุทฺโธ ภควา”ติ. ธมฺเม ฯเปฯ. สํเฆ ฯเปฯ หาสปญฺโญ, ชวนปญฺโญ, วิมุตฺติยา จ \csromanfolio{330}สมนฺนาคโต, โส อาสวานํ ขยา อนาสวํ เจโตวิมุตฺติํ ปญฺญาวิมุตฺติํ ทิฏฺเฐว ธมฺเม สยํ อภิญฺญา สจฺฉิกตฺวา อุปสมฺปชฺช วิหรติ, อยมฺปิ โข มหานาม ปุคฺคโล ปริมุตฺโต นิรยา, ปริมุตฺโต ติรจฺฉาน\-โยนิยา, ปริมุตฺโต เปตฺติวิสยา, ปริมุตฺโต อปาย\-ทุคฺคติ\-วินิปาตา.}

\prose{อิธ ปน มหานาม เอกจฺโจ ปุคฺคโล พุทฺเธ เอกนฺตคโต โหติ อภิปฺปสนฺโน “อิติปิ โส ภควา ฯเปฯ สตฺถา เทวมนุสฺสานํ พุทฺโธ ภควา”ติ. ธมฺเม ฯเปฯ. สํเฆ ฯเปฯ หาสปญฺโญ, ชวนปญฺโญ, น จ วิมุตฺติยา สมนฺนาคโต, โส ปญฺจนฺนํ โอรมฺภาคิยานํ สํโยชนานํ ปริกฺขยา อนฺตรา\-ปรินิพฺพายี โหติ, อุปหจฺจ\-ปรินิพฺพายี โหติ, อสงฺขารปรินิพฺพายี โหติ, สสงฺขาร\-ปรินิพฺพายี โหติ, อุทฺธํโสโต โหติ อกนิฏฺฐคามี. อยมฺปิ โข มหานาม ปุคฺคโล ปริมุตฺโต นิรยา, ปริมุตฺโต ติรจฺฉาน\-โยนิยา, ปริมุตฺโต เปตฺติวิสยา, ปริมุตฺโต อปาย\-ทุคฺคติ\-วินิปาตา.}

\prose{อิธ ปน มหานาม เอกจฺโจ ปุคฺคโล พุทฺเธ เอกนฺตคโต โหติ อภิปฺปสนฺโน “อิติปิ โส ภควา ฯเปฯ สตฺถา เทวมนุสฺสานํ พุทฺโธ ภควา”ติ. ธมฺเม ฯเปฯ. สํเฆ ฯเปฯ น หาสปญฺโญ, น ชวนปญฺโญ, น จ วิมุตฺติยา สมนฺนาคโต, โส ติณฺณํ สํโยชนานํ ปริกฺขยา ราค\-โทส\-โมหานํ ตนุตฺตา สกทาคามี โหติ, สกิเทว อิมํ โลกํ อาคนฺตฺวา ทุกฺขสฺสนฺตํ กโรติ. อยมฺปิ โข มหานาม ปุคฺคโล ปริมุตฺโต นิรยา, ปริมุตฺโต ติรจฺฉาน\-โยนิยา, ปริมุตฺโต เปตฺติวิสยา, ปริมุตฺโต อปาย\-ทุคฺคติ\-วินิปาตา.}

\prose{อิธ ปน มหานาม เอกจฺโจ ปุคฺคโล พุทฺเธ เอกนฺตคโต โหติ อภิปฺปสนฺโน “อิติปิ โส ภควา ฯเปฯ สตฺถา เทวมนุสฺสานํ พุทฺโธ ภควา”ติ. ธมฺเม ฯเปฯ. สํเฆ ฯเปฯ น หาสปญฺโญ, น ชวนปญฺโญ, น จ วิมุตฺติยา สมนฺนาคโต. โส ติณฺณํ สํโยชนานํ ปริกฺขยา โสตาปนฺโน โหติ อวินิปาต\-ธมฺโม นิยโต สมฺโพธิ\-ปรายโณ, อยมฺปิ โข มหานาม ปุคฺคโล ปริมุตฺโต นิรยา, ปริมุตฺโต ติรจฺฉาน\-โยนิยา, ปริมุตฺโต เปตฺติวิสยา, ปริมุตฺโต อปาย\-ทุคฺคติ\-วินิปาตา.}

\prose{อิธ ปน มหานาม เอกจฺโจ ปุคฺคโล น เหว โข พุทฺเธ เอกนฺตคโต โหติ อภิปฺปสนฺโน ฯเปฯ. น ธมฺเม ฯเปฯ. น สํเฆ ฯเปฯ น หาสปญฺโญ, น ชวนปญฺโญ, น จ วิมุตฺติยา สมนฺนาคโต, อปิ จสฺส อิเม ธมฺมา โหนฺติ}

\prose{\csromanfolio{331}สทฺธินฺทฺริยํ ฯเปฯ ปญฺญินฺทฺริยํ, ตถาคต\-ปฺปเวทิตา จสฺส ธมฺมา ปญฺญาย มตฺตโส นิชฺฌานํ ขมนฺติ. อยมฺปิ โข มหานาม ปุคฺคโล อคนฺตา นิรยํ อคนฺตา ติรจฺฉาน\-โยนิํ อคนฺตา เปตฺติวิสยํ อคนฺตา อปายํ ทุคฺคติํ วินิปาตํ.}

\prose{อิธ ปน มหานาม เอกจฺโจ ปุคฺคโล น เหว โข พุทฺเธ เอกนฺตคโต โหติ อภิปฺปสนฺโน. น ธมฺเม ฯเปฯ. น สํเฆ ฯเปฯ น หาสปญฺโญ, น ชวนปญฺโญ, น จ วิมุตฺติยา สมนฺนาคโต. อปิ จสฺส อิเม ธมฺมา โหนฺติ สทฺธินฺทฺริยํ ฯเปฯ ปญฺญินฺทฺริยํ, ตถาคเต จสฺส สทฺธามตฺตํ โหติ เปมมตฺตํ. อยมฺปิ โข มหานาม ปุคฺคโล อคนฺตา นิรยํ อคนฺตา ติรจฺฉาน\-โยนิํ อคนฺตา เปตฺติวิสยํ อคนฺตา อปายํ ทุคฺคติํ วินิปาตํ.}

\prose{เสยฺยถาปิ มหานาม ทุกฺเขตฺตํ ทุพฺภูมํ อวิหต\-ขาณุกํ, พีชานิ จสฺสุ ขณฺฑานิ ปูตีนิ วาตาตปหตานิ อสาราทานิ อสุขสยิตานิ\csromansharedfootnote{296e1680254537fb}{อสุขาปสฺสยิตานิ (ก)}, เทโว จ น สมฺมา\csromansharedfootnote{8b02074a4be903ba}{เทโว ปน สมฺมา (สฺยา, กํ), เทโว น สมฺมา (ก) ที 2. 279 ปิฏฺเฐปิ.} ธารํ อนุปฺปเวจฺเฉยฺย, อปิ นุ ตานิ พีชานิ วุทฺธิํ วุรูฬฺหิํ เวปุลฺลํ อาปชฺเชยฺยุนฺติ. โน เหตํ ภนฺเต. เอวเมว โข มหานาม อิธ ธมฺโม ทุรกฺขาโต\csromansharedfootnote{759cab869384ceaa}{ทฺวากฺขาโต (อิ, ก)} โหติ ทุปฺปเวทิโต อนิยฺยานิโก อนุปสมสํวตฺตนิโก อสมฺมาสมฺพุทฺธปฺปเวทิโต. อิทมหํ ทุกฺเขตฺตสฺมิํ วทามิ, ตสฺมิํ จ ธมฺเม สาวโก วิหรติ ธมฺมา\-นุธมฺม\-ปฺปฏิปนฺโน สามีจิ\-ปฺปฏิปนฺโน อนุธมฺมจารี, อิทมหํ ทุพฺพีชสฺมิํ วทามิ.}

\prose{เสยฺยถาปิ มหานาม สุเขตฺตํ สุภูมํ สุวิหต\-ขาณุกํ, พีชานิ จสฺสุ อขณฺฑานิ อปูตีนิ อวาตาตปหตานิ สาราทานิ สุขสยิตานิ, เทโว จ\csromansharedfootnote{36f0eaa7c9b7601f}{เทโว จสฺส (สฺยา, กํ, ก)} สมฺมา ธารํ อนุปฺปเวจฺเฉยฺย. อปิ นุ ตานิ พีชานิ วุทฺธิํ วิรูฬฺหิํ เวปุลฺลํ อาปชฺเชยฺยุนฺติ. เอวํ ภนฺเต. เอวเมว โข มหานาม อิธ ธมฺโม สฺวากฺขาโต โหติ สุปฺปเวทิโต นิยฺยานิโก อุปสม\-สํวตฺตนิโก สมฺมาสมฺพุทฺธ\-ปฺปเวทิโต, อิทมหํ สุเขตฺตสฺมิํ วทามิ, ตสฺมิํ จ ธมฺเม สาวโก วิหรติ ธมฺมา\-นุธมฺม\-ปฺปฏิปนฺโน สามีจิ\-ปฺปฏิปนฺโน อนุธมฺมจรี, อิทมหํ สุพีชสฺมิํ วทามิ, กิมงฺคํ ปน สรณานิํ สกฺกํ, สรณานิ มหานาม สกฺโก มรณกาเล สิกฺขาย ปริปูรการี อโหสีติ.  ปญฺจมํ.}

\csromanmatikaanchor{mtk.240}
\csromantocmarkref{subparagraph}{6-7. อนาถปิณฺฑิกสุตฺตานิ}{mtk.240}
\bookmark[dest={mtk.240},level=5]{6-7. อนาถปิณฺฑิกสุตฺตานิ}
\prose{\csromanfolio{332}6. ปฐม\-อนาถปิณฺฑิก\-สุตฺต}

\proseitem{1022}{สาวตฺถิ\-นิทานํ. เตน โข ปน สมเยน อนาถปิณฺฑิโก คหปติ อาพาธิโก โหติ ทุกฺขิโต พาฬฺหคิลาโน. อถ โข อนาถปิณฺฑิโก คหปติ อญฺญตรํ ปุริสํ อามนฺเตสิ “เอหิ ตฺวํ อมฺโภ ปุริส เยนายสฺมา สาริปุตฺโต เตนุปสงฺกม, อุปสงฺกมิตฺวา มม วจเนน อายสฺมโต สาริปุตฺตสฺส ปาเท สิรสา วนฺท, อนาถปิณฺฑิโก ภนฺเต คหปติ อาพาธิโก ทุกฺขิโต พาฬฺหคิลาโน, โส อายสฺมโต สาริปุตฺตสฺส ปาเท สิรสา วนฺทตี”ติ, เอวญฺจ วเทหิ “สาธุ กิร ภนฺเต อายสฺมา สาริปุตฺโต เยน อนาถ\-ปิณฺฑิกสฺส คหปติสฺส นิเวสนํ เตนุ\-ปสงฺกมตุ อนุกมฺปํ อุปาทายา”ติ.}

\prose{“เอวํ ภนฺเต”ติ โข โส ปุริโส อนาถ\-ปิณฺฑิกสฺส คหปติสฺส ปฏิสฺสุตฺวา เยนายสฺมา สาริปุตฺโต เตนุปสงฺกมิ, อุปสงฺกมิตฺวา อายสฺมนฺตํ สาริปุตฺตํ อภิวาเทตฺวา เอกมนฺตํ นิสีทิ, เอกมนฺตํ นิสินฺโน โข โส ปุริโส อายสฺมนฺตํ สาริปุตฺตํ เอตทโวจ\csromandash{}}

\prose{อนาถิปิณฺฑิโก ภนฺเต คหปติ อาพาธิโก ทุกฺขิโต พาฬฺหคิลาโน, โส อายสฺมโต สาริปุตฺตสฺส ปาเท สิรสา วนฺทติ, เอวญฺจ วทติ “สาธุ กิร ภนฺเต อายสฺมา สาริปุตฺโต เยน อนาถ\-ปิณฺฑิกสฺส คหปติสฺส นิเวสนํ เตนุ\-ปสงฺกมตุ อนุกมฺปํ อุปาทายา”ติ. อธิวาเสสิ โข อายสฺมา สาริปุตฺโต ตุณฺหีภาเวน.}

\prose{อถ โข อายสฺมา สาริปุตฺโต ปุพฺพณฺหสมยํ นิวาเสตฺวา ปตฺต\-จีวรมา\-ทาย อายสฺมตา อานนฺเทน ปจฺฉาสมเณน เยน อนาถ\-ปิณฺฑิกสฺส คหปติสฺส นิเวสนํ เตนุปสงฺกมิ, อุปสงฺกมิตฺวา ปญฺญตฺเต อาสเน นิสีทิ, นิสชฺช โข อายสฺมา สาริปุตฺโต อนาถปิณฺฑิกํ คหปติํ เอตทโวจ “กจฺจิ เต คหปติ ขมนียํ, กจฺจิ ยาปนียํ, กจฺจิ ทุกฺขา เวทนา ปฏิกฺกมนฺติ, โน อภิกฺกมนฺติ, ปฏิกฺกโมสานํ ปญฺญายติ, โน อภิกฺกโม”ติ. น เม ภนฺเต ขมนียํ, น ยาปนียํ, พาฬฺหา เม ทุกฺขา เวทนา อภิกฺกมนฺติ, โน ปฏิกฺกมนฺติ, อภิกฺกโมสานํ ปญฺญายติ, โน ปฏิกฺกโมติ.}

\prose{ยถารูเปน โข คหปติ พุทฺเธ อปฺปสาเทน สมนฺนาคโต อสฺสุตวา ปุถุชฺชโน กายสฺส เภทา ปรํ มรณา อปายํ ทุคฺคติํ วินิปาตํ นิรยํ}

\prose{\csromanfolio{333}อุปปชฺชติ, ตถารูโป เต พุทฺเธ อปฺปสาโท นตฺถิ, อตฺถิ จ โข เต คหปติ พุทฺเธ อเวจฺจปฺปสาโท “อิติปิ โส ภควา ฯเปฯ สตฺถา เทวมนุสฺสานํ พุทฺโธ ภควา”ติ, ตญฺจ ปน เต พุทฺเธ อเวจฺจปฺปสาทํ อตฺตนิ สมนุปสฺสโต ฐานโส เวทนา ปฏิปฺปสฺสมฺเภยฺย.}

\prose{ยถารูเปน โข คหปติ ธมฺเม อปฺปสาเทน สมนฺนาคโต อสฺสุตวา ปุถุชฺชโน กายสฺส เภทา ปรํ มรณา อปายํ ทุคฺคติํ วินิปาตํ นิรยํ อุปปชฺชติ, ตถารูโป เต ธมฺเม อปฺปสาโท นตฺถิ, อตฺถิ จ โข เต คหปติ ธมฺเม อเวจฺจปฺปสาโท “สฺวากฺขาโต ภควตา ธมฺโม ฯเปฯ ปจฺจตฺตํ เวทิตพฺโพ วิญฺญูหี”ติ, ตญฺจ ปน เต ธมฺเม อเวจฺจปฺปสาทํ อตฺตนิ สมนุปสฺสโต ฐานโส เวทนา ปฏิปฺปสฺสมฺเภยฺย.}

\prose{ยถารูเปน โข คหปติ สํเฆ อปฺปสาเทน สมนฺนาคโต อสฺสุตวา ปุถุชฺชโน กายสฺส เภทา ปรํ มรณา อปายํ ทุคฺคติํ วินิปาตํ นิรยํ อุปปชฺชติ, ตถารูโป เต สํเฆ อปฺปสาโท นตฺถิ, อตฺถิ จ โข เต คหปติ สํเฆ อเวจฺจปฺปสาโท “สุปฺปฏิปนฺโน ภควโต สาวกสํโฆ ฯเปฯ อนุตฺตรํ ปุญฺญกฺเขตฺตํ โลกสฺสา”ติ, ตญฺจ ปน เต สํเฆ อเวจฺจปฺปสาทํ อตฺตนิ สมนุปสฺสโต ฐานโส เวทนา ปฏิปฺปสฺสมฺเภยฺย.}

\prose{ยถารูเปน โข คหปติ ทุสฺสีเลฺยน สมนฺนาคโต อสฺสุตวา ปุถุชฺชโน กายสฺส เภทา ปรํ มรณา อปายํ ทุคฺคติํ วินิปาตํ นิรยํ อุปปชฺชติ, ตถารูปํ เต ทุสฺสีลฺยํ นตฺถิ. อตฺถิ จ โข เต คหปติ อริยกนฺตานิ สีลานิ ฯเปฯ สมาธิ\-สํวตฺตนิกานิ, ตานิ จ ปน เต อริยกนฺตานิ สีลานิ อตฺตนิ สมนุปสฺสโต ฐานโส เวทนา ปฏิปฺปสฺสมฺเภยฺย.}

\prose{ยถารูปาย โข คหปติ มิจฺฉาทิฏฺฐิยา สมนฺนาคโต อสฺสุตวา ปุถุชฺชโน กายสฺส เภทา ปรํ มรณา อปายํ ทุคฺคติํ วินิปาตํ นิรยํ อุปปชฺชติ, ตถารูปา เต มิจฺฉาทิฏฺฐิ นตฺถิ. อตฺถิ จ โข เต คหปติ สมฺมาทิฏฺฐิ, ตญฺจ ปน สมฺมาทิฏฺฐิํ อตฺตนิ สมนุปสฺสโต ฐานโส เวทนา ปฏิปฺปสฺสมฺเภยฺย.}

\prose{ยถารูเปน โข คหปติ มิจฺฉา\-สงฺกปฺเปน สมนฺนาคโต อสฺสุตวา ปุถุชฺชโน กายสฺส เภทา ปรํ มรณา อปายํ ทุคฺคติํ วินิปาตํ นิรยํ}

\prose{\csromanfolio{334}อุปปชฺชติ, ตถารูโป เต มิจฺฉาสงฺกปฺโป นตฺถิ. อตฺถิ จ โข เต คหปติ สมฺมาสงฺกปฺโป, ตญฺจ ปน เต สมฺมาสงฺกปฺปํ อตฺตนิ สมนุปสฺสโต ฐานโส เวทนา ปฏิปฺปสฺสมฺเภยฺย.}

\prose{ยถารูปาย โข คหปติ มิจฺฉาวาจาย สมนฺนาคโต อสฺสุตวา ปุถุชฺชโน กายสฺส เภทา ปรํ มรณา อปายํ ทุคฺคติํ วินิปาตํ นิรยํ อุปปชฺชติ, ตถารูปา เต มิจฺฉาวาจา นตฺถิ. อตฺถิ จ โข เต คหปติ สมฺมาวาจา, ตญฺจ ปน เต สมฺมาวาจํ อตฺตนิ สมนุปสฺสโต ฐานโส เวทนา ปฏิปฺปสฺสมฺเภยฺย.}

\prose{ยถารูเปน โข คหปติ มิจฺฉา\-กมฺมนฺเตน สมนฺนาคโต อสฺสุตวา ปุถุชฺชโน กายสฺส เภทา ปรํ มรณา อปายํ ทุคฺคติํ วินิปาตํ นิรยํ อุปปชฺชติ, ตถารูโป เต มิจฺฉากมฺมนฺโต นตฺถิ. อตฺถิ จ โข เต คหปติ สมฺมากมฺมนฺโต, ตญฺจ ปน เต สมฺมากมฺมนฺตํ อตฺตนิ สมนุปสฺสโต ฐานโส เวทนา ปฏิปฺปสฺสมฺเภยฺย.}

\prose{ยถารูเปน โข คหปติ มิจฺฉาอาชีเวน สมนฺนาคโต อสฺสุตวา ปุถุชฺชโน กายสฺส เภทา ปรํ มรณา อปายํ ทุคฺคติํ วินิปาตํ นิรยํ อุปปชฺชติ, ตถารูโป เต มิจฺฉาอาชีโว นตฺถิ. อตฺถิ จ โข เต คหปติ สมฺมาอาชีโว, ตญฺจ ปน เต สมฺมาอาชีวํ อตฺตนิ สมนุปสฺสโต ฐานโส เวทนา ปฏิปฺปสฺสมฺเภยฺย.}

\prose{ยถารูเปน โข คหปติ มิจฺฉาวายาเมน สมนฺนาคโต อสฺสุตวา ปุถุชฺชโน กายสฺส เภทา ปรํ มรณา อปายํ ทุคฺคติํ วินิปาตํ นิรยํ อุปปชฺชติ, ตถารูโป เต มิจฺฉาวายาโม นตฺถิ. อตฺถิ จ โข เต คหปติ สมฺมาวายาโม, ตญฺจ ปน เต สมฺมาวายามํ อตฺตนิ สมนุปสฺสโต ฐานโส เวทนา ปฏิปฺปสฺสมฺเภยฺย.}

\prose{ยถารูเปน โข คหปติ มิจฺฉาสติยา สมนฺนาคโต อสฺสุตวา ปุถุชฺชโน กายสฺส เภทา ปรํ มรณา อปายํ ทุคฺคติํ วินิปาตํ นิรยํ อุปปชฺชติ, ตถารูปา เต มิจฺฉาสติ นตฺถิ. อตฺถิ จ โข เต คหปติ สมฺมาสติ, ตญฺจ ปน เต สมฺมาสติํ อตฺตนิ สมนุปสฺสโต ฐานโส เวทนา ปฏิปฺปสฺสมฺเภยฺย.}

\prose{ยถารูเปน โข คหปติ มิจฺฉา\-สมาธินา สมนฺนาคโต อสฺสุตวา ปุถุชฺชโน กายสฺส เภทา ปรํ มรณา อปายํ ทุคฺคติํ วินิปาตํ นิรยํ}

\prose{\csromanfolio{335}อุปปชฺชติ, ตถารูโป เต มิจฺฉาสมาธิ นตฺถิ. อตฺถิ จ โข เต คหปติ สมฺมาสมาธิ, ตญฺจ ปน เต สมฺมาสมาธิํ อตฺตนิ สมนุปสฺสโต ฐานโส เวทนา ปฏิปฺปสฺสมฺเภยฺย.}

\prose{ยถารูเปน โข คหปติ มิจฺฉาญาเณน สมนฺนาคโต อสฺสุตวา ปุถุชฺชโน กายสฺส เภทา ปรํ มรณา อปายํ ทุคฺคติํ วินิปาตํ นิรยํ อุปปชฺชติ, ตถารูปํ เต มิจฺฉาญาณํ นตฺถิ. อตฺถิ จ โข เต คหปติ สมฺมาญาณํ, ตญฺจ ปน เต สมฺมาญาณํ อตฺตนิ สมนุปสฺสโต ฐานโส เวทนา ปฏิปฺปสฺสมฺเภยฺย.}

\prose{ยถารูปาย โข คหปติ มิจฺฉา\-วิมุตฺติยา สมนฺนาคโต อสฺสุตวา ปุถุชฺชโน กายสฺส เภทา ปรํ มรณา อปายํ ทุคฺคติํ วินิปาตํ นิรยํ อุปปชฺชติ, ตถารูปา เต มิจฺฉาวิมุตฺติ นตฺถิ. อตฺถิ จ โข เต คหปติ สมฺมาวิมุตฺติ, ตญฺจ ปน เต สมฺมาวิมุตฺติํ อตฺตนิ สมนุปสฺสโต ฐานโส เวทนา ปฏิปฺปสฺสมฺเภยฺยาติ.}

\prose{อถ โข อนาถ\-ปิณฺฑิกสฺส คหปติสฺส ฐานโส เวทนา ปฏิปฺปสฺสมฺภิํสุ. อถ โข อนาถปิณฺฑิโก คหปติ อายสฺมนฺตญฺจ สาริปุตฺตํ อายสฺมนฺตญฺจ อานนฺทํ สเกเนว ถาลิปาเกน ปริวิสิ. อถ โข อนาถปิณฺฑิโก คหปติ อายสฺมนฺตํ สาริปุตฺตํ ภุตฺตาวิํ โอนีต\-ปตฺต\-ปาณิํ อญฺญตรํ นีจาสนํ คเหตฺวา เอกมนฺตํ นิสีทิ, เอกมนฺตํ นิสินฺนํ โข อนาถปิณฺฑิกํ คหปติํ อายสฺมา สาริปุตฺโต อิมาหิ คาถาหิ อนุโมทิ\csromandash{}}

\setlength{\csromangathapagewidth}{0pt}
\setlength{\csromangathawakwidth}{0pt}
\csromangathameasuregroup{“ยสฺส สทฺธา ตถาคเต, \\ สีลญฺจ ยสฺส กลฺยาณํ, \\ สํเฆ ปสาโท ยสฺสตฺถิ, \\ อทลิทฺโทติ\textsuperscript{1} ตํ อาหุ, \\ ตสฺมา สทฺธญฺจ สีลญฺจ, \\ อนุยุญฺเชถ เมธาวี,}{\csromangathabat{“ยสฺส สทฺธา ตถาคเต,}{อจลา สุปฺปติฏฺฐิตา.} \\ \csromangathabat{สีลญฺจ ยสฺส กลฺยาณํ,}{อริยกนฺตํ ปสํสิตํ.} \\ \csromangathabat{สํเฆ ปสาโท ยสฺสตฺถิ,}{อุชุภูตญฺจ ทสฺสนํ.} \\ \csromangathabat{อทลิทฺโทติ\textsuperscript{1} ตํ อาหุ,}{อโมฆํ ตสฺส ชีวิตํ.} \\ \csromangathabat{ตสฺมา สทฺธญฺจ สีลญฺจ,}{ปสาทํ ธมฺมทสฺสนํ.} \\ \csromangathabat{อนุยุญฺเชถ เมธาวี,}{สรํ พุทฺธานสาสนนฺ”ติ.}}
\csromangathasetleft{“ยสฺส สทฺธา ตถาคเต, \\ สีลญฺจ ยสฺส กลฺยาณํ, \\ สํเฆ ปสาโท ยสฺสตฺถิ, \\ อทลิทฺโทติ\textsuperscript{1} ตํ อาหุ, \\ ตสฺมา สทฺธญฺจ สีลญฺจ, \\ อนุยุญฺเชถ เมธาวี,}
\csromangathagroup{\csromangathastanza{\csromangathabat{“ยสฺส สทฺธา ตถาคเต,}{อจลา สุปฺปติฏฺฐิตา.} \\ \csromangathabat{สีลญฺจ ยสฺส กลฺยาณํ,}{อริยกนฺตํ ปสํสิตํ.}}\csromangathastanza{\csromangathabat{สํเฆ ปสาโท ยสฺสตฺถิ,}{อุชุภูตญฺจ ทสฺสนํ.} \\ \csromangathabat{อทลิทฺโทติ\csromansharedfootnote{baa3fbcbd14e307d}{อทฬิทฺโทติ (สี, สฺยา, กํ)} ตํ อาหุ,}{อโมฆํ ตสฺส ชีวิตํ.}}\csromangathastanza{\csromangathabat{ตสฺมา สทฺธญฺจ สีลญฺจ,}{ปสาทํ ธมฺมทสฺสนํ.} \\ \csromangathabat{อนุยุญฺเชถ เมธาวี,}{สรํ พุทฺธานสาสนนฺ”ติ.}}}

\prose{อถ โข อายสฺมา สาริปุตฺโต อนาถปิณฺฑิกํ คหปติํ อิมาหิ คาถาหิ อนุโมทิตฺวา อุฏฺฐายาสนา ปกฺกามิ.}

\prose{\csromanfolio{336}อถ โข อายสฺมา อานนฺโท เยน ภควา เตนุปสงฺกมิ, อุปสงฺกมิตฺวา ภควนฺตํ อภิวาเทตฺวา เอกมนฺตํ นิสีทิ, เอกมนฺตํ นิสินฺนํ โข อายสฺมนฺตํ อานนฺทํ ภควา เอตทโวจ “หนฺท กุโต นุ ตฺวํ อานนฺท อาคจฺฉสิ ทิวา ทิวสฺสาติ. อายสฺมตา ภนฺเต สาริปุตฺเตน อนาถปิณฺฑิโก คหปติ อิมินา จ อิมินา จ โอวาเทน โอวทิโตติ. ปณฺฑิโต อานนฺท สาริปุตฺโต, มหาปญฺโญ อานนฺท สาริปุตฺโต, ยตฺร หิ นาม จตฺตาริ โสตาปตฺติยงฺคานิ ทสหากาเรหิ วิภชิสฺสตี”ติ.  ฉฏฺฐํ.}

\prose{7. ทุติย\-อนาถปิณฺฑิก\-สุตฺต}

\proseitem{1023}{สาวตฺถิ\-นิทานํ. เตน โข ปน สมเยน อนาถปิณฺฑิโก คหปติ อาพาธิโก โหติ ทุกฺขิโต พาฬฺหคิลาโน. อถ โข อนาถปิณฺฑิโก คหปติ อญฺญตรํ ปุริสํ อามนฺเตสิ “เอหิ ตฺวํ อมฺโภ ปุริส, เยนายสฺมา อานนฺโท เตนุปสงฺกม, อุปสงฺกมิตฺวา มม วจเนน อายสฺมโต อานนฺทสฺส ปาเท สิรสา วนฺท, อนาถปิณฺฑิโก ภนฺเต คหปติ อาพาธิโก ทุกฺขิโต พาฬฺหคิลาโน, โส อายสฺมโต อานนฺทสฺส ปาเท สิรสา วนฺทตีติ, เอวญฺจ วเทหิ ‘สาธุ กิร ภนฺเต อายสฺมา อานนฺโท เยน อนาถ\-ปิณฺฑิกสฺส คหปติสฺส นิเวสนํ เตนุ\-ปสงฺกมตุ, อนุกมฺปํ อุปาทายา’ติ”.}

\prose{“เอวํ ภนฺเต”ติ โข โส ปุริโส อนาถ\-ปิณฺฑิกสฺส คหปติสฺส ปฏิสฺสุตฺวา เยนายสฺมา อานนฺโท เตนุปสงฺกมิ, อุปสงฺกมิตฺวา อายสฺมานฺตํ อานนฺทํ อภิวาเทตฺวา เอกมนฺตํ นิสีทิ, เอกมนฺตํ นิสินฺโน โข โส ปุริโส อายสฺมนฺตํ อานนฺทํ เอตทโวจ “อนาถปิณฺฑิโก ภนฺเต คหปติ อาพาธิโก ทุกฺขิโต พาฬฺหคิลาโน, โส อายสฺมโต อานนฺทสฺส ปาเท สิรสา วนฺทติ, เอวญฺจ วทติ ‘สาธุ กิร ภนฺเต อายสฺมา อานนฺโท เยน อนาถ\-ปิณฺฑิกสฺส คหปติสฺส นิเวสนํ เตนุ\-ปสงฺกมตุ อนุกมฺปํ อุปาทายา’ติ”. อธิวาเสสิ โข อายสฺมา อานนฺโท ตุณฺหีภาเวน.}

\prose{อถ โข อายสฺมา อานนฺโท ปุพฺพณฺหสมยํ นิวาเสตฺวา ปตฺต\-จีวรมา\-ทาย เยน อนาถ\-ปิณฺฑิกสฺส คหปติสฺส นิเวสนํ เตนุปสงฺกมิ, อุปสงฺกมิตฺวา ปญฺญตฺเต อาสเน นิสีทิ, นิสชฺช โข อายสฺมา อานนฺโท}

\prose{\csromanfolio{337}อนาถปิณฺฑิกํ คหปติํ เอตทโวจ “กจฺจิ เต คหปติ ขมนียํ, กจฺจิ ยาปนียํ, กจฺจิ ทุกฺขา เวทนา ปฏิกฺกมนฺติ, โน อภิกฺกมนฺติ, ปฏิกฺกโมสานํ ปญฺญายติ, โน อภิกฺกโม”ติ. น เม ภนฺเต ขมนียํ, น ยาปนียํ, พาฬฺหา เม ทุกฺขา เวทนา อภิกฺกมนฺติ, โน ปฏิกฺกมนฺติ, อภิกฺกโมสานํ ปญฺญายติ, โน ปฏิกฺกโมติ.}

\prose{จตูหิ โข คหปติ ธมฺเมหิ สมนฺนาคตสฺส อสฺสุตวโต ปุถุชฺชนสฺส โหติ อุตฺตาโส, โหติ ฉมฺภิตตฺตํ, โหติ สมฺปรายิกํ มรณภยํ. กตเมหิ จตูหิ, อิธ คหปติ อสฺสุตวา ปุถุชฺชโน พุทฺเธ อปฺปสาเทน สมนฺนาคโต โหติ, ตญฺจ ปนสฺส พุทฺเธ อปฺปสาทํ อตฺตนิ สมนุปสฺสโต โหติ อุตฺตาโส, โหติ ฉมฺภิตตฺตํ, โหติ สมฺปรายิกํ มรณภยํ.}

\prose{ปุน จปรํ คหปติ อสฺสุตวา ปุถุชฺชโน ธมฺเม อปฺปสาเทน สมนฺนาคโต โหติ, ตญฺจ ปนสฺส ธมฺเม อปฺปสาทํ อตฺตนิ สมนุปสฺสโต โหติ อุตฺตาโส, โหติ ฉมฺภิตตฺตํ, โหติ สมฺปรายิกํ มรณภยํ.}

\prose{ปุน จปรํ คหปติ อสฺสุตวา ปุถุชฺชโน สํเฆ อปฺปสาเทน สมนฺนาคโต โหติ, ตญฺจ ปนสฺส สํเฆ อปฺปสาทํ อตฺตนิ สมนุปสฺสโต โหติ อุตฺตาโส, โหติ ฉมฺภิตตฺตํ, โหติ สมฺปรายิกํ มรณภยํ.}

\prose{ปุน จปรํ คหปติ อสฺสุตวา ปุถุชฺชโน ทุสฺสีเลฺยน สมนฺนาคโต โหติ, ตญฺจ ปนสฺส ทุสฺสีลฺยํ อตฺตนิ สมนุปสฺสโต โหติ อุตฺตาโส, โหติ ฉมฺภิตตฺตํ, โหติ สมฺปรายิกํ มรณภยํ. อิเมหิ โข คหปติ จตูหิ ธมฺเมหิ สมนฺนาคตสฺส อสฺสุตวโต ปุถุชฺชนสฺส โหติ อุตฺตาโส, โหติ ฉมฺภิตตฺตํ, โหติ สมฺปรายิกํ มรณภยํ.}

\prose{จตูหิ โข คหปติ ธมฺเมหิ สมนฺนาคตสฺส สุตวโต อริยสาวกสฺส น โหติ อุตฺตาโส, น โหติ ฉมฺภิตตฺตํ, น โหติ สมฺปรายิกํ มรณภยํ. กตเมหิ จตูหิ, อิธ คหปติ สุตวา อริยสาวโก พุทฺเธ อเวจฺจ\-ปฺปสาเทน สมนฺนาคโต โหติ “อิติปิ โส ภควา ฯเปฯ สตฺถา เทวมนุสฺสานํ พุทฺโธ ภควา”ติ, ตญฺจ ปนสฺส พุทฺเธ อเวจฺจปฺปสาทํ อตฺตนิ สมนุปสฺสโต น โหติ อุตฺตาโส, น โหติ ฉมฺภิตตฺตํ, น โหติ สมฺปรายิกํ มรณภยํ.}

\prose{ปุน จปรํ คหปติ สุตวา อริยสาวโก ธมฺเม อเวจฺจ\-ปฺปสาเทน สมนฺนาคโต โหติ “สฺวากฺขาโต ภควตา ธมฺโม ฯเปฯ ปจฺจตฺตํ}

\csromanmatikaanchor{mtk.241}
\csromantocmarkref{subparagraph}{8-9. ภยเวรูปสนฺตสุตฺตานิ}{mtk.241}
\bookmark[dest={mtk.241},level=5]{8-9. ภยเวรูปสนฺตสุตฺตานิ}
\prose{\csromanfolio{338}เวทิตพฺโพ วิญฺญูหี”ติ, ตญฺจ ปนสฺส ธมฺเม อเวจฺจปฺปสาทํ อตฺตนิ สมนุปสฺสโต น โหติ อุตฺตาโส, น โหติ ฉมฺภิตตฺตํ, น โหติ สมฺปรายิกํ มรณภยํ.}

\prose{ปุน จปรํ คหปติ สุตวา อริยสาวโก สํเฆ อเวจฺจ\-ปฺปสาเทน สมนฺนาคโต โหติ “สุปฺปฏิปนฺโน ภควโต สาวกสํโฆ ฯเปฯ อนุตฺตรํ ปุญฺญกฺเขตฺตํ โลกสฺสา”ติ. ตญฺจ ปนสฺส สํเฆ อเวจฺจปฺปสาทํ อตฺตนิ สมนุปสฺสโต น โหติ อุตฺตาโส, น โหติ ฉมฺภิตตฺตํ, น โหติ สมฺปรายิกํ มรณภยํ.}

\prose{ปุน จปรํ คหปติ สุตวา อริยสาวโก อริยกนฺเตหิ สีเลหิ สมนฺนาคโต โหติ อขณฺเฑหิ ฯเปฯ สมาธิ\-สํวตฺตนิเกหิ, ตานิ จ ปนสฺส อริยกนฺตานิ สีลานิ อตฺตนิ สมนุปสฺสโต น โหติ อุตฺตาโส, น โหติ ฉมฺภิตตฺตํ, น โหติ สมฺปรายิกํ มรณภยํ. อิเมหิ โข คหปติ จตูหิ ธมฺเมหิ สมนฺนาคตสฺส สุตวโต อริยสาวกสฺส น โหติ อุตฺตาโส, น โหติ ฉมฺภิตตฺตํ, น โหติ สมฺปรายิกํ มรณภยนฺติ.}

\prose{นาหํ ภนฺเต อานนฺท ภายามิ, กฺยาหํ ภายิสฺสามิ. อหํ หิ ภนฺเต พุทฺเธ อเวจฺจ\-ปฺปสาเทน สมนฺนาคโต โหมิ “อิติปิ โส ภควา ฯเปฯ สตฺถา เทวมนุสฺสานํ พุทฺโธ ภควา”ติ. ธมฺเม ฯเปฯ. สํเฆ อเวจฺจ\-ปฺปสาเทน สมนฺนาคโต โหมิ “สุปฺปฏิปนฺโน ภควโต สาวกสํโฆ ฯเปฯ อนุตฺตรํ ปุญฺญกฺเขตฺตํ โลกสฺสา”ติ, ยานิ จิมานิ ภนฺเต ภควตา คิหิสามีจิกานิ สิกฺขาปทานิ เทสิตานิ, นาหํ เตสํ กิญฺจิ อตฺตนิ ขณฺฑํ สมนุปสฺสามีติ. ลาภา เต คหปติ, สุลทฺธํ เต คหปติ, โสตาปตฺติผลํ ตยา คหปติ พฺยากตนฺติ.  สตฺตมํ.}

\csromanheader{2}{8. ปฐม\-ภย\-เวรู\-ปสนฺต\-สุตฺต}
\proseitem{1024}{สาวตฺถิ\-นิทานํ. เอกมนฺตํ นิสินฺนํ โข อนาถปิณฺฑิกํ คหปติํ ภควา เอตทโวจ\csromandash{}ยโต โข คหปติ อริยสาวกสฺส ปญฺจ ภยานิ เวรานิ วูปสนฺตานิ จ โหนฺติ, จตูหิ จ โสตาปตฺติยงฺเคหิ สมนฺนาคโต โหติ, อริโย จสฺส ญาโย ปญฺญาย สุทิฏฺโฐ โหติ สุปฺปฏิวิทฺโธ, โส อากงฺขมาโน อตฺตนาว อตฺตานํ พฺยากเรยฺย, ขีณนิรโยมฺหิ}

\prose{\csromanfolio{339}ขีณติรจฺฉานโยนิ\csromansharedfootnote{fd417d673540c829}{ขีณติรจฺฉานโยนิโย (สพฺพตฺถ)} ขีณาเปตฺติวิสโย ขีณา\-ปายทุคฺคติวินิปาโต, โสตาปนฺโนหมสฺมิ อวินิปาต\-ธมฺโม นิยโต สมฺโพธิ\-ปรายโณ.}

\prose{กตมานิ ปญฺจ ภยานิ เวรานิ วูปสนฺตานิ โหนฺติ, ยํ คหปติ ปาณาติปาตี ปาณาติปาต\-ปฺปจฺจยา ทิฏฺฐ\-ธมฺมิกมฺปิ ภยํ เวรํ ปสวติ, สมฺปรายิกมฺปิ ภยํ เวรํ ปสวติ, เจตสิกมฺปิ ทุกฺขํ โทมนสฺสํ ปฏิสํเวทิยติ. ปาณาติปาตา ปฏิวิรตสฺส เอวํ ตํ ภยํ เวรํ วูปสนฺตํ โหติ. ยํ คหปติ อทินฺนาทายี ฯเปฯ. ยํ คหปติ กาเมสุ\-มิจฺฉาจารี ฯเปฯ. ยํ คหปติ มุสาวาที ฯเปฯ. ยํ คหปติ สุรา\-เมรย\-มชฺช\-ปฺปมาท\-ฏฺฐายี สุรา\-เมรย\-มชฺช\-ปฺปมาทฏฺฐาน\-ปฺปจฺจยา ทิฏฺฐ\-ธมฺมิกมฺปิ ภยํ เวรํ ปสวติ, สมฺปรายิกมฺปิ ภยํ เวรํ ปสวติ, เจตสิกมฺปิ ทุกฺขํ โทมนสฺสํ ปฏิสํเวทิยติ, สุรา\-เมรย\-มชฺช\-ปฺปมาทฏฺฐานา ปฏิวิรตสฺส เอวํ ตํ ภยํ เวรํ วูปสนฺตํ โหติ. อิมานิ ปญฺจ ภยานิ เวรานิ วูปสนฺตานิ โหนฺติ.}

\prose{กตเมหิ จตูหิ โสตาปตฺติยงฺเคหิ สมนฺนาคโต โหติ, อิธ คหปติ อริยสาวโก พุทฺเธ อเวจฺจ\-ปฺปสาเทน สมนฺนาคโต โหติ “อิติปิ โส ภควา ฯเปฯ สตฺถา เทวมนุสฺสานํ พุทฺโธ ภควา”ติ. ธมฺเม ฯเปฯ. สํเฆ ฯเปฯ. อริยกนฺเตหิ สีเลหิ สมนฺนาคโต โหติ อขณฺเฑหิ ฯเปฯ สมาธิ\-สํวตฺตนิเกหิ. อิเมหิ จตูหิ โสตาปตฺติยงฺเคหิ สมนฺนาคโต โหติ.}

\prose{กตโม จสฺส อริโย ญาโย ปญฺญาย สุทิฏฺโฐ โหติ สุปฺปฏิวิทฺโธ, อิธ คหปติ อริยสาวโก ปฏิจฺจ\-สมุปฺปาทญฺเญว สาธุกํ โยนิโส มนสิ กโรติ “อิติ อิมสฺมิํ สติ อิทํ โหติ, อิมสฺสุปฺปาทา อิทํ อุปฺปชฺชติ. อิติ อิมสฺมิํ อสติ อิทํ น โหติ, อิมสฺส นิโรธา อิทํ นิรุชฺฌติ, ยทิทํ อวิชฺชาปจฺจยา สงฺขารา, สงฺขาร\-ปจฺจยา วิญฺญาณํ ฯเปฯ เอวเมตสฺส เกวลสฺส ทุกฺข\-กฺขนฺธสฺส สมุทโย โหติ. อวิชฺชาย เตฺวว อเสส\-วิราค\-นิโรธา สงฺขาร\-นิโรโธ ฯเปฯ เอวเมตสฺส เกวลสฺส ทุกฺข\-กฺขนฺธสฺส นิโรโธ โหติ. อยมสฺส อริโย ญาโย ปญฺญาย สุทิฏฺโฐ โหติ สุปฺปฏิวิทฺโธ.}

\prose{ยโต โข คหปติ อริยสาวกสฺส อิมานิ ปญฺจ ภยานิ เวรานิ วูปสนฺตานิ โหนฺติ, อิเมหิ จตูหิ โสตาปตฺติยงฺเคหิ สมนฺนาคโต โหติ, อยญฺจสฺส อริโย ญาโย ปญฺญาย สุทิฏฺโฐ โหติ}

\prose{\csromanfolio{340}สุปฺปฏิวิทฺโธ, โส อากงฺขมาโน อตฺตนาว อตฺตานํ พฺยากเรยฺย “ขีณนิรโยมฺหิ ขีณติรจฺฉานโยนิ ขีณเปตฺติวิสโย ขีณา\-ปายทุคฺคติวินิปาโต, โสตาปนฺโนหมสฺมิ อวินิปาต\-ธมฺโม นิยโต สมฺโพธิ\-ปรายโณ”ติ.  อฏฺฐมํ.}

\csromanheader{2}{9. ทุติย\-ภย\-เวรู\-ปสนฺต\-สุตฺต}
\proseitem{1025}{สาวตฺถิ\-นิทานํ ฯเปฯ. ยโต โข ภิกฺขเว อริยสาวกสฺส อิมานิ ปญฺจ ภยานิ เวรานิ วูปสนฺตานิ โหนฺติ, อิเมหิ จตูหิ โสตาปตฺติยงฺเคหิ สมนฺนาคโต โหติ, อยญฺจสฺส อริโย ญาโย ปญฺญาย สุทิฏฺโฐ โหติ สุปฺปฏิวิทฺโธ, โส อากงฺขมาโน อตฺตนาว อตฺตานํ พฺยากเรยฺย “ขีณนิรโยมฺหิ ขีณติรจฺฉานโยนิ ขีณเปตฺติวิสโย ขีณา\-ปายทุคฺคติวินิปาโต, โสตาปนฺโนหมสฺมิ อวินิปาต\-ธมฺโม นิยโต สมฺโพธิ\-ปรายโณ”ติ.  นวมํ.}

\csromanmatikaanchor{mtk.242}
\csromantocmarkref{subparagraph}{10. นนฺทกลิจฺฉวิสุตฺต}{mtk.242}
\bookmark[dest={mtk.242},level=5]{10. นนฺทกลิจฺฉวิสุตฺต}
\csromanheader{6}{10. นนฺทก\-ลิจฺฉวิ\-สุตฺต}
\proseitem{1026}{เอกํ สมยํ ภควา เวสาลิยํ วิหรติ มหาวเน กูฏาคาร\-สาลายํ. อถ โข นนฺทโก ลิจฺฉาวิมหามตฺโต เยน ภควา เตนุปสงฺกมิ, อุปสงฺกมิตฺวา ภควนฺตํ อภิวาเทตฺวา เอกมนฺตํ นิสีทิ, เอกมนฺตํ นิสินฺนํ โข นนฺทกํ ลิจฺฉวิ\-มหามตฺตํ ภควา เอตทโวจ\csromandash{}}

\prose{จตูหิ โข นนฺทก ธมฺเมหิ สมนฺนาคโต อริยสาวโก โสตาปนฺโน โหติ อวินิปาต\-ธมฺโม นิยโต สมฺโพธิ\-ปรายโณ. กตเมหิ จตูหิ, อิธ นนฺทก อริยสาวโก พุทฺเธ อเวจฺจ\-ปฺปสาเทน สมนฺนาคโต โหติ “อิติปิ โส ภควา ฯเปฯ สตฺถา เทวมนุสฺสานํ พุทฺโธ ภควา”ติ. ธมฺเม ฯเปฯ. สํเฆ ฯเปฯ. อริยกนฺเตหิ สีเลหิ สมนฺนาคโต โหติ อขณฺเฑหิ ฯเปฯ สมาธิ\-สํวตฺตนิเกหิ. อิเมหิ โข นนฺทก จตูหิ ธมฺเมหิ สมนฺนาคโต อริยสาวโก โสตาปนฺโน โหติ อวินิปาต\-ธมฺโม นิยโต สมฺโพธิ\-ปรายโณ.}

\prose{อิเมหิ จ ปน นนฺทก จตูหิ ธมฺเมหิ สมนฺนาคโต อริยสาวโก อายุนา สํยุตฺโต โหติ ทิพฺเพนปิ มานุเสนปิ, วณฺเณน สํยุตฺโต โหติ ทิพฺเพนปิ มานุเสนปิ, สุเขน สํยุตฺโต โหติ ทิพฺเพนปิ มานุเสนปิ, ยเสน สํยุตฺโต โหติ ทิพฺเพนปิ มานุเสนปิ, อาธิปเตยฺเยน สํยุตฺโต โหติ ทิพฺเพนปิ มานุเสนปิ. ตํ โข ปนาหํ}

\csromanmatikaanchor{mtk.243}
\csromanmatikaanchor{mtk.244}
\csromantocmarknopage{subparagraph}{4. ปุญฺญาภิสนฺทวคฺค}{mtk.243}
\bookmark[dest={mtk.243},level=5]{4. ปุญฺญาภิสนฺทวคฺค}
\csromantocmarkref{subparagraph}{1-3. ปุญฺญาภิสนฺทสุตฺตานิ}{mtk.244}
\bookmark[dest={mtk.244},level=5]{1-3. ปุญฺญาภิสนฺทสุตฺตานิ}
\prose{\csromanfolio{341}นนฺทก นาญฺญสฺส สมณสฺส วา พฺราหฺมณสฺส วา สุตฺวา วทามิ, อปิ จ ยเทว มยา สามํ ญาตํ สามํ ทิฏฺฐํ สามํ วิทิตํ, ตเทวาหํ วทามีติ.}

\prosewithclosermidruled{เอวํ วุตฺเต อญฺญตโร ปุริโส นนฺทกํ ลิจฺฉวิ\-มหามตฺตํ เอตทโวจ “นหานกาโล ภนฺเต”ติ. อลํ ทานิ ภเณ เอเตน พาหิเรน นหาเนน, อลมิทํ อชฺฌตฺตํ นหานํ ภวิสฺสติ ยทิทํ ภควติ ปสาโทติ.  ทสมํ.}{สรณานิวคฺโค ตติโย.}
\csromancenter{\textbf{ตสฺสุทฺทานํ}}
\setlength{\csromangathapagewidth}{0pt}
\setlength{\csromangathawakwidth}{0pt}
\csromangathameasuregroup{มหานาเมน เทฺว วุตฺตา, \\ ทุเว อนาถปิณฺฑิกา,}{\csromangathabat{มหานาเมน เทฺว วุตฺตา,}{โคธา จ สรณา ทุเว.} \\ \csromangathabat{ทุเว อนาถปิณฺฑิกา,}{ทุเว เวรภเยน จ.} \\ ลิจฺฉวี ทสโม วุตฺโต, วคฺโค เตน ปวุจฺจตีติ.}
\csromangathameasurewak{ลิจฺฉวี ทสโม วุตฺโต, วคฺโค เตน ปวุจฺจตีติ.}
\csromangathasetleft{มหานาเมน เทฺว วุตฺตา, \\ ทุเว อนาถปิณฺฑิกา,}
\csromangathagroup{\csromangathastanza{\csromangathabat{มหานาเมน เทฺว วุตฺตา,}{โคธา จ สรณา ทุเว.} \\ \csromangathabat{ทุเว อนาถปิณฺฑิกา,}{ทุเว เวรภเยน จ.}}\csromangathastanza{\csromangathawak{ลิจฺฉวี ทสโม วุตฺโต, วคฺโค เตน ปวุจฺจตีติ.}}}

\csromanheader{6}{4. ปุญฺญาภิสนฺท\-วคฺค}
\csromanheader{2}{1. ปฐม\-ปุญฺญาภิสนฺท\-สุตฺต}
\proseitem{1027}{สาวตฺถิ\-นิทานํ. จตฺตาโรเม ภิกฺขเว ปุญฺญาภิสนฺทา กุสลาภิสนฺทา สุขสฺสาหารา. กตเม จตฺตาโร, อิธ ภิกฺขเว อริยสาวโก พุทฺเธ อเวจฺจ\-ปฺปสาเทน สมนฺนาคโต โหติ “อิติปิ โส ภควา ฯเปฯ สตฺถา เทวมนุสฺสานํ พุทฺโธ ภควา”ติ, อยํ ปฐโม ปุญฺญาภิสนฺโท กุสลาภิสนฺโท สุขสฺสาหาโร.}

\prose{ปุน จปรํ ภิกฺขเว อริยสาวโก ธมฺเม อเวจฺจ\-ปฺปสาเทน สมนฺนาคโต โหติ “สฺวากฺขาโต ภควตา ธมฺโม ฯเปฯ ปจฺจตฺตํ เวทิตพฺโพ วิญฺญูหี”ติ, อยํ ทุติโย ปุญฺญาภิสนฺโท กุสลาภิสนฺโท สุขสฺสาหาโร.}

\prose{ปุน จปรํ ภิกฺขเว อริยสาวโก สํเฆ อเวจฺจ\-ปฺปสาเทน สมนฺนาคโต โหติ “สุปฺปฏิปนฺโน ภควโต สาวกสํโฆ ฯเปฯ อนุตฺตรํ ปุญฺญกฺเขตฺตํ โลกสฺสา”ติ, อยํ ตติโย ปุญฺญาภิสนฺโท กุสลาภิสนฺโท สุขสฺสาหาโร.}

\prose{\csromanfolio{342}ปุน จปรํ ภิกฺขเว อริยสาวโก อริยกนฺเตหิ สีเลหิ สมนฺนาคโต โหติ อขณฺเฑหิ ฯเปฯ สมาธิ\-สํวตฺตนิเกหิ, อยํ จตุตฺโถ ปุญฺญาภิสนฺโท กุสลาภิสนฺโท สุขสฺสาหาโร. อิเม โข ภิกฺขเว จตฺตาโร ปุญฺญาภิสนฺทา กุสลาภิสนฺทา สุขสฺสาหาราติ.}

\csromanheader{2}{2. ทุติย\-ปุญฺญาภิสนฺท\-สุตฺต}
\proseitem{1028}{จตฺตาโรเม ภิกฺขเว ปุญฺญาภิสนฺทา กุสลาภิสนฺทา สุขสฺสาหารา. กตเม จตฺตาโร, อิธ ภิกฺขเว อริยสาวโก พุทฺเธ อเวจฺจ\-ปฺปสาเทน สมนฺนาคโต โหติ “อิติปิ โส ภควา ฯเปฯ สตฺถา เทวมนุสฺสานํ พุทฺโธ ภควา”ติ, อยํ ปฐโม ปุญฺญาภิสนฺโท กุสลาภิสนฺโท สุขสฺสาหาโร.}

\prose{ปุน จปรํ ภิกฺขเว อริยสาวโก ธมฺเม ฯเปฯ. สํเฆ ฯเปฯ.}

\prose{ปุน จปรํ ภิกฺขเว อริยสาวโก วิคต\-มล\-มจฺเฉเรน เจตสา อคารํ อชฺฌาวสติ มุตฺตจาโค ปยตปาณิ โวสฺสคฺครโต ยาจโยโค ทาน\-สํวิภาค\-รโต, อยํ จตุตฺโถ ปุญฺญาภิสนฺโท กุสลาภิสนฺโท สุขสฺสาหาโร. อิเม โข ภิกฺขเว จตฺตาโร ปุญฺญาภิสนฺทา กุสลาภิสนฺทา สุขสฺสาหาราติ.  ทุติยํ.}

\csromanheader{2}{3. ตติย\-ปุญฺญาภิสนฺท\-สุตฺต}
\proseitem{1029}{จตฺตาโรเม ภิกฺขเว ปุญฺญาภิสนฺทา กุสลาภิสนฺทา สุขสฺสาหารา. กตเม จตฺตาโร, อิธ ภิกฺขเว อริยสาวโก พุทฺเธ อเวจฺจ\-ปฺปสาเทน สมนฺนาคโต โหติ “อิติปิ โส ภควา ฯเปฯ สตฺถา เทวมนุสฺสานํ พุทฺโธ ภควา”ติ, อยํ ปฐโม ปุณฺณาภิสนฺโท กุลสาภิสนฺโท สุขสฺสาหาโร.}

\prose{ปุน จปรํ ภิกฺขเว อริยสาวโก ธมฺเม ฯเปฯ. สํเฆ ฯเปฯ.}

\prose{ปุน จปรํ ภิกฺขเว อริยสาวโก ปญฺญวา โหติ อุทย\-ตฺถคามินิยา ปญฺญาย สมนฺนาคโต อริยาย นิพฺเพธิกาย สมฺมา ทุกฺข\-กฺขย\-คามินิยา, อยํ จตุตฺโถ ปุญฺญาภิสนฺโท กุสลาภิสนฺโท สุขสฺสาหาโร. อิเม โข ภิกฺขเว จตฺตาโร ปุญฺญาภิสนฺทา กุสลาภิสนฺทา สุขสฺสาหาราติ.  ตติยํ.}

\csromanheader{2}{4. ปฐม\-เทวปท\-สุตฺต}
\csromanmatikaanchor{mtk.245}
\csromantocmarkref{subparagraph}{4-6. เทวปทสุตฺตานิ}{mtk.245}
\bookmark[dest={mtk.245},level=5]{4-6. เทวปทสุตฺตานิ}
\proseitem{1030}{\csromanfolio{343}สาวตฺถิ\-นิทานํ. จตฺตาริมานิ ภิกฺขเว เทวานํ เทวปทานิ อวิสุทฺธานํ สตฺตานํ วิสุทฺธิยา อปริโยทาตานํ สตฺตานํ ปริโยทปนาย.}

\prose{กตมานิ จตฺตาริ, อิธ ภิกฺขเว อริยสาวโก พุทฺเธ อเวจฺจ\-ปฺปสาเทน สมนฺนาคโต โหติ “อิติปิ โส ภควา ฯเปฯ สตฺถา เทวมนุสฺสานํ พุทฺโธ ภควา”ติ, อิทํ ปฐมํ เทวปทํ อวิสุทฺธานํ สตฺตานํ วิสุทฺธิยา อปริโยทาตานํ สตฺตานํ ปริโยทปนาย.}

\prose{ปุน จปรํ ภิกฺขเว อริยสาวโก ธมฺเม ฯเปฯ. สํเฆ ฯเปฯ.}

\prose{ปุน จปรํ ภิกฺขเว อริยสาวโก อริยกนฺเตหิ สีเลหิ สมนฺนาคโต โหติ อขณฺเฑหิ ฯเปฯ สมาธิ\-สํวตฺตนิเกหิ. อิทํ จตุตฺถํ เทวานํ เทวปทํ อวิสุทฺธานํ สตฺตานํ วิสุทฺธิยา อปริโยทาตานํ สตฺตานํ ปริโยทปนาย. อิมานิ โข ภิกฺขเว จตฺตาริ เทวานํ เทวปทานิ อวิสุทฺธานํ สตฺตานํ วิสุทฺธิยา อปริโยทาตานํ สตฺตานํ ปริโยทปนายาติ.  จตุตฺถํ.}

\csromanheader{2}{5. ทุติย\-เทวปท\-สุตฺต}
\proseitem{1031}{จตฺตาริมานิ ภิกฺขเว เทวานํ เทวปทานิ อวิสุทฺธานํ สตฺตานํ วิสุทฺธิยา อปริโยทาตานํ สตฺตานํ ปริโยทปนาย.}

\prose{กตมานิ จตฺตาริ, อิธ ภิกฺขเว อริยสาวโก พุทฺเธ อเวจฺจ\-ปฺปสาเทน สมนฺนาคโต โหติ “อิติปิ โส ภควา ฯเปฯ สตฺถา เทวมนุสฺสานํ พุทฺโธ ภควา”ติ. โส อิติ ปฏิสญฺจิกฺขติ “กิํ นุ โข เทวานํ เทวปทนฺ”ติ. โส เอวํ ปชานาติ “อพฺยาพชฺฌ\-ปรเม ขฺวาหํ เอตรหิ เทเว สุณามิ, น จ โข ปนาหํ กิญฺจิ พฺยาพาเธมิ ตสํ วา ถาวรํ วา, อทฺธาหํ เทวปท\-ธมฺม\-สมนฺนาคโต วิหรามี”ติ, อิทํ ปฐมํ เทวานํ เทวปทํ อวิสุทฺธานํ สตฺตานํ วิสุทฺธิยา อปริโยทาตานํ สตฺตานํ ปริโยทปนาย.}

\prose{ปุน จปรํ ภิกฺขเว อริยสาวโก ธมฺเม ฯเปฯ. สํเฆ ฯเปฯ.}

\prose{ปุน จปรํ ภิกฺขเว อริยสาวโก อริยกนฺเตหิ สีเลหิ สมนฺนาคโต โหติ อขณฺเฑหิ ฯเปฯ สมาธิ\-สํวตฺตนิเกหิ, โส อิติ ปฏิสญฺจิกฺขติ “กิํ นุ โข เทวานํ เทวปทนฺ”ติ. โส เอวํ ปชานาติ \csromanfolio{344}“อพฺยาพชฺฌ\-ปรเม ขฺวาหํ เอตรหิ เทเว สุณามิ, น โข ปนาหํ กิญฺจิ พฺยาพาเธมิ ตสํ วา ถาวรํ วา, อทฺธาหํ เทวปท\-ธมฺม\-สมนฺนาคโต วิหรามี”ติ, อิทํ จตุตฺถํ เทวานํ เทวปทํ อวิสุทฺธานํ สตฺตานํ วิสุทฺธิยา อปริโยทาตานํ สตฺตานํ ปริโยทปนาย. อิมานิ โข ภิกฺขเว จตฺตาริ เทวานํ เทวปาทานิ อวิสุทฺธานํ สตฺตานํ วิสุทฺธิยา อปริโยทาตานํ สตฺตานํ ปริโยทปนายาติ.  ปญฺจมํ.}

\csromanheader{2}{6. เทว\-สภาค\-สุตฺต}
\proseitem{1032}{จตูหิ ภิกฺขเว ธมฺเมหิ สมนฺนาคตํ อตฺตมนา เทวา สภาคตํ กเถนฺติ. กตเมหิ จตูหิ, อิธ ภิกฺขเว อริยสาวโก พุทฺเธ อเวจฺจ\-ปฺปสาเทน สมนฺนาคโต โหติ “อิติปิ โส ภควา ฯเปฯ สตฺถา เทวมนุสฺสานํ พุทฺโธ ภควา”ติ. ยา ตา เทวตา พุทฺเธ อเวจฺจ\-ปฺปสาเทน สมนฺนาคตา อิโต จุตา ตตฺรูปปนฺนา, ตาสํ เอวํ โหติ “ยถารูเปน โข มยํ พุทฺเธ อเวจฺจ\-ปฺปสาเทน สมนฺนาคตา ตโต จุตา อิธูปปนฺนา อริยสาวโกปิ ตถารูเปน พุทฺเธ อเวจฺจ\-ปฺปสาเทน สมนฺนาคโต เอหีติ เทวานํ สนฺติเก”ติ.}

\prose{ปุน จปรํ ภิกฺขเว อริยสาวโก ธมฺเม ฯเปฯ. สํเฆ ฯเปฯ. อริยกนฺเตหิ สีเลหิ สมนฺนาคโต โหติ อขณฺเฑหิ ฯเปฯ สมาธิ\-สํวตฺตนิเกหิ, ยา ตา เทวตา อริยกนฺเตหิ สีเลหิ สมนฺนาคตา อิโต จุตา ตตฺรูปปนฺนา, ตาสํ เอวํ โหติ “ยถารูเปหิ โข มยํ อริยกนฺเตหิ สีเลหิ สมนฺนาคตา ตโต จุตา อิธูปปนฺนา อริยสาวโกปิ ตถารูเปหิ อริยกนฺเตหิ สีเลหิ สมนฺนาคโต เอหีติ เทวานํ สนฺติเก”ติ. อิเมหิ โข ภิกฺขเว จตูหิ สมนฺนาคตํ อตฺตมนา เทวา สภาคตํ กเถนฺตีติ.  ฉฏฺฐํ.}

\csromanmatikaanchor{mtk.246}
\csromantocmarkref{subparagraph}{7. มหานามสุตฺต}{mtk.246}
\bookmark[dest={mtk.246},level=5]{7. มหานามสุตฺต}
\csromanheader{6}{7. \textbf{มหานามสุตฺต}}
\proseitem{1033}{เอกํ สมยํ ภควา สกฺเกสุ วิหรติ กปิล\-วตฺถุสฺมิํ นิโคฺรธาราเม. อถ โข มหานาโม สกฺโก เยน ภควา เตนุปสงฺกมิ, อุปสงฺกมิตฺวา ภควนฺตํ อภิวาเทตฺวา เอกมนฺตํ นิสีทิ, เอกมนฺตํ นิสินฺโน โข มหานาโม สกฺโก ภควนฺตํ เอตทโวจ\csromandash{}}

\prose{กิตฺตาวตา นุ โข ภนฺเต อุปาสโก โหตีติ. ยโต โข มหานาม พุทฺธํ สรณํ คโต โหติ, ธมฺมํ สรณํ คโต โหติ,}

\prose{\csromanfolio{345}สํฆํ สรณํ คโต โหติ. เอตฺตาวตา โข มหานาม อุปาสโก โหตีติ.}

\prose{กิตฺตาวตา ปน ภนฺเต อุปาสโก สีลสมฺปนฺโน โหตีติ. ยโต โข มหานาม อุปาสโก ปาณาติปาตา ปฏิวิรโต โหติ, อทินฺนาทานา ปฏิวิรโต โหติ, กาเมสุ\-มิจฺฉาจารา ปฏิวิรโต โหติ, มุสาวาทา ปฏิวิรโต โหติ, สุรา\-เมรย\-มชฺช\-ปฺปมาทฏฺฐานา ปฏิวิรโต โหติ. เอตฺตาวตา โข มหานาม อุปาสโก สีลสมฺปนฺโน โหตีติ.}

\prose{กิตฺตาวตา ปน ภนฺเต อุปาสโก สทฺธาสมฺปนฺโน โหตีติ, อิธ มหานาม อุปาสโก สทฺโธ โหติ, สทฺทหติ ตถาคตสฺส โพธิํ “อิติปิ โส ภควา ฯเปฯ สตฺถา เทวมนุสฺสานํ พุทฺโธ ภควา”ติ. เอตฺตาวตา โข มหานาม อุปาสโก สทฺธาสมฺปนฺโน โหตีติ.}

\prose{กิตฺตาวตา ปน ภนฺเต อุปาสโก จาคสมฺปนฺโน โหตีติ. อิธ มหานาม อุปาสโก วิคต\-มล\-มจฺเฉเรน เจตสา อคารํ อชฺฌาวสติ มุตฺตจาโค ปยตปาณิ โวสฺสคฺครโต ยาจโยโค ทาน\-สํวิภาค\-รโต. เอตฺตาวตา โข มหานาม อุปาสโก จาคสมฺปนฺโน โหตีติ.}

\prose{กิตฺตาวตา ปน ภนฺเต อุปาสโก ปญฺญาสมฺปนฺโน โหตีติ. อิธ มหานาม อุปาสโก ปญฺญวา โหติ, อุทย\-ตฺถคามินิยา ปญฺญาย สมนฺนาคโต อริยาย นิพฺเพธิกาย สมฺมา ทุกฺข\-กฺขย\-คามินิยา. เอตฺตาวตา โข มหานาม อุปาสโก ปญฺญาสมฺปนฺโน โหตีติ.  สตฺตมํ.}

\csromanmatikaanchor{mtk.247}
\csromantocmarkref{subparagraph}{8. วสฺสสุตฺต}{mtk.247}
\bookmark[dest={mtk.247},level=5]{8. วสฺสสุตฺต}
\csromanheader{6}{8. \textbf{วสฺสสุตฺต}}
\proseitem{1034}{เสยฺยถาปิ ภิกฺขเว อุปริปพฺพเต ถุลฺล\-ผุสิตเก เทเว วสฺสนฺเต ตํ อุทกํ ยถานินฺนํ ปวตฺตมานํ ปพฺพต\-กนฺทร\-ปทร\-สาขา ปริปูเรติ, ปพฺพต\-กนฺทร\-ปทร\-สาขา ปริปูรา กุโสพฺเภ ปริปูเรนฺติ, กุโสพฺภา ปริปูรา มหาโสพฺเภ ปริปูเรนฺติ, มหาโสพฺภา ปริปูรา กุนฺนทิโย ปริปูเรนฺติ, กุนฺนทิโย ปริปูรา มหานทิโย ปริปูเรนฺติ, มหานทิโย ปริปูรา มหาสมุทฺทํ\csromansharedfootnote{dfca408eecbd6db5}{มหาสมุทฺทสาครํ (สพฺพตฺถ) สํ 1. 269 ปิฏฺเฐปิ.} ปริปูเรนฺติ. เอวเมว โข ภิกฺขเว อริยสาวกสฺส โย จ พุทฺเธ อเวจฺจปฺปสาโท โย จ ธมฺเม อเวจฺจปฺปสาโท โย จ \csromanfolio{346}สํเฆ อเวจฺจปฺปสาโท ยานิ จ อริยกนฺตานิ สีลานิ, อิเม ธมฺมา สนฺทมานา ปารํ คนฺตฺวา อาสวานํ ขยาย สํวตฺตนฺตีติ.  อฏฺฐมํ.}

\csromanmatikaanchor{mtk.248}
\csromantocmarkref{subparagraph}{9. กาฬิโคธสุตฺต}{mtk.248}
\bookmark[dest={mtk.248},level=5]{9. กาฬิโคธสุตฺต}
\csromanheader{6}{9. \textbf{กาฬิโคธสุตฺต}}
\proseitem{1035}{เอกํ สมยํ ภควา สกฺเกสุ วิหรติ กปิล\-วตฺถุสฺมิํ นิโคฺรธาราเม. อถ โข ภควา ปุพฺพณฺหสมยํ นิวาเสตฺวา ปตฺต\-จีวรมา\-ทาย เยน กาฬิโคธาย สากิยานิยา นิเวสนํ เตนุปสงฺกมิ, อุปสงฺกมิตฺวา ปญฺญตฺเต อาสเน นิสีทิ. อถ โข กาฬิโคธา สากิยานี เยน ภควา เตนุปสงฺกมิ, อุปสงฺกมิตฺวา ภควนฺตํ อภิวาเทตฺวา เอกมนฺตํ นิสีทิ, เอกมนฺตํ นิสินฺนํ โข กาฬิโคธํ สากิยานิํ ภควา เอตทโวจ\csromandash{}}

\prose{จตูหิ โข โคเธ ธมฺเมหิ สมนฺนาคตา อริยสาวิกา โสตาปนฺนา โหติ อวินิปาต\-ธมฺมา นิยตา สมฺโพธิ\-ปรายณา. กตเมหิ จตูหิ, อิธ โคเธ อริยสาวิกา พุทฺเธ อเวจฺจ\-ปฺปสาเทน สมนฺนาคตา โหติ “อิติปิ โส ภควา ฯเปฯ สตฺถา เทวมนุสฺสานํ พุทฺโธ ภควา”ติ. ธมฺเม ฯเปฯ. สํเฆ ฯเปฯ. วิคต\-มล\-มจฺเฉเรน เจตสา อคารํ อชฺฌาวสติ มุตฺตจาคา ปยตปาณินี\csromansharedfootnote{f97d9c3ac14e9f84}{ปยตปาณี (สพฺพตฺถ) 3. 30 โมคฺคลฺลานสุตฺตํ โอโลเกตพฺพํ.} โวสฺสคฺครตา ยาจโยคา ทาน\-สํวิภาค\-รตา. อิเมหิ โข โคเธ จตูหิ ธมฺเมหิ สมนฺนาคตา อริยสาวิกา โสตาปนฺนา โหติ อวินิปาต\-ธมฺมา นิยตา สมฺโพธิ\-ปรายณาติ.}

\prose{ยานิมานิ ภนฺเต ภควตา จตฺตาริ โสตาปตฺติยงฺคานิ เทสิตานิ, สํวิชฺชนฺเตเต ธมฺมา มยิ, อหญฺจ เตสุ ธมฺเมสุ สนฺทิสฺสามิ, อหํ หิ ภนฺเต พุทฺเธ อเวจฺจ\-ปฺปสาเทน สมนฺนาคตา “อิติปิ โส ภควา ฯเปฯ สตฺถา เทวมนุสฺสานํ พุทฺโธ ภควา”ติ. ธมฺเม ฯเปฯ. สํเฆ ฯเปฯ. ยํ โข ปน กิญฺจิ กุเล เทยฺยธมฺมํ, สพฺพํ ตํ อปฺปฏิวิภตฺตํ สีลวนฺเตหิ กลฺยาณธมฺเมหีติ. ลาภา เต โคเธ, สุลทฺธํ เต โคเธ, โสตาปตฺติผลํ ตยา โคเธ พฺยากตนฺติ.  นวมํ.}

\csromanmatikaanchor{mtk.249}
\csromantocmarkref{subparagraph}{10. นนฺทิยสกฺกสุตฺต}{mtk.249}
\bookmark[dest={mtk.249},level=5]{10. นนฺทิยสกฺกสุตฺต}
\csromanheader{6}{10. นนฺทิย\-สกฺก\-สุตฺต}
\proseitem{1036}{เอกํ สมยํ ภควา สกฺเกสุ วิหรติ กปิล\-วตฺถุสฺมิํ นิโคฺรธาราเม. อถ โข นนฺทิโย สกฺโก เยน ภควา เตนุปสงฺกมิ, อุปสงฺกมิตฺวา ภควนฺตํ อภิวาเทตฺวา เอกมนฺตํ นิสีทิ, เอกมนฺตํ นิสินฺโน}

\prose{\csromanfolio{347}โข นนฺทิโย สกฺโก ภควนฺตํ เอตทโวจ “ยสฺเสว นุ โข ภนฺเต อริยสาวกสฺส จตฺตาริ โสตาปตฺติยงฺคานิ สพฺเพน สพฺพํ สพฺพถา สพฺพํ นตฺถิ, เสฺวว นุ โข ภนฺเต อริยสาวโก ปมาทวิหารี”ติ.}

\prose{ยสฺส โข นนฺทิย จตฺตาริ โสตาปตฺติยงฺคานิ สพฺเพน สพฺพํ สพฺพถา สพฺพํ นตฺถิ, ตมหํ พาหิโร ปุถุชฺชน\-ปกฺเข ฐิโตติ วทามิ, อปิ จ นนฺทิย ยถา อริยสาวโก ปมาทวิหารี เจว โหติ อปฺปมาทวิหารี จ, ตํ สุณาหิ สาธุกํ มนสิ กโรหิ, ภาสิสฺสามีติ. “เอวํ ภนฺเต”ติ โข นนฺทิโย สกฺโก ภควโต ปจฺจสฺโสสิ. ภควา เอตทโวจ\csromandash{}}

\prose{กถญฺจ นนฺทิย อริยสาวโก ปมาทวิหารี โหติ, อิธ นนฺทิย อริยสาวโก พุทฺเธ อเวจฺจ\-ปฺปสาเทน สมนฺนาคโต โหติ “อิติปิ โส ภควา ฯเปฯ สตฺถา เทวมนุสฺสานํ พุทฺโธ ภควา”ติ. โส เตน พุทฺเธ อเวจฺจ\-ปฺปสาเทน สนฺตุฏฺโฐ น อุตฺตริ วายมติ ทิวา ปวิเวกาย รตฺติํ ปฏิสลฺลานาย, ตสฺส เอวํ ปมตฺตสฺส วิหรโต ปาโมชฺชํ น โหติ, ปาโมชฺเช อสติ ปีติ น โหติ, ปีติยา อสติ ปสฺสทฺธิ น โหติ, ปสฺสทฺธิยา อสติ ทุกฺขํ วิหรติ, ทุกฺขิโน จิตฺตํ น สมาธิยติ, อสมาหิเต จิตฺเต ธมฺมา น ปาตุภวนฺติ, ธมฺมานํ อปาตุภาวา ปมาทวิหารีเตฺวว สงฺขฺยํ คจฺฉติ.}

\prose{ปุน จปรํ นนฺทิย อริยสาวโก ธมฺเม ฯเปฯ. สํเฆ ฯเปฯ. อริยกนฺเตหิ สีเลหิ สมนฺนาคโต โหติ อขณฺเฑหิ ฯเปฯ สมาธิ\-สํวตฺตนิเกหิ, โส เตหิ อริยกนฺเตหิ สีเลหิ สนฺตุฏฺโฐ น อุตฺตริ วายมติ ทิวา ปวิเวกาย รตฺติํ ปฏิสลฺลานาย, ตสฺส เอวํ ปมตฺตสฺส วิหรโต ปาโมชฺชํ น โหติ, ปาโมชฺเช อสติ ปีติ น โหติ, ปีติยา อสติ ปสฺสทฺธิ น โหติ, ปสฺสทฺธิยา อสติ ทุกฺขํ วิหรติ, ทุกฺขิโน จิตฺตํ น สมาธิยติ, อสมาหิเต จิตฺเต ธมฺมา น ปาตุภวนฺติ, ธมฺมานํ อปาตุภาวา ปมาทวิหารีเตฺวว สงฺขฺยํ คจฺฉติ. เอวํ โข นนฺทิย อริยสาวโก ปมาทวิหารี โหติ.}

\prose{กถญฺจ นนฺทิย อริยสาวโก อปฺปมาทวิหารี โหติ, อิธ นนฺทิย อริยสาวโก พุทฺเธ อเวจฺจ\-ปฺปสาเทน สมนฺนาคโต โหติ “อิติปิ โส ภควา ฯเปฯ สตฺถา เทวมนุสฺสานํ พุทฺโธ ภควา”ติ, โส เตน พุทฺเธ อเวจฺจ\-ปฺปสาเทน อสนฺตุฏฺโฐ อุตฺตริ วายมติ ทิวา ปวิเวกาย รตฺติํ ปฏิสลฺลานาย, ตสฺส เอวํ อปฺปมตฺตสฺส วิหรโต ปาโมชฺชํ ชายติ,}

\csromanmatikaanchor{mtk.250}
\csromanmatikaanchor{mtk.251}
\csromantocmarknopage{subparagraph}{5. สคาถกปุญฺญาภิสนฺทวคฺค}{mtk.250}
\bookmark[dest={mtk.250},level=5]{5. สคาถกปุญฺญาภิสนฺทวคฺค}
\csromantocmarkref{subparagraph}{1-3. อภิสนฺทสุตฺตานิ}{mtk.251}
\bookmark[dest={mtk.251},level=5]{1-3. อภิสนฺทสุตฺตานิ}
\prose{\csromanfolio{348}ปมุทิตสฺส ปีติ ชายติ, ปีติมนสฺส กาโย ปสฺสมฺภติ, ปสฺสทฺธกาโย สุขํ เวทิยติ, สุขิโน จิตฺตํ สมาธิยติ, สมาหิเต จิตฺเต ธมฺมา ปาตุภวนฺติ, ธมฺมานํ ปาตุภาวา อปฺปมาทวิหารีเตฺวว สงฺขฺยํ คจฺฉติ.}

\prosewithclosermidruled{ปุน จปรํ นนฺทิย อริยสาวโก ธมฺเม ฯเปฯ. สํเฆ ฯเปฯ. อริยกนฺเตหิ สีเลหิ สมนฺนาคโต โหติ อขณฺเฑหิ ฯเปฯ สมาธิ\-สํวตฺตนิเกหิ. โส เตหิ อริยกนฺเตหิ สีเลหิ อสนฺตุฏฺโฐ อุตฺตริ วายมติ ทิวา ปวิเวกาย รตฺติํ ปฏิสลฺลานาย, ตสฺส เอวํ อปฺปมตฺตสฺส วิหรโต ปาโมชฺชํ ชายติ, ปมุทิตสฺส ปีติ ชายติ, ปีติมนสฺส กาโย ปสฺสมฺภติ, ปสฺสทฺธกาโย สุขํ เวทิยติ, สุขิโน จิตฺตํ สมาธิยติ, สมาหิเต จิตฺเต ธมฺมา ปาตุภวนฺติ, ธมฺมานํ ปาตุภาวา อปฺปมาทวิหารีเตฺวว สงฺขฺยํ คจฺฉติ. เอวํ โข นนฺทิย อริยสาวโก อปฺปมาทวิหารี โหตีติ.  ทสมํ.}{ปุญฺญาภิสนฺท\-วคฺโค จตุตฺโถ.}
\csromancenter{\textbf{ตสฺสุทฺทานํ}}
\setlength{\csromangathapagewidth}{0pt}
\setlength{\csromangathawakwidth}{0pt}
\csromangathameasuregroup{อภิสนฺทา ตโย วุตฺตา, \\ สภาคตํ มหานาโม,}{\csromangathabat{อภิสนฺทา ตโย วุตฺตา,}{ทุเว เทวปทานิ จ.} \\ \csromangathabat{สภาคตํ มหานาโม,}{วสฺสํ กาฬี จ นนฺทิยาติ.}}
\csromangathasetleft{อภิสนฺทา ตโย วุตฺตา, \\ สภาคตํ มหานาโม,}
\csromangathagroup{\csromangathastanza{\csromangathabat{อภิสนฺทา ตโย วุตฺตา,}{ทุเว เทวปทานิ จ.} \\ \csromangathabat{สภาคตํ มหานาโม,}{วสฺสํ กาฬี จ นนฺทิยาติ.}}}

\csromanheader{6}{5. สคาถก\-ปุญฺญาภิสนฺท\-วคฺค}
\prose{1. ปฐม\-อภิสนฺท\-สุตฺต}

\proseitem{1037}{จตฺตาโรเม ภิกฺขเว ปุญฺญาภิสนฺทา กุสลาภิสนฺทา สุขสฺสาหารา. กตเม จตฺตาโร, อิธ ภิกฺขเว อริยสาวโก พุทฺเธ อเวจฺจ\-ปฺปสาเทน สมนฺนาคโต โหติ “อิติปิ โส ภควา ฯเปฯ สตฺถา เทวมนุสฺสานํ พุทฺโธ ภควา”ติ, อยํ ปฐโม ปุญฺญาภิสนฺโท กุสลาภิสนฺโท สุขสฺสาหาโร.}

\prose{ปุน จปรํ ภิกฺขเว อริยสาวโก ธมฺเม ฯเปฯ. สํเฆ ฯเปฯ.}

\prose{ปุน จปรํ ภิกฺขเว อริยสาวโก อริยกนฺเตหิ สีเลหิ สมนฺนาคโต โหติ อขณฺเฑหิ ฯเปฯ สมาธิ\-สํวตฺตนิเกหิ, อยํ}

\prose{\csromanfolio{349}จตุตฺโถ ปุญฺญาภิสนฺโท กุสลาภิสนฺโท สุขสฺสาหาโร. อิเม โข ภิกฺขเว จตฺตาโร ปุญฺญาภิสนฺทา กุสลาภิสนฺทา สุขสฺสาหารา.}

\prose{อิเมหิ โข ภิกฺขเว จตูหิ ปุญฺญา\-ภิสนฺเทหิ กุสลา\-ภิสนฺเทหิ สมนฺนาคตสฺส อริยสาวกสฺส น สุกรํ ปุญฺญสฺส ปมาณํ คเณตุํ “เอตฺตโก ปุญฺญาภิสนฺโท กุสลาภิสนฺโท สุขสฺสาหาโร”ติ. อถ โข อสงฺเขฺยยฺโย อปฺปเมยฺโย มหาปุญฺญกฺขนฺโธ\-เตฺวว สงฺขฺยํ คจฺฉติ.}

\prose{เสยฺยถาปิ ภิกฺขเว มหาสมุทฺเท น สุกรํ อุทกสฺส ปมาณํ คเณตุํ “เอตฺตกานิ อุทกาฬฺหกานี”ติ วา “เอตฺตกานิ อุทกาฬฺหกสตานี”ติ วา “เอตฺตกานิ อุทกาฬฺหกสหสฺสานี”ติ วา “เอตฺตกานิ อุทกาฬฺหกสตสหสฺสานี”ติ วาติ, อถ โข อสงฺเขฺยยฺโย อปฺปเมยฺโย มหาอุทกกฺขนฺโธเตฺวว สงฺขฺยํ คจฺฉติ. เอวเมว โข ภิกฺขเว อิเมหิ จตูหิ ปุญฺญา\-ภิสนฺเทหิ กุสลา\-ภิสนฺเทหิ สมนฺนาคตสฺส อริยสาวกสฺส น สุกรํ ปุญฺญสฺส ปมาณํ คเณตุํ “เอตฺตโก ปุญฺญาภิสนฺโท กุสลาภิสนฺโท สุขสฺสาหาโร”ติ. อถ โข อสงฺเขฺยยฺโย อปฺปเมยฺโย มหาปุญฺญกฺขนฺโธ\-เตฺวว สงฺขฺยํ คจฺฉตีติ.}

\prose{อิทมโวจ ภควา, อิทํ วตฺวาน สุคโต อถาปรํ เอตทโวจ สตฺถา\csromandash{}}

\setlength{\csromangathapagewidth}{0pt}
\setlength{\csromangathawakwidth}{0pt}
\csromangathameasuregroup{}{“มโหทธิํ อปริมิตํ มหาสรํ, \\ พหุเภรวํ รตนคณานมาลยํ. \\ นชฺโช ยถา นรคณสํฆเสวิตา, \\ ปุถู สวนฺตี อุปยนฺติ สาครํ. \\ เอวํ นรํ อนฺนปานวตฺถททํ, \\ เสยฺยานิ ปจฺจตฺถรณสฺส\textsuperscript{1} ทายกํ. \\ ปุญฺญสฺส ธารา อุปยนฺติ ปณฺฑิตํ,}
\csromangathameasurewak{“มโหทธิํ อปริมิตํ มหาสรํ, \\ พหุเภรวํ รตนคณานมาลยํ. \\ นชฺโช ยถา นรคณสํฆเสวิตา, \\ ปุถู สวนฺตี อุปยนฺติ สาครํ. \\ เอวํ นรํ อนฺนปานวตฺถททํ, \\ เสยฺยานิ ปจฺจตฺถรณสฺส\textsuperscript{1} ทายกํ. \\ ปุญฺญสฺส ธารา อุปยนฺติ ปณฺฑิตํ,}
\setlength{\csromangathaleftcol}{0pt}
\csromangathagroup{\csromangathastanza{\csromangathawak{“มโหทธิํ อปริมิตํ มหาสรํ, \\ พหุเภรวํ รตน\-คณานมาลยํ. \\ นชฺโช ยถา นรคณสํฆเสวิตา, \\ ปุถู สวนฺตี อุปยนฺติ สาครํ.}}\csromangathastanza{\csromangathawak{เอวํ นรํ อนฺนปานวตฺถททํ, \\ เสยฺยานิ ปจฺจตฺถรณสฺส\csromansharedfootnote{fc1c0b2e0e2bcf69}{สชฺชตฺถรณสฺส (สี, สฺยา, กํ, อิ)} ทายกํ. \\ ปุญฺญสฺส ธารา อุปยนฺติ ปณฺฑิตํ,}}}

\prose{นชฺโช ยถา วาริวหาว สาครนฺ”ติ. ปฐมํ.}

\prose{\csromanfolio{350}2. ทุติย\-อภิสนฺท\-สุตฺต}

\proseitem{1038}{จตฺตาโรเม ภิกฺขเว ปุญฺญาภิสนฺทา กุสลาภิสนฺทา สุขสฺสาหารา. กตเม จตฺตาโร, อิธ ภิกฺขเว อริยสาวโก พุทฺเธ อเวจฺจ\-ปฺปสาเทน สมนฺนาคโต โหติ “อิติปิ โส ภควา ฯเปฯ สตฺถา เทวมนุสฺสานํ พุทฺโธ ภควา”ติ, อยํ ปฐโม ปุญฺญาภิสนฺโท กุสลาภิสนฺโท สุขสฺสาหาโร.}

\prose{ปุน จปรํ ภิกฺขเว อริยสาวโก ธมฺเม ฯเปฯ. สํเฆ ฯเปฯ.}

\prose{ปุน จปรํ ภิกฺขเว อริยสาวโก วิคต\-มล\-มจฺเฉเรน เจตสา อคารํ อชฺฌาวสติ มุตฺตจาโค ปยตปาณิ โวสฺสคฺครโต ยาจโยโค ทาน\-สํวิภาค\-รโต, อยํ จตุตฺโถ ปุญฺญาภิสนฺโท กุสลาภิสนฺโท สุขสฺสาหาโร. อิเม โข ภิกฺขเว จตฺตาโร ปุญฺญาภิสนฺทา กุสลาภิสนฺทา สุขสฺสาหารา.}

\prose{อิเมหิ โข ภิกฺขเว จตูหิ ปุญฺญา\-ภิสนฺเทหิ กุสลา\-ภิสนฺเทหิ สมนฺนาคตสฺส อริยสาวกสฺส น สุกรํ ปุญฺญสฺส ปมาณํ คเณตุํ “เอตฺตโก ปุญฺญาภิสนฺโท กุสลาภิสนฺโท สุขสฺสาหาโร”ติ. อถ โข อาสงฺเขฺยยฺโย อปฺปเมยฺโย มหาปุญฺญกฺขนฺโธ\-เตฺวว สงฺขฺยํ คจฺฉติ.}

\prose{เสยฺยถาปิ ภิกฺขเว ยตฺถิมา มหานทิโย สํสนฺทนฺติ สเมนฺติ. เสยฺยถิทํ, คงฺคา ยมุนา อจิรวตี สรภู มหี, ตตฺถ น สุกรํ อุทกสฺส ปมาณํ คเณตุํ “เอตฺตกานิ อุทกาฬฺหกานี”ติ วา “เอตฺตกานิ อุทกาฬฺหกสตานี”ติ วา “เอตฺตกานิ อุทกาฬฺหกสหสฺสานี”ติ วา “เอตฺตกานิ อุทกาฬฺหกสตสหสฺสานี”ติ วาติ. อถ โข อสงฺเขฺยยฺโย อปฺปเมยฺโย มหาอุทกกฺขนฺโธเตฺวว สงฺขฺยํ คจฺฉติ. เอวเมว โข ภิกฺขเว อิเมหิ จตูหิ ปุญฺญา\-ภิสนฺเทหิ กุสลา\-ภิสนฺเทหิ สมนฺนาคตสฺส อริยสาวกสฺส น สุกรํ ปุญฺญสฺส ปมาณํ คเณตุํ “เอตฺตโก ปุญฺญาภิสนฺโท กุสลาภิสนฺโท สุขสฺสาหาโร”ติ. อถ โข อสงฺเขฺยยฺโย อปฺปเมยฺโย มหาปุญฺญกฺขนฺโธ\-เตฺวว สงฺขฺยํ คจฺฉตีติ. อิทมโวจ ภควา ฯเปฯ สตฺถา\csromandash{}}

\setlength{\csromangathapagewidth}{0pt}
\setlength{\csromangathawakwidth}{0pt}
\csromangathameasuregroup{}{“มโหทธิํ อปริมิตํ มหาสรํ, \\ พหุเภรวํ รตนคณานมาลยํ. \\ นชฺโช ยถา นรคณสํฆเสวิตา, \\ ปุถู สวนฺตี อุปยนฺติ สาครํ. \\ เอวํ นรํ อนฺนปานวตฺถททํ, \\ เสยฺยานิ ปจฺจตฺถรณสฺส ทายกํ. \\ ปุญฺญสฺส ธารา อุปยนฺติ ปณฺฑิตํ,}
\csromangathameasurewak{“มโหทธิํ อปริมิตํ มหาสรํ, \\ พหุเภรวํ รตนคณานมาลยํ. \\ นชฺโช ยถา นรคณสํฆเสวิตา, \\ ปุถู สวนฺตี อุปยนฺติ สาครํ. \\ เอวํ นรํ อนฺนปานวตฺถททํ, \\ เสยฺยานิ ปจฺจตฺถรณสฺส ทายกํ. \\ ปุญฺญสฺส ธารา อุปยนฺติ ปณฺฑิตํ,}
\setlength{\csromangathaleftcol}{0pt}
\csromangathagroup{\csromangathastanza{\csromangathawak{“มโหทธิํ อปริมิตํ มหาสรํ, \\ พหุเภรวํ รตน\-คณานมาลยํ. \\ นชฺโช ยถา นรคณสํฆเสวิตา, \\ ปุถู สวนฺตี อุปยนฺติ สาครํ.}}\csromangathastanza{\csromangathawak{\csromanfolio{351}เอวํ นรํ อนฺนปานวตฺถททํ, \\ เสยฺยานิ ปจฺจตฺถรณสฺส ทายกํ. \\ ปุญฺญสฺส ธารา อุปยนฺติ ปณฺฑิตํ,}}}

\prose{นชฺโช ยถา วาริวหาว สาครนฺ”ติ. ทุติยํ.}

\prose{3. ตติย\-อภิสนฺท\-สุตฺต}

\proseitem{1039}{จตฺตาโรเม ภิกฺขเว ปุญฺญาภิสนฺทา กุสลาภิสนฺทา สุขสฺสาหารา. กตเม จตฺตาโร, อิธ ภิกฺขเว อริยสาวโก พุทฺเธ อเวจฺจ\-ปฺปสาเทน สมนฺนาคโต โหติ “อิติปิ โส ภควา ฯเปฯ สตฺถา เทวมนุสฺสานํ พุทฺโธ ภควา”ติ, อยํ ปฐโม ปุญฺญาภิสนฺโท กุสลาภิสนฺโท สุขสฺสาหาโร.}

\prose{ปุน จปรํ ภิกฺขเว อริยสาวโก ธมฺเม ฯเปฯ. สํเฆ ฯเปฯ.}

\prose{ปุน จปรํ ภิกฺขเว อริยสาวโก ปญฺญวา โหติ, อุทย\-ตฺถคามินิยา ปญฺญาย สมนฺนาคโต อริยาย นิพฺเพธิกาย สมฺมา ทุกฺข\-กฺขย\-คามินิยา, อยํ จตุตฺโถ ปุญฺญาภิสนฺโท กุสลาภิสนฺโท สุขสฺสาหาโร. อิเม โข ภิกฺขเว จตฺตาโร ปุญฺญาภิสนฺทา กุสลาภิสนฺทา สุขสฺสาหารา.}

\prose{อิเมหิ โข ภิกฺขเว จตูหิ ปุญฺญา\-ภิสนฺเทหิ กุสลา\-ภิสนฺเทหิ สมนฺนาคตสฺส อริยสาวกสฺส น สุกรํ ปุญฺญสฺส ปมาณํ คเณตุํ “เอตฺตโก ปุญฺญาภิสนฺโท กุสลาภิสนฺโท สุขสฺสาหาโร”ติ. อถ โข อสงฺเขฺยยฺโย อปฺปเมยฺโย มหาปุญฺญกฺขนฺโธ\-เตฺวว สงฺขฺยํ คจฺฉตีติ. อิทมโวจ ภควา ฯเปฯ สตฺถา\csromandash{}}

\setlength{\csromangathapagewidth}{0pt}
\setlength{\csromangathawakwidth}{0pt}
\csromangathameasuregroup{}{“โย ปุญฺญกาโม กุสเล ปติฏฺฐิโต, \\ ภาเวติ มคฺคํ อมตสฺส ปตฺติยา. \\ โส ธมฺมสาราธิคโม ขเย รโต,}
\csromangathameasurewak{“โย ปุญฺญกาโม กุสเล ปติฏฺฐิโต, \\ ภาเวติ มคฺคํ อมตสฺส ปตฺติยา. \\ โส ธมฺมสาราธิคโม ขเย รโต,}
\setlength{\csromangathaleftcol}{0pt}
\csromangathagroup{\csromangathastanza{\csromangathawak{“โย ปุญฺญกาโม กุสเล ปติฏฺฐิโต, \\ ภาเวติ มคฺคํ อมตสฺส ปตฺติยา. \\ โส ธมฺมสาราธิคโม ขเย รโต,}}}

\vspace{\tipitakaparskip}
\csromancenter{น เวธติ มจฺจุราชาคมนสฺมินฺ”ติ\csromansharedfootnote{ca4c3291264be3b2}{มจฺจุราชาคมิสฺสตีติ (สี, อิ), มจฺจุชรากมฺปิสฺมินฺติ (ก)}. ตติยํ.}
\csromanheader{2}{4. ปฐม\-มหทฺธน\-สุตฺต}
\csromanmatikaanchor{mtk.252}
\csromantocmarkref{subparagraph}{4-10. ปฐมมหทฺธนสุตฺตาทีนิ}{mtk.252}
\bookmark[dest={mtk.252},level=5]{4-10. ปฐมมหทฺธนสุตฺตาทีนิ}
\proseitem{1040}{\csromanfolio{352}จตูหิ ภิกฺขเว ธมฺเมหิ สมนฺนาคโต อริยสาวโก อฑฺโฒ มหทฺธโน มหาโภโคติ\csromansharedfootnote{0d8f5e838489c4d6}{มหาโภโค มหายโสติ (สฺยา, อิ, ก)} วุจฺจติ.}

\prose{กตเมหิ จตูหิ, อิธ ภิกฺขเว อริยสาวโก พุทฺเธ อเวจฺจ\-ปฺปสาเทน สมนฺนาคโต โหติ “อิติปิ โส ภควา ฯเปฯ สตฺถา เทวมนุสฺสานํ พุทฺโธ ภควา”ติ. ธมฺเม ฯเปฯ. สํเฆ ฯเปฯ. อริยกนฺเตหิ สีเลหิ สมนฺนาคโต โหติ อขณฺเฑหิ ฯเปฯ สมาธิ\-สํวตฺตนิเกหิ. อิเมหิ โข ภิกฺขเว จตูหิ ธมฺเมหิ สมนฺนาคโต อริยสาวโก อฑฺโฒ มหทฺธโน มหาโภโคติ วุจฺจตีติ.  จตุตฺถํ.}

\csromanheader{2}{5. ทุติย\-มหทฺธน\-สุตฺต}
\proseitem{1041}{จตูหิ ภิกฺขเว ธมฺเมหิ สมนฺนาคโต อริยสาวโก อฑฺโฒ มหทฺธโน มหาโภโค มหายโสติ วุจฺจติ.}

\prose{กตเมหิ จตูหิ, อิธ ภิกฺขเว อริยสาวโก พุทฺเธ อเวจฺจ\-ปฺปสาเทน สมนฺนาคโต โหติ “อิติปิ โส ภควา ฯเปฯ สตฺถา เทวมนุสฺสานํ พุทฺโธ ภควา”ติ. ธมฺเม ฯเปฯ. สํเฆ ฯเปฯ. อริยกนฺเตหิ สีเลหิ สมนฺนาคโต โหติ อขณฺเฑหิ ฯเปฯ. สมาธิ\-สํวตฺตนิเกหิ. อิเมหิ โข ภิกฺขเว จตูหิ ธมฺเมหิ สมนฺนาคโต อริยสาวโก อฑฺโฒ มหทฺธโน มหาโภโค มหายโสติ วุจฺจตีติ.  ปญฺจมํ.}

\csromanheader{2}{6. \textbf{สุทฺธกสุตฺต}}
\proseitem{1042}{จตูหิ ภิกฺขเว ธมฺเมหิ สมนฺนาคโต อริยสาวโก โสตาปนฺโน โหติ อวินิปาต\-ธมฺโม นิยโต สมฺโพธิ\-ปรายโณ.}

\prose{กตเมหิ จตูหิ, อิธ ภิกฺขเว อริยสาวโก พุทฺเธ อเวจฺจ\-ปฺปสาเทน สมนฺนาคโต โหติ “อิติปิ โส ภควา ฯเปฯ สตฺถา เทวมนุสฺสานํ พุทฺโธ ภควา”ติ. ธมฺเม ฯเปฯ. สํเฆ ฯเปฯ. อริยกนฺเตหิ สีเลหิ สมนฺนาคโต โหติ อขณฺเฑหิ ฯเปฯ สมาธิ\-สํวตฺตนิเกหิ. อิเมหิ โข ภิกฺขเว จตูหิ ธมฺเมหิ สมนฺนาคโต อริยสาวโก โสตาปนฺโน โหติ อวินิปาต\-ธมฺโม นิยโต สมฺโพธิ\-ปรายโณติ.  ฉฏฺฐํ.}

\csromanheader{2}{7. \textbf{นนฺทิยสุตฺต}}
\proseitem{1043}{\csromanfolio{353}กปิลวตฺถุ\-นิทานํ. เอกมนฺตํ นิสินฺนํ โข นนฺทิยํ สกฺกํ ภควา เอตทโวจ “จตูหิ โข นนฺทิย ธมฺเมหิ สมนฺนาคโต อริยสาวโก โสตาปนฺโน โหติ อวินิปาต\-ธมฺโม นิยโต สมฺโพธิ\-ปรายโณ.}

\prose{กตเมหิ จตูหิ, อิธ นนฺทิย อริยสาวโก พุทฺเธ อเวจฺจ\-ปฺปสาเทน สมนฺนาคโต โหติ “อิติปิ โส ภควา ฯเปฯ สตฺถา เทวมนุสฺสานํ พุทฺโธ ภควา”ติ. ธมฺเม ฯเปฯ. สํเฆ ฯเปฯ. อริยกนฺเตหิ สีเลหิ สมนฺนาคโต โหติ อขณฺเฑหิ ฯเปฯ สมาธิ\-สํวตฺตนิเกหิ. อิเมหิ โข นนฺทิย จตูหิ ธมฺเมหิ สมนฺนาคโต อริยสาวโก โสตาปนฺโน โหติ อวินิปาต\-ธมฺโม นิยโต สมฺโพธิ\-ปรายโณ”ติ.  สตฺตมํ.}

\csromanheader{2}{8. \textbf{ภทฺทิยสุตฺต}}
\proseitem{1044}{กปิลวตฺถุ\-นิทานํ. เอกมนฺตํ นิสินฺนํ โข ภทฺทิยํ สกฺกํ ภควา เอตทโวจ “จตูหิ โข ภทฺทิย ธมฺเมหิ สมนฺนาคโต อริยสาวโก โสตาปนฺโน โหติ อวินิปาต\-ธมฺโม นิยโต สมฺโพธิ\-ปรายโณ.}

\prose{กตเมหิ จตูหิ, อิธ ภทฺทิย อริยสาวโก พุทฺเธ ฯเปฯ. ธมฺเม ฯเปฯ. สํเฆ ฯเปฯ. อริยกนฺเตหิ สีเลหิ สมนฺนาคโต โหติ อขณฺเฑหิ ฯเปฯ สมาธิ\-สํวตฺตนิเกหิ. อิเมหิ โข ภทฺทิย จตูหิ ธมฺเมหิ สมนฺนาคโต อริยสาวโก โสตาปนฺโน โหติ, อวินิปาต\-ธมฺโม นิยโต สมฺโพธิ\-ปรายโณ”ติ.  อฏฺฐมํ.}

\csromanheader{2}{9. \textbf{มหานามสุตฺต}}
\proseitem{1045}{กปิลวตฺถุ\-นิทานํ. เอกมนฺตํ นิสินฺนํ โข มหานามํ สกฺกํ ภควา เอตทโวจ “จตูหิ โข มหานาม ธมฺเมหิ สมนฺนาคโต อริยสาวโก โสตาปนฺโน โหติ ฯเปฯ สมฺโพธิ\-ปรายโณ.}

\prose{กตเมหิ จตูหิ, อิธ มหานาม อริยสาวโก พุทฺเธ ฯเปฯ. ธมฺเม ฯเปฯ. สํเฆ ฯเปฯ. อริยกนฺเตหิ สีเลหิ สมนฺนาคโต โหติ อขณฺเฑหิ ฯเปฯ สมาธิ\-สํวตฺตนิเกหิ. อิเมหิ โข มหานาม จตูหิ ธมฺเมหิ สมนฺนาคโต อริยสาวโก โสตาปนฺโน โหติ, อวินิปาต\-ธมฺโม นิยโต สมฺโพธิ\-ปรายโณ”ติ.  นวมํ.}

\csromanheader{2}{10. \textbf{องฺคสุตฺต}}
\proseitemwithclosermidruled{1046}{\csromanfolio{354}จตฺตาริมานิ ภิกฺขเว โสตาปตฺติยงฺคานิ. กตมานิ จตฺตาริ, สปฺปุริส\-สํเสโว สทฺธมฺม\-สฺสวนํ โยนิโส\-มนสิกาโร ธมฺมา\-นุธมฺม\-ปฺปฏิปตฺติ. อิมานิ โข ภิกฺขเว จตฺตาริ โสตาปตฺติยงฺคานีติ.  ทสมํ.}{สคาถก\-ปุญฺญาภิสนฺท\-วคฺโค ปญฺจโม.}
\csromancenter{\textbf{ตสฺสุทฺทานํ}}
\setlength{\csromangathapagewidth}{0pt}
\setlength{\csromangathawakwidth}{0pt}
\csromangathameasuregroup{อภิสนฺทา ตโย วุตฺตา, \\ สุทฺธํ นนฺทิยํ ภทฺทิยํ,}{\csromangathabat{อภิสนฺทา ตโย วุตฺตา,}{ทุเว มหทฺธเนน จ.} \\ \csromangathabat{สุทฺธํ นนฺทิยํ ภทฺทิยํ,}{มหานามงฺเคน เต ทสาติ.}}
\csromangathasetleft{อภิสนฺทา ตโย วุตฺตา, \\ สุทฺธํ นนฺทิยํ ภทฺทิยํ,}
\csromangathagroup{\csromangathastanza{\csromangathabat{อภิสนฺทา ตโย วุตฺตา,}{ทุเว มหทฺธเนน จ.} \\ \csromangathabat{สุทฺธํ นนฺทิยํ ภทฺทิยํ,}{มหานามงฺเคน เต ทสาติ.}}}

\csromanmatikaanchor{mtk.253}
\csromantocmarknopage{subparagraph}{6. สปฺปญฺญวคฺค}{mtk.253}
\bookmark[dest={mtk.253},level=5]{6. สปฺปญฺญวคฺค}
\csromanheader{6}{6. \textbf{สปฺปญฺญวคฺค}}
\csromanmatikaanchor{mtk.254}
\csromantocmarkref{subparagraph}{1. สคาถกสุตฺต}{mtk.254}
\bookmark[dest={mtk.254},level=5]{1. สคาถกสุตฺต}
\csromanheader{6}{1. \textbf{สคาถกสุตฺต}}
\proseitem{1047}{จตูหิ ภิกฺขเว ธมฺเมหิ สมนฺนาคโต อริยสาวโก โสตาปนฺโน โหติ อวินิปาต\-ธมฺโม นิยโต สมฺโพธิ\-ปรายโณ.}

\prose{กตเมหิ จตูหิ, อิธ ภิกฺขเว อริยสาวโก พุทฺเธ อเวจฺจ\-ปฺปสาเทน สมนฺนาคโต โหติ “อิติปิ โส ภควา ฯเปฯ สตฺถา เทวมนุสฺสานํ พุทฺโธ ภควา”ติ. ธมฺเม ฯเปฯ. สํเฆ ฯเปฯ. อริยกนฺเตหิ สีเลหิ สมนฺนาคโต โหติ อขณฺเฑหิ ฯเปฯ สมาธิ\-สํวตฺตนิเกหิ. อิเมหิ โข ภิกฺขเว จตูหิ ธมฺเมหิ สมนฺนาคโต อริยสาวโก โสตาปนฺโน โหติ อวินิปาต\-ธมฺโม นิยโต สมฺโพธิ\-ปรายโณติ. อิทมโวจ ภควา, อิทํ วตฺวาน สุคโต อถาปรํ เอตทโวจ สตฺถา\csromandash{}}

\setlength{\csromangathapagewidth}{0pt}
\setlength{\csromangathawakwidth}{0pt}
\csromangathameasuregroup{“ยสฺส สทฺธา ตถาคเต, \\ สีลญฺจ ยสฺส กลฺยาณํ, \\ สํเฆ ปสาโท ยสฺสตฺถิ, \\ อทลิทฺโทติ ตํ อาหุ, \\ ตสฺมา สทฺธญฺจ สีลญฺจ,}{\csromangathabat{“ยสฺส สทฺธา ตถาคเต,}{อจลา สุปฺปติฏฺฐิตา.} \\ \csromangathabat{สีลญฺจ ยสฺส กลฺยาณํ,}{อริยกนฺตํ ปสํสิตํ.} \\ \csromangathabat{สํเฆ ปสาโท ยสฺสตฺถิ,}{อุชุภูตญฺจ ทสฺสนํ.} \\ \csromangathabat{อทลิทฺโทติ ตํ อาหุ,}{อโมฆํ ตสฺส ชีวิตํ.} \\ \csromangathabat{ตสฺมา สทฺธญฺจ สีลญฺจ,}{ปสาทํ ธมฺมทสฺสนํ.}}
\csromangathasetleft{“ยสฺส สทฺธา ตถาคเต, \\ สีลญฺจ ยสฺส กลฺยาณํ, \\ สํเฆ ปสาโท ยสฺสตฺถิ, \\ อทลิทฺโทติ ตํ อาหุ, \\ ตสฺมา สทฺธญฺจ สีลญฺจ,}
\csromangathagroup{\csromangathastanza{\csromangathabat{“ยสฺส สทฺธา ตถาคเต,}{อจลา สุปฺปติฏฺฐิตา.} \\ \csromangathabat{สีลญฺจ ยสฺส กลฺยาณํ,}{อริยกนฺตํ ปสํสิตํ.}}\csromangathastanza{\csromangathabat{สํเฆ ปสาโท ยสฺสตฺถิ,}{อุชุภูตญฺจ ทสฺสนํ.} \\ \csromangathabat{อทลิทฺโทติ ตํ อาหุ,}{อโมฆํ ตสฺส ชีวิตํ.} \\ \csromangathabat{ตสฺมา สทฺธญฺจ สีลญฺจ,}{ปสาทํ ธมฺมทสฺสนํ.}}}

\vspace{\tipitakaparskip}
\csromancenter{อนุยุญฺเชถ เมธาวี, สรํ พุทฺธานสาสนนฺ”ติ. ปฐมํ.}
\csromanmatikaanchor{mtk.255}
\csromantocmarkref{subparagraph}{2. วสฺสํวุตฺถสุตฺต}{mtk.255}
\bookmark[dest={mtk.255},level=5]{2. วสฺสํวุตฺถสุตฺต}
\csromanheader{6}{2. วสฺสํวุตฺถ\-สุตฺต}
\proseitem{1048}{\csromanfolio{355}เอกํ สมยํ ภควา สาวตฺถิยํ วิหรติ เชตวเน อนาถ\-ปิณฺฑิกสฺส อาราเม. เตน โข ปน สมเยน อญฺญตโร ภิกฺขุ สาวตฺถิยํ วสฺสํวุตฺโถ กปิลวตฺถุํ อนุปฺปตฺโต โหติ เกนจิเทว กรณีเยน. อสฺโสสุํ โข กาปิลวตฺถวา สกฺยา “อญฺญตโร กิร ภิกฺขุ สาวตฺถิยํ วสฺสํวุตฺโถ กปิลวตฺถุํ อนุปฺปตฺโต”ติ.}

\prose{อถ โข กาปิลวตฺถวา สกฺยา เยน โส ภิกฺขุ เตนุ\-ปสงฺกมิํสุ, อุปสงฺกมิตฺวา ตํ ภิกฺขุํ อภิวาเทตฺวา เอกมนฺตํ นิสีทิํสุ, เอกมนฺตํ นิสินฺนา โข กาปิลวตฺถวา สกฺยา ตํ ภิกฺขุํ เอตทโวจุํ “กจฺจิ ภนฺเต ภควา อโรโค เจว พลวา จาติ. อโรโค จาวุโส ภควา พลวา จาติ. กจฺจิ ปน ภนฺเต สาริปุตฺต\-โมคฺคลฺลานา อโรคา เจว พลวนฺโต จาติ. สาริปุตฺต\-โมคฺคลฺลานาปิ โข อาวุโส อโรคา เจว พลวนฺโต จาติ. กจฺจิ ปน ภนฺเต ภิกฺขุสํโฆ อโรโค จ พลวา จาติ. ภิกฺขุ\-สํโฆปิ โข อาวุโส อโรโค จ พลวา จาติ. อตฺถิ ปน เต ภนฺเต กิญฺจิ อิมินา อนฺตรวสฺเสน ภควโต สมฺมุขา สุตํ, สมฺมุขา ปฏิคฺคหิตนฺติ. สมฺมุขา เมตํ อาวุโส ภควโต สุตํ, สมฺมุขา ปฏิคฺคหิตํ “อปฺปกา เต ภิกฺขเว ภิกฺขู, เย อาสวานํ ขยา อนาสวํ เจโตวิมุตฺติํ ปญฺญาวิมุตฺติํ ทิฏฺเฐว ธมฺเม สยํ อภิญฺญา สจฺฉิกตฺวา อุปสมฺปชฺช วิหรนฺติ, อถ โข เอเตว พหุตรา ภิกฺขู, เย ปญฺจนฺนํ โอรมฺภาคิยานํ สํโยชนานํ ปริกฺขยา โอปปาติกา, ตตฺถ ปรินิพฺพายิโน อนาวตฺติธมฺมา ตสฺมา โลกา”ติ.}

\prose{อปรมฺปิ โข เม อาวุโส ภควโต สมฺมุขา สุตํ, สมฺมุขา ปฏิคฺคหิตํ “อปฺปกา เต ภิกฺขเว ภิกฺขู, เย ปญฺจนฺนํ โอรมฺภาคิยานํ สํโยชนานํ ปริกฺขยา โอปปาติกา, ตตฺถ ปรินิพฺพายิโน อนาวตฺติธมฺมา ตสฺมา โลกา, อถ โข เอเตว พหุตรา ภิกฺขู, เย ติณฺณํ สํโยชนานํ ปริกฺขยา ราค\-โทส\-โมหานํ ตนุตฺตา สกทาคามิโน สกิเทว อิมํ โลกํ อาคนฺตฺวา ทุกฺขสฺสนฺตํ กริสฺสนฺตี”ติ.}

\prose{อปรมฺปิ โข เม อาวุโส ภควโต สมฺมุขา สุตํ, สมฺมุขา ปฏิคฺคหิตํ “อปฺปกา เต ภิกฺขเว ภิกฺขู, เย ติณฺณํ สํโยชนานํ ปริกฺขยา ราค\-โทส\-โมหานํ ตนุตฺตา สกทาคามิโน สกิเทว อิมํ โลกํ \csromanfolio{356}อาคนฺตฺวา ทุกฺขสฺสนฺตํ กริสฺสนฺติ, อถ โข เอเตว พหุตรา ภิกฺขู, เย ติณฺณํ สํโยชนานํ ปริกฺขยา โสตาปนฺนา อวินิปาต\-ธมฺมา นิยตา สมฺโพธิ\-ปรายณาติ.  ทุติยํ.}

\csromanmatikaanchor{mtk.256}
\csromantocmarkref{subparagraph}{3. ธมฺมทินฺนสุตฺต}{mtk.256}
\bookmark[dest={mtk.256},level=5]{3. ธมฺมทินฺนสุตฺต}
\csromanheader{6}{3. ธมฺมทินฺน\-สุตฺต}
\proseitem{1049}{เอกํ สมยํ ภควา พาราณสิยํ วิหรติ อิสิปตเน มิคทาเย. อถ โข ธมฺมทินฺโน อุปาสโก ปญฺจหิ อุปาสกสเตหิ สทฺธิํ เยน ภควา เตนุปสงฺกมิ, อุปสงฺกมิตฺวา ภควนฺตํ อภิวาเทตฺวา เอกมนฺตํ นิสีทิ, เอกมนฺตํ นิสินฺโน โข ธมฺมทินฺโน อุปาสโก ภควนฺตํ เอตทโวจ “โอวทตุ โน ภนฺเต ภควา, อนุสาสตุ โน ภนฺเต ภควา, ยํ อมฺหากํ อสฺส ทีฆรตฺตํ หิตาย สุขายา”ติ.}

\prose{ตสฺมาติห โว ธมฺมทินฺน เอวํ สิกฺขิตพฺพํ “เย เต สุตฺตนฺตา ตถาคต\-ภาสิตา คมฺภีรา คมฺภีรตฺถา โลกุตฺตรา สุญฺญต\-ปฏิสํยุตฺตา, เต กาเลน กาลํ อุปสมฺปชฺช วิหริสฺสามา”ติ, เอวํ หิ โว ธมฺมทินฺน สิกฺขิตพฺพนฺติ. น โข เนตํ ภนฺเต สุกรํ อเมฺหหิ ปุตฺต\-สมฺพาธ\-สยนํ อชฺฌาวสนฺเตหิ กาสิกจนฺทนํ ปจฺจนุโภนฺเตหิ มาลา\-คนฺธ\-วิเลปนํ ธารยนฺเตหิ ชาตรูป\-รชตํ สาทิยนฺเตหิ, เย เต สุตฺตนฺตา ตถาคต\-ภาสิตา คมฺภีรา คมฺภีรตฺถา โลกุตฺตรา สุญฺญต\-ปฏิสํยุตฺตา, เต กาเลน กาลํ อุปสมฺปชฺช วิหริตุํ, เตสํ โน ภนฺเต ภควา อมฺหากํ ปญฺจสุ สิกฺขาปเทสุ ฐิตานํ อุตฺตริธมฺมํ เทเสตูติ.}

\prose{ตสฺมาติห โว ธมฺมทินฺน เอวํ สิกฺขิตพฺพํ “พุทฺเธ อเวจฺจ\-ปฺปสาเทน สมนฺนาคตา ภวิสฺสาม “อิติปิ โส ภควา ฯเปฯ สตฺถา เทวมนุสฺสานํ พุทฺโธ ภควา”ติ. ธมฺเม ฯเปฯ. สํเฆ ฯเปฯ. อริยกนฺเตหิ สีเลหิ สมนฺนาคตา ภวิสฺสาม อขณฺเฑหิ ฯเปฯ สมาธิสํวตฺตนิเกหีติ. เอวํ หิ โว ธมฺมทินฺน สิกฺขิตพฺพนฺติ.}

\prose{ยานิมานิ ภนฺเต ภควตา จตฺตาริ โสตาปตฺติยงฺคานิ เทสิตานิ, สํวิชฺชนฺเตเต ธมฺมา อเมฺหสุ, มยญฺจ เตสุ ธมฺเมสุ สนฺทิสฺสาม. มยํ หิ ภนฺเต พุทฺเธ อเวจฺจ\-ปฺปสาเทน สมนฺนาคตา “อิติปิ โส ภควา ฯเปฯ สตฺถา เทวมนุสฺสานํ พุทฺโธ ภควา”ติ. ธมฺเม ฯเปฯ. สํเฆ ฯเปฯ. อริยกนฺเตหิ สีเลหิ สมนฺนาคตา อขณฺเฑหิ ฯเปฯ สมาธิสํวตฺตนิเกหีติ. ลาภา}

\prose{\csromanfolio{357}โว ธมฺมทินฺน, สุลทฺธํ โว ธมฺมทินฺน, โสตาปตฺติผลํ ตุเมฺหหิ พฺยากตนฺติ.  ตติยํ.}

\csromanmatikaanchor{mtk.257}
\csromantocmarkref{subparagraph}{4. คิลานสุตฺต}{mtk.257}
\bookmark[dest={mtk.257},level=5]{4. คิลานสุตฺต}
\csromanheader{6}{4. \textbf{คิลานสุตฺต}}
\proseitem{1050}{เอกํ สมยํ ภควา สกฺเกสุ วิหรติ กปิล\-วตฺถุสฺมิํ นิโคฺรธาราเม. เตน โข ปน สมเยน สมฺพหุลา ภิกฺขู ภควโต จีวรกมฺมํ กโรนฺติ “นิฏฺฐิตจีวโร ภควา เตมาสจฺจเยน จาริกํ ปกฺกมิสฺสตี”ติ. อสฺโสสิ โข มหานาโม สกฺโก “สมฺพหุลา กิร ภิกฺขู ภควโต จีวรกมฺมํ กโรนฺติ ‘นิฏฺฐิตจีวโร ภควา เตมาสจฺจเยน จาริกํ ปกฺกมิสฺสตี’ติ”. อถ โข มหานาโม สกฺโก เยน ภควา เตนุปสงฺกมิ, อุปสงฺกมิตฺวา ภควนฺตํ อภิวาเทตฺวา เอกมนฺตํ นิสีทิ, เอกมนฺตํ นิสินฺโน โข มหานาโม สกฺโก ภควนฺตํ เอตทโวจ “สุตเมตํ ภนฺเต สมฺพหุลา กิร ภิกฺขู ภควโต จีวรกมฺมํ กโรนฺติ ‘นิฏฺฐิตจีวโร ภควา เตมาสจฺจเยน จาริกํ ปกฺกมิสฺสตี’ติ. น โข เนตํ\csromansharedfootnote{7fda7c43ca191972}{น โข เต เอตํ (สี, อิ)} ภนฺเต ภควโต สมฺมุขา สุตํ, สมฺมุขา ปฏิคฺคหิตํ ‘สปฺปญฺเญน อุปาสเกน สปฺปญฺโญ อุปาสโก อาพาธิโก ทุกฺขิโต พาฬฺหคิลาโน โอวทิตพฺโพ’ติ”.}

\prose{สปฺปญฺเญน มหานาม อุปาสเกน สปฺปญฺโญ อุปาสโก อาพาธิโก ทุกฺขิโต พาฬฺหคิลาโน จตูหิ อสฺสาสนีเยหิ ธมฺเมหิ อสฺสาเสตพฺโพ “อสฺสาสตายสฺมา, อตฺถายสฺมโต พุทฺเธ อเวจฺจปฺปสาโท ‘อิติปิ โส ภควา ฯเปฯ สตฺถา เทวมนุสฺสานํ พุทฺโธ ภควา’ติ. อสฺสาสตายสฺมา, อตฺถายสฺมโต ธมฺเม ฯเปฯ. สํเฆ ฯเปฯ. อริยกนฺตานิ สีลานิ อขณฺฑานิ ฯเปฯ สมาธิสํวตฺตนิกานี”ติ.}

\prose{สปฺปญฺเญน มหานาม อุปาสเกน สปฺปญฺโญ อุปาสโก อาพาธิโก ทุกฺขิโต พาฬฺหคิลาโน อิเมหิ จตูหิ อสฺสาสนีเยหิ ธมฺเมหิ อสฺสาเสตฺวา เอวมสฺส วจนีโย “อตฺถายสฺมโต มาตาปิตูสุ อเปกฺขา”ติ. โส เจ เอวํ วเทยฺย “อตฺถิ เม มาตาปิตูสุ อเปกฺขา”ติ. โส เอวมสฺส วจนีโย “อายสฺมา โข มาริโส มรณธมฺโม, สเจปายสฺมา มาตาปิตูสุ อเปกฺขํ กริสฺสติ มริสฺสเตว, โน \csromanfolio{358}เจปายสฺมา มาตาปิตูสุ อเปกฺขํ กริสฺสติ มริสฺสเตว. สาธายสฺมา ยา เต มาตาปิตูสุ อเปกฺขา, ตํ ปชหา”ติ.}

\prose{โส เจ เอวํ วเทยฺย “ยา เม มาตาปิตูสุ อเปกฺขา, สา ปหีนา”ติ, โส เอวมสฺส วจนีโย “อตฺถิ ปนายสฺมโต ปุตฺตทาเรสุ อเปกฺขา”ติ. โส เจ เอวํ วเทยฺย “อตฺถิ เม ปุตฺตทาเรสุ อเปกฺขา”ติ. โส เอวมสฺส วจนีโย “อายสฺมา โข มาริโส มรณธมฺโม, สเจปายสฺมา ปุตฺตทาเรสุ อเปกฺขํ กริสฺสติ มริสฺสเตว, โน เจปายสฺมา ปุตฺตทาเรสุ อเปกฺขํ กริสฺสติ มริสฺสเตว. สาธายสฺมา ยา เต ปุตฺตทาเรสุ อเปกฺขา, ตํ ปชหา”ติ.}

\prose{โส เจ เอวํ วเทยฺย “ยา เม ปุตฺตทาเรสุ อเปกฺขา, สา ปหีนาติ, โส เอวมสฺส วจนีโย “อตฺถิ ปนายสฺมโต มานุสเกสุ ปญฺจสุ กามคุเณสุ อเปกฺขา”ติ, โส เจ เอวํ วเทยฺย “อตฺถิ เม มานุสเกสุ ปญฺจสุ กามคุเณสุ อเปกฺขา”ติ, โส เอวมสฺส วจนีโย “มานุสเกหิ โข อาวุโส กาเมหิ ทิพฺพา กามา อภิกฺกนฺตตรา จ ปณีตตรา จ. สาธายสฺมา มานุสเกหิ กาเมหิ จิตฺตํ วุฏฺฐาเปตฺวา จาตุมหาราชิเกสุ\csromansharedfootnote{d07697775d53ad57}{จาตุมฺมหาราชิเกสุ (สี, สฺยา, กํ, อิ)} เทเวสุ จิตฺตํ อธิโมเจหี”ติ.}

\prose{โส เจ เอวํ วเทยฺย “มานุสเกหิ เม กาเมหิ จิตฺตํ วุฏฺฐิตํ, จาตุมหาราชิเกสุ เทเวสุ จิตฺตํ อธิโมจิตนฺ”ติ, โส เอวมสฺส วจนีโย “จาตุมหาราชิเกหิ โข อาวุโส เทเวหิ ตาวติํสา เทวา อภิกฺกนฺตตรา จ ปณีตตรา จ. สาธายสฺมา จาตุมหาราชิเกหิ เทเวหิ จิตฺตํ วุฏฺฐาเปตฺวา ตาวติํเสสุ เทเวสุ จิตฺตํ อธิโมเจหี”ติ.}

\prose{โส เจ เอวํ วเทยฺย “จาตุมหาราชิเกหิ เม เทเวหิ จิตฺตํ วุฏฺฐิตํ, ตาวติํเสสุ เทเวสุ จิตฺตํ อธิโมจิตนฺ”ติ, โส เอวมสฺส วจนีโย “ตาวติํเสหิ โข อาวุโส เทเวหิ ยามา เทวา ฯเปฯ ตุสิตา เทวา ฯเปฯ นิมฺมานรตี เทวา ฯเปฯ ปรนิมฺมิต\-วส\-วตฺตี เทวา ฯเปฯ ปรนิมฺมิต\-วส\-วตฺตีหิ โข อาวุโส เทเวหิ พฺรหฺมโลโก อภิกฺกนฺตตโร จ ปณีตตโร จ. สาธายสฺมา ปรนิมฺมิต\-วส\-วตฺตีหิ เทเวหิ จิตฺตํ วุฏฺฐาเปตฺวา พฺรหฺมโลเก จิตฺตํ อธิโมเจหี”ติ.}

\csromanmatikaanchor{mtk.258}
\csromantocmarkref{subparagraph}{5-11. โสตาปตฺติผลสุตฺตาทีนิ}{mtk.258}
\bookmark[dest={mtk.258},level=5]{5-11. โสตาปตฺติผลสุตฺตาทีนิ}
\prose{\csromanfolio{359}โส เจ เอวํ วเทยฺย “ปรนิมฺมิต\-วส\-วตฺตีหิ เม เทเวหิ จิตฺตํ วุฏฺฐิตํ, พฺรหฺมโลเก จิตฺตํ อธิโมจิตนฺ”ติ, โส เอวมสฺส วจนีโย “พฺรหฺมโลโกปิ โข อาวุโส อนิจฺโจ อทฺธุโว สกฺกาย\-ปริยาปนฺโน. สาธายสฺมา พฺรหฺมโลกา จิตฺตํ วุฏฺฐาเปตฺวา สกฺกายนิโรเธ จิตฺตํ อุปสํหราหี”ติ.}

\prose{โส เจ เอวํ วเทยฺย “พฺรหฺมโลกา เม จิตฺตํ วุฏฺฐิตํ, สกฺกายนิโรเธ จิตฺตํ อุปสํหรามี”ติ. เอวํ วิมุตฺต\-จิตฺตสฺส โข มหานาม อุปาสกสฺส อาสวา\csromansharedfootnote{32f8863c6563e04a}{วสฺสสต (สี, สฺยา)} วิมุตฺตจิตฺเตน ภิกฺขุนา น กิญฺจิ นานากรณํ วทามิ, ยทิทํ วิมุตฺติยา วิมุตฺตนฺติ.  จตุตฺถํ.}

\csromanheader{2}{5. โสตาปตฺติผล\-สุตฺต}
\proseitem{1051}{จตฺตาโรเม ภิกฺขเว ธมฺมา ภาวิตา พหุลีกตา โสตาปตฺติผล\-สจฺฉิกิริยาย สํวตฺตนฺติ. กตเม จตฺตาโร, สปฺปุริส\-สํเสโว สทฺธมฺม\-สฺสวนํ โยนิโส\-มนสิกาโร ธมฺมา\-นุธมฺม\-ปฺปฏิปตฺติ. อิเม โข ภิกฺขเว จตฺตาโร ธมฺมา ภาวิตา พหุลีกตา โสตาปตฺติผล\-สจฺฉิกิริยาย สํวตฺตนฺตีติ.  ปญฺจมํ.}

\csromanheader{2}{6. สกทาคามิ\-ผลสุตฺต}
\proseitem{1052}{จตฺตาโรเม ภิกฺขเว ธมฺมา ภาวิตา พหุลีกตา สกทาคามิ\-ผล\-สจฺฉิกิริยาย สํวตฺตนฺติ. กตเม จตฺตาโร ฯเปฯ สํวตฺตนฺตีติ.  ฉฏฺฐํ.}

\csromanheader{2}{7. อนาคามิผล\-สุตฺต}
\vspace{\tipitakaparskip}
\csromancenter{อนาคามิ\-ผล\-สจฺฉิกิริยาย สํวตฺตนฺตีติ.  สตฺตมํ.}
\csromanheader{2}{8. อรหตฺตผล\-สุตฺต}
\vspace{\tipitakaparskip}
\csromancenter{อรหตฺตผล\-สจฺฉิกิริยาย สํวตฺตนฺตีติ.  อฏฺฐมํ.}
\csromanheader{2}{9. ปญฺญาปฏิลาภ\-สุตฺต}
\vspace{\tipitakaparskip}
\csromancenter{ปญฺญา\-ปฏิลาภาย สํวตฺตนฺตีติ.  นวมํ.}
\csromanheader{2}{10. ปญฺญาวุทฺธิ\-สุตฺต}
\vspace{\tipitakaparskip}
\csromancenter{ปญฺญาวุทฺธิยา สํวตฺตนฺตีติ.  ทสมํ.}
\csromanheader{2}{11. ปญฺญาเวปุลฺล\-สุตฺต}
\csromanmatikaanchor{mtk.259}
\csromanmatikaanchor{mtk.260}
\csromantocmarknopage{subparagraph}{7. มหาปญฺญวคฺค}{mtk.259}
\bookmark[dest={mtk.259},level=5]{7. มหาปญฺญวคฺค}
\csromantocmarkref{subparagraph}{1-13. มหาปญฺญาสุตฺตาทีนิ}{mtk.260}
\bookmark[dest={mtk.260},level=5]{1-13. มหาปญฺญาสุตฺตาทีนิ}
\nitthitammidruled{ปญฺญาเวปุลฺลาย สํวตฺตนฺตีติ.  เอกาทสมํ. สปฺปญฺญวคฺโค ฉฏฺโฐ.}
\csromancenter{\textbf{ตสฺสุทฺทานํ}}
\setlength{\csromangathapagewidth}{0pt}
\setlength{\csromangathawakwidth}{0pt}
\csromangathameasuregroup{สคาถกํ วสฺสํวุตฺถํ, \\ จตุโร ผลา ปฏิลาโภ,}{\csromangathabat{สคาถกํ วสฺสํวุตฺถํ,}{ธมฺมทินฺนญฺจ คิลานํ.} \\ \csromangathabat{จตุโร ผลา ปฏิลาโภ,}{วุทฺธิ เวปุลฺลตาย จาติ.}}
\csromangathasetleft{สคาถกํ วสฺสํวุตฺถํ, \\ จตุโร ผลา ปฏิลาโภ,}
\csromangathagroup{\csromangathastanza{\csromanfolio{360}\csromangathabat{สคาถกํ วสฺสํวุตฺถํ,}{ธมฺมทินฺนญฺจ คิลานํ.} \\ \csromangathabat{จตุโร ผลา ปฏิลาโภ,}{วุทฺธิ เวปุลฺลตาย จาติ.}}}

\csromanheader{6}{7. \textbf{มหาปญฺญวคฺค}}
\csromanheader{2}{1. \textbf{มหาปญฺญาสุตฺต}}
\proseitem{1058}{จตฺตาโรเม ภิกฺขเว ธมฺมา ภาวิตา พหุลีกตา มหาปญฺญตาย สํวตฺตนฺติ. กตเม จตฺตาโร, สปฺปุริส\-สํเสโว สทฺธมฺม\-สฺสวนํ โยนิโส\-มนสิกาโร ธมฺมา\-นุธมฺม\-ปฺปฏิปตฺติ. อิเม โข ภิกฺขเว จตฺตาโร ธมฺมา ภาวิตา พหุลีกตา มหาปญฺญตาย สํวตฺตนฺตีติ.  ปฐมํ.}

\csromanheader{2}{2. ปุถุปญฺญา\-สุตฺต}
\vspace{\tipitakaparskip}
\csromancenter{ปุถุปญฺญตาย สํวตฺตนฺตีติ.  ทุติยํ.}
\csromanheader{2}{3. วิปุลปญฺญา\-สุตฺต}
\vspace{\tipitakaparskip}
\csromancenter{วิปุล\-ปญฺญตาย สํวตฺตนฺตีติ.  ตติยํ.}
\csromanheader{2}{4. คมฺภีรปญฺญา\-สุตฺต}
\vspace{\tipitakaparskip}
\csromancenter{คมฺภีร\-ปญฺญตาย สํวตฺตนฺตีติ.  จตุตฺถํ.}
\csromanheader{2}{5. อปฺปมตฺตปญฺญา\-สุตฺต}
\vspace{\tipitakaparskip}
\csromancenter{อปฺปมตฺต\-ปญฺญตาย สํวตฺตนฺตีติ.  ปญฺจมํ.}
\csromanheader{2}{6. ภูริปญฺญา\-สุตฺต}
\vspace{\tipitakaparskip}
\csromancenter{ภูรปญฺญตาย สํวตฺตนฺตีติ.  ฉฏฺฐํ.}
\csromanheader{2}{7. ปญฺญาพาหุลฺล\-สุตฺต}
\vspace{\tipitakaparskip}
\csromancenter{ปญฺญาพาหุลฺลาย สํวตฺตนฺตีติ.  สตฺตมํ.}
\csromanheader{2}{8. สีฆปญฺญา\-สุตฺต}
\vspace{\tipitakaparskip}
\csromancenter{สีฆปญฺญตาย สํวตฺตนฺตีติ.  อฏฺฐมํ.}
\csromanheader{2}{9. \textbf{ลหุปญฺญาสุตฺต}}
\vspace{\tipitakaparskip}
\csromancenter{ลหุปญฺญตาย สํวตฺตนฺตีติ.  นวมํ.}
\csromanheader{2}{10. \textbf{หาสปญฺญาสุตฺต}}
\vspace{\tipitakaparskip}
\csromancenter{หาสปญฺญตาย สํวตฺตนฺตีติ.  ทสมํ.}
\csromanheader{2}{11. ชวนปญฺญา\-สุตฺต}
\vspace{\tipitakaparskip}
\csromancenter{ชวน\-ปญฺญตาย สํวตฺตนฺตีติ.  เอกาทสมํ.}
\csromanheader{2}{12. ติกฺขปญฺญา\-สุตฺต}
\vspace{\tipitakaparskip}
\csromancenter{ติกฺข\-ปญฺญตาย สํวตฺตนฺตีติ. ทฺวาทสมํ.}
\csromanheader{2}{13. นิพฺเพธิก\-ปญฺญาสุตฺต}
\proseitemwithclosermidruled{1070}{\csromanfolio{361}นิพฺเพธิก\-ปญฺญตาย สํวตฺตนฺติ. กตเม จตฺตาโร สปฺปุริสํเสโว สทฺธมฺม\-สฺสวนํ โยนิโส\-มนสิกาโร ธมฺมา\-นุธมฺม\-ปฺปฏิปตฺติ. \csromanfolio{362}อิเม โข ภิกฺขเว จตฺตาโร ธมฺมา ภาวิตา พหุลีกตา นิพฺเพธิก\-ปญฺญตาย สํวตฺตนฺตีติ.  เตเรสมํ.}{มหาปญฺญวคฺโค สตฺตโม.}
\csromancenter{\textbf{ตสฺสุทฺทานํ}}
\setlength{\csromangathapagewidth}{0pt}
\setlength{\csromangathawakwidth}{0pt}
\csromangathameasuregroup{มหา ปุถุ วิปุล คมฺภีรํ, \\ สีฆลหุ หาส ชวน,}{\csromangathabat{มหา ปุถุ วิปุล คมฺภีรํ,}{อปฺปมตฺต ภูริ พาหุลฺลํ.} \\ \csromangathabat{สีฆลหุ หาส ชวน,}{ติกฺข นิพฺเพธิกาย จาติ.} \\ \textbf{โสตาปตฺติสํยุตฺตํ เอกาทสมํ.}}
\csromangathameasurewak{\textbf{โสตาปตฺติสํยุตฺตํ เอกาทสมํ.}}
\csromangathasetleft{มหา ปุถุ วิปุล คมฺภีรํ, \\ สีฆลหุ หาส ชวน,}
\csromangathagroup{\csromangathastanza{\csromangathabat{มหา ปุถุ วิปุล คมฺภีรํ,}{อปฺปมตฺต ภูริ พาหุลฺลํ.} \\ \csromangathabat{สีฆลหุ หาส ชวน,}{ติกฺข นิพฺเพธิกาย จาติ.}}\csromangathastanza{\csromangathawak{\textbf{โสตาปตฺติสํยุตฺตํ เอกาทสมํ.}}}}

\csromanmatikaanchor{mtk.261}
\csromanmatikaanchor{mtk.262}
\csromantocmarknopage{subparagraph}{12. สจฺจสํยุตฺต}{mtk.261}
\bookmark[dest={mtk.261},level=5]{12. สจฺจสํยุตฺต}
\csromantocmarknopage{subparagraph}{1. สมาธิวคฺค}{mtk.262}
\bookmark[dest={mtk.262},level=5]{1. สมาธิวคฺค}
\csromanheader{6}{1. \textbf{สมาธิวคฺค}}
\csromanmatikaanchor{mtk.263}
\csromantocmarkref{subparagraph}{1. สมาธิสุตฺต}{mtk.263}
\bookmark[dest={mtk.263},level=5]{1. สมาธิสุตฺต}
\csromanheader{6}{1. \textbf{สมาธิสุตฺต}}
\proseitem{1071}{\csromanfolio{363}สาวตฺถิ\-นิทานํ ฯเปฯ. สมาธิํ ภิกฺขเว ภาเวถ, สมาหิโต ภิกฺขเว ภิกฺขุ ยถาภูตํ ปชานาติ. กิญฺจ ยถาภูตํ ปชานาติ, “อิทํ ทุกฺขนฺ”ติ ยถาภูตํ ปชานาติ, “อยํ ทุกฺขสมุทโย”ติ ยถาภูตํ ปชานาติ, “อยํ ทุกฺขนิโรโธ”ติ ยถาภูตํ ปชานาติ, “อยํ ทุกฺขนิโรธ\-คามินี ปฏิปทา”ติ ยถาภูตํ ปชานาติ. สมาธิํ ภิกฺขเว ภาเวถ, สมาหิโต ภิกฺขเว ภิกฺขุ ยถาภูตํ ปชานาติ.}

\prose{ตสฺมาติห ภิกฺขเว “อิทํ ทุกฺขนฺ”ติ โยโค กรณีโย, “อยํ ทุกฺขสมุทโย”ติ โยโค กรณีโย, “อยํ ทุกฺขนิโรโธ”ติ โยโค กรณีโย, “อยํ ทุกฺขนิโรธ\-คามินี ปฏิปทา”ติ โยโค กรณีโยติ.  ปฐมํ.}

\csromanmatikaanchor{mtk.264}
\csromantocmarkref{subparagraph}{2. ปฏิสลฺลานสุตฺต}{mtk.264}
\bookmark[dest={mtk.264},level=5]{2. ปฏิสลฺลานสุตฺต}
\csromanheader{6}{2. ปฏิสลฺลาน\-สุตฺต}
\proseitem{1072}{ปฏิสลฺลาเน ภิกฺขเว โยคมาปชฺชถ, ปฏิสลฺลีโน ภิกฺขเว ภิกฺขุ ยถาภูตํ ปชานาติ. กิญฺจ ยถาภูตํ ปชานาติ, “อิทํ ทุกฺขนฺ”ติ ยถาภูตํ ปชานาติ, “อยํ ทุกฺขสมุทโย”ติ ยถาภูตํ ปชานาติ, “อยํ ทุกฺขนิโรโธ”ติ ยถาภูตํ ปชานาติ, “อยํ ทุกฺขนิโรธ\-คามินี ปฏิปทา”ติ ยถาภูตํ ปชานาติ. ปฏิสลฺลาเน ภิกฺขเว โยคมาปชฺชถ, ปฏิสลฺลีโน ภิกฺขเว ภิกฺขุ ยถาภูตํ ปชานาติ.}

\prose{ตสฺมาติห ภิกฺขเว “อิทํ ทุกฺขนฺ”ติ โยโค กรณีโย, “อยํ ทุกฺขสมุทโย”ติ โยโค กรณีโย, “อยํ ทุกฺขนิโรโธ”ติ โยโค กรณีโย, “อยํ ทุกฺขนิโรธ\-คามินี ปฏิปทา”ติ โยโค กรณีโยติ.  ทุติยํ.}

\csromanheader{2}{3. ปฐม\-กุลปุตฺต\-สุตฺต}
\csromanmatikaanchor{mtk.265}
\csromantocmarkref{subparagraph}{3-4. กุลปุตฺตสุตฺตานิ}{mtk.265}
\bookmark[dest={mtk.265},level=5]{3-4. กุลปุตฺตสุตฺตานิ}
\proseitem{1073}{\csromanfolio{364}เย หิ เกจิ ภิกฺขเว อตีตมทฺธานํ กุลปุตฺตา สมฺมา อคารสฺมา อนคาริยํ ปพฺพชิํสุ, สพฺเพ เต จตุนฺนํ อริยสจฺจานํ ยถาภูตํ อภิสมยาย. เย หิ เกจิ ภิกฺขเว อนาคตมทฺธานํ กุลปุตฺตา สมฺมา อคารสฺมา อนคาริยํ ปพฺพชิสฺสนฺติ, สพฺเพ เต จตุนฺนํ อริยสจฺจานํ ยถาภูตํ อภิสมยาย. เย หิ เกจิ ภิกฺขเว เอตรหิ กุลปุตฺตา สมฺมา อคารสฺมา อนคาริยํ ปพฺพชนฺติ, สพฺเพ เต จตุนฺนํ อริยสจฺจานํ ยถาภูตํ อภิสมยาย.}

\prose{กตเมสํ จตุนฺนํ, ทุกฺขสฺส อริยสจฺจสฺส ทุกฺข\-สมุทยสฺส อริยสจฺจสฺส ทุกฺข\-นิโรธสฺส อริยสจฺจสฺส ทุกฺขนิโรธ\-คามินิยา ปฏิปทาย อริยสจฺจสฺส. เย หิ เกจิ ภิกฺขเว อตีตมทฺธานํ กุลปุตฺตา สมฺมา อคารสฺมา อนคาริยํ ปพฺพชิํสุ ฯเปฯ ปพฺพชิสฺสนฺติ ฯเปฯ ปพฺพชนฺติ, สพฺเพ เต อิเมสํเยว จตุนฺนํ อริยสจฺจานํ ยถาภูตํ อภิสมยาย.}

\prose{ตสฺมาติห ภิกฺขเว “อิทํ ทุกฺขนฺ”ติ โยโค กรณีโย, “อยํ ทุกฺขสมุทโย”ติ โยโค กรณีโย, “อยํ ทุกฺขนิโรโธ”ติ โยโค กรณีโย, “อยํ ทุกฺขนิโรธ\-คามินี ปฏิปทา”ติ โยโค กรณีโยติ.  ตติยํ.}

\csromanheader{2}{4. ทุติย\-กุลปุตฺต\-สุตฺต}
\proseitem{1074}{เย หิ เกจิ ภิกฺขเว อตีตมทฺธานํ กุลปุตฺตา สมฺมา อคารสฺมา อนคาริยํ ปพฺพชิตา ยถาภูตํ อภิสเมสุํ, สพฺเพ เต จตฺตาริ อริยสจฺจานิ ยถาภูตํ อภิสเมสุํ. เย หิ เกจิ ภิกฺขเว อนาคตมทฺธานํ กุลปุตฺตา สมฺมา อคารสฺมา อนคาริยํ ปพฺพชิตา ยถาภูตํ อภิสเมสฺสนฺติ, สพฺเพ เต จตฺตาริ อริยสจฺจานิ ยถาภูตํ อภิสเมสฺสนฺติ. เย หิ เกจิ ภิกฺขเว เอตรหิ กุลปุตฺตา สมฺมา อคารสฺมา อนคาริยํ ปพฺพชิตา ยถาภูตํ อภิสเมนฺติ, สพฺเพ เต จตฺตาริ อริยสจฺจานิ ยถาภูตํ อภิสเมนฺติ.}

\prose{กตมานิ จตฺตาริ, ทุกฺขํ อริยสจฺจํ ทุกฺข\-สมุทยํ อริยสจฺจํ ทุกฺขนิโรธํ อริยสจฺจํ ทุกฺขนิโรธ\-คามินี ปฏิปทา อริยสจฺจํ. เย หิ เกจิ ภิกฺขเว อตีตมทฺธานํ กุลปุตฺตา สมฺมา อคารสฺมา อนคาริยํ ปพฺพชิตา ยถาภูตํ}

\csromanmatikaanchor{mtk.266}
\csromantocmarkref{subparagraph}{5-6. สมณพฺราหฺมณสุตฺตานิ}{mtk.266}
\bookmark[dest={mtk.266},level=5]{5-6. สมณพฺราหฺมณสุตฺตานิ}
\prose{\csromanfolio{365}อภิสเมสุํ ฯเปฯ อภิสเมสฺสนฺติ ฯเปฯ อภิสเมนฺติ, สพฺเพ เต อิมานิ จตฺตาริ อริยสจฺจานิ ยถาภูตํ อภิสเมนฺติ.}

\prose{ตสฺมาติห ภิกฺขเว “อิทํ ทุกฺขนฺ”ติ โยโค กรณีโย ฯเปฯ “อยํ ทุกฺขนิโรธ\-คามินี ปฏิปาทา”ติ โยโค กรณีโยติ.  จตุตฺถํ.}

\csromanheader{2}{5. ปฐม\-สมณพฺราหฺมณ\-สุตฺต}
\proseitem{1075}{เย หิ เกจิ ภิกฺขเว อตีตมทฺธานํ สมณา วา พฺราหฺมณา วา ยถาภูตํ อภิสมฺ\-โพชฺฌิํสุ, สพฺเพ เต จตฺตาริ อริยสจฺจานิ ยถาภูตํ อภิสมฺ\-โพชฺฌิํสุ. เย หิ เกจิ ภิกฺขเว อนาคตมทฺธานํ สมณา วา พฺราหฺมณา วา ยถาภูตํ อภิสมฺโพชฺฌิสฺสนฺติ, สพฺเพ เต จตฺตาริ อริยสจฺจานิ ยถาภูตํ อภิสมฺโพชฺฌิสฺสนฺติ. เย หิ เกจิ ภิกฺขเว เอตรหิ สมณา วา พฺราหฺมณา วา ยถาภูตํ อภิสมฺ\-โพชฺฌนฺติ, สพฺเพ เต จตฺตาริ อริยสจฺจานิ ยถาภูตํ อภิสมฺ\-โพชฺฌนฺติ.}

\prose{กตมานิ จตฺตาริ? ทุกฺขํ อริยสจฺจํ ฯเปฯ ทุกฺขนิโรธ\-คามินี ปฏิปทา อริยสจฺจํ. เย หิ เกจิ ภิกฺขเว อตีตมทฺธานํ สมณา วา พฺราหฺมณา วา ยถาภูตํ อภิสมฺ\-โพชฺฌิํสุ ฯเปฯ อภิสมฺโพชฺฌิสฺสนฺติ ฯเปฯ อภิสมฺ\-โพชฺฌนฺติ, สพฺเพ เต อิมานิ จตฺตาริ อริยสจฺจานิ ยถาภูตํ อภิสมฺ\-โพชฺฌนฺติ.}

\prose{ตสฺมาติห ภิกฺขเว “อิทํ ทุกฺขนฺ”ติ โยโค กรณีโย ฯเปฯ “อยํ ทุกฺขนิโรธ\-คามินี ปฏิปทา”ติ โยโค กรณีโยติ.  ปญฺจมํ.}

\csromanheader{2}{6. ทุติย\-สมณพฺราหฺมณ\-สุตฺต}
\proseitem{1076}{เย หิ เกจิ ภิกฺขเว อตีตมทฺธานํ สมณา วา พฺราหฺมณาวา ยถาภูตํ อภิสมฺพุทฺธํ ปกาเสสุํ, สพฺเพ เต จตฺตาริ อริยสจฺจานิ ยถาภูตํ อภิสมฺพุทฺธํ ปกาเสสุํ. เย หิ เกจิ ภิกฺขเว อนาคตมทฺธานํ สมณา วา พฺราหฺมณา วา ยถาภูตํ อภิสมฺพุทฺธํ ปกาเสสฺสนฺติ, สพฺเพ เต จตฺตาริ อริยสจฺจานิ ยถาภูตํ อภิสมฺพุทฺธํ ปกาเสสฺสนฺติ. เย หิ เกจิ ภิกฺขเว เอตรหิ สมณา วา พฺราหฺมณา วา ยถาภูตํ อภิสมฺพุทฺธํ ปกาเสนฺติ, สพฺเพ เต จตฺตาริ อริยสจฺจานิ ยถาภูตํ อภิสมฺพุทฺธํ ปกาเสนฺติ.}

\csromanmatikaanchor{mtk.267}
\csromantocmarkref{subparagraph}{7-9. วิตกฺกสุตฺตาทีนิ}{mtk.267}
\bookmark[dest={mtk.267},level=5]{7-9. วิตกฺกสุตฺตาทีนิ}
\prose{\csromanfolio{366}กตมานิ จตฺตาริ, ทุกฺขํ อริย สจฺจํ ฯเปฯ ทุกฺขนิโรธ\-คามินี ปฏิปทา อริยสจฺจํ. เย หิ เกจิ ภิกฺขเว อตีตมทฺธานํ สมณา วา พฺราหฺมณา วา ยถาภูตํ อภิสมฺพุทฺธํ ปกาเสสุํ ฯเปฯ ปกาเสสฺสนฺติ ฯเปฯ ปกาเสนฺติ, สพฺเพ เต อิมานิ จตฺตาริ อริยสจฺจานิ ยถาภูตํ อภิสมฺพุทฺธํ ปกาเสนฺติ.}

\prose{ตสฺมาติห ภิกฺขเว “อิทํ ทุกฺขนฺ”ติ โยโค กรณีโย ฯเปฯ “อยํ ทุกฺขนิโรธ\-คามินี ปฏิปทา”ติ โยโค กรณีโยติ.  ฉฏฺฐํ.}

\csromanheader{2}{7. \textbf{วิตกฺกสุตฺต}}
\proseitem{1077}{มา ภิกฺขเว ปาปเก อกุสเล วิตกฺเก วิตกฺเกยฺยาถ\csromansharedfootnote{125d855ebf535165}{วิตกฺเกถ (สี, สฺยา, กํ)}. เสยฺยถิทํ, กามวิตกฺกํ พฺยาปาท\-วิตกฺกํ วิหิํสา\-วิตกฺกํ. ตํ กิสฺส เหตุ, เนเต ภิกฺขเว วิตกฺกา อตฺถสํหิตา, นาทิ\-พฺรหฺม\-จริยกา, น นิพฺพิทาย น วิราคาย น นิโรธาย น อุปสมาย น อภิญฺญาย น สมฺโพธาย น นิพฺพานาย สํวตฺตนฺติ.}

\prose{วิตกฺเกนฺตา จ โข ตุเมฺห ภิกฺขเว “อิทํ ทุกฺขนฺ”ติ วิตกฺเกยฺยาถ, “อยํ ทุกฺขสมุทโย”ติ วิตกฺเกยฺยาถ, “อยํ ทุกฺขนิโรโธ”ติ วิตกฺเกยฺยาถ, “อยํ ทุกฺขนิโรธ\-คามินี ปฏิปทา”ติ วิตกฺเกยฺยาถ. ตํ กิสฺส เหตุ, เอเต ภิกฺขเว วิตกฺกา อตฺถสํหิตา, เอเต อาทิ\-พฺรหฺม\-จริยกา, เอเต นิพฺพิทาย วิราคาย นิโรธาย อุปสมาย อภิญฺญาย สมฺโพธาย นิพฺพานาย สํวตฺตนฺติ.}

\prose{ตสฺมาติห ภิกฺขเว “อิทํ ทุกฺขนฺ”ติ โยโค กรณีโย ฯเปฯ “อยํ ทุกฺขนิโรธ\-คามินี ปฏิปทา”ติ โยโค กรณีโยติ.  สตฺตมํ.}

\csromanheader{2}{8. \textbf{จินฺตสุตฺต}}
\proseitem{1078}{มา ภิกฺขเว ปาปกํ อกุสลํ จิตฺตํ จินฺเตยฺยาถ\csromansharedfootnote{ecda706c7a974102}{จินฺเตถ (สี, สฺยา, กํ)}, “สสฺสโต โลโก”ติ วา “อสสฺสโต โลโก”ติ วา, “อนฺตวา โลโก”ติ วา, “อนนฺตวา โลโก”ติ วา, “ตํ ชีวํ ตํ สรีรนฺ”ติ วา, อญฺญํ ชีวํ “อญฺญํ สรีรนฺ”ติ วา, “โหติ ตถาคโต ปรํ มรณา”ติ วา, “น โหติ ตถาคโต ปรํ มรณา”ติ วา, “โหติ จ น จ}

\prose{\csromanfolio{367}โหติ ตถาคโต ปรํ มรณา”ติ วา, “เนว โหติ น น โหติ ตถาคโต ปรํ มรณา”ติ วา. ตํ กิสฺส เหตุ, เนสา ภิกฺขเว จินฺตา อตฺถสํหิตา, นาทิ\-พฺรหฺม\-จริยกา, น นิพฺพิทาย น วิราคาย น นิโรธาย น อุปสมาย น อภิญฺญาย น สมฺโพธาย น นิพฺพานาย สํวตฺตติ.}

\prose{จินฺเตนฺตา จ โข ตุเมฺห ภิกฺขเว “อิทํ ทุกฺขนฺ”ติ จินฺเตยฺยาถ, “อยํ ทุกฺขสมุทโย”ติ จินฺเตยฺยาถ, “อยํ ทุกฺขนิโรโธ”ติ จินฺเตยฺยาถ, “อยํ ทุกฺขนิโรธ\-คามินี ปฏิปทา”ติ จินฺเตยฺยาถ. ตํ กิสฺส เหตุ, เอสา ภิกฺขเว จินฺตา อตฺถสํหิตา, เอสา อาทิ\-พฺรหฺม\-จริยกา, เอสา นิพฺพิทาย วิราคาย นิโรธาย อุปสมาย อภิญฺญาย สมฺโพธาย นิพฺพานาย สํวตฺตติ.}

\prose{ตสฺมาติห ภิกฺขเว “อิทํ ทุกฺขนฺ”ติ โยโค กรณีโย ฯเปฯ “อยํ ทุกฺขนิโรธ\-คามินี ปฏิปทา”ติ โยโค กรณีโยติ.  อฏฺฐมํ.}

\csromanheader{2}{9. วิคฺคาหิกกถา\-สุตฺต}
\proseitem{1079}{มา ภิกฺขเว วิคฺคาหิกกถํ กเถยฺยาถ\csromansharedfootnote{1a7a158975a20c24}{กเถถ (สี, สฺยา, กํ)} “น ตฺวํ อิมํ ธมฺมวินยํ อาชานาสิ, อหํ อิมํ ธมฺมวินยํ อาชานามิ. กิํ ตฺวํ อิมํ ธมฺมวินยํ อาชานิสฺสสิ, มิจฺฉา\-ปฏิปนฺโน ตฺวมสิ, อหมสฺมิ สมฺมาปฏิปนฺโน. สหิตํ เม, อสหิตํ เต, ปุเรวจนียํ ปจฺฉา อวจ, ปจฺฉา\-วจนียํ ปุเร อวจ. อธิจิณฺณํ\csromansharedfootnote{46f98a1bf49b99ff}{อาจิณฺณํ (สฺยา, กํ, อิ)} เต วิปราวตฺตํ, อาโรปิโต เต วาโท, จร วาท\-ปฺปโมกฺขาย นิคฺคหิโตสิ, นิพฺเพเฐหิ วา สเจ ปโหสี”ติ. ตํ กิสฺส เหตุ, เนสา ภิกฺขเว กถา อตฺถสํหิตา, นาทิ\-พฺรหฺม\-จริยกา, น นิพฺพิทาย น วิราคาย น นิโรธาย น อุปสมาย น อภิญฺญาย น สมฺโพธาย น นิพฺพานาย สํวตฺตติ.}

\prose{กเถนฺตา จ โข ตุเมฺห ภิกฺขเว “อิทํ ทุกฺขนฺ”ติ กเถยฺยาถ, “อยํ ทุกฺขสมุทโย”ติ กเถยฺยาถ, “อยํ ทุกฺขนิโรโธ”ติ กเถยฺยาถ, “อยํ ทุกฺขนิโรธ\-คามินี ปฏิปทา”ติ กเถยฺยาถ ฯเปฯ โยโค กรณีโยติ.  นวมํ.}

\csromanmatikaanchor{mtk.268}
\csromantocmarkref{subparagraph}{10. ติรจฺฉานกถาสุตฺต}{mtk.268}
\bookmark[dest={mtk.268},level=5]{10. ติรจฺฉานกถาสุตฺต}
\csromanheader{6}{10. ติรจฺฉาน\-กถาสุตฺต}
\proseitem{1080}{\csromanfolio{368}มา ภิกฺขเว อเนกวิหิตํ ติรจฺฉานกถํ กเถยฺยาถ. เสยฺยถิทํ, ราชกถํ โจรกถํ มหามตฺตกถํ เสนากถํ ภยกถํ ยุทฺธกถํ อนฺนกถํ ปานกถํ วตฺถกถํ สยนกถํ มาลากถํ คนฺธกถํ ญาติกถํ ยานกถํ คามกถํ นิคมกถํ นครกถํ ชนปทกถํ อิตฺถิกถํ\csromansharedfootnote{7f5e17a7ecf1935b}{อิตฺถิกถํ ปุริสกถํ (สฺยา, กํ, อิ, ก)} สูรกถํ วิสิขากถํ กุมฺภ\-ฏฺฐานกถํ ปุพฺพเปตกถํ นานตฺตกถํ โลกกฺขายิกํ สมุทฺท\-กฺขายิกํ อิติภวา\-ภวกถํ อิติ วา. ตํ กิสฺส เหตุ, เนสา ภิกฺขเว กถา อตฺถสํหิตา, นาทิ\-พฺรหฺม\-จริยกา, น นิพฺพิทาย น วิราคาย น นิโรธาย น อุปสมาย น อภิญฺญาย น สมฺโพธาย น นิพฺพานาย สํวตฺตติ.}

\prose{กเถนฺตา จ โข ตุเมฺห ภิกฺขเว “อิทํ ทุกฺขนฺ”ติ กเถยฺยาถ, “อยํ ทุกฺขสมุทโย”ติ กเถยฺยาถ, “อยํ ทุกฺขนิโรโธ”ติ กเถยฺยาถ, “อยํ ทุกฺขนิโรธ\-คามินี ปฏิปทา”ติ กเถยฺยาถ. ตํ กิสฺส เหตุ, เอสา ภิกฺขเว กถา อตฺถสํหิตา, เอสา อาทิ\-พฺรหฺม\-จริยกา, เอสา นิพฺพิทาย วิราคาย นิโรธาย อุปสมาย อภิญฺญาย สมฺโพธาย นิพฺพานาย สํวตฺตติ.}

\prosewithclosermidruled{ตสฺมาติห ภิกฺขเว “อิทํ ทุกฺขนฺ”ติ โยโค กรณีโย ฯเปฯ “อยํ ทุกฺขนิโรธ\-คามินี ปฏิปทา”ติ โยโค กรณีโยติ.}{สมาธิวคฺโค ปฐโม.}
\csromancenter{\textbf{ตสฺสุทฺทานํ}}
\setlength{\csromangathapagewidth}{0pt}
\setlength{\csromangathawakwidth}{0pt}
\csromangathameasuregroup{สมาธิ ปฏิสลฺลานา, \\ สมณพฺราหฺมณา วิตกฺกํ,}{\csromangathabat{สมาธิ ปฏิสลฺลานา,}{กุลปุตฺตา อปเร ทุเว.} \\ \csromangathabat{สมณพฺราหฺมณา วิตกฺกํ,}{จินฺตา วิคฺคาหิกา กถาติ.}}
\csromangathasetleft{สมาธิ ปฏิสลฺลานา, \\ สมณพฺราหฺมณา วิตกฺกํ,}
\csromangathagroup{\csromangathastanza{\csromangathabat{สมาธิ ปฏิสลฺลานา,}{กุลปุตฺตา อปเร ทุเว.} \\ \csromangathabat{สมณพฺราหฺมณา วิตกฺกํ,}{จินฺตา วิคฺคาหิกา กถาติ.}}}

\csromanmatikaanchor{mtk.269}
\csromantocmarknopage{subparagraph}{2. ธมฺมจกฺกปฺปวตฺตนวคฺค}{mtk.269}
\bookmark[dest={mtk.269},level=5]{2. ธมฺมจกฺกปฺปวตฺตนวคฺค}
\csromanheader{6}{2. ธมฺมจกฺกปฺปวตฺตน\-วคฺค}
\csromanmatikaanchor{mtk.270}
\csromantocmarkref{subparagraph}{1. ธมฺมจกฺกปฺปวตฺตนสุตฺต}{mtk.270}
\bookmark[dest={mtk.270},level=5]{1. ธมฺมจกฺกปฺปวตฺตนสุตฺต}
\csromanheader{6}{1. ธมฺมจกฺกปฺปวตฺตน\-สุตฺต}
\proseitem{1081}{เอกํ สมยํ ภควา พาราณสิยํ วิหรติ อิสิปตเน มิคทาเย. ตตฺร โข ภควา ปญฺจวคฺคิเย ภิกฺขู อามนฺเตสิ\csromandash{}เทฺวเม ภิกฺขเว อนฺตา}

\prose{\csromanfolio{369}ปพฺพชิเตน น เสวิตพฺพา. กตเม เทฺว, โย จายํ กาเมสุ กามสุขลฺลิกา\-นุโยโค หีโน คมฺโม โปถุชฺชนิโก อนริโย อนตฺถสํหิโต, โย จายํ อตฺต\-กิลมถา\-นุโยโค ทุกฺโข อนริโย อนตฺถสํหิโต. เอเต โข ภิกฺขเว อุโภ อนฺเต อนุปคมฺม มชฺฌิมา ปฏิปทา ตถาคเตน อภิสมฺพุทฺธา จกฺขุกรณี ญาณกรณี อุปสมาย อภิญฺญาย สมฺโพธาย นิพฺพานาย สํวตฺตติ.}

\prose{กตมา จ สา ภิกฺขเว มชฺฌิมา ปฏิปทา ตถาคเตน อภิสมฺพุทฺธา จกฺขุกรณี ญาณกรณี อุปสมาย อภิญฺญาย สมฺโพธาย นิพฺพานาย สํวตฺตติ, อยเมว อริโย อฏฺฐงฺคิโก มคฺโค. เสยฺยถิทํ, สมฺมาทิฏฺฐิ สมฺมาสงฺกปฺโป สมฺมาวาจา สมฺมากมฺมนฺโต สมฺมาอาชีโว สมฺมาวายาโม สมฺมาสติ สมฺมาสมาธิ. อยํ โข สา ภิกฺขเว มชฺฌิมา ปฏิปทา ตถาคเตน อภิสมฺพุทฺธา จกฺขุกรณี ญาณกรณี อุปสมาย อภิญฺญาย สมฺโพธาย นิพฺพานาย สํวตฺตติ.}

\prose{อิทํ โข ปน ภิกฺขเว ทุกฺขํ อริยสจฺจํ, ชาติปิ ทุกฺขา, ชราปิ ทุกฺขา, พฺยาธิปิ ทุกฺโข, มรณมฺปิ ทุกฺขํ, อปฺปิเยหิ สมฺปโยโค ทุกฺโข, ปิเยหิ วิปฺปโยโค ทุกฺโข, ยมฺปิจฺฉํ น ลภติ, ตมฺปิ ทุกฺขํ, สํขิตฺเตน ปญฺจุ\-ปาทาน\-กฺขนฺธา\csromansharedfootnote{81c36fcbd9b6911b}{ปญฺจุปาทานกฺขนฺธาปิ (อิ, ก)} ทุกฺขา. อิทํ โข ปน ภิกฺขเว ทุกฺข\-สมุทยํ อริยสจฺจํ, ยายํ ตณฺหา โปโนพฺภวิกา\csromansharedfootnote{3459575e13e816f9}{โปโนภวิกา (สี, อิ)} นนฺทิ\-ราค\-สหคตา ตตฺร\-ตตฺรา\-ภินนฺทินี. เสยฺยถิทํ\csromansharedfootnote{8ad10b4e4b08c49a}{เสยฺยถีทํ (สี, สฺยา, กํ, อิ)}, กามตณฺหา ภวตณฺหา วิภวตณฺหา. อิทํ โข ปน ภิกฺขเว ทุกฺขนิโรธํ อริยสจฺจํ, โย ตสฺสาเยว ตณฺหาย อเสส\-วิราค\-นิโรโธ จาโค ปฏินิสฺสคฺโค มุตฺติ อนาลโย. อิทํ โข ปน ภิกฺขเว ทุกฺขนิโรธ\-คามินี ปฏิปทา อริยสจฺจํ, อยเมว อริโย อฏฺฐงฺคิโก มคฺโค. เสยฺยถิทํ, สมฺมาทิฏฺฐิ ฯเปฯ สมฺมาสมาธิ.}

\prose{“อิทํ ทุกฺขํ อริยสจฺจนฺ”ติ เม ภิกฺขเว ปุพฺเพ อนนุสฺสุเตสุ ธมฺเมสุ จกฺขุํ อุทปาทิ, ญาณํ อุทปาทิ, ปญฺญา อุทปาทิ, วิชฺชา อุทปาทิ, อาโลโก อุทปาทิ. “ตํ โข ปนิทํ ทุกฺขํ อริยสจฺจํ ปริญฺเญยฺยนฺ”ติ เม ภิกฺขเว ปุพฺเพ ฯเปฯ อุทปาทิ. “ตํ โข ปนิทํ ทุกฺขํ อริยสจฺจํ ปริญฺญาตนฺ”ติ เม ภิกฺขเว ปุพฺเพ อนนุสฺสุเตสุ ธมฺเมสุ จกฺขุํ อุทปาทิ, ญาณํ อุทปาทิ, ปญฺญา อุทปาทิ วิชฺชา อุทปาทิ, อาโลโก อุทปาทิ.}

\prose{\csromanfolio{370}“อิทํ ทุกฺข\-สมุทยํ อริยสจฺจนฺ”ติ เม ภิกฺขเว ปุพฺเพ อนนุสฺสุเตสุ ธมฺเมสุ จกฺขุํ อุทปาทิ, ญาณํ อุทปาทิ, ปญฺญา อุทปาทิ, วิชฺชา อุทปาทิ, อาโลโก อุทปาทิ. “ตํ โข ปนิทํ ทุกฺข\-สมุทยํ อริยสจฺจํ ปหาตพฺพนฺ”ติ เม ภิกฺขเว ปุพฺเพ ฯเปฯ อุทปาทิ. “ตํ โข ปนิทํ ทุกฺข\-สมุทยํ อริยสจฺจํ ปหีนนฺ”ติ เม ภิกฺขเว ปุพฺเพ อนนุสฺสุเตสุ ธมฺเมสุ จกฺขุํ อุทปาทิ, ญาณํ อุทปาทิ, ปญฺญา อุทปาทิ, วิชฺชา อุทปาทิ, อาโลโก อุทปาทิ.}

\prose{“อิทํ ทุกฺขนิโรธํ อริยสจฺจนฺ”ติ เม ภิกฺขเว ปุพฺเพ อนนุสฺสุเตสุ ธมฺเมสุ จกฺขุํ อุทปาทิ, ญาณํ อุทปาทิ, ปญฺญา อุทปาทิ, วิชฺชา อุทปาทิ, อาโลโก อุทปาทิ. “ตํ โข ปนิทํ ทุกฺขนิโรธํ อริยสจฺจํ สจฺฉิกาตพฺพนฺ”ติ เม ภิกฺขเว ปุพฺเพ ฯเปฯ อุทปาทิ. “ตํ โข ปนิทํ ทุกฺขนิโรธํ อริยสจฺจํ สจฺฉิกตนฺ”ติ เม ภิกฺขเว ปุพฺเพ อนนุสฺสุเตสุ ธมฺเมสุ จกฺขุํ อุทปาทิ, ญาณํ อุทปาทิ, ปญฺญา อุทปาทิ, วิชฺชา อุทปาทิ, อาโลโก อุทปาทิ.}

\prose{“อิทํ ทุกฺขนิโรธ\-คามินี ปฏิปทา อริยสจฺจนฺ”ติ เม ภิกฺขเว ปุพฺเพ อนนุสฺสุเตสุ ธมฺเมสุ จกฺขุํ อุทปาทิ, ญาณํ อุทปาทิ, ปญฺญา อุทปาทิ, วิชฺชา อุทปาทิ, อาโลโก อุทปาทิ. “ตํ โข ปนิทํ ทุกฺขนิโรธ\-คามินี ปฏิปทา อริยสจฺจํ ภาเวตพฺพนฺ”ติ เม ภิกฺขเว ปุพฺเพ ฯเปฯ อุทปาทิ. “ตํ โข ปนิทํ ทุกฺขนิโรธ\-คามินี ปฏิปทา อริยสจฺจํ ภาวิตนฺ”ติ เม ภิกฺขเว ปุพฺเพ อนนุสฺสุเตสุ ธมฺเมสุ จกฺขุํ อุทปาทิ, ญาณํ อุทปาทิ, ปญฺญา อุทปาทิ, วิชฺชา อุทปาทิ, อาโลโก อุทปาทิ.}

\prose{ยาวกีวญฺจ เม ภิกฺขเว อิเมสุ จตูสุ อริยสจฺเจสุ เอวํ ติปริวฏฺฏํ ทฺวาทสาการํ ยถาภูตํ ญาณทสฺสนํ น สุวิสุทฺธํ อโหสิ, เนว ตาวาหํ ภิกฺขเว สเทวเก โลเก สมารเก สพฺรหฺมเก สสฺสมณ\-พฺราหฺมณิยา ปชาย สเทว\-มนุสฺสาย อนุตฺตรํ สมฺมาสมฺโพธิํ อภิสมฺพุทฺโธติ ปจฺจญฺญาสิํ\csromansharedfootnote{27ebb5cb172b84b2}{อภิสมฺพุทฺโธ ปจฺจญฺญาสิํ (สี, สฺยา, กํ)}.}

\prose{ยโต จ โข เม ภิกฺขเว อิเมสุ จตูสุ อริยสจฺเจสุ เอวํ ติปริวฏฺฏํ ทฺวาทสาการํ ยถาภูตํ ญาณทสฺสนํ สุวิสุทฺธํ อโหสิ, อถาหํ ภิกฺขเว สเทวเก โลเก สมารเก สพฺรหฺมเก สสฺสมณ\-พฺราหฺมณิยา ปชาย สเทว\-มนุสฺสาย อนุตฺตรํ สมฺมาสมฺโพธิํ อภิสมฺพุทฺโธติ ปจฺจญฺญาสิํ, ญาณญฺจ ปน เม ทสฺสนํ อุทปาทิ, อกุปฺปา เม}

\csromanmatikaanchor{mtk.271}
\csromantocmarkref{subparagraph}{2-8. ตถาคตสุตฺตาทีนิ}{mtk.271}
\bookmark[dest={mtk.271},level=5]{2-8. ตถาคตสุตฺตาทีนิ}
\prose{\csromanfolio{371}วิมุตฺติ\csromansharedfootnote{510ae28ec8cfb582}{เจโตวิมุตฺติ (สี, อิ)}, อยมนฺติมา ชาติ, นตฺถิทานิ ปุนพฺภโวติ. อิทมโวจ ภควา, อตฺตมนา ปญฺจวคฺคิยา ภิกฺขู ภควโต ภาสิตํ อภินนฺทุนฺติ.}

\prose{อิมสฺมิํ จ ปน เวยฺยากรณสฺมิํ ภญฺญมาเน อายสฺมโต โกณฺฑญฺญสฺส วิรชํ วีตมลํ ธมฺมจกฺขุํ อุทปาทิ “ยํ กิญฺจิ สมุทย\-ธมฺมํ, สพฺพํ ตํ นิโรธธมฺมนฺ”ติ.}

\prose{ปวตฺติเต จ ปน ภควตา ธมฺมจกฺเก ภุมฺมา เทวา สทฺทม\-นุสฺสาเวสุํ “เอตํ ภควตา พาราณสิยํ อิสิปตเน มิคทาเย อนุตฺตรํ ธมฺมจกฺกํ ปวตฺติตํ อปฺปฏิวตฺติยํ สมเณน วา พฺราหฺมเณน วา เทเวน วา มาเรน วา พฺรหฺมุนา วา เกนจิ วา โลกสฺมินฺ”ติ. ภุมฺมานํ เทวานํ สทฺทํ สุตฺวา จาตุมหาราชิกา เทวา สทฺทม\-นุสฺสาเวสุํ “เอตํ ภควตา พาราณสิยํ อิสิปตเน มิคทาเย อนุตฺตรํ ธมฺมจกฺกํ ปวตฺติตํ อปฺปฏิวตฺติยํ สมเณน วา พฺราหฺมเณน วา เทเวน วา มาเรน วา พฺรหฺมุนา วา เกนจิ วา โลกสฺมินฺ”ติ. จาตุมหาราชิกานํ เทวานํ สทฺทํ สุตฺวา ตาวติํสา เทวา ฯเปฯ. ยามา เทวา ฯเปฯ. ตุสิตา เทวา ฯเปฯ. นิมฺมานรตี เทวา ฯเปฯ. ปรนิมฺมิต\-วส\-วตฺตี เทวา ฯเปฯ. พฺรหฺมกายิกา เทวา สทฺทม\-นุสฺสาเวสุํ “เอตํ ภควตา พาราณสิยํ อิสิปตเน มิคทาเย อนุตฺตรํ ธมฺมจกฺกํ ปวตฺติตํ อปฺปฏิวตฺติยํ สมเณน วา พฺราหฺมเณน วา เทเวน วา มาเรน วา พฺรหฺมุนา วา เกนจิ วา โลกสฺมินฺ”ติ.}

\prose{อิติห เตน ขเณน (เตน ลเยน)\csromansharedfootnote{48ea464a4f1db7c7}{( ) นตฺถิ (สี, สฺยา, กํ)} เตน มุหุตฺเตน ยาว พฺรหฺมโลกา สทฺโท อพฺภุคฺคจฺฉิ, อยญฺจ ทสสหสฺสิ\-โลกธาตุ สงฺกมฺปิ สมฺปกมฺปิ สมฺปเวธิ, อปฺปมาโณ จ อุฬาโร โอภาโส โลเก ปาตุรโหสิ อติกฺกมฺม เทวานํ เทวานุภาวนฺติ.}

\prose{อถ โข ภควา อิมํ อุทานํ อุทาเนสิ “อญฺญาสิ วต โภ โกณฺฑญฺโญ, อญฺญาสิ วต โภ โกณฺฑญฺโญ”ติ. อิติ หิทํ อายสฺมโต โกณฺฑญฺญสฺส “อญฺญาสิโกณฺฑญฺโญ”เตฺวว นามํ อโหสีติ.  ปฐมํ.}

\csromanheader{2}{2. \textbf{ตถาคตสุตฺต}}
\proseitem{1082}{“อิทํ ทุกฺขํ อริยสจฺจนฺ”ติ ภิกฺขเว ตถาคตานํ ปุพฺเพ อนนุสฺสุเตสุ ธมฺเมสุ จกฺขุํ อุทปาทิ, ญาณํ อุทปาทิ, ปญฺญา อุทปาทิ, วิชฺชา อุทปาทิ, อาโลโก อุทปาทิ. “ตํ โข ปนิทํ ทุกฺขํ อริยสจฺจํ \csromanfolio{372}ปริญฺเญยฺยนฺ”ติ ภิกฺขเว ตถาคตานํ ปุพฺเพ ฯเปฯ อุทปาทิ. “ตํ โข ปนิทํ ทุกฺขํ อริยสจฺจํ ปริญฺญาตนฺ”ติ ภิกฺขเว ตถาคตานํ ปุพฺเพ อนนุสฺสุเตสุ ธมฺเมสุ จกฺขุํ อุทปาทิ, ญาณํ อุทปาทิ, ปญฺญา อุทปาทิ, วิชฺชา อุทปาทิ, อาโลโก อุทปาทิ.}

\prose{“อิทํ ทุกฺข\-สมุทยํ อริยสจฺจนฺ”ติ ภิกฺขเว ตถาคตานํ ปุพฺเพ อนนุสฺสุเตสุ ธมฺเมสุ จกฺขุํ อุทปาทิ, ญาณํ อุทปาทิ, ปญฺญา อุทปาทิ, วิชฺชา อุทปาทิ, อาโลโก อุทปาทิ. “ตํ โข ปนิทํ ทุกฺข\-สมุทยํ อริยสจฺจํ ปหาตพฺพนฺ”ติ ภิกฺขเว ตถาคตานํ ปุพฺเพ ฯเปฯ อุทปาทิ. “ตํ โข ปนิทํ ทุกฺข\-สมุทยํ อริยสจฺจํ ปหีนนฺ”ติ ภิกฺขเว ตถาคตานํ ปุพฺเพ อนนุสฺสุเตสุ ธมฺเมสุ จกฺขุํ อุทปาทิ, ญาณํ อุทปาทิ, ปญฺญา อุทปาทิ, วิชฺชา อุทปาทิ, อาโลโก อุทปาทิ.}

\prose{“อิทํ ทุกฺขนิโรธํ อริยสจฺจนฺ”ติ ภิกฺขเว ตถาคตานํ ปุพฺเพ อนนุสฺสุเตสุ ธมฺเมสุ จกฺขุํ อุทปาทิ, ญาณํ อุทปาทิ, ปญฺญา อุทปาทิ, วิชฺชา อุทปาทิ, อาโลโก อุทปาทิ. “ตํ โข ปนิทํ ทุกฺขนิโรธํ อริยสจฺจํ สจฺฉิกาตพฺพนฺ”ติ ภิกฺขเว ตถาคตานํ ปุพฺเพ ฯเปฯ อุทปาทิ. “ตํ โข ปนิทํ ทุกฺขนิโรธํ อริยสจฺจํ สจฺฉิกตนฺ”ติ ภิกฺขเว ตถาคตานํ ปุพฺเพ อนนุสฺสุเตสุ ธมฺเมสุ จกฺขุํ อุทปาทิ, ญาณํ อุทปาทิ, ปญฺญา อุทปาทิ, วิชฺชา อุทปาทิ, อาโลโก อุทปาทิ.}

\prose{“อิทํ ทุกฺขนิโรธ\-คามินี ปฏิปทา อริยสจฺจนฺ”ติ ภิกฺขเว ตถาคตานํ ปุพฺเพ อนนุสฺสุเตสุ ธมฺเมสุ จกฺขุํ อุทปาทิ, ญาณํ อุทปาทิ, ปญฺญา อุทปาทิ, วิชฺชา อุทปาทิ, อาโลโก อุทปาทิ. “ตํ โข ปนิทํ ทุกฺขนิโรธ\-คามินี ปฏิปทา อริยสจฺจํ ภาเวตพฺพนฺ”ติ ภิกฺขเว ตถาคตานํ ปุพฺเพ ฯเปฯ อุทปาทิ. “ตํ โข ปนิทํ ทุกฺขนิโรธ\-คามินี ปฏิปทา อริยสจฺจํ ภาวิตนฺ”ติ ภิกฺขเว ตถาคตานํ ปุพฺเพ อนนุสฺสุเตสุ ธมฺเมสุ จกฺขุํ อุทปาทิ, ญาณํ อุทปาทิ, ปญฺญา อุทปาทิ, วิชฺชา อุทปาทิ, อาโลโก อุทปาทีติ.  ทุติยํ.}

\csromanheader{2}{3. \textbf{ขนฺธสุตฺต}}
\proseitem{1083}{จตฺตาริมานิ ภิกฺขเว อริยสจฺจานิ. กตมานิ จตฺตาริ, ทุกฺขํ อริยสจฺจํ ทุกฺข\-สมุทยํ อริยสจฺจํ ทุกฺขนิโรธํ อริยสจฺจํ\csromansharedfootnote{55ebdeae0a223439}{ทุกฺขสมุทโย อริยสจฺจํ ทุกฺขนิโรโธ อริยสจฺจํ (สฺยา)} ทุกฺขนิโรธ\-คามินี ปฏิปทา อริยสจฺจํ.}

\prose{\csromanfolio{373}กตมญฺจ ภิกฺขเว ทุกฺขํ อริยสจฺจํ, “ปญฺจุ\-ปาทาน\-กฺขนฺธา”ติสฺส วจนียํ. เสยฺยถิทํ\csromansharedfootnote{9a3b673f91d4ad48}{กตเม ปญฺจ (สี, สฺยา, กํ)}, รูปุปาทานกฺขนฺโท ฯเปฯ วิญฺญาณุ\-ปาทาน\-กฺขนฺโธ. อิทํ วุจฺจติ ภิกฺขเว ทุกฺขํ อริยสจฺจํ.}

\prose{กตมญฺจ ภิกฺขเว ทุกฺข\-สมุทยํ อริยสจฺจํ, ยายํ ตณฺหา โปโนพฺภวิกา นนฺทิ\-ราค\-สหคตา ตตฺร\-ตตฺรา\-ภินนฺทินี. เสยฺยถิทํ, กามตณฺหา ภวตณฺหา วิภวตณฺหา. อิทํ วุจฺจติ ภิกฺขเว ทุกฺข\-สมุทยํ อริยสจฺจํ.}

\prose{กตมญฺจ ภิกฺขเว ทุกฺขนิโรธํ อริยสจฺจํ, โย ตสฺสาเยว ตณฺหาย อเสส\-วิราค\-นิโรโธ จาโค ปฏินิสฺสคฺโค มุตฺติ อนาลโย. อิทํ วุจฺจติ ภิกฺขเว ทุกฺขนิโรธํ อริยสจฺจํ.}

\prose{กตมญฺจ ภิกฺขเว ทุกฺขนิโรธ\-คามินี ปฏิปทา อริยสจฺจํ, อยเมว อริโย อฏฺฐงฺคิโก มคฺโค. เสยฺยถิทํ, สมฺมาทิฏฺฐิ ฯเปฯ สมฺมาสมาธิ. อิทํ วุจฺจติ ภิกฺขเว ทุกฺขนิโรธ\-คามินี ปฏิปทา อริยสจฺจํ. อิมานิ โข ภิกฺขเว จตฺตาริ อริยสจฺจานิ.}

\prose{ตสฺมาติห ภิกฺขเว “อิทํ ทุกฺขนฺ”ติ โยโค กรณีโย ฯเปฯ “อยํ ทุกฺขนิโรธ\-คามินี ปฏิปทา”ติ โยโค กรณีโยติ.  ตติยํ.}

\csromanheader{2}{4. อชฺฌตฺติกายตน\-สุตฺต}
\proseitem{1084}{จตฺตาริมานิ ภิกฺขเว อริยสจฺจานิ. กตมานิ จตฺตาริ, ทุกฺขํ อริยสจฺจํ ทุกฺข\-สมุทยํ อริยสจฺจํ ทุกฺขนิโรธํ อริยสจฺจํ ทุกฺขนิโรธ\-คามินี ปฏิปทา อริยสจฺจํ.}

\prose{กตมญฺจ ภิกฺขเว ทุกฺขํ อริยสจฺจํ, “ฉ อชฺฌตฺติกานิ อายตนานี”ติสฺส วจนียํ. กตมานิ ฉ, จกฺขายตนํ ฯเปฯ มนายตนํ. อิทํ วุจฺจติ ภิกฺขเว ทุกฺขํ อริยสจฺจํ.}

\prose{กตมญฺจ ภิกฺขเว ทุกฺข\-สมุทยํ อริยสจฺจํ, ยายํ ตณฺหา โปโนพฺภวิกา นนฺทิ\-ราค\-สหคตา ตตฺร\-ตตฺรา\-ภินนฺทินี. เสยฺยถิทํ, กามตณฺหา ภวตณฺหา วิภวตณฺหา. อิทํ วุจฺจติ ภิกฺขเว ทุกฺข\-สมุทยํ อริยสจฺจํ.}

\prose{\csromanfolio{374}กตมญฺจ ภิกฺขเว ทุกฺขนิโรธํ อริยสจฺจํ, โย ตสฺสาเยว ตณฺหาย อเสส\-วิราค\-นิโรโธ จาโค ปฏินิสฺสคฺโค มุตฺติ อนาลโย. อิทํ วุจฺจติ ภิกฺขเว ทุกฺขนิโรธํ อริยสจฺจํ.}

\prose{กตมญฺจ ภิกฺขเว ทุกฺขนิโรธ\-คามินี ปฏิปทา อริยสจฺจํ, อยเมว อริโย อฏฺฐงฺคิโก มคฺโค. เสยฺยถิทํ, สมฺมาทิฏฺฐิ ฯเปฯ สมฺมาสมาธิ. อิทํ วุจฺจติ ภิกฺขเว ทุกฺขนิโรธ\-คามินี ปฏิปทา อริยสจฺจํ. อิมานิ โข ภิกฺขเว จตฺตาริ อริยสจฺจานิ.}

\prose{ตสฺมาติห ภิกฺขเว “อิทํ ทุกฺขนฺ”ติ โยโค กรณีโย ฯเปฯ “อยํ ทุกฺขนิโรธ\-คามินี ปฏิปทา”ติ โยโค กรณีโยติ.  จตุตฺถํ.}

\csromanheader{2}{5. ปฐม\-ธารณ\-สุตฺต}
\proseitem{1085}{ธาเรถ โน ตุเมฺห ภิกฺขเว “มยา จตฺตาริ อริยสจฺจานิ เทสิตานี”ติ. เอวํ วุตฺเต อญฺญตโร ภิกฺขุ ภควนฺตํ เอตทโวจ “อหํ โข ภนฺเต ธาเรมิ ภควตา จตฺตาริ อริยสจฺจานิ เทสิตานี”ติ. ยถา กถํ ปน ตฺวํ ภิกฺขุ ธาเรสิ “มยา จตฺตาริ อริยสจฺจานิ เทสิตานี”ติ. ทุกฺขํ ขฺวาหํ ภนฺเต ภควตา ปฐมํ อริยสจฺจํ เทสิตํ ธาเรมิ, ทุกฺข\-สมุทยํ ขฺวาหํ ภนฺเต ภควตา ทุติยํ อริยสจฺจํ เทสิตํ ธาเรมิ, ทุกฺขนิโรธํ ขฺวาหํ ภนฺเต ภควตา ตติยํ อริยสจฺจํ เทสิตํ ธาเรมิ, ทุกฺขนิโรธ\-คามินิํ ปฏิปทํ ขฺวาหํ ภนฺเต ภควตา จตุตฺถํ อริยสจฺจํ เทสิตํ ธาเรมิ. เอวํ ขฺวาหํ ภนฺเต ธาเรมิ “ภควตา จตฺตาริ อริยสจฺจานิ เทสิตานี”ติ.}

\prose{สาธุ สาธุ ภิกฺขุ, สาธุ โข ตฺวํ ภิกฺขุ ธาเรสิ “มยา จตฺตาริ อริยสจฺจานิ เทสิตานี”ติ. ทุกฺขํ โข ภิกฺขุ มยา ปฐมํ อริยสจฺจํ เทสิตํ, ตถา นํ ธาเรหิ, ทุกฺข\-สมุทยํ\csromansharedfootnote{92abb159022a2cc5}{ทุกฺขสมุทโย (สฺยา, กํ)} โข ภิกฺขุ มยา ทุติยํ อริยสจฺจํ, เทสิตํ, ตถา นํ ธาเรหิ, ทุกฺขนิโรธํ\csromansharedfootnote{72cf75c19b00299e}{ทุกฺขนิโรโธ (สฺยา, กํ)} โข ภิกฺขุ มยา ตติยํ อริยสจฺจํ เทสิตํ, ตถา นํ ธาเรหิ, ทุกฺขนิโรธ\-คามินี ปฏิปทา\csromansharedfootnote{bde6dccd39e149d0}{ทุกฺขนิโรธคามินิปฏิปทํ (อิ), ทุกฺขนิโรธคามินิํ ปฏิปทํ (ก)} โข ภิกฺขุ มยา จตุตฺถํ อริยสจฺจํ เทสิตํ, ตถา นํ ธาเรหิ. เอวํ โข ภิกฺขุ ธาเรหิ “มยา จตฺตาริ อริยสจฺจานิ เทสิตานี”ติ.}

\prose{\csromanfolio{375}ตสฺมาติห ภิกฺขุ “อิทํ ทุกฺขนฺ”ติ โยโค กรณีโย ฯเปฯ “อยํ ทุกฺขนิโรธ\-คามินี ปฏิปทา”ติ โยโค กรณีโยติ.  ปญฺจมํ.}

\csromanheader{2}{6. ทุติย\-ธารณ\-สุตฺต}
\proseitem{1086}{ธาเรถ โน ตุเมฺห ภิกฺขเว “มยา จตฺตาริ อริยสจฺจานิ เทสิตานี”ติ. เอวํ วุตฺเต อญฺญตโร ภิกฺขุ ภควนฺตํ เอตทโวจ “อหํ โข ภนฺเต ธาเรมิ ภควตา จตฺตาริ อริยสจฺจานิ เทสิตานี”ติ. ยถา กถํ ปน ตฺวํ ภิกฺขุ ธาเรสิ “มยา จตฺตาริ อริยสจฺจานิ เทสิตานี”ติ.}

\prose{ทุกฺขํ ขฺวาหํ ภนฺเต ภควตา ปฐมํ อริยสจฺจํ เทสิตํ ธาเรมิ. โย หิ โกจิ ภนฺเต สมโณ วา พฺราหฺมโณ วา เอวํ วเทยฺย “เนตํ ทุกฺขํ ปฐมํ อริยสจฺจํ, ยํ สมเณน โคตเมน เทสิตํ, อหเมตํ ทุกฺขํ ปฐมํ อริยสจฺจํ ปจฺจกฺขาย อญฺญํ ทุกฺขํ ปฐมํ อริยสจฺจํ ปญฺญเปสฺสามี”ติ เนตํ ฐานํ วิชฺชติ. ทุกฺข\-สมุทยํ ขฺวาหํ ภนฺเต ภควตา ฯเปฯ ทุกฺขนิโรธ\-คามินิํ ปฏิปทํ ขฺวาหํ ภนฺเต ภควตา จตุตฺถํ อริยสจฺจํ เทสิตํ ธาเรมิ, โย หิ โกจิ ภนฺเต สมโณ วา พฺราหฺมโณ วา เอวํ วเทยฺย “เนตํ ทุกฺขนิโรธ\-คามินี ปฏิปทา จตุตฺถํ อริยสจฺจํ, ยํ สมเณน โคตเมน เทสิตํ, อหเมตํ ทุกฺขนิโรธ\-คามินิํ ปฏิปทํ จตุตฺถํ อริยสจฺจํ ปจฺจกฺขาย อญฺญํ ทุกฺขนิโรธ\-คามินิํ ปฏิปทํ จตุตฺถํ อริยสจฺจํ ปญฺญเปสฺสามี”ติ เนตํ ฐานํ วิชฺชติ. เอวํ ขฺวาหํ ภนฺเต ธาเรมิ “ภควตา จตฺตาริ อริยสจฺจานิ เทสิตานี”ติ.}

\prose{สาธุ สาธุ ภิกฺขุ, สาธุ โข ตฺวํ ภิกฺขุ ธาเรสิ “มยา จตฺตาริ อริยสจฺจานิ เทสิตานี”ติ. ทุกฺขํ โข ภิกฺขุ มยา ปฐมํ อริยสจฺจํ เทสิตํ, ตถา นํ ธาเรหิ. โย หิ โกจิ ภิกฺขุ สมโณ วา พฺราหฺมโณ วา เอวํ วเทยฺย “เนตํ ทุกฺขํ ปฐมํ อริยสจฺจํ, ยํ สมเณน โคตเมน เทสิตํ, อหเมตํ ทุกฺขํ ปฐมํ อริยสจฺจํ ปจฺจกฺขาย อญฺญํ ทุกฺขํ ปฐมํ อริยสจฺจํ ปญฺญเปสฺสามี”ติ เนตํ ฐานํ วิชฺชติ. ทุกฺข\-สมุทยํ โข ภิกฺขุ ฯเปฯ. ทุกฺขนิโรธํ โข ภิกฺขุ ฯเปฯ. ทุกฺขนิโรธ\-คามินี ปฏิปทา โข ภิกฺขุ มยา จตุตฺถํ อริยสจฺจํ เทสิตํ, ตถา นํ ธาเรหิ. โย หิ โกจิ ภิกฺขุ สมโณ วา พฺราหฺมโณ วา เอวํ วเทยฺย “เนตํ ทุกฺขนิโรธ\-คามินี ปฏิปทา จตุตฺถํ อริยสจฺจํ, ยํ สมเณน โคตเมน \csromanfolio{376}เทสิตํ, อหเมตํ ทุกฺขนิโรธ\-คามินิํ ปฏิปทํ จตุตฺถํ อริยสจฺจํ ปจฺจกฺขาย อญฺญํ ทุกฺขนิโรธ\-คามินิํ ปฏิปทํ จตุตฺถํ อริยสจฺจํ ปญฺญเปสฺสามี”ติ เนตํ ฐานํ วิชฺชติ. เอวํ โข ตฺวํ ภิกฺขุ ธาเรหิ “มยา จตฺตาริ อริยสจฺจานิ เทสิตานี”ติ.}

\prose{ตสฺมาติห ภิกฺขุ “อิทํ ทุกฺขนฺ”ติ โยโค กรณีโย ฯเปฯ “อยํ ทุกฺขนิโรธ\-คามินี ปฏิปทา”ติ โยโค กรณีโยติ.  ฉฏฺฐํ.}

\csromanheader{2}{7. \textbf{อวิชฺชาสุตฺต}}
\proseitem{1087}{เอกมนฺตํ นิสินฺโน โข โส ภิกฺขุ ภควนฺตํ เอตทโวจ “อวิชฺชา อวิชฺชาติ ภนฺเต วุจฺจติ, กตมา นุ โข ภนฺเต อวิชฺชา, กิตฺตาวตา จ อวิชฺชาคโต โหตี”ติ. ยํ โข ภิกฺขุ ทุกฺเข อญฺญาณํ, ทุกฺขสมุทเย อญฺญาณํ, ทุกฺขนิโรเธ อญฺญาณํ, ทุกฺขนิโรธ\-คามินิยา ปฏิปทาย อญฺญาณํ. อยํ วุจฺจติ ภิกฺขุ อวิชฺชา, เอตฺตาวตา จ อวิชฺชาคโต โหตีติ.}

\prose{ตสฺมาติห ภิกฺขุ “อิทํ ทุกฺขนฺ”ติ โยโค กรณีโย ฯเปฯ “อยํ ทุกฺขนิโรธ\-คามินี ปฏิปทา”ติ โยโค กรณีโยติ.  สตฺตมํ.}

\csromanheader{2}{8. \textbf{วิชฺชาสุตฺต}}
\proseitem{1088}{อถ โข อญฺญตโร ภิกฺขุ เยน ภควา เตนุปสงฺกมิ, อุปสงฺกมิตฺวา ภควนฺตํ อภิวาเทตฺวา เอกมนฺตํ นิสีทิ, เอกมนฺตํ นิสินฺโน โข โส ภิกฺขุ ภควนฺตํ เอตทโวจ “วิชฺชา วิชฺชาติ ภนฺเต วุจฺจติ, กตมา นุ โข ภนฺเต วิชฺชา, กิตฺตาวตา จ วิชฺชาคโต โหตี”ติ. ยํ โข ภิกฺขุ ทุกฺเข ญาณํ, ทุกฺขสมุทเย ญาณํ, ทุกฺขนิโรเธ ญาณํ, ทุกฺขนิโรธ\-คามินิยา ปฏิปทาย ญาณํ. อยํ วุจฺจติ ภิกฺขุ วิชฺชา, เอตฺตาวตา จ วิชฺชาคโต โหตีติ.}

\prose{ตสฺมาติห ภิกฺขุ “อิทํ ทุกฺขนฺ”ติ โยโค กรณีโย ฯเปฯ “อยํ ทุกฺขนิโรธ\-คามินี ปฏิปทา”ติ โยโค กรณีโยติ.  อฏฺฐมํ.}

\csromanmatikaanchor{mtk.272}
\csromantocmarkref{subparagraph}{9. สงฺกาสนสุตฺต}{mtk.272}
\bookmark[dest={mtk.272},level=5]{9. สงฺกาสนสุตฺต}
\csromanheader{6}{9. \textbf{สงฺกาสนสุตฺต}}
\proseitem{1089}{“อิทํ ทุกฺขํ อริยสจฺจนฺ”ติ ภิกฺขเว มยา ปญฺญตฺตํ, ตตฺถ อปริมาณา วณฺณา อปริมาณา พฺยญฺชนา อปริมาณา สงฺกาสนา อิติปิทํ ทุกฺขํ}

\csromanmatikaanchor{mtk.273}
\csromanmatikaanchor{mtk.275}
\csromantocmarkref{subparagraph}{10. ตถสุตฺต}{mtk.273}
\bookmark[dest={mtk.273},level=5]{10. ตถสุตฺต}
\prose{\csromanfolio{377}อริยสจฺจนฺติ. อิทํ ทุกฺข\-สมุทยํ ฯเปฯ. อิทํ ทุกฺขนิโรธํ ฯเปฯ “อิทํ ทุกฺขนิโรธ\-คามินี ปฏิปทา อริยสจฺจนฺ”ติ ภิกฺขเว มยา ปญฺญตฺตํ, ตตฺถ อปริมาณา วณฺณา อปริมาณา พฺยญฺชนา อปริมาณา สงฺกาสนา อิติปิทํ ทุกฺขนิโรธ\-คามินี ปฏิปทา อริยสจฺจนฺติ.}

\prose{ตสฺมาติห ภิกฺขเว “อิทํ ทุกฺขํนฺ”ติ โยโค กรณีโย ฯเปฯ “อยํ ทุกฺขนิโรธ\-คามินี ปฏิปทา”ติ โยโค กรณีโยติ.  นวมํ.}

\csromanheader{6}{10. \textbf{ตถสุตฺต}}
\proseitem{1090}{จตฺตาริมานิ ภิกฺขเว ตถานิ อวิตถานิ อนญฺญถานิ. กตมานิ จตฺตาริ, “อิทํ ทุกฺขนฺ”ติ ภิกฺขเว ตถเมตํ อวิตถเมตํ อนญฺญถเมตํ, “อยํ ทุกฺขสมุทโย”ติ ตถเมตํ อวิตถเมตํ อนญฺญถเมตํ, “อยํ ทุกฺขนิโรโธ”ติ ตถเมตํ อวิตถเมตํ อนญฺญถเมตํ, “อยํ ทุกฺขนิโรธ\-คามินี ปฏิปทา”ติ ตถเมตํ อวิตถเมตํ อนญฺญถเมตํ. อิมานิ โข ภิกฺขเว จตฺตาริ ตถานิ อวิตถานิ อนญฺญถานิ.}

\prosewithclosermidruled{ตสฺมาติห ภิกฺขเว “อิทํ ทุกฺขนฺ”ติ โยโค กรณีโย ฯเปฯ “อยํ ทุกฺขนิโรธ\-คามินี ปฏิปทา”ติ โยโค กรณีโยติ.  ทสมํ.}{ธมฺมจกฺกปฺปวตฺตน\-วคฺโค ทุติโย.}
\csromancenter{\textbf{ตสฺสุทฺทานํ}}
\setlength{\csromangathapagewidth}{0pt}
\setlength{\csromangathawakwidth}{0pt}
\csromangathameasuregroup{ธมฺมจกฺกํ ตถาคตํ, \\ ธารณา จ เทฺว อวิชฺชา,}{\csromangathabat{ธมฺมจกฺกํ ตถาคตํ,}{ขนฺธา อายตเนน จ.} \\ \csromangathabat{ธารณา จ เทฺว อวิชฺชา,}{วิชฺชา สงฺกาสนา ตถาติ.}}
\csromangathasetleft{ธมฺมจกฺกํ ตถาคตํ, \\ ธารณา จ เทฺว อวิชฺชา,}
\csromangathagroup{\csromangathastanza{\csromangathabat{ธมฺมจกฺกํ ตถาคตํ,}{ขนฺธา อายตเนน จ.} \\ \csromangathabat{ธารณา จ เทฺว อวิชฺชา,}{วิชฺชา สงฺกาสนา ตถาติ.}}}

\csromanmatikaanchor{mtk.274}
\csromantocmarknopage{subparagraph}{3. โกฏิคามวคฺค}{mtk.274}
\bookmark[dest={mtk.274},level=5]{3. โกฏิคามวคฺค}
\csromantocmarkref{subparagraph}{1-2. โกฏิคามสุตฺตานิ}{mtk.275}
\bookmark[dest={mtk.275},level=5]{1-2. โกฏิคามสุตฺตานิ}
\csromanheader{6}{3. \textbf{โกฏิคามวคฺค}}
\csromanheader{2}{1. ปฐม\-โกฏิคาม\-สุตฺต}
\proseitem{1091}{เอกํ สมยํ ภควา วชฺชีสุ วิหรติ โกฏิคาเม. ตตฺร โข ภควา ภิกฺขู อามนฺเตสิ “จตุนฺนํ ภิกฺขเว อริยสจฺจานํ อนนุโพธา อปฺปฏิเวธา เอวมิทํ ทีฆมทฺธานํ สนฺธาวิตํ สํสริตํ มมญฺเจว ตุมฺหากญฺจ.}

\prose{\csromanfolio{378}กตเมสํ จตุนฺนํ, ทุกฺขสฺส ภิกฺขเว อริยสจฺจสฺส อนนุโพธา อปฺปฏิเวธา เอวมิทํ ทีฆมทฺธานํ สนฺธาวิตํ สํสริตํ มมญฺเจว ตุมฺหากญฺจ. ทุกฺข\-สมุทยสฺส อริยสจฺจสฺส ฯเปฯ. ทุกฺข\-นิโรธสฺส อริยสจฺจสฺส ฯเปฯ. ทุกฺขนิโรธ\-คามินิยา ปฏิปทาย อริยสจฺจสฺส อนนุโพธา อปฺปฏิเวธา เอวมิทํ ทีฆมทฺธานํ สนฺธาวิตํ สํสริตํ มมญฺเจว ตุมฺหากญฺจ. ตยิทํ ภิกฺขเว ทุกฺขํ อริยสจฺจํ อนุพุทฺธํ ปฏิวิทฺธํ, ทุกฺข\-สมุทยํ อริยสจฺจํ อนุพุทฺธํ ปฏิวิทฺธํ, ทุกฺขนิโรธํ อริยสจฺจํ อนุพุทฺธํ ปฏิวิทฺธํ, ทุกฺขนิโรธ\-คามินี ปฏิปทา อริยสจฺจํ อนุพุทฺธํ ปฏิวิทฺธํ. อุจฺฉินฺนา ภวตณฺหา, ขีณา ภวเนตฺติ, นตฺถิทานิ ปุนพฺภโว”ติ.}

\prose{อิทมโวจ ภควา, อิทํ วตฺวาน สุคโต อถาปรํ เอตทโวจ สตฺถา\csromandash{}}

\setlength{\csromangathapagewidth}{0pt}
\setlength{\csromangathawakwidth}{0pt}
\csromangathameasuregroup{“จตุนฺนํ อริยสจฺจานํ, \\ สํสิตํ\textsuperscript{1} ทีฆมทฺธานํ, \\ ตานิ\textsuperscript{1} เอตานิ ทิฏฺฐานิ,}{\csromangathabat{“จตุนฺนํ อริยสจฺจานํ,}{ยถาภูตํ อทสฺสนา.} \\ \csromangathabat{สํสิตํ\textsuperscript{1} ทีฆมทฺธานํ,}{ตาสุ ตาเสฺวว ชาติสุ.} \\ \csromangathabat{ตานิ\textsuperscript{1} เอตานิ ทิฏฺฐานิ,}{ภวเนตฺติ สมูหตา.}}
\csromangathasetleft{“จตุนฺนํ อริยสจฺจานํ, \\ สํสิตํ\textsuperscript{1} ทีฆมทฺธานํ, \\ ตานิ\textsuperscript{1} เอตานิ ทิฏฺฐานิ,}
\csromangathagroup{\csromangathastanza{\csromangathabat{“จตุนฺนํ อริยสจฺจานํ,}{ยถาภูตํ อทสฺสนา.} \\ \csromangathabat{สํสิตํ\csromansharedfootnote{70fc911cceccd738}{สํสริตํ (สฺยา, กํ, ก) ที 2. 77 ปิฏฺเฐปิ.} ทีฆมทฺธานํ,}{ตาสุ ตาเสฺวว ชาติสุ.} \\ \csromangathabat{ตานิ\csromansharedfootnote{70fc911cceccd738}{สํสริตํ (สฺยา, กํ, ก) ที 2. 77 ปิฏฺเฐปิ.} เอตานิ ทิฏฺฐานิ,}{ภวเนตฺติ สมูหตา.}}}

\prose{อุจฺฉินฺนํ มูลํ ทุกฺขสฺส, นตฺถิทานิ ปุนพฺภโว”ติ. \textbf{ปฐมํ.}}

\csromanheader{2}{2. ทุติย\-โกฏิคาม\-สุตฺต}
\proseitem{1092}{เย หิ เกจิ ภิกฺขเว สมณา วา พฺราหฺมณา วา “อิทํ ทุกฺขนฺ”ติ ยถาภูตํ นปฺปชานนฺติ, “อยํ ทุกฺขสมุทโย”ติ ยถาภูตํ นปฺปชานนฺติ, “อยํ ทุกฺขนิโรโธ”ติ ยถาภูตํ นปฺปชานนฺติ, “อยํ ทุกฺขนิโรธ\-คามินี ปฏิปทา”ติ ยถาภูตํ นปฺปชานนฺติ, น เม เต ภิกฺขเว สมณา วา พฺราหฺมณา วา สมเณสุ วา สมณสมฺมตา, พฺราหฺมเณสุ วา พฺราหฺมณ\-สมฺมตา, น จ ปน เต อายสฺมนฺโต สามญฺญตฺถํ วา พฺรหฺมญฺญตฺถํ วา ทิฏฺเฐว ธมฺเม สยํ อภิญฺญา สจฺฉิกตฺวา อุปสมฺปชฺช วิหรนฺติ.}

\prose{เย จ โข เกจิ ภิกฺขเว สมณา วา พฺราหฺมณา วา “อิทํ ทุกฺขนฺ”ติ ยถาภูตํ ปชานนฺติ, “อยํ ทุกฺขสมุทโย”ติ ยถาภูตํ ปชานนฺติ, “อยํ ทุกฺขนิโรโธ”ติ ยถาภูตํ ปชานนฺติ, “อยํ ทุกฺขนิโรธ\-คามินี ปฏิปทา”ติ ยถาภูตํ ปชานนฺติ. เต โข เม ภิกฺขเว สมณา วา}

\prose{\csromanfolio{379}พฺราหฺมณา วา สมเณสุ เจว สมณสมฺมตา, พฺราหฺมเณสุ จ พฺราหฺมณ\-สมฺมตา, เต จ ปนายสฺมนฺโต สามญฺญตฺถญฺจ พฺรหฺมญฺญตฺถญฺจ ทิฏฺเฐว ธมฺเม สยํ อภิญฺญา สจฺฉิกตฺวา อุปสมฺปชฺช วิหรนฺตีติ.}

\prose{อิทมโวจ ภควา, อิทํ วตฺวาน สุคโต อถาปรํ เอตทโวจ สตฺถา\csromandash{}}

\setlength{\csromangathapagewidth}{0pt}
\setlength{\csromangathawakwidth}{0pt}
\csromangathameasuregroup{“เย ทุกฺขํ นปฺปชานนฺติ, \\ ยตฺถ จ สพฺพโส ทุกฺขํ, \\ ตญฺจ มคฺคํ น ชานนฺติ, \\ เจโตวิมุตฺติหีนา เต, \\ อภพฺพา เต อนฺตกิริยาย, \\ เย จ ทุกฺขํ ปชานนฺติ, \\ ยตฺถ จ สพฺพโส ทุกฺขํ, \\ ตญฺจ มคฺคํ ปชานนฺติ, \\ เจโตวิมุตฺติสมฺปนฺนา,}{\csromangathabat{“เย ทุกฺขํ นปฺปชานนฺติ,}{อโถ ทุกฺขสฺส สมฺภวํ.} \\ \csromangathabat{ยตฺถ จ สพฺพโส ทุกฺขํ,}{อเสสํ อุปรุชฺฌติ.} \\ \csromangathabat{ตญฺจ มคฺคํ น ชานนฺติ,}{ทุกฺขูปสมคามินํ.} \\ \csromangathabat{เจโตวิมุตฺติหีนา เต,}{อโถ ปญฺญาวิมุตฺติยา.} \\ \csromangathabat{อภพฺพา เต อนฺตกิริยาย,}{เต เว ชาติชรูปคา.} \\ \csromangathabat{เย จ ทุกฺขํ ปชานนฺติ,}{อโถ ทุกฺขสฺส สมฺภวํ.} \\ \csromangathabat{ยตฺถ จ สพฺพโส ทุกฺขํ,}{อเสสํ อุปรุชฺฌติ.} \\ \csromangathabat{ตญฺจ มคฺคํ ปชานนฺติ,}{ทุกฺขูปสมคามินํ.} \\ \csromangathabat{เจโตวิมุตฺติสมฺปนฺนา,}{อโถ ปญฺญาวิมุตฺติยา.}}
\csromangathasetleft{“เย ทุกฺขํ นปฺปชานนฺติ, \\ ยตฺถ จ สพฺพโส ทุกฺขํ, \\ ตญฺจ มคฺคํ น ชานนฺติ, \\ เจโตวิมุตฺติหีนา เต, \\ อภพฺพา เต อนฺตกิริยาย, \\ เย จ ทุกฺขํ ปชานนฺติ, \\ ยตฺถ จ สพฺพโส ทุกฺขํ, \\ ตญฺจ มคฺคํ ปชานนฺติ, \\ เจโตวิมุตฺติสมฺปนฺนา,}
\csromangathagroup{\csromangathastanza{\csromangathabat{“เย ทุกฺขํ นปฺปชานนฺติ,}{อโถ ทุกฺขสฺส สมฺภวํ.} \\ \csromangathabat{ยตฺถ จ สพฺพโส ทุกฺขํ,}{อเสสํ อุปรุชฺฌติ.}}\csromangathastanza{\csromangathabat{ตญฺจ มคฺคํ น ชานนฺติ,}{ทุกฺขู\-ปสม\-คามินํ.} \\ \csromangathabat{เจโตวิมุตฺติ\-หีนา เต,}{อโถ ปญฺญาวิมุตฺติยา.}}\csromangathastanza{\csromangathabat{อภพฺพา เต อนฺตกิริยาย,}{เต เว ชาติชรูปคา.} \\ \csromangathabat{เย จ ทุกฺขํ ปชานนฺติ,}{อโถ ทุกฺขสฺส สมฺภวํ.}}\csromangathastanza{\csromangathabat{ยตฺถ จ สพฺพโส ทุกฺขํ,}{อเสสํ อุปรุชฺฌติ.} \\ \csromangathabat{ตญฺจ มคฺคํ ปชานนฺติ,}{ทุกฺขู\-ปสม\-คามินํ.} \\ \csromangathabat{เจโตวิมุตฺติ\-สมฺปนฺนา,}{อโถ ปญฺญาวิมุตฺติยา.}}}

\vspace{\tipitakaparskip}
\csromancenter{ภพฺพา เต อนฺตกิริยาย, น เต ชาติชรูปคา”ติ. ทุติยํ.}
\csromanmatikaanchor{mtk.276}
\csromantocmarkref{subparagraph}{3. สมฺมาสมฺพุทฺธสุตฺต}{mtk.276}
\bookmark[dest={mtk.276},level=5]{3. สมฺมาสมฺพุทฺธสุตฺต}
\csromanheader{6}{3. สมฺมาสมฺพุทฺธ\-สุตฺต}
\proseitem{1093}{สาวตฺถิ\-นิทานํ. จตฺตาริมานิ ภิกฺขเว อริยสจฺจานิ. กตมานิ จตฺตาริ, ทุกฺขํ อริยสจฺจํ ฯเปฯ. ทุกฺขนิโรธ\-คามินี ปฏิปทา อริยสจฺจํ. อิมานิ โข ภิกฺขเว จตฺตาริ อริยสจฺจานิ. อิเมสํ โข ภิกฺขเว จตุนฺนํ อริยสจฺจานํ ยถาภูตํ อภิสมฺพุทฺธตฺตา ตถาคโต อรหํ “สมฺมาสมฺพุทฺโธ”ติ วุจฺจติ.}

\prose{ตสฺมาติห ภิกฺขเว “อิทํ ทุกฺขนฺ”ติ โยโค กรณีโย ฯเปฯ “อยํ ทุกฺขนิโรธ\-คามินี ปฏิปทา”ติ โยโค กรณีโยติ.  ตติยํ.}

\csromanmatikaanchor{mtk.277}
\csromantocmarkref{subparagraph}{4. อรหนฺตสุตฺต}{mtk.277}
\bookmark[dest={mtk.277},level=5]{4. อรหนฺตสุตฺต}
\csromanheader{6}{4. \textbf{อรหนฺตสุตฺต}}
\csromanmatikaanchor{mtk.278}
\csromanmatikaanchor{mtk.279}
\csromantocmarkref{subparagraph}{5. อาสวกฺขยสุตฺต}{mtk.278}
\bookmark[dest={mtk.278},level=5]{5. อาสวกฺขยสุตฺต}
\csromantocmarkref{subparagraph}{6-8. มิตฺตสุตฺตาทีนิ}{mtk.279}
\bookmark[dest={mtk.279},level=5]{6-8. มิตฺตสุตฺตาทีนิ}
\proseitem{1094}{สาวตฺถิ\-นิทานํ. เย หิ เกจิ ภิกฺขเว อตีตมทฺธานํ อรหนฺโต สมฺมาสมฺพุทฺธา ยถาภูตํ อภิสมฺ\-พุชฺฌิํสุ, สพฺเพ เต จตฺตาริ อริยสจฺจานิ \csromanfolio{380}ยถาภูตํ อภิสมฺ\-พุชฺฌิํสุ. เย หิ\csromansharedfootnote{e9bb429f2c987552}{เยปิ หิ (พหูสุ)} เกจิ ภิกฺขเว อนาคตมทฺธานํ อรหนฺโต สมฺมาสมฺพุทฺธา ยถาภูตํ อภิสมฺพุชฺฌิสฺสนฺติ, สพฺเพ เต จตฺตาริ อริยสจฺจานิ ยถาภูตํ อภิสมฺพุชฺฌิสฺสนฺติ. เย หิ เกจิ ภิกฺขเว เอตรหิ อรหนฺโต สมฺมาสมฺพุทฺธา ยถาภูตํ อภิสมฺ\-พุชฺฌนฺติ, สพฺเพ เต จตฺตาริ อริยสจฺจานิ ยถาภูตํ อภิสมฺ\-พุชฺฌนฺติ.}

\prose{กตมานิ จตฺตาริ, ทุกฺขํ อริยสจฺจํ ทุกฺข\-สมุทยํ อริยสจฺจํ ทุกฺขนิโรธํ อริยสจฺจํ ทุกฺขนิโรธ\-คามินี ปฏิปทา อริยสจฺจํ. เย หิ เกจิ ภิกฺขเว อตีตมทฺธานํ อรหนฺโต สมฺมาสมฺพุทฺธา ยถาภูตํ อภิสมฺ\-พุชฺฌิํสุ ฯเปฯ อภิสมฺพุชฺฌิสฺสนฺติ. อภิสมฺ\-พุชฺฌนฺติ, สพฺเพ เต อิมานิ จตฺตาริ อริยสจฺจานิ ยถาภูตํ อภิสมฺ\-พุชฺฌนฺติ.}

\prose{ตสฺมาติห ภิกฺขเว “อิทํ ทุกฺขนฺ”ติ โยโค กรณีโย ฯเปฯ “อยํ ทุกฺขนิโรธ\-คามินี ปฏิปทา”ติ โยโค กรณีโยติ.  จตุตฺถํ.}

\csromanheader{6}{5. อาสวกฺขย\-สุตฺต}
\proseitem{1095}{ชานโตหํ ภิกฺขเว ปสฺสโต อาสวานํ ขยํ วทามิ, โน อชานโต อปสฺสโต. กิญฺจ ภิกฺขเว ชานโต ปสฺสโต อาสวานํ ขโย โหติ, “อิทํ ทุกฺขนฺ”ติ ภิกฺขเว ชานโต ปสฺสโต อาสวานํ ขโย โหติ, “อยํ ทุกฺขสมุทโย”ติ ชานโต ปสฺสโต อาสวานํ ขโย โหติ, “อยํ ทุกฺขนิโรโธ”ติ ชานโต ปสฺสโต อาสวานํ ขโย โหติ, “อยํ ทุกฺขนิโรธ\-คามินี ปฏิปทา”ติ ชานโต ปสฺสโต อาสวานํ ขโย โหติ. เอวํ โข ภิกฺขเว ชานโต เอวํ ปสฺสโต อาสวานํ ขโย โหติ.}

\prose{ตสฺมาติห ภิกฺขเว “อิทํ ทุกฺขนฺ”ติ โยโค กรณีโย ฯเปฯ “อยํ ทุกฺขนิโรธ\-คามินี ปฏิปทา”ติ โยโค กรณีโยติ.  ปญฺจมํ.}

\csromanheader{2}{6. \textbf{มิตฺตสุตฺต}}
\proseitem{1096}{เย หิ เกจิ ภิกฺขเว อนุกมฺเปยฺยาถ, เย จ โสตพฺพํ มญฺเญยฺยุํ มิตฺตา วา อมจฺจา วา ญาตี วา สาโลหิตา วา, เต โว ภิกฺขเว จตุนฺนํ อริยสจฺจานํ ยถาภูตํ อภิสมยาย สมาทเปตพฺพา นิเวเสตพฺพา ปติฏฺฐาเป\-ตพฺพา.}

\prose{\csromanfolio{381}กตเมสํ จตุนฺนํ, ทุกฺขสฺส อริยสจฺจสฺส ทุกฺข\-สมุทยสฺส อริยสจฺจสฺส ทุกฺข\-นิโรธสฺส อริยสจฺจสฺส ทุกฺขนิโรธ\-คามินิยา ปฏิปทาย อริยสจฺจสฺส. เย หิ เกจิ ภิกฺขเว อนุกมฺเปยฺยาถ, เย จ โสตพฺพํ มญฺเญยฺยุํ มิตฺตา วา อมจฺจา วา ญาตี วา สาโลหิตา วา, เต โว ภิกฺขเว อิเมสํ จตุนฺนํ อริยสจฺจานํ ยถาภูตํ อภิสมยาย สมาทเปตพฺพา นิเวเสตพฺพา ปติฏฺฐาเป\-ตพฺพา.}

\prose{ตสฺมาติห ภิกฺขเว “อิทํ ทุกฺขนฺ”ติ โยโค กรณีโย ฯเปฯ “อยํ ทุกฺขนิโรธ\-คามินี ปฏิปาทา”ติ โยโค กรณีโยติ.  ฉฏฺฐํ.}

\csromanheader{2}{7. \textbf{ตถสุตฺต}}
\proseitem{1097}{จตฺตาริมานิ ภิกฺขเว อริยสจฺจานิ. กตมานิ จตฺตาริ, ทุกฺขํ อริยสจฺจํ ทุกฺข\-สมุทยํ อริยสจฺจํ ทุกฺขนิโรธํ อริยสจฺจํ ทุกฺขนิโรธ\-คามินี ปฏิปทา อริยสจฺจํ. อิมานิ โข ภิกฺขเว จตฺตาริ อริยสจฺจานิ ตถานิ อวิตถานิ อนญฺญถานิ, ตสฺมา “อริยสจฺจานี”ติ วุจฺจนฺติ.}

\prose{ตสฺมาติห ภิกฺขเว “อิทํ ทุกฺขนฺ”ติ โยโค กรณีโย ฯเปฯ “อยํ ทุกฺขนิโรธ\-คามินี ปฏิปทา”ติ โยโค กรณีโยติ.  สตฺตมํ.}

\csromanheader{2}{8. \textbf{โลกสุตฺต}}
\proseitem{1098}{จตฺตาริมานิ ภิกฺขเว อริยสจฺจานิ. กตมานิ จตฺตาริ, ทุกฺขํ อริยสจฺจํ ทุกฺข\-สมุทยํ อริยสจฺจํ ทุกฺขนิโรธํ อริยสจฺจํ ทุกฺขนิโรธ\-คามินี ปฏิปทา อริยสจฺจํ. สเทวเก โลเก สมารเก สพฺรหฺมเก สสฺสมณ\-พฺราหฺมณิยา ปชาย สเทว\-มนุสฺสาย ตถาคโต อริโย, ตสฺมา “อริยสจฺจานี”ติ วุจฺจนฺติ.}

\prose{ตสฺมาติห ภิกฺขเว “อิทํ ทุกฺขนฺ”ติ โยโค กรณีโย ฯเปฯ “อยํ ทุกฺขนิโรธ\-คามินี ปฏิปทา”ติ โยโค กรณีโยติ.  อฏฺฐมํ.}

\csromanmatikaanchor{mtk.280}
\csromantocmarkref{subparagraph}{9. ปริญฺเญยฺยสุตฺต}{mtk.280}
\bookmark[dest={mtk.280},level=5]{9. ปริญฺเญยฺยสุตฺต}
\csromanheader{6}{9. ปริญฺเญยฺย\-สุตฺต}
\proseitem{1099}{จตฺตาริมานิ ภิกฺขเว อริยสจฺจานิ. กตมานิ จตฺตาริ, ทุกฺขํ อริยสจฺจํ ทุกฺข\-สมุทยํ อริยสจฺจํ ทุกฺขนิโรธํ อริยสจฺจํ ทุกฺขนิโรธ\-คามินี ปฏิปาทา อริยสจฺจํ. อิมานิ โข ภิกฺขเว จตฺตาริ อริยสจฺจานิ. \csromanfolio{382}อิเมสํ โข ภิกฺขเว จตุนฺนํ อริยสจฺจานํ อตฺถิ อริยสจฺจํ ปริญฺเญยฺยํ, อตฺถิ อริยสจฺจํ ปหาตพฺพํ, อตฺถิ อริยสจฺจํ สจฺฉิกาตพฺพํ, อตฺถิ อริยสจฺจํ ภาเวตพฺพํ.}

\prose{กตมญฺจ ภิกฺขเว อริยสจฺจํ ปริญฺเญยฺยํ, ทุกฺขํ ภิกฺขเว อริยสจฺจํ ปริญฺเญยฺยํ, ทุกฺข\-สมุทยํ อริยสจฺจํ ปหาตพฺพํ, ทุกฺขนิโรธํ อริยสจฺจํ สจฺฉิกาตพฺพํ, ทุกฺขนิโรธ\-คามินี ปฏิปทา อริยสจฺจํ ภาเวตพฺพํ.}

\prose{ตสฺมาติห ภิกฺขเว “อิทํ ทุกฺขนฺ”ติ โยโค กรณีโย ฯเปฯ “อยํ ทุกฺขนิโรธ\-คามินี ปฏิปทา”ติ โยโค กรณีโยติ.  นวมํ.}

\csromanmatikaanchor{mtk.281}
\csromantocmarkref{subparagraph}{10. ควมฺปติสุตฺต}{mtk.281}
\bookmark[dest={mtk.281},level=5]{10. ควมฺปติสุตฺต}
\csromanheader{6}{10. \textbf{ควมฺปติสุตฺต}}
\proseitem{1100}{เอกํ สมยํ สมฺพหุลา เถรา ภิกฺขู เจเตสุ\csromansharedfootnote{0db9372569e563f5}{เจติเยสุ (สฺยา)} วิหรนฺติ สหญฺจนิเก\csromansharedfootnote{817c4bf30f834c56}{สหชนิเย (สี, สฺยา, กํ)}. เตน โข ปน สมเยน สมฺพหุลานํ เถรานํ ภิกฺขูนํ ปจฺฉาภตฺตํ ปิณฺฑปาต\-ปฏิกฺกนฺตานํ มณฺฑลมาเฬ สนฺนิสินฺนานํ สนฺนิปติตานํ อยม\-นฺตรากถา อุทปาทิ “โย นุ โข อาวุโส ทุกฺขํ ปสฺสติ, ทุกฺข\-สมุทยมฺปิ โส ปสฺสติ, ทุกฺข\-นิโรธมฺปิ ปสฺสติ, ทุกฺขนิโรธ\-คามินิํ ปฏิปทมฺปิ ปสฺสตี”ติ.}

\prosewithclosermidruled{เอวํ วุตฺเต อายสฺมา ควมฺปติ เถโร\csromansharedfootnote{c11c52f05fae1562}{ควมฺปติตฺเถโร (สฺยา, กํ)} ภิกฺขู เอตทโวจ “สมฺมุขา เมตํ อาวุโส ภควโต สุตํ, สมฺมุขา ปฏิคฺคหิตํ โย ภิกฺขเว ทุกฺขํ ปสฺสติ, ทุกฺข\-สมุทยมฺปิ โส ปสฺสติ, ทุกฺข\-นิโรธมฺปิ ปสฺสติ, ทุกฺขนิโรธ\-คามินิํ ปฏิปทมฺปิ ปสฺสติ. โย ทุกฺข\-สมุทยํ ปสฺสติ, ทุกฺขมฺปิ โส ปสฺสติ, ทุกฺข\-นิโรธมฺปิ ปสฺสติ, ทุกฺขนิโรธ\-คามินิํ ปฏิปทมฺปิ ปสฺสติ. โย ทุกฺขนิโรธํ ปสฺสติ, ทุกฺขมฺปิ โส ปสฺสติ, ทุกฺข\-สมุทยมฺปิ ปสฺสติ, ทุกฺขนิโรธ\-คามินิํ ปฏิปทมฺปิ ปสฺสติ. โย ทุกฺขนิโรธ\-คามินิํ ปฏิปทํ ปสฺสติ, ทุกฺขมฺปิ โส ปสฺสติ, ทุกฺข\-สมุทยมฺปิ ปสฺสติ, ทุกฺข\-นิโรธมฺปิ ปสฺสตี”ติ.  ทสมํ.}{โกฏิคามวคฺโค ตติโย.}
\csromancenter{\textbf{ตสฺสุทฺทานํ}}
\setlength{\csromangathapagewidth}{0pt}
\setlength{\csromangathawakwidth}{0pt}
\csromangathameasuregroup{เทฺว วชฺชี สมฺมาสมฺพุทฺโธ, \\ มิตฺตํ ตถา จ โลโก จ,}{\csromangathabat{เทฺว วชฺชี สมฺมาสมฺพุทฺโธ,}{อรหํ อาสวกฺขโย.} \\ \csromangathabat{มิตฺตํ ตถา จ โลโก จ,}{ปริญฺเญยฺยํ ควมฺปตีติ.}}
\csromangathasetleft{เทฺว วชฺชี สมฺมาสมฺพุทฺโธ, \\ มิตฺตํ ตถา จ โลโก จ,}
\csromangathagroup{\csromangathastanza{\csromangathabat{เทฺว วชฺชี สมฺมาสมฺพุทฺโธ,}{อรหํ อาสวกฺขโย.} \\ \csromangathabat{มิตฺตํ ตถา จ โลโก จ,}{ปริญฺเญยฺยํ ควมฺปตีติ.}}}

\csromanmatikaanchor{mtk.282}
\csromantocmarknopage{subparagraph}{4. สีสปาวนวคฺค}{mtk.282}
\bookmark[dest={mtk.282},level=5]{4. สีสปาวนวคฺค}
\csromanheader{6}{4. สีสปาวน\-วคฺค}
\csromanmatikaanchor{mtk.283}
\csromantocmarkref{subparagraph}{1. สีสปาวนสุตฺต}{mtk.283}
\bookmark[dest={mtk.283},level=5]{1. สีสปาวนสุตฺต}
\csromanheader{6}{1. สีสปาวน\-สุตฺต}
\proseitem{1101}{\csromanfolio{383}เอกํ สมยํ ภควา โกสมฺพิยํ วิหรติ สีสปาวเน\csromansharedfootnote{50fbae39639ea517}{สิํสปาวเน (สี, อิ)}. อถ โข ภควา ปริตฺตานิ สีสปาปณฺณานิ ปาณินา คเหตฺวา ภิกฺขู อามนฺเตสิ “ตํ กิํ มญฺญถ ภิกฺขเว, กตมํ นุ โข พหุตรํ ยานิ วา มยา ปริตฺตานิ สีสปาปณฺณานิ ปาณินา คหิตานิ ยทิทํ อุปริ สีสปาวเน”ติ. อปฺปมตฺตกานิ ภนฺเต ภควตา ปริตฺตานิ สีสปาปณฺณานิ ปาณินา คหิตานิ, อถ โข เอตาเนว พหุตรานิ ยทิทํ อุปริ สีสปาวเนติ. เอวเมว โข ภิกฺขเว เอตเทว พหุตรํ ยํ โว มยา อภิญฺญาย อนกฺขาตํ. กสฺมา เจตํ ภิกฺขเว มยา อนกฺขาตํ, น เหตํ ภิกฺขเว อตฺถสํหิตํ นาทิ\-พฺรหฺม\-จริยกํ น นิพฺพิทาย น วิราคาย น นิโรธาย น อุปสมาย น อภิญฺญาย น สมฺโพธาย น นิพฺพานาย สํวตฺตติ, ตสฺมา ตํ มยา อนกฺขาตํ.}

\prose{กิญฺจ ภิกฺขเว มยา อกฺขาตํ, “อิทํ ทุกฺขนฺ”ติ ภิกฺขเว มยา อกฺขาตํ, “อยํ ทุกฺขสมุทโย”ติ มยา อกฺขาตํ, “อยํ ทุกฺขนิโรโธ”ติ มยา อกฺขาตํ, “อยํ ทุกฺขนิโรธ\-คามินี ปฏิปทา”ติ มยา อกฺขาตํ.}

\prose{กสฺมา เจตํ ภิกฺขเว มยา อกฺขาตํ, เอตํ หิ ภิกฺขเว อตฺถสํหิตํ เอตํ อาทิ\-พฺรหฺม\-จริยกํ เอตํ นิพฺพิทาย วิราคาย นิโรธาย อุปสมาย อภิญฺญาย สมฺโพธาย นิพฺพานาย สํวตฺตติ, ตสฺมา ตํ มยา อกฺขาตํ.}

\prose{ตสฺมาติห ภิกฺขเว “อิทํ ทุกฺขนฺ”ติ โยโค กรณีโย ฯเปฯ “อยํ ทุกฺขนิโรธ\-คามินี ปฏิปทา”ติ โยโค กรณีโยติ.  ปฐมํ.}

\csromanmatikaanchor{mtk.284}
\csromantocmarkref{subparagraph}{2. ขทิรปตฺตสุตฺต}{mtk.284}
\bookmark[dest={mtk.284},level=5]{2. ขทิรปตฺตสุตฺต}
\csromanheader{6}{2. ขทิรปตฺต\-สุตฺต}
\proseitem{1102}{โย ภิกฺขเว เอวํ วเทยฺย “อหํ ทุกฺขํ อริยสจฺจํ ยถาภูตํ อนภิสเมจฺจ ทุกฺข\-สมุทยํ อริยสจฺจํ ยถาภูตํ อนภิสเมจฺจ ทุกฺขนิโรธํ อริยสจฺจํ ยถาภูตํ อนภิสเมจฺจ ทุกฺขนิโรธ\-คามินิํ ปฏิปทํ อริยสจฺจํ ยถาภูตํ อนภิสเมจฺจ สมฺมา ทุกฺขสฺสนฺตํ กริสฺสามี”ติ เนตํ ฐานํ วิชฺชติ.}

\prose{\csromanfolio{384}เสยฺยถาปิ ภิกฺขเว โย เอวํ วเทยฺย “อหํ ขทิร\-ปตฺตานํ วา สรลปตฺตานํ\csromansharedfootnote{f663c8c3fdcd7bbd}{ปลาสปตฺตานํ (สี, สฺยา, กํ, อิ)} วา อามลก\-ปตฺตานํ วา ปุฏํ กริตฺวา อุทกํ วา ตาลปตฺตํ วา อาหริสฺสามี”ติ เนตํ ฐานํ วิชฺชติ. เอวเมว โข ภิกฺขเว โย เอวํ วเทยฺย “อหํ ทุกฺขํ อริยสจฺจํ ยถาภูตํ อนภิสเมจฺจ ฯเปฯ ทุกฺขนิโรธ\-คามินิํ ปฏิปทํ อริยสจฺจํ ยถาภูตํ อนภิสเมจฺจ สมฺมา ทุกฺขสฺสนฺตํ กริสฺสามี”ติ เนตํ ฐานํ วิชฺชติ.}

\prose{โย จ โข ภิกฺขเว เอวํ วเทยฺย “อหํ ทุกฺขํ อริยสจฺจํ ยถาภูตํ อภิสเมจฺจ ทุกฺข\-สมุทยํ อริยสจฺจํ ยถาภูตํ อภิสเมจฺจ ทุกฺขนิโรธํ อริยสจฺจํ ยถาภูตํ อภิสเมจฺจ ทุกฺขนิโรธ\-คามินิํ ปฏิปทํ อริยสจฺจํ ยถาภูตํ อภิสเมจฺจ สมฺมา ทุกฺขสฺสนฺตํ กริสฺสามี”ติ ฐานเมตํ วิชฺชติ.}

\prose{เสยฺยถาปิ ภิกฺขเว โย เอวํ วเทยฺย “อหํ ปทุมปตฺตานํ วา ปลาสปตฺตานํ วา มาลุวปตฺตานํ วา ปุฏํ กริตฺวา อุทกํ วา ตาลปตฺตํ วา อาหริสฺสามี”ติ ฐานเมตํ วิชฺชติ. เอวเมว โข ภิกฺขเว โย เอวํ วเทยฺย “อหํ ทุกฺขํ อริยสจฺจํ ยถาภูตํ อภิสเมจฺจ ฯเปฯ ทุกฺขนิโรธ\-คามินิํ ปฏิปทํ อริยสจฺจํ ยถาภูตํ อภิสเมจฺจ สมฺมา ทุกฺขสฺสนฺตํ กริสฺสามี”ติ ฐานเมตํ วิชฺชติ.}

\prose{ตสฺมาติห ภิกฺขเว “อิทํ ทุกฺขนฺ”ติ โยโค กรณีโย ฯเปฯ “อยํ ทุกฺขนิโรธ\-คามินี ปฏิปทา”ติ โยโค กรณีโยติ.  ทุติยํ.}

\csromanmatikaanchor{mtk.285}
\csromantocmarkref{subparagraph}{3. ทณฺฑสุตฺต}{mtk.285}
\bookmark[dest={mtk.285},level=5]{3. ทณฺฑสุตฺต}
\csromanheader{6}{3. \textbf{ทณฺฑสุตฺต}}
\proseitem{1103}{เสยฺยถาปิ ภิกฺขเว ทณฺโฑ อุปริเวหาสํ ขิตฺโต สกิมฺปิ มูเลน นิปตติ, สกิมฺปิ อคฺเคน นิปตติ. เอวเมว โข ภิกฺขเว อวิชฺชานีวรณา สตฺตา ตณฺหาสํโยชนา สนฺธาวนฺตา สํสรนฺตา\csromansharedfootnote{ad4806134bc449b1}{ตณฺหาสํโยชนพนฺธา สนฺธาวตา (ก)} สกิมฺปิ อสฺมา โลกา ปรํ โลกํ คจฺฉนฺติ, สกิมฺปิ ปรสฺมา โลกา อิมํ โลกํ อาคจฺฉนฺติ. ตํ กิสฺส เหตุ, อทิฏฺฐตฺตา ภิกฺขเว จตุนฺนํ อริยสจฺจานํ. กตเมสํ จตุนฺนํ, ทุกฺขสฺส อริยสจฺจสฺส ฯเปฯ ทุกฺขนิโรธ\-คามินิยา ปฏิปทาย อริยสจฺจสฺส.}

\prose{ตสฺมาติห ภิกฺขเว “อิทํ ทุกฺขนฺ”ติ โยโค กรณีโย ฯเปฯ “อยํ ทุกฺขนิโรธ\-คามินี ปฏิปทา”ติ โยโค กรณีโยติ.  ตติยํ.}

\csromanheader{6}{12. สจฺจสํยุตฺต}
\csromanmatikaanchor{mtk.286}
\csromantocmarkref{subparagraph}{4. เจลสุตฺต}{mtk.286}
\bookmark[dest={mtk.286},level=5]{4. เจลสุตฺต}
\csromanheader{6}{4. \textbf{เจลสุตฺต}}
\proseitem{1104}{\csromanfolio{385}อาทิตฺเต ภิกฺขเว เจเล วา สีเส วา กิมสฺส กรณียนฺติ. อาทิตฺเต ภนฺเต เจเล วา สีเส วา ตสฺเสว เจลสฺส วา สีสสฺส วา นิพฺพาปนาย อธิมตฺโต ฉนฺโท จ วายาโม จ อุสฺสาโห จ อุสฺโสฬฺหี จ อปฺปฏิวานี จ สติ จ สมฺปชญฺญญฺจ กรณียนฺติ.}

\prose{อาทิตฺตํ ภิกฺขเว เจลํ วา สีสํ วา อชฺฌุเปกฺขิตฺวา อมนสิกริตฺวา อนภิสเมตานํ จตุนฺนํ อริยสจฺจานํ ยถาภูตํ อภิสมยาย อธิมตฺโต ฉนฺโท จ วายาโม จ อุสฺสาโห จ อุสฺโสฬฺหี จ อปฺปฏิวานี จ สติ จ สมฺปชญฺญญฺจ กรณียํ. กตเมสํ จตุนฺนํ, ทุกฺขสฺส อริยสจฺจสฺส ฯเปฯ ทุกฺขนิโรธ\-คามินิยา ปฏิปทาย อริยสจฺจสฺส.}

\prose{ตสฺมาติห ภิกฺขเว “อิทํ ทุกฺขนฺ”ติ โยโค กรณีโย ฯเปฯ “อยํ ทุกฺขนิโรธ\-คามินี ปฏิปทา”ติ โยโค กรณีโยติ.  จตุตฺถํ.}

\csromanmatikaanchor{mtk.287}
\csromantocmarkref{subparagraph}{5. สตฺติสตสุตฺต}{mtk.287}
\bookmark[dest={mtk.287},level=5]{5. สตฺติสตสุตฺต}
\csromanheader{6}{5. \textbf{สตฺติสตสุตฺต}}
\proseitem{1105}{เสยฺยถาปิ ภิกฺขเว ปุริโส วสฺสสตายุโก วสฺสสตชีวี, ตเมนํ เอวํ วเทยฺย “เอหมฺโภ ปุริส, ปุพฺพณฺหสมยํ ตํ สตฺติสเตน หนิสฺสนฺติ, มชฺฌนฺหิกสมยํ สตฺติสเตน หนิสฺสนฺติ, สายนฺหสมยํ สตฺติสเตน หนิสฺสนฺติ. โส โข ตฺวํ อมฺโภ ปุริส ทิวเส ทิวเส ตีหิ ตีหิ สตฺติสเตหิ หญฺญมาโน วสฺสสตายุโก วสฺสสตชีวี วสฺสสตสฺส อจฺจเยน อนภิสเมตานิ จตฺตาริ อริยสจฺจานิ อภิสเมสฺสสี”ติ.}

\prose{อตฺถวสิเกน ภิกฺขเว กุลปุตฺเตน อลํ อุปคนฺตุํ. ตํ กิสฺส เหตุ, อนมตคฺโคยํ ภิกฺขเว สํสาโร, ปุพฺพา โกฏิ นปฺปญฺญายติ สตฺติ\-ปฺปหารานํ อสิปฺปหารานํ อุสุปฺปหารานํ ผรสุ\-ปฺปหารานํ\csromansharedfootnote{058995384153c6d7}{อสิปฺปหารานํ ผรสุปฺปหารานํ (ก)}. เอวญฺเจตํ ภิกฺขเว อสฺส. น โข ปนาหํ ภิกฺขเว สห ทุกฺเขน สห โทมนสฺเสน จตุนฺนํ อริยสจฺจานํ อภิสมยํ วทามิ, อปิ จาหํ ภิกฺขเว สหาว สุเขน สหาว โสมนสฺเสน จตุนฺนํ อริยสจฺจานํ อภิสมยํ วทามิ. กตเมสํ จตุนฺนํ, ทุกฺขสฺส อริยสจฺจสฺส ฯเปฯ ทุกฺขนิโรธ\-คามินิยา ปฏิปทาย อริยสจฺจสฺส.}

\csromanmatikaanchor{mtk.288}
\csromanmatikaanchor{mtk.289}
\csromantocmarkref{subparagraph}{6. ปณสุตฺต}{mtk.288}
\bookmark[dest={mtk.288},level=5]{6. ปณสุตฺต}
\csromantocmarkref{subparagraph}{7-8. สูริยสุตฺตานิ}{mtk.289}
\bookmark[dest={mtk.289},level=5]{7-8. สูริยสุตฺตานิ}
\prose{\csromanfolio{386}ตสฺมาติห ภิกฺขเว “อิทํ ทุกฺขนฺ”ติ โยโค กรณีโย ฯเปฯ “อยํ ทุกฺขนิโรธ\-คามินี ปฏิปทา”ติ โยโค กรณีโยติ.  ปญฺจมํ.}

\csromanheader{2}{6. \textbf{ปาณสุตฺต}}
\proseitem{1106}{เสยฺยถาปิ ภิกฺขเว ปุริโส ยํ อิมสฺมิํ ชมฺพุทีเป ติณกฏฺฐ\-สาขา\-ปลาสํ ตจฺเฉตฺวา เอกชฺฌํ สํหเรยฺย, เอกชฺฌํ สํหริตฺวา สูลํ กเรยฺย, สูลํ กริตฺวา เย มหาสมุทฺเท มหนฺตกา ปาณา, เต มหนฺตเกสุ สูเลสุ อาวุเนยฺย. เย มหาสมุทฺเท มชฺฌิมกา ปาณา, เต มชฺฌิมเกสุ สูเลสุ อาวุเนยฺย. เย มหาสมุทฺเท สุขุมกา ปาณา, เต สุขุมเกสุ สูเลสุ อาวุเนยฺย, อปริยาทินฺนา จ ภิกฺขเว มหาสมุทฺเท โอฬาริกา ปาณา อสฺสุ.}

\prose{อถ อิมสฺมิํ ชมฺพุทีเป ติณกฏฺฐ\-สาขา\-ปลาสํ ปริกฺขยํ ปริยาทานํ คจฺเฉยฺย. อิโต พหุตรา โข ภิกฺขเว มหาสมุทฺเท สุขุมกา ปาณา, เย น สุกรา สูเลสุ อาวุนิตุํ. ตํ กิสฺส เหตุ, สุขุมตฺตา ภิกฺขเว อตฺตภาวสฺส. เอวํ มหา โข ภิกฺขเว อปาโย, เอวํ มหนฺตสฺมา โข ภิกฺขเว อปายสฺมา ปริมุตฺโต ทิฏฺฐิสมฺปนฺโน ปุคฺคโล “อิทํ ทุกฺขนฺ”ติ ยถาภูตํ ปชานาติ ฯเปฯ “อยํ ทุกฺขนิโรธ\-คามินี ปฏิปทา”ติ ยถาภูตํ ปชานาติ.}

\prose{ตสฺมาติห ภิกฺขเว “อิทํ ทุกฺขนฺ”ติ โยโค กรณีโย ฯเปฯ “อยํ ทุกฺขนิโรธ\-คามินี ปฏิปทา”ติ โยโค กรณีโยติ.  ฉฏฺฐํ.}

\csromanheader{2}{7. ปฐม\-สูริย\-สุตฺต}
\proseitem{1107}{สูริยสฺส\csromansharedfootnote{c9f7a8b25f5a3849}{สุริยสฺส (สี, สฺยา, กํ, อิ)} ภิกฺขเว อุทยโต เอตํ ปุพฺพงฺคมํ เอตํ ปุพฺพนิมิตฺตํ, ยทิทํ อรุณุคฺคํ. เอวเมว โข ภิกฺขเว ภิกฺขุโน จตุนฺนํ อริยสจฺจานํ ยถาภูตํ อภิสมยาย เอตํ ปุพฺพงฺคมํ เอตํ ปุพฺพนิมิตฺตํ, ยทิทํ สมฺมาทิฏฺฐิ. ตสฺเสตํ ภิกฺขเว ภิกฺขุโน ปาฏิกงฺขํ “อิทํ ทุกฺขนฺ”ติ ยถาภูตํ ปชานิสฺสติ ฯเปฯ “อยํ ทุกฺขนิโรธ\-คามินี ปฏิปทา”ติ ยถาภูตํ ปชานิสฺสติ.}

\prose{ตสฺมาติห ภิกฺขเว “อิทํ ทุกฺขนฺ”ติ โยโค กรณีโย ฯเปฯ “อยํ ทุกฺขนิโรธ\-คามินี ปฏิปทา”ติ โยโค กรณีโยติ.  สตฺตมํ.}

\csromanheader{2}{8. ทุติย\-สูริย\-สุตฺต}
\proseitem{1108}{\csromanfolio{387}ยาวกีวญฺจ ภิกฺขเว จนฺทิมสูริยา โลเก นุปฺปชฺชนฺติ, เนว ตาว มหโต อาโลกสฺส ปาตุภาโว โหติ มหโต โอภาสสฺส, อนฺธตมํ ตทา โหติ อนฺธ\-การ\-ติมิสา, เนว ตาว รตฺตินฺทิวา\csromansharedfootnote{b5825402cbce0a9f}{รตฺติทิวา (ก)} ปญฺญายนฺติ, น มาสทฺธมาสา ปญฺญายนฺติ, น อุตุสํวจฺฉรา ปญฺญายนฺติ.}

\prose{ยโต จ โข ภิกฺขเว จนฺทิมสูริยา โลเก อุปฺปชฺชนฺติ, อถ มหโต อาโลกสฺส ปาตุภาโว โหติ มหโต โอภาสสฺส, เนว อนฺธการตมํ ตทา โหติ น อนฺธ\-การ\-ติมิสา, อถ รตฺตินฺทิวา ปญฺญายนฺติ, มาสทฺธมาสา ปญฺญายนฺติ, อุตุสํวจฺฉรา ปญฺญายนฺติ. เอวเมว โข ภิกฺขเว ยาวกีวญฺจ ตถาคโต โลเก นุปฺปชฺชติ อรหํ สมฺมาสมฺพุทฺโธ, เนว ตาว มหโต อาโลกสฺส ปาตุภาโว โหติ มหโต โอภาสสฺส, อนฺธตมํ ตทา โหติ อนฺธ\-การ\-ติมิสา, เนว ตาว จตุนฺนํ อริยสจฺจานํ อาจิกฺขณา โหติ เทสนา ปญฺญาปนา ปฏฺฐปนา วิวรณา วิภชนา อุตฺตานีกมฺมํ.}

\prose{ยโต จ โข ภิกฺขเว ตถาคโต โลเก อุปฺปชฺชติ อรหํ สมฺมาสมฺพุทฺโธ, อถ มหโต อาโลกสฺส ปาตุภาโว โหติ มหโต โอภาสสฺส, เนว อนฺธตมํ ตทา โหติ น อนฺธ\-การ\-ติมิสา, อถ โข จตุนฺนํ อริยสจฺจานํ อาจิกฺขณา โหติ เทสนา ปญฺญาปนา ปฏฺฐปนา วิวรณา วิภชนา อุตฺตานีกมฺมํ. กตเมสํ จตุนฺนํ, ทุกฺขสฺส อริยสจฺจสฺส ฯเปฯ ทุกฺขนิโรธ\-คามินิยา ปฏิปทาย อริยสจฺจสฺส.}

\prose{ตสฺมาติห ภิกฺขเว “อิทํ ทุกฺขนฺ”ติ โยโค กรณีโย ฯเปฯ “อยํ ทุกฺขนิโรธ\-คามินี ปฏิปทา”ติ โยโค กรณีโยติ.  ฉฏฺฐํ.}

\csromanmatikaanchor{mtk.290}
\csromantocmarkref{subparagraph}{9. อินฺทขีลสุตฺต}{mtk.290}
\bookmark[dest={mtk.290},level=5]{9. อินฺทขีลสุตฺต}
\csromanheader{6}{9. \textbf{อินฺทขีลสุตฺต}}
\proseitem{1109}{เย หิ เกจิ ภิกฺขเว สมณา วา พฺราหฺมณา วา “อิทํ ทุกฺขนฺ”ติ ยถาภูตํ นปฺปชานนฺติ ฯเปฯ “อยํ ทุกฺขนิโรธ\-คามินี ปฏิปทา”ติ ยถาภูตํ นปฺปชานนฺติ, เต อญฺญสฺส สมณสฺส วา พฺราหฺมณสฺส วา มุขํ อุลฺโลเกนฺติ\csromansharedfootnote{c2cadabbba626274}{โอโลเกนฺติ (สี, สฺยา)} “อยํ นูน ภวํ ชานํ ชานาติ, ปสฺสํ ปสฺสตี”ติ.}

\prose{\csromanfolio{388}เสยฺยถาปิ ภิกฺขเว ตูลปิจุ วา กปฺปาสปิจุ วา ลหุโก วาตูปาทาโน สเม ภูมิภาเค นิกฺขิตฺโต, ตเมนํ ปุรตฺถิโม วาโต ปจฺฉิเมน สํหเรยฺย, ปจฺฉิโม วาโต ปุรตฺถิเมน สํหเรยฺย, อุตฺตโร วาโต ทกฺขิเณน สํหเรยฺย, ทกฺขิโณ วาโต อุตฺตเรน สํหเรยฺย. ตํ กิสฺส เหตุ, ลหุกตฺตา ภิกฺขเว กปฺปาสปิจุโน. เอวเมว โข ภิกฺขเว เย หิ เกจิ สมณา วา พฺราหฺมณา วา “อิทํ ทุกฺขนฺ”ติ ยถาภูตํ นปฺปชานนฺติ ฯเปฯ “อยํ ทุกฺขนิโรธ\-คามินี ปฏิปทา”ติ ยถาภูตํ นปฺปชานนฺติ, เต อญฺญสฺส สมณสฺส วา พฺราหฺมณสฺส วา มุขํ อุลฺโลเกนฺติ “อยํ นุน ภวํ ชานํ ชานาติ, ปสฺสํ ปสฺสตี”ติ. ตํ กิสฺส เหตุ, อทิฏฺฐตฺตา ภิกฺขเว จตุนฺนํ อริยสจฺจานํ.}

\prose{เย จ โข เกจิ ภิกฺขเว สมณา วา พฺราหฺมณา วา “อิทํ ทุกฺขํนฺ”ติ ยถาภูตํ ปชานนฺติ ฯเปฯ “อยํ ทุกฺขนิโรธ\-คามินี ปฏิปทา”ติ ยถาภูตํ ปชานนฺติ, เต น อญฺญสฺส สมณสฺส วา พฺราหฺมณสฺส วา มุขํ อุลฺโลเกนฺติ “อยํ นุน ภวํ ชานํ ชานาติ, ปสฺสํ ปสฺสตี”ติ.}

\prose{เสยฺยถาปิ ภิกฺขเว อโยขีโล วา อินฺทขีโล วา คมฺภีรเนโม สุนิขาโต อจโล อสมฺปกมฺปี, ปุรตฺถิมาย เจปิ ทิสาย อาคจฺเฉยฺย ภุสา วาตวุฏฺฐิ, เนว สงฺกมฺเปยฺย\csromansharedfootnote{98c6d53629625c44}{เนว นํ สงฺกมฺเปยฺย (สี, อิ)} น สมฺปกมฺเปยฺย น สมฺปจาเลยฺย. ปจฺฉิมาย เจปิ ทิสาย ฯเปฯ. อุตฺตราย เจปิ ทิสาย ฯเปฯ. ทกฺขิณาย เจปิ ทิสาย อาคจฺเฉยฺย ภุสา วาตวุฏฺฐิ, เนว สงฺกมฺเปยฺย น สมฺปกมฺเปยฺย น สมฺปจาเลยฺย. ตํ กิสฺส เหตุ, คมฺภีรตฺตา ภิกฺขเว เนมสฺส สุนิขาตตฺตา อินฺทขีลสฺส. เอวเมว โข ภิกฺขเว เย จ โข เกจิ สมณา วา พฺราหฺมณา วา “อิทํ ทุกฺขนฺ”ติ ยถาภูตํ ปชานนฺติ ฯเปฯ “อยํ ทุกฺขนิโรธ\-คามินี ปฏิปทา”ติ ยถาภูตํ ปชานนฺติ, เต น อญฺญสฺส สมณสฺส วา พฺราหฺมณสฺส วา มุขํ อุลฺโลเกนฺติ “อยํ นูน ภวํ ชานํ ชานาติ, ปสฺสํ ปสฺสตี”ติ. ตํ กิสฺส เหตุ, สุทิฏฺฐตฺตา ภิกฺขเว จตุนฺนํ อริยสจฺจานํ. กตเมสํ จตุนฺนํ, ทุกฺขสฺส อริยสจฺจสฺส ฯเปฯ ทุกฺขนิโรธ\-คามินิยา ปฏิปทาย อริยสจฺจสฺส.}

\prose{ตสฺมาติห ภิกฺขเว “อิทํ ทุกฺขนฺ”ติ โยโค กรณีโย ฯเปฯ “อยํ ทุกฺขนิโรธ\-คามินี ปฏิปทา”ติ โยโค กรณีโยติ.  นวมํ.}

\csromanmatikaanchor{mtk.291}
\csromantocmarkref{subparagraph}{10. วาทตฺถิกสุตฺต}{mtk.291}
\bookmark[dest={mtk.291},level=5]{10. วาทตฺถิกสุตฺต}
\csromanheader{6}{10. วาทตฺถิก\-สุตฺต}
\proseitem{1110}{\csromanfolio{389}โย หิ โกจิ ภิกฺขเว ภิกฺขุ “อิทํ ทุกฺขนฺ”ติ ยถาภูตํ ปชานาติ ฯเปฯ “อยํ ทุกฺขนิโรธ\-คามินี ปฏิปทา”ติ ยถาภูตํ ปชานาติ. ปุรตฺถิมาย เจปิ ทิสาย อาคจฺเฉยฺย สมโณ วา พฺราหฺมโณ วา วาทตฺถิโก วาทคเวสี “วาทมสฺส อาโรเปสฺสามี”ติ, ตํ วต สหธมฺเมน สงฺกมฺเปสฺสติ วา สมฺปกมฺเปสฺสติ วา สมฺปจาเลสฺสติ วาติ เนตํ ฐานํ วิชฺชติ. ปจฺฉิมาย เจปิ ทิสาย ฯเปฯ. อุตฺตราย เจปิ ทิสาย ฯเปฯ. ทกฺขิณาย เจปิ ทิสาย อาคจฺเฉยฺย สมโณ วา พฺราหฺมโณ วา วาทตฺถิโก วาทคเวสี “วาทมสฺส อาโรเปสฺสามี”ติ, ตํ วต สหธมฺเมน สงฺกมฺเปสฺสติ วา สมฺปกมฺเปสฺสติ วา สมฺปจาเลสฺสติ วาติ เนตํ ฐานํ วิชฺชติ.}

\prose{เสยฺยถาปิ ภิกฺขเว สิลายูโป โสฬสกุกฺกุโก, ตสฺสสฺสุ อฏฺฐ กุกฺกุ เหฏฺฐา เนมงฺคมา, อฏฺฐ กุกฺกุ อุปริเนมสฺส, ปุรตฺถิมาย เจปิ ทิสาย อาคจฺเฉยฺย ภุสา วาตวุฏฺฐิ, เนว สงฺกมฺเปยฺย น สมฺปกมฺเปยฺย น สมฺปจาเลยฺย. ปจฺฉิมาย เจปิ ทิสาย ฯเปฯ. อุตฺตราย เจปิ ทิสาย ฯเปฯ. ทกฺขิณาย เจปิ ทิสาย อาคจฺเฉยฺย ภุสา วาตวุฏฺฐิ, เนว สงฺกมฺเปยฺย น สมฺปกมฺเปยฺย น สมฺปจาเลยฺย. ตํ กิสฺส เหตุ, คมฺภีรตฺตา ภิกฺขเว เนมสฺส สุนิขาตตฺตา สิลายูปสฺส. เอวเมว โข ภิกฺขเว โย หิ โกจิ ภิกฺขุ “อิทํ ทุกฺขนฺ”ติ ยถาภูตํ ปชานาติ ฯเปฯ “อยํ ทุกฺขนิโรธ\-คามินี ปฏิปทา”ติ ยถาภูตํ ปชานาติ. ปุรตฺถิมาย เจปิ ทิสาย อาคจฺเฉยฺย สมโณ วา พฺราหฺมโณ วา วาทตฺถิโก วาทคเวสี “วาทมสฺส อาโรเปสฺสามี”ติ, ตํ วต สหธมฺเมน สงฺกมฺเปสฺสติ วา สมฺปกมฺเปสฺสติ วา สมฺปจาเลสฺสติ วาติ เนตํ ฐานํ วิชฺชติ. ปจฺฉิมาย เจปิ ทิสาย ฯเปฯ. อุตฺตราย เจปิ ทิสาย ฯเปฯ. ทกฺขิณาย เจปิ ทิสาย อาคจฺเฉยฺย สมโณ วา พฺราหฺมโณ วา วาทตฺถิโก วาทคเวสี “วาทมสฺส อาโรเปสฺสามี”ติ, ตํ วต สหธมฺเมน สงฺกมฺเปสฺสติ วา สมฺปกมฺเปสฺสติ วา สมฺปจาเลสฺสติ วาติ เนตํ ฐานํ วิชฺชติ. ตํ กิสฺส เหตุ, สุทิฏฺฐตฺตา ภิกฺขเว จตุนฺนํ อริยสจฺจานํ. กตเมสํ จตุนฺนํ, ทุกฺขสฺส อริยสจฺจสฺส ฯเปฯ ทุกฺขนิโรธ\-คามินิยา ปฏิปทาย อริยสจฺจสฺส.}

\prosewithclosermidruled{ตสฺมาติห ภิกฺขเว “อิทํ ทุกฺขนฺ”ติ โยโค กรณีโย ฯเปฯ “อยํ ทุกฺขนิโรธ\-คามินี ปฏิปทา”ติ โยโค กรณีโยติ.  ทสมํ.}{สีสปาวน\-วคฺโค จตุตฺโถ.}
\csromancenter{\textbf{ตสฺสุทฺทานํ}}
\setlength{\csromangathapagewidth}{0pt}
\setlength{\csromangathawakwidth}{0pt}
\csromangathameasuregroup{สีสปา ขทิโร ทณฺโฑ, \\ ปาณา สูริยูปมา เทฺวธา,}{\csromangathabat{สีสปา ขทิโร ทณฺโฑ,}{เจลา สตฺติสเตน จ.} \\ \csromangathabat{ปาณา สูริยูปมา เทฺวธา,}{อินฺทขีโล จ วาทิโนติ.}}
\csromangathasetleft{สีสปา ขทิโร ทณฺโฑ, \\ ปาณา สูริยูปมา เทฺวธา,}
\csromangathagroup{\csromangathastanza{\csromanfolio{390}\csromangathabat{สีสปา ขทิโร ทณฺโฑ,}{เจลา สตฺติสเตน จ.} \\ \csromangathabat{ปาณา สูริยูปมา เทฺวธา,}{อินฺทขีโล จ วาทิโนติ.}}}

\csromanmatikaanchor{mtk.292}
\csromantocmarknopage{subparagraph}{5. ปปาตวคฺค}{mtk.292}
\bookmark[dest={mtk.292},level=5]{5. ปปาตวคฺค}
\csromanheader{6}{5. \textbf{ปปาตวคฺค}}
\csromanmatikaanchor{mtk.293}
\csromantocmarkref{subparagraph}{1. โลกจินฺตาสุตฺต}{mtk.293}
\bookmark[dest={mtk.293},level=5]{1. โลกจินฺตาสุตฺต}
\csromanheader{6}{1. \textbf{โลกจินฺตาสุตฺต}}
\proseitem{1111}{เอกํ สมยํ ภควา ราชคเห วิหรติ เวฬุวเน กลนฺทก\-นิวาเป. ตตฺร โข ภควา ภิกฺขู อามนฺเตสิ “ภูตปุพฺพํ ภิกฺขเว อญฺญตโร ปุริโส ราชคหา นิกฺขมิตฺวา ‘โลกจินฺตํ จินฺเตสฺสามี’ติ เยน สุมาคธา โปกฺขรณี เตนุปสงฺกมิ, อุปสงฺกมิตฺวา สุมาคธาย โปกฺขรณิยา ตีเร นิสีทิ โลกจินฺตํ จินฺเตนฺโต, อทฺทสา โข ภิกฺขเว โส ปุริโส สุมาคธาย โปกฺขรณิยา ตีเร จตุรงฺคินิํ เสนํ\csromansharedfootnote{84bec219107c4d97}{จตุรงฺคินิเสนํ (ก)} ภิสมุฬาลํ\csromansharedfootnote{26aca1df5920fe75}{ภิสมูลาลํ (อิ, ก)} ปวิสนฺตํ, ทิสฺวานสฺส เอตทโหสิ ‘อุมฺมตฺโตสฺมิ นามาหํ, วิเจโตสฺมิ นามาหํ, ยํ โลเก นตฺถิ ตํ มยา ทิฏฺฐนฺ’ติ. อถ โข โส ภิกฺขเว ปุริโส นครํ ปวิสิตฺวา มหา\-ชนกายสฺส อาโรเจสิ ‘อุมฺมตฺโตสฺมิ นามาหํ ภนฺเต, วิเจโตสฺมิ นามาหํ ภนฺเต, ยํ โลเก นตฺถิ ตํ มยา ทิฏฺฐนฺ’ติ”.}

\prose{กถํ ปน ตฺวํ อมฺโภ ปุริส อุมฺมตฺโต, กถํ วิเจโต, กิญฺจ โลเก นตฺถิ ยํ ตยา ทิฏฺฐนฺติ. อิธาหํ ภนฺเต ราชคหา นิกฺขมิตฺวา “โลกจินฺตํ จินฺเตสฺสามี”ติ เยน สุมาคธา โปกฺขรณี เตนุปสงฺกมิํ, อุปสงฺกมิตฺวา สุมาคธาย โปกฺขรณิยา ตีเร นิสีทิํ โลกจินฺตํ จินฺเตนฺโต, อทฺทสํ ขฺวาหํ ภนฺเต สุมาคธาย โปกฺขรณิยา ตีเร จตุรงฺคินิํ เสนํ ภิสมุฬาลํ ปวิสนฺตํ, เอวํ ขฺวาหํ ภนฺเต อุมฺมตฺโต, เอวํ วิเจโต, อิทญฺจ โลเก นตฺถิ ยํ มยา ทิฏฺฐนฺติ. ตคฺฆ ตฺวํ อมฺโภ ปุริส อุมฺมตฺโต, ตคฺฆ วิเจโต, อิทญฺจ โลเก นตฺถิ ยํ ตยา ทิฏฺฐนฺติ.}

\prose{ตํ โข ปน ภิกฺขเว โส ปุริโส ภูตํเยว อทฺทส, โน อภูตํ. ภูตปุพฺพํ ภิกฺขเว เทวา\-สุร\-สงฺคาโม สมุปพฺยูโฬฺห อโหสิ, ตสฺมิํ โข ปน ภิกฺขเว สงฺคาเม เทวา ชินิํสุ, อสุรา ปราชินิํสุ.}

\prose{\csromanfolio{391}ปราชิตา จ โข ภิกฺขเว อสุรา ภีตา ภิสมุฬาเลน อสุรปุรํ ปวิสิํสุ เทวานํเยว โมหยมานา.}

\prose{ตสฺมาติห ภิกฺขเว มา โลกจินฺตํ จินฺเตถ, “สสฺสโต โลโก”ติ วา, “อสสฺสโต โลโก”ติ วา, “อนฺตวา โลโก”ติ วา, “อนนฺตวา โลโก”ติ วา, “ตํ ชีวํ ตํ สรีรนฺ”ติ วา, “อญฺญํ ชีวํ อญฺญํ สรีรนฺ”ติ วา, “โหติ ตถาคโต ปรํ มรณา”ติ วา, “น โหติ ตถาคโต ปรํ มรณา”ติ วา, “โหติ จ น จ โหติ ตถาคโต ปรํ มรณา”ติ วา, “เนว โหติ น น โหติ ตถาคโต ปรํ มรณา”ติ วา. ตํ กิสฺส เหตุ, เนสา ภิกฺขเว จินฺตา อตฺถสํหิตา, นาทิ\-พฺรหฺม\-จริยกา, น นิพฺพิทาย น วิราคาย น นิโรธาย น อุปสมาย น อภิญฺญาย น สมฺโพธาย น นิพฺพานาย สํวตฺตติ.}

\prose{จินฺเตนฺตา โข ตุเมฺห ภิกฺขเว “อิทํ ทุกฺขนฺ”ติ จินฺเตยฺยาถ ฯเปฯ “อยํ ทุกฺขนิโรธ\-คามินี ปฏิปทา”ติ จินฺเตยฺยาถ. ตํ กิสฺส เหตุ, เอสา ภิกฺขเว จินฺตา อตฺถสํหิตา, เอสา อาทิ\-พฺรหฺม\-จริยกา, เอสา นิพฺพิทาย วิราคาย นิโรธาย อุปสมาย อภิญฺญาย สมฺโพธาย นิพฺพานาย สํวตฺตติ.}

\prose{ตสฺมาติห ภิกฺขเว “อิทํ ทุกฺขนฺ”ติ โยโค กรณีโย ฯเปฯ “อยํ ทุกฺขนิโรธ\-คามินี ปฏิปทา”ติ โยโค กรณีโยติ.  ปฐมํ.}

\csromanmatikaanchor{mtk.294}
\csromantocmarkref{subparagraph}{2. ปปาตสุตฺต}{mtk.294}
\bookmark[dest={mtk.294},level=5]{2. ปปาตสุตฺต}
\csromanheader{6}{2. \textbf{ปปาตสุตฺต}}
\proseitem{1112}{เอกํ สมยํ ภควา ราชคเห วิหรติ คิชฺฌกูเฏ ปพฺพเต. อถ โข ภควา ภิกฺขู อามนฺเตสิ “อายาม ภิกฺขเว เยน ปฏิภานกูโฏ, เตนุปสงฺกมิสฺสาม ทิวาวิหารายา”ติ. “เอวํ ภนฺเต”ติ โข เต ภิกฺขู ภควโต ปจฺจสฺโสสุํ. อถ โข ภควา สมฺพหุเลหิ ภิกฺขูหิ สทฺธิํ เยน ปฏิภานกูโฏ เตนุปสงฺกมิ. อทฺทสา โข อญฺญตโร ภิกฺขุ ปฏิภานกูเฏ มหนฺตํ ปปาตํ, ทิสฺวาน ภควนฺตํ เอตทโวจ “มหา วตายํ ภนฺเต ปปาโต, สุภยานโก ภนฺเต ปปาโต, อตฺถิ นุ โข ภนฺเต อิมมฺหา ปปาตา อญฺโญ ปปาโต มหนฺตตโร จ ภยานกตโร จา”ติ. อตฺถิ โข ภิกฺขุ อิมมฺหา ปปาตา อญฺโญ ปปาโต มหนฺตตโร จ ภยานกตโร จาติ.}

\prose{\csromanfolio{392}กตโม ปน ภนฺเต อิมมฺหา ปปาตา อญฺโญ ปปาโต มหนฺตตโร จ ภยานกตโร จาติ. เย หิ เกจิ ภิกฺขเว สมณา วา พฺราหฺมณา วา “อิทํ ทุกฺขนฺ”ติ ยถาภูตํ นปฺปชานนฺติ, “อยํ ทุกฺขสมุทโย”ติ ยถาภูตํ นปฺปชานนฺติ, “อยํ ทุกฺขนิโรโธ”ติ ยถาภูตํ นปฺปชานนฺติ, “อยํ ทุกฺขนิโรธ\-คามินี ปฏิปทา”ติ ยถาภูตํ นปฺปชานนฺติ. เต ชาติ\-สํวตฺตนิเกสุ สงฺขาเรสุ อภิรมนฺติ, ชรา\-สํวตฺตนิเกสุ สงฺขาเรสุ อภิรมนฺติ, มรณ\-สํวตฺตนิเกสุ สงฺขาเรสุ อภิรมนฺติ, โสก\-ปริเทว\-ทุกฺข\-โทมนสฺสุ\-ปายาส\-สํวตฺตนิเกสุ สงฺขาเรสุ อภิรมนฺติ. เต ชาติ\-สํวตฺตนิเกสุ สงฺขาเรสุ อภิรตา ชรา\-สํวตฺตนิเกสุ สงฺขาเรสุ อภิรตา มรณ\-สํวตฺตนิเกสุ สงฺขาเรสุ อภิรตา โสก\-ปริเทว\-ทุกฺข\-โทมนสฺสุ\-ปายาส\-สํวตฺตนิเกสุ สงฺขาเรสุ อภิรตา ชาติ\-สํวตฺตนิเกปิ สงฺขาเร อภิส\-งฺขโรนฺติ, ชรา\-สํวตฺตนิเกปิ สงฺขาเร อภิส\-งฺขโรนฺติ, มรณ\-สํวตฺตนิเกปิ สงฺขาเร อภิส\-งฺขโรนฺติ โสกปริเทวทุกฺขโทนมสฺสุปายาสสํวตฺตนิเกปิ สงฺขาเร อภิส\-งฺขโรนฺติ. เต ชาติ\-สํวตฺตนิเกปิ สงฺขาเร อภิส\-งฺขริ\-ตฺวา ชรา\-สํวตฺตนิเกปิ สงฺขาเร อภิส\-งฺขริ\-ตฺวา มรณ\-สํวตฺตนิเกปิ สงฺขาเร อภิส\-งฺขริ\-ตฺวา โสก\-ปริเทว\-ทุกฺข\-โทมนสฺสุ\-ปายาส\-สํวตฺตนิเกปิ สงฺขาเร อภิส\-งฺขริ\-ตฺวา ชาติปปาตมฺปิ ปปตนฺติ, ชราปปาตมฺปิ ปปตนฺติ, มรณ\-ปปาตมฺปิ ปปตนฺติ, โสก\-ปริเทว\-ทุกฺข\-โทมนสฺสุ\-ปายาส\-ปปาตมฺปิ ปปตนฺติ. เต น ปริมุจฺจนฺติ ชาติยา ชราย มรเณน โสเกหิ ปริเทเวหิ ทุกฺเขหิ โทมนสฺเสหิ อุปายาเสหิ, น ปริมุจฺจนฺติ ทุกฺขสฺมาติ วทามิ.}

\prose{เย จ โข เกจิ ภิกฺขเว สมณา วา พฺราหฺมณา วา “อิทํ ทุกฺขนฺ”ติ ยถาภูตํ ปชานนฺติ ฯเปฯ “อยํ ทุกฺขนิโรธ\-คามินี ปฏิปทา”ติ ยถาภูตํ ปชานนฺติ. เต ชาติ\-สํวตฺตนิเกสุ สงฺขาเรสุ นาภิรมนฺติ, ชรา\-สํวตฺตนิเกสุ สงฺขาเรสุ นาภิรมนฺติ, มรณ\-สํวตฺตนิเกสุ สงฺขาเรสุ นาภิรมนฺติ, โสก\-ปริเทว\-ทุกฺข\-โทมนสฺสุ\-ปายาส\-สํวตฺตนิเกสุ สงฺขาเรสุ นาภิรมนฺติ, เต ชาติ\-สํวตฺตนิเกสุ สงฺขาเรสุ อนภิรตา ชรา\-สํวตฺตนิเกสุ สงฺขาเรสุ อนภิรตา มรณ\-สํวตฺตนิเกสุ สงฺขาเรสุ อนภิรตา โสก\-ปริเทว\-ทุกฺข\-โทมนสฺสุ\-ปายาส\-สํวตฺตนิเกสุ สงฺขาเรสุ อนภิรตา ชาติ\-สํวตฺตนิเกปิ สงฺขาเร นาภิสงฺขาโรนฺติ, ชรา\-สํวตฺตนิเกปิ สงฺขาเร นาภิ\-ส\-งฺขโรนฺติ, มรณ\-สํวตฺตนิเกปิ สงฺขาเร นาภิ\-ส\-งฺขโรนฺติ,}

\prose{\csromanfolio{393}โสก\-ปริเทว\-ทุกฺข\-โทมนสฺสุ\-ปายาส\-สํวตฺตนิเกปิ สงฺขาเร นาภิ\-ส\-งฺขโรนฺติ, เต ชาติ\-สํวตฺตนิเกปิ สงฺขาเร อนภิสงฺขริตฺวา ชรา\-สํวตฺตนิเกปิ สงฺขาเร อนภิสงฺขริตฺวา มรณ\-สํวตฺตนิเกปิ สงฺขาเร อนภิสงฺขริตฺวา โสก\-ปริเทว\-ทุกฺข\-โทมนสฺสุ\-ปายาส\-สํวตฺตนิเกปิ สงฺขาเร อนภิสงฺขริตฺวา ชาติปปตมฺปิ นปฺปปตนฺติ, ชราปปาตมฺปิ นปฺปปตนฺติ, มรณ\-ปปาตมฺปิ นปฺปปตนฺติ, โสก\-ปริเทว\-ทุกฺข\-โทมนสฺสุ\-ปายาส\-ปปาตมฺปิ นปฺปปตนฺติ, เต ปริมุจฺจนฺติ ชาติยา ชราย มรเณน โสเกหิ ปริเทเวหิ ทุกฺเขหิ โทมนสฺเสหิ อุปายาเสหิ, ปริมุจฺจนฺติ ทุกฺขสฺมาติ วทามิ.}

\prose{ตสฺมาติห ภิกฺขเว “อิทํ ทุกฺขนฺ”ติ โยโค กรณีโย ฯเปฯ “อยํ ทุกฺขนิโรธ\-คามินี ปฏิปทา”ติ โยโค กรณีโยติ.  ทุติยํ.}

\csromanmatikaanchor{mtk.295}
\csromantocmarkref{subparagraph}{3. มหาปริฬาหสุตฺต}{mtk.295}
\bookmark[dest={mtk.295},level=5]{3. มหาปริฬาหสุตฺต}
\csromanheader{6}{3. มหาปริฬาห\-สุตฺต}
\proseitem{1113}{อตฺถิ ภิกฺขเว มหาปริฬาโห นาม นิรโย, ตตฺถ ยํ กิญฺจิ จกฺขุนา รูปํ ปสฺสติ, อนิฏฺฐ\-รูปญฺเญว ปสฺสติ โน อิฏฺฐรูปํ, อกนฺต\-รูปญฺเญว ปสฺสติ โน กนฺตรูปํ, อมนาป\-รูปญฺเญว ปสฺสติ โน มนาปรูปํ. ยํ กิญฺจิ โสเตน สทฺทํ สุณาติ ฯเปฯ. ยํ กิญฺจิ กาเยน โผฏฺฐพฺพํ ผุสติ ฯเปฯ. ยํ กิญฺจิ มนสา ธมฺมํ วิชานาติ, อนิฏฺฐ\-รูปญฺเญว วิชานาติ โน อิฏฺฐรูปํ, อกนฺต\-รูปญฺเญว วิชานาติ, โน กนฺตรูปํ, อมนาป\-รูปญฺเญว วิชานาติ, โน มนาปรูปนฺติ.}

\prose{เอวํ วุตฺเต อญฺญตโร ภิกฺขุ ภควนฺตํ เอตทโวจ “มหา วต โส ภนฺเต ปริฬาโห, สุมหา วต โส ภนฺเต ปริฬาโห, อตฺถิ นุ โข ภนฺเต เอตมฺหา ปริฬาหา อญฺโญ ปริฬาโห มหนฺตตโร เจว ภยานกตโร จา”ติ. อตฺถิ โข ภิกฺขุ เอตมฺหา ปริฬาหา อญฺโญ ปริฬาโห มหนฺตตโร จ ภยานกตโร จาติ.}

\prose{กตโม ปน ภนฺเต เอตมฺหา ปริฬาหา อญฺโญ ปริฬาโห มหนฺตตโร จ ภยานกตโร จาติ. เยหิ เกจิ ภิกฺขเว สมณา วา พฺราหฺมณา วา “อิทํ ทุกฺขนฺ”ติ ยถาภูตํ นปฺปชานนฺติ ฯเปฯ “อยํ ทุกฺขนิโรธ\-คามินี ปฏิปทา”ติ ยถาภูตํ นปฺปชานนฺติ. เต ชาติ\-สํวตฺตนิเกสุ สงฺขาเรสุ อภิรมนฺติ ฯเปฯ อภิรตา ฯเปฯ อภิส\-งฺขโรนฺติ ฯเปฯ อภิส\-งฺขริ\-ตฺวา ชาติปริฬาเหนปิ ปริฑยฺหนฺติ, ชราปริฬาเหนปิ ปริฑยฺหนฺติ, มรณปริฬาเหนปิ ปริฑยฺหนฺติ, โสกปริเทวทุกฺขโทมนสฺสุปายาสปริฬาเหนปิ \csromanfolio{394}ปริฑยฺหนฺติ, เต น ปริมุจฺจนฺติ ชาติยา ชราย มรเณน โสเกหิ ปริเทเวหิ ทุกฺเขหิ โทมนสฺเสหิ อุปายาเสหิ, น ปริมุจฺจนฺติ ทุกฺขสฺมาติ วทามิ.}

\prose{เย จ โข เกจิ ภิกฺขเว สมณา วา พฺราหฺมณา วา “อิทํ ทุกฺขนฺ”ติ ยถาภูตํ ปชานนฺติ ฯเปฯ “อยํ ทุกฺขนิโรธ\-คามินี ปฏิปทา”ติ ยถาภูตํ ปชานนฺติ. เต ชาติ\-สํวตฺตนิเกสุ สงฺขาเรสุ นาภิรมนฺติ ฯเปฯ อนภิรตา ฯเปฯ นาภิ\-ส\-งฺขโรนฺติ ฯเปฯ อนภิสงฺขริตฺวา ชาติปริฬาเหนปิ น ปริฑยฺหนฺติ, ชราปริฬาเหนปิ น ปริฑยฺหนฺติ, มรณปริฬาเหนปิ น ปริฑยฺหนฺติ, โสกปริเทว ทุกฺข โทมนสฺสุ\-ปายา\-สปริฬาเหนปิ น ปริฑยฺหนฺติ, เต ปริมุจฺจนฺติ ชาติยา ชราย มรเณน โสเกหิ ปริเทเวหิ ทุกฺเขหิ โทมนสฺเสหิ อุปายาเสหิ, ปริมุจฺจนฺติ ทุกฺขสฺมาติ วทามิ.}

\prose{ตสฺมาติห ภิกฺขเว “อิทํ ทุกฺขนฺ”ติ โยโค กรณีโย ฯเปฯ “อยํ ทุกฺขนิโรธ\-คามินี ปฏิปทา”ติ โยโค กรณีโยติ.  ตติยํ.}

\csromanmatikaanchor{mtk.296}
\csromantocmarkref{subparagraph}{4. กูฏาคารสุตฺต}{mtk.296}
\bookmark[dest={mtk.296},level=5]{4. กูฏาคารสุตฺต}
\csromanheader{6}{4. \textbf{กูฏาคารสุตฺต}}
\proseitem{1114}{โย หิ ภิกฺขเว\csromansharedfootnote{bd57be30369f4734}{โย จ โข ภิกฺขเว (สฺยา, ก)} เอวํ วเทยฺย “อหํ ทุกฺขํ อริยสจฺจํ ยถาภูตํ อนภิสเมจฺจ ฯเปฯ ทุกฺขนิโรธ\-คามินิํ ปฏิปทํ อริยสจฺจํ ยถาภูตํ อนภิสเมจฺจ สมฺมา ทุกฺขสฺสนฺตํ กริสฺสามี”ติ เนตํ ฐานํ วิชฺชติ.}

\prose{เสยฺยถาปิ ภิกฺขเว โย เอวํ วเทยฺย “อหํ กูฏาคารสฺส เหฏฺฐิมํ ฆรํ อกริตฺวา อุปริมํ ฆรํ อาโรเปสฺสามี”ติ เนตํ ฐานํ วิชฺชติ. เอวเมว โข ภิกฺขเว โย เอวํ วเทยฺย “อหํ ทุกฺขํ อริยสจฺจํ ยถาภูตํ อนภิสเมจฺจ ฯเปฯ ทุกฺขนิโรธ\-คามินิํ ปฏิปทํ อริยสจฺจํ ยถาภูตํ อนภิสเมจฺจ สมฺมา ทุกฺขสฺสนฺตํ กริสฺสามี”ติ เนตํ ฐานํ วิชฺชติ.}

\prose{โย จ โข ภิกฺขเว เอวํ วเทยฺย “อหํ ทุกฺขํ อริยสจฺจํ ยถาภูตํ อภิสเมจฺจ ฯเปฯ ทุกฺขนิโรธ\-คามินิํ ปฏิปทํ อริยสจฺจํ ยถาภูตํ อภิสเมจฺจ สมฺมา ทุกฺขสฺสนฺตํ กริสฺสามี”ติ ฐานเมตํ วิชฺชติ.}

\prose{เสยฺยถาปิ ภิกฺขเว โย เอวํ วเทยฺย “อหํ กูฏาคารสฺส เหฏฺฐิมํ ฆรํ กริตฺวา อุปริมํ ฆรํ อาโรเปสฺสามี”ติ ฐานเมตํ วิชฺชติ. เอวเมว}

\prose{\csromanfolio{395}โข ภิกฺขเว โย เอวํ วเทยฺย “อหํ ทุกฺขํ อริยสจฺจํ ยถาภูตํ อภิสเมจฺจ ฯเปฯ ทุกฺขนิโรธ\-คามินิํ ปฏิปทํ อริยสจฺจํ ยถาภูตํ อภิสเมจฺจ สมฺมา ทุกฺขสฺสนฺตํ กริสฺสามี”ติ ฐานเมตํ วิชฺชติ.}

\prose{ตสฺมาติห ภิกฺขเว “อิทํ ทุกฺขนฺ”ติ โยโค กรณีโย ฯเปฯ “อยํ ทุกฺขนิโรธ\-คามินี ปฏิปทา”ติ โยโค กรณีโยติ.  จตุตฺถํ.}

\csromanmatikaanchor{mtk.297}
\csromantocmarkref{subparagraph}{5. วาลสุตฺต}{mtk.297}
\bookmark[dest={mtk.297},level=5]{5. วาลสุตฺต}
\csromanheader{6}{5. \textbf{วาลสุตฺต}}
\proseitem{1115}{เอกํ สมยํ ภควา เวสาลิยํ วิหรติ มหาวเน กูฏาคาร\-สาลายํ. อถ โข อายสฺมา อานนฺโท ปุพฺพณฺหสมยํ นิวาเสตฺวา ปตฺต\-จีวรมา\-ทาย เวสาลิํ ปิณฺฑาย ปาวิสิ. อทฺทสา โข อายสฺมา อานนฺโท สมฺพหุเล ลิจฺฉวิ\-กุมารเก สนฺถาคาเร อุปาสนํ กโรนฺเต ทูรโตว สุขุเมน ตาฬ\-จฺฉิคฺคเฬน อสนํ อติปาเตนฺเต โปงฺขานุโปงฺขํ\csromansharedfootnote{a35790556c52f4ab}{โปขานุโปขํ (สฺยา, กํ)} อวิราธิตํ, ทิสฺวานสฺส เอตทโหสิ “สิกฺขิตา วติเม ลิจฺฉวิ\-กุมารกา, สุสิกฺขิตา วติเม ลิจฺฉวิ\-กุมารกา, ยตฺร หิ นาม ทูรโตว สุขุเมน ตาฬ\-จฺฉิคฺคเฬน อสนํ อติปาเตสฺสนฺติ โปงฺขานุโปงฺขํ อวิราธิตนฺ”ติ.}

\prose{อถ โข อายสฺมา อานนฺโท เวสาลิํ ปิณฺฑาย จริตฺวา ปจฺฉาภตฺตํ ปิณฺฑปาต\-ปฏิกฺกนฺโต เยน ภควา เตนุปสงฺกมิ, อุปสงฺกมิตฺวา ภควนฺตํ อภิวาเทตฺวา เอกมนฺตํ นิสีทิ, เอกมนฺตํ นิสินฺโน โข อายสฺมา อานนฺโท ภควนฺตํ เอตทโวจ “อิธาหํ ภนฺเต ปุพฺพณฺหสมยํ นิวาเสตฺวา ปตฺต\-จีวรมา\-ทาย เวสาลิํ ปิณฺฑาย ปาวิสิํ, อทฺทสํ ขฺวาหํ ภนฺเต สมฺพหุเล ลิจฺฉวิ\-กุมารเก สนฺถาคาเร อุปาสนํ กโรนฺเต ทูรโตว สุขุเมน ตาฬ\-จฺฉิคฺคเฬน อสนํ อติปาเตนฺเต โปงฺขานุโปงฺขํ อวิราธิตํ, ทิสฺวาน เม เอตทโหสิ ‘สิกฺขิตา วติเม ลิจฺฉวิ\-กุมารกา, สุสิกฺขิตา วติเม ลิจฺฉวิ\-กุมารกา, ยตฺร หิ นาม ทูรโตว สุขุเมน ตาฬ\-จฺฉิคฺคเฬน อสนํ อติปาเตสฺสนฺติ โปงฺขานุโปงฺขํ อวิราธิตนฺ’ติ”.}

\csromanmatikaanchor{mtk.298}
\csromantocmarkref{subparagraph}{6. อนฺทการสุตฺต}{mtk.298}
\bookmark[dest={mtk.298},level=5]{6. อนฺทการสุตฺต}
\prose{ตํ กิํ มญฺญสิ อานนฺท, กตมํ นุ โข ทุกฺกรตรํ วา ทุรภิสมฺภวตรํ วา, โย ทูรโตว สุขุเมน ตาฬ\-จฺฉิคฺคเฬน อสนํ อติปาเตยฺย โปงฺขานุโปงฺขํ อวิราธิตํ, โย วา สตฺตธา ภินฺนสฺส วาลสฺส โกฏิยา โกฏิํ ปฏิวิชฺเฌยฺยาติ. เอตเทว ภนฺเต ทุกฺกรตร\-ญฺเจว ทุรภิสมฺภวตรญฺจ, \csromanfolio{396}โย วา\csromansharedfootnote{665710d80bd591cc}{โย (สี)} สตฺตธา ภินฺนสฺส วาลสฺส โกฏิยา โกฏิํ ปฏิวิชฺเฌยฺยาติ. อถ โข\csromansharedfootnote{e6768b0e778e72a7}{อถ โข เต (สฺยา, กํ)} อานนฺท ทุปฺปฏิวิชฺฌตรํ ปฏิวิชฺฌนฺติ, เย “อิทํ ทุกฺขนฺ”ติ ยถาภูตํ ปฏิวิชฺฌนฺติ ฯเปฯ “อยํ ทุกฺขนิโรธ\-คามินี ปฏิปทา”ติ ยถาภูตํ ปฏิวิชฺฌนฺติ.}

\prose{ตสฺมาติหานนฺท “อิทํ ทุกฺขนฺ”ติ โยโค กรณีโย ฯเปฯ “อยํ ทุกฺขนิโรธ\-คามินี ปฏิปทา”ติ โยโค กรณีโยติ.  ปญฺจมํ.}

\csromanheader{2}{6. \textbf{อนฺธการสุตฺต}}
\proseitem{1116}{อตฺถิ ภิกฺขเว โลกนฺตริกา อฆา อสํวุตา อนฺธการา อนฺธ\-การ\-ติมิสา, ยตฺถมิเมสํ จนฺทิม\-สูริยานํ เอวํ\-มหิทฺธิกานํ เอวํ\-มหา\-นุภาวานํ อาภาย นานุโภนฺตีติ.}

\prose{เอวํ วุตฺเต อญฺญตโร ภิกฺขุ ภควนฺตํ เอตทโวจ “มหา วต โส ภนฺเต อนฺธกาโร, สุมหา วต โส ภนฺเต อนฺธกาโร, อตฺถิ นุ โข ภนฺเต เอตมฺหา อนฺธการา อญฺโญ อนฺธกาโร มหนฺตตโร จ ภยานกตโร จา”ติ. อตฺถิ โข ภิกฺขุ เอตมฺหา อนฺธการา อญฺโญ อนฺธกาโร มหนฺตตโร จ ภยานกตโร จาติ.}

\prose{กตโม ปน ภนฺเต เอตมฺหา อนฺธการา อญฺโญ อนฺธกาโร มหนฺตตโร จ ภยานกตโร จาติ. เย หิ เกจิ ภิกฺขุ สมณา วา พฺราหฺมณา วา “อิทํ ทุกฺขนฺ”ติ ยถาภูตํ นปฺปชานนฺติ ฯเปฯ “อยํ ทุกฺขนิโรธ\-คามินี ปฏิปทา”ติ ยถาภูตํ นปฺปชานนฺติ. เต ชาติ\-สํวตฺตนิเกสุ สงฺขาเรสุ อภิรมนฺติ ฯเปฯ อภิรตา ฯเปฯ อภิส\-งฺขโรนฺติ ฯเปฯ อภิส\-งฺขริ\-ตฺวา ชาต\-นฺธการมฺปิ ปปตนฺติ, ชรนฺธการมฺปิ ปปตนฺติ, มรณ\-นฺธการมฺปิ ปปตนฺติ, โสกปริเทวทุกฺขโทมนสฺสุปายาส\-นฺธการมฺปิ ปปตนฺติ. เต น ปริมุจฺจนฺติ ชาติยา ชราย มรเณน โสเกหิ ปริเทเวหิ ทุกฺเขหิ โทมนสฺเสหิ อุปายาเสหิ, น ปริมุจฺจนฺติ ทุกฺขสฺมาติ วทามิ.}

\prose{เย จ โข เกจิ ภิกฺขุ สมณา วา พฺราหฺมณา วา “อิทํ ทุกฺขนฺ”ติ ยถาภูตํ ปชานนฺติ ฯเปฯ “อยํ ทุกฺขนิโรธ\-คามินี ปฏิปทา”ติ ยถาภูตํ ปชานนฺติ. เต ชาติ\-สํวตฺตนิเกสุ สงฺขาเรสุ นาภิรมนฺติ ฯเปฯ อนภิรตา ฯเปฯ นาภิ\-ส\-งฺขโรนฺติ ฯเปฯ อนภิสงฺขริตฺวา ชาต\-นฺธการมฺปิ}

\csromanmatikaanchor{mtk.299}
\csromantocmarkref{subparagraph}{7-8. ฉิคฺคฬยุคสุตฺตานิ}{mtk.299}
\bookmark[dest={mtk.299},level=5]{7-8. ฉิคฺคฬยุคสุตฺตานิ}
\prose{\csromanfolio{397}นปฺปปตนฺติ, ชรนฺธการมฺปิ นปฺปปตนฺติ, มรณ\-นฺธการมฺปิ นปฺปปตนฺติ, โสกปริเทวทุกฺขโทมนสฺสุปายาส\-นฺธการมฺปิ นปฺปปตนฺติ. เต ปริมุจฺจนฺติ ชาติยา ชราย มรเณน โสเกหิ ปริเทเวหิ ทุกฺเขหิ โทมนสฺเสหิ อุปายาเสหิ, ปริมุจฺจนฺติ ทุกฺขสฺมาติ วทามิ.}

\prose{ตสฺมาติห ภิกฺขเว “อิทํ ทุกฺขนฺ”ติ โยโค กรณีโย ฯเปฯ “อยํ ทุกฺขนิโรธ\-คามินี ปฏิปทา”ติ โยโค กรณีโยติ.  ฉฏฺฐํ.}

\csromanheader{2}{7. ปฐม\-ฉิคฺคฬ\-ยุค\-สุตฺต}
\proseitem{1117}{เสยฺยถาปิ ภิกฺขเว ปุริโส มหาสมุทฺเท เอกจฺฉิคฺคฬํ ยุคํ ปกฺขิเปยฺย, ตตฺราปิสฺส กาโณ กจฺฉโป, โส วสฺสสตสฺส วสฺสสตสฺส อจฺจเยน สกิํ สกิํ อุมฺมุชฺเชยฺย. ตํ กิํ มญฺญถ ภิกฺขเว, อปิ นุ โข กาโณ กจฺฉโป วสฺสสตสฺส วสฺสสตสฺส อจฺจเยน สกิํ สกิํ อุมฺมุชฺชนฺโต อมุสฺมิํ เอกจฺฉิคฺคเฬ ยุเค คีวํ ปเวเสยฺยาติ. ยทิ นูน ภนฺเต กทาจิ กรหจิ ทีฆสฺส อทฺธุโน อจฺจเยนาติ.}

\prose{ขิปฺปตรํ โข โส ภิกฺขเว กาโณ กจฺฉโป วสฺสสตสฺส วสฺสสตสฺส อจฺจเยน สกิํ สกิํ อุมฺมุชฺชนฺโต อมุสฺมิํ เอกจฺฉิคฺคเฬ ยุเค คีวํ ปเวเสยฺย, น เตฺววาหํ ภิกฺขเว สกิํ วินิปาตคเตน พาเลน\csromansharedfootnote{be5003e3e68e62e7}{วินีตคเตน พหุเลน (ก)} มนุสฺสตฺตํ วทามิ.}

\prose{ตํ กิสฺส เหตุ, น เหตฺถ ภิกฺขเว อตฺถิ ธมฺมจริยา สมจริยา กุสลกิริยา ปุญฺญกิริยา, อญฺญมญฺญ\-ขาทิกา เอตฺถ ภิกฺขเว วตฺตติ ทุพฺพลขาทิกา. ตํ กิสฺส เหตุ, อทิฏฺฐตฺตา ภิกฺขเว จตุนฺนํ อริยสจฺจานํ. กตเมสํ จตุนฺนํ, ทุกฺขสฺส อริยสจฺจสฺส ฯเปฯ ทุกฺขนิโรธ\-คามินิยา ปฏิปทาย อริยสจฺจสฺส.}

\prose{ตสฺมาติห ภิกฺขเว “อิทํ ทุกฺขนฺ”ติ โยโค กรณีโย ฯเปฯ “อยํ ทุกฺขนิโรธ\-คามินี ปฏิปทา”ติ โยโค กรณีโยติ.  สตฺตมํ.}

\csromanheader{2}{8. ทุติย\-ฉิคฺคฬ\-ยุค\-สุตฺต}
\csromanmatikaanchor{mtk.300}
\csromantocmarkref{subparagraph}{9-10. สิเนรุปพฺพตราชสุตฺตานิ}{mtk.300}
\bookmark[dest={mtk.300},level=5]{9-10. สิเนรุปพฺพตราชสุตฺตานิ}
\proseitem{1118}{\csromanfolio{398}เสยฺยถาปิ ภิกฺขเว อยํ มหาปถวี เอโกทกา อสฺส, ตตฺร ปุริโส เอกจฺฉิคฺคฬํ ยุคํ ปกฺขิเปยฺย, ตเมนํ ปุรตฺถิโม วาโต ปจฺฉิเมน สํหเรยฺย. ปจฺฉิโม วาโต ปุรตฺถิเมน สํหเรยฺย. อุตฺตโร วาโต ทกฺขิเณน สํหเรยฺย. ทกฺขิโณ วาโต อุตฺตเรน สํหเรยฺย. ตตฺรสฺส กาโณ กจฺฉโป, โส วสฺสสตสฺส วสฺสสตสฺส อจฺจเยน สกิํ สกิํ อุมฺมุชฺเชยฺย. ตํ กิํ มญฺญถ ภิกฺขเว, อปิ นุ โข กาโณ กจฺฉโป วสฺสสตสฺส วสฺสสตสฺส อจฺจเยน สกิํ สกิํ อุมฺมุชฺชนฺโต อมุสฺมิํ เอกจฺฉิคฺคเฬ ยุเค คีวํ ปเวเสยฺยาติ. อธิจฺจมิทํ ภนฺเต, ยํ โส กาโณ กจฺฉโป วสฺสสตสฺส วสฺสสตสฺส อจฺจเยน สกิํ สกิํ อุมฺมุชฺชนฺโต อมุสฺมิํ เอกจฺฉิคฺคเฬ ยุเค คีวํ ปเวเสยฺยาติ.}

\prose{เอวํ อธิจฺจมิทํ ภิกฺขเว ยํ มนุสฺสตฺตํ ลภติ, เอวํ อธิจฺจมิทํ ภิกฺขเว ยํ ตถาคโต โลเก อุปฺปชฺชติ อรหํ สมฺมาสมฺพุทฺโธ, เอวํ อธิจฺจมิทํ ภิกฺขเว ยํ ตถาคต\-ปฺปเวทิโต ธมฺมวินโย โลเก ทิพฺพติ, ตสฺสิทํ\csromansharedfootnote{25e5c6f528d80563}{ตยิทํ (?)} ภิกฺขเว มนุสฺสตฺตํ ลทฺธํ ตถาคโต โลเก อุปฺปนฺโน อรหํ สมฺมาสมฺพุทฺโธ, ตถาคต\-ปฺปเวทิโต จ ธมฺมวินโย โลเก ทิพฺพติ.}

\prose{ตสฺมาติห ภิกฺขเว “อิทํ ทุกฺขนฺ”ติ โยโค กรณีโย ฯเปฯ “อยํ ทุกฺขนิโรธ\-คามินี ปฏิปทา”ติ โยโค กรณีโยติ.  อฏฺฐมํ.}

\csromanheader{2}{9. ปฐม\-สิเนรุปพฺพตราช\-สุตฺต}
\proseitem{1119}{เสยฺยถาปิ ภิกฺขเว ปุริโส สิเนรุสฺส ปพฺพตราชสฺส สตฺต มุคฺคมตฺติโย ปาสาณสกฺขรา อุปนิกฺขิเปยฺย. ตํ กิํ มญฺญถ ภิกฺขเว, กตมํ นุ โข พหุตรํ, ยา วา\csromansharedfootnote{3225b06e7e15fb1d}{ยา จ} สตฺต มุคฺคมตฺติโย ปาสาณสกฺขรา อุปนิกฺขิตฺตา, โย วา\csromansharedfootnote{fe9a8a7f9c1f24c6}{โย จ (สฺยา, กํ, อิ, ก) สํ 1. 356 ปิฏฺเฐปิ.} สิเนรุ\-ปพฺพตราชาติ. เอตเทว ภนฺเต พหุตรํ ยทิทํ สิเนรุ\-ปพฺพตราชา, อปฺปมตฺติกา สตฺต มุคฺคมตฺติโย ปาสาณสกฺขรา อุปนิกฺขิตฺตา, สงฺขมฺปิ น อุเปนฺติ อุปนิธมฺปิ น อุเปนฺติ กลภาคมฺปิ น อุเปนฺติ, สิเนรุ\-ปพฺพตราชานํ อุปนิธาย สตฺต มุคฺคมตฺติโย ปาสาณสกฺขรา อุปนิกฺขิตฺตาติ. เอวเมว โข ภิกฺขเว อริยสาวกสฺส ทิฏฺฐิ\-สมฺปนฺนสฺส}

\prose{\csromanfolio{399}ปุคฺคลสฺส อภิสเมตาวิโน เอตเทว พหุตรํ ทุกฺขํ, ยทิทํ ปริกฺขีณํ ปริยาทินฺนํ อปฺปมตฺตกํ อวสิฏฺฐํ, สงฺขมฺปิ น อุเปติ อุปนิธมฺปิ น อุเปติ กลภาคมฺปิ น อุเปติ, ปุริมํ ทุกฺข\-กฺขนฺธํ ปริกฺขีณํ ปริยาทินฺนํ อุปนิธาย, ยทิทํ สตฺตกฺขตฺตุ\-ปรมตา. โย “อิทํ ทุกฺขนฺ”ติ ยถาภูตํ ปชานาติ. “อยํ ทุกฺขนิโรธ\-คามินี ปฏิปทา”ติ ยถาภูตํ ปชานาติ.}

\prose{ตสฺมาติห ภิกฺขเว “อิทํ ทุกฺขนฺ”ติ โยโค กรณีโย ฯเปฯ “อยํ ทุกฺขนิโรธ\-คามินี ปฏิปทา”ติ โยโค กรณีโยติ.  นวมํ.}

\csromanheader{2}{10. ทุติย\-สิเนรุปพฺพตราช\-สุตฺต}
\proseitem{1120}{เสยฺยถาปิ ภิกฺขเว สิเนรุ\-ปพฺพตราชายํ ปริกฺขยํ ปริยาทานํ คจฺเฉยฺย ฐเปตฺวา สตฺต มุคฺคมตฺติโย ปาสาณสกฺขรา. ตํ กิํ มญฺญถ ภิกฺขเว, กตมํ นุ โข พหุตรํ, ยํ วา สิเนรุสฺส ปพฺพตราชสฺส ปริกฺขีณํ ปริยาทินฺนํ, ยา วา สตฺต มุคฺคมตฺติโย ปาสาณสกฺขรา อวสิฏฺฐาติ. เอตเทว ภนฺเต พหุตรํ สิเนรุสฺส ปพฺพตราชสฺส, ยทิทํ ปริกฺขีณํ ปริยาทินฺนํ, อปฺปมตฺติกา สตฺต มุคฺคมตฺติโย ปาสาณสกฺขรา อวสิฏฺฐา, สงฺขมฺปิ น อุเปนฺติ อุปนิธมฺปิ น อุเปนฺติ กลภาคมฺปิ น อุเปนฺติ, สิเนรุสฺส ปพฺพตราชสฺส ปริกฺขีณํ ปริยาทินฺนํ, อุปนิธาย สตฺต มุคฺคมตฺติโย ปาสาณสกฺขรา อวสิฏฺฐาติ. เอวเมว โข ภิกฺขเว อริยสาวกสฺส ทิฏฺฐิ\-สมฺปนฺนสฺส ปุคฺคลสฺส อภิสเมตาวิโน เอตเทว พหุตรํ ทุกฺขํ, ยทิทํ ปริกฺขีณํ ปริยาทินฺนํ อปฺปมตฺตกํ อวสิฏฺฐํ, สงฺขมฺปิ น อุเปติ อุปนิธมฺปิ น อุเปติ กลภาคมฺปิ น อุเปติ, ปุริมํ ทุกฺข\-กฺขนฺธํ ปริกฺขีณํ ปริยาทินฺนํ อุปนิธาย, ยทิทํ สตฺตกฺขตฺตุ\-ปรมตา. โย “อิทํ ทุกฺขนฺ”ติ ยถาภูตํ ปชานาติ ฯเปฯ “อยํ ทุกฺขนิโรธ\-คามินี ปฏิปทา”ติ ยถาภูตํ ปชานาติ.}

\prosewithclosermidruled{ตสฺมาติห ภิกฺขเว “อิทํ ทุกฺขนฺ”ติ โยโค กรณีโย ฯเปฯ “อยํ ทุกฺขนิโรธ\-คามินี ปฏิปทา”ติ โยโค กรณีโยติ.  ทสมํ.}{ปปาตวคฺโค ปญฺจโม.}
\csromancenter{\textbf{ตสฺสุทฺทานํ}}
\setlength{\csromangathapagewidth}{0pt}
\setlength{\csromangathawakwidth}{0pt}
\csromangathameasuregroup{จินฺตา ปปาโต ปริฬาโห, \\ ฉิคฺคเฬน จ เทฺว วุตฺตา,}{\csromangathabat{จินฺตา ปปาโต ปริฬาโห,}{กูฏํ วาลนฺธกาโร จ.} \\ \csromangathabat{ฉิคฺคเฬน จ เทฺว วุตฺตา,}{สิเนรุ อปเร ทุเวติ.}}
\csromangathasetleft{จินฺตา ปปาโต ปริฬาโห, \\ ฉิคฺคเฬน จ เทฺว วุตฺตา,}
\csromangathagroup{\csromangathastanza{\csromangathabat{จินฺตา ปปาโต ปริฬาโห,}{กูฏํ วาลนฺธกาโร จ.} \\ \csromangathabat{ฉิคฺคเฬน จ เทฺว วุตฺตา,}{สิเนรุ อปเร ทุเวติ.}}}

\csromanmatikaanchor{mtk.301}
\csromantocmarknopage{subparagraph}{6. อภิสมยวคฺค}{mtk.301}
\bookmark[dest={mtk.301},level=5]{6. อภิสมยวคฺค}
\csromanheader{6}{6. อภิสมย\-วคฺค}
\csromanmatikaanchor{mtk.302}
\csromantocmarkref{subparagraph}{1. นขสิขสุตฺต}{mtk.302}
\bookmark[dest={mtk.302},level=5]{1. นขสิขสุตฺต}
\csromanheader{6}{1. นขสิข\-สุตฺต}
\proseitem{1121}{\csromanfolio{400}อถ โข ภควา ปริตฺตํ นขสิขายํ ปํสุํ อาโรเปตฺวา ภิกฺขู อามนฺเตสิ “ตํ กิํ มญฺญถ ภิกฺขเว, กตมํ นุ โข พหุตรํ โย วายํ มยา ปริตฺโต นขสิขายํ ปํสุ อาโรปิโต, อยํ วา มหาปถวี”ติ. เอตเทว ภนฺเต พหุตรํ ยทิทํ มหาปถวี, อปฺปมตฺตกายํ ภควตา ปริตฺโต นขสิขายํ ปํสุ อาโรปิโต, สงฺขมฺปิ น อุเปติ อุปนิธมฺปิ น อุเปติ กลภาคมฺปิ น อุเปติ, มหาปถวิํ อุปนิธาย ภควตา ปริตฺโต นขสิขายํ ปํสุ อาโรปิโตติ. เอวเมว โข ภิกฺขเว อริยสาวกสฺส ทิฏฺฐิ\-สมฺปนฺนสฺส ปุคฺคลสฺส อภิสเมตาวิโน เอตเทว พหุตรํ ทุกฺขํ, ยทิทํ ปริกฺขีณํ ปริยาทินฺนํ, อปฺปมตฺตกํ อวสิฏฺฐํ, สงฺขมฺปิ น อุเปติ อุปนิธมฺปิ น อุเปติ กลภาคมฺปิ น อุเปติ, ปุริมํ ทุกฺข\-กฺขนฺธํ ปริกฺขีณํ ปริยาทินฺนํ อุปนิธาย, ยทิทํ สตฺตกฺขตฺตุ\-ปรมตา. โย “อิทํ ทุกฺขนฺ”ติ ยถาภูตํ ปชานาติ ฯเปฯ “อยํ ทุกฺขนิโรธ\-คามินี ปฏิปทา”ติ ยถาภูตํ ปชานาติ.}

\prose{ตสฺมาติห ภิกฺขเว “อิทํ ทุกฺขนฺ”ติ โยโค กรณีโย ฯเปฯ “อยํ ทุกฺขนิโรธ\-คามินี ปฏิปทา”ติ โยโค กรณีโย.  ปฐมํ.}

\csromanmatikaanchor{mtk.303}
\csromantocmarkref{subparagraph}{2. โปกฺขรณีสุตฺต}{mtk.303}
\bookmark[dest={mtk.303},level=5]{2. โปกฺขรณีสุตฺต}
\csromanheader{6}{2. โปกฺขรณี\-สุตฺต}
\proseitem{1122}{เสยฺยถาปิ ภิกฺขเว โปกฺขรณี ปญฺญาส\-โยชนานิ อายาเมน ปญฺญาส\-โยชนานิ วิตฺถาเรน ปญฺญาส\-โยชนานิ อุพฺเพเธน ปุณฺณา อุทกสฺส สมติตฺติกา กากเปยฺยา, ตโต ปุริโส กุสคฺเคน อุทกํ อุทฺธเรยฺย. ตํ กิํ มญฺญถ ภิกฺขเว, กตมํ นุ โข พหุตรํ, ยํ วา กุสคฺเคน อุพฺภตํ, ยํ วา โปกฺขรณิยา อุทกนฺติ. เอตเทว ภนฺเต พหุตรํ ยทิทํ โปกฺขรณิยา อุทกํ, อปฺปมตฺตกํ กุสคฺเคน อุทกํ อุพฺภตํ, สงฺขมฺปิ น อุเปติ อุปนิธมฺปิ น อุเปติ กลภาคมฺปิ น อุเปติ, โปกฺขรณิยา อุทกํ อุปนิธาย กุสคฺเคน อุทกํ อุพฺภตนฺติ. เอวเมว โข ภิกฺขเว อริยสาวกสฺส ฯเปฯ โยโค กรณีโยติ.  ทุติยํ.}

\csromanheader{2}{3. ปฐม\-สํเภชฺช\-สุตฺต}
\csromanmatikaanchor{mtk.304}
\csromanmatikaanchor{mtk.305}
\csromantocmarkref{subparagraph}{3-4. สํเภชฺชสุตฺตานิ}{mtk.304}
\bookmark[dest={mtk.304},level=5]{3-4. สํเภชฺชสุตฺตานิ}
\csromantocmarkref{subparagraph}{5-6. มหาปถวีสุตฺตานิ}{mtk.305}
\bookmark[dest={mtk.305},level=5]{5-6. มหาปถวีสุตฺตานิ}
\proseitem{1123}{\csromanfolio{401}เสยฺยถาปิ ภิกฺขเว ยตฺถิมา มหานทิโย สํสนฺทนฺติ สเมนฺติ. เสยฺยถิทํ, คงฺคา ยมุนา อจิรวตี สรภู มหี, ตโต ปุริโส เทฺว วา ตีณิ วา อุทกผุสิตานิ อุทฺธเรยฺย. ตํ กิํ มญฺญถ ภิกฺขเว, กตมํ นุ โข พหุตรํ, ยานิ เทฺว วา ตีณิ วา อุทกผุสิตานิ อุพฺภตานิ, ยํ วา สํเภชฺชอุทกนฺติ. เอตเทว ภนฺเต พหุตรํ ยทิทํ สํเภชฺชอุทกํ, อปฺปมตฺตกานิ เทฺว วา ตีณิ วา อุทกผุสิตานิ อุพฺภตานิ, สงฺขมฺปิ น อุเปนฺติ อุปนิธมฺปิ น อุเปนฺติ กลภาคมฺปิ น อุเปนฺติ, สํเภชฺชอุทกํ อุปนิธาย เทฺว วา ตีณิ วา อุทกผุสิตานิ อุพฺภตานีติ. เอวเมว โข ภิกฺขเว อริยสาวกสฺส ฯเปฯ โยโค กรณีโยติ.  ตติยํ.}

\csromanheader{2}{4. ทุติย\-สํเภชฺช\-สุตฺต}
\proseitem{1124}{เสยฺยถาปิ ภิกฺขเว ยตฺถิมา มหานทิโย สํสนฺทนฺติ สเมนฺติ. เสยฺยถิทํ, คงฺคา ยมุนา อจิรวตี สรภู มหี, ตํ อุทกํ ปริกฺขยํ ปริยาทานํ คจฺเฉยฺย ฐเปตฺวา เทฺว วา ตีณิ วา อุทกผุสิตานิ. ตํ กิํ มญฺญถ ภิกฺขเว, กตมํ นุ โข พหุตรํ, ยํ วา สํเภชฺชอุทกํ ปริกฺขีณํ ปริยาทินฺนํ, ยานิ เทฺว วา ตีณิ วา อุทกผุสิตานิ อวสิฏฺฐานีติ. เอตเทว ภนฺเต พหุตรํ สํเภชฺชอุทกํ, ยทิทํ ปริกฺขีณํ ปริยาทินฺนํ, อปฺปมตฺตกานิ เทฺว วา ตีณิ วา อุทกผุสิตานิ อวสิฏฺฐานิ, สงฺขมฺปิ น อุเปนฺติ อุปนิธมฺปิ น อุเปนฺติ กลภาคมฺปิ น อุเปนฺติ, สํเภชฺชอุทกํ ปริกฺขีณํ ปริยาทินฺนํ อุปนิธาย เทฺว วา ตีณิ วา อุทกผุสิตานิ อวสิฏฺฐานีติ. เอวเมว โข ภิกฺขเว อริยสาวกสฺส ฯเปฯ โยโค กรณีโยติ.  จตุตฺถํ.}

\csromanheader{2}{5. ปฐม\-มหาปถวี\-สุตฺต}
\proseitem{1125}{เสยฺยถาปิ ภิกฺขเว ปุริโส มหาปถวิยา สตฺต โกลฏฺฐิ\-มตฺติโย คุฬิกา อุปนิกฺขิเปยฺย. ตํ กิํ มญฺญถ ภิกฺขเว, กตมํ นุ โข พหุตรํ, ยา วา สตฺต โกลฏฺฐิ\-มตฺติโย คุฬิกา อุปนิกฺขิตฺตา, อยํ วา มหาปถวีติ. เอตเทว ภนฺเต พหุตรํ ยทิทํ มหาปถวี, อปฺปมตฺติกา สตฺต โกลฏฺฐิ\-มตฺติโย คุฬิกา อุปนิกฺขิตฺตา, สงฺขมฺปิ น อุเปนฺติ อุปนิธมฺปิ น อุเปนฺติ กลภาคมฺปิ น อุเปนฺติ, มหาปถวิํ อุปนิธาย สตฺต โกลฏฺฐิ\-มตฺติโย คุฬิกา อุปนิกฺขิตฺตาติ. เอวเมว โข ภิกฺขเว อริยสาวกสฺส ฯเปฯ โยโค กรณีโยติ.  ปญฺจมํ.}

\csromanheader{2}{6. ทุติย\-มหาปถวี\-สุตฺต}
\csromanmatikaanchor{mtk.306}
\csromantocmarkref{subparagraph}{7-8. มหาสมุทฺทสุตฺตานิ}{mtk.306}
\bookmark[dest={mtk.306},level=5]{7-8. มหาสมุทฺทสุตฺตานิ}
\proseitem{1126}{\csromanfolio{402}เสยฺยถาปิ ภิกฺขเว มหาปถวี ปริกฺขยํ ปริยาทานํ คจฺเฉยฺย ฐเปตฺวา สตฺต โกลฏฺฐิ\-มตฺติโย คุฬิกา. ตํ กิํ มญฺญถ ภิกฺขเว, กตมํ นุ โข พหุตรํ, ยํ วา มหาปถวิยา ปริกฺขีณํ ปริยาทินฺนํ, ยา วา สตฺต โกลฏฺฐิ\-มตฺติโย คุฬิกา อวสิฏฺฐาติ. เอตเทว ภนฺเต พหุตรํ มหาปถวิยา ยทิทํ ปริกฺขีณํ ปริยาทินฺนํ, อปฺปมตฺติกา สตฺต โกลฏฺฐิ\-มตฺติโย คุฬิกา อวสิฏฺฐา, สงฺขมฺปิ น อุเปนฺติ อุปนิธมฺปิ น อุเปนฺติ กลภาคมฺปิ น อุเปนฺติ, มหาปถวี ปริกฺขีณํ ปริยาทินฺนํ อุปนิธาย สตฺต โกลฏฺฐิ\-มตฺติโย คุฬิกา อวสิฏฺฐาติ. เอวเมว โข ภิกฺขเว อริยสาวกสฺส ฯเปฯ โยโค กรณีโยติ.  ฉฏฺฐํ.}

\csromanheader{2}{7. ปฐม\-มหาสมุทฺท\-สุตฺต}
\proseitem{1127}{เสยฺยถาปิ ภิกฺขเว ปุริโส มหาสมุทฺทโต เทฺว วา ตีณิ วา อุทกผุสิตานิ อุทฺธริตานิ. ตํ กิํ มญฺญถ ภิกฺขเว, กตมํ นุ โข พหุตรํ, ยานิ เทฺว วา ตีณิ วา อุทกผุสิตานิ อุพฺภตานิ, ยํ วา มหาสมุทฺเท อุทกนฺติ. เอตเทว ภนฺเต พหุตรํ ยทิทํ มหาสมุทฺเท อุทกํ, อปฺปมตฺตกานิ เทฺว วา ตีณิ วา อุทกผุสิตานิ อุพฺภตานิ, สงฺขมฺปิ น อุเปนฺติ อุปนิธมฺปิ น อุเปนฺติ กลภาคมฺปิ น อุเปนฺติ, มหาสมุทฺเท อุทกํ อุปนิธาย เทฺว วา ตีณิ วา อุทกผุสิตานิ อุพฺภตานีติ. เอวเมว โข ภิกฺขเว อริยสาวกสฺส ฯเปฯ โยโค กรณีโยติ.  สตฺตมํ.}

\csromanheader{2}{8. ทุติย\-มหาสมุทฺท\-สุตฺต}
\proseitem{1128}{เสยฺยถาปิ ภิกฺขเว มหาสมุทฺเท อุทกํ ปริกฺขยํ\csromansharedfootnote{b99628064fdd170a}{มหาสมุทฺโท ปริกฺขยํ (สี, สฺยา, กํ) สํ 1. 355 ปิฏฺเฐปิ.} ปริยาทานํ คจฺเฉยฺย ฐเปตฺวา เทฺว วา ตีณิ วา อุทกผุสิตานิ. ตํ กิํ มญฺญถ ภิกฺขเว, กตมํ นุ โข พหุตรํ, ยํ วา มหาสมุทฺเท อุทกํ ปริกฺขีณํ ปริยาทินฺนํ, ยานิ เทฺว วา ตีณิ วา อุทกผุสิตานิ อวสิฏฺฐานีติ. เอตเทว ภนฺเต พหุตรํ มหาสมุทฺเท อุทกํ ยทิทํ ปริกฺขีณํ ปริยาทินฺนํ, อปฺปมตฺตกานิ เทฺว วา ตีณิ วา อุทกผุสิตานิ อวสิฏฺฐานิ, สงฺขมฺปิ น อุเปนฺติ อุปนิธมฺปิ น อุเปนฺติ กลภาคมฺปิ น อุเปนฺติ, มหาสมุทฺเท อุทกํ}

\csromanmatikaanchor{mtk.307}
\csromantocmarkref{subparagraph}{9-10. ปพฺพตูปมสุตฺตานิ}{mtk.307}
\bookmark[dest={mtk.307},level=5]{9-10. ปพฺพตูปมสุตฺตานิ}
\csromantocmarknopage{subparagraph}{7. ปฐมอามกธญฺญเปยฺยาลวคฺค}{mtk.308}
\bookmark[dest={mtk.308},level=5]{7. ปฐมอามกธญฺญเปยฺยาลวคฺค}
\prose{\csromanfolio{403}ปริกฺขีณํ ปริยาทินฺนํ อุปนิธาย เทฺว วา ตีณิ วา อุทกผุสิตานิ อวสิฏฺฐานีติ. เอวเมว โข ภิกฺขเว อริยสาวกสฺส ฯเปฯ โยโค กรณีโยติ.  อฏฺฐมํ.}

\csromanheader{2}{9. ปฐม\-ปพฺพตูปม\-สุตฺต}
\proseitem{1129}{เสยฺยถาปิ ภิกฺขเว ปุริโส หิมวโต ปพฺพตราชสฺส สตฺต สาสปมตฺติโย ปาสาณสกฺขรา อุปนิกฺขิเปยฺย. ตํ กิํ มญฺญถ ภิกฺขเว, กตมํ นุ โข พหุตรํ, ยา วา สตฺต สาสปมตฺติโย ปาสาณสกฺขรา อุปนิกฺขิตฺตา, อยํ วา หิมวา ปพฺพตราชาติ. เอตเทว ภนฺเต พหุตรํ ยทิทํ หิมวา ปพฺพตราชา, อปฺปมตฺติกา สตฺต สาสปมตฺติโย ปาสาณสกฺขรา อุปนิกฺขิตฺตา, สงฺขมฺปิ น อุเปนฺติ อุปนิธมฺปิ น อุเปนฺติ กลภาคมฺปิ น อุเปนฺติ, หิมวนฺตํ ปพฺพตราชานํ อุปนิธาย สตฺต สาสปมตฺติโย ปาสาณสกฺขรา อุปนิกฺขิตฺตาติ. เอวเมว โข ภิกฺขเว อริยสาวกสฺส ฯเปฯ โยโค กรณีโยติ.  นวมํ.}

\csromanheader{2}{10. ทุติย\-ปพฺพตูปม\-สุตฺต}
\proseitem{1130}{เสยฺยถาปิ ภิกฺขเว หิมวา ปพฺพตราชา ปริกฺขยํ ปริยาทานํ คจฺเฉยฺย ฐเปตฺวา สตฺต สาสปมตฺติโย ปาสาณสกฺขรา. ตํ กิํ มญฺญถ ภิกฺขเว, กตมํ นุ โข พหุตรํ, ยํ วา หิมวโต ปพฺพตราชสฺส ปริกฺขีณํ ปริยาทินฺนํ, ยา วา สตฺต สาสปมตฺติโย ปาสาณสกฺขรา อวสิฏฺฐาติ. เอตเทว ภนฺเต พหุตรํ หิมวโต ปพฺพตราชสฺส ยทิทํ ปริกฺขีณํ ปริยาทินฺนํ, อปฺปมตฺติกา สตฺต สาสปมตฺติโย ปาสาณสกฺขรา อวสิฏฺฐา, สงฺขมฺปิ น อุเปนฺติ อุปนิธมฺปิ น อุเปนฺติ กลภาคมฺปิ น อุเปนฺติ, หิมวโต ปพฺพตราชสฺส ปริกฺขีณํ ปริยาทินฺนํ, อุปนิธาย สตฺต สาสปมตฺติโย ปาสาณสกฺขรา อวสิฏฺฐาติ. เอวเมว โข ภิกฺขเว อริยสาวกสฺส ทิฏฺฐิ\-สมฺปนฺนสฺส ปุคฺคลสฺส อภิสเมตาวิโน เอตเทว พหุตรํ ทุกฺขํ, ยทิทํ ปริกฺขีณํ ปริยาทินฺนํ, อปฺปมตฺตกํ อวสิฏฺฐํ, สงฺขมฺปิ น อุเปติ อุปนิธมฺปิ น อุเปติ กลภาคมฺปิ น อุเปติ, ปุริมํ ทุกฺข\-กฺขนฺธํ ปริกฺขีณํ ปริยาทินฺนํ อุปนิธาย, ยทิทํ สตฺตกฺขตฺตุ\-ปรมตา. “โย อิทํ ทุกฺขนฺ”ติ ยถาภูตํ ปชานาติ ฯเปฯ “อยํ ทุกฺขนิโรธ\-คามินี ปฏิปทา”ติ ยถาภูตํ ปชานาติ.}

\csromanmatikaanchor{mtk.309}
\csromantocmarkref{subparagraph}{1-10. อญฺญตฺรสุตฺตาทีนิ}{mtk.309}
\bookmark[dest={mtk.309},level=5]{1-10. อญฺญตฺรสุตฺตาทีนิ}
\csromantocmarknopage{subparagraph}{8. ทุติยอามกธญฺญเปยฺยาลวคฺค}{mtk.310}
\bookmark[dest={mtk.310},level=5]{8. ทุติยอามกธญฺญเปยฺยาลวคฺค}
\prose{\csromanfolio{404}ตสฺมาติห ภิกฺขเว “อิทํ ทุกฺขนฺ”ติ โยโค กรณีโย ฯเปฯ “อยํ ทุกฺขนิโรธ\-คามินี ปฏิปทา”ติ โยโค กรณีโยติ.  ทสมํ. อภิสมย\-วคฺโค ฉฏฺโฐ.}

\vspace{\tipitakaparskip}
\csromancenter{\textbf{ตสฺสุทฺทานํ}}
\setlength{\csromangathapagewidth}{0pt}
\setlength{\csromangathawakwidth}{0pt}
\csromangathameasuregroup{นขสิขา โปกฺขรณี, \\ ปถวี เทฺว สมุทฺทา เทฺว,}{\csromangathabat{นขสิขา โปกฺขรณี,}{สํเภชฺช อปเร ทุเว.} \\ \csromangathabat{ปถวี เทฺว สมุทฺทา เทฺว,}{เทฺวมา จ ปพฺพตูปมาติ.}}
\csromangathasetleft{นขสิขา โปกฺขรณี, \\ ปถวี เทฺว สมุทฺทา เทฺว,}
\csromangathagroup{\csromangathastanza{\csromangathabat{นขสิขา โปกฺขรณี,}{สํเภชฺช อปเร ทุเว.} \\ \csromangathabat{ปถวี เทฺว สมุทฺทา เทฺว,}{เทฺวมา จ ปพฺพตูปมาติ.}}}

\chapterhead{7. ปฐม\-อามกธญฺญ\-เปยฺยาล\-วคฺค}
\csromanheader{2}{1. \textbf{อญฺญตฺรสุตฺต}}
\proseitem{1131}{อถ โข ภควา ปริตฺตํ นขสิขายํ ปํสุํ อาโรเปตฺวา ภิกฺขู อามนฺเตสิ “ตํ กิํ มญฺญถ ภิกฺขเว, กตมํ นุ โข พหุตรํ, โย วายํ มยา ปริตฺโต นขสิขายํ ปํสุ อาโรปิโต, อยํ วา มหาปถวี”ติ. เอตเทว ภนฺเต พหุตรํ ยทิทํ มหาปถวี, อปฺปมตฺตกายํ ภควตา ปริตฺโต นขสิขายํ ปํสุ อาโรปิโต, สงฺขมฺปิ น อุเปติ อุปนิธมฺปิ น อุเปติ กลภาคมฺปิ น อุเปติ, มหาปถวิํ อุปนิธาย ภควตา ปริตฺโต นขสิขายํ ปํสุ อาโรปิโตติ. เอวเมว โข ภิกฺขเว อปฺปมตฺตกา เต สตฺตา, เย มนุสฺเสสุ ปจฺจาชายนฺติ. อถ โข เอเตว พหุตรา สตฺตา เย อญฺญตฺร มนุสฺเสหิ\csromansharedfootnote{949fdf7844381e0b}{มนุสฺเสสุ (อิ, ก)} ปจฺจาชายนฺติ. ตํ กิสฺส เหตุ, อทิฏฺฐตฺตา ภิกฺขเว จตุนฺนํ อริยสจฺจานํ. กตเมสํ จตุนฺนํ, ทุกฺขสฺส อริยสจฺจสฺส ฯเปฯ ทุกฺขนิโรธ\-คามินิยา ปฏิปทาย อริยสจฺจสฺส.}

\prose{ตสฺมาติห ภิกฺขเว “อิทํ ทุกฺขนฺ”ติ โยโค กรณีโย ฯเปฯ “อยํ ทุกฺขนิโรธ\-คามินี ปฏิปทา”ติ โยโค กรณีโยติ.  ปฐมํ.}

\csromanheader{2}{2. \textbf{ปจฺจนฺตสุตฺต}}
\proseitem{1132}{อถ โข ภควา ปริตฺตํ นขสิขายํ ปํสุํ อาโรเปตฺวา ภิกฺขู อามนฺเตสิ “ตํ กิํ มญฺญถ ภิกฺขเว, กตมํ นุ โข พหุตรํ, โย}

\prose{\csromanfolio{405}วายํ มยา ปริตฺโต นขสิขายํ ปํสุ อาโรปิโต, อยํ วา มหาปถวี”ติ. เอตเทว ภนฺเต พหุตรํ ยทิทํ มหาปถวี, อปฺปมตฺตกายํ ภควตา ปริตฺโต นขสิขายํ ปํสุ อาโรปิโต, สงฺขมฺปิ น อุเปติ อุปนิธมฺปิ น อุเปติ กลภาคมฺปิ น อุเปติ, มหาปถวิํ อุปนิธาย ภควตา ปริตฺโต นขสิขายํ ปํสุ อาโรปิโตติ.}

\prose{เอวเมว โข ภิกฺขเว อปฺปมตฺตกา เต สตฺตา, เย มชฺฌิเมสุ ชนปเทสุ ปจฺจาชายนฺติ, อถ โข เอเตว พหุตรา สตฺตา, เย ปจฺจนฺติเมสุ ชนปเทสุ ปจฺจาชายนฺติ อวิญฺญาตาเรสุ มิลกฺเขสุ\csromansharedfootnote{7282586249ff41c5}{มิลกฺขูสุ (สฺยา, กํ, ก)} ฯเปฯ.  ทุติยํ.}

\csromanheader{2}{3. \textbf{ปญฺญาสุตฺต}}
\proseitem{1133}{เอวเมว โข ภิกฺขเว อปฺปกา เต สตฺตา, เย ปน อริเยน ปญฺญาจกฺขุนา สมนฺนาคตา, อถ โข เอเตว พหุตรา สตฺตา, เย อวิชฺชาคตา สมฺมุฬฺหา ฯเปฯ.  ตติยํ.}

\csromanheader{2}{4. สุราเมรย\-สุตฺต}
\proseitem{1134}{เอวเมว โข ภิกฺขเว อปฺปกา เต สตฺตา, เย สุรา\-เมรย\-มชฺช\-ปฺปมาทฏฺฐานา ปฏิวิรตา, อถ โข เอเตว พหุตรา สตฺตา, เย สุรา\-เมรย\-มชฺช\-ปฺปมาทฏฺฐานา ปฏิวิรตา ฯเปฯ.  จตุตฺถํ.}

\csromanheader{2}{5. \textbf{โอทกสุตฺต}}
\proseitem{1135}{เอวเมว โข ภิกฺขเว อปฺปกา เต สตฺตา, เย ถลชา, อถ โข เอเตว พหุตรา สตฺตา, เย อุทกชา. ตํ กิสฺส เหตุ ฯเปฯ.  ปญฺจมํ.}

\csromanheader{2}{6. \textbf{มตฺเตยฺยสุตฺต}}
\proseitem{1136}{เอวเมว โข ภิกฺขเว อปฺปกา เต สตฺตา, เย มตฺเตยฺยา, อถ โข เอเตว พหุตรา สตฺตา, เย อมตฺเตยฺยา ฯเปฯ.  ฉฏฺฐํ.}

\csromanheader{2}{7. \textbf{เปตฺเตยฺยสุตฺต}}
\csromanmatikaanchor{mtk.311}
\csromantocmarkref{subparagraph}{1-10. ปาณาติปาตสุตฺตาทีนิ}{mtk.311}
\bookmark[dest={mtk.311},level=5]{1-10. ปาณาติปาตสุตฺตาทีนิ}
\csromantocmarknopage{subparagraph}{9. ตติยอามกธญฺญเปยฺยาลวคฺค}{mtk.312}
\bookmark[dest={mtk.312},level=5]{9. ตติยอามกธญฺญเปยฺยาลวคฺค}
\proseitem{1137}{\csromanfolio{406}เอวเมว โข ภิกฺขเว อปฺปกา เต สตฺตา, เย เปตฺเตยฺยา, อถ โข เอเตว พหุตรา สตฺตา, เย อเปตฺเตยฺยา ฯเปฯ.  สตฺตมํ.}

\csromanheader{2}{8. \textbf{สามญฺญสุตฺต}}
\proseitem{1138}{เอวเมว โข ภิกฺขเว อปฺปกา เต สตฺตา, เย สามญฺญา, อถ โข เอเตว พหุตรา สตฺตา, เย อสามญฺญา ฯเปฯ.  อฏฺฐมํ.}

\csromanheader{2}{9. \textbf{พฺรหฺมญฺญสุตฺต}}
\proseitem{1139}{เอวเมว โข ภิกฺขเว อปฺปกา เต สตฺตา, เย พฺรหฺมญฺญา อถ โข เอเตว พหุตรา สตฺตา, เย อพฺรหฺมญฺญา ฯเปฯ.  นวมํ.}

\csromanheader{2}{10. \textbf{ปจายิกสุตฺต}}
\proseitemwithclosermidruled{1140}{เอวเมว โข ภิกฺขเว อปฺปกา เต สตฺตา, เย กุเล เชฏฺฐาปจายิโน. อถ โข เอเตว พหุตรา สตฺตา, เย กุเล อเชฏฺฐาปจายิโนติ\csromansharedfootnote{cf3715f85b171609}{อกุเล เชฏฺฐาปจายิโนติ (สฺยา, กํ)}.  ทสมํ.}{ปฐม\-อามกธญฺญ\-เปยฺยาล\-วคฺโค สตฺตโม.}
\csromancenter{\textbf{ตสฺสุทฺทานํ}}
\setlength{\csromangathapagewidth}{0pt}
\setlength{\csromangathawakwidth}{0pt}
\csromangathameasuregroup{อญฺญตฺร ปจฺจนฺตํ ปญฺญา, \\ มตฺเตยฺย เปตฺเตยฺยา จาปิ,}{\csromangathabat{อญฺญตฺร ปจฺจนฺตํ ปญฺญา,}{สุราเมรโยอทกา.} \\ \csromangathabat{มตฺเตยฺย เปตฺเตยฺยา จาปิ,}{สามญฺญํ พฺรหฺมปจายิกนฺติ.}}
\csromangathasetleft{อญฺญตฺร ปจฺจนฺตํ ปญฺญา, \\ มตฺเตยฺย เปตฺเตยฺยา จาปิ,}
\csromangathagroup{\csromangathastanza{\csromangathabat{อญฺญตฺร ปจฺจนฺตํ ปญฺญา,}{สุราเมรโยอทกา.} \\ \csromangathabat{มตฺเตยฺย เปตฺเตยฺยา จาปิ,}{สามญฺญํ พฺรหฺมปจายิกนฺติ.}}}

\chapterhead{8. ทุติย\-อามกธญฺญ\-เปยฺยาล\-วคฺค}
\csromanheader{2}{1. ปาณาติปาต\-สุตฺต}
\proseitem{1141}{เอวเมว โข ภิกฺขเว อปฺปกา เต สตฺตา, เย ปาณาติปาตา ปฏิวิรตา. อถ โข เอเตว พหุตรา สตฺตา, เย ปาณาติปาตา อปฺปฏิวิรตา. ตํ กิสฺส เหตุ ฯเปฯ.  ปฐมํ.}

\csromanheader{2}{2. อทินฺนาทาน\-สุตฺต}
\proseitem{1142}{\csromanfolio{407}เอวเมว โข ภิกฺขเว อปฺปกา เต สตฺตา, เย อทินฺนาทานา ปฏิวิรตา. อถ โข เอเตว พหุตรา สตฺตา, เย อทินฺนาทานา อปฺปฏิวิรตา ฯเปฯ.  ทุติยํ.}

\csromanheader{2}{3. กาเมสุมิจฺฉาจาร\-สุตฺต}
\proseitem{1143}{เอวเมว โข ภิกฺขเว อปฺปกา เต สตฺตา, เย กาเมสุ\-มิจฺฉาจารา ปฏิวิรตา. อถ โข เอเตว พหุตรา สตฺตา, เย กาเมสุ มิจฺฉาจารา อปฺปฏิวิรตา ฯเปฯ.  ตติยํ.}

\csromanheader{2}{4. \textbf{มุสาวาทสุตฺต}}
\proseitem{1144}{เอวเมว โข ภิกฺขเว อปฺปกา เต สตฺตา, เย มุสาวาทา ปฏิวิรตา. อถ โข เอเตว พหุตรา สตฺตา, เย มุสาวาทา อปฺปฏิวิรตา ฯเปฯ.  จตุตฺถํ.}

\csromanheader{2}{5. \textbf{เปสุญฺญสุตฺต}}
\proseitem{1145}{เอวเมว โข ภิกฺขเว อปฺปกา เต สตฺตา, เย ปิสุณาย วาจาย ปฏิวิรตา. อถ โข เอเตว พหุตรา สตฺตา, เย ปิสุณาย วาจาย อปฺปฏิวิรตา ฯเปฯ.  ปญฺจมํ.}

\csromanheader{2}{6. ผรุสวาจา\-สุตฺต}
\proseitem{1146}{เอวเมว โข ภิกฺขเว อปฺปกา เต สตฺตา, เย ผรุสาย วาจาย ปฏิวิรตา. อถ โข เอเตว พหุตรา สตฺตา, เย ผรุสาย วาจาย อปฺปฏิวิรตา ฯเปฯ.  ฉฏฺฐํ.}

\csromanheader{2}{7. สมฺผปฺปลาป\-สุตฺต}
\proseitem{1147}{เอวเมว โข ภิกฺขเว อปฺปกา เต สตฺตา, เย สมฺผปฺปลาปา ปฏิวิรตา. อถ โข เอเตว พหุตรา สตฺตา, เย สมฺผปฺปลาปา อปฺปฏิวิรตา ฯเปฯ.  สตฺตมํ.}

\csromanheader{2}{8. \textbf{พีชคามสุตฺต}}
\csromanmatikaanchor{mtk.313}
\csromantocmarkref{subparagraph}{1-10. นจฺจคีตสุตฺตาทีนิ}{mtk.313}
\bookmark[dest={mtk.313},level=5]{1-10. นจฺจคีตสุตฺตาทีนิ}
\proseitem{1148}{\csromanfolio{408}เอวเมว โข ภิกฺขเว อปฺปกา เต สตฺตา, เย พีชคาม\-ภูตคาม\-สมารมฺภา\csromansharedfootnote{2fbc91b9c8ed4836}{พีชคามภูตคามสมารพฺภา (ก)} ปฏิวิรตา. อถ โข เอเตว พหุตรา สตฺตา, เย พีชคาม\-ภูตคาม\-สมารมฺภา อปฺปฏิวิรตา ฯเปฯ.  อฏฺฐมํ.}

\csromanheader{2}{9. วิกาลโภชน\-สุตฺต}
\proseitem{1149}{เอวเมว โข ภิกฺขเว อปฺปกา เต สตฺตา, เย วิกาลโภชนา ปฏิวิรตา. อถ โข เอเตว พหุตรา สตฺตา, เย วิกาลโภชนา อปฺปฏิวิรตา ฯเปฯ.  นวมํ.}

\csromanheader{2}{10. คนฺธวิเลปน\-สุตฺต}
\proseitem{1150}{เอวเมว โข ภิกฺขเว อปฺปกา เต สตฺตา, เย มาลา\-คนฺธ\-วิเลปน\-ธารณ\-มณฺฑน\-วิภูสน\-ฏฺฐานา ปฏิวิรตา. อถ โข เอเตว พหุตรา สตฺตา, เย มาลาคนฺธวิเลปนธารมณฺฑนวิภูสนฏฺฐานา อปฺปฏิวิรตา. ทสมํ. ทุติย\-อามกธญฺญ\-เปยฺยาล\-วคฺโค อฏฺฐโม.}

\vspace{\tipitakaparskip}
\csromancenter{\textbf{ตสฺสุทฺทานํ}}
\setlength{\csromangathapagewidth}{0pt}
\setlength{\csromangathawakwidth}{0pt}
\csromangathameasuregroup{ปาณํ อทินฺนํ กาเมสุ, \\ ผรุสํ สมฺผปฺปลาปํ,}{\csromangathabat{ปาณํ อทินฺนํ กาเมสุ,}{มุสาวาทญฺจ เปสุญฺญํ.} \\ \csromangathabat{ผรุสํ สมฺผปฺปลาปํ,}{พีชญฺจ วิกาลํ คนฺธนฺติ.}}
\csromangathasetleft{ปาณํ อทินฺนํ กาเมสุ, \\ ผรุสํ สมฺผปฺปลาปํ,}
\csromangathagroup{\csromangathastanza{\csromangathabat{ปาณํ อทินฺนํ กาเมสุ,}{มุสาวาทญฺจ เปสุญฺญํ.} \\ \csromangathabat{ผรุสํ สมฺผปฺปลาปํ,}{พีชญฺจ วิกาลํ คนฺธนฺติ.}}}

\chapterhead{9. ตติย\-อามกธญฺญ\-เปยฺยาล\-วคฺค}
\csromanheader{2}{1. \textbf{นจฺจคีตสุตฺต}}
\proseitem{1151}{เอวเมว โข ภิกฺขเว อปฺปกา เต สตฺตา, เย นจฺจ\-คีต\-วาทิต\-วิสูก\-ทสฺสนา ปฏิวิรตา. อถ โข เอเตว พหุตรา สตฺตา, เย นจฺจ คีต วาทิต วิสูก ทสฺสนา อปฺปฏิวิรตา. ตํ กิสฺส เหตุ ฯเปฯ.  ปฐมํ.}

\csromanheader{2}{2. อุจฺจาสยน\-สุตฺต}
\proseitem{1152}{\csromanfolio{409}เอวเมว โข ภิกฺขเว อปฺปกา เต สตฺตา, เย อุจฺจาสยน\-มหาสยนา ปฏิวิรตา. อถ โข เอเตว พหุตรา สตฺตา, เย อุจฺจาสยน\-มหาสยนา อปฺปฏิวิรตา ฯเปฯ.  ทุติยํ.}

\csromanheader{2}{3. ชาตรูปรชต\-สุตฺต}
\proseitem{1153}{เอวเมว โข ภิกฺขเว อปฺปกา เต สตฺตา, เย ชาตรูป\-รชต\-ปฏิคฺคหณา ปฏิวิรตา. อถ โข เอเตว พหุตรา สตฺตา, เย ชาตรูป\-รชต\-ปฏิคฺคหณา อปฺปฏิวิรตา ฯเปฯ.  ตติยํ.}

\csromanheader{2}{4. อามกธญฺญ\-สุตฺต}
\proseitem{1154}{เอวเมว โข ภิกฺขเว อปฺปกา เต สตฺตา, เย อามก\-ธญฺญ\-ปฏิคฺคหณา ปฏิวิรตา. อถ โข เอเตว พหุตรา สตฺตา, เย อามก\-ธญฺญ\-ปฏิคฺคหณา อปฺปฏิวิรตา ฯเปฯ.  จตุตฺถํ.}

\csromanheader{2}{5. อามกมํส\-สุตฺต}
\proseitem{1155}{เอวเมว โข ภิกฺขเว อปฺปกา เต สตฺตา, เย อามก\-มํส\-ปฏิคฺคหณา ปฏิวิรตา. อถ โข เอเตว พหุตรา สตฺตา, เย อามก\-มํส\-ปฏิคฺคหณา อปฺปฏิวิรตา ฯเปฯ.  ปญฺจมํ.}

\csromanheader{2}{6. \textbf{กุมาริกสุตฺต}}
\proseitem{1156}{เอวเมว โข ภิกฺขเว อปฺปกา เต สตฺตา, เย อิตฺถิ\-กุมาริก\-ปฏิคฺคหณา\csromansharedfootnote{4058a4cf107a5407}{อิตฺถิกุมาริกาปฏิคฺคหณา (ก)} ปฏิวิรตา. อถ โข เอเตว พหุตรา สตฺตา, เย อิตฺถิ\-กุมาริก\-ปฏิคฺคหณา อปฺปฏิวิรตา ฯเปฯ.  ฉฏฺฐํ.}

\csromanheader{2}{7. \textbf{ทาสิทาสสุตฺต}}
\proseitem{1157}{เอวเมว โข ภิกฺขเว อปฺปกา เต สตฺตา, เย ทาสิ\-ทาส\-ปฏิคฺคหณา\csromansharedfootnote{3122376578ccf466}{ทาสีทาสปฏิคฺคหณา (สฺยา, กํ, อิ)} ปฏิวิรตา. อถ โข เอเตว พหุตรา สตฺตา, เย ทาสิ\-ทาส\-ปฏิคฺคหณา อปฺปฏิวิรตา ฯเปฯ.  สตฺตมํ.}

\csromanheader{2}{8. \textbf{อเชฬกสุตฺต}}
\csromanmatikaanchor{mtk.314}
\csromanmatikaanchor{mtk.315}
\csromantocmarknopage{subparagraph}{10. จตฺตุตฺถอามกธญฺญเปยฺยาลวคฺค}{mtk.314}
\bookmark[dest={mtk.314},level=5]{10. จตฺตุตฺถอามกธญฺญเปยฺยาลวคฺค}
\csromantocmarkref{subparagraph}{1-11. เขตฺตวตฺถุสุตฺตาทีนิ}{mtk.315}
\bookmark[dest={mtk.315},level=5]{1-11. เขตฺตวตฺถุสุตฺตาทีนิ}
\proseitem{1158}{\csromanfolio{410}เอวเมว โข ภิกฺขเว อปฺปกา เต สตฺตา, เย อเช\-ฬก\-ปฏิคฺคหณา ปฏิวิรตา. อถ โข เอเตว พหุตรา สตฺตา, เย อเช\-ฬก\-ปฏิคฺคหณา อปฺปฏิวิรตา ฯเปฯ.  อฏฺฐมํ.}

\csromanheader{2}{9. กุกฺกุฏสูกร\-สุตฺต}
\proseitem{1159}{เอวเมว โข ภิกฺขเว อปฺปกา เต สตฺตา, เย กุกฺกุฏ\-สูกร\-ปฏิคฺคหณา ปฏิวิรตา. อถ โข เอเตว พหุตรา สตฺตา, เย กุกฺกุฏ\-สูกร\-ปฏิคฺคหณา อปฺปฏิวิรตา ฯเปฯ.  นวมํ.}

\csromanheader{2}{10. หตฺถิควสฺส\-สุตฺต}
\proseitemwithclosermidruled{1160}{เอวเมว โข ภิกฺขเว อปฺปกา เต สตฺตา, เย หตฺถิ\-คว\-สฺส\-วฬว\-ปฏิคฺคหณา\csromansharedfootnote{bf9cd2b03f78e9f7}{หตฺถิควสฺสวฬวาปฏิคฺคหณา (สฺยา, กํ, อิ, ก)} ปฏิวิรตา. อถ โข เอเตว พหุตรา สตฺตา, เย หตฺถิควสฺสวฬาวปฏิคฺคหณา อปฺปฏิวิรตา.  ทสมํ.}{ตติย\-อามกธญฺญ\-เปยฺยาล\-วคฺโค นวโม.}
\csromancenter{\textbf{ตสฺสุทฺทานํ}}
\setlength{\csromangathapagewidth}{0pt}
\setlength{\csromangathawakwidth}{0pt}
\csromangathameasuregroup{นจฺจํ สยนํ รชตํ, \\ ทาสี อเชฬกญฺเจว,}{\csromangathabat{นจฺจํ สยนํ รชตํ,}{ธญฺญํ มํสํ กุมาริกา.} \\ \csromangathabat{ทาสี อเชฬกญฺเจว,}{กุกฺกุฏสูกรหตฺถีติ.}}
\csromangathasetleft{นจฺจํ สยนํ รชตํ, \\ ทาสี อเชฬกญฺเจว,}
\csromangathagroup{\csromangathastanza{\csromangathabat{นจฺจํ สยนํ รชตํ,}{ธญฺญํ มํสํ กุมาริกา.} \\ \csromangathabat{ทาสี อเชฬกญฺเจว,}{กุกฺกุฏสูกรหตฺถีติ.}}}

\chapterhead{10. จตุตฺถ\-อามกธญฺญ\-เปยฺยาล\-วคฺค}
\csromanheader{2}{1. เขตฺตวตฺถุ\-สุตฺต}
\proseitem{1161}{เอวเมว โข ภิกฺขเว อปฺปกา เต สตฺตา, เย เขตฺตวตฺถุ\-ปฏิคฺคหณา ปฏิวิรตา. อถ โข เอเตว พหุตรา สตฺตา, เย เขตฺตวตฺถุ\-ปฏิคฺคหณา อปฺปฏิวิรตา ฯเปฯ.  ปฐมํ.}

\csromanheader{2}{2. กยวิกฺกย\-สุตฺต}
\proseitem{1162}{\csromanfolio{411}เอวเมว โข ภิกฺขเว อปฺปกา เต สตฺตา, เย กยวิกฺกยา ปฏิวิรตา. อถ โข เอเตว พหุตรา สตฺตา, เย กยวิกฺกยา อปฺปฏิวิรตา ฯเปฯ.  ทุติยํ.}

\csromanheader{2}{3. \textbf{ทูเตยฺยสุตฺต}}
\proseitem{1163}{เอวเมว โข ภิกฺขเว อปฺปกา เต สตฺตา, เย ทูเตยฺยปหินคมนานุโยคา ปฏิวิรตา. อถ โข เอเตว พหุตรา สตฺตา, เย ทูเตยฺยปหินคมนานุโยคา อปฺปฏิวิรตา ฯเปฯ.  ตติยํ.}

\csromanheader{2}{4. \textbf{ตุลากูฏสุตฺต}}
\proseitem{1164}{เอวเมว โข ภิกฺขเว อปฺปกา เต สตฺตา, เย ตุลา\-กูฏ\-กํส\-กูฏ\-มาน\-กูฏา ปฏิวิรตา. อถ โข เอเตว พหุตรา สตฺตา, เย ตุลา\-กูฏ\-กํส\-กูฏ\-มาน\-กูฏา อปฺปฏิวิรตา ฯเปฯ.  จตุตฺถํ.}

\csromanheader{2}{5. \textbf{อุกฺโกฏนสุตฺต}}
\proseitem{1165}{เอวเมว โข ภิกฺขเว อปฺปกา เต สตฺตา, เย อุกฺโกฏน\-วญฺจน\-นิกติ\-สาจิโยคา\csromansharedfootnote{63ba1a1068af8b3a}{อุกฺโกฏนวญฺจนนิกติสาวิโยคา (สฺยา, กํ, อิ, ก)} ปฏิวิรตา, อถ โข เอเตว พหุตรา สตฺตา, เย อุกฺโกฏน\-วญฺจน\-นิกติ\-สาจิโยคา อปฺปฏิวิรตา ฯเปฯ.  ปญฺจมํ.}

\csromanheader{2}{6-11. เฉทนาทิสุตฺต}
\proseitem{1166-1171}{เอวเมว โข ภิกฺขเว อปฺปกา เต สตฺตา, เย เฉทน\-วธ\-พนฺธน\-วิปราโมส\-อาโลป\-สหสาการา\csromansharedfootnote{3511a3ba078f619b}{สาหสาการา (ก)} ปฏิวิรตา. อถ โข เอเตว พหุตรา สตฺตา, เย เฉทน\-วธ\-พนฺธน\-วิปราโมส\-อาโลป\-สหสาการา อปฺปฏิวิรตา. ตํ กิสฺส เหตุ, อทิฏฺฐตฺตา ภิกฺขเว จตุนฺนํ อริยสจฺจานํ. กตเมสํ จตุนฺนํ, ทุกฺขสฺส อริยสจฺจสฺส ฯเปฯ ทุกฺขนิโรธ\-คามินิยา ปฏิปทาย อริยสจฺจสฺส.}

\prosewithclosermidruled{ตสฺมาติห ภิกฺขเว “อิทํ ทุกฺขนฺ”ติ โยโค กรณีโย ฯเปฯ “อยํ ทุกฺขนิโรธ\-คามินี ปฏิปทา”ติ โยโค กรณีโยติ.  เอกาทสมํ.}{จตุตฺถ\-อามกธญฺญ\-เปยฺยาล\-วคฺโค ทสโม.}
\csromancenter{\textbf{ตสฺสุทฺทานํ}}
\csromanmatikaanchor{mtk.316}
\csromanmatikaanchor{mtk.317}
\csromantocmarknopage{subparagraph}{11. ปญฺจคติเปยฺยาลวคฺค}{mtk.316}
\bookmark[dest={mtk.316},level=5]{11. ปญฺจคติเปยฺยาลวคฺค}
\csromantocmarkref{subparagraph}{1-30. มนุสฺสจุตินิรยสุตฺตาทีนิ}{mtk.317}
\bookmark[dest={mtk.317},level=5]{1-30. มนุสฺสจุตินิรยสุตฺตาทีนิ}
\setlength{\csromangathapagewidth}{0pt}
\setlength{\csromangathawakwidth}{0pt}
\csromangathameasuregroup{เขตฺตํ กายํ ทูเตยฺยญฺจ, \\ เฉทนํ วธพนฺธนํ,}{\csromangathabat{เขตฺตํ กายํ ทูเตยฺยญฺจ,}{ตุลากูฏํ อุกฺโกฏนํ.} \\ \csromangathabat{เฉทนํ วธพนฺธนํ,}{วิปราโลปํ สาหสนฺติ.}}
\csromangathasetleft{เขตฺตํ กายํ ทูเตยฺยญฺจ, \\ เฉทนํ วธพนฺธนํ,}
\csromangathagroup{\csromangathastanza{\csromanfolio{412}\csromangathabat{เขตฺตํ กายํ ทูเตยฺยญฺจ,}{ตุลากูฏํ อุกฺโกฏนํ.} \\ \csromangathabat{เฉทนํ วธพนฺธนํ,}{วิปราโลปํ สาหสนฺติ.}}}

\csromanheader{6}{11. ปญฺจ\-คติ\-เปยฺยาล\-วคฺค}
\csromanheader{2}{1. มนุสฺส\-จุติ\-นิรย\-สุตฺต}
\proseitem{1172}{อถ โข ภควา ปริตฺตํ นขสิขายํ ปํสุํ อาโรเปตฺวา ภิกฺขู อามนฺเตสิ “ตํ กิํ มญฺญถ ภิกฺขเว, กตมํ นุ โข พหุตรํ, โย วายํ มยา ปริตฺโต นขสิขายํ ปํสุ อาโรปิโต, อยํ วา มหาปถวี”ติ. เอตเทว ภนฺเต พหุตรํ ยทิทํ มหาปถวี, อปฺปมตฺตกายํ ภควตา ปริตฺโต นขสิขายํ ปํสุ อาโรปิโต, สงฺขมฺปิ น อุเปติ อุปนิธมฺปิ น อุเปติ กลภาคมฺปิ น อุเปติ, มหาปถวิํ อุปนิธาย ภควตา ปริตฺโต นขสิขายํ ปํสุ อาโรปิโตติ. เอวเมว โข ภิกฺขเว อปฺปกา เต สตฺตา, เย มนุสฺสา จุตา มนุสฺเสสุ ปจฺจาชายนฺติ. อถ โข เอเตว พหุตรา สตฺตา, เย มนุสฺสา จุตา นิรเย ปจฺจาชายนฺติ ฯเปฯ.  ปฐมํ.}

\csromanheader{2}{2. \textbf{มนุสฺสจุติติรจฺฉนสุตฺต}}
\proseitem{1173}{เอวเมว โข ภิกฺขเว อปฺปกา เต สตฺตา, เย มนุสฺสา จุตา มนุสฺเสสุ ปจฺจาชายนฺติ. อถ โข เอเตว พหุตรา สตฺตา, เย มนุสฺสา จุตา ติรจฺฉนโยนิยา ปจฺจาชายนฺติ ฯเปฯ.  ทุติยํ.}

\csromanheader{2}{3. มนุสฺส\-จุติ\-เปตฺติวิสย\-สุตฺต}
\proseitem{1174}{เอวเมว โข ภิกฺขเว อปฺปกา เต สตฺตา, เย มนุสฺสา จุตา มนุสฺเสสุ ปจฺจาชายนฺติ. อถ โข เอเตว พหุตรา สตฺตา, เย มนุสฺสา จุตา เปตฺติวิสเย ปจฺจาชายนฺติ ฯเปฯ.  ตติยํ.}

\vspace{\tipitakaparskip}
\csromancenter{4-5-6. มนุสฺส\-จุติ\-เทว\-นิรยา\-ทิ\-สุตฺต}
\proseitem{1175-1177}{เอวเมว โข ภิกฺขเว อปฺปกา เต สตฺตา, เย มนุสฺสา จุตา เทเวสุ ปจฺจาชายนฺติ. อถ โข เอเตว พหุตรา สตฺตา,}

\prose{\csromanfolio{413}เย มนุสฺสา สุตา นิรเย ปจฺจชายนฺติ ฯเปฯ ติรจฺฉาน\-โยนิยา ปจฺจชายนฺติ ฯเปฯ เปตฺติวิสเย ปจฺจาชายนฺติ ฯเปฯ. ฉฏฺฐํ.}

\csromanheader{2}{7-9. เทว\-จุติ\-นิรยา\-ทิ\-สุตฺต}
\proseitem{1178-1180}{เอวเมว โข ภิกฺขเว อปฺปกา เต สตฺตา, เย เทวา จุตา เทเวสุ ปจฺจชายนฺติ. อถ โข เอเตว พหุตรา สตฺตา, เย เทวา จุตา นิรเย ปจฺจาชายนฺติ ฯเปฯ ติรจฺฉาน\-โยนิยา ปจฺจาชายนฺติ ฯเปฯ เปตฺติวิสเย ปจฺจาชายนฺติ ฯเปฯ.  นวมํ.}

\csromanheader{2}{10-12. เทวมนุสฺสนิรยสุตฺต}
\proseitem{1181-1183}{เอวเมว โข ภิกฺขเว อปฺปกา เต สตฺตา, เย เทวา จุตา มนุสฺเสสุ ปจฺจาชายนฺติ. อถ โข เอเตว พหุตรา สตฺตา, เย เทวา จุตา นิรเย ปจฺจาชายนฺติ ฯเปฯ ติรจฺฉาน\-โยนิยา ปจฺจาชายนฺติ ฯเปฯ เปตฺติวิสเย ปจฺจาชายนฺติ.  ทฺวาทสมํ.}

\csromanheader{2}{13-15. นิรย\-มนุสฺส\-นิรยา\-ทิ\-สุตฺต}
\proseitem{1184-1186}{เอวเมว โข ภิกฺขเว อปฺปกา เต สตฺตา, เย นิรยา จุตา มนุสฺเสสุ ปจฺจาชายนฺติ. อถ โข เอเตว พหุตรา สตฺตา, เย นิรยา จุตา นิรเย ปจฺจาชายนฺติ ฯเปฯ ติรจฺฉาน\-โยนิยา ปจฺจาชายนฺติ ฯเปฯ เปตฺติวิสเย ปจฺจาชายนฺติ ฯเปฯ.  ปนฺนรสมํ.}

\csromanheader{2}{16-18. นิรย\-เทว\-นิรยา\-ทิ\-สุตฺต}
\proseitem{1187-1179}{เอวเมว โข ภิกฺขเว อปฺปกา เต สตฺตา, เย นิรยา จุตา เทเวสุ ปจฺจาชายนฺติ. อถ โข เอเตว พหุตรา สตฺตา, เย นิรยา จุตา นิรเย ปจฺจาชายนฺติ ฯเปฯ ติรจฺฉาน\-โยนิยา ปจฺจาชายนฺติ ฯเปฯ เปตฺติวิสเย ปจฺจาชายนฺติ ฯเปฯ.  อฏฺฐารสมํ.}

\csromanheader{2}{19-21. ติรจฺฉาน\-มนุสฺส\-นิรยา\-ทิ\-สุตฺต}
\proseitem{1190-1192}{เอวเมว โข ภิกฺขเว อปฺปกา เต สตฺตา, เย ติรจฺฉาน\-โยนิยา จุตา มนุสฺเสสุ ปจฺจาชายนฺติ. อถ โข เอเตว พหุตรา สตฺตา, เย ติรจฺฉาน\-โยนิยา จุตา นิรเย ปจฺจาชายนฺติ ฯเปฯ \csromanfolio{414}ติรจฺฉาน\-โยนิยา ปจฺจาชายนฺติ ฯเปฯ เปตฺติวิสเย ปจฺจาชายนฺติ ฯเปฯ.  เอกวีสติมํ.}

\csromanheader{2}{22-24. ติรจฺฉาน\-เทว\-นิรยา\-ทิ\-สุตฺต}
\proseitem{1193-1195}{เอวเมว โข ภิกฺขเว อปฺปกา เต สตฺตา, เย ติรจฺฉาน\-โยนิยา จุตา เทเวสุ ปจฺจาชายนฺติ. อถ โข เอเตว พหุตรา สตฺตา, เย ติรจฺฉาน\-โยนิยา จุตา นิรเย ปจฺจาชายนฺติ ฯเปฯ ติรจฺฉาน\-โยนิยา ปจฺจาชายนฺติ ฯเปฯ เปตฺติวิสเย ปจฺจาชายนฺติ ฯเปฯ. จตุวีสติมํ.}

\csromanheader{2}{25-27. เปตฺติ\-มนุสฺส\-นิรยา\-ทิ\-สุตฺต}
\proseitem{1196-1198}{เอวเมว โข ภิกฺขเว อปฺปกา เต สตฺตา, เย เปตฺติวิสยา จุตา มนุสฺเสสุ ปจฺจาชายนฺติ. อถ โข เอเตว พหุตรา สตฺตา, เย เปตฺติวิสยา จุตา นิรเย ปจฺจาชายนฺติ ฯเปฯ ติรจฺฉาน\-โยนิยา ปจฺจาชายนฺติ ฯเปฯ เปตฺติวิสเย ปจฺจาชายนฺติ ฯเปฯ.  สตฺตวีสติมํ.}

\csromanheader{2}{28-29. เปตฺติ\-เทว\-นิรยา\-ทิ\-สุตฺต}
\proseitem{1199-1200}{เอวเมว โข ภิกฺขเว อปฺปกา เต สตฺตา, เย เปตฺติวิสยา จุตา เทเวสุ ปจฺจาชายนฺติ. อถ โข เอเตว พหุตรา สตฺตา, เย เปตฺติวิสยา จุตา นิรเย ปจฺจาชายนฺติ ฯเปฯ. เอวเมว โข ภิกฺขเว อปฺปกา เต สตฺตา, เย เปตฺติวิสยา จุตา เทเวสุ ปจฺจาชายนฺติ. อถ โข เอเตว พหุตรา สตฺตา, เย เปตฺติวิสยา จุตา ติรจฺฉาน\-โยนิยา ปจฺจาชายนฺติ ฯเปฯ.  เอกูน\-ติํสติมํ.}

\csromanheader{2}{30. เปตฺติ\-เทว\-เปตฺติวิสย\-สุตฺต}
\proseitem{1201}{เอวเมว โข ภิกฺขเว อปฺปกา เต สตฺตา, เย เปตฺติวิสยา จุตา เทเวสุ ปจฺจาชายนฺติ. อถ โข เอเตว พหุตรา สตฺตา, เย เปตฺติวิสยา จุตา เปตฺติวิสเย ปจฺจาชายนฺติ. ตํ กิสฺส เหตุ, อทิฏฺฐตฺตา ภิกฺขเว จตุนฺนํ อริยสจฺจานํ. กตเมสํ จตุนฺนํ, ทุกฺขสฺส อริยสจฺจสฺส ทุกฺข\-สมุทยสฺส อริยสจฺจสฺส ทุกฺข\-นิโรธสฺส อริยสจฺจสฺส ทุกฺขนิโรธ\-คามินิยา ปฏิปทาย อริยสจฺจสฺส.}

\prosewithclosermidruled{\csromanfolio{415}ตสฺมาติห ภิกฺขเว “อิทํ ทุกฺขนฺ”ติ โยโค กรณีโย, “อยํ ทุกฺขสมุทโย”ติ โยโค กรณีโย, “อยํ ทุกฺขนิโรโธ”ติ โยโค กรณีโย, “อยํ ทุกฺขนิโรธ\-คามินี ปฏิปทา”ติ โยโค กรณีโยติ. อิทมโวจ ภควา. อตฺตมนา เต ภิกฺขู ภควโต ภาสิตํ อภินนฺทุนฺติ.  ติํสติมํ.}{ปญฺจ\-คติ\-เปยฺยาล\-วคฺโค เอกาทสโม.}
\csromancenter{\textbf{ตสฺสุทฺทานํ}}
\setlength{\csromangathapagewidth}{0pt}
\setlength{\csromangathawakwidth}{0pt}
\csromangathameasuregroup{มนุสฺสโต จุตา ฉาปิ, \\ ติรจฺฉาเปตฺติวิสยา,}{\csromangathabat{มนุสฺสโต จุตา ฉาปิ,}{เทวา จุตา นิรยโต.} \\ \csromangathabat{ติรจฺฉาเปตฺติวิสยา,}{ติํสมตฺโต คติวคฺโคติ.} \\ สจฺจสํยุตฺตํ ทฺวาทสมํ.}
\csromangathameasurewak{สจฺจสํยุตฺตํ ทฺวาทสมํ.}
\csromangathasetleft{มนุสฺสโต จุตา ฉาปิ, \\ ติรจฺฉาเปตฺติวิสยา,}
\csromangathagroupwithclosermidruled{\csromangathastanza{\csromangathabat{มนุสฺสโต จุตา ฉาปิ,}{เทวา จุตา นิรยโต.} \\ \csromangathabat{ติรจฺฉาเปตฺติวิสยา,}{ติํสมตฺโต คติวคฺโคติ.}}\csromangathastanza{\csromangathawak{สจฺจสํยุตฺตํ ทฺวาทสมํ.}}}{\textbf{มหาวคฺโค ปญฺจโม.}}
\csromancenter{\textbf{ตสฺสุทฺทานํ}}
\setlength{\csromangathapagewidth}{0pt}
\setlength{\csromangathawakwidth}{0pt}
\csromangathameasuregroup{มคฺคโพชฺฌงฺคํ สติยา, \\ พลิทฺธิปาทานุรุทฺธา, \\ โสตาปตฺติ สจฺจญฺจาติ,}{\csromangathabat{มคฺคโพชฺฌงฺคํ สติยา,}{อินฺทฺริยํ สมฺมปฺปธานํ.} \\ \csromangathabat{พลิทฺธิปาทานุรุทฺธา,}{ฌานานาปาน สํยุตฺตํ.} \\ \csromangathabat{โสตาปตฺติ สจฺจญฺจาติ,}{มหาวคฺโคติ วุจฺจตีติ.}}
\csromangathasetleft{มคฺคโพชฺฌงฺคํ สติยา, \\ พลิทฺธิปาทานุรุทฺธา, \\ โสตาปตฺติ สจฺจญฺจาติ,}
\csromangathagroupwithclosermajor{\csromangathastanza{\csromangathabat{มคฺคโพชฺฌงฺคํ สติยา,}{อินฺทฺริยํ สมฺมปฺปธานํ.} \\ \csromangathabat{พลิทฺธิปาทานุรุทฺธา,}{ฌานานาปาน สํยุตฺตํ.} \\ \csromangathabat{โสตาปตฺติ สจฺจญฺจาติ,}{มหาวคฺโคติ วุจฺจตีติ.}}}{มหาวคฺค\-สํยุตฺต\-ปาฬิ นิฏฺฐิตา.}
